\paragraph{跨接口终端结论：奇偶权模态、同余 Fourier 扭曲与对称幂不交核的全显式封口}
本段收集三类可直接复核的“常数/闭式”条目：其一记录黄金均值稳定词 \(X_m\) 的奇偶权交错计数，
该对象对应一个隐含的 \(6\)-次割圆模态；其二把 Fold 纤维多重度按 \(|x|_1\) 的根单位扭曲提升为双变量配分函数 \(Z_m(y)\)，
由其四阶可积递推与显式有理生成函数抽取出 Tribonacci 主振荡模、Lee--Yang 型谱分歧边界、推前重量统计的
大数律/中心极限定律与全局大偏差轮廓；并进一步把纤维幂和矩提升为重量—碰撞配分函数 \(Z_m^{(q)}(y)\)（\(q=2,3,4\)），
给出其低阶常系数递推与有理生成函数、\(y=1\) 投影下的幽灵模态严格约去，以及由判别式分解所诱导的谱分歧量子化；
同时并联给出模 \(2,3,4\) 同余通道的谱半径与混合层级；
其三把不交图 \(\mathcal D_q\) 的 \(S_q\)-对称压缩矩阵 \(B_q\) 识别为 \(\mathrm{Sym}^q(M)\)，
并由此得到幺模性、显式逆矩阵与 \(n\)-步传输系数的 Fibonacci--二项式闭式。

\begin{conclusion}[黄金均值稳定词的奇偶权交错计数：严格 \(6\)-周期与有理生成函数]\label{con:xi-terminal-xm-parity-weight-6period}
令
\[
X_m:=\{x\in\{0,1\}^m:\ \forall i,\ x_ix_{i+1}=0\},
\qquad
A_m^{\mathrm{par}}:=\sum_{x\in X_m}(-1)^{|x|_1}.
\]
则对一切 \(m\ge 3\) 有二阶递推
\[
A_m^{\mathrm{par}}=A_{m-1}^{\mathrm{par}}-A_{m-2}^{\mathrm{par}},
\qquad
A_1^{\mathrm{par}}=0,\ A_2^{\mathrm{par}}=-1,
\]
从而 \(\bigl(A_m^{\mathrm{par}}\bigr)_{m\ge 1}\) 为严格 \(6\)-周期，并且
\[
A_m^{\mathrm{par}}=0\iff m\equiv 1,4\pmod 6,\qquad
A_m^{\mathrm{par}}=-1\iff m\equiv 2,3\pmod 6,\qquad
A_m^{\mathrm{par}}=1\iff m\equiv 0,5\pmod 6.
\]
若约定 \(A_0^{\mathrm{par}}:=1\)（空字的权重为 \(1\)），则生成函数满足闭式
\[
\sum_{m\ge 0}A_m^{\mathrm{par}}t^m=\frac{1-t}{1-t+t^2}.
\]
\end{conclusion}
\begin{proof}[证明]
按首段分解 \(X_m=\{0u:u\in X_{m-1}\}\sqcup \{10v:v\in X_{m-2}\}\)（\(m\ge 3\)）。
前一类不改变 \(|x|_1\)，后一类额外引入一个 \(1\)，故权重贡献多一因子 \((-1)\)。
于是得到递推
\(
A_m^{\mathrm{par}}=A_{m-1}^{\mathrm{par}}-A_{m-2}^{\mathrm{par}}
\)
及初值（由 \(X_1=\{0,1\},X_2=\{00,01,10\}\) 直接核算）。
特征方程为 \(\lambda^2-\lambda+1=0\)，根为 \(e^{\pm i\pi/3}\)，故序列严格 \(6\)-周期；
逐项展开即可得到模 \(6\) 的取值分类。
生成函数由标准递推求和（或直接对 \((1-t+t^2)\sum_{m\ge 0}A_m^{\mathrm{par}}t^m\) 展开并代入初值）得到。证毕。
\end{proof}

\begin{conclusion}[纤维—重量双变量配分函数的四阶可积性：有理生成函数与常系数递推]\label{con:xi-terminal-zm-order4-rational}
令
\[
X_m:=\{x\in\{0,1\}^m:\ \forall i,\ x_ix_{i+1}=0\},
\qquad
\Fold_m:\Omega_m=\{0,1\}^m\to X_m
\]
为标准折叠映射，并记纤维多重度 \(d_m(x)=|\Fold_m^{-1}(x)|\)。
对任意复参数 \(y\in\CC\)，定义纤维—重量双变量配分函数
\[
Z_m(y):=\sum_{x\in X_m} y^{|x|_1}\,d_m(x)
=\sum_{a\in\Omega_m} y^{|\Fold_m(a)|_1}.
\]
则对一切 \(m\ge 4\) 成立严格四阶递推
\[
\boxed{\;
Z_m(y)=Z_{m-1}(y)+(2y+1)Z_{m-2}(y)-Z_{m-3}(y)-y(y+1)Z_{m-4}(y)\;}
\]
并且生成函数在 \(\ZZ[y][[t]]\) 中具有显式有理闭式
\[
\boxed{\;
\sum_{m\ge 0} Z_m(y)\,t^m
=\frac{1+yt-yt^2-y^2t^3}{1-t-(2y+1)t^2+t^3+y(y+1)t^4}\; }.
\]
因此，对固定 \(y\) 而言，序列 \(m\mapsto Z_m(y)\) 的全部指数模态由四次特征方程
\[
\Pi(\lambda,y)=0,
\qquad
\Pi(\lambda,y):=\lambda^4-\lambda^3-(2y+1)\lambda^2+\lambda+y(y+1)
\]
完全控制。
\end{conclusion}
\begin{proof}[证明要点]
由定义 \(Z_m(y)\) 是长度 \(m\) 的微态路径权重和，其中每个微态 \(a\in\Omega_m\) 的权重仅通过其稳定输出 \(\Fold_m(a)\) 的 \(1\)-计数进入。
由于 \(\Fold_m\) 可由有限状态规约器实现（全文对 Zeckendorf/窗口折叠的固定口径），
\(Z_m(y)\) 可写为某个固定维数带权转移矩阵 \(T(y)\) 的矩阵系数；
从而 \(\sum_{m\ge 0}Z_m(y)t^m\) 为有理函数，分母次数等于最小状态数。
对本 Fold 实例，约分后的分母次数稳定为 \(4\)，从而得到四阶递推；
将 \(t\)-生成函数在 \(t=0\) 处展开并与初值匹配，即得到所示闭式分子与分母；
上述恒等式与后续的判别式/统计常数可由脚本 \texttt{scripts/exp\_fold\_zm\_bivariate\_partition\_audit.py} 一键复核。证毕。
\end{proof}

\begin{conclusion}[\texorpdfstring{$Z_m(y)$}{Zm(y)} 的次数、首项系数与分歧参数处的闭式特值]\label{con:xi-terminal-zm-degree-leading-special-values}
沿用结论 \ref{con:xi-terminal-zm-order4-rational} 的记号。
则对所有 \(m\ge 0\)：
\begin{enumerate}
  \item 次数精确为
  \[
  \deg_y Z_m(y)=\left\lceil\frac m2\right\rceil.
  \]
  \item 设 \(m=2k+1\) 为奇数，则 \(Z_{2k+1}(y)\) 为首一多项式；设 \(m=2k\) 为偶数，则其首项系数为
  \[
  [y^{k}]\,Z_{2k}(y)=1+\frac{k(k+1)}{2}.
  \]
  \item 在两条有理分歧参数处有闭式特值
  \[
  Z_m(0)=\left\lfloor\frac m2\right\rfloor+1,\qquad
  Z_m(1)=2^m.
  \]
  \item 若将 \(y\) 特化到椭圆轨道 \(y(nP)\in\QQ\)（其中 \(\mathrm{den}(y(nP))=v_n^4\)，见结论 \ref{con:xi-terminal-zm-elliptic-rational-points-denominator-law}），则对任意 \(m,n\ge 0\) 有
  \[
  Z_m\bigl(y(nP)\bigr)\in v_n^{-4\lceil m/2\rceil}\,\ZZ.
  \]
\end{enumerate}
\end{conclusion}
\begin{proof}[证明要点]
由 \(X_m\) 的语法约束 \(x_ix_{i+1}=0\) 可知 \(\max_{x\in X_m}|x|_1=\lceil m/2\rceil\)，且存在取等词；
又 \(Z_m(y)=\sum_{x\in X_m}d_m(x)y^{|x|_1}\) 且 \(d_m(x)\in\ZZ_{>0}\)，故次数结论成立。
首项系数由最大重量层的谱与纤维闭式给出（结论 \ref{con:xi-terminal-max-weight-layer-fiber-spectrum-closed}，亦即结论 \ref{con:xi-terminal-zq-top-coefficient-faulhaber} 在 \(q=1\) 的特例）。
\(Z_m(0)=d_m(0^m)=\lfloor(m+2)/2\rfloor=\lfloor m/2\rfloor+1\)（结论 \ref{con:xi-terminal-zero-weight-fiber-closed}），而 \(Z_m(1)=\sum_x d_m(x)=|\Omega_m|=2^m\)。
最后一条由 \(Z_m\in\ZZ[y]\) 且 \(\deg_y Z_m=\lceil m/2\rceil\)，结合 \(\mathrm{den}(y(nP))=v_n^4\) 直接推出。证毕。
\end{proof}

\begin{conclusion}[有限窗 Sturm 实根交错：\texorpdfstring{$m\le 100$}{m<=100} 时 \texorpdfstring{$Z_m(y)$}{Z_m(y)} 的负实根性]\label{con:xi-terminal-zm-sturm-real-roots-interlacing-m100}
对每个整数 \(1\le m\le 100\)，多项式 \(Z_m(y)\in\ZZ[y]\) 具有 \(\lceil m/2\rceil\) 个互异实根，且全部位于 \((-\infty,0)\)。
并且对每个 \(0\le m\le 98\)，\(Z_m\) 与 \(Z_{m+2}\) 的零点集合严格交错；因此固定奇偶类的两条子序列 \(\{Z_{2k}\}_{k\ge 0}\) 与 \(\{Z_{2k+1}\}_{k\ge 0}\)（在该核验窗内）各自形成一条 Sturm 实根链。
\par
上述核验可由脚本 \texttt{scripts/exp\_fold\_zm\_sturm\_real\_roots\_audit.py} 复算，其输出证书为 \texttt{artifacts/export/fold\_zm\_sturm\_real\_roots\_audit.json}。
\end{conclusion}
\begin{proof}[证明要点]
由结论 \ref{con:xi-terminal-zm-order4-rational} 的四阶递推可在 \(\ZZ[y]\) 中逐步生成 \(Z_m(y)\)。
脚本对核验窗内每个 \(Z_m\) 作高精度求根并检查：
（i）全部根的虚部在数值容许阈值内消失，从而根全为实数；
（ii）由于系数非负且常数项非零，实根必落在 \((-\infty,0)\)，并进一步核验互异性；
（iii）对同奇偶类的 \(m\mapsto m+2\) 根集逐项检验严格交错不等式链。证毕。
\end{proof}

\begin{conjecture}[负实轴完全性与全尺度交错]\label{conj:xi-terminal-zm-sturm-real-roots-interlacing}
对任意 \(m\ge 1\)，多项式 \(Z_m(y)\) 的零点全部为互异负实根，并且对任意 \(m\ge 0\)，\(Z_m\) 与 \(Z_{m+2}\) 的零点集合严格交错。
\end{conjecture}

\begin{remark}[一种结构性证明路线]\label{rem:xi-terminal-zm-sturm-proof-route}
记系数表 \(a_{m,k}:=[y^k]Z_m(y)\)。
若猜想 \ref{conj:xi-terminal-zm-sturm-real-roots-interlacing} 成立，则对每个 \(m\) 系数列 \(\bigl([y^k]Z_m(y)\bigr)_k\) 严格对数凹，因而严格单峰；若进一步可提升为 P\'olya--frequency（PF）序列，则全体阶 Hankel 行列式非负并蕴含多重总正性。结合结论 \ref{con:xi-terminal-zm-double-root-diffusive-sinc-limit} 与结论 \ref{con:xi-terminal-zm-diffusive-edge-weyl-spectral-determinant} 的边缘扩散谱刚性，可把“近零量子化”与“全局负实轴束缚”视为同一 Sturm/全正性机制在边缘与全域的两种投影。
本体系中已核验的二维局域递推
\[
a_{m,k}=a_{m-1,k}+a_{m-2,k}+2a_{m-2,k-1}-a_{m-3,k}-a_{m-4,k-1}-a_{m-4,k-2}
\]
为构造上述网络/矩阵实现提供直接代数入口。
\end{remark}

\begin{conclusion}[最小状态复杂度与退化点：除 \(y\in\{0,-1,1\}\) 外 Hankel 秩恰为 \(4\)]\label{con:xi-terminal-zm-hankel-rank4}
沿用结论 \ref{con:xi-terminal-zm-order4-rational} 的记号，并记
\(
\Pi(\lambda,y)=\lambda^4-\lambda^3-(2y+1)\lambda^2+\lambda+y(y+1)
\)。
若 \(y\in\CC\setminus\{0,-1,1\}\)，则
序列 \(\{Z_m(y)\}_{m\ge 0}\) 的最小湮灭多项式恰为 \(\Pi(\lambda,y)\)；
等价地，\(\{Z_m(y)\}\) 的 Hankel 秩为 \(4\)，不存在三阶及以下的常系数递推能够对所有 \(m\) 成立。
\par
当 \(y\in\{0,-1\}\) 时，系数 \(y(y+1)=0\) 使四阶递推在结构上退化为三阶递推，最小阶严格下降；
当 \(y=1\) 时，归一化恒等式 \(Z_m(1)=\sum_x d_m(x)=2^m\) 使生成函数发生完全抵消，从而 Hankel 秩塌缩为 \(1\)。
\end{conclusion}
\begin{proof}[证明要点]
由结论 \ref{con:xi-terminal-zm-order4-rational}，\(\sum_{m\ge 0}Z_m(y)t^m\) 的分母为
\(
D(t)=1-t-(2y+1)t^2+t^3+y(y+1)t^4
\)。
当 \(y\in\{0,-1\}\) 时，\(D(t)\) 的四次项消失，递推阶数相应降低。
当 \(y=1\) 时，分子与分母共享因子 \((1-t)(1+t)^2\)，从而生成函数约化为 \(\sum_{m\ge 0}2^m t^m=(1-2t)^{-1}\)，最小阶降为 \(1\)。
对任意 \(y\in\CC\setminus\{0,-1,1\}\)，\(\deg_t D=4\) 且分子分母无非平凡公因子，
故约分后分母次数仍为 \(4\)，从而最小递推阶数为 \(4\)；Hankel 秩表述与之等价。证毕。
\end{proof}

% (User input window)

\begin{conclusion}[\texorpdfstring{$y=0$}{y=0} 的扩散 Lee--Yang 边缘：\texorpdfstring{$\operatorname{sinc}$}{sinc} 解析核与近零根量子化]\label{con:xi-terminal-zm-double-root-diffusive-sinc-limit}
沿用结论 \ref{con:xi-terminal-zm-order4-rational} 的多项式族 \(Z_m(y)\)。
等价地，\(\{Z_m(y)\}_{m\ge 0}\) 可由初值
\[
Z_0(y)=1,\quad
Z_1(y)=1+y,\quad
Z_2(y)=2+2y,\quad
Z_3(y)=2+5y+y^2
\]
及四阶递推
\[
Z_m(y)
=Z_{m-1}(y)+(2y+1)Z_{m-2}(y)-Z_{m-3}(y)-y(y+1)Z_{m-4}(y)
\qquad(m\ge 4)
\]
唯一确定。
记
\[
N_m:=Z_m(0)=\left\lfloor\frac m2\right\rfloor+1
\]
（见结论 \ref{con:xi-terminal-zm-degree-leading-special-values}）。
定义
\[
\operatorname{sinc}(z):=
\begin{cases}
\dfrac{\sin z}{z}, & z\neq 0,\\[1mm]
1, & z=0.
\end{cases}
\qquad
S(u):=\operatorname{sinc}\!\left(\frac{\sqrt u}{2}\right).
\]
则 \(S\) 由其幂级数展开唯一确定，因而作为 \(u\) 的整函数无需选择平方根分支。
并且对任意紧集 \(K\subset\CC\)，当 \(m\to\infty\) 时成立一致估计
\[
\boxed{\;
\sup_{u\in K}\left\lvert
\frac{Z_m\!\left(-\frac{u}{m^{2}}\right)}{N_m}-S(u)
\right\rvert=O\!\left(\frac1m\right)\;}
\]
从而 \(y=0\) 处的二重特征根退化在 Lee--Yang 方向的扩散标度 \(y=-u/m^2\) 下由显式整函数核 \(S(u)\) 统一控制。
\par
由于 \(N_m=\frac m2+O(1)\)，上述一致估计亦可改写为：对任意紧集 \(K\subset\CC\)，有
\[
\sup_{u\in K}\left\lvert
\frac{1}{m}Z_m\!\left(-\frac{u}{m^{2}}\right)-\frac12 S(u)
\right\rvert=O\!\left(\frac1m\right)
\qquad(m\to\infty),
\]
其中
\(
\frac12 S(u)=\sum_{k\ge 0}\frac{(-u)^k}{2^{2k+1}(2k+1)!}
=\frac{\sin(\sqrt{u}/2)}{\sqrt{u}}
\)
以连续延拓在 \(u=0\) 处取值 \(1/2\)。

特别地，令 \(j\ge 1\) 固定，并记 \(\rho_{m,1}>\rho_{m,2}>\cdots\) 为 \(Z_m\) 在 \((-\infty,0)\) 上按从 \(0\) 向左排序的零点（若存在多重根则按重数计）。
则对每个固定 \(j\)，当 \(m\) 足够大时 \(\rho_{m,j}\) 存在且满足量子化定标
\[
\boxed{\;
m^{2}\bigl(-\rho_{m,j}\bigr)\longrightarrow 4\pi^{2}j^{2},
\qquad
\rho_{m,j}=-\frac{4\pi^{2}j^{2}}{m^{2}}+O\!\left(\frac1{m^{3}}\right)\;}
\qquad(m\to\infty),
\]
并且对任意 \(\varepsilon>0\)，当 \(m\) 充分大时，\(Z_m\) 在圆盘
\[
\left|y+\frac{4\pi^2 j^2}{m^2}\right|<\frac{\varepsilon}{m^2}
\]
内恰有一个零点（计重数）。
等价地，定义 \(\sigma_{m,j}:=\frac{m}{2}\sqrt{-\rho_{m,j}}\)，则 \(\sigma_{m,j}\to \pi j\)，体现出边缘零点在 \(\sqrt{-y}\) 坐标下的等间距“时钟过程”主阶。
\end{conclusion}
\begin{proof}[证明要点]
由结论 \ref{con:xi-terminal-zm-order4-rational} 的有理生成函数，
在 \(y=0\) 处分母因式分解为 \(D(t,0)=(1-t)^2(1+t)\)，故 \(t=1\) 为二重极点，
对应于递推特征方程在 \(\lambda=1\) 的二重特征根退化。
当 \(y=-u/m^2\) 时，与该二重极点相邻的两条简单极点对以间距 \(O(m^{-1})\) 分裂；
对其作局部留数展开并利用极限
\(
(1\pm i\alpha/m)^{-m}\to e^{\mp i\alpha}
\)
得到主项为 \(\operatorname{sinc}(\sqrt{u}/2)\) 的尺度极限；
一致性来自于对 \(u\) 在紧集上的统一估计。
归一化 \(Z_m(-u/m^2)/N_m\) 与 \((2/m)Z_m(-u/m^2)\) 的差异由 \(N_m=\frac m2+O(1)\) 吸收，并给出所示 \(O(m^{-1})\) 余项。
零点定标与局部唯一性由 Hurwitz 定理对上述一致估计与 \(S(u)\) 的简单零点 \(u=4\pi^{2}j^{2}\) 应用而得；
进一步由零点的非退化性与余项 \(O(m^{-1})\) 推出 \(u\)-坐标的偏移为 \(O(m^{-1})\)，从而 \(\rho_{m,j}=-u_{m,j}/m^2\) 的偏移为 \(O(m^{-3})\)。
可复现审计脚本：\texttt{scripts/exp\_fold\_zm\_double\_root\_diffusive\_sinc\_limit\_audit.py}。
\end{proof}

\begin{conclusion}[边缘零点计数的 Weyl 型律与谱行列式识别]\label{con:xi-terminal-zm-diffusive-edge-weyl-spectral-determinant}
在结论 \ref{con:xi-terminal-zm-double-root-diffusive-sinc-limit} 的记号下，对任意固定 \(U>0\) 定义边缘零点计数函数
\[
\mathcal N_m(U):=\#\Bigl\{j:\ \rho_{m,j}\in\bigl[-U/m^2,\,0\bigr)\Bigr\}.
\]
则当 \(m\to\infty\) 时有
\[
\boxed{\;
\mathcal N_m(U)=\left\lfloor \sqrt{\frac{U}{4\pi^2}}\right\rfloor+o(1)\;}
\]
从而在长度尺度 \(m^{-2}\) 的窗口内，零点计数呈 \(\sqrt{U}\) 增长的 Weyl 型指数。
此外，极限核 \(S(u)=\operatorname{sinc}(\sqrt{u}/2)\) 具有 Weierstrass 乘积分解
\[
\boxed{\;
S(u)=\prod_{j=1}^{\infty}\left(1-\frac{u}{4\pi^2 j^2}\right)\;}
\]
其零点集合 \(\{4\pi^2 j^2\}_{j\ge 1}\) 与单位区间 Dirichlet 拉普拉斯算子 \(-\frac{\dd^2}{\dd x^2}\) 的本征值 \(\{\pi^2 j^2\}_{j\ge 1}\) 仅差一个固定尺度因子 \(4\)。
因此，扩散标度下的归一化极限可视为一个显式谱行列式对象，从而把离散折叠—投影体系在 \(y=0\) 边缘的零点刚性与连续二阶算子谱结构对齐。
\end{conclusion}
\begin{proof}[证明要点]
由结论 \ref{con:xi-terminal-zm-double-root-diffusive-sinc-limit} 的零点定标，固定 \(U\) 时 \(\mathcal N_m(U)\) 收敛到 \(S(u)\) 在区间 \([0,U]\) 内零点数目，而 \(S(u)=0\) 当且仅当 \(u=4\pi^2 j^2\)，故得到所示计数律。
Weierstrass 乘积由经典恒等式 \(\sin z/z=\prod_{j\ge 1}(1-z^2/(\pi^2 j^2))\) 取 \(z=\sqrt{u}/2\) 直接推出；谱字典来自 Dirichlet 拉普拉斯算子的本征值闭式 \(\pi^2 j^2\)。证毕。
\end{proof}

\begin{conclusion}[\texorpdfstring{$y=0$}{y=0} 的临界窗宽为 \texorpdfstring{$m^{-2}$}{m^{-2}}：线性增长与指数增长之间的唯一尺度]\label{con:xi-terminal-zm-critical-window-m2}
在 \(y=0\) 处，结论 \ref{con:xi-terminal-zm-hankel-rank4} 所述二重特征根使 \(Z_m(0)\) 仅呈线性增长；
而对任意固定 \(y>0\)，谱半径严格大于 \(1\)，故 \(Z_m(y)\) 具有严格指数增长率。
结论 \ref{con:xi-terminal-zm-double-root-diffusive-sinc-limit} 的扩散标度极限表明，
两种增长机制之间的非平凡过渡被刚性锁定在唯一临界窗
\[
y=-\frac{u}{m^{2}},
\qquad
u\asymp 1,
\]
并在该窗内由整函数 \(u\mapsto S(u)\) 统一控制。
相应地，靠近 \(0\) 的零点簇在该窗内呈现 \(m^{-2}\) 级别的量子化定标
（结论 \ref{con:xi-terminal-zm-double-root-diffusive-sinc-limit}），
其主阶间距由 \(\sin\) 的整数倍零点完全决定。
\end{conclusion}



\begin{conclusion}[谱分歧轨迹的完全因式分解：Lee--Yang 边界由一元三次给出]\label{con:xi-terminal-zm-discriminant-leyang}
视 \(\Pi(\lambda,y)\) 为关于 \(\lambda\) 的四次多项式，则其判别式在 \(\ZZ[y]\) 中精确分解为
\[
\boxed{\;
\mathrm{Disc}_{\lambda}\bigl(\Pi(\lambda,y)\bigr)
=-y\,(y-1)\,\bigl(256y^3+411y^2+165y+32\bigr)\; }.
\]
为后续引用，记
\(
\mathrm{disc}^{\ast}_{\lambda}\Pi:=\mathrm{Disc}_{\lambda}\bigl(\Pi(\lambda,y)\bigr)
\)。
在该记号下，Lee--Yang 型边界即为 \(\mathrm{disc}^{\ast}_{\lambda}\Pi\) 在 \(\mathbb P^1_y\) 上的零点集合。
因此，\(\lambda\)-谱作为 \(y\) 的代数函数仅在
\(
y\in\{0,1\}
\)
与三次方程
\(
256y^3+411y^2+165y+32=0
\)
的根处发生分歧。
其中该三次多项式恰有一个实根
\[
y_{\mathrm{LY}}\approx-1.1344520030,
\]
并给出双变量自由能
\[
f(y):=\lim_{m\to\infty}\frac{1}{m}\log Z_m(y)
\]
在实轴上的最近边界奇点；从而 \(y\) 作为复变量时，\(f\) 的最小非平凡解析半径由 \(y_{\mathrm{LY}}\) 所决定（Lee--Yang 型边界机制；参见 \cite{XuZamolodchikov2022IsingYangLee}）。
\end{conclusion}
\begin{proof}[证明要点]
对四次多项式 \(\Pi(\lambda,y)\) 计算 \(\lambda\)-判别式并在 \(\ZZ[y]\) 中因式分解即可得到所示公式。
判别式为零当且仅当存在重根，亦即代数分支在该参数处发生分歧；三次因子的实根可由标准数值求根获得。证毕。
\end{proof}

\subsubsection{Lee--Yang 零点的射影消混合：交比不变量与 crossing 渐近}\label{subsubsec:xi-leyang-projective-demixing-crossing}
在有限体积配分函数的 Lee--Yang 零点分析中，需要区分两类互补机制。
一类来自有限维谱分解的主导谱竞争与判别式分歧：当两条主导谱支在等模曲线上发生竞争时，零点在法向方向出现 \(\frac{\log n}{n}\) 量级的定量位移；当谱曲线在临界点具有分歧指标 \(m\) 时，零点在尺度 \(n^{-m}\) 上聚束到该分歧点并呈 Puiseux 型准晶格。
另一类来自解析背景与变量混合：允许的复仿射重参数化 \(z\mapsto az+b(\tau)\) 会引入共同平移与整体缩放，使得直接以零点坐标作比值或作体积交叉的定位量对混合参数敏感。
以下首先记录可直接用于“判别式分解 \(\Rightarrow\) 零点标度反演”的谱侧定量律，随后给出射影化的消混合接口。

\begin{theorem}[两相并存临界窗的尺度极限与零点凝聚]\label{thm:xi-godel-leyang-two-phase-coexistence-window}
设 \(\mathcal C\) 为可数概念类，并给定前缀自由码长 \(\ell:\mathcal C\to\NN\)，权重 \(w(C):=2^{-\ell(C)}\)。
对观测前缀长度 \(n\) 与参数 \(\beta\in\CC\)，定义有限体积配分函数
\[
Z_n(\beta):=\sum_{C\in\mathcal C}w(C)\exp\!\bigl(-\beta E_C(n)\bigr),
\]
其中 \(E_C(n)\in\RR\) 为可加失配能量。
假设存在两概念 \(C_\pm\in\mathcal C\) 与常数 \(e_\pm,\lambda_\pm\in\RR\) 使
\[
\frac{E_{C_\pm}(n)}{n}\to e_\pm,\qquad \frac{\ell(C_\pm)}{n}\to \lambda_\pm,
\qquad
\Delta e:=e_+-e_-\neq 0,
\]
并且存在实数 \(\beta^\star\) 使两相自由能并存且严格主导：对任意其余概念 \(C\notin\{C_+,C_-\}\)，若极限
\(
e(C):=\lim_{n\to\infty}E_C(n)/n
\) 与
\(
\lambda(C):=\lim_{n\to\infty}\ell(C)/n
\)
存在，则成立严格不等式
\[
(\ln 2)\lambda_+ + \beta^\star e_+
=
(\ln 2)\lambda_- + \beta^\star e_-
\ <\ (\ln 2)\lambda(C)+\beta^\star e(C).
\]
对 \(z\in\CC^\times\)，定义临界窗重标度
\[
\beta_{n}(z):=\beta^\star-\frac{1}{n\Delta e}\log z,
\]
并定义归一化局部配分函数
\[
\mathcal Z_n(z):=
\frac{Z_n\!\bigl(\beta_n(z)\bigr)}
{w(C_-)\exp\!\bigl(-\beta_n(z)E_{C_-}(n)\bigr)}.
\]
则对任意紧集 \(K\Subset\CC\setminus\{0\}\)，有一致收敛
\[
\mathcal Z_n(z)\ \longrightarrow\ 1+z
\qquad(n\to\infty,\ z\in K).
\]
因此在该尺度下，\(\mathcal Z_n\) 的零点在 \(K\) 上趋于 \(\{-1\}\)；并且最靠近 \(\beta^\star\) 的 Lee--Yang 零点满足
\[
\beta_{n,k}
\;=\;
\beta^\star-\frac{(2k+1)\pi i}{n\Delta e}+o(n^{-1}),
\qquad k\in\ZZ.
\]
\end{theorem}
\begin{proof}
将 \(Z_n\) 分解为两主导项与余项：
\[
Z_n(\beta)=w(C_-)\,e^{-\beta E_{C_-}(n)}+w(C_+)\,e^{-\beta E_{C_+}(n)}+R_n(\beta).
\]
把 \(\beta=\beta_n(z)\) 代入并按 \(w(C_-)\,e^{-\beta E_{C_-}(n)}\) 归一化，得到
\[
\mathcal Z_n(z)=1+\frac{w(C_+)}{w(C_-)}\exp\!\bigl(-\beta_n(z)\,(E_{C_+}(n)-E_{C_-}(n))\bigr)
\;+\;\frac{R_n(\beta_n(z))}{w(C_-)\,e^{-\beta_n(z)E_{C_-}(n)}}.
\]
由 \(w(C_\pm)=2^{-\ell(C_\pm)}\) 与假设给出的并存等式，可得
\[
\frac{w(C_+)}{w(C_-)}\exp\!\bigl(-\beta^\star(E_{C_+}(n)-E_{C_-}(n))\bigr)=1+o(1).
\]
并且由 \(\beta_n(z)=\beta^\star-\frac{1}{n\Delta e}\log z\) 与 \((E_{C_+}(n)-E_{C_-}(n))/n\to \Delta e\)，有
\[
\exp\!\left(+\frac{\log z}{n\Delta e}\,(E_{C_+}(n)-E_{C_-}(n))\right)=z^{1+o(1)}.
\]
合并得到第二项趋于 \(z\)，从而给出 \(1+z\) 的主极限。
余项 \(\frac{R_n(\beta_n(z))}{w(C_-)\,e^{-\beta_n(z)E_{C_-}(n)}}\) 由严格主导条件在 \(K\) 上以指数速率趋于 \(0\)，从而得到一致收敛。
零点渐近式由 \(1+z=0\) 与 \(z=\exp(-n\Delta e(\beta-\beta^\star))\) 反解，并把误差吸收进 \(o(n^{-1})\) 得到。证毕。
\end{proof}

\begin{theorem}[两谱竞争下零点的对数位移定量律]\label{thm:xi-leyang-logshift-two-spectrum}
设 \(U\subset\CC\) 为开集，\(\lambda_1,\dots,\lambda_s:U\to\CC\setminus\{0\}\) 为解析函数。
对 \(n\in\NN\) 定义
\[
P_n(z):=\sum_{j=1}^{s}\alpha_j(n;z)\,\lambda_j(z)^n,
\qquad
\alpha_j(n;z)=\sum_{r=0}^{d_j}p_{j,r}(z)\,n^r,
\]
其中各 \(p_{j,r}\) 解析且 \(p_{j,d_j}\not\equiv 0\)。
定义严格两谱等模集
\[
\Gamma:=\Bigl\{z\in U:\ |\lambda_1(z)|=|\lambda_2(z)|>\max_{3\le j\le s}|\lambda_j(z)|\Bigr\},
\]
并取紧集 \(K\Subset \Gamma\)。
假设存在邻域 \(V\Subset U\) 与常数 \(0<\theta<1\) 使得对一切 \(z\in V\) 有
\[
\max_{3\le j\le s}\frac{|\lambda_j(z)|}{|\lambda_1(z)|}\le \theta,
\]
并且 \(\lambda_1/\lambda_2\) 在 \(V\) 上无零点，从而可取解析对数分支
\[
\mathcal{L}(z):=\log\!\Bigl(\frac{\lambda_1(z)}{\lambda_2(z)}\Bigr),
\qquad
\mathcal{L}'(z)\neq 0\ \ (z\in K).
\]
再设 \(p_{1,d_1}(z)p_{2,d_2}(z)\neq 0\) 于 \(K\) 上成立，并记
\[
\beta_n(z):=\frac{\alpha_2(n;z)}{\alpha_1(n;z)},
\qquad
D:=\max_{1\le j\le s}d_j.
\]
则存在邻域 \(W\Subset V\)（\(K\subset W\)）与 \(n_0\in\NN\) 使得对所有 \(n\ge n_0\)：
\begin{enumerate}
  \item \(P_n\) 在 \(W\) 内的全部零点皆为单零点；
  \item 若 \(z\in W\) 为零点，则存在整数 \(k\in\ZZ\) 使
  \begin{equation}\label{eq:xi-leyang-logshift-quant}
  \boxed{\;
  n\,\mathcal{L}(z)=\log\!\bigl(-\beta_n(z)\bigr)+2\pi \mathrm{i}k+O\!\bigl(n^{D}\theta^n\bigr)\;}
  \end{equation}
  其中对数取任意在 \(W\) 上连续的分支（分支变更等价于 \(k\) 的平移）。
\end{enumerate}
并且若 \(z=z_{n,k}\) 取遍上述零点，则有一致渐近展开
\[
\Re \mathcal{L}(z_{n,k})
=\frac{(d_2-d_1)\log n+\log\left|\frac{p_{2,d_2}(z_{n,k})}{p_{1,d_1}(z_{n,k})}\right|}{n}
+O\!\left(\frac{\log n}{n^2}\right),
\]
\[
\Im \mathcal{L}(z_{n,k})
=\frac{2\pi k+\arg\!\left(-\frac{p_{2,d_2}(z_{n,k})}{p_{1,d_1}(z_{n,k})}\right)}{n}
+O\!\left(\frac{\log n}{n^2}\right).
\]
由于 \(K\) 上 \(\mathcal{L}'\) 不消失，\(\Re\mathcal{L}=0\) 在 \(K\) 附近给出 \(\Gamma\) 的一条光滑实解析支；
从而零点到该曲线的欧氏距离满足
\[
\mathrm{dist}(z_{n,k},\Gamma)=\Theta\!\left(\frac{\log n}{n}\right),
\]
其中常数仅依赖于 \(K\) 上的 \(\inf|\mathcal{L}'|\) 与 \(|d_2-d_1|\)。
\end{theorem}
\begin{proof}
在 \(V\) 上写
\[
P_n(z)=\alpha_1(n;z)\lambda_1(z)^n\Bigl(1+\beta_n(z)e^{-n\mathcal{L}(z)}+R_n(z)\Bigr),
\]
其中
\[
R_n(z):=\sum_{j=3}^{s}\frac{\alpha_j(n;z)}{\alpha_1(n;z)}\Bigl(\frac{\lambda_j(z)}{\lambda_1(z)}\Bigr)^n.
\]
由 \(\alpha_j(n;z)=O(n^{d_j})\) 与谱分离假设可得在 \(V\) 上一致估计
\(
\sup_{z\in V}|R_n(z)|=O(n^{D}\theta^n)
\)。
取 \(W\Subset V\) 使 \(\alpha_1,\alpha_2\) 在 \(W\) 上对充分大 \(n\) 不为零，且 \(\mathcal{L}'\) 在 \(W\) 上无零点。
\par
取 \(z\in W\) 满足 \(P_n(z)=0\)。
由于 \(\alpha_1(n;z)\lambda_1(z)^n\neq 0\) 于 \(W\) 成立，上式等价于
\[
1+\beta_n(z)e^{-n\mathcal{L}(z)}=-R_n(z).
\]
由 \(\sup_W|R_n|=O(n^{D}\theta^n)\) 得
\[
\beta_n(z)e^{-n\mathcal{L}(z)}=-1+O(n^{D}\theta^n),
\qquad
e^{n\mathcal{L}(z)}=-\beta_n(z)\bigl(1+O(n^{D}\theta^n)\bigr)^{-1}.
\]
对两侧取对数并吸收分支不定性，得到存在 \(k\in\ZZ\) 使 \eqref{eq:xi-leyang-logshift-quant} 成立。
\par
为证单零点性，记
\[
H_n(z):=1+\beta_n(z)e^{-n\mathcal{L}(z)}+R_n(z),
\qquad
P_n(z)=\alpha_1(n;z)\lambda_1(z)^n\,H_n(z).
\]
对 \(H_n\) 求导并在零点处代入 \(\beta_n e^{-n\mathcal{L}}=-1+o(1)\)，同时注意 \(\beta_n\) 及其导数至多多项式增长，而 \(e^{-n\mathcal{L}}\) 在零点处与 \(\beta_n^{-1}\) 同阶，可得
\[
H_n'(z)=\beta_n'(z)e^{-n\mathcal{L}(z)}-n\,\beta_n(z)\mathcal{L}'(z)e^{-n\mathcal{L}(z)}+R_n'(z)
=n\,\mathcal{L}'(z)+O(1),
\]
其中余项在 \(W\) 上一致；当 \(n\) 足够大时 \(H_n'(z)\neq 0\)，故零点为单零点。
\par
由多项式前因子展开
\[
\beta_n(z)=n^{d_2-d_1}\frac{p_{2,d_2}(z)}{p_{1,d_1}(z)}\Bigl(1+O\!\left(\frac1n\right)\Bigr)
\]
（一致于 \(W\)），代入 \eqref{eq:xi-leyang-logshift-quant} 并分离实部与虚部即得所示渐近展开。
距离估计由 \(\Re\mathcal{L}\) 在 \(K\) 上的横截性与 \(\Re\mathcal{L}(z_{n,k})=\Theta((\log n)/n)\) 推出。证毕。
\end{proof}

\begin{theorem}[Jordan 调制导致的对数漂移系数]\label{thm:xi-leyang-jordan-logshift}
设 \(T:U\to\Mat_d(\CC)\) 为解析矩阵族，\(u,v\in\CC^d\) 固定，定义
\[
Z_n(z):=u^{\mathsf T}T(z)^n v\qquad(n\in\NN).
\]
设存在邻域 \(V\Subset U\) 与常数 \(0<\theta<1\)，使得 \(T(z)\) 的谱可分解为两条主导谱分支 \(\lambda_1(z),\lambda_2(z)\) 与其余谱满足
\[
\max_{j\ge 3}|\lambda_j(z)|\le \theta\,\max\bigl(|\lambda_1(z)|,|\lambda_2(z)|\bigr)\qquad(z\in V),
\]
并且 \(\lambda_1/\lambda_2\) 在 \(V\) 上可取解析对数 \(\mathcal{L}=\log(\lambda_1/\lambda_2)\)，且 \(\mathcal{L}'\) 在局部等模集上不消失。
设 \(\ell_i\ge 1\) 为与 \(\lambda_i\) 相关、且在读出 \(u^{\mathsf T}(\cdot)v\) 中出现的最大 Jordan 块尺寸，并在 \(V\) 上常值。
则在 \(V\) 的适当缩小邻域内，\(Z_n\) 的零点在等模曲线附近满足
\[
\Re\mathcal{L}(z)=\frac{(\ell_2-\ell_1)\log n+\log|c_2(z)/c_1(z)|}{n}+O\!\left(\frac1n\right),
\]
其中 \(c_1,c_2\) 为解析且在等模集上不为零的系数函数。
特别地，若 \(\ell_1\neq \ell_2\)，则 \(\frac{\ell_2-\ell_1}{n}\log n\) 项在零点法向漂移中不可消去，其系数仅由 Jordan 非半单缺陷所决定。
\end{theorem}
\begin{proof}
对每个 \(z\in V\) 取 \(T(z)=S(z)J(z)S(z)^{-1}\) 的 Jordan 标准形。
对单个 Jordan 块 \(J=\lambda(I+N)\)（\(N\) 幂零、块大小 \(\ell\)）有恒等式
\[
J^n=\lambda^n\sum_{r=0}^{\ell-1}\binom{n}{r}N^r,
\]
从而 \(Z_n(z)\) 展开为有限个项的线性组合，每项形如 \(n^r\lambda_j(z)^n\)，系数解析。
按可见最大 Jordan 阶数取出 \(\lambda_1,\lambda_2\) 的首项，得到
\[
Z_n(z)=\bigl(c_1(z)\,n^{\ell_1-1}+O(n^{\ell_1-2})\bigr)\lambda_1(z)^n
\;+\;\bigl(c_2(z)\,n^{\ell_2-1}+O(n^{\ell_2-2})\bigr)\lambda_2(z)^n
\;+\;O(n^{D}\theta^n\Lambda(z)^n),
\]
其中 \(\Lambda(z)=\max(|\lambda_1(z)|,|\lambda_2(z)|)\)。
将该表示代入定理 \ref{thm:xi-leyang-logshift-two-spectrum}（取 \(d_i=\ell_i-1\)）即得结论。证毕。
\end{proof}

\begin{theorem}[谱分歧点的 Puiseux 标度与零点聚束]\label{thm:xi-leyang-puiseux-zero-scaling}
设 \(z_c\in\CC\)。
在 \(z_c\) 的穿孔邻域内设存在两条主导谱支 \(\lambda_{\pm}(z)\neq 0\)，并存在整数 \(m\ge 2\) 与常数 \(a\neq 0\) 使得当 \(z\to z_c\) 时有 Puiseux 展开
\[
\log\!\Bigl(\frac{\lambda_+(z)}{\lambda_-(z)}\Bigr)
=2a\,(z-z_c)^{1/m}+O\!\bigl((z-z_c)^{2/m}\bigr),
\]
其中 \((z-z_c)^{1/m}\) 为某一固定的局部 Puiseux 分支。
设
\[
P_n(z)=A_n(z)\lambda_+(z)^n+B_n(z)\lambda_-(z)^n+R_n(z),
\]
其中 \(A_n,B_n\) 在 \(z_c\) 的邻域内解析且关于 \(n\) 仅含有次数有界的多项式前因子，而 \(R_n\) 为严格次主导余项：存在 \(0<\theta<1\) 与 \(D\ge 0\) 使
\[
|R_n(z)|\le C\,n^D\,\theta^n\,\max(|\lambda_+(z)|,|\lambda_-(z)|)^n
\]
在某邻域内一致成立。
并假设 \(-B_n(z_c)/A_n(z_c)\) 不以指数速率趋于 \(0\) 或 \(\infty\)。
则存在 \(n_0\) 与常数 \(C_1>0\)，使得对一切 \(n\ge n_0\)，落在 \(|z-z_c|\le C_1 n^{-m}\) 内的全部零点可表为
\[
z_{n,k}
=z_c+\Bigl(\frac{1}{2a}\Bigr)^m\Bigl(\frac{\log\!\bigl(-B_n(z_c)/A_n(z_c)\bigr)+2\pi \mathrm{i}k}{n}\Bigr)^m
+O\!\left(\frac{(\log n)^{m+1}}{n^{m+1}}\right),
\qquad k\in\ZZ,
\]
其中对数取任意连续分支（分支变更等价于 \(k\) 的平移）。
特别地，零点对 \(z_c\) 的收敛标度为 \(n^{-m}\)。
\end{theorem}
\begin{proof}
忽略余项 \(R_n\) 时，零点满足
\[
\Bigl(\frac{\lambda_+(z)}{\lambda_-(z)}\Bigr)^n=-\frac{B_n(z)}{A_n(z)}.
\]
取对数并吸收分支不定性得
\[
n\,\log\!\Bigl(\frac{\lambda_+(z)}{\lambda_-(z)}\Bigr)=\log\!\Bigl(-\frac{B_n(z)}{A_n(z)}\Bigr)+2\pi\mathrm{i}k.
\]
由于 \(B_n/A_n\) 在 \(z_c\) 的邻域内解析且非零，其对数在固定分支下满足
\[
\log\!\Bigl(-\frac{B_n(z)}{A_n(z)}\Bigr)
=\log\!\Bigl(-\frac{B_n(z_c)}{A_n(z_c)}\Bigr)+O(|z-z_c|),
\]
代入并注意 \(|z-z_c|=O(n^{-m})\) 可把该误差吸收进最终的 \(O((\log n)^{m+1}/n^{m+1})\) 项。
由 Puiseux 展开存在局部坐标 \(\eta(z)=a(z-z_c)^{1/m}+O((z-z_c)^{2/m})\) 使
\(
\log(\lambda_+/\lambda_-)=2\eta(z)
\)，并且 \(z-z_c=(\eta/a)^m+O(\eta^{m+1})\)。
由于右端仅具对数级增长，可得 \(\eta(z)=O((\log n)/n)\) 并代回得到所示展开。
最后由 \(|R_n|\ll |A_n\lambda_+^n|+|B_n\lambda_-^n|\) 的一致控制与 Rouch\'e 定理可知，加入 \(R_n\) 不改变上述小邻域内零点计数与展开。证毕。
\end{proof}

\begin{corollary}[简单判别式点的二次聚束标度]\label{cor:xi-leyang-puiseux-m2-parabola}
\leanverified{paper\_xi\_leyang\_puiseux\_m2\_parabola}
在定理 \ref{thm:xi-leyang-puiseux-zero-scaling} 的设定下，若分歧指标 \(m=2\)，则临界邻域内零点满足
\[
z_{n,k}
=z_c+\frac{1}{4a^2}\Bigl(\frac{\log\!\bigl(-B_n(z_c)/A_n(z_c)\bigr)+2\pi \mathrm{i}k}{n}\Bigr)^2
+O\!\left(\frac{(\log n)^3}{n^3}\right),
\]
从而典型间距为 \(\Theta(n^{-2})\)，并在 \(z\)-平面上形成以 \(z_c\) 为端点的二次聚束族。
\end{corollary}

\begin{corollary}[Jordan 指数差的零点漂移识别]\label{cor:xi-leyang-identify-jordan-gap-from-zeros}
在定理 \ref{thm:xi-leyang-jordan-logshift} 的设定下，设 \(\gamma\) 为等模曲线的一段紧弧，且
\[
\inf_{z\in\gamma}|\mathcal{L}'(z)|>0.
\]
对任意零点 \(z=z_{n,k}\) 取其到 \(\gamma\) 的最近点投影 \(z_\gamma\in\gamma\)，并以 \(\Re\mathcal{L}\) 的符号规定有符号法向距离 \(\delta_n(z)\)。
则
\[
\delta_n(z)=\frac{\ell_2-\ell_1}{|\mathcal{L}'(z_\gamma)|}\cdot\frac{\log n}{n}+O\!\left(\frac{1}{n}\right),
\]
其中余项在上述投影落入 \(\gamma\) 的紧子弧时一致成立。
从而 Jordan 指数差可由单尺度斜率读出：
\[
\ell_2-\ell_1=\lim_{n\to\infty}\frac{n\,|\mathcal{L}'(z_\gamma)|\,\delta_n(z)}{\log n}.
\]
\end{corollary}

\begin{remark}[与谱判别式的对齐]\label{rem:xi-leyang-logshift-puiseux-discriminant-alignment}
在本文的谱方程口径中，分歧点等价于主导谱支发生重根并合，因而由判别式为零刻画（例如 \(\Pi(\lambda,y)\) 的 \(\mathrm{Disc}_\lambda\) 分解给出 Lee--Yang 边界）。
定理 \ref{thm:xi-leyang-puiseux-zero-scaling} 将该判别式层的分歧指标 \(m\) 直接译码为零点云的收敛标度 \(n^{-m}\)，从而把“分歧几何”提升为“有限体积零点”可核验的量化证书。
\end{remark}

在允许解析背景与变量混合的口径下，零点坐标将对复仿射重参数化敏感。
下述陈述给出一种射影化的消混合接口：以四点交比构造对复仿射重参数化不变的可比量，并由此得到 crossing 定位临界点的普适渐近律；同时给出由两枚有限体积零点反演 \(A(0)\) 与 \(z_0(0)\) 的闭式估计。

\begin{theorem}[Lee--Yang 零点的射影消混合不变量与交叉点渐近]\label{thm:xi-leyang-crossratio-demixing-crossing}
\leanverified{paper\_xi\_leyang\_crossratio\_demixing\_crossing}
设 \(L\in\RR_{>0}\) 为系统线性尺度，\(\tau\in\RR\) 为距临界点的温度型标度变量（或其解析重参数化），\(z\in\CC\) 为待复化的外参。
记有限体积配分函数 \(Z_L(z;\tau)\) 关于 \(z\) 的 Lee--Yang 零点为 \(\{z_j(L,\tau)\}_{j\ge 1}\subset\CC\)（允许按某一固定约定排序）。
假设存在指数 \(y_h>0,y_t>0,\omega>0\)，解析函数 \(z_0(\tau)\) 与 \(A(\tau)\neq 0\)，以及与普适类仅相关的常数序列 \(\{\zeta_j\}_{j\ge 1}\subset\CC\) 与 \(\{\kappa_j\}_{j\ge 1},\{\xi_j\}_{j\ge 1}\subset\CC\)，使得当 \(L\to\infty\) 且 \(|\tau|\) 足够小时，对每个固定 \(j\) 有展开
\begin{equation}\label{eq:xi-leyang-fss-analytic-mixing}
z_j(L,\tau)
=
z_0(\tau)
+
A(\tau)L^{-y_h}\Bigl(
\zeta_j
+
\kappa_j\,\tau L^{y_t}
+
\xi_j\,L^{-\omega}
+
O\!\bigl(\tau^2L^{2y_t}+\tau L^{y_t-\omega}+L^{-2\omega}\bigr)
\Bigr).
\end{equation}
对任意互异指标 \(i,j,k,\ell\) 定义四点交比
\begin{equation}\label{eq:xi-leyang-crossratio-def}
\mathrm{CR}_{i,j;k,\ell}(L,\tau)
:=
\frac{\bigl(z_i-z_j\bigr)\bigl(z_k-z_\ell\bigr)}{\bigl(z_i-z_k\bigr)\bigl(z_j-z_\ell\bigr)}\Bigg|_{(L,\tau)}.
\end{equation}
则成立：
\begin{enumerate}
  \item 存在仅由 \(\{\zeta_a\}\) 决定的普适常数
  \begin{equation}\label{eq:xi-leyang-crossratio-universal}
  \mathrm{CR}_{i,j;k,\ell}^{\ast}
  :=
  \frac{(\zeta_i-\zeta_j)(\zeta_k-\zeta_\ell)}{(\zeta_i-\zeta_k)(\zeta_j-\zeta_\ell)},
  \end{equation}
  使得在临界切片 \(\tau=0\) 有
  \begin{equation}\label{eq:xi-leyang-crossratio-critical-limit}
  \mathrm{CR}_{i,j;k,\ell}(L,0)=\mathrm{CR}_{i,j;k,\ell}^{\ast}+O(L^{-\omega}).
  \end{equation}
  并且 \(\mathrm{CR}_{i,j;k,\ell}\) 对任意复仿射重参数化 \(z\mapsto az+b(\tau)\)（\(a\in\CC^{\times}\)，\(b\) 解析）严格不变。
  \item 固定 \(s>1\)，并令 \(\tau_{L,s}\) 为方程
  \begin{equation}\label{eq:xi-leyang-crossratio-crossing-eq}
  \mathrm{CR}_{i,j;k,\ell}(L,\tau)=\mathrm{CR}_{i,j;k,\ell}(sL,\tau)
  \end{equation}
  在 \(\tau=0\) 的小邻域内的解（当 \(L\) 充分大且非退化条件满足时唯一）。
  则存在常数 \(C_{i,j;k,\ell}(s)\in\CC\) 使得
  \begin{equation}\label{eq:xi-leyang-crossratio-crossing-asymp}
  \tau_{L,s}
  =
  C_{i,j;k,\ell}(s)\,L^{-(y_t+\omega)}+o\!\bigl(L^{-(y_t+\omega)}\bigr).
  \end{equation}
  更具体地，令
  \begin{equation}\label{eq:xi-leyang-crossratio-BC}
  \mathcal B:=
  \Bigl(\frac{\kappa_i-\kappa_j}{\zeta_i-\zeta_j}+\frac{\kappa_k-\kappa_\ell}{\zeta_k-\zeta_\ell}
  -\frac{\kappa_i-\kappa_k}{\zeta_i-\zeta_k}-\frac{\kappa_j-\kappa_\ell}{\zeta_j-\zeta_\ell}\Bigr),
  \qquad
  \mathcal C:=
  \Bigl(\frac{\xi_i-\xi_j}{\zeta_i-\zeta_j}+\frac{\xi_k-\xi_\ell}{\zeta_k-\zeta_\ell}
  -\frac{\xi_i-\xi_k}{\zeta_i-\zeta_k}-\frac{\xi_j-\xi_\ell}{\zeta_j-\zeta_\ell}\Bigr),
  \end{equation}
  若 \(\mathcal B\neq 0\)，则
  \begin{equation}\label{eq:xi-leyang-crossratio-crossing-constant}
  C_{i,j;k,\ell}(s)=\frac{1-s^{-\omega}}{s^{y_t}-1}\cdot\frac{\mathcal C}{\mathcal B}.
  \end{equation}
\end{enumerate}
\end{theorem}
\begin{proof}
由 \eqref{eq:xi-leyang-fss-analytic-mixing} 得对任意 \(a\neq b\),
\[
z_a-z_b
=
A(\tau)L^{-y_h}\Bigl((\zeta_a-\zeta_b)+(\kappa_a-\kappa_b)\tau L^{y_t}+(\xi_a-\xi_b)L^{-\omega}+O(\cdots)\Bigr),
\]
其中共同平移 \(z_0(\tau)\) 在差分中完全消失，且公共因子 \(A(\tau)L^{-y_h}\) 在交比 \eqref{eq:xi-leyang-crossratio-def} 中整体约去，从而交比对 \(z\mapsto az+b(\tau)\) 严格不变。
令 \(\varepsilon:=\tau L^{y_t}\)、\(\delta:=L^{-\omega}\)。
将 \eqref{eq:xi-leyang-fss-analytic-mixing} 代入 \eqref{eq:xi-leyang-crossratio-def} 并在 \((\varepsilon,\delta)=(0,0)\) 处作一阶展开得到
\[
\mathrm{CR}_{i,j;k,\ell}(L,\tau)
=
\mathrm{CR}_{i,j;k,\ell}^{\ast}
\Bigl(1+\mathcal B\,\varepsilon+\mathcal C\,\delta+O(\varepsilon^{2}+\varepsilon\delta+\delta^{2})\Bigr),
\]
其中 \(\mathcal B,\mathcal C\) 如 \eqref{eq:xi-leyang-crossratio-BC}。
取 \(\tau=0\) 即得 \eqref{eq:xi-leyang-crossratio-critical-limit}。
\par
对固定 \(s>1\)，将同一展开应用于尺度 \(L\) 与 \(sL\) 并代入 \eqref{eq:xi-leyang-crossratio-crossing-eq}，约去 \(\mathrm{CR}^{\ast}\) 后主项为
\[
\mathcal B\,(1-s^{y_t})\,\varepsilon+\mathcal C\,(1-s^{-\omega})\,\delta
+O(\varepsilon^{2}+\varepsilon\delta+\delta^{2})
=0.
\]
当 \(\mathcal B\neq 0\) 且 \(L\) 充分大时，隐函数定理给出唯一解满足 \(\varepsilon=O(\delta)\)，即
\(
\tau=O(L^{-(y_t+\omega)})
\)。
由主项平衡得到
\(
\varepsilon=\frac{1-s^{-\omega}}{s^{y_t}-1}\cdot\frac{\mathcal C}{\mathcal B}\,\delta+o(\delta)
\)，
从而推出 \eqref{eq:xi-leyang-crossratio-crossing-asymp} 与 \eqref{eq:xi-leyang-crossratio-crossing-constant}。
\end{proof}

\begin{corollary}[复混合下比值法可失稳而交比法仍严格不变]\label{cor:xi-leyang-imag-ratio-unstable-crossratio-invariant}
在定理 \ref{thm:xi-leyang-crossratio-demixing-crossing} 的语境中，考虑一类复混合：实验/格点参数 \(z_{\mathrm{phys}}\) 与标度变量 \(\tau\) 的关系满足
\(
z=z_{\mathrm{phys}}+\beta\tau+O(\tau^2)
\)，其中 \(\beta\in\CC\) 一般不为实数。
令
\[
r_{m,n}(L,\tau):=\frac{\Im z_m(L,\tau)}{\Im z_n(L,\tau)}.
\]
则存在情形使 \(r_{m,n}(L,\tau)\) 在 \(\tau\to 0\) 时对 \(\Im(\beta)\) 的一阶敏感性不为零，从而其“不同体积交叉点”不再为坐标不变量；相反，任意交比 \(\mathrm{CR}_{i,j;k,\ell}(L,\tau)\) 对 \(z\mapsto z+\beta\tau\) 严格不变，且其交叉点渐近仍满足 \eqref{eq:xi-leyang-crossratio-crossing-asymp}。
\end{corollary}
\begin{proof}
取局部线性化 \(z_m(\tau)=z_m(0)+\beta\tau+o(\tau)\) 且 \(z_m(0)\notin\RR\)。
则 \(\Im z_m(\tau)=\Im z_m(0)+\Im(\beta)\tau+o(\tau)\)，从而比值 \(r_{m,n}\) 的一阶导数一般包含 \(\Im(\beta)\) 项而不消失。
另一方面，交比仅由差分构成，公共平移 \(+\beta\tau\) 被完全消去；结论遂由定理 \ref{thm:xi-leyang-crossratio-demixing-crossing} 的不变性与 crossing 渐近立即推出。
\end{proof}
\leanverified{paper\_xi\_leyang\_imag\_ratio\_unstable\_crossratio\_invariant}

\begin{proposition}[变量混合修正的 \texorpdfstring{$1$}{1}-余边界结构]\label{prop:xi-leyang-variable-mixing-coboundary}
\leanverified{paper\_xi\_leyang\_variable\_mixing\_coboundary}
在 Lee--Yang 零点的有限尺度缩放理论中，除定理 \ref{thm:xi-leyang-crossratio-demixing-crossing} 的射影化消混合接口外，文献中亦常使用比值观测量的变量混合展开来刻画首阶偏差。
为给出一个可复用的代数结构陈述，设对某一临界切片 \(\tau=\tau_c\) 与某一指数差
\(
\bar y:=y_t-y_h<0
\)
存在如下渐近表示：对任意 \(n\neq m\)，
\begin{equation}\label{eq:xi-leyang-imag-ratio-variable-mixing-expansion}
R_{nm}(L):=\frac{\Im z_n(L,\tau_c)}{\Im z_m(L,\tau_c)}
=r_{nm}\Bigl(1+D_{nm}\,L^{2\bar y}+O(L^{4\bar y})\Bigr),
\qquad L\to\infty,
\end{equation}
其中 \(r_{nm}=X_n/X_m\) 为普适常数比值，而主导混合系数满足
\begin{equation}\label{eq:xi-leyang-variable-mixing-Dnm-square-difference}
D_{nm}=-\kappa\,(Y_n^2-Y_m^2),
\qquad
\kappa:=\frac{a_{12}^2}{a_{22}^2}.
\end{equation}
则 \(D_{nm}\) 具有严格的余边界表述：存在势函数 \(U_n\) 使得
\begin{equation}\label{eq:xi-leyang-variable-mixing-Dnm-coboundary}
D_{nm}=U_n-U_m,
\qquad
U_n:=-\kappa\,Y_n^2.
\end{equation}
因而对任意互异指标 \(n,m,p,q\) 有恒等式
\begin{equation}\label{eq:xi-leyang-variable-mixing-four-term-identity}
D_{nm}+D_{pq}-D_{nq}-D_{pm}=0,
\end{equation}
并且对任意闭路 \(\gamma:n_0\to n_1\to\cdots\to n_k=n_0\)（允许重复顶点）有
\[
\sum_{j=0}^{k-1}D_{n_jn_{j+1}}=0.
\]
\end{proposition}
\begin{proof}
由 \eqref{eq:xi-leyang-variable-mixing-Dnm-square-difference} 直接取 \(U_n:=-\kappa Y_n^2\) 即得 \eqref{eq:xi-leyang-variable-mixing-Dnm-coboundary}。
四项恒等式 \eqref{eq:xi-leyang-variable-mixing-four-term-identity} 与闭路求和为零皆为望远镜消去。
\end{proof}

\begin{proposition}[普适类的零点射影指纹：有限交比在射影等价下完备]\label{prop:xi-leyang-crossratio-projective-fingerprint-complete}
设 \(\{\zeta_1,\dots,\zeta_M\}\subset\widehat{\CC}\)（\(M\ge 4\)）两两互异，并定义交比集合
\[
\mathfrak{I}_M:=\Bigl\{\mathrm{CR}^{\ast}_{i,j;k,\ell}:\ 1\le i<j<k<\ell\le M\Bigr\},
\qquad
\mathrm{CR}^{\ast}_{i,j;k,\ell}
:=
\frac{(\zeta_i-\zeta_j)(\zeta_k-\zeta_\ell)}{(\zeta_i-\zeta_k)(\zeta_j-\zeta_\ell)}.
\]
则 \(\mathfrak{I}_M\) 在射影等价（\(\mathrm{PGL}_2(\CC)\) 作用）意义下为完备不变量：若另一配置 \(\{\zeta_1',\dots,\zeta_M'\}\subset\widehat{\CC}\)（允许重标记）满足其交比集合与 \(\mathfrak{I}_M\) 相同，则存在 \(g\in\mathrm{PGL}_2(\CC)\) 使得 \(\zeta_m'=g(\zeta_m)\) 对所有 \(m\) 成立。
在定理 \ref{thm:xi-leyang-crossratio-demixing-crossing} 的解析混合口径（仅含复仿射自由度）下，\(\mathfrak{I}_M\) 因而给出对混合免疫的普适类“零点指纹”。
\end{proposition}
\begin{proof}
\(\mathrm{PGL}_2(\CC)\) 在 \(\widehat{\CC}\) 上三重可递。
取三个互异点作归一化，可令 \(g\) 把 \((\zeta_1,\zeta_2,\zeta_3)\) 送到 \((0,1,\infty)\)。
在该归一化下，任意第四点 \(\zeta_m\)（\(m\ge 4\)）由相应交比值唯一确定，从而整组点由交比集合确定到唯一的射影轨道。
对带重标记的情形，先以重标记对齐三点归一化，再用同一论证完成。证毕。
\end{proof}

\begin{theorem}[由两枚有限体积零点反演度量因子与混合位移]\label{thm:xi-leyang-invert-metric-and-shift-from-two-zeros}
在定理 \ref{thm:xi-leyang-crossratio-demixing-crossing} 的假设下，取任意 \(i\neq j\) 且 \(\zeta_i\neq\zeta_j\)，并假设指数 \(y_h\) 已由独立机制给定。
定义
\begin{equation}\label{eq:xi-leyang-invert-twozeros-def}
A_L:=L^{y_h}\frac{z_i(L,0)-z_j(L,0)}{\zeta_i-\zeta_j},
\qquad
z_{0,L}:=z_i(L,0)-A_L L^{-y_h}\zeta_i.
\end{equation}
则当 \(L\to\infty\) 时，
\begin{equation}\label{eq:xi-leyang-invert-twozeros-asymp}
A_L=A(0)+O(L^{-\omega}),\qquad
z_{0,L}=z_0(0)+O\!\bigl(L^{-(y_h+\omega)}\bigr).
\end{equation}
因而仅用两枚有限体积零点，即可在误差 \(O(L^{-\omega})\) 的意义下反演出非普适度量因子 \(A(0)\) 与混合位移 \(z_0(0)\)。
\end{theorem}
\begin{proof}
将 \eqref{eq:xi-leyang-fss-analytic-mixing} 特化到 \(\tau=0\)，并对 \(z_i-z_j\) 作差分，得
\[
z_i(L,0)-z_j(L,0)
=
A(0)L^{-y_h}\Bigl((\zeta_i-\zeta_j)+(\xi_i-\xi_j)L^{-\omega}+O(L^{-2\omega})\Bigr),
\]
代入 \eqref{eq:xi-leyang-invert-twozeros-def} 得 \(A_L=A(0)\bigl(1+O(L^{-\omega})\bigr)\)。
再将 \(z_i(L,0)=z_0(0)+A(0)L^{-y_h}\bigl(\zeta_i+\xi_iL^{-\omega}+O(L^{-2\omega})\bigr)\) 代入 \(z_{0,L}\) 的定义，得到
\(
z_{0,L}=z_0(0)+O(L^{-(y_h+\omega)})
\)。
\end{proof}

\begin{remark}[混合位移与度量归一化的可观测反演接口]\label{rem:xi-leyang-twozero-calibration-interface}
当 \(z\) 取为复逸度/复化学势等外参时，\(z_0(0)\) 可解释为“物理参数 \(\mapsto\) 标度场坐标”映射在临界点处的解析混合位移，而 \(A(0)\) 给出标度场的非普适度量归一化。
定理 \ref{thm:xi-leyang-invert-metric-and-shift-from-two-zeros} 表明二者可由两枚有限体积零点在 \(O(L^{-\omega})\) 误差内反演，因而可作为后续零点几何不变量（例如交比/比值接口）的先验校准量；在需要数值重加权或受符号问题影响的设定中，该校准可减弱对混合参数与坐标选择的敏感性。
\end{remark}

在具有两相关标度场与变量混合的临界设定中，Lee--Yang 零点的体积交汇定位通常需要控制混合参数对坐标选择的敏感性。
下述陈述给出一种射影不变量接口：以两枚有限体积零点在两种体积下形成的四点交比作为观测量，从而将非普适的线性混合消去至无关指数控制的修正项，并给出临界点与相关指数的反演口径。

\begin{theorem}[两体积两零点交比的不变性与临界定位]\label{thm:xi-leyang-two-scale-crossratio-invariant}
设 \(\{Z_L(T,\mu)\}_{L\ge 1}\) 为一族有限体积配分函数，\(T\in\RR\) 为控制参量，\(\mu\in\CC\) 为可复化的外参（外场/化学势等）。
在临界点 \((T_c,\mu_c)\) 的邻域内，设存在两相关标度场 \((t,h)\) 的解析坐标化，使得线性化满足
\begin{equation}\label{eq:xi-leyang-mixing-linearization}
\binom{t}{h}
=
M\binom{T-T_c}{\mu-\mu_c}
+O\!\bigl(\lvert(T-T_c,\mu-\mu_c)\rvert^2\bigr),
\qquad
M=\begin{pmatrix}\alpha_t&\beta_t\\ \alpha_h&\beta_h\end{pmatrix},
\quad \beta_h\neq 0.
\end{equation}
设关于标度场 \(h\) 的 Lee--Yang 零点（在固定 \(T\) 下、按某一固定约定对 \(j\ge 1\) 编号）满足有限尺度标度：存在指数 \(y_t>0,y_h>0\) 与 \(\omega>0\)，以及普适标度函数族 \(\{\xi_j(\cdot)\}_{j\ge 1}\)，使得对每个固定 \(j\) 有一致渐近
\begin{equation}\label{eq:xi-leyang-hzero-fss}
h_j(T;L)=L^{-y_h}\,\xi_j\!\bigl((T-T_c)L^{y_t}\bigr)+O(L^{-y_h-\omega}).
\end{equation}
令 \(\mu_j(T;L)\) 表示在同一编号约定下、以 \(\mu\) 为变量得到的第 \(j\) 个 Lee--Yang 零点（即 \(Z_L(T,\mu_j(T;L))=0\)）。
对任意 \(s>1\) 与互异指标 \(j\neq k\)，定义两体积两零点四点交比
\begin{equation}\label{eq:xi-leyang-two-scale-crossratio-def}
\mathcal Q_{jk}(T;L,s)
:=
\operatorname{cr}\bigl(\mu_j(T;L),\mu_k(T;L),\mu_j(T;sL),\mu_k(T;sL)\bigr),
\qquad
\operatorname{cr}(a,b,c,d):=\frac{(a-c)(b-d)}{(a-d)(b-c)}.
\end{equation}
则存在仅依赖于普适标度零点函数 \(\xi_j,\xi_k\) 与 \(s\) 的函数 \(\Psi_{jk}(\cdot;s)\)，使得当 \(L\to\infty\) 时成立一致渐近展开
\begin{equation}\label{eq:xi-leyang-two-scale-crossratio-scaling}
\mathcal Q_{jk}(T;L,s)=\Psi_{jk}\!\bigl((T-T_c)L^{y_t};s\bigr)+O(L^{-\omega}).
\end{equation}
特别地，在临界点处有
\begin{equation}\label{eq:xi-leyang-two-scale-crossratio-critical}
\mathcal Q_{jk}(T_c;L,s)=\Psi_{jk}(0;s)+O(L^{-\omega}),
\end{equation}
并且 \(\mathcal Q_{jk}\) 对线性混合矩阵 \(M\) 的非普适参数（包括 \(\alpha_h/\beta_h\) 等混合比）不敏感。
\leanverified{paper\_xi\_leyang\_two\_scale\_crossratio\_invariant}
\end{theorem}
\begin{proof}
由 \eqref{eq:xi-leyang-mixing-linearization} 可得在固定 \(T\) 下 \(h\) 与 \(\mu\) 的关系为局部解析函数，其一阶主部为仿射变换
\[
h=\alpha_h(T-T_c)+\beta_h(\mu-\mu_c)+O\!\bigl(\lvert(T-T_c,\mu-\mu_c)\rvert^2\bigr).
\]
将零点 \(\mu=\mu_j(T;L)\) 代入并忽略可并入无关修正的高阶项，可写为
\[
\mu_j(T;L)=\mu_c-\frac{\alpha_h}{\beta_h}(T-T_c)+\frac{1}{\beta_h}h_j(T;L)+O(L^{-y_h-\omega}).
\]
因此 \(\{\mu_j(T;L)\}\) 由 \(\{h_j(T;L)\}\) 经仿射映射 \(\mu=\lambda h+\delta\)（\(\lambda=\beta_h^{-1}\neq 0\)，\(\delta\) 含平移项）得到。
交比 \(\operatorname{cr}\) 在任意 Möbius 变换下不变，故在仿射变换下严格不变；从而
\[
\mathcal Q_{jk}(T;L,s)
=
\operatorname{cr}\bigl(h_j(T;L),h_k(T;L),h_j(T;sL),h_k(T;sL)\bigr)+O(L^{-\omega}).
\]
代入 \eqref{eq:xi-leyang-hzero-fss} 并注意交比对公共乘法因子 \(L^{-y_h}\) 的整体约去性，即得到 \eqref{eq:xi-leyang-two-scale-crossratio-scaling}；取 \(T=T_c\) 得 \eqref{eq:xi-leyang-two-scale-crossratio-critical}。
\end{proof}

\begin{corollary}[交比交汇点的漂移阶]\label{cor:xi-leyang-two-scale-crossratio-crossing-drift}
在定理 \ref{thm:xi-leyang-two-scale-crossratio-invariant} 的假设下，令 \(T_{jk}(L,s)\) 为方程
\begin{equation}\label{eq:xi-leyang-two-scale-crossratio-crossing-eq}
\mathcal Q_{jk}(T;L,s)=\mathcal Q_{jk}(T;sL,s)
\end{equation}
在 \(T_c\) 的小邻域内确定的解，并假设 \(\partial_x\Psi_{jk}(0;s)\neq 0\) 以保证局部唯一性。
则存在常数 \(C_{jk}(s)\) 使得
\[
T_{jk}(L,s)-T_c
=
C_{jk}(s)\,L^{-(y_t+\omega)}+o\!\bigl(L^{-(y_t+\omega)}\bigr).
\]
\leanverified{paper\_xi\_leyang\_two\_scale\_crossratio\_crossing\_drift}
\end{corollary}
\begin{proof}
由 \eqref{eq:xi-leyang-two-scale-crossratio-scaling} 得
\(
\mathcal Q_{jk}(T;L,s)=\Psi_{jk}(x;s)+L^{-\omega}R_{jk}(x;s)+o(L^{-\omega})
\)
（\(x=(T-T_c)L^{y_t}\)），并对 \(sL\) 同理。
将两式在 \(x=0\) 作泰勒展开并令其相等，可解得 \(x=O(L^{-\omega})\)，从而 \(T-T_c=x\,L^{-y_t}=O(L^{-(y_t+\omega)})\)；系数由 \(\partial_x\Psi_{jk}(0;s)\) 与 \(R_{jk}(0;s)\) 的组合给出。
\end{proof}

\begin{proposition}[由交比斜率反演相关指数]\label{prop:xi-leyang-two-scale-crossratio-slope-exponent}
\leanverified{paper\_xi\_leyang\_two\_scale\_crossratio\_slope\_exponent}
在定理 \ref{thm:xi-leyang-two-scale-crossratio-invariant} 的假设下，若 \(\partial_x\Psi_{jk}(0;s)\neq 0\)，则
\[
\left.\frac{\dd}{\dd T}\mathcal Q_{jk}(T;L,s)\right|_{T=T_c}
=
\partial_x\Psi_{jk}(0;s)\,L^{y_t}+O(L^{y_t-\omega}).
\]
因此对任意两体积 \(L_1,L_2\) 有
\[
y_t
=
\lim_{L_1,L_2\to\infty}
\frac{\log\bigl|\partial_T\mathcal Q_{jk}(T_c;L_2,s)\bigr|-\log\bigl|\partial_T\mathcal Q_{jk}(T_c;L_1,s)\bigr|}
{\log L_2-\log L_1},
\qquad
\nu=1/y_t.
\]
\end{proposition}
\begin{proof}
对 \eqref{eq:xi-leyang-two-scale-crossratio-scaling} 关于 \(T\) 求导并在 \(T=T_c\) 处代入即可。
\end{proof}

\begin{proposition}[零点轨迹给出临界切向混合比]\label{prop:xi-leyang-two-scale-crossratio-mixing-slope}
\leanverified{paper\_xi\_leyang\_two\_scale\_crossratio\_mixing\_slope}
在定理 \ref{thm:xi-leyang-two-scale-crossratio-invariant} 的设定下，进一步假设 \(y_h>y_t\)，并令 \(\mu_j(T;L)\) 为固定 \(T\) 下第 \(j\) 个 Lee--Yang 零点。
则极限
\[
\lim_{L\to\infty}\left.\frac{\dd}{\dd T}\Re\,\mu_j(T;L)\right|_{T=T_c}
=
-\frac{\alpha_h}{\beta_h}
\]
存在且与 \(j\) 无关；并且收敛速率满足
\[
\left.\frac{\dd}{\dd T}\Re\,\mu_j(T;L)\right|_{T=T_c}
=
-\frac{\alpha_h}{\beta_h}
+O(L^{y_t-y_h})+O(L^{-\omega}).
\]
\end{proposition}
\begin{proof}
由上一证明中的仿射近似
\(
\mu_j(T;L)=\mu_c-\frac{\alpha_h}{\beta_h}(T-T_c)+\frac{1}{\beta_h}h_j(T;L)+O(L^{-y_h-\omega})
\)
与 \eqref{eq:xi-leyang-hzero-fss} 得
\[
\frac{\dd}{\dd T}\mu_j(T;L)\Big|_{T=T_c}
=
-\frac{\alpha_h}{\beta_h}
+\frac{1}{\beta_h}\,L^{-y_h}\,\xi_j'(0)\,L^{y_t}
+O(L^{y_t-y_h-\omega}),
\]
其中第二项在 \(y_h>y_t\) 时趋于 \(0\)。取实部即得结论。
\end{proof}

\begin{proposition}[二阶非线性混合在交比中的额外压制]\label{prop:xi-leyang-two-scale-crossratio-quadratic-mixing-suppression}
\leanverified{paper\_xi\_leyang\_two\_scale\_crossratio\_quadratic\_mixing\_suppression}
在定理 \ref{thm:xi-leyang-two-scale-crossratio-invariant} 的设定下，进一步设在 \((T_c,\mu_c)\) 处具有二阶展开
\[
h=\alpha_h\Delta T+\beta_h\Delta\mu+\tfrac12\gamma_h(\Delta\mu)^2+O(\lvert(\Delta T,\Delta\mu)\rvert^3),
\qquad
\Delta T:=T-T_c,\ \Delta\mu:=\mu-\mu_c,
\]
其中 \(\gamma_h\in\RR\)。
则两体积交比 \(\mathcal Q_{jk}(T;L,s)\) 的主导非线性混合修正满足
\[
\mathcal Q_{jk}(T;L,s)
=
\Psi_{jk}\!\bigl((T-T_c)L^{y_t};s\bigr)
+O(L^{-\omega})
+O(L^{-y_h}).
\]
\end{proposition}
\begin{proof}
在固定 \(T\) 下把 \(\mu\mapsto h\) 视作一元解析函数
\(
h(\mu)=A+B(\mu-\mu_c)+C(\mu-\mu_c)^2+O((\mu-\mu_c)^3)
\)
（\(B=\beta_h\neq 0\), \(C=\gamma_h/2\)）。
令 \(\tilde\mu_j\) 为忽略二次项后的仿射近似零点，则 \(\tilde\mu_j-\mu_c=\Theta(L^{-y_h})\)，并由反函数的二阶稳定性得 \(\mu_j-\tilde\mu_j=O(L^{-2y_h})\)。
交比作为四点的光滑函数，其对点位扰动的一阶变分由四点差商的有界线性组合控制；在当前尺度下差商的典型量级为 \(L^{y_h}\)，故交比差异为 \(O(L^{y_h})\cdot O(L^{-2y_h})=O(L^{-y_h})\)。
与定理 \ref{thm:xi-leyang-two-scale-crossratio-invariant} 的 \(O(L^{-\omega})\) 无关修正合并即得结论。
\end{proof}

\begin{conjecture}[交比标度函数的局部单叶性与唯一交汇]\label{conj:xi-leyang-two-scale-crossratio-univalence}
在三维 Ising 普适类中，固定 \(s>1\) 与任意 \(j\neq k\)，定理 \ref{thm:xi-leyang-two-scale-crossratio-invariant} 中的普适函数 \(\Psi_{jk}(x;s)\) 在 \(x=0\) 的某个实邻域内为单叶解析函数，且其对实轴的限制严格单调。
若该性质成立，则对充分大的 \(L\)，方程 \eqref{eq:xi-leyang-two-scale-crossratio-crossing-eq} 在 \(T_c\) 邻域内具有唯一解，并对有限体积数据的统计扰动保持稳定。
该猜想与将 \(\mathcal Q_{jk}\) 视为临界附近的普适共形坐标一致：其逆映射可在误差可控的意义下重建标度变量 \(x=(T-T_c)L^{y_t}\)，从而为比值法之外提供更强的数值判据与误差控制接口（参见 \cite{WadaKitazawaKanaya2025LYZR} 对 LYZR 框架的动机与混合误差讨论）。
\end{conjecture}


\subsubsection{一站点重--稠密 \texorpdfstring{$\mathrm{SU}(N)$}{SU(N)} toy 模型的 Lee--Yang 零点谱、解析半径与有限阶可复原性}\label{subsubsec:xi-hdqcd1-leyang-radius-susceptibility-cw}
为给出一类“零点位置—解析半径—低阶涨落”之间关系可被完全闭式核验的原型，考虑一站点重--稠密（heavy--dense）\(\mathrm{SU}(N)\) toy 模型。
该模型的配分函数在 \(\mu\) 的复平面上具有显式 Lee--Yang 零点谱，从而可把“最邻近零点（edge）”与 \(\mu=0\) 的低阶易感精确对应；同时，在引入均匀平均场耦合后，局部权重可诱导一个三态 Curie--Weiss 模型，其临界/三临界判据在颜色数 \(N\) 上出现严格分岔。

\begin{theorem}[HDQCD\(_1\)：Lee--Yang 零点谱与 \(\mu=0\) 的泰勒解析半径闭式]\label{thm:xi-hdqcd1-leyang-spectrum-radius}
\leanverified{paper\_xi\_hdqcd1\_leyang\_spectrum\_radius}
设 \(N\ge 2\) 为整数，\(h>0\)，并令 \(\dd U\) 为 \(\mathrm{SU}(N)\) 上归一化 Haar 测度。
定义
\begin{equation}\label{eq:xi-hdqcd1-partition-function}
Z_N(h,\mu)
:=
\int_{\mathrm{SU}(N)}
\det\!\bigl(I_N+h e^{\mu}U\bigr)\,
\det\!\bigl(I_N+h e^{-\mu}U^{-1}\bigr)\,
\dd U.
\end{equation}
记
\begin{equation}\label{eq:xi-hdqcd1-CN-aN}
C_N(h):=\sum_{k=0}^{N}h^{2k},
\qquad
a_N(h):=\frac{C_N(h)}{2h^{N}}.
\end{equation}
则 \(Z_N(h,\mu)\) 满足：
\begin{enumerate}
  \item \textbf{闭式表达。}对一切 \(\mu\in\CC\)，有
  \begin{equation}\label{eq:xi-hdqcd1-closed-form}
  Z_N(h,\mu)=C_N(h)+2h^{N}\cosh(N\mu).
  \end{equation}
  \item \textbf{Lee--Yang 零点谱。}由于 \(a_N(h)>1\)，方程 \(Z_N(h,\mu)=0\) 的全部解为
  \begin{equation}\label{eq:xi-hdqcd1-zeros}
  \mu_{\pm,m}(h)=\frac{1}{N}\Bigl(\pm \operatorname{arcosh}(a_N(h))+\mathrm{i}(2m+1)\pi\Bigr),
  \qquad m\in\ZZ.
  \end{equation}
  \item \textbf{解析半径。}令 \(p_N(\mu):=\log Z_N(h,\mu)\) 取与 \(\mu=0\) 附近一致的解析分支，则其在 \(\mu=0\) 的泰勒展开收敛半径为
  \begin{equation}\label{eq:xi-hdqcd1-radius}
  R_N(h)=\min_{\pm,m}\bigl|\mu_{\pm,m}(h)\bigr|
  =\frac{1}{N}\sqrt{\operatorname{arcosh}(a_N(h))^{2}+\pi^{2}}.
  \end{equation}
\end{enumerate}
\end{theorem}

\begin{proof}
令 \(z:=he^{\mu}\)、\(w:=he^{-\mu}\)，则 \(zw=h^{2}\)、\(z^{N}=h^{N}e^{N\mu}\)、\(w^{N}=h^{N}e^{-N\mu}\)。
利用外幂表示（以 \(\chi_{\wedge^{k}}\) 表示第 \(k\) 外幂表示的特征标）：
\[
\det(I_N+zU)=\sum_{k=0}^{N}z^{k}\chi_{\wedge^{k}}(U),
\qquad
\det(I_N+wU^{-1})=\sum_{\ell=0}^{N}w^{\ell}\chi_{\wedge^{\ell}}(U^{-1}).
\]
在 \(\mathrm{SU}(N)\) 上有 \(\chi_{\wedge^{\ell}}(U^{-1})=\overline{\chi_{\wedge^{\ell}}(U)}\)，并且不可约特征标满足正交性
\(
\int_{\mathrm{SU}(N)}\chi_{\rho}(U)\overline{\chi_{\rho'}(U)}\,\dd U=\delta_{\rho,\rho'}.
\)
需要注意的是：\(\wedge^{0}\) 为平凡表示，而 \(\wedge^{N}\) 为行列式表示，其在 \(\mathrm{SU}(N)\) 上亦为平凡表示，故
\(
\chi_{\wedge^{0}}=\chi_{\wedge^{N}}=1
\)
在 \(\mathrm{SU}(N)\) 上恒等。
因此积分在 \((k,\ell)\) 的求和中仅保留以下贡献：
\begin{itemize}
  \item \((k,\ell)=(0,0),(0,N),(N,0),(N,N)\) 皆给出常数 \(1\) 的积分，贡献为 \(1+w^{N}+z^{N}+z^{N}w^{N}\)；
  \item 对 \(1\le k\le N-1\)，仅 \(\ell=k\) 给出非零积分，贡献为 \(\sum_{k=1}^{N-1}(zw)^{k}\)。
\end{itemize}
从而
\[
Z_N(h,\mu)
=(1+z^{N})(1+w^{N})+\sum_{k=1}^{N-1}(zw)^{k}
=\sum_{k=0}^{N}h^{2k}+h^{N}(e^{N\mu}+e^{-N\mu}),
\]
即得到 \eqref{eq:xi-hdqcd1-closed-form}。
\par
由 \eqref{eq:xi-hdqcd1-CN-aN} 得
\[
\frac{C_N(h)}{h^{N}}=\sum_{k=0}^{N}h^{2k-N}
=
\sum_{j\in\{-N,-N+2,\dots,N-2,N\}}h^{j}
\ge N+1,
\]
其中用到对称配对 \(h^{j}+h^{-j}\ge 2\)（\(j>0\)）以及当 \(N\) 为偶数时的中心项 \(h^{0}=1\)；
因此 \(a_N(h)=C_N(h)/(2h^{N})\ge (N+1)/2>1\)。
零点条件 \(Z_N(h,\mu)=0\) 等价于
\[
\cosh(N\mu)=-a_N(h).
\]
由于 \(a_N(h)>1\)，解为
\(
N\mu=\mathrm{i}(2m+1)\pi\pm \operatorname{arcosh}(a_N(h))
\)
（\(m\in\ZZ\)），即 \eqref{eq:xi-hdqcd1-zeros}。
最后，\(\log Z_N(h,\mu)\) 在 \(\mu=0\) 的解析半径由到最近零点的距离决定；
最近零点对应 \(m=0\) 或 \(m=-1\)，从而得到 \eqref{eq:xi-hdqcd1-radius}。证毕。
\end{proof}

\begin{corollary}[大 \texorpdfstring{$N$}{N} 极限的实轴夹逼阈值]\label{cor:xi-hdqcd1-largeN-pinch}
在定理 \ref{thm:xi-hdqcd1-leyang-spectrum-radius} 的设定下，取 \(m=0\)。
则最近两条零点分支为
\[
\mu_{\pm,0}(h)=\frac{1}{N}\Bigl(\pm \operatorname{arcosh}(a_N(h))+\mathrm{i}\pi\Bigr).
\]
当 \(N\to\infty\) 时，
\[
\Im \mu_{\pm,0}(h)=\frac{\pi}{N}\longrightarrow 0,
\qquad
\Re \mu_{\pm,0}(h)=\pm \frac{1}{N}\operatorname{arcosh}(a_N(h))\longrightarrow \pm |\log h|.
\]
特别地，若写 \(h=e^{-m/T}\)，则 \(\Re \mu_{\pm,0}(h)\to \pm m/T\)，从而在 \(N\to\infty\) 的极限中出现一条沿实轴的夹逼阈值 \(|\mu|=m/T\)。
\end{corollary}

\begin{proof}
若 \(h\neq 1\)，有闭式
\[
C_N(h)=\frac{h^{2(N+1)}-1}{h^{2}-1}.
\]
当 \(h>1\) 时，
\(
C_N(h)\sim \frac{h^{2(N+1)}}{h^{2}-1}
\)，从而
\[
a_N(h)=\frac{C_N(h)}{2h^{N}}
\sim \frac{h^{N+2}}{2(h^{2}-1)},
\qquad
\operatorname{arcosh}(a_N(h))
=\log\!\bigl(a_N(h)+\sqrt{a_N(h)^{2}-1}\bigr)
\sim \log(2a_N(h))
\sim N\log h,
\]
故 \(\frac{1}{N}\operatorname{arcosh}(a_N(h))\to \log h\)。
当 \(0<h<1\) 时，由 \(C_N(h)=h^{2N}C_N(1/h)\) 得 \(a_N(h)=a_N(1/h)\)，从而同样得到极限为 \(|\log h|\)。
虚部 \(\pi/N\to 0\) 由显式表达直接得到。证毕。
\end{proof}

\begin{proposition}[有限阶可复原性：二阶易感精确重构最近 Lee--Yang 边缘]\label{prop:xi-hdqcd1-edge-from-chi2}
\leanverified{paper\_xi\_hdqcd1\_edge\_from\_chi2}
在定理 \ref{thm:xi-hdqcd1-leyang-spectrum-radius} 的设定下，定义二阶（重子数）易感
\begin{equation}\label{eq:xi-hdqcd1-chi2-def}
\chi_2(h):=\left.\frac{\dd^{2}}{\dd\mu^{2}}\log Z_N(h,\mu)\right|_{\mu=0}.
\end{equation}
则最近的 Lee--Yang 边缘零点满足精确公式
\begin{equation}\label{eq:xi-hdqcd1-edge-from-chi2}
\mu_{\mathrm{edge}}^{\pm}(h)
=\frac{1}{N}\Bigl(\mathrm{i}\pi\pm \operatorname{arcosh}\Bigl(\frac{N^{2}}{\chi_2(h)}-1\Bigr)\Bigr),
\end{equation}
该表达式不依赖任何高阶易感或级数外推。
\end{proposition}

\begin{proof}
由 \eqref{eq:xi-hdqcd1-closed-form}（记 \(C=C_N(h)\)）得
\[
Z_N(h,0)=C+2h^{N},\qquad
Z_N'(h,0)=0,\qquad
Z_N''(h,0)=2h^{N}N^{2}.
\]
因此
\[
\chi_2(h)
=\left.\left(\frac{Z_N''}{Z_N}-\Bigl(\frac{Z_N'}{Z_N}\Bigr)^{2}\right)\right|_{\mu=0}
=\frac{Z_N''(h,0)}{Z_N(h,0)}
=\frac{2h^{N}N^{2}}{C+2h^{N}}.
\]
解得
\[
\frac{C}{2h^{N}}=\frac{N^{2}}{\chi_2(h)}-1.
\]
零点条件 \(\cosh(N\mu)=-C/(2h^{N})\) 因而等价于
\(
N\mu=\mathrm{i}\pi\pm \operatorname{arcosh}\bigl(N^{2}/\chi_2(h)-1\bigr)
\)，除以 \(N\) 即得 \eqref{eq:xi-hdqcd1-edge-from-chi2}。证毕。
\end{proof}

\begin{theorem}[HDQCD\(_{\mathrm{CW}}\)：三态 Curie--Weiss 诱导模型的临界/三临界判据与颜色数分岔]\label{thm:xi-hdqcd-cw-tricritical}
\leanverified{paper\_xi\_hdqcd\_cw\_tricritical}
固定整数 \(N\ge 2\)。令 \(A>0\)、\(B>0\)，并对体积 \(V\) 定义三态 Curie--Weiss 模型（\(s_i\in\{-1,0,1\}\)）：
\begin{equation}\label{eq:xi-hdqcd-cw-partition}
\mathcal Z_{V}(A,B;\mu,J)
:=
\sum_{(s_1,\dots,s_V)\in\{-1,0,1\}^{V}}
\exp\!\Bigl(\frac{J}{2V}\Bigl(\sum_{i=1}^{V}s_i\Bigr)^{2}+N\mu\sum_{i=1}^{V}s_i\Bigr)
\prod_{i=1}^{V} w_{s_i},
\qquad
w_0=A,\ w_{+1}=w_{-1}=B,
\end{equation}
其中 \(J>0\)。
则在 \(\mu=0\) 下成立：
\begin{enumerate}
  \item \textbf{热力学极限与变分表示。}自由能密度
  \[
  f(A,B;J):=\lim_{V\to\infty}-\frac{1}{V}\log\mathcal Z_V(A,B;0,J)
  \]
  存在，并满足
  \begin{equation}\label{eq:xi-hdqcd-cw-variational}
  f(A,B;J)=\inf_{x\in\RR}\left[\frac{x^{2}}{2J}-\log\bigl(A+2B\cosh x\bigr)\right].
  \end{equation}
  \item \textbf{对称相的 spinodal。}令
  \(
  \phi_J(x):=\frac{x^{2}}{2J}-\log(A+2B\cosh x)
  \)。
  则 \(x=0\) 的失稳点为
  \begin{equation}\label{eq:xi-hdqcd-cw-spinodal}
  J_{\mathrm{sp}}(A,B)=\frac{A+2B}{2B}.
  \end{equation}
  \item \textbf{Landau 四阶判据与三临界条件。}
  令 \(\Phi_J(m):=\phi_J(Jm)\) 为对称 Landau 有效势（均匀磁化参数 \(m\)）。
  在 \(J=J_{\mathrm{sp}}(A,B)\) 处有展开
  \begin{equation}\label{eq:xi-hdqcd-cw-landau-quartic}
  \Phi_{J_{\mathrm{sp}}}(m)=\Phi_{J_{\mathrm{sp}}}(0)
  +\frac{(A+2B)^{2}(4B-A)}{192\,B^{3}}\,m^{4}
  +O(m^{6}).
  \end{equation}
  因而：
  \begin{itemize}
    \item 若 \(4B>A\)，则在 \(J=J_{\mathrm{sp}}\) 发生二阶相变；
    \item 若 \(4B=A\)，则出现三临界点（四阶项消失且六阶项主导稳定）；
    \item 若 \(4B<A\)，则二阶相变被排除，且相变（若存在）为一阶并满足转变点 \(J_{\mathrm{t}}<J_{\mathrm{sp}}\)。
  \end{itemize}
  \item \textbf{颜色数分岔（由 \(A=C_N(h),B=h^{N}\) 诱导）。}
  取
  \(
  A=C_N(h)=\sum_{k=0}^{N}h^{2k}
  \)、
  \(
  B=h^{N}
  \)
  （与定理 \ref{thm:xi-hdqcd1-leyang-spectrum-radius} 一致）。
  则三临界条件 \(4B=A\) 等价于
  \begin{equation}\label{eq:xi-hdqcd-cw-tricritical-eq}
  \sum_{k=0}^{N}h^{2k}=4h^{N}
  \qquad\Longleftrightarrow\qquad
  \sum_{k=0}^{N}h^{2k-N}=4.
  \end{equation}
  并且对任意 \(h>0\) 有普适下界
  \begin{equation}\label{eq:xi-hdqcd-cw-universal-lb}
  \sum_{k=0}^{N}h^{2k-N}\ge N+1,
  \end{equation}
  从而三临界仅可能在 \(N\le 3\) 出现；更精确地：
  \begin{itemize}
    \item \(N=3\)：唯一正根为 \(h=1\)，对应三临界点 \((h,J)=(1,3)\)；
    \item \(N=2\)：存在两正根 \(h=\frac{\sqrt5\pm 1}{2}\)（互为倒数）；
    \item \(N\ge 4\)：无正根，因而不存在三临界点，且在 spinodal 点四阶系数必为负。
  \end{itemize}
\end{enumerate}
\end{theorem}

\begin{proof}
\textbf{(1)} 对 \eqref{eq:xi-hdqcd-cw-partition} 在 \(\mu=0\) 处作 Hubbard--Stratonovich 变换：令 \(S=\sum_{i=1}^{V}s_i\)，则
\[
\exp\!\Bigl(\frac{J}{2V}S^{2}\Bigr)
=\sqrt{\frac{V}{2\pi J}}\int_{\RR}\exp\!\Bigl(-\frac{Vx^{2}}{2J}+xS\Bigr)\,\dd x.
\]
代入并交换求和与积分，得到
\[
\mathcal Z_V(A,B;0,J)
=\sqrt{\frac{V}{2\pi J}}\int_{\RR}
\exp\!\Bigl(-V\Bigl[\frac{x^{2}}{2J}-\log(A+Be^{x}+Be^{-x})\Bigr]\Bigr)\,\dd x,
\]
即 \eqref{eq:xi-hdqcd-cw-variational}。
\par
\textbf{(2)} 直接计算
\[
\phi_J''(0)=\frac{1}{J}-\frac{2B}{A+2B},
\]
故 \(x=0\) 的局部稳定性在 \(\phi_J''(0)=0\) 处改变，从而得到 \eqref{eq:xi-hdqcd-cw-spinodal}。
\par
\textbf{(3)} 在 \(x=0\) 处对 \(\phi_J(x)\) 作 Taylor 展开并以 \(x=Jm\) 重标度得到 \(\Phi_J(m)=\phi_J(Jm)\) 的 Landau 展开。
在 \(J=J_{\mathrm{sp}}\) 处二次项系数消失，四次项系数的直接展开给出 \eqref{eq:xi-hdqcd-cw-landau-quartic}。
四阶系数符号决定在 \(J\) 穿越 \(J_{\mathrm{sp}}\) 附近的分岔型：系数为正对应连续分岔（典型二阶）；系数为负对应在 \(J<J_{\mathrm{sp}}\) 已出现与 \(m=0\) 共存的非零极小并发生全局极小交换（典型一阶）；系数为零则需用下一非零偶次项决定稳定性并给出三临界结构。
\par
\textbf{(4)} 由指数集合对称性，令 \(\mathcal J=\{-N,-N+2,\dots,N-2,N\}\)，则
\[
\sum_{k=0}^{N}h^{2k-N}=\sum_{j\in\mathcal J}h^{j}.
\]
将 \(j>0\) 与 \(-j\) 配对并用 AM--GM 得 \(h^{j}+h^{-j}\ge 2\)；若 \(N\) 为偶数，另有中心项 \(j=0\) 等于 \(1\)。
故得到普适下界 \eqref{eq:xi-hdqcd-cw-universal-lb}，并且等号当且仅当 \(h=1\)。
由 \eqref{eq:xi-hdqcd-cw-tricritical-eq}，若存在正根则必有 \(4=\sum h^{2k-N}\ge N+1\)，从而 \(N\le 3\)。
当 \(N=3\) 时下界即为 \(4\)，唯一达到等号点为 \(h=1\)；
当 \(N=2\) 时方程为 \(h^{-2}+1+h^{2}=4\)，等价于 \(h^{4}-3h^{2}+1=0\)，从而
\(
h=\frac{\sqrt5\pm 1}{2}
\)；
当 \(N\ge 4\) 时 \(\sum h^{2k-N}\ge N+1\ge 5\) 排除解的存在。
证毕。
\end{proof}

\begin{corollary}[\texorpdfstring{$\mathrm{SU}(3)$}{SU(3)} 三临界点的临界指数]\label{cor:xi-hdqcd-su3-tricritical-exponent}
在定理 \ref{thm:xi-hdqcd-cw-tricritical} 的诱导情形 \(N=3\) 且 \(h=1\) 下，三临界点为 \((J,\mu)=(3,0)\)。
令 \(m\) 为对称相序参量（均匀磁化），则在 \(J\downarrow 3\) 的邻域存在常数 \(\alpha,\beta>0\) 使
\[
\Phi_J(m)-\Phi_J(0)=\alpha(J-3)m^{2}+\beta m^{6}+O(m^{8}),
\]
从而在自发对称破缺相（\(J>3\)）有
\[
m(J)\asymp (J-3)^{1/4}.
\]
\end{corollary}

\begin{proof}
当 \(N=3\)、\(h=1\) 时有 \(A=C_3(1)=4\)、\(B=1\)，由 \eqref{eq:xi-hdqcd-cw-spinodal} 得 \(J_{\mathrm{sp}}=3\)，并且 \(4B-A=0\) 使四阶项消失（见 \eqref{eq:xi-hdqcd-cw-landau-quartic}）。
继续对 \(\Phi_J\) 作六阶展开可得六阶系数为正，因而在 \(J-3\) 小量下得到所示形态。
对 \(\alpha(J-3)m^{2}+\beta m^{6}\) 取极小，非零极小满足 \(m^{4}\asymp (J-3)\)，即 \(m(J)\asymp (J-3)^{1/4}\)。证毕。
\end{proof}

\endinput


\begin{conjecture}[交比常数的数域有限生成与稳定性]\label{conj:xi-leyang-crossratio-arithmetic-rigidity}
在带代数可积结构的普适类中，设标度零点配置 \(\{\zeta_j\}\) 的有限阶交比集合 \(\mathfrak{I}_M\) 定义如命题 \ref{prop:xi-leyang-crossratio-projective-fingerprint-complete}。
猜想由 \(\mathfrak{I}_M\) 生成的数域
\[
K_M:=\QQ\bigl(\mathfrak{I}_M\bigr)
\]
为有限扩张，并在 \(M\to\infty\) 时稳定于某个固定的数域 \(K_\infty\)。
若该猜想成立，则普适类的连续零点几何可被一个离散算术对象（数域及其伽罗瓦数据）所锚定，从而为“零点指纹的算术刚性”提供可核验的判别接口。
\end{conjecture}


\subsubsection{谱曲线的椭圆化与分歧单值化}\label{subsubsec:xi-spectral-elliptic-ramification-monodromy}
沿用谱四次 \(\Pi(\lambda,y)\) 及其 Lee--Yang 型分歧的判别式设定。令 \(C\) 为仿射曲线 \(\Pi(\lambda,y)=0\) 的正规化，并将 \(C\) 通过 \(y\)-投影视为 \(\mathbb P^1_y\) 上的四层覆盖。其有限分歧像由
\(
\mathrm{Disc}_{\lambda}(\Pi)=-y(y-1)P_{\mathrm{LY}}(y)
\)
的零点给出，其中
\[
P_{\mathrm{LY}}(y):=256y^3+411y^2+165y+32
\]
为 Lee--Yang 三次因子。引入坐标
\[
u:=y-\lambda^2\qquad(\text{等价地 }w:=2u+1),
\]
则同一方程在 \(\QQ\) 上双有理同构于极小椭圆曲线
\[
E_{\min}:\quad u^2+u=\lambda^3-\lambda
\]
（见定理 \ref{thm:xi-terminal-zm-leyang-elliptic-structure} 与结论 \ref{con:xi-terminal-zm-elliptic-normalization}），从而 \(C\) 为属 \(1\) 的光滑射影曲线。等价地，令
\[
X:=\lambda,\qquad Y:=y-\lambda^2+\tfrac12,
\]
则可在短 Weierstrass 模型
\(
E:\ Y^2=X^3-X+\tfrac14
\)
上实现函数域同一；其在本文约定下满足 \(\Delta(E)=37\) 且 \(j(E)=110592/37\)，并由于 \(j(E)\) 非代数整数而不具有复乘。并且由平移 \(Y\mapsto Y-\tfrac12\) 回到整系数最小模型 \(E_{\min}:u^2+u=X^3-X\)（Cremona/LMFDB 标识 \(37\mathrm a1\)，见注记 \ref{rem:xi-terminal-elliptic-37a1-modular-anchor}）。
\par
在上述椭圆化下，重量参数成为 \(E\) 上的有理函数
\(
y=X^2+Y-\tfrac12\in H^{0}\!\bigl(E,\mathcal O_{E}(4O)\bigr)
\)；
其极点除子为 \((y)_\infty=4O\)，并且两条有理分歧纤维 \(y=0\) 与 \(y=1\) 具有完全显式的主除子分解（结论 \ref{con:xi-terminal-zm-elliptic-normalization}）。因此作为次数 \(4\) 的覆盖 \(\pi_y:E\to\mathbb P^1_y\)，纤维 \(y^{-1}(0)\) 与 \(y^{-1}(1)\) 各包含唯一一个 \(\QQ\)-有理简单分歧点，分歧指数均为 \(2\)（定理 \ref{thm:xi-terminal-zm-leyang-elliptic-structure}）。
\par
有限分歧点由 \(\dd y=0\) 刻画并分解为两条有理分歧点与一条非有理三点轨道：除去两条有理解外，其余三个分歧点的 \(X\)-坐标为三次
\(
16X^3-9X^2+1=0
\)
的三根，而对应 \(Y\)-坐标由切触线性方程
\(
4XY+3X^2-1=0
\)
唯一确定（定理 \ref{thm:xi-terminal-zm-leyang-elliptic-structure}）。该三次的判别式为 \(-2^{2}\cdot 3^{3}\cdot 37\)，与 Lee--Yang 三次 \(P_{\mathrm{LY}}(y)\) 的判别式 \(-3^{9}\cdot 31^{2}\cdot 37\) 共享平方自由核 \(-111\)，从而二次伴随域 \(\QQ(\sqrt{-111})\) 同时支配“非有理分歧点的三次坐标域”与“非有理分歧值的三次坐标域”（参见结论 \ref{con:xi-terminal-zm-leyang-cubic-arithmetic}），并构成该谱覆盖的稳定算术指纹。
\par
将 \(\Pi(\lambda,y)\) 视为 \(\QQ(y)\) 上关于 \(\lambda\) 的四次多项式，则其泛型伽罗瓦群与几何单值群为 \(S_4\simeq W(A_3)\)（结论 \ref{con:xi-terminal-zm-generic-galois-s4} 与定理 \ref{thm:xi-terminal-zm-leyang-elliptic-structure}），从而覆盖的几何单值化达到四叶覆盖所允许的最大非交换对称。以 Hitchin--cameral 的观点看，该 Weyl 单值化可作为秩 \(4\) 光谱覆盖的结构性数据，使 Lee--Yang 分歧层、Hurwitz 护照与固定椭圆曲线在同一代数对象上可计算对齐；在几何 Langlands 的物理实现中，这类谱/cameral 数据亦可作为 Hitchin 纤维化下的最小可审计原型之一。\cite{KapustinWitten2006GeometricLanglands}
\par
与该 \(y\)-投影的四层覆盖相对，\(\Pi\) 作为关于 \(y\) 的二次多项式满足 \(\mathrm{Disc}_{y}(\Pi)=4\lambda^3-4\lambda+1\)，其零点与椭圆曲线的 \(2\)-division 三次一致，故 \(y\)-方向的二次分歧由 \(E[2]\) 的坐标域精确控制。与此相并行，二阶碰撞三次
\(
P_2(\Lambda)=\Lambda^3-2\Lambda^2-2\Lambda+2
\)
所生成的三次数域与椭圆曲线 \(E\) 的 \(2\)-division 三次（非平凡 \(E[2]\) 的 \(X\)-坐标域）在 \(\QQ\) 上同构，且可由显式的二次 Tschirnhaus 变换互相恢复（结论 \ref{con:xi-terminal-zm-collision-s2-e2-field-isomorphism}）。因此“碰撞动力学的三次骨架”与“谱椭圆曲线的模 \(2\) 结构”在算术层面被识别为同一三次域对象。
\par
\begin{conjecture}[Langlands--Lee--Yang 对应的可计算约束：Hankel 导体的分歧支撑]\label{conj:xi-langlands-leyang-hankel-conductor-ramification}
沿用本段的谱四次覆盖 \(\pi_y:E\to\PP^1_y\) 与其驯/野分歧素数集合的记号。
对固定的 Rényi 阶 \(q\ge 2\)，令
\[
S_q(m):=\sum_{x\in X_m} d_m(x)^q
\]
为碰撞矩，并在“纤维谱有限支撑”口径下以推论 \ref{cor:xi-fiber-spectrum-hankel-conductor-prony-padic} 的定义构造 Hankel 导体
\[
\mathfrak N_{q,m}:=\mathrm{rad}\bigl(\det H_0^{(q,m)}\bigr),
\qquad
H_0^{(q,m)}:=\bigl(S_{q(i+j)}(m)\bigr)_{0\le i,j\le D-1},
\]
其中 \(D\) 取为该口径下的纤维谱支撑大小（使得 \(H_0^{(q,m)}\) 为首主块）。
则存在一族以 \(\pi_y\) 的几何分歧为核心的可计算“参数对象” \(\mathcal L_{q}\)，使得下述三条约束同时成立：
\begin{enumerate}
  \item \textbf{（支撑一致）}\(\mathfrak N_{q,m}\) 的素支撑在 \(m\to\infty\) 的稳定意义下与 \(\mathcal L_q\) 的坏约化/分歧素数集合一致；
  等价地，碰撞矩核在 \(p\)-进通道上的可恢复性障碍素数恰对应于几何对象的坏素数。
  \item \textbf{（好素数局部因子对齐）}对一切好素数 \(p\nmid \mathfrak N_{q,m}\)，由 Hankel 最小实现（Prony--Hankel）所诱导的局部 \(\zeta\)-因子与 \(\mathcal L_q\) 的 Hasse--Weil/Artin 型局部因子在某个完成化切片上相符。
  \item \textbf{（单位圆塔的导体约束）}上述局部因子对齐在 Joukowsky 完成化后推出一条受导体支撑控制的 Lee--Yang 单位圆零点塔：好素数通道对应的扭曲谱落在单位圆上，而坏素数通道的偏离仅由 \(\mathfrak N_{q,m}\) 的素支撑所允许。
\end{enumerate}
\end{conjecture}



对良素 \(p\notin\{2,37\}\)，令
\(
N_p(y):=\card{\{\lambda\in\FF_p:\ \Pi(\lambda,y)=0\}}
\)。
则椭圆化给出严格点计数恒等式
\(
\#E(\FF_p)=1+\sum_{y\in\FF_p}N_p(y)
\)，从而 Frobenius 迹
\(
a_p:=p+1-\#E(\FF_p)
\)
满足
\(
\sum_{y\in\FF_p}(N_p(y)-1)=-a_p
\)
（结论 \ref{con:xi-terminal-zm-fiber-root-counting-hecke} 与结论 \ref{con:xi-terminal-zm-hecke-eigenvalue-fiber-deviation}）；在模性锚定下 \(a_p\) 同时等于导子 \(37\) 新形式的 Hecke 本征值。进一步，对驯约化素数 \(p\notin\{2,3,31,37\}\)，函数域 Chebotarev 在固定单值群 \(\mathrm{Mon}(\pi_y)\cong S_4\) 上给出有效的 Chebotarev--Sato--Tate 型等分布：未分歧参数的纤维根数 \(N_p(y)\) 的极限比例由 \(S_4\) 在四点作用下的不动点分布决定，并具有 \(O(\sqrt p)\) 的平方根误差；其中
\(
\pi_0=\tfrac38,\ \pi_1=\tfrac13,\ \pi_2=\tfrac14,\ \pi_4=\tfrac1{24}
\)
且 \(N_p(y)=3\) 仅可能发生在分歧集合的模 \(p\) 化上（结论 \ref{con:xi-terminal-zm-s4-modp-fiber-rootcount-fixedpoints}）。由此，谱四次在有限域上的分解统计、Hecke 迹与单值群的共轭类分布在同一椭圆覆盖框架中被同时锁定。




\begin{theorem}[Lee--Yang 特征谱曲线的椭圆化与单值护照：固定椭圆曲线与 \(S_4\) 分支结构]\label{thm:xi-terminal-zm-leyang-elliptic-structure}
沿用谱四次多项式
\[
\Pi(\lambda,y):=\lambda^4-\lambda^3-(2y+1)\lambda^2+\lambda+y(y+1)\in\ZZ[\lambda,y].
\]
记仿射谱曲线 \(\mathcal C:=\{(\lambda,y)\in\mathbb A^2:\Pi(\lambda,y)=0\}\)。
定义平面多项式自同构
\[
T:\mathbb A^2\to\mathbb A^2,\qquad (\lambda,y)\longmapsto(x,u):=(\lambda,y-\lambda^2),
\qquad
T^{-1}:(x,u)\longmapsto(\lambda,y)=(x,u+x^2).
\]
则直接代换给出
\[
\Pi(\lambda,y)=u^2+u-(x^3-x)\qquad(x=\lambda,\ u=y-\lambda^2),
\]
从而 \(T\) 将 \(\mathcal C\) 精确送到极小椭圆曲线
\[
E_{\min}:\quad u^2+u=x^3-x,
\]
并在 \(\QQ\) 上诱导仿射曲线同构 \(\mathcal C\cong E_{\min}\)。
因此 \(\Pi(\lambda,y)=0\) 的规范化为属 \(1\) 的光滑射影曲线，并在 \(\QQ\) 上与该极小椭圆曲线同构；其分支结构因而由一条与参数 \(y\) 无关的固定有理椭圆曲线所锚定，而非任意四次覆盖的偶然现象。
等价地，若以 \(U:=\lambda\)、\(V:=y-\lambda^{2}\) 作为谱坐标，则归一化模型亦可写为
\[
V^{2}+V=U^{3}-U.
\]
为与短 Weierstrass 口径对齐，令
\[
w:=2u+1=2y-2\lambda^2+1\quad(\text{亦记 }S:=w),
\]
则恒有恒等式
\[
4\Pi(\lambda,y)=w^{2}-\bigl(4x^{3}-4x+1\bigr)=w^{2}-4x^{3}+4x-1.
\]
将 \(\Pi(\lambda,y)\) 视为关于 \(y\) 的二次多项式，则其 \(y\)-判别式满足
\(
\mathrm{Disc}_{y}(\Pi)=4\lambda^{3}-4\lambda+1
\)，
上式即等价写为 \(\,w^{2}=\mathrm{Disc}_{y}(\Pi)\,\) 在谱曲线 \(\mathcal C\) 上的实现。
从而 \((\lambda,y)\mapsto(x,w)\) 给出 \(\mathcal C\) 与短 Weierstrass 椭圆曲线
\[
\mathcal E:\quad w^{2}=4x^{3}-4x+1\qquad(x=\lambda)
\]
之间的显式仿射同构；其逆映射由
\[
y=x^{2}-\frac12+\frac{w}{2}\qquad\Bigl(\text{亦即 }y=\lambda^{2}-\frac12+\frac{w}{2}\Bigr)
\]
给出。
在该极小模型下，非平凡 \(2\)-挠点满足 \(u=-\tfrac12\)（等价 \(w=0\)），其 \(x\)-坐标由
\[
4x^{3}-4x+1=0
\]
给出；因此 \(\mathrm{Disc}_{y}(\Pi)=4\lambda^{3}-4\lambda+1\) 恰为 \(E_{\min}\) 的 \(2\)-除法多项式。
等价地，自然投影 \(\pi_x:E_{\min}\to\mathbb P^1_x\) 为次数 \(2\) 的覆盖，其分歧点的 \(x\)-坐标恰为三次方程 \(4x^{3}-4x+1=0\) 的三根，并与 \(E_{\min}(\overline{\QQ})\) 的非平凡 \(2\)-挠点一致。
将 \(E_{\min}\) 配方为 \((2u+1)^2=4x^3-4x+1\) 并以 \(w=2u+1\) 记之，则得到等价的短 Weierstrass 模型
\[
\mathcal E:\quad w^2=4x^3-4x+1\qquad(g_2=4,\ g_3=-1),
\]
其判别式仍为 \(\Delta(\mathcal E)=g_2^3-27g_3^2=37\)。
进一步以伸缩 \(X:=4x,\ Y:=4w\) 得整系数短 Weierstrass 模型
\[
Y^2=X^3-16X+16,
\]
其在 Cremona/LMFDB 的标准标识下对应同源类 \(37\mathrm a\) 的强 Weil 曲线 \(37\mathrm a1\)（注记 \ref{rem:xi-terminal-elliptic-37a1-modular-anchor}）。
等价地，谱支 \(\lambda(y)\) 可视为椭圆曲线 \(E_{\min}\) 上有理函数 \(y=x^2+u=x^2-\tfrac12+\tfrac{w}{2}\) 的局部反演，从而原先的多分支谱反演被提升为椭圆曲线上的有理函数反演问题。
在复解析层面，\(\mathcal E(\CC)\cong \CC/\Lambda\) 的 Weierstrass 一致化可取为
\[
x=\wp(z;\,g_2=4,g_3=-1),\qquad w=\wp'(z;\,g_2=4,g_3=-1),\qquad
y=\wp(z)^2-\frac12+\frac{\wp'(z)}{2},
\]
因而 Lee--Yang 分析中的四支代数谱函数 \(\{\lambda_j(y)\}_{j=1}^{4}\) 可理解为椭圆函数 \(z\mapsto\wp(z)\) 的代数反演之分支（见注记 \ref{rem:xi-terminal-zm-lambda-plus-wp-uniformization}）。
该椭圆曲线满足
\[
c_4(E_{\min})=48,\qquad c_6(E_{\min})=-216,\qquad
\Delta(E_{\min})=37,\qquad j(E_{\min})=\frac{c_4(E_{\min})^{3}}{\Delta(E_{\min})}=\frac{110592}{37},
\]
从而仅在素数 \(37\) 处坏约化且导子为 \(37\)；由此 Lee--Yang 型解析边界可在同一算术椭圆对象上与周期格/模性锚点统一表述。
由椭圆曲线模性定理，存在权 \(2\)、level \(37\) 的新形式 \(f\) 使 \(L(E_{\min},s)=L(f,s)\)；
其外部命名与交叉核验接口见注记 \ref{rem:xi-terminal-elliptic-37a1-modular-anchor}。
\par
令 \(\pi_y:E_{\min}\to\mathbb P^1_y\) 为有理函数
\(
y=x^2+u=\frac{w+2x^2-1}{2}
\)
所诱导的次数 \(4\) 态射，则 \(y\) 在无穷远点 \(O\in E_{\min}\) 处具有唯一四阶极点，因而 \(\infty\in\mathbb P^1_y\) 为全分歧分歧值。
等价地，绕 \(y=\infty\) 的几何单值化为一条 \(4\)-循环（其局部一致化参数与 \(\lambda\) 的 Puiseux 展开见定理 \ref{thm:xi-terminal-zm-puiseux-infty-4cycle}）。
并且在 \((x,w)\)-模型 \(\mathcal E:w^2=4x^3-4x+1\) 上，两个有理分歧值的纤维除子满足完全分解
\[
(y)_0=[(0,1)]+[(-1,-1)]+2[(1,-1)],
\qquad
(y-1)_0=[(2,-5)]+[(1,1)]+2[(-1,1)],
\]
其中点以 \((x,w)\)-坐标表示。
因此作为次数 \(4\) 的覆盖 \(\pi_y\)，纤维 \(y^{-1}(0)\) 与 \(y^{-1}(1)\) 各包含唯一的 \(\QQ\)-有理简单分歧点（对应于上式中系数为 \(2\) 的点），且分歧指数均为 \(2\)。
其有限临界点由 \(\dd y=0\) 刻画。

\begin{proposition}[分歧除子与 Abel--Jacobi 零和约束]\label{prop:xi-terminal-zm-leyang-ramification-divisor-stokes}
沿用定理的记号，将 \(\pi_y:E_{\min}\to\PP^1_y\) 视为次数 \(4\) 的有理映射，并令其分歧除子为
\[
R_{\pi_y}:=\sum_{P\in E_{\min}(\CC)}(e_P-1)\,[P],
\]
其中 \(e_P\) 为 \(\pi_y\) 在 \(P\) 处的分歧指数。
则在除子类群 \(\mathrm{Pic}(E_{\min})\) 中成立
\[
R_{\pi_y}\sim 2\,\pi_y^{-1}(\infty).
\]
由于 \(\pi_y^{-1}(\infty)=4O\)，从而
\[
R_{\pi_y}\sim 8O,\qquad \deg R_{\pi_y}=8.
\]
并且在识别 \(\mathrm{Pic}^0(E_{\min})\cong E_{\min}\) 的群律下，分歧点的加和满足严格零和约束
\[
\sum_{P\in E_{\min}(\CC)}(e_P-1)\,P=O.
\]
\end{proposition}
\leanverified{paper\_xi\_terminal\_zm\_leyang\_ramification\_divisor\_stokes}
\begin{proof}
取 \(\PP^1\) 的仿射坐标 \(z\)（无穷远点为 \(\infty\)），则 \((\dd z)=-2(\infty)\) 作为除子恒成立。
对非恒定有理函数 \(f\) 的一般公式给出
\[
(\dd f)=f^\ast(\dd z)+R_f.
\]
取 \(f=\pi_y\) 即得
\[
(\dd \pi_y)=-2\,\pi_y^{-1}(\infty)+R_{\pi_y}.
\]
另一方面，\(E_{\min}\) 为属 \(1\) 曲线，其典范丛平凡；任一亚纯 \(1\)-形式的除子为主除子类。
故 \((\dd \pi_y)\) 在 \(\mathrm{Pic}(E_{\min})\) 中为零，从而 \(R_{\pi_y}-2\pi_y^{-1}(\infty)\) 为主除子，得到 \(R_{\pi_y}\sim 2\pi_y^{-1}(\infty)\)。
再由 \(\pi_y^{-1}(\infty)=4O\) 得 \(R_{\pi_y}\sim 8O\) 与 \(\deg R_{\pi_y}=8\)。
将主除子在 \(\mathrm{Pic}^0(E_{\min})\cong E_{\min}\) 下对应为群和 \(O\)，即得到分歧点零和约束。证毕。
\end{proof}

将短 Weierstrass 模型 \(\mathcal E:w^2=4x^3-4x+1\) 以 \(X=x,\ Y=w/2\) 配方为
\(
Y^2=X^3-X+\tfrac14
\)，则
\(
y=X^2+Y-\tfrac12
\)。
故对任意 \(c\in\mathbb P^1\)，纤维 \(y=c\) 等价于抛物线
\[
Y=-X^2+\Bigl(c+\frac12\Bigr)
\]
与该椭圆曲线的交点除子，有限分歧点即为该抛物线族与椭圆曲线的切触点。
由 \(2Y\,\dd Y=(3X^2-1)\dd X\) 与 \(\dd y=2X\dd X+\dd Y\) 得
\[
\dd y=\frac{4XY+3X^2-1}{2Y}\,\dd X,
\]
从而有限临界点由 \(4XY+3X^2-1=0\) 刻画；在 \((x,w)\)-模型 \(\mathcal E\) 上等价写为
\[
2xw+3x^2-1=0
\qquad\bigl(\text{等价地 }4\lambda w+6\lambda^2-2=0,\ \lambda=x\bigr),
\]
并恰分为两类：
\begin{itemize}
  \item 两个有理临界点
  \[
  P_0:\ (x,u)=(1,-1)\mapsto y=0,\qquad P_1:\ (x,u)=(-1,0)\mapsto y=1;
  \]
  \item 三个非有理临界点 \(P_2,P_3,P_4\)，其 \(x\)-坐标为三次方程
  \[
  16x^3-9x^2+1=0
  \]
  的三根；相应的 \(u\)-坐标由
  \[
  u=\frac{1-3x^2-2x}{4x}
  \]
  唯一确定，并因此
  \[
  w=2u+1=-\frac{3x^2-1}{2x};
  \]
  从而临界值满足显式表达式
  \[
  y=x^2+u=\frac{4x^3-3x^2-2x+1}{4x}
  =\frac{(x-1)(4x^2+x-1)}{4x}
  =x^{2}-\frac34x-\frac12+\frac{1}{4x}.
  \]
\end{itemize}
并且消元恒等式给出
\[
\mathrm{Res}_x\!\left(16x^3-9x^2+1,\ 4xy-\bigl(4x^3-3x^2-2x+1\bigr)\right)=-64\,P_{\mathrm{LY}}(y),
\]
由结式的乘法性亦可将两条有理分歧值 \(y=0,1\) 同时并入消元，从而得到全有限分歧值集合的证书
\[
\mathrm{Res}_x\!\Bigl((x-1)(x+1)\bigl(16x^3-9x^2+1\bigr),\ 4xy-(x-1)(4x^2+x-1)\Bigr)
=1024\,y(y-1)P_{\mathrm{LY}}(y).
\]
从而上述参数化与判别式分解中的 Lee--Yang 三次因子严格一致。
上述三个非平凡有限临界值恰为不可约三次多项式
\[
P_{\mathrm{LY}}(y):=256y^3+411y^2+165y+32
\]
的三根，其中唯一负实根即 \(y_{\mathrm{LY}}\)。
记其余两根为共轭对 \(y_{\pm}\)，并记三次方程 \(16\lambda^3-9\lambda^2+1=0\) 的根为
\(\lambda_{\mathrm{LY}}\in\RR\) 与共轭对 \(\lambda_{\pm}\)（其中 \(y_{\pm}=y(\lambda_{\pm})\)）。
数值上，
\[
y_{\mathrm{LY}}\approx -1.13445200300831188727,\qquad
y_{\pm}\approx -0.235508373495844\pm 0.233925552127514\,i,
\]
\[
\lambda_{\mathrm{LY}}\approx -0.27343481065245424608,\qquad
\lambda_{\pm}\approx 0.417967405326227\pm 0.232114033943254\,i.
\]
因此 \(\pi_y\) 的有限分支值集合恰为
\[
y\in\{0,1\}\ \cup\ \{P_{\mathrm{LY}}(y)=0\},
\]
并且在每个有限分支值上均仅出现一处简单分歧（惯性生成元为换位）。
并且判别式满足
\[
\Delta(\mathcal E)=37,
\qquad
\mathrm{Disc}\bigl(4x^3-4x+1\bigr)=2^{4}\cdot 37,
\qquad
\mathrm{Disc}(16x^3-9x^2+1)=-2^{2}\cdot 3^{3}\cdot 37,
\qquad
\mathrm{Disc}\bigl(P_{\mathrm{LY}}(y)\bigr)=-3^{9}\cdot 31^{2}\cdot 37,
\]
因此 \(\mathrm{Disc}(16x^3-9x^2+1)\) 与 \(\mathrm{Disc}\bigl(P_{\mathrm{LY}}(y)\bigr)\) 的平方自由部分同为 \(-111=-3\cdot 37\)，两者分裂域共享同一二次分辨子域 \(\QQ(\sqrt{-111})\)。
该虚二次子域因而可视为该四次谱覆盖的内禀算术指纹，并与所选椭圆化坐标归一化无关。
从而素数 \(37\) 同时出现在临界点与临界值多项式的平方自由部分中，并与 \(\Delta(E_{\min})=37\) 的坏约化素点对齐。
因而判别式分解所识别的分歧集合可被解释为 \(\pi_y\) 的临界值结构，而非 \(\ZZ[y]\) 中的偶发因式分裂。
\par
对任一有限临界值 \(y_*\in\{0,1\}\cup\{P_{\mathrm{LY}}(y)=0\}\)，记 \(\lambda_*\) 为 \(\Pi(\lambda,y_*)=0\) 的并合谱值（即相应二重根）。
则在 \(y\to y_*\) 时并合的两条谱支具有平方根型 Puiseux 展开
\[
\lambda(y)=\lambda_*\pm \kappa_*(y-y_*)^{1/2}+\beta_*(y-y_*)\pm\gamma_*(y-y_*)^{3/2}+O\bigl((y-y_*)^{2}\bigr),
\]
其中首项系数满足显式代数公式
\[
\kappa_*^{2}
=-\,\frac{2\,\partial_y\Pi(\lambda_*,y_*)}{\partial_{\lambda}^{2}\Pi(\lambda_*,y_*)}
=\frac{3\lambda_*^{2}-1}{8\lambda_*^{3}-3\lambda_*^{2}-1}.
\]
在非平凡 Lee--Yang 分歧支（即 \(\lambda_*\) 满足 \(16\lambda_*^3-9\lambda_*^2+1=0\)）上，分母可化简为
\(
8\lambda_*^{3}-3\lambda_*^{2}-1=\frac{3}{2}(\lambda_*^{2}-1)
\)，因此
\[
\kappa_*^{2}=\frac{2(3\lambda_*^{2}-1)}{3(\lambda_*^{2}-1)}.
\]
并且线性项与三阶项系数可由更高阶 Taylor 系数显式给出
\[
\beta_*
=-\,\frac{\partial_{\lambda y}\Pi(\lambda_*,y_*)}{\partial_{\lambda}^{2}\Pi(\lambda_*,y_*)}
+\frac{\partial_{\lambda}^{3}\Pi(\lambda_*,y_*)\,\partial_y\Pi(\lambda_*,y_*)}{3\bigl(\partial_{\lambda}^{2}\Pi(\lambda_*,y_*)\bigr)^{2}}.
\]
在非平凡 Lee--Yang 分歧支上进一步可化简为
\[
\beta_*=-\frac{\lambda_*(3\lambda_*^{2}+16\lambda_*+5)}{18(\lambda_*^{2}-1)^{2}},
\qquad
\gamma_*=-\frac{\lambda_*^{2}\bigl(1152\lambda_*^{7}-1152\lambda_*^{6}-280\lambda_*^{5}+723\lambda_*^{4}-312\lambda_*^{3}+30\lambda_*^{2}+16\lambda_*-5\bigr)}{2\bigl(8\lambda_*^{3}-3\lambda_*^{2}-1\bigr)^{4}}\cdot \kappa_*^{-1}.
\]
特别地，在 \(y_*=0\) 与 \(y_*=1\) 处分别得到
\[
\lambda(y)=1\pm \sqrt{\frac{y}{2}}+O(y),
\qquad
\lambda(y)=-1\pm \frac{i}{\sqrt6}(y-1)^{1/2}+O(y-1),
\]
而在实 Lee--Yang 边界 \(y_*=y_{\mathrm{LY}}\) 处有
\[
y_{\mathrm{LY}}\approx -1.13445200300831188727,\qquad
\lambda_*=\lambda_{\mathrm{LY}}\approx-0.27343481065245424608,\qquad
\kappa_{\mathrm{LY}}:=\kappa_*\approx 0.74761098459045973850,\qquad
\beta_{\mathrm{LY}}:=\beta_*\approx 0.0150716815102551,\qquad
\gamma_{\mathrm{LY}}:=\gamma_*\approx -0.0418222862574037,\qquad
\kappa_{\mathrm{LY}}^{2}\approx 0.55892218428031662857,\qquad
\frac{\kappa_{\mathrm{LY}}}{\lambda_{\mathrm{LY}}}\approx -2.73414706345016524243,\qquad
\frac{\lambda_{\mathrm{LY}}^{2}}{\kappa_{\mathrm{LY}}^{2}}\approx 0.1337692397606565,\qquad
\pi^{2}\frac{\lambda_{\mathrm{LY}}^{2}}{\kappa_{\mathrm{LY}}^{2}}\approx 1.3202494774721529.
\]
从而任何由该并合谱支构造的解析量（例如选定分支的 \(\log\lambda\)）在 \(y_{\mathrm{LY}}\) 处均呈严格平方根型 Lee--Yang 奇性，其幅度由 \(\kappa_{\mathrm{LY}}\) 在代数上完全锁定。
\par
映射 \(\pi_y\) 的分歧型为五个有限简单分歧与一处无穷远全分歧（Hurwitz 护照 \((2,2,2,2,2,4)\)），故其几何单值群包含一个 \(4\)-循环与一个换位，从而
\par
该 Hurwitz 护照在 Nielsen 类意义下提供覆盖在 Hurwitz 空间中的离散同伦不变量，使分歧循环型可作为跨模型比较的稳定指纹。
\[
\mathrm{Mon}(\pi_y)\cong S_4.
\]
等价地，函数域扩张 \(\QQ(y)\subset \QQ(y,\lambda)\) 的伽罗瓦闭包之伽罗瓦群为 \(S_4\)。
并且谱四次的 \(\lambda\)-判别式满足
\[
\mathrm{Disc}_{\lambda}\bigl(\Pi(\lambda,y)\bigr)=-y(y-1)P_{\mathrm{LY}}(y),
\]
从而该 \(S_4\)-扩张的唯一指数 \(2\) 正规子域由判别式平方根生成：
\[
\QQ(y)\bigl(\sqrt{\mathrm{Disc}_{\lambda}(\Pi(\lambda,y))}\bigr)
=\QQ(y)\Bigl(\sqrt{-y(y-1)P_{\mathrm{LY}}(y)}\Bigr).
\]
该二次子扩张亦可等价表述为 \(A_4\triangleleft S_4\) 的不动域，并对应于属 \(2\) 的超椭圆曲线
\[
C:\quad \zeta^{2}=-y(y-1)P_{\mathrm{LY}}(y),
\]
其分歧点为 \(y\in\{0,1\}\cup\{P_{\mathrm{LY}}(y)=0\}\cup\{\infty\}\) 六点。
记 \(\widetilde{C}\to\mathbb P^1_y\) 为该四次覆盖的 \(S_4\)-伽罗瓦闭包，则由惰性群阶
\[
e_\infty=4,\qquad e_0=e_1=e_{y_1}=e_{y_2}=e_{y_3}=2
\qquad\bigl(P_{\mathrm{LY}}(y_j)=0\bigr)
\]
与 Riemann--Hurwitz 公式得到
\[
g(\widetilde{C})
=1+\frac{|S_4|}{2}\Bigl(-2+\Bigl(1-\frac14\Bigr)+5\Bigl(1-\frac12\Bigr)\Bigr)
=16.
\]
作为等价的代数判别，可取 \(\Pi(\lambda,y)\) 的一条标准 resolvent cubic
\[
R_y(z)=64 z^3-(176+256y) z^2+(76+128y)z-(64y^2-48y+9),
\]
其满足 \(\mathrm{Disc}_z(R_y)=2^{24}\mathrm{Disc}_{\lambda}(\Pi)\)，并在 \(\QQ(y)[z]\) 中不可约；结合 \(\mathrm{Disc}_{\lambda}(\Pi)\) 在 \(\QQ(y)\) 中非平方，即得 \(\mathrm{Gal}(\Pi/\QQ(y))\cong S_4\)。
因此任何试图将其分歧结构压缩到 \(A_4\)、\(D_4\) 或 \(V_4\) 等较小单值群的策略均与分歧循环型不相容。
由于 \(S_4\cong W(A_3)\)，上述单值化同时给出一条由简单分歧换位（局部反射）与无穷远全分歧的 \(4\)-循环（Coxeter 元）所生成的 \(A_3\)–Weyl 群最小实现。
该最大单值性在代数层面表明谱四次覆盖达到最不退化的对称复杂度，并为谱支的全局解析延拓、剪切线选择与零点凝聚结构的不可约性分析提供群论骨架。
另一方面，由标准同构 \(S_4\simeq \mathrm{PGL}_2(\FF_3)\) 可将其识别为八面体有限单值群像；
在去除分歧值的开曲线 \(U:=\mathbb P^1_y\setminus\{0,1,P_{\mathrm{LY}}(y)=0,\infty\}\) 上，
上述单值化等价给出一个有限像的 \(\mathrm{PGL}_2(\FF_3)\)-投影局部系统，其存在性与分歧集合完全由 \(\Pi(\lambda,y)\) 的判别式机制所固定。
\leanverified{paper\_xi\_terminal\_zm\_leyang\_elliptic\_structure}
\end{theorem}

\par
\begin{theorem}[\texorpdfstring{$\abs{4O}$}{|4O|} 线性扭曲的一参数四次谱模型]\label{thm:xi-terminal-zm-leyang-linear-twist-quartic-family}
\leanverified{paper\_xi\_terminal\_zm\_leyang\_linear\_twist\_quartic\_family}
设短 Weierstrass 模型
\[
E/\QQ:\quad Y^2=X^3-X+\tfrac14
\]
与无穷远点 \(O\in E\)。
令
\[
y:=X^2+Y-\tfrac12\in \QQ(E),
\qquad
y_t:=y+tX=X^2+tX+Y-\tfrac12\in \QQ(E)\qquad (t\in\QQ).
\]
则 \(y_t\) 在 \(O\) 处具有极除子 \(4O\)，从而诱导次数 \(4\) 的覆盖 \(y_t:E\to\PP^1\)。
并且在仿射 \((X,y)\)-坐标下，\(y_t\) 的像曲线由显式四次方程给出：
\[
\Pi_t(X,y):=X^4+(2t-1)X^3+(t^2-2y-1)X^2+(1-t-2ty)X+y(y+1)=0.
\]
\end{theorem}
\begin{proof}
由定义 \(y_t=X^2+tX+Y-\tfrac12\) 可解得
\[
Y=y+\tfrac12-X^2-tX.
\]
代入 \(E\) 的方程 \(Y^2=X^3-X+\tfrac14\) 并展开整理即得 \(\Pi_t(X,y)=0\)。
反之，若 \((X,y)\) 满足 \(\Pi_t(X,y)=0\)，令 \(Y:=y+\tfrac12-X^2-tX\)，则立得 \((X,Y)\in E\) 且 \(y_t(X,Y)=y\)，从而 \(\Pi_t=0\) 的归一化同构于 \(E\)。
又 \(X\) 在 \(O\) 处极点阶为 \(2\)、\(Y\) 为 \(3\)，故 \(y=X^2+Y-\tfrac12\) 在 \(O\) 处极点阶为 \(4\)，而 \(tX\) 仅为 \(2\)；从而 \((y_t)_\infty=4O\)，即 \(\deg(y_t)=4\)。
证毕。
\end{proof}

\begin{proposition}[有限分歧除子的显式生成元与群和约束]\label{prop:xi-terminal-zm-leyang-linear-twist-ramification-divisor}
对任意 \(t\in\QQ\)，设
\[
f_t(X,Y):=(4X+2t)Y+3X^2-1\in \QQ(E)^\times.
\]
则 \(y_t:E\to\PP^1\) 的全部有限分歧点（按重数计）恰为 \(f_t\) 的零点，并且在 \(E\) 上有主除子恒等式
\[
\operatorname{div}(f_t)=\sum_{i=1}^5 P_i(t)-5O,
\]
其中 \(P_1(t),\dots,P_5(t)\in E(\overline{\QQ})\) 为 \(y_t\) 的五个有限分歧点（总重数 \(5\)）。
特别地，在 \(E\) 的群结构下有
\[
P_1(t)+\cdots+P_5(t)=O.
\]
\end{proposition}
\begin{proof}
对 \(E:Y^2=X^3-X+\tfrac14\) 微分得
\[
2Y\,\dd Y=(3X^2-1)\,\dd X.
\]
因此
\[
\dd y_t=\dd(X^2+tX+Y-\tfrac12)
=\Bigl(2X+t+\frac{3X^2-1}{2Y}\Bigr)\dd X
=\frac{(4X+2t)Y+3X^2-1}{2Y}\,\dd X.
\]
令 \(\omega:=\dd X/(2Y)\)，则 \(\omega\) 为 \(E\) 上处处正则且无零极的不变微分，故
\(
\operatorname{div}(\dd y_t)=\operatorname{div}(f_t)
\)。
另一方面，\(\dd y_t=y_t^*(\dd u)\)（其中 \(u\) 为 \(\PP^1\) 上仿射坐标），且 \(\operatorname{div}(\dd u)=-2\infty\)。
由微分拉回公式得
\[
\operatorname{div}(\dd y_t)=R_t-2(y_t)^*\infty.
\]
由定理 \ref{thm:xi-terminal-zm-leyang-linear-twist-quartic-family}，\((y_t)_\infty=4O\)，故 \((y_t)^*\infty=4O\)。
并且 \(y_t\) 在 \(O\) 处全分歧，分歧指数为 \(4\)，从而 \(R_t\) 在 \(O\) 处系数为 \(3\)。
由 Hurwitz 公式（\(g(E)=1\)、\(\deg(y_t)=4\)）可得剩余有限分歧权重总和为 \(5\)，因而存在五个有限分歧点 \(P_1(t),\dots,P_5(t)\)（按重数计）使
\[
R_t=3O+\sum_{i=1}^5 P_i(t).
\]
于是
\[
\operatorname{div}(\dd y_t)=\Bigl(3O+\sum_{i=1}^5 P_i(t)\Bigr)-8O=\sum_{i=1}^5 P_i(t)-5O,
\]
从而 \(\operatorname{div}(f_t)=\sum_{i=1}^5 P_i(t)-5O\)。
由于右侧为主除子，其在 \(\operatorname{Pic}^0(E)\cong E\) 中类为零，遂得群和约束 \(\sum_i P_i(t)=O\)。
证毕。
\end{proof}

\begin{theorem}[判别式曲线的椭圆化：\texorpdfstring{$\Disc_X(\Pi_t)=0$}{Disc\_X(Pi\_t)=0} 的归一化]\label{thm:xi-terminal-zm-leyang-linear-twist-discriminant-elliptification}
\leanverified{paper\_xi\_terminal\_zm\_leyang\_linear\_twist\_discriminant\_elliptification}
令
\[
D(t,y):=\Disc_{X}\!\bigl(\Pi_t(X,y)\bigr)\in \ZZ[t,y],
\]
并令平面曲线 \(\mathcal D\subset\mathbb A^2_{t,y}\) 由 \(D(t,y)=0\) 给出。
则 \(D(t,y)\) 在 \(\QQ[t,y]\) 中不可约，且 \(\mathcal D\) 的归一化为一条椭圆曲线，并与 \(E\) 双有理等价。
更精确地，存在显式双有理互逆映射 \(\Phi:E\dashrightarrow\mathcal D\)、\(\Psi:\mathcal D\dashrightarrow E\)，其中
\[
\Phi(X,Y)=\Bigl(t=\frac{1-3X^2-4XY}{2Y},\ \ y=X^2+tX+Y-\tfrac12\Bigr),
\]
并且 \(\Psi\) 可由重根消元给出：令
\begin{align*}
A_0(t,y)=&\,4t^5-t^4+28t^3y+10t^3-32t^2y-40t^2+48ty^2+38ty+31t+2y^2-34y-4,\\
C(t,y)=&\,-8t^4y-4t^4+2t^3y+3t^3-48t^2y^2-32t^2y+14t^2+50ty^2+22ty-21t\\
&\,-64y^3-41y^2+25y+8,
\end{align*}
则在开集 \(A_0(t,y)\neq 0\) 上
\[
\Psi(t,y)=\Bigl(X=\frac{-C(t,y)}{2A_0(t,y)},\ \ Y=y+\tfrac12-X^2-tX\Bigr).
\]
此外，\(D(t,y)\) 具有如下闭式展开：
\begin{align*}
D(t,y)=\,&-4 t^7 +16 t^6 y +9 t^6 -48 t^5 y -24 t^5 -16 t^4 y^3 +136 t^4 y^2 +180 t^4 y +90 t^4\\
&+4 t^3 y^3 -318 t^3 y^2 -462 t^3 y -156 t^3 -128 t^2 y^4 +256 t^2 y^3 +640 t^2 y^2 +528 t^2 y +117 t^2\\
&+144 t y^4 -112 t y^3 -618 t y^2 -246 t y -32 t -256 y^5 -155 y^4 +246 y^3 +133 y^2 +32 y .
\end{align*}
\end{theorem}
\begin{proof}
由判别式的定义，\(D(t,y)=0\) 当且仅当存在 \(X\) 使 \(\Pi_t(X,y)=\partial_X\Pi_t(X,y)=0\)，亦即 \(\Pi_t(\cdot,y)\) 在某点出现重根。
对任意 \((X,Y)\in E\) 且 \(Y\neq 0\)，置
\[
t=\frac{1-3X^2-4XY}{2Y}.
\]
则
\[
f_t(X,Y)=(4X+2t)Y+3X^2-1=0,
\]
从而由命题 \ref{prop:xi-terminal-zm-leyang-linear-twist-ramification-divisor} 可知 \(X\) 为关于 \(y\) 的临界根。
再令 \(y=X^2+tX+Y-\tfrac12\)，即有 \(\Pi_t(X,y)=0\)，故 \(\Phi(E)\subset \mathcal D\)。
\par
反向构造中，取 \((t,y)\in\mathcal D\) 的一般点，则 \(\Pi_t(X,y)\) 具有唯一的二重根 \(\xi\)。
对多项式对 \((\Pi_t,\partial_X\Pi_t)\) 作一次子结式计算可得线性子结式
\(
S_1(X)=2A_0(t,y)\,X+C(t,y)
\)，其零点即 \(\xi\)。
因此 \(X=-C/(2A_0)\)。
令 \(Y:=y+\tfrac12-X^2-tX\)，则由 \(\Pi_t(X,y)=0\) 立刻推出 \(Y^2=X^3-X+\tfrac14\)，从而 \(\Psi(t,y)\in E\)。
直接检验可得 \(\Psi\circ\Phi=\mathrm{id}\) 与 \(\Phi\circ\Psi=\mathrm{id}\) 在各自稠密开集上成立，从而 \(\mathcal D\) 的函数域与 \(\QQ(E)\) 同构；于是 \(D(t,y)\) 不可约且 \(\mathcal D\) 的归一化同构于 \(E\)。
闭式展开由 \(D(t,y)=\Disc_X(\Pi_t)\) 的显式计算得到。
证毕。
\end{proof}

\begin{theorem}[分歧参数映射的满对称单值化：\texorpdfstring{$\Gal\cong S_5$}{Gal = S5}]\label{thm:xi-terminal-zm-leyang-linear-twist-parameter-map-s5}
\leanverified{paper\_xi\_terminal\_zm\_leyang\_linear\_twist\_parameter\_map\_s5}
令
\[
\tau:=\frac{1-3X^2-4XY}{2Y}\in\QQ(E),
\]
则 \(\tau:E\to\PP^1_t\) 为次数 \(5\) 的非恒等覆盖。
其函数域扩张 \(\QQ(E)/\QQ(t)\) 的伽罗瓦闭包之伽罗瓦群同构于 \(S_5\)。
等价地，\(X\) 在 \(\QQ(t)\) 上满足五次方程
\[
P(X,t)=16X^5+(16t-9)X^4+(4t^2-16)X^3+(-16t+10)X^2+(-4t^2+4t)X+(t^2-1)=0,
\]
且
\[
\Gal\bigl(P(X,t)/\QQ(t)\bigr)\cong S_5.
\]
\end{theorem}
\begin{proof}
由 \(\tau=(1-3X^2-4XY)/(2Y)\) 解得
\[
Y=\frac{1-3X^2}{2(2X+t)}.
\]
代入 \(E:Y^2=X^3-X+\tfrac14\) 并清分母即得 \(P(X,t)=0\)，从而 \([\QQ(E):\QQ(t)]\le 5\)。
观察 \(\tau\) 的极除子：在 \(O\) 处 \(\operatorname{ord}_O(X)=-2\)、\(\operatorname{ord}_O(Y)=-3\)，故 \(\operatorname{ord}_O(\tau)=-2\)；
在 \(Y=0\) 的三点处，分母为一阶零而分子与之互素，故各给出一阶极点。
因此 \(\tau^*\infty\) 的次数为 \(2+3=5\)，从而 \(\deg(\tau)=5\) 且 \([\QQ(E):\QQ(t)]=5\)。
\par
为判定伽罗瓦群，取特化 \(t=2\) 得
\[
P(X,2)=16X^5+23X^4-22X^2-8X+3.
\]
其模 \(5\) 仍为不可约五次多项式，故伽罗瓦群 \(G\subset S_5\) 传递且含有 \(5\)-循环。
并且其判别式
\[
\Delta(P(X,2))=-2^{16}\cdot 3^3\cdot 23\cdot 37^2
\]
非平方，故 \(G\nsubseteq A_5\)。
又在素数 \(83\) 下有分解型
\[
P(X,2)\equiv 16\,(X^2-7X+28)\,(X^3+24X^2-26X-15)\pmod{83},
\]
对应 Frobenius 循环型 \((2)(3)\)，从而 \(G\) 含有阶 \(3\) 元。
由含 \(5\)-循环的传递子群分类可知 \(G\in\{C_5,D_5,F_{20},A_5,S_5\}\)；
其中 \(C_5,D_5,A_5\subset A_5\) 已由判别式非平方排除，\(F_{20}\) 不含阶 \(3\) 元亦被排除，故 \(G=S_5\)。
最后由特化同态的基本性质，\(\Gal(P(X,2)/\QQ)\) 为 \(\Gal(P(X,t)/\QQ(t))\) 的子群；既然前者已为 \(S_5\)，则后者必为 \(S_5\)。
证毕。
\end{proof}

\begin{corollary}[判别式五次核的不可解性]\label{cor:xi-terminal-zm-leyang-linear-twist-discriminant-quintic-unsolvable}
\leanverified{paper\_xi\_terminal\_zm\_leyang\_linear\_twist\_discriminant\_quintic\_unsolvable}
将 \(D(t,y)\) 视为 \(\QQ(t)\) 上关于 \(y\) 的五次多项式，则其伽罗瓦闭包的伽罗瓦群同构于 \(S_5\)。
特别地，\(D(t,y)\) 在 \(\QQ(t)\) 上一般不可由根式求解。
\end{corollary}
\begin{proof}
由定理 \ref{thm:xi-terminal-zm-leyang-linear-twist-discriminant-elliptification}，曲线 \(\mathcal D:D(t,y)=0\) 的归一化同构于 \(E\)，而投影 \(\pi_t:\mathcal D\to\PP^1_t\) 在函数域层面即为 \(\tau:E\to\PP^1_t\)。
故 \(\QQ(\mathcal D)=\QQ(E)\)，并且 \(\QQ(E)/\QQ(t)\) 的伽罗瓦闭包群为 \(S_5\)（定理 \ref{thm:xi-terminal-zm-leyang-linear-twist-parameter-map-s5}）。
因此 \(D(t,y)\) 作为生成元的最小多项式，其伽罗瓦群亦为 \(S_5\)。
证毕。
\end{proof}

\begin{corollary}[无穷多个有理参数产生有理分歧值]\label{cor:xi-terminal-zm-leyang-linear-twist-infinitely-many-rational-branch-values}
若 \(E(\QQ)\) 为无限群（例如该曲线为正秩曲线，见 \cite{LMFDBEllipticCurve37a1}），则存在无穷多个 \(t\in\QQ\) 使得 \(D(t,y)\in\QQ[y]\) 具有有理根；
等价地，存在无穷多个 \(t\in\QQ\) 使 \(y_t:E\to\PP^1\) 至少包含一个有理分歧值。
\end{corollary}
\begin{proof}
取任意 \(P=(X,Y)\in E(\QQ)\) 且 \(Y\neq 0\)，置 \(t=\tau(P)\in\QQ\)，并令 \(y=X^2+tX+Y-\tfrac12\in\QQ\)。
由定理 \ref{thm:xi-terminal-zm-leyang-linear-twist-discriminant-elliptification} 的映射 \(\Phi\)，得到 \((t,y)\in\mathcal D(\QQ)\)，即 \(D(t,y)=0\)。
由于 \(E(\QQ)\) 无穷且仅有限多个点满足 \(Y=0\)，从而得到无穷多个这样的有理 \(t\)。
证毕。
\end{proof}


\begin{corollary}[判别式二次分辨率在椭圆归一化上的基变：属 \(6\) 的双覆盖]\label{cor:xi-terminal-zm-leyang-discriminant-square-root-basechange-genus6}
沿用定理 \ref{thm:xi-terminal-zm-leyang-elliptic-structure} 的椭圆曲线 \(E_{\min}\) 与次数 \(4\) 覆盖 \(\pi_y:E_{\min}\to\PP^{1}_{y}\)。
令
\[
D:=-y(y-1)P_{\mathrm{LY}}(y)\in \QQ(E_{\min})^{\times}.
\]
记 \(C_{E}\) 为函数域扩张
\(
\QQ(C_{E})=\QQ(E_{\min})(\sqrt{D})
\)
所对应的光滑射影曲线，并令 \(\pi:C_{E}\to E_{\min}\) 为相应的次数 \(2\) 覆盖。
则：
\begin{enumerate}
  \item \(\pi\) 的分歧支撑由 \(D\) 的奇阶零极点精确刻画；并且 \(\pi\) 在无穷远点 \(O\in E_{\min}\) 上方不分歧；
  \item \(\pi\) 的分歧点数为 \(10\)，其支撑由下述点组成：
  \begin{itemize}
    \item 纤维 \(y^{-1}(0)\) 中除去唯一简单分歧点之外的两点；
    \item 纤维 \(y^{-1}(1)\) 中除去唯一简单分歧点之外的两点；
    \item 对每个满足 \(P_{\mathrm{LY}}(y_0)=0\) 的分歧值 \(y_0\)，纤维 \(y^{-1}(y_0)\) 中除去唯一简单分歧点之外的两点（共 \(3\times 2\) 点）。
  \end{itemize}
  \item \(C_{E}\) 的亏格为 \(g(C_{E})=6\)。
\end{enumerate}
\end{corollary}
\begin{proof}
由 \((y)_\infty=4O\) 得 \(\mathrm{ord}_{O}(y)=-4\)，从而 \(\mathrm{ord}_{O}(y-1)=-4\) 且 \(\mathrm{ord}_{O}(P_{\mathrm{LY}}(y))=-12\)，故
\(
\mathrm{ord}_{O}(D)=-4-4-12=-20
\)
为偶数，\(\pi\) 在 \(O\) 上方不分歧。
\par
对有限点，定理 \ref{thm:xi-terminal-zm-leyang-elliptic-structure} 表明：\(y=0\)、\(y=1\) 以及 \(P_{\mathrm{LY}}(y)=0\) 的三根均为 \(\pi_y\) 的有限分歧值，且每个有限分歧值之上恰出现一处简单分歧点（\(\mathrm{ord}(y-y_0)=2\)），其余两点为非分歧点（\(\mathrm{ord}(y-y_0)=1\)）。
因此在每条分歧纤维中，因子 \(y-y_0\) 的奇阶零点恰为“除去分歧点之外的两点”，从而 \(D\) 的奇阶零点总数为
\(
2+2+3\cdot 2=10
\)，并与陈述一致。
\par
对次数 \(2\) 覆盖应用 Riemann--Hurwitz 公式：
\[
2g(C_{E})-2
=2(2g(E_{\min})-2)+10
=10,
\]
由于 \(g(E_{\min})=1\)，得到 \(g(C_{E})=6\)。证毕。
\end{proof}
\begin{corollary}[有理特化的谱可约性：由椭圆曲线像集刻画]\label{cor:xi-terminal-zm-pi-rational-root-specialization-elliptic-image}
\leanverified{paper\_xi\_terminal\_zm\_pi\_rational\_root\_specialization\_elliptic\_image}
设 \(y_0\in\QQ\)。则以下命题等价：
\begin{enumerate}
  \item 四次多项式 \(\Pi(\lambda,y_0)\in\QQ[\lambda]\) 在 \(\QQ\) 上存在有理根；
  \item 存在 \(P\in E_{\min}(\QQ)\) 使得 \(\pi_y(P)=y_0\)，并且该有理根可取为 \(\lambda_0=x(P)\in\QQ\)。
\end{enumerate}
\end{corollary}
\begin{proof}
由定理 \ref{thm:xi-terminal-zm-leyang-elliptic-structure} 的仿射同构 \(\mathcal C\cong E_{\min}\) 直接得到。
若 \(\Pi(\lambda_0,y_0)=0\) 且 \(\lambda_0\in\QQ\)，取 \(u_0:=y_0-\lambda_0^2\)，则 \((x,u)=(\lambda_0,u_0)\in E_{\min}(\QQ)\) 且 \(\pi_y(x,u)=y_0\)。
反之，若 \(P\in E_{\min}(\QQ)\) 满足 \(\pi_y(P)=y_0\)，则 \((\lambda,y)=(x(P),\pi_y(P))\) 给出 \(\Pi(\lambda,y)=0\) 的一组有理解，从而 \(\Pi(\lambda,y_0)\) 存在有理根。证毕。
\end{proof}

\begin{corollary}[有理特化的异常集合之薄集刻画：有理根与符号二次分辨率]\label{cor:xi-terminal-zm-specialization-thin-exceptional-geometry}
沿用定理 \ref{thm:xi-terminal-zm-leyang-elliptic-structure} 的四次谱多项式 \(\Pi(\lambda,y)\) 及 Lee--Yang 三次因子
\[
P_{\mathrm{LY}}(y):=256y^{3}+411y^{2}+165y+32.
\]
令判别式二次脊曲线为
\[
C:\quad \zeta^{2}=-y(y-1)P_{\mathrm{LY}}(y).
\]
则对任意 \(y_0\in\QQ\)：
\begin{enumerate}
  \item 若 \(\Pi(\lambda,y_0)\) 在 \(\QQ[\lambda]\) 上存在有理根，则 \(y_0\in \pi_y\bigl(E_{\min}(\QQ)\bigr)\)（推论 \ref{cor:xi-terminal-zm-pi-rational-root-specialization-elliptic-image}）。
  \item 设 \(\Pi(\lambda,y_0)\) 在 \(\QQ[\lambda]\) 上可分且不可约。
  则
  \[
  \mathrm{Gal}\bigl(\Pi(\cdot,y_0)/\QQ\bigr)\subseteq A_{4}
  \quad\Longleftrightarrow\quad
  -y_0(y_0-1)P_{\mathrm{LY}}(y_0)\in \QQ^{\times 2}
  \quad\Longleftrightarrow\quad
  \exists\,\zeta_0\in\QQ\ \text{使}\ (y_0,\zeta_0)\in C(\QQ).
  \]
\end{enumerate}
因此，“有理可约性”与“符号二次分辨率退化（\(A_4\) 型）”两类异常机制分别被 \(\pi_y\bigl(E_{\min}(\QQ)\bigr)\) 与 \(C(\QQ)\) 的 \(y\)-投影显式几何化；
并且由 Hilbert 不可约定理可知：除去一个薄集参数族后，仍存在无穷多个 \(y_0\in\QQ\) 使得 \(\Pi(\lambda,y_0)\) 在 \(\QQ[\lambda]\) 上不可约且其伽罗瓦群为 \(S_4\)。
\end{corollary}
\begin{proof}
第一条即推论 \ref{cor:xi-terminal-zm-pi-rational-root-specialization-elliptic-image}。
对第二条，当 \(\Pi(\lambda,y_0)\) 可分且不可约时，其伽罗瓦群落入 \(A_4\) 当且仅当其判别式在 \(\QQ\) 中为平方。
而定理 \ref{thm:xi-terminal-zm-leyang-elliptic-structure} 已给出
\(
\mathrm{Disc}_{\lambda}\bigl(\Pi(\lambda,y)\bigr)=-y(y-1)P_{\mathrm{LY}}(y)
\)，
故得到与曲线 \(C\) 的等价刻画。
最后的“薄集外仍有无穷多个满 \(S_4\) 特化”由该正则 \(S_4\)-扩张的 Hilbert 不可约性直接推出。证毕。
\end{proof}

\begin{corollary}[Lee--Yang 四次椭圆覆盖的本原性]\label{cor:xi-terminal-zm-leyang-cover-primitive}
\leanverified{paper\_xi\_terminal\_zm\_leyang\_cover\_primitive}
设 \(\pi_y:E_{\min}\to\mathbb P^1_y\) 为定理 \ref{thm:xi-terminal-zm-leyang-elliptic-structure} 中的次数 \(4\) 态射。
则不存在曲线 \(X\) 与非平凡分解
\[
E_{\min}\xrightarrow{\deg 2}X\xrightarrow{\deg 2}\mathbb P^1_y
\]
使 \(\pi_y\) 等于上述复合。
\end{corollary}
\begin{proof}
若存在上述分解，则 \(\pi_y\) 的几何单值群在四张纸上保留一个 \(2+2\) 的分块结构，从而为不本原子群。
但定理 \ref{thm:xi-terminal-zm-leyang-elliptic-structure} 给出 \(\mathrm{Mon}(\pi_y)\cong S_4\)，而 \(S_4\) 在四点作用下为本原群，矛盾。证毕。
\end{proof}

\begin{corollary}[分歧除子的主化与有限分歧点的群律守恒]\label{cor:xi-terminal-zm-leyang-ramification-sum-zero}
\leanverified{paper\_xi\_terminal\_zm\_leyang\_ramification\_sum\_zero}
沿用定理 \ref{thm:xi-terminal-zm-leyang-elliptic-structure} 的记号，记 \(O\) 为 \(E_{\min}\) 的无穷远点，并令 \(P_0,P_1,P_2,P_3,P_4\) 为 \(\pi_y\) 的五个有限分歧点（即 \(\dd y=0\) 的有限零点）。
则在 \(\mathrm{Pic}^0(E_{\min})\cong E_{\min}\) 中有
\[
[P_0]+[P_1]+[P_2]+[P_3]+[P_4]-5[O]\sim 0,
\]
等价地，在 \(E_{\min}(\overline{\QQ})\) 的群律下满足
\[
P_0+P_1+P_2+P_3+P_4=O.
\]
并且在 \(E_{\min}:u^2+u=x^3-x\) 的 \((x,u)\)-坐标下有
\[
P_0+P_1=\Bigl(\frac14,-\frac38\Bigr),\qquad
P_2+P_3+P_4=\Bigl(\frac14,-\frac58\Bigr)\in E_{\min}(\QQ).
\]
\end{corollary}
\begin{proof}
由定理 \ref{thm:xi-terminal-zm-leyang-elliptic-structure} 知 \(y\) 在 \(O\) 处有唯一四阶极点，故 \((y)_\infty=4O\)。
设 \(R\) 为 \(\pi_y\) 的分歧除子，则 Hurwitz 公式给出
\(
K_{E_{\min}}\sim \pi_y^*K_{\mathbb P^1}+R
\)。
由于 \(K_{E_{\min}}\sim 0\) 且 \(K_{\mathbb P^1}\sim-2\infty\)，得到
\(
R\sim 2(y)_\infty=8O
\)。
另一方面，\(O\) 为全分歧点（指数 \(4\)），故 \(R\) 在 \(O\) 处系数为 \(3\)；其余分歧点为五个有限单分歧点 \(P_0,\dots,P_4\)。
因此 \(R=3O+\sum_{i=0}^{4}P_i\)，代入 \(R\sim 8O\) 即得 \(\sum_{i=0}^{4}P_i-5O\sim 0\)。
在 \(\mathrm{Pic}^0(E_{\min})\cong E_{\min}\) 下该类对应群律和，故 \(P_0+P_1+P_2+P_3+P_4=O\)。
\par
最后，把 \(E_{\min}\) 配方为 \(Y^2=X^3-X+\tfrac14\)（其中 \(X=x,\ Y=u+\tfrac12\)），则
\[
P_0=(1,-\tfrac12),\qquad P_1=(-1,\tfrac12),
\]
由加法公式得 \(P_0+P_1=(\tfrac14,\tfrac18)\)，变换回 \(u\)-坐标即为 \((\tfrac14,-\tfrac38)\)。
再由 \(P_2+P_3+P_4=-(P_0+P_1)\) 与 \(u\mapsto -u-1\) 的取负法则即得 \((\tfrac14,-\tfrac58)\)。证毕。
\end{proof}

\begin{corollary}[倍点拉回下的分歧稳定性：分歧像不生新参数且分歧点按 \texorpdfstring{$E_{\min}[2^n]$}{E[2^n]} 平移复制]\label{cor:xi-terminal-zm-leyang-ramification-pullback-2n}
对整数 \(n\ge 0\) 令 \([2^n]:E_{\min}\to E_{\min}\) 为倍点同态，并定义
\[
y_n:=y\circ[2^n]\in \QQ(E_{\min}).
\]
则 \(y_n:E_{\min}\to\PP^1_y\) 的分歧像集合与 \(y\) 完全一致，仍为
\[
\{\infty,0,1\}\ \cup\ \{P_{\mathrm{LY}}(y)=0\}.
\]
更精确地，若记 \(R(y)\) 为 \(y\) 的分歧除子，则有严格恒等式
\[
R(y_n)=[2^n]^*R(y)
=3\sum_{Q\in E_{\min}[2^n]}[Q]\;+\;\sum_{Q\in E_{\min}[2^n]}\ \sum_{i=0}^{4}[P_i+Q].
\]
因此 \(y_n\) 的五个有限分歧点在 \([2^n]\) 的拉回下被平移复制为 \(5\cdot 4^n\) 个简单分歧点，而无穷远处的全分歧点被复制为 \(4^n\) 个全分歧点。
\end{corollary}
\begin{proof}
在特征 \(0\) 下倍点同态 \([2^n]\) 处处非分歧（\'etale），故对任意有理函数 \(f\) 有 \(R(f\circ[2^n])=[2^n]^*R(f)\)。
将结论 \ref{cor:xi-terminal-zm-leyang-ramification-sum-zero} 的 \(R(y)=3O+\sum_{i=0}^{4}P_i\) 代入并展开：
由于 \([2^n]^{-1}(O)=E_{\min}[2^n]\)。
对每个 \(i\)，任选一点 \(P_{i,n}\in[2^n]^{-1}(P_i)\)，则 \([2^n]^{-1}(P_i)=P_{i,n}+E_{\min}[2^n]\) 为 \(E_{\min}[2^n]\) 的平移余类；
将 \([2^n]^*[P_i]=\sum_{Q\in E_{\min}[2^n]}[P_{i,n}+Q]\) 代入并对 \(i\) 求和，即得所示除子等式（以余类记号书写时可简记为 \(\sum_{Q\in E_{\min}[2^n]}[P_i+Q]\)）。
分歧像陈述由“分歧点在 \'etale 拉回下取原像，而分歧值为其像”直接推出。证毕。
\end{proof}

\begin{corollary}[dyadic 分歧点逆极限的 Tate 扭子分解]\label{cor:xi-terminal-zm-leyang-branchpoint-tate-torsor-decomposition}
沿用推论 \ref{cor:xi-terminal-zm-leyang-ramification-pullback-2n} 的记号。
对 \(i\in\{0,1,2,3,4\}\) 与 \(n\ge 0\) 定义分歧点纤维
\[
\Pi_{i,n}:=[2^n]^{-1}(P_i)\subset E_{\min}(\overline{\QQ}),
\qquad
\Pi_{i,\infty}:=\varprojlim_n\bigl(\Pi_{i,n},[2]\bigr),
\qquad
\Pi_\infty:=\bigsqcup_{i=0}^{4}\Pi_{i,\infty}.
\]
则对每个 \(i\)，\(\Pi_{i,\infty}\) 为 \(T_2(E_{\min})=\varprojlim_n E_{\min}[2^n]\) 的主齐次空间（torsor）。
并且一旦选取一条相容提升 \(\mathbf P_i=(P_{i,n})_{n\ge 0}\in \Pi_{i,\infty}\)（即 \([2]P_{i,n+1}=P_{i,n}\) 且 \(P_{i,0}=P_i\)），便得到一族自然同构
\[
\Pi_{i,n}\ \cong\ P_{i,n}+E_{\min}[2^n],
\qquad
\Pi_{i,\infty}\ \cong\ \mathbf P_i+T_2(E_{\min}),
\]
其中右端表示在相应层次上的平移余类/扭子模型。
\end{corollary}

\begin{proof}
对每个 \(n\)，集合 \(\Pi_{i,n}\) 非空且为 \(E_{\min}[2^n]\) 的平移余类：任取 \(P_{i,n}\in\Pi_{i,n}\)，则
\[
\Pi_{i,n}=P_{i,n}+E_{\min}[2^n],
\]
且平移作用 \(E_{\min}[2^n]\curvearrowright \Pi_{i,n}\) 自由且传递。
结构映射 \([2]:\Pi_{i,n+1}\to\Pi_{i,n}\) 为满射并与该作用相容：
若 \(Q\in E_{\min}[2^{n+1}]\) 则 \([2](x+Q)=[2]x+[2]Q\)。
因此逆极限 \(\Pi_{i,\infty}=\varprojlim_n(\Pi_{i,n},[2])\) 上诱导出 \(T_2(E_{\min})=\varprojlim_n E_{\min}[2^n]\) 的自然平移作用，且该作用自由且传递，从而 \(\Pi_{i,\infty}\) 为 \(T_2(E_{\min})\)-torsor。
\par
若进一步给定一条相容提升 \(\mathbf P_i=(P_{i,n})_{n\ge 0}\in\Pi_{i,\infty}\)，则逐层平移给出同构
\[
E_{\min}[2^n]\stackrel{\sim}{\to}\Pi_{i,n},\quad Q\mapsto P_{i,n}+Q,
\qquad
T_2(E_{\min})\stackrel{\sim}{\to}\Pi_{i,\infty},\quad \mathbf Q\mapsto \mathbf P_i+\mathbf Q,
\]
从而得到所述“余类/扭子”模型表示。
\end{proof}

\begin{corollary}[有限 Lee--Yang 分歧塔的正则四叉地址结构]\label{cor:xi-terminal-zm-leyang-finite-branch-regular-4ary-address}
\leanverified{paper\_xi\_terminal\_zm\_leyang\_finite\_branch\_regular\_4ary\_address}
沿用推论 \ref{cor:xi-terminal-zm-leyang-branchpoint-tate-torsor-decomposition} 的记号，并定义
\[
\Pi_n^{\mathrm{fin}}:=\bigsqcup_{i=0}^{4}\Pi_{i,n},
\qquad
\Pi_\infty^{\mathrm{fin}}:=\varprojlim_n \Pi_n^{\mathrm{fin}}
=\bigsqcup_{i=0}^{4}\Pi_{i,\infty}.
\]
则成立：
\begin{enumerate}
  \item 对每个 \(i\in\{0,1,2,3,4\}\) 与每个 \(n\ge 0\)，映射
  \[
  [2]:\Pi_{i,n+1}\longrightarrow \Pi_{i,n}
  \]
  是正则四重覆盖；更精确地，每个纤维在平移作用下都是 \(E_{\min}[2]\) 的主齐次空间。
  \item 更一般地，对任意 \(r\ge 1\)，映射
  \[
  [2^r]:\Pi_{i,n+r}\longrightarrow \Pi_{i,n}
  \]
  是正则 \(4^r\)-重覆盖；其每个纤维在平移作用下都是 \(E_{\min}[2^r]\) 的主齐次空间。
  \item 一旦对每个 \(i\) 选取一条相容提升 \(\mathbf P_i\in\Pi_{i,\infty}\)，便得到同构
  \[
  \Pi_\infty^{\mathrm{fin}}
  \cong
  \bigsqcup_{i=0}^{4}\bigl(\mathbf P_i+T_2(E_{\min})\bigr)
  \cong
  \bigsqcup_{i=0}^{4}T_2(E_{\min}),
  \]
  其中第二个同构依赖于所选相容提升。
\end{enumerate}
因此，有限 Lee--Yang 分歧塔在 dyadic 方向上由五片彼此平行的 \(2\)-进二维地址片组成。
\end{corollary}
\begin{proof}
由推论 \ref{cor:xi-terminal-zm-leyang-branchpoint-tate-torsor-decomposition}，对每个 \(i\) 与 \(n\) 都可写成
\[
\Pi_{i,n}=P_{i,n}+E_{\min}[2^n]
\]
的平移余类。
若 \(x\in\Pi_{i,n}\)，任选 \(\tilde x\in\Pi_{i,n+1}\) 满足 \([2]\tilde x=x\)，则
\[
[2]^{-1}(x)=\tilde x+E_{\min}[2].
\]
由于 \(E_{\min}[2]\cong(\ZZ/2\ZZ)^2\) 的基数为 \(4\)，且其平移作用自由且传递，故 \([2]:\Pi_{i,n+1}\to\Pi_{i,n}\) 为正则四重覆盖。

同理，对任意 \(r\ge 1\)，若 \(x\in\Pi_{i,n}\) 且 \(\tilde x\in\Pi_{i,n+r}\) 满足 \([2^r]\tilde x=x\)，则
\[
[2^r]^{-1}(x)=\tilde x+E_{\min}[2^r].
\]
而
\[
|E_{\min}[2^r]|=4^r,
\]
故 \([2^r]:\Pi_{i,n+r}\to\Pi_{i,n}\) 为正则 \(4^r\)-重覆盖，并且纤维是 \(E_{\min}[2^r]\) 的主齐次空间。

第三条则由推论 \ref{cor:xi-terminal-zm-leyang-branchpoint-tate-torsor-decomposition} 对各个 \(i\) 的逐分支描述直接得到；把五个分支族作不交并即可。
\end{proof}


\paragraph{平行族截断塔上的相干停时、信息增益与 \(4\)-进读出}\label{par:xi-terminal-zm-leyang-coherence-stopping-time}
本段把 Lee--Yang 分歧塔的“平行族并/截断语义”抽象为一个 \(4\)-叉纤维的有限群塔接口，并给出相干停时的精确律、分辨率信息恒等式以及与 \(4\)-进赋值的严格同构。

\begin{theorem}[平行族 \(4\)-叉群塔上的二元相干停时律]\label{thm:xi-coherence-stopping-time-k-parallel-4ary}
\leanverified{paper\_xi\_coherence\_stopping\_time\_k\_parallel\_4ary}
设 \((G_n)_{n\ge 1}\) 为有限阿贝尔群塔，满足 \(|G_n|=4^n\)。
对每个 \(n\ge 1\) 给定满射
\(
\pi_{n+1\to n}:G_{n+1}\twoheadrightarrow G_n
\)
且每个纤维基数恒为 \(4\)。
固定整数 \(k\ge 1\)，并对每个 \(n\ge 1\) 取集合
\[
A_n:=\bigsqcup_{i=1}^{k}A_{n,i},
\qquad
A_{n,i}\cong G_n,
\]
其中每个 \(A_{n,i}\) 是一个 \(G_n\)-挠空间（主齐次空间）。
假设存在兼容的截断族 \(\Pi_{n+1\to n}:A_{n+1}\to A_n\)，使得对每个 \(i\)：
\[
\Pi_{n+1\to n}(A_{n+1,i})\subseteq A_{n,i},
\qquad
\text{且在同一标号 \(i\) 上与 }\pi_{n+1\to n}\text{ 的作用共轭。}
\]
令 \(\Pi_{N\to n}:=\Pi_{n+1\to n}\circ\cdots\circ\Pi_{N\to N-1}\)（\(1\le n\le N\)）。
取 \(X_N,Y_N\) 为 \(A_N\) 上独立均匀随机元，并定义相干停时（带终止约定）
\[
\tau:=\min\{1\le n\le N:\ \Pi_{N\to n}(X_N)\neq \Pi_{N\to n}(Y_N)\},
\qquad
\text{若不存在则置 }\tau:=N+1.
\]
则对任意整数 \(0\le n\le N\) 有精确尾分布
\begin{equation}\label{eq:xi-coherence-stopping-time-tail}
\PP(\tau>n)=\frac{1}{k\,4^n}.
\end{equation}
从而点分布为
\[
\PP(\tau=1)=1-\frac{1}{4k},\qquad
\PP(\tau=n)=\frac{3}{k\,4^n}\ (2\le n\le N),\qquad
\PP(\tau=N+1)=\frac{1}{k\,4^N}.
\]
并且期望停时满足闭式
\begin{equation}\label{eq:xi-coherence-stopping-time-mean}
\EE[\tau]=1+\frac{1}{3k}\Bigl(1-4^{-N}\Bigr).
\end{equation}
\end{theorem}
\begin{proof}
事件 \(\{\tau>n\}\) 等价于 \(\Pi_{N\to n}(X_N)=\Pi_{N\to n}(Y_N)\)。
由 \(A_N\) 的 \(k\) 分量直和结构，该事件必先要求两次抽样落在同一分量（概率 \(1/k\)）。
在该分量内，由截断共轭性与均匀性，\(\Pi_{N\to n}(X_N)\) 在 \(A_{n,i}\) 上均匀；而 \(A_{n,i}\) 作为 \(G_n\)-挠空间与 \(G_n\) 等势，故两次独立均匀抽样相等的概率为 \(1/|G_n|=4^{-n}\)。
相乘得到 \eqref{eq:xi-coherence-stopping-time-tail}。
点分布由差分 \(\PP(\tau=n)=\PP(\tau>n-1)-\PP(\tau>n)\) 给出；期望由
\(
\EE[\tau]=\sum_{n\ge 0}\PP(\tau>n)=1+\sum_{n=1}^{N}(k4^n)^{-1}
\)
求和得到 \eqref{eq:xi-coherence-stopping-time-mean}。证毕。
\end{proof}

\begin{corollary}[Lee--Yang 五平行族的首层区分律与几何尾]\label{cor:xi-coherence-stopping-time-leyang-k5}
\leanverified{paper\_xi\_coherence\_stopping\_time\_leyang\_k5}
在推论 \ref{cor:xi-terminal-zm-leyang-ramification-pullback-2n} 的 Lee--Yang 分歧塔上，
将五个有限分歧点 \(P_0,\dots,P_4\) 的平移族
\(
P_i+E_{\min}[2^n]
\)
视为 \(k=5\) 个互不相交的 \(E_{\min}[2^n]\)-挠分量，并在各分量内用倍点截断诱导兼容的 \(4\)-叉截断塔。
则对任意固定深度 \(N\ge 1\)，相应相干停时 \(\tau\) 满足
\[
\PP(\tau=1)=\frac{19}{20},\qquad
\PP(\tau=n)=\frac{3}{5\cdot 4^{n}}\ (2\le n\le N),\qquad
\PP(\tau=N+1)=\frac{1}{5\cdot 4^{N}},
\]
并且
\[
\EE[\tau]=1+\frac{1}{15}\Bigl(1-4^{-N}\Bigr).
\]
\end{corollary}
\begin{proof}
将定理 \ref{thm:xi-coherence-stopping-time-k-parallel-4ary} 取 \(k=5\) 并用 \(|E_{\min}[2^n]|=4^n\) 代入即可。证毕。
\end{proof}

\begin{proposition}[分辨率读出与信息增益的严格恒等式]\label{prop:xi-coherence-resolution_information_identity}
\leanverified{paper\_xi\_coherence\_resolution\_information\_identity}
在定理 \ref{thm:xi-coherence-stopping-time-k-parallel-4ary} 的设定下，固定 \(N\ge 1\)。
在每个分量 \(A_{n,i}\) 上选择一次性标架同构 \(\theta_{n,i}:A_{n,i}\to G_n\)，并令坐标读出
\[
O_n:=\theta_{n,I_n}\!\bigl(\Pi_{N\to n}(X_N)\bigr)\in G_n,\qquad 1\le n\le N,
\]
其中 \(I_n\) 为 \(\Pi_{N\to n}(X_N)\) 所在分量标号。
则以 \(\log_2\) 为底的 Shannon 量度下，对任意 \(1\le n\le N\) 有
\[
I(X_N;O_n)=H(O_n)=\log_2|G_n|=2n,
\]
并且条件熵满足闭式
\[
H(X_N\mid O_n)=\log_2(k4^{N-n})=\log_2 k+2(N-n).
\]
\end{proposition}
\begin{proof}
由于 \(O_n\) 为 \(X_N\) 的确定函数，故 \(I(X_N;O_n)=H(O_n)\)。
由截断共轭性与分量均匀性，\(O_n\) 在 \(G_n\) 上均匀，故 \(H(O_n)=\log_2|G_n|=2n\)。
另一方面，给定 \(O_n=o\)，其原像由 \(k\) 个分量的均匀纤维组成，每个分量纤维大小为 \(|\ker(\pi_{N\to n})|=4^{N-n}\)，故条件分布为等概率于 \(k4^{N-n}\) 个点上，得到所示条件熵。证毕。
\end{proof}

\begin{theorem}[相干等价类的仿射子立方结构]\label{thm:xi-coherence-affine-subcube}
在 Lee--Yang 分歧塔的任一固定分量 \(P_i+E_{\min}[2^N]\) 内，选定 \(T_2(E_{\min})\) 的 \(\ZZ_2\)-基并据此得到同构
\(
E_{\min}[2^N]\cong(\ZZ/2^N\ZZ)^2
\)
与地址标识
\(
\Gamma_N:P_i+E_{\min}[2^N]\to\{0,1,2,3\}^N
\)
（将每一层的两位二进地址合并为一位 \(4\)-进数字）。
对 \(1\le m\le N\) 定义“分辨率 \(m\)”相干关系为
\[
x\sim_m y
\quad:\!\!\Longleftrightarrow\quad
\Gamma_N(x)\equiv \Gamma_N(y)\pmod{4^m},
\]
其中同余按 \(\sum_{t=0}^{N-1}\Gamma_N(\cdot)_t4^t\in\ZZ/4^N\ZZ\) 的自然解释给出。
则对每个 \(m\)：
\begin{enumerate}
  \item 等价类 \([x]_{\sim_m}\) 在 \(2N\)-维超立方体顶点模型中为一个维数 \(2(N-m)\) 的仿射子立方，故
  \[
  |[x]_{\sim_m}|=4^{N-m};
  \]
  \item 该分量被分割为 \(4^m\) 个两两不交的上述仿射子立方。
\end{enumerate}
\leanverified{paper\_xi\_coherence\_affine\_subcube}
\end{theorem}
\begin{proof}
在地址模型 \(\{0,1,2,3\}^N\) 中，同余 \(\bmod 4^m\) 等价于前 \(m\) 个 \(4\)-进位（即 \(t=0,\dots,m-1\) 的低位数字）一致，而余下 \(N-m\) 个数字可自由取值。
将每个 \(4\)-进位展开为两位二进坐标后，“固定若干坐标而其余坐标自由”即为仿射子立方的定义；
自由坐标数为 \(2(N-m)\)，顶点数为 \(2^{2(N-m)}=4^{N-m}\)，并且可能的固定低位共有 \(4^m\) 种。证毕。
\end{proof}

\begin{theorem}[Gödel \(4\)-进编码与相干停时的 \(4\)-进指数表述]\label{thm:xi-coherence-godel-4adic-valuation}
在定理 \ref{thm:xi-coherence-affine-subcube} 的地址标识下，定义 Gödel 型编码
\[
\mathbf G_N(x):=\sum_{t=0}^{N-1}\Gamma_N(x)_t\,4^t\in \ZZ/4^N\ZZ.
\]
则对任意 \(x,y\) 与任意 \(1\le m\le N\) 有严格等价
\[
x\sim_m y
\quad\Longleftrightarrow\quad
\mathbf G_N(x)\equiv \mathbf G_N(y)\pmod{4^m}.
\]
并且若定义停时
\[
\tau(x,y):=\min\{1\le m\le N:\ x\not\sim_m y\},\qquad
\text{若不存在则置 }\tau(x,y):=N+1,
\]
则
\[
\tau(x,y)=1+\nu_4\!\bigl(\mathbf G_N(x)-\mathbf G_N(y)\bigr),
\]
其中 \(\nu_4(t):=\max\{r\in\{0,1,\dots,N\}: t\equiv 0\ (\mathrm{mod}\ 4^r)\}\) 为截断到深度 \(N\) 的 \(4\)-进指数。
\leanverified{paper\_xi\_coherence\_godel\_4adic\_valuation}
\end{theorem}
\begin{proof}
由 \(\mathbf G_N\) 的定义，模 \(4^m\) 的同余恰等价于低位 \(m\) 个数字一致，即 \(\sim_m\) 的定义。
停时 \(\tau(x,y)\) 是使低位前缀首次失配的层数，等价于“最大相同低位前缀长度”，而该长度正由 \(\nu_4(\mathbf G_N(x)-\mathbf G_N(y))\) 给出。证毕。
\end{proof}

\endinput


\par
\begin{conclusion}[Lee--Yang 谱四次的单参数椭圆化：二次配方与双覆盖模型]\label{con:xi-terminal-zm-leyang-spectral-quartic-a-elliptification}
为刻画谱四次椭圆化桥梁在系数方向上的稳定性，考虑有理参数 \(a\in\QQ\) 的单参数族
\[
\Pi_a(\lambda,y):=\lambda^4-\lambda^3-(ay+1)\lambda^2+\lambda+y(y+1)\in\QQ[a,\lambda,y],
\qquad
\mathcal C_a:=\{(\lambda,y)\in\mathbb A^2:\ \Pi_a(\lambda,y)=0\}.
\]
令
\[
w_a:=2y+1-a\lambda^2,
\qquad
D_a(\lambda):=(a^2-4)\lambda^4+4\lambda^3+(4-2a)\lambda^2-4\lambda+1.
\]
则在 \(\mathcal C_a\) 上恒有平方差恒等式
\[
4\Pi_a(\lambda,y)=w_a^2-D_a(\lambda),
\]
从而函数域满足
\(
\QQ(\mathcal C_a)=\QQ(\lambda,y)=\QQ(\lambda,w_a)
\)，并且
\[
(\lambda,y)\longmapsto(\lambda,w_a)
\]
给出 \(\mathcal C_a\) 与双覆盖模型
\[
\mathcal E_a:\quad w^2=D_a(\lambda)
\]
之间的显式双有理对应，其逆映射由
\(
y=(w_a+a\lambda^2-1)/2
\)
给出。
若
\[
(a-1)\bigl(16a^3-13a^2-78a+43\bigr)\neq 0,
\]
则 \(D_a(\lambda)\) 无重根，故 \(\mathcal E_a\) 的光滑射影闭包为属 \(1\) 的曲线；并且 \(\mathcal C_a\) 的规范化与该椭圆曲线双有理同构。
在 \(a=2\) 的 Lee--Yang 点上，上述恒等式退化为文内既定的椭圆化桥梁
\(
4\Pi=w^2-(4\lambda^3-4\lambda+1)
\)。
\end{conclusion}
\begin{proof}
将 \(\Pi_a(\lambda,y)\) 视为关于 \(y\) 的二次多项式：
\[
\Pi_a(\lambda,y)=y^2+(1-a\lambda^2)y+\bigl(\lambda^4-\lambda^3-\lambda^2+\lambda\bigr).
\]
配方得
\[
(2y+1-a\lambda^2)^2=(1-a\lambda^2)^2-4(\lambda^4-\lambda^3-\lambda^2+\lambda)=D_a(\lambda),
\]
从而 \(4\Pi_a=w_a^2-D_a(\lambda)\)，并推出 \(\QQ(\lambda,y)=\QQ(\lambda,w_a)\) 以及所述互逆代换。
当 \(D_a\) 无重根时，\(\mathcal E_a:w^2=D_a(\lambda)\) 为 \(\PP^1_\lambda\) 的光滑双覆盖，故其光滑射影模型为亏格 \(1\) 曲线；由于 \(\mathcal C_a\) 与 \(\mathcal E_a\) 具有同一函数域，\(\mathcal E_a\) 即给出 \(\mathcal C_a\) 的规范化。
最后，直接计算得到
\(
\mathrm{Disc}_{\lambda}(D_a)=-2^8(a-1)^2(16a^3-13a^2-78a+43)
\)，故所示条件等价于 \(D_a\) 无重根。证毕。
\end{proof}

\begin{conclusion}[归一化族的 Weierstrass 模型、判别式与 \texorpdfstring{$j$}{j}-不变量闭式]\label{con:xi-terminal-zm-leyang-spectral-quartic-a-weierstrass-j}
在结论 \ref{con:xi-terminal-zm-leyang-spectral-quartic-a-elliptification} 的非奇异条件下，\(\mathcal E_a:w^2=D_a(\lambda)\) 的 Jacobian 椭圆曲线可取为短 Weierstrass 模型
\[
W_a:\quad v^2=x^3-432(a^2-a+1)\,x+432(8a^3+15a^2-66a+35).
\]
其判别式与 \(j\)-不变量分别为
\[
\Delta(W_a)=-2^{12}3^{12}(a-1)^2\bigl(16a^3-13a^2-78a+43\bigr),
\qquad
j(W_a)= -\frac{4096\,(a^2-a+1)^3}{(a-1)^2(16a^3-13a^2-78a+43)}.
\]
特别地，当 \(a=2\) 时有 \(j(W_2)=110592/37\)，与定理 \ref{thm:xi-terminal-zm-leyang-elliptic-structure} 的固定椭圆锚点一致。
\end{conclusion}
\begin{proof}
把 \(D_a(\lambda)\) 视为二元四次 \(f_a(\lambda,1)\)，记
\(
f_a(\lambda,z)=A\lambda^4+B\lambda^3 z+C\lambda^2 z^2+D\lambda z^3+Ez^4
\)。
对本族有
\[
(A,B,C,D,E)=(a^2-4,\ 4,\ 4-2a,\ -4,\ 1).
\]
二元四次不变量
\[
I:=12AE-3BD+C^2,\qquad
J:=72ACE+9BCD-27AD^2-27EB^2-2C^3
\]
直接化简得
\[
I(a)=16(a^2-a+1),\qquad J(a)=-16(8a^3+15a^2-66a+35).
\]
取其 Jacobian 的标准短 Weierstrass 形式 \(v^2=x^3-27I(a)x-27J(a)\)，即得到所示 \(W_a\)。
短 Weierstrass 曲线的判别式与 \(j\)-不变量分别为
\(
\Delta=-16(4A^3+27B^2)
\)
与 \(j=c_4^3/\Delta\)（其中 \(c_4=-48A\)），代入 \(A=-432(a^2-a+1)\)、\(B=432(8a^3+15a^2-66a+35)\) 即得闭式。
令 \(a=2\) 立得 \(j(W_2)=110592/37\)。证毕。
\end{proof}

\begin{theorem}[归一化族 \texorpdfstring{$(W_a)$}{(Wa)} 的 \texorpdfstring{$j$}{j}-映射之全分歧谱与 \texorpdfstring{$\QQ(\sqrt{89})$}{Q(sqrt(89))} 分歧值]\label{thm:xi-terminal-zm-leyang-spectral-quartic-a-j-ramification-spectrum}
沿用结论 \ref{con:xi-terminal-zm-leyang-spectral-quartic-a-weierstrass-j} 的 \(W_a\) 与
\[
j(a):=j(W_a)= -\frac{4096\,(a^2-a+1)^3}{(a-1)^2(16a^3-13a^2-78a+43)}\in\QQ(a).
\]
则 \(j:\PP^1_a\to\PP^1_j\) 为次数 \(6\) 的有理覆盖；其全部分歧点（含无穷远处的极点分歧）恰为
\[
\{a:\ a^2-a+1=0\}\ \cup\ \{a:\ 8a^3+15a^2-66a+35=0\}\ \cup\ \{a:\ 2a^2-9a-1=0\}\ \cup\ \{a=1\},
\]
并且对应的分歧值为
\[
j=0,\qquad j=1728,\qquad j=-931\pm 89\sqrt{89},\qquad j=\infty.
\]
更精确地：
\begin{itemize}
  \item 若 \(a^2-a+1=0\)，则 \(j(a)=0\)，且该分歧点的局部分歧指数为 \(3\)；
  \item 若 \(8a^3+15a^2-66a+35=0\)，则 \(j(a)=1728\)，且局部分歧指数为 \(2\)；
  \item 若 \(2a^2-9a-1=0\)，则 \(j(a)=-931\mp 89\sqrt{89}\)，且局部分歧指数为 \(2\)；
  \item 在 \(a=1\) 处，\(j\) 具有二阶极点（从而分歧指数为 \(2\)），分歧值为 \(\infty\)。
\end{itemize}
\end{theorem}
\begin{proof}
由 \(j(a)\) 的显式式子可见：分子次数为 \(6\)、分母次数为 \(5\)，且 \(a\to\infty\) 时 \(j(a)\sim -256\,a\)，故 \(j\) 的次数为 \(6\)。
\par
对 \(j(a)\) 求导并在 \(\QQ(a)\) 中约分，可得
\[
j'(a)= -\frac{4096\,(a^2-a+1)^2\,(2a^2-9a-1)\,(8a^3+15a^2-66a+35)}
{(a-1)^3\,(16a^3-13a^2-78a+43)^2}.
\]
故除去极点结构外，全部有限临界点恰为上述三类多项式零点。
\par
对分歧值：
若 \(a^2-a+1=0\) 则显然 \(j(a)=0\)，且零点阶为 \(3\)，故分歧指数为 \(3\)。
并且存在恒等式
\[
j(a)-1728
= -\frac{64\,(8a^3+15a^2-66a+35)^2}{(a-1)^2(16a^3-13a^2-78a+43)},
\]
从而当 \(8a^3+15a^2-66a+35=0\) 时 \(j(a)=1728\) 且为二重零，故分歧指数为 \(2\)。
若 \(2a^2-9a-1=0\)，则 \(a=(9\pm\sqrt{89})/4\in\QQ(\sqrt{89})\)；将该二次关系代入 \(j(a)\) 并在 \(\QQ(\sqrt{89})\) 中化简得到
\[
j\!\left(\frac{9+\sqrt{89}}{4}\right)=-931-89\sqrt{89},\qquad
j\!\left(\frac{9-\sqrt{89}}{4}\right)=-931+89\sqrt{89},
\]
且由 \(j'(a)\) 在此处为一阶零知分歧指数为 \(2\)。
最后，\(a=1\) 处分母含 \((a-1)^2\) 而分子非零，故 \(j\) 在 \(a=1\) 处为二阶极点并给出分歧值 \(\infty\)。
\par
由 Riemann--Hurwitz 公式对次数 \(6\) 的覆盖得总分歧量为 \(2\cdot 6-2=10\)；
而上述点集的分歧贡献依次为：两点三重零贡献 \(4\)、三点二重分歧贡献 \(3\)、两点二重分歧贡献 \(2\)、以及 \(a=1\) 的二重极点贡献 \(1\)，总计 \(10\)，故无其余分歧点。证毕。
\end{proof}

\begin{corollary}[\texorpdfstring{$j=0$}{j=0} 与 \texorpdfstring{$j=1728$}{j=1728} 的有理参数阻断]\label{cor:xi-terminal-zm-leyang-spectral-quartic-a-j-0-1728-noQ}
\leanverified{paper\_xi\_terminal\_zm\_leyang\_spectral\_quartic\_a\_j\_0\_1728\_noq}
在定理 \ref{thm:xi-terminal-zm-leyang-spectral-quartic-a-j-ramification-spectrum} 的 \(j(a)\) 下，方程 \(j(a)=0\) 与 \(j(a)=1728\) 在 \(\QQ\) 上均无解。
\end{corollary}
\begin{proof}
由 \(j(a)=0\iff a^2-a+1=0\)，其判别式为 \(-3\)，故无 \(\QQ\)-解。
另一方面，由上式的恒等式
\(
j(a)-1728\propto (8a^3+15a^2-66a+35)^2
\)
知 \(j(a)=1728\iff 8a^3+15a^2-66a+35=0\)。
对该三次多项式应用有理根判据即可排除有理根，从而无 \(\QQ\)-解。证毕。
\end{proof}

\begin{proposition}[\texorpdfstring{$\QQ(\sqrt{89})$}{Q(sqrt(89))} 分歧值的最小多项式与判别式]\label{prop:xi-terminal-zm-leyang-spectral-quartic-a-j-fib89-minpoly}
令
\[
J_\pm:=-931\pm 89\sqrt{89}.
\]
则 \(J_\pm\) 的最小多项式为
\[
T^2+1862T+161792\in\ZZ[T],
\]
其判别式为 \(4\cdot 89^3\)；
从而 \(\QQ(J_\pm)=\QQ(\sqrt{89})\)。
\end{proposition}
\begin{proof}
直接计算
\[
(T-J_+)(T-J_-)
=(T+931)^2-(89\sqrt{89})^2
=T^2+1862T+161792,
\]
其判别式为
\(
1862^2-4\cdot 161792=4\cdot 89^3
\)，
故分裂域为 \(\QQ(\sqrt{89})\)。证毕。
\end{proof}
\leanverified{paper\_xi\_terminal\_zm\_leyang\_spectral\_quartic\_a\_j\_fib89\_minpoly}

\begin{conjecture}[\texorpdfstring{$j(a)$}{j(a)} 覆盖的泛型单值化群与唯一实二次分歧子域]\label{conj:xi-terminal-zm-leyang-spectral-quartic-a-j-generic-s6}
设 \(t\) 为超越元，并令 \(P_t(a)\in\QQ(t)[a]\) 为方程 \(j(a)=t\) 的分子化简所得次数 \(6\) 多项式。
猜想其在 \(\QQ(t)\) 上的伽罗瓦群为 \(S_6\)；
并且其分歧域的唯一实二次子域为 \(\QQ(\sqrt{89})\)，对应定理 \ref{thm:xi-terminal-zm-leyang-spectral-quartic-a-j-ramification-spectrum} 的两条非平凡有限分歧值。
\end{conjecture}

\begin{remark}[分歧谱作为 Hurwitz 护照输入与压力--迹桥包的可计算化]\label{rem:xi-terminal-zm-leyang-spectral-quartic-a-j-outlook}
定理 \ref{thm:xi-terminal-zm-leyang-spectral-quartic-a-j-ramification-spectrum} 给出 \(j:\PP^1_a\to\PP^1_j\) 的全分歧谱，因此可把覆盖的分歧循环型写成一份显式 Hurwitz 护照，并据此在 Nielsen 类意义下确定其单值化群与对应的 Hurwitz 空间连通分量。
另一方面，定理 \ref{thm:ostrowski-quadratic-irrational-pressure-matrix-algebraicity} 的有限状态矩阵化把压力函数压缩为主特征值数据，而定理 \ref{thm:fold-groupoid-af-inductive-limit-holographic-trace} 的 AF 极限与迹把纤维推前权重对象化为可配对的迹态。
二者联合为“压力（谱半径）\(\leftrightarrow\) 迹（维数群）”的 \(K\)-理论编码提供一条可计算桥梁：在固定折叠语法与细化映射的前提下，维数群与其迹锥可由有限矩阵数据逐层逼近并产生独立审计接口。
\end{remark}

\begin{conclusion}[\texorpdfstring{$\lambda$}{lambda}-判别式的普适分解与无穷远分歧型突变]\label{con:xi-terminal-zm-leyang-spectral-quartic-a-discriminant-branch}
将 \(\Pi_a(\lambda,y)\) 视为 \(\QQ(a,y)\) 上关于 \(\lambda\) 的四次多项式，则其 \(\lambda\)-判别式在 \(\QQ[a,y]\) 中恒分解为
\[
\mathrm{Disc}_{\lambda}\bigl(\Pi_a(\lambda,y)\bigr)=y\cdot B_a(y),
\]
其中 \(B_a(y)\) 为五次多项式
\[
\begin{aligned}
B_a(y)=&\ (16a^4-128a^2+256)y^5+(16a^4+68a^3-256a^2-400a+768)y^4\\
&+(68a^3+60a^2-800a+661)y^3+(4a^3+188a^2-290a+42)y^2\\
&+(13a^2+110a-139)y+32(a-1).
\end{aligned}
\]
令 \(\widetilde{\mathcal C}_a\) 表示 \(\mathcal C_a\) 的规范化，并令投影 \(\pi_y:\widetilde{\mathcal C}_a\to\PP^1_y\) 由坐标 \(y\) 给出。
若 \(a^2\neq 4\)，则几何纤维 \(\pi_y^{-1}(\infty)\) 由两点 \(P_\infty^\pm\) 构成，并且
\(
\operatorname{ord}_{P_\infty^\pm}(y)=-2
\)，从而 \(\infty\) 处分歧型为 \((2,2)\)；
有限分歧值集合由
\(
\{y=0\}\cup\{B_a(y)=0\}
\)
给出（一般为 \(6\) 个互异点）。
若 \(a^2=4\)，则 \(\pi_y^{-1}(\infty)\) 退化为单点 \(P_\infty\)，且
\(
\operatorname{ord}_{P_\infty}(y)=-4
\)，于是 \(\infty\) 处分歧型为 \((4)\)，并且有限分歧值数目相应降为 \(5\)。
在 \(a=2\) 时，
\[
B_2(y)=-(y-1)\bigl(256y^3+411y^2+165y+32\bigr),
\]
从而与结论 \ref{con:xi-terminal-zm-discriminant-leyang} 的 Lee--Yang 三次分歧因子完全对齐。
\end{conclusion}
\begin{proof}
判别式分解由结果式计算给出：
\(
\mathrm{Disc}_{\lambda}(\Pi_a)=\mathrm{Res}_{\lambda}(\Pi_a,\partial_\lambda\Pi_a)
\)；
在 \(\QQ[a,y]\) 中因式分解即得 \(y\cdot B_a(y)\)。
\par
对无穷远分歧型，使用结论 \ref{con:xi-terminal-zm-leyang-spectral-quartic-a-elliptification} 的双覆盖模型
\(
\mathcal E_a:w^2=D_a(\lambda)
\)
以及
\(
y=(w+a\lambda^2-1)/2
\)。
当 \(a^2\neq 4\) 时，\(\mathcal E_a\) 为四次双覆盖模型；在其光滑射影闭包的两点无穷远处，\(\lambda\) 与 \(w\) 的极点阶分别为 \(1\) 与 \(2\)，故 \(y\) 亦具有二阶极点，从而得到 \(\operatorname{ord}_{P_\infty^\pm}(y)=-2\) 与分歧型 \((2,2)\)。
当 \(a^2=4\) 时，\(\mathcal E_a\) 退化为三次 Weierstrass 模型；在唯一无穷远点处 \(\operatorname{ord}(\lambda)=-2\)、\(\operatorname{ord}(w)=-3\)，故 \(\operatorname{ord}(y)=-4\)，给出分歧型 \((4)\)。
最后代入 \(a=2\) 得到 \(B_2\) 的显式分解。证毕。
\end{proof}

\begin{conclusion}[有理 \texorpdfstring{$2$}{2}-挠点存在性的完全参数化]\label{con:xi-terminal-zm-leyang-spectral-quartic-a-2torsion}
在 \(\Delta(W_a)\neq 0\) 时，短 Weierstrass 曲线 \(W_a/\QQ\) 含非平凡有理 \(2\)-挠点当且仅当存在 \(m\in\QQ\) 使
\[
a=\frac{m^3+36m^2-18576}{(m-12)^2(m+24)}.
\]
在此情形下，
\[
x_2(m):=\frac{24\,(m^3+18m^2-702m-1728)}{(m-12)^2(m+24)}\in\QQ
\]
满足
\[
x_2(m)^3-432(a^2-a+1)x_2(m)+432(8a^3+15a^2-66a+35)=0,
\]
从而 \(T_m:=(x_2(m),0)\in W_a(\QQ)\) 为阶 \(2\) 的点。
此外，双覆盖模型 \(\mathcal E_a:w^2=D_a(\lambda)\) 含有有理分歧点（即存在 \((\lambda_0,0)\in\mathcal E_a(\QQ)\)）当且仅当存在 \(t\in\QQ^\times\setminus\{\pm 1\}\) 使
\[
a=\frac{t^4+4t^3-2t^2-2t+1}{t^4}
\quad\text{或}\quad
a=\frac{t^4-4t^3-2t^2+2t+1}{t^4},
\qquad
\lambda_0=\frac{t^2}{1-t^2}.
\]
\end{conclusion}
\begin{proof}
对 \(W_a:v^2=x^3+Ax+B\) 而言，非平凡 \(2\)-挠点等价于右端三次在 \(\QQ\) 上有根，亦即存在 \(x\in\QQ\) 使
\[
F(a,x):=x^3-432(a^2-a+1)x+432(8a^3+15a^2-66a+35)=0.
\]
该方程在 \((a,x)\)-平面上定义一条三次曲线；直接核验可知其在 \((a,x)=(1,-12)\) 处存在有理奇点（并与 \(\Delta(W_a)\) 的因子 \((a-1)^2\) 对齐），故该三次曲线为有理曲线。
取过该奇点的直线族
\[
x=m(a-1)-12\qquad(m\in\QQ),
\]
代入得
\[
F(a,m(a-1)-12)=(a-1)^2\bigl(a(m^3-432m+3456)-(m^3+36m^2-18576)\bigr).
\]
当 \(a\neq 1\) 时由线性因子解得
\[
a=\frac{m^3+36m^2-18576}{m^3-432m+3456}
=\frac{m^3+36m^2-18576}{(m-12)^2(m+24)},
\]
并由 \(x=m(a-1)-12\) 得到对应根 \(x_2(m)\) 的显式表达。
因此 \(T_m=(x_2(m),0)\) 为有理 \(2\)-挠点；反之，任一满足 \(\Delta(W_a)\neq 0\) 的有理解 \((a,x)\) 均确定斜率 \(m=(x+12)/(a-1)\)，从而落入上述参数化。
\par
最后，\(\mathcal E_a\) 存在有理分歧点当且仅当 \(D_a(\lambda)\) 在 \(\QQ\) 上有根。
将 \(D_a(\lambda)=0\) 视为关于 \(a\) 的二次方程并计算判别式
\(
\Delta_a=16\lambda^5(\lambda-1)^2(\lambda+1)
\)
可知该条件等价于 \(\lambda(\lambda+1)\) 为有理平方；令
\[
\lambda=\frac{t^2}{1-t^2},\qquad \sqrt{\lambda(\lambda+1)}=\frac{t}{1-t^2}
\]
并代入求根公式即得两族 \(a(t)\)，从而得到所示 \((\lambda_0,0)\in\mathcal E_a(\QQ)\)。证毕。
\end{proof}

\begin{conclusion}[分歧五次多项式的泛型最大对称性：\texorpdfstring{$\mathrm{Gal}(B_a/\QQ(a))\cong S_5$}{Gal(Ba/Q(a))=S5}]\label{con:xi-terminal-zm-leyang-spectral-quartic-a-branch-quintic-s5}
设 \(K=\QQ(a)\)。则 \(B_a(y)\in K[y]\) 在 \(K\) 上不可约，且其分裂域的伽罗瓦群同构于 \(S_5\)。
\end{conclusion}
\begin{proof}
取特化 \(a=3\)，得到
\[
B_3(y)=400y^5+396y^4+637y^3+972y^2+308y+64\in\ZZ[y].
\]
其在 \(\FF_7\) 上保持不可约（分解型为 \([5]\)），从而 \(B_3\) 在 \(\QQ\) 上不可约并对应一条 \(5\)-循环。
又在 \(\FF_{19}\) 上分解型为 \([2,3]\)，故其伽罗瓦群含有 \(3\)-循环，从而排除仅含 \(5\)-循环的可解候选（\(C_5,D_5,F_{20}\)）。
另一方面，\(\mathrm{Disc}(B_3)\) 含素因子 \(31\) 的奇次幂，故判别式非平方，从而排除 \(A_5\)；综上
\(
\mathrm{Gal}(B_3/\QQ)\cong S_5
\)。
\par
若 \(B_a\) 在 \(K\) 上可约，则清分母后可得在 \(\QQ[a,y]\) 中的非平凡因式分解，从而除去有限多个退化点后几乎所有有理特化皆可约，这与 \(a=3\) 特化矛盾；故 \(B_a\) 在 \(K\) 上不可约。
并且对除去有限薄集的特化，特化伽罗瓦群嵌入泛型伽罗瓦群；由于存在特化 \(a=3\) 给出 \(S_5\)，故 \(\mathrm{Gal}(B_a/K)=S_5\)。证毕。
\end{proof}


\begin{remark}[八面体单值化的射影模 \(3\) 读出与权一提升的可检验接口]\label{rem:xi-terminal-zm-octahedral-mod3-langlands}
标准同构 \(S_4\simeq \mathrm{PGL}_2(\FF_3)\)（例如 \cite{MITSmallGroupsPDF}）把定理 \ref{thm:xi-terminal-zm-leyang-elliptic-structure} 的单值群识别为八面体射影线性群像；
从而在 \(U\) 上得到一个有限像的 \(\mathrm{PGL}_2(\FF_3)\)-投影局部系统，可视为射影模 \(3\) 的几何 Langlands 参数。
另一方面，由结论 \ref{con:xi-terminal-zm-generic-galois-s4} 与结论 \ref{con:xi-terminal-zm-specialization-discriminant-ab} 的泛型 \(S_4\) 性与 Hilbert 不可约性可知：对除去薄集的有理特化 \(y_0\in\QQ\)，所得数域的伽罗瓦闭包给出一个八面体型（\(S_4\)）射影表示 \(G_\QQ\to \mathrm{PGL}_2(\FF_3)\)；
在满足 oddness 等算术条件的情形下，其二维提升落入 Langlands--Tunnell 的权 \(1\) 自守范畴，从而把“分歧—单值”数据与可计算的权一模形式检验接口对接。\cite{Tunnell1981OctahedralArtin,BuzzardPotentialModularitySurvey}
\end{remark}

\begin{remark}[零点凝聚层级的椭圆化入口]\label{rem:xi-terminal-zm-leyang-zero-measure-outlook}
在椭圆单值化框架下，有限体积多项式族 \(\{Z_m(y)\}\) 的零点集合可被转写为 \(\pi_y\) 的某类椭圆函数序列之水平集/相位锁定集；
据此可在严格代数意义上刻画零点凝聚曲线，并进一步寻求相应密度公式与推前测度的解析表达。
\end{remark}

\begin{conjecture}[Lee--Yang 椭圆谱曲线的 Stokes 常数与调和子接口]\label{conj:xi-leyang-elliptic-regulator-stokes-hub}
沿用定理 \ref{thm:xi-terminal-zm-leyang-elliptic-structure} 的椭圆谱曲线 \(E_{\min}/\QQ\) 及其模性锚点 \(L(E_{\min},s)=L(f,s)\)（权 \(2\)、level \(37\) 新形式 \(f\)）。
设 \(\alpha\) 为 \(E_{\min}\) 上由 Lee--Yang 端点/分歧数据自然诱导的一条亚纯 \(1\)-形式，其极点除子由分歧簇与端点簇决定且满足椭圆对称化的局部展开约束。
则存在 \(u,v\in\QQ(E_{\min})^\times\) 与一条由扭点 Stokes 缺陷所确定的规范 \(1\)-循环 \(\gamma\in H_1(E_{\min}(\CC),\ZZ)\)，使得 Beilinson 调和子满足
\[
\mathrm{Reg}\bigl(\{u,v\}\bigr)=\Re\int_{\gamma}\alpha,
\qquad
\mathrm{Reg}\bigl(\{u,v\}\bigr)\ \asymp_{\mathrm{norm}}\ L'(f,0),
\]
其中 \(\{u,v\}\in K_2(E_{\min})\) 为 Milnor \(K\)-群元，\(\mathrm{Reg}\) 为 Beilinson 调和子映射，而 \(\asymp_{\mathrm{norm}}\) 表示两端仅相差由归一化选择所决定的非零比例常数。
\par
在此意义下，端点审计所出现的 Stokes 常数应被视为 \(E_{\min}\) 的算术不变量（调和子/模性侧的共同投影），而非仅为渐近分析的附属系数。
\end{conjecture}

\endinput



\begin{conjecture}[金属均值族的 Lee--Yang 分歧几何与模性锚定]\label{conj:xi-metallic-mean-leyang-modular-anchor}
对每个整数 \(k\ge 1\)，令 \(\alpha_k=[k;k,k,\dots]\) 为金属均值。
假设存在一族以 \((\alpha_k,\mathrm{Fold}_m,\xi)\) 口径内生确定的二变量谱多项式
\[
\Pi_k(\lambda,y)\in \ZZ[\lambda,y],
\]
其在 \(k=1\) 时退化为本文的 \(\Pi(\lambda,y)\)。
并假设其 \(\lambda\)-判别式满足分解
\[
\mathrm{Disc}_{\lambda}\bigl(\Pi_k(\lambda,y)\bigr)
=\pm\,y(y-1)\,P_{\mathrm{LY},k}(y)\cdot Q_k(y)^{2},
\]
其中 \(P_{\mathrm{LY},k}\in\ZZ[y]\) 不可约且平方自由，\(Q_k\in\ZZ[y]\)。
令 \(S_k\) 为 \(P_{\mathrm{LY},k}\) 的判别式之平方自由部分所支撑的坏素集合（并并入 \(y(y-1)\) 的显式坏素）。
\par
则存在以下模性锚定结构：
\begin{enumerate}
  \item 存在整数 \(N_k\)（其素因子集合包含于 \(S_k\)）与一条模曲线 \(X_0(N_k)\) 的有限商，使得 \(\Pi_k(\lambda,y)=0\) 的归一化曲线在 \(\QQ\) 上支配该模曲线，并且其 Jacobian 在 \(\QQ\) 上分解出一个由 \(S_k\) 控制导子的阿贝尔因子；
  \item 在存在椭圆因子 \(E_k\) 的情形下，其 \(2\)-进 Tate 模 \(T_2(E_k)\) 给出 dyadic 分歧点逆极限的 torsor 参数化（与推论 \ref{cor:xi-terminal-zm-leyang-branchpoint-tate-torsor-decomposition} 的 \(k=1\) 情形一致），从而把“跨尺度一致的分支选择自由度”压缩为一个 \(2\)-进证书；
  \item 在全素门允许的极限口径下，\(\Sigma_{S_k}\) 的核 \(K_{S_k}\cong \prod_{p\in S_k}\ZZ_p\)（命题 \ref{prop:cdim-solenoid-kernel-product-zp}）给出相位系统的不可压缩证书轴，而圆因子 \(\TT\) 仅承载相干部分。
\end{enumerate}
\end{conjecture}

\endinput


\begin{theorem}[扭点平行四边形的 Stokes 缺陷与 Weil 配对接口]\label{thm:xi-elliptic-torsion-stokes-weil-pairing}
\leanverified{paper\_xi\_elliptic\_torsion\_stokes\_weil\_pairing}
设 \(E=\CC/\Lambda\) 为复椭圆曲线，\(\pi:\CC\to E\) 为自然投影。
取非零亚纯函数 \(f\in\CC(E)^\times\)，其除子记为
\[
\mathrm{div}(f)=\sum_{q\in E} n_q\,[q],\qquad n_q\in\ZZ,\ \text{有限支撑}.
\]
令 \(k\ge 1\)，并取两条 \(2^k\)-扭点 \(T_1,T_2\in E[2^k]\)。
对每个 \(P\in E[2^k]\) 选取提升 \(\widetilde P\in \frac{1}{2^k}\Lambda\) 使 \(\pi(\widetilde P)=P\)，并同样选取提升 \(\widetilde T_i\in \frac{1}{2^k}\Lambda\) 使 \(\pi(\widetilde T_i)=T_i\)。
对 \(i\in\{1,2\}\) 定义取值于 \(\CC/2\pi\mathrm{i}\ZZ\) 的离散 \(1\)-余链
\begin{equation}\label{eq:xi-elliptic-discrete-log-derivative-cochain}
\Omega_f(P\to P+T_i)
~:=~
\int_{\gamma(\widetilde P,\widetilde P+\widetilde T_i)} \dd\log(f\circ\pi)
\quad \in\ \CC/2\pi\mathrm{i}\ZZ,
\end{equation}
其中路径 \(\gamma(\widetilde P,\widetilde P+\widetilde T_i)\subset\CC\) 连接两端点并避开 \(\pi^{-1}(\mathrm{supp}\,\mathrm{div}(f))\)。
令 \(\square(P;T_1,T_2)\) 表示由 \(T_1,T_2\) 生成的扭点方形
\[
P\to P+T_1\to P+T_1+T_2\to P+T_2\to P.
\]
则在 \(\CC/2\pi\mathrm{i}\ZZ\) 中成立严格恒等式
\begin{equation}\label{eq:xi-elliptic-stokes-defect-divisor-index}
(d\Omega_f)\bigl(\square(P;T_1,T_2)\bigr)
\equiv
2\pi\mathrm{i}\sum_{q} n_q\,\mathrm{Ind}\bigl(q;P,T_1,T_2\bigr)
\pmod{2\pi\mathrm{i}\ZZ},
\end{equation}
其中 \(\mathrm{Ind}(q;P,T_1,T_2)\in\ZZ\) 为点 \(q\) 相对以 \(\widetilde P,\widetilde T_1,\widetilde T_2\) 确定之平行四边形围道的环绕指数。
特别地，若对所有 \(P\in E[2^k]\) 上述平行四边形围道均不包围任何 \(f\) 的零极点，则 \(d\Omega_f\equiv 0\)。
\par
进一步，若 \(f\) 为特征 theta 函数之比（从而 \(\mathrm{div}(f)\) 为扭点差的平移不变除子），则存在 \(S\in E[2^k]\) 使得由 \(\exp(\Omega_f)\) 所诱导的乘法 \(1\)-余圈与 Weil 配对同构：
\begin{equation}\label{eq:xi-elliptic-theta-weil-pairing-character}
\exp\!\bigl(\Omega_f(P\to P+T)\bigr)=e_{2^k}(T,S)
\qquad(\forall\,P,T\in E[2^k]).
\end{equation}
\end{theorem}
\begin{proof}
在 \(\CC\setminus\pi^{-1}(\mathrm{supp}\,\mathrm{div}(f))\) 上，\(\dd\log(f\circ\pi)\) 为闭 \(1\)-形式，其沿不同路径的积分差仅可能来自绕过零极点的回路，其值属于 \(2\pi\mathrm{i}\ZZ\)；故 \eqref{eq:xi-elliptic-discrete-log-derivative-cochain} 在 \(\CC/2\pi\mathrm{i}\ZZ\) 中良定。
\par
取 \(\square(P;T_1,T_2)\) 在 \(\CC\) 中的提升边界曲线 \(\partial\mathcal P\)（\(\mathcal P\) 为对应平行四边形域），则
\[
(d\Omega_f)\bigl(\square(P;T_1,T_2)\bigr)\equiv \oint_{\partial\mathcal P}\dd\log(f\circ\pi)\pmod{2\pi\mathrm{i}\ZZ}.
\]
由复分析的辩值原理/留数定理（等价于对亚纯微分的广义 Stokes 公式），右端等于 \(2\pi\mathrm{i}\) 乘以 \(\mathcal P\) 内零极点的代数计数，并以环绕指数加权：
\[
\oint_{\partial\mathcal P}\dd\log(f\circ\pi)
=2\pi\mathrm{i}\sum_q n_q\,\mathrm{Ind}(q;P,T_1,T_2),
\]
从而得到 \eqref{eq:xi-elliptic-stokes-defect-divisor-index}。
若域内无零极点，则右端为 \(0\)，故 \(d\Omega_f\equiv 0\)。
\par
最后一条关于 theta 函数与 Weil 配对的陈述，是 Mumford theta 群对 \(E[2^k]\) 的平移线性化之标准结论：由扭点差除子所定义的线丛在 \(E[2^k]\) 上的平移自同构因子给出一条双角色，而 Weil 配对刻画了该双角色的规范模型；因此 \(\exp(\Omega_f)\) 必与某个 \(S\in E[2^k]\) 的 \(e_{2^k}(\,\cdot\,,S)\) 同构。证毕。
\end{proof}

\endinput

\begin{conclusion}[Lee--Yang 特征谱曲线的椭圆化与极小模型辨识]\label{con:xi-terminal-zm-elliptic-normalization}
沿用结论 \ref{con:xi-terminal-zm-order4-rational} 的记号，并令
\[
\mathcal C:=\bigl\{(\lambda,y)\in\mathbb A^2:\ \Pi(\lambda,y)=0\bigr\}.
\]
定义新变量
\[
w:=2y+1-2\lambda^{2}.
\]
并令
\[
\eta:=y-\lambda^{2}\qquad(\text{亦记 }u).
\]
则在 \((x,u)\) 坐标（其中 \(x=\lambda\)）下有恒等式
\[
\Pi\bigl(x,\ u+x^{2}\bigr)=u^{2}+u-(x^{3}-x),
\]
从而在 \(\QQ\) 上得到整系数极小模型
\[
E_{\min}:\quad u^{2}+u=x^{3}-x\qquad(\text{亦即 }\eta^{2}+\eta=x^{3}-x).
\]
等价地，若将谱坐标重命名为 \(U:=\lambda\)、\(V:=y-\lambda^{2}\)，则该极小模型亦可写为 \(V^{2}+V=U^{3}-U\)。
并且 \((x,u)\leftrightarrow(\lambda,y)\) 为互逆的多项式代换，从而在仿射层面给出 \(\mathcal C\) 与 \(E_{\min}\) 的 \(\QQ\)-同构；
若取射影闭包，则 \(\Pi(\lambda,y)=0\) 在无穷远处可以出现奇点，但规范化后即为光滑椭圆曲线。
该同构可由显式代换
\[
(x,u)\longmapsto(\lambda,y)=(x,\ u+x^{2}),
\qquad
(\lambda,y)\longmapsto(x,u)=(\lambda,\ y-\lambda^{2})
\]
直接给出。
进一步引入 \(w=2u+1\)（等价于上式定义的 \(w=2y+1-2\lambda^{2}\)）后，得到同构
\[
\mathcal C\ \dashrightarrow\ \mathcal E,\qquad
(\lambda,y)\longmapsto(\lambda,w),
\qquad
\mathcal E:\ w^2=4\lambda^{3}-4\lambda+1,
\]
其逆变换由
\[
y=\frac{2\lambda^{2}+w-1}{2}
\]
给出。
令 \(X=\lambda,\ Y=w/2\)，则 \(\mathcal E\) 的射影闭包为短 Weierstrass 曲线
\[
E:\ Y^2=X^3-X+\frac14,
\]
并且在平面坐标 \((\lambda,y)\) 上存在完全显式的平方差分解
\[
\Pi(\lambda,y)=\Bigl(y-\lambda^2+\tfrac12\Bigr)^2-\Bigl(\lambda^3-\lambda+\tfrac14\Bigr).
\]
因此令 \(X:=\lambda\)、\(Y:=y-\lambda^2+\tfrac12\)（亦即 \(Y=w/2\)）可得到函数域同一
\[
\QQ(\mathcal C)\cong \QQ(E),\qquad \lambda=X,\quad Y=y-\lambda^2+\tfrac12.
\]
更精确地，存在完全显式的 \(\QQ\)-双有理互逆
\[
\varphi:\ \mathcal C\dashrightarrow E,\quad (\lambda,y)\longmapsto (X,Y):=\bigl(\lambda,\ y+\tfrac12-\lambda^2\bigr),
\qquad
\psi:\ E\dashrightarrow \mathcal C,\quad (X,Y)\longmapsto (\lambda,y):=\bigl(X,\ X^2-\tfrac12+Y\bigr),
\]
从而 \(\mathcal C\) 的规范化同构于 \(E\)，并且 \(\QQ(\mathcal C)\cong \QQ(E)\)。
其判别式与 \(j\)-不变量分别为
\[
\Delta(E)=37,\qquad j(E)=\frac{48^{3}}{37}=\frac{110592}{37}.
\]
进一步令 \(y_{\mathrm{int}}:=Y-\tfrac12\)（亦即 \(y_{\mathrm{int}}=u=y-\lambda^{2}\)），则得到整系数最小模型
\[
E_{\min}:\quad y_{\mathrm{int}}^{2}+y_{\mathrm{int}}=X^{3}-X,
\]
其 \(\Delta(E_{\min})=37\) 且 \(j(E_{\min})=48^{3}/37\)。
因此 \((\lambda,y)\)-谱曲线的代数分支结构可等价地提升为椭圆曲线 \(E\) 上的有理坐标结构，从而把分歧点、解析延拓与单值化问题转写为 \(E\) 上的代数几何数据。
特别地，\(\Delta(E)=37\) 表明该椭圆曲线在除 \(37\) 外处处良好约化；并且在 Cremona/LMFDB 标识下对应同源类 \(37\mathrm a\) 的强 Weil 曲线 \(37\mathrm a1\)（导子 \(37\)，见注记 \ref{rem:xi-terminal-elliptic-37a1-modular-anchor}）。
\par
在上述同构下，重量参数作为 \(E\) 上的有理函数可写为
\[
\boxed{\;
y=\frac{2\lambda^{2}+w-1}{2}=X^2+Y-\frac12\ \in\ \QQ(E)\; } ,
\qquad (O\text{ 为 }E\text{ 的无穷远点}).
\]
并记
\[
P_{0}=\Bigl(0,\frac12\Bigr),\quad P_{-1}=\Bigl(-1,-\frac12\Bigr),\quad P_{1}=\Bigl(1,-\frac12\Bigr)\in E(\QQ),
\]
\[
Q_{2}=\Bigl(2,-\frac52\Bigr),\quad Q_{1}=\Bigl(1,\frac12\Bigr),\quad Q_{-1}=\Bigl(-1,\frac12\Bigr)\in E(\QQ).
\]
由 Weierstrass 形式知 \(\operatorname{ord}_{O}(X)=-2\)、\(\operatorname{ord}_{O}(Y)=-3\)，从而 \(\operatorname{ord}_{O}(y)=-4\) 且 \((y)_\infty=4O\)，故 \(y:E\to\PP^1_y\) 为次数 \(4\) 的覆盖映射。
零点由 \(y=0\iff Y=-X^2+\tfrac12\) 与曲线方程联立得到：
\[
(-X^2+\tfrac12)^2=X^3-X+\tfrac14\iff X(X-1)^2(X+1)=0,
\]
从而主除子分解为
\[
\mathrm{div}(y)=P_{0}+P_{-1}+2P_{1}-4O.
\]
同理，由 \(y=1\iff Y=\tfrac32-X^2\) 得
\[
(\tfrac32-X^2)^2=X^3-X+\tfrac14\iff (X-2)(X-1)(X+1)^2=0,
\]
从而
\[
\mathrm{div}(y-1)=Q_{2}+Q_{1}+2Q_{-1}-4O.
\]
两式皆为主除子，故其在 \(\mathrm{Pic}^0(E)\cong E\) 中对应零类；等价地在 \(E(\QQ)\) 的群结构中成立
\[
P_{0}+P_{-1}+2P_{1}=O,
\qquad
Q_{2}+Q_{1}+2Q_{-1}=O.
\]
结合结论 \ref{con:xi-terminal-zm-elliptic-mw-generator-alignment} 的生成元 \(R=(0,\tfrac12)\) 与低倍点对齐，上述两条纤维分解亦可在除子意义下写为
\[
y^{-1}(0)=\{R\}+2\{-2R\}+\{3R\},\qquad
y^{-1}(1)=\{2R\}+2\{-3R\}+\{4R\},
\]
其中各点按重数计。
因此 \(y\) 在 \(y=0\) 与 \(y=1\) 处各出现一处二重纤维点（分别位于 \(\lambda=1\) 与 \(\lambda=-1\)），而其余分歧值由判别式分解（结论 \ref{con:xi-terminal-zm-discriminant-leyang} 或等价的结果式刻画，见结论 \ref{con:xi-terminal-zm-leyang-cubic-branch-image}）所完全决定；这给出一个纯代数几何的“除子—分歧”证书。
\end{conclusion}
\begin{proof}[证明要点]
将 \(\Pi(\lambda,y)=0\) 视为关于 \(y\) 的二次方程，得
\[
y^2+(1-2\lambda^2)y+\bigl(\lambda^4-\lambda^3-\lambda^2+\lambda\bigr)=0.
\]
配方可化为
\[
\bigl(2y+1-2\lambda^2\bigr)^2=4\lambda^3-4\lambda+1.
\]
令 \(X=\lambda\)、\(Y=y+\tfrac12-\lambda^2\)，即得 \(Y^2=X^3-X+\tfrac14\)，从而得到 \(\varphi\)。
反向代入 \(y=X^2-\tfrac12+Y\) 立即复原 \(\Pi(X,y)=0\)，故 \(\varphi,\psi\) 互为双有理逆。
Weierstrass 形式下取 \(a=-1,b=1/4\)，则
\(
\Delta=-16(4a^3+27b^2)=37
\)，且 \(j=c_4^3/\Delta\)（其中 \(c_4=-48a\)）给出 \(j(E)=110592/37\)。
可复现核验脚本：\texttt{scripts/exp\_fold\_zm\_elliptic\_leyang\_arithmetic\_audit.py}。证毕。
\end{proof}

\begin{conclusion}[有限域点计数的谱解释：仿射谱曲线点数与 Frobenius 迹的严格对应]\label{con:xi-terminal-zm-elliptic-finite-field-point-count}
沿用结论 \ref{con:xi-terminal-zm-elliptic-normalization} 的谱四次 \(\Pi(\lambda,y)\) 与极小椭圆模型
\[
E_0:\ v^{2}+v=x^{3}-x.
\]
取素数 \(p\neq 37\)，并记在 \(\FF_p\) 上的仿射解集
\[
C_p(\FF_p):=\{(\lambda,y)\in\FF_p^{2}:\ \Pi(\lambda,y)=0\},
\qquad
E_{0,p}(\FF_p):=\{(x,v)\in\FF_p^{2}:\ v^{2}+v=x^{3}-x\}\cup\{O\}.
\]
则由整系数多项式互逆代换
\[
(x,v)\longmapsto(\lambda,y)=(x,\ v+x^{2}),
\qquad
(\lambda,y)\longmapsto(x,v)=(\lambda,\ y-\lambda^{2})
\]
在 \(\FF_p\) 上诱导双射
\[
E_{0,p}(\FF_p)\setminus\{O\}\ \longrightarrow\ C_p(\FF_p),
\]
从而
\[
\#C_p(\FF_p)=\#E_{0,p}(\FF_p)-1.
\]
令
\(
a_p:=p+1-\#E_{0,p}(\FF_p)
\)
为 \(E_0\) 在良素 \(p\neq 37\) 处的 Frobenius 迹，则等价地有
\[
\#C_p(\FF_p)=p-a_p.
\]
并且对奇素数 \(p\neq 37\) 成立二次特征和公式
\[
a_p=-\sum_{x\in\FF_p}\left(\frac{4x^{3}-4x+1}{p}\right),
\]
其中 \(\bigl(\tfrac{\cdot}{p}\bigr)\) 为 Legendre 符号（约定 \(\bigl(\tfrac{0}{p}\bigr)=0\)）。
\end{conclusion}
\begin{proof}[证明要点]
变换 \((\lambda,y)\leftrightarrow(x,v)\) 由结论 \ref{con:xi-terminal-zm-elliptic-normalization} 给出，且在 \(\ZZ\) 上为多项式互逆；故对任意素数 \(p\) 均在 \(\FF_p\) 上诱导仿射点集的双射。
由于 Weierstrass 模型 \(E_{0,p}\) 的射影完备化仅多出唯一无穷远点 \(O\)，得 \(\#C_p(\FF_p)=\#E_{0,p}(\FF_p)-1\)，从而 \(\#C_p(\FF_p)=p-a_p\)。
\par
设 \(p\neq 2,37\)。
对每个 \(x\in\FF_p\)，方程 \(v^{2}+v=x^{3}-x\) 的解数等于
\[
1+\left(\frac{1+4(x^{3}-x)}{p}\right)=1+\left(\frac{4x^{3}-4x+1}{p}\right),
\]
因此
\[
\#E_{0,p}(\FF_p)
=1+\sum_{x\in\FF_p}\left(1+\left(\frac{4x^{3}-4x+1}{p}\right)\right)
=p+1+\sum_{x\in\FF_p}\left(\frac{4x^{3}-4x+1}{p}\right).
\]
代入 \(a_p=p+1-\#E_{0,p}(\FF_p)\) 得所示公式。证毕。
\end{proof}

\begin{proposition}[有限域纤维计数与 Frobenius 迹的均值恒等式]\label{prop:xi-terminal-zm-elliptic-fiber-count-mean-trace}
沿用结论 \ref{con:xi-terminal-zm-elliptic-finite-field-point-count} 的记号，并取良素 \(p\neq 37\)。
对每个 \(y\in\FF_p\) 定义纤维计数
\[
N_p(y):=\card{\{\lambda\in\FF_p:\ \Pi(\lambda,y)=0\}}.
\]
则成立：
\[
\#E_{0,p}(\FF_p)=1+\sum_{y\in\FF_p}N_p(y),
\qquad
\sum_{y\in\FF_p}\bigl(N_p(y)-1\bigr)=-a_p.
\]
\leanverified{paper\_xi\_terminal\_zm\_elliptic\_fiber\_count\_mean\_trace}
\end{proposition}

\begin{proof}
按定义，
\(
C_p(\FF_p)=\{(\lambda,y)\in\FF_p^2:\ \Pi(\lambda,y)=0\}
\)
的点数可按 \(y\)-纤维分解为
\[
\#C_p(\FF_p)=\sum_{y\in\FF_p}N_p(y).
\]
而结论 \ref{con:xi-terminal-zm-elliptic-finite-field-point-count} 已给出 \(\#C_p(\FF_p)=\#E_{0,p}(\FF_p)-1\)，故得第一式。
再由同一结论的 \(\#C_p(\FF_p)=p-a_p\) 得
\[
\sum_{y\in\FF_p}\bigl(N_p(y)-1\bigr)
=\#C_p(\FF_p)-p
=(p-a_p)-p
=-a_p.
\]
证毕。
\end{proof}

\begin{corollary}[Hecke 迹的随机纤维估计：无偏性与方差公式]\label{cor:xi-terminal-zm-elliptic-mc-trace-estimator}
\leanverified{paper\_xi\_terminal\_zm\_elliptic\_mc\_trace\_estimator}
在命题 \ref{prop:xi-terminal-zm-elliptic-fiber-count-mean-trace} 的设定下，
令 \(Y_1,\dots,Y_n\) 为 \(\FF_p\) 上独立均匀样本，并令
\[
Z_i:=N_p(Y_i)-1,
\qquad
\widehat a_p:=-p\cdot \frac{1}{n}\sum_{i=1}^{n} Z_i.
\]
则 \(\widehat a_p\) 为 \(a_p\) 的无偏估计，且
\[
\EE[\widehat a_p]=a_p,
\qquad
\Var(\widehat a_p)=\frac{p^{2}}{n}\Var(Z_1).
\]
\end{corollary}

\begin{proof}
由命题 \ref{prop:xi-terminal-zm-elliptic-fiber-count-mean-trace}，
\[
\EE[Z_1]
=\frac{1}{p}\sum_{y\in\FF_p}\bigl(N_p(y)-1\bigr)
=-\frac{a_p}{p},
\]
故 \(\EE[\widehat a_p]=a_p\)。
方差式由独立同分布与缩放恒等式直接推出：
\(
\Var(\widehat a_p)
=p^2\Var(\frac1n\sum_i Z_i)
=p^2\cdot \frac{1}{n}\Var(Z_1)
\)。
证毕。
\end{proof}

\begin{conclusion}[谱曲线在无穷远处的唯一奇点：\texorpdfstring{$A_4$}{A4} 型单分支尖点与 \texorpdfstring{$\delta=2$}{delta=2}]\label{con:xi-terminal-zm-spectral-curve-infty-a4-cusp}
沿用结论 \ref{con:xi-terminal-zm-elliptic-normalization} 的谱四次
\[
\Pi(\lambda,y)=\lambda^4-\lambda^3-(2y+1)\lambda^2+\lambda+y(y+1)\in\QQ[\lambda,y],
\]
并令 \(\overline{\mathcal C}\subset \mathbb P^1_\lambda\times \mathbb P^1_y\) 表示平面仿射曲线 \(\Pi(\lambda,y)=0\) 的闭包（视为 \((4,2)\)-型曲线）。
则成立如下事实：
\begin{enumerate}
  \item \(\overline{\mathcal C}\) 仅在 \((\lambda,y)=(\infty,\infty)\) 处与无穷远边界相交，并且该点为 \(\overline{\mathcal C}\) 的唯一奇点。
  \item 在 \((\infty,\infty)\) 的邻域取局部坐标
  \[
  \mu:=\frac{1}{\lambda},\qquad v:=\frac{1}{y},
  \]
  则 \(\overline{\mathcal C}\) 的局部方程可写为关于 \(v\) 的二次式，其判别式满足严格分解
  \[
  \boxed{\;\Delta_v(\mu)=\mu^5\cdot(\mu^3-4\mu^2+4)\;}
  \]
  从而该奇点在解析等价意义下同胚于标准形式
  \[
  v^2=u^5,
  \]
  即 \(\,A_4\)（ramphoid cusp）型单分支尖点；因此其 \(\delta\)-不变量为
  \[
  \boxed{\;\delta=\frac{(2-1)(5-1)}{2}=2\; }.
  \]
  \item 光滑 \((4,2)\)-型曲线的算术属为 \((4-1)(2-1)=3\)；由上述唯一奇点贡献 \(\delta=2\)，故几何归一化的几何属为 \(3-2=1\)，与结论 \ref{con:xi-terminal-zm-elliptic-normalization} 中“归一化为椭圆曲线”的识别一致。
\end{enumerate}
该 \(A_4\) 尖点亦可视为几何上对结论 \ref{con:xi-terminal-zm-elliptic-y-hurwitz-passport} 所述 \(y=\infty\) 处四重全分歧极点结构的平面压缩编码。
\end{conclusion}
\begin{proof}[证明要点]
结论 \ref{con:xi-terminal-zm-elliptic-normalization} 已给出：\(\Pi(\lambda,y)=0\) 的仿射部分与光滑椭圆曲线同构，
故其有限点处无奇点；因而若存在奇点，只能落在无穷远边界。
把 \(\Pi\) 视为 \(\mathbb P^1_\lambda\times \mathbb P^1_y\) 上的双齐次方程，直接检验 \(\lambda=\infty\) 且 \(y\neq\infty\)（或 \(y=\infty\) 且 \(\lambda\neq\infty\)）均不可能满足该方程，
从而边界交点只能为 \((\infty,\infty)\)。
\par
在 \((\infty,\infty)\) 附近取 \(\mu=1/\lambda\)、\(v=1/y\)，则
\[
\Pi\!\left(\frac{1}{\mu},\frac{1}{v}\right)=0
\quad\Longleftrightarrow\quad
F(\mu,v)=0,
\]
其中乘以 \(\mu^4v^2\) 清分母得到
\[
F(\mu,v)
=(1-\mu-\mu^2+\mu^3)v^2+\mu^2(\mu^2-2)v+\mu^4.
\]
视 \(F\) 为关于 \(v\) 的二次式，其判别式为
\[
\Delta_v(\mu)=\mu^4(\mu^2-2)^2-4\mu^4(1-\mu-\mu^2+\mu^3)
=\mu^5(\mu^3-4\mu^2+4),
\]
给出所示分解。
并且完成平方得到
\[
\bigl(2(1-\mu-\mu^2+\mu^3)v+\mu^2(\mu^2-2)\bigr)^2=\Delta_v(\mu).
\]
由于 \(\mu^3-4\mu^2+4\) 在 \(\mu=0\) 处为单位（取值 \(4\)），故在解析坐标变换下可将右端单位因子吸收，
从而奇点与 \(v^2=u^5\) 解析等价；\(\delta\) 由平面分支 \(v^2=u^5\) 的标准公式给出。
最后属计算由算术属与 \(\delta\) 的差给出。证毕。
\end{proof}

\begin{theorem}[无穷远处的四支 Puiseux 展开与四循环单值化]\label{thm:xi-terminal-zm-puiseux-infty-4cycle}
\leanverified{paper\_xi\_terminal\_zm\_puiseux\_infty\_4cycle}
沿用结论 \ref{con:xi-terminal-zm-spectral-curve-infty-a4-cusp} 的记号。
则在无穷远尖点 \((\lambda,y)=(\infty,\infty)\) 的穿孔邻域内存在局部参数 \(t\) 使得
\[
y=t^{-4},
\]
并且相应的一条谱分支 \(\lambda=\lambda(t)\) 具有 Puiseux 展开
\[
\lambda
=t^{-2}+\frac12 t^{-1}+\frac14+\frac{7}{64}t+\frac{37}{128}t^{2}-\frac{729}{4096}t^{3}+O(t^{4}).
\]
对同一 \(y\) 的其余三条局部分支由替换 \(t\mapsto i^{k}t\)（\(k=1,2,3\)）生成；
因此绕 \(y=\infty\) 一周的局部单值化在四张纸上的作用为一个 \(4\)-循环置换。
\end{theorem}
\begin{proof}[证明要点]
由于结论 \ref{con:xi-terminal-zm-elliptic-y-hurwitz-passport} 给出 \(y\) 在椭圆归一化上于无穷远点具有唯一四阶极点，
故可取局部一致化参数 \(t\) 使 \(y=t^{-4}\)。
另一方面，\(\lambda\) 在该点具有二阶极点，因而可设
\(
\lambda=t^{-2}+a t^{-1}+b+c t+d t^{2}+e t^{3}+O(t^{4})
\)。
将 \(y=t^{-4}\) 与该形式代入谱方程 \(\Pi(\lambda,y)=0\) 并逐阶比较系数，可得 \(a=\pm\frac12\)；
选定一条分支后，后续系数 \(b,c,d,e\) 被递推唯一确定，从而得到所示展开。
由于 \(y=t^{-4}\) 在 \(t\mapsto i t\) 下不变，解析延拓绕 \(y=\infty\) 一周对应于 \(t\mapsto i t\)，
从而四条局部分支在该作用下循环置换，给出四循环单值化。
可复现核验脚本：\texttt{scripts/exp\_fold\_zm\_elliptic\_leyang\_cover\_geometry\_audit.py}。证毕。
\end{proof}


\begin{conclusion}[次数 \(4\) 线性系 \texorpdfstring{$|4O|$}{|4O|} 的射影正规模型：两二次曲面的显式交与谱四次的线性投影复原]\label{con:xi-terminal-zm-elliptic-degree4-projective-model}
沿用结论 \ref{con:xi-terminal-zm-elliptic-normalization} 的整系数极小模型
\(
E_0:\ \eta^{2}+\eta=x^{3}-x
\)，记 \(O\) 为无穷远点，并令重量函数
\(
y:=x^{2}+\eta\in\QQ(E_0)
\)。
则
\(
1,x,x^{2},\eta\in H^{0}\!\bigl(E_0,\mathcal O_{E_0}(4O)\bigr)
\)
且构成 \(H^{0}(E_0,\mathcal O_{E_0}(4O))\) 的一组基。
由完全线性系 \(|4O|\) 诱导的态射
\[
\kappa:E_0\hookrightarrow \mathbb P^{3},\qquad
(x,\eta)\longmapsto [U:V:W:Z]=[1:x:x^{2}:\eta]
\]
为射影嵌入，其像为次数 \(4\) 的椭圆正规曲线，并在 \(\mathbb P^{3}\) 中由两条二次方程刻画为完全交：
\[
UW-V^{2}=0,
\qquad
Z^{2}+UZ-VW+VU=0.
\]
在该模型中有恒等式 \(y=(W+Z)/U\)，从而覆盖 \(y:E_0\to\mathbb P^{1}_{y}\) 可视为由二超平面丛 \(\langle U,\ W+Z\rangle\subset H^{0}(E_0,\mathcal O(4O))\) 所定义的次数 \(4\) 线性投影；在仿射片 \(U=1\) 上消去 \((W,Z)\)（等价于代入 \(\eta=y-x^{2}\)）即严格复原谱方程 \(\Pi(x,y)=0\)。
\end{conclusion}
\begin{proof}[证明要点]
在 \(E_0\) 上有 \(\mathrm{ord}_{O}(x)=-2\)、\(\mathrm{ord}_{O}(\eta)=-3\)，故 \(1,x,x^{2},\eta\in H^{0}(E_0,\mathcal O_{E_0}(4O))\)。
若 \(a_{0}+a_{1}x+a_{2}x^{2}+a_{3}\eta=0\)，比较 \(O\) 处极阶先得 \(a_{3}=0\)，继而 \(a_{2}=a_{1}=a_{0}=0\)，从而四者线性无关；由于 \(\dim H^{0}(E_0,\mathcal O_{E_0}(4O))=4\)，它们构成一组基。
当 \(\deg(4O)=4\ge 3\) 时 \(|4O|\) 为 very ample，故 \(\kappa\) 为嵌入；其像为次数 \(4\) 的椭圆正规曲线，因而为两二次曲面的完全交。
两条二次方程由代入直接核验：第一式 \(UW-V^{2}=0\) 来自 \(W=x^{2}\)；第二式在仿射图 \(U=1\) 上即为 \(\eta^{2}+\eta=x^{3}-x\)，写作 \(Z^{2}+Z=VW-V\) 并齐次化即得。
最后 \(y=(W+Z)/U\) 由定义推出，而在仿射片 \(U=1\) 上代入 \(\eta=y-x^{2}\) 即将二次关系化为 \(\Pi(x,y)=0\)。证毕。
\end{proof}

\begin{conclusion}[共轭重量支的主除子：低倍有理点轨道上的完全写法]\label{con:xi-terminal-zm-elliptic-conjugate-weight-principal-divisors}
沿用结论 \ref{con:xi-terminal-zm-elliptic-normalization} 的记号，并令椭圆反演为
\(
\iota:(X,Y)\mapsto(X,-Y)
\)。
定义共轭重量函数
\[
y^{\iota}:=\iota(y)=X^{2}-Y-\frac12\in\QQ(E).
\]
记 \(R:=(0,\tfrac12)\in E(\QQ)\) 为结论 \ref{con:xi-terminal-zm-elliptic-mw-generator-alignment} 的自由生成元对应点。
则 \(y^{\iota}\) 与 \(y^{\iota}-1\) 的主除子满足
\[
\operatorname{div}(y^{\iota})=(-R)+(-3R)+2(2R)-4O,
\qquad
\operatorname{div}(y^{\iota}-1)=(-4R)+(-2R)+2(3R)-4O.
\]
因此两条特殊纤维（\(y^{\iota}=0\) 与 \(y^{\iota}=1\)）的零点（含重数）均支撑在由 \(R\) 生成的 Mordell--Weil 轨道上；
并且 \(y^{\iota}\)（分别 \(y^{\iota}-1\)）作为 \(|4O|\) 的截面，其零除子在该轨道上的分布被上述主除子恒等式唯一固定。
\end{conclusion}
\begin{proof}[证明要点]
由 \(y^{\iota}=X^{2}-Y-\tfrac12\) 与结论 \ref{con:xi-terminal-zm-elliptic-normalization} 的除子证书，
对 \(\operatorname{div}(y)\)、\(\operatorname{div}(y-1)\) 施加 \(\iota\)（注意 \(\iota(O)=O\)）即可得到
\[
\operatorname{div}(y^{\iota})=\Bigl(0,-\tfrac12\Bigr)+\Bigl(-1,\tfrac12\Bigr)+2\Bigl(1,\tfrac12\Bigr)-4O,
\qquad
\operatorname{div}(y^{\iota}-1)=\Bigl(2,\tfrac52\Bigr)+\Bigl(1,-\tfrac12\Bigr)+2\Bigl(-1,-\tfrac12\Bigr)-4O.
\]
再由结论 \ref{con:xi-terminal-zm-elliptic-mw-generator-alignment} 的低倍点对齐
\(
[2]R=(1,\tfrac12),\ [3]R=(-1,-\tfrac12),\ [4]R=(2,-\tfrac52)
\)
即可把上述坐标表示改写为所示 \(\langle R\rangle\) 轨道形式。证毕。
\end{proof}

\begin{conclusion}[有理耦合的 Mordell--Weil 参数化：\texorpdfstring{$\Pi(\lambda,y_0)$}{Pi(lambda,y0)} 在 \texorpdfstring{$\mathbb Q$}{Q} 上有根当且仅当 \(y_0\in y\bigl(E(\QQ)\bigr)\)]\label{con:xi-terminal-zm-pi-rational-root-iff-y-image}
沿用结论 \ref{con:xi-terminal-zm-elliptic-normalization} 的谱四次
\(
\Pi(\lambda,y)=0
\)
及其椭圆化同构。
则对任意 \(y_{0}\in\QQ\)，以下两条件等价：
\begin{enumerate}
  \item 存在 \(\lambda\in\QQ\) 使得 \(\Pi(\lambda,y_{0})=0\)；
  \item 存在 \(P\in E(\QQ)\) 使得 \(y_{0}=y(P)\)。
\end{enumerate}
并且在上述条件成立时，可取 \(\lambda=X(P)\) 作为对应的有理谱根；在极小整模型 \(E_0:\eta^2+\eta=x^3-x\)（其中 \(x=X,\ \eta=Y-\tfrac12\)）下等价地有 \(y_0=x(P)^2+\eta(P)\) 且 \(\lambda_0=x(P)\)。
特别地，此时 \(\Pi(\lambda,y_{0})\) 在 \(\QQ[\lambda]\) 中含有理线性因子 \(\lambda-X(P)\)；并由结论 \ref{con:xi-terminal-zm-elliptic-mw-generator-alignment}（\(E_0(\QQ)\) 无挠且 \(\mathrm{rank}=1\)，可取生成元 \(R_0=(0,0)\)）知 \(y\bigl(E(\QQ)\bigr)\subset\QQ\) 为无限集。
若将 \(X(P)\) 写为既约分式 \(X(P)=U/V^{2}\)（\(U,V\in\ZZ,\ \gcd(U,V)=1,\ V>0\)），则存在三次多项式 \(G(\lambda)\in\QQ[\lambda]\) 使得
\[
\Pi(\lambda,y_{0})=(V^{2}\lambda-U)\,G(\lambda),
\]
并且既约分母满足
\[
\mathrm{den}(y_{0})=\mathrm{den}\bigl(X(P)\bigr)^{2}=V^{4}.
\]
\end{conclusion}
\begin{proof}[证明要点]
若存在 \(\lambda\in\QQ\) 使 \(\Pi(\lambda,y_{0})=0\)，由结论 \ref{con:xi-terminal-zm-elliptic-normalization} 的恒等式
\(
4\Pi(\lambda,y)=w^{2}-(4\lambda^{3}-4\lambda+1)
\)
（其中 \(w=2y+1-2\lambda^{2}\)）可知 \((\lambda,w)\in\mathcal E(\QQ)\)，从而 \(P=(X,Y)=(\lambda,w/2)\in E(\QQ)\)，并满足 \(y(P)=y_{0}\)。
反之，若 \(P\in E(\QQ)\)，则 \((\lambda,y)=(X(P),y(P))\) 满足 \(\Pi(\lambda,y)=0\)，从而 \(\lambda=X(P)\) 给出 \(\Pi(\lambda,y_{0})\) 的有理根。
\par
线性因子与清分母因子 \((V^{2}\lambda-U)\) 由 \(\lambda=X(P)=U/V^{2}\) 为根立即推出。
分母平方关系可直接由曲线方程与 \(y=X^{2}+Y-\tfrac12\) 推出：写 \(X=U/V^{2}\)（既约），则由
\(
Y^{2}=X^{3}-X+\tfrac14
\)
得 \(Y=W/(2V^{3})\)（\(W\in\ZZ\)）且 \(W^{2}=4U^{3}-4UV^{4}+V^{6}\)；
从而
\[
y(P)=\frac{U^{2}}{V^{4}}+\frac{W}{2V^{3}}-\frac12=\frac{2U^{2}+WV-V^{4}}{2V^{4}}.
\]
注意 \(2U^{2}+WV-V^{4}\) 恒为偶数，故 \(y(P)=N/V^{4}\)（\(N\in\ZZ\)）。
若奇素数 \(p\mid V\)，则 \(WV\equiv V^{4}\equiv 0\pmod p\) 且 \(2N\equiv 2U^{2}\not\equiv 0\pmod p\)，从而 \(p\nmid N\)；
若 \(2\mid V\)，则 \(U\) 为奇数；
并且由 \(W^{2}=4U^{3}-4UV^{4}+V^{6}\) 在模 \(8\) 下推出 \(W^{2}\equiv 4\pmod 8\)，故 \(W=2W_{1}\) 且 \(W_{1}\) 为奇数。
于是 \(WV\equiv 0\pmod 4\)、\(V^{4}\equiv 0\pmod 4\)，而 \(2U^{2}\equiv 2\pmod 4\)，从而 \(2N\equiv 2\pmod 4\)，即 \(N\) 为奇数；
综上 \(\gcd(N,V)=1\)。
因此 \(\mathrm{den}(y(P))=V^{4}=\mathrm{den}(X(P))^{2}\)。证毕。
\end{proof}

\begin{remark}[双判别桥：\texorpdfstring{$\Pi$}{Pi} 的两条判别式方向]\label{rem:xi-terminal-zm-pi-double-discriminant-bridge}
将 \(\Pi(\lambda,y)\) 视为关于 \(y\) 的二次多项式，则其 \(y\)-判别式为
\[
\mathrm{Disc}_{y}\bigl(\Pi(\lambda,y)\bigr)=4\lambda^{3}-4\lambda+1.
\]
等价地，令 \(w:=2y+1-2\lambda^{2}\)，则
\[
4\Pi(\lambda,y)=w^{2}-\bigl(4\lambda^{3}-4\lambda+1\bigr).
\]
从而同一方程在 \(y\)-方向上的判别式给出椭圆曲线模型 \(\mathcal E:w^{2}=4\lambda^{3}-4\lambda+1\)（见结论 \ref{con:xi-terminal-zm-elliptic-normalization}）；其中三次式 \(4\lambda^{3}-4\lambda+1\) 亦即 \(E_{\min}\) 的 \(2\)-division 多项式（非平凡 \(2\)-挠点的 \(x\)-坐标所满足之方程）。
另一方面，将 \(\Pi(\lambda,y)\) 视为关于 \(\lambda\) 的四次多项式，则其 \(\lambda\)-判别式在 \(\ZZ[y]\) 中分解为
\[
\mathrm{Disc}_{\lambda}\bigl(\Pi(\lambda,y)\bigr)=-y(y-1)P_{\mathrm{LY}}(y),
\qquad
P_{\mathrm{LY}}(y):=256y^{3}+411y^{2}+165y+32,
\]
从而得到属 \(2\) 的判别式曲线 \(C:\ W^{2}=-y(y-1)P_{\mathrm{LY}}(y)\)（此处记 \(W\) 以避免与椭圆坐标 \(w\) 混淆）。
该“二次—四次”的双判别结构把椭圆归一化、分歧像以及四次分裂域的二次分辨率统一编码在同一 \(\Pi\) 之中。
\end{remark}

\begin{remark}[外部命名与交叉核验接口：\texorpdfstring{$37\mathrm a1$}{37a1}、\texorpdfstring{$X_0^+(37)$}{X0+(37)} 与 level \(37\) 新形式]\label{rem:xi-terminal-elliptic-37a1-modular-anchor}
本段的椭圆曲线 \(E: Y^2=X^3-X+\tfrac14\) 在 \(\QQ\) 上与整系数模型
\[
E_0:\quad \eta^2+\eta=x^3-x
\]
同构；显式同构由
\(
X=x,\ Y=\eta+\tfrac12
\)
给出；等价地亦可取
\(
X=x,\ Y=-\eta-\tfrac12
\)
（二者相差椭圆反演）。
本文统一采用 Weierstrass 不变量的标准约定
\[
c_4^3-c_6^2=1728\Delta,
\]
并在一般 Weierstrass 形式
\(
\eta^2+a_1x\eta+a_3\eta=x^3+a_2x^2+a_4x+a_6
\)
的不变量记号下对 \(E_0\) 取
\(
(a_1,a_2,a_3,a_4,a_6)=(0,0,1,-1,0)
\)，
从而 \(b_2=0\)、\(b_4=-2\)、\(b_6=1\)、\(b_8=-1\)，并由标准判别式公式得到 \(\Delta(E_0)=64-27=37\)，与上文短 Weierstrass 模型 \(E\) 的计算一致。
若在外部文献中以极小判别式 \(-37\) 或等价的相反号判别式记之，则对应于对本文 \(\Delta\) 采用不同的符号约定，并不改变本文以 \(\Delta(E)=37\) 为锚定的后续论证。
据 \cite{LMFDBEllipticCurve37a1} 的记录，在 Cremona/LMFDB 的标准标识下，\(E\) 为导子 \(37\) 的同源类 \(37\mathrm a\) 之强 Weil 曲线（记号 \(37\mathrm a1\)）；
其极小整系数模型 \(E_0:\eta^2+\eta=x^3-x\) 亦可与 Cremona 在 \(N=37\) 的算例对照。\cite{CremonaExamplesPDF}
\par
上述外部标识虽不参与本文的代数推导，但给出一条可直接核验的 Langlands 定量锚点：与该同源类对应的 weight \(2\)、level \(37\) 新形式轨道为 \(37.2.a.a\)（系数域 \(\QQ\)，解析秩 \(1\)），相应自守表示记为 \(\pi_{37}\)，并且对任意良素 \(p\neq 37\) 有
\(
a_p=p+1-\#E(\FF_p)
\)
与该新形式的 Hecke 本征值一致，从而 \(L(E,s)=L(f,s)=L(\pi_{37},s)\)（模性定理；亦可在 Langlands 口径下视为与 \(\mathrm{GL}_2(\mathbb{A}_\QQ)\) 的相应自守对象对接）。\cite{LMFDBNewform372aa}
\par
因此对每个素数 \(\ell\)，椭圆曲线 \(E\) 的 \(\ell\)-进 Tate 模诱导二维伽罗瓦表示
\[
\rho_{E,\ell}:\mathrm{Gal}(\overline{\QQ}/\QQ)\longrightarrow \mathrm{GL}_2(\ZZ_\ell),
\]
其仅在 \(\{37,\ell\}\) 处可能分歧；并且对任意良素 \(p\neq 37,\ell\)，Frobenius 元的特征多项式为 \(T^2-a_pT+p\)，亦即 \(T_p\) 的 Hecke 特征多项式，从而局部 Euler 因子与 Hecke 数据一致。
\par
并且由 Sturm 有界性可得该外部锚定具有有限证书：对 \(k=2,N=37\) 有 \([SL_2(\ZZ):\Gamma_0(37)]=38\)，故 Sturm 界
\(
B=\bigl\lfloor \frac{k}{12}[SL_2(\ZZ):\Gamma_0(N)]\bigr\rfloor
\)
等于 \(6\)；从而仅需核验至 \(q^6\) 的有限多个 Fourier/Hecke 系数即可唯一确定该新形式轨道。\cite{LMFDBSturmBound}
在模曲线口径下，\(X_0(37)\) 为属 \(2\) 曲线；LMFDB 记录 \(E_0\) 正可作为其 Fricke involution 的商
\(
X_0^+(37):=X_0(37)/\langle w_{37}\rangle
\)
的代数模型之一。\cite{LMFDBEllipticCurve37a1}
\par
令 \(q=e^{2\pi i\tau}\)。
在上述外部锚定下，可取规范化新形式 \(f(\tau)=\sum_{n\ge 1}a_n q^n\in S_2(\Gamma_0(37))\) 使得 \(E_0\) 的 N\'eron 微分 \(\omega=\dd x/(2\eta+1)\) 在模参数化下满足
\(
\omega=2\pi i\, f(q)\,\dd q/q
\)；
并且在谱坐标下有 \(2\eta+1=2y+1-2x^2=\partial_y\Pi(x,y)\)，从而 \(\omega=\dd x/\partial_y\Pi(x,y)\)。
于是重量函数 \(y=x^{2}+\eta\in\QQ(E_0)\)（等价地 \(y=x^{2}-\eta-1\)）作为 \(X_0^+(37)\) 上的模函数在无穷远尖点具有主部 \(q^{-4}\)，其截断展开与对应新形式的 \(q\)-展开见下式：
% Auto-generated by scripts/exp_fold_zm_elliptic_modular_y_qexp_audit.py
\[
y(\tau)=q^{-4}+5q^{-3}+17q^{-2}+47q^{-1}+117+266q+571q^{2}+1162q^{3}+O\!\left(q^{4}\right).
\]
\[
f(q)=q-2q^{2}-3q^{3}+2q^{4}-2q^{5}+6q^{6}-q^{7}+O\!\left(q^{8}\right).
\]

可复现核验脚本：\texttt{scripts/exp\_fold\_zm\_elliptic\_modular\_y\_qexp\_audit.py}。
\end{remark}

\begin{conclusion}[\texorpdfstring{$j$}{j}-不变量的标准基本域反演：纯虚模参数 \texorpdfstring{$\tau$}{tau} 的唯一性与数值证书]\label{con:xi-terminal-zm-elliptic-j-inversion-tau}
结论 \ref{con:xi-terminal-zm-elliptic-normalization} 给出
\(
j(E)=110592/37
\)。
则在 \(SL_{2}(\ZZ)\) 的标准基本域内，满足 \(j(\tau)=j(E)\) 的模参数 \(\tau\) 唯一，并且为纯虚数
\[
\tau=i\cdot 1.221127360764627\ldots.
\]
相应的
\(
q
\)-参数为
\[
q=e^{2\pi i\tau}=0.000465420392294374\ldots.
\]
因此 \(E(\CC)\) 可被显式均匀化为 \(\CC/(\ZZ+\tau\ZZ)\) 的适当复伸缩；并且权函数 \(y\) 在该均匀化下为阶 \(4\) 的椭圆函数，其解析尺度与判别式素数 \(37\) 的模性锚点一致。
\end{conclusion}
\begin{proof}[证明要点]
由于 \(j(E)\in\RR\) 且 \(j(E)>1728=j(i)\)，标准模论给出在基本域内的唯一解必落在纯虚轴上，即 \(\tau=it\)（\(t>1\)）。
对函数 \(t\mapsto j(it)\) 进行数值反演即可恢复 \(t\) 与 \(q=e^{-2\pi t}\)。
该反演可在高精度下通过 Eisenstein 级数的 \(q\)-展开实现，并由脚本 \texttt{scripts/exp\_fold\_zm\_elliptic\_modular\_y\_qexp\_audit.py} 复现核验。证毕。
\end{proof}

\begin{conclusion}[模素纤维根数恒等式：由 \texorpdfstring{$\Pi(\lambda,y)$}{Pi(lambda,y)} 的分解统计恢复 \texorpdfstring{$a_p$}{a\_p}]\label{con:xi-terminal-zm-fiber-root-counting-hecke}
取素数 \(p\notin\{2,37\}\)，并在 \(\FF_p\) 上考虑仿射曲线
\[
\mathcal C_p:=\{(\lambda,y)\in\FF_p^2:\ \Pi(\lambda,y)=0\}.
\]
对每个 \(y\in\FF_p\) 定义纤维根数
\[
N_p(y):=\card{\{\lambda\in\FF_p:\ \Pi(\lambda,y)=0\}}.
\]
则有严格点计数恒等式
\[
\boxed{\;
\card{E(\FF_p)}=1+\sum_{y\in\FF_p}N_p(y)\;}
\]
其中右端的 \(1\) 对应无穷远点 \(O\)。
令 \(a_p:=p+1-\card{E(\FF_p)}\) 为 Frobenius 迹（见注记 \ref{rem:xi-terminal-elliptic-37a1-modular-anchor}），则等价地
\[
\boxed{\;
a_p=p-\sum_{y\in\FF_p}N_p(y)\;}
\]
并由 Hasse 界得到有效估计
\[
\left|\frac{1}{p}\sum_{y\in\FF_p}N_p(y)-1\right|
=\left|\frac{a_p}{p}\right|
\le \frac{2}{\sqrt p}.
\]
更一般地，对任意 \(n\ge 1\)，上述双有理变换在 \(\FF_{p^n}\) 上同样给出
\[
\card{\{(\lambda,y)\in\FF_{p^n}^2:\ \Pi(\lambda,y)=0\}}=\card{E(\FF_{p^n})}-1.
\]
因此若定义仿射谱曲线的局部 \(\zeta\)-函数为
\[
Z(\mathcal C/\FF_p,T):=\exp\!\left(\sum_{n\ge 1}\card{\mathcal C(\FF_{p^n})}\frac{T^n}{n}\right),
\]
则有闭式
\[
\boxed{\; Z(\mathcal C/\FF_p,T)=\frac{1-a_pT+pT^2}{1-pT}\; }.
\]
\end{conclusion}
\begin{proof}[证明要点]
当 \(p\notin\{2,37\}\) 时，结论 \ref{con:xi-terminal-zm-elliptic-normalization} 的双有理变换
\[
(\lambda,y)\longmapsto (\lambda,w),\qquad w:=2y+1-2\lambda^{2}
\]
在 \(\FF_p\) 上仍可逆（因为 \(2\) 可逆），并把 \(\mathcal C_p\) 双射到仿射椭圆曲线
\(
\mathcal E_p:\ w^2=4\lambda^3-4\lambda+1
\)
的 \(\FF_p\)-点集；该仿射模型与射影椭圆曲线 \(E\) 的差别仅为无穷远点 \(O\)。
因此
\[
\sum_{y\in\FF_p}N_p(y)=\card{\mathcal C_p}=\card{E(\FF_p)}-1,
\]
从而得到所示恒等式。
将其代入 \(a_p:=p+1-\card{E(\FF_p)}\) 立得 \(a_p=p-\sum_{y\in\FF_p} N_p(y)\)。
最后一式由 Hasse 界 \(|a_p|\le 2\sqrt p\) 给出。
可复现核验脚本：\texttt{scripts/exp\_fold\_zm\_elliptic\_fiber\_root\_counting\_audit.py}。证毕。
\end{proof}

\begin{conclusion}[驯约化素数上的分解型等分布：\texorpdfstring{$S_4$}{S4} 共轭类密度与平方根误差]\label{con:xi-terminal-zm-s4-modp-splitting-density}
沿用定理 \ref{thm:xi-terminal-zm-leyang-elliptic-structure} 与结论 \ref{con:xi-terminal-zm-generic-galois-s4} 的几何单值群 \(\mathrm{Mon}(\pi_y)\cong S_4\)，
并令
\(
B:=\{0,1,\infty\}\cup\{P_{\mathrm{LY}}(y)=0\}\subset \mathbb P^1_y
\)
为覆盖 \(\pi_y:E\to\mathbb P^1_y\) 的分歧值集合（见结论 \ref{con:xi-terminal-zm-discriminant-leyang}）。
取素数 \(p\notin\{2,3,31,37\}\)，并在 \(\FF_p\) 上考虑未分歧参数集合
\[
U_p:=\FF_p\setminus\bigl(B\cap \FF_p\bigr)
=\Bigl\{y_0\in\FF_p:\ y_0(y_0-1)P_{\mathrm{LY}}(y_0)\neq 0\Bigr\}.
\]
则当 \(y_0\) 在 \(U_p\) 上作均匀分布时，
四次多项式 \(\Pi(\lambda,y_0)\in\FF_p[\lambda]\) 的分解型按 \(S_4\) 的共轭类密度等分布，并具有 \(O(p^{-1/2})\) 的函数域 Weil 型误差；
更具体地，以“分解型 \(\leftrightarrow\) Frobenius 在四根上的循环型”对应，则五种分解型的极限密度为
\begin{itemize}
  \item 不可约四次（循环型 \((4)\)）：密度 \(6/24=1/4\)；
  \item 一次 \(\times\) 不可约三次（循环型 \((3+1)\)）：密度 \(8/24=1/3\)；
  \item 两个不可约二次（循环型 \((2+2)\)）：密度 \(3/24=1/8\)；
  \item 两个一次 \(\times\) 一个不可约二次（循环型 \((2+1+1)\)）：密度 \(6/24=1/4\)；
  \item 全分裂（循环型 \((1+1+1+1)\)）：密度 \(1/24\)。
\end{itemize}
该结论与结论 \ref{con:xi-terminal-zm-fiber-root-counting-hecke} 的平均纤维根数恒等式相容：在未分歧纤维上，\(\Pi(\lambda,y_0)\) 的 \(\FF_p\)-根数均值趋于 \(1\)，其全局偏差由 \(a_p\) 控制。
\end{conclusion}
\begin{proof}[证明要点]
对 \(p\notin\{2,3,31,37\}\)，结论 \ref{con:xi-terminal-zm-discriminant-leyang} 保证 \(B\) 在模 \(p\) 下仍为六个互异分歧值（含 \(\infty\)），且分歧指数均与 \(p\) 互素，从而覆盖在 \(U_p\) 上驯分歧。
由结论 \ref{con:xi-terminal-zm-generic-galois-s4} 的 \(S_4\) 单值性，取其在 \(\FF_p\) 上的 \(S_4\)-伽罗瓦闭包并对 \(U_p\) 的有理点应用函数域 Chebotarev 定理，
可得 Frobenius 共轭类在 \(S_4\) 中按类大小 \(|C|/24\) 等分布；误差项由相应 Galois 覆叠曲线的 Weil 界给出，量级为 \(O(\sqrt p)\)，折算到比例即为 \(O(p^{-1/2})\)。
循环型与 \(\FF_p[\lambda]\) 分解型的对应为标准事实。
可复现核验脚本：\texttt{scripts/exp\_fold\_zm\_elliptic\_fiber\_root\_counting\_audit.py}（输出中包含分解型频率审计）。证毕。
\end{proof}

\begin{conclusion}[未分歧纤维根数的固定点分布：\texorpdfstring{$S_4$}{S4} 四点作用的不动点密度与平方根误差]\label{con:xi-terminal-zm-s4-modp-fiber-rootcount-fixedpoints}
沿用结论 \ref{con:xi-terminal-zm-s4-modp-splitting-density} 的记号。
对 \(y_0\in\FF_p\) 定义纤维根数
\[
N_p(y_0):=\card{\{\lambda\in\FF_p:\ \Pi(\lambda,y_0)=0\}}.
\]
则存在与 \(p\) 无关的常数 \(C>0\)，使得对 \(r\in\{0,1,2,4\}\) 成立
\[
\Bigl|\card{\{y_0\in\FF_p:\ N_p(y_0)=r\}}-\pi_r\,p\Bigr|\le C\sqrt p,
\]
其中 \(\pi_r\) 为 \(S_4\) 在自然四点作用下“恰有 \(r\) 个不动点”的元素比例：
\[
\pi_0=\frac{3}{8},\qquad
\pi_1=\frac{1}{3},\qquad
\pi_2=\frac{1}{4},\qquad
\pi_4=\frac{1}{24}.
\]
此外，
\[
\card{\{y_0\in\FF_p:\ N_p(y_0)=3\}}=O(1),
\]
且任一满足 \(N_p(y_0)=3\) 的参数必落在分歧集合的模 \(p\) 化 \(B\cap\FF_p\) 上（对应 \(\mathrm{Disc}_\lambda(\Pi(\lambda,y_0))=0\) 的特化并出现重根）。
\end{conclusion}
\begin{proof}[证明要点]
对 \(y_0\in U_p\)，四次多项式 \(\Pi(\lambda,y_0)\) 可分且无重根，其 Frobenius 置换型在四根上给出 \(S_4\) 中的共轭类；
并且 \(N_p(y_0)\) 等于该置换在四点集合上的不动点数。
因此该结论可由结论 \ref{con:xi-terminal-zm-s4-modp-splitting-density} 的分解型等分布经“取线性因子数（不动点数）”投影得到，误差项仍为 \(O(\sqrt p)\)，从而给出所示的 \(O(\sqrt p)\) 有效界。
由于 \(\card{\FF_p\setminus U_p}=\card{B\cap\FF_p}\le 6\)，将计数从 \(U_p\) 扩展到 \(\FF_p\) 仅引入 \(O(1)\) 的差异，可吸收进 \(C\sqrt p\)。
\(\pi_r\) 的数值由 \(S_4\) 中不动点数的枚举得到（等价地，把分解型密度按不动点数归并）。
对 \(N_p(y_0)=3\)：
若 \(y_0\notin B\cap\FF_p\)，则 \(\Pi(\lambda,y_0)\) 无重根且次数为 \(4\)，不可能在 \(\FF_p\) 上出现恰三根情形；
故 \(N_p(y_0)=3\) 只能发生在 \(B\cap\FF_p\) 上，而 \(\card{B\cap\FF_p}\) 由 \(\deg(y(y-1)P_{\mathrm{LY}})\) 有界，遂为 \(O(1)\)。
可复现核验脚本：\texttt{scripts/exp\_fold\_zm\_s4\_chebotarev\_splitting\_distribution\_audit.py}。证毕。
\end{proof}

\begin{conclusion}[覆盖的最小坏约化素集与分歧值碰撞的精确控制]\label{con:xi-terminal-zm-cover-minimal-bad-reduction-set}
沿用结论 \ref{con:xi-terminal-zm-discriminant-leyang} 的分歧值集合 \(B\subset\mathbb P^1_y\)，其整系数模型由
\[
y(y-1)P_{\mathrm{LY}}(y)=0
\]
给出（其中 \(P_{\mathrm{LY}}(y)\) 为 \(\mathrm{Disc}_\lambda(\Pi)\) 的三次因子）。
则素数集合 \(\{2,3,31,37\}\) 构成该谱覆盖的最小坏约化集，意义如下：
\begin{itemize}
  \item 若 \(p\notin\{2,3,31,37\}\)，则 \(B\) 在模 \(p\) 下保持互异，且覆盖在 \(\FF_p\) 上为驯分歧并保持 Hurwitz 护照 \((2,2,2,2,2,4)\) 不变，从而给出良约化；
  \item 若 \(p=3\)，则 \(P_{\mathrm{LY}}(y)\equiv (y-1)^3\ (\mathrm{mod}\ 3)\)，三条非平凡分歧值在模 \(3\) 下三重碰撞，从而覆盖分歧数据退化；
  \item 若 \(p=31\)，则
  \[
  P_{\mathrm{LY}}(y)\equiv 8\,(y-12)^{2}(y-6)\pmod{31},
  \]
  从而在模 \(31\) 下出现二重根并导致分歧值二重碰撞，覆盖分歧数据退化；
  \item 若 \(p=37\)，则
  \[
  P_{\mathrm{LY}}(y)\equiv -3\,(y-14)(y-6)^{2}\pmod{37},
  \]
  并且椭圆曲线 \(E\) 发生乘法型坏约化（\(\Delta(E)=37\)）；因此既有分歧值碰撞又有基曲线坏约化，覆盖随之退化。
\end{itemize}
因此在所有 \(p\notin\{2,3,31,37\}\) 处，该谱覆盖处于“结构稳定”的驯约化情形。
\end{conclusion}
\begin{proof}[证明要点]
由结论 \ref{con:xi-terminal-zm-discriminant-leyang}，
\(
\mathrm{Disc}_\lambda(\Pi)=-y(y-1)P_{\mathrm{LY}}(y)
\)，
故分歧值的合并当且仅当 \(P_{\mathrm{LY}}\) 的模 \(p\) 约化不再可分。
又由定理 \ref{thm:xi-terminal-zm-leyang-elliptic-structure} 的判别式恒等式
\(
\mathrm{Disc}(P_{\mathrm{LY}})=-3^{9}\cdot 31^{2}\cdot 37
\)，
可知仅在 \(p\in\{3,31,37\}\) 处可能发生分歧值碰撞；并由直接模约化可区分三重与二重重根情形。
另一方面，\(E\) 的极小模型判别式为 \(\Delta(E)=37\)（结论 \ref{con:xi-terminal-zm-elliptic-normalization}），故 \(p=37\) 为唯一坏约化素数；
而 \(p=2\) 时变换 \(w=2y+1-2\lambda^2\) 不再可逆并引入特征 \(2\) 的不可分退化。
除去上述有限素集后，分歧指数与 \(p\) 互素且分歧值互异，从而覆盖保持驯分歧并具有良约化。证毕。
\end{proof}

\begin{conclusion}[模形式 Hecke 特征值的纤维统计实现：总偏差恒等式]\label{con:xi-terminal-zm-hecke-eigenvalue-fiber-deviation}
对素数 \(p\notin\{2,37\}\)，沿用结论 \ref{con:xi-terminal-zm-fiber-root-counting-hecke} 的纤维根数 \(N_p(y)\)，并定义纤维差分函数
\[
\delta_p(y):=N_p(y)-1.
\]
则有严格恒等式
\[
\boxed{\;\sum_{y\in\FF_p}\delta_p(y)=-a_p\;}
\]
其中 \(a_p=p+1-\#E(\FF_p)\) 为椭圆曲线 \(E\) 的 Frobenius 迹。
在注记 \ref{rem:xi-terminal-elliptic-37a1-modular-anchor} 的模性锚定下，\(a_p\) 同时等于对应权 \(2\)、level \(37\) 新形式的 \(T_p\)-Hecke 本征值；
从而该本征值可被解释为四次谱方程模 \(p\) 的“纤维根数相对均值 \(1\)”之总偏差。
\end{conclusion}
\begin{proof}[证明要点]
由结论 \ref{con:xi-terminal-zm-fiber-root-counting-hecke} 的恒等式
\(
a_p=p-\sum_{y\in\FF_p}N_p(y)
\)，
移项即得
\(
\sum_{y\in\FF_p}(N_p(y)-1)=-a_p
\)。
Hecke 本征值解释由模性锚定给出（见注记 \ref{rem:xi-terminal-elliptic-37a1-modular-anchor}）。证毕。
\end{proof}

\begin{conclusion}[落入 \texorpdfstring{$A_4$}{A4} 的专化判别：判别式平方条件与属 \texorpdfstring{$2$}{2} 曲线的有理点约束]\label{con:xi-terminal-zm-specialization-a4-square-criterion}
设 \(K\) 为特征 \(0\) 的域，并取 \(y_{0}\in K\setminus\bigl(\{0,1\}\cup\{P_{\mathrm{LY}}(y)=0\}\bigr)\)。
假设专化四次 \(\Pi(\lambda,y_{0})\in K[\lambda]\) 在 \(K\) 上不可约。
则成立判别准则
\[
\boxed{\;
\mathrm{Gal}\bigl(\Pi(\lambda,y_{0})/K\bigr)\subseteq A_{4}
\quad\Longleftrightarrow\quad
-y_{0}(y_{0}-1)P_{\mathrm{LY}}(y_{0})\in K^{\times 2}\;}
\]
等价地，存在 \(w_{0}\in K\) 使
\[
w_{0}^{2}=-y_{0}(y_{0}-1)P_{\mathrm{LY}}(y_{0}),
\]
亦即 \((y_{0},w_{0})\in C(K)\)，其中属 \(2\) 曲线
\[
C:\quad w^{2}=-y(y-1)P_{\mathrm{LY}}(y)
\]
为结论 \ref{con:xi-terminal-zm-discriminant-leyang} 的判别式二次脊曲线。
因此，当 \(K=\QQ\) 时，“落入交错群”的全部有理专化参数被严格嵌入到一条属 \(2\) 曲线的有理点集合中，并因而在算术意义上受到强刚性约束（例如可由 Faltings 定理推出其有理点集合有限）。
\end{conclusion}
\begin{proof}[证明要点]
在 \(y_{0}\) 避开分歧集合的假设下，\(\Pi(\lambda,y_{0})\) 的四根互异，故其专化伽罗瓦群可视为 \(S_{4}\) 的一条传递子群。
对不可约四次而言，伽罗瓦群落入交错群当且仅当判别式在底域中为平方。
而由结论 \ref{con:xi-terminal-zm-discriminant-leyang}，
\(
\mathrm{Disc}_{\lambda}\bigl(\Pi(\cdot,y_{0})\bigr)=-y_{0}(y_{0}-1)P_{\mathrm{LY}}(y_{0})
\)，
从而得到所示等价；曲线 \(C\) 的刻画为直接改写。证毕。
\end{proof}

\begin{conclusion}[判别式平方约束沿椭圆有理轨道的有限性：属 \texorpdfstring{$6$}{6} 纤维积曲线与 Mordell--Faltings 刚性]\label{con:xi-terminal-zm-discriminant-square-elliptic-orbit-finiteness}
沿用定理 \ref{thm:xi-terminal-zm-leyang-elliptic-structure} 的椭圆曲线 \(E/\QQ\) 与次数 \(4\) 映射 \(\pi_y:E\to\PP^{1}_{y}\)，并沿用结论 \ref{con:xi-terminal-zm-specialization-a4-square-criterion} 的属 \(2\) 判别式二次脊曲线
\[
C:\quad w^{2}=-y(y-1)P_{\mathrm{LY}}(y).
\]
令 \(D\) 为纤维积 \(E\times_{\PP^{1}_{y}} C\) 的正规化（等价地，\(\QQ(D)=\QQ(E)(w)\) 且 \(w^{2}=-y(y-1)P_{\mathrm{LY}}(y)\)）。
则 \(D\) 为属 \(6\) 的光滑射影曲线，从而 \(D(\QQ)\) 为有限集。
因此集合
\[
\Bigl\{\,Q\in E(\QQ):\ -y(Q)\bigl(y(Q)-1\bigr)P_{\mathrm{LY}}\bigl(y(Q)\bigr)\in \QQ^{\times 2}\Bigr\}
\]
为有限集；特别地，对任意无穷阶点 \(Q_{0}\in E(\QQ)\)，整数集合
\[
\Bigl\{\,n\in\ZZ:\ -y(nQ_{0})\bigl(y(nQ_{0})-1\bigr)P_{\mathrm{LY}}\bigl(y(nQ_{0})\bigr)\in \QQ^{\times 2}\Bigr\}
\]
亦为有限集。
\end{conclusion}
\begin{proof}[证明要点]
由于 \(D\) 的函数域由 \(E\) 的函数域加入平方根 \(w^{2}=-y(y-1)P_{\mathrm{LY}}(y)\) 得到，投影 \(D\to E\) 为次数 \(2\) 覆盖。
其分歧点恰为函数 \(-y(y-1)P_{\mathrm{LY}}(y)\in\QQ(E)^{\times}\) 的零极点中具有奇数阶的支撑点。
由结论 \ref{con:xi-terminal-zm-elliptic-y-hurwitz-passport} 的主除子分解可知：在纤维 \(y=0\) 与 \(y=1\) 上各有两点以奇数重数出现（分歧点本身以重数 \(2\) 出现，故不贡献奇数阶），从而给出 \(4\) 个分歧点。
另一方面，由定理 \ref{thm:xi-terminal-zm-leyang-elliptic-structure}，对每个分歧值 \(y_{*}\) 满足 \(P_{\mathrm{LY}}(y_{*})=0\)，纤维 \(y=y_{*}\) 的型为 \((2,1,1)\)，因此在 \(E\) 上的除子 \(\pi_y^{*}(y-y_{*})\) 中恰有两点以奇数重数出现；三条 Lee--Yang 分歧值合计给出 \(6\) 个分歧点。
此外，由 \((y)_\infty=4O\) 得 \(y\) 在 \(O\) 处极阶为 \(4\)，故 \(-y(y-1)P_{\mathrm{LY}}(y)\) 在 \(O\) 处的极阶为 \(4+4+12=20\)（偶数），无穷远处不分歧。
综上，\(D\to E\) 的分歧点总数为 \(4+6=10\)，且皆为单分歧。
由 Riemann--Hurwitz 公式（\(g(E)=1\)）得到
\[
2g(D)-2=2(2g(E)-2)+10=10,
\]
故 \(g(D)=6\)。
由 Mordell--Faltings 定理，属 \(>1\) 的曲线 \(D/\QQ\) 满足 \(D(\QQ)\) 有限。
将 \(D(\QQ)\) 通过投影映到 \(E(\QQ)\) 即得所述有限性；而对无穷阶点 \(Q_0\) 的倍点轨道，因 \(\{nQ_0\}_{n\in\ZZ}\) 互异，故其与上述有限集的交亦为有限。证毕。
\end{proof}

\begin{conclusion}[\texorpdfstring{$A_4$}{A4} 特化的符号通道平凡化：符号扭曲 Ihara--Artin 因子与跨支比值的退化]\label{con:xi-terminal-zm-a4-specialization-sign-channel-trivialization}
沿用结论 \ref{con:xi-terminal-zm-specialization-a4-square-criterion} 的设定，并令 \(\beta_{\mathrm{LY}}\) 表示由符号表示 \(\mathrm{sgn}:S_4\to\{\pm1\}\) 在专化伽罗瓦群 \(\mathrm{Gal}(\Pi(\lambda,y_0)/K)\) 上诱导的秩 \(1\) \(\ZZ_2\) 局部系统。
记 \(\zeta(u,t)\) 与 \(\zeta_{\beta_{\mathrm{LY}}}(u,t)\) 分别为未扭曲与 \(\beta_{\mathrm{LY}}\)-扭曲的 Ihara--Artin 轨道乘积（例如可取结论 \ref{con:xi-split-tunneling-zeta-detect-aflip} 的 Euler 乘积口径）。
若
\[
\mathrm{Gal}\bigl(\Pi(\lambda,y_{0})/K\bigr)\subseteq A_{4},
\]
则对所有 primitive 轨道 \(p\) 有 \(\chi_{\beta_{\mathrm{LY}}}(p)=+1\)，从而
\[
\zeta_{\beta_{\mathrm{LY}}}(u,t)\equiv \zeta(u,t),
\qquad
R(u,t):=\frac{\zeta_{\beta_{\mathrm{LY}}}(u,t)}{\zeta(u,t)}\equiv 1.
\]
在结论 \ref{con:xi-split-tunneling-zeta-detect-aflip} 的记号下，这等价于 \(A_{\mathrm{flip}}=+\infty\)（即不存在任何 \(\chi_{\beta_{\mathrm{LY}}}=-1\) 的 Ihara prime 环）。
\end{conclusion}
\begin{proof}[证明要点]
当 \(\mathrm{Gal}(\Pi(\lambda,y_0)/K)\subseteq A_4\) 时，符号角色在该群上恒为平凡，故任意由该单值化诱导的 primitive 轨道之 \(\ZZ_2\) holonomy 皆为 \(+1\)，从而扭曲与未扭曲 Euler 因子逐项相同并给出 \(\zeta_{\beta_{\mathrm{LY}}}\equiv \zeta\)。
按结论 \ref{con:xi-split-tunneling-zeta-detect-aflip} 中 \(\min\varnothing:=+\infty\) 的约定，此时 \(A_{\mathrm{flip}}=+\infty\)。证毕。
\end{proof}

\begin{conclusion}[Hitchin--cameral 取向通道：\texorpdfstring{$W(A_3)\cong S_4$}{W(A3)=S4} 的符号特征给出 \texorpdfstring{$\beta_{\mathrm{LY}}$}{beta_LY} 与行列式比]\label{con:xi-terminal-zm-hitchin-cameral-orientation-channel}
沿注 \ref{rem:xi-terminal-zm-hitchin-cameral-cover} 的解释，把四次谱覆盖 \(\Pi(\lambda,y)=0\) 的伽罗瓦闭包视为 \(S_4\simeq W(A_3)\) 的 cameral 结构。
则 Weyl 群的符号角色 \(\mathrm{sgn}\) 为其取向特征，并诱导一条规范的秩 \(1\) 取向局部系统 \(\mathcal L_{\mathrm{or}}\)。
在附录命题 \ref{prop:ihara-artin-factorization} 的群扩张 Hashimoto 口径下，\(\rho=\mathrm{sgn}\) 的扭曲通道即为该取向通道，亦即
\[
B_{\mathcal L_{\mathrm{or}}}(u)=B_{\mathrm{sgn}}(u).
\]
因此若定义
\[
\zeta(u,t):=\det\bigl(I-tB_{\mathbf 1}(u)\bigr)^{-1},
\qquad
\zeta_{\beta_{\mathrm{LY}}}(u,t):=\det\bigl(I-tB_{\mathcal L_{\mathrm{or}}}(u)\bigr)^{-1}
=\det\bigl(I-tB_{\mathrm{sgn}}(u)\bigr)^{-1},
\]
则跨支比值可提升为行列式比
\[
\boxed{\;
\frac{\zeta_{\beta_{\mathrm{LY}}}(u,t)}{\zeta(u,t)}
=
\frac{\det\bigl(I-tB_{\mathbf 1}(u)\bigr)}{\det\bigl(I-tB_{\mathcal L_{\mathrm{or}}}(u)\bigr)}
=
\frac{\det\bigl(I-tB_{\mathbf 1}(u)\bigr)}{\det\bigl(I-tB_{\mathrm{sgn}}(u)\bigr)}\; }.
\]
从而 cyclotomic 共振与端点通量等可审计指标可在取向通道内独立抽取（例如定理 \ref{thm:conclusion75-ihara-cyclotomic-factor-classification}、结论 \ref{con:xi-endpoint-flux-operator-trace-identity} 与结论 \ref{con:xi-endpoint-flux-artin-additivity}）。
\end{conclusion}
\begin{proof}[证明要点]
由注 \ref{rem:xi-terminal-zm-hitchin-cameral-cover}，\(S_4\simeq W(A_3)\) 的 Weyl 对称为谱覆盖的内禀单值化；其符号角色即取向特征并给出 \(\mathcal L_{\mathrm{or}}\)。
对群扩张 Hashimoto 算子取 \(\rho=\mathrm{sgn}\) 的不可约分量即实现该取向扭曲，故 \(B_{\mathcal L_{\mathrm{or}}}=B_{\mathrm{sgn}}\)。
最后一式由 Ihara--Artin 因子定义
\(
\zeta_{\rho}(u,t)=\det(I-tB_\rho(u))^{-1}
\)
取比值得到。证毕。
\end{proof}

\begin{conclusion}[判别式脊点诱导的循环三次分量：由三次预解式生成 \texorpdfstring{$C_3$}{C3} 专化]\label{con:xi-terminal-zm-specialization-a4-c3-resolvent}
设 \(K\) 为特征 \(0\) 的域，并取 \((y_{0},w_{0})\in C(K)\) 满足
\[
y_{0}\notin \{0,1\}\cup\{P_{\mathrm{LY}}(y)=0\}.
\]
假设专化四次 \(\Pi(\lambda,y_{0})\in K[\lambda]\) 在 \(K\) 上不可约。
令 \(\mathscr R(z,y)\in\ZZ[z,y]\) 为结论 \ref{con:xi-terminal-zm-generic-galois-s4} 的三次预解式，并记 \(\mathscr R_{y_{0}}(z):=\mathscr R(z,y_{0})\in K[z]\)。
则：
\begin{enumerate}
  \item \(\mathrm{Disc}_{z}\bigl(\mathscr R_{y_{0}}\bigr)=w_{0}^{2}\in K^{\times 2}\)，从而 \(\mathrm{Gal}\bigl(\mathscr R_{y_{0}}/K\bigr)\subseteq A_{3}\cong C_{3}\)；
  \item 若 \(\mathscr R_{y_{0}}(z)\) 在 \(K\) 上不可约，则其分裂域为循环三次扩张 \(K_{3}/K\)，并且
  \[
  \mathrm{Gal}\bigl(\Pi(\lambda,y_{0})/K\bigr)\cong A_{4}.
  \]
  在此情形下，\(K_{3}\) 可识别为该 \(A_{4}\)-扩张中对应正规 Klein 四元子群 \(V_{4}\triangleleft A_{4}\) 的不动域（即 \(A_{4}/V_{4}\cong C_{3}\) 的循环三次分量）。
\end{enumerate}
因此，在判别式平方约束下，属 \(2\) 判别式脊曲线 \(C\) 的 \(K\)-有理点为从专化四次谱方程中抽取循环三次扩张的显式接口。
\end{conclusion}
\begin{proof}[证明要点]
由结论 \ref{con:xi-terminal-zm-generic-galois-s4} 的判别式恒等式
\(
\mathrm{Disc}_{z}\bigl(\mathscr R(\cdot,y)\bigr)=-y(y-1)P_{\mathrm{LY}}(y)
\)，
代入 \(y=y_{0}\) 并利用 \((y_{0},w_{0})\in C(K)\) 得
\[
\mathrm{Disc}_{z}\bigl(\mathscr R_{y_{0}}\bigr)=-y_{0}(y_{0}-1)P_{\mathrm{LY}}(y_{0})=w_{0}^{2}\in K^{\times 2},
\]
故 \(\mathrm{Gal}(\mathscr R_{y_{0}}/K)\subseteq A_{3}\cong C_{3}\)；若 \(\mathscr R_{y_{0}}\) 在 \(K\) 上不可约，则必有 \(\mathrm{Gal}(\mathscr R_{y_{0}}/K)\cong C_{3}\)，从而其分裂域为循环三次扩张。
\par
另一方面，由结论 \ref{con:xi-terminal-zm-specialization-a4-square-criterion}，判别式为平方给出
\(
\mathrm{Gal}(\Pi(\lambda,y_{0})/K)\subseteq A_{4}
\)。
对不可约四次而言，三次预解式在底域上不可约将伽罗瓦群约束为 \(\{A_{4},S_{4}\}\)（同结论 \ref{con:xi-terminal-zm-generic-galois-s4} 的判别）；而此处已知 \(\mathrm{Gal}(\Pi(\lambda,y_{0})/K)\subseteq A_{4}\)，故必有
\(
\mathrm{Gal}(\Pi(\lambda,y_{0})/K)\cong A_{4}
\)。
最后，三次预解式的构造使其分裂域与 \(V_{4}\) 的不动域一致：其三根对应四根之三种二二配对不变量，从而给出商 \(A_{4}/V_{4}\cong C_{3}\) 的循环三次分量。证毕。
\end{proof}



\begin{conclusion}[Lee--Yang 三次多项式的分歧像刻画：结果式给出的最小整系数证书]\label{con:xi-terminal-zm-leyang-cubic-branch-image}
把 \(\Pi(\lambda,y)=0\) 视为关于 \(\lambda\) 的四次方程，并令投影
\(
\pi:\mathcal C\to\mathbb A^1_y
\)
由 \((\lambda,y)\mapsto y\) 给出。
则 \(\pi\) 的分歧值集合在 \(\ZZ[y]\) 中由结果式精确刻画：
\[
\boxed{\;
\mathrm{Res}_{\lambda}\!\bigl(\Pi(\lambda,y),\partial_{\lambda}\Pi(\lambda,y)\bigr)
=-y(y-1)\bigl(256y^3+411y^2+165y+32\bigr)\; }.
\]
因此除平凡分歧值 \(y\in\{0,1\}\) 外，全部非平凡分歧值恰为三次多项式
\(
256y^3+411y^2+165y+32
=0
\)
的三根；其中唯一负实根即结论 \ref{con:xi-terminal-zm-discriminant-leyang} 的 \(y_{\mathrm{LY}}\)。
该三次因子因而并非经验拟合对象，而是投影 \(\pi\) 的分歧像在整系数层面的最小消元证书。
\end{conclusion}
\begin{proof}[证明要点]
分歧点当且仅当存在 \(\lambda\) 使得 \(\Pi(\lambda,y)=\partial_\lambda\Pi(\lambda,y)=0\)。
对 \(\lambda\) 作消元即可得到结果式分解；由于 \(\Pi\) 首一，该结果式与 \(\mathrm{Disc}_{\lambda}(\Pi)\) 一致，从而与结论 \ref{con:xi-terminal-zm-discriminant-leyang} 的判别式因式分解完全相容。证毕。
\end{proof}

\begin{conclusion}[分歧点的 \texorpdfstring{$\lambda$}{lambda}-参数化与 Lee--Yang 实根的同一化]\label{con:xi-terminal-zm-leyang-lambda-cubic}
设 \((\lambda,y)\) 满足谱曲线 \(\Pi(\lambda,y)=0\)。
则全部有限分歧点（即存在二重根之 \(y\)）由联立方程
\[
\Pi(\lambda,y)=0,\qquad \partial_{\lambda}\Pi(\lambda,y)=0
\]
所刻画，并且消元得到分歧所发生之谱坐标必满足
\[
\boxed{\;(\lambda-1)(\lambda+1)\bigl(16\lambda^3-9\lambda^2+1\bigr)=0\;}
\]
且相应分歧值由显式有理式给出
\[
\boxed{\;
y=\frac{4\lambda^3-3\lambda^2-2\lambda+1}{4\lambda}
=\frac{(\lambda-1)(4\lambda^2+\lambda-1)}{4\lambda}\; }.
\]
其中 \(\lambda=1\) 与 \(\lambda=-1\) 分别对应分歧值 \(y=0\) 与 \(y=1\)；
其余三根对应的分歧值恰为三次方程 \(256y^3+411y^2+165y+32=0\) 的三根（结论 \ref{con:xi-terminal-zm-leyang-cubic-branch-image}）。
特别地，结论 \ref{con:xi-terminal-zm-discriminant-leyang} 的唯一负实分歧值 \(y_{\mathrm{LY}}\) 对应于三次方程 \(16\lambda^3-9\lambda^2+1=0\) 的唯一实根 \(\lambda_{\mathrm{LY}}<0\)。
从而获得对 \(y_{\mathrm{LY}}\) 的等价刻画：
\[
\lambda_{\mathrm{LY}}\approx -0.2734348106524542,
\qquad y_{\mathrm{LY}}=y(\lambda_{\mathrm{LY}}),
\qquad 16\lambda_{\mathrm{LY}}^3-9\lambda_{\mathrm{LY}}^2+1=0.
\]
因此 \(y_{\mathrm{LY}}\) 可等价刻画为有理映射 \(y=y(\lambda)\) 在三次证书 \(16\lambda^3-9\lambda^2+1=0\) 的唯一实根处的取值。
\end{conclusion}
\begin{proof}[证明要点]
由 \(\partial_{\lambda}\Pi(\lambda,y)=0\) 可解出
\(
y=(4\lambda^3-3\lambda^2-2\lambda+1)/(4\lambda)
\)（\(\lambda\neq 0\)）。
代回 \(\Pi(\lambda,y)=0\) 并清分母得到
\(
(\lambda-1)(\lambda+1)(16\lambda^3-9\lambda^2+1)=0
\)；
其中 \(\lambda=\pm 1\) 分别对应分歧值 \(y=0,1\)；
其余三根由三次因子刻画，并经 \(y=y(\lambda)\) 映射落入结论 \ref{con:xi-terminal-zm-leyang-cubic-branch-image} 的三次分歧像，从而得到 Lee--Yang 实分歧值与伴随谱根的同一化。证毕。
\end{proof}

\begin{conclusion}[Lee--Yang 三次场的判别式与二次子域：\texorpdfstring{$\QQ(\sqrt{-111})$}{Q(sqrt(-111))}]\label{con:xi-terminal-zm-leyang-cubic-arithmetic}
令
\(
g(y):=256y^3+411y^2+165y+32
\)。
则其判别式满足严格分解
\[
\boxed{\;\mathrm{Disc}(g)=-3^{9}\cdot 31^{2}\cdot 37\;}
\]
且 \(g\) 在 \(\QQ[y]\) 上不可约，从而其伽罗瓦群为 \(S_3\)。
记 \(L\) 为 \(g\) 的分裂域，则其唯一二次子域为
\[
\boxed{\ \QQ(\sqrt{\mathrm{Disc}(g)})=\QQ(\sqrt{-111})\ },
\]
其中 \(-111\) 为 \(\mathrm{Disc}(g)\) 的平方自由部分，亦为该虚二次域的基本判别式。
\end{conclusion}
\begin{proof}[证明要点]
判别式因子分解由对三次多项式的标准判别式公式直接计算并在 \(\ZZ\) 中分解得到。
由有理根判据可排除 \(g\) 的有理根，从而 \(g\) 不可约；对不可约三次而言，判别式非平方当且仅当伽罗瓦群为 \(S_3\)。
二次子域由 \(\QQ(\sqrt{\mathrm{Disc}(g)})\) 给出；由于 \(\mathrm{Disc}(g)\) 的平方自由部分为 \(-111\)，故该二次域为 \(\QQ(\sqrt{-111})\)。
可复现核验脚本：\texttt{scripts/exp\_fold\_zm\_elliptic\_leyang\_arithmetic\_audit.py}。证毕。
\end{proof}

\begin{conclusion}[Lee--Yang 三次在 \texorpdfstring{$\FF_p$}{Fp} 上的分解型密度律：由二次特征 \texorpdfstring{$\left(\frac{-111}{p}\right)$}{(-111/p)} 完全控制“恰有一根”情形]\label{con:xi-terminal-zm-leyang-cubic-modp-splitting-density}
令
\[
g(y):=256y^3+411y^2+165y+32\in\ZZ[y],
\]
并取素数 \(p\nmid 3\cdot 31\cdot 37\)。
记 \(\varepsilon_p:=\left(\frac{-111}{p}\right)\in\{\pm1\}\) 为二次特征（Legendre 符号），并把 \(g\) 模 \(p\) 化为 \(g_p\in\FF_p[y]\)。
则 \(g_p\) 的分解型满足：
\begin{itemize}
  \item 若 \(\varepsilon_p=-1\)，则 \(g_p\) 必为 \((1)(2)\) 型，从而在 \(\FF_p\) 上恰有一个根；
  \item 若 \(\varepsilon_p=+1\)，则 \(g_p\) 要么全分裂（\((1)(1)(1)\) 型，三根），要么不可约（\((3)\) 型，无根）。
\end{itemize}
并且上述三类素数集合的自然密度分别为
\[
\delta\bigl((1)(2)\bigr)=\frac12,\qquad
\delta\bigl((1)(1)(1)\bigr)=\frac16,\qquad
\delta\bigl((3)\bigr)=\frac13.
\]
因此 \(g_p\) 在 \(\FF_p\) 上具有至少一个根（即分解型为 \((1)(2)\) 或 \((1)(1)(1)\)）的素数集合密度为
\[
\delta\bigl(g_p\ \text{在}\ \FF_p\ \text{上有根}\bigr)=\frac12+\frac16=\frac23.
\]
对判别式二次脊曲线
\(
C_p:\ w^{2}=-y(y-1)g_p(y)
\)
而言，这等价于其分歧点集合在 \(\FF_p\) 上除 \(\{y=0,1,\infty\}\) 之外还出现额外的 \(\FF_p\)-Weierstrass 点。
\end{conclusion}
\begin{proof}[证明要点]
由结论 \ref{con:xi-terminal-zm-leyang-cubic-arithmetic}，\(g\) 在 \(\QQ\) 上不可约且其分裂域的伽罗瓦群为 \(S_3\)，并具有唯一二次分辨子域 \(\QQ(\sqrt{-111})\)。
对 \(p\nmid \mathrm{Disc}(g)\) 的素数，Frobenius 元 \(\mathrm{Frob}_p\in S_3\) 的共轭类与 \(g_p\) 的分解型一一对应：
恒等类对应全分裂 \((1)(1)(1)\)，换位类对应 \((1)(2)\)，而三循环类对应不可约 \((3)\)。
另一方面，符号同态 \(\mathrm{sgn}:S_3\to\{\pm1\}\) 的固定域正是上述二次子域；
因而 \(\mathrm{sgn}(\mathrm{Frob}_p)=\left(\frac{-111}{p}\right)=\varepsilon_p\)。
于是 \(\varepsilon_p=-1\) 强制 \(\mathrm{Frob}_p\) 落入换位共轭类，从而 \(g_p\) 必为 \((1)(2)\) 型；而 \(\varepsilon_p=+1\) 等价于 \(\mathrm{Frob}_p\) 为偶置换，从而仅可能为恒等或三循环，给出 \((1)(1)(1)\) 或 \((3)\) 两型。
密度由 Chebotarev 定理给出，为相应共轭类大小除以 \(|S_3|=6\)，即分别为 \(3/6,1/6,2/6\)。
可复现核验脚本：\texttt{scripts/exp\_fold\_zm\_elliptic\_leyang\_cubic\_modp\_splitting\_audit.py}。证毕。
\end{proof}

\begin{conclusion}[判别式的素因子骨架与不可约性：素数 \texorpdfstring{$37$}{37} 在三条互异分歧机制中以一次幂强制出现]\label{con:xi-terminal-37-squarefree-coupling}
设 \(E\) 取短 Weierstrass 模型 \(Y^2=X^3-X+\tfrac14\)，并记
\[
\psi_2(X):=4X^3-4X+1\quad(\text{\(E\) 的 \(2\)-division 多项式}),
\qquad
S(X):=16X^3-9X^2+1\quad(\text{结论 \ref{con:xi-terminal-zm-elliptic-y-hurwitz-passport} 的分歧三次式}),
\]
\[
P_{\mathrm{LY}}(y):=256y^3+411y^2+165y+32\quad(\text{Lee--Yang 三次因子}).
\]
则三者判别式与 \(E\) 的判别式之间满足严格恒等式
\[
\Delta(E)=37,\qquad
\mathrm{Disc}(\psi_2)=2^4\cdot 37,\qquad
\mathrm{Disc}(S)=-2^2\cdot 3^3\cdot 37,\qquad
\mathrm{Disc}(P_{\mathrm{LY}})=-3^9\cdot 31^2\cdot 37.
\]
从而 \(\psi_2\) 的分裂域之唯一二次分辨子域为 \(\QQ(\sqrt{37})\)。
从而素数 \(37\) 以一次幂同时出现在：
\begin{itemize}
  \item 椭圆曲线 \(E\) 的最小坏约化处；
  \item \(y\)-投影分歧点的 \(X\)-坐标三次域（由 \(S\) 控制）的基本判别式因子中；
  \item Lee--Yang 三次因子 \(P_{\mathrm{LY}}(y)\) 的分裂判别式中。
\end{itemize}
特别地，对任意素数 \(p\notin\{2,3,31,37\}\)，上述三多项式在 \(\FF_p\) 上均无重根，且 \(y=0,1\) 与 \(P_{\mathrm{LY}}(y)=0\) 的三点在模 \(p\) 下保持互异；
从而该 \(S_4\)-覆盖塔在此集合之外呈现纯粹驯分歧并具有稳定约化型。
\end{conclusion}
\begin{proof}[证明要点]
\(\Delta(E)\) 由短 Weierstrass 形式的标准公式 \(\Delta(E)=-16(4a^3+27b^2)\)（对 \(Y^2=X^3+aX+b\)）直接计算（此处 \(a=-1,b=\tfrac14\)）得到 \(\Delta(E)=37\)。
其余三个判别式为一元判别式公式的直接代数计算，并在 \(\ZZ\) 中分解。
最后的约化断言来自：若 \(p\nmid \mathrm{Disc}(f)\) 则 \(f\) 模 \(p\) 无重根；
并且当 \(p\notin\{2,3\}\) 时 \(P_{\mathrm{LY}}(0)=32\)、\(P_{\mathrm{LY}}(1)=864\) 皆非零，从而 \(0,1\) 不与 \(P_{\mathrm{LY}}\) 的根碰撞；
再结合 \(p\nmid |S_4|=24\) 得驯分歧与稳定约化结论。证毕。
\end{proof}

\begin{conclusion}[二阶碰撞三次与 \texorpdfstring{$E[2]$}{E[2]} 的三次坐标域：Tschirnhaus 生成元变换与 Möbius 共轭的显式同构]\label{con:xi-terminal-zm-collision-s2-e2-field-isomorphism}
记二阶碰撞幂和 \(S_2(m)=\sum_{x\in X_m}d_m(x)^2\) 的最小递推特征多项式为 \(P_2(\Lambda)\)：
\[
P_2(\Lambda):=\Lambda^3-2\Lambda^2-2\Lambda+2\in\ZZ[\Lambda].
\]
另一方面，在结论 \ref{con:xi-terminal-zm-elliptic-normalization} 的短 Weierstrass 模型
\(
E:\ Y^2=X^3-X+\tfrac14
\)
上，非平凡 \(2\)-挠点的 \(X\)-坐标由 \(2\)-division 三次
\[
D_2(X):=4X^3-4X+1\in\ZZ[X]
\]
刻画。
取 \(\theta\) 为 \(D_2(X)=0\) 的任一根，并令 \(\beta:=2\theta\)。
则 \(\beta\) 为代数整数，满足首一三次
\[
\beta^3-4\beta+2=0,
\qquad
\mathrm{Disc}\bigl(\beta^3-4\beta+2\bigr)=148=4\cdot 37.
\]
定义 Tschirnhaus 变换
\[
\boxed{\;\Lambda:=-2+2\theta+4\theta^2=\beta^2+\beta-2\;}
\]
则 \(P_2(\Lambda)=0\) 在 \(\QQ(\beta)=\QQ(\theta)\) 中成立，从而诱导域同构
\[
\boxed{\ \QQ(\Lambda)\cong \QQ(\beta)=\QQ(\theta)=:\QQ(E[2])\ }.
\]
因此二阶碰撞谱的三次代数数域与重量谱椭圆曲线的 \(2\)-挠坐标域在算术上被严格识别为同一对象；
并且两条首一整式 \(P_2(\Lambda)\) 与 \(\beta^3-4\beta+2\) 的判别式相同，给出该三次数域的域判别式 \(148\)。
等价地，若以非首一三次 \(D_2(X)\) 描述 \(X\)-坐标，则其多项式判别式为 \(\mathrm{Disc}(D_2)=2^{4}\cdot 37\)，与 \(148\) 仅相差平方因子 \(2^{2}\)。
\par
更进一步地，上述二次 Tschirnhaus 变换在同一三次数域内可显式反演：在 \(\QQ(\theta)=\QQ(\Lambda)\) 中严格成立
\[
\boxed{\;\theta=\frac{\Lambda^{2}-\Lambda-2}{2}\; }.
\]
等价地，直接以 \((\theta,\Lambda)\) 表述，则互逆二次变换为
\[
\Lambda=4\theta^{2}+2\theta-2,
\qquad
\theta=\frac{\Lambda^{2}-\Lambda-2}{2}.
\]
并且两幂基 \(\{1,\theta,\theta^{2}\}\) 与 \(\{1,\Lambda,\Lambda^{2}\}\) 的判别式满足严格比例
\[
\mathrm{Disc}(D_2)=4\cdot \mathrm{Disc}(P_2)
\qquad
\bigl(\mathrm{Disc}(D_2)=16\cdot 37,\ \mathrm{Disc}(P_2)=4\cdot 37\bigr),
\]
从而对应基变换行列式的绝对值为 \(2\)。
\par
等价地，同一三次数域亦可由一次分式变换给出：令
\[
\phi(\mu):=\frac{\mu}{2(\mu-1)}.
\]
则严格恒等式成立
\[
\boxed{\;D_2\bigl(\phi(\mu)\bigr)\cdot\bigl(2(\mu-1)\bigr)^3=-4\,P_2(\mu)\;}
\]
从而 \(P_2(\mu)=0\) 当且仅当 \(D_2(\phi(\mu))=0\)；其逆变换为 \(\mu=\frac{2X}{2X-1}\)（其中 \(X=\phi(\mu)\)）。
\end{conclusion}
\begin{proof}[证明要点]
将 \(\Lambda=-2+2X+4X^2\) 代入 \(P_2(\Lambda)\) 并在 \(\QQ[X]\) 中对 \(D_2(X)=4X^3-4X+1\) 作带余除法，可得余式恒为零，从而当 \(D_2(\theta)=0\) 时必有 \(P_2(-2+2\theta+4\theta^2)=0\)。
等价地，结果式给出严格消元证书
\[
\mathrm{Res}_X\!\bigl(D_2(X),\ \Lambda-(-2+2X+4X^2)\bigr)=16\,P_2(\Lambda),
\]
因此 \(\Lambda\) 在 \(\QQ(\theta)\) 中满足同一三次极小多项式，诱导出所述三次域同构。
反解式 \(\theta=(\Lambda^{2}-\Lambda-2)/2\) 为在商环 \(\QQ[X]/(D_2)\) 中对 \(\Lambda=-2+2X+4X^2\) 作代数化简所得之恒等式。
一次分式变换 \(\phi(\mu)=\mu/(2(\mu-1))\) 的恒等式可由直接代入并清分母逐项展开验证。
可复现核验脚本：\texttt{scripts/exp\_fold\_zm\_collision\_s2\_e2\_field\_isomorphism\_audit.py}。证毕。
\end{proof}

\begin{conclusion}[二阶碰撞分裂型与 \texorpdfstring{$E[2]$}{E[2]} 的 Frobenius 作用：模 \texorpdfstring{$p$}{p} 的可检验等价]\label{con:xi-terminal-zm-collision-s2-splitting-e2-torsion}
沿用结论 \ref{con:xi-terminal-zm-collision-s2-e2-field-isomorphism} 的记号，并取素数 \(p\neq 2,37\)。
则以下条件等价：
\begin{enumerate}
  \item 二阶碰撞三次 \(P_2(\Lambda)\) 在 \(\FF_p[\Lambda]\) 上完全分裂；
  \item \(2\)-division 三次 \(D_2(X)=4X^3-4X+1\) 在 \(\FF_p[X]\) 上完全分裂；
  \item 椭圆曲线 \(E/\FF_p\) 的三条非平凡 \(2\)-挠点全在 \(\FF_p\) 上有理，即
  \[
  E(\FF_p)[2]\cong (\ZZ/2\ZZ)^2.
  \]
\end{enumerate}
更进一步，\(P_2\)（等价地 \(D_2\)）在 \(\FF_p\) 上的分解型与 Frobenius 元在
\(
S_3\cong \mathrm{Aut}\bigl(E[2]\setminus\{O\}\bigr)
\)
中的共轭类一致，从而二阶碰撞谱的模素分裂行为可被精确译码为 \(E[2]\) 上的伽罗瓦表示型。
\par
与此相对，结论 \ref{con:xi-terminal-zm-leyang-ramification-critical-values} 的 Lee--Yang 非有理分歧轨道属于另一条 \(S_3\)-型扩张，其判别式二次子域为 \(\QQ(\sqrt{-111})\)（见结论 \ref{con:xi-terminal-zm-leyang-cubic-arithmetic} 与结论 \ref{con:xi-terminal-zm-leyang-cubic-field-isomorphism}）。
由于 \(\QQ(\sqrt{-111})\neq \QQ(\sqrt{37})\)，两条 \(S_3\)-伽罗瓦闭包在 \(\QQ\) 上线性互不相交；其模素分解型的统计独立性由结论 \ref{con:xi-terminal-zm-collision-s2-leyang-linear-disjointness} 给出。
\end{conclusion}
\begin{proof}[证明要点]
由结论 \ref{con:xi-terminal-zm-collision-s2-e2-field-isomorphism} 的 Tschirnhaus 同构，可知 \(P_2\) 与 \(D_2\) 在 \(\QQ\) 上生成同一三次数域；
对 \(p\nmid 2\cdot 37\) 的素数，该数域在 \(p\) 处无分歧，因而 \(p\) 的分解型可由任一整系数生成多项式在 \(\FF_p\) 上的分解型读取，得到 \((1)\Longleftrightarrow(2)\)。
另一方面，在特征不为 \(2\) 的短 Weierstrass 形式下，非平凡 \(2\)-挠点恰对应 \(Y=0\)，其 \(X\)-坐标即为 \(D_2(X)=0\) 的三根；
故 \(D_2\) 在 \(\FF_p\) 上完全分裂当且仅当三条非平凡 \(2\)-挠点全为 \(\FF_p\)-点，即 \((2)\Longleftrightarrow(3)\)。
最后，Frobenius 在三条非平凡 \(2\)-挠点集合上的置换型与三次扩张的分解型之间的对应为标准事实。
可复现核验脚本：\texttt{scripts/exp\_fold\_zm\_collision\_s2\_e2\_field\_isomorphism\_audit.py}。证毕。
\end{proof}

\begin{conclusion}[两条 \texorpdfstring{$S_3$}{S3} 扩张的线性互不相交：\texorpdfstring{$\QQ(\sqrt{37})$}{Q(sqrt(37))} 与 \texorpdfstring{$\QQ(\sqrt{-111})$}{Q(sqrt(-111))} 的二次脊柱分离]\label{con:xi-terminal-zm-collision-s2-leyang-linear-disjointness}
设 \(L_{2}:=\mathrm{Spl}\bigl(P_2(\Lambda)\bigr)\) 为二阶碰撞三次 \(P_2(\Lambda)\) 的分裂域，并设 \(L_{\mathrm{LY}}:=\mathrm{Spl}\bigl(P_{\mathrm{LY}}(y)\bigr)\) 为 Lee--Yang 三次
\[
P_{\mathrm{LY}}(y)=256y^{3}+411y^{2}+165y+32\in\ZZ[y]
\]
的分裂域。
则两者均为 \(S_3\)-伽罗瓦扩张，且其唯一二次子域分别为
\[
L_2^{A_3}=\QQ(\sqrt{37}),\qquad L_{\mathrm{LY}}^{A_3}=\QQ(\sqrt{-111}).
\]
因此 \(L_2\cap L_{\mathrm{LY}}=\QQ\)，两扩张线性互不相交，并有
\[
\mathrm{Gal}(L_2L_{\mathrm{LY}}/\QQ)\cong S_3\times S_3.
\]
对所有素数 \(p\nmid 2\cdot 3\cdot 37\)，Frobenius 在两因子上的共轭类独立等分布；特别地，“\(P_2\) 与 \(P_{\mathrm{LY}}\) 同时在 \(\FF_p\) 上完全分裂”的素数集合自然密度为 \(1/36\)。
\end{conclusion}
\begin{proof}[证明要点]
由结论 \ref{con:xi-terminal-zm-collision-s2-e2-field-isomorphism} 的判别式计算，\(P_2\) 的判别式平方自由核为 \(37\)，从而其分裂域的唯一二次子域为 \(\QQ(\sqrt{37})\)。
另一方面，结论 \ref{con:xi-terminal-zm-leyang-cubic-arithmetic} 给出 Lee--Yang 三次 \(P_{\mathrm{LY}}\) 的判别式平方自由核为 \(-111\)，其分裂域的唯一二次子域为 \(\QQ(\sqrt{-111})\)。
\par
由于 \(L_2\) 与 \(L_{\mathrm{LY}}\) 均为伽罗瓦扩张，交域 \(L_2\cap L_{\mathrm{LY}}\) 亦为伽罗瓦扩张。
而 \(S_3\) 的真伽罗瓦子扩张仅可能为其唯一二次子域；由 \(\QQ(\sqrt{37})\neq\QQ(\sqrt{-111})\) 即得 \(L_2\cap L_{\mathrm{LY}}=\QQ\)。
因此两扩张线性互不相交，复合域的伽罗瓦群为直积 \(S_3\times S_3\)。
\par
对 \(p\nmid 2\cdot 3\cdot 37\)，两域在 \(p\) 处均未分歧；Chebotarev 密度定理在直积群上的等分布给出联合分布的乘积分解。
完全分裂对应恒等共轭类，其在每个 \(S_3\) 中所占比例为 \(1/6\)，故同时完全分裂的自然密度为 \((1/6)^2=1/36\)。证毕。
\end{proof}

\begin{conclusion}[Fibonacci cube 的 \(f\)-多项式分裂域：全实阿贝尔性与导数约束]\label{con:xi-terminal-fibcube-splitting-field-real-abelian-conductor}
对 \(n\ge 0\) 令 \(K_n:=\mathrm{Spl}\bigl(C_n(u)\bigr)\) 为 Fibonacci cube 的 \(f\)-多项式 \(C_n(u)\)（定义 \ref{def:pom-fibcube-fvector}）在 \(\QQ\) 上的分裂域。
则：
\begin{enumerate}
  \item \(K_n\subset \QQ(\zeta_{2(n+2)})^+\)，其中 \(\QQ(\zeta_{2(n+2)})^+\) 为 \(2(n+2)\)-次单位根域的最大实子域；因而 \(K_n/\QQ\) 为阿贝尔且全实的伽罗瓦扩张；
  \item \(K_n\) 的任意二次子域的导数（conductor）整除 \(2(n+2)\)。
\end{enumerate}
\end{conclusion}
\begin{proof}[证明要点]
由结论 \ref{con:xi-terminal-fibcube-fpoly-real-negative-roots-param} 的显式根式，
\[
u_{n,k}=-\frac{5+\tan^2\!\left(\frac{k\pi}{n+2}\right)}{4}.
\]
而
\(
\tan^2\theta=\frac{1-\cos(2\theta)}{1+\cos(2\theta)}
\)
并且
\(
\cos\!\left(\frac{2k\pi}{n+2}\right)\in \QQ(\zeta_{n+2})^+\subset \QQ(\zeta_{2(n+2)})^+
\)。
故每个根 \(u_{n,k}\) 均属于 \(\QQ(\zeta_{2(n+2)})^+\)，从而由其生成的分裂域 \(K_n\) 亦包含于该最大实子域，得 (1)。
\par
(2) 由 (1) 得 \(K_n\) 为导数整除 \(2(n+2)\) 的实阿贝尔扩张的子域；二次子域对应二次 Dirichlet 特征，其导数必整除同一模数，故得到所述约束。证毕。
\end{proof}

\begin{conclusion}[与 Lee--Yang 三次分裂域的线性互素性]\label{con:xi-terminal-fibcube-leyang-linear-disjointness}
令 \(L_{\mathrm{LY}}:=\mathrm{Spl}\bigl(P_{\mathrm{LY}}\bigr)\) 为 Lee--Yang 三次 \(P_{\mathrm{LY}}(y)\) 的分裂域。
则对一切 \(n\ge 0\) 有
\[
K_n\cap L_{\mathrm{LY}}=\QQ,
\qquad\text{从而}\qquad
\mathrm{Gal}(K_nL_{\mathrm{LY}}/\QQ)\cong \mathrm{Gal}(K_n/\QQ)\times S_3.
\]
\end{conclusion}
\begin{proof}[证明要点]
由结论 \ref{con:xi-terminal-fibcube-splitting-field-real-abelian-conductor}，\(K_n/\QQ\) 为全实伽罗瓦扩张。
由于 \(L_{\mathrm{LY}}/\QQ\) 为 \(S_3\)-伽罗瓦扩张，其伽罗瓦子扩张仅可能为 \(\QQ\) 或唯一二次子域 \(\QQ(\sqrt{-111})\)（结论 \ref{con:xi-terminal-zm-leyang-cubic-arithmetic}）。
交域 \(K_n\cap L_{\mathrm{LY}}\) 作为两个伽罗瓦扩张的交必为 \(\QQ\) 上伽罗瓦扩张，故只能为上述两者之一；
但 \(\QQ(\sqrt{-111})\) 为虚二次域，不可能嵌入全实域 \(K_n\)，故交域只能为 \(\QQ\)。
两边皆为伽罗瓦扩张且交域平凡，得到线性互素与伽罗瓦群直积分解。证毕。
\end{proof}

\begin{conclusion}[三通道 Chebotarev 乘积律：二阶碰撞、Lee--Yang 与 Fibonacci cube 的联合完全分裂密度]\label{con:xi-terminal-zm-collision-leyang-fibcube-chebotarev-triple-product}
沿用结论 \ref{con:xi-terminal-zm-collision-s2-leyang-linear-disjointness} 的 \(L_2:=\mathrm{Spl}(P_2)\) 与 \(L_{\mathrm{LY}}:=\mathrm{Spl}(P_{\mathrm{LY}})\)，并令 \(K_n:=\mathrm{Spl}(C_n)\) 如结论 \ref{con:xi-terminal-fibcube-splitting-field-real-abelian-conductor}。
若 \(37\nmid (n+2)\)，则三域两两线性互素：
\[
L_2\cap K_n=\QQ,\qquad
L_{\mathrm{LY}}\cap K_n=\QQ,\qquad
L_2\cap L_{\mathrm{LY}}=\QQ,
\]
从而
\[
\mathrm{Gal}(L_2L_{\mathrm{LY}}K_n/\QQ)\cong S_3\times S_3\times \mathrm{Gal}(K_n/\QQ).
\]
并且对除去有限坏素集合后的素数 \(p\)，满足下列三条件同时成立的素数集合具有自然密度
\[
\boxed{\;\delta_n=\frac{1}{36\,[K_n:\QQ]}\;}
\]
其中三条件为：
\begin{enumerate}
  \item \(P_2(\Lambda)\) 在 \(\FF_p\) 上完全分裂；
  \item \(P_{\mathrm{LY}}(y)\) 在 \(\FF_p\) 上完全分裂；
  \item \(C_n(u)\) 在 \(\FF_p\) 上完全分裂。
\end{enumerate}
\end{conclusion}
\begin{proof}[证明要点]
由结论 \ref{con:xi-terminal-zm-collision-s2-leyang-linear-disjointness} 已有 \(L_2\cap L_{\mathrm{LY}}=\QQ\)。
又由结论 \ref{con:xi-terminal-fibcube-leyang-linear-disjointness} 得 \(K_n\cap L_{\mathrm{LY}}=\QQ\)。
剩余只需证明 \(L_2\cap K_n=\QQ\)。
\par
由于 \(L_2/\QQ\) 为 \(S_3\)-伽罗瓦扩张，其非平凡伽罗瓦子扩张只能为唯一二次子域 \(\QQ(\sqrt{37})\)（结论 \ref{con:xi-terminal-zm-collision-s2-leyang-linear-disjointness}）。
而由结论 \ref{con:xi-terminal-fibcube-splitting-field-real-abelian-conductor}，\(K_n\) 的任意二次子域导数整除 \(2(n+2)\)；
若 \(\QQ(\sqrt{37})\subseteq K_n\)，则其导数 \(37\) 必整除 \(2(n+2)\)，从而 \(37\mid(n+2)\)，与假设矛盾。
故 \(L_2\cap K_n=\QQ\)。
\par
三域皆为 \(\QQ\) 上伽罗瓦扩张且两两交平凡，复合域的伽罗瓦群为直积。
对不分歧素数 \(p\)，“在 \(\FF_p\) 上完全分裂”等价于 Frobenius 落在对应因子群的恒等共轭类；
直积结构给出独立性，Chebotarev 密度定理遂推出联合事件密度为三边际密度之乘积。
其中 \(L_2\) 与 \(L_{\mathrm{LY}}\) 同时完全分裂的密度为 \(1/36\)（结论 \ref{con:xi-terminal-zm-collision-s2-leyang-linear-disjointness}），而 \(K_n\) 中完全分裂素数的密度为 \(1/[K_n:\QQ]\)。
相乘即得 \(\delta_n=1/(36[K_n:\QQ])\)。证毕。
\end{proof}




\begin{theorem}[Lee--Yang 特征曲线的椭圆化与 \(37\)-模性]\label{thm:xi-terminal-zm-leyang-elliptic-modularity-37}\label{thm:xi-terminal-zm-leyang-elliptic-37-modularity-monodromy-langlands}
设
\[
\Pi(\lambda,y):=\lambda^{4}-\lambda^{3}-(2y+1)\lambda^{2}+\lambda+y(y+1)\in\ZZ[\lambda,y],
\]
并令仿射曲线 \(\mathcal C\subset\mathbb A^2_{(\lambda,y)}\) 由 \(\Pi(\lambda,y)=0\) 给出。
令
\[
x:=\lambda,\qquad v:=y-x^{2}.
\]
其中 \(v\) 与前文记号 \(u:=y-\lambda^2\) 一致。
则 \(\mathcal C\) 与椭圆曲线
\[
E_0:\quad v^{2}+v=x^{3}-x
\]
在 \(\QQ\) 上同构；可取显式互逆同构
\[
\phi:E_0\to \mathcal C,\quad (x,v)\longmapsto(\lambda,y)=(x,v+x^{2}),
\qquad
\phi^{-1}:\mathcal C\to E_0,\quad (\lambda,y)\longmapsto (x,v)=(\lambda,y-\lambda^{2}).
\]
并且 \(\Delta(E_0)=37\)、\(j(E_0)=110592/37\)；在 Cremona/LMFDB 记号下 \(E_0\) 对应强 Weil 曲线 \(37\mathrm a1\)，并可作为
\(
X_0^+(37):=X_0(37)/\langle w_{37}\rangle
\)
的代数模型之一。\cite{LMFDBEllipticCurve37a1}
\end{theorem}

\leanverified{paper\_xi\_terminal\_zm\_leyang\_elliptic\_37\_modularity\_monodromy\_langlands}

\begin{proof}
将 \(y=v+x^2\) 与 \(\lambda=x\) 代入得到恒等式
\[
\Pi(x,v+x^2)=v^2+v-x^3+x,
\]
从而 \(\Pi(\lambda,y)=0\) 与 \(v^2+v=x^3-x\) 等价；并且 \((x,v)\leftrightarrow(\lambda,y)\) 的互逆性由 \(y=v+x^2\) 在仿射坐标环中的可逆性直接给出，故 \(\phi\) 为 \(\QQ\)-同构。
不变量计算与外部标识见注记 \ref{rem:xi-terminal-elliptic-37a1-modular-anchor}。证毕。
\end{proof}
\par

\begin{proposition}[重量参数作为椭圆曲线上的四重函数：分歧点、分歧值与群律约束]\label{prop:xi-terminal-zm-leyang-f-ramification}
令 \(E_0\) 为定理 \ref{thm:xi-terminal-zm-leyang-elliptic-modularity-37} 的椭圆曲线，\(O\) 为无穷远点。
设
\[
f:=v+x^{2}=\phi^{*}(y)\in\QQ(E_0),
\]
并以同一符号记由 \(f\) 诱导的有理映射 \(f:E_0\to\PP^{1}\)。
则：
\begin{enumerate}
  \item \(\QQ(E_0)/\QQ(f)\) 为次数 \(4\) 的函数域扩张，并且 \(x\) 在 \(\QQ(f)\) 上的极小多项式为
  \[
  \Pi(x,f)=x^{4}-x^{3}-(2f+1)x^{2}+x+f(f+1)\in\QQ(f)[x].
  \]
  \item \(f\) 在 \(O\) 处具有唯一四阶极点，且 \((f)_\infty=4O\)。因此 \(\deg(f)=4\)，并在 \(O\) 上方全分歧（分歧指数 \(e_O=4\)）。
  \item \(f\) 的有限分歧点恰为 \(\dd f=0\) 的点；其 \(x\)-坐标满足
  \[
  \dd f=0\quad\Longleftrightarrow\quad (x^{2}-1)\bigl(16x^{3}-9x^{2}+1\bigr)=0.
  \]
  更具体地，有限分歧点由五点组成：两条有理点
  \[
  P_0=(1,-1),\qquad P_1=(-1,0),
  \]
  以及三点 \(P_2,P_3,P_4\)，其中 \(x(P_i)\) 为三次方程 \(16x^{3}-9x^{2}+1=0\) 的三根，且
  \[
  v(P_i)=\frac{1-3x(P_i)^2-2x(P_i)}{4x(P_i)}\qquad(i=2,3,4).
  \]
  \item 令 \(B\subset\PP^1\) 为覆盖 \(f:E_0\to\PP^1\) 的分歧值集合，则
  \[
  B=\{\infty,0,1\}\ \cup\ \{t\in\overline{\QQ}:\ P_{\mathrm{LY}}(t)=0\},
  \qquad
  P_{\mathrm{LY}}(t):=256t^3+411t^2+165t+32.
  \]
  并且 \(B\) 的每个有限分歧值之上恰出现一处简单分歧（分歧指数 \(2\)），其余点不分歧。
  \item 作为 \(E_0\) 上的亚纯微分，\(\dd f\) 的零极点除子满足
  \[
  \mathrm{div}(\dd f)=P_0+P_1+P_2+P_3+P_4-5O,
  \]
  从而在 \(E_0(\overline{\QQ})\) 的群律下有
  \[
  P_0+P_1+P_2+P_3+P_4=O,
  \qquad
  P_2+P_3+P_4=-(P_0+P_1)=\Bigl(\frac14,-\frac58\Bigr)\in E_0(\QQ).
  \]
\end{enumerate}
\end{proposition}
\begin{proof}
由 \(f=v+x^2\) 得 \(v=f-x^2\)。代回 \(E_0:v^2+v=x^3-x\) 并展开，得到
\[
(f-x^2)^2+(f-x^2)=x^3-x
\quad\Longleftrightarrow\quad
x^{4}-x^{3}-(2f+1)x^{2}+x+f(f+1)=0,
\]
即 \(\Pi(x,f)=0\)。因此 \([\QQ(E_0):\QQ(f)]\le 4\)。
\par
把 \(E_0\) 配方为短 Weierstrass 形式 \(Y^2=X^3-X+\tfrac14\)（取 \(X=x\)、\(Y=v+\tfrac12\)），则 \(x\) 与 \(v\) 在 \(O\) 处的极点阶分别为 \(2\) 与 \(3\)，故 \(f=v+x^2\) 在 \(O\) 处极点阶为 \(4\)，从而 \((f)_\infty=4O\) 且 \(\deg(f)=4\)。由次数等于函数域扩张次数可得 \([\QQ(E_0):\QQ(f)]=4\)，并且 \(\Pi(x,f)\) 为 \(x\) 的极小多项式。
\par
对 \(E_0:v^2+v=x^3-x\) 微分得到
\[
(2v+1)\,\dd v=(3x^2-1)\,\dd x.
\]
于是
\[
\dd f=\dd(v+x^2)=\dd v+2x\,\dd x=\left(\frac{3x^2-1}{2v+1}+2x\right)\dd x.
\]
因此 \(\dd f=0\) 等价于 \(3x^2-1+2x(2v+1)=0\)，亦即
\[
4xv+3x^2+2x-1=0.
\]
在 \(x\neq 0\) 上可解得 \(v=(1-3x^2-2x)/(4x)\)；
将其代回 \(v^2+v=x^3-x\) 并化简即得 \((x^2-1)(16x^3-9x^2+1)=0\)，从而得到五个有限分歧点的显式刻画。
\par
记 \(E_0\) 的标准正则微分
\[
\omega:=\frac{\dd x}{2v+1}.
\]
由上式可得 \(\dd f=(4xv+3x^2+2x-1)\,\omega\)。
由于 \(g(E_0)=1\)，\(\omega\) 在 \(E_0\) 上处处正则且无零极点；因此 \(\dd f\) 在五个有限分歧点处各有一个简单零点。
另一方面，由 \((f)_\infty=4O\) 可取局部参数 \(t\) 使 \(f\sim t^{-4}\)，则 \(\dd f\sim -4t^{-5}\dd t\)，故 \(\dd f\) 在 \(O\) 处为五阶极点，从而 \(\mathrm{div}(\dd f)=\sum_{i=0}^{4}P_i-5O\)。
该除子为典范类，而椭圆曲线的典范类在 \(\mathrm{Pic}^0(E_0)\) 中为零；在同一识别 \(E_0\simeq \mathrm{Pic}^0(E_0)\) 下得到群律恒等式 \(\sum_{i=0}^4 P_i=O\)。
最后，在短 Weierstrass 坐标 \((X,Y)\) 中，
\[
P_0=(1,-\tfrac12),\qquad P_1=(-1,\tfrac12),
\]
故弦割斜率为 \(-\tfrac12\)，从而 \(P_0+P_1=(\tfrac14,\tfrac18)\)，变回 \(v=Y-\tfrac12\) 得
\(
P_0+P_1=(\tfrac14,-\tfrac38)
\)；
取负元 \(v\mapsto -v-1\) 即得
\(
-(P_0+P_1)=(\tfrac14,-\tfrac58)
\)，从而 \(P_2+P_3+P_4=-(P_0+P_1)\)。
\par
最后，由 \(\Pi(x,f)=0\) 作为 \(\QQ(f)\) 上的四次方程，分歧值由判别式为零刻画；结合结论 \ref{con:xi-terminal-zm-discriminant-leyang} 的分解
\(
\mathrm{Disc}_{x}\bigl(\Pi(x,f)\bigr)=-f(f-1)P_{\mathrm{LY}}(f)
\)
得到分歧值集合 \(B=\{\infty,0,1\}\cup\{P_{\mathrm{LY}}(t)=0\}\)。
证毕。
\end{proof}
\par

\begin{theorem}[分歧型 \texorpdfstring{$(4;2^{5})$}{(4;2^5)} 覆盖的全对称单值化：\texorpdfstring{$\mathrm{Mon}(f)\cong S_4$}{Mon(f)=S4}]\label{thm:xi-terminal-zm-leyang-monodromy-s4}
\leanverified{paper\_xi\_terminal\_zm\_leyang\_monodromy\_s4}
沿用命题 \ref{prop:xi-terminal-zm-leyang-f-ramification} 的 \(f:E_0\to\PP^{1}\)。
则 \(f\) 的几何单值群同构于 \(S_4\)。
等价地，将 \(\Pi(\lambda,y)\) 视为 \(\QQ(y)\) 上关于 \(\lambda\) 的四次多项式，则其函数域伽罗瓦群为 \(S_4\)。
\end{theorem}
\begin{proof}
由命题 \ref{prop:xi-terminal-zm-leyang-f-ramification}，覆盖 \(f:E_0\to\PP^1\) 在 \(O\) 上方全分歧，故绕 \(f=\infty\) 的局部单值为一个 \(4\)-循环；而任一有限分歧值处均为简单分歧，其局部单值为换位。
一个 \(4\)-循环与一个换位生成 \(S_4\)，故 \(\mathrm{Mon}(f)\cong S_4\)。
四次多项式的函数域伽罗瓦群与几何单值群之等价表述见结论 \ref{con:xi-terminal-zm-generic-galois-s4}。证毕。
\end{proof}
\par

\begin{proposition}[素数 \(37\) 的双重显影：\(2\)-除子域与 Lee--Yang 分歧域的统一分歧控制]\label{prop:xi-terminal-37-double-ramification-fingerprint}
沿用定理 \ref{thm:xi-terminal-zm-leyang-elliptic-modularity-37} 与命题 \ref{prop:xi-terminal-zm-leyang-f-ramification} 的记号，并记 \(E_0[2]\) 为 \(E_0\) 的 \(2\)-挠子群。
则存在如下算术分歧指纹的一致性：
\begin{enumerate}
  \item 令 \(Z:=2v+1\)，则 \(E_0\) 等价配方为
  \[
  E_0:\quad Z^{2}=4x^{3}-4x+1.
  \]
  其 \(2\)-除法三次 \(4x^{3}-4x+1\) 的判别式为 \(\mathrm{Disc}(4x^{3}-4x+1)=2^{4}\cdot 37\)，从而 \(E_0[2]\) 的分裂域仅在素数 \(2\) 与 \(37\) 处可能分歧。
  \item Lee--Yang 三次
  \[
  P_{\mathrm{LY}}(t)=256t^{3}+411t^{2}+165t+32
  \]
  的判别式满足 \(\mathrm{Disc}(P_{\mathrm{LY}})=-3^{9}\cdot 31^{2}\cdot 37\)，从而其分裂域仅在素数 \(3,31,37\) 处可能分歧，且 \(37\) 以一次幂出现。
  \item 因而同一素数 \(37\) 同时作为 \(E_0\) 的极小判别式（亦即唯一坏约化素数）与 Lee--Yang 分歧域的必分歧素数出现；从分歧数据的观点看，Lee--Yang 分歧集合与椭圆曲线的坏素数在 \(37\) 处发生不可约耦合。
\end{enumerate}
\end{proposition}
\begin{proof}
对三次多项式 \(aX^{3}+bX^{2}+cX+d\) 的判别式取标准公式
\[
\mathrm{Disc}(aX^{3}+bX^{2}+cX+d)=b^{2}c^{2}-4ac^{3}-4b^{3}d-27a^{2}d^{2}+18abcd.
\]
对 \(4x^{3}-4x+1\) 代入 \((a,b,c,d)=(4,0,-4,1)\) 得 \(\mathrm{Disc}=1024-432=592=2^{4}\cdot 37\)，从而 \(E_0[2]\) 的分裂域仅在 \(\{2,37\}\) 处可能分歧。
第二条的判别式分解见结论 \ref{con:xi-terminal-zm-leyang-cubic-arithmetic}。
最后一条由“分裂域仅在判别式素因子处可能分歧”的一般事实立即推出。证毕。
\end{proof}
\par

\begin{proposition}[唯一坏素数与 \texorpdfstring{$\mathrm{GL}_2$}{GL2} 局部数据的锁定]\label{prop:xi-prime-bad-set-determines-unramified-data}
设 \(E/\QQ\) 为椭圆曲线，且其极小判别式满足 \(|\Delta(E)|=p\) 为素数。
则 \(E\) 的坏约化素数集合包含于 \(\{p\}\)，且 \(E\) 在 \(p\) 处为乘法约化，从而导子 \(N(E)=p\)。
若进一步满足 \(\dim S_2(\Gamma_0(p))=1\)，则存在唯一规范化新形式 \(g\in S_2(\Gamma_0(p))\) 使
\[
L(E,s)=L(g,s),
\]
并且对每个良素 \(\ell\neq p\)，与 \(E\) 对应的 \(\ell\)-进 Galois 表示在所有非分歧处的 Frobenius 特征多项式（等价地 Hecke 本征值序列 \((a_\ell)_{\ell\neq p}\)）由该单一算术输入 \(p\) 唯一确定。
特别地，\(p=37\) 时可取 \(E=E_0: v^2+v=x^3-x\)，其满足 \(\Delta(E_0)=37\) 且 \(\dim S_2(\Gamma_0(37))=1\)，故其全部非分歧局部数据由坏素数集合 \(\{37\}\) 锁定。
\end{proposition}
\begin{proof}
\(|\Delta(E)|=p\) 蕴含极小判别式仅在 \(p\) 处有正赋值，故坏约化素数集合包含于 \(\{p\}\)。
并且 \(\ord_p(\Delta(E))=1\) 强制 \(E\) 在 \(p\) 处为乘法约化，因而导子指数为 \(1\)，得到 \(N(E)=p\)。
由模性定理存在 \(g\in S_2(\Gamma_0(p))\) 满足 \(L(E,s)=L(g,s)\)。
当 \(\dim S_2(\Gamma_0(p))=1\) 时，规范化条件 \(a_1=1\) 迫使 \(g\) 唯一，从而其 Hecke 本征值序列唯一；
这等价于所有 \(\ell\neq p\) 处非分歧局部因子（Frobenius 特征多项式）被唯一锁定。
对 \(p=37\) 的锚定由定理 \ref{thm:xi-terminal-zm-leyang-elliptic-modularity-37} 与推论 \ref{cor:xi-terminal-zm-leyang-langlands-level37} 的外部核验接口给出。证毕。
\end{proof}

\begin{conjecture}[判别式的 gcd--导子原则：Lee--Yang 分歧素数与坏约化素数的可计算接口]\label{conj:xi-leyang-discriminant-gcd-conductor-principle}
设由某一 Lee--Yang 特征谱曲线的椭圆化得到椭圆曲线 \(E/\QQ\)，并设模型内禀地给出一整数
\(\mathfrak D_{\mathrm{LY}}\in\ZZ\)
（可取为控制临界值/分歧的判别式、结式或其平方自由部分的规范化代表）。
记 \(\mathrm{sqf}(\cdot)\) 为平方自由部分算子。
则预期存在同一模型类中不依赖坐标选择的对应，使得
\[
\mathrm{sqf}\!\Bigl(\gcd\bigl(\Delta(E),\mathfrak D_{\mathrm{LY}}\bigr)\Bigr)
=
\mathrm{sqf}\!\bigl(N(E)\bigr),
\]
其中 \(\Delta(E)\) 为极小判别式、\(N(E)\) 为导子。
在 \(\Delta(E)\) 平方自由的情形，该关系等价于：\(\gcd(\Delta(E),\mathfrak D_{\mathrm{LY}})\) 的素因子集合与坏约化素数集合一致。
\end{conjecture}

\begin{corollary}[Langlands--\texorpdfstring{$\mathrm{GL}_2$}{GL2} 的可审计实现：导子 \(37\) 的模性对象与八面体有限参数]\label{cor:xi-terminal-zm-leyang-langlands-level37}
令 \(\mathcal C\) 为定理 \ref{thm:xi-terminal-zm-leyang-elliptic-modularity-37} 的特征曲线（亦即 \(\Pi(\lambda,y)=0\) 的正规化），并令 \(E_0\) 为相应椭圆曲线。
并沿用命题 \ref{prop:xi-terminal-zm-leyang-f-ramification} 的四重覆盖 \(f:E_0\to\PP^1\) 及其分歧值集合 \(B\subset\PP^1\)，令 \(U:=\PP^1\setminus B\)。
则：
\begin{enumerate}
  \item 存在唯一的权 \(2\)、level \(37\) 的规范化新形式 \(g\in S_2(\Gamma_0(37))\) 及对应自守表示 \(\pi_{37}\)，并相应存在与 \(E_0\) 相容的 \(\ell\)-进 Galois 表示族 \(\{\rho_{E_0,\ell}\}_{\ell}\)，使得
  \[
  L(\mathcal C,s)=L(E_0,s)=L(g,s)=L(\pi_{37},s).
  \]
  并且对任意良素 \(p\neq 37\) 有点计数恒等式
  \[
  \#E_0(\FF_p)=p+1-a_p,
  \]
  其中 \(a_p\) 为 \(g\) 的 Hecke 本征值（亦即 \(E_0\) 的 Frobenius 迹）。
  \item 推前得到的秩 \(4\) 局部系统
  \[
  \mathcal L:=\left.(f_\ast\underline{\QQ})\right|_{U}
  \]
  具有几何单值像 \(\mathrm{Mon}(\mathcal L)\cong S_4\)，并经标准同构 \(S_4\simeq \mathrm{PGL}_2(\FF_3)\) 可识别为八面体有限像投影局部系统。\cite{MITSmallGroupsPDF}
  此外，命题 \ref{prop:xi-terminal-37-double-ramification-fingerprint} 表明控制 \(E_0[2]\) 与 \(B\setminus\{\infty,0,1\}\) 的两条三次分裂域仅在 \(\{2,3,31,37\}\) 处可能分歧。
\end{enumerate}
由此，谱四次 \(\Pi(\lambda,y)\) 的正规化与分歧单值数据在算术层面内生锁定了一个 level \(37\) 的 \(\mathrm{GL}_2(\mathbb A_\QQ)\) 模性对象（由 \(L\)-函数与 \(\{a_p\}\) 可计算核验）以及一个八面体有限像投影参数（由 \(\mathrm{Mon}(f)\cong S_4\simeq \mathrm{PGL}_2(\FF_3)\) 读出）。
\end{corollary}
\begin{proof}
第一条由定理 \ref{thm:xi-terminal-zm-leyang-elliptic-modularity-37} 的同构 \(\mathcal C\simeq E_0\) 与椭圆曲线模性定理推出；对 level \(37\) 的外部命名与交叉核验接口见注记 \ref{rem:xi-terminal-elliptic-37a1-modular-anchor}。\cite{LMFDBEllipticCurve37a1,LMFDBNewform372aa,ConradTSWFinal}
该新形式的唯一性可由 \(\dim S_2(\Gamma_0(37))=1\) 的标准维数计算或数据库条目直接读出。\cite{LMFDBNewform372aa}
\(\#E_0(\FF_p)=p+1-a_p\)（\(p\neq 37\)）为 Eichler--Shimura 对应与模性表述下的标准点计数公式。
\par
第二条中 \(\mathrm{Mon}(\mathcal L)\cong \mathrm{Mon}(f)\cong S_4\) 由定理 \ref{thm:xi-terminal-zm-leyang-monodromy-s4} 直接给出；八面体识别由 \(S_4\simeq \mathrm{PGL}_2(\FF_3)\) 的标准同构得到。分歧素数集合由命题 \ref{prop:xi-terminal-37-double-ramification-fingerprint} 立即推出。
证毕。
\end{proof}
\par

\subsubsection{Lee--Yang 零因子张量函子与未分歧参数}\label{subsubsec:xi-leyang-zero-divisor-tensor-functor-unramified}
\noindent
在未分歧局部 Langlands 参数的常用约定下，半单共轭类 \([g]_{\mathrm{ss}}\) 可视为参数的代数群部分。
下述构造把该共轭类在各有限维表示中的特征多项式零点组织为一个对称张量函子，从而得到可传输且可逆读出的 Lee--Yang 零因子数据。
\par

\begin{definition}[Lee--Yang 零因子与乘法卷积半群]\label{def:xi-leyang-zero-divisor-convolution}
记 \(\mathrm{Div}^{+}(\CC^{\times})\) 为 \(\CC^{\times}\) 上有效除子所成的交换幺半群：
\[
\mathrm{Div}^{+}(\CC^{\times})
=\left\{\sum_{i=1}^{n}m_i[z_i]:\ n<\infty,\ m_i\in\ZZ_{\ge 1},\ z_i\in\CC^{\times}\right\},
\]
其加法为除子相加。
定义乘法卷积
\begin{equation}\label{eq:xi-divisor-multiplicative-convolution}
\boxed{\;
\boxtimes:\ \mathrm{Div}^{+}(\CC^{\times})\times \mathrm{Div}^{+}(\CC^{\times})\to \mathrm{Div}^{+}(\CC^{\times}),
\qquad
\Bigl(\sum_i m_i[z_i]\Bigr)\boxtimes\Bigl(\sum_j n_j[w_j]\Bigr)
:=\sum_{i,j}m_in_j[z_iw_j]\; }.
\end{equation}
等价地，\(\boxtimes\) 为乘法映射 \(m:\CC^{\times}\times\CC^{\times}\to\CC^{\times}\)（\(m(z,w)=zw\)）下的推前：\(D_1\boxtimes D_2=m_{\ast}(D_1\times D_2)\)。
\end{definition}

\begin{definition}[半单参数的 Lee--Yang 零因子张量函子]\label{def:xi-leyang-zero-divisor-functor}
设 \({}^{L}G\) 为复连通约化群，并记 \(\mathrm{Rep}({}^{L}G)\) 为其有限维复表示所成的对称张量范畴。
取半单元 \(g\in {}^{L}G(\CC)\) 与表示 \(V\in\mathrm{Rep}({}^{L}G)\)，定义一元多项式
\begin{equation}\label{eq:xi-leyang-characteristic-polynomial}
\boxed{\;
P_g(V;z):=\det\!\bigl(I-z\cdot V(g)\bigr)\in\CC[z]\; }.
\end{equation}
由于 \(P_g(V;0)=1\)，其全部零点位于 \(\CC^{\times}\)。
记 \(\operatorname{div}_0(P_g(V;z))\in\mathrm{Div}^{+}(\CC^{\times})\) 为该多项式在 \(\CC^{\times}\) 上的零点有效除子（按重数计），并定义
\begin{equation}\label{eq:xi-leyang-zero-divisor}
\boxed{\;
\mathrm{LY}_g(V):=\operatorname{div}_0\!\bigl(P_g(V;z)\bigr)\in\mathrm{Div}^{+}(\CC^{\times})\; }.
\end{equation}
\end{definition}

\begin{theorem}[Lee--Yang 零因子张量函子的单射性与有限生成读出]\label{thm:xi-leyang-zero-divisor-tensor-functor-faithful}
\leanverified{paper\_xi\_leyang\_zero\_divisor\_tensor\_functor\_faithful}
设 \({}^{L}G\) 为复连通约化群，\(g\in{}^{L}G(\CC)\) 为半单元。
则由 \eqref{eq:xi-leyang-zero-divisor} 定义的赋值 \(V\mapsto \mathrm{LY}_g(V)\) 满足：
\begin{enumerate}
  \item \textbf{（对称幺半张量函子）}对任意 \(V,W\in\mathrm{Rep}({}^{L}G)\)，有
  \begin{equation}\label{eq:xi-leyang-monoidal-add}
  \boxed{\;
  \mathrm{LY}_g(V\oplus W)=\mathrm{LY}_g(V)+\mathrm{LY}_g(W)\;}
  \end{equation}
  以及
  \begin{equation}\label{eq:xi-leyang-monoidal-tensor}
  \boxed{\;
  \mathrm{LY}_g(V\otimes W)=\mathrm{LY}_g(V)\boxtimes \mathrm{LY}_g(W)\;}
  \end{equation}
  其中 \(\boxtimes\) 由定义 \ref{def:xi-leyang-zero-divisor-convolution} 给出。
  \item \textbf{（半单共轭类的完备性）}若 \(g,g'\in{}^{L}G(\CC)\) 为半单且对一切 \(V\in\mathrm{Rep}({}^{L}G)\) 有 \(\mathrm{LY}_g(V)=\mathrm{LY}_{g'}(V)\)，则 \(g\) 与 \(g'\) 在 \({}^{L}G(\CC)\) 中半单共轭。
  \item \textbf{（有限生成读出）}取极大环面 \(T\subset {}^{L}G\) 与 Weyl 群 \(W\)。
  若存在有限表示族 \(\{V_i\}_{i=1}^{M}\subset\mathrm{Rep}({}^{L}G)\) 使其角色在 \(T/W\) 上生成坐标环 \(\CC[T]^W\)，则由有限数据 \(\{\mathrm{LY}_g(V_i)\}_{i=1}^{M}\) 即可确定 \(g\) 的半单共轭类。
  当 \({}^{L}G\) 半单且单连通时，可取对应基本权的基本表示为此生成系。
\end{enumerate}
\end{theorem}
\begin{proof}
\textbf{(1)} 由特征多项式的乘法性，
\(
P_g(V\oplus W;z)=P_g(V;z)\,P_g(W;z)
\)，
零点除子相加即得 \eqref{eq:xi-leyang-monoidal-add}。
另一方面，\(V(g)\otimes W(g)\) 的特征值为 \(V(g)\) 与 \(W(g)\) 特征值的两两乘积；
故 \(P_g(V\otimes W;z)\) 的零点集合（按重数）为两集合在 \(\CC^{\times}\) 上的乘法卷积推前，从而得到 \eqref{eq:xi-leyang-monoidal-tensor}。
\par
\textbf{(2)} 由 \(\mathrm{LY}_g(V)=\mathrm{LY}_{g'}(V)\) 得对每个 \(V\) 有 \(P_g(V;z)=P_{g'}(V;z)\)，从而 \(\chi_V(g)=\tr(V(g))=\chi_V(g')\)。
由于半单共轭类与 \(T/W\) 的点对应，而有限维表示的角色生成并分离 \(\CC[T]^W\) 上的正则函数，遂 \(g,g'\) 对应同一点，故半单共轭。
\par
\textbf{(3)} 在假设下，由 \(\mathrm{LY}_g(V_i)\) 可恢复 \(P_g(V_i;z)\)（例如由零点及重数确定多项式），从而恢复 \(\chi_{V_i}(g)\)。
于是 \(g\) 在所有生成元上的角色值确定其在 \(\CC[T]^W\) 上的取值，进而确定其半单共轭类。
最后一条为基本表示生成 \(\CC[T]^W\) 的标准事实。证毕。
\end{proof}

\begin{theorem}[有限个 Lee--Yang 圆周条件刻画紧实半单类]\label{thm:xi-leyang-compactness-by-fundamental-circle}
\leanverified{paper\_xi\_leyang\_compactness\_by\_fundamental\_circle}
设 \({}^{L}G\) 半单且单连通，\(K\subset {}^{L}G(\CC)\) 为极大紧子群。
令 \(\{\omega_1,\dots,\omega_r\}\) 为基本权，\(V_{\omega_i}\) 为相应基本表示。
对半单元 \(g\in {}^{L}G(\CC)\)，下列命题等价：
\begin{enumerate}
  \item \(g\) 共轭于 \(K\) 中某元；
  \item 对每个 \(i=1,\dots,r\)，有效除子 \(\mathrm{LY}_g(V_{\omega_i})\) 的支撑包含于单位圆 \(\{z\in\CC^{\times}:\ |z|=1\}\)。
\end{enumerate}
\end{theorem}
\begin{proof}
(i)\(\Rightarrow\)(ii) 若 \(g\) 共轭于 \(K\)，则 \(V_{\omega_i}(g)\) 可取为酉算子，故其全部特征值模长为 \(1\)；
于是 \(P_g(V_{\omega_i};z)\) 的零点（特征值的倒数）亦全在单位圆上。
\par
(ii)\(\Rightarrow\)(i) 取极大环面 \(T\subset {}^{L}G\) 使 \(g\) 共轭于某 \(t\in T(\CC)\)。
将 \(t\) 作极分解 \(t=\exp(\mu)\cdot u\)（\(u\in T\cap K\)），并以 Weyl 共轭令 \(\mu\) 落入闭支配 Weyl 室。
对基本表示 \(V_{\omega_i}\)，其谱半径满足
\(
\rho\bigl(V_{\omega_i}(\exp(\mu))\bigr)=\exp\langle\omega_i,\mu\rangle
\)
（\(\langle\cdot,\cdot\rangle\) 为权与 \(\mathfrak t_{\RR}\) 的自然配对）。
由 (ii) 得 \(V_{\omega_i}(t)\) 的特征值模长全为 \(1\)，故谱半径为 \(1\)，从而 \(\langle\omega_i,\mu\rangle=0\) 对所有 \(i\) 成立。
基本权在 \(\mathfrak t_{\RR}^{\ast}\) 中成基，遂 \(\mu=0\)，即 \(t\in T\cap K\)。
因此 \(g\) 共轭于 \(K\) 中元素。证毕。
\end{proof}

\input{para__xi-leyang-solenoidal-lift-tensor-functor}

\begin{proposition}[函子性下 Lee--Yang 零因子数据的传输与缺陷上界]\label{prop:xi-leyang-functoriality-defect-bound}
\leanverified{paper\_xi\_leyang\_functoriality\_defect\_bound}
令 \(\phi:{}^{L}H\to {}^{L}G\) 为复约化群同态，并取半单元 \(h\in {}^{L}H(\CC)\)。
则对任意 \(V\in\mathrm{Rep}({}^{L}G)\) 有严格等式
\begin{equation}\label{eq:xi-leyang-functoriality-restriction}
\boxed{\;
\mathrm{LY}^{G}_{\phi(h)}(V)
=\mathrm{LY}^{H}_{h}\!\left(\mathrm{Res}^{\,{}^{L}G}_{{}^{L}H}V\right)\; }.
\end{equation}
进一步，设 \({}^{L}H\) 半单且单连通，\(\{\omega_i^{H}\}_{i=1}^{r_H}\) 为其基本权。
定义 Lee--Yang 缺陷量
\begin{equation}\label{eq:xi-leyang-defect}
\boxed{\;
\mathrm{Def}_{\mathrm{LY}}(h)
:=\max_{1\le i\le r_H}\log \rho\!\bigl(V_{\omega_i^{H}}(h)\bigr)\ \in\ [0,\infty)\; }.
\end{equation}
则 \(\mathrm{Def}_{\mathrm{LY}}(h)=0\) 当且仅当 \(h\) 共轭于 \({}^{L}H(\CC)\) 的极大紧子群。
并且对任意固定 \(V\in\mathrm{Rep}({}^{L}G)\)，存在常数 \(C_{V,\phi}\in\NN\) 使对所有半单 \(h\) 成立
\begin{equation}\label{eq:xi-leyang-defect-upperbound}
\boxed{\;
\log \rho\!\bigl(V(\phi(h))\bigr)\ \le\ C_{V,\phi}\cdot \mathrm{Def}_{\mathrm{LY}}(h)\; }.
\end{equation}
\end{proposition}
\begin{proof}
等式 \eqref{eq:xi-leyang-functoriality-restriction} 由表示限制的自然性给出：\(V(\phi(h))\) 与 \(\bigl(\mathrm{Res}V\bigr)(h)\) 为同一线性变换，故其特征多项式与零点除子一致。
\par
\(\mathrm{Def}_{\mathrm{LY}}(h)=0\) 与紧实型等价性由定理 \ref{thm:xi-leyang-compactness-by-fundamental-circle}（应用于 \({}^{L}H\)）立得。
\par
为证 \eqref{eq:xi-leyang-defect-upperbound}，将 \(h\) 共轭到极大环面并取极分解 \(h\sim \exp(\mu)\cdot u\)，令 \(\mu\) 支配。
任取 \({}^{L}H\) 的支配权 \(\lambda=\sum_i a_i\omega_i^{H}\)（\(a_i\in\ZZ_{\ge 0}\)），则
\[
\langle\lambda,\mu\rangle=\sum_i a_i\langle\omega_i^{H},\mu\rangle
\le \Bigl(\sum_i a_i\Bigr)\cdot \mathrm{Def}_{\mathrm{LY}}(h).
\]
将 \(\mathrm{Res}^{\,{}^{L}G}_{{}^{L}H}V\) 分解为不可约直和 \(\bigoplus_j m_j W_{\lambda_j}\)，其中 \(W_{\lambda_j}\) 最高权为 \(\lambda_j=\sum_i a_{j,i}\omega_i^{H}\)。
则
\[
\rho\!\bigl(V(\phi(h))\bigr)=\rho\!\bigl((\mathrm{Res}V)(h)\bigr)
=\max_j \rho\!\bigl(W_{\lambda_j}(h)\bigr)
=\exp\!\Bigl(\max_j\langle\lambda_j,\mu\rangle\Bigr).
\]
令 \(C_{V,\phi}:=\max_j\sum_i a_{j,i}\)，即得 \(\max_j\langle\lambda_j,\mu\rangle\le C_{V,\phi}\mathrm{Def}_{\mathrm{LY}}(h)\)，从而得到 \eqref{eq:xi-leyang-defect-upperbound}。证毕。
\end{proof}

\leanverified{paper\_xi\_leyang\_fock\_trace\_euler\_factor}
\begin{proposition}[玻色 Fock 迹公式与未分歧 Euler 因子的实现]\label{prop:xi-leyang-fock-trace-euler-factor}
令 \(V\) 为有限维复向量空间，并设 \(\mathcal F(V):=\mathrm{Sym}^{\bullet}(V)=\bigoplus_{n\ge 0}\mathrm{Sym}^{n}(V)\)。
对满足 \(\rho(A)<1\) 的线性算子 \(A\in\mathrm{End}(V)\)，令 \(\Gamma(A)\) 表示其在 \(\mathcal F(V)\) 上的二次量子化作用：\(\Gamma(A)\) 在 \(\mathrm{Sym}^{n}(V)\) 上的限制为对称张量幂 \(A^{\otimes n}\) 的诱导算子。
则有严格恒等式
\begin{equation}\label{eq:xi-leyang-fock-trace-determinant}
\boxed{\;
\mathrm{Tr}_{\mathcal F(V)}\bigl(\Gamma(A)\bigr)=\det(I-A)^{-1}\; }.
\end{equation}
从而对任意半单 \(g\in{}^{L}G(\CC)\) 与 \(V\in\mathrm{Rep}({}^{L}G)\)，当 \(|z|\) 充分小时有
\[
\mathrm{Tr}_{\mathcal F(V)}\bigl(\Gamma(z\,V(g))\bigr)=\det(I-z\,V(g))^{-1}.
\]
并且对任意有限指标集 \(S\) 与赋值 \(q_v>1\)（\(v\in S\)），以及任意半单族 \((g_v)_{v\in S}\subset {}^{L}G(\CC)\)，定义
\[
\mathcal H_S(V):=\bigotimes_{v\in S}\mathcal F(V),
\qquad
Z_S(V,(g_v);s):=\mathrm{Tr}_{\mathcal H_S(V)}\Bigl(\bigotimes_{v\in S}\Gamma(q_v^{-s}\,V(g_v))\Bigr),
\]
则对 \(\Re(s)\) 充分大有乘积恒等式
\begin{equation}\label{eq:xi-leyang-fock-euler-product}
\boxed{\;
Z_S(V,(g_v);s)=\prod_{v\in S}\det\!\bigl(I-q_v^{-s}V(g_v)\bigr)^{-1}\; }.
\end{equation}
\end{proposition}
\begin{proof}
取 \(A\) 的特征值为 \(\alpha_1,\dots,\alpha_d\)。
则在 \(\mathrm{Sym}^n(V)\) 上，\(\Gamma(A)\) 的迹为所有次数为 \(n\) 的单项式 \(\alpha_1^{k_1}\cdots\alpha_d^{k_d}\)（\(\sum k_i=n\)）之和；
从而
\[
\mathrm{Tr}_{\mathcal F(V)}(\Gamma(A))
=\prod_{j=1}^{d}\left(\sum_{k\ge 0}\alpha_j^{k}\right)
=\prod_{j=1}^{d}(1-\alpha_j)^{-1}
=\det(I-A)^{-1},
\]
其中收敛由 \(\rho(A)<1\) 保证。
张量积情形下迹按因子分解，得 \eqref{eq:xi-leyang-fock-euler-product}。证毕。
\end{proof}

\begin{theorem}[Lee--Yang 边界指数与 Weyl 多面体：由非解析集合重构根系]\label{thm:xi-leyang-boundary-support-function-root-system}
令 \({}^{L}G\) 为复半单群，取极大环面 \(T\subset {}^{L}G\) 与 Weyl 群 \(W\)。
记
\(
\mathfrak a:=X_{\ast}(T)\otimes_{\ZZ}\RR
\)
并令 \(\overline{\mathfrak a^{+}}\subset\mathfrak a\) 为闭支配 Weyl 室。
对半单 \(g\in{}^{L}G(\CC)\) 的极分解 \(g\sim k\exp(\mu)k^{-1}\)（\(\mu\in\overline{\mathfrak a^{+}}\)），对任意支配最高权 \(\lambda\) 与相应最高权表示 \(V_{\lambda}\)，定义 Lee--Yang 边界半径
\[
R_{\lambda}(g):=\min\{|z|:\ P_g(V_{\lambda};z)=0\}
=\rho\!\bigl(V_{\lambda}(g)\bigr)^{-1},
\]
并定义边界指数函数
\begin{equation}\label{eq:xi-leyang-boundary-exponent}
\boxed{\;
\mathcal E^{\mathrm{LY}}_{\lambda}(\mu)
:=-\log R_{\lambda}(\exp(\mu))
=\log \rho\!\bigl(V_{\lambda}(\exp(\mu))\bigr)\; }.
\end{equation}
令 \(P_{\lambda}:=\mathrm{conv}(W\cdot \lambda)\subset \mathfrak a^{\ast}\) 为 Weyl 多面体。
则对所有 \(\mu\in\overline{\mathfrak a^{+}}\) 有恒等式
\begin{equation}\label{eq:xi-leyang-boundary-support-function}
\boxed{\;
\mathcal E^{\mathrm{LY}}_{\lambda}(\mu)=h_{P_{\lambda}}(\mu)
:=\max_{x\in P_{\lambda}}\langle x,\mu\rangle\; }.
\end{equation}
因此 \(\mathcal E^{\mathrm{LY}}_{\lambda}\) 唯一确定凸体 \(P_{\lambda}\)。
若 \(\lambda\) 为正则权，则 \(P_{\lambda}\) 的正规扇等于 Weyl 扇；从而 \(\mathcal E^{\mathrm{LY}}_{\lambda}\) 的不可微集合恰为 Weyl 壁之并，并唯一确定 \({}^{L}G\) 的根系（因而根数据，至多差一个同源类型）。
\end{theorem}
\begin{proof}
在权分解下，\(V_{\lambda}(\exp(\mu))\) 对权 \(\nu\) 取特征值 \(\exp\langle \nu,\mu\rangle\)；
故谱半径为
\[
\rho\!\bigl(V_{\lambda}(\exp(\mu))\bigr)
=\max_{\nu\in\mathrm{Wt}(V_{\lambda})}\exp\langle\nu,\mu\rangle
=\exp\!\Bigl(\max_{\nu\in\mathrm{Wt}(V_{\lambda})}\langle\nu,\mu\rangle\Bigr).
\]
权多面体为权集合的凸包；最大值在紧凸集上取于极点。
对最高权表示 \(V_{\lambda}\)，其权多面体的极点为 \(W\cdot\lambda\)，从而
\[
\max_{\nu\in\mathrm{Wt}(V_{\lambda})}\langle\nu,\mu\rangle
=\max_{x\in\mathrm{conv}(W\cdot\lambda)}\langle x,\mu\rangle
=h_{P_{\lambda}}(\mu),
\]
取对数即得 \eqref{eq:xi-leyang-boundary-support-function}。
支撑函数唯一确定紧凸体，故 \(P_{\lambda}\) 可由 \(\mathcal E^{\mathrm{LY}}_{\lambda}\) 恢复。
最后，当 \(\lambda\) 正则时，\(P_{\lambda}\) 的法锥分解与 Weyl 室分解一致，等价地其正规扇为 Weyl 扇；
支撑函数在各室内线性而在壁上出现不可微性，故不可微集合为 Weyl 壁之并，从而可重构根系超平面排列与根数据。证毕。
\end{proof}

\endinput

\par

\begin{conjecture}[Kapustin--Witten 物理几何朗兰兹中的 Lee--Yang 分歧壁：范畴化对应]\label{conj:xi-terminal-zm-leyang-kapustin-witten-wall}
沿用推论 \ref{cor:xi-terminal-zm-leyang-langlands-level37} 第二条的分歧集 \(B\)、开曲线 \(U=\PP^1\setminus B\) 与局部系统 \(\mathcal L\)。
猜想在 Kapustin--Witten 所述的 \(4\mathrm d\ \mathcal N=4\) 超对称 Yang--Mills 的拓扑扭曲与边界/线算符框架下，可构造一类范畴对象 \(\mathfrak E(\mathcal L)\)，使其在几何朗兰兹等价中成为 Hecke 特征对象，特征值局部系统为 \(\mathcal L\)，并且 \(B\) 中的有限分歧值对应于耦合常数复延拓中的壁穿越与非正则极限。\cite{KapustinWitten2006GeometricLanglands}
\end{conjecture}

\begin{conjecture}[谱—外置坏素一致性：坏约化素点为谱兼容 prime-register 的不可规避轴]\label{conj:xi-terminal-zm-spectral-compatible-externalization-badprime-minimal}
设 \(\mathcal C/\QQ\) 为一条由参数无关的代数谱方程所给的特征曲线，并在 \(\QQ\) 上与某条极小椭圆曲线 \(E/\QQ\) 双有理等价。
考虑一类\textbf{谱兼容的外置制度} \(\mathcal S\)：其以 prime-register 纤维 \(\mathcal P\)（定义 \ref{def:pom-prime-fiber}）作为残差账本，使得谱反演/分歧护照的单值化选择在外置上可追溯且可比较；并要求存在有限素数集
\(\mathfrak O(\mathcal S)\subset\mathbb{P}\)
使得对任意 \(p\notin\mathfrak O(\mathcal S)\)，该制度在模 \(p\) 约化下满足“扩张保持”意义的稳定性：覆盖的分歧型与外置比较规则在 \(\FF_p\) 上不发生退化或碰撞。
则猜想不可规避的障碍素集必包含坏约化素点：
\[
\boxed{\ \{p:\ p\mid \Delta(E)\}\ \subseteq\ \mathfrak O(\mathcal S)\ }.
\]
在 Lee--Yang 原型中，由定理 \ref{thm:xi-terminal-zm-leyang-elliptic-structure} 的固定椭圆锚点知 \(\Delta(E)=37\)，故任何谱兼容外置都应满足 \(37\in\mathfrak O(\mathcal S)\)；
并且在“最小障碍谱”口径下，奇素部分可取为 \(\mathfrak O(\mathcal S)\cap(2\mathbb{Z}+1)=\{37\}\)。
\end{conjecture}
\begin{remark}
该猜想与本文关于 \(p=37\) 的多重退化证书一致：
坏约化（定理 \ref{thm:xi-terminal-zm-leyang-elliptic-structure}），
Latt\`es 稳定约化降阶与 Chebyshev 共轭（结论 \ref{con:xi-terminal-zm-elliptic-lattes-badprime-37-degree-drop}），
以及重量像的高重坍缩（例如结论 \ref{con:xi-terminal-zm-elliptic-3torsion-y-image-mod37-collapse}）
均在同一素数处同步出现。
\end{remark}

\endinput

\begin{conclusion}[时间纤维标尺与模尖点增长的同一：\texorpdfstring{$q^{-4}$}{q^{-4}} 主部锁定对数深度]\label{con:xi-terminal-zm-timefiber-cusp-depth}
沿用注记 \ref{rem:xi-terminal-elliptic-37a1-modular-anchor} 的模参数化。
令 \(q=e^{2\pi i\tau}\)，并令 \(y(\tau)\) 表示该参数化下的重量函数。
则在无穷远尖点具有主部 \(q^{-4}\)（同注记 \ref{rem:xi-terminal-elliptic-37a1-modular-anchor} 的截断展开），从而沿纯虚轴 \(\tau=iT\to i\infty\) 有
\[
\boxed{\;
\log y(\tau)=8\pi\,\Im\tau+O(1)\;}
\qquad(T\to+\infty).
\]
特别地，若以 \(b_{-1}:=\log B^{-1}\) 表示端点时间纤维的对数深度坐标，则 Cayley 标尺变换只改变其加性零点（备注 \ref{rem:dephys-cayley-scale-gauge-four}）；
因此在模尖点处，\(q^{-4}\) 主部把该深度与 \(8\pi\,\Im\tau\) 在同一标尺类中锁定到有界误差。
\end{conclusion}
\begin{proof}[证明要点]
由注记 \ref{rem:xi-terminal-elliptic-37a1-modular-anchor} 的 \(q\)-展开可写
\[
y(\tau)=q^{-4}\bigl(1+O(q)\bigr)\qquad(\Im\tau\to+\infty).
\]
沿 \(\tau=iT\) 有 \(q=e^{-2\pi T}\)，故
\[
\log y(iT)=-4\log q+\log(1+O(q))=8\pi T+O(1),
\]
其中 \(\log(1+O(q))\) 在 \(T\to\infty\) 时有界（事实上趋于 \(0\)）。
最后，\(b_{-1}=\log B^{-1}\) 的标尺变换律由备注 \ref{rem:dephys-cayley-scale-gauge-four} 给出，因而上述 \(O(1)\) 误差与标尺常数可在同一加性规范下合并吸收。证毕。
\end{proof}


\begin{conclusion}[重量谱椭圆曲线上的倍乘—Latt\`es 自映射：\texorpdfstring{$x([2]P)$}{x([2]P)} 的闭式四次有理函数]\label{con:xi-terminal-zm-elliptic-lattes-doubling}\label{con:xi-terminal-zm-lattes-4division-certificate}
在结论 \ref{con:xi-terminal-zm-elliptic-normalization} 的椭圆模型
\[
E:\ Y^2=X^3-X+\frac14
\]
上，记倍乘映射 \([2]:E\to E\)，并令 \(x:E\to\mathbb P^1\) 为 \(x\)-坐标函数 \(x(P)=X\)。
则 \(x\circ[2]\) 在仿射坐标 \(X\) 上诱导一个四次 Latt\`es 自映射 \(\Phi\in\QQ(X)\)：
\[
\boxed{\;
\Phi(X):=x([2]P)=\frac{X^4+2X^2-2X+1}{4X^3-4X+1}\; }.
\]
其分母
\(
4X^3-4X+1
=4\bigl(X^3-X+\tfrac14\bigr)
=4Y^2
\)
等价于 \(E\) 的 \(2\)-division 条件 \(Y=0\)（从而恰为 \(x\)-坐标的 \(2\)-division 多项式）。
此外，\(\Phi\) 的临界点可由 \(4\)-division 数据完全刻画：记 \(E[2]\)、\(E[4]\) 分别为 \(2\)-挠与 \(4\)-挠子群，则对仿射坐标 \(X\) 有严格等价
\[
\boxed{\;
\Phi'(X)=0
\ \Longleftrightarrow\
\exists\,Q\in E[4]\setminus E[2]\ \text{使得}\ X=x(Q)
\ \Longleftrightarrow\
2X^{6}-10X^{4}+10X^{3}-10X^{2}+2X+1=0\; }.
\]
从而该六次多项式给出 \(x\)-坐标图上的 \(4\)-division 证书（等价于经典 \(4\)-division 多项式 \(\psi_4\) 除去因子 \(2Y\) 后的 \(X\)-部分），倍乘 Latt\`es 动力学的临界结构在 \(\QQ[X]\) 层面完全可审计。
\par
更进一步，记
\[
D(X):=4X^{3}-4X+1,\qquad
C(X):=2X^{6}-10X^{4}+10X^{3}-10X^{2}+2X+1.
\]
则 \(2\)-division 多项式在 \(\Phi\) 下满足严格平方拉回恒等式
\[
\boxed{\;
D\bigl(\Phi(X)\bigr)=\frac{C(X)^{2}}{D(X)^{3}}\ \in\ \QQ(X)\; }.
\]
并且对任意 \(P=(X,Y)\in E(\overline{\QQ})\) 且 \(Y\neq 0\)，倍点的 \(Y\)-坐标满足同一六次多项式的闭式表达
\[
\boxed{\;
Y([2]P)=\frac{C(X)}{16\,Y^{3}}\ \in\ \QQ(E)\; }.
\]
\end{conclusion}
\begin{proof}[证明要点]
对短 Weierstrass 曲线 \(Y^2=X^3+aX+b\)（此处 \(a=-1,b=1/4\)），倍点公式给出
\[
x([2]P)=\left(\frac{3X^2+a}{2Y}\right)^{\!2}-2X.
\]
代入并用 \(4Y^2=4X^3-4X+1\) 化简，即得所示 \(\Phi(X)\)。
对 \(\Phi=N/D\)（\(N=X^4+2X^2-2X+1\)、\(D=4X^3-4X+1\)）直接求导并清分母，
\(
\Phi'(X)=0\iff N'D-ND'=0
\)，化简得到所示六次多项式。
另一方面，对一般短 Weierstrass 曲线 \(Y^2=X^3+aX+b\)，四除多项式满足
\[
\psi_4=4Y\bigl(X^6+5aX^4+20bX^3-5a^2X^2-4abX-8b^2-a^3\bigr).
\]
代入 \((a,b)=(-1,\tfrac14)\) 得
\[
\psi_4(X)=2Y\bigl(2X^{6}-10X^{4}+10X^{3}-10X^{2}+2X+1\bigr),
\]
从而该六次因子刻画 \(x\bigl(E[4]\setminus E[2]\bigr)\)，与 Latt\`es 临界集合一致。
又由 division 多项式的标准恒等式
\(
Y([2]P)=\psi_4(P)\big/ \bigl(2\,\psi_2(P)^{4}\bigr)
\)
（其中 \(\psi_2(P)=2Y\)），代入上式 \(\psi_4(P)=2Y\cdot C(X)\) 即得
\(
Y([2]P)=C(X)/(16Y^{3})
\)。
\par
最后，将 \(D(X)=4X^{3}-4X+1\) 视为 \(2\)-division 三次，则
\(
D(\Phi(X))=4\Phi(X)^{3}-4\Phi(X)+1
\)；
写 \(\Phi=N/D\) 并清分母得到
\[
D(\Phi(X))=\frac{4N(X)^{3}-4N(X)D(X)^{2}+D(X)^{3}}{D(X)^{3}}.
\]
对分子直接化简与因式分解可得
\(
4N^{3}-4ND^{2}+D^{3}=C(X)^{2}
\)，从而得到 \(D(\Phi(X))=C(X)^{2}/D(X)^{3}\)。
可复现核验脚本：\texttt{scripts/exp\_fold\_zm\_elliptic\_lattes\_rational\_points\_audit.py}。证毕。
\end{proof}

\begin{conclusion}[四次 Latt\`es 自映射的临界值消去恒等式与后临界肖像的闭式描述]\label{con:xi-terminal-zm-elliptic-lattes-pcf-portrait}
沿用结论 \ref{con:xi-terminal-zm-elliptic-lattes-doubling} 的记号，令
\[
\Phi(X)=\frac{N(X)}{D(X)},\qquad
N(X)=X^{4}+2X^{2}-2X+1,\qquad
D(X)=4X^{3}-4X+1,
\]
并令临界多项式
\[
C(X)=2X^{6}-10X^{4}+10X^{3}-10X^{2}+2X+1.
\]
则对任意形式变量 \(T\) 存在严格消去恒等式
\[
\boxed{\;
\mathrm{Res}_{X}\!\bigl(C(X),\,N(X)-T\,D(X)\bigr)
=37^{3}\bigl(4T^{3}-4T+1\bigr)^{2}\; }.
\]
因此 \(\Phi\) 的临界值集合在 \(\overline{\QQ}\) 中严格等于
\[
\Phi\bigl(\mathrm{Crit}(\Phi)\bigr)
=\{\,T\in\overline{\QQ}:\ 4T^{3}-4T+1=0\,\}
=x\bigl(E[2]\setminus\{O\}\bigr),
\]
并且对每一临界值 \(\alpha\)（即 \(D(\alpha)=0\) 的根），方程 \(\Phi(X)=\alpha\) 在 \(\overline{\QQ}\) 中恰有两解，且两解皆为二重根；
从而每一临界值纤维中恰含两个单纯临界点，总临界度数 \(2\deg\Phi-2=6\) 精确分配为 \(3\times 2\)。
\par
又由于 \(D(\alpha)=0\) 等价于 \(\alpha\) 为 \(\Phi\) 的极点，必有
\[
\alpha\in \Phi\bigl(\mathrm{Crit}(\Phi)\bigr)\ \Longrightarrow\ \Phi(\alpha)=\infty,
\qquad
\Phi(\infty)=\infty.
\]
从而后临界集合为
\[
\mathrm{PostCrit}(\Phi)=\{\infty\}\ \cup\ \{\,T\in\overline{\QQ}:\ 4T^{3}-4T+1=0\,\},
\]
并且对任意有限临界点 \(c\in\mathrm{Crit}(\Phi)\) 必有 \(\Phi(c)\in x\bigl(E[2]\setminus\{O\}\bigr)\) 且 \(\Phi^{2}(c)=\infty\)，从而后临界肖像由扭点塔 \(E[4]\to E[2]\to O\) 的 \(x\)-像完全决定。
该后临界肖像在纯代数意义下完全由上述消去恒等式锁定，因而可逐项审计。
\end{conclusion}
\begin{proof}[证明要点]
由结论 \ref{con:xi-terminal-zm-elliptic-lattes-doubling}，有限临界点集合为
\(
\mathrm{Crit}(\Phi)=x\bigl(E[4]\setminus E[2]\bigr)
\)，且临界多项式 \(C(X)\) 刻画 \(\Phi'(X)=0\)。
对变量 \(X\) 消去临界方程 \(C(X)=0\) 与像约束 \(N(X)-T D(X)=0\)，可得结果式恒等式；其右端因子 \(\bigl(4T^{3}-4T+1\bigr)^{2}\) 表明临界值集合恰为 \(D(T)=0\) 的根集，并且每个临界值在消去意义下以二次重数出现。
\par
对任意 \(\alpha\) 满足 \(D(\alpha)=0\)，在 \(\QQ(\alpha)[X]\) 中有严格平方分解
\[
N(X)-\alpha D(X)=\bigl(X^{2}-2\alpha X-2\alpha^{2}+1\bigr)^{2},
\]
从而 \(\Phi(X)-\alpha\) 在 \(\alpha\) 的纤维上呈完全平方，纤维中两点均为二重前像点，因而为单纯临界点。
\par
最后，当 \(D(\alpha)=0\) 时 \(\alpha\) 为 \(\Phi\) 的极点，故 \(\Phi(\alpha)=\infty\)；
而由 \(\deg N=4\)、\(\deg D=3\) 得 \(X\to\infty\) 时 \(\Phi(X)\sim \tfrac14 X\)，从而 \(\Phi(\infty)=\infty\)。
可复现核验脚本：\texttt{scripts/exp\_fold\_zm\_elliptic\_lattes\_rational\_points\_audit.py}。证毕。
\end{proof}



\begin{conclusion}[Latt\`es 周期点方程的整式分解：\texorpdfstring{$\Phi^n(X)=X$}{Phi^n(X)=X} 的分子由 \texorpdfstring{$(2^n\pm1)$}{(2^n±1)}-division 乘积给出]\label{con:xi-terminal-zm-elliptic-lattes-periodic-division-factorization}
沿用结论 \ref{con:xi-terminal-zm-elliptic-lattes-doubling} 的 Latt\`es 自映射
\[
\Phi(X)=\frac{X^4+2X^2-2X+1}{4X^3-4X+1}\in\QQ(X).
\]
对每个整数 \(n\ge 1\)，将 \(\Phi^n(X)-X\) 写为既约分式
\[
\Phi^n(X)-X=\frac{A_n(X)}{B_n(X)},
\qquad
A_n(X),B_n(X)\in\ZZ[X],\ \gcd(A_n,B_n)=1.
\]
对任意奇数 \(m\ge 1\)，定义 \(x\)-坐标的 \(m\)-division 整式
\[
\Psi_m(X):=\prod_{Q\in E[m]\setminus\{O\}/\{\pm1\}}\bigl(X-x(Q)\bigr)\ \in\ \ZZ[X],
\]
其零点集合按重数计恰为非平凡 \(m\)-扭点的 \(x\)-坐标像（次数为 \((m^2-1)/2\)）。
则对所有 \(n\ge 1\) 成立严格分解律
\[
\boxed{\;
A_n(X)=\pm\,\Psi_{2^n-1}(X)\,\Psi_{2^n+1}(X)\; }.
\]
特别地：
\begin{enumerate}
  \item \(n=1\)（不动点方程）：
  \[
  \Phi(X)-X=-\frac{3X^4-6X^2+3X-1}{4X^3-4X+1},
  \qquad
  A_1(X)=-(3X^4-6X^2+3X-1).
  \]
  \item \(n=2\)（二周期点方程）：
  \[
  A_2(X)=-(3X^4-6X^2+3X-1)\cdot Q_{12}(X),
  \]
  其中
  \[
  \begin{aligned}
  Q_{12}(X)=\;&5X^{12}-62X^{10}+95X^{9}-105X^{8}-60X^{7}+285X^{6}-174X^{5}\\
  &-5X^{4}-5X^{3}+35X^{2}-15X+2.
  \end{aligned}
  \]
  并且 \(Q_{12}(X)\) 的零点集合正对应于 \(5\)-扭点在 \(x\)-坐标上的像（其次数 \(12=(5^2-1)/2\) 与计数完全吻合）。
\end{enumerate}
\end{conclusion}
\begin{proof}[证明要点]
由 \(\Phi(X)=x([2]P)\)（其中 \(X=x(P)\)）可得迭代满足 \(\Phi^n(X)=x([2^n]P)\)。
因此 \(\Phi^n(X)=X\) 等价于 \(x([2^n]P)=x(P)\)，亦即 \([2^n]P=\pm P\)。
于是
\[
[2^n]P=P\ \Longrightarrow\ (2^n-1)P=O,
\qquad
[2^n]P=-P\ \Longrightarrow\ (2^n+1)P=O.
\]
反之，若 \(P\) 为 \((2^n-1)\)-扭或 \((2^n+1)\)-扭点，则显然有 \([2^n]P=\pm P\)，从而 \(X=x(P)\) 为 \(\Phi^n(X)=X\) 的解。
由于 \(2^n\pm1\) 为奇数，上述两类扭点在 \(x\)-坐标下自然按 \(\pm\) 识别，故对应的一元整式恰为 \(\Psi_{2^n-1}\) 与 \(\Psi_{2^n+1}\) 的乘积；与 \(A_n\) 的差别至多为一个整体符号。
\par
两条显式多项式分解（\(n=1,2\)）可由直接代入化简与多项式除法核验得到；并且次数计数 \(\deg \Psi_3=4\)、\(\deg \Psi_5=12\) 与 \(A_1,A_2\) 的次数相容，从而 \(Q_{12}\) 可识别为 \(\Psi_5\)（至单位因子）。
证毕。
\end{proof}

\begin{conclusion}[临界点分裂域的带符号置换结构：\texorpdfstring{$C_{2}\wr S_{3}$}{C2 wr S3} 的显式呈现与判别式素数 \(37\)]\label{con:xi-terminal-zm-elliptic-lattes-critical-sextic-galois}
沿用结论 \ref{con:xi-terminal-zm-elliptic-lattes-doubling} 的记号，并记
\[
C(X):=2X^{6}-10X^{4}+10X^{3}-10X^{2}+2X+1\in\QQ[X],
\qquad
D(T):=4T^{3}-4T+1\in\QQ[T].
\]
设 \(\alpha_{1},\alpha_{2},\alpha_{3}\) 为 \(D(T)=0\) 的三根。
则在其分裂域中，临界六次多项式发生完全分解并具有闭式乘积结构
\[
\boxed{\;
C(X)=2\prod_{i=1}^{3}\bigl(X^{2}-2\alpha_{i}X-2\alpha_{i}^{2}+1\bigr)\; }.
\]
特别地，对每一 \(\alpha_i\)，临界点方程在 \(\Phi(X)=\alpha_i\) 的纤维上等价于二次方程
\[
X^{2}-2\alpha_{i}X-2\alpha_{i}^{2}+1=0,
\]
从而对应的两临界点为
\[
X_{i,\pm}=\alpha_{i}\pm\sqrt{\,3\alpha_{i}^{2}-1\,}.
\]
并且其范数满足严格恒等式
\[
\boxed{\;
\operatorname{N}_{\QQ(\alpha_{i})/\QQ}\!\bigl(3\alpha_{i}^{2}-1\bigr)=-\frac{37}{16}\; }.
\]
\par
此外，\(C(X)\) 在 \(\QQ\) 上不可约，其判别式为
\[
\boxed{\ \mathrm{Disc}(C)=-2^{8}\cdot 37^{5}=-2^{8}\cdot \Delta(E)^{5}\ }.
\]
其中 \(\Delta(E)=37\) 为椭圆曲线 \(E:\ Y^2=X^3-X+\tfrac14\) 的判别式。
因而临界点分裂域必含二次子域 \(\QQ(\sqrt{-37})\)，且除 \(2\) 与 \(37\) 外无其它素分歧。
更精确地，设 \(L=\mathrm{Spl}(C)\) 为 \(C\) 的分裂域，则其伽罗瓦群同构于三维带符号置换群（Coxeter 群 \(B_{3}\)）
\[
\boxed{\;
\mathrm{Gal}(L/\QQ)\ \cong\ C_{2}^{3}\rtimes S_{3}\ \cong\ C_{2}\wr S_{3}\ \cong\ S_{4}\times C_{2}\; } ,
\]
其阶为 \(48\)。
该结构可解释为：三次多项式 \(D(T)\) 的分裂域提供 \(S_{3}\)-对称的“临界值三重体”，而三组二次分支 \(\sqrt{3\alpha_{i}^{2}-1}\) 叠加给出 \(C_{2}^{3}\) 的符号翻转层；两者半直积给出全域对称性。
\end{conclusion}
\begin{proof}[证明要点]
若 \(\alpha\) 满足 \(D(\alpha)=0\)，则在 \(\QQ(\alpha)[X]\) 中存在严格恒等式
\[
\bigl(X^{4}+2X^{2}-2X+1\bigr)-\alpha\bigl(4X^{3}-4X+1\bigr)=\bigl(X^{2}-2\alpha X-2\alpha^{2}+1\bigr)^{2},
\]
由此得到所示二次因子与 \(C(X)\) 的三组分裂；并由二次判别式立即得到临界点闭式
\(
X_{i,\pm}=\alpha_i\pm\sqrt{3\alpha_i^{2}-1}
\)。
范数恒等式可由结果式计算
\(
\operatorname{N}_{\QQ(\alpha)/\QQ}(3\alpha^{2}-1)=\mathrm{Res}_{\alpha}\bigl(\alpha^{3}-\alpha+\tfrac14,\ 3\alpha^{2}-1\bigr)
\)
直接得到。
\par
不可约性、判别式分解与伽罗瓦群阶为直接符号计算结论；
其中 \(\mathrm{Gal}(L/\QQ)\cong S_{4}\times C_{2}\) 与 \(C_{2}^{3}\rtimes S_{3}\) 的同构为全八面体群的标准等价表述。
可复现核验脚本：\texttt{scripts/exp\_fold\_zm\_elliptic\_lattes\_rational\_points\_audit.py}。证毕。
\end{proof}

\begin{conclusion}[八面体 \texorpdfstring{$4$}{4}-挠射影场的 \texorpdfstring{$S_3$}{S3} 商：临界六次分裂域强制包含 \texorpdfstring{$E[2]$}{E[2]} 的伽罗瓦闭包]\label{con:xi-terminal-zm-elliptic-lattes-critical-sextic-s3-quotient}
沿用结论 \ref{con:xi-terminal-zm-elliptic-lattes-critical-sextic-galois} 的记号，设
\[
L:=\mathrm{Spl}\bigl(C(X)\bigr),\qquad L_{2}:=\mathrm{Spl}\bigl(D(T)\bigr).
\]
则有包含关系 \(L_2\subset L\)，并且由 \(C_2^3\rtimes S_3\) 的半直积结构给出规范满射
\[
\boxed{\;
\mathrm{Gal}(L/\QQ)\twoheadrightarrow \mathrm{Gal}(L_2/\QQ)\cong S_3\;}
\]
其核同构于 \(C_2^3\)（对应三组平方根 \(\sqrt{3\alpha_i^2-1}\) 的独立符号翻转）。
因此 \(L_2=\mathrm{Spl}\bigl(D_2(X)\bigr)\) 可识别为 \(\QQ(E[2])\) 的 \(S_3\)-伽罗瓦闭包（结论 \ref{con:xi-terminal-zm-collision-s2-e2-field-isomorphism}），从而 \(L\) 在群论意义上不可避免地支配 \(E[2]\) 的 \(S_3\)-闭包。
特别地，由结论 \ref{con:xi-terminal-zm-collision-s2-e2-field-isomorphism}，
\(
\QQ(\Lambda)\cong \QQ(E[2])
\)
为该 \(S_3\)-闭包的三次子域，因而亦嵌入 \(L\)。
\end{conclusion}
\begin{proof}[证明要点]
由结论 \ref{con:xi-terminal-zm-elliptic-lattes-critical-sextic-galois} 的分解式，
\(
C(X)=2\prod_{i=1}^{3}(X^{2}-2\alpha_{i}X-2\alpha_{i}^{2}+1)
\)，
可知 \(L\) 由 \(D(T)\) 的分裂域与三组二次分支生成：
\[
L=L_2\bigl(\sqrt{3\alpha_1^2-1},\sqrt{3\alpha_2^2-1},\sqrt{3\alpha_3^2-1}\bigr).
\]
对每个 \(i\)，平方根的符号翻转给出一个阶 \(2\) 自同构；三者生成的正规子群同构于 \(C_2^3\)，并在商上诱导对三元组 \((\alpha_1,\alpha_2,\alpha_3)\) 的置换作用。
因此存在短正合列
\(
1\to C_2^3\to \mathrm{Gal}(L/\QQ)\to \mathrm{Gal}(L_2/\QQ)\to 1
\)，
得到所示满射。
最后，\(D(T)\) 与结论 \ref{con:xi-terminal-zm-collision-s2-e2-field-isomorphism} 的 \(D_2(X)\) 同型（仅变量记号不同），其分裂域即为 \(E[2]\) 的 \(S_3\)-伽罗瓦闭包。
由该结论的三次域同构 \(\QQ(\Lambda)\cong\QQ(E[2])\)，立得嵌入断言。证毕。
\end{proof}



\begin{conclusion}[重量参数的二次扩张与判别式的椭圆化：\(y\) 在 \texorpdfstring{$\QQ(X)$}{Q(X)} 上的最小多项式]\label{con:xi-terminal-zm-y-quadratic-extension-minpoly}\label{con:xi-terminal-zm-elliptic-quadratic-discriminant}
沿用结论 \ref{con:xi-terminal-zm-elliptic-normalization} 的坐标约定（\(X=\lambda,\ Y=w/2\)，从而 \(y=(2X^2+w-1)/2=X^2+Y-\tfrac12\)）。
令物理主支的重量参数取为
\[
\boxed{\;
y:=X^2+Y-\frac12\ \in\ \QQ(E)\; }.
\]
则 \(y\) 在 \(\QQ(X)\) 上满足严格二次最小多项式
\[
\boxed{\;
y^2-(2X^2-1)\,y+X(X-1)^2(X+1)=0\; }.
\]
其判别式恰为
\[
\boxed{\;
(2X^2-1)^2-4X(X-1)^2(X+1)
=4X^3-4X+1
=4Y^2\; }.
\]
因此重量参数的双支共轭由椭圆反演 \(\iota:(X,Y)\mapsto(X,-Y)\) 给出：
\[
y^{\iota}=(2X^2-1)-y
=X^2-Y-\frac12,
\qquad
yy^{\iota}=X(X-1)^2(X+1).
\]
从而“重量变量—谱变量”的非线性耦合在函数域层面被压缩为一个二次扩张与其椭圆判别式，结构上完全刚性。
\end{conclusion}
\begin{proof}[证明要点]
由定义 \(Y=y-X^2+\frac12\)，代入曲线方程 \(Y^2=X^3-X+\frac14\) 并展开化简，得到所示二次多项式。
判别式恒等式由直接展开验证，并用 \(4Y^2=4X^3-4X+1\) 改写即可。
共轭与乘积恒等式来自二次方程两根之和与两根之积。证毕。
\end{proof}

\begin{conclusion}[倍点拉回下重量参数的线性变换：临界六次作为线性系数]\label{con:xi-terminal-zm-elliptic-y-doubling-linear}
在函数域 \(\QQ(E)=\QQ(X,y)\) 中，倍乘自同态 \([2]:E\to E\) 对物理主支重量参数 \(y\) 的拉回满足严格恒等式
\[
\boxed{\;
[2]^*(y)=\frac{C(X)\,y+B(X)}{(4X^3-4X+1)^2}\; },
\]
其中
\[
C(X)=2X^{6}-10X^{4}+10X^{3}-10X^{2}+2X+1,
\]
\[
B(X)=-X^{8}+7X^{6}-14X^{5}+27X^{4}-9X^{3}-6X^{2}+X+1.
\]
特别地，在二次扩张 \(\QQ(E)/\QQ(X)\) 的自然基 \(\{1,y\}\) 下，\([2]^*\) 在 \(y\) 上严格呈线性；
其线性系数恰为 Latt\`es 临界六次 \(C(X)\)。
因此
\[
\boxed{\;
C(X)=0\quad\Longleftrightarrow\quad [2]^*(y)\in \QQ(X)\; } ,
\]
即倍乘 Latt\`es 的临界性可被解释为“重量二次分歧依赖在倍点拉回下发生坍缩”的代数证书。
\end{conclusion}
\begin{proof}[证明要点]
由倍点公式写出 \([2](X,Y)=(X_2,Y_2)\) 的闭式有理表达，再将 \(y_2=X_2^2+Y_2-\tfrac12\) 视为 \(\QQ(X)\)-线性函数于 \(y\)，并在四次证书 \(F(X,y)=0\) 下化简即可得到所示系数 \(C(X)\)、\(B(X)\) 与分母 \((4X^3-4X+1)^2\)。
可复现核验脚本：\texttt{scripts/exp\_fold\_zm\_elliptic\_lattes\_rational\_points\_audit.py}。证毕。
\end{proof}

\begin{conclusion}[倍点拉回在权平面四次模型中的三次/三次有理表达]\label{con:xi-terminal-zm-elliptic-y-doubling-cubic-fraction}
在结论 \ref{con:xi-terminal-zm-elliptic-quadratic-discriminant} 的首一四次表示
\[
\QQ(E)\cong \QQ(y)(X)/(F(X,y)),
\qquad
F(X,y):=X^{4}-X^{3}-(2y+1)X^{2}+X+y(y+1),
\]
下，记 \(y_2:=[2]^*(y)\in\QQ(E)\)。
则 \(y_2\) 在扩张基 \(\{1,X,X^{2},X^{3}\}\) 上可写为关于 \(X\) 的三次/三次分式：
\[
\boxed{\;
y_2=\frac{\mathcal N(X,y)}{\mathcal D(X,y)}\in \QQ(y)(X)/(F)\;},
\]
其中 \(\mathcal N,\mathcal D\in\ZZ[y][X]\) 均为关于 \(X\) 的三次多项式，并可取为
\[
\begin{aligned}
\mathcal N(X,y)=\;&(-2y^{2}-10y+1)X^{3}+(2y^{3}-6y^{2}+34y+27)X^{2}\\
&\quad+(2y^{3}+10y^{2}+12y-23)X-(y^{4}+23y^{2}+23y-1),\\[1mm]
\mathcal D(X,y)=\;&(64y+8)X^{3}+(48y^{2}+16y)X^{2}-(16y^{2}+48y+8)X-(32y^{3}+32y^{2}-1).
\end{aligned}
\]
该表示把倍点拉回的非线性压缩为 \(\QQ(y)\)-线性基上的有限阶分式表达，从而为分歧/极点的局部诊断提供直接的代数接口。
\end{conclusion}
\begin{proof}[证明要点]
由结论 \ref{con:xi-terminal-zm-elliptic-y-doubling-linear}，\(y_2=[2]^*(y)\) 在 \(\QQ(E)=\QQ(X,y)\) 中给出为 \(\QQ(X)\)-线性分式。
将该表达视为 \(\QQ(y)(X)/(F)\) 中的元素，并以 \(F(X,y)=0\) 对 \(X\) 的高次幂作降次归约后清分母，即可得到一个关于 \(X\) 次数不超过 \(3\) 的分子分母表示；
对所示 \(\mathcal N,\mathcal D\) 的逐项代入核验可由脚本 \texttt{scripts/exp\_fold\_zm\_elliptic\_lattes\_rational\_points\_audit.py} 完成。证毕。
\end{proof}

\begin{conclusion}[重量椭圆曲线的有理点结构：\texorpdfstring{$E(\QQ)_{\mathrm{tors}}=\{O\}$}{E(Q)\_tors={O}}，且物理点为无限阶]\label{con:xi-terminal-zm-elliptic-rational-points-denominator-law}
在 \(E:\ Y^2=X^3-X+\frac14\) 上，有理挠点群为平凡群
\[
\boxed{\;E(\QQ)_{\mathrm{tors}}=\{O\}\; }.
\]
将 \(E\) 通过整变换
\[
(\widetilde X,\widetilde Y)=(4X,8Y)
\]
写成积分模型
\[
\widetilde E:\ \widetilde Y^{2}=\widetilde X^{3}-16\widetilde X+16,
\qquad
\Delta(\widetilde E)=2^{12}\cdot 37,
\]
则在 \(p\neq 2,37\) 处具有良好约化；并且点计数给出
\[
\card{\widetilde E(\FF_{5})}=8,
\qquad
\card{\widetilde E(\FF_{7})}=9,
\qquad
\gcd(8,9)=1,
\]
从而推出 \(E(\QQ)_{\mathrm{tors}}=\{O\}\)。
特别地，点
\[
P=\left(2,-\frac52\right)\in E(\QQ)
\qquad(\text{等价地 }-P=(2,\tfrac52))
\]
为无限阶点，并满足 \(y(P)=1\)（其中 \(y\) 为结论 \ref{con:xi-terminal-zm-y-quadratic-extension-minpoly} 的物理主支）。
此外直接代入得 \(y(-P)=6\)。
并且其在 \(\FF_{11}\) 上的约化满足
\[
\mathrm{ord}\bigl(P\bmod 11\bigr)=17,
\]
因而 \(P\) 不可能为挠点；于是 \(\mathrm{rank}\,E(\QQ)\ge 1\)，且 \(E(\QQ)\) 为无挠的自由阿贝尔群。
写
\(
x(nP)=u_n/v_n^2
\)
为既约形式（\(v_n\in\ZZ_{>0}\)），则 \(v_n\) 构成一条椭圆分母序列，并且由
\(
y=X^2+Y-\tfrac12
\)
得到严格整性尺度律：
\[
\boxed{\;
y(nP)\in\QQ,\qquad \mathrm{den}\bigl(y(nP)\bigr)=v_n^{4}\; }.
\]
前六项可取为（列出 \(v_n\) 与 \(y(nP)\)）：
\[
\begin{array}{c|cccccc}
n & 1 & 2 & 3 & 4 & 5 & 6\\ \hline
v_n & 1 & 5 & 29 & 65 & 16264 & 2382785\\
y(nP)
& 1
& \dfrac{96}{625}
& \dfrac{2679201}{707281}
& \dfrac{252115165246}{17850625}
& \dfrac{162231294668143481}{69969611497148416}
& \dfrac{2533325779202487675774201}{32235872541947843696250625}
\end{array}
\]
相应的谱根（亦即 \(x(nP)\) 的有理值）前六项为
\[
x(P)=2,\quad
x(2P)=\dfrac{21}{25},\quad
x(3P)=\dfrac{1357}{841},\quad
x(4P)=\dfrac{480106}{4225},\quad
x(5P)=\dfrac{683916417}{264517696},\quad
x(6P)=\dfrac{4881674119706}{5677664356225}.
\]
从而“重量自由能的椭圆化”进一步提升为一条可直接进行局部—整体诊断的有理点动力学：同一物理点的倍乘轨道在分母上携带严格的四次幂尺度。
\end{conclusion}
\begin{proof}[证明要点]
（i）挠点平凡性可由有限多个良约化素数的模 \(p\) 群阶作 gcd 审计：由约化注入
\(
E(\QQ)_{\mathrm{tors}}\hookrightarrow E(\mathbb F_p)
\)
知 \(|E(\QQ)_{\mathrm{tors}}|\) 整除 \(\gcd_p \#E(\mathbb F_p)\)；
在积分模型 \(\widetilde E\) 上选取良约化素数 \(p=5,7\) 即得 \(\gcd(\#\widetilde E(\FF_5),\#\widetilde E(\FF_7))=\gcd(8,9)=1\)，从而 \(E(\QQ)_{\mathrm{tors}}=\{O\}\)。
（ii）由 \(\mathrm{ord}(P\bmod 11)=17\) 可直接排除 \(P\) 为挠点，从而 \(P\) 为无限阶点，\(\mathrm{rank}\,E(\QQ)\ge 1\)。
（iii）令 \(x(nP)=u_n/v_n^2\)、\(Y(nP)=r_n/v_n^3\) 为既约表示，则
\(
y(nP)=x(nP)^2+Y(nP)-\tfrac12
\)
的分母整除 \(v_n^4\)；
对该轨道可核验得到精确等号 \(\mathrm{den}(y(nP))=v_n^4\)，并可回收所示前六项。
可复现核验脚本：\texttt{scripts/exp\_fold\_zm\_elliptic\_lattes\_rational\_points\_audit.py}。证毕。
\end{proof}

% (User input window)

\begin{conclusion}[最小模型的 Mordell--Weil 自由生成元与分支纤维的低倍点对齐：\texorpdfstring{$(R,Q_0,Q_0',P)$}{(R,Q0,Q0',P)} 的统一表示]\label{con:xi-terminal-zm-elliptic-mw-generator-alignment}
考虑最小整 Weierstrass 模型
\[
E_0:\ y_0^2+y_0=x^3-x,
\]
其与结论 \ref{con:xi-terminal-zm-elliptic-normalization} 的短模型 \(E\) 由 \(X=x\)、\(Y=y_0+\tfrac12\) 互相对应。
由该曲线的标准 Mordell--Weil 数据（Cremona 标记 \(37\mathrm{a}1\)，见 \cite{LMFDBEllipticCurve37a1}），
\(E_0(\QQ)\) 无挠且 \(\mathrm{rank}=1\)，可取自由生成元 \(R_0=(0,0)\)。
记其在 \(E(\QQ)\) 中对应点为
\[
R=(0,\tfrac12).
\]
则存在完全显式的低倍关系
\[
[2]R=\left(1,\tfrac12\right)=-Q_0,\qquad
[3]R=\left(-1,-\tfrac12\right)=-Q_0',\qquad
[4]R=\left(2,-\tfrac52\right)=P,
\]
其中 \(Q_0=(1,-\tfrac12)\)、\(Q_0'=(-1,\tfrac12)\)，并且 \(P=(2,-\tfrac52)\) 为结论 \ref{con:xi-terminal-zm-elliptic-rational-points-denominator-law} 的物理基点。
由于 \(E(\QQ)\) 无挠，方程 \([4]G=P\) 在 \(E(\QQ)\) 中仅有唯一解 \(G=R\)。
\par
由结论 \ref{con:xi-terminal-zm-elliptic-lattes-doubling} 的倍点 Latt\`es 自映射 \(\Phi\)（以 \(x\)-坐标为变量）满足对任意 \(G\in E(\overline{\QQ})\) 有
\(
x([2]G)=\Phi\bigl(x(G)\bigr)
\)；
从而
\[
x(P)=x([4]R)=\Phi^{\circ 2}\!\bigl(x(R)\bigr),
\qquad
x(R)=0,\ \Phi(0)=1,\ \Phi(1)=2.
\]
并且
\(
y(Q_0)=0,\ y(Q_0')=1
\)，
且 \(-Q_0=(1,\tfrac12)\)、\(-Q_0'=(-1,-\tfrac12)\) 分别给出 \(y=1\) 与 \(y=0\) 纤维上的另一条有理点。
结合结论 \ref{con:xi-terminal-zm-elliptic-normalization} 中 \(\mathrm{div}(y)\)、\(\mathrm{div}(y-1)\) 的二重纤维点可见，
\(Q_0\) 与 \(Q_0'\) 正是覆盖 \(y:E\to\mathbb P^1_y\) 在 \(y=0,1\) 处的有理分支点。
因此 \((R,Q_0,Q_0',P)\) 作为同一自由生成元轨道上的低倍点组，给出“生成元—分支纤维—物理基点”的统一表示。
\end{conclusion}
\begin{proof}[证明要点]
倍点与加法公式可直接核对所示坐标恒等式；分支点身份由结论 \ref{con:xi-terminal-zm-elliptic-normalization} 的除子证书（\(y=0\) 在 \(X=1\) 处二重、\(y=1\) 在 \(X=-1\) 处二重）给出。
可复现核验脚本：\texttt{scripts/exp\_fold\_zm\_elliptic\_mw\_height\_denominator\_growth\_audit.py}。证毕。
\end{proof}

\begin{conclusion}[有理耦合的无穷性：\texorpdfstring{$\Pi(\lambda,y_0)$}{Pi(lambda,y0)} 在 \texorpdfstring{$\QQ$}{Q} 上可解的特化参数无限]\label{con:xi-terminal-zm-pi-rational-root-infinite-y}
沿用结论 \ref{con:xi-terminal-zm-elliptic-mw-generator-alignment} 的记号，并令 \(y=X^2+Y-\tfrac12\in\QQ(E)\)。
则存在无穷多个 \(y_0\in\QQ\) 使得 \(\Pi(\lambda,y_0)\in\QQ[\lambda]\) 含一次因子；
等价地，集合 \(y\bigl(E(\QQ)\bigr)\subset\QQ\) 为无限集。
\end{conclusion}
\begin{proof}[证明要点]
由结论 \ref{con:xi-terminal-zm-elliptic-mw-generator-alignment}，\(E(\QQ)\) 无挠且 \(\mathrm{rank}=1\)，故 \(E(\QQ)\) 为无限集。
另一方面，结论 \ref{con:xi-terminal-zm-elliptic-normalization} 给出 \(y:E\to\PP^1_y\) 为次数 \(4\) 的非常数有理函数，
从而对任意 \(c\in\PP^1(\overline{\QQ})\) 纤维 \(y^{-1}(c)\) 的点数（按重数）至多为 \(4\)，尤其 \(y^{-1}(c)\cap E(\QQ)\) 有限。
若 \(y\bigl(E(\QQ)\bigr)\) 为有限集，则 \(E(\QQ)\) 为有限多个有限集之并，从而有限，矛盾。
最后由结论 \ref{con:xi-terminal-zm-pi-rational-root-iff-y-image} 的充要条件得到“可解特化”等价表述。证毕。
\end{proof}

\begin{conclusion}[物理基点的二分解与 \(2\)-可除性：有理半点唯一]\label{con:xi-terminal-zm-elliptic-physical-basepoint-halving-uniqueness}
沿用结论 \ref{con:xi-terminal-zm-elliptic-mw-generator-alignment} 的记号，并记物理基点 \(P=(2,-\tfrac52)\)。
则存在唯一 \(Q\in E(\QQ)\) 满足 \([2]Q=P\)，且
\[
Q=\Bigl(1,\frac12\Bigr)=-Q_0.
\]
等价地，
\[
[2]^{-1}(\{\pm P\})\cap E(\QQ)=\{\pm Q\}.
\]
因此 \(P\in 2E(\QQ)\)，并且在 \(E(\QQ)\) 的 Mordell--Weil 群中并非原始点。
由于 \(E(\QQ)_{\mathrm{tors}}=\{O\}\)（结论 \ref{con:xi-terminal-zm-elliptic-rational-points-denominator-law}），点 \(Q\) 为无限阶。
\end{conclusion}
\begin{proof}[证明要点]
由结论 \ref{con:xi-terminal-zm-elliptic-mw-generator-alignment}，取 \(Q:=[2]R=(1,\tfrac12)\) 且 \(P=[4]R\)，则 \([2]Q=[4]R=P\)。
若 \(Q'\in E(\QQ)\) 亦满足 \([2]Q'=P\)，则 \(Q'-Q\in E[2](\QQ)\subseteq E(\QQ)_{\mathrm{tors}}\)，从而 \(Q'=Q\)。
对 \([2]Q'=-P\) 的情形同理得到 \(Q'=-Q\)。
最后，由 \(E(\QQ)_{\mathrm{tors}}=\{O\}\) 推出 \(Q\neq O\) 必为无限阶，且 \(P\) 的 \(2\)-可除性随即给出其非原始性。证毕。
\end{proof}

\begin{conclusion}[标准 N\'eron--Tate 高度的高精度常数与严格倍乘标度：可引用的数值证书]\label{con:xi-terminal-zm-elliptic-nt-height-scaling}
沿用结论 \ref{con:xi-terminal-zm-elliptic-mw-generator-alignment} 的记号。
在 LMFDB 的标准 N\'eron--Tate 规范化下（见 \cite{LMFDBEllipticCurve37a1}），秩一调节子由高度配对给出，且满足
\[
\langle R_0,R_0\rangle=0.05111140815411797\ldots.
\]
相应的 N\'eron--Tate 高度（自然对数归一）为
\[
\boxed{\;
\hat h(R)=\hat h(R_0)=\frac12\langle R_0,R_0\rangle
=0.025555704077058985\ldots\;} .
\]
由二次性 \(\hat h([n]G)=n^2\hat h(G)\) 及 \(Q_0=-[2]R\)、\(Q_0'=-[3]R\)、\(P=[4]R\) 得
\[
\hat h(Q_0)=4\hat h(R)=0.10222281630823594\ldots,
\]
\[
\hat h(Q_0')=9\hat h(R)=0.23000133669353087\ldots,
\]
\[
\hat h(P)=16\hat h(R)=0.40889126523294376\ldots.
\]
上述常数为后续分母增长、局部高度与轨道尺度化计算提供统一标尺。
\end{conclusion}

\begin{conclusion}[算术“隧穿降阶”的高度刚性：\texorpdfstring{$4{:}3$}{4:3} 的极点阶比]\label{con:xi-terminal-zm-elliptic-height-tunneling-degree-rigidity}
在 Lee--Yang 椭圆化坐标中取
\[
E:\ w^{2}=4x^{3}-4x+1,
\qquad
y=\frac{w+2x^{2}-1}{2}\in\QQ(E),
\]
并令 \(\iota:E\to E\) 为椭圆反演 \(P\mapsto -P\)。
令 \(\hat h(\cdot)\) 为 \(E(\QQ)\) 上的 N\'eron--Tate 高度，\(h(\cdot)\) 为 \(\PP^{1}(\QQ)\) 上的对数朴素高度。
则对全体 \(P\in E(\QQ)\) 成立
\[
\boxed{\;
h\bigl(y(P)\bigr)=4\hat h(P)+O(1),
\qquad
h\bigl(y(P)-y(\iota P)\bigr)=3\hat h(P)+O(1)\; }.
\]
从而具有比例律
\[
h\bigl(y(P)-y(\iota P)\bigr)=\frac{3}{4}h\bigl(y(P)\bigr)+O(1).
\]
\end{conclusion}
\begin{proof}[证明要点]
对椭圆曲线上的有理函数 \(f\in\QQ(E)\)，高度比较定理给出
\(
h(f(P))=\deg(f)\,\hat h(P)+O(1)
\)，
其中 \(\deg(f)\) 为 \(f:E\to\PP^1\) 的态射次数（等价于在无穷远点 \(O\) 处的极点阶）。
在短 Weierstrass 模型中，\(x\) 于 \(O\) 处极点阶为 \(2\)，\(w\) 极点阶为 \(3\)，而 \(y=(w+2x^2-1)/2\) 极点阶为 \(4\)（亦见定理 \ref{thm:xi-terminal-zm-leyang-elliptic-structure} 关于 \(\pi_y\) 为次数 \(4\) 覆盖的刻画）。
故 \(h(y(P))=4\hat h(P)+O(1)\) 与 \(h(w(P))=3\hat h(P)+O(1)\)。
又由 \(\iota:(x,w)\mapsto(x,-w)\) 得 \(y(P)-y(\iota P)=w(P)\)，代入即得第二式与比例律。
\end{proof}

\begin{conclusion}[分母序列的二次标度与重量分母的四次提升：以 \texorpdfstring{$\hat h(P)$}{hhat(P)} 为极限常数的增长律]\label{con:xi-terminal-zm-elliptic-denominator-growth-nt-height}
对任意 \(n\ge 1\) 将
\(
x(nP)=u_n/v_n^2
\)
写为既约形式（\(v_n\in\ZZ_{>0}\)），并沿用结论 \ref{con:xi-terminal-zm-elliptic-rational-points-denominator-law} 的重量坐标 \(y(nP)\)。
则仍有严格同位
\[
\mathrm{den}\bigl(y(nP)\bigr)=v_n^4.
\]
并且由标准高度与对数高度的比较理论，分母增长满足二次尺度律（以自然对数计）
\[
\log v_n=\hat h(P)\,n^{2}+o(n^{2}),\qquad
\log\bigl(\mathrm{den}(y(nP))\bigr)=4\hat h(P)\,n^{2}+o(n^{2}).
\]
在 \(n=1,\dots,12\) 时，\(v_n\) 的显式值为
\[
( v_n )_{n=1}^{12}=(1,\ 5,\ 29,\ 65,\ 16264,\ 2382785,\ 398035821,\ 43244638645,\ 124106986093951,\ 541051130050800400,\ 2591758672670554328449,\ 10358960321661880987253845),
\]
并且 \(\log(v_n)/n^2\) 在 \(n=8,9,10,11,12\) 上分别为
\[
0.3826584234\ldots,\ 0.4006440135\ldots,\ 0.4083229017\ldots,\ 0.4074927572\ldots,\ 0.3999992646\ldots,
\]
已围绕极限常数 \(\hat h(P)=0.40889126523294376\ldots\) 振荡收敛，形成“标准高度—分母增长”的可审计数值证书。
\end{conclusion}
\begin{proof}[证明要点]
极限律来自 N\'eron--Tate 高度与对数 Weil 高度的二次等价；显式分母与比值可由倍乘轨道直接计算并复核。
可复现核验脚本：\texttt{scripts/exp\_fold\_zm\_elliptic\_mw\_height\_denominator\_growth\_audit.py}。证毕。
\end{proof}

\begin{conclusion}[分母序列的 Ward 型双线性递推：符号因子 \texorpdfstring{$(-1)^{\lfloor n/2\rfloor}$}{(-1)^{floor(n/2)}} 的闭式控制]\label{con:xi-terminal-zm-elliptic-denominator-ward-bilinear-recurrence}
沿用结论 \ref{con:xi-terminal-zm-elliptic-denominator-growth-nt-height} 的记号。对 \(n\ge 1\) 将
\[
x(nP)=\frac{u_n}{v_n^2}\qquad(\gcd(u_n,v_n)=1,\ v_n>0)
\]
写为既约形式，并约定 \(v_0=0\)。
则 \((v_n)_{n\ge 0}\) 满足完全闭合的双线性递推：对一切 \(n\ge 2\),
\[
\boxed{\;
v_{n+2}v_{n-2}=29\,v_n^2+(-1)^{\lfloor n/2\rfloor}\,25\,v_{n+1}v_{n-1}\; }.
\]
并且由初值
\[
v_0=0,\qquad v_1=1,\qquad v_2=5,\qquad v_3=29,\qquad v_4=65
\]
唯一确定整个分母序列。
因此，物理点倍乘轨道上的分母信息可由一条四阶（Ward/Somos 型）可积双线性递推在 \(\ZZ\) 内自洽生成，而无需显式调用椭圆曲线加法公式。
\end{conclusion}
\begin{proof}[证明要点]
由结论 \ref{con:xi-terminal-zm-elliptic-rational-points-denominator-law}，点 \(P\) 为无限阶；相应的 \((v_n)\) 构成一条椭圆可除序列。
该类序列满足 Ward 型双线性递推；在将 division 多项式的符号约定推送到“分母取正”的规范化后，会出现一个仅依赖 \(n\) 的显式符号因子 \( (-1)^{\lfloor n/2\rfloor}\)。
在本文的记号下，取 \(m=2\) 的特化给出所示四阶双线性闭式，其系数由 \(v_2=5\)、\(v_3=29\) 固定。
可复现核验脚本：\texttt{scripts/exp\_fold\_zm\_elliptic\_mw\_height\_denominator\_growth\_audit.py}。证毕。
\end{proof}



\begin{conclusion}[物理基点诱导的椭圆可除序列：强可除性与低阶原始素因子]\label{con:xi-terminal-zm-elliptic-eds-strong-divisibility}
在结论 \ref{con:xi-terminal-zm-elliptic-rational-points-denominator-law} 的积分模型 \(\widetilde E\) 上，令
\(
\widetilde P:=(8,20)\in\widetilde E(\QQ)
\)
对应于 \(P\)。
对 \(n\ge 1\) 将 \(\widetilde X\bigl(n\widetilde P\bigr)\) 写成既约分式
\[
\widetilde X\bigl(n\widetilde P\bigr)=\frac{A_n}{B_n^2},
\qquad
A_n,B_n\in\ZZ,\quad \gcd(A_n,B_n)=1,\quad B_n>0,
\]
则 \(\bigl(B_n\bigr)_{n\ge 1}\) 构成椭圆可除序列，并满足强可除性
\[
\gcd(B_m,B_n)=B_{\gcd(m,n)},
\qquad
m\mid n\Longrightarrow B_m\mid B_n.
\]
其前十项为
\[
\begin{aligned}
B_1&=1,&
B_2&=5,&
B_3&=29,&
B_4&=65,&
B_5&=8132,\\
B_6&=2382785,&
B_7&=398035821,&
B_8&=43244638645,&
B_9&=124106986093951,&
B_{10}&=270525565025400200.
\end{aligned}
\]
并且在低阶上已出现若干原始分母素因子，例如
\[
29\mid B_3,\qquad 13\mid B_4,\qquad 19\mid B_5,\qquad 17\mid B_9.
\]
其中 \(\{13,17,19,29\}\) 同时属于 \(E_8\) Coxeter 指数集合 \(\{1,7,11,13,17,19,23,29\}\) 的素数子集，从而在极低阶即可给出一个可审计的指数素数共现样本。
\end{conclusion}
\begin{proof}[证明要点]
强可除性与可除性为椭圆可除序列的一般性质；
上述前十项与所示整除关系可由对 \(\widetilde P\) 的倍乘轨道直接计算并在 \(\ZZ\) 中分解验证。
可复现核验脚本：\texttt{scripts/exp\_fold\_zm\_elliptic\_lattes\_rational\_points\_audit.py}。证毕。
\end{proof}

\begin{conclusion}[分母素因子与模 \texorpdfstring{$p$}{p} 阶的严格等价：\texorpdfstring{$p\mid B_n$}{p|Bn} 的群论解码]\label{con:xi-terminal-zm-elliptic-eds-modp-order-equivalence}
设 \(p\neq 2,37\) 为 \(\widetilde E\) 的良好约化素数。
则对任意 \(n\ge 1\) 有严格等价
\[
p\mid B_n
\ \Longleftrightarrow\
n\widetilde P\equiv O\pmod p
\ \Longleftrightarrow\
\mathrm{ord}_{\widetilde E(\FF_p)}(\widetilde P)\mid n.
\]
因此 \(\bigl(B_n\bigr)\) 的素因子集合精确记录了 \(\widetilde P\) 在各有限域点群 \(\widetilde E(\FF_p)\) 中的阶分布；任何能够检测 \(p\mid B_n\) 的局部“模 \(p\) 指纹”均可等价转写为一个纯群论命题“\(\mathrm{ord}(\widetilde P\bmod p)\) 是否整除 \(n\)”。
\end{conclusion}
\begin{proof}[证明要点]
在良好约化处，\(\widetilde E(\QQ)\to\widetilde E(\FF_p)\) 的约化映射对分母的 \(p\)-进信息具有精确判别：\(p\mid B_n\) 当且仅当 \(\widetilde X(n\widetilde P)\) 在 \(p\)-进赋值下出现极点，这等价于 \(n\widetilde P\) 约化为无穷远点 \(O\)。
最后一条等价为有限阿贝尔群中元素阶的定义。证毕。
\end{proof}

% (User input window)

\begin{conclusion}[坏素数 \texorpdfstring{$37$}{37} 处的非分裂乘法约化与分母的精确周期律：首次触发指数与周期的闭式]\label{con:xi-terminal-zm-elliptic-badprime-37-nonsplit-periodicity}
考虑与结论 \ref{con:xi-terminal-zm-elliptic-normalization} 同一 \(\QQ\)-同构类的最小整模型
\[
E_0:\ y_0^{2}+y_0=x^{3}-x.
\]
其 Weierstrass 不变量满足
\(
c_6=-216
\)，从而
\[
-c_6\equiv 31\ (\mathrm{mod}\ 37),
\qquad
\Bigl(\frac{31}{37}\Bigr)=-1,
\]
因此 \(p=37\) 处为\emph{非分裂}乘法约化（Kodaira 型 \(I_1\) 的非分裂情形）。
\par
令 \(R_0=(0,0)\in E_0(\QQ)\)，并令 \(P_0=(2,-3)\in E_0(\QQ)\) 为结论 \ref{con:xi-terminal-zm-elliptic-rational-points-denominator-law} 的物理基点 \(P=(2,-\tfrac52)\) 在该最小模型下的对应点。
对任意 \(n\ge 1\) 将 \(x\)-坐标写为既约形式
\[
x(nR_0)=\frac{A_n}{B_n^{2}},
\qquad
x(nP_0)=\frac{U_n}{V_n^{2}},
\qquad
B_n,V_n\in\ZZ_{>0}.
\]
则在素数 \(37\) 处存在完全精确的 \(37\)-进分母律
\[
v_{37}(B_n)=
\begin{cases}
0,& 38\nmid n,\\
v_{37}(n)+1,& 38\mid n,
\end{cases}
\qquad
v_{37}(V_n)=
\begin{cases}
0,& 19\nmid n,\\
v_{37}(n)+1,& 19\mid n.
\end{cases}
\]
等价地，由 \(\mathrm{den}(x(nR_0))=B_n^2\)、\(\mathrm{den}(x(nP_0))=V_n^2\) 得
\[
v_{37}\!\bigl(\mathrm{den}(x(nR_0))\bigr)=
\begin{cases}
0,& 38\nmid n,\\
2\bigl(v_{37}(n)+1\bigr),& 38\mid n,
\end{cases}
\qquad
v_{37}\!\bigl(\mathrm{den}(x(nP_0))\bigr)=
\begin{cases}
0,& 19\nmid n,\\
2\bigl(v_{37}(n)+1\bigr),& 19\mid n.
\end{cases}
\]
特别地，\(R_0\) 的分母在 \(n=38\) 首次出现素因子 \(37\)，而 \(P_0\) 的分母在 \(n=19\) 首次出现 \(37\)；其后在每个 \(38\)（分别 \(19\)）的整倍处出现，并且在进一步乘以 \(37\) 的幂时严格按提升指数律累加 \(37\)-进指数。
\end{conclusion}
\begin{proof}[证明要点]
\(-c_6\) 的同余与勒让德符号给出非分裂乘法约化判据。
对非分裂乘法约化，非奇点约化点群 \(E_0^{\mathrm{ns}}(\FF_{37})\) 为一维非分裂代数环面之 \(\FF_{37}\)-点，
从而
\(
\#E_0^{\mathrm{ns}}(\FF_{37})=37+1=38
\)。
并且约化映射的核由形式群 \(E_0^{1}(\QQ_{37})\) 给出；
因此对任意 \(Q\in E_0(\QQ)\)，有
\[
[n]Q\in E_0^{1}(\QQ_{37})
\quad\Longleftrightarrow\quad
n\cdot \overline{Q}=O\ \text{于}\ E_0^{\mathrm{ns}}(\FF_{37}),
\]
其中 \(\overline{Q}\) 为模 \(37\) 约化。
对 \(R_0\) 与 \(P_0=[4]R_0\)，由结点三次曲线的环面正规形（Tate/结点参数化）可将
\(
E_0^{\mathrm{ns}}(\FF_{37})
\)
同构到
\(
\{u\in \FF_{37^2}^\times:\mathrm N_{\FF_{37^2}/\FF_{37}}(u)=1\}
\)
之乘法群（阶 \(38\)），并由直接代入得到 \(\overline{R_0}\) 的阶为 \(38\)，从而 \(\overline{P_0}\) 的阶为 \(19\)。
故 \(nR_0\)（分别 \(nP_0\)）落入形式群当且仅当 \(38\mid n\)（分别 \(19\mid n\)），这给出分母出现的必要充分条件。
\par
令 \(m:=\mathrm{ord}(\overline{Q})\)，则对任意整数 \(n\ge 1\) 有 \([mn]Q\in E_0^{1}(\QQ_{37})\)，并且其局部参数的 \(37\)-进赋值由乘法型形式群的提升指数律控制；
等价地，在 Tate 曲线（或环面）参数 \(u\) 下有
\(
v_{37}(u^{mn}-1)=v_{37}(u^{m}-1)+v_{37}(n)
\)
（LTE，因 \(37\nmid m\)），并且首次触发处 \(v_{37}(u^m-1)=1\)。
由 \(x\)-坐标在形式群参数下的主部 \(x\sim (u-1)^{-2}\) 得到
\(
v_{37}(B_{mn})=v_{37}(n)+1
\)；
对 \(m\nmid n\) 则 \([n]Q\notin E_0^{1}(\QQ_{37})\) 从而 \(v_{37}(B_n)=0\)。
将 \(Q=R_0\)（\(m=38\)）与 \(Q=P_0\)（\(m=19\)）代入即得所示两式。
可复现核验脚本：\texttt{scripts/exp\_fold\_zm\_elliptic\_mw\_height\_denominator\_growth\_audit.py}（其中对首次触发与周期作精确审计）。证毕。
\end{proof}

\begin{proposition}[坏素数 \texorpdfstring{$37$}{37} 分母指数的极限分布]\label{prop:xi-terminal-zm-elliptic-badprime-37-denominator-valuation-limit-law}
\leanverified{paper\_xi\_terminal\_zm\_elliptic\_badprime\_37\_denominator\_valuation\_limit\_law}
沿用结论 \ref{con:xi-terminal-zm-elliptic-badprime-37-nonsplit-periodicity} 的记号，并令
\[
V_N(n):=v_{37}\!\bigl(\mathrm{den}(x(nP_0))\bigr),
\]
其中 \(n\) 在 \(\{1,\dots,N\}\) 上均匀取样。
则当 \(N\to\infty\) 时，随机变量 \(V_N\) 的极限分布存在，且由
\[
\PP(V=0)=\frac{18}{19},
\qquad
\PP\!\bigl(V=2(k+1)\bigr)=\frac{1}{19}\cdot\frac{36}{37^{k+1}}
\quad(k\ge 0)
\]
给出。
等价地，尾分布满足
\[
\PP\!\bigl(V\ge 2(k+1)\bigr)=\frac{1}{19\cdot 37^{k}}
\quad(k\ge 0).
\]
\end{proposition}

\begin{proof}
由结论 \ref{con:xi-terminal-zm-elliptic-badprime-37-nonsplit-periodicity}，
\[
V_N(n)=
\begin{cases}
0,&19\nmid n,\\
2\bigl(v_{37}(n)+1\bigr),&19\mid n.
\end{cases}
\]
因此 \(\PP(V=0)=\lim_{N\to\infty}\#\{n\le N:19\nmid n\}/N=18/19\)。
在条件 \(19\mid n\) 下，事件 \(v_{37}(n)=k\) 的自然密度为
\[
\lim_{N\to\infty}\frac{\#\{n\le N:\ 19\mid n,\ 37^k\mid n,\ 37^{k+1}\nmid n\}}{\#\{n\le N:\ 19\mid n\}}
=\frac{36}{37^{k+1}},
\]
故
\(
\PP(V=2(k+1))=(1/19)\cdot 36/37^{k+1}
\)。
尾分布由几何级数求和得到：
\[
\PP\!\bigl(V\ge 2(k+1)\bigr)=\sum_{j\ge k}\PP\!\bigl(V=2(j+1)\bigr)=\frac{1}{19\cdot 37^{k}}.
\]
证毕。
\end{proof}

\begin{conclusion}[\texorpdfstring{$\PP^{1}(\FF_{37})$}{P1(F37)} 上的稳定约化动力学：周期谱、基点锁定与 Artin--Mazur \texorpdfstring{$\zeta$}{zeta} 函数]\label{con:xi-terminal-zm-elliptic-badprime-37-cycle-spectrum}
令 \(\overline{\Phi}\) 为结论 \ref{con:xi-terminal-zm-elliptic-lattes-badprime-37-degree-drop} 的稳定约化二次映射
\[
\overline{\Phi}(X)=\frac{X^{2}+10X+3}{4X+3}\in \FF_{37}(X)
\]
并视其为有限集 \(\PP^{1}(\FF_{37})\) 上的自映射。
则 \(\overline{\Phi}\) 的周期结构由五个互不相交的周期完全给出：
\begin{itemize}
  \item 一个 \(9\)-周期：
  \[
  1\mapsto 2\mapsto 26\mapsto 15\mapsto 6\mapsto 16\mapsto 30\mapsto 17\mapsto 31\mapsto 1;
  \]
  \item 一个 \(3\)-周期：
  \[
  3\mapsto 25\mapsto 29\mapsto 3;
  \]
  \item 三个不动点：\(\{5\},\{22\},\{\infty\}\)。
\end{itemize}
此外，结论 \ref{con:xi-terminal-zm-elliptic-badprime-37-nonsplit-periodicity} 的物理基点 \(P_{0}=(2,-3)\) 满足 \(x(P_{0})\equiv 2\pmod{37}\)，从而必落在上述唯一 \(9\)-周期之中；
并且该事实与 \(\overline{P_{0}}\in E_{0}^{\mathrm{ns}}(\FF_{37})\) 之阶为 \(19\) 的结构严格相容（\(x\)-坐标在 \(\pm\) 商下使倍乘轨道长度减半）。
\par
在 \(p=37\) 下，两类非平凡周期长度由
\[
\operatorname{ord}_{9}(2)=6,\qquad \operatorname{ord}_{19}(2)=18
\]
所控制，并在 \(\pm\)-识别后给出长度
\[
\frac{1}{2}\operatorname{ord}_{9}(2)=3,\qquad \frac{1}{2}\operatorname{ord}_{19}(2)=9.
\]
因此 \(\overline{\Phi}\) 的周期谱在同一有限动力系统中同时编码了 \((p-1)_{\mathrm{odd}}=9\) 与 \((p+1)_{\mathrm{odd}}=19\) 两类环面之奇部分。
其 Artin--Mazur \(\zeta\)-函数因而具有闭式
\[
\zeta_{\overline{\Phi},37}(t)=\frac{1}{(1-t)^{3}(1-t^{3})(1-t^{9})}.
\]
\end{conclusion}
\begin{proof}[证明要点]
周期分解可由对 \(\PP^{1}(\FF_{37})\) 的直接枚举得到；
\(\operatorname{ord}_{9}(2)\)、\(\operatorname{ord}_{19}(2)\) 的控制来自结论 \ref{con:xi-terminal-zm-elliptic-lattes-badprime-37-degree-drop} 的 Chebyshev 共轭与 Latt\`es \(\pm\)-商结构。
可复现核验脚本：\texttt{scripts/exp\_fold\_zm\_elliptic\_lattes\_rational\_points\_audit.py}。
最后，Artin--Mazur \(\zeta\)-函数由“每个长度 \(k\) 的周期贡献因子 \((1-t^{k})^{-1}\)”的有限集动力系统恒等式直接得到。证毕。
\end{proof}


\begin{conclusion}[非复乘与原始素因子定理：椭圆可除序列的原始素因子无穷性]\label{con:xi-terminal-zm-elliptic-eds-primitive-primes-noncm}
结论 \ref{con:xi-terminal-zm-elliptic-normalization} 给出
\(
j(E)=110592/37
\)
非代数整数，故 \(E\) 无复乘。
因此由非复乘椭圆可除序列的原始素因子定理，存在常数 \(n_0\) 使得对一切 \(n>n_0\)，\(B_n\) 含有原始素因子 \(p_n\)（即 \(p_n\mid B_n\) 且 \(p_n\nmid\prod_{m<n}B_m\)），并由结论 \ref{con:xi-terminal-zm-elliptic-eds-modp-order-equivalence} 得
\[
\mathrm{ord}_{\widetilde E(\FF_{p_n})}(\widetilde P)=n.
\]
从而该椭圆结构在严格意义上产生无穷多个原始素因子素数，其上 \(\widetilde P\) 具有预设的精确阶，为后续一切模 \(p\) 指纹检验提供可控的“阶—素数”供给机制。
\end{conclusion}

\begin{conclusion}[重量配分函数的域迹表示：分母仅由 Lee--Yang 三次因子控制]\label{con:xi-terminal-zm-field-trace-representation}
沿用结论 \ref{con:xi-terminal-zm-order4-rational} 的记号，令
\[
\Pi(\lambda,y):=\lambda^4-\lambda^3-(2y+1)\lambda^2+\lambda+y(y+1),
\qquad
P_{\mathrm{LY}}(y):=256y^3+411y^2+165y+32.
\]
记 \(K:=\QQ(y)(\lambda)\)（其中 \(\lambda\) 为 \(\Pi(\lambda,y)=0\) 的一根），并令
\(
s_m(y):=\mathrm{Tr}_{K/\QQ(y)}(\lambda^m)
=\sum_{i=1}^{4}\lambda_i(y)^m
\)
为四根的幂和（Newton 幂和）。
则存在唯一的元素 \(u(\lambda)\in K\)，其在基 \(\{1,\lambda,\lambda^2,\lambda^3\}\) 下的代表可取为三次多项式
\[
u(\lambda)=a_0(y)+a_1(y)\lambda+a_2(y)\lambda^2+a_3(y)\lambda^3,
\]
使得对一切 \(m\ge 0\) 有严格域迹表示
\[
\boxed{\;
Z_m(y)=\mathrm{Tr}_{K/\QQ(y)}\!\bigl(u(\lambda)\lambda^m\bigr)\; }.
\]
并且 \(u\) 的系数可显式取为
\[
\boxed{\;
a_0(y)=\frac{-19y^3-75y^2-42y-8}{P_{\mathrm{LY}}(y)},\qquad
a_2(y)=\frac{87y^2+48y+9}{P_{\mathrm{LY}}(y)}\; } ,
\]
\[
\boxed{\;
a_1(y)=\frac{80y^4-20y^3-82y^2-42y-8}{y\,P_{\mathrm{LY}}(y)},\qquad
a_3(y)=\frac{-16y^3+49y^2+31y+8}{y\,P_{\mathrm{LY}}(y)}\; } .
\]
因此尽管 \(Z_m(y)\in\ZZ[y]\) 为整系数多项式，其“谱投影权”在函数域上只可能在
\(
y=0
\)
与
\(
P_{\mathrm{LY}}(y)=0
\)
处出现极点；这为 Lee--Yang 分歧点在代数层面的不可避免性给出一条直接证书。
\end{conclusion}
\begin{proof}[证明要点]
对未知 \(u=b_0+b_1\lambda+b_2\lambda^2+b_3\lambda^3\in K\)，有恒等式
\[
\mathrm{Tr}(u\lambda^m)=b_0\,s_m+b_1\,s_{m+1}+b_2\,s_{m+2}+b_3\,s_{m+3}.
\]
由于 \(m\mapsto Z_m(y)\) 与 \(m\mapsto s_m(y)\) 均满足由 \(\Pi(\lambda,y)\) 诱导的同一四阶递推，故只需对 \(m=0,1,2,3\) 匹配即可唯一确定 \((b_0,\dots,b_3)\)；
解线性系统并化简即可得到所示 \(a_i(y)\)。
可复现核验脚本：\texttt{scripts/exp\_fold\_zm\_trace\_galois\_audit.py}。证毕。
\end{proof}

\begin{conclusion}[清分母后的可审计整式恒等式：\texorpdfstring{$yP_{\mathrm{LY}}(y)Z_m$}{yP(y)Z\_m} 为幂和的整系数组合]\label{con:xi-terminal-zm-integral-identity-power-sums}
在结论 \ref{con:xi-terminal-zm-field-trace-representation} 的记号下，乘以 \(yP_{\mathrm{LY}}(y)\) 可将域迹公式提升为整式恒等式：对一切 \(m\ge 0\) 有
\[
\boxed{\;
yP_{\mathrm{LY}}(y)\,Z_m(y)
=A_0(y)\,s_m(y)+A_1(y)\,s_{m+1}(y)+A_2(y)\,s_{m+2}(y)+A_3(y)\,s_{m+3}(y)\;}
\]
其中系数多项式为
\[
\boxed{\;
\begin{aligned}
A_0(y)&=-19y^4-75y^3-42y^2-8y,\\
A_1(y)&=80y^4-20y^3-82y^2-42y-8,\\
A_2(y)&=87y^3+48y^2+9y,\\
A_3(y)&=-16y^3+49y^2+31y+8.
\end{aligned}}
\]
该恒等式把“重量计数”（左端）与“谱幂和”（右端）在整系数层面直接对接，并将所有潜在的解析障碍压缩为显式代数集合 \(yP_{\mathrm{LY}}(y)=0\)。
\end{conclusion}
\begin{proof}
将结论 \ref{con:xi-terminal-zm-field-trace-representation} 的 \(a_i(y)\) 代入
\(
\mathrm{Tr}(u\lambda^m)=a_0 s_m+a_1 s_{m+1}+a_2 s_{m+2}+a_3 s_{m+3}
\)，并清分母即可得到。证毕。
\end{proof}

\begin{conclusion}[迹函数解释：四阶递推的几何来源与奇点集合的约束]\label{con:xi-terminal-zm-trace-pushforward-origin}
在结论 \ref{con:xi-terminal-zm-elliptic-normalization} 的椭圆化识别下，谱曲线的归一化可视为固定椭圆曲线 \(E\) 上的次数 \(4\) 有理映射
\(
y\in\QQ(E)
\)。
对任意 \(y_0\in\CC\setminus\{0,1\}\) 且 \(P_{\mathrm{LY}}(y_0)\neq 0\)，纤维 \(y^{-1}(y_0)\subset E(\CC)\) 恰含四个点；
其 \(X\)-坐标集合与四次方程 \(\Pi(\lambda,y_0)=0\) 的四根集合一致。
因此 \(\lambda\) 作为 \(y\) 的四值代数函数之全分支 Riemann 面即为 \(E(\CC)\)（等价于结论 \ref{con:xi-terminal-zm-elliptic-normalization} 的归一化描述），并由椭圆判别式 \(\Delta(E)=37\) 的统一化所锚定。
\par
在函数域层面，扩张 \(\QQ(E)/\QQ(y)\) 为次数 \(4\) 的有限可分扩张，其原始元可取为 \(X=\lambda\in\QQ(E)\)。
对每个整数 \(m\ge 0\) 定义迹函数
\[
T_m(y):=\Tr_{\QQ(E)/\QQ(y)}(X^m)\in\QQ(y).
\]
则 \((T_m)_{m\ge 0}\) 满足由 \(\Pi(\lambda,y)\) 决定的四阶线性递推（其特征多项式即为 \(\Pi\)）：
\[
T_m(y)=T_{m-1}(y)+(2y+1)T_{m-2}(y)-T_{m-3}(y)-y(y+1)T_{m-4}(y)\qquad(m\ge 4).
\]
此外，作为 \(\mathbb P^1_y\) 上的有理函数，\(T_m\) 的可能极点仅位于 \(y=\infty\)；
因而其奇点集合包含于分歧值集合
\(
\{\infty,0,1\}\cup\{P_{\mathrm{LY}}(y)=0\}
\)（见结论 \ref{con:xi-terminal-zm-elliptic-y-hurwitz-passport}）。
结合结论 \ref{con:xi-terminal-zm-field-trace-representation} 的域迹表示，四阶递推可解释为次数 \(4\) 覆叠 \(\pi_y:E\to\mathbb P^1_y\) 上的迹推送结构，而 Lee--Yang 三次因子则在该覆叠的分歧几何中被内在锁定。
\end{conclusion}
\begin{proof}[证明要点]
由 \(\Pi(X,y)=0\) 在 \(\QQ(E)\) 中成立，乘以 \(X^{m-4}\) 并取 \(\Tr_{\QQ(E)/\QQ(y)}\) 得到所示递推。
又由于 \(X\) 在 \(E\) 上仅在无穷远点 \(O\) 处有极点，而 \(y\) 在 \(O\) 处具有四阶极点（结论 \ref{con:xi-terminal-zm-elliptic-normalization}），故 \(T_m\) 在有限 \(y\)-点处不出现极点，从而其可能极点仅在 \(y=\infty\)。
可复现核验脚本：\texttt{scripts/exp\_fold\_zm\_elliptic\_audit.py}。证毕。
\end{proof}

\begin{conclusion}[谱的最大代数复杂度：\texorpdfstring{$\mathrm{Gal}_{\QQ(y)}\!\bigl(\Pi(\lambda,y)\bigr)\cong S_4$}{Gal(Q(y))(Pi)=S4}]\label{con:xi-terminal-zm-generic-galois-s4}
把 \(\Pi(\lambda,y)\) 视为关于 \(\lambda\) 的四次多项式。
其三次分解式（cubic resolvent）可取为
\[
\boxed{\;
\mathscr R(z,y)
=z^3+(1+2y)z^2-(1+4y+4y^2)z-(1+5y+13y^2+8y^3)\; }.
\]
其判别式满足严格恒等式
\[
\boxed{\;
\mathrm{Disc}_{z}\bigl(\mathscr R(\cdot,y)\bigr)
=-y(y-1)\,P_{\mathrm{LY}}(y)\; } ,
\]
并且与四次谱方程的判别式完全一致：
\[
\boxed{\;
\mathrm{Disc}_{z}\bigl(\mathscr R(\cdot,y)\bigr)
=\mathrm{Disc}_{\lambda}\bigl(\Pi(\cdot,y)\bigr)\; } .
\]
因此在 \(y\)-线上，四次谱覆盖的分歧值集合恰为 \(\{0,1\}\) 与 \(P_{\mathrm{LY}}(y)=0\) 的三根；Lee--Yang 三次因子因而给出判别式的不可约分歧核。
在任一良特化（例如 \(y=2\)）下，
\[
\mathscr R(z,2)=z^{3}+5z^{2}-25z-127
\]
在 \(\QQ\) 上不可约且其判别式为 \(-8108\)，因而非平方，
从而 \(\Pi(\lambda,y)\) 在 \(\QQ(y)\) 上的泛型伽罗瓦群为
\[
\boxed{\;\mathrm{Gal}\bigl(\Pi/\QQ(y)\bigr)\cong S_4.\;}
\]
因此 \(\Pi(\lambda,y)\) 在 \(\QQ(y)\) 上给出一个正则 \(S_4\)-扩张（常数域为 \(\QQ\)），而其唯一二次分辨率由 \(-y(y-1)P_{\mathrm{LY}}(y)\) 锁定，直接连接到 Lee--Yang 分歧集合。
等价地，谱支 \(\lambda_i(y)\) 的单值解析延拓所诱导的单值群在四个谱分支上给出满的 \(S_4\)-作用；因而不存在任何将谱覆盖全局分解为两个互不混合的二次子覆盖的代数约化能够成立。
\end{conclusion}
\begin{proof}[证明要点]
分解式 \(\mathscr R(z,y)\) 与判别式恒等式可由直接符号计算得到，并与结论 \ref{con:xi-terminal-zm-discriminant-leyang} 的判别式分解一致。
若在某特化 \(y=y_0\) 下 \(\mathscr R(z,y_0)\) 在 \(\QQ[z]\) 上不可约，则 \(\mathscr R\) 在 \(\QQ(y)[z]\) 上不可约；
对不可约四次而言，三次分解式不可约将伽罗瓦群约束为 \(\{A_4,S_4\}\)；
而判别式非平方排除 \(A_4\)，故必为 \(S_4\)。
可复现核验脚本：\texttt{scripts/exp\_fold\_zm\_trace\_galois\_audit.py}。证毕。
\end{proof}

\begin{remark}[Hitchin 谱覆盖与 \texorpdfstring{$A_3$}{A3} cameral 覆叠：\texorpdfstring{$S_4\simeq W(A_3)$}{S4=W(A3)} 的结构性解释]\label{rem:xi-terminal-zm-hitchin-cameral-cover}
结论 \ref{con:xi-terminal-zm-generic-galois-s4} 给出的泛型对称群 \(S_4\) 可被理解为一种 Weyl 单值化：在代数几何层面，四次谱覆盖 \(\Pi(\lambda,y)=0\) 的伽罗瓦闭包给出一个 \(S_4\)-cameral 覆叠，其群作用与 \(A_3\) 型 Weyl 群 \(W(A_3)\cong S_4\) 同构。
若将 \(\Pi(\lambda,y)=0\) 解释为某一秩 \(4\) Higgs 场（例如 \(G=\mathrm{SL}_4\) 或 \(\mathrm{PGL}_4\)）的特征方程在 Hitchin 基底上的一维切片，则该 \(S_4\) 单值化即对应 cameral cover 的 Weyl 对称；
而同一判别式核 \(-y(y-1)P_{\mathrm{LY}}(y)\) 则给出该切片与 Hitchin 判别式除子的交，从而把 Lee--Yang 分歧集合内在地识别为谱覆盖奇纤维的显式坐标化。
几何朗兰兹的物理实现中，Hitchin 纤维化与其谱/cameral 数据是刻画拓扑扭曲后真空结构与缺陷算子的重要几何输入。\cite{KapustinWitten2006GeometricLanglands}
\end{remark}

\begin{conclusion}[有理参数特化的四次整式模型与判别式：分歧素集的显式上界]\label{con:xi-terminal-zm-specialization-discriminant-ab}
取有理参数 \(y_0=a/b\in\QQ\)（\(\gcd(a,b)=1,\ b>0\)），并考虑清分母后的四次整式
\[
P_{a,b}(\lambda):=b^{2}\Pi(\lambda,a/b)
=b^{2}\lambda^{4}-b^{2}\lambda^{3}-(2ab+b^{2})\lambda^{2}+b^{2}\lambda+a(a+b)\in\ZZ[\lambda].
\]
则其关于 \(\lambda\) 的判别式满足严格闭式
\[
\boxed{\;
\mathrm{Disc}_{\lambda}\bigl(P_{a,b}\bigr)
=-a(a-b)\bigl(256a^{3}+411a^{2}b+165ab^{2}+32b^{3}\bigr)\,b^{7}\; }.
\]
因而当且仅当 \(a=0\)、\(a=b\) 或 \(256a^{3}+411a^{2}b+165ab^{2}+32b^{3}=0\) 时出现重根；
在其余情形，\(P_{a,b}\) 在 \(\overline\QQ\) 上可分（等价地：特化覆叠在该参数处 \'{e}tale）。
此外，若素数 \(p\) 在 \(P_{a,b}\) 的分裂域中分歧，则必整除
\[
2ab(a-b)\bigl(256a^{3}+411a^{2}b+165ab^{2}+32b^{3}\bigr),
\]
从而分歧素集被上述显式整数所包含。
结合结论 \ref{con:xi-terminal-zm-generic-galois-s4} 的泛型 \(S_4\) 性，由 Hilbert 不可约定理可知：除去薄集参数外，特化多项式 \(P_{a,b}\) 的分裂域伽罗瓦群仍为 \(S_4\)，
从而得到一族由 \(y_0\in\QQ\) 参数化的 \(S_4\) 数域族，其分歧由同一 Lee--Yang 三次核在清分母后完全控制。
\end{conclusion}
\begin{proof}[证明要点]
由结论 \ref{con:xi-terminal-zm-discriminant-leyang}，
\(
\mathrm{Disc}_{\lambda}\bigl(\Pi(\cdot,y)\bigr)=-y(y-1)P_{\mathrm{LY}}(y)
\)。
代入 \(y=a/b\) 并利用四次情形的齐次性
\(
\mathrm{Disc}_{\lambda}(c\cdot f)=c^{6}\mathrm{Disc}_{\lambda}(f)
\)
得到
\(
\mathrm{Disc}_{\lambda}\bigl(P_{a,b}\bigr)=b^{12}\mathrm{Disc}_{\lambda}\bigl(\Pi(\cdot,a/b)\bigr)
\)，
化简即得闭式公式。
分歧素集包含性来自一般判据“分裂域的分歧素数必整除多项式判别式”。
薄集断言由泛型伽罗瓦群与 Hilbert 不可约定理推出。
可复现核验脚本：\texttt{scripts/exp\_fold\_zm\_pi\_specialization\_discriminant\_ab\_audit.py}。证毕。
\end{proof}

\begin{conclusion}[“椭圆—\texorpdfstring{$S_4$}{S4}—Lee--Yang”三重对偶：同一判别式同时支配归一化分歧与分裂域二次脊柱]\label{con:xi-terminal-zm-elliptic-s4-leyang-triple-duality}
曲线 \(\Pi(\lambda,y)=0\) 的几何归一化为椭圆曲线 \(E\)（结论 \ref{con:xi-terminal-zm-elliptic-normalization}），
而同一四次谱方程在 \(\QQ(y)\) 上具有 \(S_4\) 伽罗瓦群（结论 \ref{con:xi-terminal-zm-generic-galois-s4}）。
两者在判别式层面被同一多项式
\[
-y(y-1)\,P_{\mathrm{LY}}(y)
\]
同时支配：
\begin{enumerate}
  \item 作为几何对象，它给出投影 \(\pi:E\to\mathbb P^1_y\) 的分歧除子（见结论 \ref{con:xi-terminal-zm-leyang-cubic-branch-image} 与结论 \ref{con:xi-terminal-zm-discriminant-leyang}）；
  \item 作为代数对象，它给出 \(S_4\) 分裂域的唯一二次分辨率子域 \(\QQ\bigl(\sqrt{-y(y-1)P_{\mathrm{LY}}(y)}\bigr)\)。
\end{enumerate}
因此 Lee--Yang 三次因子 \(P_{\mathrm{LY}}(y)\) 并非“热力学支的偶发奇点”，而是一个同时控制椭圆归一化的分歧几何与四次分裂域二次脊柱的共同判别式核；
从而把“解析临界”与“伽罗瓦分辨率”在结构上锁定为同一条不可约代数数据。
\end{conclusion}

\begin{conclusion}[分支点倍乘的 Lee--Yang 近平方特化：\texorpdfstring{$y=6,21$}{y=6,21} 的平方因子与四次谱方程线性裂解]\label{con:xi-terminal-zm-branchpoint-doubling-leyang-nearsquare-specializations}
令 Lee--Yang 三次因子
\[
P_{\mathrm{LY}}(y):=256y^{3}+411y^{2}+165y+32.
\]
沿用结论 \ref{con:xi-terminal-zm-elliptic-mw-generator-alignment} 的记号（\(Q_0=[2]R\)、\(Q_0'=[3]R\) 为 \(y=0,1\) 上的有理分支点），则其倍点在重量坐标上的取值满足
\[
y([2]Q_0)=6,\qquad y([2]Q_0')=21.
\]
并且出现完全显式的“近平方”特化
\[
P_{\mathrm{LY}}(6)=2\cdot 37\cdot 31^{2},\qquad P_{\mathrm{LY}}(21)=4\cdot 11\cdot 241^{2}.
\]
等价地，判别式核在这两点处可抽出完整平方因子：
\[
-6\cdot 5\cdot P_{\mathrm{LY}}(6)= -555\cdot(2\cdot 31)^{2},\qquad
-21\cdot 20\cdot P_{\mathrm{LY}}(21)= -1155\cdot(4\cdot 241)^{2}.
\]
同时，四次谱方程 \(\Pi(\lambda,y)=0\) 在对应纤维上发生一次线性裂解：
\[
\Pi(\lambda,6)=(\lambda-2)(\lambda^{3}+\lambda^{2}-11\lambda-21),\qquad
\Pi(\lambda,21)=(\lambda-6)(\lambda^{3}+5\lambda^{2}-13\lambda-77).
\]
上述恒等式把“分支点倍乘”与“Lee--Yang 判别式核的平方分解”以两组完全显式整数分解连接起来，为后续局部域分析（尤其是坏素数与平方类控制）提供最短算术入口。
\end{conclusion}
\begin{proof}[证明要点]
由结论 \ref{con:xi-terminal-zm-elliptic-mw-generator-alignment} 的低倍点坐标直接计算 \(y([4]R)\)、\(y([6]R)\) 即得 \(6,21\)；其余分解由在 \(\ZZ\) 中代入并因式分解得到。
可复现核验脚本：\texttt{scripts/exp\_fold\_zm\_elliptic\_mw\_height\_denominator\_growth\_audit.py}。证毕。
\end{proof}

\begin{conclusion}[判别式双覆叠的亏格 \(2\) 模型与 resolvent 的无分歧三重提升]\label{con:xi-terminal-zm-discriminant-ridge-etale-c3-explicit-model}
沿用结论 \ref{con:xi-terminal-zm-leyang-discriminant-ridge-polydisc-badprimes} 的属 \(2\) 曲线
\[
C:\quad w^{2}=-y(y-1)P_{\mathrm{LY}}(y),
\]
并令 \(\mathscr R(z,y)\in\ZZ[y][z]\) 为结论 \ref{con:xi-terminal-zm-generic-galois-s4} 的三次分解式，满足
\[
\mathrm{Disc}_{z}\bigl(\mathscr R(\cdot,y)\bigr)=-y(y-1)P_{\mathrm{LY}}(y).
\]
令 \(C_{6}\) 为仿射曲线
\[
w^{2}=-y(y-1)P_{\mathrm{LY}}(y),\qquad \mathscr R(z,y)=0
\]
的规范化之光滑射影模型，并记自然投影 \(\pi:C_{6}\to C\) 由 \((y,w,z)\mapsto(y,w)\) 给出。
则：
\begin{enumerate}
  \item \(\pi\) 为循环 Galois 三重覆盖，且 \(\mathrm{Gal}(C_{6}/C)\cong C_{3}\)；
  \item \(\pi\) 处处非分歧（即为 \'etale 覆盖）；
  \item 因而由 Riemann--Hurwitz 公式得到 \(g(C_{6})=4\)。
\end{enumerate}
进一步地，复合映射 \(C_{6}\to C\to\mathbb P^1_y\) 为正规 \(S_3\)-Galois 覆叠；
其分歧值集合恰为 \(\{0,1,\infty\}\cup\{P_{\mathrm{LY}}(y)=0\text{ 的三根}\}\)，并且在上述每个分歧值处的惰性群阶均为 \(2\)。
换言之，\(\QQ(C_{6})=\QQ(y,w,z)\) 为由三次分解式 \(\mathscr R(z,y)=0\) 所定义的三次扩张 \(\QQ(y)\subset\QQ(y,z)\) 的 \(S_3\)-伽罗瓦闭包，而 \(C\) 即其符号表示对应之唯一二次分辨子域的规范化。
该显式模型与结论 \ref{con:xi-terminal-zm-s4-galois-closure-genus-tower-holomorphic-differentials} 中由 \(S_4\) 子群格诱导的中间商曲线 \(C_6=\widetilde X/V_4\) 相一致。
\end{conclusion}
\begin{proof}[证明要点]
设 \(k=\QQ(y)\)，并以 \(\mathscr R(z,y)\) 定义三次扩张 \(k\subset k(z)\)。
由结论 \ref{con:xi-terminal-zm-generic-galois-s4}，\(\mathscr R\) 的判别式给出其 \(S_3\)-伽罗瓦闭包的唯一二次子扩张，
而该二次子扩张正是 \(k(w)\)（其中 \(w^2=\mathrm{Disc}_z(\mathscr R)\)），对应于曲线 \(C\) 的函数域。
因此在底变 \(k\subset k(w)\) 后，伽罗瓦群由 \(S_3\) 收缩为 \(A_3\cong C_3\)，从而 \(\pi\) 为循环三重覆盖。
\par
另一方面，\(k\subset k(z)\) 的全部惯性群（对应 \(\mathrm{Disc}_z(\mathscr R)=0\) 的分歧点）由换位生成，
其与 \(A_3\) 的交为空；故在二次底变后，循环三重商覆盖上不再出现任何惰性群，从而 \(\pi\) 处处非分歧。
最后，由 \(g(C)=2\) 与无分歧次数 \(3\) 得
\(
2g(C_{6})-2=3(2g(C)-2)=6
\)，故 \(g(C_{6})=4\)。证毕。
\end{proof}

\input{thm__xi-terminal-zm-s3-closure-jacobian-prym-a2}
\input{thm__xi-terminal-zm-cubic-cofactor-discriminant-kernel-s3-rigidity-height}

\begin{conclusion}[判别式二次脊曲线的五次判别式与坏约化素集：\texorpdfstring{$\{2,3,31,37\}$}{\{2,3,31,37\}} 的严格分解]\label{con:xi-terminal-zm-leyang-discriminant-ridge-polydisc-badprimes}
令
\[
P_{\mathrm{LY}}(y):=256y^{3}+411y^{2}+165y+32,\qquad
f(y):=-y(y-1)P_{\mathrm{LY}}(y)\in\ZZ[y],
\]
并考虑属 \(2\) 超椭圆曲线
\[
C:\ w^{2}=f(y).
\]
则多项式判别式满足严格分解
\[
\boxed{\ \mathrm{Disc}_{y}(f)=-2^{20}\cdot 3^{15}\cdot 31^{2}\cdot 37\ }.
\]
因此对任意素数 \(p\notin\{2,3,31,37\}\)，右端 \(f(y)\) 在 \(\FF_p[y]\) 中为平方自由多项式，曲线 \(C\) 在 \(p\) 处良好约化。
结合结论 \ref{con:xi-terminal-zm-elliptic-normalization} 的 \(\Delta(E)=37\) 可知：在判别式二次脊的坏约化素集 \(\{2,3,31,37\}\) 中，\(37\) 为唯一同时出现在谱椭圆曲线 \(E\) 与判别式脊曲线 \(C\) 的共同坏素数。
\end{conclusion}
\begin{proof}[证明要点]
多项式判别式的因式分解由直接符号计算得到。
对奇素数 \(p\)，模型 \(w^2=f(y)\) 在 \(\FF_p\) 上光滑当且仅当 \(f\) 在 \(\FF_p\) 上无重根，等价于 \(\gcd(f,f')=1\)，亦等价于 \(p\nmid \mathrm{Disc}_{y}(f)\)。
可复现核验脚本：\texttt{scripts/exp\_fold\_zm\_leyang\_discriminant\_ridge\_difference\_square\_audit.py}。证毕。
\end{proof}

\input{con__xi-terminal-zm-leyang-discriminant-ridge-bad-reduction-nodal-31-37}

\begin{conclusion}[判别式二次脊曲线雅可比簇的有理挠子上界：良约化注入与 \texorpdfstring{$\gcd_{p}\#J(\FF_{p})=12$}{gcd\#J(Fp)=12}]\label{con:xi-terminal-zm-discriminant-ridge-jacobian-torsion-upper-bound}
沿用结论 \ref{con:xi-terminal-zm-leyang-discriminant-ridge-polydisc-badprimes} 的属 \(2\) 曲线
\[
C:\quad w^{2}=-y(y-1)P_{\mathrm{LY}}(y),
\qquad
J:=\operatorname{Jac}(C).
\]
则：
\begin{enumerate}
  \item 由有理分歧点 \((y,w)=(0,0),(1,0)\) 与无穷远点 \(\infty\) 得到两条非零 \(2\)-挠类
  \[
  D_{0}:=[(0,0)-\infty],\qquad D_{1}:=[(1,0)-\infty]\in J(\QQ)[2],
  \]
  且 \(D_{0}\neq D_{1}\)，因此 \((\ZZ/2\ZZ)^{2}\hookrightarrow J(\QQ)_{\mathrm{tors}}\)。
  \item 对任意好素数 \(p\notin\{2,3,31,37\}\)，N\'eron 模型的专门化单射给出
  \[
  J(\QQ)_{\mathrm{tors}}\hookrightarrow J(\FF_{p}),
  \]
  从而 \(\#J(\QQ)_{\mathrm{tors}}\mid \gcd_{p}\#J(\FF_{p})\)。
  以 \(p=5,7,11,13,17,19,23,29\) 为一组审计样本，有
  \input{sections/generated/eq_fold_zm_discriminant_ridge_jacobian_torsion_upper_bound_audit}
  因而 \(\#J(\QQ)_{\mathrm{tors}}\mid 12\)。
\end{enumerate}
\end{conclusion}
\begin{proof}[证明要点]
对超椭圆曲线 \(w^{2}=f(y)\)（\(\deg f=5\)）而言，Weierstrass 点与无穷远点之差在 \(\mathrm{Pic}^{0}\) 中均为 \(2\)-挠；
此处 \(y=0,1\) 为有理分歧点，故 \(D_{0},D_{1}\in J(\QQ)[2]\) 且非零。
两者互异可由支撑不同或等价地由分歧点差结构的基本性质推出，从而得到 \((\ZZ/2\ZZ)^2\) 的注入。
\par
当 \(p\) 为好素数且 \(p\neq 2,3\) 时，挠子在专门化下保持单射，故有 \(J(\QQ)_{\mathrm{tors}}\hookrightarrow J(\FF_p)\)。
而 \(\#J(\FF_p)=L_p(1)\)，其中 \(L_p(T)\) 为 \(C/\FF_p\) 的局部 \(L\)-多项式；
上列数值可由点数 \(\#C(\FF_p),\#C(\FF_{p^2})\) 回收 \(L_p\) 后严格计算得到。
可复现核验脚本：\texttt{scripts/exp\_fold\_zm\_discriminant\_ridge\_jacobian\_torsion\_upper\_bound\_audit.py}。证毕。
\end{proof}

\begin{conclusion}[判别式二次脊曲线的 \(2\)-挠分裂域与模 \(2\) 伽罗瓦像]\label{con:xi-terminal-zm-discriminant-ridge-jacobian-2torsion-field}
沿用结论 \ref{con:xi-terminal-zm-leyang-discriminant-ridge-polydisc-badprimes} 的属 \(2\) 曲线
\(
C:\ w^{2}=-y(y-1)P_{\mathrm{LY}}(y)
\)
及其雅可比簇 \(J=\operatorname{Jac}(C)\)。
则：
\begin{enumerate}
  \item \(2\)-挠分裂域满足
  \[
  \QQ(J[2])=\QQ\bigl(\text{\(P_{\mathrm{LY}}(y)\) 的全部根}\bigr),
  \qquad
  \mathrm{Gal}\bigl(\QQ(J[2])/\QQ\bigr)\cong S_{3}.
  \]
  \item 该 \(S_{3}\)-扩张的唯一二次子域为
  \[
  \QQ(\sqrt{-111}).
  \]
  \item 模 \(2\) 伽罗瓦表示
  \[
  \rho_{J,2}:\mathrm{Gal}(\overline{\QQ}/\QQ)\longrightarrow \mathrm{Sp}_{4}(\FF_{2})
  \]
  的像同构于 \(S_{3}\)，并可解释为对三个位于 \(P_{\mathrm{LY}}(y)=0\) 的非有理 Weierstrass 点的自然置换作用（而 \(y=0,1,\infty\) 对应的三点均为有理 Weierstrass 点）。
  \item 特别地，
  \[
  J(\QQ)[2]\ \cong\ (\ZZ/2\ZZ)^2,
  \]
  且由结论 \ref{con:xi-terminal-zm-discriminant-ridge-jacobian-torsion-upper-bound} 的 \(D_0,D_1\) 生成。
\end{enumerate}
\end{conclusion}
\begin{proof}[证明要点]
对奇次数超椭圆模型 \(w^2=f(y)\)（此处 \(f(y)=-y(y-1)P_{\mathrm{LY}}(y)\)）而言，\(\overline{\QQ}\)-Weierstrass 点由 \(f(y)=0\) 的五个有限根与无穷远点给出，共 \(6\) 个；并且 \(J[2]\) 由 Weierstrass 点差的类生成。
因此 \(\QQ(J[2])\) 由 Weierstrass 点坐标生成。
由于 \(y=0,1,\infty\) 为有理 Weierstrass 点，非平凡生成仅来自 \(P_{\mathrm{LY}}(y)=0\) 的三根，从而 \(\QQ(J[2])\) 即为 \(P_{\mathrm{LY}}\) 的分裂域。
由结论 \ref{con:xi-terminal-zm-leyang-cubic-arithmetic}，\(P_{\mathrm{LY}}\) 在 \(\QQ\) 上不可约且其分裂域伽罗瓦群为 \(S_3\)，并具有唯一二次子域 \(\QQ(\sqrt{-111})\)，遂得前两点。
\par
模 \(2\) 表示 \(\rho_{J,2}\) 与 Galois 对六个 Weierstrass 点差类的作用相容；由于三个位于 \(P_{\mathrm{LY}}(y)=0\) 的 Weierstrass 点构成一个 \(S_3\)-轨道而其余三点固定，可见 \(\rho_{J,2}\) 的像同构于该置换群 \(S_3\)。
最后，由上述 \(S_3\)-作用的 Galois 不变元仅来自有理 Weierstrass 点之差，可知 \(J(\QQ)[2]\) 不含额外元素，故 \(J(\QQ)[2]\cong(\ZZ/2\ZZ)^2\)，并由 \(D_0,D_1\) 生成。证毕。
\end{proof}

\begin{conclusion}[判别式二次脊曲线在 \texorpdfstring{$\QQ$}{Q} 上的自同构刚性：除超椭圆对合外无额外对称]\label{con:xi-terminal-zm-discriminant-ridge-autq-rigidity}
仍设 \(C:w^{2}=-y(y-1)P_{\mathrm{LY}}(y)\)。
则
\[
\mathrm{Aut}_{\QQ}(C)\ \cong\ \ZZ/2\ZZ,
\]
其生成元为超椭圆对合 \((y,w)\mapsto (y,-w)\)。
\end{conclusion}
\begin{proof}[证明要点]
任意 \(\QQ\)-自同构在底曲线 \(\mathbb{P}^{1}_{y}\) 上诱导一条 Möbius 变换并置换分歧集合
\(
\{0,1,\infty\}\cup\{P_{\mathrm{LY}}(y)=0\text{ 的三根}\}
\)。
由于 \(\{0,1,\infty\}\) 为 \(\QQ\)-有理分歧点集合，诱导的 Möbius 变换必为稳定 \(\{0,1,\infty\}\) 的六个标准变换之一。
逐一代入并比较 \(P_{\mathrm{LY}}\) 的变换后系数（等价地，比对根集像）可知：除恒等以外，均不保持 \(P_{\mathrm{LY}}\) 的根集。
故诱导变换只能为恒等，从而自同构在坐标 \(y\) 上恒等，因而仅可能为超椭圆对合。证毕。
\end{proof}

由此可见，由谱四次多项式 \(\Pi(\lambda,y)\)（定理 \ref{thm:xi-terminal-zm-leyang-elliptic-structure}）所诱导的判别式二次脊曲线 \(C\) 在算术几何层面同时呈现可见的 \(2\)-挠与严格的有理挠子上界（结论 \ref{con:xi-terminal-zm-discriminant-ridge-jacobian-torsion-upper-bound}），并携带一条规范的循环三次 \'etale 覆叠（结论 \ref{con:xi-terminal-zm-discriminant-ridge-etale-c3-explicit-model}），其在 \(K=\QQ(\sqrt{-3})\) 上可 Kummer 化并进一步导出 \(\QQ\)-定义的 \(3\)-同源（结论 \ref{con:xi-terminal-zm-discriminant-ridge-3isogeny-descent}）。
结合 Lee--Yang 三次的 Chebotarev 密度律（结论 \ref{con:xi-terminal-zm-leyang-cubic-modp-splitting-density}），模约化处的分歧点可见性可被转写为分解型统计，从而把分歧谱数据与挠子/覆盖塔结构以可审计方式对接。

\begin{conclusion}[判别式二次脊上的差平方范数证书：\texorpdfstring{$Q(y)^2+27w^2=(16A(y)^3)^2$}{Q(y)^2+27w^2=(16A(y)^3)^2}]\label{con:xi-terminal-zm-leyang-discriminant-ridge-diff-square-norm}
沿用结论 \ref{con:xi-terminal-zm-leyang-discriminant-ridge-polydisc-badprimes} 的曲线 \(C\) 及其仿射坐标 \((y,w)\)。
令
\[
A(y):=1+2y,\qquad Q(y):=16+69y+219y^{2}+128y^{3}.
\]
则在 \(\QQ[y]\)（事实上在 \(\ZZ[y]\)）中同时成立两条一次恒等式
\[
\boxed{\ Q(y)-16A(y)^3=27y(y-1)\ },\qquad
\boxed{\ Q(y)+16A(y)^3=P_{\mathrm{LY}}(y)\ },
\]
其中 \(P_{\mathrm{LY}}(y)=256y^3+411y^2+165y+32\) 为结论 \ref{con:xi-terminal-zm-discriminant-leyang} 的 Lee--Yang 三次因子。
从而自动推出差平方分解
\[
\boxed{\ Q(y)^2-256A(y)^6=27y(y-1)P_{\mathrm{LY}}(y)\ }.
\]
等价地，在函数域 \(\QQ(C)=\QQ(y,w)\) 中成立范数形式
\[
\boxed{\ Q(y)^2+27w^2=(16A(y)^3)^2\ }.
\]
从而在二次扩张 \(\QQ(C)(\sqrt{-3})/\QQ(C)\) 上，元素
\[
U:=\frac{Q(y)+3w\sqrt{-3}}{16A(y)^3}
\]
满足严格范数恒等式
\[
\operatorname{N}_{\QQ(C)(\sqrt{-3})/\QQ(C)}(U)=1.
\]
该范数--\(1\) 证书把判别式核 \(-y(y-1)P_{\mathrm{LY}}(y)\) 结构化为 \(\QQ(\sqrt{-3})\) 上的二次范数表达，从而为结论 \ref{con:xi-terminal-zm-s4-galois-closure-genus-tower-holomorphic-differentials} 的循环三次无分歧覆盖给出一条 Kummer 型描述的显式代数入口。
\end{conclusion}
\begin{proof}[证明要点]
由
\(
16A(y)^3=16(1+2y)^3=16+96y+192y^2+128y^3
\)
与 \(Q(y)\) 逐项比较可得
\(
Q(y)-16A(y)^3=27y^2-27y=27y(y-1)
\)
与
\(
Q(y)+16A(y)^3=256y^3+411y^2+165y+32=P_{\mathrm{LY}}(y)
\)。
两式相乘即得差平方分解。
再将曲线方程 \(w^2=-y(y-1)P_{\mathrm{LY}}(y)\) 代入即可得到范数形式。
范数恒等式由
\(
\operatorname{N}(Q+3w\sqrt{-3})=Q^2+27w^2
\)
与分母 \((16A^3)^2\) 的显式平方性立即推出。
可复现核验脚本：\texttt{scripts/exp\_fold\_zm\_leyang\_discriminant\_ridge\_difference\_square\_audit.py}。证毕。
\end{proof}

\input{con__xi-terminal-zm-leyang-compression-chebyshev-conjugacy}

\begin{conclusion}[Chebyshev 纤维积结构：二次脊曲线上的 \(S_3\) 六次闭包与循环三重 \'etale 覆叠]\label{con:xi-terminal-zm-discriminant-ridge-chebyshev-fiber-product}
沿用结论 \ref{con:xi-terminal-zm-leyang-discriminant-ridge-diff-square-norm} 的记号，并沿用结论 \ref{con:xi-terminal-zm-leyang-compression-chebyshev-conjugacy-explicit} 中的有理函数 \(a(y)=\frac{Q(y)}{8A(y)^3}\in\QQ(y)\)。
考虑 Chebyshev 型三次映射
\[
\psi:\ \PP^1_s\longrightarrow\PP^1_a,\qquad a=s^3-3s,
\]
其关于变量 \(s\) 的判别式满足
\[
\boxed{\ \mathrm{Disc}_{s}\bigl(s^3-3s-a\bigr)=-27\,(a^2-4)\ }.
\]
由结论 \ref{con:xi-terminal-zm-leyang-discriminant-ridge-diff-square-norm} 的两条一次恒等式立得
\[
a(y)-2=\frac{Q(y)-16A(y)^3}{8A(y)^3}=\frac{27y(y-1)}{8A(y)^3},
\qquad
a(y)+2=\frac{Q(y)+16A(y)^3}{8A(y)^3}=\frac{P_{\mathrm{LY}}(y)}{8A(y)^3}.
\]
从而 \(a(y)\) 在 \(y\)-线上把分歧集
\(
\{0,1,\infty\}\cup\{P_{\mathrm{LY}}(y)=0\}
\)
压缩为 \(\{2,-2,\infty\}\subset\PP^1_a\)：
\(a^{-1}(2)=\{y=0,1,\infty\}\)，\(a^{-1}(-2)=\{P_{\mathrm{LY}}(y)=0\}\)，而 \(a^{-1}(\infty)=\{A(y)=0\}=\{y=-\tfrac12\}\)。
\par
令 \(C_6\) 为纤维积 \(C\times_{\PP^1_a}\PP^1_s\) 的正规化（其中 \(C\to\PP^1_a\) 由 \(C\to\PP^1_y\) 与 \(a(y)\) 复合得到）。
等价地，\(C_6\) 可取为仿射曲线
\[
w^{2}=-y(y-1)P_{\mathrm{LY}}(y),\qquad s^{3}-3s=a(y)=\frac{Q(y)}{8A(y)^{3}}
\]
的规范化之光滑射影模型。
则：
\begin{enumerate}
  \item 复合映射 \(C_6\to\PP^1_y\) 为次数 \(6\) 的 \(S_3\)-Galois 覆叠，其惰性群由换位（阶 \(2\) 元）生成；
  \item 诱导的次数 \(3\) 覆叠 \(C_6\to C\) 为循环 Galois 覆叠，且处处无分歧（\'etale）；
  \item 因而 \(g(C_6)=3(g(C)-1)+1=4\)。
\end{enumerate}
并且该 \(C_6\to C\) 与结论 \ref{con:xi-terminal-zm-discriminant-ridge-etale-c3-explicit-model} 中由三次分解式 \(\mathscr R(z,y)=0\) 所刻画的循环三次 \'etale 覆叠一致。
\end{conclusion}
\begin{proof}[证明要点]
\(\psi\) 在 \(a=\pm 2\) 处具有指数 \(2\) 的驯分歧，在 \(a=\infty\) 处具有指数 \(3\) 的全分歧；
另一方面，由上述分解可见：在 \(C\to\PP^1_y\) 的全部分歧点（\(y=0,1,\infty\) 与 \(P_{\mathrm{LY}}(y)=0\) 的三根）上，复合映射 \(C\to\PP^1_a\) 的分歧指数均为 \(2\)，而在 \(A(y)=0\) 的两点 \(P_\pm\) 上分歧指数为 \(3\)。
因此在每个分歧点处，复合映射 \(C\to\PP^1_a\) 的分歧指数与 \(\psi\) 的分歧指数相同；由 Abhyankar 引理（驯分歧情形）知，拉回后的三次覆盖在 \(C\) 上无分歧，从而 \(C_6\to C\) 为 \'etale 循环三重覆盖。
其与 \(S_3\)-闭包的对应由判别式恒等式
\(
\mathrm{Disc}_{s}(s^3-3s-a(y))\sim -y(y-1)P_{\mathrm{LY}}(y)
\)
给出（常数与 \(A(y)^6\) 仅贡献平方因子），遂得到 \(C_6\to\PP^1_y\) 的 \(S_3\)-Galois 性。
最后由 Riemann--Hurwitz 公式得 \(g(C_6)=4\)。证毕。
\end{proof}

\begin{conclusion}[循环三次无分歧覆盖的 Kummer--Cardano 正规形与二面体提升：显式生成元与分解式同构]\label{con:xi-terminal-zm-discriminant-ridge-kummer-dihedral}
令 \(K:=\QQ(\sqrt{-3})\)，并取
\[
\zeta_{3}:=\frac{-1+\sqrt{-3}}{2}.
\]
沿用结论 \ref{con:xi-terminal-zm-discriminant-ridge-etale-c3-explicit-model} 的循环三次 \'etale 覆盖 \(\pi:C_{6}\to C\)，并在结论 \ref{con:xi-terminal-zm-leyang-discriminant-ridge-diff-square-norm} 的记号下令
\[
u:=\frac{Q(y)+3w\sqrt{-3}}{16A(y)^3}\in K(C)^{\times}.
\]
则 \(\operatorname{N}_{K(C)/\QQ(C)}(u)=1\)，且在 \(K\) 上该循环三次覆盖可 Kummer 化为
\[
C_{6,K}:\quad t^{3}=u,
\]
其阶 \(3\) 的 deck 变换由 \(t\mapsto \zeta_{3}t\) 给出。
进一步地，置
\[
t_{0}:=\frac{2A(y)}{3}\,t
\qquad
\Bigl(\text{从而 }t_{0}^{3}=\frac{Q(y)+3w\sqrt{-3}}{54}\Bigr),
\]
并定义
\[
z:=t_{0}+\frac{4A(y)^{2}}{9t_{0}}-\frac{A(y)}{3}\ \in\ K(C_{6,K}).
\]
则 \(z\) 满足分解式方程 \(\mathscr R(z,y)=0\)，从而给出函数域的显式同构
\[
K(C_{6})\otimes_{\QQ}K
=K(y,w,t)\ \cong\ K(y,w,z),
\]
并与结论 \ref{con:xi-terminal-zm-discriminant-ridge-etale-c3-explicit-model} 中由 \(\mathscr R(z,y)=0\) 定义的 \(C_{6}\) 模型一致。
此外，超椭圆反演 \(\iota:(y,w)\mapsto(y,-w)\) 在该 Kummer 模型上提升为
\[
(y,w,t)\longmapsto (y,-w,t^{-1}),
\]
其对 \(t_0\) 的作用等价写为 \(t_{0}\mapsto \frac{4A(y)^{2}}{9t_{0}}\)，并使 \(z\) 不变；因此
\[
K(C_{6,K})^{\langle t\mapsto \zeta_{3}t\rangle}=K(y,w)=K(C)\otimes_{\QQ}K,
\qquad
K(C_{6,K})^{\langle \iota\rangle}=K(y,z).
\]
从而在 \(K\) 上得到一个显式的半直积作用
\[
\bigl\langle\ t\mapsto \zeta_{3}t,\ (y,w,t)\mapsto (y,-w,t^{-1})\ \bigr\rangle
\ \cong\ C_{3}\rtimes C_{2}\ \cong\ D_{6}\ \cong\ S_{3}
\]
于 \(C_{6,K}\) 上；该作用固定底坐标 \(y\)，并具体实现 \(C_{6,K}\to\mathbb P^{1}_{y}\) 的 \(S_{3}\)-正规单值群。
\end{conclusion}
\begin{proof}[证明要点]
由于 \(K\) 含有三次单位根 \(\zeta_3\)，任意次数 \(3\) 的循环函数域扩张在 \(K\) 上均可由 Kummer 方程 \(t^{3}=u\) 给出（其中 \(u\in K(C)^{\times}\) 唯一到乘以立方为止）。
结论 \ref{con:xi-terminal-zm-leyang-discriminant-ridge-diff-square-norm} 已构造出满足 \(\operatorname{N}_{K(C)/\QQ(C)}(u)=1\) 的显式生成元 \(u\)，其与 \(\pi:C_6\to C\) 的 \(A_3\)-子扩张一致，因而给出所示 Kummer 正规形。
最后，\(\iota\) 作用下 \(w\mapsto -w\) 使
\(
u\mapsto u^{-1}
\)，
故 \(t\mapsto t^{-1}\) 给出其提升。
对 \(t_0=\frac{2A(y)}{3}t\) 与 \(z=t_0+\frac{4A(y)^2}{9t_0}-\frac{A(y)}{3}\) 的陈述由 Cardano 三次解式的直接代入核验得到，且该构造把 Kummer 立方根与分解式根 \(z\) 的同构显式化。
可复现核验脚本：\texttt{scripts/exp\_fold\_zm\_trace\_galois\_audit.py}。
令
\(
\sigma:t\mapsto \zeta_3 t
\)
与
\(
\tau:(y,w,t)\mapsto (y,-w,t^{-1})
\)，
则满足关系 \(\tau\sigma\tau=\sigma^{-1}\)，从而生成二面体群 \(D_6\)（阶 \(6\)）。
证毕。
\end{proof}

\input{con__xi-terminal-zm-discriminant-ridge-dihedral-fixedpoints-elliptic-quotient}

\begin{conclusion}[Kummer 生成元的除子完全局部化与显式 \(3\)-挠点]\label{con:xi-terminal-zm-discriminant-ridge-kummer-divisor-3torsion}
在结论 \ref{con:xi-terminal-zm-discriminant-ridge-kummer-dihedral} 的记号下，令
\[
P_{\pm}:=\left(y=-\frac12,\ w=\pm\frac{9}{4}\sqrt{-3}\right)\in C(K),
\qquad
u:=\frac{Q(y)+3w\sqrt{-3}}{16A(y)^3}\in K(C)^{\times},
\qquad
v:=\frac{Q(y)+3w\sqrt{-3}}{Q(y)-3w\sqrt{-3}}\in K(C)^{\times}.
\]
则 \(v=u^{2}\)，并且主除子满足严格恒等式
\[
\operatorname{div}(u)=3(P_{+}-P_{-}),
\qquad
\operatorname{div}(v)=6(P_{+}-P_{-}),
\]
从而 \(u\) 与 \(v\) 的零点与极点均仅支撑在 \(\{P_{+},P_{-}\}\) 上。
因此
\[
\tau:=[P_{+}-P_{-}]\in \mathrm{Pic}^{0}(C_{K})\simeq J(C)(K)
\]
满足 \(3\tau=0\) 且 \(\tau\neq 0\)，从而 \(\mathrm{ord}(\tau)=3\)。
并且由于 \(v=u^{2}\) 且 \(\gcd(2,3)=1\)，Kummer 扩张 \(t^{3}=u\) 与 \(s^{3}=v\) 给出同一 \(\mu_{3}\)-扭子，
因而与结论 \ref{con:xi-terminal-zm-discriminant-ridge-kummer-dihedral} 所给出的循环三次覆盖在函数域意义下同构。
\end{conclusion}
\begin{proof}[证明要点]
由 \(P_{\mathrm{LY}}(-\tfrac12)=\tfrac{81}{4}\) 得
\(
w^{2}\bigl\vert_{y=-1/2}=-(-\tfrac12)(-\tfrac32)\cdot\tfrac{81}{4}=-\tfrac{243}{16}
\)，
故在 \(K=\QQ(\sqrt{-3})\) 上纤维 \(y=-\tfrac12\) 含两点 \(P_{\pm}\) 且 \(w(P_{\pm})=\pm\frac{9}{4}\sqrt{-3}\)。
另一方面，由结论 \ref{con:xi-terminal-zm-leyang-discriminant-ridge-diff-square-norm} 的范数恒等式可得
\[
\bigl(Q(y)+3w\sqrt{-3}\bigr)\bigl(Q(y)-3w\sqrt{-3}\bigr)=(16A(y)^{3})^{2}.
\]
在点 \(P_{+}\) 处有 \(Q(P_{+})+3w(P_{+})\sqrt{-3}=0\) 而 \(Q(P_{+})-3w(P_{+})\sqrt{-3}\neq 0\)，
从而后者在局部环 \(\mathcal O_{C,P_{+}}\) 中为单位，遂
\(
\mathrm{ord}_{P_{+}}(Q+3w\sqrt{-3})=\mathrm{ord}_{P_{+}}\!\bigl((16A(y)^{3})^{2}\bigr)=6
\)
（因 \(A(y)=1+2y\) 在 \(y=-\tfrac12\) 处为单零且该纤维不分歧）。
同理在 \(P_{-}\) 处得 \(\mathrm{ord}_{P_{-}}(Q-3w\sqrt{-3})=6\) 且 \(\mathrm{ord}_{P_{-}}(Q+3w\sqrt{-3})=0\)。
又 \(u=(Q+3w\sqrt{-3})/(16A(y)^{3})\)，而 \(\mathrm{ord}_{P_{\pm}}(A(y))=1\)，故
\(
\operatorname{div}(u)=3(P_{+}-P_{-})
\)；
由 \(v=(Q+3w\sqrt{-3})/(Q-3w\sqrt{-3})=u^{2}\) 得 \(\operatorname{div}(v)=6(P_{+}-P_{-})\)。
于是 \(3(P_{+}-P_{-})\) 为主除子，故 \(3\tau=0\)。
若 \(\tau=0\)，则 \(P_{+}-P_{-}\) 为主除子，推出存在仅有一次极点的非常值函数 \(C\to\PP^1\)，从而 \(C\) 必为有理曲线；
但结论 \ref{con:xi-terminal-zm-leyang-discriminant-ridge-polydisc-badprimes} 已给出 \(g(C)=2\)，矛盾，故 \(\tau\neq 0\)。
最后，由 \(v=u^{2}\) 且 \(\gcd(2,3)=1\) 得由 \(t^{3}=u\) 与 \(s^{3}=v\) 所给出的 \(\mu_3\)-扭子相同。证毕。
\end{proof}

\begin{conclusion}[显式 \(3\)-循环子群的下降与 \texorpdfstring{$\QQ$}{Q}-定义 \(3\)-同源]\label{con:xi-terminal-zm-discriminant-ridge-3isogeny-descent}
令 \(K=\QQ(\sqrt{-3})\)，并取结论 \ref{con:xi-terminal-zm-discriminant-ridge-kummer-divisor-3torsion} 的 \(3\)-挠点 \(\tau\in J(C)(K)[3]\setminus\{0\}\) 与循环子群
\[
G:=\langle \tau\rangle=\{0,\pm \tau\}\subset J(C)[3].
\]
由于 \(\mathrm{Gal}(K/\QQ)\) 的非平凡元交换 \(P_{+}\leftrightarrow P_{-}\) 从而诱导 \(\tau\mapsto -\tau\)，可知子群 \(G\) 在 \(\mathrm{Gal}(\overline{\QQ}/\QQ)\) 作用下稳定。
因此存在定义于 \(\QQ\) 的度 \(3\) 同源
\[
\varphi:\ J(C)\longrightarrow J(C)/G,
\]
其核恰为 \(G\)。
\end{conclusion}
\begin{proof}[证明要点]
对任意域 \(F\) 上的阿贝尔簇 \(A/F\)，若有限子群方案 \(G\subset A\) 在 \(\mathrm{Gal}(\overline F/F)\) 作用下稳定，则商阿贝尔簇 \(A/G\) 与自然同源 \(A\to A/G\) 均可在 \(F\) 上下降。
此处 \(F=\QQ\)、\(A=J(C)\)，而上式 \(G\) 的伽罗瓦稳定性由 \(\tau\mapsto -\tau\) 立即给出；由于 \(|G|=3\)，同源次数亦为 \(3\)。证毕。
\end{proof}

\begin{conclusion}[分裂型控制下的 \(3\)-扭专门化：良约化素数上的可见性二分律]\label{con:xi-terminal-zm-discriminant-ridge-3torsion-specialization}
设 \(p\notin\{2,3,31,37\}\) 为素数，从而曲线 \(C\) 在 \(p\) 处良好约化（结论 \ref{con:xi-terminal-zm-leyang-discriminant-ridge-polydisc-badprimes}），并且 \(p\neq 3\)。
令 \(\tau\in J(C)(\QQ(\sqrt{-3}))[3]\setminus\{0\}\) 为结论 \ref{con:xi-terminal-zm-discriminant-ridge-kummer-divisor-3torsion} 的 \(3\)-挠点。
则：
\begin{enumerate}
  \item 若 \(p\equiv 1\pmod 3\)，则 \(\sqrt{-3}\in\FF_{p}\)，从而点 \(P_{\pm}\) 与 \(\tau\) 均在 \(\FF_{p}\) 上可专门化，并给出非零点
  \[
  \tau_{p}\in J(C)(\FF_{p})[3]\setminus\{0\}.
  \]
  \item 若 \(p\equiv 2\pmod 3\)，则 \(\sqrt{-3}\notin\FF_{p}\) 但 \(\sqrt{-3}\in\FF_{p^{2}}\)，因此 \(P_{\pm}\) 与 \(\tau\) 在 \(\FF_{p^{2}}\) 上专门化，并给出非零点
  \[
  \tau_{p^{2}}\in J(C)(\FF_{p^{2}})[3]\setminus\{0\}.
  \]
\end{enumerate}
从而该 \(3\)-扭结构在有限域上的可见性由二次扩张 \(\QQ(\sqrt{-3})/\QQ\) 的分裂型（亦即 \(p\bmod 3\)）完全决定。
\end{conclusion}
\begin{proof}[证明要点]
对良约化素数 \(p\neq 3\)，N\'eron 模型的基本性质保证专门化映射在素因子与 \(p\) 互素的挠子上为单射；因此 \(\tau\neq 0\) 的 \(3\)-挠性在约化后保持非平凡。
当 \(p\equiv 1\pmod 3\) 时，\(-3\) 在 \(\FF_{p}\) 中为平方，故 \(\sqrt{-3}\) 与点 \(P_{\pm}\) 在 \(\FF_{p}\) 上可定义；
当 \(p\equiv 2\pmod 3\) 时，\(-3\) 在 \(\FF_p\) 中为非平方而在 \(\FF_{p^2}\) 中必为平方，故所需定义域扩张恰为 \(\FF_{p^2}\)。证毕。
\end{proof}

\input{con__xi-terminal-zm-discriminant-ridge-3torsion-frobenius-mod3}

\begin{conclusion}[判别式二次脊曲线的雅可比簇在 \texorpdfstring{$\QQ$}{Q} 上为简单阿贝尔曲面：不可约 Frobenius 多项式证书]\label{con:xi-terminal-zm-discriminant-ridge-jacobian-qsimple}
沿用结论 \ref{con:xi-terminal-zm-leyang-discriminant-ridge-polydisc-badprimes} 的属 \(2\) 曲线 \(C\)。
对好素数 \(p\notin\{2,3,31,37\}\) 记其局部 \(L\)-多项式为
\[
L_{p}(T)=1-a_{1}T+a_{2}T^{2}-a_{1}pT^{3}+p^{2}T^{4}.
\]
则以下局部数据成立（给出 \(\#C(\FF_p)\)、\(\#C(\FF_{p^{2}})\) 与分解形态）：
\begin{itemize}
  \item \(p=5\)：\(\#C(\FF_{5})=5\)、\(\#C(\FF_{25})=33\)，且 \(L_{5}(T)=(5T^{2}-3T+1)(5T^{2}+2T+1)\)；
  \item \(p=7\)：\(\#C(\FF_{7})=5\)、\(\#C(\FF_{49})=61\)，且 \(L_{7}(T)=(7T^{2}-4T+1)(7T^{2}+T+1)\)；
  \item \(p=11\)：\(\#C(\FF_{11})=6\)、\(\#C(\FF_{121})=130\)，且 \(L_{11}(T)=(11T^{2}+1)(11T^{2}-6T+1)\)；
  \item \(p=13\)：\(\#C(\FF_{13})=14\)、\(\#C(\FF_{169})=166\)，且 \(L_{13}(T)=169T^{4}-2T^{2}+1\) 在 \(\ZZ[T]\) 中不可约；
  \item \(p=17\)：\(\#C(\FF_{17})=15\)、\(\#C(\FF_{289})=337\)，且 \(L_{17}(T)=289T^{4}-51T^{3}+28T^{2}-3T+1\) 在 \(\ZZ[T]\) 中不可约。
\end{itemize}
由 \(p=13\)（亦可取 \(p=17\)）处 \(L_{p}(T)\) 在 \(\ZZ[T]\) 不可约可推出：
雅可比簇 \(J(C)\) 在 \(\QQ\) 上不存在非平凡分解（不含任何 \(\QQ\)-椭圆因子），即 \(J(C)\) 为 \(\QQ\)-简单阿贝尔曲面。
\end{conclusion}
\begin{proof}[证明要点]
由点数关系
\(
a_1=p+1-\#C(\FF_p)
\)
与
\(
\#C(\FF_{p^2})=p^2+1-(a_1^2-2a_2)
\)
可由 \(\#C(\FF_p)\)、\(\#C(\FF_{p^2})\) 唯一回收 \(L_p(T)\)；
若 \(J(C)\) 在 \(\QQ\) 上分解，则对几乎所有好素数 \(p\) 其约化的 Frobenius 多项式必相应分解，故任一好素数处出现不可约 \(L_p(T)\) 即排除 \(\QQ\)-分解。
可复现核验脚本：\texttt{scripts/exp\_fold\_zm\_discriminant\_ridge\_jacobian\_frobenius\_audit.py}。证毕。
\end{proof}

\begin{conclusion}[判别式二次脊曲线雅可比簇的绝对单性与端同态环刚性]\label{con:xi-terminal-zm-discriminant-ridge-jacobian-absolute-simplicity-endomorphism}
仍设结论 \ref{con:xi-terminal-zm-leyang-discriminant-ridge-polydisc-badprimes} 的属 \(2\) 曲线 \(C\) 及其雅可比簇 \(J=\operatorname{Jac}(C)\)。
则：
\begin{enumerate}
  \item \(J\) 在 \(\overline{\QQ}\) 上绝对单（absolutely simple）。
  \item \[
  \mathrm{End}_{\overline{\QQ}}(J)\ \cong\ \ZZ,
  \qquad\text{等价地}\qquad
  \mathrm{End}^0_{\overline{\QQ}}(J):=\mathrm{End}_{\overline{\QQ}}(J)\otimes_{\ZZ}\QQ\ \cong\ \QQ.
  \]
  \item 几何 N\'eron--Severi 群满足
  \[
  \mathrm{NS}_{\overline{\QQ}}(J)\ \cong\ \ZZ,
  \]
  且由主极化（theta 除子类）生成。
\end{enumerate}
\end{conclusion}
\begin{proof}[证明要点]
取良好约化素数 \(p=17,29\)（结论 \ref{con:xi-terminal-zm-leyang-discriminant-ridge-polydisc-badprimes}）。
记 \(P_p(x)\) 为 \(J\) 在 \(p\) 处的 Frobenius 特征多项式（次数 \(4\) 的 Weil 多项式），并令 \(\pi_p\) 为其任一根。
脚本 \texttt{scripts/exp\_fold\_zm\_discriminant\_ridge\_jacobian\_endomorphism\_ring\_audit.py} 给出如下可复现证书：
\input{sections/generated/eq_fold_zm_discriminant_ridge_jacobian_endomorphism_ring_audit}
其中 \(P_{17},P_{29}\) 在 \(\QQ\) 上不可约，且其伽罗瓦群均为二面体群 \(D_4\)（阶 \(8\)）。
由不可约性知 \(J\) 在 \(p=17,29\) 处的约化 \(J_p\) 皆为简单阿贝尔曲面；
由 Honda--Tate 理论（此处 \(\dim J=2\) 且 \(\deg P_p=4\)）得
\[
\mathrm{End}^0_{\overline{\FF}_{p}}(J_p)\ \cong\ \QQ(\pi_p)\qquad (p=17,29),
\]
为四次 CM 域。
\par
进一步令
\(
u_p:=\pi_p+p/\pi_p
\)。
则 \(u_p\) 的极小多项式为 \(u^2-a_1u+(a_2-2p)=0\)（其中 \(a_1,a_2\) 由 \(P_p\) 给出），从而 \(\QQ(u_{17})=\QQ(\sqrt{33})\)、\(\QQ(u_{29})=\QQ(\sqrt{73})\)；
两者互异。
又由于 \(P_{17},P_{29}\) 的伽罗瓦闭包均为 \(D_4\)，对应四次 CM 域各自仅含唯一二次子域，且该子域即其极大实子域；
故 \(\QQ(\pi_{17})\cap \QQ(\pi_{29})=\QQ\)。
\par
由良好约化处端同态专门化映射的单射性（Tate--Faltings 型结果），\(\mathrm{End}^0_{\overline{\QQ}}(J)\) 可嵌入到每个 \(\mathrm{End}^0_{\overline{\FF}_{p}}(J_p)\) 中；
于是
\[
\mathrm{End}^0_{\overline{\QQ}}(J)\ \hookrightarrow\ \QQ(\pi_{17})\cap \QQ(\pi_{29})=\QQ,
\]
从而 \(\mathrm{End}^0_{\overline{\QQ}}(J)=\QQ\)，即 \(\mathrm{End}_{\overline{\QQ}}(J)\cong\ZZ\)。
该端同态代数无非平凡幂等元，故 \(J\) 在 \(\overline{\QQ}\) 上绝对单。
最后，\(\mathrm{End}^0_{\overline{\QQ}}(J)=\QQ\) 的 Rosati 不变部分为 \(\QQ\)，从而 \(\mathrm{NS}_{\overline{\QQ}}(J)\) 的秩为 \(1\)；
而 \(J\) 作为 Jacobian 携带自然主极化，遂 \(\mathrm{NS}_{\overline{\QQ}}(J)\cong\ZZ\) 由 theta 除子类生成。证毕。
\end{proof}

\begin{conclusion}[\texorpdfstring{$S_4$}{S4} 伽罗瓦闭包曲线的属、子群塔与全纯微分的表示分解：\texorpdfstring{$g=16$}{g=16} 与 \texorpdfstring{$H^{0}(\Omega^{1})$}{H0(Omega1)} 的不可约谱]\label{con:xi-terminal-zm-s4-galois-closure-genus-tower-holomorphic-differentials}
设 \(\widetilde{X}\to\mathbb{P}^{1}_{y}\) 为四次覆盖 \(\Pi(\lambda,y)=0\) 的 \(S_{4}\) 伽罗瓦闭包（\(|S_{4}|=24\)）。
其分歧循环类型为：在 \(y=\infty\) 处惯性阶 \(4\)，在 \(y=0,1\) 及 \(P_{\mathrm{LY}}(y)=0\) 的三根处惯性阶均为 \(2\)。
等价地，其 Hurwitz 分歧护照可记为 \((4)\cdot(2)^5\)。
由 Riemann--Hurwitz 公式得到闭式属数
\[
g(\widetilde{X})=16.
\]
并且存在三条结构刚性的中间商曲线
\[
\widetilde{X}/A_{4}=C\ \ (g=2),\qquad \widetilde{X}/V_{4}=C_{6}\ \ (g=4),\qquad \widetilde{X}/S_{3}=E\ \ (g=1),
\]
其中 \(C\) 的函数域即为 \(\QQ\bigl(y,\sqrt{-y(y-1)P_{\mathrm{LY}}(y)}\bigr)\)，可取为上式判别式二次脊曲线；
而 \(E\) 正是结论 \ref{con:xi-terminal-zm-elliptic-normalization} 的椭圆归一化曲线。
此外，取任意 Sylow \(2\)-子群 \(D_{4}\le S_{4}\)（阶 \(8\) 的二面体子群），则对应中间商 \(\widetilde{X}/D_{4}\) 亦为属 \(1\) 曲线：
\[
g(\widetilde{X}/D_{4})=1,
\]
从而在同一分歧护照下出现另一条椭圆中间商（不要求该子群为正规子群）。
由于全部惯性阶均与 \(3\) 互素，诱导出一条无分歧的循环三次覆盖
\[
\pi:\ C_{6}\longrightarrow C,\qquad \mathrm{Gal}(C_{6}/C)\cong C_{3}.
\]
进一步地，\(\widetilde{X}\) 的全纯微分空间作为 \(S_{4}\)-表示具有完全显式分解
\[
H^{0}(\widetilde{X},\Omega^{1})\ \cong\ 2\cdot \mathrm{sgn}\ \oplus\ V_{2}\ \oplus\ V_{3}\ \oplus\ 3V_{3}',
\]
其中 \(\mathrm{sgn}\) 为符号表示，\(V_{2}\) 为唯一二维不可约表示，\(V_{3},V_{3}'\) 为两条三维不可约表示。
于是雅可比簇在 \(\QQ\) 上诱导出分解（至 \(\QQ\)-同源意义）
\[
J(\widetilde{X})\ \sim\ A_{\mathrm{sgn}}\times A_{2}\times A_{3}\times A_{3'},
\]
其维数分别为 \((2,2,3,9)\)（总维数 \(16\)）。
并且 \(\mathrm{sgn}\)-等型因子满足 \(A_{\mathrm{sgn}}\sim J(C)\)，同时 \(J(C_{6})\sim J(C)\times A_{2}\)。
从而“判别式二次脊的简单阿贝尔曲面”被结构性嵌入到 \(S_{4}\) 伽罗瓦闭包的表示谱之中，给出从 Lee--Yang 判别核通向高属 \(S_{4}\) 动力学闭包的全局算术通道。
\end{conclusion}
\begin{proof}[证明要点]
属数由
\(
g=1+\frac{|S_4|}{2}\bigl(-2+\sum(1-1/e_i)\bigr)
\)
直接代入惯性阶集合 \(e_\infty=4\)、其余五点 \(e_i=2\) 得到。
中间商的函数域由 \(S_4\) 的标准子群格（\(A_4\lhd S_4\)、\(V_4\lhd A_4\)）与判别式二次子域的唯一性给出，且 \(C_6\to C\) 的无分歧性由惯性阶与 \(3\) 互素推出。
对 \(D_4\) 的属数，可由 \(S_4\) 在陪集 \(S_4/D_4\) 上的次数 \(3\) 置换表示中计算由 \((4)\) 与五个 \((2)\) 所诱导的分歧型，并用 Riemann--Hurwitz 公式得到 \(g(\widetilde{X}/D_4)=1\)。
表示分解可由等式 \(\dim H^{0}(\widetilde{X},\Omega^{1})^{H}=g(\widetilde{X}/H)\)（对 \(H=S_4,A_4,V_4,S_3\)）结合 \(S_4\) 角色表解出不可约重数，因而完全刚性。
可复现核验脚本：\texttt{scripts/exp\_fold\_zm\_s4\_galois\_closure\_representation\_audit.py}。证毕。
\end{proof}




\begin{conclusion}[Lee--Yang 六次闭包与 Latt\`es 八面体闭包的 Chebotarev 独立：完全分裂素数密度 \texorpdfstring{$1/288$}{1/288}]\label{con:xi-terminal-zm-leyang-lattes-chebotarev-independence}
令
\[
P_{\mathrm{LY}}(y):=256y^{3}+411y^{2}+165y+32\in\ZZ[y],
\qquad
C(X):=2X^{6}-10X^{4}+10X^{3}-10X^{2}+2X+1\in\ZZ[X].
\]
记 \(L_{\mathrm{LY}}:=\mathrm{Spl}(P_{\mathrm{LY}})\) 与 \(L_{4}:=\mathrm{Spl}(C)\) 为其分裂域。
则：
\begin{enumerate}
  \item \(\mathrm{Gal}(L_{\mathrm{LY}}/\QQ)\cong S_{3}\)，其唯一二次子域为 \(\QQ(\sqrt{-111})\)，并且 \(3\) 在该二次域中分歧（结论 \ref{con:xi-terminal-zm-leyang-cubic-field-minimal-discriminant}）。
  \item \(\mathrm{Gal}(L_{4}/\QQ)\cong S_{4}\times C_{2}\)（阶 \(48\)），并且 \(L_{4}/\QQ\) 的分歧素数集合包含于 \(\{2,37\}\)（结论 \ref{con:xi-terminal-zm-elliptic-lattes-critical-sextic-galois}）。
  \item 因而两闭包在 \(\QQ\) 上线性互不相交：\(L_{\mathrm{LY}}\cap L_{4}=\QQ\)。
  等价地，其复合扩张满足
  \[
  \mathrm{Gal}(L_{\mathrm{LY}}L_{4}/\QQ)\ \cong\ (S_{4}\times C_{2})\times S_{3},
  \qquad [L_{\mathrm{LY}}L_{4}:\QQ]=48\cdot 6=288.
  \]
\end{enumerate}
因此对任意素数 \(p\notin\{2,3,31,37\}\)，满足
\[
P_{\mathrm{LY}}(y)\ \text{与}\ C(X)\ \text{同时在}\ \FF_{p}\ \text{上完全分裂}
\]
的素数集合具有自然密度
\[
\boxed{\;\delta=\frac{1}{288}\;}
\]
（Chebotarev 意义下的“同时落在恒等共轭类”事件）。
该密度律把 Lee--Yang 三分歧值的全体有理性与 Latt\`es 临界六次的全体可见性在素数层面严格解耦，并给出一条可复核的统计预言。
\end{conclusion}
\begin{proof}[证明要点]
由结论 \ref{con:xi-terminal-zm-leyang-cubic-field-minimal-discriminant}，\(L_{\mathrm{LY}}/\QQ\) 为 \(S_3\)-伽罗瓦扩张且其唯一二次子域为 \(\QQ(\sqrt{-111})\)，其中素数 \(3\) 分歧。
而由结论 \ref{con:xi-terminal-zm-elliptic-lattes-critical-sextic-galois}，\(L_4/\QQ\) 仅在 \(\{2,37\}\) 处分歧，故 \(3\) 在 \(L_4\) 的任一子扩张中亦不分歧。
由于 \(L_{\mathrm{LY}}\) 与 \(L_4\) 均为 \(\QQ\) 上的伽罗瓦扩张，交域 \(L_{\mathrm{LY}}\cap L_4\) 亦为 \(\QQ\) 上伽罗瓦扩张，故作为 \(L_{\mathrm{LY}}\) 的伽罗瓦子扩张只能为 \(\QQ\) 或 \(\QQ(\sqrt{-111})\)；
但后者会引入素数 \(3\) 的分歧，与 \(L_4\) 的分歧素集矛盾，遂 \(L_{\mathrm{LY}}\cap L_4=\QQ\)。
因此复合伽罗瓦群为直积并给出次数 \(288\)。
\par
对 \(p\notin\{2,3,31,37\}\)，两扩张在 \(p\) 处均未分歧。
“在 \(\FF_p\) 上完全分裂”等价于 Frobenius 落在恒等共轭类；
在线性互不相交下，复合扩张的 Frobenius 共轭类由两边的 Frobenius 独立确定，故同时全分裂的密度为
\(
1/|\,\mathrm{Gal}(L_{\mathrm{LY}}L_4/\QQ)\,|=1/288
\)。
可复现核验脚本：\texttt{scripts/exp\_fold\_zm\_leyang\_lattes\_joint\_splitting\_density\_audit.py}（输出证书：\texttt{artifacts/export/fold\_zm\_leyang\_lattes\_joint\_splitting\_density\_audit.json}）。
证毕。
\end{proof}


\begin{conclusion}[三次预解式曲线的平面三次结构与无穷远纤维的分歧型 \texorpdfstring{$(2,1)$}{(2,1)}]\label{con:xi-terminal-zm-resolvent-cubic-plane-cubic-geometry}
沿用结论 \ref{con:xi-terminal-zm-generic-galois-s4} 的三次预解式（cubic resolvent）
\[
\mathscr R(z,y)
:=z^3+(1+2y)z^2-(1+4y+4y^2)z-(1+5y+13y^2+8y^3)\in\ZZ[z,y],
\]
并令 \(\mathcal R\subset \mathbb A^2_{z,y}\) 为仿射曲线 \(\mathscr R(z,y)=0\)，令 \(\overline{\mathcal R}\subset \mathbb P^2_{[Z:Y:W]}\) 为其齐次化闭包（\(z=Z/W,\ y=Y/W\)）。
记 \(\Pi(\lambda,y)\) 为结论 \ref{con:xi-terminal-zm-order4-rational} 的谱四次多项式，\(\lambda_1,\dots,\lambda_4\) 为 \(\Pi(\lambda,y)=0\) 的四根（于代数闭包中）。
定义三组配对不变量
\[
\lambda_1\lambda_2+\lambda_3\lambda_4,\qquad
\lambda_1\lambda_3+\lambda_2\lambda_4,\qquad
\lambda_1\lambda_4+\lambda_2\lambda_3.
\]
则它们恰为 \(\mathscr R(z,y)\) 的三根；因而 \(\mathscr R\) 作为 Ferrari 预解三次在代数层面压缩了四次谱覆盖的配对张量不变量。
并且在判别式层面存在严格同一性
\[
\mathrm{Disc}_{z}\bigl(\mathscr R(\cdot,y)\bigr)=\mathrm{Disc}_{\lambda}\bigl(\Pi(\cdot,y)\bigr),
\]
从而 Lee--Yang 型分歧集合亦可等价刻画为该三次预解式退化（出现重根）的参数集合。
则：
\begin{enumerate}
  \item \(\overline{\mathcal R}\) 为非奇异平面三次曲线，且含有有理点 \((y,z)=(0,1)\)（因为 \(\mathscr R(1,0)=0\)），从而 \(\overline{\mathcal R}\) 本身即为一条定义在 \(\QQ\) 上的椭圆曲线（无需再作归一化）。
  \item 投影
  \[
  \pi^{(3)}_y:\ \overline{\mathcal R}\longrightarrow \mathbb P^1_y,\qquad [Z:Y:W]\longmapsto [Y:W]
  \]
  为次数 \(3\) 的覆盖。其分歧像由判别式严格刻画：视 \(\mathscr R\) 为关于 \(z\) 的三次多项式，则
  \[
  \mathrm{Disc}_{z}\bigl(\mathscr R(\cdot,y)\bigr)=-y(y-1)P_{\mathrm{LY}}(y),
  \]
  因而在 \(y=0\)、\(y=1\) 与 \(P_{\mathrm{LY}}(y)=0\) 的三点处发生简单分歧（分歧指数 \(2\)）。
  \item 在无穷远纤维 \(y=\infty\) 上，\(\pi^{(3)}_y\) 的纤维由两个无穷远点构成，且极点阶数分布为 \(2+1\)（分歧型 \((2,1)\)）。
  更精确地，直线 \(W=0\) 与 \(\overline{\mathcal R}\) 的交点为
  \[
  [2:1:0]\ \text{（交数 \(1\)）},\qquad [-2:1:0]\ \text{（交数 \(2\)）},
  \]
  故 \([-2:1:0]\) 为 \(y=\infty\) 上唯一的分歧点（分歧指数 \(2\)），而 \([2:1:0]\) 为非分歧点。
\end{enumerate}
\end{conclusion}
\begin{proof}[证明要点]
注意到 \(\mathscr R(1,0)=0\)，故 \((y,z)=(0,1)\in \mathcal R(\QQ)\)，从而 \(\overline{\mathcal R}(\QQ)\neq\varnothing\)。
齐次化取
\[
\mathscr R_h(Z,Y,W)
=Z^3+(W+2Y)Z^2-(W^2+4YW+4Y^2)Z-(W^3+5YW^2+13Y^2W+8Y^3).
\]
在无穷远直线 \(W=0\) 上，有因式分解
\[
\mathscr R_h(Z,Y,0)=Z^3+2YZ^2-4Y^2Z-8Y^3=(Z-2Y)(Z+2Y)^2,
\]
从而得到两点 \([2:1:0]\)、\([-2:1:0]\) 及其交数 \(1,2\)。
由 \(y=Y/W\) 可知该交数即为 \(y\) 在相应点处的极点阶数，
故 \(y=\infty\) 的纤维型为 \(2+1\)，且仅 \([-2:1:0]\) 对应分歧指数 \(2\)。
\par
作为关于 \(z\) 的三次多项式，其判别式恒等式已在结论 \ref{con:xi-terminal-zm-generic-galois-s4} 给出，故有限分歧值集合与分歧指数结论随之成立。
最后，对 \(\overline{\mathcal R}\) 的光滑性可在仿射图 \(W=1\) 与无穷远图 \(Y=1\) 上分别检验梯度不同时为零；
等价地，对 \(\mathscr R_h=0\) 的射影曲线，联立 \(\mathscr R_h=\partial_Z\mathscr R_h=\partial_Y\mathscr R_h=\partial_W\mathscr R_h=0\) 在代数闭包中无解。
证毕。
\end{proof}

\begin{conclusion}[三次预解式曲线的 Weierstrass 化与显式双有理同一]\label{con:xi-terminal-zm-resolvent-cubic-weierstrass-birational}
沿用结论 \ref{con:xi-terminal-zm-resolvent-cubic-plane-cubic-geometry} 的记号。
则平面三次曲线 \(\overline{\mathcal R}/\QQ\) 与下述椭圆曲线 \(E_{\mathscr R}/\QQ\) 双有理同构：
\[
E_{\mathscr R}:\quad Y^2=X^3-37X^2+307X+361.
\]
并且可取一组完全显式的双有理互逆变换：在 \(E_{\mathscr R}\) 上令
\[
y=-\frac{8X}{Y+19+14X-X^2},
\qquad
z=1-\frac{2X(7-X)}{Y+19+14X-X^2},
\]
则 \((y,z)\) 满足 \(\mathscr R(z,y)=0\)，从而给出 \(E_{\mathscr R}\dashrightarrow \overline{\mathcal R}\)；
反之，在开子集 \(y\neq 0\) 上，取
\[
X=\frac{7y+4-4z}{y},
\qquad
Y=X^2-14X-19-\frac{8X}{y},
\]
则 \((X,Y)\) 满足 \(E_{\mathscr R}\) 的 Weierstrass 方程，从而给出 \(\overline{\mathcal R}\dashrightarrow E_{\mathscr R}\)。
\par
此外，\(E_{\mathscr R}\) 的基本不变量可显式写为
\[
c_4=7168,\quad c_6=-341504,\quad \Delta=2^{12}\cdot 31^{2}\cdot 37,\quad
j(E_{\mathscr R})=\frac{2^{18}\cdot 7^{3}}{31^{2}\cdot 37}.
\]
\end{conclusion}
\begin{proof}[证明要点]
由所给双有理变换将 \((X,Y)\) 代入 \((y,z)\) 并清分母，可验证 \(\mathscr R(z,y)=0\) 在理想 \((Y^2-(X^3-37X^2+307X+361))\) 中为零；
反向同理由直接代入并清分母得到。
判别式、\(c_4,c_6\) 与 \(j\) 的显式表达来自对该 Weierstrass 模型套用标准不变量公式。
可复现核验脚本：\texttt{scripts/exp\_fold\_zm\_resolvent\_cubic\_weierstrass\_audit.py}。证毕。
\end{proof}

\begin{conclusion}[分解式椭圆因子的双二次模型与短 Weierstrass 正规形]\label{con:xi-terminal-zm-resolvent-cubic-biquadratic-invariants-short-weierstrass}
沿用结论 \ref{con:xi-terminal-zm-resolvent-cubic-plane-cubic-geometry} 的 \((y,z)\) 坐标，并在开子集 \(z\neq 1\) 上置
\[
x:=\frac{y}{z-1}.
\]
则对过基点 \((y,z)=(0,1)\) 的直线族 \(z=1+\frac{y}{x}\) 代入 \(\mathscr R(z,y)=0\) 并约去因子 \(y\)，可得关于 \(y\) 的二次方程
\[
(1+2x-4x^2-8x^3)y^2+(4x-17x^3)y+(4x^2-7x^3)=0.
\]
令
\[
u:=\frac{2(1+2x-4x^2-8x^3)y+(4x-17x^3)}{x},
\]
则 \((x,u)\) 满足双二次模型
\[
\mathcal Q:\quad
u^2=x\bigl(65x^3+16x^2-16x-4\bigr)
=65x^4+16x^3-16x^2-4x,
\]
并与平面三次 \(\overline{\mathcal R}\) 在 \(\QQ\) 上双有理等价。
对二元四次 \(65x^4+16x^3-16x^2-4x\) 的经典不变量，有
\[
I=448,\qquad J=-10672,\qquad 4I^3-J^2=2^8\cdot 3^3\cdot 31^2\cdot 37,
\]
从而其 Jacobian 与 \(E_{\mathscr R}\) 在 \(\QQ\) 上同构，并同具
\[
j=\frac{2^{18}\cdot 7^3}{31^2\cdot 37}.
\]
此外，定义短 Weierstrass 曲线
\[
E_{\mathrm{res}}:\quad Y^2=X^3-12096X+288144,
\]
则 \(E_{\mathrm{res}}\cong_{\QQ}E_{\mathscr R}\)；例如由变换
\[
X_{\mathrm{res}}=9X-111,\qquad Y_{\mathrm{res}}=27Y
\]
可将 \(E_{\mathscr R}:\ Y^2=X^3-37X^2+307X+361\) 送至短模型 \(E_{\mathrm{res}}:\ Y_{\mathrm{res}}^2=X_{\mathrm{res}}^3-12096X_{\mathrm{res}}+288144\)。
\end{conclusion}
\begin{proof}[证明要点]
将 \(z=1+\frac{y}{x}\) 代入 \(\mathscr R(z,y)=0\) 并清分母，出现因子 \(y\) 对应基点 \((0,1)\)；其余二次因子即给出所示关于 \(y\) 的二次方程。
令该二次方程的系数为 \(Ay^2+By+C=0\)，则恒等式 \((2Ay+B)^2=B^2-4AC\) 导出
\[
u^2=\frac{B^2-4AC}{x^2}=x\bigl(65x^3+16x^2-16x-4\bigr),
\]
从而得到 \(\mathcal Q\)。
二元四次不变量 \(I,J\) 的计算为经典公式的直接代入；由
\(
c_4=2^4I,\ c_6=2^5J
\)
可回收 \(j=c_4^3/\Delta\) 的闭式。
最后，将 \(X=\frac{X_{\mathrm{res}}+111}{9}\)、\(Y=\frac{Y_{\mathrm{res}}}{27}\) 代回 \(E_{\mathscr R}\) 即得短 Weierstrass 方程，并给出 \(E_{\mathrm{res}}\cong_{\QQ}E_{\mathscr R}\)。
可复现核验脚本：\texttt{scripts/exp\_fold\_zm\_resolvent\_cubic\_weierstrass\_audit.py}。证毕。
\end{proof}

\begin{proposition}[无复乘与几何端同态代数的极大性]\label{prop:xi-terminal-zm-resolvent-cubic-no-cm-max-end}
\leanverified{paper\_xi\_terminal\_zm\_resolvent\_cubic\_no\_cm\_max\_end}
令 \(E_{\mathscr R}\) 为结论 \ref{con:xi-terminal-zm-resolvent-cubic-weierstrass-birational} 的 Weierstrass 模型。
则 \(E_{\mathscr R}\) 不具有复乘；因此
\[
\mathrm{End}_{\overline{\QQ}}(E_{\mathscr R})=\ZZ,
\qquad
\mathrm{End}^0_{\overline{\QQ}}(E_{\mathscr R})=\QQ.
\]
从而
\[
\mathrm{End}^0_{\overline{\QQ}}(E_{\mathscr R}^{2})\ \cong\ M_2(\QQ).
\]
特别地，若 \(A\) 为任意满足 \(A\sim_{\QQ}E_{\mathscr R}^{2}\) 的二维阿贝尔簇（例如结论 \ref{con:xi-terminal-zm-s4-jacobian-a2-isog-r2} 的 \(A_2\)，或定理 \ref{thm:xi-terminal-zm-s3-closure-jacobian-splitting} 的标准表示分量），则
\(
\mathrm{End}^0_{\overline{\QQ}}(A)\cong M_2(\QQ)
\)。
\end{proposition}
\begin{proof}
若 \(E_{\mathscr R}\) 具有复乘，则其 \(j\)-不变量为代数整数；而结论 \ref{con:xi-terminal-zm-resolvent-cubic-weierstrass-birational} 给出
\[
j(E_{\mathscr R})=\frac{2^{18}\cdot 7^{3}}{31^{2}\cdot 37}\in\QQ\setminus\ZZ,
\]
其在 \(\QQ\) 中非整，因而不可能为代数整数，矛盾。
故 \(E_{\mathscr R}\) 无复乘，得到 \(\mathrm{End}^0_{\overline{\QQ}}(E_{\mathscr R})=\QQ\)。
最后由一般恒等式
\(
\mathrm{End}^0_{\overline{\QQ}}(E^{2})\cong M_2(\mathrm{End}^0_{\overline{\QQ}}(E))
\)
得到 \(\mathrm{End}^0_{\overline{\QQ}}(E_{\mathscr R}^{2})\cong M_2(\QQ)\)，并由同源不变性推广到任意 \(A\sim_{\QQ}E_{\mathscr R}^{2}\)。证毕。
\end{proof}

\begin{conclusion}[\texorpdfstring{$\pi^{(3)}_y$}{pi\_y(3)} 的零极点分解与三条有理纤维的分歧剖面]\label{con:xi-terminal-zm-resolvent-cubic-y-divisor-rational-fibers}
通过结论 \ref{con:xi-terminal-zm-resolvent-cubic-weierstrass-birational} 的同构，把投影
\(
\pi^{(3)}_y:\overline{\mathcal R}\to\mathbb P^1_y
\)
识别为 \(E_{\mathscr R}\) 上的有理函数
\[
y=-\frac{8X}{Y+19+14X-X^2}\in\QQ(E_{\mathscr R}).
\]
则其极点、零点除子在 \(E_{\mathscr R}(\overline{\QQ})\) 上满足精确分解
\[
(y)_\infty=2(15,-4)+(-1,-4),
\qquad
(y)_0=2O+(0,19),
\]
其中 \(O\) 为 \(E_{\mathscr R}\) 的无穷远点。
因此，\(\pi^{(3)}_y\) 的三条有理分歧纤维在 \(E_{\mathscr R}\) 上可具体写为
\[
(\pi^{(3)}_y)^{-1}(0)=2O+(0,19),
\qquad
(\pi^{(3)}_y)^{-1}(1)=2(23,4)+(-1,4),
\qquad
(\pi^{(3)}_y)^{-1}(\infty)=2(15,-4)+(-1,-4),
\]
从而在 \(y=0,1,\infty\) 三处均呈现“一点二次分歧、其余不分歧”的纤维型。
\end{conclusion}
\begin{proof}[证明要点]
由
\(
y=-8X/(Y+19+14X-X^2)
\)
可知零点来自 \(X=0\) 与无穷远点处的阶数比较，而极点来自分母与曲线的交点及其重数；
将 \(Y=X^2-14X-19\) 代入 \(E_{\mathscr R}\) 的方程并因式分解可得分母零点的 \(X\)-坐标分解，从而读出 \((y)_\infty\)。
对 \(y-1\) 同理得到 \(y=1\) 纤维的重数分解。
可复现核验脚本：\texttt{scripts/exp\_fold\_zm\_resolvent\_cubic\_weierstrass\_audit.py}。证毕。
\end{proof}

\begin{conclusion}[非有理分歧值的 \texorpdfstring{$X$}{X}-坐标消元证书与 \texorpdfstring{$P_{\mathrm{LY}}$}{P\_LY} 的再出现]\label{con:xi-terminal-zm-resolvent-cubic-nonrational-branch-elimination}
考虑 \(\pi^{(3)}_y\) 的非有理分歧点（即不在 \(y=0,1,\infty\) 之上者）。
则其在 \(E_{\mathscr R}\) 上的 \(X\)-坐标必满足不可约三次方程
\[
g_{\mathscr R}(X):=4X^3+3X^2+114X-269=0,
\]
且相应分歧值 \(y\) 可纯以 \(X\) 表示为
\[
y=\frac{16(X^2+19)}{(X-15)(2X^2+X+15)}.
\]
对 \(X\) 作结式消元可得显式证书
\[
\boxed{\;
\mathrm{Res}_X\!\Bigl(g_{\mathscr R}(X),\ (X-15)(2X^2+X+15)y-16(X^2+19)\Bigr)
=-68460544\,P_{\mathrm{LY}}(y)\; }.
\]
从而上述三条非有理二次分歧值的极小多项式恰为
\[
P_{\mathrm{LY}}(y)=256y^3+411y^2+165y+32.
\]
\end{conclusion}
\begin{proof}[证明要点]
对 \(y\) 的显式表达式计算 \(\dd y\) 并在 \(E_{\mathscr R}\) 上消去 \(Y\) 得到分歧点的 \(X\)-坐标多项式；
对 \(X\) 作结式消元即可回收 \(P_{\mathrm{LY}}(y)\)。
可复现核验脚本：\texttt{scripts/exp\_fold\_zm\_resolvent\_cubic\_weierstrass\_audit.py}。证毕。
\end{proof}

\begin{conclusion}[\texorpdfstring{$\QQ(\sqrt{-111})$}{Q(sqrt(-111))} 的几何实现：由 \texorpdfstring{$\mathrm{Disc}(g_{\mathscr R})$}{Disc(g\_R)} 直接读出]\label{con:xi-terminal-zm-resolvent-cubic-sqrt-111-geometric}
结论 \ref{con:xi-terminal-zm-resolvent-cubic-nonrational-branch-elimination} 的分歧控制多项式 \(g_{\mathscr R}(X)\) 的判别式满足
\[
\mathrm{Disc}(g_{\mathscr R})=-(2^3\cdot 3\cdot 31)^2\cdot 111,
\]
故其分裂域的二次分辨子域必为
\[
\QQ\!\left(\sqrt{\mathrm{Disc}(g_{\mathscr R})}\right)=\QQ(\sqrt{-111}).
\]
因此，二次域 \(\QQ(\sqrt{-111})\) 在此获得一条完全显式的几何实现：它即为 \(\pi^{(3)}_y\) 于 \(y\neq 0,1,\infty\) 的三条二次分歧值之定义域的公共二次分辨域（由 \(g_{\mathscr R}\) 的判别式直接给出），并与 \(P_{\mathrm{LY}}\) 的判别式平方自由部分精确一致。
\end{conclusion}
\begin{proof}[证明要点]
判别式分解可由对三次多项式的标准判别式公式计算并在 \(\ZZ\) 中分解得到。
由不可约三次分裂域的标准结构，二次分辨子域由 \(\sqrt{\mathrm{Disc}}\) 给出。
可复现核验脚本：\texttt{scripts/exp\_fold\_zm\_resolvent\_cubic\_weierstrass\_audit.py}。证毕。
\end{proof}

\begin{conclusion}[\texorpdfstring{$E_{\mathscr R}(\QQ)$}{E\_R(Q)} 的扭点平凡性与可计算的无穷阶点]\label{con:xi-terminal-zm-resolvent-cubic-mw-point}
\(E_{\mathscr R}(\QQ)\) 的扭点子群为平凡群
\[
E_{\mathscr R}(\QQ)_{\mathrm{tors}}=\{O\}.
\]
特别地，任一 \(\QQ\)-有理点只要不同于 \(O\) 即为无穷阶点；例如
\[
P:=(23,4)\in E_{\mathscr R}(\QQ)
\]
为无穷阶点。
并且在结论 \ref{con:xi-terminal-zm-resolvent-cubic-biquadratic-invariants-short-weierstrass} 的短 Weierstrass 模型
\[
E_{\mathrm{res}}:\quad Y^2=X^3-12096X+288144
\]
上，点
\[
P_{\mathrm{res}}:=(25,37)\in E_{\mathrm{res}}(\QQ)
\]
亦为无穷阶点，从而 \(\operatorname{rank}E_{\mathrm{res}}(\QQ)\ge 1\)（等价于 \(\operatorname{rank}E_{\mathscr R}(\QQ)\ge 1\)）。
在结论 \ref{con:xi-terminal-zm-resolvent-cubic-weierstrass-birational} 的平面模型坐标下，
\[
\bigl(y(P),z(P)\bigr)=(1,-3),
\]
从而沿 \((nP)\)（\(n\ge 1\)）取值可产生无穷多个有理参数
\[
y_n:=y(nP)\in\QQ,
\]
使得三次预解式 \(\mathscr R(z,y_n)=0\) 在 \(\QQ\) 上具有有理根 \(z_n:=z(nP)\in\QQ\)，并因此在每个 \(y_n\) 处发生线性因子化。
\end{conclusion}
\begin{proof}[证明要点]
扭点注入良约化的有限域点群，故扭点阶整除任意良约化素数 \(p\nmid \Delta\) 处的 \(\#E_{\mathscr R}(\FF_p)\)；
例如在 \(p=3,5\) 处良约化且 \(\#E_{\mathscr R}(\FF_3)=5\)、\(\#E_{\mathscr R}(\FF_5)=8\)，其最大公因子为 \(1\) 即推出扭点平凡性。
点 \(P_{\mathrm{res}}=(25,37)\in E_{\mathrm{res}}(\QQ)\) 的代入验证为直接计算；又由 \(E_{\mathrm{res}}\cong_{\QQ}E_{\mathscr R}\) 知其对应点在 \(E_{\mathscr R}(\QQ)\) 中非零，因扭点平凡而必为无穷阶点。
\((y(P),z(P))\) 的取值由结论 \ref{con:xi-terminal-zm-resolvent-cubic-weierstrass-birational} 的显式公式直接代入得到。
可复现核验脚本：\texttt{scripts/exp\_fold\_zm\_resolvent\_cubic\_weierstrass\_audit.py}。证毕。
\end{proof}

\paragraph{二面体特化机制：由预解根到单二次域闭合}
在结论 \ref{con:xi-terminal-zm-s4-resolvent-cubic-discriminant-equals-delta} 与结论
\ref{con:xi-terminal-zm-resolvent-cubic-weierstrass-birational} 的同一记号体系下，
平面三次预解曲线 \(\mathscr R(z,y)=0\) 的 \(\QQ\)-有理点提供了一个可计算的特化机制：当 \(z\) 在基域中出现时，四次谱覆盖的伽罗瓦群从泛型的 \(S_4\) 下降到配对稳定子，且 Ferrari 分解所需的两处平方根可在同一二次扩张中闭合。

\begin{theorem}[函数域上的二面体特化：预解根强制 \texorpdfstring{$D_4$}{D4} 伽罗瓦群]\label{thm:xi-terminal-zm-dihedral-d4-specialization-functionfield}
令 \(K:=\QQ(E_{\mathscr R})\) 为椭圆曲线 \(E_{\mathscr R}\) 的函数域（结论 \ref{con:xi-terminal-zm-resolvent-cubic-weierstrass-birational}），并在 \(K\) 中取
\[
y=-\frac{8X}{Y+19+14X-X^2},
\qquad
z=1-\frac{2X(7-X)}{Y+19+14X-X^2}.
\]
则 \((y,z)\) 满足 \(\mathscr R(z,y)=0\)。
将四次谱多项式
\[
\Pi(\lambda,y)=\lambda^{4}-\lambda^{3}-(2y+1)\lambda^{2}+\lambda+y(y+1)
\in K[\lambda]
\]
视为 \(K\) 上的可分四次多项式，则其伽罗瓦群满足
\[
\mathrm{Gal}\bigl(\Pi(\lambda,y)/K\bigr)\ \subseteq\ D_4.
\]
此外，存在 \(P\in E_{\mathscr R}(\QQ)\) 使得特化多项式 \(\Pi(\lambda,y(P))\in\QQ[\lambda]\) 的伽罗瓦群同构于 \(D_4\)；
从而
\[
\mathrm{Gal}\bigl(\Pi(\lambda,y)/K\bigr)\ \cong\ D_4.
\]
因此存在无穷多个 \(P\in E_{\mathscr R}(\QQ)\) 使得 \(\Pi(\lambda,y(P))\) 的伽罗瓦群同构于 \(D_4\)。
\end{theorem}

\begin{proof}[证明要点]
由 \(\mathscr R(z,y)=0\) 可知，\(\Pi(\lambda,y)\) 的 Ferrari 预解三次在 \(K\) 上具有根 \(z\)。
四次多项式的经典判别指出：若预解三次在基域上有根，则伽罗瓦群稳定某一根配对，因而包含于配对稳定子 \(D_4\)。
进一步，取一个使特化伽罗瓦群为 \(D_4\) 的有理点 \(P\)，即可推出函数域伽罗瓦群必须达到 \(D_4\)。
最后由曲线覆盖的 Hilbert 不可约机制得到无穷多个保持该伽罗瓦群的有理点特化。证毕。
\end{proof}

\begin{proposition}[预解曲线上的二判别式锁定恒等式]\label{prop:xi-terminal-zm-dihedral-double-discriminant-lock}
\leanverified{paper\_xi\_terminal\_zm\_dihedral\_double\_discriminant\_lock}
在仿射曲线 \(\mathscr R(z,y)=0\) 上恒成立
\[
\boxed{\;
\bigl(z^{2}-4y(y+1)\bigr)\,\bigl(8y+5+4z\bigr)=(z+2)^{2}\; }.
\]
等价地，在多项式环 \(\ZZ[y,z]\) 中有恒等式
\[
\bigl(z^{2}-4y(y+1)\bigr)\,\bigl(8y+5+4z\bigr)-(z+2)^{2}=4\,\mathscr R(z,y).
\]
\end{proposition}

\begin{proof}
将左端展开并与 \((z+2)^2\) 相减，直接整理可得差为 \(4\mathscr R(z,y)\)。
因此在理想 \((\mathscr R)\) 上为零，等价于在曲线 \(\mathscr R=0\) 上恒成立。证毕。
\end{proof}

\begin{corollary}[Ferrari 分解的单二次域闭合与显式二次因子]\label{cor:xi-terminal-zm-dihedral-ferrari-single-quadratic-closure}
\leanverified{paper\_xi\_terminal\_zm\_dihedral\_ferrari\_single\_quadratic\_closure}
设 \(F\) 为特征不为 \(2\) 的域。
取 \(y,z\in F\) 满足 \(\mathscr R(z,y)=0\)，并在某个二次扩张中取
\[
s^2=z^2-4y(y+1).
\]
则在 \(F(s)[\lambda]\) 中存在显式因子分解
\[
\boxed{\;
\Pi(\lambda,y)=Q_{+}(\lambda)\,Q_{-}(\lambda)\;}
\]
其中
\[
Q_{\pm}(\lambda)=
\lambda^2-\frac12\left(1\pm\frac{z+2}{s}\right)\lambda+\frac12\bigl(z\pm s\bigr).
\]
并且由命题 \ref{prop:xi-terminal-zm-dihedral-double-discriminant-lock}，有
\[
\left(\frac{z+2}{s}\right)^2=8y+5+4z\in F,
\]
从而 Ferrari 过程中出现的两处平方根属于同一二次扩张 \(F(s)/F\)。
\end{corollary}

\begin{proof}[证明要点]
令 \(\lambda_1,\dots,\lambda_4\) 为 \(\Pi(\lambda,y)\) 的四根。
当 \(z\) 为预解三次的一根时，对应某一配对使得
\(
\lambda_1\lambda_2+\lambda_3\lambda_4=z
\)
且
\(
\lambda_1\lambda_2\lambda_3\lambda_4=y(y+1)
\)。
于是 \(p_{\pm}:=\lambda_1\lambda_2,\lambda_3\lambda_4\) 为二次方程 \(t^2-zt+y(y+1)=0\) 的两根，即 \(p_{\pm}=\frac12(z\pm s)\)。
另一方面，利用 \(\lambda_1+\cdots+\lambda_4=1\) 与二次对称函数关系可得两组和 \(\lambda_1+\lambda_2,\lambda_3+\lambda_4\) 的判别式为 \(8y+5+4z\)；
命题 \ref{prop:xi-terminal-zm-dihedral-double-discriminant-lock} 表明该判别式与 \(s^2\) 属于同一平方类，故同属 \(F(s)\)。
据此写出两二次因子并整理即得所示 \(Q_\pm\)。证毕。
\end{proof}

\begin{proposition}[\texorpdfstring{$D_4$}{D4} 特化时的二次子域生成元可计算性]\label{prop:xi-terminal-zm-dihedral-d4-quadratic-subfields}
\leanverified{paper\_xi\_terminal\_zm\_dihedral\_d4\_quadratic\_subfields}
设 \(y_0\in\QQ\) 使存在 \(z_0\in\QQ\) 满足 \(\mathscr R(z_0,y_0)=0\)，并且 \(\Pi(\lambda,y_0)\) 在 \(\QQ[\lambda]\) 上不可约且判别式非平方。
记其分裂域为 \(L/\QQ\)。
则 \(\mathrm{Gal}(L/\QQ)\cong D_4\)，并且 \(L\) 的二次子域可由两条平方根生成：
\[
K_1=\QQ\!\left(\sqrt{z_0^2-4y_0(y_0+1)}\right),
\qquad
K_2=\QQ\!\left(\sqrt{\mathrm{Disc}_\lambda\Pi(\lambda,y_0)}\right),
\]
第三个二次子域为 \(K_1K_2\) 的另一二次子域。
\end{proposition}

\begin{proof}[证明要点]
由 \(z_0\) 为预解根得 \(\mathrm{Gal}(L/\QQ)\subseteq D_4\)。
不可约性给出传递性，判别式非平方排除 \(V_4\)。
另一方面，推论 \ref{cor:xi-terminal-zm-dihedral-ferrari-single-quadratic-closure} 给出 \(K_1\subset L\)，而判别式二次分辨子域 \(K_2\subset L\) 恒成立。
若 \(\mathrm{Gal}\cong C_4\)，则分裂域仅含唯一二次子域，与 \(K_1\neq K_2\) 矛盾；故只能为 \(D_4\)。
二次子域格结构由 \(D_4\) 的指数 \(2\) 子群分类推出。证毕。
\end{proof}

\begin{conjecture}[沿无穷阶轨道的无例外二面体性]\label{conj:xi-terminal-zm-dihedral-d4-no-exception-orbit}
取结论 \ref{con:xi-terminal-zm-resolvent-cubic-mw-point} 的无穷阶点 \(P=(23,4)\in E_{\mathscr R}(\QQ)\)，并记 \(y_n:=y(nP)\in\QQ\)。
猜想对所有 \(n\ge 2\)，四次多项式 \(\Pi(\lambda,y_n)\) 在 \(\QQ\) 上不可约且
\[
\mathrm{Gal}\bigl(\Pi(\lambda,y_n)/\QQ\bigr)\cong D_4.
\]
\end{conjecture}


\begin{conclusion}[\texorpdfstring{$31$}{31} 与 \texorpdfstring{$37$}{37} 处的约化型与分裂性判据]\label{con:xi-terminal-zm-resolvent-cubic-reduction-31-37}
对结论 \ref{con:xi-terminal-zm-resolvent-cubic-weierstrass-birational} 的 \(E_{\mathscr R}\)，由
\(
v_{31}(c_4)=v_{37}(c_4)=0
\)
且
\[
v_{31}(\Delta)=2,\qquad v_{37}(\Delta)=1,
\]
可知其在 \(p=31,37\) 处均为乘性约化。
进一步由 \(-c_6\) 的二次特征判定分裂性：
\[
\left(\frac{-c_6}{31}\right)=+1
\quad\Longrightarrow\quad
\text{\(31\) 处分裂乘性约化（Kodaira 型 \(I_2\)）},
\]
\[
\left(\frac{-c_6}{37}\right)=-1
\quad\Longrightarrow\quad
\text{\(37\) 处非分裂乘性约化（Kodaira 型 \(I_1\)）}.
\]
因此，素数 \(31\) 与 \(37\) 的局部约化行为在 \(E_{\mathscr R}\) 的 Weierstrass 模型上被完全显式地控制，并与三次预解式所诱导的分歧结构在判别式层面严格对齐。
\end{conclusion}
\begin{proof}[证明要点]
由 \(v_p(c_4)=0\) 且 \(v_p(\Delta)>0\) 的 Tate 判据知为乘性约化；
Kodaira 型由 \(v_p(\Delta)\) 决定，而分裂性由 \(\left(\frac{-c_6}{p}\right)\) 判别。
可复现核验脚本：\texttt{scripts/exp\_fold\_zm\_resolvent\_cubic\_weierstrass\_audit.py}。证毕。
\end{proof}

\begin{conclusion}[三次预解式在坏素数 \texorpdfstring{$37$}{37} 处的约化奇点：唯一非分裂结点与有理归一化]\label{con:xi-terminal-zm-resolvent-cubic-badprime-37-nonsplit-node}
令 \(\overline{\mathcal R}\subset \PP^{2}_{[Z:Y:W]}\) 为结论 \ref{con:xi-terminal-zm-resolvent-cubic-plane-cubic-geometry} 的平面三次曲线，
并记其在 \(\FF_{37}\) 上的约化为 \(\overline{\mathcal R}_{37}\)。
则 \(\overline{\mathcal R}_{37}\) 的射影闭包仅有一个奇点，且在仿射图 \(W=1\) 上可取为
\[
(z,y)\equiv (29,6)\pmod{37}.
\]
该奇点为非分裂结点；因此 \(\overline{\mathcal R}_{37}\) 的归一化为 \(\PP^{1}_{\FF_{37}}\)。
并且该退化由判别式核 \(P_{\mathrm{LY}}(y)\) 在 \(y\equiv 6\pmod{37}\) 的二重消失所触发，
从而与判别式二次脊曲线在该素处的结点机制保持一致（参见结论 \ref{con:xi-terminal-zm-elliptic-lattes-badprime-37-degree-drop}）。
\end{conclusion}
\begin{proof}[证明要点]
在 \(\FF_{37}\) 上联立
\(
\mathscr R(z,y)=\partial_{z}\mathscr R(z,y)=\partial_{y}\mathscr R(z,y)=0
\)
得到唯一解 \((z,y)=(29,6)\)，从而约化曲线仅有该一处奇点。
将坐标平移到该点附近，最低齐次项为二次型，其判别式为 \(2\in\FF_{37}^{\times}\) 的非平方元，
故切锥在 \(\FF_{37}\) 上不分裂而在二次扩张上分裂，给出非分裂结点。
平面三次曲线在出现结点时的归一化为 \(\PP^{1}\) 为标准结论。
可复现核验脚本：\texttt{scripts/exp\_fold\_zm\_resolvent\_cubic\_weierstrass\_audit.py}。证毕。
\end{proof}

\begin{conclusion}[Lee--Yang 三次因子的几何起源：重量映射 \texorpdfstring{$\pi_y$}{pi\_y} 的临界值消元与切触包络]\label{con:xi-terminal-zm-leyang-ramification-critical-values}
在结论 \ref{con:xi-terminal-zm-elliptic-normalization} 的椭圆模型
\[
E:\ Y^2=X^3-X+\frac14
\]
上，沿用结论 \ref{con:xi-terminal-zm-y-quadratic-extension-minpoly} 的物理主支
\[
y:=X^2+Y-\frac12\in\QQ(E),
\qquad
\omega:=\frac{\dd X}{2Y}.
\]
则 \(y\) 诱导次数 \(4\) 的有限态射 \(\pi_y:E\to\mathbb P^1_y\)，并满足 \((y)_\infty=4O\)（见结论 \ref{con:xi-terminal-zm-elliptic-y-hurwitz-passport}）。
则有严格微分恒等式
\[
\boxed{\;
\dd y=\bigl(4XY+3X^2-1\bigr)\,\omega
=\frac{4X\eta+3X^2+2X-1}{2\eta+1}\,\dd X\qquad(\eta:=Y-\tfrac12)\; }.
\]
并在仿射开集 \(Y\neq 0\) 上等价写为 \(\dd y=\frac{4XY+3X^2-1}{2Y}\,\dd X\)。
因此映射 \(\pi_y:E\to\mathbb P^1_y\)、\(P\mapsto y(P)\) 的有限分歧点由显式代数曲线
\[
4XY+3X^2-1=0
\]
在 \(E\) 上的交点给出；消去 \(Y\) 后得到其 \(X\)-坐标满足完全分解
\[
\boxed{\;
(X-1)(X+1)\bigl(16X^3-9X^2+1\bigr)=0\; }.
\]
等价地，其五次展开式为
\[
16X^{5}-9X^{4}-16X^{3}+10X^{2}-1=0.
\]
更精确地：
\begin{enumerate}
  \item 当 \(X=1\) 时，唯一有限分歧点为 \(\bigl(1,-\tfrac12\bigr)\)，其分歧值为 \(y=0\)；
  \item 当 \(X=-1\) 时，唯一有限分歧点为 \(\bigl(-1,\tfrac12\bigr)\)，其分歧值为 \(y=1\)；
  \item 若 \(\alpha\) 为不可约三次方程 \(16X^{3}-9X^{2}+1=0\) 的任一根，取
  \[
  Y=\frac{1-3\alpha^{2}}{4\alpha},
  \qquad
  P_{\alpha}:=(\alpha,Y)\in E(\overline{\QQ}),
  \]
  则 \(P_{\alpha}\) 为 \(\pi_y\) 的简单分歧点，且其分歧值为
  \[
  y_{\alpha}
  =\alpha^{2}+\frac{1-3\alpha^{2}}{4\alpha}-\frac12
  =\frac{(\alpha-1)\bigl(4\alpha^{2}+\alpha-1\bigr)}{4\alpha}.
  \]
\end{enumerate}
从而 \(\pi_y\) 的有限分歧值集合恰为 \(y=0\)、\(y=1\) 与 \(\{y_{\alpha}\}\)。
其中 Lee--Yang 三次多项式
\[
P_{\mathrm{LY}}(y):=256y^3+411y^2+165y+32
\]
可由纯消元证书刻画为严格结果式恒等式
\[
\boxed{\;
\mathrm{Res}_X\!\left(16X^3-9X^2+1,\ 4Xy-\bigl(4X^3-3X^2-2X+1\bigr)\right)
=-64\,P_{\mathrm{LY}}(y)\; }.
\]
等价地，交换消元顺序（此处两式关于 \(X\) 的次数皆为 \(3\)）得到
\[
\boxed{\;
P_{\mathrm{LY}}(y)=\frac1{64}\,
\mathrm{Res}_X\!\left(4Xy-\bigl(4X^3-3X^2-2X+1\bigr),\ 16X^3-9X^2+1\right)\; }.
\]
因此上述三点给出的三个分歧值 \(\{y_{\alpha}\}\) 恰为 \(P_{\mathrm{LY}}(y)=0\) 的三根。
因而 Lee--Yang 三次因子不仅是谱四次的判别式核（结论 \ref{con:xi-terminal-zm-discriminant-leyang}），亦等价给出 \(\pi_y\) 的三条非有理临界值在 \(\QQ[y]\) 中的极小多项式（结论 \ref{con:xi-terminal-zm-leyang-cubic-arithmetic}）；
在几何上，这等价于一族抛物线
\[
Y=y-X^2+\frac12
\]
与 \(E\) 的切触参数在消元意义下的像。
\end{conclusion}
\begin{proof}[证明要点]
由曲线方程微分关系 \(2Y\,\dd Y=(3X^2-1)\dd X\) 得 \(\dd Y=(3X^2-1)\omega\)，从而
\[
\dd y=2X\dd X+\dd Y=\bigl(4XY+3X^2-1\bigr)\omega.
\]
因此有限分歧点满足 \(4XY+3X^2-1=0\)。
由此解出
\(
Y=(1-3X^2)/(4X)
\)
并代入 \(E\) 的方程，化简即得
\(
(X-1)(X+1)(16X^3-9X^2+1)=0
\)。
当 \(X=1\) 与 \(X=-1\) 时分别得到 \(y=0\) 与 \(y=1\)。
对满足 \(16X^3-9X^2+1=0\) 的非平凡分歧点，将 \(y=X^2+Y-\frac12\) 与 \(Y=(1-3X^2)/(4X)\) 联立得到线性约束
\[
4Xy=4X^3-3X^2-2X+1.
\]
对变量 \(X\) 作消元（结果式）得到严格恒等式
\(
\mathrm{Res}_X(16X^3-9X^2+1,\ 4Xy-(4X^3-3X^2-2X+1))=-64\,P_{\mathrm{LY}}(y)
\)，
从而非平凡有限分歧值由 \(P_{\mathrm{LY}}(y)=0\) 精确刻画，并与结论 \ref{con:xi-terminal-zm-discriminant-leyang} 的判别式分解一致。
可复现核验脚本：\texttt{scripts/exp\_fold\_zm\_elliptic\_leyang\_cover\_geometry\_audit.py}。证毕。
\end{proof}

\begin{conclusion}[权平面四次模型的结式分裂：Lee--Yang 三次因子与共轭纤维三次因子的对偶消元]\label{con:xi-terminal-zm-leyang-conjugate-fiber-cubic}
沿用结论 \ref{con:xi-terminal-zm-leyang-ramification-critical-values} 的记号，并由
\(
Y=y-X^{2}+\tfrac12
\)
消去 \(Y\)。
则由下式所定义之首一四次多项式 \(F(X,y)\) 在 \(\QQ(E)\) 中满足 \(F(X,y)=0\)，并且由 \(Y=y-X^{2}+\tfrac12\) 可逆重建 \(Y\)，从而 \(\QQ(E)=\QQ(y)(X)\) 且 \([\QQ(E):\QQ(y)]=4\)；等价地，\(E\) 的函数域具有首一四次表示
\[
\QQ(E)\cong \QQ(X,y)/(F),
\qquad
F(X,y):=X^{4}-X^{3}-(2y+1)X^{2}+X+y(y+1).
\]
该表示与结论 \ref{con:xi-terminal-zm-elliptic-normalization} 的规范化同构一致，并使谱变量在规范化层面可识别为椭圆坐标 \(X\)。
在该平面四次模型下，\(\pi_y:E\to\mathbb P^{1}_{y}\) 的有限分歧点等价于方程组
\[
F(X,y)=0,\qquad \partial_{X}F(X,y)=0,
\]
并且对 \(y\) 消元得到严格结果式恒等式
\[
\boxed{\;
\mathrm{Res}_{y}\bigl(F(X,y),\partial_{X}F(X,y)\bigr)=-(X-1)(X+1)\bigl(16X^{3}-9X^{2}+1\bigr)\; }.
\]
\par
令
\(
g(X):=16X^{3}-9X^{2}+1
\)，并定义与 \(P_{\mathrm{LY}}\) 共轭的三次因子
\[
Q_{\mathrm{LY}}(y):=256y^{3}+195y^{2}+93y-48\in\ZZ[y].
\]
则对变量 \(X\) 消元给出另一条严格结式分裂
\[
\boxed{\;
\mathrm{Res}_{X}\bigl(F(X,y),g(X)\bigr)=P_{\mathrm{LY}}(y)\,Q_{\mathrm{LY}}(y)\; }.
\]
因此，对任意满足 \(g(X)=0\) 的根 \(\alpha\)：
在分歧约束 \(\partial_XF=0\) 下的重量值为
\[
y=\frac{4\alpha^{3}-3\alpha^{2}-2\alpha+1}{4\alpha},
\qquad
P_{\mathrm{LY}}(y)=0,
\]
而其在椭圆反演 \(\iota:(X,Y)\mapsto (X,-Y)\) 下的共轭点给出第二张（非分歧）纤维，
其重量值满足 \(Q_{\mathrm{LY}}(y)=0\)；
两组三次根由 \(\iota\) 交换。
从而 Lee--Yang 因子 \(P_{\mathrm{LY}}(y)\) 被识别为 \(\pi_y\) 的非平凡分歧值集合之不可约消元核，而 \(Q_{\mathrm{LY}}(y)\) 刻画同一 \(g(X)=0\) 之上的共轭非分歧纤维。
\end{conclusion}
\begin{proof}[证明要点]
由 \(Y=y-X^{2}+\tfrac12\) 代回 \(E\) 的方程并展开，即得 \(F(X,y)=0\) 的首一四次表示。
在平面模型中，有限分歧点满足 \(F=\partial_XF=0\)；
对 \(y\) 作结果式消元并在 \(\ZZ[X]\) 中因式分解，即得 \(\mathrm{Res}_{y}(F,\partial_XF)\) 的分裂恒等式。
\par
对 \(g(X)=0\) 的三根，方程 \(F(X,y)=0\) 给出对应的两张重量纤维（由 \(Y\mapsto -Y\) 交换）；
对 \(X\) 作结果式消元得到 \(y\) 的一元消元多项式，直接分解为两条三次因子。
其中分歧点的重量值由结论 \ref{con:xi-terminal-zm-leyang-ramification-critical-values} 的线性约束
\(
4Xy=4X^{3}-3X^{2}-2X+1
\)
确定，故其对应因子必为 \(P_{\mathrm{LY}}\)，另一因子即定义为 \(Q_{\mathrm{LY}}\) 并对应共轭非分歧纤维。
可复现核验脚本：\texttt{scripts/exp\_fold\_zm\_elliptic\_leyang\_cover\_geometry\_audit.py}。证毕。
\end{proof}

\begin{conclusion}[四次重量覆叠的分歧全刻画与三点消元恒等式]\label{con:xi-terminal-zm-elliptic-y-hurwitz-passport}
设
\[
E:\quad Y^2=X^3-X+\frac14,\qquad
y:=X^2+Y-\frac12\in\QQ(E),
\]
并记
\[
P_{\mathrm{LY}}(y):=256y^3+411y^2+165y+32.
\]
则有：
\begin{enumerate}
  \item \(y:E\to\PP^1_y\) 为次数 \(4\) 覆叠，且
  \[
  y^{-1}(\infty)=\{O\},\qquad e_O(y)=4.
  \]
  换言之，\(\infty\) 为全分歧点。

  \item 有限分歧值恰为
  \[
  y=0,\qquad y=1,\qquad P_{\mathrm{LY}}(y)=0\ \text{的三根},
  \]
  并且每个有限分歧值的分歧指数均为 \(2\)。

  \item 有限分歧点的 \(X\)-坐标恰为
  \[
  X=\pm1,\qquad 16X^3-9X^2+1=0\ \text{的三根}.
  \]
  对后者三根，分歧值可写为
  \[
  y=\frac{4X^3-3X^2-2X+1}{4X},
  \]
  且成立消元恒等式
  \[
  \operatorname{Res}_X\!\Bigl(16X^3-9X^2+1,\ 4Xy-(4X^3-3X^2-2X+1)\Bigr)
  =-64\,P_{\mathrm{LY}}(y).
  \]
\end{enumerate}
\end{conclusion}
\begin{proof}[证明要点]
在无穷远点 \(O\) 处，\(X\) 与 \(Y\) 的极点阶分别为 \(2,3\)，故 \(y=X^2+Y-\tfrac12\) 在 \(O\) 处极点阶为 \(4\)，并且 \(E\) 上无其他极点；于是 \(\deg(y)=4\) 且 \(e_O=4\)。

由隐函数求导 \(2Y\,\mathrm dY=(3X^2-1)\,\mathrm dX\)，得
\[
\frac{\mathrm dy}{\mathrm dX}
=2X+\frac{3X^2-1}{2Y}
=\frac{4XY+3X^2-1}{2Y}.
\]
有限分歧点满足 \(4XY+3X^2-1=0\)，即
\[
Y=\frac{1-3X^2}{4X}.
\]
代回 \(E\) 的方程得到
\[
16X^5-9X^4-16X^3+10X^2-1=0
=(X^2-1)(16X^3-9X^2+1),
\]
从而有限分歧点的 \(X\)-坐标正是 \(X=\pm1\) 与三次式 \(16X^3-9X^2+1\) 的三根。

在临界关系 \(4XY+3X^2-1=0\) 上，直接代入 \(y=X^2+Y-\tfrac12\) 得
\[
y=\frac{4X^3-3X^2-2X+1}{4X}.
\]
据此对 \(X\) 作结果式消元即可得到
\[
\operatorname{Res}_X\!\Bigl(16X^3-9X^2+1,\ 4Xy-(4X^3-3X^2-2X+1)\Bigr)
=-64\,P_{\mathrm{LY}}(y).
\]
并且代入 \(X=1,-1\) 分别得到 \(y=0,1\)。

最后由 Riemann--Hurwitz 公式（\(g(E)=1,\deg(y)=4\)）知总分歧度为 \(8\)。
无穷远点贡献 \(e_O-1=3\)，其余有限分歧总贡献为 \(5\)，故五个有限分歧均为单分歧（分歧指数 \(2\)）。证毕。
\end{proof}


\begin{conclusion}[分歧除子的主化与群律迹：Lee--Yang 三分歧点的加法迹坍缩为显式有理点]\label{con:xi-terminal-zm-leyang-ramification-divisor-principalization-trace}
沿用结论 \ref{con:xi-terminal-zm-elliptic-y-hurwitz-passport} 的记号，并令
\[
f:=4XY+3X^2-1\in\QQ(E)^\times.
\]
设 \(O\) 为无穷远点，记两条有理分歧点
\[
Q_0:=\Bigl(1,-\frac12\Bigr),\qquad Q_0':=\Bigl(-1,\frac12\Bigr),
\]
并令 \(Q_1,Q_2,Q_3\in E(\overline{\QQ})\) 为满足 \(16X^3-9X^2+1=0\) 的三条非有理分歧点（其 \(Y\)-坐标由 \(4XY+3X^2-1=0\) 唯一确定）。
则有严格主除子恒等式
\[
\boxed{\;
\mathrm{div}(f)=Q_0+Q_0'+Q_1+Q_2+Q_3-5O\; }.
\]
从而在椭圆曲线群律意义下成立
\[
Q_1+Q_2+Q_3=-(Q_0+Q_0')=\Bigl(\frac14,-\frac18\Bigr)\in E(\QQ),
\]
并且
\[
Q_0+Q_0'=\Bigl(\frac14,\frac18\Bigr).
\]
因此，位于 Lee--Yang 三个分歧值之上的三条分歧点（其定义域为 \(\QQ(\alpha)\)，\(\alpha\) 为 \(16X^3-9X^2+1=0\) 的根）在群律意义下满足一条全局加法约束：其加法迹严格落入固定的有理点 \((\tfrac14,-\tfrac18)\)。
结合结论 \ref{con:xi-terminal-zm-elliptic-mw-generator-alignment} 的自由生成元 \(R=(0,\tfrac12)\) 及 \(Q_0=-[2]R\)、\(Q_0'=-[3]R\)，可将上述迹恒等式改写为
\[
Q_1+Q_2+Q_3=5R.
\]
该恒等式亦可视为 Riemann--Hurwitz 所给出的线性等价关系在 \(\mathrm{Pic}^0(E)\cong E\) 下的具象化。
由于结论 \ref{con:xi-terminal-zm-elliptic-rational-points-denominator-law} 已给出 \(E(\QQ)_{\mathrm{tors}}=\{O\}\)，且 \((\tfrac14,-\tfrac18)\neq O\)，故该有理迹点为无限阶。
\end{conclusion}
\begin{proof}[证明要点]
由结论 \ref{con:xi-terminal-zm-leyang-ramification-critical-values} 的微分恒等式 \(\dd y=f\,\omega\) 可知：
\(f\) 的有限零点恰为 \(\pi_y\) 的有限分歧点，且由结论 \ref{con:xi-terminal-zm-elliptic-y-hurwitz-passport} 的 Hurwitz 护照可知其均为一阶零。
在无穷远点 \(O\) 处取均匀化参数 \(t\) 使 \(X\sim t^{-2}\)、\(Y\sim t^{-3}\)，则
\[
f=4XY+3X^2-1\sim 4t^{-5},
\]
故 \(f\) 在 \(O\) 处有五阶极点且无其他极点，从而得到所示主除子恒等式。
由于 \(\mathrm{div}(f)\) 在 \(\mathrm{Pic}^{0}(E)\cong E\) 中为零类，推出 \(Q_0+Q_0'+Q_1+Q_2+Q_3=O\)。
最后对有理点 \(Q_0,Q_0'\) 施用弦割加法公式：斜率
\[
m=\frac{\frac12-(-\frac12)}{-1-1}=-\frac12
\]
给出
\[
x(Q_0+Q_0')=m^2-1-(-1)=\frac14,\qquad
y(Q_0+Q_0')=m\Bigl(1-\frac14\Bigr)-\Bigl(-\frac12\Bigr)=\frac18,
\]
从而 \(Q_0+Q_0'=(\tfrac14,\tfrac18)\) 且 \(-(Q_0+Q_0')=(\tfrac14,-\tfrac18)\)，并得到结论。
可复现核验脚本：\texttt{scripts/exp\_fold\_zm\_elliptic\_leyang\_cover\_geometry\_audit.py}。证毕。
\end{proof}

\begin{conclusion}[Lee--Yang 判别式的范数恒等式：\texorpdfstring{$-y(y-1)P_{\mathrm{LY}}(y)$}{-y(y-1)P(y)} 为不同生成元的相对范数]\label{con:xi-terminal-zm-leyang-discriminant-norm-identity}
设
\[
k:=\QQ(y),\qquad K:=\QQ(E)=\QQ(X,Y),
\]
则 \(K/k\) 为次数 \(4\) 的函数域扩张。
该扩张可分，且其不同理想在 \(K\) 中为主理想。
令
\[
\delta:=4Xy-\bigl(4X^3-3X^2-2X+1\bigr)\in K.
\]
在 \(E\) 上利用 \(Y=y-X^2+\tfrac12\) 可化为
\[
\delta=4XY+3X^2-1,
\]
并且 \(\delta\) 在 \(\pi_y\) 的有限分歧点处恰为一阶零。
因此
\[
\mathfrak D_{K/k}=(\delta)=\bigl(1-3X^2-4XY\bigr).
\]
则其相对范数满足严格恒等式
\[
\boxed{\;
\operatorname{N}_{K/k}(\delta)=-y(y-1)\,P_{\mathrm{LY}}(y)\; }.
\]
\end{conclusion}
\begin{proof}[证明要点]
由 \(Y=y-X^2+\tfrac12\) 消去 \(Y\) 得 \(X\) 在 \(k\) 上满足的四次方程
\[
F(X,y)=X^4-X^3-(2y+1)X^2+X+y(y+1)=0,
\]
并且直接计算给出 \(\delta=-\partial_XF(X,y)\)。
由于 \(F\) 为关于 \(X\) 的首一可分多项式，标准理论给出不同理想
\(
\mathfrak D_{K/k}=\bigl(\partial_XF(X,y)\bigr)
\)；
从而 \(\mathfrak D_{K/k}=(\delta)\)。
再将 \(y=X^2+Y-\tfrac12\) 代入 \(\partial_XF(X,y)\)，即得
\(
\partial_XF(X,y)=1-3X^2-4XY
\)，故 \((\delta)=\bigl(1-3X^2-4XY\bigr)\)。
由于 \(F\) 关于 \(X\) 首一，范数 \(\operatorname{N}_{K/k}(\delta)\) 等于结果式 \(\mathrm{Res}_X(F,\delta)\)（等价于 \(\mathrm{Disc}_X(F)\) 的标准表示）。
而该结果式恰计算为 \(-y(y-1)P_{\mathrm{LY}}(y)\)，与结论 \ref{con:xi-terminal-zm-discriminant-leyang} 的判别式分解一致。
可复现核验脚本：\texttt{scripts/exp\_fold\_zm\_elliptic\_leyang\_cover\_geometry\_audit.py}。证毕。
\end{proof}

\begin{conclusion}[\texorpdfstring{$2$}{2}-division 多项式在 \texorpdfstring{$y$}{y}-投影下的范数平方化：隐藏三次控制多项式]\label{con:xi-terminal-zm-2division-norm-square}
沿用结论 \ref{con:xi-terminal-zm-leyang-discriminant-norm-identity} 的记号，仍令 \(k=\QQ(y)\)、\(K=\QQ(E)\)。
记 \(2\)-division 三次
\[
D(X):=4X^3-4X+1=4Y^2\in K.
\]
并记其在 \(y\)-坐标下的像控制多项式
\[
N_2(y):=16y^3-8y^2-4y+1\qquad(\text{亦记 }S(y)).
\]
则其在扩张 \(K/k\) 下的相对范数满足严格恒等式
\[
\boxed{\;
\operatorname{N}_{K/k}\bigl(D(X)\bigr)=N_2(y)^2\; }.
\]
等价地，由 \(2Y=2y-2X^2+1\)（见结论 \ref{con:xi-terminal-zm-y-quadratic-extension-minpoly}）得到
\[
\boxed{\;
\operatorname{N}_{K/k}(2Y)=-N_2(y)\; }.
\]
因此三次多项式 \(N_2(y)\) 精确刻画 \(E[2]\setminus\{O\}\) 在 \(\pi_y\) 下的像；并且由于 \(D(X)=(2Y)^2\) 在 \(K\) 中为平方，其相对范数在 \(k\) 中必为完全平方，上式给出该平方化的最小整系数证书。
\par
等价地，由 \(E[2]\setminus\{O\}\) 的特征方程 \(Y=0\) 与 \(D(X)=0\) 得 \(y=X^2-\tfrac12\)，从而存在严格消元证书
\[
\boxed{\;
\mathrm{Res}_X\!\left(4X^3-4X+1,\ y-\Bigl(X^2-\frac12\Bigr)\right)=N_2(y)\; }.
\]
\end{conclusion}
\begin{proof}[证明要点]
仍以四次方程 \(F(X,y)=0\)（见结论 \ref{con:xi-terminal-zm-leyang-discriminant-norm-identity}）作为 \(K/k\) 的首一表示，则
\(
\operatorname{N}_{K/k}(D(X))=\mathrm{Res}_X(F,D)
\)
与
\(
\operatorname{N}_{K/k}(2Y)=\mathrm{Res}_X(F,2y-2X^2+1)
\)
为标准结果式计算；
直接化简即得到所示恒等式。
可复现核验脚本：\texttt{scripts/exp\_fold\_zm\_elliptic\_leyang\_cover\_geometry\_audit.py}。证毕。
\end{proof}

\begin{conclusion}[Lee--Yang 三次域与临界三次域同构：显式生成元对应与判别式的 \texorpdfstring{$37$}{37} 不可消性]\label{con:xi-terminal-zm-leyang-cubic-field-isomorphism}
令 \(\alpha\) 为三次方程
\[
16X^3-9X^2+1=0
\]
的任一根，并定义
\[
\beta:=\frac{4\alpha^3-3\alpha^2-2\alpha+1}{4\alpha}
=\alpha^{2}-\frac34\alpha-\frac12+\frac{1}{4\alpha}
=-3\alpha^2+\frac32\alpha-\frac12.
\]
则：
\begin{enumerate}
  \item \(\beta\) 满足 Lee--Yang 三次方程 \(P_{\mathrm{LY}}(\beta)=0\)。
  \item 因此诱导出立方域同一性 \(\QQ(\alpha)=\QQ(\beta)\)。
  反向生成元可取为如下两支之一（对应该三次域的两种 \(\QQ\)-嵌入选择）：
  \[
  \alpha=\frac{-512\beta^2-518\beta-5}{279}\quad\text{或}\quad \alpha=\frac{512\beta^2-970\beta-739}{279}.
  \]
  \item 两个表示的多项式判别式满足
  \[
  \boxed{\ \mathrm{Disc}(16X^3-9X^2+1)=-2^2\cdot 3^3\cdot 37\ },
  \qquad
  \boxed{\ \mathrm{Disc}(P_{\mathrm{LY}})=-3^9\cdot 31^2\cdot 37\ }.
  \]
  并且二者严格相差一个平方因子：
  \[
  \boxed{\ \mathrm{Disc}(P_{\mathrm{LY}})=\mathrm{Disc}(16X^3-9X^2+1)\cdot\Bigl(\frac{837}{2}\Bigr)^{2}\ }.
  \]
  因而两者判别式的平方自由部分同为 \(-111=-3\cdot 37\)，对应的唯一二次分辨子域皆为 \(\QQ(\sqrt{-111})\)；
  从而两条 \(S_3\)-扩张的分裂域共享该虚二次子域；并且三次数域 \(K=\QQ(\alpha)=\QQ(\beta)\) 的域判别式为 \(\mathrm{Disc}(K)=-999=-3^{3}\cdot 37\)（见结论 \ref{con:xi-terminal-zm-leyang-cubic-field-minimal-discriminant}），其中 \(-111\) 为其二次分辨子域 \(\QQ(\sqrt{-111})\) 的基本判别式。
  同时在该三次域中，素数 \(37\) 必然分歧（其在上述判别式中以奇数次幂出现，无法被整数指标平方消去），从而 Lee--Yang 立方域在算术上被椭圆判别式素数 \(37\) 牢固锚定。
  \item 由于 \(16X^3-9X^2+1\) 与 \(P_{\mathrm{LY}}(y)\) 皆在 \(\QQ\) 上不可约且判别式非平方，故其分裂域伽罗瓦群皆同构于 \(S_3\)；
  并由 \(\QQ(\alpha)=\QQ(\beta)\) 得两者的最小分裂域一致。
\end{enumerate}
\end{conclusion}
\begin{proof}[证明要点]
由结论 \ref{con:xi-terminal-zm-leyang-ramification-critical-values} 的临界条件得到
\(
4\alpha\beta=4\alpha^3-3\alpha^2-2\alpha+1
\)，故 \(\beta\in\QQ(\alpha)\)。
将该线性关系与 \(\alpha\) 的三次方程作消元可得 \(P_{\mathrm{LY}}(\beta)=0\)，从而 \(\QQ(\beta)\subseteq\QQ(\alpha)\)。
两者均为三次扩张，故必相等。
判别式分解为直接代数计算。
可复现核验脚本：\texttt{scripts/exp\_fold\_zm\_elliptic\_leyang\_cover\_geometry\_audit.py}。证毕。
\end{proof}

\begin{conclusion}[\texorpdfstring{$3$}{3}-division 判别式的平方自由核：\texorpdfstring{$\QQ(\sqrt{-111})$}{Q(sqrt(-111))} 的再现]\label{con:xi-terminal-zm-elliptic-3division-discriminant}
沿用结论 \ref{con:xi-terminal-zm-elliptic-normalization} 的短 Weierstrass 模型
\(
E:\ Y^{2}=X^{3}-X+\tfrac14
\)。
令 \(\psi_{3}(X)\) 表示 \(E\) 的 \(3\)-division 多项式（短 Weierstrass 形式下），则
\[
\psi_{3}(X)=3X^{4}-6X^{2}+3X-1\in\ZZ[X],
\qquad
\mathrm{Disc}\bigl(\psi_{3}\bigr)=-3^{3}\cdot 37^{2}.
\]
因此 \(\mathrm{Disc}(\psi_3)\) 的平方自由部分为 \(-111\)，从而 \(\psi_3\) 的分裂域的唯一二次子域为
\[
\QQ\!\left(\sqrt{\mathrm{Disc}(\psi_3)}\right)=\QQ(\sqrt{-111}).
\]
该二次域与结论 \ref{con:xi-terminal-zm-leyang-cubic-field-isomorphism} 中由 Lee--Yang 非平凡分歧轨道所强制出现的二次预解域一致，从而把 \(y\)-映射的有限分歧算术与 \(E[3]\) 的投影伽罗瓦作用在常数层面锁定到同一虚二次背景。
\end{conclusion}
\begin{proof}[证明要点]
\(\psi_3\) 的标准闭式来自短 Weierstrass 形式的 division 多项式公式（在 \(a=-1,b=1/4\) 的代入下即得所示四次）；其判别式可由四次判别式公式或符号计算直接得到。
唯一二次子域由 \(\QQ(\sqrt{\mathrm{Disc}(\psi_3)})\) 给出，而 \(\mathrm{Disc}(\psi_3)\) 的平方自由核为 \(-111\) 即得结论。
可复现核验脚本：\texttt{scripts/exp\_fold\_zm\_elliptic\_3torsion\_y\_image\_octic\_audit.py}。证毕。
\end{proof}

\begin{conclusion}[\texorpdfstring{$E[3]$}{E[3]} 在 \texorpdfstring{$y$}{y}-映射下的消元刻画：八次整系数方程与 \texorpdfstring{$37^{10}$}{37^10} 的高次嵌入]\label{con:xi-terminal-zm-elliptic-3torsion-y-image-octic}
仍沿用结论 \ref{con:xi-terminal-zm-elliptic-normalization} 的 \(E\) 与有理函数
\[
y:=X^{2}+Y-\tfrac12\in\QQ(E).
\]
令
\[
L_{3}(y):=\mathrm{Res}_{X}\bigl(\psi_{3}(X),\Pi(X,y)\bigr)\in\ZZ[y],
\qquad
\Pi(X,y):=X^{4}-X^{3}-(2y+1)X^{2}+X+y(y+1).
\]
则 \(L_{3}(y)\) 为次数 \(8\) 的整系数多项式，并可取其闭式表达为：
\[
\boxed{\;
L_{3}(y)=81 y^{8}-324 y^{7}+297 y^{6}+729 y^{5}-216 y^{4}-1395 y^{3}-480 y^{2}-120 y+7\; }.
\]
并且 \(E[3]\setminus\{O\}\) 在重量投影下的有限像满足严格刻画
\[
y\bigl(E[3]\setminus\{O\}\bigr)=\{\,u\in\overline{\QQ}:\ L_{3}(u)=0\,\}.
\]
该八次多项式的判别式分解为
\[
\boxed{\;
\mathrm{Disc}=-3^{31}\cdot 37^{10}\cdot 48973^{2}\; }.
\]
因此，\(y(E[3]\setminus\{O\})\) 的消元描述在判别式层面不可避免地携带导子坏素数 \(37\) 的高次嵌入；该现象为后续研究 \(37\)-进局部分歧、导子指数与扭点像空间的细分解提供了直接的整系数证书入口。
\par
此外，结论 \ref{con:xi-terminal-zm-elliptic-weight-tripling-modular-equation} 的三倍乘模方程 \(M_{3}(y,y_{3})\) 的首项系数满足 \(d_{4}(y)=L_{3}(y)^{4}\)，从而 \([3]\) 的极点支撑像与三除点重量像在同一八次消元数据上对齐。
\end{conclusion}
\begin{proof}[证明要点]
由 \(y=X^{2}+Y-\tfrac12\) 得
\(
Y=y-X^{2}+\tfrac12
\)，代入曲线方程 \(Y^{2}=X^{3}-X+\tfrac14\) 即得谱四次关系 \(\Pi(X,y)=0\)。
另一方面，非平凡三除点满足 \(\psi_3(X)=0\)。
因此对 \(X\) 作消元（取结果式 \(\mathrm{Res}_X(\psi_3(X),\Pi(X,y))\)）即可得到所示八次多项式 \(L_3(y)\)，其零点集合恰刻画 \(y(E[3]\setminus\{O\})\) 的有限像；
判别式分解为该多项式在 \(\ZZ[y]\) 中的直接代数计算。
可复现核验脚本：\texttt{scripts/exp\_fold\_zm\_elliptic\_3torsion\_y\_image\_octic\_audit.py}。证毕。
\end{proof}



\begin{conclusion}[Lee--Yang 三次数域的域判别式与分支判别式：坏素数 \texorpdfstring{$37$}{37} 的同步标记]\label{con:xi-terminal-zm-leyang-cubic-field-minimal-discriminant}
沿用结论 \ref{con:xi-terminal-zm-leyang-cubic-field-isomorphism} 的记号，令 \(\beta\) 为 Lee--Yang 三次
\[
P_{\mathrm{LY}}(y)=256y^{3}+411y^{2}+165y+32\in\ZZ[y]
\]
的任一根，并记 \(K:=\QQ(\beta)\)（等价地 \(K=\QQ(\alpha)\)，其中 \(\alpha\) 为 \(16X^{3}-9X^{2}+1=0\) 的根）。
则：
\begin{enumerate}
  \item 多项式判别式与域判别式分别满足
  \[
  \boxed{\ \mathrm{Disc}\bigl(P_{\mathrm{LY}}\bigr)=-3^{9}\cdot 31^{2}\cdot 37\ },
  \qquad
  \boxed{\ \mathrm{Disc}(K)=-999=-3^{3}\cdot 37\ }.
  \]
  并且二者满足指数平方关系
  \[
  \mathrm{Disc}\bigl(P_{\mathrm{LY}}\bigr)=(3^{3}\cdot 31)^2\,\mathrm{Disc}(K).
  \]
  因而素数 \(31\) 仅通过指数 \([\mathcal O_K:\ZZ[\beta]]\) 的平方因子进入 \(\mathrm{Disc}\bigl(P_{\mathrm{LY}}\bigr)\)，并不必作为 \(K/\QQ\) 的分歧素数出现。
  \item 因而扩张 \(K/\QQ\) 的分歧素集恰为 \(\{3,37\}\)；
  其中 \(3\) 为野分歧且全分歧，而 \(37\) 为驯分歧且分解为
  \[
  \boxed{\ 3\mathcal O_K=\mathfrak p_3^{3},\qquad 37\mathcal O_K=\mathfrak p_{37}^{2}\mathfrak q_{37}\ }.
  \]
  特别地，素数 \(37\) 在 \(K\) 中必分歧；结合结论 \ref{con:xi-terminal-zm-elliptic-normalization} 与注记 \ref{rem:xi-terminal-elliptic-37a1-modular-anchor}，其亦为谱椭圆曲线 \(E\) 的导子与 level \(37\) 的模性锚点坏素数，从而在素数层面把 Lee--Yang 分支算术与导出自守对象的最小坏素数严格对齐。
  \item \(K\) 的 \(S_3\)-分裂域之唯一二次分辨子域为
  \[
  \boxed{\ \QQ(\sqrt{-111})\ }.
  \]
  其中 \(-111=-3\cdot 37\) 为 \(\mathrm{Disc}\bigl(P_{\mathrm{LY}}\bigr)\) 的平方自由部分，因此该虚二次域同时携带 Lee--Yang 分支三次与坏素数 \(37\) 的共同判别式指纹。
\end{enumerate}
\end{conclusion}
\begin{proof}[证明要点]
取整生成元 \(\theta:=256\beta\)（从而 \(\QQ(\theta)=\QQ(\beta)\)），则 \(\theta\) 满足首一整式
\(
\theta^{3}+411\theta^{2}+42240\theta+2097152=0
\)。
对该整模型求最大整环并计算域判别式得到 \(\mathrm{Disc}(K)=-999\)，并由判别式判据确定仅在 \(3\) 与 \(37\) 处分歧；
再对 \(p\in\{3,37\}\) 进行素理想分解得到所示分解型。
另一方面，\(\mathrm{Disc}\bigl(P_{\mathrm{LY}}\bigr)\) 可由多项式判别式公式直接计算；并由一般恒等式 \(\mathrm{Disc}\bigl(P_{\mathrm{LY}}\bigr)=\mathrm{Disc}(K)\,[\mathcal O_K:\ZZ[\beta]]^{2}\) 得到上式指数平方关系。
可复现核验脚本：\texttt{scripts/exp\_fold\_zm\_elliptic\_leyang\_arithmetic\_audit.py}。证毕。
\end{proof}

\begin{conjecture}[射线类场刻画：\texorpdfstring{$\QQ(\sqrt{-111})$}{Q(sqrt(-111))} 上的循环三次分裂域]\label{conj:xi-terminal-zm-leyang-ray-class-field}
令 \(L\) 为 Lee--Yang 三次 \(P_{\mathrm{LY}}(y)\) 的分裂域，并令
\(
K_{\Delta}:=\QQ(\sqrt{-111})
\)
为其二次分辨子域（等价地，\(K_{\Delta}=\QQ(\sqrt{\mathrm{Disc}(P_{\mathrm{LY}})})\)，见结论 \ref{con:xi-terminal-zm-leyang-cubic-arithmetic}）。
则 \(L/K_{\Delta}\) 为循环三次扩张。
进一步猜想 \(L\) 恰为 \(K_{\Delta}\) 的某一射线类域，并且其导子可取为
\[
\mathfrak f=3^{4}\cdot 31
\qquad(\text{按 }K_{\Delta}\text{ 的理想分解理解}).
\]
等价地，记 \(y_1,y_2,y_3\) 为重量覆盖 \(\pi_y:E\to\mathbb P^1_y\) 的三条非平凡有限分歧值（即 \(P_{\mathrm{LY}}(y)=0\) 的三根），则
\[
L=\QQ(y_1,y_2,y_3)=K_{\Delta}(y_1,y_2,y_3).
\]
\end{conjecture}

\begin{remark}[判别式分解与导子候选]\label{rem:xi-terminal-zm-leyang-rayclass-discriminant-motivation}
由定理 \ref{thm:xi-terminal-zm-leyang-elliptic-structure} 的判别式分解
\[
\mathrm{Disc}\bigl(P_{\mathrm{LY}}(y)\bigr)=-3^{9}\cdot 31^{2}\cdot 37
=-111\cdot(3^{4}\cdot 31)^{2}
\]
可见二次分辨子域的基本判别式 \(-111\) 被平方因子精确剥离；在类场论的导子--判别式对照中，这一平方结构为 \(\mathfrak f\) 的候选提供了自然的可审计接口。
\end{remark}

\begin{conclusion}[分歧特征值与分歧参数的二次互换：缩放生成元的同域性与指数公式]\label{con:xi-terminal-zm-leyang-uv-quadratic-exchange-index}
仍取非平凡分歧值 \(y_0\)（即 \(P_{\mathrm{LY}}(y_0)=0\)）与其对应的并合谱值 \(\lambda_0\)（即 \(\Pi(\lambda_0,y_0)=\partial_\lambda\Pi(\lambda_0,y_0)=0\)，等价于 \(16\lambda_0^3-9\lambda_0^2+1=0\) 且 \(y_0=\frac{4\lambda_0^3-3\lambda_0^2-2\lambda_0+1}{4\lambda_0}\)，见结论 \ref{con:xi-terminal-zm-leyang-ramification-critical-values}）。
定义缩放生成元
\[
u:=16\lambda_0,\qquad v:=256y_0.
\]
则 \(u\) 与 \(v\) 生成同一三次域 \(K\)，并且满足显式二次互换关系
\begin{equation}\label{eq:xi-terminal-zm-leyang-uv-quadratic-exchange}
\boxed{\;
3u^2+(v+128)u-768=0\;}
\end{equation}
从而 \(u\in\QQ(v)\) 且 \(v\in\QQ(u)\)，并且 \(\QQ(u)=\QQ(v)=K\)。
等价地，利用 \(u\neq 0\) 可写出
\[
v=u^2-12u-128+\frac{1024}{u}=-3u^2+24u-128\in\ZZ[u],
\]
因而 \(\ZZ[v]\subseteq \ZZ[u]\)。
\par
记
\[
f_u(t):=t^3-9t^2+256,\qquad
f_v(t):=t^3+411t^2+42240t+2^{21}.
\]
则 \(u\) 为 \(f_u\) 的根且 \(v\) 为 \(f_v\) 的根；并且二者的判别式满足
\[
\mathrm{Disc}(f_u)=-2^{10}\cdot 3^{3}\cdot 37,\qquad
\mathrm{Disc}(f_v)=-2^{16}\cdot 3^{9}\cdot 31^{2}\cdot 37.
\]
从而两单生成整环的指数满足严格公式
\begin{equation}\label{eq:xi-terminal-zm-leyang-uv-index}
\boxed{\;
[\ZZ[u]:\ZZ[v]]=2^3\cdot 3^3\cdot 31=6696.\;}
\end{equation}
\end{conclusion}
\begin{proof}[证明要点]
由结论 \ref{con:xi-terminal-zm-leyang-ramification-critical-values} 的参数化
\(
y_0=\frac{4\lambda_0^3-3\lambda_0^2-2\lambda_0+1}{4\lambda_0}
\)
得
\[
v
=256y_0
=\frac{64}{\lambda_0}\bigl(4\lambda_0^3-3\lambda_0^2-2\lambda_0+1\bigr)
=u^2-12u-128+\frac{1024}{u}.
\]
又由 \(16\lambda_0^3-9\lambda_0^2+1=0\) 得 \(u\) 满足 \(u^3-9u^2+256=0\)，从而 \(\frac{1024}{u}=-4u^2+36u\)，得到 \(v=-3u^2+24u-128\in\ZZ[u]\) 及包含 \(\ZZ[v]\subseteq\ZZ[u]\)。
将 \(u^3=9u^2-256\) 代入恒等式 \(vu=u^3-12u^2-128u+1024\) 并整理即得 \eqref{eq:xi-terminal-zm-leyang-uv-quadratic-exchange}，从而 \(\QQ(u)=\QQ(v)\)。
\par
由 \(\lambda_0=u/16\) 代入三次方程 \(16\lambda^3-9\lambda^2+1=0\) 得 \(f_u(u)=0\)；
由 \(y_0=v/256\) 代入 \(P_{\mathrm{LY}}(y)=0\) 并清分母得 \(f_v(v)=0\)。
判别式由缩放公式与 \(\mathrm{Disc}(16\lambda^3-9\lambda^2+1)\)、\(\mathrm{Disc}(P_{\mathrm{LY}})\) 的闭式（定理 \ref{thm:xi-terminal-zm-leyang-elliptic-structure}）直接推出。
最后，由 \(\mathrm{Disc}(\ZZ[\theta])=\mathrm{Disc}(m_\theta)\)（\(\theta\) 为代数整数且 \(m_\theta\) 为其首一极小多项式）与
\(
\mathrm{Disc}(\ZZ[v])=\mathrm{Disc}(\ZZ[u])\,[\ZZ[u]:\ZZ[v]]^{2}
\)
得
\(
[\ZZ[u]:\ZZ[v]]=\sqrt{\mathrm{Disc}(f_v)/\mathrm{Disc}(f_u)}=2^3\cdot 3^3\cdot 31
\)。
证毕。
\end{proof}

\begin{remark}[二次分辨子域的规范特征与权一自守对接]\label{rem:xi-terminal-zm-leyang-artin-weight1-interface}
令 \(L\) 为 \(P_{\mathrm{LY}}\) 的分裂域。
由 \(\mathrm{Disc}\bigl(P_{\mathrm{LY}}\bigr)=-3^{9}\cdot 31^{2}\cdot 37\) 的非平方性可知 \(\mathrm{Gal}(L/\QQ)\cong S_3\)，并且其唯一二次分辨子域为
\(
\QQ(\sqrt{-111})
\)
（结论 \ref{con:xi-terminal-zm-leyang-cubic-field-minimal-discriminant}）。
取 \(S_3\hookrightarrow \mathrm{GL}_2(\CC)\) 为标准二维不可约表示，则复合
\[
\rho_{\mathrm{LY}}:\mathrm{Gal}(\overline{\QQ}/\QQ)\twoheadrightarrow \mathrm{Gal}(L/\QQ)\hookrightarrow \mathrm{GL}_2(\CC)
\]
给出一个二维不可约 Artin 表示，其行列式等于符号特征，从而与二次特征 \(\chi_{-111}\) 对齐。
该表示的分歧素集包含于 \(\{3,31,37\}\)，其中 \(31\) 由 \(\mathrm{Disc}\bigl(P_{\mathrm{LY}}\bigr)\) 的指数平方因子进入；而在域判别式口径下仅由 \(\{3,37\}\) 触发（\(\mathrm{Disc}(K)=-999\)）。
由于 \(S_3\) 为可解群像，\(\rho_{\mathrm{LY}}\) 属于经典 dihedral/可解像 reciprocity 的已知范畴；
从而其 Artin \(L\)-函数可与权 \(1\) 模形式的 \(L\)-函数对接，并据此将 Lee--Yang 分歧三次转写为一条可计算的权一自守数据通道。\cite{DeligneSerre1974WeightOne}
\end{remark}

\begin{remark}[\texorpdfstring{$\sqrt{-3}$}{sqrt(-3)} 自反通道与 \texorpdfstring{$\sqrt{37}$}{sqrt(37)} 模性锚点的二次扭转：\texorpdfstring{$\QQ(\sqrt{-111})$}{Q(sqrt(-111))} 的出现]\label{rem:xi-terminal-zm-leyang-sqrtminus111-twist}
结论 \ref{con:xi-terminal-zm-self-inversive-unique} 表明自反性仅在 \(y^2+y+1=0\) 时发生，从而唯一非平凡自反通道由虚二次域
\(
K_{-3}:=\QQ(\sqrt{-3})
\)
标记。
另一方面，结论 \ref{con:xi-terminal-37-squarefree-coupling} 给出椭圆锚点的 \(2\)-division 伴随二次域为
\(
K_{37}:=\QQ(\sqrt{37})
\)。
在双二次扩张 \(K_{37,-3}:=\QQ(\sqrt{37},\sqrt{-3})\) 中，第三个二次子域为
\[
\QQ(\sqrt{37}\sqrt{-3})=\QQ(\sqrt{-111}),
\]
其与 Lee--Yang 三次的二次分辨子域一致（结论 \ref{con:xi-terminal-zm-leyang-cubic-field-minimal-discriminant}）。
等价地，对任意 \(p\nmid 3\cdot 37\) 有 Kronecker 符号的乘法恒等式
\[
\left(\frac{-111}{p}\right)=\left(\frac{37}{p}\right)\left(\frac{-3}{p}\right),
\]
从而 Lee--Yang 三次分歧的二次特征可视为在 level \(37\) 椭圆锚点之上施加的唯一 \((-3)\) 扭转读出。
\end{remark}



\begin{conclusion}[Lee--Yang 分歧多项式的判别式与模 \texorpdfstring{$p$}{p} 退化：\texorpdfstring{$\{3,31,37\}$}{3,31,37} 处的分歧值碰撞]\label{con:xi-terminal-zm-leyang-branchpolynomial-discriminant-modp-collisions}
令 Lee--Yang 三次多项式
\[
P_{\mathrm{LY}}(y)=256y^3+411y^2+165y+32\in\ZZ[y],
\]
并记整体分歧多项式
\(
B(y):=y(y-1)P_{\mathrm{LY}}(y)\in\ZZ[y]
\)。
则：
\begin{enumerate}
  \item \(P_{\mathrm{LY}}\) 在 \(\QQ\) 上不可约。
  \item \(P_{\mathrm{LY}}\) 的判别式为
  \[
  \boxed{\ \mathrm{Disc}(P_{\mathrm{LY}})=-3^9\cdot 31^2\cdot 37\ }.
  \]
  \item 因而 \(P_{\mathrm{LY}}\) 的分裂域仅可能在素数 \(\{3,31,37\}\) 处分歧，并且 \(37\) 必分歧。
  \item \(B(y)\) 的判别式为
  \[
  \boxed{\ \mathrm{Disc}(B)=-2^{20}\cdot 3^{15}\cdot 31^2\cdot 37\ }.
  \]
  因此分歧值在特征 \(p\) 下发生合并的可能性仅出现在 \(p\in\{2,3,31,37\}\)；
  其中对 \(P_{\mathrm{LY}}\) 的实际碰撞严格发生在 \(p\in\{3,31,37\}\)。
  \item 上述退化在模素数下具有完全显式的分解型：
  \[
  P_{\mathrm{LY}}(y)\equiv (y-1)^3 \pmod 3,\qquad
  P_{\mathrm{LY}}(y)\equiv 8\,(y-12)^2(y-6)\pmod{31},
  \]
  \[
  P_{\mathrm{LY}}(y)\equiv -3\,(y-14)(y-6)^2\pmod{37}.
  \]
  因而在特征 \(3\) 下三条 Lee--Yang 分歧值全体塌缩为单点（对应三重根），而在特征 \(31\) 与 \(37\) 下发生二重碰撞（一次根与一次二次根并存）。
  \item 素数 \(37\) 同时作为结论 \ref{con:xi-terminal-zm-elliptic-lattes-critical-sextic-galois} 的 Latt\`es 临界六次判别式分歧素数之一，亦作为 \(\Delta(E)=37\) 的唯一奇素坏约化点（结论 \ref{con:xi-terminal-zm-elliptic-normalization}），从而该动力系统的临界结构与 Lee--Yang 分歧几何在算术敏感性上发生对齐。
\end{enumerate}
\end{conclusion}
\begin{proof}[证明要点]
不可约性可由模 \(5\) 的化简给出：在 \(\FF_5[y]\) 中有
\[
P_{\mathrm{LY}}(y)\equiv y^{3}+y^{2}+2\pmod 5,
\]
而对 \(y\in\FF_5\) 逐点代入可知其无根；由于三次多项式在域上可约当且仅当有一次因子，故 \(P_{\mathrm{LY}}\) 在 \(\FF_5\) 上不可约，从而在 \(\QQ\) 上不可约。
\par
对分歧素数集合，令 \(\beta\) 为 \(P_{\mathrm{LY}}\) 的任一根并记 \(L:=\QQ(\beta)\)，再记 \(K\) 为 \(P_{\mathrm{LY}}\) 的分裂域。
由判别式与整指数的标准关系
\[
\mathrm{Disc}(P_{\mathrm{LY}})=\mathrm{Disc}(L)\cdot[\mathcal O_L:\ZZ[\beta]]^{2}
\]
可知：\(\mathrm{Disc}(P_{\mathrm{LY}})\) 的素因子给出 \(L/\QQ\) 的分歧素数上界，因而分裂域 \(K\) 的分歧素数也只能来自 \(\{3,31,37\}\)。
并且由于指数项为平方，\(37\) 在 \(\mathrm{Disc}(P_{\mathrm{LY}})\) 中以奇数次幂出现不能被吸收，故 \(37\mid \mathrm{Disc}(L)\)，从而 \(37\) 在 \(L\) 中分歧；因此在包含 \(L\) 的分裂域 \(K\) 中亦必分歧。
\par
判别式分解与模 \(p\) 因式分解均为 \(\ZZ[y]\) 与 \((\ZZ/p\ZZ)[y]\) 内的有限代数计算；
\(\mathrm{Disc}(B)\) 亦可由 \(\mathrm{Disc}(P_{\mathrm{LY}})\) 与结果式 \(\mathrm{Res}(P_{\mathrm{LY}},y(y-1))\) 的乘法公式回收。
可复现核验脚本：\texttt{scripts/exp\_fold\_zm\_elliptic\_audit.py}。证毕。
\end{proof}



\begin{conclusion}[坏素数 \texorpdfstring{$37$}{37} 处的同根塌缩：三条三次控制多项式的判别式与同余分解]\label{con:xi-terminal-zm-leyang-badprime-37-common-double-root}
沿用结论 \ref{con:xi-terminal-zm-leyang-conjugate-fiber-cubic} 的共轭三次
\(
Q_{\mathrm{LY}}(y)=256y^{3}+195y^{2}+93y-48
\)
与 Lee--Yang 三次
\(
P_{\mathrm{LY}}(y)=256y^{3}+411y^{2}+165y+32
\)，
并记二分点像控制三次（结论 \ref{con:xi-terminal-zm-2division-norm-square}）
\[
C(y):=16y^{3}-8y^{2}-4y+1
\qquad(\text{亦即结论 \ref{con:xi-terminal-zm-2division-norm-square} 中的 }S(y)).
\]
则三者的判别式均携带素因子 \(37\)，且满足严格分解：
\[
\mathrm{Disc}\bigl(C(y)\bigr)=2^{8}\cdot 37,\qquad
\mathrm{Disc}\bigl(P_{\mathrm{LY}}(y)\bigr)=-3^{9}\cdot 31^{2}\cdot 37,\qquad
\mathrm{Disc}\bigl(Q_{\mathrm{LY}}(y)\bigr)=-3^{3}\cdot 37\cdot 2677^{2}.
\]
更进一步地，在 \(\FF_{37}[y]\) 中三者发生同一参数值的双根塌缩：
\[
\begin{aligned}
C(y)&\equiv 16\,(y-7)(y-6)^{2}\ (\mathrm{mod}\ 37),\\
P_{\mathrm{LY}}(y)&\equiv -3\,(y-14)(y-6)^{2}\ (\mathrm{mod}\ 37),\\
Q_{\mathrm{LY}}(y)&\equiv -3\,(y-16)(y-6)^{2}\ (\mathrm{mod}\ 37).
\end{aligned}
\]
因此 \(y\equiv 6\ (\mathrm{mod}\ 37)\) 同时为三条三次控制多项式的共同重根。
与此同时，对物理基点 \(P=(2,-\tfrac52)\in E(\QQ)\)（结论 \ref{con:xi-terminal-zm-elliptic-rational-points-denominator-law}），其逆元满足 \(y(-P)=6\)，
从而该塌缩点在整数层面与一条显式有理点标记相一致。
\end{conclusion}
\begin{proof}[证明要点]
判别式与模 \(37\) 因式分解为 \(\ZZ[y]\) 与 \((\ZZ/37\ZZ)[y]\) 内的有限代数计算。
\(y(-P)=6\) 由 \(y=X^{2}+Y-\tfrac12\) 在点 \(-P=(2,\tfrac52)\) 的直接代入得到。
可复现核验脚本：\texttt{scripts/exp\_fold\_zm\_elliptic\_leyang\_cover\_geometry\_audit.py}。证毕。
\end{proof}

\begin{conclusion}[坏素数 \(37\) 处的 Latt\`es 稳定约化：公共平方因子、Chebyshev 共轭与结点退化的同步]\label{con:xi-terminal-zm-elliptic-lattes-badprime-37-degree-drop}
沿用结论 \ref{con:xi-terminal-zm-elliptic-lattes-doubling} 的倍点 Latt\`es 自映射
\[
\Phi(X)=\frac{X^4+2X^2-2X+1}{4X^3-4X+1}\in\QQ(X).
\]
则在 \(\FF_{37}[X]\) 中出现非平凡公共因子，且满足显式同余分解
\[
X^4+2X^2-2X+1\equiv (X-5)^2\,(X^2+10X+3)\pmod{37},
\]
\[
4X^3-4X+1\equiv (X-5)^2\,(4X+3)\pmod{37}.
\]
因此在稳定约化意义下（约去公共因子）得到次数严格下降的约化映射
\[
\boxed{\;
\overline{\Phi}(X)=\frac{X^2+10X+3}{4X+3}\in\FF_{37}(X),\qquad \deg\overline{\Phi}=2<\deg\Phi=4\; } .
\]
\par
并且在 \(\PP^{1}/\FF_{37}\) 上存在显式 Möbius 变换
\[
M(X)=\frac{2X+13}{X-5}\in \mathrm{PGL}_{2}(\FF_{37})
\]
使得
\[
\boxed{\;
M\circ \overline{\Phi}\circ M^{-1}(X)=X^{2}-2\; } .
\]
因而 \(\overline{\Phi}\) 在 \(\FF_{37}\) 上与归一化 Chebyshev 多项式 \(T_{2}(X)=X^{2}-2\) 共轭。
在该约化动力系统中，\(X=5\) 满足 \(\overline{\Phi}(5)=5\) 且 \(\overline{\Phi}'(5)=0\)，从而为唯一的临界不动点；
同时 \(M\) 在 \(X=5\) 处具有极点并将其送至 \(\infty\)，从而把 \(\overline{\Phi}\) 刚性化为多项式正规形。
\par
另一方面，Lee--Yang 分歧点的三次控制多项式（结论 \ref{con:xi-terminal-zm-leyang-ramification-critical-values}）
\[
g(X):=16X^3-9X^2+1
\]
在 \(\FF_{37}[X]\) 中发生双根塌缩：
\[
g(X)\equiv 16\,(X-16)(X-5)^2\pmod{37}.
\]
并且 Lee--Yang 三次因子在 \(\FF_{37}[y]\) 上同样出现重根融合（结论 \ref{con:xi-terminal-zm-leyang-badprime-37-common-double-root}）：
\[
P_{\mathrm{LY}}(y)\equiv -3\,(y-14)(y-6)^2\pmod{37}.
\]
\par
因此判别式二次脊曲线
\[
C:\quad W^{2}=-y(y-1)P_{\mathrm{LY}}(y)
\]
在 \(p=37\) 处出现由二重根 \(y\equiv 6\) 所触发的半稳定退化：
其特殊纤维满足
\[
C_{\FF_{37}}:\quad W^{2}=3\,y(y-1)(y-14)(y-6)^{2},
\]
为不可约结点曲线。
令 \(Z=W/(y-6)\)，则归一化为椭圆曲线
\[
E_{14}:\quad Z^{2}=3\,y(y-1)(y-14)
\]
并且该结点为非分裂结点（两支仅在 \(\FF_{37^{2}}\) 上分离）。
从而对任意 \(n\ge 1\) 有点数恒等式
\[
\card{C(\FF_{37^{n}})}=\card{E_{14}(\FF_{37^{n}})}+(-1)^{n+1},
\]
并得到局部 \(\zeta\)-函数的显式分解
\[
Z_{C,37}(T)=(1+T)\,Z_{E_{14},37}(T),\qquad
Z_{E_{14},37}(T)=\frac{1-2T+37T^{2}}{(1-T)(1-37T)}.
\]
\par
综上，素数 \(37\) 同时承载：倍点 Latt\`es 自映射的稳定约化降阶与 Chebyshev 共轭、Lee--Yang 三次分歧的重根碰撞、以及判别式二次脊曲线的非分裂结点退化；
这些同步现象共同刻画了该椭圆化机制在唯一坏素处的算术奇异性。
\end{conclusion}
\begin{proof}[证明要点]
模 \(37\) 的因式分解与约去为有限整系数计算。
Möbius 共轭可在 \(\FF_{37}(X)\) 中直接代入化简核验。
结点曲线的归一化由代换 \(Z=W/(y-6)\) 给出，而非分裂性等价于归一化纤维在 \(y=6\) 处的两点仅在二次扩张上可见，从而点数差异呈 \((-1)^{n+1}\) 交替并推出所示 \(\zeta\)-因子。
可复现核验脚本：\texttt{scripts/exp\_fold\_zm\_elliptic\_lattes\_rational\_points\_audit.py}（\(\overline{\Phi}\) 的约化、Chebyshev 共轭与周期枚举），\texttt{scripts/exp\_fold\_zm\_elliptic\_leyang\_cover\_geometry\_audit.py}（\(E_{14}\) 点数与 \((-1)^{n+1}\) 修正）。
最后一条重根分解由结论 \ref{con:xi-terminal-zm-leyang-badprime-37-common-double-root} 给出。证毕。
\end{proof}

\begin{conclusion}[权坐标除法结果式的整系数性、根像刻画与次数律]\label{con:xi-terminal-zm-elliptic-division-resultant-integrality-degree}
在椭圆曲线
\[
E:\ Y^{2}=X^{3}-X+\frac14,
\qquad
y:=X^{2}+Y-\frac12,
\]
与四次关系
\[
\Pi(X,y):=X^{4}-X^{3}-(2y+1)X^{2}+X+y(y+1)=0
\]
的设定下，对任意奇整数 \(n\ge 3\)，令 \(\psi_{n}(X)\in\ZZ[X]\) 为 \(E\) 的 \(n\)-division 多项式，并定义
\[
L_{n}(y):=\operatorname{Res}_{X}\!\bigl(\psi_{n}(X),\Pi(X,y)\bigr)\in\ZZ[y].
\]
则成立：
\begin{enumerate}
  \item \(L_{n}(y)\) 具有整系数，且 \(\deg L_{n}=n^{2}-1\)；
  \item 根集满足
  \[
  \{\,\alpha\in\overline{\QQ}:L_{n}(\alpha)=0\,\}
  =\{\,y(P):P\in E[n]\setminus\{O\}\,\},
  \]
  并且在特征 \(0\) 下 \(L_n\) 可分。
\end{enumerate}
\end{conclusion}
\begin{proof}[证明要点]
由 \(\psi_n(X)\in\ZZ[X]\) 与 \(\Pi(X,y)\in\ZZ[X,y]\)，结果式 \(L_n(y)\) 直接属于 \(\ZZ[y]\)。
对任意 \(\psi_n(x_0)=0\)，多项式 \(\Pi(x_0,y)\) 关于 \(y\) 为首一二次式，故每个 \(x_0\) 贡献 \(2\) 次 \(y\)-度；
结合奇数 \(n\) 的标准次数公式 \(\deg_X\psi_n=(n^2-1)/2\)，得到 \(\deg_y L_n=n^2-1\)。
\par
若 \(P=(x_0,Y_0)\in E[n]\setminus\{O\}\)，则 \(\psi_n(x_0)=0\)，且 \(y(P)=x_0^2+Y_0-\tfrac12\) 满足 \(\Pi(x_0,y(P))=0\)，故 \(L_n(y(P))=0\)。
反之若 \(L_n(\alpha)=0\)，则存在 \(x_0\) 使 \(\psi_n(x_0)=0\) 且 \(\Pi(x_0,\alpha)=0\)；
令 \(Y_0:=\alpha+\tfrac12-x_0^2\)，代回即得 \((x_0,Y_0)\in E(\overline{\QQ})\)，并由 \(\psi_n(x_0)=0\) 知其为非平凡 \(n\)-挠点。
因此 \(\alpha=y(P)\) 对某 \(P\in E[n]\setminus\{O\}\) 成立。
在 \(\mathrm{char}=0\) 下，\(E[n]\) 为 étale 群簇，故对应消元方程无不可分重根。证毕。
\end{proof}

\begin{conclusion}[三除点重量像在 \texorpdfstring{$\FF_{37}$}{F37} 上的六重坍缩]\label{con:xi-terminal-zm-elliptic-3torsion-y-image-mod37-collapse}
仍沿用结论 \ref{con:xi-terminal-zm-elliptic-3torsion-y-image-octic} 的八次消元多项式
\[
L_{3}(y)=81 y^{8}-324 y^{7}+297 y^{6}+729 y^{5}-216 y^{4}-1395 y^{3}-480 y^{2}-120 y+7\in\ZZ[y].
\]
则在 \(\FF_{37}[y]\) 中发生高重塌缩分解
\[
L_{3}(y)\equiv 7\,(y-6)^{6}\,(y^{2}-5y-1)\pmod{37}.
\]
特别地，\(y\equiv 6\) 为 \(y\bigl(E[3]\setminus\{O\}\bigr)\) 的模 \(37\) 主聚集像，其重数为 \(6=3^{2}-3\)。
\end{conclusion}
\begin{proof}[证明要点]
对 \(L_{3}(y)\) 作模 \(37\) 约化并在 \(\FF_{37}[y]\) 中因式分解即可得到所示同余分解。
可复现核验脚本：\texttt{scripts/exp\_fold\_zm\_elliptic\_odd\_torsion\_mod37\_collapse\_audit.py}。证毕。
\end{proof}

\begin{conclusion}[五除点重量消元多项式的显式证书、模 \(37\) 分解与判别式平方类]\label{con:xi-terminal-zm-elliptic-5torsion-y-image-degree24-mod37-collapse}
仍在椭圆曲线
\(
E:\ Y^{2}=X^{3}-X+\tfrac14
\)
与重量函数
\(
y=X^{2}+Y-\tfrac12
\)
的设定下，取 \(5\)-division 多项式（\(X\)-坐标形式）
\[
\psi_{5}(X)=5X^{12}-62X^{10}+95X^{9}-105X^{8}-60X^{7}+285X^{6}-174X^{5}-5X^{4}-5X^{3}+35X^{2}-15X+2\in\ZZ[X].
\]
令
\[
L_{5}(y):=\operatorname{Res}_{X}\bigl(\psi_{5}(X),\Pi(X,y)\bigr)\in\ZZ[y],
\qquad
\Pi(X,y):=X^{4}-X^{3}-(2y+1)X^{2}+X+y(y+1).
\]
则 \(L_{5}(y)\) 在 \(\QQ\) 上不可约，且 \(\deg L_{5}=24\)。
其显式展开、模 \(37\) 分解及判别式平方类分解如下：
\input{sections/generated/eq_fold_zm_elliptic_5torsion_discriminant_quadratic_character_audit}
特别地，\(\operatorname{Disc}(L_5)\) 的平方自由核为 \(5\)。
\end{conclusion}
\begin{proof}[证明要点]
由结果式的基本性质，\(L_{5}(y)\) 的根集合恰为 \(y\bigl(E[5]\setminus\{O\}\bigr)\) 的消元像。
不可约性与模 \(37\) 分解可由脚本 \texttt{scripts/exp\_fold\_zm\_elliptic\_odd\_torsion\_mod37\_collapse\_audit.py} 核验；
判别式分解与平方类由脚本 \texttt{scripts/exp\_fold\_zm\_elliptic\_5torsion\_discriminant\_quadratic\_character\_audit.py} 核验。证毕。
\end{proof}

\begin{corollary}[五次权除法分裂域的符号二次特征]\label{cor:xi-terminal-zm-elliptic-5torsion-sign-quadratic-field}
\leanverified{paper\_xi\_terminal\_zm\_elliptic\_5torsion\_sign\_quadratic\_field}
设 \(K\) 为 \(L_{5}\) 的分裂域，\(G:=\mathrm{Gal}(K/\QQ)\)，并记
\[
\rho:G\longrightarrow S_{24}
\]
为 \(G\) 在 \(L_{5}\) 的 \(24\) 个根上的自然置换表示。
则符号特征 \(\mathrm{sgn}\circ\rho\) 所对应的二次子域为
\[
\QQ(\sqrt5),
\]
等价地，
\[
\QQ\!\bigl(\sqrt{\operatorname{Disc}(L_{5})}\bigr)=\QQ(\sqrt5).
\]
\end{corollary}
\begin{proof}
对任意可分多项式 \(f\in\QQ[y]\)，\(\operatorname{Disc}(f)\) 的平方类对应于分裂域中由符号特征切出的二次子扩张。
由结论 \ref{con:xi-terminal-zm-elliptic-5torsion-y-image-degree24-mod37-collapse}，
\(\operatorname{Disc}(L_{5})\) 的平方自由核为 \(5\)，故结论成立。证毕。
\end{proof}

\begin{conclusion}[七除点重量像在模 \(37\) 下的四十二重坍缩]\label{con:xi-terminal-zm-elliptic-7torsion-y-image-mod37-collapse}
令 \(\psi_{7}(X)\in\ZZ[X]\) 为同一曲线 \(E\) 的 \(7\)-division 多项式（\(X\)-坐标形式，次数 \(24\)）。
在有限域 \(\FF_{37}\) 上定义
\[
L_{7}^{(37)}(y):=\operatorname{Res}_{X}\bigl(\psi_{7}(X),\Pi(X,y)\bigr)\in\FF_{37}[y],
\]
其中 \(\Pi(X,y)\) 如上。
则在 \(\FF_{37}[y]\) 中存在精确分解
\[
L_{7}^{(37)}(y)= -4\,(y-6)^{42}\,\bigl(y^{6}-11y^{5}+18y^{4}+7y^{3}-6y^{2}-13y-9\bigr),
\]
其中指数 \(42=7^{2}-7\)，剩余因子次数 \(6=7-1\)。
\end{conclusion}
\begin{proof}[证明要点]
在 \(\FF_{37}[X,y]\) 中计算结果式并因式分解即可得到所示恒等式；见脚本 \texttt{scripts/exp\_fold\_zm\_elliptic\_odd\_torsion\_mod37\_collapse\_audit.py}。证毕。
\end{proof}

\begin{theorem}[坏素数 \(37\) 处全阶 torsion 重量像坍缩定理]\label{con:xi-terminal-zm-elliptic-odd-torsion-mod37-collapse-exponent}\label{thm:xi-terminal-zm-elliptic-full-order-torsion-mod37-collapse}
设
\[
E:\ Y^{2}=X^{3}-X+\tfrac14,
\qquad
y:=X^{2}+Y-\tfrac12,
\]
并令
\[
\Pi(X,y):=X^{4}-X^{3}-(2y+1)X^{2}+X+y(y+1)=0.
\]
对任意整数 \(n\ge 2\) 且 \(\gcd(n,37)=1\)，记 \(\psi_n(X)\) 为 \(E\) 的 \(n\)-division 多项式的 \(X\)-坐标形式（偶数 \(n\) 时取经典除点多项式去除因子 \(2Y\) 后所得的纯 \(X\) 多项式），并定义
\[
L_n(y):=\operatorname{Res}_{X}\!\bigl(\psi_n(X),\Pi(X,y)\bigr)\in\ZZ[y].
\]
则在 \(\FF_{37}[y]\) 中存在 \(c_n\in\FF_{37}^{\times}\) 与 \(R_n(y)\in\FF_{37}[y]\)，满足
\[
L_n(y)\equiv c_n\,(y-6)^{n^{2}-n}\,R_n(y)\pmod{37},
\qquad
\deg R_n=n-1,
\qquad
R_n(6)\neq 0.
\]
因此，对所有 \(\gcd(n,37)=1\) 的 \(n\ge 2\)，重量像在模 \(37\) 下具有精确主聚集重数 \(n^{2}-n\)，并保留严格 \(n-1\) 次的非坍缩自由度。
\end{theorem}
\begin{proof}
在 \(p=37\) 处，\(E\) 具有非分裂乘法约化。令 \(K:=\QQ_{37}\)，取其唯一不分歧二次扩张 \(K^{(2)}\)，则 \(E/K^{(2)}\) 与某 Tate 曲线 \(E_q\) 同构。于是对任意 \(\gcd(n,37)=1\)，有标准分解
\[
E[n]\cong \mu_n\times \ZZ/n\ZZ
\]
（作为 \(K^{(2)}\)-点集的抽象群），并与结点参数化兼容。

在 Tate 参数 \(u\in \overline{K^{(2)}}^{\times}/q^{\ZZ}\) 下，\(E[n]\) 的点写作 \(u=\zeta\,q^{k/n}\)（\(\zeta\in\mu_n,\ k\in\ZZ/n\ZZ\)）。当 \(k=0\) 时，\(|u|=1\)，约化落在特殊纤维的光滑分支；当 \(k\neq 0\) 时，\(|u|\neq 1\)，约化落在非分裂结点。故 \(k\neq 0\) 的点总数 \(n(n-1)=n^{2}-n\) 在特殊纤维上汇聚到同一结点。

在本模型中，结点的重量坐标为 \(y\equiv 6\pmod{37}\)，故全部 \(k\neq 0\) 点对 \(L_n(y)\) 贡献因子 \((y-6)^{n^{2}-n}\)。而 \(k=0\) 的非平凡点仅有 \(n-1\) 个，且不约化到结点，故对应剩余因子在 \(y=6\) 处不消失。再由结果式次数恒等式 \(\deg L_n=n^{2}-1\)，得到
\[
\deg R_n=(n^{2}-1)-(n^{2}-n)=n-1,
\qquad
R_n(6)\neq 0.
\]
结论成立。
\end{proof}

\begin{theorem}[倍点 Latt\`es 逆像纤维的 \(2k\)-超立方体实现]\label{thm:xi-terminal-zm-lattes-preimage-hypercube-2k}
设 \(E\) 为特征 \(\neq 2\) 的域上的椭圆曲线，\(x:E\to\PP^{1}\) 为商映射 \(P\mapsto x(P)\)，并令 \(\Phi\) 为倍点 \([2]:E\to E\) 在 \(\PP^{1}\) 上诱导的 Latt\`es 自映射，使
\[
x([2]P)=\Phi(x(P)).
\]
取 \(a\in\PP^{1}\) 使其不落在 \(\Phi\) 的后临界集，取 \(P\in E\) 满足 \(x(P)=a\) 且 \(P\notin E[2]\)。则对每个 \(k\ge 1\)：
\begin{enumerate}
  \item 存在自然双射
  \[
  \Phi^{-k}(a)\cong [2^{k}]^{-1}(P),
  \]
  从而 \(\Phi^{-k}(a)\) 是 \(E[2^k]\) 的 torsor，且
  \[
  |\Phi^{-k}(a)|=4^k=2^{2k}.
  \]
  \item \(2\)-幂挠群滤过
  \[
  E[2]\subset E[4]\subset\cdots\subset E[2^k]
  \]
  的关联分次给出
  \[
  \mathrm{gr}(E[2^k])\cong(\ZZ/2\ZZ)^{2k},
  \]
  因而在选定 \(E[2]\) 基后，\(\Phi^{-k}(a)\) 在集合层面诱导为 \(2k\) 维超立方体图 \(Q_{2k}\)。
\end{enumerate}
在曲线
\[
E:\ Y^{2}=X^{3}-X+\tfrac14
\]
上，对应的倍点 Latt\`es 映射为
\[
\Phi(X)=\frac{X^{4}+2X^{2}-2X+1}{4X^{3}-4X+1},
\]
故 \(k=3\) 时逆像纤维 \(\Phi^{-3}(a)\) 具有 \(64\) 顶点的 \(Q_{6}\) 结构。
\leanverified{paper\_xi\_terminal\_zm\_lattes\_preimage\_hypercube\_2k}
\end{theorem}
\begin{proof}
取 \(Q_0\in E\) 使 \([2^k]Q_0=P\)。定义
\[
E[2^k]\longrightarrow \Phi^{-k}(a),
\qquad
T\longmapsto x(Q_0+T).
\]
若 \(x(Q_0+T)=x(Q_0+T')\)，则 \(Q_0+T'=\pm(Q_0+T)\)。施加 \([2^k]\) 得
\[
P=[2^k](Q_0+T')=\pm[2^k](Q_0+T)=\pm P.
\]
因 \(P\notin E[2]\)，只能取正号，故 \(T'=T\)，映射单射。两侧基数同为 \(4^k\)，故为双射，得到第 1 点。

再由短正合列
\[
0\to E[2^{j-1}]\to E[2^j]\xrightarrow{[2^{j-1}]}E[2]\to 0
\]
迭代，得到 \(E[2^k]\) 的二进层位编码，并将该编码平移到 torsor \([2^k]^{-1}(P)\)，遂得 \(Q_{2k}\) 的顶点坐标与单比特翻转边关系。结论成立。
\end{proof}

\begin{corollary}[坏素数 \(37\) 处的平方根全息压缩律]\label{cor:xi-terminal-zm-badprime37-square-root-holographic-compression}
沿用定理 \ref{thm:xi-terminal-zm-elliptic-full-order-torsion-mod37-collapse} 的 \((E,y)\) 与定理 \ref{thm:xi-terminal-zm-lattes-preimage-hypercube-2k} 的 \(k\)-步 dyadic 逆像结构。令 \(n=2^k\)，并将 \(E[n]\)（等价地 \(\Phi^{-k}(a)\)）视为大小 \(4^k=2^{2k}\) 的体相空间。则
\[
E[n]\setminus\{O\}\xrightarrow{\ y\ }\PP^1\xrightarrow{\bmod\,37}\PP^1(\overline{\FF}_{37})
\]
的像集大小满足
\[
\#\bigl(y(E[n]\setminus\{O\})\bmod 37\bigr)\le n=2^k=\sqrt{|E[n]|}.
\]
且 \(y\equiv 6\pmod{37}\) 的纤维重数精确为
\[
n^2-n=2^{2k}-2^k.
\]
\end{corollary}
\begin{proof}
由定理 \ref{thm:xi-terminal-zm-elliptic-full-order-torsion-mod37-collapse}，
\[
L_n(y)\equiv c_n\,(y-6)^{n^2-n}R_n(y)\pmod{37},
\qquad
\deg R_n=n-1.
\]
因此除 \(y=6\) 外，模 \(37\) 可出现的其余 \(y\)-值至多为 \(n-1\) 个；连同 \(y=6\) 本身，总像集基数至多 \(n\)。重数结论直接来自 \((y-6)^{n^2-n}\) 的指数。证毕。
\end{proof}

\leanverified{paper\_xi\_terminal\_zm\_noncollapsed\_boundary\_frobenius\_period}
\begin{proposition}[非坍缩边界的 Frobenius 周期与边界时间上界]\label{prop:xi-terminal-zm-noncollapsed-boundary-frobenius-period}
在定理 \ref{thm:xi-terminal-zm-elliptic-full-order-torsion-mod37-collapse} 的设定下，取 \(n\ge 2\) 且 \(\gcd(n,37)=1\)。令 \(K^{(2)}/\QQ_{37}\) 为使 \(E\) 分裂乘法约化的不分歧二次扩张，其剩余域为 \(\FF_{37^2}\)。记
\[
f:=\operatorname{ord}_{n}(37^2).
\]
则：
\begin{enumerate}
  \item 定理 \ref{thm:xi-terminal-zm-elliptic-full-order-torsion-mod37-collapse} 中剩余因子 \(R_n(y)\) 的全部根都落在 \(\FF_{37^{2f}}\)；
  \item Frobenius \(\mathrm{Fr}_{37^2}\) 对这些根的轨道长度均整除 \(f\)。
\end{enumerate}
\end{proposition}
\begin{proof}
在 \(K^{(2)}\) 上的 Tate 统一化下，分量群坐标 \(0\) 的 \(n\)-torsion 对应 \(\mu_n\)。Frobenius \(\mathrm{Fr}_{37^2}\) 在 \(\mu_n\) 上作用为 \(u\mapsto u^{37^2}\)，故任一 \(u\in\mu_n\) 的轨道长度整除 \(\operatorname{ord}_n(37^2)=f\)，且 \(\mu_n\subset \FF_{37^{2f}}\)。

非坍缩 \(y\)-值来自该分量上的参数代入，因而属于同一有限扩张并继承同样的 Frobenius 轨道整除关系。证毕。
\end{proof}

\begin{theorem}[Tate--Hypercube 全息分解与坏素视界单层可见性]\label{thm:xi-terminal-zm-tate-hypercube-single-layer-visibility}
取 \(p=37\) 与定理 \ref{thm:xi-terminal-zm-elliptic-full-order-torsion-mod37-collapse} 中的曲线 \(E\)。令 \(n=2^k\)（\(k\ge 1\)）。在 \(K^{(2)}\) 上有 Tate 分解
\[
E[n]\cong \mu_n\times \ZZ/n\ZZ.
\]
选取 \(\mu_n\) 与 \(\ZZ/n\ZZ\) 的生成元，并以 \(2\)-进数位分别识别为 \(\{0,1\}^k\)，得到
\[
E[n]\cong \{0,1\}^k\times \{0,1\}^k\cong Q_k\times Q_k\cong Q_{2k}.
\]
模 \(37\) 约化后，重量读出满足：
\begin{enumerate}
  \item 对所有 \((\zeta,m)\in \mu_n\times(\ZZ/n\ZZ)\) 且 \(m\neq 0\)，有
  \[
  y(\zeta,m)\equiv 6\pmod{37};
  \]
  \item 因而 \(y\bmod 37\) 在 \(Q_k\times (Q_k\setminus\{0\})\) 上为常数 \(6\)，仅在零层 \(Q_k\times\{0\}\) 上保留非平凡变化。
\end{enumerate}
特别地，\(k=3\) 时 \(E[8]\)（等价于 \(\Phi^{-3}(a)\)）的 \(Q_6\) 顶点可分为 \(8\) 个互不相交的 \(3\)-维子立方体，其中 \(7\) 层在模 \(37\) 下塌缩到 \(6\)，仅 \(1\) 层承载非坍缩边界信息。
\end{theorem}
\begin{proof}
由 Tate 分解，点 \((\zeta,m)\) 对应参数 \(u=\zeta q^{m/n}\)。当 \(m\neq 0\) 时，\(u\) 赋值非零，约化落入非分裂结点，故重量约化恒为结点值 \(6\)。这与定理 \ref{thm:xi-terminal-zm-elliptic-full-order-torsion-mod37-collapse} 中因子 \((y-6)^{n^2-n}\) 的来源一致。

当 \(n=2^k\) 时，两因子分别以 \(k\) 比特编码，得到超立方体直积表示；零层即 \(m=0\)。从而得到单层可见性。证毕。
\end{proof}
\leanverified{paper\_xi\_terminal\_zm\_tate\_hypercube\_single\_layer\_visibility}

\begin{proposition}[Gödel 素数编码下的视界因子化与信息忘却]\label{prop:xi-terminal-zm-godel-horizon-factorization-forgetting}
\leanverified{paper\_xi\_terminal\_zm\_godel\_horizon\_factorization\_forgetting}
令 \(p_1<p_2<\cdots\) 为素数序列。对 \(k\ge 1\) 定义平方自由 Gödel 编码
\[
\mathrm{G}_k:\{0,1\}^k\to \NN_{\mathrm{sf}},
\qquad
\mathrm{G}_k(b_1,\dots,b_k):=\prod_{i=1}^k p_i^{\,b_i}.
\]
在定理 \ref{thm:xi-terminal-zm-tate-hypercube-single-layer-visibility} 的识别
\(
Q_{2k}\cong Q_k\times Q_k
\)
下，定义
\[
\mathrm{G}_{2k}(b,c):=\mathrm{G}_k(b)\cdot \prod_{i=1}^k p_{k+i}^{\,c_i}.
\]
则模 \(37\) 的重量读出塌缩等价于如下忘却因子化：存在函数 \(\mathcal B_k\) 使
\[
y\bmod 37=\mathcal B_k\circ \pi_k,
\qquad
\pi_k:\ \mathrm{G}_{2k}(b,c)\longmapsto \mathrm{G}_k(b),
\]
其中 \(\pi_k\) 即去除后 \(k\) 个素数块的投影。并且每条纤维
\[
\pi_k^{-1}\!\bigl(\mathrm{G}_k(b)\bigr)
\]
大小为 \(2^k\)，而坏素视界值 \(6\) 的最大坍缩重数为
\[
2^{2k}-2^k.
\]
\end{proposition}
\begin{proof}
由定理 \ref{thm:xi-terminal-zm-tate-hypercube-single-layer-visibility}，在 \(Q_{2k}\cong Q_k\times Q_k\) 中，除零层外第二因子 \(c\) 的全部取值在模 \(37\) 下均映到 \(6\)。故边界读出不区分后 \(k\) 个比特，只依赖前 \(k\) 个比特，即可写为 \(\mathcal B_k\circ\pi_k\)。

Gödel 编码把比特翻转转写为素数幂的乘除，去除后 \(k\) 比特即对应去除后 \(k\) 个素数块。纤维大小为第二因子基数 \(2^k\)；最大坍缩重数由定理 \ref{thm:xi-terminal-zm-elliptic-full-order-torsion-mod37-collapse} 的指数 \(n^2-n\) 在 \(n=2^k\) 时给出 \(2^{2k}-2^k\)。证毕。
\end{proof}

\input{con__xi-terminal-zm-leyang-elliptic-prime-ledger-synthesis}


\begin{conclusion}[重量二倍对应的四次模方程：双次数曲线与最小多项式证书]\label{con:xi-terminal-zm-elliptic-weight-doubling-modular-equation}
设
\[
E_0:\ \eta^{2}+\eta=x^{3}-x
\]
并定义其函数域中的次数 \(4\) 有理函数
\(
y:=x^{2}-\eta-1\in\QQ(E_0)
\)。
在同构
\(
X=x,\ Y=-\eta-\tfrac12
\)
下，\(E_0\) 等价配方为短 Weierstrass 模型
\(
E:\ Y^{2}=X^{3}-X+\tfrac14
\)，并且该重量函数等价写为 \(y=X^{2}+Y-\tfrac12\)。
记 \([2]:E_0\to E_0\) 为倍点算子，令 \(y_{2}:=y([2]P)\)。
在短 Weierstrass 坐标下，谱坐标与椭圆 \(X\)-坐标同一（结论 \ref{con:xi-terminal-zm-elliptic-normalization}），
并且 \(X([2]P)\) 由 \(X\) 的四次 Latt\`es 自映射 \(\Phi\) 给出（结论 \ref{con:xi-terminal-zm-lattes-4division-certificate}）。
考虑在 \(\PP^{1}_{y}\times \PP^{1}_{y_{2}}\) 中的 Zariski 闭包
\[
C_{2}:=\overline{\{(y(P),y([2]P)):\ P\in E_0(\overline{\QQ})\}}\subset \PP^{1}_{y}\times \PP^{1}_{y_{2}}.
\]
则 \(C_{2}\) 为一条双次数 \((16,4)\) 的不可约曲线，并存在唯一（除整体非零常数因子）整系数多项式
\[
\Phi_{2}(y,y_{2})=A_{4}(y)y_{2}^{4}+A_{3}(y)y_{2}^{3}+A_{2}(y)y_{2}^{2}+A_{1}(y)y_{2}+A_{0}(y)\in\ZZ[y,y_{2}]
\]
使 \(C_{2}=\{\Phi_{2}=0\}\)。
其中 \(\Phi_{2}\) 的显式系数证书可由对模型 \(E\) 与重量函数 \(y\) 作消元得到。
具体地，可取 \(\Phi_{2}=\mathcal H\)，其闭式展开为
\input{sections/generated/eq_fold_zm_elliptic_weight_doubling_H}
并且其首项系数为完全四次幂
\[
A_{4}(y)=\bigl(16y^{3}-8y^{2}-4y+1\bigr)^{4},
\]
其中三次因子刻画 \(E[2]\setminus\{O\}\) 在 \(y\)-坐标下的像（见结论 \ref{con:xi-terminal-zm-2division-norm-square}）。
\par
因此对任意 \(P\in E_0(\overline{\QQ})\) 均有 \(\Phi_{2}(y(P),y([2]P))=0\)，且对一般 \(y_{0}\in\overline{\QQ}\)，方程 \(\Phi_{2}(y_{0},y_{2})=0\) 在 \(y_{2}\) 中恰为四次，从而给出倍点对应在重量坐标下的四值性。
\par
在两个有理分歧值 \(y\in\{0,1\}\) 上，四次对应在第二变量中发生退化分解：
\[
\Phi_{2}(0,y_{2})=(y_{2}-50)(y_{2}-6)^{2}(y_{2}-1),\qquad
\Phi_{2}(1,y_{2})=(y_{2}-21)^{2}(y_{2}-1)(625y_{2}-96).
\]
因此在纤维 \(y=0\) 上，倍点像的重量值（计重数）构成多重集合 \(\{50,6,6,1\}\)；在纤维 \(y=1\) 上则为 \(\{21,21,1,\dfrac{96}{625}\}\)。
平方因子对应于纤维型 \((2,1,1)\) 中唯一分歧点在倍点推前下诱导的根并合。
\par
进一步，若以谱四次平面曲线
\[
\Pi(X,y):=X^{4}-X^{3}-(2y+1)X^{2}+X+y(y+1)
\]
表征 \(\QQ(E_0)/\QQ(y)\)（结论 \ref{con:xi-terminal-zm-elliptic-normalization}），则 \(\Phi_{2}\) 可由对 \(X\) 的纯代数消元得到，并满足结果式分解恒等式
\par
并且由 \(y=X^{2}+Y-\tfrac12\) 得
\(
2y+1-2X^{2}=2Y
\)，
从而在同一首一表示 \(\Pi(X,y)=0\) 下存在严格消元恒等式
\[
\boxed{\;
\mathrm{Res}_{X}\bigl(\Pi(X,y),\,2y+1-2X^{2}\bigr)
=-\bigl(16y^{3}-8y^{2}-4y+1\bigr)\;}
\]
因此三次多项式 \(16y^{3}-8y^{2}-4y+1\) 的零点集合恰为 \(y\bigl(E[2]\setminus\{O\}\bigr)\)。
等价地，在次数 \(4\) 的函数域扩张 \(\QQ(E)/\QQ(y)\) 中满足
\[
\operatorname{N}_{\QQ(E)/\QQ(y)}(2Y)=-(16y^{3}-8y^{2}-4y+1),
\qquad
\operatorname{N}_{\QQ(E)/\QQ(y)}(Y)=-\frac{16y^{3}-8y^{2}-4y+1}{16}.
\]
并且其判别式满足 \(\mathrm{Disc}(16y^{3}-8y^{2}-4y+1)=2^{8}\cdot 37\)，与 \(\mathrm{Disc}(P_{\mathrm{LY}})=-3^{9}\cdot 31^{2}\cdot 37\) 共享素因子 \(37\)（结论 \ref{con:xi-terminal-zm-leyang-badprime-37-common-double-root}），从而对任意 \(p\neq 37\) 两者在 \(\FF_{p}\) 上均可分离。
\[
\mathrm{Res}_{X}\bigl(\Pi(X,y),\,T-y([2]P)\bigr)
=-16\cdot\bigl(16y^{3}-8y^{2}-4y+1\bigr)^{3}\cdot \Phi_{2}(y,T).
\]
\par
此外，在谱平面曲线 \(\Pi(\lambda,y)=0\) 的函数域中，重量倍点像 \(y([2]P)\) 可写为关于谱坐标 \(\lambda=X\) 与重量 \(y\) 的显式有理提升：
\input{sections/generated/eq_fold_zm_elliptic_weight_doubling_psi}
\end{conclusion}
\begin{proof}[证明要点]
在模型 \(E:\ Y^{2}=X^{3}-X+\tfrac14\) 与重量函数 \(y=X^{2}+Y-\tfrac12\) 下，由倍点公式把 \([2](X,Y)=(X_{2},Y_{2})\) 写成 \(X,Y\) 的有理函数；
进而把 \(y_{2}=X_{2}^{2}+Y_{2}-\tfrac12\) 写成 \(\QQ(X)\)-线性函数于 \(\{1,y\}\)。
在 \(\Pi(X,y)=0\) 下化简并对 \(X\) 作消元，即得平面消元方程与结果式分解；不可约性与系数核验可由脚本 \texttt{scripts/exp\_fold\_zm\_elliptic\_weight\_doubling\_audit.py} 复核。证毕。
\end{proof}

\begin{conclusion}[特殊值 \(0\) 与 \(1\) 的全逆像范数分解：\texorpdfstring{$\operatorname{N}(y_2)$}{N(y2)} 与 \texorpdfstring{$\operatorname{N}(1-y_2)$}{N(1-y2)} 的整因子结构]\label{con:xi-terminal-zm-elliptic-weight-doubling-norms-0-1}
沿用结论 \ref{con:xi-terminal-zm-elliptic-weight-doubling-modular-equation} 的记号，并令
\[
k:=\QQ(y),\qquad K:=\QQ(E).
\]
记
\[
Q(y):=16y^{3}-8y^{2}-4y+1,
\qquad
A_{4}(y)=Q(y)^{4}.
\]
则在次数 \(4\) 的函数域扩张 \(K/k\) 中有严格相对范数恒等式：
\[
\boxed{\;
\operatorname{N}_{K/k}(y_{2})
=\frac{\Phi_{2}(y,0)}{A_{4}(y)}
=\frac{(y+1)^{2}(y^{3}-7y^{2}+13y-5)^{2}(y^{4}-6y^{3}+66y^{2}+28y+9)(y^{4}+2y^{3}+8y^{2}+8y+8)}{Q(y)^{4}}\; }.
\]
并且
\[
\boxed{\;
\operatorname{N}_{K/k}(1-y_{2})
=\frac{\Phi_{2}(y,1)}{A_{4}(y)}
=\frac{y(y-1)(y^{3}-37y^{2}+3y+25)(y^{3}-14y^{2}-14y+2)(y^{4}-14y^{3}+46y^{2}+22y+14)^{2}}{Q(y)^{4}}\; }.
\]
\end{conclusion}
\begin{proof}[证明要点]
由结论 \ref{con:xi-terminal-zm-elliptic-weight-doubling-modular-equation}，\(y_{2}\) 在 \(k\) 上满足四次关系 \(\Phi_{2}(y,y_{2})=0\)，其首项系数为 \(A_{4}(y)\)。
因此 \(y_{2}\) 的首一极小多项式可取为 \(A_{4}(y)^{-1}\Phi_{2}(y,T)\)，从而 \(\operatorname{N}_{K/k}(y_{2})=\Phi_{2}(y,0)/A_{4}(y)\)。
对 \(1-y_{2}\) 同理取 \(T=1\) 得 \(\operatorname{N}_{K/k}(1-y_{2})=\Phi_{2}(y,1)/A_{4}(y)\)。
所示因式分解为 \(\ZZ[y]\) 中的直接分解计算。证毕。
\end{proof}

\begin{conclusion}[倍乘模方程的判别式平方类与 Lee--Yang 分歧核完全一致]\label{con:xi-terminal-zm-elliptic-weight-doubling-disc-squareclass}
把结论 \ref{con:xi-terminal-zm-elliptic-weight-doubling-modular-equation} 的 \(\Phi_{2}(y,T)\) 视为关于 \(T\) 的四次多项式，则其判别式在 \(\QQ(y)^{\times}/(\QQ(y)^{\times})^{2}\) 中的平方类满足
\[
\mathrm{Disc}_{T}\bigl(\Phi_{2}(y,\cdot)\bigr)\equiv -y(y-1)P_{\mathrm{LY}}(y)\pmod{(\QQ(y)^{\times})^{2}},
\]
其中 \(P_{\mathrm{LY}}(y)=256y^{3}+411y^{2}+165y+32\) 为 Lee--Yang 三次因子（结论 \ref{con:xi-terminal-zm-leyang-ramification-critical-values}）。
更强地，存在 \(A(y)\in\ZZ[y]\)（\(\deg A=12\)）与 \(B_{+}(y),B_{-}(y)\in\ZZ[y]\)（\(\deg B_{\pm}=26\)），并记
\[
S_{12}(y):=A(y),
\qquad
S_{26}(y):=B_{+}(y),
\qquad
\widetilde S_{26}(y):=B_{-}(y).
\]
则在 \(\ZZ[y]\) 中严格成立判别式分解
\[
\boxed{\;
\mathrm{Disc}_{T}\bigl(\Phi_{2}(y,\cdot)\bigr)=-y(y-1)P_{\mathrm{LY}}(y)\cdot S_{12}(y)^{2}\cdot S_{26}(y)^{2}\; }.
\]
其中 \(A,B_{+},B_{-}\) 的显式展开见下列闭式证书；同一证书亦表明共轭重量支下的四次对应在判别式层面共享相同的平方自由核：
\input{sections/generated/eq_fold_zm_elliptic_weight_doubling_discriminant}
因此倍乘诱导的重量四次模方程之唯一二次分辨率子域由
\[
\QQ(y)\Bigl(\sqrt{-y(y-1)P_{\mathrm{LY}}(y)}\Bigr)
\]
锁定；Lee--Yang 三次因子在平方类意义下对倍乘动力学保持刚性不变。
\end{conclusion}
\begin{proof}[证明要点]
按定义计算 \(\mathrm{Disc}_{T}(\Phi_2)\) 并在 \(\ZZ[y]\) 中因式分解即可得到所示恒等式；平方类结论随即推出。
可复现核验脚本：\texttt{scripts/exp\_fold\_zm\_elliptic\_weight\_doubling\_audit.py}。证毕。
\end{proof}

\begin{conclusion}[原始元的判别式平方自由类不变性：二次分辨率在二倍重整化塔中恒定]\label{con:xi-terminal-zm-discriminant-squarefree-class-primitive-element}
令 \(k:=\QQ(y)\) 且 \(K:=\QQ(E)\)。
则对任意满足 \([k(\alpha):k]=4\) 的原始元 \(\alpha\in K\)，其首一最小多项式 \(m_{\alpha}(T)\in k[T]\) 的判别式在平方类意义下满足
\[
\boxed{\;
\mathrm{Disc}_{T}\bigl(m_{\alpha}(T)\bigr)\equiv -y(y-1)P_{\mathrm{LY}}(y)\pmod{(k^{\times})^{2}}\; }.
\]
特别地，对每个整数 \(n\ge 1\)，令 \(y_{n}:=y\circ[2^{n}]\in K\)；
若 \([k(y_{n}):k]=4\)，则其最小多项式判别式同样满足
\[
\mathrm{Disc}_{T}\bigl(m_{y_{n}}(T)\bigr)\equiv -y(y-1)P_{\mathrm{LY}}(y)\pmod{(k^{\times})^{2}}.
\]
从而 Lee--Yang 判别式核 \(-y(y-1)P_{\mathrm{LY}}(y)\) 在任意二倍迭代层级上均不给出新的二次判别式数据。
\end{conclusion}
\begin{proof}[证明要点]
由于 \(K/k\) 的伽罗瓦闭包之唯一二次分辨率子域由 \(\sqrt{\mathrm{Disc}_{\lambda}(\Pi)}\) 生成（结论 \ref{con:xi-terminal-zm-generic-galois-s4} 与结论 \ref{con:xi-terminal-zm-leyang-discriminant-norm-identity}），并且
\(
\mathrm{Disc}_{\lambda}(\Pi)=-y(y-1)P_{\mathrm{LY}}(y)
\)
（结论 \ref{con:xi-terminal-zm-discriminant-leyang}），故该二次子域恒为
\(
k\bigl(\sqrt{-y(y-1)P_{\mathrm{LY}}(y)}\bigr)
\)。
另一方面，对任意原始元 \(\alpha\)，\(\sqrt{\mathrm{Disc}_{T}(m_{\alpha})}\) 在分裂域上实现同一符号表示的生成元，因而生成相同的二次分辨率子域；于是 \(\mathrm{Disc}_{T}(m_{\alpha})\) 与 \(-y(y-1)P_{\mathrm{LY}}(y)\) 在 \(k^{\times}/(k^{\times})^{2}\) 中同类。
对 \(\alpha=y_{n}\) 的情形仅需附加假设 \([k(y_{n}):k]=4\)（即 \(y_{n}\) 为原始元），结论随即得到。证毕。
\end{proof}

\begin{conclusion}[\texorpdfstring{$2$}{2}-扭平移不变纤维的参数约束：平方型碰撞因子与判别式平方因子对齐]\label{con:xi-terminal-zm-elliptic-weight-doubling-2torsion-translation-collision-parameters}
令 \(E[2]\subset E(\overline{\QQ})\) 为二阶扭点子群，并定义参数集合
\[
\Lambda_{2}:=\Bigl\{\,y_{0}\in\overline{\QQ}:\ \exists\,P\in E(\overline{\QQ}),\ \exists\,T\in E[2]\setminus\{O\}\ \text{s.t.}\ y(P)=y(P+T)=y_{0}\,\Bigr\}.
\]
记
\[
Q(y):=16y^{3}-8y^{2}-4y+1,
\]
并令 \(A(y)\in\ZZ[y]\) 为结论 \ref{con:xi-terminal-zm-elliptic-weight-doubling-disc-squareclass} 的判别式分解中出现的次数 \(12\) 因子（\(\mathrm{Disc}_{T}(\Phi_{2})\) 的平方因子之一）。
则 \(\Lambda_{2}\) 的 Zariski 闭包满足
\[
\Lambda_{2}\subseteq \{\,y\in\overline{\QQ}:\ Q(y)\,A(y)=0\,\}.
\]
其中 \(Q(y)=0\) 的零点集合恰为 \(y\bigl(E[2]\setminus\{O\}\bigr)\)，而其余的平移碰撞参数被 \(A(y)=0\) 约束；
并且 \(A(y)\) 在 \(\mathrm{Disc}_{T}(\Phi_{2})\) 中以平方方式出现，与平移碰撞仅贡献偶分歧的机制一致。
\end{conclusion}
\begin{proof}[证明要点]
取短 Weierstrass 形式下的非平凡二扭点 \(T=(a,0)\)，满足 \(4a^{3}-4a+1=0\)；
将加法公式写出 \(P+T\) 并把条件 \(y(P)=y(P+T)\) 视为关于 \((X,a,y)\) 的代数系统，随后对 \(X\) 与 \(a\) 作消元即可得到 \(y\) 的必要约束多项式。
该消元多项式的平方自由部分严格等于 \(Q(y)A(y)\)，从而得到所示包含关系；可复现核验脚本：\texttt{scripts/exp\_fold\_zm\_elliptic\_weight\_doubling\_2torsion\_translation\_collision\_audit.py}。证毕。
\end{proof}

\begin{conclusion}[倍乘模方程的泛型伽罗瓦群为 \texorpdfstring{$S_{4}$}{S4}]\label{con:xi-terminal-zm-elliptic-weight-doubling-galois-s4}
设 \(K:=\QQ(y)\)，令 \(L\) 为结论 \ref{con:xi-terminal-zm-elliptic-weight-doubling-modular-equation} 的四次多项式 \(\Phi_{2}(y,T)\) 在 \(K\) 上的分裂域。
则泛型情形下
\[
\mathrm{Gal}(L/K)\cong S_{4}.
\]
\end{conclusion}
\begin{proof}[证明要点]
由结论 \ref{con:xi-terminal-zm-elliptic-weight-doubling-modular-equation}，\(\Phi_{2}(y,\cdot)\) 在 \(K\) 上不可约；
而结论 \ref{con:xi-terminal-zm-elliptic-weight-doubling-disc-squareclass} 给出 \(\mathrm{Disc}_{T}(\Phi_{2})\) 的平方类为 \(-y(y-1)P_{\mathrm{LY}}(y)\)，因而在 \(K\) 中非平方，排除 \(\mathrm{Gal}\subset A_{4}\)。
再对 \(y=2\) 作一次专门化，得到
\[
\Phi_{2}(2,T)=62742241T^{4}-225603363T^{3}+534557358T^{2}-28061748T+235224\in\QQ[T],
\]
其专门化伽罗瓦群可由可审计算法判定为 \(S_{4}\)；
由 Hilbert 不可约性定理，泛型伽罗瓦群亦为 \(S_{4}\)。
可复现核验脚本：\texttt{scripts/exp\_fold\_zm\_elliptic\_weight\_doubling\_audit.py}。证毕。
\end{proof}

\begin{conclusion}[用 \texorpdfstring{$(y,y\circ[2])$}{(y,y∘[2])} 重构函数域：倍乘模曲线的归一化双有理同构于 \(E\)]\label{con:xi-terminal-zm-elliptic-weight-doubling-function-field-reconstruction}
设 \(C_{2}\subset\mathbb P^{1}_{y}\times\mathbb P^{1}_{y_{2}}\) 为由方程 \(\Phi_{2}(y,y_{2})=0\) 定义的平面代数曲线。
则：
\begin{enumerate}
  \item \(\QQ(E)=\QQ(y,y_{2})\)，即 \(y\) 与 \(y_{2}=y\circ[2]\) 共同生成 \(E\) 的函数域；
  \item \(C_{2}\) 的归一化与 \(E\) 双有理同构；等价地，态射
  \[
  E\longrightarrow \mathbb P^{1}_{y}\times\mathbb P^{1}_{y_{2}},\qquad
  P\longmapsto\bigl(y(P),y([2]P)\bigr)
  \]
  在几何泛型点处为一一对应。
\end{enumerate}
\end{conclusion}
\begin{proof}[证明要点]
由结论 \ref{con:xi-terminal-zm-elliptic-weight-doubling-modular-equation}，\(y([2]P)\) 在 \(\QQ(y)\) 上满足次数为 \(4\) 的代数关系 \(\Phi_{2}(y,y_{2})=0\)，
故 \([\QQ(y,y_{2}):\QQ(y)]=4=[\QQ(E):\QQ(y)]\)，从而 \(\QQ(E)=\QQ(y,y_{2})\)。
双有理同构可由显式逆映射证书（重构 \(X,Y\)）给出，见脚本 \texttt{scripts/exp\_fold\_zm\_elliptic\_weight\_doubling\_audit.py} 的输出。证毕。
\end{proof}

\begin{theorem}[Lee--Yang 权函数同态对应曲线的双次数与奇点总 \texorpdfstring{$\delta$}{delta}-不变量]\label{thm:xi-terminal-zm-elliptic-weight-n-correspondence-bidegree-delta}
\leanverified{paper\_xi\_terminal\_zm\_elliptic\_weight\_n\_correspondence\_bidegree\_delta}
沿用结论 \ref{con:xi-terminal-zm-elliptic-weight-doubling-modular-equation} 的记号，令
\[
E_0:\ \eta^2+\eta=x^3-x,\qquad
y:=x^2-\eta-1.
\]
在定理 \ref{thm:xi-terminal-zm-leyang-elliptic-structure} 的记号下，\(E_0\) 与 \(E_{\min}\) 为同一同构类，以下不再区分。
对任意整数 \(n\ge 2\)，记 \([n]:E_0\to E_0\) 为乘法同态，并定义
\[
y^{(n)}:=y\circ [n].
\]
令
\[
C_n:=\overline{\{(y(P),y^{(n)}(P)):\ P\in E_0(\overline{\QQ})\}}\subset \PP^1_y\times\PP^1_{y^{(n)}}.
\]
则：
\begin{enumerate}
  \item \(C_n\) 不可约，且其正规化与 \(E_0\) 双有理同构；
  \item 在 \(\PP^1\times\PP^1\) 中，\(C_n\) 的双次数为
  \[
  \deg_{(1,0)}(C_n)=4n^2,\qquad \deg_{(0,1)}(C_n)=4;
  \]
  \item \(C_n\) 的算术亏格满足
  \[
  p_a(C_n)=(4n^2-1)(4-1)=3(4n^2-1),
  \]
  且
  \[
  \sum_{x\in \mathrm{Sing}(C_n)}\delta_x
  =p_a(C_n)-g(E_0)
  =12n^2-4.
  \]
\end{enumerate}
\end{theorem}
\begin{proof}
映射
\[
\Phi_n:E_0\to \PP^1_y\times\PP^1_{y^{(n)}},\qquad P\mapsto (y(P),y^{(n)}(P))
\]
的像闭包即 \(C_n\)；因 \(E_0\) 不可约，故 \(C_n\) 不可约。
设 \(\nu:C_n^\nu\to C_n\) 为正规化，并把 \(\Phi_n\) 分解为
\[
E_0\xrightarrow{\psi} C_n^\nu \xrightarrow{\nu} C_n.
\]
记 \(m:=\deg(\psi)\)。
令 \(p_1,p_2\) 为两坐标投影，则在函数域层面有
\[
y=(p_1\circ\nu)\circ\psi,\qquad y^{(n)}=(p_2\circ\nu)\circ\psi.
\]
故
\[
4=\deg(y)=m\cdot \deg(p_1\circ\nu),\qquad
4n^2=\deg(y^{(n)})=m\cdot \deg(p_2\circ\nu),
\]
从而 \(m\mid 4\)。

若 \(m=2\)，则 \(y\) 分解为次数 \(2\) 与次数 \(2\) 的复合，和推论 \ref{cor:xi-terminal-zm-leyang-cover-primitive} 的本原性矛盾。
若 \(m=4\)，则 \(\deg(p_1\circ\nu)=1\)，故 \(k(C_n^\nu)=\QQ(y)\)，从而 \(y^{(n)}\in\QQ(y)\)。
但
\[
(y)_\infty=4O,\qquad
\bigl(y^{(n)}\bigr)_\infty=[n]^*(y)_\infty
=4\sum_{Q\in E_0[n]}[Q],
\]
其极点支撑含 \(n^2\) 个点（\(n\ge2\)）。
而任意 \(\QQ(y)\) 中函数在 \(E_0\) 上的极点只能落在 \(y^{-1}(\infty)=\{O\}\)，矛盾。
故 \(m=1\)，于是 \(\psi\) 双有理，第一条成立，并得到
\[
\deg(p_1\circ\nu)=4,\qquad
\deg(p_2\circ\nu)=4n^2.
\]
按本节双次数记号，即
\[
\deg_{(0,1)}(C_n)=4,\qquad
\deg_{(1,0)}(C_n)=4n^2.
\]

对 \((a,b)\)-型不可约曲线（此处 \(a=4n^2,b=4\)），算术亏格为 \((a-1)(b-1)\)，故
\[
p_a(C_n)=3(4n^2-1).
\]
由 \(g(C_n^\nu)=g(E_0)=1\) 与恒等式
\[
p_a(C_n)=g(C_n^\nu)+\sum_{x\in\mathrm{Sing}(C_n)}\delta_x
\]
即得
\[
\sum_{x\in\mathrm{Sing}(C_n)}\delta_x=12n^2-4.
\]
证毕。
\end{proof}

\begin{corollary}[二幂层级的闭式 \texorpdfstring{$\delta$}{delta}-账本]\label{cor:xi-terminal-zm-elliptic-weight-2k-delta-ledger}
\leanverified{paper\_xi\_terminal\_zm\_elliptic\_weight\_2k\_delta\_ledger}
在定理 \ref{thm:xi-terminal-zm-elliptic-weight-n-correspondence-bidegree-delta} 的设定下，取 \(n=2^k\)（\(k\ge1\)），则
\[
\sum_{x\in \mathrm{Sing}(C_{2^k})}\delta_x=12\cdot 4^k-4=3\cdot 4^{k+1}-4.
\]
\end{corollary}
\begin{proof}
把 \(n=2^k\) 代入定理 \ref{thm:xi-terminal-zm-elliptic-weight-n-correspondence-bidegree-delta} 的闭式
\(
\sum\delta_x=12n^2-4
\)
即可。证毕。
\end{proof}

\paragraph{重量坐标下的三倍乘对应：四值模方程与 Lee--Yang 判别式核的层级保持}
沿用结论 \ref{con:xi-terminal-zm-elliptic-quadratic-discriminant} 的椭圆曲线 \(\mathcal{E}\) 与重量函数
\[
y:=\lambda^{2}+Y-\tfrac12\in\QQ(\mathcal{E}),
\qquad
\mathcal{E}:\ Y^{2}=\lambda^{3}-\lambda+\tfrac14.
\]
令 \(y_{3}:=y\circ[3]\in\QQ(\mathcal{E})\)。
由于 \(y:\mathcal{E}\to\PP^{1}_{y}\) 为次数 \(4\) 的覆盖，三倍乘 \([3]\) 在重量线上同样诱导一个四值对应；
其平面消元闭包可写为一条双次数 \((36,4)\) 的不可约对应曲线。

\begin{conclusion}[三倍乘在重量线上诱导的四次模方程：\texorpdfstring{$(36,4)$}{(36,4)}-对应与 \texorpdfstring{$T(y)^{4}$}{T(y)^4} 极点簇的必然剥离]\label{con:xi-terminal-zm-elliptic-weight-tripling-modular-equation}
存在原始整系数多项式
\[
M_{3}(y,y_{3})\in\ZZ[y,y_{3}],
\qquad
\deg_{y_{3}}M_{3}=4,\ \deg_{y}M_{3}=36,
\]
使得对一切 \(P\in\mathcal{E}(\overline{\QQ})\) 恒有 \(M_{3}\bigl(y(P),y([3]P)\bigr)=0\)。
并且可取表示
\[
M_{3}(y,y_{3})=d_{4}(y)y_{3}^{4}+d_{3}(y)y_{3}^{3}+d_{2}(y)y_{3}^{2}+d_{1}(y)y_{3}+d_{0}(y),
\qquad d_{j}(y)\in\ZZ[y],
\]
其显式系数证书为：
\input{sections/generated/eq_fold_zm_elliptic_weight_tripling_M3}
\par
并且其首项系数为完全四次幂
\[
d_{4}(y)=L_{3}(y)^{4},
\]
其中 \(L_{3}(y)\in\ZZ[y]\) 为 \(\mathcal{E}[3]\setminus\{O\}\) 在重量投影下的八次消元多项式（结论 \ref{con:xi-terminal-zm-elliptic-3torsion-y-image-octic}）。
\par
其常数项亦在 \(\ZZ[y]\) 中完全分解：
\[
\boxed{\;
\begin{aligned}
d_{0}(y)=\;&y\,(y^{8}-39y^{7}+1461y^{6}+1422y^{5}+1890y^{4}+2019y^{3}-93y^{2}-627y-147)\\
&\cdot (y^{9}-39y^{8}+237y^{7}+762y^{6}-498y^{5}-141y^{4}+3495y^{3}+1725y^{2}+309y-46)^{2}\\
&\cdot (y^{9}+15y^{8}+201y^{7}+786y^{6}+2928y^{5}+1206y^{4}-2583y^{3}-1431y^{2}-342y-45).
\end{aligned}
\; }.
\]
\par
令
\[
T(y):=16y^{3}-8y^{2}-4y+1,
\]
则对消去构造所得的全结果式可作结构性分解
\[
\operatorname{Res}_{\lambda}\Bigl(\Pi(\lambda,y),\ y_{3}-y([3]P)\Bigr)
=16\cdot T(y)^{4}\cdot M_{3}(y,y_{3}),
\]
其中 \(\Pi(\lambda,y)=\lambda^{4}-\lambda^{3}-(2y+1)\lambda^{2}+\lambda+y(y+1)\) 为结论 \ref{con:xi-terminal-zm-elliptic-quadratic-discriminant} 的谱四次方程。
此处 \(T(y)^{4}\) 的出现编码了 \([3]\) 的中间分式表达在 \(\mathcal{E}[2]\) 上引入的极点簇；而 \(M_{3}\) 为剥离该极点簇后所剩的唯一动力学本体对应。
\end{conclusion}
\begin{proof}[证明要点]
将 \([3]\) 的有理公式与 \(y=\lambda^{2}+Y-\tfrac12\) 联立，并对 \((\lambda,Y)\) 作消元可得到结果式；
其在 \(\ZZ[y,y_{3}]\) 中的本原部分给出 \(M_{3}\)，而由 \(y\) 在 \(\mathcal{E}[2]\) 上的像三次 \(T(y)=0\)（结论 \ref{con:xi-terminal-zm-2division-norm-square}）可解释极点簇的 \(T(y)^{4}\) 剥离项。
可复现核验脚本：\texttt{scripts/exp\_fold\_zm\_elliptic\_weight\_tripling\_audit.py}。证毕。
\end{proof}

\begin{conclusion}[三级模方程的判别式平方类：平方自由核仍由 \texorpdfstring{$-y(y-1)P_{\mathrm{LY}}(y)$}{-y(y-1)P} 唯一锁定]\label{con:xi-terminal-zm-elliptic-weight-tripling-disc-squareclass}
把结论 \ref{con:xi-terminal-zm-elliptic-weight-tripling-modular-equation} 的 \(M_{3}(y,y_{3})\) 视为关于 \(y_{3}\) 的四次多项式，则其判别式在 \(\ZZ[y]\) 中发生严格平方分解：
\[
\mathrm{Disc}_{y_{3}}\bigl(M_{3}(y,\cdot)\bigr)
=-y(y-1)P_{\mathrm{LY}}(y)\cdot A_{32}(y)^{2}\cdot B_{66}(y)^{2},
\]
其中 \(P_{\mathrm{LY}}(y)=256y^{3}+411y^{2}+165y+32\) 为 Lee--Yang 三次因子（结论 \ref{con:xi-terminal-zm-discriminant-leyang}），并且 \(A_{32}(y)\in\ZZ[y]\)、\(B_{66}(y)\in\ZZ[y]\) 分别为次数 \(32\) 与次数 \(66\) 的整式。
因而在 \(\QQ(y)^{\times}/(\QQ(y)^{\times})^{2}\) 中有平方类恒等式
\[
\mathrm{Disc}_{y_{3}}\bigl(M_{3}(y,\cdot)\bigr)\equiv -y(y-1)P_{\mathrm{LY}}(y).
\]
其显式分解证书为：
\input{sections/generated/eq_fold_zm_elliptic_weight_tripling_discriminant}
\end{conclusion}
\begin{proof}[证明要点]
按定义计算四次判别式并在 \(\ZZ[y]\) 中分解即可得到所示恒等式；平方类陈述随即推出。
可复现核验脚本：\texttt{scripts/exp\_fold\_zm\_elliptic\_weight\_tripling\_audit.py}。证毕。
\end{proof}

\begin{conclusion}[原始生成元的判别式平方类刚性：Lee--Yang 二次脊柱在任一四次生成元下保持]\label{con:xi-terminal-zm-elliptic-primitive-generator-discriminant-squareclass}
设 \(K:=\QQ(\mathcal{E})\)、\(k:=\QQ(y)\)，则 \(K/k\) 为次数 \(4\) 的可分扩张。
对任意原始生成元 \(\theta\in K\)（即 \([k(\theta):k]=4\)），记 \(m_{\theta}(T)\in k[T]\) 为其首一极小多项式，则
\[
\mathrm{Disc}\bigl(m_{\theta}\bigr)\equiv -y(y-1)P_{\mathrm{LY}}(y)\pmod{k^{\times 2}}.
\]
特别地，\(\theta=y\circ[2]\)（结论 \ref{con:xi-terminal-zm-elliptic-weight-doubling-disc-squareclass}）与 \(\theta=y\circ[3]\)（结论 \ref{con:xi-terminal-zm-elliptic-weight-tripling-disc-squareclass}）均为原始生成元，且二者的判别式平方类与谱四次 \(\Pi(\lambda,y)\) 的 Lee--Yang 判别式核完全一致。
因此，四次覆盖的 \(S_{4}\)-伽罗瓦闭包所对应的唯一二次分辨率子域在所有原始生成元下均可一致刻画为
\[
\QQ(y)\Bigl(\sqrt{-y(y-1)P_{\mathrm{LY}}(y)}\Bigr).
\]
\end{conclusion}
\begin{proof}[证明要点]
扩张 \(K/k\) 的伽罗瓦闭包为 \(S_{4}\)-扩张（结论 \ref{con:xi-terminal-zm-generic-galois-s4}），其指数 \(2\) 的唯一正规子群为 \(A_{4}\)，对应的二次子域由任意原始生成元之极小多项式判别式的平方类给出。
由于闭包与其二次子域均与生成元选择无关，故 \(\mathrm{Disc}(m_{\theta})\) 在 \(k^{\times}/k^{\times 2}\) 中的平方类为不变量；
再由 \(\Pi(\lambda,y)\) 的判别式恒等式 \(\mathrm{Disc}_{\lambda}(\Pi)=-y(y-1)P_{\mathrm{LY}}(y)\)（结论 \ref{con:xi-terminal-zm-discriminant-leyang}）确定该不变量即得。
证毕。
\end{proof}

\input{sections/body/zeta_finite_part/xi/thm__xi-terminal-zm-elliptic-weight-tripling-c3-singularities}


\begin{conclusion}[\texorpdfstring{$(4)$}{(4)}-division 的重量像：临界集合的 \(y\)-值由不可约十二次多项式支配]\label{con:xi-terminal-zm-4division-weight-image}
在同一椭圆曲线模型
\[
E:\ Y^{2}=X^{3}-X+\frac14,
\qquad
y:=X^{2}+Y-\frac12
\]
及其诱导的谱四次平面曲线
\[
\Pi(X,y):=X^{4}-X^{3}-(2y+1)X^{2}+X+y(y+1)=0
\]
的背景下，令
\[
N_{6}(X):=2X^{6}-10X^{4}+10X^{3}-10X^{2}+2X+1
\]
（其零点给出 \(x\bigl(E[4]\setminus E[2]\bigr)\)，等价于结论 \ref{con:xi-terminal-zm-lattes-4division-certificate} 的 \(4\)-division 证书多项式 \(\psi_{4}/(2Y)\) 在 \(X\)-坐标上的表示）。
定义
\[
T_{4}(y):=\operatorname{Res}_{X}\bigl(N_{6}(X),\Pi(X,y)\bigr)\in\ZZ[y].
\]
则 \(T_{4}\) 为不可约的十二次整系数多项式，并且
\[
y\bigl(E[4]\setminus E[2]\bigr)=\{\,u\in\overline{\QQ}:\ T_{4}(u)=0\,\}.
\]
此外，\(T_{4}\) 与其判别式具有强烈的算术刚性；其展开式与判别式素分解如下：
\input{sections/generated/eq_fold_zm_elliptic_4division_weight_image_audit}
特别地，\(\operatorname{Disc}(T_{4})\) 中素数 \(37\) 的指数为 \(23\)，且其平方自由部分等于 \(37\)；从而该重量像所张成的数域在 \(37\) 处发生高阶分歧。
\end{conclusion}
\begin{proof}[证明要点]
由 \(P\in E[4]\setminus E[2]\) 当且仅当 \([2]P\in E[2]\setminus\{O\}\)，可知其 \(X\)-坐标满足 \(4\)-division 证书六次 \(N_{6}(X)=0\)。
另一方面，由 \(y=X^{2}+Y-\tfrac12\) 消去 \(Y\) 得 \(\Pi(X,y)=0\)（结论 \ref{con:xi-terminal-zm-elliptic-normalization}）。
因此 \(y(E[4]\setminus E[2])\) 正是消去 \(X\) 后的参数像，结果式 \(\operatorname{Res}_{X}(N_{6},\Pi)\) 给出其最小整系数消元证书；
不可约性与判别式分解可由脚本 \texttt{scripts/exp\_fold\_zm\_elliptic\_4division\_weight\_image\_audit.py} 复核。证毕。
\end{proof}


\paragraph{四次判别脊柱、Lee--Yang 三次分支层与 \texorpdfstring{$4$}{4}-division 重量像的代数耦合}
在四次覆盖的原始生成元判别式平方类刚性、Lee--Yang 分歧支平方根系数的三次数域刻画以及 \(4\)-division 重量像的十二次不可约消元证书均已建立之后，还可进一步抽出下列四条后果。它们分别刻画二次分辨率对子任意原始坐标的严格不变性、三次分支层与四次二次脊柱在 \(k=\QQ(y)\) 上的线性无交、素数 \(37\) 在两端诱导的相反符号二次影，以及 \(4\)-division 十二个重量值的单轨道刚性。

\begin{theorem}[四次 Lee--Yang 覆盖的二次分辨率对一切原始生成元严格不变]\label{thm:xi-terminal-zm-leyang-quadratic-resolvent-primitive-generator-invariance}
沿用结论 \ref{con:xi-terminal-zm-elliptic-primitive-generator-discriminant-squareclass} 的记号，记
\[
K:=\QQ(E),\qquad k:=\QQ(y),
\]
并把 \(K/k\) 视为次数 \(4\) 的可分扩张。对任意原始生成元 \(\theta\in K\)（即 \([k(\theta):k]=4\)），记 \(m_\theta(T)\in k[T]\) 为其首一极小多项式。则对任意两原始生成元 \(\theta_1,\theta_2\in K\) 都有
\[
\frac{\mathrm{Disc}(m_{\theta_1})}{\mathrm{Disc}(m_{\theta_2})}\in k^{\times 2}.
\]
特别地，由任意原始生成元恢复出的二次分辨率子域都相同，并统一等于
\[
k^{(2)}:=k\!\left(\sqrt{-y(y-1)P_{\mathrm{LY}}(y)}\right).
\]
\end{theorem}
\begin{proof}
结论 \ref{con:xi-terminal-zm-elliptic-primitive-generator-discriminant-squareclass} 已表明：对任意原始生成元 \(\theta\in K\)，都有
\[
\mathrm{Disc}(m_\theta)\equiv -y(y-1)P_{\mathrm{LY}}(y)\pmod{k^{\times 2}}.
\]
因此对任意两原始生成元 \(\theta_1,\theta_2\) 有
\[
\mathrm{Disc}(m_{\theta_1})\equiv \mathrm{Disc}(m_{\theta_2})\pmod{k^{\times 2}},
\]
移项即得
\[
\frac{\mathrm{Disc}(m_{\theta_1})}{\mathrm{Disc}(m_{\theta_2})}\in k^{\times 2}.
\]
同一结论同时说明：该平方类所决定的二次子域正是四次覆盖 \(K/k\) 的唯一二次分辨率子域，故其统一闭式即为
\[
k^{(2)}=k\!\left(\sqrt{-y(y-1)P_{\mathrm{LY}}(y)}\right).
\]
证毕。
\end{proof}

\begin{theorem}[Lee--Yang 三次分支层与四次二次脊柱在 \texorpdfstring{$k=\QQ(y)$}{k=Q(y)} 上线性无交]\label{thm:xi-terminal-zm-leyang-cubic-layer-quadratic-spine-linear-disjoint-over-k}
沿用定理 \ref{thm:xi-terminal-zm-leyang-quadratic-resolvent-primitive-generator-invariance} 的记号，并令
\[
u:=\kappa_\ast^2
\]
为定理 \ref{thm:xi-terminal-zm-kappa-square-cubic-field-s3} 中的 Lee--Yang 分歧平方根系数。记
\[
L_3:=k(u),\qquad L_2:=k^{(2)}.
\]
则有
\[
L_2\cap L_3=k,
\qquad
[L_2L_3:k]=6.
\]
\end{theorem}
\begin{proof}
由定理 \ref{thm:xi-terminal-zm-kappa-square-cubic-field-s3}，\(u\) 满足在 \(\QQ\) 上不可约的三次多项式
\[
324u^3-540u^2+333u-74=0.
\]
由于 \(k=\QQ(y)\) 是 \(\QQ\) 的纯超越扩张，故 \(k/\QQ\) 为 regular 扩张，从而与任意有限代数扩张线性无交。于是
\[
[L_3:k]=[\QQ(u):\QQ]=3.
\]
另一方面，由定理 \ref{thm:xi-terminal-zm-leyang-quadratic-resolvent-primitive-generator-invariance}，
\[
L_2=k^{(2)}=k\!\left(\sqrt{-y(y-1)P_{\mathrm{LY}}(y)}\right)
\]
是 \(K/k\) 的唯一二次分辨率子域，故 \([L_2:k]=2\)。

因此交域次数 \([L_2\cap L_3:k]\) 同时整除 \(2\) 与 \(3\)，只能等于 \(1\)。于是
\[
L_2\cap L_3=k.
\]
再由次数乘法公式，
\[
[L_2L_3:k]
=
\frac{[L_2:k][L_3:k]}{[L_2\cap L_3:k]}
=
\frac{2\cdot 3}{1}
=6.
\]
证毕。
\end{proof}

\begin{theorem}[素数 \texorpdfstring{$37$}{37} 在 \texorpdfstring{$4$}{4}-division 权重像与 Lee--Yang 三次分支层中诱导相反符号的二次影]\label{thm:xi-terminal-zm-prime37-opposite-sign-quadratic-shadows}
令 \(M_{T_4}\) 为 \(T_4(y)\) 的分裂域，并令 \(M_u\) 为定理 \ref{thm:xi-terminal-zm-kappa-square-cubic-field-s3} 中三次多项式
\[
324u^3-540u^2+333u-74
\]
的分裂域。则 \(M_{T_4}\) 所对应的符号二次子域与 \(M_u\) 的唯一二次子域分别为
\[
\QQ\!\left(\sqrt{\mathrm{Disc}(T_4)}\right)=\QQ(\sqrt{37}),
\qquad
M_u^{A_3}=\QQ(\sqrt{-37}).
\]
从而
\[
\QQ(\sqrt{37})\cap \QQ(\sqrt{-37})=\QQ,
\qquad
[\QQ(\sqrt{37},\sqrt{-37}):\QQ]=4.
\]
并且这两个二次域的域判别式分别为 \(148\) 与 \(-148\)，因而具有共同奇分歧素数 \(37\)，但在无穷位上符号相反：前者为实二次域，后者为虚二次域。
\end{theorem}
\begin{proof}
由结论 \ref{con:xi-terminal-zm-4division-weight-image} 的显式判别式证书，
\[
\mathrm{Disc}(T_4)=2^{56}\cdot 37^{23}\cdot 59^{6}\cdot(652225136287)^2,
\]
故其平方自由部分为 \(37\)。于是 \(T_4\) 分裂域所携带的符号二次子域为
\[
\QQ\!\left(\sqrt{\mathrm{Disc}(T_4)}\right)=\QQ(\sqrt{37}).
\]

另一方面，定理 \ref{thm:xi-terminal-zm-kappa-square-cubic-field-s3} 已给出
\[
\mathrm{Disc}(324u^3-540u^2+333u-74)=-2^{6}\cdot 3^{9}\cdot 37,
\]
其平方自由部分为 \(-37\)。又 \(M_u/\QQ\) 的 Galois 群同构于 \(S_3\)，故其唯一二次子域为
\[
M_u^{A_3}=\QQ(\sqrt{-37}).
\]

由于 \(\QQ(\sqrt{37})\neq \QQ(\sqrt{-37})\)，两个二次域的交只能为 \(\QQ\)。因此
\[
\QQ(\sqrt{37})\cap \QQ(\sqrt{-37})=\QQ,
\]
并由次数乘法公式得到
\[
[\QQ(\sqrt{37},\sqrt{-37}):\QQ]=4.
\]
最后，二次域 \(\QQ(\sqrt{37})\) 与 \(\QQ(\sqrt{-37})\) 的域判别式分别为 \(148\) 与 \(-148\)，故二者具有共同奇分歧素数 \(37\)，同时在无穷位上分别对应实与虚符号。证毕。
\end{proof}

\begin{theorem}[\texorpdfstring{$4$}{4}-division 权重像的十二个值构成单一 Galois 轨道并避开有理层]\label{thm:xi-terminal-zm-4division-weight-image-single-galois-orbit}
记
\[
\Omega_4:=y\bigl(E[4]\setminus E[2]\bigr)\subset\overline{\QQ}.
\]
则
\[
\Omega_4=\{\,u\in\overline{\QQ}:T_4(u)=0\,\},
\qquad
|\Omega_4|=12,
\qquad
\Omega_4\cap\QQ=\varnothing.
\]
此外，绝对 Galois 群 \(\Gal(\overline{\QQ}/\QQ)\) 在 \(\Omega_4\) 上的作用是传递的；等价地，\(\Omega_4\) 的十二个点构成单一 Galois 轨道。
\end{theorem}
\begin{proof}
结论 \ref{con:xi-terminal-zm-4division-weight-image} 已证明
\[
\Omega_4=\{\,u\in\overline{\QQ}:T_4(u)=0\,\},
\]
且 \(T_4\in\ZZ[y]\) 为次数 \(12\) 的不可约多项式。又由同一结论中的判别式公式可知 \(\mathrm{Disc}(T_4)\neq 0\)，故 \(T_4\) 有 \(12\) 个互异根，从而 \(|\Omega_4|=12\)。

若 \(\Omega_4\cap\QQ\neq\varnothing\)，则 \(T_4\) 存在有理根，与其在 \(\QQ\) 上不可约矛盾，故
\[
\Omega_4\cap\QQ=\varnothing.
\]
最后，不可约多项式的全部根在绝对 Galois 群作用下必构成单一轨道，故 \(\Gal(\overline{\QQ}/\QQ)\) 在 \(\Omega_4\) 上作用传递。证毕。
\end{proof}

因此，四次 Lee--Yang 覆盖的二次判别脊柱固定在
\[
k\!\left(\sqrt{-y(y-1)P_{\mathrm{LY}}(y)}\right),
\]
而 Lee--Yang 三次分支层作为一条三次常数扩张，只在次数 \(6\) 的复合层与之首次耦合。与此同时，在 \(\QQ\) 上由 \(4\)-division 权重像与分支三次判别式所切出的两条二次影
\[
\QQ(\sqrt{37}),\qquad \QQ(\sqrt{-37})
\]
分别承载实符号与虚符号；全部 \(4\)-division 重量值则作为单一绝对 Galois 轨道，构成该代数壳层的不可再分分量。


\begin{conclusion}[所有分歧点处的 Puiseux 统一化：平方根临界指数与首项系数的显式公式]\label{con:xi-terminal-zm-branch-puiseux-uniformization}
将 \(F(X,y)=0\) 视为定义代数函数 \(X=X(y)\) 的四次方程（见结论 \ref{con:xi-terminal-zm-leyang-discriminant-norm-identity}）。
则对每一有限分歧值
\(
y_0\in\{0,1\}\cup\{y_1,y_2,y_3\}
\)
（其中 \(P_{\mathrm{LY}}(y_i)=0\)），存在局部参数 \(t\) 使 \(y-y_0=t^2\)，并且恰有两条分支在该点发生二重并合，其 Puiseux 展开具有普适平方根型
\[
X=\alpha_0\pm c_0\,t+O(t^2).
\]
其中：
\begin{enumerate}
  \item 在 \(y_0=0\)（并合点 \(X=1\)）处，可取 \(y=t^2\)，则两条并合分支满足
  \[
  X=1\pm \frac{1}{\sqrt2}\,t+\frac58\,t^2\mp \frac{43}{64\sqrt2}\,t^3-\frac{1}{16}t^4+O(t^5),
  \]
  即
  \[
  X=1\pm \frac{1}{\sqrt2}\,y^{1/2}+\frac58\,y\mp \frac{43}{64\sqrt2}\,y^{3/2}-\frac{1}{16}y^{2}+O(y^{5/2}).
  \]
  \item 在 \(y_0=1\)（并合点 \(X=-1\)）处，可取 \(y-1=t^2\)，则两条并合分支满足
  \[
  X=-1\pm \frac{i}{\sqrt6}\,t+O(t^2),
  \]
  即
  \[
  X=-1\pm \frac{i}{\sqrt6}\,(y-1)^{1/2}+O(y-1).
  \]
  \item 在任一 Lee--Yang 分歧值 \(y_0=y_i\)（并合点 \(X=\alpha_i\)，其中 \(\alpha_i\) 为 \(16X^3-9X^2+1=0\) 的根）处，首项系数满足
  \[
  \boxed{\;
  c_0^2=\frac{2}{3}\cdot\frac{3\alpha_i^2-1}{\alpha_i^2-1}\ \in\ \QQ(\alpha_i)=\QQ(y_i)\; }.
  \]
\end{enumerate}
\par
取局部参数 \(t=\sqrt{y-y_0}\) 的任一分支，则两条并合谱支可写为
\[
\lambda_{\pm}(y)=\alpha_0\pm c_0\sqrt{y-y_0}+O(y-y_0).
\]
因此对任一对数分支定义边界自由能 \(f_{\mathrm{edge}}(y):=\log\lambda_{\pm}(y)\)，则在 \(y\to y_0\) 时具有标准 \(1/2\) 阶分歧
\[
f_{\mathrm{edge}}(y)=\log\alpha_0\pm \frac{c_0}{\alpha_0}\sqrt{y-y_0}+O(y-y_0),
\qquad
\frac{\dd}{\dd y}f_{\mathrm{edge}}(y)=\pm \frac{c_0}{2\alpha_0}(y-y_0)^{-1/2}+O(1).
\]
在 \(y_0=y_{\mathrm{LY}}\) 的情形下，\(\alpha_0=\lambda_{\mathrm{LY}}\) 为结论 \ref{con:xi-terminal-zm-leyang-lambda-cubic} 的唯一实根，故
\[
c_0^2=\frac{2(3\lambda_{\mathrm{LY}}^{2}-1)}{3(\lambda_{\mathrm{LY}}^{2}-1)},
\qquad
c_0\approx 0.7476109845904598\ldots.
\]
\end{conclusion}
\begin{proof}[证明要点]
在分歧点处有 \(F(\alpha_0,y_0)=0\)、\(\partial_XF(\alpha_0,y_0)=0\)，并且 \(\partial_{XX}F(\alpha_0,y_0)\neq 0\)。
对 \(F\) 作二元 Taylor 展开并令 \(y-y_0=t^2\) 即得平方根型 Puiseux。
一般地，首项系数由二阶展开给出 \(c_0^2=-2(\partial_yF)/(\partial_{XX}F)\) 在 \((\alpha_0,y_0)\) 处的取值；
在原谱坐标 \((\lambda,y)\) 下（其中 \(F(X,y)=\Pi(X,y)\)），同一公式亦可写为
\(
c_0^2=-2\,\partial_y\Pi(\alpha_0,y_0)/\partial_{\lambda}^{2}\Pi(\alpha_0,y_0)
\)。
对 \(y_0=y_i\) 的情形再利用 \(\alpha_i\) 的三次证书化简，即得到所示闭式。
对 \(y_0=0,1\) 可进一步递推系数得到所列展开。
可复现核验脚本：\texttt{scripts/exp\_fold\_zm\_elliptic\_leyang\_cover\_geometry\_audit.py}。证毕。
\end{proof}

\begin{conclusion}[分歧值处并合谱根的线性闭式：一次消元子给出的全局刚性]\label{con:xi-terminal-zm-branch-double-root-linear-closure}
对任意分歧值
\(
y_0\in\{0,1\}\cup\{P_{\mathrm{LY}}(y)=0\}
\)，
令 \(\lambda_*=\lambda_*(y_0)\) 表示谱四次 \(\Pi(\lambda,y_0)=0\) 的唯一二重根（等价地，\(\Pi(\lambda_*,y_0)=\partial_{\lambda}\Pi(\lambda_*,y_0)=0\) 的唯一解）。
则并合谱根满足如下线性闭式
\[
\boxed{\;
2232\,\lambda_*+16384y_0^{4}+8128y_0^{3}-14525y_0^{2}-5523y_0-2232=0\; }.
\]
并且在 Lee--Yang 三次分歧支 \(P_{\mathrm{LY}}(y_0)=0\) 上，上式等价化简为
\[
\boxed{\;
279\,\lambda_*+512y_0^{2}+518y_0+5=0\; }.
\]
因此在全部分歧集合上，并合谱值 \(\lambda_*(y_0)\) 可由线性闭式直接读出，无需在每个 \(y_0\) 处再行求解四次方程。
\end{conclusion}
\begin{proof}[证明要点]
对理想 \((\Pi,\partial_{\lambda}\Pi)\subset\QQ[\lambda,y]\) 作 Gröbner 消元可得到一条关于 \(\lambda\) 的一次消元子，因而在分歧集合上以线性形式唯一刻画共同根（二重根）；
其 \(y\)-消元子即为 \(\mathrm{Disc}_{\lambda}(\Pi)=-y(y-1)P_{\mathrm{LY}}(y)\)。
在模 \(P_{\mathrm{LY}}(y)\) 的三次商环中将该一次消元子约化即可得到化简式。
可复现核验脚本：\texttt{scripts/exp\_fold\_zm\_elliptic\_leyang\_resy\_spectral\_decomposition\_audit.py}。证毕。
\end{proof}

\begin{conclusion}[\texorpdfstring{$S_4$}{S4} 分裂域的几何化：\(24\)-重 \texorpdfstring{$S_4$}{S4}-伽罗瓦闭包曲线属 \(16\)，并伴随属 \(2\) 的二次分辨曲线与属 \(1\) 的三次预解式曲线]\label{con:xi-terminal-zm-s4-splitting-curve-genus16}
令 \(\widetilde{K}/\QQ(y)\) 为四次方程 \(F(X,y)\) 的分裂域，并令 \(\widetilde{C}\) 为其函数域对应的光滑射影曲线。
在结论 \ref{con:xi-terminal-zm-generic-galois-s4} 的泛型 \(S_4\) 性下，\(\widetilde{C}\to\mathbb P^1_y\) 为 \(S_4\)-伽罗瓦覆叠，且在六个分歧值
\[
\{\infty,0,1,y_1,y_2,y_3\}
\]
处的惰性群阶分别为
\[
e_\infty=4,\qquad e_0=e_1=e_{y_1}=e_{y_2}=e_{y_3}=2.
\]
其中 \(e_\infty=4\) 对应于 \(y\) 在无穷远点 \(O\) 处的全分歧极点 \((y)_\infty=4O\)（结论 \ref{con:xi-terminal-zm-elliptic-y-hurwitz-passport}）。
由 Riemann--Hurwitz 公式得到
\[
\boxed{\ g(\widetilde{C})=16\ }.
\]

此外，二次分辨子域与三次预解式分别给出两条自然的中间曲线：
\begin{enumerate}
  \item \textbf{符号二次商（固定子群 \(A_4\)）。}
  对应曲线为双覆叠
  \[
  C_2:\ w^2=-y(y-1)P_{\mathrm{LY}}(y),
  \]
  其分歧点为上述六点，故 \(g(C_2)=2\)。
  \item \textbf{三次预解式曲线（cubic resolvent）。}
  令 \(\overline{\mathcal R}\subset\PP^{2}\) 为三次预解式 \(\mathscr R(z,y)=0\) 的射影闭包（结论 \ref{con:xi-terminal-zm-resolvent-cubic-plane-cubic-geometry}）。
  则 \(\overline{\mathcal R}\) 为光滑射影平面三次曲线，且含有有理点 \((y,z)=(0,1)\)，从而 \(g(\overline{\mathcal R})=1\) 并可视为定义在 \(\QQ\) 上的椭圆曲线。
\end{enumerate}
因此，尽管根域归一化为属 \(1\) 的椭圆曲线 \(E\)，其 \(S_4\)-分裂域的伽罗瓦闭包在几何上提升为属 \(16\) 的高属曲线；
与此同时，二次分辨曲线与三次预解式曲线分别落在属 \(2\) 与属 \(1\)，为该 \(S_4\) 结构提供两条可见的低属投影。
\end{conclusion}
\begin{proof}[证明要点]
对 \(S_4\)-伽罗瓦覆叠使用 Riemann--Hurwitz 公式：
若在分歧值 \(b\) 处惰性群阶为 \(e_b\)，则总分歧贡献为 \(|S_4|\,(1-1/e_b)\)。
代入 \(|S_4|=24\) 与 \((e_\infty,e_0,e_1,e_{y_1},e_{y_2},e_{y_3})=(4,2,2,2,2,2)\) 得 \(g(\widetilde{C})=16\)。
二次商 \(C_2\) 的属由同一公式在次数 \(2\) 上计算得到 \(g(C_2)=2\)。
三次预解式曲线 \(\overline{\mathcal R}\) 的光滑性与属计算见结论 \ref{con:xi-terminal-zm-resolvent-cubic-plane-cubic-geometry}。
可复现核验脚本：\texttt{scripts/exp\_fold\_zm\_elliptic\_leyang\_cover\_geometry\_audit.py}。证毕。
\end{proof}

\begin{remark}[进一步结构化方向：\texorpdfstring{$S_4$}{S4} 伽罗瓦闭包上的表示与周期数据]\label{rem:xi-terminal-zm-s4-closure-representation-periods}
由结论 \ref{con:xi-terminal-zm-s4-splitting-curve-genus16} 的 \(S_4\)-伽罗瓦覆叠 \(\widetilde{C}\to\mathbb P^1_y\)（属 \(16\)），
自然得到群 \(S_4\) 在 \(H^{0}(\widetilde{C},\Omega^{1})\) 上的表示。
对该表示作不可约分解并与椭圆曲线 \(E\) 的模性锚点（判别式 \(37\)，见注记 \ref{rem:xi-terminal-elliptic-37a1-modular-anchor}）对接，
可在同一框架下统一刻画分支单值化、谱支置换以及相关的周期积分与椭圆对数数据。
\end{remark}


\begin{conclusion}[纤维主除子诱导的群律迹：固定重量纤维的四点零和]\label{con:xi-terminal-zm-elliptic-weight-fiber-sum-zero}
沿用结论 \ref{con:xi-terminal-zm-elliptic-y-hurwitz-passport} 的次数 \(4\) 重量态射
\(
\pi_y:E\to\PP^1_y
\)
与极点结构 \((y)_\infty=4O\)。
对任意非分歧值 \(y_0\in\overline{\QQ}\setminus\bigl(\{\infty,0,1\}\cup\{P_{\mathrm{LY}}(y)=0\}\bigr)\)，
记纤维为
\(
\pi_y^{-1}(y_0)=\{P_1,P_2,P_3,P_4\}
\)
（四点互异）。
则在椭圆曲线群律下成立严格恒等式
\[
\boxed{\;P_1+P_2+P_3+P_4=O\;}
\]
并且更一般地，对任意有限值 \(y_0\in\overline{\QQ}\)（允许分歧值）有主除子表达
\[
\mathrm{div}(y-y_0)=\sum_{P\in \pi_y^{-1}(y_0)}\operatorname{ord}_{P}(y-y_0)\,P-4O,
\]
从而在 \(\mathrm{Pic}^0(E)\cong E\) 的识别下给出加权群律迹
\(
\sum_{P}\operatorname{ord}_{P}(y-y_0)\,P=O
\)。
无穷远纤维由 \((y)_\infty=4O\) 给出。
若 \(Q\in E(\QQ)\) 且 \(y(Q)=y_0\) 为非分歧值，则其余三点（一般位于三次扩张中）之和严格等于 \(-Q\)。
\end{conclusion}
\begin{proof}[证明要点]
对非分歧值 \(y_0\)，函数 \(y-y_0\in\QQ(E)^\times\) 的零点恰为 \(P_1,\dots,P_4\) 且均为单零；
又由 \((y)_\infty=4O\) 得 \((y-y_0)_\infty=4O\)，从而
\(
\mathrm{div}(y-y_0)=P_1+P_2+P_3+P_4-4O
\)。
由于主除子在 \(\mathrm{Pic}^0(E)\cong E\) 中对应单位元，遂得 \(P_1+\cdots+P_4=O\)。
分歧值情形以 \(\operatorname{ord}_{P}(y-y_0)\) 计入重数后同理。证毕。
\end{proof}

\begin{conclusion}[有理点特化的整化分解：线性因子与三次余因子的整系数闭式]\label{con:xi-terminal-zm-elliptic-rational-specialization-integral-factorization}
沿用结论 \ref{con:xi-terminal-zm-order4-rational} 的四次谱方程 \(\Pi(\lambda,y)\) 及椭圆化识别。
令 \(Q\in E(\QQ)\) 为有理点，并写其谱坐标为既约形式
\[
X(Q)=\frac{u}{v^{2}},
\qquad u\in\ZZ,\ v\in\ZZ_{>0},\ \gcd(u,v)=1.
\]
再将重量值写为既约形式
\[
y(Q)=\frac{a}{v^{4}},
\qquad a\in\ZZ,\ \gcd(a,v)=1
\]
（分母精确为 \(v^{4}\) 见结论 \ref{con:xi-terminal-zm-pi-rational-root-iff-y-image}）。
则在 \(\ZZ[\lambda]\) 中成立完全显式分解恒等式
\[
\boxed{\;
v^{8}\Pi\Bigl(\lambda,\frac{a}{v^{4}}\Bigr)=\bigl(v^{2}\lambda-u\bigr)\,G_{Q}(\lambda)\;}
\]
其中三次余因子 \(G_Q(\lambda)\in\ZZ[\lambda]\) 可取为
\[
\boxed{\;
\begin{aligned}
G_Q(\lambda)=\;&v^{6}\lambda^{3}
+v^{4}(u-v^{2})\lambda^{2}
+v^{2}\bigl(u^{2}-uv^{2}-2a-v^{4}\bigr)\lambda\\
&+\Bigl(v^{6}+u\bigl(u^{2}-uv^{2}-2a-v^{4}\bigr)\Bigr).
\end{aligned}}
\]
因此，特化多项式 \(\Pi(\lambda,y(Q))\) 必含有理线性因子 \(\lambda-X(Q)\)，而三次因子 \(G_Q\) 的三根精确对应纤维 \(\pi_y^{-1}(y(Q))\) 中除 \(Q\) 外三点之 \(X\)-坐标。
\end{conclusion}
\begin{proof}[证明要点]
由结论 \ref{con:xi-terminal-zm-pi-rational-root-iff-y-image}，\(\lambda=X(Q)=u/v^{2}\) 为 \(\Pi(\lambda,a/v^{4})\) 的有理根；
两边乘以 \(v^{8}\) 得 \(v^{8}\Pi(\lambda,a/v^{4})\in\ZZ[\lambda]\) 并含因子 \((v^{2}\lambda-u)\)。
在 \(\QQ[\lambda]\) 中以 \((v^{2}\lambda-u)\) 作多项式除法，并逐项比较系数即可回收所示三次商 \(G_Q\) 的闭式；
该闭式的整系数性由 \(u,v,a\in\ZZ\) 直接推出。
根的几何解释由 \(\QQ(E)=\QQ(y)(X)\) 及 \(\Pi(X,y)=0\) 的归一化识别给出：对固定 \(y_0\)，方程 \(\Pi(\lambda,y_0)=0\) 的根集即为纤维点的 \(X\)-坐标多重集。
可复现核验脚本：\texttt{scripts/exp\_fold\_zm\_elliptic\_rational\_fiber\_cubic\_factor\_audit.py}。证毕。
\end{proof}

\begin{conjecture}[倍乘轨道诱导的三次余因子之满对称伽罗瓦性]\label{conj:xi-terminal-zm-elliptic-rational-fiber-cubic-s3-galois}
取物理基点 \(P=(2,-\tfrac52)\in E(\QQ)\)（结论 \ref{con:xi-terminal-zm-elliptic-rational-points-denominator-law}），并令 \(Q_n:=[n]P\)。
对每个 \(n\ge 1\) 将
\[
X(Q_n)=\frac{u_n}{v_n^2},
\qquad
y(Q_n)=\frac{a_n}{v_n^4}
\]
写为既约形式（\(u_n,a_n\in\ZZ\)、\(v_n\in\ZZ_{>0}\)）。
令 \(G_{Q_n}(\lambda)\in\ZZ[\lambda]\) 表示结论 \ref{con:xi-terminal-zm-elliptic-rational-specialization-integral-factorization} 的三次余因子。
猜想对所有 \(n\ge 2\)，多项式 \(G_{Q_n}(\lambda)\) 在 \(\QQ\) 上不可约且其判别式不为 \(\QQ^{\times 2}\) 中的平方，
从而
\[
\mathrm{Gal}\bigl(G_{Q_n}/\QQ\bigr)\cong S_3.
\]
此外，预期其三次分裂域除有限例外外两两不相同，并在合成意义下呈现近乎线性无关的增长行为（可通过判别式的平方自由核与分歧素数的稀疏互素性形式化）。
\end{conjecture}


\subsubsection{\texorpdfstring{$S_4$}{S4} 伽罗瓦闭包 Jacobian 的可见椭圆--Prym 因子分解：由子群格到等型因子之几何识别}\label{subsubsec:xi-s4-visible-jacobian-prym-factorization}
以下内容在结论 \ref{con:xi-terminal-zm-s4-galois-closure-genus-tower-holomorphic-differentials} 的表示分解与同源分解框架上，
把其中的 \(V_2,V_3,V_3'\) 等型因子分别几何化为三条具体的中间椭圆曲线与一条 Prym 三维体，从而得到 \(J(X_e)\) 的四因子可见分解。

\begin{conclusion}[\texorpdfstring{$D_4$}{D4}-固定子域的几何化：三次预解式椭圆曲线即为 \texorpdfstring{$X_e/D_4$}{Xe/D4}]\label{con:xi-terminal-zm-s4-d4fixedfield-resolvent-elliptic}
设 \(X_e\to\mathbb P^1_y\) 为谱四次覆盖 \(\Pi(\lambda,y)=0\) 的 \(S_4\) 伽罗瓦闭包曲线（结论 \ref{con:xi-terminal-zm-s4-galois-closure-genus-tower-holomorphic-differentials}；亦即结论 \ref{con:xi-terminal-zm-s4-splitting-curve-genus16} 的几何闭包曲线）。
取任意 Sylow \(2\)-子群 \(D_{4}\le S_{4}\)（阶 \(8\) 的二面体子群，亦常记为 \(D_{8}\)），并记
\[
R:=X_e/D_{4}.
\]
则 \(R\) 与结论 \ref{con:xi-terminal-zm-resolvent-cubic-plane-cubic-geometry} 的三次预解式平面三次曲线 \(\overline{\mathcal R}\) 双有理同一；因而 \(R\) 本身为椭圆曲线，并且（在 \(\QQ\) 上）满足
\[
R\ \cong\ \overline{\mathcal R}\ \cong\ E_{\mathscr R},
\]
其中 \(E_{\mathscr R}\) 为结论 \ref{con:xi-terminal-zm-resolvent-cubic-weierstrass-birational} 的 Weierstrass 模型。
\end{conclusion}
\begin{proof}[证明要点]
在代数闭包中记谱四次的四根为 \(\lambda_1,\dots,\lambda_4\)。
由结论 \ref{con:xi-terminal-zm-resolvent-cubic-plane-cubic-geometry}，配对不变量
\(
z_{12|34}:=\lambda_1\lambda_2+\lambda_3\lambda_4
\)
等满足三次预解式 \(\mathscr R(z,y)=0\)，并且 \(S_4\) 在三组配对 \(\{12|34,13|24,14|23\}\) 上诱导出次数 \(3\) 的自然置换作用。
其中稳定某一配对（例如 \(12|34\)）的稳定子恰为阶 \(8\) 的二面体子群 \(D_4\)。
因此固定子域 \(\QQ(X_e)^{D_4}\) 由 \((y,z_{12|34})\) 生成，并与 \(\mathscr R(z,y)=0\) 的函数域同一，从而得到 \(X_e/D_4\) 与 \(\overline{\mathcal R}\) 的双有理同一。
\end{proof}

\begin{conclusion}[\texorpdfstring{$V_2$}{V2}-等型因子的椭圆平方化：\texorpdfstring{$A_2\sim R^2$}{A2~R^2}]\label{con:xi-terminal-zm-s4-jacobian-a2-isog-r2}
沿用结论 \ref{con:xi-terminal-zm-s4-galois-closure-genus-tower-holomorphic-differentials} 的记号：
\[
J(X_e)\ \sim\ A_{\mathrm{sgn}}\times A_{2}\times A_{3}\times A_{3'}.
\]
令 \(R:=X_e/D_4\) 为结论 \ref{con:xi-terminal-zm-s4-d4fixedfield-resolvent-elliptic} 的椭圆商。
则 \(V_2\)-等型因子满足
\[
\boxed{\ A_{2}\ \sim\ R^{2}\ }.
\]
\end{conclusion}
\begin{proof}[证明要点]
由 \(S_4\) 在三组配对上的次数 \(3\) 置换作用可知：指数为 \(3\) 的共轭子群族 \(\{gD_4g^{-1}\}\) 与三条中间商曲线 \(\{X_e/(gD_4g^{-1})\}\) 一一对应。
其置换表示同构于 \(\mathrm{Ind}_{D_4}^{S_4}\mathbf{1}\)（维数 \(3\)），并分解为
\[
\mathrm{Ind}_{D_4}^{S_4}\mathbf{1}\ \cong\ \mathbf{1}\ \oplus\ V_2.
\]
另一方面，\(X_e/S_4\cong\mathbb P^1\) 故 \(H^0(X_e,\Omega^1)\) 中不含平凡表示分量；
因此由三条椭圆商 \(R\) 的 pullback 在 \(J(X_e)\) 中生成的阿贝尔子簇恰对应于 \(V_2\) 分量，
且其“去掉对角平凡分量”后得到两份独立拷贝，从而给出同源 \(A_2\sim R^2\)。
\end{proof}

\begin{conclusion}[\texorpdfstring{$V_3$}{V3}-等型因子的椭圆立方化：\texorpdfstring{$A_3\sim E^3$}{A3~E^3}]\label{con:xi-terminal-zm-s4-jacobian-a3-isog-e3}
令
\[
E:=X_e/S_3
\]
为结论 \ref{con:xi-terminal-zm-s4-galois-closure-genus-tower-holomorphic-differentials} 的椭圆中间商（亦即结论 \ref{con:xi-terminal-zm-elliptic-normalization} 的椭圆归一化曲线）。
则 \(V_3\)-等型因子满足
\[
\boxed{\ A_{3}\ \sim\ E^{3}\ }.
\]
\end{conclusion}
\begin{proof}[证明要点]
四次根域对应子群为点稳定子 \(S_3\le S_4\)，其共轭类对应于四个根的选择。
于是 \(\mathrm{Ind}_{S_3}^{S_4}\mathbf{1}\) 给出次数 \(4\) 的置换表示，并分解为
\(
\mathrm{Ind}_{S_3}^{S_4}\mathbf{1}\cong \mathbf{1}\oplus V_3
\)。
同理，由于平凡分量不出现在 \(H^0(X_e,\Omega^1)\) 中，
四条共轭椭圆商 \(X_e/(gS_3g^{-1})\cong E\) 的 pullback 在 Jacobian 中生成的非平凡部分恰对应 \(V_3\) 等型因子，
从而给出同源 \(A_3\sim E^3\)。
\end{proof}

\begin{conclusion}[\texorpdfstring{$V_3'$}{V3'}-等型因子的 Prym 三维体识别：\texorpdfstring{$A_{3'}\sim P^3$}{A3'~P^3}]\label{con:xi-terminal-zm-s4-jacobian-a3p-isog-prym3-cubed}
取结论 \ref{con:xi-terminal-zm-s4-d4fixedfield-resolvent-elliptic} 中的二面体子群 \(D_4\) 的旋转子群
\(
C_4\lhd D_4
\)
（阶 \(4\) 的循环子群），并令
\[
X_4:=X_e/C_4,
\qquad
R:=X_e/D_4.
\]
则自然映射 \(X_4\to R\) 为次数 \(2\) 覆盖。
定义 Prym 三维体
\[
P:=\operatorname{Prym}(X_4/R):=\ker\!\bigl(\operatorname{Nm}_{X_4/R}:J(X_4)\to J(R)\bigr)^{0}.
\]
则 \(V_3'\)-等型因子满足
\[
\boxed{\ A_{3'}\ \sim\ P^{3}\ } .
\]
\end{conclusion}
\begin{proof}[证明要点]
对二次覆盖 \(X_4\to R\)，经典 Prym 分解给出
\(
J(X_4)\sim J(R)\times P
\)；
其中 \(\dim P=g(X_4)-g(R)\)。
\par
另一方面，\(C_4\) 在 \(S_4\) 中的陪集置换表示为 \(\mathrm{Ind}_{C_4}^{S_4}\mathbf{1}\)（维数 \(6\)），并满足分解
\[
\mathrm{Ind}_{C_4}^{S_4}\mathbf{1}\ \cong\ \mathbf{1}\ \oplus\ V_2\ \oplus\ V_3'.
\]
结合结论 \ref{con:xi-terminal-zm-s4-galois-closure-genus-tower-holomorphic-differentials} 的
\(
H^0(X_e,\Omega^1)\cong 2\cdot\mathrm{sgn}\oplus V_2\oplus V_3\oplus 3V_3'
\)
可见
\(
g(X_4)=\dim H^0(X_e,\Omega^1)^{C_4}=1+3=4
\)
且 \(g(R)=1\)，从而 \(\dim P=3\)。
\par
再由同源分解中 \(V_3'\) 的重数为 \(3\)，并注意 \(C_4\) 的共轭类恰给出三条同构的二次塔 \(X_e/C_4\to X_e/D_4\)，
其 Prym 三维体在同源意义下同构；由此把 \(V_3'\) 的三份拷贝几何化为 \(P^3\)，即得到 \(A_{3'}\sim P^3\)。
\end{proof}

\begin{conclusion}[\texorpdfstring{$S_4$}{S4}-闭包曲线 Jacobian 的四因子可见分解]\label{con:xi-terminal-zm-s4-jacobian-visible-factorization}
沿用上述记号，并记 \(C:=X_e/A_4\)（属 \(2\) 的判别式二次脊曲线，结论 \ref{con:xi-terminal-zm-s4-galois-closure-genus-tower-holomorphic-differentials}）。
则在 \(\QQ\) 上成立同源分解
\[
\boxed{\;
J(X_e)\ \sim\ J(C)\times R^{2}\times E^{3}\times P^{3}\; }.
\]
\end{conclusion}
\begin{proof}[证明要点]
由结论 \ref{con:xi-terminal-zm-s4-galois-closure-genus-tower-holomorphic-differentials} 已知
\(
J(X_e)\sim J(C)\times A_2\times A_3\times A_{3'}
\)。
再分别代入结论 \ref{con:xi-terminal-zm-s4-jacobian-a2-isog-r2}、
\ref{con:xi-terminal-zm-s4-jacobian-a3-isog-e3} 与
\ref{con:xi-terminal-zm-s4-jacobian-a3p-isog-prym3-cubed} 的同源识别，即得所示四因子可见分解。
\end{proof}

\begin{conclusion}[\texorpdfstring{$C_2$}{C2}-换位商的属 \(6\) 中间覆盖与 \texorpdfstring{$A_{2}\times A_{3}\times A_{3'}$}{A2xA3xA3'} 因子剥离]\label{con:xi-terminal-zm-s4-transposition-quotient-genus6}
沿用结论 \ref{con:xi-terminal-zm-s4-galois-closure-genus-tower-holomorphic-differentials} 的 \(S_4\)-伽罗瓦闭包曲线 \(X_e\to\PP^{1}_{y}\)（\(|S_4|=24\)）。
取任意由换位生成的子群 \(H\simeq C_{2}\subset S_4\)，并令
\[
X_{H}:=X_e/H.
\]
则商映射 \(X_{H}\to\PP^{1}_{y}\) 为次数 \(12\) 的有限态射，其分歧型由 \(S_4\) 在左陪集集合 \(S_4/H\) 上的作用给出，并满足：
\begin{enumerate}
  \item 在每个惯性为换位的分歧值 \(y\in\{0,1\}\cup\{P_{\mathrm{LY}}(y)=0\}\) 上，惯性元在 \(12\) 个陪集上的循环型为 \(1^{2}2^{5}\)；
  等价地，该纤维上恰有 \(5\) 个分歧指数 \(2\) 的点与 \(2\) 个不分歧点。
  \item 在无穷远分歧值 \(y=\infty\) 上，惯性 \(4\)-循环在 \(12\) 个陪集上的循环型为 \(4^{3}\)，即该纤维上恰有 \(3\) 个分歧指数 \(4\) 的点。
  \item 因而由 Riemann--Hurwitz 公式得到
  \[
  \boxed{\ g(X_H)=6\ }.
  \]
\end{enumerate}
并且在结论 \ref{con:xi-terminal-zm-s4-galois-closure-genus-tower-holomorphic-differentials} 的同源分解记号下，
\[
J(X_{H})\ \sim_{\QQ}\ A_{2}\times A_{3}\times A_{3'},
\]
亦即 \(X_H\) 的 Jacobian 恰对应于 \(J(X_e)\) 中去掉 \(\mathrm{sgn}\)-等型因子后的剩余三分量。
\end{conclusion}
\begin{proof}[证明要点]
分歧循环型：对任意 \(g\in S_4\)，其在陪集集 \(S_4/H\) 上的固定点个数为
\(
\#\{xH:\ gxH=xH\}=\#\{x\in S_4:\ x^{-1}gx\in H\}/|H|
\)。
当 \(g\) 为换位时，\(H\cap g^{S_4}\) 恰含一元素，且 \(|C_{S_4}(g)|=2\cdot 2!=4\)，故固定点数为 \(4/2=2\)，其余 \(10\) 点成 \(5\) 个 \(2\)-轮换，得循环型 \(1^{2}2^{5}\)。
当 \(g\) 为 \(4\)-循环时，\(H\cap g^{S_4}=\varnothing\)，故无固定点。
又 \(g^{2}\) 为双换位，仍有 \(H\cap (g^{2})^{S_4}=\varnothing\)，从而 \(g^{2}\) 在陪集集上亦无固定点；
因此 \(g\) 的循环分解不含 \(1\)-循环或 \(2\)-循环，只能为 \(4\)-循环，且 \(12\) 点必分解为 \(4^{3}\)。
\par
属数：对次数 \(12\) 覆盖应用 Riemann--Hurwitz 公式
\[
2g(X_H)-2=-2\cdot 12+\sum_{b}\sum_{P\mid b}(e_P-1).
\]
五个换位分歧值处，每个分歧值的总贡献为 \(5\cdot(2-1)=5\)；
无穷远处总贡献为 \(3\cdot(4-1)=9\)。
总分歧贡献 \(5\cdot 5+9=34\)，代入得 \(2g(X_H)-2=-24+34=10\)，故 \(g(X_H)=6\)。
\par
雅可比因子：由于 \(X_e\to X_H\) 为 Galois 商映射，有
\(
H^{0}(X_H,\Omega^{1})\cong H^{0}(X_e,\Omega^{1})^{H}
\)。
由结论 \ref{con:xi-terminal-zm-s4-galois-closure-genus-tower-holomorphic-differentials} 的表示分解
\(
H^{0}(X_e,\Omega^{1})\cong 2\cdot \mathrm{sgn}\oplus V_2\oplus V_3\oplus 3V_3'
\)
以及 \(S_4\) 角色表在换位处的迹值可知：\(\mathrm{sgn}\) 在 \(H\) 上无不变元，而 \(V_2,V_3,V_3'\) 的不变元维数分别为 \(1,2,1\)。
故
\(
\dim H^{0}(X_e,\Omega^{1})^{H}=1+2+3\cdot 1=6=g(X_H)
\)，并且 \(H\)-不变部分恰由 \(A_2,A_3,A_{3'}\) 三个等型因子贡献，从而得到 \(J(X_H)\sim_{\QQ}A_2\times A_3\times A_{3'}\)。证毕。
\end{proof}

\endinput

\paragraph{共轭重量支下的闭式审计链：临界域、分歧域与倍乘对应的算术分解}
本段在既有椭圆统一化（结论 \ref{con:xi-terminal-zm-elliptic-normalization}）与 Lee--Yang 分歧刻画（结论 \ref{con:xi-terminal-zm-leyang-ramification-critical-values}）的基础上，
将 Latt\`es 临界多项式、共轭重量覆盖的分歧几何与倍乘诱导的四值对应统一为一组可直接引用的闭式证书。

\begin{conclusion}[临界六次的 \(4\)-分点解释与分裂域刻画]\label{con:xi-terminal-zm-lattes-critical-sextic-4division-splitting-field}
沿用结论 \ref{con:xi-terminal-zm-elliptic-lattes-doubling} 的椭圆曲线
\(
E/\QQ:\ Y^{2}=X^{3}-X+\tfrac14
\)
与倍点 Latt\`es 自映射
\(
\Phi(X)=\frac{X^{4}+2X^{2}-2X+1}{4X^{3}-4X+1}
\)。
令
\[
\Psi(X):=2X^{6}-10X^{4}+10X^{3}-10X^{2}+2X+1\in\ZZ[X].
\]
则：
\begin{enumerate}
  \item \(\Psi\) 恰为 \(4\)-division 多项式 \(\psi_{4}\) 去除 \(2\)-division 因子后的本原因子；等价地，有限临界集合满足
  \[
  \Phi'(X)=0
  \ \Longleftrightarrow\
  \exists\,Q\in E[4]\setminus E[2]\ \text{使得}\ X=x(Q)
  \ \Longleftrightarrow\
  \Psi(X)=0.
  \]
  因而 \(\Phi\) 的六个有限临界点正是 \(x\bigl(E[4]\setminus E[2]\bigr)\)（在 \(\pm\) 识别下）的像。
  \item \(\Psi\) 在 \(\QQ\) 上不可约，且其判别式满足严格恒等式
  \[
  \mathrm{Disc}(\Psi)=-2^{8}\cdot 37^{5}.
  \]
  从而 \(\Psi\) 的分裂域 \(K_{4}/\QQ\) 的分歧素数集合包含于 \(\{2,37\}\)。
  \item 伽罗瓦群为阶 \(48\) 的传递群，且同构于
  \[
  \mathrm{Gal}(K_{4}/\QQ)\cong S_{4}\times C_{2}.
  \]
  特别地，六个有限临界点在 \(\mathrm{Gal}(\overline{\QQ}/\QQ)\) 作用下构成单一轨道。
\end{enumerate}
\end{conclusion}
\begin{proof}[证明要点]
第一点由结论 \ref{con:xi-terminal-zm-elliptic-lattes-doubling} 的 \(\psi_{4}/(2Y)\) 恒等式给出；
第二、三点由结论 \ref{con:xi-terminal-zm-elliptic-lattes-critical-sextic-galois} 的不可约性、判别式与伽罗瓦群计算推出；
传递性等价于“临界点构成单一伽罗瓦轨道”。
可复现核验脚本：\texttt{scripts/exp\_fold\_zm\_elliptic\_lattes\_rational\_points\_audit.py}。证毕。
\end{proof}

\begin{conclusion}[谱四次的显式椭圆化与共轭重量覆盖的分歧剖面]\label{con:xi-terminal-zm-elliptic-normalization-y-iota-ramification-portrait}
令平面曲线
\[
C:\quad \Pi(\lambda,y)=\lambda^{4}-\lambda^{3}-(2y+1)\lambda^{2}+\lambda+y(y+1)=0.
\]
在结论 \ref{con:xi-terminal-zm-elliptic-normalization} 的椭圆统一化下，\(C\) 与
\(
E:\ Y^{2}=X^{3}-X+\tfrac14
\)
之间存在完全显式的坐标同一（在仿射开集上为同构）：
\[
E\dashrightarrow C,\qquad (X,Y)\longmapsto (\lambda,y)=(X,\ X^{2}+Y-\tfrac12),
\]
\[
C\dashrightarrow E,\qquad (\lambda,y)\longmapsto (X,Y)=(\lambda,\ y+\tfrac12-\lambda^{2}).
\]
记物理主支重量函数与其共轭支
\[
y:=X^{2}+Y-\frac12\in\QQ(E),
\qquad
y^{\iota}:=X^{2}-Y-\frac12\in\QQ(E),
\qquad
P_{\mathrm{LY}}(t):=256t^{3}+411t^{2}+165t+32\in\ZZ[t].
\]
则 \(y:E\to\mathbb P^{1}\) 为次数 \(4\) 的覆盖，并满足以下全分歧剖面：
\begin{enumerate}
  \item 在无穷远点 \(O\) 处，\(y\) 具有唯一极点且极点阶为 \(4\)，从而在 \(y=\infty\) 处全分歧（\(e_{O}=4\)）。
  \item 有限分歧值集合恰为
  \[
  y\in\{0,1\}\ \cup\ V\!\bigl(P_{\mathrm{LY}}\bigr),
  \]
  并且每个有限分歧值均为单分歧（上方存在唯一分歧点，分歧指标 \(2\)）。
  \item 两个有理分歧值的纤维分解完全显式：
  \[
  (y)^{-1}(0)=\Bigl\{\Bigl(0,\frac12\Bigr),\ \Bigl(-1,-\frac12\Bigr),\ \underbrace{\Bigl(1,-\frac12\Bigr)}_{e=2}\Bigr\},
  \qquad
  (y)^{-1}(1)=\Bigl\{\Bigl(1,\frac12\Bigr),\ \Bigl(2,-\frac52\Bigr),\ \underbrace{\Bigl(-1,\frac12\Bigr)}_{e=2}\Bigr\}.
  \]
\end{enumerate}
\end{conclusion}
\begin{proof}[证明要点]
由结论 \ref{con:xi-terminal-zm-y-quadratic-extension-minpoly}，物理主支 \(y=X^{2}+Y-\tfrac12\) 与共轭支 \(y^{\iota}=X^{2}-Y-\tfrac12\) 互为共轭，且由椭圆反演 \(\iota:(X,Y)\mapsto(X,-Y)\) 交换。
结论 \ref{con:xi-terminal-zm-elliptic-y-hurwitz-passport} 给出 \(y\) 的 Hurwitz 护照与两个有理分歧纤维。
有限分歧值集合与 \(P_{\mathrm{LY}}\) 的出现由判别式分解（结论 \ref{con:xi-terminal-zm-discriminant-leyang}，等价地见结论 \ref{con:xi-terminal-zm-leyang-cubic-branch-image}）决定。证毕。
\end{proof}

\begin{conclusion}[Lee--Yang 三次因子的分歧消元产生式：共轭重量支的临界点立方刻画]\label{con:xi-terminal-zm-leyang-branch-elimination-origin-y-iota}
在结论 \ref{con:xi-terminal-zm-elliptic-normalization-y-iota-ramification-portrait} 的记号下，令
\[
g(X):=16X^{3}-9X^{2}+1.
\]
则 \(y:E\to\mathbb P^{1}\) 在 \(V\!\bigl(P_{\mathrm{LY}}\bigr)\) 上方的三条非有理分歧点恰为满足
\[
4XY+3X^{2}-1=0,\qquad X\neq \pm 1
\]
的三点；等价地可由
\[
g(X)=0,\qquad Y=\frac{1-3X^{2}}{4X}
\]
刻画。
更进一步，分歧值与上述立方分歧点之间存在严格消元恒等式：在分歧约束下有线性关系
\[
4Xy=4X^{3}-3X^{2}-2X+1,
\]
并且满足结果式证书
\[
\mathrm{Res}_{X}\!\Bigl(g(X),\ 4Xy-\bigl(4X^{3}-3X^{2}-2X+1\bigr)\Bigr)=-64\,P_{\mathrm{LY}}(y).
\]
因此 \(P_{\mathrm{LY}}\)（除去单位因子）可被解释为共轭重量覆盖的“非有理分歧值集合”在基底上的最小整系数消元像。
\end{conclusion}
\begin{proof}[证明要点]
与结论 \ref{con:xi-terminal-zm-leyang-ramification-critical-values} 的推导同理：由
\(
y=X^{2}+Y-\tfrac12
\)
与曲线微分关系 \(2Y\,\dd Y=(3X^{2}-1)\dd X\) 得
\(
\dd y=(4XY+3X^{2}-1)\,\omega
\)（\(\omega=\dd X/(2Y)\)）。
因此有限分歧点满足 \(4XY+3X^{2}-1=0\)；
消去 \(Y\) 得
\(
(X-1)(X+1)g(X)=0
\)，其中 \(X=\pm 1\) 对应有理分歧值 \(0,1\)，其余三点由 \(g(X)=0\) 给出。
线性关系 \(4Xy=4X^{3}-3X^{2}-2X+1\) 由代回 \(Y=(1-3X^{2})/(4X)\) 得到；
对变量 \(X\) 作消元即得结果式恒等式。
可复现核验脚本：\texttt{scripts/exp\_fold\_zm\_elliptic\_leyang\_cover\_geometry\_audit.py}。证毕。
\end{proof}

\begin{conclusion}[倍点的雅可比共周期：\texorpdfstring{$Y([2]P)$}{Y([2]P)} 与 \texorpdfstring{$\Phi'(X)$}{Phi'(X)} 的严格乘法公式]\label{con:xi-terminal-zm-elliptic-y-doubling-jacobian-cocycle}
沿用结论 \ref{con:xi-terminal-zm-elliptic-lattes-doubling} 的椭圆曲线 \(E\) 与倍点 Latt\`es 自映射 \(\Phi\)。
对任意 \(P=(X,Y)\in E(\overline{\QQ})\) 且 \(Y\neq 0\)，成立严格恒等式
\[
\boxed{\;
Y([2]P)=\frac{\Phi'(X)}{2}\,Y\;}
\]
进一步，对任意整数 \(n\ge 1\)，有
\[
\boxed{\;
Y([2^{n}]P)=\frac{(\Phi^{n})'(X)}{2^{n}}\,Y,
\qquad\text{等价地}\qquad
(\Phi^{n})'(X)=2^{n}\frac{Y([2^{n}]P)}{Y}\;}
\]
其中 \(\Phi^{n}\) 为 \(\Phi\) 的 \(n\) 次迭代。
\end{conclusion}
\begin{proof}[证明要点]
取 \(E\) 的不变微分 \(\omega=\dd X/(2Y)\)。
倍点映射满足 \([2]^{\ast}\omega=2\omega\)。
另一方面，由 \(x([2]P)=\Phi(X)\) 得
\[
[2]^{\ast}\omega=\frac{\dd(\Phi(X))}{2\,Y([2]P)}
=\frac{\Phi'(X)}{2\,Y([2]P)}\,\dd X.
\]
与 \(2\omega=\dd X/Y\) 比较即得 \(Y([2]P)=\frac{\Phi'(X)}{2}Y\)。
对 \(n\ge 1\) 迭代同理：由 \([2^{n}]^{\ast}\omega=2^{n}\omega\) 与 \(x([2^{n}]P)=\Phi^{n}(X)\) 得
\(
\frac{(\Phi^{n})'(X)}{2\,Y([2^{n}]P)}=\frac{2^{n}}{2Y}
\)，从而推出所示恒等式。证毕。
\end{proof}

\begin{conclusion}[重量参数的仿射共周期：\texorpdfstring{$y([2]P)$}{y([2]P)} 对 \texorpdfstring{$y(P)$}{y(P)} 的线性--平移分解]\label{con:xi-terminal-zm-elliptic-y-doubling-affine-cocycle-derivative}
沿用结论 \ref{con:xi-terminal-zm-elliptic-normalization-y-iota-ramification-portrait} 的重量函数
\(
y=X^{2}+Y-\tfrac12
\)。
对任意 \(P=(X,Y)\in E(\overline{\QQ})\) 且 \(Y\neq 0\)，有严格仿射恒等式
\[
\boxed{\;
y([2]P)
=\Phi(X)^{2}-\frac12+\frac{\Phi'(X)}{2}\Bigl(y-\bigl(X^{2}-\tfrac12\bigr)\Bigr)\;}
\]
亦即
\[
y([2]P)=\frac{\Phi'(X)}{2}\,y
+\Bigl(\Phi(X)^{2}-\frac12-\frac{\Phi'(X)}{2}\bigl(X^{2}-\tfrac12\bigr)\Bigr),
\]
其中括号项仅依赖于 \(X\)。
\par
更强地，记 \(2\)-division 多项式
\(
D(X)=4X^{3}-4X+1
\)，
则存在唯一的八次整式
\[
B_{8}(X)=X^{8}-7X^{6}+14X^{5}-27X^{4}+9X^{3}+6X^{2}-X-1\in\ZZ[X]
\]
使得
\[
\boxed{\;
y([2]P)=\frac{\Phi'(X)}{2}\,y-\frac{B_{8}(X)}{D(X)^{2}}\;}
\]
\end{conclusion}
\begin{proof}[证明要点]
由定义 \(y([2]P)=X([2]P)^{2}+Y([2]P)-\tfrac12=\Phi(X)^{2}+Y([2]P)-\tfrac12\)。
代入结论 \ref{con:xi-terminal-zm-elliptic-y-doubling-jacobian-cocycle} 的
\(
Y([2]P)=\frac{\Phi'(X)}{2}Y
\)
与恒等式 \(Y=y-(X^{2}-\tfrac12)\)，即得第一组公式。
最后的分式表示由把平移项化为分母 \(D(X)^{2}\) 的既约形式得到；
其显式分子与唯一性可由结论 \ref{con:xi-terminal-zm-elliptic-y-doubling-linear} 的整式证书直接读取（取 \(B_{8}=-B\)）。证毕。
\end{proof}

\begin{conclusion}[倍点拉回在重量二次扩张上的 Möbius 共周期：\texorpdfstring{$[2]^*y$}{[2]*y} 的 \texorpdfstring{$\mathrm{PGL}_2(\QQ(X))$}{PGL2(Q(X))} 表示]\label{con:xi-terminal-zm-elliptic-y-doubling-mobius-cocycle}
沿用结论 \ref{con:xi-terminal-zm-elliptic-normalization-y-iota-ramification-portrait} 的椭圆曲线
\(
E:\ Y^{2}=X^{3}-X+\tfrac14
\)
与重量函数
\(
y=X^{2}+Y-\tfrac12
\)。
记
\(
D(X):=4X^{3}-4X+1
\)。
则存在唯一的元素 \([M(X)]\in\mathrm{PGL}_{2}(\QQ(X))\)，使得在 \(\QQ(E)\) 中成立严格恒等式
\begin{equation}\label{eq:xi-elliptic-y-doubling-mobius}
\boxed{\;
[2]^{\ast}(y)=y([2]P)=\frac{a(X)\,y+b(X)}{c(X)\,y+d(X)}\;}
\end{equation}
其中可取代表矩阵
\[
M(X)=
\begin{pmatrix}
a(X) & b(X)\\
c(X) & d(X)
\end{pmatrix}\in\mathrm{Mat}_{2}(\QQ(X)),
\]
并且
\[
\begin{aligned}
a(X)&:=2\bigl(2X^{8}-8X^{6}-8X^{5}+44X^{4}-24X^{3}+1\bigr),\\
b(X)&:=-2\bigl(2X^{10}-4X^{9}-9X^{8}+16X^{7}+27X^{6}-20X^{5}-X^{4}-15X^{3}+10X^{2}+X-1\bigr),\\
c(X)&:=4\,D(X)^{2},\qquad
d(X):=-2\,(2X^{2}-1)\,D(X)^{2}.
\end{aligned}
\]
并且该归一化代表满足完全因子分解
\[
\boxed{\;
\det M(X)=-4\,D(X)^{3}\cdot\bigl(2X^{6}-10X^{4}+10X^{3}-10X^{2}+2X+1\bigr)\;}
\]
以及
\[
\boxed{\;
\tr M(X)=-4\cdot \psi_{3}(X)\cdot \bigl(5X^{4}-2X^{2}-X+1\bigr),\qquad
\psi_{3}(X)=3X^{4}-6X^{2}+3X-1\;}
\]
其中 \(\psi_{3}\) 为该短 Weierstrass 模型的三除多项式。
特别地，\(\det M(X)\) 的零除子（除去由 \(D(X)=0\) 产生的极点层）恰截获倍点 Latt\`es 自映射的临界六次多项式 \(2X^{6}-10X^{4}+10X^{3}-10X^{2}+2X+1\)；
从而重量二次扩张在 \([2]^{\ast}\) 下的 Möbius 退化与 \(x\)-坐标动力学的临界结构在同一显式除子上对齐。
\end{conclusion}
\begin{proof}[证明要点]
在短 Weierstrass 模型上使用标准倍点公式
\(
X([2]P)=\bigl(\frac{3X^{2}-1}{2Y}\bigr)^{2}-2X
\)
与
\(
Y([2]P)=\frac{3X^{2}-1}{2Y}\bigl(X-X([2]P)\bigr)-Y
\)，
并以 \(Y=y-(X^{2}-\tfrac12)\) 消去 \(Y\)。
再在重量二次扩张的最小多项式（结论 \ref{con:xi-terminal-zm-elliptic-normalization-y-iota-ramification-portrait}）下对 \(y^{2}\) 降次，即可得到 \eqref{eq:xi-elliptic-y-doubling-mobius} 的线性分式表示。
\(\det M\) 与 \(\tr M\) 的因子分解由直接计算并在 \(\QQ[X]\) 中因式分解得到，其中 \(\psi_{3}\) 为标准三除多项式。
可复现核验脚本：\texttt{scripts/exp\_fold\_zm\_elliptic\_y\_doubling\_mobius\_cocycle\_audit.py}。证毕。
\end{proof}

\begin{conclusion}[全二幂迭代的闭式仿射化：\texorpdfstring{$y([2^{n}]P)$}{y([2^n]P)} 由 \texorpdfstring{$\Phi^{n}$}{Phi^n} 与 \texorpdfstring{$(\Phi^{n})'$}{(Phi^n)'} 控制]\label{con:xi-terminal-zm-elliptic-y-doubling-iterate-affine-closed}
沿用结论 \ref{con:xi-terminal-zm-elliptic-normalization-y-iota-ramification-portrait} 的记号。
对任意整数 \(n\ge 1\) 与任意 \(P=(X,Y)\in E(\overline{\QQ})\) 且 \(Y\neq 0\)，成立严格闭式表示
\[
\boxed{\;
y([2^{n}]P)
=\Phi^{n}(X)^{2}-\frac12+\frac{(\Phi^{n})'(X)}{2^{n}}\Bigl(y-\bigl(X^{2}-\tfrac12\bigr)\Bigr)\;}
\]
等价地，
\[
y([2^{n}]P)
=\frac{(\Phi^{n})'(X)}{2^{n}}\,y
+\Bigl(\Phi^{n}(X)^{2}-\frac12-\frac{(\Phi^{n})'(X)}{2^{n}}\bigl(X^{2}-\tfrac12\bigr)\Bigr).
\]
\end{conclusion}
\begin{proof}[证明要点]
由定义 \(y([2^{n}]P)=X([2^{n}]P)^{2}+Y([2^{n}]P)-\tfrac12=\Phi^{n}(X)^{2}+Y([2^{n}]P)-\tfrac12\)。
代入结论 \ref{con:xi-terminal-zm-elliptic-y-doubling-jacobian-cocycle} 的
\(
Y([2^{n}]P)=\frac{(\Phi^{n})'(X)}{2^{n}}Y
\)，再用 \(Y=y-(X^{2}-\tfrac12)\) 消去 \(Y\) 即得。证毕。
\end{proof}

\begin{conclusion}[重量变量上的四次倍乘对应：\texorpdfstring{$y([2]P)$}{y([2]P)} 的整系数最小多项式]\label{con:xi-terminal-zm-elliptic-weight-doubling-quartic-minpoly}
沿用结论 \ref{con:xi-terminal-zm-elliptic-normalization-y-iota-ramification-portrait} 的重量函数 \(y\)，并令
\[
y_0:=y(P),\qquad y_1:=y([2]P)\qquad(P\in E(\overline{\QQ})).
\]
则 \(y_1\) 在基域 \(\QQ(y_0)\) 上满足唯一的原始四次整系数关系
\[
\mathcal H(y_0,y_1)=0,
\qquad
\mathcal H\in\ZZ[y_0,y_1],
\]
并且 \(\mathcal H\) 在 \(\QQ(y_0)[y_1]\) 中不可约；从而 \(\mathcal H\) 即为 \(y_1\) 在 \(\QQ(y_0)\) 上的最小多项式。
其显式系数证书为：
% Auto-generated by scripts/exp_fold_zm_elliptic_weight_doubling_audit.py
\[
\mathcal H(y_{0},y_{1})=A_{4}(y_{0})y_{1}^{4}+A_{3}(y_{0})y_{1}^{3}+A_{2}(y_{0})y_{1}^{2}+A_{1}(y_{0})y_{1}+A_{0}(y_{0}).
\]
\[
\begin{aligned}
A_4(y_0)&=65536 y_{0}^{12} - 131072 y_{0}^{11} + 32768 y_{0}^{10} + 81920 y_{0}^{9} - 45056 y_{0}^{8} - 16384 y_{0}^{7} + 13824 y_{0}^{6} + 512 y_{0}^{5} - 1664 y_{0}^{4} + 192 y_{0}^{3} + 64 y_{0}^{2} - 16 y_{0} + 1\\
A_3(y_0)&=- 16384 y_{0}^{13} + 61440 y_{0}^{12} - 454656 y_{0}^{11} + 82944 y_{0}^{10} + 1247744 y_{0}^{9} - 57600 y_{0}^{8} - 906688 y_{0}^{7} - 167648 y_{0}^{6} + 174368 y_{0}^{5} + 23020 y_{0}^{4} - 12596 y_{0}^{3} - 1198 y_{0}^{2} + 346 y_{0} - 63\\
A_2(y_0)&=1536 y_{0}^{14} - 9216 y_{0}^{13} + 81792 y_{0}^{12} - 118976 y_{0}^{11} + 294336 y_{0}^{10} + 641648 y_{0}^{9} - 232314 y_{0}^{8} - 391280 y_{0}^{7} - 430462 y_{0}^{6} - 198332 y_{0}^{5} + 280857 y_{0}^{4} + 263078 y_{0}^{3} + 108044 y_{0}^{2} + 14594 y_{0} + 698\\
A_1(y_0)&=- 64 y_{0}^{15} + 608 y_{0}^{14} - 5376 y_{0}^{13} + 15052 y_{0}^{12} - 37980 y_{0}^{11} + 75342 y_{0}^{10} + 75922 y_{0}^{9} - 208389 y_{0}^{8} - 257076 y_{0}^{7} - 198536 y_{0}^{6} + 143116 y_{0}^{5} + 199086 y_{0}^{4} - 37204 y_{0}^{3} - 57694 y_{0}^{2} - 26364 y_{0} - 2436\\
A_0(y_0)&=y_{0}^{16} - 16 y_{0}^{15} + 158 y_{0}^{14} - 872 y_{0}^{13} + 2243 y_{0}^{12} - 3190 y_{0}^{11} + 7024 y_{0}^{10} - 4728 y_{0}^{9} - 23757 y_{0}^{8} - 1652 y_{0}^{7} + 35994 y_{0}^{6} + 65440 y_{0}^{5} + 497 y_{0}^{4} - 35454 y_{0}^{3} - 2792 y_{0}^{2} + 1640 y_{0} + 1800
\end{aligned}
\]

其中首项系数为完全四次幂
\[
A_{4}(y_0)=\bigl(16(y_0)^{3}-8(y_0)^{2}-4y_0+1\bigr)^{4},
\]
而常数项具有非平凡因式分解
\[
\boxed{\;
A_{0}(y_0)
=(y_0+1)^{2}\bigl((y_0)^{3}-7(y_0)^{2}+13y_0-5\bigr)^{2}
\bigl((y_0)^{4}-6(y_0)^{3}+66(y_0)^{2}+28y_0+9\bigr)
\bigl((y_0)^{4}+2(y_0)^{3}+8(y_0)^{2}+8y_0+8\bigr)\; }.
\]
\end{conclusion}
\begin{proof}[证明要点]
由结论 \ref{con:xi-terminal-zm-elliptic-normalization-y-iota-ramification-portrait}，覆盖 \(y:E\to\mathbb P^{1}\) 的次数为 \(4\)，从而在函数域层面
\(
[\QQ(E):\QQ(y_0)]=4
\)。
另一方面，\(y_1=[2]^*(y)\in\QQ(E)\)。
将倍点公式与 \(y=X^{2}+Y-\tfrac12\) 联立并对 \((X,Y)\) 消元，可得到 \((y_0,y_1)\) 的平面整系数关系；
取其在 \(\ZZ[y_0,y_1]\) 中的既约本原部分即为 \(\mathcal H\)。
再取一处无分歧专门化（例如 \(y_0=2\)，此时判别式非零）并核验所得四次多项式在 \(\QQ[y_1]\) 中不可约，即推出 \(\mathcal H\) 在 \(\QQ(y_0)[y_1]\) 中不可约；
由于扩张次数为 \(4\)，不可约四次必为最小多项式。
可复现核验脚本：\texttt{scripts/exp\_fold\_zm\_elliptic\_weight\_doubling\_audit.py}。证毕。
\end{proof}

\begin{conclusion}[判别式的平方自由核刚性与泛型 \texorpdfstring{$S_{4}$}{S4}：Lee--Yang 二次脊柱不增生]\label{con:xi-terminal-zm-elliptic-weight-doubling-discriminant-squareclass-s4}\label{con:xi-terminal-zm-elliptic-weight-doubling-quartic-elimination-disc}
令 \(\mathcal H\) 如结论 \ref{con:xi-terminal-zm-elliptic-weight-doubling-quartic-minpoly}。
将 \(\mathcal H(y,v)\) 视为关于 \(v\) 的四次多项式，则其判别式具有完全可审计的显式分解：
\[
\boxed{\;
\mathrm{Disc}_{v}\bigl(\mathcal H(y,v)\bigr)=-y(y-1)P_{\mathrm{LY}}(y)\cdot(\text{完全平方})\;}
\]
更精确地，对共轭重量支的另一条不可约四次对应 \(R_{-}\) 一并有
\[
\mathrm{Disc}_{v}\bigl(R_{\pm}(y,v)\bigr)=-y(y-1)P_{\mathrm{LY}}(y)\,A(y)^{2}\,B_{\pm}(y)^{2},
\]
其中 \(A\in\ZZ[y]\) 为次数 \(12\) 因子，\(B_{\pm}\in\ZZ[y]\) 为次数 \(26\) 因子（其显式展开见下列闭式证书）：
% Auto-generated by scripts/exp_fold_zm_elliptic_weight_doubling_audit.py
\[
P_{\mathrm{LY}}(y)=256y^{3}+411y^{2}+165y+32.
\]
\[
A(y)=64 y^{12} - 128 y^{11} - 2576 y^{10} - 2160 y^{9} + 10892 y^{8} + 32064 y^{7} + 28873 y^{6} - 11139 y^{5} - 31715 y^{4} - 8333 y^{3} - 958 y^{2} - 100 y + 8.
\]
\[
B_{+}(y)=262144 y^{26} - 10747904 y^{25} + 191954944 y^{24} - 1897332736 y^{23} + 13439238144 y^{22} - 47170043904 y^{21} + 30127935488 y^{20} + 55661785088 y^{19} - 112465192960 y^{18} + 135823686656 y^{17} + 205179341056 y^{16} - 455994979584 y^{15} - 230261903040 y^{14} + 622458964864 y^{13} + 231233432960 y^{12} - 451887225664 y^{11} - 230891455552 y^{10} + 102570870816 y^{9} + 95472617744 y^{8} + 15921249600 y^{7} - 7132590928 y^{6} - 3074856528 y^{5} - 21077156 y^{4} + 227572860 y^{3} + 2463184 y^{2} - 12256842 y - 4594975.
\]
\[
B_{-}(y)=191102976 y^{26} - 3248750592 y^{25} + 15336013824 y^{24} + 3965386752 y^{23} - 141567492096 y^{22} - 71797653504 y^{21} + 1007302348800 y^{20} + 2608289648640 y^{19} + 2522478716928 y^{18} - 366569892864 y^{17} - 3997689244416 y^{16} - 5518986465024 y^{15} - 4422184238528 y^{14} - 2328409689344 y^{13} - 764216979776 y^{12} - 98954676928 y^{11} + 41010765504 y^{10} + 26703630624 y^{9} + 8069517360 y^{8} + 2341777728 y^{7} + 1038528048 y^{6} + 413681776 y^{5} + 105306364 y^{4} + 14201380 y^{3} + 624020 y^{2} - 56922 y - 4725.
\]
\[
\mathrm{Disc}_{v}\bigl(R_{\pm}(y,v)\bigr)=-y(y-1)P_{\mathrm{LY}}(y)\,A(y)^{2}B_{\pm}(y)^{2}.
\]

因此
\[
\boxed{\;
\bigl(\mathrm{Disc}_{v}(\mathcal H)\bigr)_{\mathrm{sf}}=-y(y-1)P_{\mathrm{LY}}(y)\;}
\]
并且在 \(\QQ(y)\) 上，分裂域的唯一二次分辨率子域恒等于
\[
\QQ(y)\Bigl(\sqrt{-y(y-1)P_{\mathrm{LY}}(y)}\Bigr).
\]
此外，作为 \(\QQ(y)\) 上的四次多项式，\(R_{+}(y,\cdot)=\mathcal H(y,\cdot)\) 与 \(R_{-}(y,\cdot)\) 的泛型伽罗瓦群均同构于 \(S_{4}\)。
\end{conclusion}
\begin{proof}[证明要点]
判别式分解由对 \(R_{\pm}\) 直接计算 \(\mathrm{Disc}_{v}\) 并在 \(\ZZ[y]\) 中因式分解得到；其平方自由部分随即给出唯一二次分辨率子域。
\par
对泛型伽罗瓦群，注意 \(\mathrm{Disc}_{v}(\mathcal H)\) 在 \(\QQ(y)\) 中非平方，从而排除 \(A_{4}\)。
再取一处无分歧专门化（例如 \(y=2\)）并核验所得四次多项式在 \(\QQ\) 上的伽罗瓦群为 \(S_{4}\)；
由于专门化伽罗瓦群为泛型伽罗瓦群之子群，而 \(S_{4}\) 已为四次情形的最大可能，故泛型必为 \(S_{4}\)。
上述核验可由脚本 \texttt{scripts/exp\_fold\_zm\_elliptic\_weight\_doubling\_audit.py} 复核。证毕。
\end{proof}

由此，重量参数在倍点动力学下的非线性被 \(\Phi\) 的导数共周期严格线性化（结论 \ref{con:xi-terminal-zm-elliptic-y-doubling-jacobian-cocycle}--\ref{con:xi-terminal-zm-elliptic-y-doubling-iterate-affine-closed}），
而消元后产生的四次倍乘对应在判别式层面把唯一二次障碍刚性锁定为同一 Lee--Yang 数据 \(-y(y-1)P_{\mathrm{LY}}(y)\)（结论 \ref{con:xi-terminal-zm-elliptic-weight-doubling-discriminant-squareclass-s4}）。

\begin{conclusion}[Lee--Yang 立方域上分歧点的显式坐标化与隐藏平方结构]\label{con:xi-terminal-zm-leyang-branchpoint-explicit-coordinates-hidden-square}
令 \(y_0\) 为 \(P_{\mathrm{LY}}(y)=0\) 的任一根，并记 \(K:=\QQ(y_0)\)。
则在 \(K\) 上存在唯一分歧点 \(Q\in E(K)\) 满足
\[
y(Q)=y_0,\qquad \dd y(Q)=0.
\]
其坐标可由 \(y_0\) 显式表示为
\[
\boxed{\;
X(Q)=\frac{-512y_0^{2}-518y_0-5}{279},\qquad
Y(Q)=\frac{512y_0^{2}+1262y_0+377}{558}
=\frac{(512y_0^{2}+518y_0+5)^{2}-25947}{372\,\bigl(512y_0^{2}+518y_0+5\bigr)}\; }.
\qquad (25947=3^{3}\cdot 31^{2})
\]
特别地，\(X(Q)\) 满足立方证书 \(g(X(Q))=0\)，其中 \(g(X)=16X^{3}-9X^{2}+1\)（见结论 \ref{con:xi-terminal-zm-leyang-branch-elimination-origin-y-iota}）。
\par
此外，域 \(K\) 中存在一条非显然的“平方”恒等式：定义
\[
s:=\frac{-512y_0^{2}+226y_0+367}{93}\in K,
\]
则严格成立
\[
\boxed{\;s^{2}=64y_0^{2}+64y_0+25\ \text{于 }K\text{ 中}\; }.
\]
该恒等式给出一条显式平方证书，表明 Lee--Yang 立方域中出现的分歧点坐标可在不引入额外二次扩张的前提下完成闭式表达。
\end{conclusion}
\begin{proof}[证明要点]
在商环 \(\QQ[y,X]/(P_{\mathrm{LY}}(y))\) 中，对系统 \(\Pi(X,y)=0\)、\(\partial_{X}\Pi(X,y)=0\) 作 Gröbner 消元可得到一次消元子
\(
279X+512y^{2}+518y+5=0
\)，从而给出 \(X(Q)\) 的显式表达；
再由 \(y=X^{2}+Y-\tfrac12\) 回代得到 \(Y(Q)\)。
隐藏平方恒等式为同一三次商环中的代数恒等式检验。
可复现核验脚本：\texttt{scripts/exp\_fold\_zm\_elliptic\_leyang\_resy\_spectral\_decomposition\_audit.py}。证毕。
\end{proof}

\begin{conclusion}[Lee--Yang 分歧值处谱四次的精确退化：二重根与剩余二次因子的闭式分解]\label{con:xi-terminal-zm-leyang-branch-spectral-quartic-double-root-factorization}
令 \(y_0\) 为 \(P_{\mathrm{LY}}(y)=0\) 的任一根，并记 \(K:=\QQ(y_0)\)。
考虑专门化谱四次
\[
\Pi(\lambda,y_0)
=\lambda^{4}-\lambda^{3}-(2y_0+1)\lambda^{2}+\lambda+y_0(y_0+1)\ \in\ K[\lambda].
\]
则存在严格分解
\[
\boxed{\;
\Pi(\lambda,y_0)=\bigl(\lambda-X(Q)\bigr)^{2}\cdot Q_{y_0}^{(2)}(\lambda)\; },
\]
其中 \(Q\in E(K)\) 为结论 \ref{con:xi-terminal-zm-leyang-branchpoint-explicit-coordinates-hidden-square} 的唯一分歧点，且剩余二次因子可取为
\[
\boxed{\;
Q_{y_0}^{(2)}(\lambda)=\lambda^{2}
-\frac{1024y_0^{2}+1036y_0+289}{279}\,\lambda
+\frac{256y_0^{2}-578y_0-416}{279}\; }.
\]
该分解把 Lee--Yang 分歧值处 \(\Pi(\lambda,y_0)=0\) 的谱退化精确量化为：一处由分歧点给出的二重根与两条非分歧谱支所决定的剩余二次片。
\end{conclusion}
\begin{proof}[证明要点]
在 \(P_{\mathrm{LY}}(y_0)=0\) 的 Lee--Yang 立方域上，系统
\(
\Pi(\lambda,y_0)=\partial_{\lambda}\Pi(\lambda,y_0)=0
\)
在 \(K\) 中具有唯一解 \(\lambda=X(Q)\)，其中 \(Q\) 为结论 \ref{con:xi-terminal-zm-leyang-branchpoint-explicit-coordinates-hidden-square} 的分歧点。
因此 \(\lambda=X(Q)\) 为 \(\Pi(\lambda,y_0)\) 的二重根；并且由单分歧性可知其重数恰为 \(2\)。
将 \(\Pi(\lambda,y_0)\) 除以 \((\lambda-X(Q))^{2}\) 并在 \(K\) 中化简即可得到剩余二次因子的闭式系数。
可复现核验脚本：\texttt{scripts/exp\_fold\_zm\_elliptic\_leyang\_resy\_spectral\_decomposition\_audit.py}。证毕。
\end{proof}

\begin{conclusion}[两类二次伴随场的分化：临界域与 Lee--Yang 分歧域的算术阴影]\label{con:xi-terminal-zm-quadratic-companion-fields-separation}
该模型中出现两条彼此不同但同由椭圆判别式素数 \(37=\Delta(E)\) 驱动的二次伴随场：
\begin{enumerate}
  \item 对临界六次多项式 \(\Psi\)，其判别式平方自由核为 \(-37\)（结论 \ref{con:xi-terminal-zm-lattes-critical-sextic-4division-splitting-field}），因此
  \[
  \QQ(\sqrt{-37})\ \subset\ K_4
  \]
  必嵌入临界点分裂域 \(K_4/\QQ\)。
  \item 对 Lee--Yang 三次分歧（等价地，对 \(P_{\mathrm{LY}}\) 或结论 \ref{con:xi-terminal-zm-leyang-cubic-arithmetic} 的 \(S_3\) 分裂域），其判别式平方自由核为 \(-111=-3\cdot 37\)，因此
  \[
  \QQ(\sqrt{-111})\ \subset\ L_{\mathrm{LY}}
  \]
  必嵌入 Lee--Yang 分歧域的伽罗瓦闭包 \(L_{\mathrm{LY}}/\QQ\)。
\end{enumerate}
两者分别对应 \(-\Delta(E)\) 与 \(-3\Delta(E)\) 的二次阴影，
从而在最低阶二次信息上，动力学临界域（由 \(\Phi\) 的临界点控制）与分歧/Lee--Yang 域（由 \(y:E\to\mathbb P^{1}\) 的分歧控制）已发生不可约的算术分化。
\end{conclusion}
\begin{proof}[证明要点]
第一点由结论 \ref{con:xi-terminal-zm-elliptic-lattes-critical-sextic-galois} 的判别式平方类立即推出；
第二点由结论 \ref{con:xi-terminal-zm-leyang-cubic-arithmetic} 的判别式平方自由部分与不可约性推出，并由不可约三次的标准结构得到其唯一二次子域。
二次阴影的对比仅为代入 \(\Delta(E)=37\) 的算术重写。证毕。
\end{proof}

\begin{conclusion}[二次判别角色的最小同时消障场：\texorpdfstring{$\QQ(\sqrt{-111},\sqrt{-37})$}{Q(sqrt(-111),sqrt(-37))}]\label{con:xi-terminal-zm-quadratic-simultaneous-obstruction-field}
令 \(L_{\mathrm{LY}}\) 为 Lee--Yang 三次 \(P_{\mathrm{LY}}(y)\) 的分裂域，并令 \(K_4\) 为临界六次 \(\Psi(X)\)（结论 \ref{con:xi-terminal-zm-lattes-critical-sextic-4division-splitting-field}）的分裂域。
则由结论 \ref{con:xi-terminal-zm-quadratic-companion-fields-separation} 可知
\[
\QQ(\sqrt{-111})\subset L_{\mathrm{LY}},
\qquad
\QQ(\sqrt{-37})\subset K_4.
\]
定义双二次扩张
\[
F:=\QQ(\sqrt{-111},\sqrt{-37}).
\]
则 \(F/\QQ\) 为 \(V_4\)-伽罗瓦扩张，且为同时包含上述两条二次子域的最小数域：任何同时包含 \(\QQ(\sqrt{-111})\) 与 \(\QQ(\sqrt{-37})\) 的数域必包含 \(F\)。
因此，Lee--Yang 分歧域与临界域在二次层级上出现的两条判别角色在 \(F\) 上同时平凡化；其二次障碍被统一压缩为一个可枚举的双二次框架。
并且 \(F/\QQ\) 的分歧素数集合包含于
\[
\{2,3,37\}.
\]
\end{conclusion}
\begin{proof}[证明要点]
由结论 \ref{con:xi-terminal-zm-leyang-cubic-field-minimal-discriminant}，\(\QQ(\sqrt{-111})\) 为 \(L_{\mathrm{LY}}\) 的二次分辨子域；
由结论 \ref{con:xi-terminal-zm-elliptic-lattes-critical-sextic-galois}，\(\QQ(\sqrt{-37})\) 嵌入 \(K_4\)。
二者的合成域 \(F=\QQ(\sqrt{-111})\QQ(\sqrt{-37})\) 即为最小同时容纳两子域的扩张，且为双二次（故伽罗瓦群同构于 \(V_4\)）。
\par
分歧素数集合的包含性由二次域判别式给出：
\(\mathrm{Disc}\bigl(\QQ(\sqrt{-111})\bigr)=-111\) 与 \(\mathrm{Disc}\bigl(\QQ(\sqrt{-37})\bigr)=-148\)，
故二者分歧素数分别包含于 \(\{3,37\}\) 与 \(\{2,37\}\)；
而合成域的分歧素数只能来自二者分歧素数之并集，因而包含于 \(\{2,3,37\}\)。证毕。
\end{proof}

\begin{conclusion}[不同元 \texorpdfstring{$\delta$}{delta} 的迹对偶与一行反演：四次幂基上的系数可审计重构]\label{con:xi-terminal-zm-delta-trace-dual-inversion}
设 \(k:=\QQ(y)\) 且 \(K:=\QQ(E)\)。
在首一四次表示
\[
K\cong k[X]/(F),
\qquad
F(X,y):=X^4-X^3-(2y+1)X^2+X+y(y+1)
\]
中，令 \(X\in K\) 为 \(X\) 的像，并令
\[
\delta:=4XY+3X^2-1\in K
\]
为结论 \ref{con:xi-terminal-zm-leyang-ramification-critical-values} 的 Lee--Yang 雅可比元（从而 \(\dd y=\delta\,\omega\)），并记
\[
\delta_{\mathrm{diff}}:=\partial_XF(X,y)=1-3X^2-4XY=-\delta.
\]
则在扩张 \(K/k\) 中成立严格迹恒等式
\[
\Tr_{K/k}\!\left(\frac{X^m}{\delta_{\mathrm{diff}}}\right)=0\quad(m=0,1,2),
\qquad
\Tr_{K/k}\!\left(\frac{X^3}{\delta_{\mathrm{diff}}}\right)=1.
\]
写
\[
F(X,y)=X^4+a_3X^3+a_2X^2+a_1X+a_0
\qquad
(a_3=-1,\ a_2=-(2y+1),\ a_1=1,\ a_0=y(y+1)),
\]
并定义 Horner 截断多项式
\[
q_0:=1,\qquad
q_1:=X+a_3=X-1,\qquad
q_2:=X^2+a_3X+a_2=X^2-X-(2y+1),
\]
\[
q_3:=X^3+a_3X^2+a_2X+a_1=X^3-X^2-(2y+1)X+1.
\]
则两组 \(k\)-基
\[
\{1,X,X^2,X^3\},
\qquad
\left\{\frac{q_3}{\delta_{\mathrm{diff}}},\frac{q_2}{\delta_{\mathrm{diff}}},\frac{q_1}{\delta_{\mathrm{diff}}},\frac{q_0}{\delta_{\mathrm{diff}}}\right\}
\]
互为迹对偶基。
特别地，对任意 \(a\in K\) 有一行反演公式
\[
a=\sum_{i=0}^{3}\Tr_{K/k}\!\left(\frac{aq_i}{\delta_{\mathrm{diff}}}\right)X^{3-i}.
\]
由于 \(\delta_{\mathrm{diff}}=-\delta\)，亦可等价写成
\(
a=-\sum_{i=0}^{3}\Tr_{K/k}\!\left(\frac{aq_i}{\delta}\right)X^{3-i}
\)；
该重构通道表明幂基系数的分母结构完全由同一不同理想生成元（Lee--Yang 雅可比元）控制。
\end{conclusion}
\begin{proof}[证明要点]
由于 \(F\) 在 \(k\) 上可分，\(K/k\) 为可分扩张，且 \(\delta_{\mathrm{diff}}=\partial_XF(X,y)\) 生成不同理想（结论 \ref{con:xi-terminal-zm-leyang-discriminant-norm-identity}）。
设 \(X_1,\dots,X_4\) 为 \(F(\cdot,y)\) 在代数闭包中的四根，则对任意 \(m\ge 0\) 有
\[
\Tr_{K/k}\!\left(\frac{X^{m}}{\delta_{\mathrm{diff}}}\right)=\sum_{j=1}^{4}\frac{X_j^{m}}{F'(X_j)}.
\]
将部分分式分解
\(
1/F(T)=\sum_{j}(F'(X_j))^{-1}(T-X_j)^{-1}
\)
在 \(T=\infty\) 展开并比较 \(T^{-1}\) 的系数，即得
\(
\sum_j X_j^{m}/F'(X_j)=0
\)（\(m\le 2\)）与
\(
\sum_j X_j^{3}/F'(X_j)=1
\)。
将 \(k\)-基 \(\{1,X,X^2,X^3\}\) 与候选族 \(\{X^3/\delta_{\mathrm{diff}},X^2/\delta_{\mathrm{diff}},X/\delta_{\mathrm{diff}},1/\delta_{\mathrm{diff}}\}\) 的迹配对写成矩阵，可得其为对角为 \(1\)、上三角为 \(0\) 的单位下三角矩阵；
其逆矩阵给出迹对偶基在上述候选族中的唯一线性校正。
对本四次模型，该校正恰由 Horner 截断 \(q_0,q_1,q_2,q_3\) 给出，从而得到所示迹对偶基与反演公式。
可复现核验脚本：\texttt{scripts/exp\_fold\_zm\_elliptic\_trace\_dual\_delta\_audit.py}。证毕。
\end{proof}

\begin{conclusion}[有理点平移在重量坐标上的双四次对应]\label{con:xi-terminal-zm-elliptic-translation-q-biquartic-correspondence}
沿用结论 \ref{con:xi-terminal-zm-elliptic-normalization-y-iota-ramification-portrait} 的椭圆曲线
\[
E/\QQ:\quad Y^{2}=X^{3}-X+\tfrac14,
\qquad
y:=X^{2}+Y-\tfrac12\in\QQ(E),
\]
并取有理点
\[
Q:=\Bigl(1,-\tfrac12\Bigr)\in E(\QQ),
\qquad
\tau_{Q}:E\to E,\ \ P\longmapsto P+Q
\]
为平移自同构。
对任意点 \(P\in E(\overline{\QQ})\)（使下式有意义），记
\[
u:=y(P),\qquad v:=y\bigl(\tau_{Q}(P)\bigr)=\tau_Q^{*}y(P).
\]
则 \((u,v)\in\PP^{1}\times\PP^{1}\) 落在一条既约 \((4,4)\)-曲线 \(C_{Q}\) 上；并且 \(C_{Q}\) 由如下双四次关系式在 \(\QQ[u,v]\) 中本原刻画：
\[
\begin{aligned}
0=\ &u^{4}v^{4}-4u^{3}v^{4}-37u^{3}v^{3}-8u^{3}v^{2}+6u^{2}v^{4}-129u^{2}v^{3}+36u^{2}v^{2}+20u^{2}v+16u^{2} \\
&-4uv^{4}-79uv^{3}+44uv^{2}+79uv+9u+v^{4}-11v^{3}+24v^{2}+36v.
\end{aligned}
\]
等价地，作为扩张 \(\QQ(E)/\QQ(u)\) 中的元素，函数 \(v=\tau_Q^*y\) 在 \(\QQ(u)\) 上满足上述四次整依赖关系。
\end{conclusion}
\begin{proof}[证明要点]
将 \(u=y(P)\) 与 \(v=y(P+Q)\) 视为 \(P\) 的有理函数，并在函数域 \(\QQ(E)\) 中对 \((X,Y)\) 消元即可得到 \((u,v)\) 的本原消元多项式；
该多项式在 \(\QQ[u,v]\) 中既约，故其零点集的 Zariski 闭包即为 \(C_Q\)。
可复现核验脚本：\texttt{scripts/exp\_fold\_zm\_elliptic\_translation\_q\_biquartic\_audit.py}。证毕。
\end{proof}

\begin{conclusion}[双四次图像曲线的归一化与椭圆坐标的显式回收]\label{con:xi-terminal-zm-elliptic-translation-q-biquartic-birational-recovery}
令 \(C_Q\subset \PP^{1}\times\PP^{1}\) 为结论 \ref{con:xi-terminal-zm-elliptic-translation-q-biquartic-correspondence} 的双四次曲线，并在其光滑部分上保持坐标 \((u,v)\)。
则 \(C_Q\) 的归一化与 \(E\) 在 \(\QQ\) 上双有理同构：在开集 \(\{u\neq 0\}\cap\{D(u,v)\neq 0\}\) 上，
由下述显式公式可恢复椭圆坐标
\[
X=\frac{N(u,v)}{D(u,v)},
\qquad
Y=u-X^{2}+\tfrac12,
\]
其中
\[
\begin{aligned}
N(u,v)=\ &2u^{3}v^{3}+23u^{2}v^{3}+4u^{2}v^{2}+8u^{2}v+28uv^{3}+17uv^{2}+3uv-u \\
&\quad+11v^{3}+7v^{2}-4v,\\[2mm]
D(u,v)=\ &13u^{2}v^{3}+9u^{2}v^{2}+38uv^{3}+6uv^{2}-7uv-4u \\
&\quad+13v^{3}-3v^{2}-16v.
\end{aligned}
\]
因此，平移对应的图像曲线 \(C_Q\) 的归一化即为 \(E\)，从而 \((u,v)\) 给出 \(E\) 的一组低次数双射模型。
\end{conclusion}
\begin{proof}[证明要点]
由消元得到的双四次模型与 \(E\) 的函数域同构均为显式可核验陈述：在 \(\QQ(E)\) 中以 \(u=y\)、\(v=\tau_Q^*y\) 生成同一子域，并可解出 \(X\) 为 \(\QQ(u,v)\) 的有理函数，从而得到所示反解式。
对开集上的互逆性可通过代回 \(E\) 的 Weierstrass 方程与双四次约束作一致性检验。
可复现核验脚本：\texttt{scripts/exp\_fold\_zm\_elliptic\_translation\_q\_biquartic\_audit.py}。证毕。
\end{proof}

\begin{conclusion}[平移双四次对应的判别式分解与 Lee--Yang 分歧集合的一致性]\label{con:xi-terminal-zm-elliptic-translation-q-biquartic-discriminant}
令 \(R(u,v)\in\QQ[u,v]\) 表示结论 \ref{con:xi-terminal-zm-elliptic-translation-q-biquartic-correspondence} 中定义 \(C_Q\) 的双四次多项式，并将其视为关于 \(v\) 的四次多项式 \(R\in\QQ[u][v]\)。
再记 Lee--Yang 三次
\[
P_{\mathrm{LY}}(u):=256u^{3}+411u^{2}+165u+32\in\ZZ[u].
\]
则在 \(\QQ[u]\) 中，判别式具有严格分解
\[
\mathrm{Disc}_{v}\bigl(R(u,v)\bigr)
=-u(u-1)\,P_{\mathrm{LY}}(u)\cdot S(u)^{2},
\]
其中
\[
S(u)=1024u^{6}+20543u^{5}+92840u^{4}+180215u^{3}+167685u^{2}+75958u+13671\in\ZZ[u].
\]
因此，该平移对应作为 \(u\)-直线上的四次覆叠，仅在 \(u=0,1\) 及 \(P_{\mathrm{LY}}(u)=0\) 的三根处发生分歧；
并且其判别式平方自由核严格等于 \(-u(u-1)P_{\mathrm{LY}}(u)\)，与重量覆盖 \(\pi_y:E\to\PP^{1}\) 的 Lee--Yang 分歧因子完全一致。
\end{conclusion}
\begin{proof}[证明要点]
将 \(R(u,v)\) 视为 \(v\) 的四次多项式并在 \(\ZZ[u]\) 内计算判别式，再作因式分解即可得到所示恒等式；
其中平方因子以 \(S(u)^2\) 的形式出现，故平方自由部分仅由 \(-u(u-1)P_{\mathrm{LY}}(u)\) 控制。
可复现核验脚本：\texttt{scripts/exp\_fold\_zm\_elliptic\_translation\_q\_biquartic\_audit.py}。证毕。
\end{proof}


\begin{theorem}[双参数四次模型的分歧判别式闭式与 \texorpdfstring{$C_{3}$}{C3} 协变律]\label{thm:xi-terminal-zm-translation-t-branch-discriminant-c3-closed}
\leanverified{paper\_xi\_terminal\_zm\_translation\_t\_branch\_discriminant\_c3\_closed}
定义
\[
F_t(\mu,u):=\mu^4-\mu^3-(tu+1)\mu^2+\mu(tu-u^3)+tu.
\]
令
\[
\operatorname{Disc}_{\mu}\!\bigl(F_t(\mu,u)\bigr)=u\,D_t(u),
\qquad D_t(u)\in\ZZ[t,u].
\]
则 \(\deg_u D_t=11\)，并且
\[
\boxed{
\begin{aligned}
D_t(u)=\;&-27u^{11}+90t\,u^9+(4t^3-22)u^8-239t^2u^7-(8t^4+70t)u^6\\
&+(236t^3+5)u^5+(20t^5-164t^2)u^4-(120t^4+108t)u^3\\
&+220t^3u^2-120t^2u+20t.
\end{aligned}}
\]
此外，对任意满足 \(\zeta^3=1\) 的三次单位根 \(\zeta\)，有严格协变关系
\[
\boxed{\,D_{\zeta t}(u)=\zeta\,D_t(\zeta u)\, }.
\]
\end{theorem}
\begin{proof}
由判别式定义
\[
\operatorname{Disc}_{\mu}(F)=(-1)^{n(n-1)/2}\operatorname{Res}_{\mu}(F,\partial_{\mu}F),\qquad n=4,
\]
对 \(F_t\) 直接消元可得
\[
\operatorname{Disc}_{\mu}(F_t)=u\,D_t(u),
\]
并得到上述闭式系数。
\par
对协变性，先验证恒等式
\[
F_{\zeta t}(\mu,\zeta^{-1}u)=F_t(\mu,u),
\]
因为 \(u^3\) 在 \(u\mapsto \zeta^{-1}u\) 下不变，且 \(tu\) 在 \((t,u)\mapsto(\zeta t,\zeta^{-1}u)\) 下不变。
对两边取 \(\mu\)-判别式：
\[
u\,D_t(u)=\zeta^{-1}u\,D_{\zeta t}(\zeta^{-1}u).
\]
将 \(u\) 替换为 \(\zeta u\) 即得
\[
D_{\zeta t}(u)=\zeta D_t(\zeta u).
\]
可复现核验脚本：\texttt{scripts/exp\_fold\_zm\_elliptic\_translation\_t\_discriminant\_c3\_audit.py}。证毕。
\end{proof}

\begin{theorem}[\texorpdfstring{$\mu_3$}{mu3} 等变作用与判别式权重律]\label{thm:xi-terminal-zm-translation-t-branch-discriminant-c3-mu3-weight}
令 \(\zeta\) 为原始三次单位根，并定义
\[
\gamma:\ (t,u)\longmapsto(\zeta^2 t,\zeta u).
\]
则有严格恒等式
\[
F_{\zeta^2 t}(\mu,\zeta u)=F_t(\mu,u),
\]
从而
\[
D_{\zeta^2 t}(\zeta u)=\zeta^2D_t(u).
\]
记
\[
\Delta(t):=\operatorname{Disc}_u\!\bigl(D_t(u)\bigr).
\]
则
\[
\operatorname{Disc}_{u}\!\bigl(D_{\zeta^2 t}(u)\bigr)=\zeta^2\,\operatorname{Disc}_{u}\!\bigl(D_t(u)\bigr),
\qquad\text{即}\qquad
\Delta(\zeta^2 t)=\zeta^2\Delta(t).
\]
并且存在 \(\Phi(s)\in\ZZ[s]\) 使得
\[
\Delta(t)=t^4\,\Phi(t^3).
\]
\end{theorem}
\begin{proof}
在变换 \(\gamma\) 下，
\[
(\zeta^2 t)(\zeta u)=tu,\qquad (\zeta u)^3=u^3.
\]
代入 \(F_t(\mu,u)\) 的显式表达即得 \(F_{\zeta^2 t}(\mu,\zeta u)=F_t(\mu,u)\)。
对两边取 \(\mu\)-判别式，得到
\[
u\,D_t(u)=\operatorname{Disc}_{\mu}\!\bigl(F_t\bigr)
=\operatorname{Disc}_{\mu}\!\bigl(F_{\zeta^2 t}(\cdot,\zeta u)\bigr)
=(\zeta u)\,D_{\zeta^2 t}(\zeta u),
\]
故 \(D_{\zeta^2 t}(\zeta u)=\zeta^2D_t(u)\)。
将 \(u\mapsto \zeta^2u\) 代回，得
\[
D_{\zeta^2 t}(u)=\zeta^2D_t(\zeta^2u).
\]
记 \(n:=\deg_u D_t=11\)。判别式缩放律
\[
\operatorname{Disc}_u(cP)=c^{2n-2}\operatorname{Disc}_u(P),\qquad
\operatorname{Disc}_u(P(au))=a^{n(n-1)}\operatorname{Disc}_u(P)
\]
给出
\[
\Delta(\zeta^2 t)
=(\zeta^2)^{2n-2}(\zeta^2)^{n(n-1)}\Delta(t)
=\zeta^{4n-4+2n(n-1)}\Delta(t)
=\zeta^{260}\Delta(t)
=\zeta^2\Delta(t).
\]
写 \(\Delta(t)=\sum_k a_k t^k\)。由 \(\Delta(\zeta^2 t)=\zeta^2\Delta(t)\) 可知 \(a_k\neq 0\Rightarrow k\equiv 1\pmod 3\)，故
\[
\Delta(t)=t\,\Psi(t^3),\qquad \Psi\in\ZZ[s].
\]
再结合定理 \ref{thm:xi-terminal-zm-translation-t-branch-discriminant-c3-delta-factor} 的显式分解中 \(t^4\) 的精确因子，得 \(\Delta(t)=t^4\Phi(t^3)\)。证毕。
\end{proof}

\begin{theorem}[二次判别式 \texorpdfstring{$\Delta(t)=\operatorname{Disc}_u(D_t)$}{Delta(t)=Disc_u(Dt)} 的完全闭式分解]\label{thm:xi-terminal-zm-translation-t-branch-discriminant-c3-delta-factor}
令
\[
\Delta(t):=\operatorname{Disc}_u\!\bigl(D_t(u)\bigr)\in\ZZ[t].
\]
记 \(s:=t^3\)，并定义
\[
A(s):=80s^2+16424s-84375,
\]
\[
B(s):=256s^2-1312s+3125,
\]
\[
C(s):=8000s^4+429800s^3+96378087s^2-2034256s-64.
\]
则有严格恒等式
\[
\boxed{
\Delta(t)
=2^{20}3^{9}5^{2}\,t^4\,A(t^3)\,B(t^3)^2\,C(t^3)^3.}
\]
\end{theorem}
\begin{proof}
由 \(\deg_u D_t=11\) 与判别式定义
\[
\operatorname{Disc}_u(D_t)=(-1)^{m(m-1)/2}\operatorname{Res}_u(D_t,\partial_u D_t),\qquad m=11,
\]
对 \(D_t\) 与 \(\partial_u D_t\) 进行 Sylvester 消元并在 \(\ZZ[t]\) 中作不可约分解，得到
\[
\Delta(t)=515978035200\;t^4\,A(t^3)\,B(t^3)^2\,C(t^3)^3.
\]
再由
\[
515978035200=2^{20}3^{9}5^{2}
\]
即得所述闭式。该分解与定理
\ref{thm:xi-terminal-zm-translation-t-branch-discriminant-c3-mu3-weight}
中的 \(\mu_3\) 权重律完全一致。证毕。
\end{proof}

\begin{corollary}[正实参数半轴上的分歧碰撞仅发生于两处]\label{cor:xi-terminal-zm-translation-t-branch-discriminant-c3-real-collision}
\leanverified{paper\_xi\_terminal\_zm\_translation\_t\_branch\_discriminant\_c3\_real\_collision}
记
\[
\Delta(t):=\operatorname{Disc}_u\!\bigl(D_t(u)\bigr).
\]
则对 \(t>0\)，有
\[
\Delta(t)=0
\quad\Longleftrightarrow\quad
t=\sqrt[3]{s_A}\ \text{或}\ t=\sqrt[3]{s_C},
\]
其中 \(s_A>0\) 为 \(A(s)=0\) 的唯一正根，\(s_C>0\) 为 \(C(s)=0\) 的唯一正根。特别地，
\[
t>0\ \Longrightarrow\ B(t^3)\neq 0.
\]
\end{corollary}
\begin{proof}
由定理 \ref{thm:xi-terminal-zm-translation-t-branch-discriminant-c3-delta-factor}，
\[
\Delta(t)=0
\iff
t=0\ \text{或}\ A(t^3)=0\ \text{或}\ B(t^3)=0\ \text{或}\ C(t^3)=0.
\]
在 \(t>0\) 条件下，仅需讨论后三者。对 \(B(s)\) 计算判别式：
\[
(-1312)^2-4\cdot 256\cdot 3125=-1478656<0,
\]
故 \(B(s)\) 在 \(\RR\) 上无根，从而 \(B(t^3)\neq 0\) 对一切 \(t\in\RR\) 成立。
再看 \(A(s)=80s^2+16424s-84375\)，其符号序列为 \((+,+,-)\)，由 Descartes 符号法则知正根至多一个；且 \(A(0)=-84375<0\)、\(A(s)\to+\infty\)（\(s\to+\infty\)），故恰有一正根 \(s_A\)。
类似地，\(C(s)=8000s^4+429800s^3+96378087s^2-2034256s-64\) 的符号序列为 \((+,+,+,-,-)\)，故正根至多一个；又 \(C(0)=-64<0\)、\(C(1)>0\)，故恰有一正根 \(s_C\)。
因 \(t\mapsto t^3\) 在 \(\RR_{>0}\) 上为双射，结论成立。证毕。
\end{proof}

\begin{proposition}[\texorpdfstring{$u=1$}{u=1} 进入分歧值的完整代数条件]\label{prop:xi-terminal-zm-translation-t-branch-discriminant-c3-u1}
\leanverified{paper\_xi\_terminal\_zm\_translation\_t\_branch\_discriminant\_c3\_u1}
对 \(u=1\)，有
\[
\boxed{
D_t(1)=0\iff (t-2)\bigl(20t^4-88t^3+284t^2+45t+22\bigr)=0.}
\]
因此在 \(\overline{\QQ}\) 上恰有 \(5\) 个参数使 \(u=1\) 为分歧值：一实根 \(t=2\) 与该四次多项式的四个根。
\end{proposition}
\begin{proof}
将 \(u=1\) 代入定理 \ref{thm:xi-terminal-zm-translation-t-branch-discriminant-c3-closed} 的闭式表达并因式分解，立即得到结论。
可复现核验脚本：\texttt{scripts/exp\_fold\_zm\_elliptic\_translation\_t\_discriminant\_c3\_audit.py}。证毕。
\end{proof}

\begin{theorem}[四次因子的严格正性与唯一实 \texorpdfstring{$u=1$}{u=1} 共振]\label{thm:xi-terminal-zm-translation-t-branch-discriminant-c3-u1-positivity}
\leanverified{paper\_xi\_terminal\_zm\_translation\_t\_branch\_discriminant\_c3\_u1\_positivity}
定义
\[
q(t):=20t^4-88t^3+284t^2+45t+22.
\]
则有严格平方和分解
\[
\boxed{
q(t)=20\Bigl(t^2-\frac{11}{5}t\Bigr)^2+\frac{936}{5}\Bigl(t+\frac{25}{208}\Bigr)^2+\frac{8027}{416}.}
\]
从而 \(q(t)>0\) 对一切 \(t\in\RR\) 成立。
因此
\[
t\in\RR,\ D_t(1)=0\ \Longrightarrow\ t=2.
\]
\end{theorem}
\begin{proof}
上式直接展开可核验。右端三项均非负且常数项严格为正，故 \(q(t)\) 在 \(\RR\) 上严格正。
结合命题 \ref{prop:xi-terminal-zm-translation-t-branch-discriminant-c3-u1} 即得唯一实解 \(t=2\)。证毕。
\end{proof}

\begin{proposition}[判别曲线奇点的消元刻画与唯一性]\label{prop:xi-terminal-zm-translation-t-branch-discriminant-c3-singular-elimination}
设
\[
\Gamma\subset\mathbb{A}^2_{t,u},\qquad \Gamma:\ D_t(u)=0.
\]
则 \(\Gamma\) 的奇点参数满足：
\begin{enumerate}
  \item \(\Gamma\) 存在奇点当且仅当
  \[
  \boxed{
  \bigl(256t^6-1312t^3+3125\bigr)\,
  \bigl(8000t^{12}+429800t^9+96378087t^6-2034256t^3-64\bigr)=0.}
  \]
  \item 等价地，在 \(t\)-线上由消元多项式
  \[
  \boxed{
  \bigl(256t^6-1312t^3+3125\bigr)\,
  \bigl(8000t^{12}+429800t^9+96378087t^6-2034256t^3-64\bigr)^2=0
  }
  \]
  刻画；其中平方指数给出消元重数。
  \item 对每个满足上述条件的 \(t\)，奇点的 \(u\)-坐标唯一，且由
  \[
  \boxed{
  u=\frac{t^2\,P_9(t^3)}{C}}
  \]
  给出，其中
  \[
  C=
  1146919918223400883707609447433426476251203385787440489213053015719484911581940154368000,
  \]
  且
  \[
  \boxed{
  \begin{aligned}
  P_9(s)=\;&
  17864502295709707825431356157528873283413020214750253582630010725020746370305048576000000\,s^9\\
  &+1827985552391858929673625577699869358461325615692919356153396939222537398020953119539200000\,s^8\\
  &+472380765502860761689023714837178856458962674487907763054197839296312918358467091455103488000\,s^7\\
  &+20669312921272651046479667688582870966622217807989339629125906175063486375887028759114636876800\,s^6\\
  &+2479717207964943525348394675715513569066639511983825407996293452948613349769192552883221946978496\,s^5\\
  &-13112777363415463937040364405348120900527714677007038191152130648141031999868546590632040898213488\,s^4\\
  &+32206159583214755787286512483786868328764234040970721269159053665093307376534469693668131746127017\,s^3\\
  &-1341350108695509062380855870071374542799221268567807990255827240403545308460434755576190878625288\,s^2\\
  &+14031773541681724880812976025532230691811442528685420734972982181307286334826738788279124405440\,s\\
  &+1142522353125628145251947755351541775770680303528808392151252348593203216967111461525699072.
  \end{aligned}}
  \]
\end{enumerate}
\end{proposition}
\begin{proof}
对理想
\[
I:=\langle D_t(u),\partial_u D_t(u),\partial_t D_t(u)\rangle\subset\QQ[u,t]
\]
取字典序 \(u\succ t\) 的 Gröbner 基。
其 \(t\)-消元子恰为
\[
\bigl(256t^6-1312t^3+3125\bigr)\bigl(8000t^{12}+429800t^9+96378087t^6-2034256t^3-64\bigr)^2.
\]
故奇点参数集合由其平方自由部分给出，即第 1 条。
\par
同一 Gröbner 基含一条关于 \(u\) 的一次关系
\[
Cu-t^2P_9(t^3)=0,
\]
其中 \(C\neq 0\) 且 \(P_9\) 如上，故每个可行 \(t\) 对应唯一 \(u\)。
可复现核验脚本：\texttt{scripts/exp\_fold\_zm\_elliptic\_translation\_t\_discriminant\_c3\_audit.py}。证毕。
\end{proof}

\begin{theorem}[\texorpdfstring{$t=1$}{t=1} 处的分支多项式具有满对称伽罗瓦群 \texorpdfstring{$S_{11}$}{S11}]\label{thm:xi-terminal-zm-translation-t-branch-discriminant-c3-t1-galois-s11}
取 \(t=1\)，令
\[
\begin{aligned}
P(u):=D_1(u)=\;&-27u^{11}+90u^9-18u^8-239u^7-78u^6+241u^5\\
&-144u^4-228u^3+220u^2-120u+20.
\end{aligned}
\]
则 \(P\) 在 \(\QQ\) 上不可约，且其分裂域伽罗瓦群满足
\[
\operatorname{Gal}(P/\QQ)\cong S_{11}.
\]
\end{theorem}
\begin{proof}
记 \(G:=\operatorname{Gal}(P/\QQ)\le S_{11}\)。
\begin{enumerate}
  \item \textbf{不可约性与原始性。}
  将 \(P\) 在 \(p=19\) 下约化，得到 \(\overline P\in\FF_{19}[u]\) 不可约（分解型为单一 \(11\) 次因子）。
  由 Gauss 引理与约化判据，\(P\) 在 \(\QQ[u]\) 不可约，故 \(G\) 在 \(11\) 个根上是传递群；又 \(11\) 为素数，故该传递作用自动原始。
  \item \textbf{三循环证书。}
  由 \(\operatorname{Disc}(P)\not\equiv 0\pmod{11}\)，\(p=11\) 为非分歧素数。
  在 \(p=11\) 下有分解型
  \[
  \overline P(u)=\overline f_3(u)\,\overline f_8(u),
  \]
  其中 \(\deg\overline f_3=3\)、\(\deg\overline f_8=8\) 且两因子均不可约。
  因而对应 Frobenius 元素循环型为 \((3)(8)\)，其八次幂给出一个三循环，故 \(G\) 含 \(3\)-循环。
  \item \textbf{Jordan 推论。}
  原始子群 \(G\le S_{11}\) 含 \(3\)-循环，且 \(3\le 11-3\)。
  由 Jordan 定理，\(A_{11}\subseteq G\)。
  \item \textbf{判别式非平方排除 \(A_{11}\)。}
  直接计算
  \[
  \operatorname{Disc}(P)=-127646260798715703506840120525857458033760665600<0.
  \]
  该数在 \(\QQ\) 中非平方，故 \(G\nsubseteq A_{11}\)。
  结合 \(A_{11}\subseteq G\le S_{11}\)，得 \(G=S_{11}\)。
\end{enumerate}
模 \(19\)/模 \(11\) 分解型与判别式值可由脚本
\texttt{scripts/exp\_fold\_zm\_elliptic\_translation\_t\_discriminant\_c3\_audit.py}
独立复核。
证毕。
\end{proof}

\begin{conjecture}[泛型 \texorpdfstring{$S_{11}$}{S11} 稳定性]\label{conj:xi-terminal-zm-translation-t-branch-discriminant-c3-galois-s11}
设 \(t\in\QQ\setminus\{0,2\}\)，并满足 \(\Delta(t)\neq 0\)（即 \(D_t(u)\) 无重根）。
猜想：
\[
\boxed{
D_t(u)\ \text{在}\ \QQ[u]\ \text{上不可约，且}\ \operatorname{Gal}\bigl(D_t/\QQ\bigr)\cong S_{11}.}
\]
定理 \ref{thm:xi-terminal-zm-translation-t-branch-discriminant-c3-t1-galois-s11} 给出该猜想的一个严格基点 \(t=1\)。
若该猜想成立，则除显式可描述的异常参数族外，\(D_t(u)\) 在 \(\QQ\)-参数线上保持满对称群复杂度。
\end{conjecture}



\begin{theorem}[Lee--Yang 根处平方根系数 \texorpdfstring{$\kappa_*^2$}{kappa*^2} 的三次数域刻画与判别式]\label{thm:xi-terminal-zm-kappa-square-cubic-field-s3}\label{prop:xi-terminal-zm-branch-puiseux-coefficient-cubic}
\leanverified{paper\_xi\_terminal\_zm\_kappa\_square\_cubic\_field\_s3}
\leanverified{paper\_xi\_terminal\_zm\_branch\_puiseux\_coefficient\_cubic}
沿用谱四次
\[
\Pi(\lambda,y):=\lambda^4-\lambda^3-(2y+1)\lambda^2+\lambda+y(y+1)
\]
及其分歧集合 \(B_{\mathrm{fin}}=\{0,1\}\cup\{P_{\mathrm{LY}}(y)=0\}\) 的记号（结论 \ref{con:xi-terminal-zm-discriminant-leyang}）。
取任一非平凡分歧值 \(y_0\in\CC\) 满足 \(P_{\mathrm{LY}}(y_0)=0\)，并令 \(\lambda_0\) 为 \(\Pi(\lambda,y_0)=0\) 的唯一二重根（等价地 \(\Pi(\lambda_0,y_0)=\partial_\lambda\Pi(\lambda_0,y_0)=0\)）。
则存在局部参数 \(t\) 使 \(y-y_0=t^2\)，并且并合的两条谱支满足平方根型 Puiseux 展开
\[
\lambda=\lambda_0\pm c_0\,t+O(t^2),
\qquad
c_0^{2}=-\frac{2\,\partial_{y}\Pi(\lambda_0,y_0)}{\partial_{\lambda}^{2}\Pi(\lambda_0,y_0)}.
\]
定义
\[
u:=c_0^2=\kappa_*^2.
\]
在该 Lee--Yang 分歧支上进一步有
\[
\lambda_0\ \text{满足}\ 16\lambda_0^3-9\lambda_0^2+1=0
\qquad(\text{结论 \ref{con:xi-terminal-zm-leyang-lambda-cubic}}),
\]
并且 \(u\) 可化简为
\[
\boxed{\;
u=\frac{2}{3}\cdot\frac{3\lambda_0^2-1}{\lambda_0^2-1}\;}
\]
从而
\[
\boxed{\;324u^{3}-540u^{2}+333u-74=0\; }.
\]
该三次在 \(\QQ\) 上不可约，且其判别式为
\[
\mathrm{Disc}(324u^{3}-540u^{2}+333u-74)=-2^{6}\cdot 3^{9}\cdot 37.
\]
因而分裂域伽罗瓦群同构于 \(S_3\)。
此外
\[
\lambda_0=\frac{3(2u-1)}{4(3u-2)},
\]
故 \(\QQ(u)=\QQ(\lambda_0)\)。
\end{theorem}
\begin{proof}[证明要点]
平方根型 Puiseux 与系数公式由 \(\Pi(\lambda_0,y_0)=\partial_\lambda\Pi(\lambda_0,y_0)=0\) 处的二元 Taylor 展开得到，见结论 \ref{con:xi-terminal-zm-branch-puiseux-uniformization} 与定理 \ref{thm:xi-terminal-zm-leyang-elliptic-structure} 中的一般公式。
在 Lee--Yang 分歧支 \(P_{\mathrm{LY}}(y_0)=0\) 上，由结论 \ref{con:xi-terminal-zm-leyang-lambda-cubic} 得 \(\lambda_0\) 满足 \(16\lambda^3-9\lambda^2+1=0\)，并可将
\(
u=-2\,\partial_y\Pi/\partial_\lambda^2\Pi
\)
化简为所示有理式。
\par
由
\(
u=\frac{2}{3}\frac{3\lambda_0^2-1}{\lambda_0^2-1}
\)
可直接解得
\[
\lambda_0^2=\frac{3u-2}{3(u-2)}.
\]
又因 \(\lambda_0\neq 0\)，将 \(16\lambda_0^3-9\lambda_0^2+1=0\) 除以 \(\lambda_0^2\) 得
\[
\lambda_0=\frac{9\lambda_0^2-1}{16\lambda_0^2}
=\frac{3(2u-1)}{4(3u-2)}.
\]
将其代回 \(16\lambda_0^3-9\lambda_0^2+1=0\) 并清分母，即得
\[
324u^{3}-540u^{2}+333u-74=0.
\]
有理根检验排除一次因子，故该三次在 \(\QQ\) 上不可约。
其判别式按三次判别式公式计算为
\(
-2^{6}\cdot 3^{9}\cdot 37
\)，非平方。
对不可约三次，判别式非平方当且仅当分裂域伽罗瓦群为 \(S_3\)，故结论成立。
最后，由上式 \(\lambda_0=\frac{3(2u-1)}{4(3u-2)}\) 得 \(\lambda_0\in\QQ(u)\)，而 \(u\in\QQ(\lambda_0)\) 由定义即得，故 \(\QQ(u)=\QQ(\lambda_0)\)。证毕。
\end{proof}

\begin{corollary}[Lee--Yang 分歧数据在 \texorpdfstring{$\kappa_*^2$}{kappa*^2} 下的二次重构]\label{cor:xi-terminal-zm-leyang-branchdata-quadratic-reconstruction}
\leanverified{paper\_xi\_terminal\_zm\_leyang\_branchdata\_quadratic\_reconstruction}
沿用定理 \ref{thm:xi-terminal-zm-kappa-square-cubic-field-s3} 的记号，并令 \(\beta_*\) 为定理 \ref{thm:xi-terminal-zm-leyang-elliptic-structure} 中 Puiseux 展开
\[
\lambda(y)=\lambda_0\pm \kappa_*(y-y_0)^{1/2}+\beta_*(y-y_0)+O\bigl((y-y_0)^{3/2}\bigr)
\]
的线性项系数。
则在三次数域 \(K:=\QQ(u)=\QQ(\lambda_0)=\QQ(y_0)\) 内成立以下恒等式：
\[
16\lambda_0+31=108u(1-u),
\]
\[
256y_0+1091=108u(25-21u),
\]
\[
72\beta_*+22=9u(10u-1).
\]
等价地，
\[
\lambda_0=-\frac{31}{16}+\frac{27}{4}u-\frac{27}{4}u^2,\qquad
y_0=-\frac{1091}{256}+\frac{675}{64}u-\frac{567}{64}u^2,\qquad
\beta_*=-\frac{11}{36}-\frac{1}{8}u+\frac{5}{4}u^2.
\]
\end{corollary}
\begin{proof}
由定理 \ref{thm:xi-terminal-zm-kappa-square-cubic-field-s3}，\(u\) 满足三次关系 \(324u^{3}-540u^{2}+333u-74=0\)，且 \(\lambda_0\in\QQ(u)\)。
在商环 \(\QQ[u]/(324u^{3}-540u^{2}+333u-74)\) 中可直接核验
\(
16\lambda_0+31-108u(1-u)=0
\)，从而得到 \(\lambda_0\) 的二次表达式。
将其代入定理 \ref{thm:xi-terminal-zm-leyang-elliptic-structure} 的临界值显式式
\(
y_0=\frac{4\lambda_0^3-3\lambda_0^2-2\lambda_0+1}{4\lambda_0}
\)
与 \(\beta_*\) 的有理表达式（同一定理），并在同一商环内化简，即得其余两式。
\end{proof}

\begin{corollary}[\texorpdfstring{$\beta_*$}{beta*} 的极小多项式与显式整性化]\label{cor:xi-terminal-zm-beta-minpoly-integralization}
\leanverified{paper\_xi\_terminal\_zm\_beta\_minpoly\_integralization}
沿用上文记号。则 \(\beta_*\) 满足不可约三次方程
\[
419904\,\beta^{3}+93312\,\beta^{2}+72279\,\beta-1112=0.
\]
并且令
\[
T:=162u,\qquad S:=52488\,\beta_*,
\]
则 \(T,S\) 为 \(K\) 的代数整数，且分别满足首一整系数多项式
\[
T^{3}-270T^{2}+26973T-971028=0,
\]
\[
S^{3}+11664S^{2}+474222519S-382943630016=0.
\]
\end{corollary}
\begin{proof}
由推论 \ref{cor:xi-terminal-zm-leyang-branchdata-quadratic-reconstruction}，\(\beta_*\in\QQ(u)\) 且为 \(u\) 的二次多项式。
消去 \(u\)（或等价地在 \(\QQ[u]/(324u^{3}-540u^{2}+333u-74)\) 中化简）得到 \(\beta_*\) 的三次整系数关系；其不可约性由有理根判别排除一次因子。
对 \(T=162u\) 与 \(S=52488\beta_*\) 作变量伸缩并清分母，可得到所示首一整系数多项式，从而 \(T,S\) 为代数整数。
\end{proof}

\begin{proposition}[Lee--Yang 三次因子的二重覆盖与分歧核的完全拉回分解]\label{prop:xi-terminal-zm-leyang-double-cover-pullback-factorization}
\leanverified{paper\_xi\_terminal\_zm\_leyang\_double\_cover\_pullback\_factorization}
设
\[
g(u):=324u^{3}-540u^{2}+333u-74,\qquad
\iota(u):=\frac{25}{21}-u,
\]
并定义二次映射
\[
y(u):=-\frac{1091}{256}+\frac{675}{64}u-\frac{567}{64}u^2.
\]
则 \(y(\iota(u))=y(u)\)，并且 Lee--Yang 三次因子与全有限分歧核在 \(u\)-坐标下满足显式分解
\[
P_{\mathrm{LY}}(y(u))=\frac{3^{4}\,7^{3}}{2^{14}}\; g(u)\,g(\iota(u)),
\]
\[
y(u)\bigl(y(u)-1\bigr)P_{\mathrm{LY}}(y(u))
=\frac{3^{5}\,7^{3}}{2^{30}}\,
\bigl(2268u^2-2700u+1091\bigr)\,
\bigl(756u^2-900u+449\bigr)\,
g(u)\,g(\iota(u)).
\]
\end{proposition}
\begin{proof}
由推论 \ref{cor:xi-terminal-zm-leyang-branchdata-quadratic-reconstruction} 的恒等式 \(256y(u)+1091=108u(25-21u)\) 可见右端关于 \(u\mapsto\frac{25}{21}-u\) 对称，故 \(y(\iota(u))=y(u)\)。
将 \(P_{\mathrm{LY}}(y)=256y^3+411y^2+165y+32\) 代入 \(y=y(u)\) 得到 \(\QQ[u]\) 中的六次多项式。
由 \(u\in K\) 满足 \(g(u)=0\) 时 \(y(u)\) 为 Lee--Yang 分歧值（定理 \ref{thm:xi-terminal-zm-leyang-elliptic-structure}），可知 \(P_{\mathrm{LY}}(y(u))\) 被 \(g(u)\) 整除；对称性进一步给出其亦被 \(g(\iota(u))\) 整除。
两者均为三次且互素（由 \(g\) 的不可约性与 \(\iota\) 非平凡性），故
\(
P_{\mathrm{LY}}(y(u))=c\,g(u)g(\iota(u))
\)
对某 \(c\in\QQ^\times\) 成立；比较首项系数得到 \(c=\frac{3^{4}7^{3}}{2^{14}}\)。
第二式将 \(y(u)\) 与 \(y(u)-1\) 的分母清除并与第一式合并，再比较首项系数得到所示常数与二次因子。
\end{proof}

\begin{conclusion}[谱四次在实轴上的分岔型：实根数目仅在 \texorpdfstring{$y_{\mathrm{LY}},0,1$}{yLY,0,1} 处变化]\label{con:xi-terminal-zm-pi-real-root-bifurcation}
令
\[
\Pi(\lambda,y)=\lambda^4-\lambda^3-(2y+1)\lambda^2+\lambda+y(y+1),
\qquad
P_{\mathrm{LY}}(y)=256y^3+411y^2+165y+32,
\]
并令 \(y_{\mathrm{LY}}\in\RR\) 为 \(P_{\mathrm{LY}}(y)=0\) 的唯一实根（结论 \ref{con:xi-terminal-zm-discriminant-leyang}）。
则对任意实参数 \(y\in\RR\)，四次方程 \(\Pi(\lambda,y)=0\) 的实根数目满足如下分岔型（在开区间内根均互异）：
\begin{enumerate}
  \item 若 \(y<y_{\mathrm{LY}}\)，则 \(\Pi(\lambda,y)\) 在 \(\RR\) 上无实根；
  \item 若 \(y_{\mathrm{LY}}<y<0\)，则 \(\Pi(\lambda,y)\) 在 \(\RR\) 上恰有两条实根分支；
  \item 若 \(0<y<1\)，则 \(\Pi(\lambda,y)\) 在 \(\RR\) 上恰有四条实根分支；
  \item 若 \(y>1\)，则 \(\Pi(\lambda,y)\) 在 \(\RR\) 上恰有两条实根分支。
\end{enumerate}
并且在 \(y=y_{\mathrm{LY}},0,1\) 处分别出现一次二重根退化（对应结论 \ref{con:xi-terminal-zm-discriminant-leyang} 的分歧值集合）。
\end{conclusion}
\begin{proof}[证明要点]
由结论 \ref{con:xi-terminal-zm-discriminant-leyang}，
\(\mathrm{Disc}_{\lambda}(\Pi)=-y(y-1)P_{\mathrm{LY}}(y)\)，
故当 \(y\notin\{0,1\}\cup\{P_{\mathrm{LY}}(y)=0\}\) 时，四根皆为单根，实根数目在实轴上随 \(y\) 作局部常数变化。
另一方面，定理 \ref{thm:xi-terminal-zm-leyang-elliptic-structure} 给出 \(\mathrm{Disc}(P_{\mathrm{LY}})<0\)，从而 \(P_{\mathrm{LY}}\) 仅有一个实根 \(y_{\mathrm{LY}}\)；
故实轴上的分岔点严格为 \(y\in\{y_{\mathrm{LY}},0,1\}\)。
对每个开区间 \((-\infty,y_{\mathrm{LY}})\)、\((y_{\mathrm{LY}},0)\)、\((0,1)\)、\((1,\infty)\) 取代表点并用 Sturm 计数确定实根数目，即得到所示四段分岔型；在分岔点处由判别式为零得到二重根退化。
可复现核验脚本：\texttt{scripts/exp\_fold\_zm\_pi\_real\_bifurcation\_asymptotics\_audit.py}。证毕。
\end{proof}

\begin{conclusion}[Lee--Yang 实分歧点的 Cardano 根式表示：\texorpdfstring{$\sqrt{37}$}{sqrt(37)} 的判别式强制]\label{con:xi-terminal-zm-leyang-real-branchpoint-cardano}
沿用结论 \ref{con:xi-terminal-zm-leyang-lambda-cubic} 的记号，令
\[
g(\lambda):=16\lambda^3-9\lambda^2+1,
\qquad
y(\lambda):=\frac{4\lambda^3-3\lambda^2-2\lambda+1}{4\lambda}
=\lambda^2-\frac34\lambda-\frac12+\frac{1}{4\lambda}.
\]
设 \(\lambda_{\mathrm{LY}}\in\RR\) 为 \(g(\lambda)=0\) 的唯一实根，并令 \(y_{\mathrm{LY}}:=y(\lambda_{\mathrm{LY}})\)。
则 \(y_{\mathrm{LY}}\) 为 \(P_{\mathrm{LY}}(y)=0\) 的唯一实根（结论 \ref{con:xi-terminal-zm-discriminant-leyang}），并且在取实立方根分支的约定下存在根式闭式
\[
A:=\sqrt[3]{-101+16\sqrt{37}},\qquad B:=\sqrt[3]{-101-16\sqrt{37}},
\qquad
\boxed{\;\lambda_{\mathrm{LY}}=\frac{3+A+B}{16}\;}
\]
从而
\[
\boxed{\;y_{\mathrm{LY}}=\lambda_{\mathrm{LY}}^2-\frac34\lambda_{\mathrm{LY}}-\frac12+\frac{1}{4\lambda_{\mathrm{LY}}}\;}
\]
并且在 Tschirnhaus 平移 \(\lambda=t+\tfrac{3}{16}\) 下，
\[
g(\lambda)=0\ \Longleftrightarrow\ t^3-\frac{27}{256}\,t+\frac{101}{2048}=0,
\]
其 Cardano 判别式精确为
\[
\boxed{\;\Delta=\Bigl(\frac{101}{4096}\Bigr)^2+\Bigl(-\frac{9}{256}\Bigr)^3=\frac{37}{2^{16}}\;}
\]
从而 \(\sqrt{37}\) 在上述根式表达中由判别式机制强制出现。
\end{conclusion}
\begin{proof}[证明要点]
对 \(g(\lambda)=16\lambda^3-9\lambda^2+1\) 作平移 \(\lambda=t+\tfrac{3}{16}\) 得
\[
g(\lambda)=\frac{1}{128}\,\bigl(2048t^3-216t+101\bigr),
\]
从而等价于压低三次
\(
t^3-\frac{27}{256}\,t+\frac{101}{2048}=0
\)。
由 Cardano 公式，其唯一实根可写为
\[
t=\sqrt[3]{-\frac{101}{4096}+\frac{\sqrt{37}}{256}}+\sqrt[3]{-\frac{101}{4096}-\frac{\sqrt{37}}{256}}
=\frac{A+B}{16},
\]
于是 \(\lambda_{\mathrm{LY}}=t+\tfrac{3}{16}=\tfrac{3+A+B}{16}\)。
判别式 \(\Delta\) 由 \(p=-\tfrac{27}{256}\)、\(q=\tfrac{101}{2048}\) 代入
\(
\Delta=(q/2)^2+(p/3)^3
\)
得到，且等于 \(37/2^{16}>0\)，因此该三次仅有一个实根且 Cardano 根式必经 \(\sqrt{\Delta}=\sqrt{37}/2^{8}\)。
最后由结论 \ref{con:xi-terminal-zm-leyang-lambda-cubic} 的分歧参数化，非平凡分歧值满足 \(y=y(\lambda)\) 且 \(P_{\mathrm{LY}}(y)=0\)，从而 \(y_{\mathrm{LY}}=y(\lambda_{\mathrm{LY}})\) 为 \(P_{\mathrm{LY}}\) 的唯一实根。
可复现核验脚本：\texttt{scripts/exp\_fold\_zm\_pi\_real\_bifurcation\_asymptotics\_audit.py}。证毕。
\end{proof}

\begin{conclusion}[无穷远处的正实谱支分数幂渐近：\texorpdfstring{$y\to+\infty$}{y->+infty} 的 Puiseux jet]\label{con:xi-terminal-zm-pi-positive-real-branch-puiseux-infty}
设 \(y>1\)，并令 \(\lambda_{\pm}(y)\) 为 \(\Pi(\lambda,y)=0\) 在 \(\RR\) 上的两条实根分支，按大小约定
\[
0<\lambda_{-}(y)<\lambda_{+}(y).
\]
则当 \(y\to +\infty\) 沿实轴趋于无穷时，存在分数幂渐近展开
\[
\boxed{\;
\lambda_{\pm}(y)
= y^{1/2}\ \pm\ \frac{1}{2}y^{1/4}\ +\ \frac{1}{4}\ \pm\ \frac{7}{64}y^{-1/4}\ +\ \frac{37}{128}y^{-1/2}\ \mp\ \frac{729}{4096}y^{-3/4}\ +\ O(y^{-1})\;}
\]
特别地，
\[
\lambda_{+}(y)=y^{1/2}+\frac12 y^{1/4}+\frac14+O(y^{-1/4}).
\]
\end{conclusion}
\begin{proof}[证明要点（Puiseux 于无穷远）]
令 \(t=y^{-1/4}\)（故 \(y=t^{-4}\)）。
对根分支作 Puiseux 型设定
\[
\lambda(t)=t^{-2}+c_{1}t^{-1}+c_{0}+c_{-1}t+c_{-2}t^{2}+c_{-3}t^{3}+O(t^{4}),
\]
代入方程 \(\Pi(\lambda,y)=0\)，并乘以 \(t^{8}\) 以消去负幂，逐阶比较 \(t\)-系数得
\[
c_{1}=\pm\frac12,\qquad c_{0}=\frac14,\qquad c_{-1}=\pm\frac{7}{64},\qquad c_{-2}=\frac{37}{128},\qquad c_{-3}=\mp\frac{729}{4096},
\]
其中同号（分别异号）对应于 \(\lambda_{+}\)（分别 \(\lambda_{-}\)）的两支选择。
把 \(t=y^{-1/4}\) 回代即得所示展开。
可复现核验脚本：\texttt{scripts/exp\_fold\_zm\_pi\_real\_bifurcation\_asymptotics\_audit.py}。证毕。
\end{proof}

\begin{conclusion}[极端权重下 Perron 分支的双尺度 Puiseux：\texorpdfstring{$y\to0^+$}{y->0+} 与 \texorpdfstring{$y\to+\infty$}{y->+infty}]\label{con:xi-terminal-zm-perron-puiseux-two-scale}
令 \(\lambda_{+}(y)\) 表示对 \(y>0\) 的 Perron 根分支（最大正实根），并由 \(\lambda_{+}(1)=2\) 固定其与热力学主支的连通性（结论 \ref{con:xi-terminal-zm-lln-clt}）。
则 \(\lambda_{+}\) 在两端点呈现不可约分数幂尺度：
\begin{enumerate}
  \item 当 \(y\to 0^{+}\) 时，
  \[
  \lambda_{+}(y)
  =1+\frac{1}{\sqrt2}y^{1/2}+\frac58 y-\frac{43}{64\sqrt2}y^{3/2}-\frac1{16}y^{2}+O(y^{5/2})
  \qquad(\text{见结论 \ref{con:xi-terminal-zm-branch-puiseux-uniformization}}).
  \]

  \item 当 \(y\to+\infty\) 时，令 \(t=y^{-1/4}\)，则
  \[
  \lambda_{+}(y)
  =y^{1/2}+\frac12 y^{1/4}+\frac14+\frac{7}{64}y^{-1/4}+\frac{37}{128}y^{-1/2}-\frac{729}{4096}y^{-3/4}+O(y^{-1})
  \qquad(\text{见结论 \ref{con:xi-terminal-zm-pi-positive-real-branch-puiseux-infty}}).
  \]
\end{enumerate}
由此，定义密度函数
\[
\alpha(y):=\frac{y\,\lambda_{+}'(y)}{\lambda_{+}(y)}
=\left.\frac{\dd}{\dd s}\psi(s)\right|_{s=\log y},
\qquad
\psi(s):=\log\lambda_{+}(e^{s})-\log 2
\quad(\text{结论 \ref{con:xi-terminal-zm-lln-clt}}),
\]
则端点渐近为
\[
\alpha(y)=\frac{\sqrt2}{4}\,y^{1/2}+O(y)\quad(y\to0^{+}),
\qquad
\alpha(y)=\frac12-\frac18\,y^{-1/4}+O(y^{-1/2})\quad(y\to+\infty).
\]
特别地，\(\alpha(y)\) 以 \(y^{-1/4}\) 的速率逼近饱和值 \(1/2\)，该 \(1/4\) 指数由无穷远处谱方程的二重切触诱导，而非通常的线性扰动尺度。
\end{conclusion}



\begin{conclusion}[推前分布下的重量大数律与中心极限定律：均值与方差为有理常数]\label{con:xi-terminal-zm-lln-clt}
令 \(w_m(x)=d_m(x)/2^m\) 为均匀微态推前分布，并令 \(X\sim w_m\)。
记 \(H_m:=|X|_1\)，则对一切 \(y\in\CC\) 有概率母函数恒等式
\[
\EE\!\left[y^{H_m}\right]=\frac{Z_m(y)}{2^m}.
\]
令 \(\lambda_+(y)\) 表示 \(\Pi(\lambda,y)=0\) 在 \(y=1\) 的邻域内、满足 \(\lambda_+(1)=2\) 的解析特征支，则存在极限对数母函数
\[
\lim_{m\to\infty}\frac{1}{m}\log \EE\!\left(e^{sH_m}\right)
=\log\lambda_+(e^s)-\log 2.
\]
由此推出：
\begin{enumerate}
  \item \(m^{-1}H_m\to 5/18\)（依概率收敛）；
  \item \(\frac{H_m-\frac{5}{18}m}{\sqrt{m}}\Rightarrow \mathcal N\!\bigl(0,\frac{67}{972}\bigr)\)。
\end{enumerate}
上述常数由 \(\Pi(\lambda,y)\) 在点 \((\lambda,y)=(2,1)\) 的隐式微分唯一确定，并不依赖任何外部输入或额外归一化选择。
\end{conclusion}
\begin{proof}[证明要点（隐式微分）]
由定义
\(
\EE[y^{H_m}]
=\sum_{x\in X_m}y^{|x|_1}w_m(x)
=2^{-m}Z_m(y)
\)。
又由于 \(\Pi(\lambda,y)\) 在 \((2,1)\) 处满足 \(\partial_\lambda\Pi(2,1)\neq 0\)，
由隐函数定理得 \(\lambda_+(y)\) 在 \(y=1\) 的邻域内解析。
记 \(y=e^s\) 并令 \(\psi(s)=\log\lambda_+(e^s)-\log 2\)，则 \(\psi'(0)\) 与 \(\psi''(0)\) 分别给出线性漂移与方差率。
\par
对 \(\Pi(\lambda_+(y),y)=0\) 作一次与二次隐式微分，在点 \((\lambda,y)=(2,1)\) 代入可得
\[
\lambda_+'(1)=\frac{5}{9},
\qquad
\lambda_+''(1)=-\frac{64}{243}.
\]
于是
\(
\psi'(0)=\lambda_+'(1)/\lambda_+(1)=5/18
\)，并且
\[
\psi''(0)=\frac{\lambda_+'(1)}{\lambda_+(1)}+\frac{\lambda_+''(1)}{\lambda_+(1)}-\frac{\lambda_+'(1)^2}{\lambda_+(1)^2}
=\frac{67}{972}.
\]
大数律与中心极限定律由有限状态加权转移矩阵的解析谱扰动与标准准幂型定理推出。证毕。
\end{proof}

\begin{conclusion}[推前重量的一阶矩闭式：线性主项、常数漂移与交替相位]\label{con:xi-terminal-zm-m1-closed}
沿用结论 \ref{con:xi-terminal-zm-lln-clt} 的记号，并定义一阶矩分子
\[
M_1(m):=\sum_{x\in X_m}|x|_1\,d_m(x)=Z_m'(1).
\]
则对一切整数 \(m\ge 1\) 有严格恒等式
\[
\boxed{\;
M_1(m)=\Bigl(\frac{5}{18}m-\frac{1}{27}\Bigr)2^m+\frac14-\Bigl(\frac{m}{18}+\frac{23}{108}\Bigr)(-1)^m\;}
\]
从而均值具有精确展开
\[
\boxed{\;
\EE[H_m]=\frac{M_1(m)}{2^m}
=\frac{5}{18}m-\frac{1}{27}
+2^{-m}\Bigl(\frac14-\Bigl(\frac{m}{18}+\frac{23}{108}\Bigr)(-1)^m\Bigr)\;}
\]
其中有限尺寸偏离同时包含一个非零常数漂移与携带严格交替相位的 \(2^{-m}\) 层修正。
\end{conclusion}
\begin{proof}[证明要点（对 \(Z_m(y)\) 递推作一次微分）]
由结论 \ref{con:xi-terminal-zm-order4-rational}，\(Z_m(y)\) 满足四阶递推。
对该递推关于 \(y\) 微分并在 \(y=1\) 处取值，结合恒等式 \(Z_m(1)=2^m\)，得到 \(M_1(m)=Z_m'(1)\) 的常系数递推；
其特征多项式由 \(\Pi(\lambda,1)=(\lambda-2)(\lambda-1)(\lambda+1)^2\) 决定，从而解必可写为 \(2^m\) 通道与 \(\pm 1\) 通道的有限阶 Jordan 叠加。
用小 \(m\) 初值定出系数即可得到所示闭式；可由脚本 \texttt{scripts/exp\_fold\_zm\_bivariate\_partition\_audit.py} 复核。证毕。
\end{proof}

\begin{conclusion}[推前重量的二阶矩与方差闭式：线性涨落系数、常数项与双尺度交替修正]\label{con:xi-terminal-zm-m2-var-closed}
沿用结论 \ref{con:xi-terminal-zm-lln-clt} 的记号，并定义二阶矩分子
\[
M_2(m):=\sum_{x\in X_m}|x|_1^2\,d_m(x).
\]
则对一切整数 \(m\ge 1\) 有严格恒等式
\[
\boxed{\;
M_2(m)=\Bigl(\frac{25}{324}m^2+\frac{47}{972}m+\frac{113}{729}\Bigr)2^m+\frac{m}{8}
+\Bigl(\frac{m^3}{324}-\frac{m^2}{54}-\frac{31m}{216}-\frac{113}{729}\Bigr)(-1)^m\;}
\]
并因此得到方差的全量闭式分解
\[
\Var(H_m)=\frac{67}{972}m+\frac{112}{729}
+2^{-m}\Bigl(\frac{1}{54}-\frac{m}{72}+\Bigl(\frac{m^3}{324}+\frac{m^2}{81}-\frac{19m}{648}-\frac{83}{486}\Bigr)(-1)^m\Bigr)
\]
\[
\qquad\qquad
-2^{-2m}\Bigl(\frac{1}{16}-\Bigl(\frac{m}{36}+\frac{23}{216}\Bigr)(-1)^m+\Bigl(\frac{m}{18}+\frac{23}{108}\Bigr)^2\Bigr),
\]
其中主导涨落项为 \(\frac{67}{972}m+\frac{112}{729}\)，而全部有限尺寸偏离被压缩为两条严格指数尺度 \(2^{-m}\) 与 \(2^{-2m}\) 并携带交替相位。
\end{conclusion}
\begin{proof}[证明要点（对 \(Z_m(y)\) 递推作二次微分）]
注意 \(M_2(m)=Z_m''(1)+Z_m'(1)\)。
对结论 \ref{con:xi-terminal-zm-order4-rational} 的递推作二次微分并在 \(y=1\) 处取值，即得到 \(Z_m''(1)\) 的常系数递推；
再结合结论 \ref{con:xi-terminal-zm-m1-closed} 给出的 \(Z_m'(1)\) 闭式，可推出 \(M_2(m)\) 与 \(\Var(H_m)\) 的所示表达。
上述恒等式可由脚本 \texttt{scripts/exp\_fold\_zm\_bivariate\_partition\_audit.py} 复核。证毕。
\end{proof}

\begin{conclusion}[任意阶原始矩的湮灭律：\texorpdfstring{$(E-2)^{r+1}(E-1)^r(E+1)^{2r}$}{(E-2)^{r+1}(E-1)^r(E+1)^{2r}}]\label{con:xi-terminal-zm-mr-annihilator}
对整数 \(r\ge 0\)，定义原始矩分子
\[
M_r(m):=\sum_{x\in X_m}|x|_1^r\,d_m(x)\in\ZZ\qquad(m\ge 0).
\]
则序列 \(m\mapsto M_r(m)\) 满足常系数齐次递推，其最小湮灭多项式整除
\[
(E-2)^{r+1}(E-1)^r(E+1)^{2r},
\]
其中 \(E\) 为右移算子 \((Ea)(m)=a(m+1)\)。
等价地，存在多项式 \(P_r,Q_r,R_r\in\QQ[m]\) 使得对所有 \(m\ge 0\)
\[
M_r(m)=2^m P_r(m)+Q_r(m)+(-1)^m R_r(m),
\]
并且次数界为
\[
\deg P_r=r,\qquad \deg Q_r\le r-1,\qquad \deg R_r\le 2r-1.
\]
该结构把全部有限尺寸复杂性压缩为三条谱通道 \(2,1,-1\) 的有限阶 Jordan 叠加，从而任意阶矩在 \(m\) 上不仅存在极限律，而且具有全阶有限闭式展开。
\end{conclusion}
\begin{proof}[证明要点（谱通道与有限维扰动）]
由结论 \ref{con:xi-terminal-zm-order4-rational} 的“有限维带权转移矩阵”表示，
\(\{Z_m(y)\}\) 可写为固定维数矩阵 \(T(y)\) 的矩阵元；于是 \(M_r(m)\) 是该表示在 \(y=1\) 的有限阶导数所产生的线性组合。
在 \(y=1\) 时主多项式分解 \(\Pi(\lambda,1)=(\lambda-2)(\lambda-1)(\lambda+1)^2\) 给出三条谱通道；
有限阶导数仅提升相应 Jordan 阶数，从而得到所示湮灭多项式整除关系与三通道分解。
对 \(r\le 3\) 的直接代入核验与符号回归可由脚本 \texttt{scripts/exp\_fold\_zm\_bivariate\_partition\_audit.py} 完成。证毕。
\end{proof}

\begin{conclusion}[单位长度累积量的有理累积塔：所有极限累积量皆属 \texorpdfstring{$\QQ$}{Q}]\label{con:xi-terminal-zm-cumulant-rational}
令 \(\kappa_r(H_m)\) 为 \(H_m\) 的第 \(r\) 阶累积量（cumulant）。
则极限
\[
\kappa_r^\infty:=\lim_{m\to\infty}\frac{1}{m}\kappa_r(H_m)
\]
对每个整数 \(r\ge 1\) 存在，且满足 \(\kappa_r^\infty\in\QQ\)。
其生成机制可取为：令 \(\lambda_+(y)\) 表示特征方程 \(\Pi(\lambda,y)=0\) 在 \(y=1\) 的邻域内满足 \(\lambda_+(1)=2\) 的解析特征支，
并定义
\[
\psi(s):=\log\lambda_+(e^s)-\log 2.
\]
则 \(\psi\) 在 \(s=0\) 的泰勒系数全为有理数，且
\(
\kappa_r^\infty=\psi^{(r)}(0)
\)。
前四阶显式为
\[
\kappa_1^\infty=\frac{5}{18},\qquad
\kappa_2^\infty=\frac{67}{972},\qquad
\kappa_3^\infty=-\frac{22}{2187},\qquad
\kappa_4^\infty=-\frac{1763}{157464}.
\]
因此标准化极限偏度与超额峰度分别为
\[
\gamma_1=\frac{\kappa_3^\infty}{(\kappa_2^\infty)^{3/2}}\approx -0.5558544555,\qquad
\gamma_2=\frac{\kappa_4^\infty}{(\kappa_2^\infty)^2}\approx -2.3564268211.
\]
\end{conclusion}
\begin{proof}[证明要点（隐函数定理与系数域闭合）]
由结论 \ref{con:xi-terminal-zm-lln-clt}，极限对数矩母函数为 \(\psi(s)\)。
由于 \(\Pi(\lambda,y)\in\ZZ[\lambda,y]\) 且 \(\partial_\lambda\Pi(2,1)\neq 0\)，隐函数定理给出 \(\lambda_+(y)\) 在 \(y=1\) 邻域的解析展开；
逐阶比较系数可知其系数落在 \(\QQ\) 中，进而 \(\psi(s)=\log\lambda_+(e^s)-\log 2\) 的泰勒系数同属 \(\QQ\)。
累积量密度 \(\kappa_r^\infty\) 等于 \(\psi^{(r)}(0)\) 为标准结论；显式分数可由符号微分回收，
并可由脚本 \texttt{scripts/exp\_fold\_zm\_bivariate\_partition\_audit.py} 复核。证毕。
\end{proof}

\begin{conclusion}[三阶累积量的有限尺寸闭式：仿射主项与三层指数修正]\label{con:xi-terminal-zm-kappa3-finite-closed}
令 \(\kappa_3(H_m)\) 为 \(H_m\) 的第三阶累积量，则对一切整数 \(m\ge 1\) 有精确分解
\[
\kappa_3(H_m)= -\frac{22}{2187}m+\frac{52}{19683}
+2^{-m}\mathcal P_1\!\bigl(m,(-1)^m\bigr)
+2^{-2m}\mathcal P_2\!\bigl(m,(-1)^m\bigr)
+2^{-3m}\mathcal P_3\!\bigl(m,(-1)^m\bigr),
\]
其中每个 \(\mathcal P_j(\cdot,\cdot)\) 都可写成“偶/奇”二扇区可分离的形式
\(
\mathcal P_j(m,\varepsilon)=A_j(m)+\varepsilon B_j(m)
\)
（\(A_j,B_j\) 为有理系数多项式）。
三层修正可取为
\[
\mathcal P_1(m,\varepsilon)=\Bigl(\frac{m^2}{1728}+\frac{151m}{1728}-\frac{293}{10368}\Bigr)
+\varepsilon\Bigl(-\frac{m^5}{12960}+\frac{7m^4}{15552}+\frac{41m^3}{11664}-\frac{487m^2}{46656}-\frac{2699m}{77760}+\frac{2407}{279936}\Bigr),
\]
\[
\mathcal P_2(m,\varepsilon)=\Bigl(\frac{m^4}{1944}+\frac{47m^3}{11664}+\frac{35m^2}{11664}-\frac{143m}{3888}-\frac{269}{2187}\Bigr)
+\varepsilon\Bigl(-\frac{m^3}{432}-\frac{5m^2}{432}+\frac{7m}{432}+\frac{34}{243}\Bigr),
\]
\[
\mathcal P_3(m,\varepsilon)=\Bigl(\frac{m^2}{216}+\frac{23m}{648}+\frac{193}{1944}\Bigr)
+\varepsilon\Bigl(-\frac{m^3}{2916}-\frac{23m^2}{5832}-\frac{629m}{17496}-\frac{15617}{157464}\Bigr).
\]
因此三阶非高斯性在单位长度尺度上呈严格线性斜率 \(-22/2187\)，且其全部有限尺寸偏离被压缩为三层指数尺度 \(2^{-m},2^{-2m},2^{-3m}\) 的代数可审计修正。
\end{conclusion}
\begin{proof}[证明要点（由闭式矩回收累积量并分层收集）]
由结论 \ref{con:xi-terminal-zm-m1-closed} 与结论 \ref{con:xi-terminal-zm-m2-var-closed} 得到 \(\EE[H_m]\) 与 \(\EE[H_m^2]\) 的闭式；
并由结论 \ref{con:xi-terminal-zm-mr-annihilator} 在 \(r=3\) 的特例得到 \(\EE[H_m^3]\) 的三通道闭式。
将三者代入
\(\kappa_3(H_m)=\EE[H_m^3]-3\EE[H_m^2]\EE[H_m]+2\EE[H_m]^3\)，
并把 \(2^{-m}\) 的幂次与 \((-1)^m\) 的相位逐层收集，即得到所示分解与多项式系数。
恒等式可由脚本 \texttt{scripts/exp\_fold\_zm\_bivariate\_partition\_audit.py} 复核。证毕。
\end{proof}

\begin{conclusion}[重量分布的全局大偏差：速率函数为四次代数谱的 Legendre 变换]\label{con:xi-terminal-zm-ldp}
定义压力函数
\[
\psi(s):=\log\lambda_+(e^s)-\log 2.
\]
则 \(H_m/m\) 满足大偏差原理，其速率函数为
\[
I(\alpha)=\sup_{s\in\RR}\bigl(s\alpha-\psi(s)\bigr).
\]
其中 \(\psi\) 在实轴上严格凸且实解析；
其复解析延拓的最近实轴障碍由结论 \ref{con:xi-terminal-zm-discriminant-leyang} 的 \(y_{\mathrm{LY}}\) 所决定（经 \(y=e^s\) 的复对数支提升）。
因此重量统计的多分形谱可在不引入任何额外热力学形式体系的前提下，直接由 \(\Pi(\lambda,y)\) 的主特征支给出。
\end{conclusion}
\begin{proof}[证明要点]
由结论 \ref{con:xi-terminal-zm-lln-clt}，\(\psi\) 给出极限对数矩母函数。
在 \(\psi\) 的可微与严格凸条件下，G\"artner--Ellis 定理推出大偏差原理及其 Legendre 变换形式。
最近解析障碍由 \(\lambda_+(y)\) 的分歧点决定，而分歧点由判别式为零刻画，见结论 \ref{con:xi-terminal-zm-discriminant-leyang}。证毕。
\end{proof}

\begin{remark}[主特征支的椭圆一致化：\texorpdfstring{$\lambda_+(y)$}{lambda+(y)} 的 Weierstrass \texorpdfstring{$\wp$}{wp} 表示]\label{rem:xi-terminal-zm-lambda-plus-wp-uniformization}
在结论 \ref{con:xi-terminal-zm-elliptic-normalization} 的椭圆模型
\[
E:\ Y^2=X^3-X+\frac14
\]
上，复解析一致化给出存在周期格 \(\Lambda\subset\CC\) 与 Weierstrass 函数 \(\wp_\Lambda\) 使得
\[
E(\CC)\cong \CC/\Lambda,\qquad
X=\wp_\Lambda(z),\qquad
Y=\frac{\wp_\Lambda'(z)}{2},
\]
并满足微分方程
\[
\bigl(\wp_\Lambda'(z)\bigr)^2=4\wp_\Lambda(z)^3-4\wp_\Lambda(z)+1
\qquad(g_2=4,\ g_3=-1).
\]
在该一致化下，物理主支重量函数（结论 \ref{con:xi-terminal-zm-y-quadratic-extension-minpoly}）
\[
y=X^2+Y-\frac12
\]
提升为椭圆函数
\[
y(z)=\wp_\Lambda(z)^2+\frac{\wp_\Lambda'(z)-1}{2}.
\]
令 \(\lambda(z):=X(z)=\wp_\Lambda(z)\)。
由结论 \ref{con:xi-terminal-zm-elliptic-normalization} 的双有理互逆 \(\nu:E\to\mathcal C\)，
对任意 \(z\in\CC/\Lambda\) 有 \((\lambda(z),y(z))\in\mathcal C(\CC)\)，即
\[
\Pi\bigl(\lambda(z),y(z)\bigr)=0.
\]
反之，\(\Pi(\lambda,y)=0\) 的正规化在复解析意义下同构于 \(E(\CC)\cong\CC/\Lambda\)，
因此其任一解析支均可由上述椭圆函数参数化统一给出。
因此 \(\pi_y\) 的有限分歧点在一致化参数 \(z\) 下等价于临界点条件 \(\frac{\dd}{\dd z}y(z)=0\)；
而 \(y=\infty\) 对应于 \(z\to 0\) 的四阶极点结构，与 \((y)_\infty=4O\) 的全分歧极点描述一致。
令 \(P_1=(2,-\tfrac52)\in E(\QQ)\)（满足 \(y(P_1)=1\)，见结论 \ref{con:xi-terminal-zm-elliptic-rational-points-denominator-law}），并取 \(z_1\in\CC/\Lambda\) 使
\(
P_1=(\wp_\Lambda(z_1),\wp_\Lambda'(z_1)/2)
\)。
由于 \(P_1\) 不在覆盖 \(\pi_y:E\to\mathbb P^1_y\) 的分歧点上（结论 \ref{con:xi-terminal-zm-elliptic-y-hurwitz-passport}），
\(y(z)\) 在 \(z=z_1\) 的邻域内局部可逆；令 \(z=z(y)\) 表示其局部反函数，则主特征支可写为
\[
\lambda_+(y)=X\bigl(z(y)\bigr)=\wp_\Lambda\bigl(z(y)\bigr),
\qquad \lambda_+(1)=2.
\]
因此 \(\lambda_+(y)\) 在避开分歧值集合的任一路径上均可解析延拓，并可通过在 \(\CC/\Lambda\) 上延拓 \(z(y)\) 后回代 \(\wp_\Lambda\) 进行全局表示。
从而 \(\lambda_+(y)\)（以及与之伴随的 \(\log\lambda_+(y)\) 的多值解析结构）可被等价视为固定椭圆格上的反演与周期问题，而非一般代数函数分支的偶然叠加。
进一步，由结论 \ref{con:xi-terminal-zm-leyang-ramification-critical-values} 的微分恒等式
\(
\dd y=(4XY+3X^2-1)\,\omega
\)
（其中 \(\omega=\dd X/(2Y)\)）得：在任意去分歧区域内取满足 \(\pi_y(P(y))=y\) 的局部解析截面 \(P(y)\in E(\CC)\)，令
\(
X(y)=X(P(y))
\)、\(
Y(y)=Y(P(y))
\)，则
\[
\boxed{\;
\frac{\dd}{\dd y}\log \lambda_+(y)
=\frac{1}{X(y)}\frac{\dd X(y)}{\dd y}
=\frac{2Y(y)}{X(y)\,\bigl(4X(y)Y(y)+3X(y)^2-1\bigr)}\; }.
\]
等价地，在最小整模型 \(E_0:\eta^2+\eta=x^3-x\) 上令 \(x=\lambda\)、\(\eta=y-\lambda^2\)，则同一恒等式可写为
\[
\frac{\dd}{\dd y}\log \lambda(y)
=\frac{2\eta+1}{x\bigl(2x(2\eta+1)+3x^2-1\bigr)}.
\]
因此 \(\bigl(\dd(\log\lambda_+)/\dd y\bigr)\dd y=\dd X/X\) 为椭圆曲线上的第三类阿贝尔微分，其周期决定 \(\log\lambda_+\) 的多值解析延拓。
结合注记 \ref{rem:xi-terminal-elliptic-37a1-modular-anchor} 中 \(E\) 作为 \(X_0^+(37)\) 模型的外部锚定，上述周期亦可在 level \(37\) 的模形式周期口径下理解。
其分歧集合与局部型由 \(\pi_y\) 的分歧结构完全确定：有限分歧值为
\(
y\in\{0,1\}\cup\{P_{\mathrm{LY}}(y)=0\}
\)，
且在每个有限分歧值处呈平方根型 Puiseux 分歧（结论 \ref{con:xi-terminal-zm-branch-puiseux-uniformization}）；
而在 \(y=\infty\) 处由 \((y)_\infty=4O\) 给出四重全分歧（结论 \ref{con:xi-terminal-zm-elliptic-y-hurwitz-passport}）。
特别地，实轴上的最近分歧值 \(y_{\mathrm{LY}}\) 及其伴随退化特征值 \(\lambda_{\mathrm{LY}}\) 分别由
\(
P_{\mathrm{LY}}(y_{\mathrm{LY}})=0
\)
与
\(
16\lambda_{\mathrm{LY}}^{3}-9\lambda_{\mathrm{LY}}^{2}+1=0
\)
代数刻画（结论 \ref{con:xi-terminal-zm-leyang-lambda-cubic}）。
\end{remark}

\begin{conclusion}[大偏差速率函数的代数参数化与端点对数奇性]\label{con:xi-terminal-zm-ldp-rate-algebraic-param-endpoint-log}
沿用结论 \ref{con:xi-terminal-zm-ldp} 的压力函数
\(
\psi(s)=\log\lambda_{+}(e^{s})-\log2
\)
及其 Legendre--Fenchel 对偶 \(I(\alpha)=\sup_{s\in\RR}(s\alpha-\psi(s))\)。
对任意 \(y>0\)，令 \(\lambda=\lambda_{+}(y)\)。
定义
\[
\alpha(y):=\frac{y\,\lambda_{+}'(y)}{\lambda_{+}(y)}\in(0,\tfrac12).
\]
则由谱方程 \(\Pi(\lambda,y)=0\) 的隐式微分得到显式有理表达式
\[
\alpha(y)
=-\frac{y\,\partial_y\Pi(\lambda,y)}{\lambda\,\partial_{\lambda}\Pi(\lambda,y)}
=\frac{y\bigl(2\lambda^{2}-(2y+1)\bigr)}{\lambda\bigl(4\lambda^{3}-3\lambda^{2}-2(2y+1)\lambda+1\bigr)}
\Big|_{\lambda=\lambda_{+}(y)}.
\]
并且沿该参数支有
\[
I\bigl(\alpha(y)\bigr)=\alpha(y)\log y-\log\lambda_{+}(y)+\log2.
\]
在本体系中，\(\alpha:(0,\infty)\to(0,\tfrac12)\) 为严格递增的实解析双射，因此 \(I\) 在 \((0,\tfrac12)\) 上实解析且严格凸。
其端点行为由 Perron 分支的分数幂尺度锁定（结论 \ref{con:xi-terminal-zm-perron-puiseux-two-scale}）：
\begin{itemize}
  \item 当 \(\alpha\downarrow0\) 时，
  \[
  I(\alpha)
  =\log2+2\alpha\log\alpha+(\log8-2)\alpha+O\!\bigl(\alpha^{2}\log(1/\alpha)\bigr).
  \]
  \item 当 \(\alpha\uparrow\tfrac12\) 时，令 \(\delta:=\tfrac12-\alpha\downarrow0\)，则
  \[
  I\!\left(\tfrac12-\delta\right)
  =\log2+4\delta\log\delta+(12\log2-4)\delta+O\!\bigl(\delta^{2}\log(1/\delta)\bigr).
  \]
\end{itemize}
其中 \(\alpha\log\alpha\) 与 \(\delta\log\delta\) 两个非解析项分别由 \(y\to0^{+}\) 的 \(y^{1/2}\) Puiseux 与 \(y\to+\infty\) 的 \(y^{-1/4}\) Puiseux 诱导，因而构成该 Lee--Yang 权重体系在两端点的可检验奇性指纹。
\end{conclusion}

\begin{conclusion}[复 Lee--Yang 主导的解析半径与高阶累积量振荡律]\label{con:xi-terminal-zm-psi-analytic-radius-cumulant-oscillation}
考虑 \(\psi(s)=\log\lambda_{+}(e^{s})-\log2\) 在 \(s=0\) 的解析芽（结论 \ref{con:xi-terminal-zm-lln-clt}）。
由结论 \ref{con:xi-terminal-zm-discriminant-leyang} 与结论 \ref{con:xi-terminal-zm-leyang-cubic-branch-image}，
\(y\)-平面的有限分歧值除 \(\{0,1\}\) 外还包含 Lee--Yang 三次
\(
P_{\mathrm{LY}}(y)=256y^{3}+411y^{2}+165y+32
\)
的共轭复根对
\[
y_{\mathrm c}^{\pm}\approx-0.2355083735\pm0.2339255521\,i.
\]
令
\[
s_{\mathrm c}^{\pm}:=\log y_{\mathrm c}^{\pm}\quad(\text{取主值}),
\qquad
\rho:=|s_{\mathrm c}^{+}|\approx2.6045558274,
\qquad
\theta:=\arg(s_{\mathrm c}^{+})\approx2.0080026785,
\]
则 \(\psi\) 在圆盘 \(|s|<\rho\) 内单值全纯，并且该圆盘为以 \(0\) 为中心的最大解析圆盘；其边界上首先遇到的奇点即 \(s_{\mathrm c}^{\pm}\) 所对应的代数分歧。
\par
记 \(\kappa_n^\infty:=\psi^{(n)}(0)\in\QQ\) 为单位长度的第 \(n\) 阶累积量密度（结论 \ref{con:xi-terminal-zm-cumulant-rational}）。
则当 \(n\to\infty\) 时存在常数 \(B\neq 0\) 与相位 \(\varphi\in\RR\)，使得
\[
\kappa_n^\infty
=\frac{n!}{\sqrt\pi}\,|B|\,\rho^{-n+\frac12}\,n^{-\frac32}
\cos\!\Bigl((n-\tfrac12)\theta+\varphi\Bigr)
+o\!\bigl(n!\rho^{-n}n^{-\frac32}\bigr).
\]
更精确地，取临界对 \((\lambda_{\mathrm c},y_{\mathrm c})\) 满足
\[
\Pi(\lambda_{\mathrm c},y_{\mathrm c})=0,\qquad
\partial_{\lambda}\Pi(\lambda_{\mathrm c},y_{\mathrm c})=0,
\qquad (y_{\mathrm c}=y_{\mathrm c}^{+}),
\]
并令
\[
c:=\sqrt{-\,\frac{2\,\partial_{y}\Pi(\lambda_{\mathrm c},y_{\mathrm c})}{\partial_{\lambda\lambda}\Pi(\lambda_{\mathrm c},y_{\mathrm c})}},
\qquad
B:=\frac{c\sqrt{y_{\mathrm c}}}{\lambda_{\mathrm c}},
\]
则上述振荡律成立；数值上可取
\[
\lambda_{\mathrm c}\approx0.4179674053-0.2321140339\,i,
\qquad
|B|\approx0.9634999692.
\]
因此高阶累积量密度除呈现 \(\rho^{-n}\) 的指数控制与 \(n^{-3/2}\) 的普适幂律修正外，还携带由 \(\theta=\arg(s_{\mathrm c}^{+})\) 决定的确定性振荡频率，给出一条可直接数值核验的“复 Lee--Yang 振荡律”。
可复现核验脚本：\texttt{scripts/exp\_fold\_zm\_bivariate\_partition\_audit.py}。
\end{conclusion}

\begin{theorem}[Lee--Yang 主振荡角的 Gelfond--Schneider 超越性]\label{thm:xi-terminal-zm-leyang-oscillation-angle-transcendence}
沿用结论 \ref{con:xi-terminal-zm-psi-analytic-radius-cumulant-oscillation} 的记号，令
\(
s_{\mathrm c}:=s_{\mathrm c}^{+}=\log y_{\mathrm c}^{+}
\)
为主值对数，并写 \(\rho:=|s_{\mathrm c}|\)、\(\theta:=\arg(s_{\mathrm c})\in(0,\pi)\)。
则：
\begin{enumerate}
  \item \(\tan\theta\) 为超越数；
  \item 特别地，\(\theta/\pi\notin\QQ\)。
\end{enumerate}
\end{theorem}
\begin{proof}
设
\[
N:=y_{\mathrm c}^{+}\overline{y_{\mathrm c}^{+}}=|y_{\mathrm c}^{+}|^{2}\in\overline{\QQ}\cap\RR_{>0},
\qquad
r:=\frac{y_{\mathrm c}^{+}}{\overline{y_{\mathrm c}^{+}}}\in\overline{\QQ},
\qquad |r|=1.
\]
由于 \(y_{\mathrm c}^{+}\notin\RR_{\le 0}\)，主值对数与共轭相容，故
\[
\log N=s_{\mathrm c}+\overline{s_{\mathrm c}}=2\Re(s_{\mathrm c})\neq 0,
\qquad
s_{\mathrm c}-\overline{s_{\mathrm c}}=2i\,\Im(s_{\mathrm c})
\]
并且 \(\exp(s_{\mathrm c}-\overline{s_{\mathrm c}})=r\neq 1\)。
令
\[
\beta:=\frac{s_{\mathrm c}-\overline{s_{\mathrm c}}}{s_{\mathrm c}+\overline{s_{\mathrm c}}}
=i\,\frac{\Im(s_{\mathrm c})}{\Re(s_{\mathrm c})}=i\tan\theta.
\]
若 \(\tan\theta\) 为代数数，则 \(\beta\in\overline{\QQ}\setminus\QQ\)（因为 \(\beta\notin\RR\)）。
取 \(\alpha:=N\in\overline{\QQ}\setminus\{0,1\}\)，由 Gelfond--Schneider 定理可知任一值 \(\alpha^{\beta}\) 皆为超越数。\cite{Baker1975TranscendentalNumberTheory}
另一方面，由 \(\beta\log N=s_{\mathrm c}-\overline{s_{\mathrm c}}\) 得
\[
N^{\beta}=\exp(\beta\log N)=\exp(s_{\mathrm c}-\overline{s_{\mathrm c}})=r\in\overline{\QQ},
\]
矛盾。
故 \(\tan\theta\) 不可能为代数数，即 \(\tan\theta\) 为超越数。
\par
若 \(\theta/\pi\in\QQ\)，则 \(\zeta:=e^{2i\theta}\) 为单位根，从而
\(
\tan\theta=\frac{\zeta-1}{i(\zeta+1)}
\in\overline{\QQ}
\)，
与已证超越性矛盾。
故 \(\theta/\pi\notin\QQ\)。证毕。
\end{proof}

\begin{theorem}[归一化高阶累积量密度的反正弦极限律]\label{thm:xi-terminal-zm-cumulant-arcsine-limit-law}
\leanverified{paper\_xi\_terminal\_zm\_cumulant\_arcsine\_limit\_law}
沿用结论 \ref{con:xi-terminal-zm-psi-analytic-radius-cumulant-oscillation} 的振荡渐近：
存在 \(B\neq 0\)、\(\varphi\in\RR\) 使
\[
\kappa_n^\infty
=\frac{n!}{\sqrt\pi}\,|B|\,\rho^{-n+\frac12}\,n^{-\frac32}
\cos\!\Bigl((n-\tfrac12)\theta+\varphi\Bigr)
+o\!\bigl(n!\rho^{-n}n^{-\frac32}\bigr).
\]
定义归一化序列
\[
A_n:=\frac{\sqrt\pi}{|B|}\,\rho^{\,n-\frac12}\,n^{\frac32}\,\frac{\kappa_n^\infty}{n!}.
\]
则 \(A_n=\cos\!\bigl((n-\tfrac12)\theta+\varphi\bigr)+o(1)\)，并且由定理 \ref{thm:xi-terminal-zm-leyang-oscillation-angle-transcendence} 的 \(\theta/\pi\notin\QQ\) 得：\(A_n\) 的经验分布弱收敛到 \([-1,1]\) 上的反正弦分布。
具体地，对任意连续函数 \(f\in C([-1,1])\)，有
\[
\lim_{N\to\infty}\frac1N\sum_{n=1}^{N} f(A_n)
=\frac{1}{\pi}\int_{-1}^{1}\frac{f(x)}{\sqrt{1-x^2}}\,\dd x.
\]
\end{theorem}
\begin{proof}
由渐近式直接得到 \(A_n-\cos(t_n)\to 0\)，其中 \(t_n:=(n-\tfrac12)\theta+\varphi\)。
由 \(\theta/\pi\notin\QQ\) 可知旋转序列 \(t_n\ (\mathrm{mod}\ 2\pi)\) 在 \(\RR/2\pi\ZZ\) 上等分布（Kronecker--Weyl），故对任意连续 \(g\in C(\RR/2\pi\ZZ)\) 有
\[
\frac1N\sum_{n=1}^{N} g(t_n)\to \frac{1}{2\pi}\int_{0}^{2\pi} g(t)\,\dd t.
\]
取 \(g(t)=f(\cos t)\) 即得 \(\cos(t_n)\) 的经验分布为反正弦律的推前测度。\cite{KuipersNiederreiter1974}
再由 \(f\) 的一致连续性与 \(A_n-\cos(t_n)\to 0\) 控制 Ces\`aro 平均差，得到所述极限式。证毕。
\end{proof}

\begin{corollary}[符号密度、极值充满与无限次符号翻转]\label{cor:xi-terminal-zm-cumulant-sign-density-extremes}
\leanverified{paper\_xi\_terminal\_zm\_cumulant\_sign\_density\_extremes}
沿用定理 \ref{thm:xi-terminal-zm-cumulant-arcsine-limit-law} 的记号。
则集合 \(\{n:\kappa_n^\infty>0\}\) 与 \(\{n:\kappa_n^\infty<0\}\) 的自然密度均为 \(1/2\)，并且
\[
\limsup_{n\to\infty}A_n=1,\qquad \liminf_{n\to\infty}A_n=-1.
\]
从而序列 \((\kappa_n^\infty)\) 发生无限次符号翻转，且不可能最终周期。
\end{corollary}
\begin{proof}
由反正弦分布关于 \(0\) 对称且 \(\cos t\) 在一半圆周取正、一半取负，
结合 \(A_n-\cos(t_n)\to 0\) 得符号密度结论。
反正弦分布的支撑为 \([-1,1]\) 且任意开邻域均有正测度，故 \(\cos(t_n)\) 在 \([-1,1]\) 稠密；
再由 \(A_n-\cos(t_n)\to 0\) 得 \(\limsup A_n=1\)、\(\liminf A_n=-1\)。
最后，无限次正负出现推出无限次符号翻转；若 \((\kappa_n^\infty)\) 最终周期，则其主振荡相位 \((n-\tfrac12)\theta+\varphi\) 只能对应有理旋转，从而 \(\theta/\pi\in\QQ\)，与定理 \ref{thm:xi-terminal-zm-leyang-oscillation-angle-transcendence} 矛盾。证毕。
\end{proof}

\begin{proposition}[相邻阶局部增长率在扩展实轴上的稠密聚点]\label{prop:xi-terminal-zm-cumulant-local-growth-dense}
\leanverified{paper\_xi\_terminal\_zm\_cumulant\_local\_growth\_dense}
沿用结论 \ref{con:xi-terminal-zm-psi-analytic-radius-cumulant-oscillation} 的振荡渐近，并在 \(\kappa_n^\infty\neq 0\) 时定义
\[
R_n:=\rho\left(\frac{n}{n+1}\right)^{\frac32}\cdot\frac{\kappa_{n+1}^\infty}{(n+1)\kappa_n^\infty}.
\]
则 \(R_n\) 在扩展实轴 \(\widehat{\RR}:=\RR\cup\{\infty\}\) 上具有稠密的聚点集合：对任意 \(L\in\RR\)，存在无穷子列 \(n_k\to\infty\) 使 \(R_{n_k}\to L\)，并且亦存在无穷子列使 \(|R_{n_k}|\to\infty\)。
\end{proposition}
\begin{proof}
由结论 \ref{con:xi-terminal-zm-psi-analytic-radius-cumulant-oscillation}，
\[
\frac{\kappa_{n+1}^\infty}{(n+1)\kappa_n^\infty}
=\frac{1}{\rho}\left(\frac{n}{n+1}\right)^{\frac32}
\frac{\cos(t_n+\theta)}{\cos(t_n)}+o(1),
\qquad
t_n:=(n-\tfrac12)\theta+\varphi.
\]
故
\[
R_n=\frac{\cos(t_n+\theta)}{\cos(t_n)}+o(1)
=\cos\theta-\sin\theta\,\tan(t_n)+o(1),
\]
其中 \(\sin\theta\neq 0\)（因 \(\theta\in(0,\pi)\)）。
由 \(\theta/\pi\notin\QQ\) 得 \(t_n\ (\mathrm{mod}\ \pi)\) 等分布，从而 \(\tan(t_n)\) 在 \(\widehat{\RR}\) 上稠密；
再由仿射同胚 \(x\mapsto \cos\theta-\sin\theta\,x\) 与 \(o(1)\) 微扰的稳定性得到 \(R_n\) 的聚点稠密性，并可取 \(|\tan(t_n)|\to\infty\) 的子列推出 \(|R_n|\to\infty\)。证毕。
\end{proof}

\begin{theorem}[Stokes 幅值平方 \texorpdfstring{$B^2$}{B^2} 的三次最小多项式、判别式与 \texorpdfstring{$S_3$}{S3} 刚性]\label{thm:xi-terminal-zm-stokes-amplitude-square-cubic-rigidity}
沿用结论 \ref{con:xi-terminal-zm-psi-analytic-radius-cumulant-oscillation} 的临界对 \((\lambda_{\mathrm c},y_{\mathrm c})\) 与幅值常数 \(B\)，并令
\[
b:=B^{2}=\frac{c^{2}y_{\mathrm c}}{\lambda_{\mathrm c}^{2}}\in\CC,
\qquad t:=b.
\]
设 \(K:=\QQ(y_{\mathrm c})\) 为 Lee--Yang 三次域（\(P_{\mathrm{LY}}(y_{\mathrm c})=0\)）。
则在 \(K\) 内成立显式有理恒等式
\begin{equation}\label{eq:xi-terminal-zm-stokes-t-rational-in-yc}
\boxed{\;
t
=-\frac{1}{6}\cdot
\frac{440y_{\mathrm c}^{2}+47y_{\mathrm c}-64}{46y_{\mathrm c}^{2}+40y_{\mathrm c}-5}\;}
\end{equation}
并且 \(t\) 在 \(\QQ\) 上满足三次整式
\begin{equation}\label{eq:xi-terminal-zm-stokes-t-minpoly}
\boxed{\;
162t^{3}+1593t^{2}+1998t+1184=0\;}
\end{equation}
等价地，\(b\) 满足
\[
162b^{3}+1593b^{2}+1998b+1184=0.
\]
该三次在 \(\QQ\) 上不可约，且判别式为
\[
\mathrm{Disc}(162b^{3}+1593b^{2}+1998b+1184)
=-2^{2}\cdot 3^{9}\cdot 11^{2}\cdot 37\cdot 109^{2}.
\]
因判别式非平方，其分裂域伽罗瓦群同构于 \(S_3\)。
从而 \(\QQ(b)=\QQ(t)=\QQ(y_{\mathrm c})\)，亦即 \(B^2\) 本身为 \(K\) 的一个原始生成元。
\end{theorem}
\begin{proof}[证明要点]
由结论 \ref{con:xi-terminal-zm-branch-double-root-linear-closure}，在约束 \(P_{\mathrm{LY}}(y_{\mathrm c})=0\) 下二重根谱支满足线性刚性
\(
279\lambda_{\mathrm c}+512y_{\mathrm c}^{2}+518y_{\mathrm c}+5=0
\)，
故 \(\lambda_{\mathrm c}\in \QQ(y_{\mathrm c})\)。
另一方面，
\(
\partial_y\Pi(\lambda,y)=-2\lambda^2+2y+1
\)、
\(
\partial_{\lambda\lambda}\Pi(\lambda,y)=12\lambda^2-6\lambda-4y-2
\)，
从而
\(
t=-2\cdot\frac{y_{\mathrm c}\,\partial_y\Pi(\lambda_{\mathrm c},y_{\mathrm c})}{\lambda_{\mathrm c}^{2}\,\partial_{\lambda\lambda}\Pi(\lambda_{\mathrm c},y_{\mathrm c})}
\in \QQ(\lambda_{\mathrm c},y_{\mathrm c})
=\QQ(y_{\mathrm c})
\)。
将 \(\lambda_{\mathrm c}\) 以 \(y_{\mathrm c}\) 的有理式表示并在三次商环 \(\QQ[y]/(P_{\mathrm{LY}})\) 中约化，即得 \eqref{eq:xi-terminal-zm-stokes-t-rational-in-yc}。
再将 \eqref{eq:xi-terminal-zm-stokes-t-rational-in-yc} 改写为关于 \(y_{\mathrm c}\) 的二次方程并与 \(P_{\mathrm{LY}}(y_{\mathrm c})=0\) 联立作消元，即得到 \eqref{eq:xi-terminal-zm-stokes-t-minpoly}。
有理根检验排除一次因子，故该三次在 \(\QQ\) 上不可约。
其判别式按三次判别式公式计算即为
\(
-2^{2}\cdot 3^{9}\cdot 11^{2}\cdot 37\cdot 109^{2}
\)，非平方，故分裂域伽罗瓦群为 \(S_3\)。
由不可约性可知 \([\QQ(t):\QQ]=3\)，且 \(t\in\QQ(y_{\mathrm c})\) 已证，遂 \(\QQ(t)=\QQ(y_{\mathrm c})\)；
由 \(b=t\) 得 \(\QQ(b)=\QQ(t)\)。证毕。
\end{proof}

\begin{corollary}[幅值三次的二次分辨子域：\texorpdfstring{$\QQ(\sqrt{-111})$}{Q(sqrt(-111))}]\label{cor:xi-terminal-zm-stokes-b2-quadratic-resolvent}
沿用定理 \ref{thm:xi-terminal-zm-stokes-amplitude-square-cubic-rigidity} 的记号，令 \(K=\QQ(b)\) 并令 \(L\) 为 \(K/\QQ\) 的伽罗瓦闭包。
则 \(L/\QQ\) 的唯一二次子域为
\[
\boxed{\ \QQ(\sqrt{-111})\ }.
\]
特别地，该二次子域与 Lee--Yang 三次的二次分辨域一致（见结论 \ref{con:xi-terminal-zm-leyang-cubic-arithmetic}），从而 \(B^2\) 与 \(y_{\mathrm c}\) 的任一原始生成元给出同一条二次判别脊柱。
\end{corollary}
\begin{proof}
由定理 \ref{thm:xi-terminal-zm-stokes-amplitude-square-cubic-rigidity}，\(b\) 的极小多项式为不可约三次 \(f(t)=162t^{3}+1593t^{2}+1998t+1184\)，且
\[
\mathrm{Disc}(f)=-2^{2}\cdot 3^{9}\cdot 11^{2}\cdot 37\cdot 109^{2}.
\]
判别式的平方自由部分为 \(-3\cdot 37=-111\)。
对不可约三次且 \(\mathrm{Disc}(f)\) 非平方的情形，其分裂域的伽罗瓦群为 \(S_3\)，并且其唯一二次子域由 \(\QQ(\sqrt{\mathrm{Disc}(f)})\) 给出；
因此 \(L\) 的唯一二次子域为 \(\QQ(\sqrt{-111})\)。证毕。
\end{proof}

\begin{conjecture}[Stokes 幅值平方的 CM--模参数化]\label{conj:xi-terminal-zm-stokes-b2-cm-modular-parameterization}
沿用定理 \ref{thm:xi-terminal-zm-stokes-amplitude-square-cubic-rigidity} 的记号，令 \(b:=B^2\) 并记三次域 \(K:=\QQ(b)\)。
猜想存在某一层级 \(N\) 的模曲线上的模函数 \(f\)（例如 \(X_0(N)\) 的 Hauptmodul 或 \(\lambda\)-函数的有理变换）与某一 CM 点 \(\tau\in\mathbb H\)，使得
\[
b=f(\tau),
\qquad
\QQ(b)\ \text{为相应环类域的某一三次子域}.
\]
并且 \(N\) 的素因子集合可由 \(\mathrm{Disc}(162b^{3}+1593b^{2}+1998b+1184)\) 的素因子集合控制。
该猜想可通过对候选 \(N\) 与 CM 判别式的有限枚举检验：比较 \(f(\tau)\) 的极小多项式与 \eqref{eq:xi-terminal-zm-stokes-t-minpoly}。
\end{conjecture}

% Exponential period model for the Stokes amplitude via third-kind differentials.

\begin{conjecture}[Stokes 幅值平方的第三类微分指数化周期模型]\label{conj:xi-terminal-zm-stokes-b2-thirdkind-exponential-period}
沿用定理 \ref{thm:xi-terminal-zm-leyang-elliptic-structure} 的椭圆曲线 \(E_{\min}/\QQ\) 及其有理函数
\(
y= X^{2}+Y-\tfrac12
\)
（在短 Weierstrass 口径 \(Y^{2}=X^{3}-X+\tfrac14\) 下），并沿用定理 \ref{thm:xi-terminal-zm-stokes-amplitude-square-cubic-rigidity} 的幅值常数 \(B\)。
则存在一个零度因子 \(D\in\mathrm{Div}^{0}(E_{\min})\)，其支撑包含于
\[
\mathrm{supp}(D)\subseteq y^{-1}(0)\cup y^{-1}(1),
\]
以及一条由 Stokes 归一化选择所确定的基本闭路 \(\gamma\in H_{1}(E_{\min}(\CC),\ZZ)\)，使得对满足
\(
\mathrm{Res}(\eta_D)=D
\)
且 \(a\)-周期为零的第三类微分 \(\eta_D\)，成立指数化周期恒等式
\[
\boxed{\;
B^{2}=\exp\!\left(\int_{\gamma}\eta_D\right).\;}
\]
进一步地，若 \(\varphi:X_0(37)\to E_{\min}\) 为模参数化（同构到 \(37\mathrm a1\) 的模性锚点口径），并令 \(f\in S_2(\Gamma_0(37))\) 为与 \(E_{\min}\) 对应的权二新形式，则预期 \(\varphi^{*}\eta_D\) 属于由 \(f(q)\,\dd q/q\) 及其 \(q\)-导数所张成的微分域（等价地，可由该模形式及其微分闭包中的元素表示）。
\end{conjecture}

\endinput



\begin{theorem}[振幅平方 \texorpdfstring{$B^2$}{B^2} 与 \texorpdfstring{$\kappa_*^2$}{kappa*^2} 的显式有理互构]\label{thm:xi-terminal-zm-b2-kappa2-rational-interconstruction}
\leanverified{paper\_xi\_terminal\_zm\_b2\_kappa2\_rational\_interconstruction}
令
\[
u:=\kappa_*^2,
\qquad
b:=B^2.
\]
其中 \(u\) 取自定理 \ref{thm:xi-terminal-zm-kappa-square-cubic-field-s3}，\(b\) 取自定理 \ref{thm:xi-terminal-zm-stokes-amplitude-square-cubic-rigidity}。
则存在显式有理关系
\[
\boxed{\;
b
=-2\cdot
\frac{24708u^2-28628u+7733}{51084u^2-44412u+8735}\; }.
\]
因而
\[
\QQ(b)=\QQ(u)=\QQ(\lambda_*),
\]
即 \(B^2\) 与 \(\kappa_*^2\) 生成同一 Lee--Yang 三次数域。
\end{theorem}
\begin{proof}
由定理 \ref{thm:xi-terminal-zm-kappa-square-cubic-field-s3}，
\[
\lambda_*=\frac{3(2u-1)}{4(3u-2)}.
\]
由定理 \ref{thm:xi-terminal-zm-stokes-t-elliptic-branch-coordinate} 的表达式 \eqref{eq:xi-terminal-zm-stokes-t-rational-in-alpha}（取 \(\alpha=\lambda_*\) 且 \(t=b\)）：
\[
b=\frac{2(1113\lambda_*^2-2192\lambda_*+303)}{3(1831\lambda_*^2+144\lambda_*-175)}.
\]
把 \(\lambda_*=\frac{3(2u-1)}{4(3u-2)}\) 代入并在 \(\QQ(u)\) 中约化，即得所示有理式，故 \(b\in\QQ(u)\)。
另一方面，定理 \ref{thm:xi-terminal-zm-kappa-square-cubic-field-s3} 与定理 \ref{thm:xi-terminal-zm-stokes-amplitude-square-cubic-rigidity} 分别给出
\(
[\QQ(u):\QQ]=[\QQ(b):\QQ]=3
\)。
故包含关系 \(\QQ(b)\subseteq\QQ(u)\) 只能是等号；再结合
\(
\QQ(u)=\QQ(\lambda_*)
\)
即得结论。
\end{proof}

\begin{theorem}[幅值平方的椭圆分歧坐标表达]\label{thm:xi-terminal-zm-stokes-t-elliptic-branch-coordinate}
沿用定理 \ref{thm:xi-terminal-zm-stokes-amplitude-square-cubic-rigidity} 的记号，并令 \(\alpha:=\lambda_{\mathrm c}\)。
则 \(\alpha\) 满足 Lee--Yang 三次谱分歧方程 \(16\alpha^{3}-9\alpha^{2}+1=0\)（结论 \ref{con:xi-terminal-zm-leyang-ramification-critical-values}）。
并且在分歧约束下有恒等式
\[
4\alpha\,y_{\mathrm c}=4\alpha^{3}-3\alpha^{2}-2\alpha+1
\qquad(\text{结论 \ref{con:xi-terminal-zm-leyang-ramification-critical-values}}),
\]
从而 \(t=B^{2}\) 可写为 \(\QQ(\alpha)\) 中的有理函数：
\begin{equation}\label{eq:xi-terminal-zm-stokes-t-rational-in-alpha}
\boxed{\;
t
=\frac{2\,(1113\alpha^{2}-2192\alpha+303)}{3\,(1831\alpha^{2}+144\alpha-175)}\; }.
\end{equation}
\end{theorem}
\begin{proof}[证明要点]
由结论 \ref{con:xi-terminal-zm-leyang-ramification-critical-values} 的分歧恒等式可将 \(y_{\mathrm c}\) 表为 \(\alpha\) 的有理式。
将其代入定理 \ref{thm:xi-terminal-zm-stokes-amplitude-square-cubic-rigidity} 的表达式 \eqref{eq:xi-terminal-zm-stokes-t-rational-in-yc}，
并在商环 \(\QQ[\alpha]/(16\alpha^{3}-9\alpha^{2}+1)\) 中约化，即得 \eqref{eq:xi-terminal-zm-stokes-t-rational-in-alpha}。证毕。
\end{proof}

\begin{theorem}[隐藏平方与纯有理反演：由 \texorpdfstring{$t=B^{2}$}{t=B^2} 重构 \texorpdfstring{$y_{\mathrm c}$}{yc}]\label{thm:xi-terminal-zm-stokes-hidden-square-rational-inversion}
令 \(t\) 满足 \eqref{eq:xi-terminal-zm-stokes-t-minpoly}，并在三次域 \(K:=\QQ(t)\) 中定义
\begin{equation}\label{eq:xi-terminal-zm-stokes-s-def}
s:=\frac{-71045-13947t+12258t^{2}}{1199}.
\end{equation}
则在 \(K\) 内成立隐藏平方恒等式
\begin{equation}\label{eq:xi-terminal-zm-stokes-hidden-square}
\boxed{\;
s^{2}=10080t^{2}+16224t+12761\; }.
\end{equation}
并且由 \eqref{eq:xi-terminal-zm-stokes-t-rational-in-yc} 导出的二次方程在 \(K\) 内分解，其两根可写为纯有理表达
\begin{equation}\label{eq:xi-terminal-zm-stokes-yc-rational-in-t}
\boxed{\;
y^{(1)}(t)
=\frac{36774t^{2}-329601t-269488}{9592(69t+110)},
\qquad
y^{(2)}(t)
=\frac{-36774t^{2}-245919t+156782}{9592(69t+110)}\; }.
\end{equation}
其中与 \(P_{\mathrm{LY}}(y)=0\) 相容者即为 Lee--Yang 临界值 \(y_{\mathrm c}\)；在结论 \ref{con:xi-terminal-zm-psi-analytic-radius-cumulant-oscillation} 取 \(y_{\mathrm c}=y_{\mathrm c}^{+}\) 的嵌入下，对应 \(y_{\mathrm c}=y^{(1)}(t)\)。
\end{theorem}
\begin{proof}[证明要点]
由 \eqref{eq:xi-terminal-zm-stokes-t-rational-in-yc} 可得 \(y_{\mathrm c}\) 满足关于 \(y\) 的二次方程
\[
(276t+440)y^{2}+(240t+47)y-(30t+64)=0.
\]
其判别式为
\(
\Delta=9(10080t^{2}+16224t+12761)
\)。
式 \eqref{eq:xi-terminal-zm-stokes-hidden-square} 表明该判别式平方根已属于 \(K\)，故该二次方程在 \(K\) 内分解并给出两条根的纯有理表达；
化简即得 \eqref{eq:xi-terminal-zm-stokes-yc-rational-in-t}。证毕。
\end{proof}

\begin{proposition}[临界谱值的幅值重构]\label{prop:xi-terminal-zm-stokes-lambda-critical-rational-in-t}
在定理 \ref{thm:xi-terminal-zm-stokes-hidden-square-rational-inversion} 的设定下，
\(\lambda_{\mathrm c}\) 可由 \(t\) 显式重构为
\begin{equation}\label{eq:xi-terminal-zm-stokes-lambda-critical-rational-in-t}
\boxed{\;
\lambda_{\mathrm c}
=-\frac{525411t^{2}-4582128t-2625484}{1199(69t+110)^{2}}\; }.
\end{equation}
\end{proposition}
\begin{proof}
由结论 \ref{con:xi-terminal-zm-branch-double-root-linear-closure} 的线性刚性关系，
\(
\lambda_{\mathrm c}=-(512y_{\mathrm c}^{2}+518y_{\mathrm c}+5)/279
\)。
将 \eqref{eq:xi-terminal-zm-stokes-yc-rational-in-t} 中与 \(P_{\mathrm{LY}}(y)=0\) 相容的根 \(y_{\mathrm c}\) 代入并在 \(\QQ(t)\) 中约化即可得到 \eqref{eq:xi-terminal-zm-stokes-lambda-critical-rational-in-t}。证毕。
\end{proof}

\begin{theorem}[Stokes 幅值的六次代数度与极小多项式]\label{thm:xi-terminal-zm-stokes-amplitude-degree6-minpoly}
设 \(B\) 为结论 \ref{con:xi-terminal-zm-psi-analytic-radius-cumulant-oscillation} 中由临界对 \((\lambda_{\mathrm c},y_{\mathrm c})\) 所定义的 Stokes 幅值常数，并令 \(t=B^{2}\)。
则 \([\QQ(B):\QQ]=6\)，且 \(B\) 的极小多项式为
\begin{equation}\label{eq:xi-terminal-zm-stokes-B-minpoly}
\boxed{\;
162B^{6}+1593B^{4}+1998B^{2}+1184=0\; }.
\end{equation}
\end{theorem}
\begin{proof}
由 \eqref{eq:xi-terminal-zm-stokes-t-minpoly}，\(t\) 在 \(\QQ\) 上次数为 \(3\)，故 \(K:=\QQ(t)\) 为三次域且 \(K\subset \QQ(B)\)。
若 \(B\in K\)，则 \(t=B^{2}\) 在 \(K\) 中为平方，从而其范数 \(\mathrm{N}_{K/\QQ}(t)\) 必为有理数平方，因而非负。
另一方面，把 \eqref{eq:xi-terminal-zm-stokes-t-minpoly} 除以 \(162\) 化为首一多项式，可得
\[
\mathrm{N}_{K/\QQ}(t)=(-1)^{3}\cdot\frac{1184}{162}=-\frac{592}{81}<0,
\]
矛盾。
故 \(B\notin K\)，从而 \([\QQ(B):K]=2\) 并推出 \([\QQ(B):\QQ]=6\)。
式 \eqref{eq:xi-terminal-zm-stokes-B-minpoly} 由 \eqref{eq:xi-terminal-zm-stokes-t-minpoly} 代入 \(t=B^{2}\) 直接得到；
其次数为 \(6\)，与 \([\QQ(B):\QQ]=6\) 一致，故为极小多项式。证毕。
\end{proof}

\begin{conclusion}[全阶累积量密度的有理性与可审计生成方程]\label{con:xi-terminal-zm-cumulant-dalgebraic-equation}
结论 \ref{con:xi-terminal-zm-cumulant-rational} 已给出 \(\psi(s)\) 在 \(s=0\) 的 Taylor 展开属于 \(\QQ[[s]]\)，因而 \(\kappa_n^\infty=\psi^{(n)}(0)\in\QQ\) 对一切 \(n\ge 1\) 成立。
进一步地，由 \(\Pi(\lambda,y)=0\) 在代换
\(
\lambda=2e^{\psi(s)},\ y=e^{s}
\)
下得到 \(\psi\) 的显式代数生成方程
\[
16e^{4\psi(s)}-8e^{3\psi(s)}-4(2e^{s}+1)e^{2\psi(s)}+2e^{\psi(s)}+e^{2s}+e^{s}=0,
\]
从而 \(\psi\) 为 \(D\)-代数函数；这为任意阶 \(\kappa_n^\infty\) 的符号递推与误差审计提供封闭的代数约束。
除结论 \ref{con:xi-terminal-zm-cumulant-rational} 已列出的低阶项外，高阶例如
\[
\kappa_{5}^\infty=\frac{15650}{1594323},\quad
\kappa_{6}^\infty=\frac{307723}{57395628},\quad
\kappa_{7}^\infty=-\frac{2332358}{129140163},\quad
\kappa_{8}^\infty=\frac{166894943}{55788550416}.
\]
可复现核验脚本：\texttt{scripts/exp\_fold\_zm\_bivariate\_partition\_audit.py}。
\end{conclusion}

\begin{conclusion}[Abel--Jacobi 坐标下的 Lee--Yang 微分闭式：\texorpdfstring{$y(u)$}{y(u)} 的一阶代数方程]\label{con:xi-terminal-zm-leyang-abel-jacobi-ode}
沿用结论 \ref{con:xi-terminal-zm-elliptic-normalization} 的椭圆模型
\[
E:\quad Y^2=X^3-X+\frac14,
\qquad
y:=X^2+Y-\frac12\in\QQ(E),
\qquad
\omega:=\frac{\dd X}{2Y}.
\]
取 Abel--Jacobi 坐标 \(u\) 使 \(\dd u=\omega\)，并令
\[
\delta(u):=\frac{\dd y}{\dd u}\ \in\ \QQ(E).
\]
则在去分歧区域内，\((y(u),\delta(u))\) 满足完全显式的一阶代数微分闭式
\[
\boxed{\;
\delta^{4}+(8y-3)\delta^{3}+(2y^{2}-34y-4)\delta^{2}
-y(y-1)\,P_{\mathrm{LY}}(y)=0\;},
\qquad
P_{\mathrm{LY}}(y):=256y^{3}+411y^{2}+165y+32.
\]
该四次关系在 \(\QQ(y)[T]\) 中不可约，从而给出函数域扩张的首一四次本原表示
\[
\QQ(E)=\QQ(y,\delta),
\qquad
[\QQ(E):\QQ(y)]=4,
\]
并且常数项精确等于 Lee--Yang 判别核（带符号）\(-y(y-1)P_{\mathrm{LY}}(y)\)。
因此常数项（即 \(\delta=0\) 的驻值谱）严格给出
\[
\delta(u)=0\quad\Longleftrightarrow\quad
y(u)\in \{0,1\}\ \cup\ \{P_{\mathrm{LY}}(y)=0\},
\]
从而 Lee--Yang 分歧像在 Abel--Jacobi 时间下作为驻值集合内禀出现，并与结论 \ref{con:xi-terminal-zm-discriminant-leyang} 的分歧判别核一致。
\end{conclusion}
\begin{proof}[证明要点（消元）]
由结论 \ref{con:xi-terminal-zm-leyang-ramification-critical-values} 的微分恒等式
\(
\dd y=(4XY+3X^2-1)\,\omega
\)
得
\[
\delta=4XY+3X^2-1.
\]
又由 \(y=X^2+Y-\tfrac12\) 得 \(Y=y-X^2+\tfrac12\)，代入上式得到关于 \(X\) 的三次约束
\[
4X^{3}-3X^{2}-(4y+2)X+(1+\delta)=0.
\]
另一方面，消去 \(Y\) 的平面四次模型由
\[
F(X,y):=X^{4}-X^{3}-(2y+1)X^{2}+X+y(y+1)=0
\]
给出（结论 \ref{con:xi-terminal-zm-leyang-discriminant-norm-identity}）。
对变量 \(X\) 作结果式消元 \(\mathrm{Res}_{X}(F,\ 4X^{3}-3X^{2}-(4y+2)X+(1+\delta))\)，并在 \(\ZZ[y,\delta]\) 中因式整理，即得到所示四次方程；
其常数项分解为 \(-y(y-1)P_{\mathrm{LY}}(y)\) 与分歧像刻画（结论 \ref{con:xi-terminal-zm-discriminant-leyang}）一致。
不可约性可由一次良性专化读出：例如取 \(y=2\)，所得四次多项式在 \(\QQ[T]\) 上不可约，从而原四次在 \(\QQ(y)[T]\) 中不可约。
可复现核验脚本：\texttt{scripts/exp\_fold\_zm\_elliptic\_leyang\_cover\_geometry\_audit.py}。证毕。
\end{proof}

\begin{conclusion}[判别式自反与二级退化五次：\texorpdfstring{$\mathrm{Disc}_T(F_\delta)$}{Disc\_T(F\_delta)} 的平方自由部仍为 Lee--Yang 核]\label{con:xi-terminal-zm-leyang-delta-quartic-discriminant-reflexive}
令
\[
F_\delta(T;y):=T^{4}+(8y-3)T^{3}+(2y^{2}-34y-4)T^{2}-y(y-1)\,P_{\mathrm{LY}}(y)\in\ZZ[y][T],
\qquad
P_{\mathrm{LY}}(y)=256y^{3}+411y^{2}+165y+32.
\]
则其关于 \(T\) 的判别式满足严格因式分解
\[
\boxed{\;
\mathrm{Disc}_{T}\bigl(F_\delta(T;y)\bigr)
=-y(y-1)\,P_{\mathrm{LY}}(y)\cdot Q_{5}(y)^{2}\;},
\]
其中
\[
Q_{5}(y)=4096y^{5}+5376y^{4}-464y^{3}-2749y^{2}-723y+80\in\ZZ[y].
\]
因此 \(\delta\)-生成表示的全部不可约平方自由障碍仍由同一 Lee--Yang 判别核 \(-y(y-1)P_{\mathrm{LY}}(y)\) 控制；
而 \(Q_{5}(y)=0\) 给出一组与 \(y(y-1)P_{\mathrm{LY}}(y)=0\) 不同的二级退化参数集合，其上四叶结构发生额外并合。
\end{conclusion}
\begin{proof}[证明要点]
直接计算 \(\mathrm{Disc}_{T}(F_\delta)\) 并在 \(\ZZ[y]\) 中因式分解得到所示恒等式。
可复现核验脚本：\texttt{scripts/exp\_fold\_zm\_elliptic\_leyang\_cover\_geometry\_audit.py}。证毕。
\end{proof}

\paragraph{\texorpdfstring{$(y,\delta)$}{(y,delta)} 平面闭包模型：五节点、伴随二次与谱坐标反演}
沿用结论 \ref{con:xi-terminal-zm-leyang-abel-jacobi-ode} 的椭圆曲线
\[
E:\quad Y^{2}=X^{3}-X+\frac14,
\qquad
y:=X^{2}+Y-\frac12\in\QQ(E),
\qquad
\omega:=\frac{\dd X}{2Y}.
\]
令 Abel--Jacobi 坐标 \(u\) 满足 \(\dd u=\omega\)，并令
\[
\delta:=\frac{\dd y}{\dd u}\in\QQ(E).
\]
由 \(\dd y=(4XY+3X^{2}-1)\,\omega\)（结论 \ref{con:xi-terminal-zm-leyang-abel-jacobi-ode} 的证明中已给出）可得
\[
\delta=4XY+3X^{2}-1.
\]

\begin{theorem}[二阶射流的三次代数闭式]\label{thm:xi-terminal-zm-leyang-second-jet-cubic-closure}
\leanverified{paper\_xi\_terminal\_zm\_leyang\_second\_jet\_cubic\_closure}
令 \(u\) 为满足 \(\dd u=\omega\) 的 Abel--Jacobi 坐标，并令
\[
\delta=\frac{\dd y}{\dd u},
\qquad
\delta_{2}:=\frac{\dd^{2}y}{\dd u^{2}}=\frac{\dd\delta}{\dd u}.
\]
则在 \(E\) 上有显式表达式
\[
\delta_{2}=20X^{3}-12X+2+12XY.
\]
并且三元组 \((y,\delta,\delta_2)\) 满足一条三次代数关系：
\[
\begin{aligned}
0={}&64\delta^{3}+96\delta^{2}\delta_{2}+576\delta^{2}y-288\delta^{2}+48\delta\delta_{2}^{2}-144\delta\delta_{2}y-216\delta\delta_{2}\\
&+9600\delta y^{3}-5424\delta y^{2}+1200\delta y+132\delta+8\delta_{2}^{3}-216\delta_{2}^{2}y-63\delta_{2}^{2}\\
&-3200\delta_{2}y^{3}-1392\delta_{2}y^{2}+1008\delta_{2}y+128\delta_{2}+16000y^{3}-2640y^{2}-816y-16.
\end{aligned}
\]
\end{theorem}
\begin{proof}
由 \(E:Y^2=X^3-X+\tfrac14\) 得 \(\frac{\dd X}{\dd u}=2Y\) 且 \(\frac{\dd Y}{\dd u}=3X^2-1\)。
对 \(\delta=4XY+3X^2-1\) 求导得
\[
\delta_2
=4\left(\frac{\dd X}{\dd u}Y+X\frac{\dd Y}{\dd u}\right)+6X\frac{\dd X}{\dd u}
=8Y^2+12X^3-4X+12XY.
\]
再用 \(8Y^2=8(X^3-X+\tfrac14)=8X^3-8X+2\) 即得 \(\delta_2=20X^3-12X+2+12XY\)。
\par
另一方面，\((y,\delta)\) 与 \(X\) 满足三次约束
\[
4X^{3}-3X^{2}-(4y+2)X+(1+\delta)=0,
\]
其由 \(y=X^2+Y-\tfrac12\) 与 \(\delta=4XY+3X^2-1\) 消去 \(Y\) 得到（见定理 \ref{thm:xi-terminal-zm-delta-spectral-x-inversion} 的证明中的同型消元式）。
将 \(\delta_2\) 视作未知量，与上式及 \(\delta_2=20X^{3}-12X+2+12XY\) 联立，并对 \(X\) 作消元（例如取关于 \(X\) 的结果式）得到所示三次关系；该关系在 \(E\) 上恒成立，故为二阶射流的代数闭式。
\end{proof}

\begin{theorem}[\texorpdfstring{$\delta$}{delta}-闭包的平面五次与五节点归一化]\label{thm:xi-terminal-zm-delta-closure-plane-quintic}
令
\[
F_\delta(T;y):=T^{4}+(8y-3)T^{3}+(2y^{2}-34y-4)T^{2}-y(y-1)\,P_{\mathrm{LY}}(y)\in\ZZ[y][T],
\]
\[
P_{\mathrm{LY}}(y):=256y^{3}+411y^{2}+165y+32.
\]
记仿射平面曲线
\[
\mathcal C_\delta:=\{(y,T)\in\mathbb A^{2}_{(y,T)}:\ F_\delta(T;y)=0\},
\]
并记 \(\overline{\mathcal C}_\delta\subset\mathbb P^{2}\) 为将 \(F_\delta\) 以总次数 \(5\) 齐次化所得的射影闭包。
则成立：
\begin{enumerate}
  \item \(\overline{\mathcal C}_\delta\) 在仿射部分 \(\mathbb A^{2}\) 内恰有五个奇点，并且全为普通双点（节点）；
  \item \(\overline{\mathcal C}_\delta\) 在无穷远处仅有一个点且为光滑点；
  \item \(\overline{\mathcal C}_\delta\) 的归一化为亏格 \(1\) 的光滑射影曲线，并与 \(E\) 在 \(\QQ\) 上同构；
  等价地，
  \[
  \QQ(\overline{\mathcal C}_\delta)\cong \QQ(y,T)/(F_\delta)\cong \QQ(E)=\QQ(X,Y).
  \]
\end{enumerate}
\end{theorem}

\begin{proof}[证明要点]
首先，由结论 \ref{con:xi-terminal-zm-leyang-abel-jacobi-ode} 得 \(\QQ(E)=\QQ(y,\delta)\)，并且 \((y,\delta)\) 满足 \(F_\delta(\delta;y)=0\)；
因此 \(\QQ(y,T)/(F_\delta)\) 的光滑模型为亏格 \(1\) 的曲线，并与 \(E\) 同构（函数域同构确定归一化）。

其次，齐次化取总次数 \(5\)，可写为
\[
F_\delta^{\mathrm{h}}(T,y,Z)=
Z\Bigl(T^{4}+(8y-3Z)T^{3}+(2y^{2}-34yZ-4Z^{2})T^{2}\Bigr)-y(y-Z)P_{\mathrm{LY}}^{\mathrm{h}}(y,Z),
\]
其中
\(
P_{\mathrm{LY}}^{\mathrm{h}}(y,Z)=256y^{3}+411y^{2}Z+165yZ^{2}+32Z^{3}
\)。
令 \(Z=0\) 得 \(F_\delta^{\mathrm{h}}(T,y,0)=-256y^{5}\)，故无穷远处唯一点为 \([T:y:Z]=[1:0:0]\)。
并且
\(
\partial_ZF_\delta^{\mathrm{h}}(1,0,0)=1\neq 0
\)，因此该无穷远点光滑，得到第二条。

对仿射奇点，设
\[
I_{\mathrm{sing}}:=(F_\delta,\partial_TF_\delta,\partial_yF_\delta)\subset\QQ[T,y].
\]
由于奇点必满足 \(F_\delta(T;y)\) 在 \(T\) 上有重根，故其 \(y\)-坐标只能落在判别式零点上。
由结论 \ref{con:xi-terminal-zm-leyang-abel-jacobi-ode} 的判别式分解
\[
\mathrm{Disc}_{T}(F_\delta)=-y(y-1)P_{\mathrm{LY}}(y)\,Q_{5}(y)^{2},
\quad
Q_{5}(y)=4096y^{5}+5376y^{4}-464y^{3}-2749y^{2}-723y+80,
\]
奇点候选仅可能出现在 \(y\in\{0,1\}\cup\{P_{\mathrm{LY}}(y)=0\}\cup\{Q_5(y)=0\}\)。

当 \(T=0\) 时，\(F_\delta(0;y)=-y(y-1)P_{\mathrm{LY}}(y)\)。
直接计算
\[
\partial_yF_\delta(0;y)=-(2y-1)P_{\mathrm{LY}}(y)-y(y-1)P_{\mathrm{LY}}'(y),
\]
并注意 \(P_{\mathrm{LY}}\) 无重根且其根不等于 \(0,1\)，可知在 \(y=0,1\) 与 \(P_{\mathrm{LY}}(y)=0\) 的各点处 \(\partial_yF_\delta(0;y)\neq 0\)；
因此 \((y,T)=(0,0)\)、\((1,0)\) 以及 \(\{P_{\mathrm{LY}}=0\}\times\{0\}\) 仅对应分歧而非平面奇点。
从而一切平面奇点均满足 \(T\neq 0\)，并且其 \(y\)-坐标必满足 \(Q_5(y)=0\)。

在 \(I_{\mathrm{sing}}\) 上可作消元得到
\(
I_{\mathrm{sing}}\cap\QQ[y]=(Q_5(y))
\)，并且在 \(\QQ[T,y]/I_{\mathrm{sing}}\) 中有线性关系
\[
93\,T=P_4(y),
\qquad
P_4(y)=8192y^{4}+5888y^{3}-3680y^{2}-2848y+152.
\]
由于 \(\gcd(Q_5,Q_5')=1\)，故 \(Q_5\) 的五个根互异，从而得到五个互异奇点
\[
\left(y,\ T=\frac{P_4(y)}{93}\right),\qquad Q_5(y)=0.
\]

最后，为判定奇点类型，取二次切锥判别式
\[
\Delta:=(\partial_{Ty}F_\delta)^{2}-(\partial_{TT}F_\delta)(\partial_{yy}F_\delta).
\]
将 \(\Delta\) 在 \(\QQ[T,y]/I_{\mathrm{sing}}\) 中约化得到
\[
\Delta\equiv 2\cdot\bigl(19072y^{4}+51120y^{3}+45129y^{2}+14759y+1032\bigr).
\]
该四次多项式与 \(Q_5\) 互素，故在 \(Q_5(y)=0\) 的各点处 \(\Delta\neq 0\)，从而切锥分解为两条不同直线，五个奇点均为节点。
由于平面五次曲线的算术亏格为 \(\frac{(5-1)(5-2)}{2}=6\)，节点的 \(\delta\)-不变量均为 \(1\)，总和为 \(5\)，故几何亏格为 \(1\)。
与第一段的函数域同构合并，即得第三条。证毕。
\end{proof}

\begin{proposition}[一阶射流的纤维碰撞与 \texorpdfstring{$Q_5$}{Q5} 判据]\label{prop:xi-terminal-zm-leyang-first-jet-fiber-collision-q5}
令 \(\pi_y:E\to\PP^1_y\) 由函数 \(y=X^2+Y-\tfrac12\) 诱导（次数 \(4\)），并在仿射图上定义一阶射流映射
\[
j^{1}y:\ E\longrightarrow \mathbb A^{2}_{(y,T)},
\qquad
P\longmapsto \bigl(y(P),\delta(P)\bigr).
\]
对任意 \(y_0\in\overline{\QQ}\) 满足 \(y_0(y_0-1)P_{\mathrm{LY}}(y_0)\neq 0\)，下列二择一成立：
\begin{enumerate}
  \item 若 \(Q_5(y_0)\neq 0\)，则纤维 \(\pi_y^{-1}(y_0)\) 含四个不同点，且 \(\delta\) 在该纤维上取四个互异值；等价地，\(j^1y\) 在该纤维上为单射。
  \item 若 \(Q_5(y_0)=0\)，则纤维 \(\pi_y^{-1}(y_0)\) 仍含四个不同点，但其中恰有两点 \(P\neq P'\) 满足
  \[
  y(P)=y(P')=y_0,
  \qquad
  \delta(P)=\delta(P')=\frac{8192y_0^{4}+5888y_0^{3}-3680y_0^{2}-2848y_0+152}{93},
  \]
  其余两点在 \(\delta\) 下取值互异。
\end{enumerate}
\end{proposition}
\begin{proof}
固定 \(y=y_0\) 时，\(\delta\) 的可能取值为四次多项式 \(F_\delta(T;y_0)=0\) 的根。
由于 \(y_0(y_0-1)P_{\mathrm{LY}}(y_0)\neq 0\)，由定理 \ref{thm:xi-terminal-zm-leyang-elliptic-structure} 中对 \(\pi_y\) 的分歧刻画可知：\(y_0\) 不是分歧值，故纤维 \(\pi_y^{-1}(y_0)\) 由四个不同点组成。
\par
若 \(Q_5(y_0)\neq 0\)，由定理 \ref{thm:xi-terminal-zm-delta-closure-plane-quintic} 的判别式分解
\(
\Disc_T(F_\delta)=-y(y-1)P_{\mathrm{LY}}(y)\,Q_5(y)^2
\)
知 \(\Disc_T(F_\delta(\cdot;y_0))\neq 0\)，从而 \(F_\delta(\cdot;y_0)\) 有四个互异根；它们分别对应四个片层上的 \(\delta\)-取值，故 \(j^1y\) 在该纤维上单射。
\par
若 \(Q_5(y_0)=0\)，在同一非分歧假设下纤维仍由四个不同点组成；而 \(\Disc_T(F_\delta(\cdot;y_0))=0\) 表明 \(F_\delta(\cdot;y_0)\) 在 \(T\)-轴上出现重根。
由定理 \ref{thm:xi-terminal-zm-delta-closure-plane-quintic} 的消元刻画，该重根满足线性闭式
\(
93\,T=P_4(y)
\)，因此两点在同一 \((y_0,\delta_0)\) 处发生射流识别，且 \(\delta_0=P_4(y_0)/93\)。
其余两点对应的根为单根，故 \(\delta\) 取值互异。
\end{proof}

\begin{proposition}[由 \texorpdfstring{$\delta$}{delta} 乘法诱导的秩 \(4\) 谱结构]\label{prop:xi-terminal-zm-leyang-delta-multiplication-spectral-structure}
令 \(\pi_y:E\to\PP^1_y\) 如上，并令
\[
V:=\pi_{y*}\mathcal O_E.
\]
则 \(V\) 为 \(\PP^1_y\) 上秩 \(4\) 向量丛。
令 \(\Phi:V\to V\) 为“乘以 \(\delta\)”诱导的 \(\mathcal O_{\PP^1_y}\)-线性内自同态。
则在仿射图 \(\mathbb A^1_y\) 上有恒等式
\[
\det\bigl(\eta\cdot\mathrm{Id}_{V}-\Phi\bigr)=F_\delta(\eta;y),
\]
从而 \(F_\delta(\eta;y)=0\) 给出 \((V,\Phi)\) 的谱曲线；其归一化同构于 \(E\)，而平方因子 \(Q_5(y)^2\) 精确记录平面模型中由节点导致的非正则识别。
\end{proposition}
\begin{proof}
在 \(\mathbb A^1_y\) 上，\(V\) 的局部截面可识别为有限平坦代数
\[
\mathcal O_{\mathbb A^1_y}[T]\big/(F_\delta(T;y)),
\]
其中 \(T\) 对应 \(\delta\) 的代数生成元（定理 \ref{thm:xi-terminal-zm-delta-closure-plane-quintic} 的函数域同构）。
在该表示下，\(\Phi\) 恰为“乘以 \(T\)”的乘法算子。
其特征多项式等于 \(T\) 在 \(\QQ(y)\) 上的极小多项式；由于 \(F_\delta\) 为首一四次且不可约，二者一致，从而得到所示行列式恒等式。
谱曲线的归一化与平方因子的几何含义由定理 \ref{thm:xi-terminal-zm-delta-closure-plane-quintic} 与推论 \ref{cor:xi-terminal-zm-delta-discriminant-square-geometry} 给出。
\end{proof}

\begin{theorem}[五节点平面模型的广义雅可比与半阿贝尔分解]\label{thm:xi-terminal-zm-delta-generalized-jacobian-semiabelian}
令 \(\overline{\mathcal C}_\delta\subset\PP^2\) 为定理 \ref{thm:xi-terminal-zm-delta-closure-plane-quintic} 的射影闭包，并令 \(\nu:E\to\overline{\mathcal C}_\delta\) 为归一化映射。
则 \(\overline{\mathcal C}_\delta\) 为 Deligne--Mumford 稳定曲线，算术亏格为 \(6\)，并且其广义雅可比（\(\mathrm{Pic}^0(\overline{\mathcal C}_\delta)\)）为维数 \(6\) 的半阿贝尔簇，满足正合列
\[
1\longrightarrow (\mathbb G_m)^5\longrightarrow \mathrm{Pic}^0(\overline{\mathcal C}_\delta)\longrightarrow E\longrightarrow 0.
\]
更具体地，设五个节点在 \(E\) 上的两两预像为 \(\{P_i,Q_i\}\subset E(\overline{\QQ})\)（\(i=1,\dots,5\)），则该扩张的扩张类由差分除子类 \([P_i-Q_i]\in \mathrm{Pic}^0(E)\cong E\) 决定。
\end{theorem}
\begin{proof}[证明要点]
由定理 \ref{thm:xi-terminal-zm-delta-closure-plane-quintic}，\(\overline{\mathcal C}_\delta\) 仅含五个节点且归一化为椭圆曲线 \(E\)，故
\(
p_a(\overline{\mathcal C}_\delta)=g(E)+5=6
\)，并且稳定性成立。
\par
对节点曲线有标准的单位层短正合列
\[
1\to \mathcal O_{\overline{\mathcal C}_\delta}^{\times}\to \nu_{*}\mathcal O_{E}^{\times}\to \bigoplus_{i=1}^{5}\mathbb G_m\to 1,
\]
其中末项逐节点记录两支粘合的标量比。
取 \(\mathrm H^{1}\) 并取零度部分得到
\[
1\to(\mathbb G_m)^5\to \mathrm{Pic}^0(\overline{\mathcal C}_\delta)\to \mathrm{Pic}^0(E)\to 0,
\]
而 \(\mathrm{Pic}^0(E)\cong E\)，从而得到所示半阿贝尔分解。
扩张类由每个节点的两预像点对 \(\{P_i,Q_i\}\) 的差分除子类决定，这是广义雅可比的标准刻画。
\end{proof}

\begin{proposition}[伴随二次曲线的唯一性与 \texorpdfstring{$Q_5$}{Q5} 的内禀消元刻画]\label{prop:xi-terminal-zm-delta-adjoint-conic-and-elimination}
令
\[
H(T,y):=\frac{1}{T}\,\partial_TF_\delta(T;y)=4T^{2}+(24y-9)T+(4y^{2}-68y-8).
\]
则：
\begin{enumerate}
  \item \(H(T,y)=0\) 在射影意义下为 \(\overline{\mathcal C}_\delta\) 的唯一伴随二次曲线（即通过全部五个节点的二次曲线，按比例唯一）；
  \item 结果式消元满足恒等式
  \[
  \mathrm{Res}_{T}\bigl(F_\delta(T;y),\,H(T,y)\bigr)=Q_{5}(y)^{2}.
  \]
  因而 \(Q_5\) 可被内禀刻画为理想 \((F_\delta,\frac{1}{T}\partial_TF_\delta)\) 的消元因子。
\end{enumerate}
\end{proposition}

\begin{proof}[证明要点]
由于 \(F_\delta\) 不含 \(T\) 的一次项，\(\partial_TF_\delta\) 含因子 \(T\)，故 \(H\) 为二次多项式。
任一节点满足 \(F_\delta=\partial_TF_\delta=\partial_yF_\delta=0\)，并且由定理 \ref{thm:xi-terminal-zm-delta-closure-plane-quintic} 的证明知节点处 \(T\neq 0\)，故节点亦满足 \(H=0\)；
于是 \(H=0\) 通过全部五个节点。

另一方面，次数为 \(5\) 且仅含节点的平面曲线，其正则微分由次数 \(d-3=2\) 的伴随曲线给出；
在亏格 \(1\) 情形，伴随二次曲线空间维数为 \(1\)，故按比例唯一。
因此 \(H=0\) 为唯一伴随二次曲线。

对结果式，使用判别式--结果式恒等式
\[
\mathrm{Disc}_{T}(F_\delta)=\pm\,\mathrm{Res}_{T}(F_\delta,\partial_TF_\delta),
\]
并用 \(\partial_TF_\delta=T\cdot H\) 与结果式乘法性得到
\[
\mathrm{Res}_{T}(F_\delta,\partial_TF_\delta)
=\mathrm{Res}_{T}(F_\delta,T)\,\mathrm{Res}_{T}(F_\delta,H).
\]
注意 \(\mathrm{Res}_{T}(F_\delta,T)=F_\delta(0;y)=-y(y-1)P_{\mathrm{LY}}(y)\)。
与结论 \ref{con:xi-terminal-zm-leyang-abel-jacobi-ode} 的判别式分解合并，即得
\(
\mathrm{Res}_{T}(F_\delta,H)=Q_5(y)^2
\)。
证毕。
\end{proof}

\begin{theorem}[谱坐标 \texorpdfstring{$\lambda=X$}{lambda=X} 的有理反演与极点除子]\label{thm:xi-terminal-zm-delta-spectral-x-inversion}
继续沿用上述记号，并令谱坐标 \(\lambda:=X\)。
在开集
\[
U:=\{(y,T)\in \mathcal C_\delta:\ Q_{5}(y)\neq 0\}
\]
上，有显式有理反演
\[
\lambda=\frac{N(y,T)}{Q_{5}(y)},
\]
其中
\[
\begin{aligned}
N(y,T)=\;&32T^{3}y+19T^{3}+256T^{2}y^{3}+416T^{2}y^{2}+24T^{2}y-81T^{2}\\
&+512Ty^{4}+160Ty^{3}-709Ty^{2}-411Ty-20T\\
&-1792y^{5}-2528y^{4}-914y^{3}-377y^{2}-85y+80.
\end{aligned}
\]
并令
\[
Y:=y-\lambda^{2}+\frac12.
\]
则 \((\lambda,Y)\) 满足椭圆曲线方程
\[
Y^{2}=\lambda^{3}-\lambda+\frac14,
\]
且反演关系
\[
y=\lambda^{2}+Y-\frac12,
\qquad
T=4\lambda Y+3\lambda^{2}-1
\]
成立。
因此在 \(U\) 上，映射
\[
(\lambda,Y)\longmapsto (y,T)=\Bigl(\lambda^{2}+Y-\frac12,\ 4\lambda Y+3\lambda^{2}-1\Bigr)
\]
给出 \(E\dashrightarrow\mathcal C_\delta\) 的双有理参数化，而 \(\lambda=N/Q_{5}\) 给出其在 \(U\) 上的显式逆。
\end{theorem}

\begin{proof}[证明要点]
由结论 \ref{con:xi-terminal-zm-leyang-abel-jacobi-ode} 的消元链条，\(\QQ(E)=\QQ(y,T)\)，故 \(\lambda\in\QQ(y,T)\)。
另一方面，\((\lambda,y)\) 必满足谱四次
\[
\Pi(\lambda,y)=\lambda^{4}-\lambda^{3}-(2y+1)\lambda^{2}+\lambda+y(y+1)=0,
\]
并且由 \(T=4\lambda Y+3\lambda^{2}-1\) 与 \(Y=y-\lambda^{2}+\frac12\) 消去 \(Y\) 得到三次约束
\[
4\lambda^{3}-3\lambda^{2}-(4y+2)\lambda+(1+T)=0.
\]
在 \(\QQ[\lambda,y,T]\) 中对上述两式作消元，可得到线性关系
\[
Q_{5}(y)\lambda-N(y,T)=0,
\]
其中 \(Q_5\) 与 \(N\) 如上。
当 \(Q_5(y)\neq 0\) 时可解得 \(\lambda=N/Q_5\)。

取 \(Y:=y-\lambda^2+\tfrac12\)，则由定义有 \(y=\lambda^2+Y-\tfrac12\)。
将 \(\Pi(\lambda,y)=0\) 代入并整理即得
\(
Y^2=\lambda^3-\lambda+\tfrac14
\)，从而 \((\lambda,Y)\in E\)。
最后将 \((\lambda,Y)\) 代回 \(T=4\lambda Y+3\lambda^2-1\) 得恒等成立，故得双有理互逆。证毕。
\end{proof}

\begin{proposition}[平面射流坐标的显式反演]\label{prop:xi-terminal-zm-delta-jet-plane-explicit-inversion}
\leanverified{paper\_xi\_terminal\_zm\_delta\_jet\_plane\_explicit\_inversion}
在开集
\[
U_{\mathrm{inv}}:=\{(y,T)\in\mathcal C_\delta:\ \mathcal Q(y,T)\neq 0\},
\qquad
\mathcal Q(y,T):=8192y^{4}+5888y^{3}-3680y^{2}-2848y+152-93T,
\]
上，归一化同构 \(\QQ(E)\cong \QQ(y,T)/(F_\delta)\) 的谱坐标可写成显式有理式
\[
X=\frac{\mathcal N(y,T)}{\mathcal Q(y,T)},
\qquad
Y=y-X^{2}+\frac12,
\]
其中
\[
\begin{aligned}
\mathcal N(y,T)=\;&64T^{3}+512T^{2}y^{2}+528T^{2}y-219T^{2}+1024Ty^{3}-288Ty^{2}\\
&-1154Ty-131T-3584y^{4}-2928y^{3}-415y^{2}-929y+152.
\end{aligned}
\]
并且在 \(U_{\mathrm{inv}}\) 上恒等式
\[
y=X^{2}+Y-\frac12,
\qquad
T=4XY+3X^{2}-1
\]
成立。
\end{proposition}
\begin{proof}
由 \(\mathcal Q(y,T)=P_4(y)-93T\) 可知其零点恰对应平面模型中“射流分母退化”的条带；在开集 \(\mathcal Q\neq 0\) 上，可将函数域同构写为 \(X\in\QQ(y,T)\) 的显式有理表达式。
所示 \(\mathcal N/\mathcal Q\) 由对 \((F_\delta(T;y),\,4X^{3}-3X^{2}-(4y+2)X+(1+T))\) 的消元反演给出，并与定理 \ref{thm:xi-terminal-zm-delta-spectral-x-inversion} 的反演在共同定义域上一致。
令 \(Y=y-X^{2}+\tfrac12\) 后，直接代入 \(E:Y^2=X^3-X+\tfrac14\) 与 \(T=4XY+3X^2-1\) 即得两条恒等式。
\end{proof}

\begin{corollary}[判别式分解中平方因子的几何来源]\label{cor:xi-terminal-zm-delta-discriminant-square-geometry}
令 \(\widetilde{\mathcal C}_\delta\to\overline{\mathcal C}_\delta\) 为归一化，\(\pi_y:\widetilde{\mathcal C}_\delta\to\mathbb P^1_y\) 为由坐标 \(y\) 诱导的投影。
则：
\begin{enumerate}
  \item 因子 \(-y(y-1)P_{\mathrm{LY}}(y)\) 恰给出 \(\pi_y\) 的分歧像（有限分歧值），对应于 \(T=0\) 的分支；
  \item 多项式 \(Q_{5}(y)\) 的零点不对应 \(\pi_y\) 的分歧值，而对应于平面闭包 \(\overline{\mathcal C}_\delta\subset\mathbb P^{2}\) 的五个节点在 \(y\)-直线上的投影；
  \item 因而
  \[
  \mathrm{Disc}_{T}(F_\delta(T;y))
  =-\underbrace{y(y-1)P_{\mathrm{LY}}(y)}_{\text{归一化覆盖的分歧贡献}}
  \cdot
  \underbrace{Q_{5}(y)^{2}}_{\text{平面模型节点造成的平方因子}}
  \]
  给出“分歧贡献与平面自交贡献”的正交分解：平方因子完全记录平面化表示中的节点识别，而不改变归一化后的本征分歧集合。
\end{enumerate}
\end{corollary}

\begin{proof}[证明要点]
由命题 \ref{prop:xi-terminal-zm-delta-adjoint-conic-and-elimination} 的结果式分解
\[
\mathrm{Res}_{T}(F_\delta,\partial_TF_\delta)
=\mathrm{Res}_{T}(F_\delta,T)\,\mathrm{Res}_{T}(F_\delta,H),
\]
其中 \(\mathrm{Res}_{T}(F_\delta,T)=F_\delta(0;y)=-y(y-1)P_{\mathrm{LY}}(y)\) 对应 \(T=0\) 处的重根，即 \(\pi_y\) 的分歧像；
而 \(\mathrm{Res}_{T}(F_\delta,H)=Q_5(y)^2\) 对应 \(T\neq 0\) 的约化导数零点。
由定理 \ref{thm:xi-terminal-zm-delta-closure-plane-quintic}，这些点恰为平面模型的节点，其 \(y\)-坐标由 \(Q_5(y)=0\) 给出。
因此平方因子仅记录平面闭包的自交，并不引入新的归一化分歧值。证毕。
\end{proof}

\begin{theorem}[\texorpdfstring{$(5,4)$}{(5,4)} 双齐次闭包的有限节点与无穷远尖点]\label{thm:xi-terminal-zm-delta-bihomogeneous-closure-singularity-spectrum}
令 \(\overline{\mathcal C}_\delta^{\mathrm{bi}}\subset \mathbb P^1_y\times \mathbb P^1_T\) 为 \(F_\delta(T;y)=0\) 的 \((\deg_y,\deg_T)=(5,4)\) 双齐次闭包：取齐次坐标 \((y_0:y_1)\)、\((T_0:T_1)\)，并令
\[
F_\delta^{\mathrm{bi}}(T_0,T_1;y_0,y_1)
:=
y_1^5T_0^4
(8y_0y_1^4-3y_1^5)T_0^3T_1
(2y_0^2y_1^3-34y_0y_1^4-4y_1^5)T_0^2T_1^2
-y_0(y_0-y_1)P_{\mathrm{LY}}^{\mathrm{h}}(y_0,y_1)T_1^4,
\]
其中 \(P_{\mathrm{LY}}^{\mathrm{h}}(y_0,y_1)=256y_0^3+411y_0^2y_1+165y_0y_1^2+32y_1^3\)，并定义
\[
\overline{\mathcal C}_\delta^{\mathrm{bi}}:=\{(y_0:y_1;T_0:T_1):\ F_\delta^{\mathrm{bi}}(T_0,T_1;y_0,y_1)=0\}.
\]
则 \(\overline{\mathcal C}_\delta^{\mathrm{bi}}\) 的奇点谱满足：
\begin{enumerate}
  \item 在仿射图 \(y\neq\infty\)、\(T\neq\infty\) 上，\(\overline{\mathcal C}_\delta^{\mathrm{bi}}\) 恰有五个奇点，且全为普通节点；其 \(y\)-坐标恰为 \(Q_5(y)\) 的五个根，并且对应的 \(T\)-坐标为 \(\delta_\ast(y_0)\)，其中 \(\delta_\ast\) 为定理 \ref{thm:xi-fdelta-subresultant-unique-double-root} 所给出的唯一二重根闭式；
  \item 在边界 \(\{y=\infty\}\cup\{T=\infty\}\) 上，\(\overline{\mathcal C}_\delta^{\mathrm{bi}}\) 仅有一个点，且必为 \((y,T)=(\infty,\infty)\)；在局部坐标 \(z:=1/y\)、\(s:=1/T\) 下，其 Newton 主部与
  \[
  z^{5}-256\,s^{4}=0
  \]
  等价，从而该点为单分支尖点，Puiseux 对为 \((4,5)\)，并且其 \(\delta\)-不变量为 \(6\)；
  \item 除上述五个节点与一个 \((4,5)\)-尖点外，不存在任何额外奇点；
  \item \(\overline{\mathcal C}_\delta^{\mathrm{bi}}\) 的归一化为亏格 \(1\) 的光滑射影曲线，且其函数域与 \(\QQ(E)\) 同构；因而归一化与 \(E\) 在 \(\QQ\) 上同构。
\end{enumerate}
\end{theorem}

\begin{proof}[证明要点]
有限处结论由定理 \ref{thm:xi-terminal-zm-delta-closure-plane-quintic} 与定理 \ref{thm:xi-fdelta-subresultant-unique-double-root} 合并得到：判别式零点给出有限奇点候选，且在 \(Q_5(y)=0\) 处出现唯一二重根 \(\delta_\ast(y_0)\)；同时节点性可由切锥判别式的非退化性判定。

对无穷远处，双齐次方程迫使 \(y=\infty\) 与 \(T=\infty\) 同时发生，故边界候选点唯一。
在局部坐标 \(z=1/y\)、\(s=1/T\) 下，将 \(F_\delta(T;y)=0\) 写成 \(z,s\) 的方程并取最低阶主导平衡，即得 Newton 主部 \(z^5-256s^4\)。
由标准参数化 \(z=t^4\)、\(s=\tfrac{1}{4}t^5\) 可读出 Puiseux 对为 \((4,5)\)，从而 \(\delta\)-不变量为 \((4-1)(5-1)/2=6\)。

最后，\((5,4)\)-型曲线的算术亏格为 \((5-1)(4-1)=12\)；
五个节点贡献 \(\sum\delta_P=5\)，无穷远尖点贡献 \(6\)，故几何亏格为 \(12-(5+6)=1\)。
结合 \(\QQ(\overline{\mathcal C}_\delta^{\mathrm{bi}})=\QQ(y,T)/(F_\delta)\cong \QQ(E)\) 即得归一化同构断言。证毕。
\end{proof}

\begin{proposition}[整闭包导子与判别式的指数平方分解]\label{prop:xi-terminal-zm-delta-conductor-q5-and-discriminant-index-square}
令 \(k:=\QQ(y)\)、\(K:=\QQ(y,T)/(F_\delta)\cong \QQ(E)\)，并设
\[
A:=\QQ[y]\subset K,\qquad
B_0:=A[T]\big/(F_\delta(T;y)).
\]
记 \(B\) 为 \(A\) 在 \(K\) 中的整闭包（等价于 \(\mathcal C_\delta\) 的归一化在 \(y\neq\infty\) 处的坐标环）。
则 \(B_0\subset B\) 的导子理想满足
\[
\mathfrak f:=\operatorname{Ann}_{B}(B/B_0)=(Q_5(y))\subset A,
\]
并且判别式理想满足指数平方公式
\[
\operatorname{Disc}(B_0/A)=\operatorname{Disc}(B/A)\cdot (Q_5(y))^2,
\]
其中 \(\operatorname{Disc}(B_0/A)\) 由 \(F_\delta\) 的 \(T\)-判别式给出：
\[
\operatorname{Disc}(B_0/A)=\bigl(\mathrm{Disc}_T(F_\delta(T;y))\bigr)
=\bigl(-y(y-1)P_{\mathrm{LY}}(y)\,Q_5(y)^2\bigr),
\qquad
\operatorname{Disc}(B/A)=\bigl(-y(y-1)P_{\mathrm{LY}}(y)\bigr).
\]
\end{proposition}

\begin{proof}[证明要点]
由定理 \ref{thm:xi-terminal-zm-delta-closure-plane-quintic}，\(B_0\) 的整闭缺陷恰发生在五个节点之上，且其 \(y\)-坐标为 \(Q_5\) 的五个互异根。
对每个节点的完备局部模型，\(B_0\) 的局部环同构于 \(k[[u,v]]/(uv)\)，其整闭包为 \(k[[u]]\oplus k[[v]]\)，导子由极大理想生成；
在基环 \(A\) 上的像即为 \((y-y_0)\)。
对五点相乘得到 \(\mathfrak f=(Q_5(y))\)。

判别式的指数平方关系 \(\operatorname{Disc}(\text{order})=\operatorname{Disc}(\text{整闭包})\cdot (\text{指数})^2\) 在一维整环的有限扩张中逐点成立；
将五个节点处的局部指数相乘即得 \(Q_5(y)\)，从而得到所示分解。
最后两条显式判别式恒等式分别由
\(\mathrm{Disc}_T(F_\delta)=-y(y-1)P_{\mathrm{LY}}(y)\,Q_5(y)^2\)
以及归一化覆盖的分歧判别得到。证毕。
\end{proof}

\begin{proposition}[\texorpdfstring{$\delta$}{delta}-投影判别式的分歧--结点分离分解]\label{prop:xi-terminal-zm-delta-discy-factorization-a5b5}
将 \(F_\delta(T;y)\) 视为 \(\QQ[T]\) 上关于变量 \(y\) 的五次多项式。则其 \(y\)-判别式在 \(\ZZ[T]\) 中满足精确分解
\begin{equation}\label{eq:xi-terminal-zm-delta-discy-factorization-a5b5}
\boxed{\;
\mathrm{Disc}_{y}\!\bigl(F_\delta(T;y)\bigr)
=16\,(T-4)\,A_5(T)^2\,B_5(T)\; } ,
\end{equation}
其中
\[
\begin{aligned}
A_5(T)
&:=4096T^5-47360T^4+186992T^3-231939T^2-200880T+482112,\\
B_5(T)
&:=50000T^5+136112T^4+60776T^3-26712T^2+38961T+35964.
\end{aligned}
\]
因此 \(\mathrm{Disc}_{y}(F_\delta)\) 的平方因子恰为 \(A_5(T)^2\)，而其平方自由部分为 \(16\,(T-4)\,B_5(T)\)。
\end{proposition}
\begin{proof}
直接计算 \(\mathrm{Disc}_{y}(F_\delta(T;y))\) 并在 \(\ZZ[T]\) 中分解即可得到 \eqref{eq:xi-terminal-zm-delta-discy-factorization-a5b5}。
可复现核验脚本：\texttt{scripts/exp\_fold\_zm\_delta\_projection\_a5b5\_audit.py}。证毕。
\end{proof}

\begin{proposition}[结点子概形的全消元形与坐标域互换同构]\label{prop:xi-terminal-zm-delta-nodes-elimination-a5}
记 \(\mathcal C_\delta\subset\mathbb A^2_{(y,T)}\) 由 \(F_\delta(T;y)=0\) 给出，并令
\[
I_{\mathrm{sing}}:=(F_\delta,\partial_yF_\delta,\partial_TF_\delta)\subset \QQ[y,T].
\]
则 \(\mathcal C_\delta\) 的奇点子概形 \(\mathrm{Sing}(\mathcal C_\delta)\) 由理想 \(I_{\mathrm{sing}}\) 给出，并满足以下两组完全消元形：
\begin{enumerate}
  \item 在消去 \(T\) 后有
  \[
  I_{\mathrm{sing}}\cap \QQ[y]=(Q_5(y)),
  \qquad
  Q_{5}(y)=4096y^{5}+5376y^{4}-464y^{3}-2749y^{2}-723y+80,
  \]
  并且在 \(\QQ[y,T]/I_{\mathrm{sing}}\) 中出现一次关系
  \begin{equation}\label{eq:xi-terminal-zm-delta-node-linear-T}
  \boxed{\;
  93\,T=\bigl(8192y^{4}+5888y^{3}-3680y^{2}-2848y+152\bigr)\; } .
  \end{equation}
  \item 在消去 \(y\) 后有
  \[
  I_{\mathrm{sing}}\cap \QQ[T]=(A_5(T)),
  \]
  其中 \(A_5(T)\) 见命题 \ref{prop:xi-terminal-zm-delta-discy-factorization-a5b5}；并且在 \(\QQ[y,T]/I_{\mathrm{sing}}\) 中出现一次关系
  \begin{equation}\label{eq:xi-terminal-zm-delta-node-linear-y}
  \boxed{\;
  353664\,y+1863680T^{4}-16240384T^{3}+38835472T^{2}+5074731T-77031720=0\; } .
  \end{equation}
\end{enumerate}
从而结点的 \(y\)-坐标全体恰为 \(Q_5\) 的根集，结点的 \(T\)-坐标全体恰为 \(A_5\) 的根集。并且对任一结点 \((y_0,T_0)\) 成立严格域等式
\begin{equation}\label{eq:xi-terminal-zm-delta-node-coordinate-field-swap}
\QQ(y_0)=\QQ(T_0).
\end{equation}
\end{proposition}
\begin{proof}
理想刻画由雅可比判据给出。消元形 \eqref{eq:xi-terminal-zm-delta-node-linear-T} 已在定理 \ref{thm:xi-terminal-zm-delta-closure-plane-quintic} 的证明中出现；
而 \eqref{eq:xi-terminal-zm-delta-node-linear-y} 与 \(I_{\mathrm{sing}}\cap\QQ[T]=(A_5(T))\) 可由对 \(I_{\mathrm{sing}}\) 取 \(y\succ T\) 的字典序 Gröbner 消元得到。
由 \eqref{eq:xi-terminal-zm-delta-node-linear-T} 立得 \(T_0\in\QQ(y_0)\)，由 \eqref{eq:xi-terminal-zm-delta-node-linear-y} 立得 \(y_0\in\QQ(T_0)\)，从而得到 \eqref{eq:xi-terminal-zm-delta-node-coordinate-field-swap}。
可复现核验脚本：\texttt{scripts/exp\_fold\_zm\_delta\_projection\_a5b5\_audit.py}。证毕。
\end{proof}

\begin{theorem}[两个五次多项式的全对称算术单值化与分裂域一致]\label{thm:xi-terminal-zm-delta-a5-b5-galois-s5}
命题 \ref{prop:xi-terminal-zm-delta-discy-factorization-a5b5} 中的 \(A_5(T)\) 与 \(B_5(T)\) 在 \(\QQ\) 上不可约，且其伽罗瓦群均为全对称群 \(S_5\)。
并且 \(Q_5(y)\) 与 \(A_5(T)\) 具有相同分裂域；特别地，
\[
\Gal\bigl(A_5/\QQ\bigr)\cong \Gal\bigl(Q_5/\QQ\bigr)\cong S_5.
\]
\end{theorem}
\begin{proof}
先证 \(A_5\)。取素数 \(p=7\)，由 \(A_5\bmod 7\) 在 \(\FF_7[T]\) 中保持五次不可约可知 \(A_5\) 在 \(\QQ\) 上不可约，故其伽罗瓦群 \(G_A\le S_5\) 传递。
再取素数 \(p=61\)，可验证 \(A_5\bmod 61\) 的分解型为 \((2)(3)\)，故 \(G_A\) 含有一个 \(3\)-循环。
同时 \(\mathrm{Disc}(A_5)\) 在 \(\ZZ\) 中非平方，故 \(G_A\nsubseteq A_5\)。
由含 \(3\)-循环的传递子群分类可得 \(G_A=S_5\)。
\par
对 \(B_5\) 亦同理：取 \(p=59\) 得 \(B_5\bmod 59\) 不可约，故 \(G_B\) 传递；取 \(p=7\) 得分解型 \((2)(3)\) 从而含 \(3\)-循环；并由 \(\mathrm{Disc}(B_5)\) 非平方排除 \(A_5\)，遂 \(G_B=S_5\)。
\par
最后，由命题 \ref{prop:xi-terminal-zm-delta-nodes-elimination-a5} 的域互换同构 \eqref{eq:xi-terminal-zm-delta-node-coordinate-field-swap}：对任一结点 \((y_0,T_0)\) 有 \(\QQ(y_0)=\QQ(T_0)\)。
因此 \(A_5\) 的根生成的分裂域与 \(Q_5\) 的根生成的分裂域一致，从而两者的伽罗瓦群同构。
可复现核验脚本：\texttt{scripts/exp\_fold\_zm\_delta\_projection\_a5b5\_audit.py}。证毕。
\end{proof}

\begin{theorem}[\texorpdfstring{$\delta:E\to\PP^1_\delta$}{delta-map} 的分歧护照与几何单值群]\label{thm:xi-terminal-zm-delta-map-hurwitz-passport-s5}
将 \(\delta=\dd y/\dd u=4XY+3X^2-1\in\QQ(E)\) 视为 \(E\) 上的有理函数，并记由其诱导的有限态射
\[
\pi_\delta:\ E\longrightarrow \PP^1_\delta .
\]
则：
\begin{enumerate}
  \item \(\delta\) 的极点除子满足 \((\delta)_\infty=5O\)。因此 \(\deg(\pi_\delta)=5\)，并且在 \(\delta=\infty\) 处全分歧（分歧指数 \(e_O=5\)）。
  \item \(\pi_\delta\) 的有限分歧点恰由 \(\dd \delta=0\) 给出；其 \(X\)-坐标满足完全分解
  \[
  (X+1)\bigl(100X^5-136X^4+16X^3+40X^2-13X+1\bigr)=0,
  \]
  且其 \(Y\)-坐标由
  \[
  Y=\frac{-10X^3+6X-1}{6X}
  \]
  唯一确定。上述六点均为简单分歧点（分歧指数 \(2\)）。
  \item 有限分歧值集合为
  \[
  \{\delta=4\}\ \cup\ \{\delta\in\overline{\QQ}:\ B_5(\delta)=0\},
  \]
  其中 \(B_5\) 如命题 \ref{prop:xi-terminal-zm-delta-discy-factorization-a5b5}。
  更精确地，消去 \(X\) 得到严格结果式恒等式
  \begin{equation}\label{eq:xi-terminal-zm-delta-branch-resultant-b5}
  \boxed{\;
  \mathrm{Res}_X\!\left(
  100X^5-136X^4+16X^3+40X^2-13X+1,\ \ 3\delta+20X^3-9X^2-12X+5
  \right)=4860\,B_5(\delta)\; }.
  \end{equation}
  \item 该覆盖的几何单值群同构于 \(S_5\)。等价地，函数域扩张 \(\QQ(E)/\QQ(\delta)\) 的伽罗瓦闭包群为 \(S_5\)。
\end{enumerate}
\end{theorem}
\begin{proof}
由 \(\operatorname{ord}_O(X)=-2\)、\(\operatorname{ord}_O(Y)=-3\) 得
\(
\operatorname{ord}_O(\delta)=-(2+3)=-5
\)，从而 \((\delta)_\infty=5O\) 并得到第一条。
\par
取 \(\omega=\dd X/(2Y)\) 并用 \(\dd Y=(3X^2-1)\omega\)、\(\dd X=2Y\omega\) 微分 \(\delta=4XY+3X^2-1\)，得
\[
\dd\delta=2\bigl(10X^3+6XY-6X+1\bigr)\,\omega.
\]
从而有限分歧点由 \(10X^3+6XY-6X+1=0\) 给出，并解得所示 \(Y\) 的闭式；
代回 \(Y^2=X^3-X+\tfrac14\) 化简即得 \(X\)-坐标多项式的完全分解，得到第二条。
由 Riemann--Hurwitz 公式，度为 \(5\) 的覆盖 \(E\to\PP^1\) 的总分歧度为 \(2\deg(\pi_\delta)=10\)；
无穷远点贡献 \(e_O-1=4\)，故有限分歧总贡献为 \(6\)，与上述六个有限分歧点一致，从而它们均为简单分歧点。
\par
第三条中，\(\delta\) 在 \(X=-1\) 的分歧点处取值为 \(4\)；
对五次因子根消去 \(X\) 即得 \eqref{eq:xi-terminal-zm-delta-branch-resultant-b5}，从而其分歧值满足 \(B_5(\delta)=0\)。
\par
第四条由分歧循环型给出：\(\delta=\infty\) 处全分歧对应一个 \(5\)-循环；任一有限分歧值处为简单分歧，对应换位。
包含 \(5\)-循环与换位的传递子群必为 \(S_5\)。证毕。
\end{proof}

\begin{proposition}[特征 \texorpdfstring{$37$}{37} 上的分歧塌缩：\texorpdfstring{$B_5$}{B5} 的双重根]\label{prop:xi-terminal-zm-delta-badprime-37-b5-double-root}
\leanverified{paper\_xi\_terminal\_zm\_delta\_badprime\_37\_b5\_double\_root}
令 \(B_5(\delta)\) 如命题 \ref{prop:xi-terminal-zm-delta-discy-factorization-a5b5}。
则在 \(\FF_{37}[\delta]\) 中有精确分解
\[
B_5(\delta)\equiv 13\,\delta^2\bigl(\delta^3+2\delta^2-4\delta+3\bigr)\pmod{37}.
\]
因此 \(\pi_\delta\) 的有限分歧值在特征 \(37\) 的约化中发生并合，等价地，\(\delta\)-分歧多项式在 \(p=37\) 处出现重根并导致该覆盖不具良分歧约化。
\end{proposition}
\begin{proof}
对 \(B_5\) 的系数作模 \(37\) 化简并因式分解即可得到所示恒等式。
可复现核验脚本：\texttt{scripts/exp\_fold\_zm\_delta\_projection\_a5b5\_audit.py}。证毕。
\end{proof}

\begin{corollary}[有限节点集的伽罗瓦传递性]\label{cor:xi-terminal-zm-delta-nodes-single-galois-orbit}
\(\overline{\mathcal C}_\delta^{\mathrm{bi}}\) 的五个有限节点在 \(\operatorname{Gal}(\overline{\QQ}/\QQ)\) 作用下构成单一轨道；特别地，\(\overline{\mathcal C}_\delta^{\mathrm{bi}}\) 在 \(\QQ\) 上不存在有限有理奇点。
\end{corollary}

\begin{proof}
由定理 \ref{thm:xi-terminal-zm-delta-bihomogeneous-closure-singularity-spectrum}，有限节点的 \(y\)-坐标恰为 \(Q_5\) 的根；
由定理 \ref{thm:xi-q5-galois-s5}，\(\operatorname{Gal}(Q_5/\QQ)\cong S_5\) 在根集上传递，故节点集亦为单一伽罗瓦轨道。
若存在 \(\QQ\)-有理有限节点，则其 \(y\)-坐标给出 \(Q_5\) 的有理根，与不可约性矛盾。证毕。
\end{proof}

\input{thm__xi-terminal-zm-leyang-delta-node-preimage-r10-wreath}

\endinput

\begin{corollary}[完全分裂素数处的奇偶律]\label{cor:xi-terminal-zm-delta-node-tangent-parity-law}
沿用推论 \ref{cor:xi-terminal-zm-delta-node-tangent-chebotarev} 的设定，并进一步要求 \(p\neq 2\) 且在 \(L/\QQ\) 中不分歧。
在事件 \(E_p\) 下，将 \(Q_5\) 的五个根在 \(\FF_p\) 中记为 \(y_1(p),\dots,y_5(p)\)，并令
\[
\varepsilon(p):=\left(\left(\frac{D(y_1(p))}{p}\right),\dots,\left(\frac{D(y_5(p))}{p}\right)\right)\in\{\pm 1\}^5.
\]
则成立确定性恒等式
\[
\prod_{i=1}^{5}\left(\frac{D(y_i(p))}{p}\right)=\left(\frac{-3}{p}\right).
\]
等价地，若记 \(N_-(p):=\#\{i:\ (\tfrac{D(y_i(p))}{p})=-1\}\)，则
\[
(-1)^{N_-(p)}=\left(\frac{-3}{p}\right).
\]
\end{corollary}
\begin{proof}
由结论 \ref{con:xi-terminal-zm-delta-node-tangent-quadratic-subfields}，\(L\) 的“总翻转”二次特征的固定域为 \(\QQ(\sqrt{-3})\)。
在事件 \(E_p\) 下，\(\mathrm{Frob}_p\) 落在核 \((C_2)^5\) 中，其在各坐标上的符号恰为 \((\tfrac{D(y_i(p))}{p})\)；
因而该二次特征在 \(\mathrm{Frob}_p\) 处的取值为 \(\prod_i(\tfrac{D(y_i(p))}{p})\)。
另一方面，二次子域 \(\QQ(\sqrt{-3})\) 的 Frobenius 特征为 \((\tfrac{-3}{p})\)，故得到所示恒等式。证毕。
\end{proof}

\endinput

\begin{conclusion}[\texorpdfstring{$\delta$}{delta}-分支五次的判别式谱、分歧素集与二次分辨子域]\label{con:xi-terminal-zm-delta-b5-discriminant-ramification-quadratic}
沿用命题 \ref{prop:xi-terminal-zm-delta-discy-factorization-a5b5} 的 \(\delta\)-分支五次多项式 \(B_5(\delta)\in\ZZ[\delta]\)，记其分裂域为 \(L_B/\QQ\)。
则其判别式在 \(\ZZ\) 中满足素因子分解
\[
\Disc(B_5)
=-2^{16}\cdot 3^{16}\cdot 17^{3}\cdot 37\cdot 223^{3}\cdot 3739^{2}\cdot 7333^{2},
\]
从而 \(L_B/\QQ\) 的分歧素数集合包含于
\[
S_B=\{2,3,17,37,223,3739,7333\}.
\]
并且 \(\Disc(B_5)\notin(\QQ^\times)^2\)，故 \(L_B\) 的二次分辨子域为
\[
\QQ\!\left(\sqrt{\Disc(B_5)}\right)
=\QQ\!\left(\sqrt{-17\cdot 37\cdot 223}\right)
=\QQ(\sqrt{-140267}).
\]
\end{conclusion}
\begin{proof}
判别式分解由直接计算给出，可复现核验脚本为 \texttt{scripts/exp\_fold\_zm\_delta\_b5\_arithmetic\_audit.py}。
由 \(\Disc(B_5)\) 非平方，二次分辨子域由符号表征的判别式刻画得到。证毕。
\end{proof}

\begin{conclusion}[节点消元五次的判别式全素分解与二次分辨子域]\label{con:xi-terminal-zm-delta-a5-discriminant-quadratic}
沿用命题 \ref{prop:xi-terminal-zm-delta-discy-factorization-a5b5} 的节点消元五次多项式 \(A_5(T)\in\ZZ[T]\)，记其分裂域为 \(L_A/\QQ\)。
则其判别式在 \(\ZZ\) 中满足素因子分解
\[
\Disc(A_5)=-2^{42}\cdot 3^{11}\cdot 11^{2}\cdot 31\cdot 307^{2}\cdot 727,
\]
从而 \(\Disc(A_5)\notin(\QQ^\times)^2\)，并且其平方自由核满足
\[
\mathrm{sf}\bigl(\Disc(A_5)\bigr)=-3\cdot 31\cdot 727=-67611.
\]
因此当 \(\Gal(L_A/\QQ)\cong S_5\)（定理 \ref{thm:xi-terminal-zm-delta-a5-b5-galois-s5}）时，\(L_A\) 的唯一二次分辨子域为
\[
\QQ\!\left(\sqrt{\Disc(A_5)}\right)=\QQ(\sqrt{-67611}).
\]
并且由同一分裂域识别（定理 \ref{thm:xi-terminal-zm-delta-a5-b5-galois-s5}），该二次子域亦等于 \(Q_5(y)\) 的分裂域 \(L_5\) 的唯一二次子域（定理 \ref{thm:xi-leyang-ply-q5-linearly-disjoint-chebotarev} 的证明）。
\end{conclusion}
\begin{proof}
对 \(A_5\) 直接作判别式计算并在 \(\ZZ\) 中素分解得到所示分解式；可复现核验脚本为 \texttt{scripts/exp\_fold\_zm\_delta\_a5\_discriminant\_audit.py}。
平方自由核与二次分辨子域的刻画为标准伽罗瓦理论：当 \(\Gal(\mathrm{Spl}(f)/\QQ)\cong S_n\) 时，唯一指数 \(2\) 正规子群为 \(A_n\)，对应二次子域由 \(\sqrt{\Disc(f)}\) 生成。证毕。
\end{proof}

\begin{conclusion}[\texorpdfstring{$\delta$}{delta}-分支谱的坏约化素数与碰撞型的算术指纹]\label{con:xi-terminal-zm-delta-b5-bad-reduction-primes-collision-fingerprint}
令 \(S_B\) 如结论 \ref{con:xi-terminal-zm-delta-b5-discriminant-ramification-quadratic}。
则：
\begin{enumerate}
  \item 若 \(p\notin S_B\)，则 \(B_5\bmod p\) 在 \(\FF_p[\delta]\) 上可分，等价地，\(\delta\)-分支谱在特征 \(p\) 下保持分离；
  \item 若 \(p\in\{17,37,223,3739,7333\}\)，则 \(B_5\bmod p\) 恰出现一个二重根，且该二重根同余类分别为
  \[
  \delta\equiv -2\ (17),\qquad
  \delta\equiv 0\ (37),\qquad
  \delta\equiv -56\ (223),\qquad
  \delta\equiv 48\ (3739),\qquad
  \delta\equiv 3257\ (7333),
  \]
  其余根均为单根；
  \item 在特征 \(3\) 下发生更高阶塌缩：
  \[
  B_5(\delta)\equiv \delta^{3}(\delta-1)^{2}\pmod 3.
  \]
\end{enumerate}
\end{conclusion}
\begin{proof}
第一条为判别式与可分性的同值命题。
第二、三条由对 \(B_5\) 在相应有限域上的直接因式分解得到；可复现核验脚本为 \texttt{scripts/exp\_fold\_zm\_delta\_b5\_arithmetic\_audit.py}。证毕。
\end{proof}

\input{thm__xi-terminal-zm-delta-s5-tame-collision-semistable-fiber}

\begin{conclusion}[\texorpdfstring{$S_5$}{S5} 正则闭包曲线的精确属数]\label{con:xi-terminal-zm-delta-s5-regular-closure-genus-109}
设 \(k=\QQ(\delta)\)，并令 \(K=\QQ(E)\) 视为次数 \(5\) 的函数域扩张 \(K/k\)（由定理 \ref{thm:xi-terminal-zm-delta-map-hurwitz-passport-s5} 的 \(\pi_\delta:E\to\PP^1_\delta\) 给出）。
令 \(L/k\) 为该扩张的伽罗瓦闭包，\(G:=\Gal(L/k)\cong S_5\)，并令 \(X\) 为函数域满足 \(\QQ(X)=L\) 的光滑射影曲线（正则闭包曲线）。
记 \(\widehat\pi_\delta:X\to\PP^1_\delta\) 为诱导的正则闭包覆盖，则 \(\widehat\pi_\delta\) 为次数 \(120\) 的正则 \(S_5\)-伽罗瓦覆盖，且
\[
g(X)=109.
\]
\end{conclusion}
\begin{proof}[证明要点]
由定理 \ref{thm:xi-terminal-zm-delta-map-hurwitz-passport-s5}，分歧护照为一个阶 \(5\) 惯性点与六个阶 \(2\) 惯性点，即 \((5)\cdot(2)^6\)。
对正则伽罗瓦覆盖应用 Hurwitz 公式
\[
g(X)=1+\frac{|S_5|}{2}\!\left(-2+\sum_{p\in B}\Bigl(1-\frac1{e_p}\Bigr)\right),
\]
代入 \(|S_5|=120\)、\((e_p)=(5,2,2,2,2,2,2)\) 即得
\[
g(X)=1+60\left(-2+\frac45+6\cdot\frac12\right)=109.
\]
证毕。
\end{proof}

在本小节后续各处，\(\widehat\pi_\delta:X\to\PP^1_\delta\) 均表示上述正则闭包；并固定其为正则 \(S_5\)-伽罗瓦覆盖、分歧护照为 \((5)\cdot(2)^6\)、\(\Gal(L/\QQ(\delta))\cong S_5\)，以及 \(g(X)=109\)。

\begin{conclusion}[\texorpdfstring{$A_5$}{A5} 固定域的判别式模型与属 \(2\) 超椭圆商]\label{con:xi-terminal-zm-delta-a5-fixedfield-hyperelliptic-genus2}
沿用结论 \ref{con:xi-terminal-zm-delta-s5-regular-closure-genus-109} 的记号。
则唯一指数 \(2\) 的中间子域 \(L^{A_5}\) 由判别式平方根给出：
\[
L^{A_5}=\QQ\!\left(\delta,\sqrt{(\delta-4)\,B_5(\delta)}\right),
\]
其中 \(B_5\) 见命题 \ref{prop:xi-terminal-zm-delta-discy-factorization-a5b5}。
因此 \(X/A_5\) 具有仿射模型
\[
C_{A_5}:\quad s^2=(\delta-4)\,B_5(\delta),
\]
并且 \(g(C_{A_5})=2\)。
\end{conclusion}
\begin{proof}[证明要点]
将 \(F_\delta(T;y)\) 视为关于 \(y\) 的五次多项式（系数在 \(\QQ(T)\) 中），其 \(y\)-判别式满足
\(
\Disc_{y}(F_\delta(T;y))=16\,(T-4)\,A_5(T)^2B_5(T)
\)
（命题 \ref{prop:xi-terminal-zm-delta-discy-factorization-a5b5}）。
对五次可分扩张，伽罗瓦闭包群为 \(S_5\) 时，\(A_5\) 固定域由“判别式平方根”生成；此处平方自由部分等价于 \((T-4)B_5(T)\)（常数 \(16\) 为全局平方），故得所示 \(L^{A_5}\) 与超椭圆模型。
\par
又 \((T-4)B_5(T)\) 在特征 \(0\) 上可分且次数为 \(6\)，故双覆盖在 \(\overline{\QQ}\) 上有 \(6\) 个互异分歧点，从而
\(
g(C_{A_5})=(6-2)/2=2
\)。
证毕。
\end{proof}

\input{subsubsec__xi-terminal-zm-delta-s4-layer-branch-fiberproduct}

\begin{conclusion}[全纯 \(1\)-形式的 Chevalley--Weil 精确分解]\label{con:xi-terminal-zm-delta-s5-holomorphic-oneforms-chevalley-weil}
令 \(X\) 如结论 \ref{con:xi-terminal-zm-delta-s5-regular-closure-genus-109}。
将 \(S_5\) 的不可约复表示按分拆 \(\lambda\vdash 5\) 记为 \(V_\lambda\)，并记符号表示为 \(\varepsilon\)。
则 \(S_5\) 在 \(H^0(X,\Omega_X^1)\) 上的表示分解为
\[
H^0(X,\Omega_X^1)\ \cong\
\varepsilon^{\oplus 2}\ \oplus\
V_{(4,1)}^{\oplus 1}\ \oplus\
V_{(3,2)}^{\oplus 3}\ \oplus\
V_{(3,1,1)}^{\oplus 5}\ \oplus\
V_{(2,2,1)}^{\oplus 6}\ \oplus\
V_{(2,1,1,1)}^{\oplus 7},
\]
并且平凡表示 \(V_{(5)}\) 不出现。
特别地，\(\dim_\CC H^0(X,\Omega_X^1)=g(X)=109\)。
\end{conclusion}
\begin{proof}[证明要点]
本覆盖的惯性群仅有 \(C_5\) 与 \(C_2\) 两类（定理 \ref{thm:xi-terminal-zm-delta-map-hurwitz-passport-s5}）。
对 \(\rho\neq V_{(5)}\)，Chevalley--Weil 公式在本情形可写为
\[
m_\rho=-\dim\rho+\frac12\sum_{p\in B}\bigl(\dim\rho-\dim(\rho^{I_p})\bigr),
\qquad
m_{V_{(5)}}=0,
\]
其中 \(I_p\) 为分歧点的惯性群。
所需不变维数可由特征平均给出：
\[
\dim(\rho^{C_2})=\frac{\dim\rho+\chi_\rho(\text{换位})}{2},
\qquad
\dim(\rho^{C_5})=\frac{\dim\rho+4\chi_\rho(\text{(5)-循环})}{5}.
\]
对各不可约表示用 Murnaghan--Nakayama 规则计算 \(\chi_\rho\) 在换位与 \(5\)-循环处的值，得到
\[
\begin{array}{c|cccccc}
\rho & \varepsilon & V_{(4,1)} & V_{(3,2)} & V_{(3,1,1)} & V_{(2,2,1)} & V_{(2,1,1,1)}\\
\hline
\dim(\rho^{C_2}) & 0 & 3 & 3 & 3 & 2 & 1\\
\dim(\rho^{C_5}) & 1 & 0 & 1 & 2 & 1 & 0
\end{array}
\]
并代入一个 \(C_5\) 分歧点与六个 \(C_2\) 分歧点，得到重数
\(
(2,1,3,5,6,7)
\)。
维数求和即得 \(\dim H^0(X,\Omega^1)=109\)，与结论 \ref{con:xi-terminal-zm-delta-s5-regular-closure-genus-109} 一致。证毕。
\end{proof}

\input{thm__xi-terminal-zm-delta-s5-pluricanonical-character-and-artin-conductor}

\begin{conclusion}[Jacobian 的 \texorpdfstring{$\QQ$}{Q}-同源分解与属 \(2\) 判别曲线因子]\label{con:xi-terminal-zm-delta-s5-jacobian-qisogeny-decomposition}
设 \(J(X)=\mathrm{Jac}(X)\)。
则存在定义于 \(\QQ\) 上的阿贝尔簇
\[
A_\varepsilon,\ A_{(4,1)},\ A_{(3,2)},\ A_{(3,1,1)},\ A_{(2,2,1)},\ A_{(2,1,1,1)},
\]
使得在 \(\QQ\)-同源意义下
\[
J(X)\sim_{\QQ}
A_\varepsilon\times A_{(4,1)}\times A_{(3,2)}\times A_{(3,1,1)}\times A_{(2,2,1)}\times A_{(2,1,1,1)},
\]
且
\[
\dim A_\varepsilon=2,\quad
\dim A_{(4,1)}=4,\quad
\dim A_{(3,2)}=15,\quad
\dim A_{(3,1,1)}=30,\quad
\dim A_{(2,2,1)}=30,\quad
\dim A_{(2,1,1,1)}=28.
\]
此外，结论 \ref{con:xi-terminal-zm-delta-a5-fixedfield-hyperelliptic-genus2} 的属 \(2\) 曲线 \(C_{A_5}=X/A_5\) 满足
\[
A_\varepsilon\sim_{\QQ}\mathrm{Jac}(C_{A_5}).
\]
\end{conclusion}
\begin{proof}[证明要点]
由 \(S_5\) 对 \(X\) 的作用，群代数 \(\QQ[S_5]\) 经由自同构诱导自同态嵌入 \(\mathrm{End}^0(J(X))\)。
结合结论 \ref{con:xi-terminal-zm-delta-s5-holomorphic-oneforms-chevalley-weil} 中 \(H^0(X,\Omega^1)\) 的等型分解，可由 \(\QQ[S_5]\) 的中心幂等元得到 \(J(X)\) 的 \(\QQ\)-同源等型分裂；各因子维数由“重数 \(\times\) 表示维数”给出。
\par
对 \(A_\varepsilon\)：由于 \(\varepsilon\) 在 \(A_5\) 上平凡且在陪集上变号，
有 \(H^0(X,\Omega^1)^{A_5}\cong \varepsilon^{\oplus 2}\)，从而
\(
g(X/A_5)=\dim H^0(X,\Omega^1)^{A_5}=2
\)。
由正规商的函子性与 Kani--Rosen 型分解的标准对应，可识别 \(\mathrm{Jac}(X/A_5)\) 正对应于 \(\varepsilon\) 等型分量，故 \(A_\varepsilon\sim_{\QQ}\mathrm{Jac}(C_{A_5})\)。证毕。
\end{proof}

\begin{conclusion}[\texorpdfstring{$S_5$}{S5} 对称性强制的 N\'eron--Severi 秩下界]\label{con:xi-terminal-zm-delta-s5-jacobian-neron-severi-rank-lower-bound-73}
设 \(J:=J(X)=\mathrm{Jac}(X)\)。
则
\[
\rho(J):=\rank \mathrm{NS}(J)\ \ge\ 73.
\]
更精确地，\(\QQ[S_5]\subset \End^0(J)\) 在 Rosati 反演 \(\dagger\) 下的不变子空间维数为 \(73\)，从而在
\[
\mathrm{NS}(J)_{\QQ}\cong \End^0(J)^{\dagger=1}
\]
中给出一个显式 \(73\) 维子空间。
\end{conclusion}
\begin{proof}[证明要点]
群作用诱导嵌入 \(\QQ[S_5]\hookrightarrow \End^0(J)\)，且 Rosati 在群代数上对应 \(g\mapsto g^{-1}\)。
又因 \(S_5\) 的不可约特征标取整数值，\(\QQ[S_5]\) 的 Wedderburn 分解在 \(\QQ\) 上裂开，并且
\[
\QQ[S_5]\ \cong\ \QQ\ \oplus\ \QQ\ \oplus\ M_4(\QQ)\ \oplus\ M_4(\QQ)\ \oplus\ M_5(\QQ)\ \oplus\ M_5(\QQ)\ \oplus\ M_6(\QQ),
\]
其中矩阵大小对应 \(S_5\) 的不可约维数 \(1,1,4,4,5,5,6\)。
在每个 \(M_n(\QQ)\) 上，\(\dagger\) 与转置同型，故不变子空间维数为对称矩阵维数 \(n(n+1)/2\)。
因此
\[
\dim_{\QQ}\QQ[S_5]^{\dagger=1}
=1+1+\frac{4\cdot 5}{2}+\frac{4\cdot 5}{2}+\frac{5\cdot 6}{2}+\frac{5\cdot 6}{2}+\frac{6\cdot 7}{2}
=73.
\]
最后用主极化阿贝尔簇的标准同构 \(\mathrm{NS}(J)_{\QQ}\cong \End^0(J)^{\dagger=1}\) 即得结论；亦可参照 \cite{Internal2026S5JacobianPrymNSRank} 的群代数--Rosati 口径。证毕。
\end{proof}

\begin{conclusion}[Jacobian 的幂次型精化：由 \texorpdfstring{$Y_{10}$}{Y10} 与 \texorpdfstring{$Z_{28}$}{Z28} 投影抽取的乘法分量]\label{con:xi-terminal-zm-delta-s5-jacobian-power-isogeny-refinement}
沿用结论 \ref{con:xi-terminal-zm-delta-s5-jacobian-qisogeny-decomposition} 的记号。
则存在椭圆曲线 \(E\)（\(\dim E=1\)）与阿贝尔簇 \(B_3,B_5,B_6,B_7\)（分别满足 \(\dim(B_3,B_5,B_6,B_7)=(3,5,6,7)\)），使得在 \(\QQ\)-同源意义下
\[
J(X)\ \sim_{\QQ}\ A_\varepsilon\ \times\ E^{\,4}\ \times\ B_3^{\,5}\ \times\ B_5^{\,6}\ \times\ B_6^{\,5}\ \times\ B_7^{\,4}.
\]
其中 \(B_7\) 可由 \(Y_{10}:=X/A_4\) 的 Jacobian 的群代数投影切出（结论 \ref{con:xi-terminal-zm-delta-y10-jacobian-qisogeny-1-2-7}），而 \(B_3,B_6\) 可由 \(Z_{28}:=X/V_4\) 的 Jacobian 的群代数投影切出（结论 \ref{con:xi-terminal-zm-delta-z28-jacobian-qisogeny-1-2-7-6-12}）。
\end{conclusion}
\begin{proof}[证明要点]
由结论 \ref{con:xi-terminal-zm-delta-s5-holomorphic-oneforms-chevalley-weil}，对各不可约表示 \(V_\lambda\) 的出现重数记为 \(m_\lambda\)。
由于 \(\QQ[S_5]\) 在 \(\QQ\) 上裂开，各等型分量 \(A_\lambda\) 的端同态代数含有满矩阵块 \(M_{\dim V_\lambda}(\QQ)\)，从而由 Poincar\'e 完全可约性可取一在 \(\QQ\)-同源意义下唯一的阿贝尔簇 \(B_\lambda\)（\(\dim B_\lambda=m_\lambda\)），满足
\[
A_\lambda\ \sim_{\QQ}\ B_\lambda^{\,\dim V_\lambda}.
\]
对 \(\lambda=(4,1)\) 有 \(m_\lambda=1\)、\(\dim V_\lambda=4\)，故 \(A_{(4,1)}\sim_{\QQ}E^{\,4}\)（取 \(E:=B_{(4,1)}\)）。
对 \(\lambda=(3,2),(3,1,1),(2,2,1),(2,1,1,1)\) 分别取
\[
B_3:=B_{(3,2)},\quad B_5:=B_{(3,1,1)},\quad B_6:=B_{(2,2,1)},\quad B_7:=B_{(2,1,1,1)},
\]
并由 \((m_\lambda)=(3,5,6,7)\) 与 \((\dim V_\lambda)=(5,6,5,4)\) 得到相应幂次型分解。
最后将其与结论 \ref{con:xi-terminal-zm-delta-s5-jacobian-qisogeny-decomposition} 的总分解合并即可；而 \(B_7,B_3,B_6\) 的几何识别由结论 \ref{con:xi-terminal-zm-delta-y10-jacobian-qisogeny-1-2-7} 与 \ref{con:xi-terminal-zm-delta-z28-jacobian-qisogeny-1-2-7-6-12} 给出。证毕。
\end{proof}

\input{thm__xi-terminal-zm-delta-s5-tame-collision-torus-rank-signature}

\begin{conclusion}[Hurwitz--Burnside 属数公式：中间商曲线的属数由 \(5\)-循环与换位计数控制]\label{con:xi-terminal-zm-delta-s5-intermediate-genus-hurwitz-burnside}
沿用结论 \ref{con:xi-terminal-zm-delta-s5-regular-closure-genus-109} 的记号，令 \(G:=S_5\)。
对任意子群 \(H\le G\)，记中间商曲线 \(Y_H:=X/H\)。
并定义
\[
a(H):=\#\{h\in H:\ h\ \text{为}\ (5)\text{-循环}\},\qquad
b(H):=\#\{h\in H:\ h\ \text{为换位}\}.
\]
则 \(Y_H\) 的几何属数满足显式恒等式
\[
g(Y_H)=1+\frac{108-2a(H)-18b(H)}{|H|}.
\]
\end{conclusion}
\begin{proof}[证明要点]
令 \(n:=[G\!:\!H]=120/|H|\)。
由定理 \ref{thm:xi-terminal-zm-delta-map-hurwitz-passport-s5}，覆盖 \(X\to\PP^1_\delta\) 的分歧惯性共轭类为一个 \(5\)-循环与六个换位。
因此非伽罗瓦覆盖 \(Y_H\to\PP^1_\delta\) 的 Hurwitz 公式可写为
\[
2g(Y_H)-2=-2n+\Bigl(n-c_5(H)\Bigr)+6\Bigl(n-c_2(H)\Bigr),
\]
其中 \(c_\ell(H)\) 为固定的 \(\ell\)-阶惯性元在陪集集 \(G/H\) 上诱导置换的循环数（\(\ell\in\{5,2\}\)）。
\par
由于 \(\ell\) 为素数，轨道长度只能为 \(1\) 或 \(\ell\)。
若记 \(f_\ell(H)\) 为该置换在 \(G/H\) 上的不动点数，则有
\[
c_5(H)=\frac{n+4f_5(H)}5,\qquad
c_2(H)=\frac{n+f_2(H)}2.
\]
另一方面，对固定惯性元 \(u\in G\)（阶为 \(\ell\)），有
\[
f_\ell(H)=\#\{gH:\ g^{-1}ug\in H\}
      =\frac{\#\{g\in G:\ g^{-1}ug\in H\}}{|H|}
      =\frac{|C_G(u)|\cdot |H\cap \mathrm{Cl}(u)|}{|H|},
\]
其中 \(C_G(u)\) 为中心化子，\(\mathrm{Cl}(u)\) 为共轭类。
对 \(S_5\)，当 \(u\) 为 \(5\)-循环时 \(|C_G(u)|=5\) 且 \(|H\cap\mathrm{Cl}(u)|=a(H)\)；
当 \(u\) 为换位时 \(|C_G(u)|=12\) 且 \(|H\cap\mathrm{Cl}(u)|=b(H)\)。
故
\[
f_5(H)=\frac{5a(H)}{|H|},\qquad
f_2(H)=\frac{12b(H)}{|H|}.
\]
代入并化简（用 \(n=120/|H|\)）即得
\(
g(Y_H)=1+\frac{108-2a(H)-18b(H)}{|H|}.
\)
证毕。
\end{proof}

\input{parts/thm__xi-terminal-zm-delta-s5-intermediate-quotient-collision-degeneration}

\begin{conclusion}[全部中间商曲线的属数谱：\(S_5\) 子群共轭类的统一表]\label{con:xi-terminal-zm-delta-s5-intermediate-genus-spectrum-19}
在结论 \ref{con:xi-terminal-zm-delta-s5-intermediate-genus-hurwitz-burnside} 的记号下，\(S_5\) 的子群按共轭分类恰为 \(19\) 类。
对每一类取代表 \(H\)，其 \(\bigl(|H|,a(H),b(H),[S_5\!:\!H],g(Y_H)\bigr)\) 与该共轭类中子群个数如下：
\[
\begin{array}{c|c|c|c|c|c|c}
H & |H| & a(H) & b(H) & [S_5\!:\!H] & g(Y_H) & \#\text{共轭子群}\\
\hline
S_5 & 120 & 24 & 10 & 1 & 0 & 1\\
A_5 & 60 & 24 & 0 & 2 & 2 & 1\\
S_4\ (\text{一点稳定子}) & 24 & 0 & 6 & 5 & 1 & 5\\
F_{20}=C_5\rtimes C_4 & 20 & 4 & 0 & 6 & 6 & 6\\
S_3\times C_2 & 12 & 0 & 4 & 10 & 4 & 10\\
A_4 & 12 & 0 & 0 & 10 & 10 & 5\\
D_{10} & 10 & 4 & 0 & 12 & 11 & 6\\
D_{8} & 8 & 0 & 2 & 15 & 10 & 15\\
S_3 & 6 & 0 & 3 & 20 & 10 & 10\\
C_6 & 6 & 0 & 1 & 20 & 16 & 10\\
D_3\ (\subset A_5) & 6 & 0 & 0 & 20 & 19 & 10\\
C_5 & 5 & 4 & 0 & 24 & 21 & 6\\
V_4\ (\text{含换位}) & 4 & 0 & 2 & 30 & 19 & 15\\
V_4\ (\subset A_4) & 4 & 0 & 0 & 30 & 28 & 5\\
C_4 & 4 & 0 & 0 & 30 & 28 & 15\\
C_3 & 3 & 0 & 0 & 40 & 37 & 10\\
C_2\ (\text{双换位}) & 2 & 0 & 0 & 60 & 55 & 15\\
C_2\ (\text{换位}) & 2 & 0 & 1 & 60 & 46 & 10\\
\{1\} & 1 & 0 & 0 & 120 & 109 & 1
\end{array}
\]
\end{conclusion}
\begin{proof}[证明要点]
子群共轭类的枚举为有限群论事实。
在每一类上，\(|H|\) 与 \((a(H),b(H))\) 由循环型直接确定，属数由结论 \ref{con:xi-terminal-zm-delta-s5-intermediate-genus-hurwitz-burnside} 立即给出。证毕。
\end{proof}

\begin{conclusion}[分歧三分律：真子群商仅可能为纯 \(2\)-分歧、纯 \(5\)-分歧或全不分歧]\label{con:xi-terminal-zm-delta-s5-branching-trichotomy-proper-quotients}
对任意真子群 \(1\neq H\lneq S_5\)，记 \(\pi_H:X\to Y_H=X/H\)。
则：
\begin{itemize}
  \item 若 \(a(H)>0\)，则 \(b(H)=0\)，且 \(\pi_H\) 的全部分歧指数均为 \(5\)；
  \item 若 \(b(H)>0\)，则 \(a(H)=0\)，且 \(\pi_H\) 的全部分歧指数均为 \(2\)；
  \item 若 \(a(H)=b(H)=0\)，则 \(\pi_H\) 为全不分歧覆盖。
\end{itemize}
并且在分歧存在时，分歧点个数满足显式公式：
\[
r_5(H)=\frac{5}{4}\left(\frac{216}{|H|}-(2g(Y_H)-2)\right)\quad(a(H)>0),
\qquad
r_2(H)=2\left(\frac{216}{|H|}-(2g(Y_H)-2)\right)\quad(b(H)>0),
\]
其中 \(r_5(H)\)（resp.\ \(r_2(H)\)）为 \(\pi_H\) 的分歧点数，且每一点分歧指数恒为 \(5\)（resp.\ \(2\)）。
对结论 \ref{con:xi-terminal-zm-delta-s5-intermediate-genus-spectrum-19} 中出现的共轭类型，具体数值为
\[
r_5(A_5)=2,\quad r_5(F_{20})=1,\quad r_5(D_{10})=2,\quad r_5(C_5)=4,
\]
\[
r_2(S_4)=18,\quad r_2(S_3\times C_2)=24,\quad r_2(D_8)=18,\quad r_2(S_3)=36,\quad r_2(C_6)=12,\quad r_2(V_4\ (\text{含换位}))=36,\quad r_2(C_2\ (\text{换位}))=36.
\]
\end{conclusion}
\begin{proof}[证明要点]
在 \(S_5\) 中，一个 \(5\)-循环与一个换位生成 \(S_5\)，故真子群不可能同时含有换位与 \(5\)-循环，得到互斥性 \(a(H)b(H)=0\)。
\par
由结论 \ref{con:xi-terminal-zm-delta-s5-regular-closure-genus-109}，\(X\to\PP^1_\delta\) 的点稳定子仅可能为 \(C_5\) 或 \(C_2\)。
因此 \(\pi_H\) 的惯性群只能来自 \(H\) 与这些稳定子的共轭交；在真子群情形下，互斥性迫使交要么恒为平凡、要么恒为 \(C_5\)、要么恒为 \(C_2\)，从而得到三分律。
\par
当 \(\pi_H\) 为 \(H\)-伽罗瓦覆盖且分歧指数恒为 \(e\in\{5,2\}\) 时，Hurwitz 公式给出
\[
2g(X)-2=|H|\,(2g(Y_H)-2)+r\cdot |H|\left(1-\frac1e\right),
\]
其中 \(r\) 为分歧点数。
代入 \(2g(X)-2=216\) 并解出 \(r\) 即得所示 \(r_5(H),r_2(H)\)。
将结论 \ref{con:xi-terminal-zm-delta-s5-intermediate-genus-spectrum-19} 中的 \((|H|,g(Y_H))\) 逐一代入得到具体数值。证毕。
\end{proof}

\begin{conclusion}[可割度的精确值]\label{con:xi-terminal-zm-delta-s5-gonality-10}
设 \(\operatorname{gon}(X)\) 为 \(X\) 在代数闭包上的几何可割度。
则
\[
\operatorname{gon}(X)=10.
\]
\end{conclusion}
\begin{proof}[证明要点]
由结论 \ref{con:xi-terminal-zm-delta-s5-intermediate-genus-spectrum-19}，存在指数 \(5\) 的子群 \(S_4\le S_5\) 使得 \(C_{S_4}:=X/S_4\) 为属 \(1\) 曲线，且自然映射 \(\pi:X\to C_{S_4}\) 的次数为 \(5\)。
在代数闭包上任取 \(P\in C_{S_4}\)，线性系 \(|2P|\) 给出次数 \(2\) 的态射 \(C_{S_4}\to\PP^1\)，从而得到次数 \(10\) 的复合态射 \(X\to\PP^1\)，故 \(\operatorname{gon}(X)\le 10\)。
\par
若存在次数 \(d\le 9\) 的态射 \(f:X\to\PP^1\)，则 \(f\) 不可能经由 \(\pi\) 分解。
此时由 Castelnuovo--Severi 不等式（对 \(\pi\) 与 \(f\)）得到
\[
g(X)\le d\cdot g(C_{S_4})+5\cdot g(\PP^1)+(d-1)(5-1)=5d-4,
\]
与 \(g(X)=109\) 矛盾。
故不存在 \(d\le 9\) 的态射，从而 \(\operatorname{gon}(X)\ge 10\)。
综上 \(\operatorname{gon}(X)=10\)。证毕。
\end{proof}

\begin{conclusion}[次数阈值：\(\deg<23\) 的 \texorpdfstring{$\PP^1$}{P1} 态射必经由属 \(1\) 商]\label{con:xi-terminal-zm-delta-s5-degree-threshold-23}
令 \(C_{S_4}:=X/S_4\) 与 \(\pi:X\to C_{S_4}\) 如结论 \ref{con:xi-terminal-zm-delta-s5-gonality-10}。
对任意非恒等态射 \(f:X\to\PP^1\)，若 \(\deg f<23\)，则必存在态射 \(g:C_{S_4}\to\PP^1\) 使得 \(f=g\circ\pi\)。
因此在 \(1\le d\le 22\) 范围内，存在次数为 \(d\) 的态射 \(X\to\PP^1\) 当且仅当
\[
d\in\{10,15,20\}.
\]
\end{conclusion}
\begin{proof}[证明要点]
若 \(f\) 不经由 \(\pi\) 分解，则与结论 \ref{con:xi-terminal-zm-delta-s5-gonality-10} 同理，对 \(\pi\) 与 \(f\) 应用 Castelnuovo--Severi 不等式得到
\(
g(X)\le 5\deg(f)-4
\)，代入 \(g(X)=109\) 得 \(\deg(f)\ge 23\)。
故当 \(\deg f<23\) 时必有分解 \(f=g\circ\pi\)。
\par
一旦 \(f=g\circ\pi\)，则 \(\deg f=5\deg g\)。
属 \(1\) 曲线到 \(\PP^1\) 的非恒等态射次数不小于 \(2\)，并且在代数闭包上对任意 \(m\ge 2\) 均存在次数 \(m\) 的态射 \(C_{S_4}\to\PP^1\)（取任意次数 \(m\) 线丛并选取二维线性系）。
因此 \(5m\le 22\) 的可能性仅为 \(m=2,3,4\)，即 \(d\in\{10,15,20\}\)，并且这些次数均可实现。证毕。
\end{proof}

\begin{conclusion}[完全不分歧的可解塔]\label{con:xi-terminal-zm-delta-s5-etale-solvable-tower}
存在子群链
\[
C_2\ (\text{双换位})\ \trianglelefteq\ V_4\ (\subset A_4)\ \trianglelefteq\ A_4\le S_5,
\]
使得相应商曲线
\[
X\longrightarrow X/C_2\longrightarrow X/V_4\longrightarrow X/A_4
\]
构成完全不分歧的伽罗瓦覆盖塔，覆盖次数依次为 \(2,2,3\)，并且属数分别为
\[
g(X/C_2)=55,\qquad g(X/V_4)=28,\qquad g(X/A_4)=10.
\]
\end{conclusion}
\begin{proof}[证明要点]
对所列三类子群，结论 \ref{con:xi-terminal-zm-delta-s5-intermediate-genus-spectrum-19} 给出 \(a(H)=b(H)=0\)。
由结论 \ref{con:xi-terminal-zm-delta-s5-branching-trichotomy-proper-quotients}，相应商映射 \(X\to X/H\) 均为全不分歧覆盖。
覆盖次数由子群指数给出，属数由结论 \ref{con:xi-terminal-zm-delta-s5-intermediate-genus-spectrum-19} 直接读出。证毕。
\end{proof}

\begin{conclusion}[全不分歧二次商的 Jacobian 与 Prym 的同源分解]\label{con:xi-terminal-zm-delta-s5-c2-quotient-jacobian-prym-decomposition}
令 \(C_2\trianglelefteq V_4\) 为结论 \ref{con:xi-terminal-zm-delta-s5-etale-solvable-tower} 中的双换位子群，并置
\[
X_{55}:=X/C_2,\qquad J_{55}:=\Jac(X_{55}),\qquad P_{54}:=\Prym(X\to X_{55}).
\]
则在结论 \ref{con:xi-terminal-zm-delta-s5-jacobian-power-isogeny-refinement} 的 \(E,B_3,B_5,B_6,B_7\) 记号下，有 \(\QQ\)-同源分解
\[
J_{55}\ \sim_{\QQ}\ A_\varepsilon\ \times\ E^{\,2}\ \times\ B_3^{\,3}\ \times\ B_5^{\,2}\ \times\ B_6^{\,3}\ \times\ B_7^{\,2},
\]
\[
P_{54}\ \sim_{\QQ}\ E^{\,2}\ \times\ B_3^{\,2}\ \times\ B_5^{\,4}\ \times\ B_6^{\,2}\ \times\ B_7^{\,2}.
\]
并且 \(P_{54}\) 带自然主极化。
\end{conclusion}
\begin{proof}[证明要点]
由结论 \ref{con:xi-terminal-zm-delta-s5-etale-solvable-tower}，覆盖 \(X\to X_{55}\) 为全不分歧二次覆叠。
故 \(\pi^*:J_{55}\to J(X)\) 的像为 \(J(X)^{C_2}\) 的连通分量，且在 \(\QQ\)-同源意义下
\[
J(X)\ \sim_{\QQ}\ J(X)^{C_2}\times P_{54}.
\]
对任意不可约表示 \(V\)（特征标 \(\chi_V\)），由特征平均得
\[
\dim(V^{C_2})=\frac{\chi_V(1)+\chi_V((2,2))}{2}.
\]
而由结论 \ref{con:xi-terminal-zm-delta-z28-jacobian-qisogeny-1-2-7-6-12} 的 \(V_4\)-不变维数（其中 \(V_4\) 的非平凡元素皆为双换位）可解得
\[
\chi_\varepsilon((2,2))=1,\quad
\chi_{(4,1)}((2,2))=0,\quad
\chi_{(3,2)}((2,2))=1,\quad
\chi_{(3,1,1)}((2,2))=-2,\quad
\chi_{(2,2,1)}((2,2))=1,\quad
\chi_{(2,1,1,1)}((2,2))=0,
\]
从而
\[
\dim(V^{C_2})=1,2,3,2,3,2
\]
依次对应于
\(\varepsilon,V_{(4,1)},V_{(3,2)},V_{(3,1,1)},V_{(2,2,1)},V_{(2,1,1,1)}\)。
将该不变维数代入结论 \ref{con:xi-terminal-zm-delta-s5-jacobian-power-isogeny-refinement} 的幂次型分解，即得 \(J(X)^{C_2}\sim_{\QQ}J_{55}\) 的分解；反不变部分幂指数为 \(\dim V-\dim(V^{C_2})\)，从而得到 \(P_{54}\) 的分解。
全不分歧二次覆叠的 Prym 携带自然主极化为经典定理。证毕。
\end{proof}

\input{prop__xi-terminal-zm-delta-s5-tame-collision-etale-quotient-node-law}

\begin{conclusion}[素数 \texorpdfstring{$p=37$}{p=37} 处的半稳定约化与 \texorpdfstring{$j\equiv 16$}{j=16}]\label{con:xi-terminal-zm-delta-a5-curve-semistable-reduction-p37-j16}
令 \(C_{A_5}\) 如结论 \ref{con:xi-terminal-zm-delta-a5-fixedfield-hyperelliptic-genus2}，即
\(
C_{A_5}:\ s^2=(\delta-4)B_5(\delta)
\)。
则在素数 \(p=37\) 处有同余分解
\[
B_5(\delta)\equiv 13\,\delta^2(\delta^3+2\delta^2-4\delta+3)\pmod{37}
\]
（命题 \ref{prop:xi-terminal-zm-delta-badprime-37-b5-double-root}）。
因此 \(C_{A_5}\) 的特殊纤维在 \((\delta,s)=(0,0)\) 处出现唯一普通双点；其归一化为 \(\FF_{37}\) 上的椭圆曲线
\[
E_{37}:\quad v^2=(\delta-4)(\delta^3+2\delta^2-4\delta+3),
\qquad v:=s/\delta,
\]
并且其 \(j\)-不变量满足
\[
j(E_{37})\equiv 16\pmod{37}.
\]
\end{conclusion}
\begin{proof}[证明要点]
由同余分解可写
\[
s^2\equiv 13\,\delta^2(\delta-4)(\delta^3+2\delta^2-4\delta+3)\pmod{37},
\]
从而 \(\delta=0\) 为二重分支点并导致 \((0,0)\) 处节点。
作归一化，令 \(v=s/\delta\) 即得 \(E_{37}\) 的四次模型。
\par
记四次多项式
\[
f(\delta):=(\delta-4)(\delta^3+2\delta^2-4\delta+3)=\delta^4-2\delta^3-12\delta^2+19\delta-12.
\]
对四次模型 \(v^2=f(\delta)\) 的二元四次不变量，取系数 \(a=1,b=-2,c=-12,d=19,e=-12\)，则
\[
I:=12ae-3bd+c^2\equiv 3\pmod{37},\qquad
J:=72ace+9bcd-27ad^2-27b^2e-2c^3\equiv 5\pmod{37},
\]
故
\(
\Delta:=4I^3-J^2\equiv 9\not\equiv 0
\)
并且
\[
j=1728\cdot\frac{4I^3}{\Delta}\equiv 1728\cdot\frac{108}{9}\equiv 16\pmod{37}.
\]
证毕。
\end{proof}

\begin{conclusion}[\texorpdfstring{$p=37$}{p=37} 处 \texorpdfstring{$\Jac(C_{A_5})$}{Jac(CA5)} 的半稳定约化：环面秩与椭圆归一化核]\label{con:xi-terminal-zm-delta-jac-ca5-neron-model-torus-rank1}
令 \(J_{A_5}:=\Jac(C_{A_5})\)，并令 \(\mathcal J\) 表示其在 \(\ZZ_{37}\) 上的 N\'eron 模型，\(\mathcal J^0\) 为单位连通分量。
则 \(J_{A_5}\) 在 \(p=37\) 处为半稳定但非良约化，并且存在半阿贝尔精确列
\[
0\ \longrightarrow\ \mathbf G_m\ \longrightarrow\ \mathcal J^0_{\FF_{37}}\ \longrightarrow\ E_{37}\ \longrightarrow\ 0,
\]
其中 \(E_{37}\) 为结论 \ref{con:xi-terminal-zm-delta-a5-curve-semistable-reduction-p37-j16} 给出的归一化椭圆曲线。
特别地，该约化的环面秩恰为 \(1\)。
\end{conclusion}
\begin{proof}[证明要点]
由结论 \ref{con:xi-terminal-zm-delta-a5-curve-semistable-reduction-p37-j16}，\(C_{A_5}\) 在 \(p=37\) 的特殊纤维不可约且恰含一个普通双点，其归一化为椭圆曲线 \(E_{37}\)。
对离散赋值环上的稳定曲线，Raynaud 的广义 Jacobian 理论给出：若特殊纤维含 \(1\) 个节点，则广义 Jacobian 的单位连通分量为归一化 Jacobian 对 \(\mathbf G_m\) 的一次扩张，且环面秩等于节点数。
据此得到所示精确列与环面秩 \(1\)，从而约化半稳定且非良。证毕。
\end{proof}

\input{parts/prop__xi-terminal-zm-delta-ca5-node-splitting-symbol}

\input{thm__xi-terminal-zm-delta-ca5-semistable-local-factors}

\endinput

\begin{theorem}[节点坐标域上两条单生成整序的判别式与指数比值]\label{thm:xi-terminal-zm-delta-node-order-discriminant-index-ratio}
取 \(\overline{\mathcal C}_\delta\) 的任一有限节点 \((y_0,T_0)\)。
令
\[
K_0:=\QQ(y_0)=\QQ(T_0),
\qquad
\mathcal O_y:=\ZZ[y_0]\subset K_0,
\qquad
\mathcal O_T:=\ZZ[T_0]\subset K_0,
\]
并令 \(\mathcal O_{K_0}\) 为 \(K_0\) 的整数环。
则有严格恒等式
\[
\frac{[\mathcal O_{K_0}:\mathcal O_T]}{[\mathcal O_{K_0}:\mathcal O_y]}
=\frac{2^{7}\cdot 3\cdot 307}{31},
\qquad
\frac{\Disc(\mathcal O_T)}{\Disc(\mathcal O_y)}
=\left(\frac{2^{7}\cdot 3\cdot 307}{31}\right)^2.
\]
等价地，对任意素数 \(\ell\) 有逐素数赋值差分
\[
v_\ell([\mathcal O_{K_0}:\mathcal O_T]) - v_\ell([\mathcal O_{K_0}:\mathcal O_y])
=
v_\ell\!\left(\frac{2^{7}\cdot 3\cdot 307}{31}\right),
\]
从而
\[
v_2(\cdot)=7,\qquad v_3(\cdot)=1,\qquad v_{307}(\cdot)=1,\qquad v_{31}(\cdot)=-1,
\]
并且在 \(\ell\in\{11,727\}\) 处差分为 \(0\)。
\end{theorem}
\begin{proof}
由命题 \ref{prop:xi-terminal-zm-delta-nodes-elimination-a5}，节点坐标满足 \(\QQ(y_0)=\QQ(T_0)\)。
由于 \(Q_5(y)\) 与 \(A_5(T)\) 均为首一整系数多项式，\(y_0,T_0\) 均为代数整数，故 \(\mathcal O_y,\mathcal O_T\) 为 \(K_0\) 的整序。
\par
由结论 \ref{con:xi-terminal-zm-leyang-q5-discriminant-727} 与结论 \ref{con:xi-terminal-zm-delta-a5-discriminant-quadratic}，
\[
\Disc(\mathcal O_y)=\Disc(Q_5)
=-2^{28}\cdot 3^{9}\cdot 11^{2}\cdot 31^{3}\cdot 727,
\]
\[
\Disc(\mathcal O_T)=\Disc(A_5)
=-2^{42}\cdot 3^{11}\cdot 11^{2}\cdot 31\cdot 307^{2}\cdot 727.
\]
对同一数域 \(K_0\) 的任意整序 \(\mathcal O\subseteq \mathcal O_{K_0}\)，判别式满足
\(
\Disc(\mathcal O)=\Disc(\mathcal O_{K_0})\,[\mathcal O_{K_0}:\mathcal O]^2
\)。
将其分别应用于 \(\mathcal O_y,\mathcal O_T\) 并取比值得
\[
\frac{\Disc(\mathcal O_T)}{\Disc(\mathcal O_y)}
=\left(\frac{[\mathcal O_{K_0}:\mathcal O_T]}{[\mathcal O_{K_0}:\mathcal O_y]}\right)^2
=\frac{2^{42}\cdot 3^{11}\cdot 11^{2}\cdot 31\cdot 307^{2}\cdot 727}{2^{28}\cdot 3^{9}\cdot 11^{2}\cdot 31^{3}\cdot 727}
=\left(\frac{2^{7}\cdot 3\cdot 307}{31}\right)^2.
\]
取正平方根得到指数比值恒等式。
逐素数差分由对两侧取 \(v_\ell\) 并除以 \(2\) 得出。证毕。
\end{proof}
\leanverified{paper\_xi\_terminal\_zm\_delta\_node\_order\_discriminant\_index\_ratio}

\endinput


\begin{conclusion}[二级退化五次的全局算术指纹：新的平方自由敏感素数 \texorpdfstring{$727$}{727}]\label{con:xi-terminal-zm-leyang-q5-discriminant-727}
令 \(Q_5(y)\) 如结论 \ref{con:xi-terminal-zm-leyang-delta-quartic-discriminant-reflexive}。
则其判别式满足严格素因子分解
\[
\boxed{\;
\mathrm{Disc}\bigl(Q_5\bigr)=-2^{28}\cdot 3^{9}\cdot 11^{2}\cdot 31^{3}\cdot 727\; }.
\]
因此除去结论 \ref{con:xi-terminal-zm-leyang-branchpolynomial-discriminant-modp-collisions} 所限定的一级分歧素集 \(\{2,3,31,37\}\) 之外，
二级退化集合还引入新的平方自由敏感素数 \(727\)（并且 \(11\) 仅以平方因子出现）；
在该素数特征下，\(Q_5(y)=0\) 的约化不可避免地产生额外的并合/不可分现象，从而提供一条与 Lee--Yang 一级核正交的更深层算术奇异源。
\end{conclusion}
\begin{proof}[证明要点]
对 \(Q_5\) 作判别式计算并在 \(\ZZ\) 中分解即可。
可复现核验脚本：\texttt{scripts/exp\_fold\_zm\_elliptic\_leyang\_cover\_geometry\_audit.py}。证毕。
\end{proof}

\begin{conclusion}[模性--微分闭合：\texorpdfstring{$D(\tau)=\frac{q\,\partial_q y(\tau)}{f(q)}$}{D(tau)=q y'/f} 满足同一四次方程]\label{con:xi-terminal-zm-leyang-modular-differential-closure}
沿用注记 \ref{rem:xi-terminal-elliptic-37a1-modular-anchor} 的模参数化：取 \(q=e^{2\pi i\tau}\)，
并令 \(f(\tau)=\sum_{n\ge 1}a_n q^n\) 为 level \(37\) 的规范化新形式，使得 \(f(q)\,\dd q/q\) 等于椭圆曲线 \(E\) 的 N\'eron 微分 \(\omega=\dd X/(2Y)\) 在同一参数化下的拉回；
令 \(y(\tau)\) 为该参数化下的重量函数。
定义
\[
D(\tau):=\frac{q\,\frac{\dd}{\dd q}y(\tau)}{f(q)}\in \CC((q)).
\]
则 \((y(\tau),D(\tau))\) 满足完全闭合的一阶代数微分方程
\[
\boxed{\;
D^{4}+(8y-3)D^{3}+(2y^{2}-34y-4)D^{2}
-y(y-1)\,P_{\mathrm{LY}}(y)=0\;},
\qquad
P_{\mathrm{LY}}(y)=256y^{3}+411y^{2}+165y+32,
\]
其中 \(y=y(\tau)\)、\(D=D(\tau)\)。
\end{conclusion}
\begin{proof}[证明要点]
由定义 \(D(\tau)=\frac{q\,(\dd/\dd q)y(\tau)}{f(q)}\) 与 \(\dd u=\omega=f(q)\,\dd q/q\) 得 \(D=\dd y/\dd u=\delta\)。
代入结论 \ref{con:xi-terminal-zm-leyang-abel-jacobi-ode} 的四次闭式即得。证毕。
\end{proof}

\begin{conclusion}[Lee--Yang Perron 分支的中心化三阶 Fuchs 线性算子]\label{con:xi-terminal-zm-leyang-perron-centered-fuchs-operator}
令
\[
\Pi(\lambda,y):=\lambda^{4}-\lambda^{3}-(2y+1)\lambda^{2}+\lambda+y(y+1),\qquad
P_{\mathrm{LY}}(y):=256y^{3}+411y^{2}+165y+32,
\]
并记 \(\lambda_{+}(y)\) 为 \(y=1\) 邻域内满足 \(\lambda_{+}(1)=2\) 的解析分支（结论 \ref{con:xi-terminal-zm-lln-clt}）。
定义中心化分支
\[
f(y):=\lambda_{+}(y)-\frac14.
\]
则 \(f\) 满足三阶线性微分方程
\[
A_{3}(y)f^{(3)}+A_{2}(y)f^{(2)}+A_{1}(y)f'+A_{0}(y)f=0,
\]
其中 \(A_{j}(y)\in\ZZ[y]\) 显式如下：
% Auto-generated by scripts/exp_fold_zm_leyang_perron_fuchs_operator_audit.py
\[
\Pi(\lambda,y)=\lambda^{4}-\lambda^{3}-(2y+1)\lambda^{2}+\lambda+y(y+1),\qquad P_{\mathrm{LY}}(y)=256y^{3}+411y^{2}+165y+32.
\]
\[
f(y):=\lambda_{+}(y)-\frac14,\qquad A_{3}(y)f^{(3)}+A_{2}(y)f^{(2)}+A_{1}(y)f'+A_{0}(y)f=0.
\]
\[
A_{3}(y)=y(y-1)P_{\mathrm{LY}}(y)Q(y),\qquad Q(y)=168y^{5}+1260y^{4}-422y^{3}-183y^{2}-639y-76.
\]
\[
\begin{aligned}A_{2}(y)&=6\Big(17920y^{9}+192500y^{8}-12480y^{7}-94031y^{6}-139107y^{5}-71300y^{4}+80597y^{3}+41863y^{2}+6758y+608\Big),\\A_{1}(y)&=6\Big(3136y^{8}+45164y^{7}-56178y^{6}-36438y^{5}-138724y^{4}-62488y^{3}+18443y^{2}+5933y-32\Big),\\A_{0}(y)&=24\Big(56y^{7}+1036y^{6}+3844y^{5}+3740y^{4}+8288y^{3}+1417y^{2}+19y+32\Big).\end{aligned}
\]
\[
\operatorname{Sing}_{\mathrm{fin}}(L_{3})=\{A_{3}(y)=0\}=\{0,1\}\cup\{P_{\mathrm{LY}}(y)=0\}\cup\{Q(y)=0\}.
\]
\[
\operatorname{Exp}_{y_{0}}(L_{3})=\left\{0,1,2-\frac{A_{2}(y_{0})}{A_{3}'(y_{0})}\right\},
\]
\[
\operatorname{Exp}_{y_{0}}(L_{3})=\left\{0,\frac12,1\right\}\ (y_{0}\in\{0,1,P_{\mathrm{LY}}=0\}),\qquad \operatorname{Exp}_{y_{0}}(L_{3})=\{0,1,3\}\ (y_{0}\in\{Q=0\}).
\]
\[
\operatorname{Exp}_{\infty}(L_{3})=\left\{-\frac14,\frac14,\frac12\right\},\qquad 32\sigma^{3}-16\sigma^{2}-2\sigma+1=(2\sigma-1)(4\sigma-1)(4\sigma+1).
\]
\[
W(y)^{2}=C\cdot\frac{Q(y)^{2}}{y^{3}(y-1)^{3}P_{\mathrm{LY}}(y)^{3}},\qquad C\in\mathbb{C}^{\times}.
\]
\[
\Lambda(t)=2+\frac{5}{9}t-\frac{64}{243}t^{2}+\frac{2404}{6561}t^{3}-\frac{159436}{177147}t^{4}+\frac{15226688}{4782969}t^{5}+O(t^{6}).
\]
\[
\Lambda(t)\in \mathbb{Z}\!\left[\frac13\right][[t]],\qquad \kappa_{n}^{\infty}:=\left.\frac{d^{n}}{ds^{n}}\right|_{s=0}\log\frac{\Lambda(e^{s}-1)}{2}\in \mathbb{Z}\!\left[\frac16\right]\ \ (n\ge 1).
\]
\[
(p\neq 3)\ \Longrightarrow\ \overline{\Lambda}_{p}(t)\in\mathbb{F}_{p}[[t]]\ \text{is algebraic over}\ \mathbb{F}_{p}(t),\ \text{hence its coefficient sequence is \(p\)-automatic (Christol).}
\]
\[
\operatorname{Mon}(L_{3})\cong S_{4}\ \text{on}\ \{(x_{1},x_{2},x_{3},x_{4})\in\mathbb{C}^{4}:\ x_{1}+x_{2}+x_{3}+x_{4}=0\},\ \dim=3.
\]

\end{conclusion}
\begin{proof}[证明要点]
设
\(
K:=\QQ(y)[\lambda]/(\Pi(\lambda,y))
\)。
由 Vieta 关系 \(\sum_{i=1}^{4}\lambda_i=1\) 得
\(
\mathrm{Tr}_{K/\QQ(y)}(\lambda-\tfrac14)=0
\)，故
\(
g:=\lambda-\tfrac14\in K_0:=\ker\mathrm{Tr}
\)。
又 \(\dim_{\QQ(y)}K_0=3\)。
令 \(\partial:=\frac{d}{dy}\)；
可分扩张下迹与导数可交换，
\(
\mathrm{Tr}(\partial u)=\partial\mathrm{Tr}(u)
\)，故 \(\partial(K_0)\subseteq K_0\)。
因此 \(g,\partial g,\partial^2 g,\partial^3 g\) 在 \(K_0\) 中线性相关，存在 \(a_j(y)\in\QQ(y)\) 使
\[
\partial^3 g+a_2(y)\partial^2 g+a_1(y)\partial g+a_0(y)g=0.
\]
将该恒等式在 \(K\) 中按基 \(\{1,\lambda,\lambda^2,\lambda^3\}\) 消元并清分母，得到整系数多项式 \(A_3,A_2,A_1,A_0\)。
脚本 \texttt{scripts/exp\_fold\_zm\_leyang\_perron\_fuchs\_operator\_audit.py} 对上述恒等式给出代数核验。证毕。
\end{proof}

\begin{conclusion}[Riemann 方案中的真奇点与表观奇点分离]\label{con:xi-terminal-zm-leyang-perron-fuchs-riemann-apparent-split}
设 \(L_3\) 为结论 \ref{con:xi-terminal-zm-leyang-perron-centered-fuchs-operator} 的三阶算子。
其有限奇点恰为 \(A_3(y)=0\)：
\[
\{0,1\}\cup\{P_{\mathrm{LY}}(y)=0\}\cup\{Q(y)=0\}.
\]
其中 \(\{0,1\}\cup\{P_{\mathrm{LY}}=0\}\) 为 \(\Pi\) 的分歧点集；\(\{Q=0\}\) 为五个有限表观奇点。
若 \(y_0\) 为 \(A_3\) 的简单零点，则局部指标为
\[
\operatorname{Exp}_{y_0}(L_3)=\left\{0,1,2-\frac{A_2(y_0)}{A_3'(y_0)}\right\}.
\]
并且
\[
\operatorname{Exp}_{y_0}(L_3)=\left\{0,\frac12,1\right\}
\quad\bigl(y_0\in\{0,1,P_{\mathrm{LY}}=0\}\bigr),
\]
\[
\operatorname{Exp}_{y_0}(L_3)=\{0,1,3\}
\quad\bigl(y_0\in\{Q=0\}\bigr).
\]
在 \(y=\infty\) 处，指标集合为
\[
\left\{-\frac14,\frac14,\frac12\right\},
\]
其指标多项式为
\[
32\sigma^{3}-16\sigma^{2}-2\sigma+1=(2\sigma-1)(4\sigma-1)(4\sigma+1).
\]
\end{conclusion}
\begin{proof}[证明要点]
将方程写成首一形式 \(f^{(3)}+p_2f^{(2)}+p_1f'+p_0f=0\)。
在 \(A_3\) 的简单零点处 \(p_2\) 仅有一阶极点、\(p_1,p_0\) 不产生更高阶极部，故为正则奇点，
并得到指标公式
\(
r(r-1)\bigl(r-2+p_{2,-1}\bigr)=0
\)；
即
\(
\{0,1,2-A_2(y_0)/A_3'(y_0)\}
\)。
由恒等式
\[
2A_2-3A_3'\equiv 0\pmod{y(y-1)P_{\mathrm{LY}}},\qquad
A_2+A_3'\equiv 0\pmod{Q}
\]
分别得到上述两组局部指数。
又 \(\gcd\!\bigl(Q,\ y(y-1)P_{\mathrm{LY}}\bigr)=1\)，故 \(Q\) 的零点不在 \(\Pi\) 的分歧集上；
四个代数分支在这些点局部单值解析，故该五点为表观奇点。
无穷远指标由主平衡项直接给出。证毕。
\end{proof}

\begin{conclusion}[Wronskian 的平方闭式与表观奇点结构方程]\label{con:xi-terminal-zm-leyang-perron-fuchs-wronskian-square}
设 \(V\) 为算子 \(L_3\) 的解空间，\(\dim V=3\)。
取任意基 \((u_1,u_2,u_3)\)，Wronskian 记为
\[
W(y):=\det\bigl(u_i^{(j-1)}(y)\bigr)_{1\le i,j\le 3}.
\]
则存在常数 \(C\neq 0\) 使
\[
W(y)^2
=
C\cdot\frac{Q(y)^2}{y^3(y-1)^3P_{\mathrm{LY}}(y)^3}.
\]
因此 \(Q(y)=0\) 给出 \(W\) 的有限零点，而 \(P_{\mathrm{LY}}(y)=0\) 与 \(y=0,1\) 给出 \(W\) 的有限极点。
\end{conclusion}
\begin{proof}[证明要点]
对首一三阶方程有经典恒等式
\(
\frac{W'}{W}=-p_2
\)。
从而
\(
\frac{(W^2)'}{W^2}=-2p_2=-2A_2/A_3
\)。
对
\(
R(y):=\frac{Q(y)^2}{y^3(y-1)^3P_{\mathrm{LY}}(y)^3}
\)
直接计算得
\(
\frac{R'}{R}=-2A_2/A_3
\)。
故 \(W^2=CR\)。
脚本 \texttt{scripts/exp\_fold\_zm\_leyang\_perron\_fuchs\_operator\_audit.py} 已给出符号核验。证毕。
\end{proof}

\begin{conclusion}[Perron 分支 \texorpdfstring{$\Lambda(t)$}{Lambda(t)} 的全局 \(p\)-整性与仅 \(3\)-进分母]\label{con:xi-terminal-zm-leyang-perron-global-pintegrality}
令 \(t:=y-1\)，并令 \(\Lambda(t)\in\QQ[[t]]\) 满足
\[
\Pi(\Lambda(t),1+t)=0,\qquad \Lambda(0)=2.
\]
则对任意素数 \(p\neq 3\)，有
\[
\Lambda(t)\in\ZZ_p[[t]].
\]
因此 \(\Lambda(t)=2+\sum_{n\ge 1}a_n t^n\) 的每个系数满足 \(a_n\in\ZZ[1/3]\)。
\end{conclusion}
\begin{proof}[证明要点]
设 \(F(\lambda,t):=\Pi(\lambda,1+t)\in\ZZ[\lambda,t]\)。
有 \(F(2,0)=0\) 且
\(
F_{\lambda}(2,0)=\partial_{\lambda}\Pi(2,1)=9
\)。
当 \(p\neq 3\) 时 \(9\in\ZZ_p^\times\)，由 \(p\)-进隐函数定理得到唯一
\(
\Lambda_p(t)\in\ZZ_p[[t]]
\)
满足 \(F(\Lambda_p(t),t)=0\)、\(\Lambda_p(0)=2\)。
而该初值下 \(\QQ[[t]]\) 形式解唯一，故 \(\Lambda=\Lambda_p\) 视为 \(\QQ_p[[t]]\) 元成立。
于是对所有 \(p\neq 3\) 的 \(p\)-整性同时成立，分母仅可能含素因子 \(3\)。证毕。
\end{proof}

\begin{theorem}[\texorpdfstring{$p=3$}{p=3} 的 \(27\)-尺度良约化与 \(3\)-进 Hensel 可解性]\label{thm:xi-terminal-zm-leyang-perron-p3-renormalization-hensel}
\leanverified{paper\_xi\_terminal\_zm\_leyang\_perron\_p3\_renormalization\_hensel}
在结论 \ref{con:xi-terminal-zm-leyang-perron-global-pintegrality} 的设定下，素数 \(p=3\) 为唯一使 \(F_\lambda(2,0)\notin\ZZ_p^\times\) 的例外。
令
\[
t=27u,\qquad \lambda=2+3v.
\]
则恒有分解
\[
\Pi(2+3v,1+27u)=27\,H(u,v),
\]
其中
\[
H(u,v)=27u^{2}-18uv^{2}-24uv-5u+3v^{4}+7v^{3}+5v^{2}+v\in\ZZ[u,v].
\]
并且
\[
\partial_v H(0,0)=1\in\ZZ_3^\times.
\]
因此存在唯一的 \(v(u)\in\ZZ_3[[u]]\) 满足
\[
H(u,v(u))=0,\qquad v(0)=0,
\]
从而
\[
\Lambda(27u)=2+3v(u)\in 2+3\,\ZZ_3[[u]].
\]
\end{theorem}
\begin{proof}
把 \(\lambda=2+3v\)、\(t=27u\) 代入 \(F(\lambda,t)=\Pi(\lambda,1+t)\) 并在 \(\ZZ[u,v]\) 中展开整理，即得所示恒等式与多项式 \(H(u,v)\)。
对 \(H\) 求偏导并在 \((u,v)=(0,0)\) 处取值可得 \(\partial_vH(0,0)=1\)，故在完备局部环 \(\ZZ_3[[u]]\) 中为单位。
由形式隐函数定理（等价于对 \(\ZZ_3[[u]]\) 的 Hensel 型提升），存在且仅存在唯一的 \(v(u)\in\ZZ_3[[u]]\) 解 \(H(u,v)=0\) 且满足 \(v(0)=0\)。
由 \(\Pi(\Lambda(t),1+t)=0\) 的形式解唯一性知 \(t=27u\) 的代入与上述解一致，即 \(\Lambda(27u)=2+3v(u)\)。
可复现核验脚本：\texttt{scripts/exp\_fold\_zm\_leyang\_perron\_p3\_renormalization\_audit.py}。证毕。
\end{proof}

\begin{proposition}[\texorpdfstring{$27$}{27} 的最小性：\texorpdfstring{$t=9u$}{t=9u} 不产生 \texorpdfstring{$3$}{3}-进整系数展开]\label{prop:xi-terminal-zm-leyang-perron-p3-minimal-27}
\leanverified{paper\_xi\_terminal\_zm\_leyang\_perron\_p3\_minimal\_27}
不存在 \(\widetilde\Lambda(u)\in\ZZ_3[[u]]\) 使得
\[
\Pi(\widetilde\Lambda(u),1+9u)=0,\qquad \widetilde\Lambda(0)=2.
\]
等价地，\(\Lambda(9u)\notin\ZZ_3[[u]]\)。
\end{proposition}
\begin{proof}
在 \(\FF_3\) 上约化。
由于 \(9u\equiv 0\pmod 3\)，有
\[
\Pi(\lambda,1+9u)\equiv \Pi(\lambda,1)\pmod 3.
\]
又
\[
\Pi(\lambda,1)=\lambda^4-\lambda^3-3\lambda^2+\lambda+2\equiv \lambda^4-\lambda^3+\lambda-1
=(\lambda-1)(\lambda+1)^3\quad(\bmod\,3).
\]
若存在 \(\widetilde\Lambda(u)\in\ZZ_3[[u]]\) 满足命题条件，则其模 \(3\) 约化 \(\overline{\widetilde\Lambda}(u)\in\FF_3[[u]]\) 满足
\((\overline{\widetilde\Lambda}(u)-1)(\overline{\widetilde\Lambda}(u)+1)^3=0\)。
由于 \(\FF_3[[u]]\) 为整环，推出 \(\overline{\widetilde\Lambda}(u)=\pm 1\)；
而 \(\widetilde\Lambda(0)=2\equiv-1\pmod 3\)，故 \(\overline{\widetilde\Lambda}(u)\equiv-1\)，从而可写
\[
\widetilde\Lambda(u)=2+3v(u),\qquad v(u)\in\ZZ_3[[u]].
\]
将 \(\lambda=2+3v\)、\(t=9u\) 代入并在 \(\ZZ[u,v]\) 中展开得
\[
\Pi(2+3v,1+9u)=9\,H_9(u,v),
\]
其中
\[
H_9(u,v)=9u^{2}-18uv^{2}-24uv-5u+9v^{4}+21v^{3}+15v^{2}+3v\in\ZZ[u,v].
\]
模 \(3\) 约化时含 \(v\) 的各项系数均被 \(3\) 整除而消失，且 \(-5u\equiv u\pmod 3\)，故
\[
H_9(u,v)\equiv u\quad(\bmod\,3).
\]
于是 \(H_9(u,v(u))=0\) 在 \(\FF_3[[u]]\) 中强制 \(u=0\)，与 \(u\) 为形式变量矛盾。证毕。
\end{proof}

\begin{corollary}[Perron 分支系数的 \(3\)-进估值线性下界]\label{cor:xi-terminal-zm-leyang-perron-v3-linear-bound}
\leanverified{paper\_xi\_terminal\_zm\_leyang\_perron\_v3\_linear\_bound}
写
\[
\Lambda(t)=2+\sum_{n\ge 1}a_n t^n\qquad(a_n\in\QQ).
\]
则对任意 \(n\ge 1\) 有
\[
v_3(a_n)\ge 1-3n.
\]
等价地，存在 \(b_n\in\ZZ_3\) 使得
\[
a_n=\frac{b_n}{3^{3n-1}}\qquad(n\ge 1).
\]
\end{corollary}
\begin{proof}
由定理 \ref{thm:xi-terminal-zm-leyang-perron-p3-renormalization-hensel}，存在 \(v(u)=\sum_{n\ge 1}b_nu^n\in\ZZ_3[[u]]\) 使
\[
\Lambda(27u)=2+3v(u).
\]
比较系数得
\[
2+\sum_{n\ge 1}a_n(27u)^n
=2+3\sum_{n\ge 1}b_nu^n,
\]
从而 \(a_n=\frac{3b_n}{27^n}=\frac{b_n}{3^{3n-1}}\)，且 \(b_n\in\ZZ_3\)。证毕。
\end{proof}

\begin{theorem}[特征 \(3\) 的显式约化方程与自动结构]\label{thm:xi-terminal-zm-leyang-perron-p3-mod3-algebraic-automatic}
\leanverified{paper\_xi\_terminal\_zm\_leyang\_perron\_p3\_mod3\_algebraic\_automatic}
令 \(v(u)\in\ZZ_3[[u]]\) 为定理 \ref{thm:xi-terminal-zm-leyang-perron-p3-renormalization-hensel} 的唯一解，并记其逐项模 \(3\) 约化为 \(\overline v(u)\in\FF_3[[u]]\)。
则 \(\overline v(u)\) 满足显式三次代数方程
\[
\overline v(u)^3-\overline v(u)^2+\overline v(u)+u=0,
\qquad
\overline v(0)=0,
\]
并由初值条件确定唯一解。
因此 \(\overline v(u)\) 在 \(\FF_3(u)\) 上代数，其系数序列为 \(3\)-自动序列（等价地，具有有限 \(3\)-核）。
\end{theorem}
\begin{proof}
将定理 \ref{thm:xi-terminal-zm-leyang-perron-p3-renormalization-hensel} 的 \(H(u,v)\) 模 \(3\) 约化。
由于 \(27,18,24,3\equiv 0\pmod 3\)，且 \(-5\equiv 1\)、\(7\equiv 1\)、\(5\equiv -1\) 于 \(\FF_3\)，得
\[
H(u,v)\equiv u+v^3-v^2+v\quad(\bmod\,3),
\]
从而 \(\overline v\) 满足 \(\overline v^3-\overline v^2+\overline v+u=0\)。
又 \(\partial_v(u+v^3-v^2+v)|_{(0,0)}=1\neq 0\) 于 \(\FF_3\)，形式隐函数定理给出满足 \(\overline v(0)=0\) 的唯一解。
代数性推出自动性由 Christol 定理给出。证毕。
\end{proof}

\begin{theorem}[Perron 分支在 \texorpdfstring{$\QQ_3$}{Q3} 上的收敛半径]\label{thm:xi-terminal-zm-leyang-perron-p3-radius}
\leanverified{paper\_xi\_terminal\_zm\_leyang\_perron\_p3\_radius}
将 \(\Lambda(t)\) 视为 \(\QQ_3\) 上的幂级数，则其 \(3\)-进收敛半径严格等于
\[
R_3(\Lambda)=3^{-3}.
\]
\end{theorem}
\begin{proof}
由推论 \ref{cor:xi-terminal-zm-leyang-perron-v3-linear-bound}，\(v_3(a_n)\ge 1-3n\)，故
\(
|a_n|_3\le 3^{3n-1}
\)
并推出 \(R_3(\Lambda)\ge 3^{-3}\)。
另一方面，若仅有有限多个 \(n\) 使 \(b_n\not\equiv 0\pmod 3\)，则 \(\overline v(u)\) 为多项式；
但方程 \(\overline v^3-\overline v^2+\overline v+u=0\) 不存在非零多项式解（比较次数即可），矛盾。
故存在无穷多个 \(n\) 使 \(b_n\not\equiv 0\pmod 3\)，从而对这无穷多 \(n\) 有 \(v_3(a_n)=1-3n\)。
因此
\[
\limsup_{n\to\infty}|a_n|_3^{1/n}=3^{3},
\]
从而 \(R_3(\Lambda)\le 3^{-3}\)。
合并两端不等式得到 \(R_3(\Lambda)=3^{-3}\)。证毕。
\end{proof}

\begin{corollary}[\texorpdfstring{$p=3$}{p=3} 处的局部半稳定化与覆盖度降阶]\label{cor:xi-terminal-zm-leyang-perron-p3-semistable-degree-drop}
将 \(\ZZ_3\) 上的平面曲线族 \(\Pi(\lambda,1+t)=0\) 视为 \((\lambda,t)=(2,0)\) 的形式邻域内的局部模型。
在变量替换 \(t=27u\)、\(\lambda=2+3v\) 下，该形式邻域可写为 \(H(u,v)=0\) 且 \(\partial_vH(0,0)\in\ZZ_3^\times\)，因而在 \(\ZZ_3\) 上正则。
其特殊纤维（模 \(3\) 约化）由
\[
u=-v^3+v^2-v
\]
给出，从而在特征 \(3\) 上诱导出一般可分的三次覆盖 \(v\mapsto u\)。
\end{corollary}
\begin{proof}
正则性由 \(\partial_vH(0,0)\in\ZZ_3^\times\) 的隐函数口径直接得到。
模 \(3\) 约化由
\(
H(u,v)\equiv u+v^3-v^2+v
\)
立得。
最后，\(\frac{d}{dv}(-v^3+v^2-v)= -3v^2+2v-1\equiv 2v-1\) 于 \(\FF_3\) 非恒零，故该三次映射在一般点可分。证毕。
\end{proof}
\leanverified{paper\_xi\_terminal\_zm\_leyang\_perron\_p3\_semistable\_degree\_drop}

% Ramification and inertia at p=3 for the semistable special fiber.

\begin{theorem}[特征 \(3\) 特殊纤维三次覆盖的判别式与 Swan 指数]\label{thm:xi-terminal-zm-leyang-perron-p3-cubic-discriminant-s3-swan}
\leanverified{paper\_xi\_terminal\_zm\_leyang\_perron\_p3\_cubic\_discriminant\_s3\_swan}
在推论 \ref{cor:xi-terminal-zm-leyang-perron-p3-semistable-degree-drop} 的特殊纤维上，考虑三次关系
\[
u=-v^3+v^2-v,
\]
其给出 \(\PP^1_v\to\PP^1_u\) 的次数 \(3\) 的可分覆盖。
令函数域扩张为 \(\FF_3(v)/\FF_3(u)\)，并取最小多项式
\[
f(v):=v^3-v^2+v+u\in \FF_3(u)[v].
\]
则：
\begin{enumerate}
  \item \(f\) 的判别式在 \(\FF_3(u)\) 中等于 \(u\)，因而非平方；从而该三次扩张的正规闭包伽罗瓦群为 \(S_3\)；
  \item 分歧点恰为 \(u=0\) 与 \(u=\infty\)。
        在 \(u=0\) 处存在唯一分歧点，分歧指数为 \(2\)，且为驯分歧；
        在 \(u=\infty\) 处，分歧指数为 \(3\)，并必为野分歧，其 Swan 指数为 \(1\)。
\end{enumerate}
\end{theorem}
\begin{proof}
将关系式改写为 \(f(v)=0\)。
对一般三次 \(v^3+av^2+bv+c\) 的判别式公式
\[
\Delta(a,b,c)=a^2b^2-4b^3-4a^3c-27c^2+18abc
\]
代入 \(a=-1,b=1,c=u\)，并在特征 \(3\) 下化简（此时 \(27=18=0\)，且 \(4\equiv 1\)），得到
\(
\Delta=u
\)。
由于 \(u\) 在 \(u=0\) 处为一阶零点、在无穷远处为一阶极点，故其平方类不平凡，从而正规闭包伽罗瓦群为 \(S_3\)。
\par
分歧点由判别式与无穷远处的局部次数决定。
\(\Delta=u\) 强制 \(u=0\) 为分歧值。
另一方面，映射 \(u=-v^3+v^2-v\) 在 \(v=\infty\) 处具有极点阶 \(3\)，故 \(u=\infty\) 亦为分歧值。
导数为
\[
\frac{du}{dv}=-3v^2+2v-1\equiv 2v-1\quad(\bmod\,3),
\]
其唯一有限零点为 \(v=2\)，且 \(u(2)=0\)，从而 \(u=0\) 处恰有一个分歧点，分歧指数为 \(2\)，并因 \(3\nmid 2\) 为驯分歧。
\par
对可分覆盖 \(\PP^1_v\to\PP^1_u\) 应用 Riemann--Hurwitz 的不同除子形式，得到不同指数总度为 \(2\cdot 3-2=4\)。
已知 \(u=0\) 的驯分歧点贡献不同指数 \(e-1=1\)，故剩余不同指数度为 \(3\)，只能由 \(u=\infty\) 承担。
在无穷远处分歧指数为 \(3\) 且与特征相同，必为野分歧；其 Swan 指数由
\(
\mathrm{Swan}=d-(e-1)=3-(3-1)=1
\)
读出。证毕。
\end{proof}

\begin{theorem}[特征 \(3\) 的惯性作用在 \(S_4\) 中的点稳定子嵌入]\label{thm:xi-terminal-zm-leyang-perron-p3-inertia-point-stabilizer}
\leanverified{paper\_xi\_terminal\_zm\_leyang\_perron\_p3\_inertia\_point\_stabilizer}
考虑四叶覆盖 \(\Pi(\lambda,1+t)=0\) 在 \(t=0\) 的特征 \(3\) 约化。
则在 \(\FF_3\) 上有分解
\[
\Pi(\lambda,1)\equiv (\lambda-1)(\lambda+1)^3.
\]
因而在 \(t=0\) 的特殊纤维上出现一张简单片与三张三重并合片。
在推论 \ref{cor:xi-terminal-zm-leyang-perron-p3-semistable-degree-drop} 的半稳定化坐标下，三重并合片诱导的三次可分分支的正规闭包惯性群为 \(S_3\)，并在四叶单值群 \(S_4\) 中以点稳定子 \(S_3\hookrightarrow S_4\) 的方式嵌入：其作用固定简单片，并在其余三张片上给出自然三点置换表示；其中 \(3\)-Sylow 子群由一个 \(3\)-循环生成。
\end{theorem}
\begin{proof}
模 \(3\) 约化分解已在命题 \ref{prop:xi-terminal-zm-leyang-perron-p3-minimal-27} 的证明中给出。
简单因子 \((\lambda-1)\) 表明存在一张在 \(t=0\) 处保持分离的片；三重因子 \((\lambda+1)^3\) 对应三张在该点并合的片族。
推论 \ref{cor:xi-terminal-zm-leyang-perron-p3-semistable-degree-drop} 的半稳定化把三重并合片在特征 \(3\) 下的局部模型识别为三次覆盖 \(u=-v^3+v^2-v\)。
由定理 \ref{thm:xi-terminal-zm-leyang-perron-p3-cubic-discriminant-s3-swan}，其正规闭包伽罗瓦群为 \(S_3\)，且在无穷远处出现野 \(3\)-分歧，从而惯性群含有 \(3\)-循环。
因此在四叶分支的置换表示下，惯性群固定简单片并对其余三张片给出 \(S_3\) 的自然作用，亦即点稳定子嵌入。证毕。
\end{proof}

\begin{corollary}[\texorpdfstring{$3$}{3}-进估值下界的无穷次取到]\label{cor:xi-terminal-zm-leyang-perron-p3-slope-attained-infinitely-often}
沿用推论 \ref{cor:xi-terminal-zm-leyang-perron-v3-linear-bound} 的记号，写
\[
\Lambda(t)=2+\sum_{n\ge 1}a_n t^n,\qquad
b_n:=3^{3n-1}a_n\in\ZZ_3.
\]
则存在无穷多个 \(n\) 使 \(b_n\in\ZZ_3^\times\)，等价地，
\[
v_3(a_n)=1-3n
\]
无穷多次成立。
\end{corollary}
\begin{proof}
由定理 \ref{thm:xi-terminal-zm-leyang-perron-p3-mod3-algebraic-automatic}，\(\overline v(u)=\sum_{n\ge 1}(b_n\bmod 3)u^n\in\FF_3[[u]]\) 满足三次代数方程
\(
\overline v^3-\overline v^2+\overline v+u=0
\)，从而 \(\overline v\) 非多项式。
因此其系数中非零项出现无穷多次，即存在无穷多 \(n\) 使 \(b_n\bmod 3\neq 0\)，等价于 \(b_n\in\ZZ_3^\times\) 无穷多次成立。
由 \(a_n=b_n/3^{3n-1}\) 得所述估值等式无穷多次取到。证毕。
\end{proof}

\endinput


\begin{conclusion}[累积量密度序列的分母支撑仅含 \texorpdfstring{$\{2,3\}$}{\{2,3\}}]\label{con:xi-terminal-zm-leyang-perron-cumulant-denominator-2-3}
定义
\[
\psi(s):=\log\frac{\Lambda(e^s-1)}{2},\qquad
\kappa_n^\infty:=\psi^{(n)}(0)\quad(n\ge 1).
\]
则
\[
\kappa_n^\infty\in\ZZ\!\left[\frac16\right]\qquad(n\ge 1).
\]
\end{conclusion}
\begin{proof}[证明要点]
由结论 \ref{con:xi-terminal-zm-leyang-perron-global-pintegrality}，\(\Lambda(t)\) 系数属于 \(\ZZ[1/3]\)。
又
\[
(e^s-1)^k
=
k!\sum_{n\ge k}S(n,k)\frac{s^n}{n!},
\]
其中 \(S(n,k)\in\ZZ\)，故 \(\Lambda(e^s-1)\) 的各阶导数在 \(s=0\) 落在 \(\ZZ[1/3]\)。
再由 \(\psi(s)=\log(\Lambda(e^s-1))-\log 2\)，对数微分仅额外引入 \(\Lambda(0)=2\) 的分母因子，
因此最多新增素因子 \(2\)。
故 \(\kappa_n^\infty\in\ZZ[1/6]\)。证毕。
\end{proof}

\begin{conclusion}[\texorpdfstring{$p\neq 3$}{p!=3} 下的模 \(p\) 自动性与有限 \(p\)-核]\label{con:xi-terminal-zm-leyang-perron-modp-automaticity}
对任意素数 \(p\neq 3\)，记
\[
\overline{\Lambda}_p(t):=\sum_{n\ge 0}\overline{a}_n t^n\in\FF_p[[t]]
\]
为 \(\Lambda(t)\) 的逐项模 \(p\) 约化。
则 \(\overline{\Lambda}_p(t)\) 在 \(\FF_p(t)\) 上代数，且系数序列 \((\overline{a}_n)_{n\ge 0}\) 为 \(p\)-自动序列（等价地，其 \(p\)-核有限）。
\end{conclusion}
\begin{proof}[证明要点]
由结论 \ref{con:xi-terminal-zm-leyang-perron-global-pintegrality}，\(\Lambda(t)\in\ZZ_p[[t]]\)（\(p\neq 3\)），故可逐项约化得到 \(\overline{\Lambda}_p\)。
约化后仍满足
\(
\Pi(\overline{\Lambda}_p(t),1+t)=0
\)
的模 \(p\) 方程，故 \(\overline{\Lambda}_p\) 为 \(\FF_p(t)\)-代数幂级数。
由 Christol 定理，代数性当且仅当系数序列 \(p\)-自动，结论即得。证毕。
\end{proof}

\begin{conclusion}[三阶算子的单值群与 \texorpdfstring{$S_4$}{S4} 标准三维表示]\label{con:xi-terminal-zm-leyang-perron-fuchs-monodromy-s4}
设 \(\lambda_1,\dots,\lambda_4\) 为 \(\Pi(\lambda,y)=0\) 的四个解析分支，并定义
\[
f_i(y):=\lambda_i(y)-\frac14\qquad(1\le i\le 4).
\]
则
\[
V=\mathrm{span}_{\CC}\{f_1,f_2,f_3,f_4\}
=
\left\{\sum_{i=1}^4 c_i f_i:\ \sum_{i=1}^4 c_i=0\right\},
\qquad \dim_{\CC}V=3,
\]
且 \(V\) 恰为结论 \ref{con:xi-terminal-zm-leyang-perron-centered-fuchs-operator} 中 \(L_3\) 的解空间。
解析延拓在 \(V\) 上诱导的表示同构于 \(S_4\) 在
\(
\{(x_1,x_2,x_3,x_4)\in\CC^4:\sum x_i=0\}
\)
上的标准三维表示，故单值群同构于 \(S_4\)。
\end{conclusion}
\begin{proof}[证明要点]
由 Vieta 关系 \(\sum_i\lambda_i=1\) 得 \(\sum_i f_i=0\)，因此四分支张成空间至多三维。
另一方面，\(L_3\) 为三阶算子且各 \(f_i\) 均被其湮灭，故解空间维数恰为 \(3\)，于是两者一致。
由定理 \ref{thm:xi-terminal-zm-leyang-monodromy-s4}，\(\Pi\) 的四叶覆盖单值群为 \(S_4\)；
其在 \((f_1,\dots,f_4)\) 上为指标置换作用，并保持和为零超平面，得到标准三维表示。
该表示忠实，故诱导单值群同构于 \(S_4\)。证毕。
\end{proof}

\begin{conclusion}[指数对数坐标下的 Lee--Yang 微分闭合：\texorpdfstring{$\lambda(s)=\lambda_+(e^s)$}{lambda(s)} 的一阶代数方程与椭圆核一致性]\label{con:xi-terminal-zm-leyang-exponential-time-algebraic-differential-closure}
沿用谱四次
\[
\Pi(\lambda,y):=\lambda^4-\lambda^3-(2y+1)\lambda^2+\lambda+y(y+1)=0
\]
及其在 \(y=1\) 邻域满足 \(\lambda_+(1)=2\) 的解析特征支 \(\lambda_+(y)\)。
在 \(s=0\) 的邻域定义对数参数化
\[
\lambda(s):=\lambda_+(e^s),
\qquad
\dot\lambda:=\frac{\dd \lambda}{\dd s}.
\]
则 \(\lambda(s)\) 满足显式一阶代数微分闭式
\begin{equation}\label{eq:xi-terminal-zm-leyang-exponential-time-lambda-ode}
\boxed{\;
\begin{aligned}
0={}&(16\lambda^4+7\lambda^3-9\lambda^2+\lambda+1)\,\dot\lambda^{2}
(-16\lambda^5-4\lambda^4+20\lambda^3-5\lambda+1)\,\dot\lambda \\
&\qquad\qquad+(4\lambda^6-8\lambda^4+\lambda^3+4\lambda^2-\lambda).
\end{aligned}\;}
\end{equation}
并且将 \eqref{eq:xi-terminal-zm-leyang-exponential-time-lambda-ode} 视为关于 \(\dot\lambda\) 的二次方程，其判别式满足严格因式分解
\begin{equation}\label{eq:xi-terminal-zm-leyang-exponential-time-lambda-ode-discriminant}
\boxed{\;
\Delta_{\dot\lambda}(\lambda)
=\bigl(4\lambda^3-4\lambda+1\bigr)\bigl(2\lambda^3+2\lambda^2-\lambda+1\bigr)^2\; }.
\end{equation}
因此在函数域层面成立
\begin{equation}\label{eq:xi-terminal-zm-leyang-exponential-time-function-field}
\QQ(\lambda,\dot\lambda)=\QQ\!\left(\lambda,\sqrt{4\lambda^3-4\lambda+1}\right),
\end{equation}
从而指数对数坐标的微分闭合与椭圆核
\begin{equation}\label{eq:xi-terminal-zm-leyang-exponential-time-elliptic-kernel}
\mathcal E:\quad w^2=4\lambda^3-4\lambda+1
\end{equation}
处于同一函数域（与结论 \ref{con:xi-terminal-zm-elliptic-normalization} 的椭圆化一致）。
更进一步，外参 \(y=e^s\) 可由 \((\lambda,\dot\lambda)\) 有理重构为
\begin{equation}\label{eq:xi-terminal-zm-leyang-exponential-time-y-reconstruction}
\boxed{\;
y=
\frac{\dot\lambda\,(4\lambda^3-3\lambda^2-2\lambda+1)-2\lambda(\lambda-1)^2(\lambda+1)}{-2\lambda^2+4\lambda\dot\lambda+1}\; }.
\end{equation}
\end{conclusion}
\begin{proof}[证明要点（隐式微分与消元）]
对恒等式 \(\Pi(\lambda(s),e^s)\equiv0\) 微分得
\(
\partial_\lambda\Pi(\lambda,y)\,\dot\lambda+\partial_y\Pi(\lambda,y)\,\dot y=0
\)
且 \(\dot y=y\)。
由 \(\Pi(\lambda,y)=0\) 与
\(
\partial_\lambda\Pi(\lambda,y)\,\dot\lambda+y\,\partial_y\Pi(\lambda,y)=0
\)
得到关于 \(y\) 的两条二次代数约束；消去 \(y^2\) 得到 \(y\) 的有理表达式 \eqref{eq:xi-terminal-zm-leyang-exponential-time-y-reconstruction}，
代回即得 \eqref{eq:xi-terminal-zm-leyang-exponential-time-lambda-ode}。
判别式分解 \eqref{eq:xi-terminal-zm-leyang-exponential-time-lambda-ode-discriminant} 为直接计算并在 \(\ZZ[\lambda]\) 中因式分解。
函数域恒等式 \eqref{eq:xi-terminal-zm-leyang-exponential-time-function-field} 由判别式的唯一非平凡平方自由部 \((4\lambda^3-4\lambda+1)\) 立即推出；
而 \eqref{eq:xi-terminal-zm-leyang-exponential-time-elliptic-kernel} 与结论 \ref{con:xi-terminal-zm-elliptic-normalization} 的椭圆模型一致。
可复现核验脚本：\texttt{scripts/exp\_fold\_zm\_leyang\_exponential\_time\_differential\_closure\_audit.py}。证毕。
\end{proof}

\begin{conclusion}[\texorpdfstring{$\psi(s)=\log\lambda_+(e^s)-\log2$}{psi(s)} 的微分代数性与非 \texorpdfstring{$D$}{D}-有限性]\label{con:xi-terminal-zm-psi-differential-algebraic-non-dfinite}
令
\[
\psi(s):=\log\lambda_+(e^s)-\log2.
\]
则 \(\psi\) 在 \(s=0\) 的解析芽满足：
\begin{enumerate}
  \item \(\psi\) 为微分代数函数：其满足某个系数属 \(\QQ(s)\) 的非平凡代数微分方程；
  \item \(\psi\) 非 \(D\)-有限：不存在首项系数为多项式的非零线性常微分算子湮灭 \(\psi\)。
  因而泰勒系数列 \(\bigl(\psi^{(n)}(0)\bigr)_{n\ge 0}\) 非 \(P\)-递归（非 holonomic）。
\end{enumerate}
\end{conclusion}
\begin{proof}[证明要点]
第一条由结论 \ref{con:xi-terminal-zm-leyang-exponential-time-algebraic-differential-closure} 给出的 \(\lambda(s)=\lambda_+(e^s)\) 的一阶代数微分闭式推出：
\(\lambda\) 因而微分代数；又 \(\psi'=\dot\lambda/\lambda\) 为 \((\lambda,\dot\lambda)\) 的有理函数且 \(\lambda(0)=2\neq0\)，故 \(\psi\) 的解析芽同属微分代数闭包。
\par
第二条利用奇点集合的无限性。
由结论 \ref{con:xi-terminal-zm-discriminant-leyang}，\(\lambda_+(y)\) 的有限分歧值集合包含某个满足 \(y_\star\neq0\) 的代数分歧值（例如 Lee--Yang 三次 \(P_{\mathrm{LY}}(y)=0\) 的任一根）。
取任一对数支 \(s_\star:=\log y_\star\)，则指数映射的周期性给出无限前像族
\[
s_\star+2\pi i\,\ZZ.
\]
在每个点 \(s=s_\star+2\pi i k\) 处，复合函数 \(s\mapsto \lambda_+(e^s)\) 继承 \(y=y_\star\) 的代数分歧而出现不可去奇性，故 \(\psi\) 在 \(\CC\) 上具有无穷多个奇点。
另一方面，若 \(\psi\) 为 \(D\)-有限，则其满足某个首项系数为多项式的线性常微分方程，故其解析延拓的可能奇点仅能位于该首项系数的零点集合（有限集）与无穷远点。
两者矛盾，故 \(\psi\) 非 \(D\)-有限。
最后，\(D\)-有限解析芽的泰勒系数列必为 \(P\)-递归为标准等价，因而 \(\bigl(\psi^{(n)}(0)\bigr)\) 非 \(P\)-递归。证毕。
\end{proof}

\begin{conclusion}[主特征支局部喷射的 \(3\)-局部整性：\texorpdfstring{$\lambda_+(1+z)\in\ZZ[1/3][[z]]$}{lambda+(1+z) in Z[1/3]} 与 \texorpdfstring{$2^n\psi^{(n)}(0)\in\ZZ[1/3]$}{2^n psi^(n)(0) in Z[1/3]}]\label{con:xi-terminal-zm-lambda-plus-3local-denominators}
记
\[
\lambda_+(1+z)=2+\sum_{n\ge 1}a_n z^n\in\QQ[[z]],
\qquad
\psi(s)=\sum_{n\ge 1}\frac{\psi^{(n)}(0)}{n!}\,s^n.
\]
则：
\begin{enumerate}
  \item 对一切 \(n\ge 1\) 有 \(a_n\in\ZZ[1/3]\)，亦即 \(\lambda_+(1+z)\in\ZZ[1/3][[z]]\)；
  \item 对一切 \(n\ge 1\) 有 \(2^n\psi^{(n)}(0)\in\ZZ[1/3]\)，从而 \(\psi^{(n)}(0)\in\ZZ[1/6]\)。
\end{enumerate}
\end{conclusion}
\begin{proof}[证明要点（\(p\)-进隐函数与导数闭合）]
由 \(\Pi(2,1)=0\) 且
\(
\partial_\lambda\Pi(2,1)=9
\)，
对任一素数 \(p\neq3\)，\(\partial_\lambda\Pi(2,1)\) 在 \(\ZZ_p\) 中可逆。
由 \(p\)-进隐函数定理，存在唯一 \(\lambda_p(z)\in\ZZ_p[[z]]\) 使
\(
\Pi(\lambda_p(z),1+z)\equiv0
\)
且 \(\lambda_p(0)=2\)。
而 \(\lambda_+(1+z)\in\QQ[[z]]\) 亦满足同一初值条件，故在 \(\QQ_p[[z]]\) 中与 \(\lambda_p\) 一致；
于是 \(a_n\in\ZZ_p\) 对一切 \(p\neq3\) 成立，从而 \(a_n\in\ZZ[1/3]\)。
\par
再记 \(\lambda(s)=\lambda_+(e^s)\)。
由 \(e^s\) 在 \(s=0\) 处各阶导数皆为 \(1\)，复合求导（Fa\`a di Bruno）表明 \(\lambda^{(n)}(0)\) 为 \(\{\lambda_+^{(k)}(1)\}_{k\le n}\) 的整系数组合；
而 \(\lambda_+^{(k)}(1)=k!\,a_k\in\ZZ[1/3]\)，故 \(\lambda^{(n)}(0)\in\ZZ[1/3]\)。
又 \(\psi=\log(\lambda/2)\) 且 \(\lambda(0)=2\)，于是 \(\psi^{(n)}(0)\) 为 \(\{\lambda^{(k)}(0)/2\}_{k\le n}\) 的整系数多项式（来自对数导数的递推），
每项分母至多贡献 \(2^n\)，从而 \(2^n\psi^{(n)}(0)\in\ZZ[1/3]\)。证毕。
\end{proof}

\begin{conclusion}[有限阶局部喷射唯一确定谱四次：\(13\)-阶芽锁定 \texorpdfstring{$\Pi(\lambda,y)$}{Pi(lambda,y)} 与 Lee--Yang 三次]\label{con:xi-terminal-zm-local-jet-rigidly-determines-pi}
在二元多项式集合中考虑逆问题：设
\(
Q(\lambda,y)\in\QQ[\lambda,y]
\) 满足 \(\deg_\lambda Q\le4\)、\(\deg_y Q\le2\)，并且 \(\lambda^4\) 的系数归一为 \(1\)。
若已知主特征支 \(\lambda_+(y)\) 在 \(y=1\) 的 \(13\)-阶喷射（即 \(a_1,\dots,a_{13}\) 连同 \(\lambda_+(1)=2\)），
则满足
\[
Q(\lambda_+(y),y)\equiv 0 \qquad\text{于}\qquad \QQ[[y-1]]/(y-1)^{14}
\]
的上述 \(Q\) 唯一存在，且必为谱四次 \(\Pi\) 本身。
因此由恒等式
\(
\mathrm{Disc}_{\lambda}(\Pi)=-y(y-1)P_{\mathrm{LY}}(y)
\)
所定义的 Lee--Yang 三次 \(P_{\mathrm{LY}}(y)\) 亦被该有限喷射唯一决定。
\end{conclusion}
\begin{proof}[证明要点（线性化与可逆性）]
将一般 \(Q\) 写为
\[
Q(\lambda,y)=\lambda^4+\sum_{(i,j)\neq(4,0)}c_{ij}\lambda^i y^j,
\qquad 0\le i\le4,\ 0\le j\le2,
\]
未知量为 \(14\) 个系数 \(c_{ij}\in\QQ\)。
代入 \(y=1+z\) 与
\(
\lambda=\lambda_+(1+z)=2+\sum_{n=1}^{13}a_n z^n+O(z^{14})
\)，
并取模 \(z^{14}\)，条件 \(Q(\lambda_+(1+z),1+z)\equiv0\pmod{z^{14}}\) 等价于 \(z^0,\dots,z^{13}\) 的 \(14\) 个系数同时为零，
从而得到一个 \(14\times14\) 的线性方程组。
直接计算可得其系数矩阵行列式非零，故解唯一存在。
另一方面，\(\Pi\) 显然给出一个满足条件的解，故唯一解必为 \(\Pi\)。
由判别式恒等式 \(\mathrm{Disc}_{\lambda}(\Pi)=-y(y-1)P_{\mathrm{LY}}(y)\) 立得 \(P_{\mathrm{LY}}\) 的唯一性。证毕。
可复现核验脚本：\texttt{scripts/exp\_fold\_zm\_leyang\_exponential\_time\_differential\_closure\_audit.py}。
\end{proof}



\begin{conclusion}[重量大偏差的代数鞍点方程：倾斜参数的五次消元证书]\label{con:xi-terminal-zm-ldp-saddle-quintic}
对 \(\alpha\in(0,\tfrac12)\)，令 \(y(\alpha)>0\) 为满足鞍点条件
\[
\alpha=\frac{y\,\lambda_+'(y)}{\lambda_+(y)}
\]
的唯一正解（\(\lambda_+\) 为结论 \ref{con:xi-terminal-zm-lln-clt} 中 \(\Pi(\lambda,y)=0\) 的主特征支）。
则 \(y(\alpha)\) 必满足关于 \((\alpha,y)\) 的显式五次代数方程
\[
\begin{aligned}
0=&\ 256\alpha^4y^5+411\alpha^4y^4-91\alpha^4y^3-379\alpha^4y^2-165\alpha^4y-32\alpha^4\\
&-512\alpha^3y^5-566\alpha^3y^4+337\alpha^3y^3+512\alpha^3y^2+197\alpha^3y+32\alpha^3\\
&+384\alpha^2y^5+228\alpha^2y^4-298\alpha^2y^3-214\alpha^2y^2-28\alpha^2y\\
&-128\alpha y^5-8\alpha y^4+86\alpha y^3+16\alpha y^2-4\alpha y\\
&+16y^5-8y^4-4y^3+y^2,
\end{aligned}
\]
并且速率函数在该支上可写为
\[
I(\alpha)=\alpha\log y(\alpha)-\log \lambda_+\!\bigl(y(\alpha)\bigr)+\log 2.
\]
因此该 LDP 的鞍点求解与速率函数评估在严格意义下为有限代数可计算对象。
\end{conclusion}
\begin{proof}[证明要点（隐式微分 + 消元）]
鞍点条件等价于 \(\alpha=\psi'(s)\) 与 \(y=e^s\) 的联立，从而可写为
\(\alpha=-y\,\partial_y\Pi(\lambda,y)/(\lambda\,\partial_\lambda\Pi(\lambda,y))\) 与 \(\Pi(\lambda,y)=0\) 的共同解（其中 \(\lambda=\lambda_+(y)\)）。
对变量 \(\lambda\) 作消元（例如取关于 \(\lambda\) 的 resultant）即可得到所示五次多项式；
并由结论 \ref{con:xi-terminal-zm-ldp} 的 Legendre 形式推出 \(I(\alpha)\) 的鞍点表达。
上述消元证书可由脚本 \texttt{scripts/exp\_fold\_zm\_bivariate\_partition\_audit.py} 复核。证毕。
\end{proof}

\begin{conclusion}[自反型同余通道的唯一性：仅 \(y^2+y+1=0\) 触发单位圆对称]\label{con:xi-terminal-zm-self-inversive-unique}
对 \(|y|=1\)，称 \(\Pi(\lambda,y)\) 为自反（self-inversive）多项式，
若存在常数 \(c\in\CC\)（\(|c|=1\)）使得
\[
\lambda^4\,\overline{\Pi\!\left(\frac{1}{\overline{\lambda}},y\right)}=c\,\Pi(\lambda,y).
\]
则当且仅当 \(y(y+1)=-1\)，亦即 \(y\) 为本原三次单位根（满足 \(y^2+y+1=0\)），\(\Pi(\lambda,y)\) 才具有该自反性。
在且仅在此情形，\(\Pi(\lambda,y)\) 的根集关于单位圆满足严格对称：要么位于单位圆上，要么以
\(\lambda\mapsto 1/\overline{\lambda}\) 成对出现；并且必存在一对互为倒共轭的根，其模长严格大于 \(1\)。
因此“同余权重的单位圆对称”在该体系中是严格量子化的：模 \(3\) 为唯一触发自反对称的非平凡同余通道。
\end{conclusion}
\begin{proof}[证明要点]
写
\(
\Pi(\lambda,y)=\sum_{k=0}^{4}a_k(y)\lambda^k
\)，其中
\(
a_4=1,\ a_3=-1,\ a_2=-(2y+1),\ a_1=1,\ a_0=y(y+1)
\)。
自反性等价于存在 \(|c|=1\) 使 \(a_k(y)=c\,\overline{a_{4-k}(y)}\) 对所有 \(k\) 成立。
由 \(a_3=-1=c\,\overline{a_1}=c\) 得 \(c=-1\)。
再由 \(a_4=1=c\,\overline{a_0}=-\overline{y(y+1)}\) 得 \(\overline{y(y+1)}=-1\)；
当 \(|y|=1\) 时 \(\overline{y(y+1)}=(1+y)/y^2\)，故等价于 \(y^2+y+1=0\)。
根集的单位圆对称为自反多项式的标准结论。证毕。
\end{proof}

\begin{conclusion}[重量细分系数的二维线性动力：\((m,k)\) 平面上的有限邻域闭合递推]\label{con:xi-terminal-zm-2d-recurrence-amk}
令
\[
a_{m,k}:=\sum_{\substack{x\in X_m\\ |x|_1=k}} d_m(x),
\qquad
Z_m(y)=\sum_{k\in\ZZ} a_{m,k}y^k .
\]
则由结论 \ref{con:xi-terminal-zm-order4-rational} 的四阶递推在 \(y\) 上的多项式结构，得到 \(a_{m,k}\) 的闭合二维递推：
对一切 \(m\ge 4\) 与 \(k\in\ZZ\)，
\[
\boxed{\;
a_{m,k}=a_{m-1,k}+a_{m-2,k}+2a_{m-2,k-1}-a_{m-3,k}-a_{m-4,k-1}-a_{m-4,k-2}\;}
\]
并以边界条件 \(a_{m,k}=0\)（当 \(k<0\) 或 \(k>\lceil m/2\rceil\)）封闭。
该二维递推把“推前分布的重量剖分”转写为一个完全局域的线性动力系统，使得重量层（固定 \(k\)）与长度层（固定 \(m\)）之间的耦合仅通过有限邻域发生；
因此重量分布的全体精细信息可在不显式构造任何高维碰撞核的前提下，由 \((m,k)\) 层级递推直接生成并审计。
\end{conclusion}
\begin{proof}
将结论 \ref{con:xi-terminal-zm-order4-rational} 的递推式写成 \(y\)-幂级数，并比较 \(y^k\) 的系数即可得到所示二维递推；
边界条件由黄金均值合法语法的最大 \(1\)-计数 \(\max_{x\in X_m}|x|_1=\lceil m/2\rceil\) 直接给出。证毕。
\end{proof}

% (User input window)

\begin{conclusion}[固定重量层系数生成函数的分母刚性：统一闭式与二阶递推]\label{con:xi-terminal-zm-fixed-weight-generating-function-rigidity}
沿用结论 \ref{con:xi-terminal-zm-order4-rational} 与结论 \ref{con:xi-terminal-zm-2d-recurrence-amk} 的记号，并记
\[
a_{m,k}:=[y^k]\,Z_m(y),\qquad
A_k(t):=\sum_{m\ge 0}a_{m,k}t^m=[y^k]F(t,y),
\]
其中
\[
F(t,y)=\sum_{m\ge 0}Z_m(y)t^m
=\frac{1+yt-yt^2-y^2t^3}{1-t-(2y+1)t^2+t^3+y(y+1)t^4}.
\]
则对每个固定整数 \(k\ge 0\)，存在唯一多项式 \(\mathcal{P}_k(t)\in\ZZ[t]\) 使得
\[
\boxed{\;
A_k(t)=\frac{\mathcal{P}_k(t)}{(1-t)^{2k+2}(1+t)^{k+1}}\;}
\]
并且 \(\mathcal{P}_k\) 可由显式二阶递推唯一确定：令
\[
\mathcal{P}_0(t)=1,\qquad
\mathcal{P}_1(t)=-t\bigl(t^4-t^3-1\bigr),\qquad
\mathcal{P}_2(t)=t^3\bigl(t^5-2t^4-t^3+t^2+t+1\bigr),
\]
则对一切 \(k\ge 3\) 有
\[
\boxed{\;
\mathcal{P}_k(t)=(2t^2-t^4)\mathcal{P}_{k-1}(t)-t^4(1-t)^2(1+t)\mathcal{P}_{k-2}(t)\; }.
\]
并且当 \(k\ge 1\) 时成立整除性
\[
\boxed{\; t^{2k-1}\mid \mathcal{P}_k(t)\;}
\]
\end{conclusion}
\begin{proof}[证明要点]
将分母记为
\[
D(t,y):=1-t-(2y+1)t^2+t^3+y(y+1)t^4=D_0(t)+yD_1(t)+y^2D_2(t),
\]
其中
\[
D_0(t)=(1-t)^2(1+t),\qquad D_1(t)=-2t^2+t^4,\qquad D_2(t)=t^4,
\]
并将分子记为
\(
N(t,y)=1+y(t-t^2)-y^2t^3
\)。
写 \(F(t,y)=\sum_{k\ge 0}A_k(t)y^k\)，则由恒等式 \(D(t,y)F(t,y)=N(t,y)\) 比较 \(y^k\) 系数得
\[
D_0(t)A_k(t)+D_1(t)A_{k-1}(t)+D_2(t)A_{k-2}(t)=N_k(t),
\]
其中 \(N_0(t)=1\)、\(N_1(t)=t-t^2\)、\(N_2(t)=-t^3\)、且 \(N_k(t)=0\)（当 \(k\ge 3\)），并约定 \(A_{-1}=A_{-2}=0\)。
令 \(\mathcal{P}_k(t):=D_0(t)^{k+1}A_k(t)\)，则上式等价于
\[
\mathcal{P}_k(t)=\bigl(2t^2-t^4\bigr)\mathcal{P}_{k-1}(t)-t^4D_0(t)\mathcal{P}_{k-2}(t)+N_k(t)D_0(t)^k.
\]
对 \(k\ge 3\) 时 \(N_k\equiv 0\)，得到所示齐次二阶递推；
\(k=0,1,2\) 的显式初值由上式直接化简得到。
由于 \(D_0,D_1,D_2,N_k\in\ZZ[t]\) 且初值整系数，递推给出 \(\mathcal{P}_k\in\ZZ[t]\)。
\par
另一方面，\(\mathcal{P}_1(t)\) 含因子 \(t\)、\(\mathcal{P}_2(t)\) 含因子 \(t^3\)，且对 \(k\ge 3\) 的递推右端两项分别额外带因子 \(t^2\) 与 \(t^4\)，
故可归纳得到 \(t^{2k-1}\mid \mathcal{P}_k(t)\)（\(k\ge 1\)）。
最后由 \(D_0(t)^{k+1}=(1-t)^{2k+2}(1+t)^{k+1}\) 得到分母结构。证毕。
\end{proof}

\begin{conclusion}[分子多项式的黄金因子：模 \texorpdfstring{$t^2-t-1$}{t^2-t-1} 的五步标度闭合]\label{con:xi-terminal-zm-fixed-weight-numerator-golden-factor}
在结论 \ref{con:xi-terminal-zm-fixed-weight-generating-function-rigidity} 的记号下，令
\[
\Phi_{\mathrm{gold}}(t):=t^2-t-1\in\ZZ[t].
\]
则对一切整数 \(n\ge 0\)，成立严格整除性
\[
\boxed{\;\Phi_{\mathrm{gold}}(t)\mid \mathcal{P}_{5n+4}(t)\; }.
\]
\end{conclusion}
\begin{proof}[证明要点]
在商环 \(R:=\ZZ[t]/(\Phi_{\mathrm{gold}})\) 中有关系 \(t^2=t+1\)，从而
\[
(1-t)^2(1+t)=1-t-t^2+t^3\equiv 1,\qquad
2t^2-t^4\equiv -t,\qquad
t^4\equiv 3t+2
\pmod{\Phi_{\mathrm{gold}}}.
\]
将结论 \ref{con:xi-terminal-zm-fixed-weight-generating-function-rigidity} 的齐次递推在 \(R\) 中约化，得到对 \(k\ge 3\)
\[
\mathcal{P}_k\equiv -t\,\mathcal{P}_{k-1}-(3t+2)\mathcal{P}_{k-2}.
\]
令向量 \(v_k:=(\mathcal{P}_k,\mathcal{P}_{k-1})^{\mathsf T}\)，则对 \(k\ge 3\) 有
\[
v_k\equiv M v_{k-1},\qquad
M:=\begin{pmatrix}-t&-(3t+2)\\[2pt]1&0\end{pmatrix}\in \Mat_{2}(R).
\]
在 \(R\) 中对 \(M\) 作直接幂计算并约化可得
\[
M^5\equiv -(55t+34)\,I_2.
\]
因此对一切 \(k\ge 2\) 有五步标度关系
\[
v_{k+5}\equiv M^5 v_k\equiv -(55t+34)\,v_k,
\qquad\text{特别地}\quad
\mathcal{P}_{k+5}\equiv -(55t+34)\,\mathcal{P}_k.
\]
由初值在 \(R\) 中的约化计算
\(
\mathcal{P}_1\equiv -(1+t)
\)、\(
\mathcal{P}_2\equiv -(2+3t)
\)，进而由递推得到 \(\mathcal{P}_4\equiv 0\)。
于是对任意 \(n\ge 0\)，反复应用五步标度推出 \(\mathcal{P}_{5n+4}\equiv 0\)，
等价于 \(\Phi_{\mathrm{gold}}(t)\mid \mathcal{P}_{5n+4}(t)\)。
可复核脚本：\texttt{scripts/exp\_fold\_zm\_low\_weight\_quasipolynomial\_audit.py}。证毕。
\end{proof}

\begin{conclusion}[固定重量层系数的准多项式性：周期 \(2\)、次数量子化与首项系数闭式]\label{con:xi-terminal-zm-fixed-weight-quasipolynomial}
沿用结论 \ref{con:xi-terminal-zm-order4-rational} 与结论 \ref{con:xi-terminal-zm-2d-recurrence-amk} 的记号，并记
\[
a_{m,k}:=[y^k]\,Z_m(y)\qquad(m\ge 0,\ k\ge 0).
\]
则对每个固定整数 \(k\ge 0\)，存在唯一一对有理系数多项式 \(P_k,Q_k\in\QQ[m]\)，使对所有 \(m\ge 0\) 成立
\[
\boxed{\;a_{m,k}=P_k(m)+(-1)^m Q_k(m)\;}
\]
并且次数满足
\[
\deg P_k=2k+1,\qquad \deg Q_k=k.
\]
此外，\(P_k\) 与 \(Q_k\) 的首项系数具有封闭形式
\[
\bigl[m^{2k+1}\bigr]P_k(m)=\frac{1}{2^{k+1}(2k+1)!},
\qquad
\bigl[m^{k}\bigr]Q_k(m)=\frac{1}{2^{2k+2}k!}.
\]
等价地，令 \(m=2n+\varepsilon\)（\(\varepsilon\in\{0,1\}\)），则两条奇偶子序列 \(n\mapsto a_{2n,k}\) 与 \(n\mapsto a_{2n+1,k}\) 均为 \(n\) 的次数 \(2k+1\) 多项式。
并且两条奇偶分支的最高次项一致：当 \(n\to\infty\)（固定 \(k\)）时，
\[
a_{2n,k}=\frac{2^k}{(2k+1)!}\,n^{2k+1}+O_k(n^{2k}),
\qquad
a_{2n+1,k}=\frac{2^k}{(2k+1)!}\,n^{2k+1}+O_k(n^{2k}).
\]
\end{conclusion}
\begin{proof}[证明要点]
由结论 \ref{con:xi-terminal-zm-fixed-weight-generating-function-rigidity}，
对每个固定 \(k\ge 0\) 有显式闭式
\[
A_k(t)=\frac{\mathcal{P}_k(t)}{(1-t)^{2k+2}(1+t)^{k+1}}\in\QQ(t),
\]
从而 \(A_k(t)\) 仅在 \(t=\pm 1\) 处出现极点，且在 \(t=1\) 与 \(t=-1\) 处的极点阶分别恰为 \(2k+2\) 与 \(k+1\)。
对 \(A_k\) 作部分分式并逐项取系数即可得到对所有 \(m\ge 0\) 成立的分解
\(
a_{m,k}=P_k(m)+(-1)^mQ_k(m)
\)
及次数界；唯一性来自于：若 \(P(m)+(-1)^mQ(m)\equiv 0\) 对所有整数 \(m\ge 0\) 成立，则必有 \(P\equiv 0\) 且 \(Q\equiv 0\)。
\par
首项系数由 \(t=1\) 处的最高阶极点给出：由结论 \ref{con:xi-terminal-zm-fixed-weight-generating-function-rigidity} 的递推在 \(t=1\) 处代入可归纳得 \(\mathcal{P}_k(1)=1\)；并且 \((1+t)^{-(k+1)}\sim 2^{-(k+1)}\)（\(t\to 1\)），故
\(
A_k(t)\sim 2^{-(k+1)}(1-t)^{-(2k+2)}
\)，从而
\(
[t^m]A_k(t)=2^{-(k+1)}\frac{m^{2k+1}}{(2k+1)!}+O(m^{2k})
\)
给出所示首项系数。
\par
对 \(t=-1\) 的首项系数，记分母
\[
D(t,y)=D_0(t)+yD_1(t)+y^2D_2(t),
\qquad
D_0(t)=D(t,0),\ \ D_1(t)=\partial_y D(t,0)=-2t^2+t^4,
\]
则
\(
D_0(t)\sim 4(1+t)
\)
（\(t\to-1\)）且 \(D_1(-1)=-1\)。
由于 \(F(t,y)\) 的分子在 \(y=0\) 处取值为 \(1\)，且其 \(y\)-次数至多为 \(2\)，故 \(A_k(t)\) 在 \(t=-1\) 处的最高阶极点完全由 \([y^k]\bigl(1/D(t,y)\bigr)\) 决定。
写
\[
\frac{1}{D(t,y)}=\frac{1}{D_0(t)}
\left(1+\frac{yD_1(t)+y^2D_2(t)}{D_0(t)}\right)^{-1},
\]
则 \([y^k]\bigl(1/D(t,y)\bigr)\) 的最高阶项为 \((-1)^kD_1(t)^k/D_0(t)^{k+1}\)；
结合 \(D_0(t)\sim 4(1+t)\) 与 \(D_1(-1)=-1\) 得：
\[
A_k(t)=\frac{1}{4^{k+1}}(1+t)^{-(k+1)}+O\!\bigl((1+t)^{-k}\bigr)\qquad(t\to-1),
\]
从而 \(t=-1\) 的极点阶恰为 \(k+1\)，因此 \(\deg Q_k=k\)。
并且由
\(
[t^m](1+t)^{-(k+1)}=(-1)^m\binom{m+k}{k}
\)
得到
\(
[m^k]Q_k(m)=4^{-(k+1)}/k!=2^{-(2k+2)}/k!
\)。
上述分解与后续低阶闭式可由脚本 \texttt{scripts/exp\_fold\_zm\_low\_weight\_quasipolynomial\_audit.py} 复核。证毕。
\end{proof}

\begin{conclusion}[两项主部的显式化：多项式主增长与交替振荡的阶分离]\label{con:xi-terminal-zm-fixed-weight-two-term-asymptotic}
沿用结论 \ref{con:xi-terminal-zm-fixed-weight-generating-function-rigidity} 的记号，并记 \(\mathcal{P}_k(t)\in\ZZ[t]\) 为 \(A_k(t)\) 的分子。
则对一切 \(k\ge 0\) 有
\[
\boxed{\;
\mathcal{P}_k(1)=\mathcal{P}_k(-1)=1,\qquad \mathcal{P}_k'(1)=0\;}
\]
并且当 \(k\ge 1\) 时成立二阶与三阶导数的一次律：
\[
\boxed{\;
\mathcal{P}_k''(1)=4-12k,\qquad \mathcal{P}_k^{(3)}(1)=42-78k\; }.
\]
因此，对每个固定 \(k\)（令 \(m\to\infty\)）有两极点主部叠加的显式展开
\[
\boxed{\;
a_{m,k}
=\frac{1}{2^{k+1}}\binom{m+2k+1}{2k+1}
+\frac{k+1}{2^{k+2}}\binom{m+2k}{2k}
+O_k(m^{2k-1})
+(-1)^m\left(\frac{1}{2^{2k+2}}\binom{m+k}{k}+O_k(m^{k-1})\right)\; }.
\]
其中 \(O_k(m^{2k-1})\) 来自 \(t=1\) 的次高阶极点，而交替项来自 \(t=-1\) 的极点。
特别地，
\[
a_{m,k}
=\frac{m^{2k+1}}{2^{k+1}(2k+1)!}
+(-1)^m\frac{m^k}{2^{2k+2}k!}
+O_k(m^{2k}).
\]
\end{conclusion}
\begin{proof}[证明要点]
由结论 \ref{con:xi-terminal-zm-fixed-weight-generating-function-rigidity} 的递推
\[
\mathcal{P}_k(t)=(2t^2-t^4)\mathcal{P}_{k-1}(t)-t^4(1-t)^2(1+t)\mathcal{P}_{k-2}(t)
\]
在 \(t=\pm 1\) 处代入即可得 \(\mathcal{P}_k(\pm1)=\mathcal{P}_{k-1}(\pm1)=1\)。
对该递推在 \(t=1\) 处作导数并注意 \((2t^2-t^4)'|_{t=1}=0\) 且 \(t^4(1-t)^2(1+t)\) 在 \(t=1\) 处至少二阶消失，得到 \(\mathcal{P}_k'(1)=\mathcal{P}_{k-1}'(1)=0\)。
二阶与三阶导数律同理由在 \(t=1\) 处对递推作二次与三次微分并归纳得到（以 \(k=1,2\) 的直接核验作为起始）。
\par
最后由
\[
A_k(t)=\frac{\mathcal{P}_k(t)}{(1-t)^{2k+2}(1+t)^{k+1}}
\]
在 \(t=1\) 与 \(t=-1\) 处的局部展开作系数抽取：
在 \(t=1\) 处，
\(
(1+t)^{-(k+1)}=2^{-(k+1)}\bigl(1+\frac{k+1}{2}(1-t)+O((1-t)^2)\bigr)
\)
并结合 \(\mathcal{P}_k(1)=1\)、\(\mathcal{P}_k'(1)=0\) 得到所示两项主部；
在 \(t=-1\) 处，由 \(\mathcal{P}_k(-1)=1\) 得
\[
A_k(t)=\frac{1}{2^{2k+2}}(1+t)^{-(k+1)}+O\!\bigl((1+t)^{-k}\bigr)\qquad(t\to-1),
\]
从而
\(
[t^m]A_k(t)=(-1)^m\frac{1}{2^{2k+2}}\binom{m+k}{k}+(-1)^mO_k(m^{k-1})
\)。
证毕。
\end{proof}

\begin{conclusion}[固定重量层系数的最小常系数递推：阶 \texorpdfstring{$3k+3$}{3k+3} 与 \texorpdfstring{$\pm 1$}{±1} 极点分解]\label{con:xi-terminal-zm-fixed-weight-minimal-cc-recurrence-order}
固定 \(k\ge 0\)。
则序列 \(\{a_{m,k}\}_{m\ge 0}\) 在常系数递推意义下为 \(C\)-finite：
存在一个阶数为 \(3k+3\) 的齐次常系数线性递推，使得对所有充分大 \(m\) 成立；其特征多项式（以移位算子 \(E:a_m\mapsto a_{m+1}\) 记）为
\[
\boxed{\;(E-1)^{2k+2}(E+1)^{k+1}\;}
\]
并且在“充分大 \(m\)”意义下该阶数为最小：不存在阶数 \(<3k+3\) 的非平凡常系数递推能在全部充分大 \(m\) 上湮灭 \(\{a_{m,k}\}\)。
\end{conclusion}
\begin{proof}[证明要点]
由结论 \ref{con:xi-terminal-zm-fixed-weight-generating-function-rigidity}，
\(
A_k(t)=\sum_{m\ge 0}a_{m,k}t^m
\)
为有理函数，且其分母为
\(
Q_k(t)=(1-t)^{2k+2}(1+t)^{k+1}
\)
并满足 \(Q_k(0)=1\)。
由恒等式 \(Q_k(t)A_k(t)=\mathcal{P}_k(t)\) 逐项比较系数，得到由 \(Q_k\) 的系数给出的常系数线性递推；
当 \(m\) 足够大时右端系数归零，于是递推齐次化。
将 \(Q_k(t)\) 写成
\(
Q_k(t)=\sum_{j=0}^{3k+3}q_j t^j
\)
即可得到对应阶数为 \(3k+3\) 的递推；其特征多项式等价于
\(
E^{3k+3}Q_k(1/E)=(E-1)^{2k+2}(E+1)^{k+1}
\)。
\par
最小性归结为 \(\gcd(\mathcal{P}_k,Q_k)=1\)。
由于 \(Q_k\) 的全部不可约因子仅为 \((1-t)\) 与 \((1+t)\)，而结论 \ref{con:xi-terminal-zm-fixed-weight-two-term-asymptotic} 给出 \(\mathcal{P}_k(\pm1)=1\)，
故 \(\mathcal{P}_k\) 不被 \((1-t)\) 或 \((1+t)\) 整除，从而 \(Q_k\) 在既约意义下为最小湮灭分母（只差非零常数因子），递推阶数不可再降。证毕。
\end{proof}

\begin{conclusion}[低重量层系数的全量闭式：\(k\le 3\) 时对所有 \(m\ge 0\) 逐点成立]\label{con:xi-terminal-zm-fixed-weight-kle3-closed}
在结论 \ref{con:xi-terminal-zm-fixed-weight-quasipolynomial} 的记号下，令 \(n\ge 0\)。
则 \(k=0,1,2,3\) 的两条奇偶子序列 \(\{a_{2n,k}\}_{n\ge 0}\)、\(\{a_{2n+1,k}\}_{n\ge 0}\) 具有如下逐点闭式：
\[
a_{2n,0}=n+1,\qquad a_{2n+1,0}=n+1.
\]
\[
a_{2n,1}=\frac{1}{3}n^3+\frac{3}{2}n^2+\frac{1}{6}n,
\qquad
a_{2n+1,1}=\frac{1}{3}n^3+2n^2+\frac{5}{3}n+1.
\]
\[
a_{2n,2}=\frac{1}{30}n^5+\frac{3}{8}n^4-\frac{1}{12}n^3-\frac{7}{8}n^2+\frac{11}{20}n,
\]
\[
a_{2n+1,2}=\frac{1}{30}n^5+\frac{11}{24}n^4+\frac{3}{4}n^3-\frac{11}{24}n^2+\frac{13}{60}n.
\]
\[
a_{2n,3}=\frac{1}{630}n^7+\frac{1}{30}n^6+\frac{7}{360}n^5-\frac{17}{24}n^4+\frac{523}{360}n^3-\frac{33}{40}n^2+\frac{11}{420}n,
\]
\[
a_{2n+1,3}=\frac{1}{630}n^7+\frac{7}{180}n^6+\frac{23}{180}n^5-\frac{19}{36}n^4+\frac{29}{180}n^3+\frac{22}{45}n^2-\frac{61}{210}n.
\]
\end{conclusion}
\begin{proof}[证明要点]
对 \(k\le 3\)，由 \(A_k(t)=[y^k]F(t,y)\in\QQ(t)\) 且其极点仅位于 \(t=\pm1\) 可知
\(\{a_{2n,k}\}_{n\ge 0}\)、\(\{a_{2n+1,k}\}_{n\ge 0}\)
均为关于 \(n\) 的多项式序列（次数至多 \(2k+1\)）。
因此可用有限点值回收唯一多项式并与二维递推（结论 \ref{con:xi-terminal-zm-2d-recurrence-amk}）逐点核对；
亦可直接对 \(A_k(t)\) 作部分分式并回收系数。
上述闭式与其逐点一致性可由脚本 \texttt{scripts/exp\_fold\_zm\_low\_weight\_quasipolynomial\_audit.py} 复核。证毕。
\end{proof}



\paragraph{硬边界层的精确层级：二项式基整数证书与纯极点生成函数}
结论 \ref{con:xi-terminal-zm-2d-recurrence-amk} 的二维递推在硬边界方向（\(k\) 贴近 \(\lceil m/2\rceil\)）发生进一步的差分退化：
固定层深 \(j\) 后，偶长与奇长的边界系数在长度参数 \(k\) 上不仅具有渐近展开，
而是给出\emph{逐点成立}的精确多项式结构；并且其 Newton 二项式基展开自动产生可审计的整系数证书。

\begin{conclusion}[硬边界系数的精确多项式化：偶长 \texorpdfstring{$m=2k$}{m=2k} 的近最大重量层仅呈代数级增长]\label{con:xi-terminal-zm-hard-boundary-even-polynomial}
沿用结论 \ref{con:xi-terminal-zm-order4-rational} 的 \(Z_m(y)\) 记号。
对固定整数 \(j\ge 0\)，定义偶长硬边界系数
\[
A_j(k):=[y^{k-j}]\,Z_{2k}(y)\qquad(k\ge 0),
\]
其中约定当 \(k<j\) 时 \(A_j(k)=0\)。
则 \(A_j(k)\) 在所有整数 \(k\ge 0\) 上为精确多项式，其次数恰为 \(4j+2\)，且首项系数严格为 \(1/(4j+2)!\)。
更进一步，\(A_j(k)\) 存在唯一的二项式基展开
\[
A_j(k)=\sum_{r=j}^{4j+2}c_{j,r}\binom{k}{r},\qquad c_{j,r}\in\ZZ,\ \ c_{j,4j+2}=1,
\]
其中约定 \(\binom{k}{r}=0\)（当 \(r>k\)）。
因此在硬边界方向 \(k-j\) 处，推前重量剖分的每一固定层深 \(j\) 仅携带代数级前因子；
并且当 \(j\mapsto j+1\) 时，多项式次数严格增加 \(4\)。
\end{conclusion}

\begin{conclusion}[偶长硬边界系数的前三层闭式：二项式基展开给出完全可审计的整系数证书]\label{con:xi-terminal-zm-hard-boundary-even-first3}
在结论 \ref{con:xi-terminal-zm-hard-boundary-even-polynomial} 的记号下，前三个非平凡层（\(j=1,2,3\)）满足下列逐点恒等式（对所有整数 \(k\ge 0\) 成立）：
\[
\boxed{\ A_1(k)=2\binom{k}{1}+5\binom{k}{2}+9\binom{k}{3}+7\binom{k}{4}+4\binom{k}{5}+\binom{k}{6}\ }
\]
等价地，
\[
A_1(k)=\frac{k^6+9k^5+55k^4+435k^3-56k^2+996k}{720}.
\]
\[
\boxed{\ A_2(k)=3\binom{k}{2}+14\binom{k}{3}+39\binom{k}{4}+59\binom{k}{5}+63\binom{k}{6}+45\binom{k}{7}+22\binom{k}{8}+7\binom{k}{9}+\binom{k}{10}\ }
\]
\[
\boxed{\ A_3(k)=4\binom{k}{3}+30\binom{k}{4}+119\binom{k}{5}+273\binom{k}{6}+439\binom{k}{7}+507\binom{k}{8}+437\binom{k}{9}+283\binom{k}{10}+135\binom{k}{11}+46\binom{k}{12}+10\binom{k}{13}+\binom{k}{14}\ }.
\]
上述二项式基系数可由 forward-difference（Newton 级数）直接回收，因而给出无需谱计算的整数证书；
并可由脚本 \texttt{scripts/exp\_fold\_zm\_hard\_boundary\_coeffs\_audit.py} 一键复核。
\end{conclusion}

\begin{conclusion}[硬边界系数的精确多项式化：奇长 \texorpdfstring{$m=2k+1$}{m=2k+1} 的层次数严格量子化为 \texorpdfstring{$4j$}{4j}]\label{con:xi-terminal-zm-hard-boundary-odd-polynomial}
沿用结论 \ref{con:xi-terminal-zm-order4-rational} 的 \(Z_m(y)\) 记号。
对固定整数 \(j\ge 0\)，定义奇长硬边界系数
\[
B_j(k):=[y^{k+1-j}]\,Z_{2k+1}(y)\qquad(k\ge 0),
\]
其中约定当 \(k<j-1\) 时 \(B_j(k)=0\)。
则 \(B_j(k)\) 在所有整数 \(k\ge 0\) 上为精确多项式，其次数恰为 \(4j\)，且首项系数严格为 \(1/(4j)!\)。
并且存在唯一的二项式基展开
\[
B_j(k)=\sum_{r=j-1}^{4j}d_{j,r}\binom{k}{r},\qquad d_{j,r}\in\ZZ,\ \ d_{j,4j}=1,
\]
其中约定 \(\binom{k}{r}=0\)（当 \(r>k\)）。
在 \(j=1,2\) 时可给出完全闭式：
\[
\boxed{\ B_1(k)=\binom{k}{0}+4\binom{k}{1}+4\binom{k}{2}+3\binom{k}{3}+\binom{k}{4}
=\frac{k^4+6k^3+23k^2+66k+24}{24}\ }
\]
\[
\boxed{\ B_2(k)=2\binom{k}{1}+11\binom{k}{2}+23\binom{k}{3}+32\binom{k}{4}+28\binom{k}{5}+16\binom{k}{6}+6\binom{k}{7}+\binom{k}{8}\ }.
\]
\end{conclusion}

\begin{conclusion}[硬边界大偏差的闭式前因子：边界层概率质量的严格多项式—指数分解]\label{con:xi-terminal-zm-hard-boundary-ldp-prefactor}
令 \(w_m(x)=d_m(x)/2^m\) 为均匀微态推前分布，并令 \(X\sim w_m\)、\(H_m:=|X|_1\)。
对固定整数 \(j\ge 1\)，存在严格闭式：
\begin{itemize}
  \item 偶长 \(m=2k\)：
  \[
  \mathbb{P}(H_{2k}=k-j)=\frac{A_j(k)}{4^k},
  \]
  其中 \(A_j(k)\) 为结论 \ref{con:xi-terminal-zm-hard-boundary-even-polynomial} 的次数 \(4j+2\) 多项式。
  因而
  \[
  \mathbb{P}(H_{2k}=k-j)=\frac{k^{4j+2}}{(4j+2)!\,4^k}\Bigl(1+O(k^{-1})\Bigr).
  \]
  \item 奇长 \(m=2k+1\)：
  \[
  \mathbb{P}(H_{2k+1}=k+1-j)=\frac{B_j(k)}{2^{2k+1}},
  \]
  其中 \(B_j(k)\) 为结论 \ref{con:xi-terminal-zm-hard-boundary-odd-polynomial} 的次数 \(4j\) 多项式。
  因而
  \[
  \mathbb{P}(H_{2k+1}=k+1-j)=\frac{k^{4j}}{(4j)!\,2^{2k+1}}\Bigl(1+O(k^{-1})\Bigr).
  \]
\end{itemize}
上述分解表明：在最大密度边界附近，推前重量分布的指数速率保持为纯几何项 \(4^{-k}\)（或 \(2^{-(2k+1)}\)），
而全部非平凡算术结构被压缩到一个阶乘归一化的多项式前因子之中。
\end{conclusion}

\begin{conclusion}[硬边界层生成函数的纯极点结构：极点阶恰为 \texorpdfstring{$4j+3$}{4j+3} 与 \texorpdfstring{$4j+1$}{4j+1}]\label{con:xi-terminal-zm-hard-boundary-pure-pole}
令 \(A_j(k)\)、\(B_j(k)\) 如结论 \ref{con:xi-terminal-zm-hard-boundary-even-polynomial} 与结论 \ref{con:xi-terminal-zm-hard-boundary-odd-polynomial}，
并记其普通生成函数
\[
\mathcal A_j(t):=\sum_{k\ge 0}A_j(k)t^k,\qquad \mathcal B_j(t):=\sum_{k\ge 0}B_j(k)t^k.
\]
则二者均为有理函数，且极点阶严格量子化为
\[
\mathcal A_j(t)=\frac{P_j(t)}{(1-t)^{4j+3}},\qquad
\mathcal B_j(t)=\frac{Q_j(t)}{(1-t)^{4j+1}},
\]
其中 \(P_j,Q_j\in\ZZ[t]\)。
例如对 \(j=1\)（偶长第一硬边界层）有完全闭式
\[
\boxed{\ \mathcal A_1(t)=\frac{2t-5t^2+9t^3-10t^4+7t^5-2t^6}{(1-t)^7}\ }.
\]
因此硬边界层不仅在系数层面为精确多项式，其整生成函数亦呈现单一纯极点的最简解析结构。
\end{conclusion}


\begin{theorem}[Stokes 幅值三次与 Lee--Yang 三次的共同二维 Artin 表示]\label{thm:xi-terminal-zm-stokes-leyang-shared-artin-representation}
沿用定理 \ref{thm:xi-terminal-zm-b2-kappa2-rational-interconstruction} 的记号，记
\[
b:=B^2,
\qquad
u:=\kappa_*^2.
\]
再沿用结论 \ref{con:xi-terminal-zm-leyang-cubic-field-isomorphism} 的生成元 \(\alpha,\beta\)。
则有共同三次数域识别
\[
\boxed{\;
\QQ(b)=\QQ(u)=\QQ(\lambda_*)=\QQ(\alpha)=\QQ(\beta)=:K.\;}
\]
令 \(L\) 为 \(K/\QQ\) 的伽罗瓦闭包。
记 \(\rho_B\) 为由 Stokes 幅值三次
\[
162b^3+1593b^2+1998b+1184=0
\]
所定义的二维 Artin 表示；等价地，\(\rho_B\) 是 \(\Gal(L/\QQ)\cong S_3\) 的标准二维不可约表示沿自然商映射
\(\Gal(\overline\QQ/\QQ)\twoheadrightarrow \Gal(L/\QQ)\)
的拉回。
再记 \(\rho_{\mathrm{LY}}\) 为定理 \ref{thm:xi-time-part9c-leyang-cubic-artin-weight1} 中附着于 Lee--Yang 三次数域 \(K\) 的二维 Artin 表示。
则
\[
\boxed{\;\rho_B\cong \rho_{\mathrm{LY}}.\;}
\]
特别地，
\[
L(s,\rho_B)=L(s,\rho_{\mathrm{LY}}),
\qquad
\det(\rho_B)=\det(\rho_{\mathrm{LY}})=\chi_{-111},
\]
其中 \(\chi_{-111}\) 对应二次分辨子域 \(\QQ(\sqrt{-111})\)。
\leanverified{paper\_xi\_terminal\_zm\_stokes\_leyang\_shared\_artin\_representation}
\end{theorem}
\begin{proof}
由定理 \ref{thm:xi-terminal-zm-b2-kappa2-rational-interconstruction}，
\[
\QQ(b)=\QQ(u)=\QQ(\lambda_*).
\]
又由定理 \ref{thm:xi-terminal-zm-stokes-t-elliptic-branch-coordinate}，\(\lambda_*\) 满足
\[
16\lambda_*^3-9\lambda_*^2+1=0.
\]
结合结论 \ref{con:xi-terminal-zm-leyang-cubic-field-isomorphism} 的
\[
\QQ(\alpha)=\QQ(\beta),
\qquad
16\alpha^3-9\alpha^2+1=0,
\]
可知 \(\QQ(\lambda_*)=\QQ(\alpha)=\QQ(\beta)\)。
遂得第一条共同域识别。

再由定理 \ref{thm:xi-terminal-zm-stokes-amplitude-square-cubic-rigidity} 与定理 \ref{thm:xi-time-part9c-leyang-cubic-artin-weight1}，
Stokes 幅值三次与 Lee--Yang 三次都给出同一非伽罗瓦三次域 \(K\)，其伽罗瓦闭包均为 \(L\)，并满足
\[
\Gal(L/\QQ)\cong S_3.
\]
而 \(S_3\) 的二维不可约复表示只有一个同构类，即标准表示。
因此由 Stokes 幅值三次得到的 \(\rho_B\) 与 Lee--Yang 三次得到的 \(\rho_{\mathrm{LY}}\) 必同构：
\[
\rho_B\cong \rho_{\mathrm{LY}}.
\]

最后，定理 \ref{thm:xi-time-part9c-leyang-cubic-artin-weight1} 已给出
\[
\det(\rho_{\mathrm{LY}})=\chi_{-111}.
\]
另一方面，由推论 \ref{cor:xi-terminal-zm-stokes-b2-quadratic-resolvent}，Stokes 幅值三次的唯一二次分辨子域同样是 \(\QQ(\sqrt{-111})\)，故
\[
\det(\rho_B)=\chi_{-111}.
\]
表示同构即推出 Artin \(L\)-函数相同，得到
\[
L(s,\rho_B)=L(s,\rho_{\mathrm{LY}}).
\]
证毕。
\end{proof}

\begin{corollary}[Lee--Yang 三次数域的有理观测原始元刚性]\label{cor:xi-terminal-zm-leyang-cubic-rational-observable-primitivity}
\leanverified{paper\_xi\_terminal\_zm\_leyang\_cubic\_rational\_observable\_primitivity}
沿用定理 \ref{thm:xi-terminal-zm-stokes-leyang-shared-artin-representation} 的记号，令
\[
K:=\QQ(u)=\QQ(b)=\QQ(\lambda_*).
\]
则 \([K:\QQ]=3\)。若 \(R(T)\in\QQ(T)\) 在 \(u\) 处有定义且 \(R(u)\notin\QQ\)，则 \(R(u)\) 已为 \(K\) 的原始元，即
\[
\QQ(R(u))=K.
\]
同样地，只要 \(R(b)\) 或 \(R(\lambda_*)\) 有定义且不属于 \(\QQ\)，则相应元素亦为 \(K\) 的原始元，即
\[
\QQ(R(b))=K,
\qquad
\QQ(R(\lambda_*))=K.
\]
\end{corollary}
\begin{proof}
由定理 \ref{thm:xi-terminal-zm-kappa-square-cubic-field-s3}，\(u\) 满足不可约三次关系
\[
324u^3-540u^2+333u-74=0,
\]
故 \([\QQ(u):\QQ]=3\)。
再由定理 \ref{thm:xi-terminal-zm-stokes-leyang-shared-artin-representation} 的共同域识别，得 \([K:\QQ]=3\)。

若 \(R(u)\notin\QQ\)，则
\[
\QQ\subsetneq \QQ(R(u))\subseteq K.
\]
由于 \([K:\QQ]=3\) 为素数，\(K\) 不存在非平凡真中间域，故必有 \(\QQ(R(u))=K\)。
对 \(b\) 与 \(\lambda_*\) 的情形完全同理，因为 \(\QQ(b)=\QQ(\lambda_*)=K\)。证毕。
\end{proof}

\leanverified{paper\_xi\_terminal\_zm\_stokes\_leyang\_common\_quadratic\_resolvent}
\begin{corollary}[Stokes 幅值三次与 Lee--Yang 三次的共同二次分辨子域]\label{cor:xi-terminal-zm-stokes-leyang-common-quadratic-resolvent}
沿用上述记号，并令 \(L\) 为三次数域 \(K/\QQ\) 的伽罗瓦闭包。则 \(L/\QQ\) 的唯一二次子域为
\[
\boxed{\ \QQ(\sqrt{-111})\ }.
\]
等价地，\(u\) 与 \(b\) 的三次最小多项式
\[
324u^3-540u^2+333u-74,
\qquad
162b^3+1593b^2+1998b+1184
\]
的判别式分别为
\[
-2^6\cdot 3^9\cdot 37,
\qquad
-2^2\cdot 3^9\cdot 11^2\cdot 37\cdot 109^2,
\]
其平方自由部分同为
\[
-111=-3\cdot 37.
\]
\end{corollary}
\begin{proof}
由定理 \ref{thm:xi-terminal-zm-kappa-square-cubic-field-s3}，\(u\) 的极小多项式判别式为
\[
-2^6\cdot 3^9\cdot 37,
\]
故其分裂域的唯一二次子域为 \(\QQ(\sqrt{-111})\)。
另一方面，由推论 \ref{cor:xi-terminal-zm-stokes-b2-quadratic-resolvent}，\(b\) 的三次模型所对应的唯一二次分辨子域同样是 \(\QQ(\sqrt{-111})\)。

又由定理 \ref{thm:xi-terminal-zm-stokes-leyang-shared-artin-representation}，\(u\) 与 \(b\) 生成同一三次数域 \(K\)，故其伽罗瓦闭包一致，均等于 \(L\)。
因此 \(L/\QQ\) 的唯一二次子域只能是上述共同的 \(\QQ(\sqrt{-111})\)。
最后一条仅是把两判别式去除平方因子后的直接改写。证毕。
\end{proof}

\endinput

\paragraph{Lee--Yang 谱曲线在 \texorpdfstring{$u$}{u}-坐标下的统一化：分歧剖分、纤维群律与模性均匀化}
沿用
\[
\Pi(\lambda,y)=\lambda^{4}-\lambda^{3}-(2y+1)\lambda^{2}+\lambda+y(y+1),
\]
并记
\[
\Disc_{\lambda}\Pi=-y(y-1)\bigl(256y^{3}+411y^{2}+165y+32\bigr).
\]

\begin{conclusion}[Lee--Yang 谱曲线的 \texorpdfstring{$u$}{u}-坐标椭圆化与导子 \texorpdfstring{$37$}{37} 标识]\label{con:xi-terminal-zm-leyang-u-coordinate-birational-37a1}
设仿射曲线 \(C\subset\mathbb A^2_{\lambda,y}\) 由 \(\Pi(\lambda,y)=0\) 定义，并令
\[
u:=2y+1-2\lambda^{2}.
\]
则 \((\lambda,y)\mapsto(\lambda,u)\) 在 \(\QQ\) 上给出双有理同构
\[
C\ \dashrightarrow\ E_u:\quad u^{2}=4\lambda^{3}-4\lambda+1,
\]
其逆映射为
\[
(\lambda,u)\longmapsto\left(\lambda,\frac{u-1+2\lambda^2}{2}\right).
\]
进一步设 \(v:=(u-1)/2\)，则
\[
E_u\cong E_0:\quad v^2+v=\lambda^3-\lambda.
\]
因此该谱曲线归一化与导子 \(37\) 的模椭圆曲线同构（见定理 \ref{thm:xi-terminal-zm-leyang-elliptic-modularity-37}）。
\end{conclusion}
\begin{proof}[证明要点]
将 \(\Pi\) 视为关于 \(y\) 的二次式：
\[
\Pi(\lambda,y)=y^2+(1-2\lambda^2)y+\bigl(\lambda^4-\lambda^3-\lambda^2+\lambda\bigr).
\]
乘以 \(4\) 并配方，得
\[
(2y+1-2\lambda^2)^2=(1-2\lambda^2)^2-4(\lambda^4-\lambda^3-\lambda^2+\lambda)=1+4\lambda^3-4\lambda,
\]
即 \(u^2=4\lambda^3-4\lambda+1\)。由 \(u=2y+1-2\lambda^2\) 立得逆式。再令 \(u=2v+1\) 即得 \(v^2+v=\lambda^3-\lambda\)。证毕。
\end{proof}

\begin{conclusion}[分歧剖分与积分坐标下的临界五次方程]\label{con:xi-terminal-zm-leyang-u-coordinate-ramification-integral-quintic}
在结论 \ref{con:xi-terminal-zm-leyang-u-coordinate-birational-37a1} 的同构下，
\[
y=\frac{u-1+2\lambda^2}{2},\qquad y:E_u\to\PP^1
\]
为次数 \(4\) 的态射。其分歧满足：
\begin{enumerate}
  \item \(y\) 在无穷远点 \(O\) 处有唯一四阶极点，故 \(O\) 为 \(y=\infty\) 上的全分歧点（分歧指数 \(4\)）；
  \item 有限分歧值集合恰为
  \[
  \{0,1\}\cup\{\,\alpha\in\CC:\ 256\alpha^3+411\alpha^2+165\alpha+32=0\,\},
  \]
  且每个有限分歧值均对应简单分歧（分歧指数 \(2\)）；
  \item 在积分 Weierstrass 坐标
  \[
  X:=4\lambda,\qquad Y:=4u,\qquad Y^2=X^3-16X+16,
  \]
  有限临界点满足
  \[
  2XY+3X^2-16=0,
  \]
  消去 \(Y\) 得
  \[
  4X^5-9X^4-64X^3+160X^2-256=0.
  \]
\end{enumerate}
\end{conclusion}
\begin{proof}[证明要点]
由 \(y=(u-1+2\lambda^2)/2\) 得
\[
dy=\frac{du+4\lambda\,d\lambda}{2}.
\]
而 \(u^2=4\lambda^3-4\lambda+1\) 微分为 \(2u\,du=(12\lambda^2-4)\,d\lambda\)，故
\[
dy=\left(2\lambda+\frac{3\lambda^2-1}{u}\right)d\lambda.
\]
有限临界点由 \(dy=0\) 给出，即
\[
2\lambda u+3\lambda^2-1=0.
\]
换元 \(X=4\lambda,\ Y=4u\) 即得 \(2XY+3X^2-16=0\)。与 \(Y^2=X^3-16X+16\) 消元得到所示五次方程。
分歧值集合由 \(\Disc_\lambda\Pi=-y(y-1)P_{\mathrm{LY}}(y)\) 直接刻画，且有限分歧点总贡献与 Riemann--Hurwitz 公式一致，故均为简单分歧。证毕。
\end{proof}

\begin{conclusion}[非分歧纤维四点零和]\label{con:xi-terminal-zm-leyang-u-coordinate-fiber-fourpoint-zero-sum}
取 \(y_0\in\CC\) 避开分歧值集合，记纤维
\[
y^{-1}(y_0)=\{P_1,P_2,P_3,P_4\}\subset E_u(\CC).
\]
则椭圆曲线群律中恒有
\[
P_1+P_2+P_3+P_4=O.
\]
\end{conclusion}
\begin{proof}[证明要点]
定义
\[
f_{y_0}:=u+2\lambda^2-(2y_0+1).
\]
其零点恰为纤维四点，且在 \(O\) 处有四阶极点，从而
\[
\div(f_{y_0})=P_1+P_2+P_3+P_4-4O.
\]
主除子在 \(\operatorname{Pic}^0(E_u)\cong E_u\) 中对应零元，即得结论。
该陈述与结论 \ref{con:xi-terminal-zm-elliptic-weight-fiber-sum-zero} 一致。证毕。
\end{proof}

\begin{conclusion}[Néron 微分与导子 \texorpdfstring{$37$}{37} 的模性实现]\label{con:xi-terminal-zm-leyang-u-coordinate-neron-modular-pullback}
在 \(E_u:u^2=4\lambda^3-4\lambda+1\) 上定义
\[
\omega:=\frac{d\lambda}{u}.
\]
则 \(\omega\) 为 \(E_u\) 的 Néron 不变微分。并且存在权 \(2\)、水平 \(37\) 的归一化新形式
\(f_{37}\in S_2(\Gamma_0(37))\) 及模参数化
\[
\varphi:X_0(37)\to E_u
\]
使得
\[
\varphi^*\omega=c_\varphi\,2\pi i\,f_{37}(\tau)\,d\tau,\qquad c_\varphi\in\ZZ_{>0}.
\]
对最优参数化可取 \(c_\varphi=1\)。
\end{conclusion}
\begin{proof}[证明要点]
\(u^2=4\lambda^3-4\lambda+1\) 为 Weierstrass 形式，\(d\lambda/u\) 即其标准不变微分。
由模性定理，导子 \(37\) 的椭圆曲线对应 \(S_2(\Gamma_0(37))\) 新形式，且模参数化拉回 Néron 微分为新形式乘 \(2\pi i\,d\tau\)（带 Manin 常数）。证毕。
\end{proof}

\begin{conclusion}[Weierstrass 均匀化与 Lee--Yang 分支的椭圆函数参数化]\label{con:xi-terminal-zm-leyang-u-coordinate-weierstrass-uniformization}
存在格 \(\Lambda\subset\CC\) 使
\[
\lambda(z)=\wp(z;\,g_2=4,g_3=-1),\qquad
u(z)=\wp'(z;\,g_2=4,g_3=-1),
\]
并定义
\[
y(z)=\frac{\wp'(z)-1+2\wp(z)^2}{2}.
\]
则对一切 \(z\) 恒有
\[
\Pi(\lambda(z),y(z))=0.
\]
因此 \(\lambda\) 作为 \(y\) 的代数多值反演可在 \(\CC/\Lambda\) 上由平移作用实现单值化；临界值集合由 \(dy(z)=0\) 给出，恰为
\[
y\in\{0,1\}\cup\{P_{\mathrm{LY}}(y)=0\}.
\]
\end{conclusion}
\begin{proof}[证明要点]
由 \((\wp')^2=4\wp^3-4\wp+1\) 立得 \((\lambda,u)\in E_u\)，再代入
\(
y=(u-1+2\lambda^2)/2
\)
即得 \(\Pi(\lambda(z),y(z))=0\)。
临界值陈述由结论 \ref{con:xi-terminal-zm-leyang-u-coordinate-ramification-integral-quintic} 的 \(dy=0\) 判据与判别式分解一致性得到。证毕。
\end{proof}

\paragraph{\texorpdfstring{$S_4$}{S4} 闭包下的 Chevalley--Weil 重数刚性、Jacobian 群代数分解与局部算术证书}
沿用结论 \ref{con:xi-terminal-zm-s4-galois-closure-genus-tower-holomorphic-differentials} 的设定，
记 \(\widetilde{X}\to \PP^1_y\) 为四次谱覆盖 \(\Pi(\lambda,y)=0\) 的 \(S_4\)-伽罗瓦闭包曲线，且 \(g(\widetilde{X})=16\)。
并记
\[
F(\lambda,y):=\Pi(\lambda,y)
=\lambda^{4}-\lambda^{3}-(2y+1)\lambda^{2}+\lambda+y(y+1).
\]
记
\[
V:=H^0(\widetilde{X},\Omega^1_{\widetilde{X}}),
\qquad
\varepsilon:=\mathrm{sgn},\ \rho_2:=V_2,\ \rho_3:=V_3,\ \rho_3':=V_3'.
\]

\begin{conclusion}[\texorpdfstring{$S_4$}{S4} 闭包覆盖的 Hurwitz 型与子群商亏格谱]\label{con:xi-terminal-zm-s4-hurwitz-subgroup-spectrum-cycletypes}
设 \(\widetilde{X}\to\PP^1_y\) 为 \(S_4\)-伽罗瓦覆盖，其分歧惯性阶为
\[
(4,2,2,2,2,2),
\]
即 Hurwitz 型 \((0;4,2,2,2,2,2)\)。由 Riemann--Hurwitz 公式可得 \(g(\widetilde{X})=16\)，并且
\[
g(\widetilde{X}/S_3)=1,\quad
g(\widetilde{X}/A_4)=2,\quad
g(\widetilde{X}/V_4^{\mathrm{norm}})=4,\quad
g(\widetilde{X}/D_8)=1,\quad
g(\widetilde{X}/C_3)=6.
\]
进一步设 \(\tau=(12)\)、\(\rho=(1234)\in S_4\)。
对任意子群 \(H\le S_4\)，记
\[
d:=[S_4:H],\qquad
c_\tau(H):=\#\text{\(\tau\) 在 }S_4/H\text{ 上的循环个数},\qquad
c_\rho(H):=\#\text{\(\rho\) 在 }S_4/H\text{ 上的循环个数}.
\]
则有统一亏格公式
\[
g(\widetilde{X}/H)
=1-d+\frac12\Bigl(5\bigl(d-c_\tau(H)\bigr)+\bigl(d-c_\rho(H)\bigr)\Bigr).
\]
并且在各共轭类代表子群上，可得完整枚举
\[
\begin{array}{c|c|c|c|c}
H & d & \tau\text{ 在 }S_4/H\text{ 上的循环型} & \rho\text{ 在 }S_4/H\text{ 上的循环型} & g(\widetilde{X}/H)\\
\hline
S_4 & 1 & 1 & 1 & 0\\
A_4 & 2 & 2 & 2 & 2\\
S_3 & 4 & 2\cdot 1^2 & 4 & 1\\
D_8 & 3 & 2\cdot 1 & 2\cdot 1 & 1\\
V_4^{\mathrm{norm}} & 6 & 2^3 & 2^3 & 4\\
V_4^{\mathrm{pair}}=\langle(12),(34)\rangle & 6 & 2^2\cdot 1^2 & 4\cdot 2 & 2\\
C_4 & 6 & 2^3 & 4\cdot 1^2 & 4\\
C_3 & 8 & 2^4 & 4^2 & 6\\
C_2^{\mathrm{tr}} & 12 & 2^5\cdot 1^2 & 4^3 & 6\\
C_2^{\mathrm{dt}} & 12 & 2^6 & 4^2\cdot 2^2 & 8\\
\{1\} & 24 & 2^{12} & 4^6 & 16
\end{array}
\]
其中 \(D_8\) 表示阶 \(8\) 二面体子群，\(V_4^{\mathrm{norm}}\) 为正规 Klein 四元，\(V_4^{\mathrm{pair}}\) 为两两换位 Klein 四元。
\end{conclusion}
\begin{proof}[证明要点]
对 \(\widetilde{X}\to\PP^1_y\) 应用 Riemann--Hurwitz：
\[
2g(\widetilde{X})-2
=|S_4|(-2)+|S_4|\Bigl(1-\frac14\Bigr)+5|S_4|\Bigl(1-\frac12\Bigr)=30,
\]
故 \(g(\widetilde{X})=16\)。

对任意 \(H\le S_4\)，商覆盖 \(\widetilde{X}/H\to\PP^1_y\) 的总分歧由 \(\tau,\rho\) 在 \(S_4/H\) 上的循环分解给出，分别为 \(d-c_\tau(H)\)、\(d-c_\rho(H)\)。
代入 Riemann--Hurwitz 即得统一亏格公式与上表。证毕。
\end{proof}

\begin{conclusion}[指数 \texorpdfstring{$6$}{6} 配对稳定子商曲线的显式模型]\label{con:xi-terminal-zm-s4-v4pair-explicit-model}
在代数闭包中记 \(F(\lambda,y)=0\) 的四根为 \(\lambda_1,\lambda_2,\lambda_3,\lambda_4\)，并取
\[
V_4^{\mathrm{pair}}=\langle(12),(34)\rangle\le S_4.
\]
定义
\[
v:=\lambda_1\lambda_2,\quad w:=\lambda_3\lambda_4,\quad z:=v+w,\quad s:=v-w.
\]
则固定域可写为
\[
\QQ(\widetilde{X}/V_4^{\mathrm{pair}})=\QQ(y,z,s),
\]
并满足完全显式关系
\[
\mathscr R(z,y)=0,\qquad s^2=z^2-4y(y+1),
\]
其中
\[
\mathscr R(z,y)=z^3+(1+2y)z^2-(1+4y+4y^2)z-(1+5y+13y^2+8y^3).
\]
记 \(E_{\mathrm{res}}\) 为 \(\mathscr R(z,y)=0\) 的光滑射影模型，\(C_{\mathrm{pair}}\) 为上述二元方程组的光滑射影模型，则
\[
\QQ(C_{\mathrm{pair}})=\QQ(\widetilde{X}/V_4^{\mathrm{pair}}),\qquad
g(C_{\mathrm{pair}})=2,
\]
且自然投影 \(C_{\mathrm{pair}}\to E_{\mathrm{res}}\) 为次数 \(2\) 覆盖。
\end{conclusion}
\begin{proof}[证明要点]
\(z=v+w\) 为配对 \(\{12|34\}\) 的不变量，故满足三次预解式 \(\mathscr R(z,y)=0\)。
又由 \(vw=\lambda_1\lambda_2\lambda_3\lambda_4=y(y+1)\)，得
\[
s^2=(v-w)^2=(v+w)^2-4vw=z^2-4y(y+1).
\]
于是 \(\QQ(y,z,s)\subseteq \QQ(\widetilde{X}/V_4^{\mathrm{pair}})\)。
反向由 \(v=(z+s)/2,\ w=(z-s)/2\) 可恢复配对乘积信息，故函数域相等。
亏格与覆盖次数由结论 \ref{con:xi-terminal-zm-s4-hurwitz-subgroup-spectrum-cycletypes} 中 \(H=V_4^{\mathrm{pair}}\) 的条目给出。证毕。
\end{proof}

\begin{conclusion}[黄金比率分歧：\texorpdfstring{$C_{\mathrm{pair}}\to E_{\mathrm{res}}$}{Cpair to Eres} 的分歧除子由 \texorpdfstring{$y^2+y-1$}{y2+y-1} 唯一刻画]\label{con:xi-terminal-zm-s4-v4pair-golden-ratio-branch-divisor}
在结论 \ref{con:xi-terminal-zm-s4-v4pair-explicit-model} 的记号下，令
\[
f:=z^2-4y(y+1)\in \QQ(E_{\mathrm{res}})^\times.
\]
则有严格消去恒等式
\[
\operatorname{Res}_z\!\bigl(\mathscr R(z,y),\,z^2-4y(y+1)\bigr)=(y^2+y-1)^2.
\]
因此 \(f=0\) 与 \(\mathscr R=0\) 的公共点恰为
\[
P_\pm:\quad
(y,z)=\Bigl(-\frac12\pm\frac{\sqrt5}{2},\,-2\Bigr),
\]
并且二次覆盖 \(C_{\mathrm{pair}}\to E_{\mathrm{res}}\) 的分歧除子为 \(P_++P_-\)。
从而
\[
2g(C_{\mathrm{pair}})-2=\deg(P_++P_-)=2,\qquad g(C_{\mathrm{pair}})=2,
\]
且分歧点由唯一二次域 \(\QQ(\sqrt5)\) 控制。
\end{conclusion}
\begin{proof}[证明要点]
结果式计算给出仅当 \(y^2+y-1=0\) 时可能有公共点。
将该二次方程两根代回 \(\mathscr R(z,y)=0\) 与 \(z^2-4y(y+1)=0\) 可唯一得到 \(z=-2\)。
对二次覆盖 \(s^2=f\)，分歧由 \(f\) 的奇阶零点给出，故分歧除子正为 \(P_++P_-\)。
再由 Riemann--Hurwitz 公式即得亏格。证毕。
\end{proof}

\begin{conclusion}[Chevalley--Weil 型 \texorpdfstring{$S_4$}{S4} 表示分解的显式重数解]\label{con:xi-terminal-zm-s4-chevalley-weil-explicit-decomposition}
在上述记号下，\(V\) 作为 \(S_4\)-表示满足
\[
V\cong 2\varepsilon\oplus \rho_2\oplus \rho_3\oplus 3\rho_3'.
\]
等价地，\(\dim V=16\) 的不可约重数向量唯一为
\[
(m_{\varepsilon},m_{\rho_2},m_{\rho_3},m_{\rho_3'})=(2,1,1,3).
\]
\end{conclusion}
\begin{proof}[证明要点]
由结论 \ref{con:xi-terminal-zm-s4-hurwitz-subgroup-spectrum-cycletypes}，
\[
g(\widetilde{X}/S_4)=0,\quad
g(\widetilde{X}/A_4)=2,\quad
g(\widetilde{X}/S_3)=1,\quad
g(\widetilde{X}/D_8)=1.
\]
对任意 \(K\le S_4\) 有
\[
\dim V^K=g(\widetilde{X}/K).
\]
设
\[
V\cong a\mathbf 1\oplus b\varepsilon\oplus c\rho_2\oplus d\rho_3\oplus e\rho_3'.
\]
由 \(\dim V^{S_4}=0\) 得 \(a=0\)。

再由角色平均公式可得：
在 \(A_4\) 上仅 \(\varepsilon\) 贡献一维不变元，故 \(b=2\)；
在 \(S_3\) 上仅 \(\rho_3\) 贡献一维不变元，故 \(d=1\)；
在 \(D_8\) 上仅 \(\rho_2\) 贡献一维不变元，故 \(c=1\)。

最后由总维数
\[
16=\dim V=2b+2c+3d+3e
=2\cdot2+2\cdot1+3\cdot1+3e
\]
解得 \(e=3\)。
因此
\[
V\cong 2\varepsilon\oplus \rho_2\oplus \rho_3\oplus 3\rho_3',
\]
且分解唯一。证毕。
\end{proof}

\begin{conclusion}[Jacobian 的 \texorpdfstring{$S_4$}{S4} 等变同源分解与可识别等型因子]\label{con:xi-terminal-zm-s4-jacobian-group-algebra-prym-e3-b3}
记
\[
J:=J(\widetilde{X}),\qquad
J_{\mathrm{sgn}}:=J(\widetilde{X}/A_4),\qquad
E:=\widetilde{X}/S_3,\qquad
E_{\mathrm{res}}:=\widetilde{X}/D_8.
\]
则存在定义在 \(\QQ\) 上的三维阿贝尔簇 \(A,B\)，使得
\[
J\sim_{\QQ}J_{\mathrm{sgn}}\times J_{\rho_2}\times A\times B^3,
\]
其中
\[
J_{\mathrm{sgn}}\sim_{\QQ}J(\widetilde{X}/A_4),\qquad
J_{\rho_2}\sim_{\QQ}J(\widetilde{X}/D_8)^2=E_{\mathrm{res}}^2,
\]
并且 \(A\) 对应 \(\rho_3\)-等型（可取 \(A\sim_{\QQ}E^3\)），\(B\) 对应 \(\rho_3'\)-等型。

等价地，可写为
\[
J\sim_{\QQ}J(\widetilde{X}/A_4)\times \operatorname{Prym}\!\bigl((\widetilde{X}/V_4^{\mathrm{norm}})/(\widetilde{X}/A_4)\bigr)\times E^{3}\times B^{3}.
\]
并且有中间 Jacobian 分拆
\[
J(\widetilde{X}/V_4^{\mathrm{norm}})\sim_{\QQ}J_{\mathrm{sgn}}\times E_{\mathrm{res}}^{2},
\]
\[
J(C_{\mathrm{pair}})\sim_{\QQ}E_{\mathrm{res}}\times E,\qquad
C_{\mathrm{pair}}:=\widetilde{X}/V_4^{\mathrm{pair}},
\]
\[
J(\widetilde{X}/C_4)\sim_{\QQ}E_{\mathrm{res}}\times B.
\]
\end{conclusion}
\begin{proof}[证明要点]
由
\[
\QQ[S_4]\cong \QQ\oplus \QQ\oplus M_2(\QQ)\oplus M_3(\QQ)\oplus M_3(\QQ)
\]
与结论 \ref{con:xi-terminal-zm-s4-chevalley-weil-explicit-decomposition}，
\(J\) 在 \(\QQ\)-同源意义下按 \(\varepsilon,\rho_2,\rho_3,\rho_3'\) 四个等型分解，维数分别为 \(2,2,3,9\)。

\(\varepsilon\)-等型由 \(A_4\) 不变微分识别为 \(J(\widetilde{X}/A_4)\)；
\(\rho_2\)-等型由 \(D_8\) 商与结论 \ref{con:xi-terminal-zm-s4-prym-c6-over-c-d8-elliptic-square} 识别为 \(E_{\mathrm{res}}^2\)；
\(\rho_3\)-等型可由 \(S_3\) 商识别为 \(E^3\)；
\(\rho_3'\)-等型重数为 \(3\)，记为 \(B^3\)。

对中间子群 \(V_4^{\mathrm{norm}},V_4^{\mathrm{pair}},C_4\)，利用
\[
H^0(\widetilde{X}/H,\Omega^1)\cong H^0(\widetilde{X},\Omega^1)^H
\]
与角色平均公式，结合结论 \ref{con:xi-terminal-zm-s4-hurwitz-subgroup-spectrum-cycletypes} 的亏格数据，即得所列三条中间 Jacobian 同源分拆。证毕。
\end{proof}

\begin{conjecture}[\texorpdfstring{$\rho_3'$}{rho3'}-三维因子 \texorpdfstring{$B$}{B} 的绝对单纯性与端同态刚性]\label{conj:xi-terminal-zm-s4-b-absolute-simplicity-endomorphism}
在结论 \ref{con:xi-terminal-zm-s4-jacobian-group-algebra-prym-e3-b3} 的记号下，
设 \(B\) 为 \(\rho_3'\)-等型三维构件。
则 \(B\) 在 \(\overline{\QQ}\) 上绝对单纯，且
\[
\End^0_{\overline{\QQ}}(B)\cong \QQ.
\]
等价地，除由 \(\QQ[S_4]\) 的中心化子所强制的结构外，\(B\) 不携带额外的代数自同态。
\end{conjecture}

\begin{conclusion}[\texorpdfstring{$D_8$}{D8} 商椭圆曲线与二维 Prym 的椭圆平方化]\label{con:xi-terminal-zm-s4-prym-c6-over-c-d8-elliptic-square}
设 \(D_8\le S_4\) 为任一 \(4\)-循环的正规化子（阶 \(8\) 二面体子群），并定义
\[
E_8:=\widetilde{X}/D_8.
\]
并记
\[
C:=\widetilde{X}/A_4,\qquad C_6:=\widetilde{X}/V_4^{\mathrm{norm}}.
\]
则 \(g(E_8)=1\)，并且
\[
\operatorname{Prym}(C_6/C)\sim_{\QQ}E_8^2.
\]
从而
\[
\mathrm{End}^0_{\QQ}\!\bigl(\operatorname{Prym}(C_6/C)\bigr)\supset M_2(\QQ).
\]
在无分歧循环三次覆盖情形下，\(\operatorname{Prym}(C_6/C)\) 的极化型为 \((1,3)\)。
\end{conclusion}
\begin{proof}[证明要点]
由结论 \ref{con:xi-terminal-zm-s4-galois-closure-genus-tower-holomorphic-differentials}，\(g(\widetilde{X}/D_8)=1\)。
由结论 \ref{con:xi-terminal-zm-s4-jacobian-a2-isog-r2}，\(\rho_2\)-等型分量满足 \(A_2\sim R^2\)，其中 \(R=\widetilde{X}/D_8\)。
另一方面，定理 \ref{thm:xi-terminal-zm-s3-closure-jacobian-splitting} 识别 \(A_2\sim \operatorname{Prym}(C_6/C)\)。
合并即得 \(\operatorname{Prym}(C_6/C)\sim_{\QQ}E_8^2\)。
极化型由定理 \ref{thm:xi-terminal-zm-prym-polarization-a2-rigidity} 给出。证毕。
\end{proof}

\begin{conclusion}[好素数处的局部 \texorpdfstring{$L$}{L} 因子强分解与模 \texorpdfstring{$3$}{3} 点数同余证书]\label{con:xi-terminal-zm-s4-local-factorization-mod3-congruence-certificate}
设 \(p\) 为 \(\widetilde{X}\) 的好约化素数，并令 \(q=p^m\)（\(m\ge 1\)）。
沿用结论 \ref{con:xi-terminal-zm-s4-jacobian-group-algebra-prym-e3-b3} 与
\ref{con:xi-terminal-zm-s4-prym-c6-over-c-d8-elliptic-square} 的记号，并记
\[
C:=\widetilde{X}/A_4.
\]
则局部因子满足
\[
P_p(\widetilde{X},T)
=P_p(C,T)\,P_p(E,T)^3\,P_p(E_8,T)^2\,P_p(B,T)^3.
\]
进一步，对有限域点数记号 \(N_Y(q):=\#Y(\FF_q)\)，有同余
\[
N_{\widetilde{X}}(q)\equiv N_C(q)+2N_{E_8}(q)-2(q+1)\pmod 3.
\]
\end{conclusion}
\begin{proof}[证明要点]
\(\QQ\)-同源分解在好约化处与 \(\ell\)-进 \(H^1\) 分解相容，故 Frobenius 特征多项式按因子相乘，得到 \(P_p\) 强分解。
又
\[
N_Y(q)=q+1-\mathrm{Tr}\!\left(\mathrm{Frob}_q\mid H^1_{\acute et}(Y_{\overline{\FF}_q},\QQ_\ell)\right).
\]
由分解式可得
\[
\mathrm{Tr}_q(\widetilde{X})
\equiv \mathrm{Tr}_q(C)+2\mathrm{Tr}_q(E_8)\pmod 3,
\]
因为 \(3\mathrm{Tr}_q(E)\) 与 \(3\mathrm{Tr}_q(B)\) 在模 \(3\) 下消失。
代回点数公式即得所示同余。证毕。
\end{proof}

\paragraph{\texorpdfstring{$S_4$}{S4} 作用曲线的低次数映射、典范二次关系与 \texorpdfstring{$n$}{n}-典范表示}
以下统一记 \(X:=\widetilde{X}\)。
则 \(X/S_4\simeq \PP^1_y\) 的分支型为一个阶 \(4\) 分支点与五个阶 \(2\) 分支点，且 \(g(X)=16\)。

\begin{conclusion}[Castelnuovo--Severi 阻断下的低次数可约性消失]\label{con:xi-terminal-zm-s4-low-gonality-vanishing-castelnuovo-severi}
曲线 \(X\) 不存在次数为 \(2\) 或 \(3\) 的非常值态射 \(X\to \PP^1\)。
因此 \(X\) 非双曲线且非三折曲线。
\end{conclusion}
\begin{proof}[证明要点]
由结论 \ref{con:xi-terminal-zm-s4-galois-closure-genus-tower-holomorphic-differentials}，
\(E:=X/S_3\) 为亏格 \(1\) 曲线，且自然商映射 \(\pi:X\to E\) 的次数为 \(6\)。
设 \(\varphi:X\to \PP^1\) 为次数 \(d\in\{2,3\}\) 的非常值态射。
若 \(\varphi=\psi\circ\pi\)，则 \(d=6\deg(\psi)\)，矛盾。
故 \(\pi,\varphi\) 在 Castelnuovo--Severi 意义下相互独立，满足
\[
g(X)\le d\,g(E)+6\,g(\PP^1)+(d-1)(6-1)=6d-5.
\]
当 \(d=2\) 时右端为 \(7\)，当 \(d=3\) 时右端为 \(13\)，均与 \(g(X)=16\) 矛盾。
\end{proof}

\begin{conclusion}[典范模型的射影正规性与二次生成]\label{con:xi-terminal-zm-s4-canonical-projective-normality-quadratic-generation}
典范嵌入 \(X\hookrightarrow \PP^{15}\) 射影正规，且其齐次理想由二次方程生成。
等价地，
\[
\mu_2:\operatorname{Sym}^2 H^0(X,K_X)\twoheadrightarrow H^0(X,2K_X)
\]
为满射，并且
\[
\dim\ker(\mu_2)=91.
\]
\end{conclusion}
\begin{proof}[证明要点]
由结论 \ref{con:xi-terminal-zm-s4-low-gonality-vanishing-castelnuovo-severi}，\(X\) 非双曲线。
Max Noether 定理给出典范嵌入的射影正规性，特别地 \(\mu_2\) 满射。
又由同一结论知 \(X\) 非三折；且 \(g(X)=16\neq 6\)，故 \(X\) 不可能为平面五次曲线。
由 Petri 定理可得典范理想由二次方程生成。
再由 Riemann--Roch，
\[
h^0(X,2K_X)=3g-3=45,\qquad
\dim\operatorname{Sym}^2 H^0(X,K_X)=\binom{16+1}{2}=136,
\]
故 \(\dim\ker(\mu_2)=136-45=91\)。
\end{proof}

\begin{conclusion}[\texorpdfstring{$S_4$}{S4} 作用在 \texorpdfstring{$X=\widetilde{X}$}{X} 上的固定点数完全表]\label{con:xi-terminal-zm-s4-conjugacy-fixedpoints-local-eigenvalue-full-count}
设 \(g\in S_4\)。
则 \(|\mathrm{Fix}(g)|\) 仅依赖于 \(g\) 的共轭类，并满足
\[
|\mathrm{Fix}(g)|=
\begin{cases}
10,& g\text{ 为换位 }(2),\\
2,& g\text{ 为双换位 }(2,2),\\
0,& g\text{ 为三循环 }(3),\\
2,& g\text{ 为四循环 }(4).
\end{cases}
\]
此外，阶 \(2\) 元在定点处切空间特征值为 \(-1\)；
阶 \(4\) 元的两个定点切空间特征值分别为 \(i\) 与 \(-i\)。
\end{conclusion}
\begin{proof}[证明要点]
由结论 \ref{con:xi-terminal-zm-s4-hurwitz-subgroup-spectrum-cycletypes}，\(X\to\PP^1_y\) 的惯性只可能为换位（五个有限分歧值）或四循环（无穷远分歧值）。
若 \(g\) 为三循环，则其阶为 \(3\)，不可能包含于任一惯性子群，故无固定点。

若 \(g\) 为换位，则固定点仅可能来自五个有限分歧纤维。每个此类分歧值上，惯性群同构于 \(\langle\tau\rangle\cong C_2\)。给定换位 \(g\) 的固定点数为
\[
\frac{|C_{S_4}(g)|}{|\langle\tau\rangle|}=\frac{4}{2}=2,
\]
故总数为 \(5\cdot2=10\)。

若 \(g\) 为四循环，则固定点仅来自无穷远分歧纤维。设其惯性为 \(I=\langle c\rangle\cong C_4\)。固定点条件等价于 \(g\in hIh^{-1}\)，即 \(hch^{-1}=g\) 或 \(g^{-1}\)，从而
\[
|\mathrm{Fix}(g)|=\frac{|C_{S_4}(c)|+|C_{S_4}(c)|}{|I|}=\frac{4+4}{4}=2.
\]

若 \(g\) 为双换位，则只能来自无穷远纤维。在 \(I=\langle c\rangle\) 中唯一非平凡二阶元为 \(c^2\)，故
\[
|\mathrm{Fix}(g)|=\frac{|C_{S_4}(c^2)|}{|I|}=\frac{8}{4}=2.
\]
局部特征值方面，阶 \(2\) 元在定点处导数为 \(-1\)。
对阶 \(4\) 元，由结论 \ref{con:xi-terminal-zm-s4-chevalley-weil-explicit-decomposition} 的
\(\mathrm{tr}((1234)\mid H^0(X,K_X))=0\)。
并与全纯 Lefschetz 公式
\[
\mathrm{tr}\!\bigl(g\mid H^0(X,K_X)\bigr)
=1+\sum_{P\in \mathrm{Fix}(g)}\frac{\zeta_P}{1-\zeta_P},
\]
可知两定点贡献相消，故特征值为 \(i\) 与 \(-i\)。证毕。
\end{proof}

\begin{conclusion}[全体 \texorpdfstring{$n$}{n}-典范表示的显式不可约分解]\label{con:xi-terminal-zm-s4-pluricanonical-irrep-decomposition-all-n}
设 \(n\ge 2\)。
记 \(\mathbf 1,\varepsilon,\rho_2,\rho_3,\rho_3'\) 分别为 \(S_4\) 的平凡、符号、二维不可约、三维标准及其符号扭曲表示。
则 \(H^0(X,nK_X)\) 的不可约分解如下：
\[
\text{若 }n\equiv 0\!\!\pmod 4,\quad
H^0(X,nK_X)\cong
\frac{5n+4}{4}\mathbf 1\oplus
\frac{5n-8}{4}\varepsilon\oplus
\frac{5n-2}{2}\rho_2\oplus
\frac{15n-4}{4}\rho_3\oplus
\frac{15n-12}{4}\rho_3'.
\]
\[
\text{若 }n\equiv 1\!\!\pmod 4,\quad
H^0(X,nK_X)\cong
\frac{5n-9}{4}\mathbf 1\oplus
\frac{5n+3}{4}\varepsilon\oplus
\frac{5n-3}{2}\rho_2\oplus
\frac{15n-11}{4}\rho_3\oplus
\frac{15n-3}{4}\rho_3'.
\]
\[
\text{若 }n\equiv 2\!\!\pmod 4,\quad
H^0(X,nK_X)\cong
\frac{5n+2}{4}\mathbf 1\oplus
\frac{5n-6}{4}\varepsilon\oplus
\frac{5n-2}{2}\rho_2\oplus
\frac{15n-2}{4}\rho_3\oplus
\frac{15n-14}{4}\rho_3'.
\]
\[
\text{若 }n\equiv 3\!\!\pmod 4,\quad
H^0(X,nK_X)\cong
\frac{5n-7}{4}\mathbf 1\oplus
\frac{5n+1}{4}\varepsilon\oplus
\frac{5n-3}{2}\rho_2\oplus
\frac{15n-13}{4}\rho_3\oplus
\frac{15n-1}{4}\rho_3'.
\]
特别地，平凡表示重数满足
\[
\dim H^0(X,nK_X)^{S_4}
=1-2n+5\Bigl\lfloor\frac n2\Bigr\rfloor+\Bigl\lfloor\frac{3n}{4}\Bigr\rfloor.
\]
\end{conclusion}
\begin{proof}[证明要点]
当 \(n\ge 2\) 时 \(H^1(X,nK_X)=0\)。
由结论 \ref{con:xi-terminal-zm-s4-conjugacy-fixedpoints-local-eigenvalue-full-count} 的定点与局部特征值数据，全纯 Lefschetz 公式给出
\[
\mathrm{tr}\!\bigl((12)\mid H^0(X,nK_X)\bigr)=5(-1)^n,\qquad
\mathrm{tr}\!\bigl((12)(34)\mid H^0(X,nK_X)\bigr)=(-1)^n,
\]
\[
\mathrm{tr}\!\bigl((123)\mid H^0(X,nK_X)\bigr)=0,
\]
以及
\[
\mathrm{tr}\!\bigl((1234)\mid H^0(X,nK_X)\bigr)=
\begin{cases}
1,& n\equiv 0,3\pmod 4,\\
-1,& n\equiv 1,2\pmod 4.
\end{cases}
\]
另一方面
\[
\mathrm{tr}\!\bigl(1\mid H^0(X,nK_X)\bigr)=h^0(X,nK_X)=30n-15.
\]
将上述五个共轭类迹代入 \(S_4\) 特征标正交关系，逐一解得各不可约重数，即得到四个同余类分解公式。
平凡重数的地板函数表达式由同一解在平凡特征标分量上的化简得到。
\end{proof}

\begin{conclusion}[二次典范像的 \texorpdfstring{$S_4$}{S4} 等变二次关系分解]\label{con:xi-terminal-zm-s4-canonical-quadrics-equivariant-decomposition}
记
\[
V:=H^0(X,K_X),\qquad
Q:=\ker\!\Bigl(\operatorname{Sym}^2V\xrightarrow{\mu_2}H^0(X,2K_X)\Bigr).
\]
则：
\begin{enumerate}
  \item
  \[
  H^0(X,2K_X)\cong
  3\mathbf 1\oplus \varepsilon\oplus 4\rho_2\oplus 7\rho_3\oplus 4\rho_3';
  \]
  \item
  \[
  Q\cong
  8\mathbf 1\oplus 2\varepsilon\oplus 9\rho_2\oplus 13\rho_3\oplus 8\rho_3'.
  \]
\end{enumerate}
特别地，\(\dim Q^{S_4}=8\)。
\end{conclusion}
\begin{proof}[证明要点]
第 \(1\) 点是结论 \ref{con:xi-terminal-zm-s4-pluricanonical-irrep-decomposition-all-n} 在 \(n=2\) 时的特例。
由结论 \ref{con:xi-terminal-zm-s4-canonical-projective-normality-quadratic-generation}，
\[
0\to Q\to \operatorname{Sym}^2V\xrightarrow{\mu_2}H^0(X,2K_X)\to 0
\]
为 \(S_4\)-等变正合列。
又由结论 \ref{con:xi-terminal-zm-s4-chevalley-weil-explicit-decomposition} 中
\(V\cong 2\varepsilon\oplus\rho_2\oplus\rho_3\oplus 3\rho_3'\)，
可算得
\[
\operatorname{Sym}^2V\cong
11\mathbf 1\oplus 3\varepsilon\oplus 13\rho_2\oplus 20\rho_3\oplus 12\rho_3'.
\]
与第 \(1\) 点逐项相减即得第 \(2\) 点，平凡重数遂为 \(8\)。
\end{proof}

\paragraph{四次谱方程在固定记号下的 \texorpdfstring{$V_4$}{V4} 分辨率与闭包几何刚性链}
以下记号在本节固定：
\[
\Pi(\lambda,y)=\lambda^4-\lambda^3-(2y+1)\lambda^2+\lambda+y(y+1)\in \QQ(y)[\lambda],
\]
\[
C_{\mathrm{LY}}(y):=256y^3+411y^2+165y+32,\qquad
\Delta(y):=\Disc_{\lambda}\Pi(\lambda,y)=-y(y-1)C_{\mathrm{LY}}(y)\in\QQ[y].
\]
并沿用文内既有中间商记号
\[
X:=\widetilde X,\qquad C:=X/A_4,\qquad C_6:=X/V_4^{\mathrm{norm}}.
\]

\begin{conclusion}[\texorpdfstring{$V_4$}{V4} 分辨率三次多项式与判别式的完全一致]\label{con:xi-terminal-zm-v4-resolvent-discriminant-complete-coincidence}
可取 \(\Pi(\lambda,y)\) 的 \(V_4\)-分辨率三次多项式为
\[
R(u,y)=u^3+(2y+1)u^2-(2y+1)^2u-(8y^3+13y^2+5y+1)\in\QQ(y)[u],
\]
并满足严格恒等式
\[
\Disc_u R(u,y)=\Disc_{\lambda}\Pi(\lambda,y)=\Delta(y)=-y(y-1)C_{\mathrm{LY}}(y).
\]
\end{conclusion}
\begin{proof}[证明要点]
一般首一四次多项式
\(
f(x)=x^4+ax^3+bx^2+cx+d
\)
的分辨率三次可写为
\[
u^3-bu^2+(ac-4d)u+(4bd-a^2d-c^2).
\]
代入
\(
a=-1,\ b=-(2y+1),\ c=1,\ d=y(y+1)
\)
即得上式 \(R(u,y)\)。
判别式恒等式与 \(\Disc_{\lambda}\Pi\) 的一致性已在结论 \ref{con:xi-terminal-zm-generic-galois-s4} 给出，故结论成立。证毕。
\end{proof}

\begin{conclusion}[\texorpdfstring{$S_4$}{S4} 闭包曲线的分歧护照与中间商亏格谱]\label{con:xi-terminal-zm-s4-closure-passport-genus-spectrum-restated}
设 \(L/\QQ(y)\) 为 \(\Pi(\lambda,y)\) 的分裂域，\(X\) 为对应光滑射影曲线。则：
\begin{enumerate}
  \item \(\Gal(L/\QQ(y))\cong S_4\)；
  \item 覆盖 \(X\to\PP^1_y\) 的分歧惯性阶为 \((2,2,2,2,2,4)\)；
  \item \(g(X)=16\)；
  \item 曲线 \(C=X/A_4\) 亏格为 \(2\)，并可取模型
  \[
  C:\ s^2=\Delta(y)=-y(y-1)C_{\mathrm{LY}}(y);
  \]
  \item 曲线 \(C_6=X/V_4^{\mathrm{norm}}\) 满足 \(g(C_6)=4\)，且
  \[
  C_6\to \PP^1_y
  \]
  为 \(S_3\)-Galois 覆盖，六个分歧值处惯性均为换位（阶 \(2\)）。
  此外，\(C_6\) 可由
  \[
  \begin{cases}
  s^2=\Delta(y),\\
  R(u,y)=0
  \end{cases}
  \]
  的仿射簇取规范化得到。
\end{enumerate}
\end{conclusion}
\begin{proof}[证明要点]
由结论 \ref{con:xi-terminal-zm-generic-galois-s4} 得到 \(S_4\) 泛型伽罗瓦群与判别式核；
由结论 \ref{con:xi-terminal-zm-s4-galois-closure-genus-tower-holomorphic-differentials} 得到分歧护照 \((4)\cdot(2)^5\) 与 \(g(X)=16\)，并给出 \(X/A_4\) 与 \(X/V_4^{\mathrm{norm}}\) 的亏格分别为 \(2,4\)。
由 \(S_4/V_4^{\mathrm{norm}}\cong S_3\) 及结论 \ref{con:xi-terminal-zm-discriminant-ridge-etale-c3-explicit-model} 的显式模型可知，\(C_6\) 的规范化表示与所述方程组一致。证毕。
\end{proof}

\begin{conclusion}[无分歧循环三次覆盖与 \texorpdfstring{$3$}{3}-挠类]\label{con:xi-terminal-zm-c6-over-c-etale-c3-pic3-class}
存在定义在 \(\QQ\) 上的无分歧循环三次覆盖
\[
\pi:C_6\longrightarrow C
\]
且 \(\Gal(C_6/C)\cong C_3\)。
该覆盖对应于 \(\Pic^0(C)\) 的非平凡三挠类，并等价给出 \(J(C)\) 上 \(\QQ\)-定义的阶 \(3\) 循环子群方案与次数 \(3\) 同源。
\end{conclusion}
\begin{proof}[证明要点]
结论 \ref{con:xi-terminal-zm-discriminant-ridge-etale-c3-explicit-model} 已给出 \(\pi\) 的循环三次与处处非分歧性。
对光滑射影曲线，循环 \'{e}tale \(3\)-覆盖与 \(\Pic^0[3]\) 的非平凡类对应；据此得到三挠类与相应 \(3\)-同源。
该下降与核的 \(\QQ\)-定义性可由结论 \ref{con:xi-terminal-zm-discriminant-ridge-3isogeny-descent} 对照验证。证毕。
\end{proof}

\begin{conclusion}[判别式二次脊 Jacobian 的有理挠子下界]\label{con:xi-terminal-zm-discriminant-ridge-jacobian-rational-torsion-lower-bound-2x2x3}
曲线
\(
C:\ s^2=-y(y-1)C_{\mathrm{LY}}(y)
\)
的 Jacobian 满足
\[
(\ZZ/2\ZZ)^2\hookrightarrow J(C)(\QQ)_{\mathrm{tors}}.
\]
并且由三次无分歧覆盖诱导的 \(3\)-同源核提供一个 \(\QQ\)-定义阶 \(3\) 循环子群方案；
在该核分裂为常值群方案的情形下，进一步得到
\[
(\ZZ/2\ZZ)^2\times(\ZZ/3\ZZ)\hookrightarrow J(C)(\QQ)_{\mathrm{tors}}.
\]
\end{conclusion}
\begin{proof}[证明要点]
由结论 \ref{con:xi-terminal-zm-discriminant-ridge-jacobian-torsion-upper-bound}，
\[
D_0=[(0,0)-\infty],\qquad D_1=[(1,0)-\infty]
\]
给出两条线性无关的有理 \(2\)-挠元，故有 \((\ZZ/2\ZZ)^2\) 注入。
另一方面，由结论 \ref{con:xi-terminal-zm-c6-over-c-etale-c3-pic3-class} 得到阶 \(3\) 循环子群方案与度 \(3\) 同源；
若该子群在 \(\QQ\) 上常值分裂，则可取非零有理 \(3\)-挠点，与前述 \(2\)-挠按互素阶直积分解，得到所示注入。证毕。
\end{proof}

\begin{conclusion}[三阶 Sylow 自由作用与亏格 \(6\) 商曲线]\label{con:xi-terminal-zm-sylow3-free-action-genus6-quotient}
设 \(P\le S_4\) 为任一三阶 Sylow 子群（\(P\cong C_3\)）。
则 \(P\) 在 \(X\) 上无不动点；故
\[
X\to X_P:=X/P
\]
为三重 \'{e}tale 覆盖，且
\[
g(X_P)=6.
\]
\end{conclusion}
\begin{proof}[证明要点]
由结论 \ref{con:xi-terminal-zm-s4-closure-passport-genus-spectrum-restated}，\(X\to\PP^1_y\) 的惯性阶仅为 \(2\) 与 \(4\)，故任一点稳定子阶整除 \(4\)。
因此三阶元素不可能固定点，\(P\) 作用自由。
应用 Riemann--Hurwitz：
\[
2g(X)-2=|P|\bigl(2g(X_P)-2\bigr),
\]
代入 \(g(X)=16\)、\(|P|=3\) 得 \(g(X_P)=6\)。证毕。
\end{proof}

\begin{conclusion}[全纯微分空间的 \texorpdfstring{$S_4$}{S4} 表示分解刚性]\label{con:xi-terminal-zm-holomorphic-differential-s4-decomposition-rigid-restated}
记
\[
W:=H^0(X,\Omega^1_X).
\]
则 \(W\) 作为 \(\QQ[S_4]\)-模唯一分解为
\[
W\cong 2\varepsilon\oplus V_2\oplus V_3\oplus 3V_3',
\]
其中 \(\varepsilon\) 为符号表示，\(V_2\) 为二维不可约表示，\(V_3\) 为标准三维表示，\(V_3'=V_3\otimes\varepsilon\)。
\end{conclusion}
\begin{proof}[证明要点]
对任意 \(H\le S_4\) 有
\(
W^H\simeq H^0(X/H,\Omega^1)
\)，故 \(\dim W^H=g(X/H)\)。
由结论 \ref{con:xi-terminal-zm-s4-galois-closure-genus-tower-holomorphic-differentials}，
\[
\dim W^{S_4}=0,\quad
\dim W^{A_4}=2,\quad
\dim W^{V_4^{\mathrm{norm}}}=4,\quad
\dim W^{S_3}=1,
\]
并且 \(\dim W=g(X)=16\)。
将 \(S_4\) 五个不可约表示的子群不变维数代入该线性系统，得到唯一解即所示分解。证毕。
\end{proof}

\begin{conclusion}[Jacobian 的群代数同源分解与 \texorpdfstring{$(1,3)$}{(1,3)}-极化 Prym 因子]\label{con:xi-terminal-zm-jacobian-group-algebra-prym13-visible-factor-restated}
存在 \(\QQ\)-同源分解
\[
J(X)\sim_{\QQ}J(C)\times P\times A_3\times A_9,
\]
其中：
\begin{enumerate}
  \item \(P=\operatorname{Prym}(C_6/C)\) 为二维阿贝尔簇，极化型为 \((1,3)\)；
  \item \(A_3\) 与 \(A_9\) 分别对应 \(V_3\) 与 \(3V_3'\) 的等型部分（\(\dim A_3=3,\ \dim A_9=9\)）；
  \item deck 群 \(C_3\) 在 \(P\) 上诱导
  \[
  \QQ(\zeta_3)\hookrightarrow \End^0(P).
  \]
\end{enumerate}
\end{conclusion}
\begin{proof}[证明要点]
由结论 \ref{con:xi-terminal-zm-s4-jacobian-group-algebra-prym-e3-b3}，
\[
J(X)\sim_{\QQ}J(C)\times E_{\mathrm{res}}^2\times E^3\times B^3.
\]
由定理 \ref{thm:xi-terminal-zm-s3-closure-jacobian-splitting} 与定理 \ref{thm:xi-terminal-zm-prym-polarization-a2-rigidity}，
\[
P=\operatorname{Prym}(C_6/C)\sim_{\QQ}E_{\mathrm{res}}^2,\qquad \mathrm{type}(P)=(1,3).
\]
令 \(A_3:=E^3,\ A_9:=B^3\) 即得所示分解。
又 \(C_3\) 对覆盖 \(C_6\to C\) 的作用在 Prym 部分给出群代数作用
\(
\QQ[C_3]\cong \QQ\times\QQ(\zeta_3)
\)，故得到 \(\QQ(\zeta_3)\subset \End^0(P)\)。证毕。
\end{proof}

\paragraph{判别式二次分辨率与 \texorpdfstring{$C_3$}{C3} 商几何：纤维积规范化、Prym 因子与固定子商 Jacobian 分解}
沿用结论 \ref{con:xi-terminal-zm-s4-galois-closure-genus-tower-holomorphic-differentials} 的设定，记
\[
\widetilde{C}:=\widetilde{X},\qquad
L:=\QQ(\widetilde{C}),\qquad
G:=\mathrm{Gal}(L/\QQ(y))\cong S_4.
\]
并记
\[
E:=\widetilde{C}/S_3,\qquad
H_{\mathrm{LY}}:=\widetilde{C}/A_4,\qquad
C_3:=S_3\cap A_4.
\]

\begin{conclusion}[判别式二次分辨率的几何合成：\texorpdfstring{$C_3$}{C3} 商即椭圆谱曲线与 Lee--Yang 二次脊的规范化纤维积]\label{con:xi-terminal-zm-s4-c3-fiberproduct-normalization}
存在曲线同构
\[
\widetilde{C}/C_3\ \cong\ \mathrm{Norm}\!\bigl(E\times_{\PP^1_y}H_{\mathrm{LY}}\bigr),
\]
等价地，在函数域层面有
\[
\QQ(\widetilde{C}/C_3)=\QQ(E)\cdot \QQ(H_{\mathrm{LY}})\subset L.
\]
此外，令
\[
P_{\mathrm{LY}}(y):=256y^3+411y^2+165y+32,
\]
则由判别式恒等式 \(\mathrm{Disc}_{\lambda}(F)=-y(y-1)P_{\mathrm{LY}}(y)\) 可取
\[
\QQ(H_{\mathrm{LY}})=\QQ(y)\!\left(\sqrt{-y(y-1)P_{\mathrm{LY}}(y)}\right),
\]
从而
\[
\QQ(\widetilde{C}/C_3)=\QQ(E)\!\left(\sqrt{-y(y-1)P_{\mathrm{LY}}(y)}\right).
\]
\end{conclusion}
\begin{proof}[证明要点]
由伽罗瓦对应，
\[
\QQ(E)=L^{S_3},\qquad \QQ(H_{\mathrm{LY}})=L^{A_4}.
\]
由于 \(L/\QQ(y)\) 为 Galois 扩张，固定域合成满足
\[
L^{S_3}\,L^{A_4}=L^{S_3\cap A_4}=L^{C_3},
\]
即
\[
\QQ(\widetilde{C}/C_3)=\QQ(E)\cdot\QQ(H_{\mathrm{LY}}).
\]
规范化纤维积的函数域等于右侧合成域，故得到曲线同构。
二次分辨子域的显式式由结论 \ref{con:xi-terminal-zm-discriminant-leyang} 与
\ref{con:xi-terminal-zm-s4-galois-closure-genus-tower-holomorphic-differentials} 直接给出。证毕。
\end{proof}

\begin{conclusion}[\texorpdfstring{$C_3$}{C3} 商曲线 Jacobian 的三因子同源分解]\label{con:xi-terminal-zm-s4-c3-quotient-jacobian-three-factor-decomposition}
设
\[
C:=\widetilde{C}/C_3,\qquad
J_{\mathrm{LY}}:=\mathrm{Jac}(H_{\mathrm{LY}}).
\]
沿用结论 \ref{con:xi-terminal-zm-s4-jacobian-group-algebra-prym-e3-b3} 的 \(E\) 与三维阿贝尔簇 \(B\) 记号，则有
\[
\mathrm{Jac}(C)\sim_{\QQ}J_{\mathrm{LY}}\times E\times B.
\]
\end{conclusion}
\begin{proof}[证明要点]
由结论 \ref{con:xi-terminal-zm-s4-chevalley-weil-explicit-decomposition}，
\[
H^0(\widetilde{C},\Omega^1)\cong 2\varepsilon\oplus \rho_2\oplus \rho_3\oplus 3\rho_3'.
\]
对 \(C_3=\langle(123)\rangle\) 取不变子空间并用角色平均公式
\[
\dim W^{C_3}=\frac{1}{3}\bigl(\chi_W(1)+\chi_W((123))+\chi_W((132))\bigr),
\]
可得
\[
\dim(\varepsilon^{C_3})=1,\quad
\dim(\rho_2^{C_3})=0,\quad
\dim(\rho_3^{C_3})=1,\quad
\dim((\rho_3')^{C_3})=1.
\]
结合结论 \ref{con:xi-terminal-zm-s4-jacobian-group-algebra-prym-e3-b3} 的同源分解
\[
J(\widetilde{C})\sim_{\QQ}J_{\mathrm{LY}}\times E_{\mathrm{res}}^2\times E^3\times B^3,
\]
对幂等元 \(e_{C_3}=\frac1{3}\sum_{h\in C_3}h\) 取像后，分别截取到
\(J_{\mathrm{LY}}\)、\(E\)、\(B\) 各一次，且 \(E_{\mathrm{res}}^2\) 分量不贡献，遂得结论。证毕。
\end{proof}

\begin{conclusion}[Prym--Lee--Yang 等价：\texorpdfstring{$\mathrm{Prym}(C/E)\sim_{\QQ}J_{\mathrm{LY}}\times B$}{Prym(C/E) isogenous to JLY x B}]\label{con:xi-terminal-zm-s4-c3-prym-leyang-equivalence}
由包含 \(C_3\le S_3\) 诱导次数 \(2\) 映射
\[
\pi:C=\widetilde{C}/C_3\longrightarrow E=\widetilde{C}/S_3.
\]
则其 Prym 变体满足
\[
\mathrm{Prym}(C/E)\sim_{\QQ}J_{\mathrm{LY}}\times B.
\]
\end{conclusion}
\begin{proof}[证明要点]
对次数 \(2\) 覆盖 \(\pi:C\to E\)，有标准同源分解
\[
\mathrm{Jac}(C)\sim_{\QQ}\mathrm{Jac}(E)\times \mathrm{Prym}(C/E),
\]
且 \(\mathrm{Jac}(E)=E\)。
代入结论 \ref{con:xi-terminal-zm-s4-c3-quotient-jacobian-three-factor-decomposition} 的
\[
\mathrm{Jac}(C)\sim_{\QQ}J_{\mathrm{LY}}\times E\times B
\]
并消去 \(E\) 因子即得。证毕。
\end{proof}

\begin{conclusion}[Lee--Yang 非分歧六点的 Abel 和闭式与二次覆叠分歧除子的群和消去]\label{con:xi-terminal-zm-leyang-unramified-sixpoint-abel-closed-form-branch-balance}
仍在椭圆曲线
\[
E:\ Y^2=X^3-X+\frac14,\qquad
y=X^2+Y-\frac12
\]
上，记 \(P_{\mathrm{LY}}(y)=256y^3+411y^2+165y+32\)。
设 \(Q_1,Q_2,Q_3\) 为 Lee--Yang 三个分歧点（见结论 \ref{con:xi-terminal-zm-leyang-ramification-divisor-principalization-trace}），并定义有效除子 \(D_{\mathrm{unr}}\)（次数 \(6\)）满足
\[
\mathrm{div}_0\!\bigl(P_{\mathrm{LY}}(y)\bigr)=2(Q_1+Q_2+Q_3)+D_{\mathrm{unr}}.
\]
再记
\[
Q_0=\Bigl(1,-\frac12\Bigr),\qquad
Q_0'=\Bigl(-1,\frac12\Bigr)
\]
（分别位于 \(y=0,1\) 的有理分歧点）。
则在椭圆曲线群律下有
\[
\bigoplus_{P\in\mathrm{Supp}(D_{\mathrm{unr}})}P
=2(Q_0\oplus Q_0')
=\Bigl(\frac{161}{16},\,\frac{2033}{64}\Bigr)\in E(\QQ).
\]
进一步，考虑二次覆叠
\[
\pi:\ C\to E,\qquad C:\ w^2=-y(y-1)P_{\mathrm{LY}}(y),
\]
记其分歧除子为 \(B_{\pi}\)。
则
\[
\bigoplus_{P\in\mathrm{Supp}(B_{\pi})}P=O,
\]
等价地 \(B_{\pi}-10O\) 在 \(\mathrm{Pic}^0(E)\) 中为零类。
\end{conclusion}
\begin{proof}[证明要点]
由 \((y)_\infty=4O\) 可得
\[
\mathrm{div}\!\bigl(P_{\mathrm{LY}}(y)\bigr)=2(Q_1+Q_2+Q_3)+D_{\mathrm{unr}}-12O,
\]
从而
\[
2(Q_1\oplus Q_2\oplus Q_3)\oplus\!\!\bigoplus_{P\in\mathrm{Supp}(D_{\mathrm{unr}})}\!\!P=O.
\]
又由结论 \ref{con:xi-terminal-zm-leyang-ramification-divisor-principalization-trace}，
\[
Q_1\oplus Q_2\oplus Q_3=-(Q_0\oplus Q_0'),
\]
故得到
\[
\bigoplus_{P\in\mathrm{Supp}(D_{\mathrm{unr}})}P=2(Q_0\oplus Q_0').
\]
同一结论给出 \(Q_0\oplus Q_0'=(\tfrac14,\tfrac18)\)，对该点施行倍点公式即得
\((\tfrac{161}{16},\tfrac{2033}{64})\)。
\par
对分歧除子，设
\[
\mathrm{div}(y)=2Q_0+D_0-4O,\qquad
\mathrm{div}(y-1)=2Q_0'+D_1-4O,
\]
其中 \(D_0,D_1\) 分别是 \(y=0,1\) 纤维中的简单点部分（各次数 \(2\)）。
则 \(B_{\pi}=D_0+D_1+D_{\mathrm{unr}}\)。
由上述三条主除子关系在 \(\mathrm{Pic}^0(E)\cong E\) 下可得
\[
\bigoplus_{P\in\mathrm{Supp}(D_0)}P=-2Q_0,\qquad
\bigoplus_{P\in\mathrm{Supp}(D_1)}P=-2Q_0',\qquad
\bigoplus_{P\in\mathrm{Supp}(D_{\mathrm{unr}})}P=2(Q_0\oplus Q_0'),
\]
三式相加即得
\(
\bigoplus_{P\in\mathrm{Supp}(B_{\pi})}P=O
\)。
证毕。
\end{proof}

\begin{conclusion}[换位固定子商曲线的 Jacobian 分解]\label{con:xi-terminal-zm-s4-transposition-quotient-jacobian-full-decomposition}
取换位 \(\tau=(12)\in S_4\)，令
\[
C_{\tau}:=\widetilde{C}/\langle\tau\rangle.
\]
沿用 \(E_{\mathrm{res}}:=\widetilde{C}/D_8\)、\(E:=\widetilde{C}/S_3\) 与 \(B\) 的记号，则有
\[
\mathrm{Jac}(C_{\tau})\sim_{\QQ}E_{\mathrm{res}}\times E^2\times B.
\]
\end{conclusion}
\begin{proof}[证明要点]
由结论 \ref{con:xi-terminal-zm-s4-chevalley-weil-explicit-decomposition} 的表示分解，对 \(\langle\tau\rangle\) 取不变并用
\[
\dim W^{\langle\tau\rangle}=\frac12\bigl(\chi_W(1)+\chi_W(\tau)\bigr)
\]
得
\[
\dim(\varepsilon^{\langle\tau\rangle})=0,\quad
\dim(\rho_2^{\langle\tau\rangle})=1,\quad
\dim(\rho_3^{\langle\tau\rangle})=2,\quad
\dim((\rho_3')^{\langle\tau\rangle})=1.
\]
结合结论 \ref{con:xi-terminal-zm-s4-jacobian-group-algebra-prym-e3-b3} 的等型对应
\[
\varepsilon\leftrightarrow J_{\mathrm{LY}},\quad
\rho_2\leftrightarrow E_{\mathrm{res}}^2,\quad
\rho_3\leftrightarrow E^3,\quad
\rho_3'\leftrightarrow B^3,
\]
可知分别贡献 \(0,1,2,1\) 个因子，遂得所示同源分解。
其总维数为 \(1+2+3=6\)，与结论 \ref{con:xi-terminal-zm-s4-hurwitz-subgroup-spectrum-cycletypes} 的 \(g(\widetilde{C}/\langle\tau\rangle)=6\) 一致。证毕。
\end{proof}

\begin{conclusion}[双换位固定子商曲线的 Jacobian 分解]\label{con:xi-terminal-zm-s4-doubletransposition-quotient-jacobian-full-decomposition}
取双换位 \(\delta=(12)(34)\in S_4\)，令
\[
C_{\delta}:=\widetilde{C}/\langle\delta\rangle.
\]
则有
\[
\mathrm{Jac}(C_{\delta})\sim_{\QQ}J_{\mathrm{LY}}\times E_{\mathrm{res}}^2\times E\times B.
\]
\end{conclusion}
\begin{proof}[证明要点]
同理，
\[
\dim W^{\langle\delta\rangle}=\frac12\bigl(\chi_W(1)+\chi_W(\delta)\bigr),
\]
并由角色值计算得
\[
\dim(\varepsilon^{\langle\delta\rangle})=1,\quad
\dim(\rho_2^{\langle\delta\rangle})=2,\quad
\dim(\rho_3^{\langle\delta\rangle})=1,\quad
\dim((\rho_3')^{\langle\delta\rangle})=1.
\]
代入同一等型对应即得
\[
\mathrm{Jac}(C_{\delta})\sim_{\QQ}J_{\mathrm{LY}}\times E_{\mathrm{res}}^2\times E\times B.
\]
总维数 \(2+2+1+3=8\)，与
\(g(\widetilde{C}/\langle\delta\rangle)=8\)
一致。证毕。
\end{proof}

\begin{conclusion}[换位与双换位商的 Prym 极化型及 \texorpdfstring{$(\pm1)$}{(+/-1)}-本征同源分解]\label{con:xi-terminal-zm-s4-involution-prym-polarization-eigensplitting}
取换位 \(\tau\in S_4\) 与双换位 \(\delta\in S_4\)，并记
\[
C_{\tau}:=\widetilde{C}/\langle\tau\rangle,\qquad
C_{\delta}:=\widetilde{C}/\langle\delta\rangle,
\]
以及自然二次覆盖
\[
\pi_{\tau}:\widetilde{C}\to C_{\tau},\qquad
\pi_{\delta}:\widetilde{C}\to C_{\delta}.
\]
则：
\begin{enumerate}
  \item \(\pi_{\tau}\) 的分歧点数为 \(10\)，\(\dim \Prym(\pi_{\tau})=10\)，且由 \(\Jac(\widetilde{C})\) 的主极化在 \(\Prym(\pi_{\tau})\) 上诱导的极化型为
  \[
  (1^{6},2^{4}).
  \]
  并且在 \(\QQ\)-同源意义下有 \((\pm1)\)-本征分解
  \[
  \Jac(C_{\tau})\ \sim_{\QQ}\ E_{\mathrm{res}}\times E^2\times B,\qquad
  \Prym(\pi_{\tau})\ \sim_{\QQ}\ J_{\mathrm{LY}}\times E_{\mathrm{res}}\times E\times B^2.
  \]
  \item \(\pi_{\delta}\) 的分歧点数为 \(2\)，\(\dim \Prym(\pi_{\delta})=8\)，并且 \(\Prym(\pi_{\delta})\) 为主极化阿贝尔簇。
  在 \(\QQ\)-同源意义下有 \((\pm1)\)-本征分解
  \[
  \Jac(C_{\delta})\ \sim_{\QQ}\ J_{\mathrm{LY}}\times E_{\mathrm{res}}^2\times E\times B,\qquad
  \Prym(\pi_{\delta})\ \sim_{\QQ}\ E^2\times B^2.
  \]
\end{enumerate}
\end{conclusion}
\begin{proof}[证明要点]
由结论 \ref{con:xi-terminal-zm-s4-transposition-quotient-jacobian-full-decomposition} 与
\ref{con:xi-terminal-zm-s4-doubletransposition-quotient-jacobian-full-decomposition} 已知 \(g(C_{\tau})=6\)、\(g(C_{\delta})=8\)。
又由结论 \ref{con:xi-terminal-zm-s4-conjugacy-class-fixedpoints-and-cyclic-quotient-genera}，
\[
\#\mathrm{Fix}_{\widetilde C}(\tau)=10,\qquad \#\mathrm{Fix}_{\widetilde C}(\delta)=2,
\]
故两条二次覆盖的分歧点数分别为 \(10,2\)，并有
\[
\dim\Prym(\pi_{\tau})=g(\widetilde C)-g(C_{\tau})=16-6=10,\qquad
\dim\Prym(\pi_{\delta})=16-8=8.
\]
对分歧二次覆盖 \(\pi:C'\to C\)，Prym 极化型的标准公式给出
\(\mathrm{type}=(1^{g(C)},2^{r/2-1})\)，其中 \(r\) 为分歧点数。
据此代入 \((g,r)=(6,10)\) 得 \((1^6,2^4)\)，代入 \((g,r)=(8,2)\) 得主极化。
\par
同源分解由幂等元 \(e_{\pm}=(1\pm g)/2\)（\(g=\tau\) 或 \(\delta\)）作用于 \(\Jac(\widetilde C)\) 给出：
\(
e_+\Jac(\widetilde C)\sim_{\QQ}\Jac(\widetilde C/\langle g\rangle)
\)
且
\(
e_-\Jac(\widetilde C)\sim_{\QQ}\Prym(\pi_g)
\)。
将结论 \ref{con:xi-terminal-zm-s4-jacobian-group-algebra-prym-e3-b3} 的等型分解
\[
\Jac(\widetilde C)\sim_{\QQ}J_{\mathrm{LY}}\times E_{\mathrm{res}}^2\times E^3\times B^3
\]
与换位、双换位在各不可约表示上的特征值重数比较，即得所列 \((\pm1)\)-本征同源分解。证毕。
\end{proof}

\begin{conclusion}[四循环与 Klein 四群商的 Jacobian 纯分量分裂]\label{con:xi-terminal-zm-s4-c4-v4-quotient-jacobian-pure-splitting}
取无穷远惯性生成元 \(c\in S_4\)（四循环），并记
\[
C_{(4)}:=\widetilde C/\langle c\rangle,\qquad
C_{V_4}:=\widetilde C/V_4^{\mathrm{norm}}.
\]
则二者皆为属 \(4\) 曲线，且其 Jacobian 在 \(\QQ\)-同源意义下满足
\[
\Jac(C_{(4)})\ \sim_{\QQ}\ E_{\mathrm{res}}\times B,\qquad
\Jac(C_{V_4})\ \sim_{\QQ}\ J_{\mathrm{LY}}\times E_{\mathrm{res}}^2.
\]
\end{conclusion}
\begin{proof}[证明要点]
属数断言来自结论 \ref{con:xi-terminal-zm-s4-conjugacy-class-fixedpoints-and-cyclic-quotient-genera} 中
\(
g(\widetilde C/\langle c\rangle)=4
\)
及结论 \ref{con:xi-terminal-zm-s4-closure-passport-genus-spectrum-restated} 中
\(
g(\widetilde C/V_4^{\mathrm{norm}})=4
\)。
对 \(\Jac(C_{(4)})\) 的分裂由结论 \ref{con:xi-terminal-zm-s4-c4-prym-threefold-and-v3sgn-identification} 的 \(\Jac(\widetilde C/\langle c\rangle)\) 同源分解给出。
对 \(\Jac(C_{V_4})\) 的分裂由定理 \ref{thm:xi-terminal-zm-s3-closure-jacobian-splitting} 给出的
\(
\Jac(\widetilde C/V_4^{\mathrm{norm}})\sim_{\QQ}\Jac(\widetilde C/A_4)\times E_{\mathrm{res}}^2
\)
与 \(J_{\mathrm{LY}}=\Jac(\widetilde C/A_4)\)（结论 \ref{con:xi-terminal-zm-s4-c3-quotient-jacobian-three-factor-decomposition} 的记号）直接推出。证毕。
\end{proof}

\paragraph{\texorpdfstring{$A_4$}{A4}/\texorpdfstring{$S_3$}{S3} 中间覆盖的分歧刚性、标准表示因子与定点谱}
在既定记号下，记
\[
X:=\widetilde{C},\qquad
L=\QQ(X),\qquad
\mathrm{Gal}(L/\QQ(y))\cong S_4,
\]
并令
\[
X_A:=X/A_4,\qquad E:=X/S_3.
\]
其中 \(S_3\le S_4\) 为字母稳定子，且 \(X\to \PP^1_y\) 的惯性型为
\((C_4;C_2,C_2,C_2,C_2,C_2)\)。

\begin{conclusion}[\texorpdfstring{$A_4$}{A4} 子覆盖的单点分歧刚性]\label{con:xi-terminal-zm-a4-subcover-single-branch-rigidity}
自然商映射
\[
f:X\longrightarrow X_A=X/A_4
\]
为次数 \(12\) 的 \(A_4\)-Galois 覆盖。
其分歧点在 \(X_A\) 上恰有一处；该唯一分歧点的惯性群同构于 \(C_2\)，并由双换位生成。
\end{conclusion}
\begin{proof}[证明要点]
记 \(\varphi:X\to\PP^1_y\) 为 \(S_4\)-正则覆盖。对 \(P\in X\)，中间覆盖的惯性满足
\[
I_P(f)=I_P(\varphi)\cap A_4.
\]
对五个有限分歧值，\(I_P(\varphi)\) 由换位生成，属于奇置换，故
\(
I_P(\varphi)\cap A_4=\{1\}
\)；
于是 \(f\) 在有限分歧值上方处处不分歧。

在无穷远处分歧，\(I_P(\varphi)\cong C_4\) 由四循环 \(\rho\) 生成。
其平方 \(\rho^2\) 为双换位，属于 \(A_4\)；
且 \(A_4\) 中无阶 \(4\) 元，故
\[
I_P(\varphi)\cap A_4=\langle \rho^2\rangle\cong C_2.
\]
因此 \(f\) 仅可能在 \(\infty\) 上方分歧。

又 \(X_A\to \PP^1_y\) 为符号二次子覆盖，在 \(\infty\) 处分歧，故 \(\infty\) 在 \(X_A\) 上方仅有一个点。
综上，\(f\) 的分歧点唯一，惯性为双换位生成的 \(C_2\)。证毕。
\end{proof}

\begin{conclusion}[\texorpdfstring{$S_3$}{S3} 子覆盖 \texorpdfstring{$X\to E$}{X to E} 的分歧除子全刻画]\label{con:xi-terminal-zm-s3-subcover-branch-divisor-complete}
令
\[
\pi:X\longrightarrow E=X/S_3
\]
为次数 \(6\) 的 \(S_3\)-Galois 覆盖，则：
\begin{enumerate}
  \item \(\pi\) 的所有分歧均为简单二分歧（惯性同构于 \(C_2\)），不存在阶 \(3\) 惯性；
  \item \(\pi\) 在 \(E\) 上恰有 \(10\) 个分歧点；
  \item 对每个有限分歧值 \(y_0\in\PP^1_y\)，纤维 \(E_{y_0}\) 在 \(E\to\PP^1_y\) 中呈 \(2+1+1\) 型：其中唯一二分歧点上方 \(\pi\) 不分歧，而两个不分歧点上方 \(\pi\) 恰好分歧；
  \item 在 \(E\) 的无穷远分歧点上方，\(\pi\) 不分歧。
\end{enumerate}
\end{conclusion}
\begin{proof}[证明要点]
设 \(H=S_3\le G=S_4\)。则
\[
\QQ(y)\subset \QQ(E)=L^H\subset L=\QQ(X),
\]
且 \(\mathrm{Gal}(L/L^H)=H\)。

取有限分歧值 \(y_0\)。在 \(\varphi\) 中其惯性为换位 \(\tau\in S_4\)。
考察 \(\tau\) 在陪集空间 \(S_4/H\simeq\{1,2,3,4\}\) 上的作用，循环分解为 \(2+1+1\)。
这对应 \(E\to\PP^1_y\) 在 \(y_0\) 上有一个二分歧点和两个不分歧点。

若 \(R\in E_{y_0}\)，取 \(P\in X\) 在其上方，则
\[
I_P(\pi)=H\cap g\langle\tau\rangle g^{-1}
\]
（\(g\) 由点选择决定）。
当 \(R\) 对应 \(\tau\) 的不动点时，\(g^{-1}\tau g\in H\)，故 \(I_P(\pi)\cong C_2\)；
当 \(R\) 对应长度 \(2\) 轨道时，\(g^{-1}\tau g\notin H\)，故 \(I_P(\pi)=\{1\}\)。
于是每个有限分歧值贡献恰好两处 \(\pi\)-分歧点。

在无穷远处，\(\varphi\) 的惯性为 \(C_4=\langle\rho\rangle\)。
对任意 \(g\)，\(H\cap gC_4g^{-1}=\{1\}\)：因为 \(H\cong S_3\) 不含阶 \(4\) 元，且其阶 \(2\) 元均为换位，而 \(gC_4g^{-1}\) 的唯一阶 \(2\) 元是双换位。
故 \(\pi\) 在无穷远上方不分歧。

因此总分歧点数为 \(5\times2=10\)，且惯性均为 \(C_2\)。
最后由 Riemann--Hurwitz（\(g(E)=1\)）
\[
2g(X)-2=\sum_{P\in X}(e_P-1).
\]
每个分歧点在 \(E\) 上方有 \(6/2=3\) 个 ramification 点，每点贡献 \(1\)，故总贡献 \(10\times3=30\)，与 \(g(X)=16\) 一致，并排除额外分歧与阶 \(3\) 惯性。证毕。
\end{proof}

\begin{conclusion}[标准表示因子的几何实现与 \texorpdfstring{$A_{V_3}\sim_{\QQ}E^3$}{AV3 isogenous to E^3}]\label{con:xi-terminal-zm-standard-factor-geometric-realization-e-cube}
令 \(H_1,\dots,H_4\le S_4\) 为四个共轭字母稳定子（均同构于 \(S_3\)），并记
\[
E_i:=X/H_i\qquad(1\le i\le4).
\]
则各 \(E_i\) 皆为椭圆曲线，且在字母重标下两两同构。
记 \(\pi_i:X\to E_i\)，固定 \(O_i\in E_i\)，定义同态
\[
\Phi:E_1\times E_2\times E_3\times E_4\longrightarrow \operatorname{Jac}(X),\qquad
(P_1,\dots,P_4)\longmapsto \sum_{i=1}^4 \pi_i^*(P_i-O_i).
\]
令 \(A_{V_3}:=\operatorname{Im}(\Phi)\)。则：
\begin{enumerate}
  \item \(A_{V_3}\) 为 \(\operatorname{Jac}(X)\) 的 \(S_4\)-稳定阿贝尔子簇，其余切表示同构于标准不可约表示 \(V_3\)；
  \item \(\ker(\Phi)^0\) 为对角嵌入椭圆曲线
  \[
  E_\Delta=\{(P,P,P,P)\}\subset E_1\times\cdots\times E_4;
  \]
  \item 存在 \(S_4\)-等变同源
  \[
  (E_1\times\cdots\times E_4)/E_\Delta\ \sim_{\QQ}\ A_{V_3},
  \]
  从而（忽略标号同构 \(E_i\cong E\)）
  \[
  A_{V_3}\sim_{\QQ}E^3.
  \]
\end{enumerate}
\end{conclusion}
\begin{proof}[证明要点]
在余切空间上，\(\Phi\) 的微分为
\[
d\Phi:\bigoplus_{i=1}^4 H^0(E_i,\Omega^1)\longrightarrow H^0(X,\Omega^1).
\]
左端是四点置换表示，分解为 \(\mathbf 1\oplus V_3\)。
由 \(X/S_4\simeq \PP^1_y\) 得
\[
H^0(X,\Omega^1)^{S_4}\cong H^0(\PP^1,\Omega^1)=0,
\]
故像中无平凡分量。
另一方面，每个 \(\pi_i^*:H^0(E_i,\Omega^1)\to H^0(X,\Omega^1)\) 非零，故 \(d\Phi\) 像非零。
于是 \(\operatorname{Im}(d\Phi)\) 必为 \(V_3\) 分量，维数 \(3\)；从而 \(A_{V_3}\) 维数为 \(3\)，且余切表示为 \(V_3\)。

又 \(\dim(E_1\times\cdots\times E_4)=4\)，故 \(\ker(\Phi)^0\) 维数为 \(1\)。
其切空间为置换表示中的唯一 \(S_4\)-稳定一维子空间，即对角线 \(\mathbf 1\)；
因此 \(\ker(\Phi)^0=E_\Delta\)。

\(\Phi\) 于是分解为
\[
E_1\times\cdots\times E_4 \twoheadrightarrow (E_1\times\cdots\times E_4)/E_\Delta
\xrightarrow{\ \overline{\Phi}\ } A_{V_3},
\]
且 \(\overline{\Phi}\) 在余切空间上同构，故核有限，\(\overline{\Phi}\) 为同源。
识别 \(E_i\cong E\) 即得 \(A_{V_3}\sim_{\QQ}E^3\)。证毕。
\end{proof}

\begin{conclusion}[\texorpdfstring{$A_4$}{A4} 作用的定点谱：双换位固定两点而三循环无定点]\label{con:xi-terminal-zm-a4-fixedpoint-spectrum-v4-c3}
设 \(V_4\lhd A_4\) 为 Klein 四元子群。则：
\begin{enumerate}
  \item 对任意非平凡双换位 \(\delta\in V_4\setminus\{1\}\)，有
  \[
  \#\mathrm{Fix}_X(\delta)=2;
  \]
  \item 对任意三循环 \(\sigma\in A_4\)（阶 \(3\) 元），有
  \[
  \mathrm{Fix}_X(\sigma)=\varnothing.
  \]
\end{enumerate}
\end{conclusion}
\begin{proof}[证明要点]
由结论 \ref{con:xi-terminal-zm-a4-subcover-single-branch-rigidity}，\(f:X\to X_A\) 仅在一点 \(Q\in X_A\) 分歧，且惯性 \(I\cong C_2\)。
故任一非平凡群元若有定点，该定点必位于 \(f^{-1}(Q)\)。
又
\[
\#f^{-1}(Q)=\frac{|A_4|}{|I|}=\frac{12}{2}=6.
\]

取 \(\delta\in I\setminus\{1\}\)。其在陪集 \(A_4/I\) 上固定陪集数为
\[
\#\{gI:\delta gI=gI\}
=\frac{|N_{A_4}(I)|}{|I|}.
\]
而 \(N_{A_4}(I)=V_4\)，故固定陪集数为 \(4/2=2\)。
这正对应 \(\delta\) 在 \(f^{-1}(Q)\subset X\) 上的定点数；在其余点无定点，故该 \(\delta\) 满足 \(\#\mathrm{Fix}_X(\delta)=2\)。
又 \(A_4\) 中三个非平凡双换位彼此共轭，而定点数按共轭类不变，故结论对任意 \(\delta\in V_4\setminus\{1\}\) 成立。

若 \(\sigma\) 为三循环，则 \(\sigma\notin V_4\)，且不可能落入任何阶 \(2\) 惯性子群。
由于分歧仅来自阶 \(2\) 惯性，\(\sigma\) 不能稳定任何点，故 \(\mathrm{Fix}_X(\sigma)=\varnothing\)。证毕。
\end{proof}

\endinput


\begin{conclusion}[Lee--Yang 判别式双覆盖在 \(\QQ(E)\) 上的线性化表达]\label{con:xi-terminal-zm-leyang-doublecover-linearization-a-plus-by}
在椭圆曲线
\[
E/\QQ:\qquad Y^2=X^3-X+\frac14
\]
上定义
\[
y:=X^2+Y-\frac12,\qquad
P(y):=256y^3+411y^2+165y+32.
\]
于函数域 \(\QQ(E)\) 中取 \(w\) 满足
\[
w^2=-y(y-1)P(y).
\]
则存在唯一 \(A(X),B(X)\in\ZZ[X]\)，使得
\[
w^2=A(X)+B(X)Y.
\]
可取显式系数
\[
\begin{aligned}
A(X)=\,&-256X^{10}-2560X^9-795X^8+5470X^7+2321X^6-3802X^5-1671X^4\\
&\quad+986X^3+423X^2-94X-22,
\end{aligned}
\]
\[
\begin{aligned}
B(X)=\,&-1280X^8-2560X^7+1684X^6+4500X^5-380X^4-2152X^3-68X^2+212X+44.
\end{aligned}
\]
并具有共同结构分解
\[
A(X)=-(X^2-1)A_8(X),\qquad
B(X)=-4(X^2-1)B_6(X),
\]
其中
\[
A_8(X)=256X^8+2560X^7+1051X^6-2910X^5-1270X^4+892X^3+401X^2-94X-22,
\]
\[
B_6(X)=320X^6+640X^5-101X^4-485X^3-6X^2+53X+11.
\]
\end{conclusion}
\begin{proof}[证明要点]
将 \(y=X^2+Y-\frac12\) 代入 \(-y(y-1)P(y)\) 并在 \(\QQ[X,Y]/(Y^2-X^3+X-\frac14)\) 中约化。
由于 \(\QQ(E)=\QQ(X)\oplus \QQ(X)\,Y\)（作为 \(\QQ(X)\)-向量空间），任意元素可唯一写为 \(A(X)+B(X)Y\) 形，故表示唯一。
直接展开并收集项即可得到上述 \(A,B\)。
再对 \(A,B\) 作整系数因式分解，得公共因子 \((X^2-1)\) 与给定 \(A_8,B_6\)。证毕。
\end{proof}

\begin{conclusion}[曲线 \(C\) 的平面四次模型与双二次规范方程]\label{con:xi-terminal-zm-leyang-doublecover-plane-quartic-model}
令 \(C\) 为 \(\QQ(C)=\QQ(E)(w)\) 的光滑射影模型。
定义
\[
f(X):=X^3-X+\frac14,\qquad D(X):=4X^3-4X+1=4f(X).
\]
则 \(C\) 在 \((X,w)\)-平面上满足显式四次方程
\[
\bigl(w^2-A(X)\bigr)^2=B(X)^2f(X).
\]
等价地，可写作整数系数规范式
\[
\bigl(w^2+(X^2-1)A_8(X)\bigr)^2
=
4(X^2-1)^2B_6(X)^2D(X).
\]
\end{conclusion}
\begin{proof}[证明要点]
由结论 \ref{con:xi-terminal-zm-leyang-doublecover-linearization-a-plus-by}
\[
w^2=A(X)+B(X)Y
\]
可得
\[
\bigl(w^2-A(X)\bigr)^2=B(X)^2Y^2=B(X)^2f(X).
\]
再代入
\(
A=-(X^2-1)A_8,\ B=-4(X^2-1)B_6,\ D=4f
\)
并约去常数因子即得第二式。证毕。
\end{proof}

\begin{conclusion}[四次投影关于 \(w\) 的判别式完全因式分解]\label{con:xi-terminal-zm-leyang-quartic-projection-discriminant-full-factorization}
设
\[
H(w;X):=\bigl(w^2-A(X)\bigr)^2-B(X)^2f(X)\in\ZZ[X][w].
\]
则关于 \(w\) 的判别式满足
\[
\Disc_w\!\bigl(H(w;X)\bigr)
=4096\cdot \Xi(X)\cdot D(X)^2\cdot g(X)^2\cdot Q_6(X)\cdot B_6(X)^4,
\]
其中
\[
\Xi(X):=X(X-2)(X-1)^7(X+1)^7,
\]
\[
g(X):=16X^3-9X^2+1,
\]
\[
Q_6(X):=256X^6-480X^5+201X^4+750X^3-921X^2-78X+704.
\]
因此，投影 \(\pi_X:C\to \PP^1_X\) 的有限分歧 \(X\)-坐标包含于显式集合
\[
\left\{0,2,\pm1\right\}
\cup\{D(X)=0\}
\cup\{g(X)=0\}
\cup\{Q_6(X)=0\}
\cup\{B_6(X)=0\}.
\]
\end{conclusion}
\begin{proof}[证明要点]
写
\[
H(w;X)=w^4-2A(X)w^2+\bigl(A(X)^2-B(X)^2f(X)\bigr),
\]
\[
\partial_w H=4w\bigl(w^2-A(X)\bigr).
\]
由单首多项式判别式公式
\[
\Disc_w(H)=\operatorname{Res}_w\!\bigl(H,\partial_wH\bigr)
=4^4\operatorname{Res}_w\!\bigl(H,w(w^2-A)\bigr)
=256\,H(0;X)\,\operatorname{Res}_w\!\bigl(H,w^2-A\bigr).
\]
再由
\[
H(0;X)=A^2-B^2f,\qquad
\operatorname{Res}_w\!\bigl(H,w^2-A\bigr)=\bigl(-B^2f\bigr)^2,
\]
得
\[
\Disc_w(H)=256\,(A^2-B^2f)\,B^4f^2.
\]
代入 \(B=-4(X^2-1)B_6,\ f=D/4\) 可化为
\[
\Disc_w(H)=4096\,(X^2-1)^4\,B_6^4\,D^2\,(A^2-B^2f).
\]
对最后因子作完全分解
\[
A^2-B^2f
=X(X-2)(X-1)^3(X+1)^3\,g(X)^2\,Q_6(X),
\]
合并即得所述闭式分解。分歧点由判别式零点给出。证毕。
\end{proof}

\begin{conclusion}[\(A_4\)-子扩张中的单点分歧机制]\label{con:xi-terminal-zm-s4-a4-subextension-single-point-ramification}
设 \(L/\QQ(y)\) 为四次谱方程 \(F(\lambda,y)=0\) 的分裂域，且
\[
\mathrm{Gal}(L/\QQ(y))\cong S_4,\qquad
\QQ(H_{\mathrm{LY}})=L^{A_4}
=\QQ(y)\!\left(\sqrt{-y(y-1)P(y)}\right).
\]
记 \(\widetilde{C}\) 为 \(\QQ(\widetilde{C})=L\) 的光滑射影模型，则 \(\widetilde{C}\to H_{\mathrm{LY}}\) 为 \(A_4\)-伽罗瓦覆盖。
设 \(\widetilde{C}\to\PP^1_y\) 的有限分歧惯性由换位 \(\tau\) 生成，无穷远处分歧惯性由四循环 \(\rho\) 生成，则：
\begin{enumerate}
  \item 有限分歧值处，惯性与 \(A_4\) 的交为空，故在 \(L/L^{A_4}\) 中全部消失；
  \item 唯一残存分歧位于 \(H_{\mathrm{LY}}\) 上方无穷远点 \(\infty_{\mathrm{LY}}\)，其惯性群
  \[
  I_{\infty_{\mathrm{LY}}}
  =\langle \rho^2\rangle
  \cong C_2,
  \qquad \rho^2=(13)(24)\in A_4.
  \]
\end{enumerate}
故 \(\widetilde{C}\to H_{\mathrm{LY}}\) 为几何单点分歧覆盖，总分歧度为 \(6\)。
\end{conclusion}
\begin{proof}[证明要点]
由 \(A_4\triangleleft S_4\) 正规性，\(L/L^{A_4}\) 在给定基点上的惯性群为原惯性群与 \(A_4\) 的交。
对有限分歧，惯性由换位生成，换位不属于 \(A_4\)，故交为空。
对无穷远处，\(\rho\notin A_4\) 但 \(\rho^2\in A_4\) 且阶为 \(2\)，故交为 \(\langle\rho^2\rangle\)。
又 \(H_{\mathrm{LY}}\to\PP^1_y\) 在无穷远处分歧，故其上方仅有一点 \(\infty_{\mathrm{LY}}\)。
于是仅此一点分歧，且
\[
\deg\!\bigl(\mathrm{Ram}(\widetilde{C}/H_{\mathrm{LY}})\bigr)
=|A_4|\!\left(1-\frac12\right)
=12\cdot\frac12
=6.
\]
证毕。
\end{proof}

\begin{conclusion}[次数 \(4\) 非伽罗瓦覆盖的极小分歧型]\label{con:xi-terminal-zm-c3-quotient-to-hly-minimal-ramification-type}
令 \(C_3\le A_4\) 为任一阶 \(3\) 子群，记
\[
C:=\widetilde{C}/C_3.
\]
则自然映射 \(C\to H_{\mathrm{LY}}\) 为次数 \(4\) 覆盖，且：
\begin{enumerate}
  \item 除 \(\infty_{\mathrm{LY}}\) 上方外处处无分歧；
  \item 在 \(\infty_{\mathrm{LY}}\) 的纤维中，局部循环型为 \(2+2\)，即恰有两处简单分歧点；
  \item 总分歧度为 \(2\)，从而 \(g(C)=6\)。
\end{enumerate}
\end{conclusion}
\begin{proof}[证明要点]
由结论 \ref{con:xi-terminal-zm-s4-a4-subextension-single-point-ramification}，
\(\widetilde{C}\to H_{\mathrm{LY}}\) 的唯一分歧惯性由 \(\rho^2=(13)(24)\) 生成。
商覆盖 \(C=\widetilde{C}/C_3\) 的分歧由 \(\rho^2\) 在陪集集 \(A_4/C_3\)（大小 \(4\)）上的置换作用决定；
该作用同构于 \(A_4\) 的自然四点作用，故 \(\rho^2\) 的循环分解为两不交换位，即 \(2+2\)。
因此 \(\infty_{\mathrm{LY}}\) 上方有两处简单分歧点，总分歧度为 \(2\)。
由 Riemann--Hurwitz
\[
2g(C)-2
=4\bigl(2g(H_{\mathrm{LY}})-2\bigr)+2
=4(2\cdot2-2)+2
=10,
\]
得 \(g(C)=6\)。证毕。
\end{proof}

\begin{conclusion}[\texorpdfstring{$C_3$}{C3} 商双覆盖的分歧线性类型与 Prym 极化补偿]\label{con:xi-terminal-zm-prym-polarization-type-rigidity-11112}
设
\[
E=\widetilde{C}/S_3,\qquad
C=\widetilde{C}/C_3,\qquad
\pi:C\longrightarrow E
\]
为由 \(C_3\le S_3\) 诱导的二次覆盖。记 \(B_\pi\) 为其分歧除子，
\[
P:=\operatorname{Prym}(C/E),
\]
并以 \(\Xi\) 表示 \(\Jac(C)\) 主极化在 \(P\) 上诱导的 Prym 极化。
则成立：
\begin{enumerate}
  \item \(g(C)=6\)，且
  \[
  B_\pi\sim 10O\in \Pic(E);
  \]
  \item 任一满足
  \[
  L^{\otimes 2}\cong \mathcal O_E(B_\pi)
  \]
  的平方根线丛均可写为
  \[
  L\cong \mathcal O_E(5O+\eta),\qquad \eta\in E[2],
  \]
  因而这样的 \(L\) 恰有 \(4\) 个；
  \item \(\dim P=5\) 且
  \[
  \mathrm{type}(\Xi)=(1,2,2,2,2).
  \]
  特别地，若 \(L_P\) 为表示 \(\Xi\) 的充足线丛，则
  \[
  h^0(P,L_P)=16;
  \]
  \item 若在 \(\QQ\)-同源意义下 \(P\sim_{\QQ}J_{\mathrm{LY}}\times B\)，则可取与幂等分解相容的同源
  \[
  \varphi:J_{\mathrm{LY}}\times B\longrightarrow P
  \]
  使
  \[
  \varphi^*\Xi=\lambda_{\mathrm{LY}}\boxtimes \lambda_B,
  \qquad
  \deg\Xi=\deg\lambda_{\mathrm{LY}}\deg\lambda_B=16.
  \]
  特别地，若 \(\deg\lambda_{\mathrm{LY}}=4\)，则 \(\deg\lambda_B=4\)；反之亦然。
\end{enumerate}
\end{conclusion}
\begin{proof}[证明要点]
由结论 \ref{con:xi-terminal-zm-s4-c3-quotient-jacobian-three-factor-decomposition}，
\[
\Jac(C)\sim_{\QQ}J_{\mathrm{LY}}\times E\times B,
\]
故 \(g(C)=6\)。又 \(g(E)=1\)，于是 \(\dim P=g(C)-g(E)=5\)。
对二次覆盖 \(\pi:C\to E\)，设分歧点数为 \(r\)。
由 Hurwitz 公式
\[
2g(C)-2=2(2g(E)-2)+r
\]
得 \(r=10\)。
另一方面，结论 \ref{con:xi-terminal-zm-leyang-unramified-sixpoint-abel-closed-form-branch-balance} 已证明
\[
\bigoplus_{P\in\mathrm{Supp}(B_\pi)}P=O,
\]
这在椭圆曲线 \(E\simeq \Pic^0(E)\) 的群律翻译下等价于
\[
B_\pi-10O\sim 0,
\]
从而 \(B_\pi\sim 10O\)。

任一平方根线丛 \(L\) 满足
\[
L^{\otimes 2}\cong \mathcal O_E(B_\pi)\cong \mathcal O_E(10O),
\]
故
\[
L\in [2]^{-1}\!\bigl(\mathcal O_E(5O)\bigr).
\]
而 \([2]:\Pic^0(E)\to\Pic^0(E)\) 的核为 \(E[2]\)，故全部解构成一个 \(E[2]\)-torsor，亦即
\[
L\cong \mathcal O_E(5O+\eta),\qquad \eta\in E[2].
\]
对分歧二次覆盖到底曲线属 \(1\) 的 Prym，标准极化公式给出
\[
\mathrm{type}(\Xi)=\bigl(1^{g(E)},2^{r/2-1}\bigr)=(1,2,2,2,2).
\]
于是对任一表示 \(\Xi\) 的充足线丛 \(L_P\)，有
\[
h^0(P,L_P)=1\cdot 2^4=16.
\]

最后，由结论 \ref{con:xi-terminal-zm-s4-c3-prym-leyang-equivalence} 的同源分解
\[
P\sim_{\QQ}J_{\mathrm{LY}}\times B
\]
可取与幂等分解相容的同源 \(\varphi\)，使 \(\varphi^*\Xi\) 分裂为外积极化。按本文对极化次数的记法，极化型各项之积即为其次数，故
\[
\deg\Xi=1\cdot 2^4=16=\deg\lambda_{\mathrm{LY}}\deg\lambda_B.
\]
于是若一侧次数为 \(4\)，则另一侧亦必为 \(4\)。证毕。
\end{proof}

\paragraph{\texorpdfstring{$V_4$}{V4} 特征 Prym 分量与 \texorpdfstring{$C_3$}{C3} 商几何的同源细化}
以下结论均在既定 \(S_4\)-伽罗瓦闭包数据下成立：
\[
X:=\widetilde{C},\qquad
C:=X/A_4,\qquad
C_6:=X/V_4^{\mathrm{norm}},\qquad
Y:=X/C_3,\qquad
E:=X/S_3.
\]
其中 \(g(X)=16\)、\(g(C)=2\)、\(g(C_6)=4\)、\(g(E)=1\)，并记
\[
q:X\to C_6.
\]
对三个非平凡特征 \(\chi_i\in \widehat{V_4^{\mathrm{norm}}}\)（\(i=1,2,3\)），令
\[
D_i:=X/\ker(\chi_i),\qquad
P_i:=\operatorname{Prym}(D_i/C_6).
\]

\begin{conclusion}[\texorpdfstring{$V_4$}{V4} 特征 Prym 因子的 \texorpdfstring{$E\times B$}{E x B} 分裂]\label{con:xi-terminal-zm-v4-character-prym-e-b-splitting}
存在在 \(\QQ\)-同源意义下唯一确定的三维阿贝尔簇 \(B\)，使得对每个 \(i\in\{1,2,3\}\) 均有
\[
P_i\sim_{\QQ}E\times B.
\]
并且该分裂与 \(J(X)\) 的 \(\QQ[S_4]\)-模分解兼容：\(E\) 对应 \(\rho_3\)-等型分量的 \(\chi_i\)-特征部分，\(B\) 对应 \(\rho_3'\)-等型分量的 \(\chi_i\)-特征部分。
\end{conclusion}
\begin{proof}[证明要点]
由结论 \ref{con:xi-terminal-zm-s4-jacobian-group-algebra-prym-e3-b3}，
\[
J(X)\sim_{\QQ}J(C)\times E_{\mathrm{res}}^2\times E^3\times B^3.
\]
又因 \(V_4^{\mathrm{norm}}\triangleleft S_4\)，群代数幂等元 \(e_{\chi_i}\in \QQ[V_4^{\mathrm{norm}}]\) 与 \(S_4\)-中心幂等元可交换。对 \(V_4^{\mathrm{norm}}\) 限制后，\(\rho_3\) 与 \(\rho_3'\) 各含三个非平凡特征各一次，而 \(\varepsilon,\rho_2\) 不贡献非平凡特征部分，故
\[
e_{\chi_i}J(X)\sim_{\QQ}E\times B.
\]
另一方面，\(V_4^{\mathrm{norm}}\)-覆盖 \(q\) 的三个非平凡特征分量与三条双覆盖 \(D_i\to C_6\) 的 Prym 因子对应，故 \(P_i\sim_{\QQ}e_{\chi_i}J(X)\)。合并即得结论。证毕。
\end{proof}

\begin{conclusion}[Prym 的 \texorpdfstring{\(\Theta\)}{Theta} 奇点与 Andreotti--Mayer 轨道]\label{con:xi-terminal-zm-prym-theta-singular-andreotti-mayer}
对每个 \(i\)，设 \((P_i,\Xi_i)\) 为对应主极化 Prym 变体，则存在平移 \(t_i\in P_i(\CC)\) 使
\[
t_i+E\subset \operatorname{Sing}(\Xi_i).
\]
因此
\[
(P_i,\Xi_i)\in N_1\subset \mathcal A_4.
\]
\end{conclusion}
\begin{proof}[证明要点]
由结论 \ref{con:xi-terminal-zm-v4-character-prym-e-b-splitting}，\(P_i\) 在同源意义下含椭圆因子 \(E\)。
在该构造诱导的主极化下，\(E\)-方向给出 theta 除子的正维退化轨道，故奇点集含正维平移阿贝尔子簇。
据 Andreotti--Mayer 判据，这等价于落入 \(N_1\)。证毕。
\end{proof}

\begin{conclusion}[\texorpdfstring{$C_3$}{C3} 商曲线的 Jacobian 三因子分解]\label{con:xi-terminal-zm-c3-quotient-jacobian-jc-e-b}
设 \(Y:=X/C_3\)。则 \(g(Y)=6\)，且
\[
J(Y)\sim_{\QQ}J(C)\times E\times B.
\]
\end{conclusion}
\begin{proof}[证明要点]
由结论 \ref{con:xi-terminal-zm-s4-conjugacy-fixedpoints-local-eigenvalue-full-count}，三循环在 \(X\) 上无定点，故 \(X\to Y\) 为三重非分歧循环覆盖，Riemann--Hurwitz 给出 \(g(Y)=6\)。
分解式即结论 \ref{con:xi-terminal-zm-s4-c3-quotient-jacobian-three-factor-decomposition} 的同义改写（其中该结论中的 \(C\) 记号对应当前 \(Y\)）。证毕。
\end{proof}

\begin{conclusion}[三重非分歧 Prym 的 \texorpdfstring{$\ZZ[\zeta_3]$}{Z[zeta3]}-PEL 结构]\label{con:xi-terminal-zm-c3-etale-prym-pel-signature-polarization}
对三重非分歧循环覆盖
\[
\pi_3:X\to Y
\]
定义
\[
P_{(3)}:=\operatorname{Prym}(X/Y).
\]
则 \(\dim P_{(3)}=10\)，并满足：
\begin{enumerate}
  \item 存在自然嵌入
  \[
  \ZZ[\zeta_3]\hookrightarrow \operatorname{End}(P_{(3)}),
  \]
  且 Rosati 反演在该子环上对应复共轭；
  \item 在 \(H^{1,0}(P_{(3)})\) 上的两条非平凡特征重数均为 \(5\)，即 signature 为 \((5,5)\)；
  \item 诱导极化型为
  \[
  (1^5,3^5).
  \]
\end{enumerate}
\end{conclusion}
\begin{proof}[证明要点]
\(\pi_3\) 的覆盖群由阶 \(3\) 自同构 \(\sigma\) 生成，故在 Prym 部分上给出 \(\ZZ[\sigma]/(\sigma^2+\sigma+1)\cong \ZZ[\zeta_3]\) 作用；Rosati 与 \(\sigma\mapsto \sigma^{-1}\) 对应，遂为复共轭。
又 \(g(X)=16\)、\(g(Y)=6\)，故 \(\dim P_{(3)}=10\)。
不变部分维数为 \(6\)，其余 \(10\) 维按共轭成对分布，故两非平凡特征重数均为 \(5\)。
对三重非分歧循环覆盖，Prym 极化型为 \((1^{g(Y)-1},3^{g(Y)-1})\)，代入 \(g(Y)=6\) 即得 \((1^5,3^5)\)。证毕。
\end{proof}

\begin{conclusion}[\texorpdfstring{$Y$}{Y} 的双椭圆结构与双覆盖 Prym 同源型]\label{con:xi-terminal-zm-c3-quotient-bielliptic-prym-jc-b}
存在对合
\[
\iota\in\operatorname{Aut}(Y)
\]
使得
\[
Y/\langle\iota\rangle\cong E,
\]
且对应二次覆盖 \(Y\to E\) 的分歧点数为 \(10\)。
记
\[
P_{(2)}:=\operatorname{Prym}(Y/E),
\]
则
\[
P_{(2)}\sim_{\QQ}J(C)\times B.
\]
\end{conclusion}
\begin{proof}[证明要点]
由 \(N_{S_4}(C_3)=S_3\) 得
\[
N_{S_4}(C_3)/C_3\cong C_2,
\]
从而在 \(Y=X/C_3\) 上诱导对合 \(\iota\)；
商的函子性给出
\[
Y/\langle\iota\rangle\cong X/S_3=E.
\]
由 \(g(Y)=6,g(E)=1\) 的 Hurwitz 公式得分歧点数 \(10\)。
又对二次覆盖有
\[
J(Y)\sim_{\QQ}E\times P_{(2)},
\]
并结合结论 \ref{con:xi-terminal-zm-c3-quotient-jacobian-jc-e-b} 即得
\[
P_{(2)}\sim_{\QQ}J(C)\times B.
\]
证毕。
\end{proof}

\begin{conclusion}[\texorpdfstring{$S_4$}{S4} 作用强制的 N\'eron--Severi 秩下界]\label{con:xi-terminal-zm-jacobian-neron-severi-rank-lower-bound-17}
\[
\rho\!\bigl(J(X)\bigr)\ge 17.
\]
更精确地，\(\QQ[S_4]\subset \operatorname{End}^0(J(X))\) 在 Rosati 反演下的不变子空间维数为 \(17\)，从而在
\[
\mathrm{NS}(J(X))_{\QQ}\cong \operatorname{End}^0(J(X))^{\dagger=1}
\]
中给出显式 \(17\) 维子空间。
\end{conclusion}
\begin{proof}[证明要点]
群作用诱导嵌入 \(\QQ[S_4]\hookrightarrow \operatorname{End}^0(J(X))\)，且 Rosati 在群代数上对应 \(g\mapsto g^{-1}\)。
又
\[
\QQ[S_4]\cong \QQ\oplus \QQ\oplus M_2(\QQ)\oplus M_3(\QQ)\oplus M_3(\QQ).
\]
在各块上 Rosati 对应转置，故不变维数为
\[
1+1+\frac{2\cdot3}{2}+\frac{3\cdot4}{2}+\frac{3\cdot4}{2}=17.
\]
由 \(\mathrm{NS}(A)_{\QQ}\cong \operatorname{End}^0(A)^{\dagger=1}\)（\(A\) 为主极化阿贝尔簇）即得结论。证毕。
\end{proof}

\begin{conclusion}[三循环商曲线的四割性]\label{con:xi-terminal-zm-c3-quotient-gonality-four}
\[
\operatorname{gon}(Y)=4.
\]
\end{conclusion}
\begin{proof}[证明要点]
由结论 \ref{con:xi-terminal-zm-c3-quotient-bielliptic-prym-jc-b}，存在次数 \(2\) 映射 \(Y\to E\)。
再与椭圆曲线的次数 \(2\) 到 \(\PP^1\) 态射复合，得 \(\operatorname{gon}(Y)\le 4\)。
若存在次数 \(d\le 3\) 的 \(Y\to \PP^1\)，与次数 \(2\) 的 \(Y\to E\) 联用 Castelnuovo--Severi 不等式：
\[
g(Y)\le d\,g(E)+2\,g(\PP^1)+(d-1)(2-1)=2d-1\le 5,
\]
与 \(g(Y)=6\) 矛盾，故 \(\operatorname{gon}(Y)\ge 4\)。
综上 \(\operatorname{gon}(Y)=4\)。证毕。
\end{proof}


\paragraph{\texorpdfstring{$S_4$}{S4} 子群格中低亏格商与双覆盖 Prym 的补充显式化}
沿用结论 \ref{con:xi-terminal-zm-s4-jacobian-group-algebra-prym-e3-b3} 与
\ref{con:xi-terminal-zm-s4-chevalley-weil-explicit-decomposition} 的记号，记
\[
\Jac(\widetilde C)\sim_{\QQ}J_{\mathrm{LY}}\times E_{\mathrm{res}}^2\times E^3\times B^3,
\]
\[
H^0(\widetilde C,\Omega^1)\cong 2\varepsilon\oplus \rho_2\oplus \rho_3\oplus 3\rho_3',
\]
并置
\[
V_4^{\mathrm{tr}}:=\langle(12),(34)\rangle.
\]

\begin{conclusion}[子群格中三条低亏格商曲线的 Jacobian 同源型显式化]\label{con:xi-terminal-zm-s4-lowgenus-subgroup-quotients-jacobian-explicit}
记
\[
C_4:=\widetilde C/\langle(1234)\rangle,\qquad
C_{V_4^{\mathrm{norm}}}:=\widetilde C/V_4^{\mathrm{norm}},\qquad
C_{V_4^{\mathrm{tr}}}:=\widetilde C/V_4^{\mathrm{tr}}.
\]
则
\[
g(C_4)=4,\qquad \Jac(C_4)\sim_{\QQ}E_{\mathrm{res}}\times B;
\]
\[
g(C_{V_4^{\mathrm{norm}}})=4,\qquad
\Jac(C_{V_4^{\mathrm{norm}}})\sim_{\QQ}J_{\mathrm{LY}}\times E_{\mathrm{res}}^2;
\]
\[
g(C_{V_4^{\mathrm{tr}}})=2,\qquad
\Jac(C_{V_4^{\mathrm{tr}}})\sim_{\QQ}E_{\mathrm{res}}\times E.
\]
\end{conclusion}
\begin{proof}[证明要点]
对任意子群 \(H\le S_4\)，有
\[
H^0(\widetilde C/H,\Omega^1)\cong H^0(\widetilde C,\Omega^1)^H,
\]
故只需计算
\(
2\varepsilon\oplus \rho_2\oplus \rho_3\oplus 3\rho_3'
\)
在各子群下的不变向量维数。对任意不可约表示 \(W\)，由角色平均公式
\[
\dim W^H=\frac{1}{|H|}\sum_{h\in H}\chi_W(h).
\]

对 \(H=\langle(1234)\rangle\)，其元素型为
\(
1,(1234),(13)(24),(1432)
\)，从而
\[
\dim(\varepsilon^H)=0,\quad
\dim(\rho_2^H)=1,\quad
\dim(\rho_3^H)=0,\quad
\dim((\rho_3')^H)=1.
\]
于是
\[
g(C_4)=2\cdot 0+1\cdot 1+1\cdot 0+3\cdot 1=4,
\]
并由等型对应得到
\[
\Jac(C_4)\sim_{\QQ}E_{\mathrm{res}}\times B.
\]

对 \(H=V_4^{\mathrm{norm}}\)，其非平凡元均为双换位，故
\[
\dim(\varepsilon^H)=1,\quad
\dim(\rho_2^H)=2,\quad
\dim(\rho_3^H)=0,\quad
\dim((\rho_3')^H)=0.
\]
因此
\[
g(C_{V_4^{\mathrm{norm}}})=2\cdot 1+1\cdot 2=4,
\]
且
\[
\Jac(C_{V_4^{\mathrm{norm}}})\sim_{\QQ}J_{\mathrm{LY}}\times E_{\mathrm{res}}^2.
\]

对 \(H=V_4^{\mathrm{tr}}\)，其元素型为
\(
1,(12),(34),(12)(34)
\)，故
\[
\dim(\varepsilon^H)=0,\quad
\dim(\rho_2^H)=1,\quad
\dim(\rho_3^H)=1,\quad
\dim((\rho_3')^H)=0.
\]
于是
\[
g(C_{V_4^{\mathrm{tr}}})=1+1=2,
\]
并且
\[
\Jac(C_{V_4^{\mathrm{tr}}})\sim_{\QQ}E_{\mathrm{res}}\times E.
\]
证毕。
\end{proof}

\begin{conclusion}[两条属 \(4\) 双覆盖到 \texorpdfstring{$E_{\mathrm{res}}$}{Eres} 的 Prym 完全辨识]\label{con:xi-terminal-zm-s4-two-genus4-doublecovers-prym-identification}
自然商映射
\[
\pi_{C_4}:C_4\longrightarrow E_{\mathrm{res}},\qquad
\pi_{V_4}:C_{V_4^{\mathrm{norm}}}\longrightarrow E_{\mathrm{res}}
\]
均为分歧二次覆盖，并满足
\[
g(C_4)=g(C_{V_4^{\mathrm{norm}}})=4,\qquad
\deg\operatorname{Branch}(\pi_{C_4})=\deg\operatorname{Branch}(\pi_{V_4})=6.
\]
此外，
\[
\operatorname{Prym}(C_4/E_{\mathrm{res}})\sim_{\QQ}B,
\]
\[
\operatorname{Prym}(C_{V_4^{\mathrm{norm}}}/E_{\mathrm{res}})\sim_{\QQ}J_{\mathrm{LY}}\times E_{\mathrm{res}},
\]
且两条覆盖的 Prym 极化型同为
\[
(1,2,2).
\]
\end{conclusion}
\begin{proof}[证明要点]
由 \(\langle(1234)\rangle\subset D_8\) 与 \(V_4^{\mathrm{norm}}\subset D_8\) 的指数均为 \(2\)，可知两条商曲线都带有到
\(
E_{\mathrm{res}}=\widetilde C/D_8
\)
的自然二次映射。由结论
\ref{con:xi-terminal-zm-s4-lowgenus-subgroup-quotients-jacobian-explicit}，
\[
g(C_4)=g(C_{V_4^{\mathrm{norm}}})=4.
\]
再由 Hurwitz 公式
\[
2g(Y)-2=2\bigl(2g(E_{\mathrm{res}})-2\bigr)+\deg\operatorname{Branch}(Y/E_{\mathrm{res}})
\]
并代入 \(g(E_{\mathrm{res}})=1\)，得两条覆盖的分支次数均为 \(6\)。

对分歧二次覆盖的 Jacobian 标准分解
\[
\Jac(Y)\sim_{\QQ}E_{\mathrm{res}}\times \operatorname{Prym}(Y/E_{\mathrm{res}})
\]
与结论
\ref{con:xi-terminal-zm-s4-lowgenus-subgroup-quotients-jacobian-explicit}
比较，可得
\[
\operatorname{Prym}(C_4/E_{\mathrm{res}})\sim_{\QQ}B,
\qquad
\operatorname{Prym}(C_{V_4^{\mathrm{norm}}}/E_{\mathrm{res}})\sim_{\QQ}J_{\mathrm{LY}}\times E_{\mathrm{res}}.
\]
最后，对底曲线属 \(1\)、分支次数 \(6\) 的分歧二次覆盖，Prym 极化型由一般公式给出为
\[
\bigl(1^{g(E_{\mathrm{res}})},2^{6/2-1}\bigr)=(1,2,2).
\]
证毕。
\end{proof}

\begin{conclusion}[\texorpdfstring{$V_4^{\mathrm{tr}}$}{V4tr} 商 Jacobian 分裂的部分塔式可见性与剩余 endoscopic 因子]\label{con:xi-terminal-zm-s4-v4tr-splitting-partial-tower-visibility}
对
\[
C_{V_4^{\mathrm{tr}}}:=\widetilde C/V_4^{\mathrm{tr}}
\]
有
\[
\Jac(C_{V_4^{\mathrm{tr}}})\sim_{\QQ}E_{\mathrm{res}}\times E.
\]
其中 \(E_{\mathrm{res}}\) 因子可由包含链
\[
V_4^{\mathrm{tr}}\subset D_8^{\mathrm{tr}}:=N_{S_4}\!\bigl(\langle(1324)\rangle\bigr)
\]
几何实现，因为 \(\widetilde C/D_8^{\mathrm{tr}}\cong E_{\mathrm{res}}\)；然而不存在 \(S_3\le S_4\) 含 \(V_4^{\mathrm{tr}}\)，故 \(E\)-因子不能由任何属 \(1\) 子群塔商直接产生。因而该属 \(2\) Jacobian 的椭圆分裂并非完全由子群塔解释，其中 \(E\)-方向仍由 \(\QQ[S_4]\)-幂等投影强制出现。
\end{conclusion}
\begin{proof}[证明要点]
第一断言即结论
\ref{con:xi-terminal-zm-s4-lowgenus-subgroup-quotients-jacobian-explicit}
的第三个分裂式。

取
\[
D_8^{\mathrm{tr}}=\langle(1324),(12)\rangle.
\]
则
\[
(1324)^2=(12)(34)\in V_4^{\mathrm{tr}},\qquad
(12)(12)(34)=(34),
\]
从而 \(V_4^{\mathrm{tr}}=\langle(12),(34)\rangle\subset D_8^{\mathrm{tr}}\)。又所有 \(D_8\) 子群在 \(S_4\) 中彼此共轭，故
\[
\widetilde C/D_8^{\mathrm{tr}}\cong \widetilde C/D_8=E_{\mathrm{res}},
\]
于是 \(E_{\mathrm{res}}\) 因子可由子群塔商解释。

另一方面，由结论 \ref{con:xi-terminal-zm-s4-hurwitz-subgroup-spectrum-cycletypes}，属 \(1\) 的子群商仅来自 \(D_8\) 与 \(S_3\) 两类子群。任一 \(S_3\) 作为点稳定子都不可能同时含有两条互不共享固定点的换位 \((12)\) 与 \((34)\)，故 \(V_4^{\mathrm{tr}}\not\subset S_3\)。因此 \(E\)-因子不能来自任何包含 \(V_4^{\mathrm{tr}}\) 的属 \(1\) 子群商。综上，\(\Jac(C_{V_4^{\mathrm{tr}}})\) 的分裂只有 \(E_{\mathrm{res}}\) 一侧具备塔式几何来源，而 \(E\)-方向仍是表示论意义下的剩余 endoscopic 因子。证毕。
\end{proof}

\endinput

\begin{theorem}[三类换位碰撞边界下四个 Langlands 分量的环面秩与导子指数]\label{thm:xi-terminal-zm-s4-langlands-factor-torus-rank-conductor-table}
\leanverified{paper\_xi\_terminal\_zm\_s4\_langlands\_factor\_torus\_rank\_conductor\_table}
沿用定理 \ref{thm:xi-s4-isotypic-jacobian-toric-rank-by-boundary-type} 的三类换位碰撞边界记号 \(r\in\{1,2,3\}\)，以及结论 \ref{con:xi-terminal-zm-s4-jacobian-group-algebra-prym-e3-b3} 的 \(\QQ\)-同源分解
\[
\Jac(\widetilde C)\ \sim_{\QQ}\ J_{\mathrm{LY}}\times E_{\mathrm{res}}^2\times E^3\times B^3
\]
与等型对应
\(
\varepsilon\leftrightarrow J_{\mathrm{LY}},\
\rho_2\leftrightarrow E_{\mathrm{res}}^2,\
\rho_3\leftrightarrow E^3,\
\rho_3'\leftrightarrow B^3
\)。
对每个分量 \(A\in\{J_{\mathrm{LY}},E_{\mathrm{res}},E,B\}\)，记其在对应半稳定退化处的环面秩为 \(t(A)\)。
则对三类碰撞型 \(r\) 有表：
\[
\begin{array}{c|cccc}
r & t(J_{\mathrm{LY}}) & t(E_{\mathrm{res}}) & t(E) & t(B)\\
\hline
1 & 1 & 1 & 1 & 2\\
2 & 0 & 0 & 1 & 1\\
3 & 1 & 0 & 0 & 1
\end{array}
\]
并且在特征 \(0\) 的半稳定几何退化处（或更一般地在温和半稳定情形下），\(\ell\)-进 Tate 模的惯性表示无 Swan 部分，因而 Artin 导子指数等于环面秩。
\end{theorem}
\begin{proof}
由定理 \ref{thm:xi-s4-isotypic-jacobian-toric-rank-by-boundary-type}，在碰撞型 \(r\) 处四个 \(\QQ[S_4]\)-等型因子 \(A_\varepsilon,A_{\rho_2},A_{\rho_3},A_{\rho_3'}\) 的环面秩
\((t_\varepsilon,t_{\rho_2},t_{\rho_3},t_{\rho_3'})\)
由该定理中的表格确定。
另一方面，结论 \ref{con:xi-terminal-zm-s4-jacobian-group-algebra-prym-e3-b3} 的等型对应表明
\[
A_\varepsilon\sim_{\QQ} J_{\mathrm{LY}},\qquad
A_{\rho_2}\sim_{\QQ}E_{\mathrm{res}}^2,\qquad
A_{\rho_3}\sim_{\QQ}E^3,\qquad
A_{\rho_3'}\sim_{\QQ}B^3.
\]
环面秩对同源不变且对直积可加，故
\[
t(J_{\mathrm{LY}})=t_\varepsilon,\qquad
2\,t(E_{\mathrm{res}})=t_{\rho_2},\qquad
3\,t(E)=t_{\rho_3},\qquad
3\,t(B)=t_{\rho_3'}.
\]
将三类 \(r\) 的数值代入即得表格。
\par
最后，在特征 \(0\) 的半稳定几何退化处（或温和半稳定情形下），惯性表示的 Swan 部分为 \(0\)，而 Artin 导子指数等于环面秩与 Swan 部分之和，从而导子指数等于环面秩。证毕。
\end{proof}

\begin{corollary}[碰撞选择律与 \texorpdfstring{$B$}{B}-通道的普遍非良约化]\label{cor:xi-terminal-zm-s4-langlands-factor-collision-selection}
\leanverified{paper\_xi\_terminal\_zm\_s4\_langlands\_factor\_collision\_selection}
在定理 \ref{thm:xi-terminal-zm-s4-langlands-factor-torus-rank-conductor-table} 的设定下：
\begin{enumerate}
  \item 对所有三类换位碰撞边界均有 \(t(B)\ge 1\)，因此 \(B\) 在这些边界处均为半稳定但非良约化；并且在 \(r=1\) 型边界上有 \(t(B)=2\)。
  \item \(J_{\mathrm{LY}}\) 的良约化仅发生在 \(r=2\) 型边界；\(E_{\mathrm{res}}\) 的良约化发生在 \(r\in\{2,3\}\)；\(E\) 的良约化仅发生在 \(r=3\) 型边界。
\end{enumerate}
\end{corollary}
\begin{proof}
良约化等价于环面秩为 \(0\)。直接读取定理 \ref{thm:xi-terminal-zm-s4-langlands-factor-torus-rank-conductor-table} 的表格即可。证毕。
\end{proof}

\begin{theorem}[两类 Prym 的环面秩分层]\label{thm:xi-terminal-zm-s4-prym-torus-rank-layering}
\leanverified{paper\_xi\_terminal\_zm\_s4\_prym\_torus\_rank\_layering}
沿用结论 \ref{con:xi-terminal-zm-s4-involution-prym-polarization-eigensplitting} 的记号，
令 \(\pi_\tau:\widetilde C\to C_\tau\) 为换位商诱导的二次覆盖，\(\pi_\delta:\widetilde C\to C_\delta\) 为双换位商诱导的二次覆盖。
则在三类换位碰撞边界 \(r\in\{1,2,3\}\) 上，其 Prym 的环面秩分别为
\[
t\!\bigl(\Prym(\pi_\delta)\bigr)=(6,4,2),
\qquad
t\!\bigl(\Prym(\pi_\tau)\bigr)=(7,3,3),
\]
其中括号中三元组按 \(r=(1,2,3)\) 顺序排列。
\end{theorem}
\begin{proof}
由结论 \ref{con:xi-terminal-zm-s4-involution-prym-polarization-eigensplitting} 的 \(\QQ\)-同源分解，
\[
\Prym(\pi_\tau)\sim_{\QQ} J_{\mathrm{LY}}\times E_{\mathrm{res}}\times E\times B^2,
\qquad
\Prym(\pi_\delta)\sim_{\QQ} E^2\times B^2.
\]
环面秩对直积可加并对同源不变，故
\[
t\!\bigl(\Prym(\pi_\delta)\bigr)=2t(E)+2t(B),
\qquad
t\!\bigl(\Prym(\pi_\tau)\bigr)=t(J_{\mathrm{LY}})+t(E_{\mathrm{res}})+t(E)+2t(B).
\]
将定理 \ref{thm:xi-terminal-zm-s4-langlands-factor-torus-rank-conductor-table} 的三类 \(r\) 下 \((t(J_{\mathrm{LY}}),t(E_{\mathrm{res}}),t(E),t(B))\) 逐项代入即可得到所示数值。证毕。
\end{proof}

\endinput

\subsubsection{\texorpdfstring{$S_4$}{S4} 正则闭包的整体谱：属数、Chevalley--Weil 分解、正规子群商与 Jacobian 表示分裂}\label{subsubsec:xi-s4-regular-closure-global-spectrum}
沿用
\[
\Pi(\lambda,y)=\lambda^{4}-\lambda^{3}-(2y+1)\lambda^{2}+\lambda+y(y+1)\in\QQ(y)[\lambda],
\]
并记
\[
g(y):=256y^{3}+411y^{2}+165y+32,\qquad
\Delta(y):=-y(y-1)g(y).
\]
设 \(L/\QQ(y)\) 为 \(\Pi\) 的分裂域，\(G:=\mathrm{Gal}(L/\QQ(y))\cong S_4\)，并令 \(X\) 为 \(\QQ(X)=L\) 的光滑射影模型（即 \(X=\widetilde X\)）。

\begin{conclusion}[\texorpdfstring{$S_4$}{S4} 正则闭包曲线的精确属数]\label{con:xi-terminal-zm-s4-regular-closure-exact-genus-16}
覆盖 \(X\to\PP^1_y\) 为次数 \(24\) 的正则伽罗瓦覆盖，且
\[
g(X)=16.
\]
\end{conclusion}
\begin{proof}[证明要点]
由结论 \ref{con:xi-terminal-zm-s4-galois-closure-genus-tower-holomorphic-differentials}，分歧护照为一个阶 \(4\) 惯性点与五个阶 \(2\) 惯性点，即 \((4)\cdot(2)^5\)。
对正则伽罗瓦覆盖应用 Hurwitz 公式
\[
g(X)=1+\frac{|S_4|}{2}\!\left(-2+\sum_{p\in B}\Bigl(1-\frac1{e_p}\Bigr)\right),
\]
代入 \(|S_4|=24\)、\((e_p)=(4,2,2,2,2,2)\) 即得
\[
g(X)=1+12\!\left(-2+\frac34+5\cdot\frac12\right)=16.
\]
证毕。
\end{proof}

\begin{conclusion}[\texorpdfstring{$S_4$}{S4} 正则闭包中共轭类定点计数与循环商曲线属数]\label{con:xi-terminal-zm-s4-conjugacy-class-fixedpoints-and-cyclic-quotient-genera}
设
\(
\varphi:X\to \PP^1_y
\)
为结论 \ref{con:xi-terminal-zm-s4-regular-closure-exact-genus-16} 中的正则 \(S_4\)-覆盖，其分歧惯性型为一个 \(C_4\) 与五个 \(C_2\)。
对共轭类代表 \(g\in S_4\)，记 \(\mathrm{Fix}_X(g)\) 为其在 \(X\) 上的定点集，则有
\[
\begin{array}{c|c|c}
g\text{ 的循环型} & \#\mathrm{Fix}_X(g) & g(X/\langle g\rangle)\\
\hline
(2) & 10 & 6\\
(2,2) & 2 & 8\\
(4) & 2 & 4\\
(3) & 0 & 6
\end{array}
\]
\end{conclusion}
\begin{proof}[证明]
设 \(b\in \PP^1_y\) 为分歧值，\(I_b\le S_4\) 为其惯性群。由正则覆盖的陪集描述，
\(\varphi^{-1}(b)\cong S_4/I_b\)，且 \(xI_b\) 对应点的稳定子为 \(xI_bx^{-1}\)。故对任意 \(g\in S_4\) 有
\begin{equation}\label{eq:xi-terminal-zm-s4-fixedpoint-fiber-coset-formula}
\#\bigl(\mathrm{Fix}_X(g)\cap \varphi^{-1}(b)\bigr)
=\frac{\#\{x\in S_4\mid x^{-1}gx\in I_b\}}{|I_b|}.
\end{equation}

先取换位 \(\tau\)。五个有限分歧点 \(b_1,\dots,b_5\) 的惯性均为 \(C_2\)，其非平凡元均为换位。
由 \eqref{eq:xi-terminal-zm-s4-fixedpoint-fiber-coset-formula}，每个 \(b_i\) 的贡献为
\[
\frac{\#\{x\in S_4\mid x^{-1}\tau x=\tau_i\}}{2}
=\frac{|C_{S_4}(\tau)|}{2}
=\frac{4}{2}=2,
\]
故 \(\#\mathrm{Fix}_X(\tau)=10\)。对阶 \(2\) 商映射 \(X\to X/\langle\tau\rangle\) 用 Hurwitz 公式并代入 \(g(X)=16\)：
\[
2g(X)-2=2\bigl(2g(X/\langle\tau\rangle)-2\bigr)+\#\mathrm{Fix}_X(\tau),
\]
得 \(30=4g'-4+10\)，故 \(g'=6\)。

再取双换位 \(\delta\)。有限分歧惯性不含双换位，故仅无穷远分歧值 \(b_\infty\) 可能贡献。
记其惯性 \(I_{b_\infty}=\langle c\rangle\cong C_4\)，则唯一阶 \(2\) 元 \(c^2\) 为双换位。
由 \eqref{eq:xi-terminal-zm-s4-fixedpoint-fiber-coset-formula}
\[
\#\mathrm{Fix}_X(\delta)
=\frac{\#\{x\in S_4\mid x^{-1}\delta x=c^2\}}{4}
=\frac{|C_{S_4}(\delta)|}{4}
=\frac{8}{4}=2.
\]
再由 Hurwitz（阶 \(2\)）得
\[
30=2\bigl(2g(X/\langle\delta\rangle)-2\bigr)+2,
\]
故 \(g(X/\langle\delta\rangle)=8\)。

取四循环 \(\gamma\)。同理仅 \(b_\infty\) 贡献。
由于四循环同属一共轭类，可不失一般性取 \(\gamma=c\)。
此时式 \eqref{eq:xi-terminal-zm-s4-fixedpoint-fiber-coset-formula} 的条件 \(x^{-1}\gamma x\in \langle\gamma\rangle\) 等价于 \(x\in N_{S_4}(\langle\gamma\rangle)\)，
而 \(|N_{S_4}(\langle\gamma\rangle)|=8\)，故
\[
\#\mathrm{Fix}_X(\gamma)=\frac{8}{4}=2.
\]
又 \(\gamma^2\) 为双换位，已知 \(\#\mathrm{Fix}_X(\gamma^2)=2\)，且 \(\mathrm{Fix}_X(\gamma)\subseteq \mathrm{Fix}_X(\gamma^2)\)，
从而两者相等，不存在“仅被 \(\gamma^2\) 固定”之点。
故 \(X\to X/\langle\gamma\rangle\) 的分歧仅来自这 \(2\) 个点，且各分歧指数为 \(4\)。Hurwitz（阶 \(4\)）给出
\[
30=4\bigl(2g(X/\langle\gamma\rangle)-2\bigr)+2(4-1),
\]
解得 \(g(X/\langle\gamma\rangle)=4\)。

最后取三循环 \(\sigma\)。各惯性群仅含阶 \(2\) 或 \(4\) 元，故 \(\sigma\) 不可能固定分歧纤维点，从而 \(\mathrm{Fix}_X(\sigma)=\varnothing\)。
于是 \(X\to X/\langle\sigma\rangle\) 无分歧，Hurwitz（阶 \(3\)）给出
\[
30=3\bigl(2g(X/\langle\sigma\rangle)-2\bigr),
\]
故 \(g(X/\langle\sigma\rangle)=6\)。证毕。
\end{proof}

\begin{conclusion}[四循环商的 Prym 三维因子与 \texorpdfstring{$9$}{9} 维等型分量识别]\label{con:xi-terminal-zm-s4-c4-prym-threefold-and-v3sgn-identification}
取 \(c\in S_4\) 为无穷远惯性生成元（四循环），定义
\[
Y:=X/\langle c\rangle,\qquad
E_{\mathrm D}:=X/D_8
\]
（其中 \(D_8=N_{S_4}(\langle c\rangle)\)）。令
\[
\pi:Y\longrightarrow E_{\mathrm D}
\]
为诱导的次数 \(2\) 映射，并定义
\[
B:=\mathrm{Prym}(Y/E_{\mathrm D})
=\ker\!\bigl(\mathrm{Nm}_\pi:\mathrm{Jac}(Y)\to \mathrm{Jac}(E_{\mathrm D})\bigr)^0.
\]
则有：
\begin{enumerate}
  \item \(\pi\) 为分歧双覆盖，且 \(\deg\mathrm{Branch}(\pi)=6\)，从而 \(\dim B=3\)；
  \item \(\mathrm{Jac}(Y)\sim_{\QQ}E_{\mathrm D}\times B\)；
  \item 对任意换位 \(\tau\) 与双换位 \(\delta\)，有
  \[
  \mathrm{Jac}(X/\langle\tau\rangle)\sim_{\QQ}E^2\times E_{\mathrm D}\times B,
  \]
  \[
  \mathrm{Jac}(X/\langle\delta\rangle)\sim_{\QQ}\mathrm{Jac}(X_A)\times E\times E_{\mathrm D}^2\times B;
  \]
  \item 若 \(A_{V_3\otimes\mathrm{sgn}}\subset \mathrm{Jac}(X)\) 表示对应 \(V_3\otimes\mathrm{sgn}\) 的 \(9\) 维等型分量，则
  \[
  A_{V_3\otimes\mathrm{sgn}}\sim_{\QQ}B^3.
  \]
\end{enumerate}
\end{conclusion}
\begin{proof}[证明]
由子群塔 \(\langle c\rangle\subset D_8\subset S_4\) 可知 \(\deg(\pi)=2\)。
又由结论 \ref{con:xi-terminal-zm-s4-conjugacy-class-fixedpoints-and-cyclic-quotient-genera} 与
\ref{con:xi-terminal-zm-s4-hurwitz-subgroup-spectrum-cycletypes}，
\[
g(Y)=g(X/\langle c\rangle)=4,\qquad g(E_{\mathrm D})=g(X/D_8)=1.
\]
Hurwitz 公式
\[
2g(Y)-2=2\bigl(2g(E_{\mathrm D})-2\bigr)+\deg\mathrm{Branch}(\pi)
\]
给出 \(\deg\mathrm{Branch}(\pi)=6\)，并有
\(
\dim B=g(Y)-g(E_{\mathrm D})=3
\)。
这证明第 \(1\) 点。

第 \(2\) 点由双覆盖的 Prym 标准分解
\[
\mathrm{Jac}(Y)\sim_{\QQ}\pi^*\mathrm{Jac}(E_{\mathrm D})\times \mathrm{Prym}(Y/E_{\mathrm D})
\]
直接得到。

第 \(3\) 点与结论
\ref{con:xi-terminal-zm-s4-transposition-quotient-jacobian-full-decomposition}、
\ref{con:xi-terminal-zm-s4-doubletransposition-quotient-jacobian-full-decomposition}
一致：其中文中记号 \(E_{\mathrm{res}}\) 即此处 \(E_{\mathrm D}\)，且
\(
J_{\mathrm{LY}}\sim_{\QQ}\mathrm{Jac}(X_A)
\)
见结论
\ref{con:xi-terminal-zm-s4-quotient-genus-spectrum-and-jacobian-representation-splitting}。

第 \(4\) 点可由两种等价方式得到：其一，由结论
\ref{con:xi-terminal-zm-s4-jacobian-group-algebra-prym-e3-b3}
的分解
\[
\mathrm{Jac}(X)\sim_{\QQ}\mathrm{Jac}(X_A)\times E_{\mathrm D}^2\times E^3\times B^3
\]
直接识别 \(V_3\otimes\mathrm{sgn}\) 等型即为 \(B^3\)；
其二，由
\(
e_{V_3\otimes\mathrm{sgn}}\QQ[S_4]e_{V_3\otimes\mathrm{sgn}}\cong M_3(\QQ)
\)
与原始幂等分解得到同结论。证毕。
\end{proof}

\begin{conclusion}[Hasse--Weil 动机分解与良素处局部 \texorpdfstring{$L$}{L} 因子幂次刚性]\label{con:xi-terminal-zm-s4-hasse-weil-full-factorization-rigid-exponents}
在 \(\QQ\)-同源意义下有
\[
\mathrm{Jac}(X)\sim_{\QQ}\mathrm{Jac}(X_A)\times E^3\times E_{\mathrm D}^2\times B^3.
\]
因而对任一使上述各因子皆良约化的素数 \(p\)，局部 \(L\)-多项式分子满足
\[
P_p(X;T)
=P_p(X_A;T)\cdot P_p(E;T)^3\cdot P_p(E_{\mathrm D};T)^2\cdot P_p(B;T)^3,
\]
并且
\(
\deg P_p(B;T)=2\dim B=6
\)。
\end{conclusion}
\begin{proof}[证明]
第一式由结论
\ref{con:xi-terminal-zm-s4-c4-prym-threefold-and-v3sgn-identification}
第 \(4\) 点与既有 \(A_{\mathrm{sgn}}\)、\(A_{V_3}\)、\(A_{V_2}\) 识别合并得到。
第二式由“同源阿贝尔簇在良素处具有相同局部 \(L\)-因子”及乘积阿贝尔簇局部因子的乘法性直接推出。证毕。
\end{proof}

\begin{conclusion}[全纯 \(1\)-形式的 Chevalley--Weil 刚性分解]\label{con:xi-terminal-zm-s4-holomorphic-oneforms-chevalley-weil-rigid-decomposition}
记 \(S_4\) 的不可约表示为
\[
\mathbf 1,\ \mathrm{sgn},\ V_3,\ V_3\otimes\mathrm{sgn},\ V_2.
\]
则有
\[
H^0(X,\Omega_X^1)\cong
\mathrm{sgn}^{\oplus 2}\oplus V_3\oplus (V_3\otimes\mathrm{sgn})^{\oplus 3}\oplus V_2,
\]
并且平凡表示不出现。
\end{conclusion}
\begin{proof}[证明要点]
对 \(\rho\neq \mathbf 1\)，Chevalley--Weil 公式在本情形可写为
\[
m_\rho=-\dim\rho+\frac12\sum_{p\in B}\bigl(\dim\rho-\dim(\rho^{I_p})\bigr),
\]
且 \(m_{\mathbf 1}=0\)（底曲线 \(\PP^1\) 无全纯 \(1\)-形式）。
本覆盖仅需 \(C_2\) 与 \(C_4\) 两类惯性不变维数：
\[
\begin{aligned}
&\dim(\mathrm{sgn}^{C_2})=0,\ \dim(\mathrm{sgn}^{C_4})=0,\\
&\dim(V_3^{C_2})=2,\ \dim(V_3^{C_4})=0,\\
&\dim((V_3\!\otimes\!\mathrm{sgn})^{C_2})=1,\ \dim((V_3\!\otimes\!\mathrm{sgn})^{C_4})=1,\\
&\dim(V_2^{C_2})=1,\ \dim(V_2^{C_4})=1.
\end{aligned}
\]
代入五个 \(C_2\) 分歧点与一个 \(C_4\) 分歧点，得到
\[
m_{\mathrm{sgn}}=2,\quad m_{V_3}=1,\quad m_{V_3\otimes\mathrm{sgn}}=3,\quad m_{V_2}=1,\quad m_{\mathbf 1}=0.
\]
与结论 \ref{con:xi-terminal-zm-s4-chevalley-weil-explicit-decomposition} 一致。证毕。
\end{proof}

\begin{conclusion}[正规子群商曲线的显式模型与属数：\texorpdfstring{$A_4$}{A4} 商与 \texorpdfstring{$V_4$}{V4} 商]\label{con:xi-terminal-zm-s4-normal-subgroup-quotients-explicit-models-genus}
设
\[
X_A:=X/A_4,\qquad X_V:=X/V_4^{\mathrm{norm}}.
\]
则：
\begin{enumerate}
  \item \(X_A\to\PP^1_y\) 为二次覆盖，且
  \[
  \QQ(X_A)=\QQ\!\left(y,\sqrt{\Delta(y)}\right),
  \qquad
  X_A:\ t^2=\Delta(y),
  \qquad g(X_A)=2.
  \]
  \item \(X_V\to\PP^1_y\) 为次数 \(6\) 的 \(S_3\)-正则覆盖，且
  \[
  \QQ(X_V)=\QQ(y,z,\sqrt{\Delta(y)}),\qquad \mathscr R(z,y)=0,
  \]
  其中可取
  \[
  \mathscr R(z,y)=z^3+(2y+1)z^2-(4y^2+4y+1)z-(8y^3+13y^2+5y+1),
  \]
  并且 \(g(X_V)=4\)。
\end{enumerate}
\end{conclusion}
\begin{proof}[证明要点]
对 \(A_4\)：指数 \(2\) 正规子群对应符号同态核，故固定域由判别式平方根生成，即
\(
\QQ(X_A)=\QQ(y,\sqrt{\Delta(y)})
\)。
由于 \(\deg\Delta=5\) 为奇数，双覆盖在五个有限零点与无穷远点分歧，总分歧点数 \(6\)，故 \(g(X_A)=2\)。
\par
对 \(V_4^{\mathrm{norm}}\)：商群 \(S_4/V_4^{\mathrm{norm}}\cong S_3\)，固定域由三次解式闭包给出。
以 \(\mathscr R\) 为 resolvent cubic，附加其判别式平方根即得 \(S_3\)-闭包域，
从而 \(\QQ(X_V)=\QQ(y,z,\sqrt{\Delta(y)})\)。
由结论 \ref{con:xi-terminal-zm-s4-hurwitz-subgroup-spectrum-cycletypes} 已知 \(g(X_V)=4\)；等价地，用 \(|S_3|=6\) 与六个阶 \(2\) 分歧值代入 Hurwitz 公式亦得同值。证毕。
\end{proof}

\begin{conclusion}[双换位商曲线的双二次结构与纤维积刻画]\label{con:xi-terminal-zm-s4-delta-quotient-v4-fiberproduct}
沿用本小节记号，取无穷远惯性生成元 \(c\in S_4\) 为四循环，置
\[
\delta:=c^2,\qquad
Z:=X/\langle\delta\rangle,\qquad
E_{\mathrm D}:=X/D_8,\qquad
Y:=X/\langle c\rangle.
\]
则有：
\begin{enumerate}
  \item \(g(Z)=8\)；
  \item \(Z\to E_{\mathrm D}\) 为次数 \(4\) 的 Galois 覆盖，覆盖群同构于 \(V_4\)；
  \item 在函数域层面
  \[
  \QQ(Z)=\QQ(X_V)\cdot\QQ(Y)\quad\text{于}\ \QQ(E_{\mathrm D})\ \text{之上},
  \]
  且
  \[
  \QQ(X_V)\cap\QQ(Y)=\QQ(E_{\mathrm D}),
  \]
  从而 \(Z\) 同构于纤维积 \(X_V\times_{E_{\mathrm D}}Y\) 的规范化。
\end{enumerate}
\end{conclusion}
\begin{proof}
第 \(1\) 点由结论 \ref{con:xi-terminal-zm-s4-conjugacy-class-fixedpoints-and-cyclic-quotient-genera} 中
\(
\#\mathrm{Fix}_X(\delta)=2
\)
与 \(g(X)=16\) 代入二次商 Hurwitz 公式：
\[
2g(X)-2=2(2g(Z)-2)+\#\mathrm{Fix}_X(\delta)
\]
即得 \(g(Z)=8\)。

第 \(2\) 点中，\(\delta\in\langle c\rangle\subset D_8\)，且 \(V_4^{\mathrm{norm}}\subset D_8\)。
并有
\[
\langle\delta\rangle=V_4^{\mathrm{norm}}\cap\langle c\rangle,\qquad
D_8/\langle\delta\rangle\cong V_4.
\]
故 \(Z=X/\langle\delta\rangle\to X/D_8=E_{\mathrm D}\) 为次数 \(4\) 的 \(V_4\)-Galois 覆盖。

第 \(3\) 点令 \(L:=\QQ(X)\)，则
\[
\QQ(X_V)=L^{V_4^{\mathrm{norm}}},\qquad
\QQ(Y)=L^{\langle c\rangle},\qquad
\QQ(E_{\mathrm D})=L^{D_8}.
\]
由 Galois 对应性，
\[
L^{V_4^{\mathrm{norm}}}\cap L^{\langle c\rangle}
=L^{\langle V_4^{\mathrm{norm}},\langle c\rangle\rangle}
=L^{D_8}
=\QQ(E_{\mathrm D}),
\]
\[
L^{V_4^{\mathrm{norm}}}\cdot L^{\langle c\rangle}
=L^{V_4^{\mathrm{norm}}\cap\langle c\rangle}
=L^{\langle\delta\rangle}
=\QQ(Z).
\]
由“纤维积规范化的函数域等于子域合成”即得结论。证毕。
\end{proof}

\begin{conclusion}[三条二次子覆盖的分支奇偶对称差律]\label{cor:xi-terminal-zm-s4-v4-three-quad-subcovers-branch-parity}
在结论 \ref{con:xi-terminal-zm-s4-delta-quotient-v4-fiberproduct} 的 \(V_4\)-Galois 覆盖
\(
Z\to E_{\mathrm D}
\)
中，存在且仅存在三条非平凡二次子覆盖：
\[
X_V\to E_{\mathrm D},\qquad
Y\to E_{\mathrm D},\qquad
X_H\to E_{\mathrm D},
\]
其中可取 \(H=\langle\delta,\tau\rangle\le D_8\)（\(\tau\) 为 \(D_8\) 内换位），并置 \(X_H:=X/H\)。则：
\begin{enumerate}
  \item \(g(X_H)=2\)，且 \(X_H\to E_{\mathrm D}\) 为分支次数 \(2\) 的二次覆盖；
  \item 若
  \[
  \QQ(X_V)=\QQ(E_{\mathrm D})(\sqrt{f_V}),\qquad
  \QQ(Y)=\QQ(E_{\mathrm D})(\sqrt{f_Y}),
  \]
  则
  \[
  \QQ(X_H)=\QQ(E_{\mathrm D})(\sqrt{f_Vf_Y}),
  \]
  并满足
  \[
  \mathrm{Branch}(X_H/E_{\mathrm D})
  \equiv
  \mathrm{Branch}(X_V/E_{\mathrm D})\triangle \mathrm{Branch}(Y/E_{\mathrm D})
  \pmod 2.
  \]
  \item 设 \(D_V,D_Y,D_H\) 为三条二次覆盖的约化分支集合。
  已知 \(\deg D_V=\deg D_Y=6,\ \deg D_H=2\)，故
  \[
  \deg(D_V\cap D_Y)=5,
  \]
  即 \(D_V,D_Y\) 共有 \(5\) 个分支点，且各自仅有 \(1\) 个独有分支点；这两个独有点恰构成 \(D_H\)。
\end{enumerate}
\end{conclusion}
\begin{proof}
对第 \(1\) 点，取 \(H=\langle\delta,\tau\rangle=\{1,\delta,\tau,\delta\tau\}\)。
由结论 \ref{con:xi-terminal-zm-s4-conjugacy-class-fixedpoints-and-cyclic-quotient-genera}，
\[
\#\mathrm{Fix}_X(\delta)=2,\qquad
\#\mathrm{Fix}_X(\tau)=\#\mathrm{Fix}_X(\delta\tau)=10.
\]
群作用 Hurwitz 公式给出
\[
2g(X)-2=|H|\bigl(2g(X_H)-2\bigr)+\sum_{1\neq h\in H}\#\mathrm{Fix}_X(h),
\]
即
\[
30=4\bigl(2g(X_H)-2\bigr)+(2+10+10),
\]
故 \(g(X_H)=2\)。
又 \([D_8:H]=2\)，故 \(X_H\to E_{\mathrm D}\) 为二次覆盖；由 \(g(E_{\mathrm D})=1\) 得分支次数 \(2g(X_H)-2=2\)。

第 \(2\) 点是双二次扩张的局部赋值判别：在任意离散赋值 \(v\) 下，\(\QQ(E_{\mathrm D})(\sqrt{f})/\QQ(E_{\mathrm D})\) 分歧当且仅当 \(v(f)\) 为奇数。对 \(f_V,f_Y,f_Vf_Y\) 逐点比较即得模 \(2\) 对称差律。

第 \(3\) 点由第 \(1\) 点与已有 \(g(X_V)=g(Y)=4\) 立即得到 \(\deg D_V=\deg D_Y=6,\deg D_H=2\)。
在特征 \(0\) 下二次分支为简单分支，故
\[
2=\deg(D_V\triangle D_Y)=\deg D_V+\deg D_Y-2\deg(D_V\cap D_Y)=12-2\deg(D_V\cap D_Y),
\]
从而 \(\deg(D_V\cap D_Y)=5\)。证毕。
\end{proof}

\begin{conclusion}[\(\delta\)-不变/反不变分解与 \(\mathrm{Jac}(Z)\)、\(\mathrm{Prym}(X/Z)\) 同源型]\label{con:xi-terminal-zm-s4-delta-eigensplitting-jacz-prymxz}
在结论 \ref{con:xi-terminal-zm-s4-delta-quotient-v4-fiberproduct} 的记号下，存在定义于 \(\QQ\) 上的同源分解
\[
\mathrm{Jac}(Z)\sim_{\QQ}\mathrm{Jac}(X_A)\times E\times E_{\mathrm D}^2\times B,
\]
以及
\[
\mathrm{Prym}(X/Z)\sim_{\QQ}E^2\times B^2.
\]
\end{conclusion}
\begin{proof}
记 \(p:X\to Z=X/\langle\delta\rangle\)，并在 \(\QQ[\langle\delta\rangle]\) 中取幂等元
\[
e_+:=\frac{1+\delta}{2},\qquad e_-:=\frac{1-\delta}{2}.
\]
标准理论给出
\[
\mathrm{Jac}(X)\sim_{\QQ}e_+\mathrm{Jac}(X)\times e_-\mathrm{Jac}(X),
\]
且
\(
e_+\mathrm{Jac}(X)\sim_{\QQ}\mathrm{Jac}(Z)
\),
\(
e_-\mathrm{Jac}(X)\sim_{\QQ}\mathrm{Prym}(X/Z)
\)。

再由结论 \ref{con:xi-terminal-zm-s4-jacobian-group-algebra-prym-e3-b3}
\[
\mathrm{Jac}(X)\sim_{\QQ}\mathrm{Jac}(X_A)\times E^3\times E_{\mathrm D}^2\times B^3.
\]
对双换位 \(\delta\)，其在 \(\mathrm{sgn},V_2\) 上取值为 \(+1\)，在 \(V_3,V_3\otimes\mathrm{sgn}\) 上迹为 \(-1\)，故 \(+1\)-特征维数分别为
\[
1,\ 2,\ 1,\ 1.
\]
于是 \(e_+\)-像分别贡献
\[
\mathrm{Jac}(X_A),\ E_{\mathrm D}^2,\ E,\ B,
\]
而 \(e_-\)-像分别贡献
\[
0,\ 0,\ E^2,\ B^2.
\]
合并即得结论。证毕。
\end{proof}

\begin{conclusion}[属 \(2\) 中间商曲线的 Jacobian 分裂]\label{con:xi-terminal-zm-s4-xh-jacobian-splitting-ed-e}
取
\[
H=\langle\delta,\tau\rangle\le D_8,\qquad X_H:=X/H
\]
（\(\tau\) 为 \(D_8\) 中换位）。则 \(g(X_H)=2\)，且
\[
\mathrm{Jac}(X_H)\sim_{\QQ}E\times E_{\mathrm D}.
\]
特别地，\(X_H\to E_{\mathrm D}\) 为分支次数 \(2\) 的二次覆盖，且 \(X_H\) 为双椭圆型属 \(2\) 曲线。
\end{conclusion}
\begin{proof}
由结论 \ref{cor:xi-terminal-zm-s4-v4-three-quad-subcovers-branch-parity} 已知 \(g(X_H)=2\)。
令
\[
e_H:=\frac{1}{|H|}\sum_{h\in H}h\in\QQ[H].
\]
则 \(e_H\mathrm{Jac}(X)\sim_{\QQ}\mathrm{Jac}(X_H)\)。
对任一不可约表示 \(\rho\)，有
\[
\dim(\rho^H)=\frac14\!\left(\chi_\rho(1)+\chi_\rho(\delta)+2\chi_\rho(\text{换位})\right).
\]
代入 \(S_4\) 特征标值得
\[
\dim(\mathrm{sgn}^H)=0,\quad
\dim((V_3\otimes\mathrm{sgn})^H)=0,\quad
\dim(V_3^H)=1,\quad
\dim(V_2^H)=1.
\]
故 \(\mathrm{Jac}(X_A)\) 与 \(B^3\) 不贡献不变部分，而 \(E^3\) 与 \(E_{\mathrm D}^2\) 各贡献一维椭圆因子，遂有
\[
\mathrm{Jac}(X_H)\sim_{\QQ}E\times E_{\mathrm D}.
\]
其余断言由 \(g(X_H)=2\) 与该分裂立即得到。证毕。
\end{proof}

\begin{conclusion}[三次解式判别式与 \texorpdfstring{$\Delta(y)$}{Delta(y)} 的一致性]\label{con:xi-terminal-zm-s4-resolvent-cubic-discriminant-equals-delta}
在结论 \ref{con:xi-terminal-zm-s4-normal-subgroup-quotients-explicit-models-genus} 的记号下，恒有
\[
\mathrm{Disc}_z\!\bigl(\mathscr R(z,y)\bigr)=\Delta(y)=-y(y-1)g(y).
\]
\end{conclusion}
\begin{proof}[证明要点]
将 \(\mathscr R=z^3+az^2+bz+c\) 代入三次判别式闭式
\[
\mathrm{Disc}_z(\mathscr R)=a^2b^2-4b^3-4a^3c-27c^2+18abc,
\]
其中
\[
a=2y+1,\quad b=-(4y^2+4y+1),\quad c=-(8y^3+13y^2+5y+1),
\]
直接化简即得 \(-y(y-1)g(y)\)。
该恒等式亦与结论 \ref{con:xi-terminal-zm-discriminant-leyang}、\ref{con:xi-terminal-zm-s4-c3-fiberproduct-normalization} 的二次分辨数据一致。证毕。
\end{proof}

\begin{conclusion}[中间商曲线属数谱与 Jacobian 的表示论分裂]\label{con:xi-terminal-zm-s4-quotient-genus-spectrum-and-jacobian-representation-splitting}
对任意子群 \(H\le S_4\)，有
\[
g(X/H)=\dim_{\CC}H^0(X,\Omega_X^1)^H.
\]
由结论 \ref{con:xi-terminal-zm-s4-holomorphic-oneforms-chevalley-weil-rigid-decomposition} 可直接得到
\[
g(X/A_4)=2,\quad
g(X/V_4^{\mathrm{norm}})=4,\quad
g(X/S_3)=1,\quad
g(X/C_3)=6,\quad
g(X/D_8)=1.
\]
此外，存在定义在 \(\QQ\) 上的阿贝尔簇
\[
A_{\mathrm{sgn}},\ A_{V_3},\ A_{V_3\otimes\mathrm{sgn}},\ A_{V_2},
\]
使
\[
\mathrm{Jac}(X)\sim_{\QQ}
A_{\mathrm{sgn}}\times A_{V_3}\times A_{V_3\otimes\mathrm{sgn}}\times A_{V_2},
\]
并且
\[
\dim A_{\mathrm{sgn}}=2,\quad
\dim A_{V_3}=3,\quad
\dim A_{V_3\otimes\mathrm{sgn}}=9,\quad
\dim A_{V_2}=2.
\]
同时有规范同源识别
\[
A_{\mathrm{sgn}}\sim_{\QQ}\mathrm{Jac}(X/A_4),\qquad
A_{\mathrm{sgn}}\times A_{V_2}\sim_{\QQ}\mathrm{Jac}(X/V_4^{\mathrm{norm}}).
\]
\end{conclusion}
\begin{proof}[证明要点]
恒等式 \(H^0(X/H,\Omega^1)\cong H^0(X,\Omega^1)^H\) 给出第一式。
将
\[
H^0(X,\Omega_X^1)\cong
2\,\mathrm{sgn}\oplus V_3\oplus 3(V_3\otimes\mathrm{sgn})\oplus V_2
\]
限制到各子群并取不变维数，即得所列属数（与结论
\ref{con:xi-terminal-zm-s4-hurwitz-subgroup-spectrum-cycletypes} 一致）。
\par
对 Jacobian，按 \(\QQ[S_4]\) 中心幂等元作用得到 isotypic 分裂；由于 \(S_4\) 不可约特征均为有理值，各因子可在 \(\QQ\) 上定义。
维数由重数乘表示维数给出：
\[
2\cdot 1,\ 1\cdot 3,\ 3\cdot 3,\ 1\cdot 2,
\]
即 \((2,3,9,2)\)。
当 \(H\) 正规（如 \(A_4,V_4^{\mathrm{norm}}\)）时，\(\mathrm{Jac}(X/H)\) 正对应于在 \(H\) 上平凡的表示分量之直和，故得两条规范同源识别。证毕。
\end{proof}
\paragraph{\texorpdfstring{$S_4$}{S4}-等变 \texorpdfstring{$H^1$}{H1} 分解、驯导子谱与函数域 Langlands 分量化}
沿用子小节 \ref{subsubsec:xi-s4-regular-closure-global-spectrum} 的记号，令
\(
X=\widetilde X
\)
为 Lee--Yang 四次谱覆盖的 \(S_4\)-正则闭包曲线，分歧惯性型为 \((4,2,2,2,2,2)\)。
记 \(S_4\) 的四个非平凡不可约表示为
\[
\varepsilon=\mathrm{sgn},\qquad V_2,\qquad V_3,\qquad V_3':=V_3\otimes\varepsilon.
\]

\begin{conclusion}[等变 Weil--Artin 分量分解的次数与幂指数刚性]\label{con:xi-terminal-zm-s4-equivariant-weil-artin-isotypic-factorization}
设 \(p\nmid 2\cdot 3\cdot 31\cdot 37\)，并设 \(X\) 在 \(p\) 处具有光滑良约化 \(X_p/\FF_p\)。
取素数 \(\ell\neq p\)，并记几何 Frobenius 为 \(\mathrm{Frob}_p\)。
则 \(H^1_{\acute et}(X_{\overline{\FF}_p},\QQ_\ell)\) 作为 \(\QQ_\ell[S_4]\)-半单模存在唯一的等变直和分解
\[
H^1_{\acute et}(X_{\overline{\FF}_p},\QQ_\ell)
\ \cong\
(\varepsilon\otimes M_{\varepsilon})
\ \oplus\
(V_2\otimes M_{2})
\ \oplus\
(V_3\otimes M_{3})
\ \oplus\
(V_3'\otimes M_{3'}),
\]
其中
\(
M_\sigma:=\Hom_{S_4}\!\bigl(\sigma,H^1_{\acute et}(X_{\overline{\FF}_p},\QQ_\ell)\bigr)
\)
为 \(\mathrm{Gal}(\overline{\FF}_p/\FF_p)\)-表示，并满足
\[
\dim M_{\varepsilon}=4,\qquad
\dim M_{2}=2,\qquad
\dim M_{3}=2,\qquad
\dim M_{3'}=6.
\]
令
\[
P_\sigma(T):=\det\!\left(1-T\mathrm{Frob}_p\mid M_\sigma\right)\in 1+T\overline{\QQ}[T]
\qquad(\sigma\in\{\varepsilon,2,3,3'\}),
\]
并记曲线的局部因子分子
\(
P_p(X;T):=\det(1-T\mathrm{Frob}_p\mid H^1_{\acute et}(X_{\overline{\FF}_p},\QQ_\ell))
\)。
则有幂指数分解
\[
P_p(X;T)=P_{\varepsilon}(T)^{1}\cdot P_{2}(T)^{2}\cdot P_{3}(T)^{3}\cdot P_{3'}(T)^{3},
\]
并且
\[
\deg P_{\varepsilon}=4,\qquad \deg P_{2}=2,\qquad \deg P_{3}=2,\qquad \deg P_{3'}=6.
\]
因此
\[
Z(X_p,T)=\frac{P_{\varepsilon}(T)\,P_{2}(T)^{2}\,P_{3}(T)^{3}\,P_{3'}(T)^{3}}{(1-T)(1-pT)}.
\]
\end{conclusion}
\begin{proof}
对任意 \(\ell\neq p\)，有自然同构
\(
H^1_{\acute et}(X_{\overline{\FF}_p},\QQ_\ell)\cong V_\ell(\mathrm{Jac}(X))
\)
且该同构与 \(S_4\) 作用及 \(\mathrm{Frob}_p\) 作用相容。
由结论 \ref{con:xi-terminal-zm-s4-quotient-genus-spectrum-and-jacobian-representation-splitting}，
\(\mathrm{Jac}(X)\) 在 \(\QQ\)-同源意义下按 \(\varepsilon,V_2,V_3,V_3'\) 的等型分量分裂，并满足
\[
\dim A_{\varepsilon}=2,\qquad
\dim A_{V_2}=2,\qquad
\dim A_{V_3}=3,\qquad
\dim A_{V_3'}=9.
\]
因此 \(H^1_{\acute et}\) 的相应等型分量维数分别为 \(4,4,6,18\)，再除以 \(\dim\sigma\in\{1,2,3,3\}\) 得
\[
\dim M_{\varepsilon}=4,\qquad
\dim M_{2}=2,\qquad
\dim M_{3}=2,\qquad
\dim M_{3'}=6,
\]
从而得到所示等变分解与次数数据。
再由 Schur 引理，
\(
\det(1-T\mathrm{Frob}_p\mid \sigma\otimes M_\sigma)=\det(1-T\mathrm{Frob}_p\mid M_\sigma)^{\dim\sigma}
\)，
故 \(P_p(X;T)=\prod_\sigma P_\sigma(T)^{\dim\sigma}\)。
次数陈述由 \(\deg P_\sigma=\dim M_\sigma\) 得出；最后一式为曲线的标准 ζ-函数表达式。证毕。
\end{proof}

\begin{conclusion}[分歧集上的驯 Artin 导子谱与全局导子度]\label{con:xi-terminal-zm-s4-tame-artin-conductor-spectrum}
设
\[
U:=\PP^1_y\setminus\Bigl(\{0,1,\infty\}\cup\{P_{\mathrm{LY}}(y)=0\}\Bigr),
\]
并令 \(\pi_1(U)\twoheadrightarrow S_4\) 为由正则闭包 \(X\to\PP^1_y\) 给出的单值化满射。
对每个不可约表示 \(\sigma\in\{\varepsilon,V_2,V_3,V_3'\}\)，记相应驯 lisse 局部系统为 \(\mathcal L_\sigma\)。
则对每个分歧点 \(b\)（其惰性群 \(I_b\) 为 \(C_2\) 或 \(C_4\)）有驯 Artin 导子指数
\[
a_b(\sigma)=\dim\sigma-\dim\sigma^{I_b},
\]
并且在五个有限分歧点（惰性生成元可取换位 \(\tau\)）处有
\[
a_b(\varepsilon)=1,\qquad a_b(V_2)=1,\qquad a_b(V_3)=1,\qquad a_b(V_3')=2,
\]
在无穷远分歧点（惰性生成元可取四循环 \(\rho\)）处有
\[
a_\infty(\varepsilon)=1,\qquad a_\infty(V_2)=1,\qquad a_\infty(V_3)=3,\qquad a_\infty(V_3')=2.
\]
因而四个表示的全局驯导子度
\(
A(\sigma):=\sum_b a_b(\sigma)
\)
分别为
\[
A(\varepsilon)=6,\qquad A(V_2)=6,\qquad A(V_3)=8,\qquad A(V_3')=12.
\]
\end{conclusion}
\begin{proof}
有限分歧点处 \(I_b\cong\langle\tau\rangle\)（\(\tau\) 为换位），无穷远处 \(I_\infty\cong\langle\rho\rangle\)（\(\rho\) 为四循环）。
对任意循环群 \(\langle g\rangle\)（阶 \(n\)）有角色平均
\[
\dim\sigma^{\langle g\rangle}=\frac1n\sum_{k=0}^{n-1}\chi_\sigma(g^k).
\]
取 \(S_4\) 角色值（按共轭类 \(1,(12),(12)(34),(123),(1234)\)）：
\[
\chi_{\varepsilon}=(1,-1,1,1,-1),\quad
\chi_{V_2}=(2,0,2,-1,0),\quad
\chi_{V_3}=(3,1,-1,0,-1),\quad
\chi_{V_3'}=(3,-1,-1,0,1).
\]
于是
\(
\dim\sigma^{\langle\tau\rangle}=\tfrac12(\chi_\sigma(1)+\chi_\sigma(\tau))
\)
给出
\[
\dim\varepsilon^{\langle\tau\rangle}=0,\quad
\dim V_2^{\langle\tau\rangle}=1,\quad
\dim V_3^{\langle\tau\rangle}=2,\quad
\dim V_3'^{\langle\tau\rangle}=1,
\]
从而五个有限分歧点的导子指数如表。
另一方面 \(\rho\) 为四循环且 \(\rho^2\) 为双换位，故
\(
\dim\sigma^{\langle\rho\rangle}=\tfrac14(\chi_\sigma(1)+2\chi_\sigma(\rho)+\chi_\sigma(\rho^2))
\)
给出
\[
\dim\varepsilon^{\langle\rho\rangle}=0,\quad
\dim V_2^{\langle\rho\rangle}=1,\quad
\dim V_3^{\langle\rho\rangle}=0,\quad
\dim V_3'^{\langle\rho\rangle}=1,
\]
从而无穷远处导子指数如表。
对六个分歧点求和即得全局导子度。证毕。
\end{proof}

\begin{conclusion}[函数域 Langlands 的分量化实现与导子精确性]\label{con:xi-terminal-zm-s4-function-field-langlands-artin-components}
设 \(p\nmid 2\cdot 3\cdot 31\cdot 37\)，并令 \(F=\FF_p(y)\)。
对每个不可约表示 \(\sigma\in\{\varepsilon,V_2,V_3,V_3'\}\)，考虑由正则闭包单值化得到的纯权 \(0\) 驯 lisse 局部系统 \(\mathcal L_\sigma\)（定义于 \(U\subset\PP^1_{\FF_p}\)）。
则存在一个（在同构意义下）唯一的尖点自守表示
\[
\pi_{\sigma,p}\ \subset\ \mathrm{GL}_{\dim\sigma}(\mathbb A_F),
\]
使得对每个非分歧处 \(v\) 的未分歧局部 \(L\)-因子有
\[
L_v(\pi_{\sigma,p},T)
=
\det\!\left(1-T\mathrm{Frob}_v\mid (\mathcal L_{\sigma,\bar\eta})^{I_v}\right)^{-1}.
\]
并且 \(\pi_{\sigma,p}\) 的算术导子在每个分歧点处的局部指数等于结论 \ref{con:xi-terminal-zm-s4-tame-artin-conductor-spectrum} 所给的 \(a_b(\sigma)\)，从而
\[
\deg\mathfrak f(\pi_{\varepsilon,p})=6,\qquad
\deg\mathfrak f(\pi_{V_2,p})=6,\qquad
\deg\mathfrak f(\pi_{V_3,p})=8,\qquad
\deg\mathfrak f(\pi_{V_3',p})=12.
\]
\end{conclusion}
\begin{proof}[证明要点]
由于 \(\mathcal L_\sigma\) 具有有限幺半单单值群，且仅在分歧集处驯分歧，函数域全球 Langlands 对应（Lafforgue）将其与唯一尖点自守表示 \(\pi_{\sigma,p}\) 配对，并逐点匹配未分歧局部 \(L\)-因子并保持导子指数。
将结论 \ref{con:xi-terminal-zm-s4-tame-artin-conductor-spectrum} 的驯导子谱代入即可得到导子度公式。证毕。\cite{LafforgueGlobalLanglandsFunctionFields}
\end{proof}

\begin{conclusion}[分量化 Weil 圆周性质与正号函数方程]\label{con:xi-terminal-zm-s4-componentwise-rh-functional-equation-plus}
在结论 \ref{con:xi-terminal-zm-s4-equivariant-weil-artin-isotypic-factorization} 的记号下，对每个 \(\sigma\in\{\varepsilon,2,3,3'\}\)：
\begin{enumerate}
  \item \(P_\sigma(T)\) 的全部逆根 \(\alpha\) 满足 \(|\alpha|=\sqrt p\)（任一复嵌入下），等价地其全部零点位于圆周 \(|T|=p^{-1/2}\)。
  \item 记 \(m_\varepsilon=2,m_2=1,m_3=1,m_{3'}=3\) 为 \(H^0(X,\Omega^1)\) 中的重数（结论 \ref{con:xi-terminal-zm-s4-holomorphic-oneforms-chevalley-weil-rigid-decomposition}）。
  则 \(\deg P_\sigma=2m_\sigma\)，并且函数方程具有正号形式
  \[
  P_\sigma(T)=p^{m_\sigma}\,T^{2m_\sigma}\,P_\sigma\!\left(\frac{1}{pT}\right).
  \]
\end{enumerate}
\end{conclusion}
\begin{proof}[证明要点]
纯性给出圆周性质：\(H^1_{\acute et}(X_{\overline{\FF}_p},\QQ_\ell)\) 为权 \(1\) 的纯表示，故其 Frobenius 特征值绝对值为 \(\sqrt p\)，并对任一 Frobenius 稳定子表示（尤其各 \(S_4\)-等型分量 \(\sigma\otimes M_\sigma\)）保持。
\par
函数方程来自 Poincar\'e 对偶：\(H^1\) 上存在 \(\mathrm{Frob}_p\)-半线性扭转的非退化交替配对，且 \(S_4\) 作用保持该配对。
由于四个 \(\sigma\) 皆自对偶且可取 \(S_4\)-不变对称双线性形式，配对在 \(\sigma\otimes M_\sigma\) 上诱导出 \(M_\sigma\) 的交替配对；于是 \(\det(1-T\mathrm{Frob}_p\mid M_\sigma)\) 为回文多项式且符号为 \(+1\)，指数由 \(\deg P_\sigma=\dim M_\sigma=2m_\sigma\) 给出。证毕。\cite{DeligneWeilII}
\end{proof}

\endinput


\paragraph{\texorpdfstring{$V_4$}{V4} 中间塔的分支刚化、幂等分解与 Prym 位置唯一化}

\begin{conclusion}[双二次覆盖 \texorpdfstring{$Z\to E_{\mathrm D}$}{Z to ED} 的分支总数与惯性分布]\label{con:xi-terminal-zm-s4-v4-ed-branchcount-inertia-115}
在结论 \ref{con:xi-terminal-zm-s4-delta-quotient-v4-fiberproduct} 与
\ref{cor:xi-terminal-zm-s4-v4-three-quad-subcovers-branch-parity} 的记号下，记
\[
\pi:Z\longrightarrow E_{\mathrm D},\qquad \operatorname{Gal}(Z/E_{\mathrm D})\cong V_4.
\]
则成立：
\begin{enumerate}
  \item \(\pi\) 在 \(E_{\mathrm D}\) 上恰有 \(7\) 个分支点，且均为惯性阶 \(2\) 的简单分支；
  \item 存在 \(V_4\setminus\{1\}=\{a,b,\gamma\}\) 的标号，使分支集合分解为互不相交约化有效除子
  \[
  D_a\sqcup D_b\sqcup D_\gamma,\qquad
  \deg D_a=\deg D_b=1,\ \deg D_\gamma=5,
  \]
  并满足
  \[
  \operatorname{Branch}(X_V/E_{\mathrm D})=D_b+D_\gamma,\qquad
  \operatorname{Branch}(Y/E_{\mathrm D})=D_a+D_\gamma,\qquad
  \operatorname{Branch}(X_H/E_{\mathrm D})=D_a+D_b.
  \]
\end{enumerate}
\end{conclusion}
\begin{proof}
由 \(\operatorname{Gal}(Z/E_{\mathrm D})\cong V_4\) 可知每个分支点惯性群均为阶 \(2\)。
记 \(N:=\deg\operatorname{Branch}(\pi)\)。因 \(g(E_{\mathrm D})=1,\ g(Z)=8\)，Hurwitz 公式给出
\[
2g(Z)-2
=\sum_{Q\in\operatorname{Branch}(\pi)}\frac{|V_4|}{|I_Q|}\bigl(|I_Q|-1\bigr)
=\sum_{Q\in\operatorname{Branch}(\pi)}2
=2N,
\]
故 \(14=2N\)，即 \(N=7\)。

设三条二次子覆盖对应非平凡特征
\(
\chi_V,\chi_Y,\chi_H\in V_4^\vee
\)。
对分支点 \(Q\)，记其惯性生成元为 \(s_Q\in V_4\setminus\{1\}\)。
二次子覆盖在 \(Q\) 处分歧当且仅当对应特征在 \(s_Q\) 上取值 \(-1\)。
对任意 \(s\neq 1\)，三条非平凡特征中恰有一条在 \(s\) 上取值 \(+1\)，其余两条取值 \(-1\)；
因此每个分支点处，三条二次子覆盖恰有一条不分歧。

记
\[
n_V:=\#\{Q\in\operatorname{Branch}(\pi):X_V\ \text{在}\ Q\ \text{处不分歧}\},
\]
\[
n_Y:=\#\{Q\in\operatorname{Branch}(\pi):Y\ \text{在}\ Q\ \text{处不分歧}\},\qquad
n_H:=\#\{Q\in\operatorname{Branch}(\pi):X_H\ \text{在}\ Q\ \text{处不分歧}\}.
\]
由结论 \ref{cor:xi-terminal-zm-s4-v4-three-quad-subcovers-branch-parity}，
\[
\deg\operatorname{Branch}(X_V/E_{\mathrm D})=6,\quad
\deg\operatorname{Branch}(Y/E_{\mathrm D})=6,\quad
\deg\operatorname{Branch}(X_H/E_{\mathrm D})=2.
\]
又 \(\#\operatorname{Branch}(\pi)=7\)，故
\[
7-n_V=6,\qquad 7-n_Y=6,\qquad 7-n_H=2,
\]
从而 \((n_V,n_Y,n_H)=(1,1,5)\)。
令 \(D_a,D_b,D_\gamma\) 分别为使 \(X_V\)、\(Y\)、\(X_H\) 不分歧的分支点集合，即得
\[
\deg D_a=1,\ \deg D_b=1,\ \deg D_\gamma=5
\]
及三条分支公式。证毕。
\end{proof}

\begin{conclusion}[第七分支点位于 \texorpdfstring{$E_{\mathrm D}\to\PP^1_y$}{ED to P1} 的分歧原像]\label{con:xi-terminal-zm-s4-v4-ed-seventh-branchpoint-cubic-ramified-preimage}
设 \(\varpi:E_{\mathrm D}\to\PP^1_y\) 为结论
\ref{thm:xi-s3-fiberproduct-branch-parity-xa-xv-ed} 中的三次覆盖，
其分歧值集合为
\[
\{0,1,\infty\}\cup\{g(y)=0\}.
\]
则：
\begin{enumerate}
  \item \(X_V\to E_{\mathrm D}\) 的 \(6\) 个分支点正是上述六个分歧值在 \(E_{\mathrm D}\) 上的未分歧原像点；
  \item 在结论 \ref{con:xi-terminal-zm-s4-v4-ed-branchcount-inertia-115} 中使 \(X_V\) 不分歧而 \(Z\) 分歧的唯一点，必为某个分歧值的分歧原像点。
\end{enumerate}
\end{conclusion}
\begin{proof}
第 \(1\) 点即结论 \ref{thm:xi-s3-fiberproduct-branch-parity-xa-xv-ed} 的第 \(3\) 点。

对第 \(2\) 点，记 \(S\subset\PP^1_y\) 为上述六个分歧值。
若 \(y_0\notin S\)，则 \(X\to\PP^1_y\) 在 \(y_0\) 上无分歧，故其上方点稳定子平凡。
任何中间商覆盖（特别是 \(Z\to E_{\mathrm D}\)）在该纤维上亦不可能新生分歧。
因此 \(\operatorname{Branch}(Z/E_{\mathrm D})\subset\varpi^{-1}(S)\)。

由结论 \ref{thm:xi-s3-fiberproduct-branch-parity-xa-xv-ed}，对每个 \(y_0\in S\)，纤维型为 \((2,1)\)：
记分歧原像为 \(R_{y_0}\)，未分歧原像为 \(U_{y_0}\)，且
\(
X_V\to E_{\mathrm D}
\)
恰在全部 \(U_{y_0}\) 处分歧。
由结论 \ref{con:xi-terminal-zm-s4-v4-ed-branchcount-inertia-115}，存在唯一分支点 \(P\) 使 \(X_V\) 在 \(P\) 处不分歧而 \(Z\) 分歧。
该点不能属于任何 \(U_{y_0}\)，故必有 \(P=R_{y_0}\) 对某个 \(y_0\in S\)。证毕。
\end{proof}

\begin{conclusion}[\texorpdfstring{$V_4$}{V4} 幂等关系导出的 Kani--Rosen 型分解]\label{con:xi-terminal-zm-s4-v4-ed-kanirosen-jacz-decomposition}
令 \(G:=\operatorname{Gal}(Z/E_{\mathrm D})\cong V_4\)，
并令 \(H_1,H_2,H_3\le G\) 为三条阶 \(2\) 子群，使
\[
Z/H_1\cong X_V,\qquad Z/H_2\cong Y,\qquad Z/H_3\cong X_H,\qquad Z/G\cong E_{\mathrm D}.
\]
则在 \(\QQ\)-同源意义下有
\[
\Jac(Z)\times \Jac(E_{\mathrm D})^{2}
\sim_{\QQ}
\Jac(X_V)\times \Jac(Y)\times \Jac(X_H).
\]
\end{conclusion}
\begin{proof}
在 \(\QQ[G]\) 中记
\[
e_H:=\frac{1}{|H|}\sum_{h\in H}h.
\]
对 \(G=V_4\) 的三条阶 \(2\) 子群直接计算得
\[
e_{H_1}+e_{H_2}+e_{H_3}=e_{\{1\}}+2e_G.
\]
将此恒等式作用于 \(\Jac(Z)\)。
按标准范数--拉回理论，
\(
\operatorname{Im}(e_H)\sim_{\QQ}\Jac(Z/H)
\) 且
\(
\operatorname{Im}(e_G)\sim_{\QQ}\Jac(Z/G)=\Jac(E_{\mathrm D})
\)。
于是得到
\[
\operatorname{Im}(e_{H_1})\times \operatorname{Im}(e_{H_2})\times \operatorname{Im}(e_{H_3})
\sim_{\QQ}
\Jac(Z)\times \Jac(E_{\mathrm D})^2,
\]
即所述分解。证毕。
\end{proof}

\begin{conclusion}[阿贝尔三维因子 \texorpdfstring{$B$}{B} 的局部 \texorpdfstring{$L$}{L} 因子有理消去公式]\label{con:xi-terminal-zm-s4-v4-ed-b-localfactor-rational-cancellation}
仍记
\(
B:=\Prym(Y/E_{\mathrm D})
\)。
若素数 \(p\) 处 \(Z,X_V,X_H,E_{\mathrm D}\) 皆良约化，记曲线 \(C\) 的局部 \(L\)-多项式分子为 \(P_p(C;T)\)，则
\[
P_p(B;T)
=
\frac{P_p(Z;T)\,P_p(E_{\mathrm D};T)}
{P_p(X_V;T)\,P_p(X_H;T)},
\qquad
\deg P_p(B;T)=6.
\]
等价地，在 Grothendieck 群中
\[
[H^1(B)]
=[H^1(Z)]+[H^1(E_{\mathrm D})]-[H^1(X_V)]-[H^1(X_H)],
\]
故对任意 \(n\ge 1\) 有
\[
\Tr\!\bigl(\mathrm{Frob}_p^n\mid H^1(B)\bigr)
=
\Tr\!\bigl(\mathrm{Frob}_p^n\mid H^1(Z)\bigr)
+\Tr\!\bigl(\mathrm{Frob}_p^n\mid H^1(E_{\mathrm D})\bigr)
-\Tr\!\bigl(\mathrm{Frob}_p^n\mid H^1(X_V)\bigr)
-\Tr\!\bigl(\mathrm{Frob}_p^n\mid H^1(X_H)\bigr).
\]
\end{conclusion}
\begin{proof}
由结论 \ref{con:xi-terminal-zm-s4-v4-ed-kanirosen-jacz-decomposition} 与
\ref{con:xi-terminal-zm-s4-c4-prym-threefold-and-v3sgn-identification}
中的
\(
\Jac(Y)\sim_{\QQ}E_{\mathrm D}\times B
\)，得到
\[
\Jac(Z)\times E_{\mathrm D}
\sim_{\QQ}
\Jac(X_V)\times \Jac(X_H)\times B.
\]
良素处同源阿贝尔簇具有相同 Frobenius 特征多项式，乘积对应局部因子相乘，故得比值公式。
再由 \(\dim B=3\)（同一结论已给出）得 \(\deg P_p(B;T)=2\dim B=6\)。
上式对 \(H^1\) 的 Grothendieck 群表达与迹恒等式与之等价。证毕。
\end{proof}

\begin{conclusion}[\texorpdfstring{$V_4$}{V4} 覆盖构造数据在椭圆基曲线上的刚化]\label{con:xi-terminal-zm-s4-v4-ed-abelian-cover-linebundle-rigidification}
设 \(D_a,D_b,D_\gamma\) 如结论
\ref{con:xi-terminal-zm-s4-v4-ed-branchcount-inertia-115}。
则存在线丛 \(L_a,L_b,L_\gamma\in\Pic(E_{\mathrm D})\)，满足
\[
L_a^{\otimes 2}\cong \mathcal O_{E_{\mathrm D}}(D_b+D_\gamma),\qquad
L_b^{\otimes 2}\cong \mathcal O_{E_{\mathrm D}}(D_a+D_\gamma),\qquad
L_\gamma^{\otimes 2}\cong \mathcal O_{E_{\mathrm D}}(D_a+D_b),
\]
且三条二次子覆盖分别由上述三组平方根数据给出。
特别地，
\[
\deg L_a=\deg L_b=3,\qquad \deg L_\gamma=1.
\]
此外 \(L_\gamma\in\Pic^1(E_{\mathrm D})\cong E_{\mathrm D}\) 等价于一点 \(Q\in E_{\mathrm D}\) 使
\[
2Q\sim D_a+D_b.
\]
\end{conclusion}
\begin{proof}
对光滑曲线上的 \(V_4\)-阿贝尔覆盖，标准构造数据为：分支除子
\(
\{D_s\}_{s\in V_4\setminus\{1\}}
\)
与角色线丛
\(
\{L_\chi\}_{\chi\in V_4^\vee\setminus\{1\}}
\)，满足
\[
L_\chi^{\otimes 2}\cong
\mathcal O_{E_{\mathrm D}}\!\Bigl(\sum_{\chi(s)=-1}D_s\Bigr).
\]
将 \(\chi_V,\chi_Y,\chi_H\) 与结论
\ref{con:xi-terminal-zm-s4-v4-ed-branchcount-inertia-115}
中的三条分支公式逐一对应，即得三条平方根关系。
次数由
\[
\deg L_\chi=\frac12\deg\!\Bigl(\sum_{\chi(s)=-1}D_s\Bigr)
\]
立即推出。
最后利用椭圆曲线上的标准同一
\(
\Pic^1(E_{\mathrm D})\cong E_{\mathrm D}
\)，即可将 \(L_\gamma\) 写成 \(2Q\sim D_a+D_b\)。证毕。
\end{proof}

\begin{conclusion}[三条 involution 商映射的分支极端不对称]\label{con:xi-terminal-zm-s4-v4-ed-involution-quotients-branch-2210}
令
\[
\iota_V,\iota_Y,\iota_H\in \operatorname{Gal}(Z/E_{\mathrm D})\setminus\{1\}
\]
分别对应
\[
Z/\langle\iota_V\rangle\cong X_V,\qquad
Z/\langle\iota_Y\rangle\cong Y,\qquad
Z/\langle\iota_H\rangle\cong X_H.
\]
则三条二次商映射分支次数为
\[
\deg\operatorname{Branch}(Z/X_V)=2,\qquad
\deg\operatorname{Branch}(Z/Y)=2,\qquad
\deg\operatorname{Branch}(Z/X_H)=10.
\]
等价地，
\[
\#\mathrm{Fix}_Z(\iota_V)=2,\qquad
\#\mathrm{Fix}_Z(\iota_Y)=2,\qquad
\#\mathrm{Fix}_Z(\iota_H)=10.
\]
\end{conclusion}
\begin{proof}
对任意二次覆盖 \(\rho:Z\to Q\) 有
\[
2g(Z)-2=2\bigl(2g(Q)-2\bigr)+\deg\operatorname{Branch}(\rho).
\]
代入 \(g(Z)=8\)、\(g(X_V)=4\)、\(g(Y)=4\)、\(g(X_H)=2\)，分别得
\[
14=2(6)+\deg\operatorname{Branch}(Z/X_V),\qquad
14=2(6)+\deg\operatorname{Branch}(Z/Y),
\]
\[
14=2(2)+\deg\operatorname{Branch}(Z/X_H),
\]
即分支次数依次为 \(2,2,10\)。
二次商中分支点与对应 involution 的固定点一一对应，故固定点个数结论等价成立。证毕。
\end{proof}

\begin{conclusion}[三条 Prym 因子的同源辨识与 \texorpdfstring{$B$}{B} 的唯一位置]\label{con:xi-terminal-zm-s4-v4-ed-three-prym-identifications}
定义
\[
P_V:=\Prym(Z/X_V),\qquad
P_Y:=\Prym(Z/Y),\qquad
P_H:=\Prym(Z/X_H).
\]
则在 \(\QQ\)-同源意义下
\[
P_V\sim_{\QQ}E\times B,
\]
\[
P_Y\sim_{\QQ}\Jac(X_A)\times E\times E_{\mathrm D},
\]
\[
P_H\sim_{\QQ}\Jac(X_A)\times E_{\mathrm D}\times B.
\]
\end{conclusion}
\begin{proof}
由结论 \ref{con:xi-terminal-zm-s4-delta-eigensplitting-jacz-prymxz}，
\[
\Jac(Z)\sim_{\QQ}\Jac(X_A)\times E\times E_{\mathrm D}^{2}\times B.
\]
另一方面：
\begin{itemize}
  \item 由定理 \ref{thm:xi-terminal-zm-s3-closure-jacobian-splitting}，
  \[
  \Jac(X_V)\sim_{\QQ}\Jac(X_A)\times E_{\mathrm D}^{2};
  \]
  \item 由结论 \ref{con:xi-terminal-zm-s4-c4-prym-threefold-and-v3sgn-identification}，
  \[
  \Jac(Y)\sim_{\QQ}E_{\mathrm D}\times B;
  \]
  \item 由结论 \ref{con:xi-terminal-zm-s4-xh-jacobian-splitting-ed-e}，
  \[
  \Jac(X_H)\sim_{\QQ}E\times E_{\mathrm D}.
  \]
\end{itemize}
对每条二次覆盖 \(Z\to Q\) 使用
\(
\Jac(Z)\sim_{\QQ}\Jac(Q)\times \Prym(Z/Q)
\)，并代入以上三式逐一消去公共因子，即得三条同源辨识。证毕。
\end{proof}

\endinput

\paragraph{\texorpdfstring{$(4,2^5)$}{(4,2^5)} 分歧型的 Nielsen 计数、辫单值群与 Jacobian 三挠几何}
以下结论在本节既定 \(S_4\)-正则覆盖记号下成立，并固定一个阶 \(4\) 分歧点为“无穷远位”。
记 \(\mathcal C_4\subset S_4\) 为四循环共轭类，\(\mathcal C_2\subset S_4\) 为换位共轭类。

\begin{conclusion}[内 Nielsen 类基数 \(160\) 与 \(B_5\) 传递性]\label{con:xi-terminal-zm-s4-nielsen-in-160-b5-transitive}
设
\[
\mathrm{Ni}^{\mathrm{in}}(S_4;\mathcal C_4,\mathcal C_2^5)
:=\left\{
(g_\infty,g_1,\ldots,g_5)\in S_4^6\ \middle|\
\begin{array}{l}
g_\infty\in \mathcal C_4,\ g_i\in\mathcal C_2,\\
g_\infty g_1\cdots g_5=1,\\
\langle g_\infty,g_1,\ldots,g_5\rangle=S_4
\end{array}
\right\}\Big/\!\sim,
\]
其中 \(\sim\) 为同时共轭。
则
\[
\bigl|\mathrm{Ni}^{\mathrm{in}}(S_4;\mathcal C_4,\mathcal C_2^5)\bigr|=160.
\]
并且 \(5\)-股辫群 \(B_5\)（仅作用于 \((g_1,\ldots,g_5)\)）在该 \(160\) 元集合上作用传递。
\end{conclusion}
\begin{proof}
设 Hurwitz 生成元 \(\sigma_i\ (1\le i\le 4)\) 作用为
\[
(g_i,g_{i+1})\longmapsto \bigl(g_{i+1},g_{i+1}^{-1}g_i g_{i+1}\bigr).
\]
脚本
\texttt{scripts/exp\_fold\_zm\_s4\_s3\_nielsen\_monodromy\_audit.py}
对 \(S_4\) 完全枚举给出原始满足约束的六元组数为 \(3840\)，按同时共轭商去后得到
\[
3840/|S_4|=3840/24=160.
\]
同一脚本对 \(\sigma_i,\sigma_i^{-1}\) 诱导图执行轨道遍历，证得唯一轨道，故 \(B_5\) 作用传递。证毕。
\end{proof}

\begin{conclusion}[\texorpdfstring{$S_4\twoheadrightarrow S_3$}{S4 to S3} 诱导的 \(40\) 块四重分块]\label{con:xi-terminal-zm-s4-to-s3-nielsen-40-blocks-of-4}
令 \(V_4^{\mathrm{norm}}\triangleleft S_4\) 为 Klein 四元正规子群，并取商映射
\[
\pi:S_4\longtwoheadrightarrow S_4/V_4^{\mathrm{norm}}\cong S_3.
\]
设 \(\mathcal T\subset S_3\) 为换位共轭类，则 \(\pi\) 诱导映射
\[
\pi_*:\mathrm{Ni}^{\mathrm{in}}(S_4;\mathcal C_4,\mathcal C_2^5)\longrightarrow
\mathrm{Ni}^{\mathrm{in}}(S_3;\mathcal T^6).
\]
并满足：
\begin{enumerate}
  \item \(\bigl|\mathrm{Ni}^{\mathrm{in}}(S_3;\mathcal T^6)\bigr|=40\)；
  \item \(\pi_*\) 满射，且每个纤维基数恒为 \(4\)；
  \item 上述 \(160\) 元集合因此分解为 \(40\) 个 \(B_5\)-不变块，每块大小 \(4\)。
\end{enumerate}
\end{conclusion}
\begin{proof}
同一脚本在 \(S_3\) 侧完全枚举得原始六元组总数 \(240\)，按同时共轭商去得到
\[
240/|S_3|=240/6=40.
\]
对每个 \(S_4\)-内类代表施行 \(\pi\) 并计数纤维，得到纤维大小分布恒为 \(4\)。
由 Hurwitz 作用与群同态、同时共轭均可交换，\(\pi_*\) 的纤维自动组成 \(B_5\)-不变分块系统。证毕。
\end{proof}

\begin{conclusion}[块上辫单值群识别为 \(\mathrm{PSp}_4(\mathbf F_3)\)]\label{con:xi-terminal-zm-braid-monodromy-psp4f3-identification}
记 \(\Gamma_{S_3}\subset S_{40}\) 为 \(B_5\) 在
\(\mathrm{Ni}^{\mathrm{in}}(S_3;\mathcal T^6)\) 上的像，\(\Gamma_{S_4}\subset S_{160}\) 为 \(B_5\) 在
\(\mathrm{Ni}^{\mathrm{in}}(S_4;\mathcal C_4,\mathcal C_2^5)\) 上的像。
则：
\begin{enumerate}
  \item \(|\Gamma_{S_3}|=25920\)，且 \(\Gamma_{S_3}\) 中心平凡、导群等于自身；
  \item 点稳定子次轨道分解为 \((1,12,27)\)；
  \item 因而 \(\Gamma_{S_3}\cong \mathrm{PSp}_4(\mathbf F_3)\)；
  \item \(160\) 点作用为带 \(40\) 个 \(4\)-块的非本原作用，其块上商像即 \(\Gamma_{S_3}\)。
\end{enumerate}
进一步地，若记
\[
|\Gamma_{S_4}|=2^{102}\cdot 3^{44}\cdot 5,
\]
则块核满足
\[
|K|=\frac{|\Gamma_{S_4}|}{|\Gamma_{S_3}|}=2^{96}\cdot 3^{40}.
\]
\end{conclusion}
\begin{proof}
脚本证书给出
\[
|\Gamma_{S_3}|=25920,\quad Z(\Gamma_{S_3})=1,\quad [\Gamma_{S_3},\Gamma_{S_3}]=\Gamma_{S_3},
\]
并给出次轨道 \((1,12,27)\)。
另一方面，\(\mathrm{PSp}_4(\mathbf F_3)\) 的阶为 \(25920\)，其在 \(\mathbf P^3(\mathbf F_3)\) 的 \(40\) 点作用为秩 \(3\)，次轨道正为 \((1,12,27)\)。
由这些群论不变量可识别块上单值群为 \(\mathrm{PSp}_4(\mathbf F_3)\)。

块结构由结论 \ref{con:xi-terminal-zm-s4-to-s3-nielsen-40-blocks-of-4} 给出，故 \(160\) 点作用非本原且块上商像为 \(\Gamma_{S_3}\)。
关于 \(|\Gamma_{S_4}|\) 与 \(|K|\) 的数值，脚本中保留了对应 Schreier--Sims 证书值及
\[
|\Gamma_{S_4}|/|\Gamma_{S_3}|=2^{96}3^{40}
\]
的一致性核验。证毕。
\end{proof}

\begin{conclusion}[四十个分块与属 \(2\) Jacobian 的循环 \(3\)-挠子群]\label{con:xi-terminal-zm-40-blocks-equal-cyclic-3torsion-subgroups}
设
\[
\varphi:X\to\PP^1
\]
为本节固定的 \(S_4\)-Galois 覆盖，分歧型 \((4,2,2,2,2,2)\)。
记
\[
C:=X/A_4,\qquad C_6:=X/V_4^{\mathrm{norm}}.
\]
则：
\begin{enumerate}
  \item \(C\to\PP^1\) 为在六个分歧值处分歧的二次覆盖，故 \(g(C)=2\)；
  \item \(C_6\to C\) 为次数 \(3\) 的循环无分歧覆盖；
  \item 该覆盖由 \(L\in \mathrm{Pic}^0(C)[3]\setminus\{0\}\)（等价于循环子群 \(\langle L\rangle\)）确定；
  \item 因 \(\mathrm{Pic}^0(C)[3]\cong (\ZZ/3\ZZ)^4\)，非平凡循环子群个数为
  \[
  \#\mathbf P^3(\mathbf F_3)=\frac{3^4-1}{3-1}=40;
  \]
  \item 结论 \ref{con:xi-terminal-zm-s4-to-s3-nielsen-40-blocks-of-4} 的 \(40\) 个块可几何识别为这些循环 \(3\)-挠子群。
\end{enumerate}
\end{conclusion}
\begin{proof}
第 \(1\) 点由 \(S_4/A_4\cong C_2\) 及结论
\ref{con:xi-terminal-zm-s4-normal-subgroup-quotients-explicit-models-genus}
直接得到。
第 \(2\) 点由 \(A_4/V_4^{\mathrm{norm}}\cong C_3\) 与结论
\ref{con:xi-terminal-zm-c6-over-c-etale-c3-pic3-class}
给出的无分歧性得到。
第 \(3\) 点为无分歧循环 \(3\)-覆盖与 \(3\)-挠线丛的标准对应。
第 \(4\) 点为有限线性代数计数。
第 \(5\) 点由结论 \ref{con:xi-terminal-zm-s4-to-s3-nielsen-40-blocks-of-4} 的基数与纤维分块，与上述 \(40\)-点参数空间同基数同构识别给出。证毕。
\end{proof}

\begin{theorem}[块轨道的辛正交刻画与强正则结构]\label{thm:xi-terminal-zm-block-orbits-symplectic-orthogonality-srg}
在结论 \ref{con:xi-terminal-zm-40-blocks-equal-cyclic-3torsion-subgroups} 的识别下，将
\[
\mathcal B:=\mathrm{Ni}^{\mathrm{in}}(S_3;\mathcal T^6)
\]
视为 \(V:=\mathrm{Pic}^0(C)[3]\cong (\ZZ/3\ZZ)^4\) 的一维 \(\FF_3\)-子空间（即循环子群）之集合。
记主极化所诱导的 Weil 配对为 \(e_3:V\times V\to\mu_3\)，并对两条线 \(\ell_1,\ell_2\subset V\) 定义关系
\[
\ell_1\perp \ell_2
\quad:\Longleftrightarrow\quad
e_3(x,y)=1\ \text{对任意}\ x\in\ell_1\setminus\{0\},\ y\in\ell_2\setminus\{0\}.
\]
则：
\begin{enumerate}
  \item 在 \(\mathcal B\) 上存在唯一的 \(\Gamma_{S_3}\)-不变非平凡对称二元关系，其在点稳定子下的两条非平凡子轨道基数为 \((12,27)\)；并且该关系与 \(\perp\) 完全一致。
  \item 因而以 \(\mathcal B\) 为顶点、以 \(\ell\neq\ell'\) 且 \(\ell\perp\ell'\) 为边的图为强正则图，参数为
  \[
  (v,k,\lambda,\mu)=(40,12,2,4).
  \]
  \item 其复置换表示分解满足
  \[
  \CC[\mathcal B]\cong \mathbf 1\oplus W_{24}\oplus W_{15},
  \qquad
  \dim W_{24}=24,\ \dim W_{15}=15,
  \]
  等价地邻接算子谱为 \(12,2,-4\)（重数分别为 \(1,24,15\)）。
\end{enumerate}
\end{theorem}

\begin{proof}
由结论 \ref{con:xi-terminal-zm-braid-monodromy-psp4f3-identification}，\(\Gamma_{S_3}\cong \mathrm{PSp}_4(\mathbf F_3)\) 在 \(\mathcal B\) 上为秩 \(3\) 的传递作用，点稳定子次轨道分解为 \((1,12,27)\)；因此任一 \(\Gamma_{S_3}\)-不变对称二元关系由两条非平凡子轨道之一唯一确定，从而至多存在一个非平凡关系。
\par
另一方面，在四维辛空间 \(V\cong\FF_3^4\) 的射影点集 \(\PP(V)\) 上，固定一条线 \(\ell\) 后其正交补 \(\ell^\perp\) 为三维子空间，故 \(\PP(\ell^\perp)\) 含
\(
(3^3-1)/(3-1)=13
\)
个点；去除 \(\ell\) 自身尚余 \(12\) 点，正给出 \(\ell\) 的正交邻域
\(
\mathcal B_\perp(\ell):=\{\ell'\neq\ell:\ \ell'\subset\ell^\perp\}
\)。
余下 \(27\) 点为非正交类。
因此 \(\perp\) 给出一条 \((12,27)\) 子轨道划分，并与上述唯一性一致，得到第 (1) 条。
\par
对强正则参数，若 \(\ell_1\perp\ell_2\) 且 \(\ell_1\neq\ell_2\)，则 \(\ell_1^\perp\cap \ell_2^\perp\) 为二维子空间，其射影点数为 \(4\)；其中包含 \(\ell_1,\ell_2\)，故公共邻点数为 \(2\)，即 \(\lambda=2\)。
若 \(\ell_1\not\perp\ell_2\)，则 \(\ell_1,\ell_2\not\subset \ell_1^\perp\cap\ell_2^\perp\)，公共邻点数为 \(4\)，即 \(\mu=4\)。
于是得到第 (2) 条。
第 (3) 条为强正则图的一般谱结论，与 Klingen Hecke 分解中的同一谱值一致（参见定理 \ref{thm:conclusion-m2-level3-delta0-inertia-klingen-charpoly}）。证毕。
\end{proof}

\begin{proposition}[由强正则图内生重构辛极空间 \(W(3)\) 的线系]\label{prop:xi-terminal-zm-reconstruct-W3-from-srg}
设图为定理 \ref{thm:xi-terminal-zm-block-orbits-symplectic-orthogonality-srg} 的强正则图。
则其所有极大团的大小均为 \(4\)。
记 \(\mathcal L\) 为全部 \(4\)-团之集合，则：
\begin{enumerate}
  \item \(|\mathcal L|=40\)；
  \item 每个顶点恰落在 \(4\) 个 \(4\)-团中；
  \item \((\mathcal B,\mathcal L,\in)\) 构成阶 \(3\) 的广义四边形（辛极空间）\(W(3)\)：每点在 \(4\) 条线上、每线含 \(4\) 点，且满足广义四边形公理。
\end{enumerate}
\end{proposition}

\begin{proof}
在 \(\mathcal B\simeq \PP(V)\) 的识别下，\(4\)-团等价于 \(\PP(M)\)，其中 \(M\subset V\) 为二维全各向同性子空间（Lagrangian 平面）。
确有 \(|\PP(M)|=(3^2-1)/(3-1)=4\)，且任意两点正交，故给出 \(4\)-团。
反向地，若四点两两正交，则其张成的子空间必为二维各向同性子空间，从而团必为某个 \(\PP(M)\)。
\par
计数：固定一点 \(\ell\in\PP(V)\)，包含 \(\ell\) 的 Lagrangian 平面个数等于 \(\PP(\ell^\perp/\ell)\) 的点数；而 \(\ell^\perp/\ell\) 为二维 \(\FF_3\)-空间，故该点数为 \((3^2-1)/(3-1)=4\)，得到第 (2) 条。
再由总旗数计数得 \(|\mathcal L|\cdot 4=|\mathcal B|\cdot 4=160\)，故 \(|\mathcal L|=40\)，得第 (1) 条。
广义四边形公理为辛极空间的标准性质：对点 \(\ell\) 与不含 \(\ell\) 的 Lagrangian 平面 \(M\)，存在唯一 \(\ell'\in\PP(M)\) 使 \(\ell\perp\ell'\)，从而唯一线 \(\langle \ell,\ell'\rangle\) 与 \(M\) 相交于 \(\ell'\)。证毕。
\end{proof}

\begin{theorem}[\texorpdfstring{\(\pi_*\)}{pi*} 的四重纤维与“通过一点的四条线”]\label{thm:xi-terminal-zm-pistar-fiber-equals-four-lines-through-point}
\leanverified{paper\_xi\_terminal\_zm\_pistar\_fiber\_equals\_four\_lines\_through\_point}
保持结论 \ref{con:xi-terminal-zm-s4-to-s3-nielsen-40-blocks-of-4} 的四重覆盖
\(
\pi_*:\mathrm{Ni}^{\mathrm{in}}(S_4;\mathcal C_4,\mathcal C_2^5)\to \mathcal B
\)。
固定 \(b\in\mathcal B\)，令 \(B:=\pi_*^{-1}(b)\) 为其纤维。
在定理 \ref{thm:xi-terminal-zm-block-orbits-symplectic-orthogonality-srg} 与命题 \ref{prop:xi-terminal-zm-reconstruct-W3-from-srg} 的识别下，将 \(b\) 视为 \(V\) 中一条线 \(\ell\subset V\)，并记
\[
\mathcal L(\ell):=\{M\in\mathcal L:\ \ell\in M\},
\qquad |\mathcal L(\ell)|=4.
\]
则存在一个由单值作用所确定的同构类，使得纤维 \(B\) 的点稳定子单值表示与 \(\mathcal L(\ell)\) 的自然作用同构；从而在该同构类意义下可将 \(B\) 识别为“通过点 \(\ell\) 的四条线”集合 \(\mathcal L(\ell)\)。
在此识别下，\(\ell\) 的正交邻域分裂为四个互不交叠的三元组
\[
\mathcal B_\perp(\ell)
=
\bigsqcup_{M\in\mathcal L(\ell)}\bigl(\PP(M)\setminus\{\ell\}\bigr),
\qquad |\mathcal B_\perp(\ell)|=12,
\]
并且上述分裂由纤维 \(B\) 内在标记。
\end{theorem}

\begin{proof}
由结论 \ref{con:xi-terminal-zm-braid-monodromy-psp4f3-identification}，块上单值群为 \(G=\Gamma_{S_3}\cong \mathrm{PSp}_4(\mathbf F_3)\)，并且 \(G\) 在 \(\mathcal B\) 上秩为 \(3\)。
设 \(P=\mathrm{Stab}_G(\ell)\)。
在辛几何中，集合 \(\mathcal L(\ell)\) 与 \(\PP(\ell^\perp/\ell)\cong \PP^1(\FF_3)\) 等价，且 \(P\) 在其上诱导的作用同构于 \(\mathrm{PGL}_2(\FF_3)\cong S_4\) 的自然四点作用，因此该四点传递作用在同构意义下唯一。
另一方面，\(\pi_*\) 给出大小为 \(4\) 的纤维，且纤维上的单值表示由点稳定子诱导；于是纤维单值表示与 \(\mathcal L(\ell)\) 的自然表示在同构类意义下必须一致。
最后的三元组分裂由 \(\PP(M)\) 各含四点且不同 \(M\) 仅交于 \(\ell\) 直接推出。证毕。
\end{proof}

\begin{corollary}[\texorpdfstring{\((3,3)\)}{(3,3)}-同源核的四重扩充律]\label{cor:xi-terminal-zm-cyclic3-to-33-isogeny-four-extensions}
令 \(A=\mathrm{Pic}^0(C)\) 为主极化阿贝尔曲面。
对任一非平凡循环子群 \(\ell\subset A[3]\)，包含 \(\ell\) 的极大各向同性子群 \(M\subset A[3]\) 恰有 \(4\) 个，并与 \(\mathcal L(\ell)\) 等势。
对每个 \(M\) 商 \(A/M\) 继承主极化，从而诱导一条 \((3,3)\)-同源
\[
\varphi_M:A\to A/M.
\]
因此，“给定循环 \(3\)-同源核 \(\ell\)”具有且仅具有四种方式可扩充为一个 \((3,3)\)-同源核 \(M\)。
\end{corollary}

\begin{proof}
极大各向同性二维子空间与 \((3,3)\)-同源核的对应由点--线四重纤维识别给出（见定理 \ref{thm:xi-terminal-zm-pistar-fiber-equals-four-lines-through-point}）。
而包含给定 \(\ell\) 的极大各向同性平面数目为 \(|\PP(\ell^\perp/\ell)|=4\)，与命题 \ref{prop:xi-terminal-zm-reconstruct-W3-from-srg} 一致。证毕。
\end{proof}

\begin{proposition}[线系的强正则对偶与旗空间的双 \(4\)-重投影]\label{prop:xi-terminal-zm-line-srg-and-flag-double-projection}
令 \(\mathcal L\) 为命题 \ref{prop:xi-terminal-zm-reconstruct-W3-from-srg} 的线系，并在 \(\mathcal L\) 上定义邻接关系
\[
M\sim N\quad:\Longleftrightarrow\quad \card{M\cap N}=3
\quad(\text{即 } M\cap N\ \text{为一条线}).
\]
则：
\begin{enumerate}
  \item \((\mathcal L,\sim)\) 亦为参数 \((40,12,2,4)\) 的强正则图；
  \item 点–线关联图 \(\mathcal I\subset \mathcal B\times\mathcal L\)（\((\ell,M)\in\mathcal I\iff \ell\subset M\)）为 \((4,4)\)-正则二部图，边数为 \(160\)；
  \item \(\mathcal I\to\mathcal B\) 与 \(\mathcal I\to\mathcal L\) 均为四重投影，并在组合意义下互为对偶：每个点邻域的 \(12\) 个正交邻点被唯一分裂为四个三元组，对应于通过该点的四条线；线邻域亦同理。
\end{enumerate}
\end{proposition}

\begin{proof}
在辛极空间 \(W(3)\) 中点与线完全对偶。
固定 \(M\in\mathcal L\)，其包含四条线 \(\ell\subset M\)；每条 \(\ell\) 被恰好四条线包含，其中之一为 \(M\) 本身，故与 \(M\) 交于 \(\ell\) 的其他线有三条。
因此 \(M\) 的邻接度为 \(4\cdot 3=12\)，余下为 \(27\)，与点图同参。
其余参数由同样的二维交与公共邻点计数得到，与点图的 \((\lambda,\mu)=(2,4)\) 完全平行。
二部图与双四重投影为命题 \ref{prop:xi-terminal-zm-reconstruct-W3-from-srg} 的计数结论之直接改写；对偶分裂由
\(
\mathcal B_\perp(\ell)=\bigsqcup_{M\in\mathcal L(\ell)}(\PP(M)\setminus\{\ell\})
\)
给出。证毕。
\end{proof}

\begin{conjecture}[内 Hurwitz 空间的抛物型 \(3\)-水平结构模解释]\label{conj:xi-terminal-zm-hurwitz-parabolic-level3-moduli}
设 \(\mathcal H_{S_3}^{\mathrm{in}}(\mathcal T^6)\) 与 \(\mathcal H_{S_4}^{\mathrm{in}}(\mathcal C_4,\mathcal C_2^5)\) 分别为相应分歧型别的内 Hurwitz 空间连通分量。
则期望存在一条与结论 \ref{con:xi-terminal-zm-40-blocks-equal-cyclic-3torsion-subgroups} 的 Jacobian 三挠识别与定理 \ref{thm:xi-terminal-zm-pistar-fiber-equals-four-lines-through-point} 的点–线–旗几何同时相容的双有理同构链，使得：
\begin{enumerate}
  \item \(\mathcal H_{S_3}^{\mathrm{in}}(\mathcal T^6)\) 与某个抛物型 \(3\)-水平结构的 Siegel 模叠（等价地：主极化阿贝尔曲面携带一条 \(3\)-挠线 \(\ell\subset A[3]\) 的模叠）在双有理意义下同构；
  \item 覆盖 \(\mathcal H_{S_4}^{\mathrm{in}}(\mathcal C_4,\mathcal C_2^5)\to \mathcal H_{S_3}^{\mathrm{in}}(\mathcal T^6)\) 对应于进一步选择一条包含 \(\ell\) 的极大各向同性子群 \(M\subset A[3]\)（即选择 \((3,3)\)-同源核），其几何纤维与 \(\PP(\ell^\perp/\ell)\cong \PP^1(\FF_3)\) 的四点截面一致。
\end{enumerate}
该猜想的可核验截面由 \(\Gamma_{S_3}\cong \mathrm{PSp}_4(\mathbf F_3)\) 的块上单值化、强正则结构与旗空间双投影给出；其几何提升需在代数簇层面构造相应的 \(A[3]\) 局部系统并验证其与 Hurwitz 单值数据的自然性一致。
\end{conjecture}

\begin{corollary}[Hurwitz 模空间不可约性与“属 \(2\)+\(3\)-挠数据”解释]\label{cor:xi-terminal-zm-hurwitz-irreducible-and-genus2-jac3torsion-moduli}
记 \(\mathcal H_{S_4}\) 为上述 \((4,2^5)\) 型 \(S_4\)-覆盖的 Hurwitz 模空间，
\(\mathcal H_{S_3}\) 为其 \(S_3\)-商型空间。则：
\begin{enumerate}
  \item \(\mathcal H_{S_4}\) 与 \(\mathcal H_{S_3}\) 均几何不可约，维数均为 \(3\)；
  \item 存在有限 \'{e}tale 映射 \(\mathcal H_{S_4}\to \mathcal H_{S_3}\)，泛纤维基数为 \(4\)；
  \item \(\mathcal H_{S_3}\) 可解释为“带六个分歧标记的属 \(2\) 曲线 \(C\) 及其 \(J(C)[3]\) 非平凡循环子群”之粗模空间。
\end{enumerate}
\end{corollary}
\begin{proof}
由结论 \ref{con:xi-terminal-zm-s4-nielsen-in-160-b5-transitive}、
\ref{con:xi-terminal-zm-s4-to-s3-nielsen-40-blocks-of-4} 的辫群传递性，Hurwitz 覆盖在分歧点配置空间上连通，从而几何不可约。
分歧点总数为 \(6\)，故维数为 \(6-3=3\)。
第 \(2\) 点由纤维大小恒为 \(4\) 得到。
第 \(3\) 点由结论 \ref{con:xi-terminal-zm-40-blocks-equal-cyclic-3torsion-subgroups} 给出的几何识别得到。证毕。
\end{proof}

\begin{proposition}[Hurwitz 生成元循环型的离散证书]\label{prop:xi-terminal-zm-hurwitz-generator-cycle-types-on-160-and-40}
在上述置换表示中，对任一标准生成元 \(\sigma_i\ (1\le i\le 4)\)：
\begin{enumerate}
  \item 在 \(160\) 点集合 \(\mathrm{Ni}^{\mathrm{in}}(S_4;\mathcal C_4,\mathcal C_2^5)\) 上，
  \[
  \sigma_i\sim 3^{36}2^{12}1^{28},
  \]
  因而置换阶为 \(6\)；
  \item 在 \(40\) 点集合 \(\mathrm{Ni}^{\mathrm{in}}(S_3;\mathcal T^6)\) 上，
  \[
  \sigma_i\sim 3^{9}1^{13},
  \]
  因而置换阶为 \(3\)。
\end{enumerate}
\end{proposition}
\begin{proof}
由脚本对四个 Hurwitz 生成元的循环分解直接给出：
\[
\{1^{28},2^{12},3^{36}\}\quad\text{与}\quad \{1^{13},3^{9}\},
\]
且四个生成元同型。置换阶由循环长度最小公倍数得到。证毕。
\end{proof}

\begin{conclusion}[模空间像的不可约性与一般点自同构群]\label{con:xi-terminal-zm-m16-locus-irreducible-generic-aut-s4}
令 \(\mathcal L\subset \mathcal M_{16}\) 为该 Hurwitz 家族在属 \(16\) 模空间中的像。
则：
\begin{enumerate}
  \item \(\mathcal L\) 为维数 \(3\) 的不可约代数簇；
  \item \(\mathcal L\) 一般点对应曲线 \(X\) 满足
  \[
  \mathrm{Aut}(X)=S_4.
  \]
\end{enumerate}
\end{conclusion}
\begin{proof}
第 \(1\) 点由推论
\ref{cor:xi-terminal-zm-hurwitz-irreducible-and-genus2-jac3torsion-moduli}
及模映射像的不可约性得到。
第 \(2\) 点中，对一般分歧配置，基底 \(\PP^1\) 上保持分歧集合的自同构群为平凡。
任何 \(X\) 的自同构都必须保持 \(S_4\)-覆盖结构并在基底诱导该平凡自同构，故仅剩甲板变换群 \(S_4\) 本身。证毕。
\end{proof}

\begin{remark}[可复现证书路径]\label{rem:xi-terminal-zm-s4-s3-nielsen-monodromy-certificate-path}
本组结论的离散计数与置换群证书由脚本
\texttt{scripts/exp\_fold\_zm\_s4\_s3\_nielsen\_monodromy\_audit.py}
生成，输出文件为
\texttt{artifacts/export/fold\_zm\_s4\_s3\_nielsen\_monodromy\_audit.json}。
\end{remark}

\endinput


\paragraph{\texorpdfstring{$S_4$}{S4} 典范表示的共轭类谱型与低次数关系模的等变闭式}
沿用前文记号，设 \(X=\widetilde{X}\) 且
\[
V:=H^0(X,K_X),\qquad \dim V=16.
\]
并记 \(S:=\mathrm{Sym}(V)\)，\(I\subset S\) 为典范嵌入的齐次理想，\(I_m:=I\cap \mathrm{Sym}^mV\)。

\begin{conclusion}[典范表示在共轭类上的特征多项式刚性确定]\label{con:xi-terminal-zm-s4-canonical-characteristic-polynomial-conjugacy-rigidity}
对 \(S_4\) 的共轭类代表元 \(g\)（按循环型）在 \(V\) 上的特征多项式为
\[
\chi_{V,g}(T)=
\begin{cases}
(T-1)^{16}, & g=1,\\[1mm]
(T-1)^6(T+1)^{10}, & g\sim(12),\\[1mm]
(T-1)^8(T+1)^8, & g\sim(12)(34),\\[1mm]
(T-1)^6(T^2+T+1)^5, & g\sim(123),\\[1mm]
(T-1)^4(T+1)^4(T^2+1)^4, & g\sim(1234).
\end{cases}
\]
因此，典范表示在各共轭类上的谱型由上述多项式唯一确定。
\end{conclusion}
\begin{proof}[证明要点]
由结论 \ref{con:xi-terminal-zm-s4-chevalley-weil-explicit-decomposition}
\[
V\cong 2\varepsilon\oplus \rho_2\oplus \rho_3\oplus 3\rho_3'
\]
可得类函数迹
\[
\tr(g\mid V)=
\begin{cases}
16,& g=1,\\
-4,& g\sim(12),\\
0,& g\sim(12)(34),\\
1,& g\sim(123),\\
0,& g\sim(1234).
\end{cases}
\]
对阶 \(2\) 元，特征值仅为 \(\pm1\)，由迹与维数可解出重数，得到前两类式子。
对阶 \(3\) 元，特征值属于 \(\{1,\omega,\omega^2\}\)；又 \(\tr(g\mid V)=\tr(g^2\mid V)=1\)，故 \(1\) 的重数为 \(6\)，\(\omega,\omega^2\) 的重数各为 \(5\)。
对阶 \(4\) 元，特征值属于 \(\{1,-1,i,-i\}\)；由 \(\tr(g\mid V)=0\) 与 \(\tr(g^2\mid V)=\tr((12)(34)\mid V)=0\) 解得四者重数均为 \(4\)。证毕。
\end{proof}

\begin{conclusion}[对称代数的类函数型 Molien 生成函数与等型分量级数]\label{con:xi-terminal-zm-s4-symmetric-algebra-molien-isotypic-series}
对任意 \(g\in S_4\)，定义
\[
\Phi_g(t):=\sum_{m\ge0}\tr\!\bigl(g\mid \mathrm{Sym}^mV\bigr)t^m.
\]
则 \(\Phi_g(t)\) 仅依赖于 \(g\) 的共轭类，并有闭式
\[
\Phi_g(t)=
\begin{cases}
(1-t)^{-16}, & g=1,\\[1mm]
(1-t)^{-6}(1+t)^{-10}, & g\sim(12),\\[1mm]
(1-t)^{-8}(1+t)^{-8}, & g\sim(12)(34),\\[1mm]
(1-t)^{-6}(1+t+t^2)^{-5}, & g\sim(123),\\[1mm]
(1-t)^{-4}(1+t)^{-4}(1+t^2)^{-4}, & g\sim(1234).
\end{cases}
\]
进一步，对每个不可约表示
\(\sigma\in\{\mathbf 1,\varepsilon,\rho_2,\rho_3,\rho_3'\}\)，定义
\[
M_\sigma(t):=\sum_{m\ge0}\mathrm{mult}_\sigma\!\bigl(\mathrm{Sym}^mV\bigr)t^m.
\]
则有
\begin{align*}
M_{\mathbf 1}(t)
&=\frac1{24}(1-t)^{-16}
+\frac14(1-t)^{-6}(1+t)^{-10}
+\frac18(1-t)^{-8}(1+t)^{-8} \\
&\quad
+\frac13(1-t)^{-6}(1+t+t^2)^{-5}
+\frac14(1-t)^{-4}(1+t)^{-4}(1+t^2)^{-4},\\[1mm]
M_{\varepsilon}(t)
&=\frac1{24}(1-t)^{-16}
-\frac14(1-t)^{-6}(1+t)^{-10}
+\frac18(1-t)^{-8}(1+t)^{-8} \\
&\quad
+\frac13(1-t)^{-6}(1+t+t^2)^{-5}
-\frac14(1-t)^{-4}(1+t)^{-4}(1+t^2)^{-4},\\[1mm]
M_{\rho_2}(t)
&=\frac1{12}(1-t)^{-16}
+\frac14(1-t)^{-8}(1+t)^{-8}
-\frac13(1-t)^{-6}(1+t+t^2)^{-5},\\[1mm]
M_{\rho_3}(t)
&=\frac18(1-t)^{-16}
+\frac14(1-t)^{-6}(1+t)^{-10}
-\frac18(1-t)^{-8}(1+t)^{-8}
-\frac14(1-t)^{-4}(1+t)^{-4}(1+t^2)^{-4},\\[1mm]
M_{\rho_3'}(t)
&=\frac18(1-t)^{-16}
-\frac14(1-t)^{-6}(1+t)^{-10}
-\frac18(1-t)^{-8}(1+t)^{-8}
+\frac14(1-t)^{-4}(1+t)^{-4}(1+t^2)^{-4}.
\end{align*}
\end{conclusion}
\begin{proof}[证明要点]
对任意线性算子 \(A\)，恒有
\[
\sum_{m\ge0}\tr\!\bigl(\mathrm{Sym}^m(A)\bigr)t^m
=\det(I-tA)^{-1}.
\]
将结论 \ref{con:xi-terminal-zm-s4-canonical-characteristic-polynomial-conjugacy-rigidity} 的各类特征值代入即可得到 \(\Phi_g(t)\)。
再由特征标正交关系
\[
M_\sigma(t)=\frac1{|S_4|}\sum_{[g]}|[g]|\,\chi_\sigma(g)\,\Phi_g(t),
\]
代入 \(S_4\) 的共轭类大小与特征标值，化简即得各显式式。证毕。
\end{proof}

\begin{conclusion}[三次与四次对称幂及典范理想低次分量的 \texorpdfstring{$S_4$}{S4} 分解]\label{con:xi-terminal-zm-s4-sym3-sym4-canonical-ideal-low-degree-decomposition}
有如下等变分解：
\[
\mathrm{Sym}^3V
\cong
25\mathbf 1\oplus 47\varepsilon\oplus 66\rho_2\oplus 91\rho_3\oplus 113\rho_3',
\]
\[
\mathrm{Sym}^4V
\cong
198\mathbf 1\oplus 138\varepsilon\oplus 330\rho_2\oplus 508\rho_3\oplus 452\rho_3'.
\]
并且
\[
I_3
\cong
23\mathbf 1\oplus 43\varepsilon\oplus 60\rho_2\oplus 83\rho_3\oplus 102\rho_3',
\]
\[
I_4
\cong
192\mathbf 1\oplus 135\varepsilon\oplus 321\rho_2\oplus 494\rho_3\oplus 440\rho_3'.
\]
\end{conclusion}
\begin{proof}[证明要点]
由结论 \ref{con:xi-terminal-zm-s4-symmetric-algebra-molien-isotypic-series}，
\[
\mathrm{mult}_\sigma\!\bigl(\mathrm{Sym}^mV\bigr)=[t^m]M_\sigma(t),
\]
取 \(m=3,4\) 即得 \(\mathrm{Sym}^3V,\mathrm{Sym}^4V\) 分解。
另一方面，由结论 \ref{con:xi-terminal-zm-s4-canonical-projective-normality-quadratic-generation} 的射影正规性，
\[
\mathrm{Sym}^mV\twoheadrightarrow H^0(X,mK_X)\qquad (m\ge 0),
\]
故
\[
0\to I_m\to \mathrm{Sym}^mV\to H^0(X,mK_X)\to 0
\]
为 \(S_4\)-等变正合列。由结论 \ref{con:xi-terminal-zm-s4-pluricanonical-irrep-decomposition-all-n} 在
\(n=3,4\) 时
\[
H^0(X,3K_X)\cong
2\mathbf 1\oplus 4\varepsilon\oplus 6\rho_2\oplus 8\rho_3\oplus 11\rho_3',
\]
\[
H^0(X,4K_X)\cong
6\mathbf 1\oplus 3\varepsilon\oplus 9\rho_2\oplus 14\rho_3\oplus 12\rho_3',
\]
在表示环中作差即得 \(I_3,I_4\) 分解。证毕。
\end{proof}

\begin{conclusion}[二次生成下三次层线性 syzygy 模的完全分解]\label{con:xi-terminal-zm-s4-linear-syzygy-module-complete-decomposition}
记
\[
Q:=I_2,\qquad
\mathcal S_3:=\ker\!\bigl(Q\otimes V\to I_3\bigr).
\]
则
\[
\mathcal S_3
\cong
27\mathbf 1\oplus 29\varepsilon\oplus 61\rho_2\oplus 88\rho_3\oplus 91\rho_3'.
\]
\end{conclusion}
\begin{proof}[证明要点]
由结论 \ref{con:xi-terminal-zm-s4-canonical-projective-normality-quadratic-generation}，典范理想由二次式生成，故 \(Q\otimes V\twoheadrightarrow I_3\) 为 \(S_4\)-等变满射。
结合结论 \ref{con:xi-terminal-zm-s4-canonical-quadrics-equivariant-decomposition} 的
\[
Q\cong 8\mathbf 1\oplus 2\varepsilon\oplus 9\rho_2\oplus 13\rho_3\oplus 8\rho_3'
\]
与结论 \ref{con:xi-terminal-zm-s4-chevalley-weil-explicit-decomposition} 的 \(V\) 分解，按特征标点乘可得
\[
Q\otimes V
\cong
50\mathbf 1\oplus 72\varepsilon\oplus 121\rho_2\oplus 171\rho_3\oplus 193\rho_3'.
\]
再由结论 \ref{con:xi-terminal-zm-s4-sym3-sym4-canonical-ideal-low-degree-decomposition} 的 \(I_3\) 分解，在表示环中作差
\(
[\mathcal S_3]=[Q\otimes V]-[I_3]
\)
即得结论。证毕。
\end{proof}

\begin{conclusion}[四次层二次系数核表示的闭式分解]\label{con:xi-terminal-zm-s4-quadratic-coefficient-kernel-degree4-closed-decomposition}
定义
\[
\mathcal K_4:=\ker\!\bigl(Q\otimes \mathrm{Sym}^2V\to I_4\bigr).
\]
则有
\[
\mathcal K_4
\cong
375\mathbf 1\oplus 344\varepsilon\oplus 724\rho_2\oplus 1090\rho_3\oplus 1056\rho_3'.
\]
\end{conclusion}
\begin{proof}[证明要点]
仍由二次生成性，\(Q\otimes \mathrm{Sym}^2V\twoheadrightarrow I_4\) 为 \(S_4\)-等变满射。
由结论 \ref{con:xi-terminal-zm-s4-canonical-quadrics-equivariant-decomposition} 与
\ref{con:xi-terminal-zm-s4-symmetric-algebra-molien-isotypic-series} 在 \(m=2\) 的系数可得
\[
\mathrm{Sym}^2V
\cong
11\mathbf 1\oplus 3\varepsilon\oplus 13\rho_2\oplus 20\rho_3\oplus 12\rho_3'.
\]
特征标点乘后得到
\[
Q\otimes \mathrm{Sym}^2V
\cong
567\mathbf 1\oplus 479\varepsilon\oplus 1045\rho_2\oplus 1584\rho_3\oplus 1496\rho_3'.
\]
再由结论 \ref{con:xi-terminal-zm-s4-sym3-sym4-canonical-ideal-low-degree-decomposition} 的 \(I_4\) 分解，作差
\(
[\mathcal K_4]=[Q\otimes \mathrm{Sym}^2V]-[I_4]
\)
即得所示重数。证毕。
\end{proof}

\endinput


\subsubsection{谱兼容外置、半直积修复与有限/无穷二影刚性}\label{subsubsec:xi-spectrum-compatible-semidirect-two-shadow-rigidity}
本段在 Lee--Yang 四层覆盖的既有单值化框架上，将“可交换外置障碍—非交换修复—线性椭圆双影刚性”组织为同一条可检验链条。

\begin{theorem}[可交换外置寄存器对非交换单值群的不可完备性]\label{thm:xi-terminal-abelian-register-obstruction-nonabelian-monodromy}
设 \(U\) 为道路连通拓扑空间，\(\pi:X\to U\) 为连通有限分支覆盖，分支集记为 \(B\subset U\)，并取基点 \(u_0\in U\setminus B\)。
令
\[
\rho:\pi_1(U\setminus B,u_0)\longrightarrow \mathfrak S_n
\]
为单值表示，且假设 \(\rho(\pi_1(U\setminus B,u_0))\) 非交换。

称三元组 \((A,\iota,D)\) 为该覆盖的一组可交换外置寄存器数据，若：
\begin{enumerate}
  \item \(A\) 为交换群；
  \item \(\iota:\pi_1(U\setminus B,u_0)\to A\) 为群同态；
  \item \(D:\{1,\dots,n\}\times A\to\{1,\dots,n\}\) 满足对任意 \(\gamma\in\pi_1(U\setminus B,u_0)\) 与 \(i\in\{1,\dots,n\}\),
  \[
  \rho(\gamma)(i)=D\bigl(i,\iota(\gamma)\bigr).
  \]
\end{enumerate}
则此类三元组不存在。
\end{theorem}
\begin{proof}
由 \(\rho(\pi_1(U\setminus B,u_0))\) 非交换，可取 \(\alpha,\beta\in\pi_1(U\setminus B,u_0)\) 使
\[
\rho([\alpha,\beta])\neq \mathrm{id},
\qquad
[\alpha,\beta]:=\alpha\beta\alpha^{-1}\beta^{-1}.
\]
由 \(A\) 交换且 \(\iota\) 为同态，
\[
\iota([\alpha,\beta])
=\iota(\alpha)\iota(\beta)\iota(\alpha)^{-1}\iota(\beta)^{-1}
=e_A
=\iota(1).
\]
故对任意 \(i\),
\[
D\bigl(i,\iota([\alpha,\beta])\bigr)=D(i,e_A)=D\bigl(i,\iota(1)\bigr).
\]
左侧按定义等于 \(\rho([\alpha,\beta])(i)\)，右侧等于 \(\rho(1)(i)=i\)，从而 \(\rho([\alpha,\beta])(i)=i\) 对所有 \(i\) 成立，矛盾于 \(\rho([\alpha,\beta])\neq\mathrm{id}\)。
证毕。
\end{proof}

\paragraph{window-6 三轴中心约束、Pati--Salam 连通实现与标准模型包络的中心 2-挠兼容性}
在 window-6 的 boundary 读出中，三轴独立 sheet-flip 给出一个抽象中心寄存器
\[
\mathcal C_{\mathrm{bdry}}\cong (\ZZ/2\ZZ)^3.
\]
本段讨论其与紧连通 Lie 包络在中心 2-挠层面的可实现条件。

\begin{theorem}[标准模型连通包络对三轴中心寄存器的不可实现性]\label{thm:xi-window6-sm-envelope-center-2torsion-obstruction}
设 \(G\) 为紧连通 Lie 群，满足
\[
\Lie(G)\cong \mathfrak u(1)\oplus\mathfrak{su}(2)\oplus\mathfrak{su}(3).
\]
则
\[
\dim_{\FF_2} Z(G)[2]\le 2.
\]
从而不存在单同态
\[
\mathcal C_{\mathrm{bdry}}\cong (\ZZ/2\ZZ)^3\hookrightarrow Z(G)[2].
\]
\end{theorem}
\begin{proof}
由紧连通群分类，可取
\[
G\cong \widetilde G/\Gamma,\qquad
\widetilde G:=U(1)\times SU(2)\times SU(3),
\]
其中 \(\Gamma\subset Z(\widetilde G)\) 为有限中心子群。记
\[
A:=Z(\widetilde G)\cong U(1)\times \ZZ_2\times \ZZ_3,
\qquad
Z(G)\cong A/\Gamma.
\]
设 \(\pi:A\to \ZZ_2\times \ZZ_3\) 为自然投影，则有短正合列
\[
1\to U(1)\to A\to \ZZ_2\times \ZZ_3\to 1.
\]
对 \(\Gamma\) 取商后得到
\[
1\to T\to A/\Gamma\to Q\to 1,
\]
其中 \(T\cong U(1)\)（一维紧环面），\(Q\) 为 \((\ZZ_2\times \ZZ_3)\) 的商群。

记 \(E:=A/\Gamma\)。考虑 \(E[2]\to Q[2]\)。对任意 \(q\in Q[2]\)，其纤维中两点之差落在 \(T[2]\)。
故每个纤维大小至多 \(|T[2]|=2\)，并且
\[
|Q[2]|\le |(\ZZ_2\times \ZZ_3)[2]|=2.
\]
从而
\[
|E[2]|\le 2\cdot 2=4,\qquad
\dim_{\FF_2}E[2]\le 2.
\]
即 \(\dim_{\FF_2} Z(G)[2]\le 2\)。因此 \((\ZZ/2\ZZ)^3\) 不能嵌入 \(Z(G)[2]\)。证毕。
\end{proof}

\begin{theorem}[三轴 canonical sheet-flip 保真下降下的 Pati--Salam 连通实现唯一性]\label{thm:xi-window6-ps-faithful-canonical-descent-unique}
设
\[
\widetilde G_{\mathrm{PS}}:=SU(4)\times SU(2)_L\times SU(2)_R,\qquad
G:=\widetilde G_{\mathrm{PS}}/H,
\]
其中 \(H\subset Z(\widetilde G_{\mathrm{PS}})\) 为有限中心子群。记 canonical involution 子群
\[
\mathcal C_{\mathrm{can}}
:=Z(SU(4))[2]\times Z(SU(2)_L)\times Z(SU(2)_R)
\cong (\ZZ/2\ZZ)^3.
\]
若商映射 \(q:\widetilde G_{\mathrm{PS}}\to G\) 在 \(\mathcal C_{\mathrm{can}}\) 上限制为单同态
\[
q|_{\mathcal C_{\mathrm{can}}}:\mathcal C_{\mathrm{can}}\hookrightarrow Z(G)[2],
\]
则必有 \(H=\{e\}\)，因而
\[
G\cong SU(4)\times SU(2)_L\times SU(2)_R.
\]
\end{theorem}
\begin{proof}
有
\[
Z(\widetilde G_{\mathrm{PS}})
\cong \ZZ_4\times \ZZ_2\times \ZZ_2.
\]
取任意非平凡 \(H\subset Z(\widetilde G_{\mathrm{PS}})\)。任取 \(h\in H\setminus\{e\}\)。
若 \(\ord(h)=2\)，则 \(h\in \mathcal C_{\mathrm{can}}\)；
若 \(\ord(h)=4\)，则 \(h^2\neq e\) 且 \(h^2\in \mathcal C_{\mathrm{can}}\)。
故总有
\[
H\cap \mathcal C_{\mathrm{can}}\neq \{e\}.
\]
而 \(q\) 在 \(H\) 上恒等于单位元，故 \(q|_{\mathcal C_{\mathrm{can}}}\) 的核至少含
\(H\cap \mathcal C_{\mathrm{can}}\) 的非平凡元，不能单射，矛盾。
故 \(H\) 只能为平凡子群。证毕。
\end{proof}

\begin{theorem}[Pati--Salam 到标准模型包络的中心 2-挠压缩不可避免]\label{thm:xi-window6-ps-to-sm-center2-kernel}
设
\[
G_{\mathrm{PS}}:=SU(4)\times SU(2)_L\times SU(2)_R,
\]
且 \(G_{\mathrm{SM}}\) 为任意紧连通 Lie 群，满足
\(
\Lie(G_{\mathrm{SM}})\cong \mathfrak u(1)\oplus\mathfrak{su}(2)\oplus\mathfrak{su}(3)
\)。
若连续群同态
\(
f:G_{\mathrm{PS}}\to G_{\mathrm{SM}}
\)
满足中心兼容条件
\[
f\bigl(\mathcal C_{\mathrm{can}}\bigr)\subseteq Z(G_{\mathrm{SM}})[2],
\]
则限制映射
\[
f_*:\mathcal C_{\mathrm{can}}\longrightarrow Z(G_{\mathrm{SM}})[2]
\]
必有非平凡核。
\end{theorem}
\begin{proof}
由定理 \ref{thm:xi-window6-sm-envelope-center-2torsion-obstruction}，
\(
\dim_{\FF_2} Z(G_{\mathrm{SM}})[2]\le 2
\)。
而
\(
\mathcal C_{\mathrm{can}}\cong (\ZZ/2\ZZ)^3
\)
的 \(\FF_2\)-维数为 \(3\)。
故线性映射 \(f_*\) 的秩至多 \(2\)，其核维数至少 \(1\)，即核非平凡。证毕。
\end{proof}

\begin{proposition}[sheet parity 在 \(m=6\) 的严格锚点性]\label{prop:xi-window6-sheet-parity-anchor-m6}
\leanverified{paper\_xi\_window6\_sheet\_parity\_anchor\_m6}
在 dyadic boundary uplift 口径下，取
\[
\Delta_m^{\mathrm{bdry}}
:=\{n-V_m(w):\ n\in(\Fold_m^{\mathrm{bin}})^{-1}(w)\}
\]
（对 boundary \(w\) 不依赖于具体词）。
核验给出
\[
\Delta_6^{\mathrm{bdry}}=\{0,34\},\qquad
\Delta_7^{\mathrm{bdry}}=\{0,55,89\},\qquad
\Delta_8^{\mathrm{bdry}}=\{0,89,144\}.
\]
在“parity 由 uplift 数据无需额外选择即可恢复”的准则下，该结构仅在 \(m=6\) 成立。
\end{proposition}
\begin{proof}
若 parity 可由 \(\Delta_m^{\mathrm{bdry}}\) 唯一恢复，则每个 boundary 纤维必须存在唯一二元配对规则；
这要求非零 uplift 增量唯一。
当 \(m=6\) 时，\(\Delta_6^{\mathrm{bdry}}\) 仅含一个非零元 \(34\)，且 boundary 预像层数为 \(2\)，配对唯一。
当 \(m=7,8\) 时，\(\Delta_m^{\mathrm{bdry}}\) 各含两个不同非零元且预像层数为 \(3\)；
任意二元化都需额外选择（例如忽略某一层），不再由数据唯一决定。证毕。
\end{proof}

\begin{theorem}[纤维内自由 involution 的奇数障碍与最小覆盖度]\label{thm:xi-bdry-free-involution-odd-fiber-obstruction-min-cover}
\leanverified{paper\_xi\_bdry\_free\_involution\_odd\_fiber\_obstruction\_min\_cover}
设 \(f:\widetilde X\to X\) 为有限集合之间的满射。
若存在自由 involution \(\sigma:\widetilde X\to\widetilde X\) 满足
\[
\sigma^2=\mathrm{id},\qquad
\mathrm{Fix}(\sigma)=\varnothing,\qquad
f\circ\sigma=f,
\]
则对任意 \(x\in X\) 必有
\[
\bigl|f^{-1}(x)\bigr|\equiv 0\pmod 2.
\]
因此若存在 \(x\in X\) 使 \(|f^{-1}(x)|\) 为奇数，则不存在满足 \(f\circ\sigma=f\) 的自由 involution。
\par
进一步，若希望通过一个覆盖 \(\pi:Y\to\widetilde X\) 在 \(Y\) 上实现同类的纤维内自由 involution（即 \((f\circ\pi)\circ\widetilde\sigma=f\circ\pi\) 且 \(\widetilde\sigma\) 自由），
则覆盖度必须为偶数，因而 \(\deg(\pi)\ge 2\)。
\end{theorem}
\begin{proof}
若 \(\sigma\) 自由，则其轨道分解全由二元轨道组成，故 \(|S|\) 为偶数对任何 \(\sigma\)-不变有限子集 \(S\subseteq \widetilde X\) 成立。
取 \(S=f^{-1}(x)\)。由 \(f\circ\sigma=f\) 得 \(S\) 为 \(\sigma\)-不变，故 \(|f^{-1}(x)|\) 必为偶数；若存在奇数纤维则矛盾。
\par
对覆盖情形：若 \(|f^{-1}(x)|\) 为奇数，则 \(|(f\circ\pi)^{-1}(x)|=\deg(\pi)\cdot |f^{-1}(x)|\) 与 \(\deg(\pi)\) 同奇偶。
而在 \(Y\) 上实现纤维内自由 involution 仍要求每条纤维基数为偶数，故 \(\deg(\pi)\) 必为偶数，特别地 \(\deg(\pi)\ge 2\)。证毕。
\end{proof}

\begin{corollary}[boundary 三层 uplift 的二值外置下界]\label{cor:xi-bdry-uplift-three-layer-needs-one-bit}
在 dyadic boundary uplift 口径下，若 \(m\in\{7,8\}\)，则由命题 \ref{prop:xi-window6-sheet-parity-anchor-m6} 的
\(
\Delta_m^{\mathrm{bdry}}=\{0,\cdot,\cdot\}
\)
可知 boundary 预像层数为 \(3\)，从而存在 boundary 词 \(w\) 使
\[
\bigl|(\Fold_m^{\mathrm{bin}})^{-1}(w)\bigr|=3.
\]
因此不存在在 boundary uplift 纤维内闭合的自由 \(\ZZ_2\) involution；
并且任何试图在覆盖替代后实现该自由性的数据结构，其最小覆盖度满足 \(\deg(\pi)\ge 2\)，从而至少需要引入一个二值外置轴以消除奇数纤维的轨道障碍。
\end{corollary}
\begin{proof}
由命题 \ref{prop:xi-window6-sheet-parity-anchor-m6}，当 \(m=7,8\) 时 boundary 预像层数为 \(3\)，故存在边界纤维基数为 \(3\)。
将定理 \ref{thm:xi-bdry-free-involution-odd-fiber-obstruction-min-cover} 应用于 \(f=(\Fold_m^{\mathrm{bin}})|_{\widetilde X_m^{\mathrm{bdry}}}\) 即得结论。证毕。
\end{proof}

\begin{theorem}[全对称不变性强制的 \texorpdfstring{$\ZZ_2$}{Z/2} 分级平凡化]\label{thm:xi-sym-invariant-z2-grading-trivial}\leanverified{paper\_xi\_sym\_invariant\_z2\_grading\_trivial}
设 \(\Delta\) 为有限集合，\(|\Delta|=k\ge 3\)。
将一个 \(\ZZ_2\)-分级理解为选取“奇层子集” \(A\subset\Delta\)（允许 \(A=\varnothing\) 或 \(A=\Delta\) 的平凡分级）。
若要求该分级在全对称群 \(\mathrm{Sym}(\Delta)\cong S_k\) 的自然作用下不变，即对任意 \(\sigma\in S_k\) 都有 \(\sigma(A)=A\)，则该分级必为平凡分级：\(A=\varnothing\) 或 \(A=\Delta\)。
\end{theorem}
\begin{proof}
由于 \(S_k\) 在 \(\Delta\) 上传递，若 \(A\neq\varnothing\)，取 \(x\in A\)。
对任意 \(y\in\Delta\) 存在 \(\sigma\in S_k\) 使 \(\sigma(x)=y\)。
由不变性 \(\sigma(A)=A\) 得 \(y\in A\)，从而 \(A=\Delta\)。
同理，若 \(A\neq\Delta\) 则只能 \(A=\varnothing\)。
故除平凡分级外不存在 \(S_k\)-不变的 \(\ZZ_2\) 分级。证毕。
\end{proof}

\begin{theorem}[左--右交换下的规范中心模与不可下降通道]\label{thm:xi-window6-lr-swap-canonical-center-mod}
在 \(\mathcal C_{\mathrm{can}}\cong(\ZZ/2\ZZ)^3\) 中取基
\[
a:=(-I_4,1,1),\qquad b:=(1,-I_2,1),\qquad c:=(1,1,-I_2).
\]
令 \(\iota\) 为交换 \(SU(2)_L\leftrightarrow SU(2)_R\) 的自同构，则
\[
\iota(a)=a,\qquad \iota(b)=c,\qquad \iota(c)=b.
\]
定义满同态
\[
\pi_{\mathrm{LR}}:\mathcal C_{\mathrm{can}}\twoheadrightarrow \ZZ/2\ZZ,\qquad
\pi_{\mathrm{LR}}(\alpha a+\beta b+\gamma c):=\beta+\gamma.
\]
则
\[
\ker(\pi_{\mathrm{LR}})=\mathcal C_{\mathrm{can}}^{\iota}
\cong (\ZZ/2\ZZ)^2.
\]
因而任何要求“左--右交换在中心层不可见”的下降机制，必然至少压缩一个规范 \(\ZZ_2\) 模通道。
\end{theorem}
\begin{proof}
\(\pi_{\mathrm{LR}}\) 显然为群同态且满。
\[
\pi_{\mathrm{LR}}(\alpha,\beta,\gamma)=0
\iff \beta=\gamma.
\]
另一方面，
\[
\iota(\alpha,\beta,\gamma)=(\alpha,\gamma,\beta),
\]
故
\(
(\alpha,\beta,\gamma)\in\mathcal C_{\mathrm{can}}^{\iota}
\iff \beta=\gamma
\)。
两条件等价，故核恰为不动子群，且其维数为 \(2\)。证毕。
\end{proof}

\begin{conjecture}[dyadic boundary uplift 层数在 \(m=6\) 的唯一退化]\label{conj:xi-window6-bdry-uplift-unique-degeneration}
设 \(\Delta_m^{\mathrm{bdry}}\) 为 boundary uplift 的增量集合。猜想对所有 \(m\ge 7\) 有
\[
\Delta_m^{\mathrm{bdry}}=\{0,F_{m+3},F_{m+4}\},
\]
从而 boundary 覆盖层数恒为 \(3\)。
而 \(m=6\) 为唯一退化分辨率：
\[
\Delta_6^{\mathrm{bdry}}=\{0,F_9\}=\{0,34\},
\]
即唯一仅含单一非零增量的二层情形。
\end{conjecture}

\begin{conjecture}[异常与共振的通道稀疏性：由 boundary sheet parity 子群支配]\label{conj:xi-anomaly-resonance-channel-sparsity-controlled-by-boundary-parity}
在 window-\(6\) 的 boundary 读出中，三轴独立 sheet-flip 诱导规范中心寄存器
\(\mathcal C_{\mathrm{bdry}}\cong(\ZZ/2\ZZ)^3\)。
猜想：对所有足够大的分辨率 \(m\)，所有可见通道中的非平凡绕数/真共振指数（从而异常的可见非零性）只能出现在那些在 \(\mathcal C_{\mathrm{bdry}}\) 上取非平凡限制的特征标通道中。
等价地：若某一可见通道 \(\overline\chi\) 在 \(\mathcal C_{\mathrm{bdry}}\) 上平凡，则其可见散射绕数与相应的 Toeplitz 负平方指数均为零。
\par
该稀疏性一旦成立，将把“通道级异常定位/证书化”的搜索空间压缩到由 \(\mathcal C_{\mathrm{bdry}}\) 诱导的有限层筛选子族，并与定理 \ref{thm:xi-anom-winding-obstruction} 与定理 \ref{thm:xi-visible-anom-iff-radial-wallcrossing} 的可见通道判别口径形成闭合接口。
\end{conjecture}

\begin{remark}[可复现核验脚本]\label{rem:xi-window6-center-sheetflip-ps-sm-audit-script}
本段代数与组合核验由脚本
\texttt{scripts/exp\_xi\_window6\_center\_sheetflip\_ps\_sm\_audit.py}
给出，证书文件为
\texttt{artifacts/export/xi\_window6\_center\_sheetflip\_ps\_sm\_audit.json}。
\end{remark}


\begin{corollary}[Lee--Yang 四层覆盖中的可交换外置障碍]\label{cor:xi-terminal-abelian-register-obstruction-leyang-s4}
在定理 \ref{thm:xi-terminal-zm-leyang-elliptic-structure} 给出的 Lee--Yang 四层覆盖 \(\pi_y:E_{\min}\to\mathbb P^1_y\) 上，
\[
\mathrm{Mon}(\pi_y)\cong S_4
\]
为非交换群，故满足定理 \ref{thm:xi-terminal-abelian-register-obstruction-nonabelian-monodromy} 的不可完备性结论：任何要求路径拼接同态化的可交换外置寄存器都不能实现全局单值恢复。
\end{corollary}

\begin{theorem}[位移半直积 Gödel 寄存器对非交换算法轨迹的忠实外置]\label{thm:xi-terminal-semidir-godel-register-faithful}
\leanverified{paper\_xi\_terminal\_semidir\_godel\_register\_faithful}
设 \(\Sigma\) 为有限字母表（\(|\Sigma|\ge 2\)），\(\Sigma^\ast\) 为自由幺半群（串联为乘法），并取单射编码
\(
c:\Sigma\hookrightarrow \NN_{\ge 1}
\)。
记 \((p_i)_{i\ge 1}\) 为递增素数序列。

对 \(w=a_1\cdots a_n\in\Sigma^\ast\) 定义位序 Gödel 编码
\[
G(w):=\prod_{i=1}^{n}p_i^{\,c(a_i)},
\qquad
|w|=n.
\]
定义素数位移同态
\[
\mathrm{Sh}^{k}\!\left(\prod_{i\ge1}p_i^{e_i}\right):=\prod_{i\ge1}p_{i+k}^{e_i}.
\]
令
\[
\mathcal R:=\NN\times \NN_{>0},
\qquad
(n,A)\odot(m,B):=\bigl(n+m,\ A\cdot \mathrm{Sh}^{n}(B)\bigr).
\]
则映射
\[
J:\Sigma^\ast\to\mathcal R,\qquad
J(w):=(|w|,G(w))
\]
为单射幺半群同态，且 \((|w|,G(w))\) 可唯一恢复 \(w\)。
此外，对任意交换幺半群 \(M\)，不存在单射幺半群同态 \(\Sigma^\ast\hookrightarrow M\)。
\end{theorem}
\begin{proof}
\(\mathrm{Sh}^{n}\circ\mathrm{Sh}^{m}=\mathrm{Sh}^{n+m}\) 且每个 \(\mathrm{Sh}^n\) 为幺半群同态，故 \(\odot\) 结合，单位元为 \((0,1)\)。
对 \(u=a_1\cdots a_n\)、\(v=b_1\cdots b_m\)，
\[
G(uv)
=\left(\prod_{i=1}^{n}p_i^{c(a_i)}\right)\left(\prod_{j=1}^{m}p_{n+j}^{c(b_j)}\right)
=G(u)\,\mathrm{Sh}^{n}(G(v)),
\]
故 \(J(uv)=J(u)\odot J(v)\)。

单射性由唯一分解定理给出：\(J(w)=(n,G(w))\) 已给出长度 \(n\)，且
\(
v_{p_i}(G(w))=c(a_i)
\)
可逐位恢复 \(a_i\)（因 \(c\) 单射）。

若 \(f:\Sigma^\ast\to M\) 为到交换幺半群的同态，取 \(a\neq b\in\Sigma\)，则
\[
f(ab)=f(a)f(b)=f(b)f(a)=f(ba).
\]
若 \(f\) 单射则推出 \(ab=ba\)，与自由幺半群非交换性矛盾。证毕。
\end{proof}

% AUTO-GENERATED by scripts/exp_xi_prime_register_semidirect_two_shadow_audit.py
\begin{equation}\label{eq:xi_prime_register_semidirect_two_shadow_audit}
\begin{aligned}
\rho([\alpha,\beta])&=(2,3,1,4),\qquad \rho([\alpha,\beta])(1)\neq 1,\\
J(uv)&=J(u)\odot J(v)\ \text{(checked for all words with }|u|,|v|\le 4\text{)},\\
\left|O_2(\ZZ)\right|&=8,\qquad \left|SO_2(\ZZ)\right|=4.
\end{aligned}
\end{equation}


\begin{theorem}[Lee--Yang--Gödel 嵌入与素数寄存器的代数可识别性]\label{thm:xi-leyang-godel-prime-register-algebraic-identifiability}
令 \(\mathcal P\) 为素数集。记
\[
\NN^{(\mathcal P)}:=\{r:\mathcal P\to\NN:\ \#\mathrm{supp}(r)<\infty\},
\qquad
\mathrm{supp}(r):=\{p\in\mathcal P:\ r(p)\neq 0\}.
\]
并定义 Gödel 读数
\[
N(r):=\prod_{p\in\mathcal P}p^{r(p)}\in\NN_{\ge 1}.
\]
定义映射
\[
\mathcal L:\NN^{(\mathcal P)}\to \ZZ[z],
\qquad
\mathcal L(r)(z):=\prod_{p\in\mathcal P}\bigl(1+z^p\bigr)^{r(p)}.
\]
则成立：
\begin{enumerate}
  \item \(\mathcal L\) 为交换幺半群同态（以点加 \(r+s\) 与多项式乘法为幺半群运算）：
  \[
  \mathcal L(r+s)=\mathcal L(r)\mathcal L(s).
  \]
  \item \(\mathcal L\) 为单射。
  \item \(\mathcal L(r)\) 的全部零点均落在单位圆 \(S^1\)，且皆为偶阶单位根。
\end{enumerate}
\end{theorem}
\begin{proof}
第一条由指数加法与乘法分配律直接给出。
\par
第三条：任一因子 \(1+z^p\) 的零点满足 \(z^p=-1\)，从而 \(z^{2p}=1\)，故其零点均为偶阶单位根，亦在 \(S^1\) 上；乘积的零点集为各因子零点的并（按重数计），结论成立。
\par
第二条：对奇素数 \(p\) 有分解
\[
z^p+1=(z+1)\Phi_{2p}(z),
\]
且 \(\Phi_{2p}\) 在 \(\QQ\) 上不可约，并且不同 \(p\) 给出不同的不可约因子（可参见 \cite{Serre1973CourseArithmetic} 对 cyclotomic 多项式的基本性质与不可约性论述）。
对 \(p=2\)，有 \(z^2+1=\Phi_4(z)\)。
因此在 \(\mathcal L(r)\) 的不可约分解中，每个 \(\Phi_{2p}\)（\(p\) 为奇素数）出现的指数恰为 \(r(p)\)，而 \(\Phi_4\) 出现的指数为 \(r(2)\)，从而 \(r\) 由 \(\mathcal L(r)\) 唯一确定，故 \(\mathcal L\) 单射。
\end{proof}

\begin{theorem}[Gödel--Lee--Yang 对数振幅的逐素频全息反演]\label{thm:xi-leyang-logamplitude-prime-frequency-inversion}
在定理 \ref{thm:xi-leyang-godel-prime-register-algebraic-identifiability} 的记号下，对任意 \(r\in\NN^{(\mathcal P)}\) 定义单位圆上的对数振幅
\[
\mathcal A_r(\theta):=\log\abs{\mathcal L(r)(e^{i\theta})}\qquad(\theta\in[-\pi,\pi]).
\]
则 \(\mathcal A_r\in L^1([-\pi,\pi])\)。
并且对每个素数 \(p\) 成立严格反演公式
\[
\boxed{\quad
r(p)=\frac1\pi\int_{-\pi}^{\pi}\mathcal A_r(\theta)\cos(p\theta)\,\dd\theta.\quad}
\]
因此 \(\{\mathcal A_r(\theta)\}_{\theta\in[-\pi,\pi]}\) 在测度意义下逐素数唯一确定寄存器 \(r\)。
\end{theorem}
\begin{proof}
由于 \(r\) 仅有限支撑，
\[
\mathcal A_r(\theta)=\sum_{p\in\mathcal P}r(p)\log\abs{1+e^{ip\theta}}.
\]
对任意 \(\rho\in(0,1)\)，在单位圆盘内取主值对数并取实部，得到一致可积展开
\[
\Re\log(1+\rho e^{ix})
=\sum_{k\ge 1}\frac{(-1)^{k+1}}{k}\rho^k\cos(kx).
\]
令 \(\rho\uparrow 1\) 可得在 \(L^2([-\pi,\pi])\)（从而 \(L^1\)）意义下的余弦级数
\[
\log\abs{1+e^{ix}}
=\sum_{k\ge 1}\frac{(-1)^{k+1}}{k}\cos(kx)
\qquad(x\in[-\pi,\pi]\text{ a.e.}).
\]
代入 \(x=p\theta\) 并对有限个 \(p\) 求和（可交换求和次序）得到
\[
\mathcal A_r(\theta)
=\sum_{p}r(p)\sum_{k\ge 1}\frac{(-1)^{k+1}}{k}\cos(kp\theta).
\]
因此 \(\mathcal A_r\) 的第 \(n\) 个余弦系数为
\[
\frac1\pi\int_{-\pi}^{\pi}\mathcal A_r(\theta)\cos(n\theta)\,\dd\theta
=\sum_{p\mid n} r(p)\,\frac{(-1)^{n/p+1}}{n/p}.
\]
当 \(n=p\) 为素数时右端仅含 \(p\) 一项且 \(n/p=1\)，从而该余弦系数等于 \(r(p)\)，即得所述反演公式。证毕。
\end{proof}

\begin{corollary}[整数刚性阈值：\(L^2\) 误差控制下的逐素数精确恢复]\label{cor:xi-leyang-logamplitude-rounding-threshold}
在定理 \ref{thm:xi-leyang-logamplitude-prime-frequency-inversion} 的设定下，令 \(\widetilde{\mathcal A}\in L^2([-\pi,\pi])\) 为任意观测/学习得到的对数振幅近似，并对素数 \(p\) 定义
\[
\widetilde r(p):=\operatorname{Round}\!\left(\frac1\pi\int_{-\pi}^{\pi}\widetilde{\mathcal A}(\theta)\cos(p\theta)\,\dd\theta\right)\in\ZZ,
\]
其中 \(\operatorname{Round}\) 表示四舍五入到最近整数。
若满足误差阈值
\[
\norm{\widetilde{\mathcal A}-\mathcal A_r}_{L^2([-\pi,\pi])}<\frac{\sqrt\pi}{2},
\]
则对所有素数 \(p\) 均有 \(\widetilde r(p)=r(p)\)。
\end{corollary}
\begin{proof}
记
\[
c_p(F):=\frac1\pi\int_{-\pi}^{\pi}F(\theta)\cos(p\theta)\,\dd\theta.
\]
由 Cauchy--Schwarz 不等式并注意 \(\norm{\cos(p\theta)}_{L^2([-\pi,\pi])}=\sqrt\pi\)，得到
\[
\abs{c_p(\widetilde{\mathcal A})-c_p(\mathcal A_r)}
\le \frac1\pi\norm{\widetilde{\mathcal A}-\mathcal A_r}_{L^2}\,\norm{\cos(p\theta)}_{L^2}
=\frac1{\sqrt\pi}\norm{\widetilde{\mathcal A}-\mathcal A_r}_{L^2}
<\frac12.
\]
而由定理 \ref{thm:xi-leyang-logamplitude-prime-frequency-inversion}，\(c_p(\mathcal A_r)=r(p)\in\ZZ\)。
因此 \(c_p(\widetilde{\mathcal A})\) 与整数 \(r(p)\) 的距离严格小于 \(1/2\)，四舍五入必回到 \(r(p)\)。证毕。
\end{proof}

\begin{definition}[素数根均匀测度与零点经验测度]\label{def:xi-leyang-prime-root-measures}
对每个素数 \(p\)，令 \(\eta_p\) 为方程 \(z^p=-1\) 的 \(p\) 个解上的均匀概率测度：
\[
\eta_p:=\frac1p\sum_{k=0}^{p-1}\delta_{\exp\!\left(i\frac{(2k+1)\pi}{p}\right)}
\quad\text{（概率测度，支撑于 \(S^1\)）}.
\]
对任意 \(r\in\NN^{(\mathcal P)}\setminus\{0\}\)，记
\[
D(r):=\deg \mathcal L(r)=\sum_{p\in\mathcal P}r(p)\,p,
\]
并定义其零点归一化经验测度（按重数计）
\[
\mu_r:=\frac1{D(r)}\sum_{\zeta:\ \mathcal L(r)(\zeta)=0}\mathrm{mult}(\zeta)\,\delta_\zeta
\quad\text{（概率测度，支撑于 \(S^1\)）}.
\]
\end{definition}

\begin{lemma}[Fourier 系数的显式消失律]\label{lem:xi-eta-p-fourier}
\leanverified{paper\_xi\_eta\_p\_fourier}
对任意整数 \(n\neq 0\)，有
\[
\widehat{\eta_p}(n):=\int_{S^1}z^n\,\dd\eta_p(z)=
\begin{cases}
(-1)^{n/p},&p\mid n,\\
0,&p\nmid n.
\end{cases}
\]
\end{lemma}
\begin{proof}
\[
\int_{S^1}z^n\,\dd\eta_p(z)
=\frac1p\sum_{k=0}^{p-1}\exp\!\left(i n\frac{(2k+1)\pi}{p}\right)
=e^{i n\pi/p}\cdot \frac1p\sum_{k=0}^{p-1}e^{i2\pi nk/p}.
\]
若 \(p\nmid n\)，则内层为非平凡 \(p\) 次单位根之和，取值 \(0\)；若 \(p\mid n\)，则内层为 \(1\)，从而整体为 \(e^{i\pi(n/p)}=(-1)^{n/p}\)。
\end{proof}

\begin{proposition}[零点测度的素数分解与单频谱解码]\label{prop:xi-leyang-prime-measure-linear-independence-decoding}
对任意 \(r\in\NN^{(\mathcal P)}\setminus\{0\}\)，有凸组合分解
\[
\mu_r=\sum_{p\in\mathcal P}\frac{r(p)\,p}{D(r)}\,\eta_p.
\]
进一步，对任意有限支撑整数族 \((c_p)_{p\in\mathcal P}\subset\ZZ\) 与带符号测度
\(
\sigma:=\sum_pc_p\,\eta_p
\)，若 \(\sigma=0\) 则必有 \(c_p=0\)（对所有 \(p\)）。
并且对任意素数 \(p\)，有解码公式
\[
c_p=-\int_{S^1}z^p\,\dd\sigma(z).
\]
\end{proposition}
\begin{proof}
由定义 \(\mathcal L(r)=\prod_p(1+z^p)^{r(p)}\)，每个因子 \(1+z^p\) 贡献 \(p\) 个单位根零点，且零点重数按指数 \(r(p)\) 线性叠加，因此零点测度满足
\(
\mu_r=\sum_p\frac{r(p)p}{D(r)}\eta_p
\)。
\par
对线性无关性与解码，取 Fourier 矩 \(n=p\)。由引理 \ref{lem:xi-eta-p-fourier}，对 \(q\neq p\) 因 \(q\nmid p\) 有
\(
\int z^p\,\dd\eta_q=0
\)，而对 \(q=p\) 有
\(
\int z^p\,\dd\eta_p=-1
\)。
故
\[
\int_{S^1}z^p\,\dd\sigma(z)=\sum_q c_q\int z^p\,\dd\eta_q(z)=-c_p,
\]
从而若 \(\sigma=0\) 则 \(c_p=0\)（对任意 \(p\)），并得解码公式。
\end{proof}

\begin{theorem}[素数截断零点谱的 Haar 极限与显式差分界]\label{thm:xi-leyang-prime-truncation-haar-discrepancy}
令
\[
\mathcal L_{\le P}(z):=\prod_{p\le P}(1+z^p),
\qquad
D_P:=\deg\mathcal L_{\le P}=\sum_{p\le P}p,
\]
并令 \(\mu_P\) 为 \(\mathcal L_{\le P}\) 的零点归一化经验测度（按重数计）。
则当 \(P\to\infty\) 时，\(\mu_P\) 在弱-\(\ast\) 意义下收敛到单位圆上的 Haar 概率测度 \(m_{S^1}\)。
此外，存在绝对常数 \(C>0\)，使对任意弧段 \(A\subset S^1\) 有
\[
\sup_{A\subset S^1}\abs{\mu_P(A)-m_{S^1}(A)}
\le \frac{C}{\sqrt{D_P}}
=\frac{C}{\sqrt{\sum_{p\le P}p}}.
\]
\end{theorem}
\begin{proof}
由命题 \ref{prop:xi-leyang-prime-measure-linear-independence-decoding} 的分解（取 \(r(p)=\ind_{\{p\le P\}}\)），有
\[
\mu_P=\sum_{p\le P}\frac{p}{D_P}\,\eta_p.
\]
对任意整数 \(n\neq 0\)，其 Fourier 系数为
\[
\widehat{\mu_P}(n)
=\sum_{p\le P}\frac{p}{D_P}\widehat{\eta_p}(n)
=\frac1{D_P}\sum_{\substack{p\le P\\ p\mid n}}p\,(-1)^{n/p},
\]
其中最后一步用引理 \ref{lem:xi-eta-p-fourier}。
对固定 \(n\neq 0\)，分子仅为有限和且与 \(P\) 无关，而 \(D_P\to\infty\)，故 \(\widehat{\mu_P}(n)\to 0\)。
Haar 测度 \(m_{S^1}\) 的所有非零 Fourier 系数均为 \(0\)，因而 \(\mu_P\Rightarrow m_{S^1}\)。
\par
差分界由 Erd\H{o}s--Tur\'an 不等式给出（参见 \cite{KuipersNiederreiter1974}）：存在绝对常数 \(C_0>0\)，对任意 \(H\in\NN\) 有
\[
\sup_A\abs{\mu_P(A)-m_{S^1}(A)}
\le C_0\left(\frac1H+\sum_{k=1}^H\frac{|\widehat{\mu_P}(k)|}{k}\right).
\]
由上式表达与 \(\sum_{p\mid k}p\le k\) 得
\(
|\widehat{\mu_P}(k)|\le k/D_P
\)，从而
\[
\sum_{k=1}^H\frac{|\widehat{\mu_P}(k)|}{k}\le \sum_{k=1}^H\frac1{D_P}=\frac{H}{D_P}.
\]
取 \(H=\lfloor\sqrt{D_P}\rfloor\) 得
\(
\sup_A\abs{\mu_P(A)-m_{S^1}(A)}\le 2C_0/\sqrt{D_P}
\)，并将常数并入 \(C\) 即得。
\end{proof}

\begin{proposition}[Gödel 椭圆边界的推前测度与显式密度]\label{prop:xi-godel-ellipse-pushforward-density}
取 \(a,b>0\)，并定义单位圆到轴对齐椭圆边界的参数化
\[
T_{a,b}:S^1\to \partial E_{a,b},
\qquad
T_{a,b}(e^{it})=(a\cos t,\ b\sin t),
\]
其中 \(\partial E_{a,b}\) 为 \(x^2/a^2+y^2/b^2=1\) 的边界。
令
\(
\nu_{a,b}:=(T_{a,b})_*m_{S^1}
\)
为 Haar 测度的推前测度，则 \(\nu_{a,b}\) 关于椭圆弧长测度 \(ds\) 绝对连续，且 Radon--Nikodym 导数满足
\[
\frac{\dd\nu_{a,b}}{\dd s}\Bigl(T_{a,b}(e^{it})\Bigr)
=\frac1{2\pi}\cdot\frac1{\sqrt{a^2\sin^2 t+b^2\cos^2 t}}.
\]
进一步，若 \(\mu_P\) 为定理 \ref{thm:xi-leyang-prime-truncation-haar-discrepancy} 中的零点谱测度，则
\[
(T_{a,b})_*\mu_P\Rightarrow \nu_{a,b}.
\]
\end{proposition}
\begin{proof}
由弧长元素
\(
ds=\sqrt{(a\sin t)^2+(b\cos t)^2}\,\dd t
\)
可知 \(m_{S^1}\) 在参数 \(t\) 下的密度为 \(\dd t/(2\pi)\)。推前到 \(\partial E_{a,b}\) 后得到关于 \(ds\) 的密度为
\[
\frac{\dd t}{2\pi}\cdot \frac1{\dd s}
=\frac1{2\pi}\cdot\frac1{\sqrt{a^2\sin^2 t+b^2\cos^2 t}}.
\]
弱收敛在连续映射下的推前稳定性给出第二条：由 \(\mu_P\Rightarrow m_{S^1}\) 且 \(T_{a,b}\) 连续，得 \((T_{a,b})_*\mu_P\Rightarrow (T_{a,b})_*m_{S^1}=\nu_{a,b}\)。
\end{proof}

\begin{theorem}[素数寄存器序列的 Haar 光谱收敛判别]\label{thm:xi-leyang-prime-register-haar-criterion}
设 \((r_m)_{m\ge 1}\subset \NN^{(\mathcal P)}\setminus\{0\}\) 为任意序列，并令
\[
\mathcal L_m(z):=\mathcal L(r_m)(z),
\qquad
D_m:=D(r_m)=\sum_{p\in\mathcal P}r_m(p)\,p,
\qquad
\mu_m:=\mu_{r_m}.
\]
则以下三条等价：
\begin{enumerate}
  \item \(\mu_m\Rightarrow m_{S^1}\)。
  \item 对每个固定素数 \(q\)，有
  \[
  \frac{r_m(q)\,q}{D_m}\longrightarrow 0.
  \]
  \item 对每个固定整数 \(n\neq 0\)，有
  \[
  \frac{\sum_{p\mid n}r_m(p)\,p}{D_m}\longrightarrow 0.
  \]
\end{enumerate}
\end{theorem}
\begin{proof}
由定义 \(\mu_m=\sum_p \frac{r_m(p)p}{D_m}\eta_p\) 与引理 \ref{lem:xi-eta-p-fourier}，对任意整数 \(n\neq 0\) 有
\[
\widehat{\mu_m}(n)
=\int z^n\,\dd\mu_m(z)
=\frac1{D_m}\sum_{p\mid n}r_m(p)\,p\,(-1)^{n/p}.
\]
Haar 测度满足 \(\widehat{m_{S^1}}(n)=0\)（\(n\neq 0\)），故 \(\mu_m\Rightarrow m_{S^1}\) 等价于对一切 \(n\neq 0\) 有 \(\widehat{\mu_m}(n)\to 0\)；
而这与第三条等价（符号项不影响绝对值阶）。
第三条显然推出第二条（取 \(n=q\)）。
第二条推出第三条是因为 \(n\) 的素因子集合有限，\(\sum_{p\mid n}r_m(p)p/D_m\) 为有限项之和，逐项趋于 \(0\) 即整体趋于 \(0\)。
\end{proof}

\begin{theorem}[零点普适律与对数振幅保持律的严格分裂]\label{thm:xi-leyang-bulk-boundary-splitting-haar-vs-logamplitude}
设 \((r_N)_{N\ge 1}\subset \NN^{(\mathcal P)}\setminus\{0\}\)，并令 \(\mu_{r_N}\) 与
\[
D_N:=D(r_N)=\sum_{p\in\mathcal P} r_N(p)\,p
\]
如定义 \ref{def:xi-leyang-prime-root-measures} 所示。
假设 \(D_N\to\infty\)，且对每个固定素数 \(p\)，序列 \((r_N(p))_{N}\) 有界。
则：
\begin{enumerate}
  \item 零点经验测度满足 \(\mu_{r_N}\Rightarrow m_{S^1}\)（单位圆上的 Haar 概率测度）；
  \item 与此同时，若记 \(\mathcal A_{r_N}(\theta)=\log|\mathcal L(r_N)(e^{i\theta})|\)，则对每个素数 \(p\) 恒有
  \[
  \frac1\pi\int_{-\pi}^{\pi}\mathcal A_{r_N}(\theta)\cos(p\theta)\,\dd\theta=r_N(p),
  \]
  因而 \((\mathcal A_{r_N})_N\) 不可能在 \(L^2([-\pi,\pi])\) 中收敛到某个与 \(r_N\) 无关的“普适极限”（除非 \(r_N(p)\) 在每个素数处逐点稳定）。
\end{enumerate}
\end{theorem}
\begin{proof}
对任意固定素数 \(q\)，由有界性得 \(r_N(q)\,q=O_q(1)\)，而 \(D_N\to\infty\)，故
\(
\frac{r_N(q)\,q}{D_N}\to 0
\)。
由定理 \ref{thm:xi-leyang-prime-register-haar-criterion} 即得 \(\mu_{r_N}\Rightarrow m_{S^1}\)，从而得到 (1)。
\par
(2) 由定理 \ref{thm:xi-leyang-logamplitude-prime-frequency-inversion} 对每个 \(N\) 与每个素数 \(p\) 的反演恒等式直接给出。
若 \(\mathcal A_{r_N}\to \mathcal A\) 于 \(L^2\)，则其每个余弦系数必收敛，因而 \(r_N(p)\) 必对每个 \(p\) 收敛；若要极限不依赖于序列 \((r_N)\)，则只能在逐素数意义下稳定。证毕。
\end{proof}

\begin{conjecture}[外置诱导素数账本的质量不凝聚与普适差分阶]\label{conj:xi-leyang-prime-register-universal-nonconcentration}
在本文折叠/投影外置范式中，由算法轨迹自然诱导的素数寄存器序列 \((r_m)\) 满足定理 \ref{thm:xi-leyang-prime-register-haar-criterion} 的质量不凝聚条件，从而其 Lee--Yang--Gödel 多项式零点谱 \(\mu_{r_m}\) 弱收敛到 \(m_{S^1}\)。
并且在适当的规范化与均匀性假设下，弧段差分满足 \(\sup_A|\mu_{r_m}(A)-m_{S^1}(A)|=O(D(r_m)^{-1/2})\)，且该指数阶在该范式内具有普适性。
\end{conjecture}

\paragraph{Zeckendorf--Gödel 素数码集合的密度、Lee--Yang 实根结构与谱行列式重整}
设 \((p_k)_{k\ge 1}\) 为递增素数序列，并定义
\[
\mathrm{ZG}
:=\Bigl\{
n\in\NN:\ \forall p,\ v_p(n)\in\{0,1\},\ \forall k\ge 1,\ v_{p_k}(n)\,v_{p_{k+1}}(n)=0
\Bigr\}.
\]
即 \(n\) 为平方因子自由，且其素因子索引集合不含相邻整数。

\begin{theorem}[Gödel--Zeckendorf 素数约束 Dirichlet 级数的对数重整化与半临界解析性]\label{thm:xi-zg-dirichlet-log-renormalization}
\leanverified{paper\_xi\_zg\_dirichlet\_log\_renormalization}
令
\[
\Omega:=\Bigl\{\varepsilon=(\varepsilon_k)_{k\ge 1}:\ \varepsilon_k\in\{0,1\},\ \varepsilon_k\varepsilon_{k+1}=0,\ \varepsilon\ \text{仅有限支撑}\Bigr\}.
\]
对 \(\Re(s)>1\) 定义 Dirichlet 级数
\[
F(s):=\sum_{n\in \mathrm{ZG}}n^{-s}
=\sum_{\varepsilon\in\Omega}\ \prod_{k\ge 1}p_k^{-s\varepsilon_k}.
\]
并对 \(N\ge 0\) 定义截断
\[
F_N(s):=\sum_{\substack{\varepsilon\in\{0,1\}^N\\ \varepsilon_k\varepsilon_{k+1}=0}}\ \prod_{k=1}^{N}p_k^{-s\varepsilon_k},
\qquad
F_0(s)=1,\ \ F_1(s)=1+2^{-s}.
\]
令 \(r_N(s):=F_N(s)/F_{N-1}(s)\)（\(N\ge 1\)），并令
\[
\Delta_N(s):=\log r_N(s)-p_N^{-s}\qquad(N\ge 2),
\]
其中对数取为与实轴上 \(s>1\) 的正值一致的解析分支。
则：
\begin{enumerate}
  \item 对所有 \(\Re(s)>0\) 成立递推恒等式
  \[
  F_N(s)=F_{N-1}(s)+p_N^{-s}F_{N-2}(s),\qquad
  r_N(s)=1+\frac{p_N^{-s}}{r_{N-1}(s)}.
  \]
  \item 级数
  \[
  G(s):=\sum_{N\ge 2}\Delta_N(s)
  \]
  在每个半平面 \(\Re(s)\ge \sigma_0>\tfrac12\) 上绝对且一致收敛，从而在 \(\Re(s)>\tfrac12\) 上给出全纯函数。
  \item 令素数 \(\zeta\) 函数（prime zeta）
  \(
  P(s):=\sum_{k\ge 1}p_k^{-s}
  \)
  （在 \(\Re(s)>1\) 绝对收敛）。
  则对 \(\Re(s)>1\) 有严格分解
  \[
  \log F(s)=P(s)+G(s),
  \qquad\text{从而}\qquad
  F(s)=\exp\!\bigl(P(s)+G(s)\bigr).
  \]
\end{enumerate}
\end{theorem}

\begin{proof}
第一条由末位 \(\varepsilon_N\in\{0,1\}\) 分拆：\(\varepsilon_N=0\) 贡献为 \(F_{N-1}\)，而 \(\varepsilon_N=1\) 强制 \(\varepsilon_{N-1}=0\) 并贡献 \(p_N^{-s}F_{N-2}\)。
对 \(r_N=F_N/F_{N-1}\) 取比即得递推。
\par
对第二条，令 \(\sigma=\Re(s)\ge \sigma_0>\tfrac12\)。
由递推 \(r_N=1+p_N^{-s}/r_{N-1}\) 与 \(|p_N^{-s}|\le p_N^{-\sigma}\) 可在每个紧子域上固定 \(r_{N-1}(s)\) 远离 \(0\) 的一致下界，
从而存在常数 \(C(\sigma_0)\) 使
\[
|\Delta_N(s)|
\le
C(\sigma_0)\Bigl(p_N^{-2\sigma}+(p_Np_{N-1})^{-\sigma}\Bigr)
\qquad(\sigma\ge\sigma_0).
\]
由于
\(
\sum_N p_N^{-2\sigma}
\)
与
\(
\sum_N(p_Np_{N-1})^{-\sigma}
\)
在 \(\sigma>\tfrac12\) 收敛，故 \(\sum_{N\ge 2}\Delta_N(s)\) 在 \(\Re(s)\ge\sigma_0\) 上绝对一致收敛，从而 \(G\) 在 \(\Re(s)>\tfrac12\) 全纯。
\par
第三条由
\(
\log F_N(s)=\sum_{k=1}^N\log r_k(s)
\)
与 \(\Delta_k=\log r_k-p_k^{-s}\) 得
\[
\log F_N(s)=\sum_{k=1}^{N}p_k^{-s}+\sum_{k=2}^{N}\Delta_k(s).
\]
令 \(N\to\infty\) 并在 \(\Re(s)>1\) 用绝对收敛交换极限，得到 \(\log F(s)=P(s)+G(s)\)。证毕。
\end{proof}

\begin{theorem}[平方自由 Euler 因子分离与 \texorpdfstring{$\zeta$}{zeta}-因子化]\label{thm:xi-zg-zeta-factorization}
\leanverified{paper\_xi\_zg\_zeta\_factorization}
沿用定理 \ref{thm:xi-zg-dirichlet-log-renormalization} 的记号，并令
\[
B_N(s):=\prod_{k=1}^{N}\bigl(1+p_k^{-s}\bigr),\qquad
H_N(s):=\frac{F_N(s)}{B_N(s)}.
\]
则存在唯一函数 \(H\) 满足：
\begin{enumerate}
  \item \(H_N(s)\to H(s)\) 在每个紧集 \(K\subset\{s:\Re(s)>\tfrac12\}\) 上一致收敛；
  \item \(H(s)\) 在半平面 \(\Re(s)>\tfrac12\) 上全纯且处处非零；
  \item 在 \(\Re(s)>1\) 上成立严格因子化
  \[
  F(s)=H(s)\frac{\zeta(s)}{\zeta(2s)}.
  \]
\end{enumerate}
从而 \(F\) 在 \(\Re(s)>\tfrac12\) 上亚纯，且在该半平面内唯一可能的极点位于 \(s=1\)，并且为一阶极点。
\end{theorem}

\begin{proof}
由递推可写 \(F_N(s)=\prod_{k=1}^{N}r_k(s)\)，从而
\[
\log H_N(s)
=\sum_{k=1}^{N}\Bigl(\log r_k(s)-\log\bigl(1+p_k^{-s}\bigr)\Bigr).
\]
对 \(k\ge 2\) 使用 \(\Delta_k(s)=\log r_k(s)-p_k^{-s}\)，并将上式改写为
\[
\log H_N(s)
=\Bigl(\log r_1(s)-\log(1+p_1^{-s})\Bigr)
\sum_{k=2}^{N}\Delta_k(s)
\sum_{k=2}^{N}\Bigl(p_k^{-s}-\log(1+p_k^{-s})\Bigr).
\]
第一项恒为 \(0\)。
由定理 \ref{thm:xi-zg-dirichlet-log-renormalization}，\(\sum_{k\ge 2}\Delta_k(s)\) 在 \(\Re(s)>\tfrac12\) 的每个紧集上一致收敛。
另一方面，当 \(\Re(s)\ge \sigma_0>\tfrac12\) 时有
\(
\bigl|p_k^{-s}-\log(1+p_k^{-s})\bigr|
\ll p_k^{-2\sigma_0}
\),
从而 \(\sum_{k\ge 2}\bigl(p_k^{-s}-\log(1+p_k^{-s})\bigr)\) 在同一紧集上一致绝对收敛。
故 \(\log H_N\) 在 \(\Re(s)>\tfrac12\) 的紧集上一致收敛到某个全纯函数 \(\log H\)，并且 \(H=\exp(\log H)\) 全纯且不取零值。

对 \(\Re(s)>1\)，由平方自由 Dirichlet 级数的标准恒等式
\[
\prod_{p}\bigl(1+p^{-s}\bigr)
=\sum_{\substack{m\ge 1\\ m\ \text{平方因子自由}}}\frac1{m^s}
=\frac{\zeta(s)}{\zeta(2s)}
\]
（参见 \cite{Apostol1976,Titchmarsh1986RiemannZeta}），得 \(B_N(s)\to \zeta(s)/\zeta(2s)\)。
再由 \(F_N(s)=H_N(s)B_N(s)\) 并令 \(N\to\infty\)（在 \(\Re(s)>1\) 可由绝对收敛交换极限），得到 \(F(s)=H(s)\zeta(s)/\zeta(2s)\)。
最后，\(\zeta(2s)\) 在 \(\Re(s)>\tfrac12\) 上解析且无零点，而 \(\zeta(s)\) 在该半平面内唯一奇性为 \(s=1\) 的一阶极点，因而 \(F\) 在 \(\Re(s)>\tfrac12\) 上仅可能在 \(s=1\) 处有一阶极点。证毕。
\end{proof}

在平方自由 Euler 因子分离之外，\(\mathrm{ZG}\) 的 Dirichlet 级数还可进一步写成一条递归重整化的乘法链，其局部因子由相邻素数禁配约束逐步修正。

\begin{theorem}[Zeckendorf--Gödel Dirichlet 级数的 continued-Euler 型分解]\label{thm:xi-zg-continued-euler-factorization}
\leanverified{paper\_xi\_zg\_continued\_euler\_factorization}
沿用定理 \ref{thm:xi-zg-dirichlet-log-renormalization} 的记号。对实数 \(s>1\)，定义递归因子
\[
R_1(s):=p_1^{-s},
\qquad
R_n(s):=\frac{p_n^{-s}}{1+R_{n-1}(s)}
\qquad(n\ge 2).
\]
则对每个 \(n\ge 1\) 都有
\[
\boxed{\;
0<R_n(s)<p_n^{-s}.\;}
\]
并且 \(\mathrm{ZG}\) 的 Dirichlet 级数
\[
F(s)=\sum_{n\in \mathrm{ZG}}n^{-s}
\]
具有精确乘法分解
\[
\boxed{\;
F(s)=\prod_{n=1}^{\infty}\bigl(1+R_n(s)\bigr).\;}
\]
特别地，
\[
\boxed{\;
1<F(s)<\prod_{n=1}^{\infty}\bigl(1+p_n^{-s}\bigr)=\frac{\zeta(s)}{\zeta(2s)}
\qquad(s>1).\;}
\]
此外，\(F(s)\) 在区间 \((1,\infty)\) 上严格递减，并满足
\[
\boxed{\;
\lim_{s\to+\infty}F(s)=1.\;}
\]
\end{theorem}
\begin{proof}
由定理 \ref{thm:xi-zg-dirichlet-log-renormalization}，截断和
\[
F_N(s):=\sum_{\substack{\varepsilon\in\{0,1\}^N\\ \varepsilon_k\varepsilon_{k+1}=0}}
\prod_{k=1}^{N}p_k^{-s\varepsilon_k}
\]
满足递推
\[
F_N(s)=F_{N-1}(s)+p_N^{-s}F_{N-2}(s),
\qquad
F_0(s)=1,\quad F_1(s)=1+p_1^{-s}.
\]
再记
\[
r_N(s):=\frac{F_N(s)}{F_{N-1}(s)}
\qquad(N\ge 1).
\]
则同一定理给出
\[
r_1(s)=1+p_1^{-s},
\qquad
r_N(s)=1+\frac{p_N^{-s}}{r_{N-1}(s)}
\qquad(N\ge 2).
\]
定义
\[
R_N(s):=r_N(s)-1.
\]
于是
\[
R_1(s)=p_1^{-s},
\qquad
R_N(s)=\frac{p_N^{-s}}{1+R_{N-1}(s)}
\qquad(N\ge 2),
\]
即得所示递推。由 \(R_1(s)>0\) 出发归纳可知 \(R_N(s)>0\)；再因 \(1+R_{N-1}(s)>1\)，便有
\[
R_N(s)<p_N^{-s}.
\]

另一方面，
\[
F_N(s)=F_{N-1}(s)\,r_N(s)=F_{N-1}(s)\bigl(1+R_N(s)\bigr),
\]
故迭代得到
\[
F_N(s)=\prod_{n=1}^{N}\bigl(1+R_n(s)\bigr).
\]
由于 \(s>1\) 时 \(F_N(s)\to F(s)\)，令 \(N\to\infty\) 即得
\[
F(s)=\prod_{n=1}^{\infty}\bigl(1+R_n(s)\bigr).
\]

由 \(0<R_n(s)<p_n^{-s}\)，逐项比较可得
\[
1+R_n(s)<1+p_n^{-s}.
\]
因此
\[
1<F(s)=\prod_{n\ge 1}\bigl(1+R_n(s)\bigr)
<
\prod_{n\ge 1}\bigl(1+p_n^{-s}\bigr).
\]
最后一项等于 \(\zeta(s)/\zeta(2s)\) 是定理 \ref{thm:xi-zg-zeta-factorization} 中已经使用的平方自由 Euler 乘积恒等式。

再证单调性。若 \(1<s_1<s_2\)，则对每个 \(n\in\mathrm{ZG}\) 有
\[
n^{-s_1}\ge n^{-s_2},
\]
且在 \(n=2\) 处严格不等，于是
\[
F(s_1)-F(s_2)
=
\sum_{n\in\mathrm{ZG}}\bigl(n^{-s_1}-n^{-s_2}\bigr)>0.
\]
故 \(F\) 在 \((1,\infty)\) 上严格递减。

最后，由上界
\[
F(s)<\frac{\zeta(s)}{\zeta(2s)}
\]
且
\[
\frac{\zeta(s)}{\zeta(2s)}\longrightarrow 1
\qquad(s\to+\infty),
\]
并结合显然的下界 \(F(s)>1\)，由夹逼定理得到
\[
\lim_{s\to+\infty}F(s)=1.
\]
证毕。
\end{proof}


\begin{corollary}[素数路径硬核常数与 \texorpdfstring{$s=1$}{s=1} 留数]\label{cor:xi-zg-hardcore-constant-residue}
在定理 \ref{thm:xi-zg-zeta-factorization} 的记号下，记
\[
\mathfrak C_{\mathrm{HC}}:=H(1)=\lim_{N\to\infty}\frac{F_N(1)}{\prod_{k=1}^{N}\bigl(1+\frac1{p_k}\bigr)}\in(0,1).
\]
则 \(F(s)\) 在 \(s=1\) 的留数满足
\[
\operatorname*{Res}_{s=1}F(s)=\frac{\mathfrak C_{\mathrm{HC}}}{\zeta(2)}=\frac{6}{\pi^2}\,\mathfrak C_{\mathrm{HC}}.
\]
并且该留数与 \(\mathrm{ZG}\) 的 Dirichlet 密度一致；结合定理 \ref{con:xi-zg-density-closed-form-tail-bound} 的自然密度存在性，得到
\[
\delta_{\mathrm{ZG}}=\frac{\mathfrak C_{\mathrm{HC}}}{\zeta(2)}.
\]
数值核验可得到
\[
\mathfrak C_{\mathrm{HC}}=0.8461958103\ldots,
\qquad
\delta_{\mathrm{ZG}}=0.5144253666\ldots
\]
\leanverified{paper\_xi\_zg\_hardcore\_constant\_residue}
\end{corollary}

\begin{proof}
由定理 \ref{thm:xi-zg-zeta-factorization}，在 \(s=1\) 的邻域内
\(
F(s)=H(s)\zeta(s)/\zeta(2s)
\)
且 \(H\) 在 \(s=1\) 解析。
因此
\[
\operatorname*{Res}_{s=1}F(s)
=H(1)\operatorname*{Res}_{s=1}\frac{\zeta(s)}{\zeta(2s)}
=\frac{H(1)}{\zeta(2)}.
\]
留数等于 \(\mathrm{ZG}\) 的 Dirichlet 密度是推论 \ref{cor:xi-zg-dirichlet-residue} 的内容，而该推论亦指出留数与自然密度一致，从而得 \(\delta_{\mathrm{ZG}}=\mathfrak C_{\mathrm{HC}}/\zeta(2)\)。
最后，由于对实数 \(s>1\) 有 \(F_N(s)\le B_N(s)\) 且约束非平凡，故 \(0<H(s)<1\) 并推出 \(\mathfrak C_{\mathrm{HC}}\in(0,1)\)。证毕。
\end{proof}

\begin{proposition}[\texorpdfstring{$s=1$}{s=1} 的 Laurent 展开与 Euler--Kronecker 型常数]\label{prop:xi-zg-hardcore-laurent}
记
\[
\kappa_{\mathrm{HC}}:=\left.\frac{d}{ds}\log H(s)\right|_{s=1}.
\]
则 \(\kappa_{\mathrm{HC}}\) 存在，并且 \(F\) 在 \(s=1\) 的 Laurent 展开为
\[
F(s)=\frac{\mathfrak C_{\mathrm{HC}}}{\zeta(2)}\cdot\frac{1}{s-1}
\;+\;
\frac{\mathfrak C_{\mathrm{HC}}}{\zeta(2)}
\left(\gamma-\frac{2\zeta'(2)}{\zeta(2)}+\kappa_{\mathrm{HC}}\right)
 + O(s-1),
\qquad s\to 1,
\]
其中 \(\gamma\) 为 Euler--Mascheroni 常数。
数值核验表明 \(\kappa_{\mathrm{HC}}=0.40404\ldots\)。
\end{proposition}

\begin{proof}
由定理 \ref{thm:xi-zg-zeta-factorization} 的构造，\(\log H_N(s)\) 在 \(\Re(s)>\tfrac12\) 的紧集上一致收敛为 \(\log H(s)\)，并且其尾项由 \(\sum_{k\ge 2}\Delta_k(s)\) 与 \(\sum_{k\ge 2}(p_k^{-s}-\log(1+p_k^{-s}))\) 的一致收敛控制。
同样地，对 \(s\) 位于 \(s=1\) 的小邻域且满足 \(\Re(s)>\tfrac12\)，上述两级数可逐项求导，从而 \(\kappa_{\mathrm{HC}}\) 存在。

另一方面，
\(
\zeta(s)=\frac1{s-1}+\gamma+O(s-1)
\)
且
\[
\frac{1}{\zeta(2s)}
\;=\;
\frac{1}{\zeta(2)}
\left(1-\frac{2\zeta'(2)}{\zeta(2)}(s-1)+O((s-1)^2)\right).
\]
再由
\(
H(s)=H(1)\bigl(1+\kappa_{\mathrm{HC}}(s-1)+O((s-1)^2)\bigr)
\)
三者相乘即得所述展开式。证毕。
\end{proof}

% Auto-split to keep each TeX source < 600 lines.

\begin{theorem}[Zeckendorf--Gödel 素数码集合计数函数的幂节省误差]\label{thm:xi-zg-counting-power-saving-error}
记
\[
A(x):=\#\{n\le x:\ n\in\mathrm{ZG}\}
=\sum_{n\le x}\mathbf 1_{\mathrm{ZG}}(n).
\]
则对任意 \(\varepsilon>0\)，当 \(x\to\infty\) 有
\[
A(x)=\delta_{\mathrm{ZG}}\,x+O_\varepsilon\!\bigl(x^{\frac12+\varepsilon}\bigr).
\]
\end{theorem}

\begin{proof}
令 \(F(s)=\sum_{n\in\mathrm{ZG}}n^{-s}\)。
由定理 \ref{thm:xi-zg-zeta-factorization}，\(F\) 在 \(\Re(s)>\tfrac12\) 上亚纯且在 \(s=1\) 处有一阶极点；
其留数为 \(\delta_{\mathrm{ZG}}\)（推论 \ref{cor:xi-zg-hardcore-constant-residue}）。
\par
对非负系数 Dirichlet 级数应用 Perron 反演（参见 \cite{Apostol1976,Titchmarsh1986RiemannZeta}）可表示 \(A(x)\) 为沿竖线 \(\Re(s)=c>1\) 的积分。
将积分线移至 \(\Re(s)=\tfrac12+\varepsilon\) 时穿越 \(s=1\) 的留数贡献为 \(\delta_{\mathrm{ZG}}x\)。
由于 \(\zeta(2s)\) 在 \(\Re(s)\ge\tfrac12+\varepsilon\) 上解析且远离 \(0\)，并且 \(H(s)\) 在该半平面解析且由其对数展开的绝对收敛给出多项式增长界（定理 \ref{thm:xi-zg-dirichlet-log-renormalization}--\ref{thm:xi-zg-zeta-factorization}），结合 \(\zeta(s)\) 的经典竖线界即可控制移线后的积分余项，从而得到
\(
A(x)=\delta_{\mathrm{ZG}}x+O_\varepsilon(x^{\frac12+\varepsilon})
\)。
证毕。
\end{proof}

\begin{corollary}[Zeckendorf--Gödel 素数码集合的倒数调和渐近]\label{cor:xi-zg-reciprocal-harmonic-asymptotic}
在定理 \ref{thm:xi-zg-counting-power-saving-error} 的记号下，对任意 \(0<\varepsilon<\tfrac12\)，存在常数 \(C_{\mathrm{ZG}}\in\RR\) 使
\[
\sum_{\substack{n\le X\\ n\in\mathrm{ZG}}}\frac{1}{n}
=\delta_{\mathrm{ZG}}\log X+C_{\mathrm{ZG}}+O_\varepsilon\!\bigl(X^{-\frac12+\varepsilon}\bigr)
\qquad(X\to\infty).
\]
特别地，
\[
\sum_{\substack{n\le X\\ n\in\mathrm{ZG}}}\frac{1}{n}
\sim \delta_{\mathrm{ZG}}\log X.
\]
\end{corollary}

\begin{proof}
记误差项
\[
E(x):=A(x)-\delta_{\mathrm{ZG}}x.
\]
由定理 \ref{thm:xi-zg-counting-power-saving-error}，对任意 \(0<\varepsilon<\tfrac12\) 有
\[
E(x)=O_\varepsilon\!\bigl(x^{\frac12+\varepsilon}\bigr).
\]
对计数函数 \(A\) 应用 Abel 求和，得到
\[
\sum_{\substack{n\le X\\ n\in\mathrm{ZG}}}\frac{1}{n}
=\frac{A(X)}{X}+\int_{1}^{X}\frac{A(t)}{t^{2}}\,\dd t.
\]
代入 \(A(t)=\delta_{\mathrm{ZG}}t+E(t)\) 可得
\[
\sum_{\substack{n\le X\\ n\in\mathrm{ZG}}}\frac{1}{n}
=\delta_{\mathrm{ZG}}\log X+\delta_{\mathrm{ZG}}
+\frac{E(X)}{X}
+\int_{1}^{X}\frac{E(t)}{t^{2}}\,\dd t.
\]
由于
\[
\frac{E(X)}{X}=O_\varepsilon\!\bigl(X^{-\frac12+\varepsilon}\bigr),
\qquad
\frac{E(t)}{t^{2}}=O_\varepsilon\!\bigl(t^{-\frac32+\varepsilon}\bigr),
\]
且指数 \(-\tfrac32+\varepsilon<-1\)，故积分
\[
\int_{1}^{\infty}\frac{E(t)}{t^{2}}\,\dd t
\]
绝对收敛。定义
\[
C_{\mathrm{ZG}}
:=\delta_{\mathrm{ZG}}+\int_{1}^{\infty}\frac{E(t)}{t^{2}}\,\dd t.
\]
于是
\[
\int_{1}^{X}\frac{E(t)}{t^{2}}\,\dd t
=C_{\mathrm{ZG}}-\delta_{\mathrm{ZG}}+O_\varepsilon\!\bigl(X^{-\frac12+\varepsilon}\bigr),
\]
从而得到所述渐近展开。特别地，主项为 \(\delta_{\mathrm{ZG}}\log X\)，故结论成立。
\end{proof}



\begin{theorem}[Zeckendorf--Gödel 素数码集合的自然密度、闭式可计算性与尾项误差界]\label{con:xi-zg-density-closed-form-tail-bound}
对每个 \(N\ge 1\)，定义截断集合
\[
\mathrm{ZG}^{(N)}
:=\Bigl\{
n\in\NN:\ \forall 1\le k\le N,\ v_{p_k}(n)\in\{0,1\},\
\forall 1\le k\le N-1,\ v_{p_k}(n)\,v_{p_{k+1}}(n)=0
\Bigr\}.
\]
记
\[
\delta_N:=\mathrm{dens}\bigl(\mathrm{ZG}^{(N)}\bigr).
\]
则自然密度
\[
\delta_{\mathrm{ZG}}:=\mathrm{dens}(\mathrm{ZG})
\]
存在，且
\[
\delta_{\mathrm{ZG}}=\lim_{N\to\infty}\delta_N\in(0,1).
\]
进一步，\(\delta_N\) 具有显式闭式
\[
\delta_N=\Bigl(\prod_{k=1}^N\Bigl(1-\frac1{p_k}\Bigr)\Bigr)I_N,
\qquad
I_N:=\sum_{\substack{S\subseteq\{1,\dots,N\}\\ S\ \mathrm{无相邻}}}\ \prod_{k\in S}\frac1{p_k},
\]
并满足尾项误差上界
\[
0\le \delta_N-\delta_{\mathrm{ZG}}
\le
\sum_{k>N}\frac1{p_k^2}
\;+\;
\sum_{k\ge N}\frac1{p_kp_{k+1}}.
\]
特别地，\(\delta_{\mathrm{ZG}}\) 可通过生成素数并迭代计算 \(\delta_N\) 在任意精度下有效逼近。
\end{theorem}
\begin{proof}
对固定 \(N\)，成员资格仅依赖于模
\(
M_N:=\prod_{k=1}^N p_k^2
\)
的剩余类，故 \(\delta_N\) 存在。

记
\(
a_k:=1-\frac1{p_k},\ b_k:=a_k\frac1{p_k}
\)。
在前 \(N\) 个素数坐标上，“允许指数为 \(1\)”的索引集合恰为路径图 \(\{1,\dots,N\}\) 的独立集 \(S\)；
对应权重为
\(
\prod_{k\in S}b_k\prod_{k\notin S}a_k
\)。
求和得
\[
\delta_N
=\sum_{S\ \mathrm{无相邻}}\ \prod_{k\in S}b_k\prod_{k\notin S}a_k
=\Bigl(\prod_{k=1}^N a_k\Bigr)\sum_{S\ \mathrm{无相邻}}\prod_{k\in S}\frac1{p_k},
\]
即所示闭式。

又 \(\mathrm{ZG}=\bigcap_{N\ge 1}\mathrm{ZG}^{(N)}\)，故 \(\delta_N\) 单调下降且
\[
\overline{\mathrm{dens}}(\mathrm{ZG})\le \delta_N\quad(\forall N).
\]
若 \(n\in \mathrm{ZG}^{(N)}\setminus \mathrm{ZG}\)，则必发生尾部违约事件之一：
\[
\exists\,k>N,\ p_k^2\mid n
\qquad\text{或}\qquad
\exists\,k\ge N,\ p_kp_{k+1}\mid n.
\]
故
\[
\mathrm{ZG}^{(N)}\setminus \mathrm{ZG}
\subseteq
\Bigl(\bigcup_{k>N}\{p_k^2\mid n\}\Bigr)
\cup
\Bigl(\bigcup_{k\ge N}\{p_kp_{k+1}\mid n\}\Bigr).
\]
由并集上界与
\(
\mathrm{dens}\{m\mid n\}=1/m
\)
得
\[
\mathrm{dens}\bigl(\mathrm{ZG}^{(N)}\setminus \mathrm{ZG}\bigr)
\le
\sum_{k>N}\frac1{p_k^2}
+
\sum_{k\ge N}\frac1{p_kp_{k+1}},
\]
即误差估计。令 \(N\to\infty\) 得上下密度同限，故自然密度存在且等于 \(\lim\delta_N\)。

最后，
\[
\delta_{\mathrm{ZG}}
\ge 1-\sum_{k\ge 1}\frac1{p_k^2}-\sum_{k\ge 1}\frac1{p_kp_{k+1}}>0,
\]
因为两级数收敛且其和严格小于 \(1\)。证毕。
\end{proof}

\begin{corollary}[Mertens 重整常数与素数路径独立集分区函数的对数增长律]\label{cor:xi-zg-mertens-log-growth}
沿用定理 \ref{con:xi-zg-density-closed-form-tail-bound} 的记号。
则
\[
\delta_{\mathrm{ZG}}
=\lim_{N\to\infty}
\Bigl(\prod_{k=1}^N\Bigl(1-\frac1{p_k}\Bigr)\Bigr)I_N,
\]
并且
\[
I_N\sim \delta_{\mathrm{ZG}}\,e^{\gamma}\,\log p_N
\qquad (N\to\infty),
\]
其中 \(\gamma\) 为 Euler--Mascheroni 常数。
\end{corollary}
\begin{proof}
由定理 \ref{con:xi-zg-density-closed-form-tail-bound}，
\[
I_N=\frac{\delta_N}{\prod_{k\le N}(1-\frac1{p_k})}.
\]
再用 Mertens 乘积渐近
\(
\prod_{p\le y}(1-\frac1p)\sim e^{-\gamma}/\log y
\)
并令 \(y=p_N\)、\(\delta_N\to\delta_{\mathrm{ZG}}\) 即得。证毕。
\end{proof}

\begin{theorem}[素数权重 Fibonacci 路径的 Lee--Yang 实根性与 Poisson--binomial 分解]\label{con:xi-zg-weighted-path-leyang-poisson-binomial}
定义
\[
P_N(t):=\sum_{\substack{S\subseteq\{1,\dots,N\}\\ S\ \mathrm{无相邻}}}
t^{|S|}\prod_{k\in S}\frac1{p_k},
\]
则
\[
P_0(t)=1,\qquad
P_1(t)=1+\frac{t}{p_1},\qquad
P_N(t)=P_{N-1}(t)+\frac{t}{p_N}P_{N-2}(t)\quad(N\ge 2).
\]
并且成立：
\begin{enumerate}
  \item \(P_N\) 的全部零点均为互异负实根；
  \item \(P_N\) 与 \(P_{N-1}\) 的零点严格交错；
  \item 设 \(m_N:=\deg P_N=\lfloor(N+1)/2\rfloor\)，则存在 \(\alpha_{N,j}>0\) 使
  \[
  P_N(t)=\prod_{j=1}^{m_N}\Bigl(1+\frac{t}{\alpha_{N,j}}\Bigr).
  \]
\end{enumerate}
对固定 \(t>0\)，令
\[
\pi_{N,j}(t):=\frac{t}{t+\alpha_{N,j}}\in(0,1).
\]
若按 Gibbs 权重
\(
\propto t^{|S|}\prod_{k\in S}p_k^{-1}
\)
在独立集上取样，则随机变量 \(|S|\) 满足分布分解
\[
|S|\overset{d}{=}\sum_{j=1}^{m_N}B_{N,j},
\qquad
B_{N,j}\sim\mathrm{Bernoulli}\!\bigl(\pi_{N,j}(t)\bigr),
\]
其中 \((B_{N,j})\) 相互独立。
\end{theorem}
\begin{proof}
递推直接由索引 \(N\) 是否入独立集的分拆得到。

令 \(\mathcal T_{N+1}\) 为长度 \(N+1\) 的路径树，赋边权
\(
w_k:=p_k^{-1/2}
\)
（\(1\le k\le N\)），并记 \(A_{N+1}\) 为其加权邻接矩阵。
匹配多项式与独立集母函数的标准对应给出
\[
\chi_{N+1}(x):=\det(xI-A_{N+1})
=x^{N+1}P_N\!\Bigl(-\frac1{x^2}\Bigr).
\]
由于 \(A_{N+1}\) 为实对称 Jacobi 矩阵（副对角元严格正），其特征值实且互异；
故 \(P_N\) 的零点为 \(-1/\lambda_j^2\)，均为互异负实数。
交错性由主子矩阵谱交错（Cauchy interlacing）传递到 \(P_N,P_{N-1}\) 的零点。

第三条由实负根分解直接成立。
设
\[
G_{N,t}(s):=\frac{P_N(ts)}{P_N(t)}
=\prod_{j=1}^{m_N}
\frac{1+\frac{ts}{\alpha_{N,j}}}{1+\frac{t}{\alpha_{N,j}}}
=\prod_{j=1}^{m_N}\bigl(1-\pi_{N,j}(t)+\pi_{N,j}(t)s\bigr),
\]
右侧即独立 Bernoulli 和的概率母函数，故得到分布恒等式。证毕。
\end{proof}

\paragraph{素数权硬核归一化常数、Mertens--Lee--Yang 渐近与系数不等式链}
沿用定理 \ref{con:xi-zg-weighted-path-leyang-poisson-binomial} 的素数权路径独立集分区函数 \(P_N(t)\) 与其递推。
\par
\begin{theorem}[硬核归一化常数的存在性与显式尾项控制]\label{thm:xi-prime-hardcore-normalization-constant}
\leanverified{paper\_xi\_prime\_hardcore\_normalization\_constant}
对 \(t>0\) 与 \(N\ge 1\) 定义归一化比值
\begin{equation}\label{eq:xi-prime-hardcore-RN-def}
R_N(t)
:=
\frac{P_N(t)}{\prod_{k=1}^{N}\left(1+\frac{t}{p_k}\right)}\in(0,1].
\end{equation}
则：
\begin{enumerate}
  \item \((R_N(t))_{N\ge 1}\) 严格单调递减；
  \item 极限
  \begin{equation}\label{eq:xi-prime-hardcore-deltaHC-def}
  \boxed{\ \delta_{\mathrm{HC}}(t):=\lim_{N\to\infty}R_N(t)\ }
  \end{equation}
  存在且满足 \(0<\delta_{\mathrm{HC}}(t)<1\)；
  \item 对所有 \(N\ge 2\) 有尾项界
  \begin{equation}\label{eq:xi-prime-hardcore-RN-tail}
  \boxed{\;
  0\ \le\ R_N(t)-\delta_{\mathrm{HC}}(t)\ \le\ \sum_{k\ge N}\frac{t^2}{p_kp_{k+1}}\; }.
  \end{equation}
\end{enumerate}
\end{theorem}
\begin{proof}
记 \(a_k:=t/p_k\) 与 \(D_N(t):=\prod_{k=1}^{N}(1+a_k)\)。
由递推 \(P_N=P_{N-1}+a_NP_{N-2}\) 得
\[
R_N=\frac{R_{N-1}}{1+a_N}+\frac{a_N}{(1+a_{N-1})(1+a_N)}R_{N-2}.
\]
对 \(N\ge 3\) 进一步消元得到差分恒等式
\[
R_{N-1}-R_N
=\frac{a_{N-1}a_N}{(1+a_{N-2})(1+a_{N-1})(1+a_N)}\,R_{N-3}\ >\ 0,
\]
从而 \((R_N)\) 严格递减。
又由于 \(P_N\) 计数的独立集严格包含于所有子集，故 \(0<P_N\le D_N\)，从而 \(0<R_N\le 1\)。
\par
由上式立得上界 \(0<R_{N-1}-R_N\le a_{N-1}a_N=t^2/(p_{N-1}p_N)\)。
由于 \(\sum_{N\ge 2}p_{N-1}^{-1}p_N^{-1}\) 收敛，故 \(\sum_{N\ge 2}(R_{N-1}-R_N)\) 收敛，进而 \(R_N\) 收敛并满足尾项界 \eqref{eq:xi-prime-hardcore-RN-tail}。
\par
最后证明极限严格为正。
由 \(R_{N-2}\ge R_{N-1}\) 与上一差分恒等式得
\[
R_N
\ge R_{N-1}\left(1-\frac{a_{N-1}a_N}{(1+a_{N-1})(1+a_N)}\right),
\]
对 \(N\) 连乘并用 \(\sum_{N\ge 2}a_{N-1}a_N<\infty\) 得无穷乘积下界严格为正，从而 \(\delta_{\mathrm{HC}}(t)>0\)。证毕。
\end{proof}

\begin{corollary}[\texorpdfstring{$t=1$}{t=1} 与 \texorpdfstring{$\mathfrak C_{\mathrm{HC}}$}{C_HC} 的一致性]\label{cor:xi-prime-hardcore-deltaHC-matches-C_HC}
\leanverified{paper\_xi\_prime\_hardcore\_deltahc\_matches\_c\_hc}
在推论 \ref{cor:xi-zg-hardcore-constant-residue} 的记号下，有
\[
\delta_{\mathrm{HC}}(1)=\mathfrak C_{\mathrm{HC}}.
\]
从而
\(
\delta_{\mathrm{ZG}}=\delta_{\mathrm{HC}}(1)/\zeta(2)
\)
与推论 \ref{cor:xi-zg-hardcore-constant-residue} 的闭式一致。
\end{corollary}
\begin{proof}
由 \eqref{eq:xi-prime-hardcore-RN-def} 在 \(t=1\) 的定义，\(\delta_{\mathrm{HC}}(1)=\lim_{N\to\infty}P_N(1)/\prod_{k\le N}(1+1/p_k)\)。
而推论 \ref{cor:xi-zg-hardcore-constant-residue} 中 \(\mathfrak C_{\mathrm{HC}}\) 正是同一极限。证毕。
\end{proof}

\begin{theorem}[Mertens--Lee--Yang 渐近与两常数分解]\label{thm:xi-prime-hardcore-mertens-leyang-asymptotic}
令 \(\delta_{\mathrm{HC}}(t)\) 如定理 \ref{thm:xi-prime-hardcore-normalization-constant} 定义，并定义收敛欧拉乘积
\begin{equation}\label{eq:xi-prime-hardcore-Ct-euler}
C(t):=\prod_{p}\Bigl(1-\frac{1}{p}\Bigr)^{t}\Bigl(1+\frac{t}{p}\Bigr)\in(0,\infty).
\end{equation}
则当 \(N\to\infty\) 时有
\begin{equation}\label{eq:xi-prime-hardcore-asymptotic}
\boxed{\;
P_N(t)\ \sim\ \delta_{\mathrm{HC}}(t)\,C(t)\,e^{\gamma t}\,(\log p_N)^{t}\;}
\end{equation}
其中 \(\gamma\) 为 Euler--Mascheroni 常数。
\end{theorem}
\begin{proof}
由 \eqref{eq:xi-prime-hardcore-RN-def} 与 \eqref{eq:xi-prime-hardcore-deltaHC-def}，
\(
P_N(t)=R_N(t)\prod_{k\le N}(1+t/p_k)\sim \delta_{\mathrm{HC}}(t)\prod_{k\le N}(1+t/p_k)
\)。
将有限乘积分解为
\[
\prod_{k\le N}\Bigl(1+\frac{t}{p_k}\Bigr)
=
\Bigl(\prod_{k\le N}\Bigl(1-\frac{1}{p_k}\Bigr)^{-t}\Bigr)
\cdot
\Bigl(\prod_{k\le N}\Bigl(1-\frac{1}{p_k}\Bigr)^t\Bigl(1+\frac{t}{p_k}\Bigr)\Bigr).
\]
第二个括号内因子满足
\(
\bigl(1-\frac1p\bigr)^t(1+\frac{t}{p})
=1-\frac{t(t+1)}{2p^2}+O(p^{-3})
\)，
其对数为 \(O(p^{-2})\)，故乘积收敛并趋于 \(C(t)\)。
第一个括号由 Mertens 乘积渐近
\(
\prod_{p\le y}(1-\frac1p)\sim e^{-\gamma}/\log y
\)
（取 \(y=p_N\)）给出
\(
\prod_{k\le N}(1-\frac1{p_k})^{-t}\sim e^{\gamma t}(\log p_N)^t
\)。
合并即得 \eqref{eq:xi-prime-hardcore-asymptotic}。证毕。
\end{proof}

\begin{theorem}[素数权硬核 Gibbs 场的规模涨落与中心极限定理]\label{thm:xi-prime-hardcore-gibbs-clt}
\leanverified{paper\_xi\_prime\_hardcore\_gibbs\_clt}
固定 \(t>0\)。
在 \(\{1,\dots,N\}\) 的无相邻集合上按权重
\(
\propto t^{|S|}\prod_{k\in S}p_k^{-1}
\)
抽样，并记 \(K_N:=|S|\)。
则存在仅依赖于 \(t\) 的常数 \(\mu(t),\nu(t)\in\RR\)，使得
\[
\EE[K_N]=t\log\log p_N+\mu(t)+o(1),
\qquad
\mathrm{Var}(K_N)=t\log\log p_N+\nu(t)+o(1).
\]
并且令 \(V_N:=\mathrm{Var}(K_N)\)，则 \(V_N\to\infty\) 且存在绝对常数 \(C>0\) 使
\[
\sup_{x\in\RR}\left|
\PP\!\left(\frac{K_N-\EE[K_N]}{\sqrt{V_N}}\le x\right)-\Phi(x)
\right|
\le \frac{C}{\sqrt{V_N}}
=O\!\left(\frac{1}{\sqrt{\log\log p_N}}\right),
\]
从而 \(\frac{K_N-\EE[K_N]}{\sqrt{V_N}}\Rightarrow\mathcal N(0,1)\)。
\end{theorem}
\begin{proof}
由定理 \ref{thm:xi-prime-hardcore-mertens-leyang-asymptotic}，
\[
\log P_N(t)=t\log\log p_N+\log\delta_{\mathrm{HC}}(t)+\log C(t)+\gamma t+o(1).
\]
对该指数族分配，有标准恒等式
\(
\EE[K_N]=t\,\partial_t\log P_N(t)
\)
与
\(
\mathrm{Var}(K_N)=t\,\partial_t\EE[K_N]
\)，
从而得到所示两条渐近；其中常数 \(\mu(t),\nu(t)\) 由 \(\log\delta_{\mathrm{HC}}(t)+\log C(t)+\gamma t\) 的一、二阶导数给出。
\par
另一方面，由定理 \ref{con:xi-zg-weighted-path-leyang-poisson-binomial} 的 Poisson--binomial 分解，
\(
K_N\overset{d}{=}\sum_{j=1}^{m_N}B_{N,j}
\)
为独立 Bernoulli 和，且 \(\sum_j\mathrm{Var}(B_{N,j})=V_N\to\infty\)。
对 Poisson--binomial 和应用 Berry--Esseen 型界（以第三绝对中心矩上界为代价）得到误差 \(O(V_N^{-1/2})\)，从而得到中心极限定理。证毕。
\end{proof}

\begin{corollary}[单位权素数硬核 Gibbs 场的中心极限定律]\label{cor:xi-prime-hardcore-gibbs-clt-t1-explicit}
\leanverified{paper\_xi\_prime\_hardcore\_gibbs\_clt\_t1\_explicit}
在定理 \ref{thm:xi-prime-hardcore-gibbs-clt} 中取 \(t=1\)。
记
\[
I_N
:=
\sum_{\substack{S\subseteq\{1,\dots,N\}\\ S\ \mathrm{无相邻}}}\prod_{i\in S}\frac1{p_i},
\qquad
\PP_N(S):=\frac{1}{I_N}\prod_{i\in S}\frac1{p_i},
\]
并令 \(\omega_N:=|S|\)。
则
\[
\frac{\omega_N-\mu_N}{\sigma_N}\Longrightarrow \mathcal N(0,1),
\]
其中
\[
\mu_N=\sum_{i=1}^{N}\frac{1}{p_i+1}+O(1),
\qquad
\sigma_N^2=\sum_{i=1}^{N}\frac{p_i}{(p_i+1)^2}+O(1),
\]
并且
\[
\sigma_N^2\sim \sum_{i\le N}\frac1{p_i}\sim \log\log p_N.
\]
\end{corollary}
\begin{proof}
定理 \ref{thm:xi-prime-hardcore-gibbs-clt} 在 \(t=1\) 时给出
\[
\EE[\omega_N]=\log\log p_N+O(1),
\qquad
\mathrm{Var}(\omega_N)=\log\log p_N+O(1),
\]
并且
\(
(\omega_N-\EE[\omega_N])/\sqrt{\mathrm{Var}(\omega_N)}
\Rightarrow
\mathcal N(0,1)
\)。

另一方面，
\[
\frac1{p_i}-\frac{1}{p_i+1}=\frac{1}{p_i(p_i+1)}=O(p_i^{-2}),
\]
以及
\[
\frac1{p_i}-\frac{p_i}{(p_i+1)^2}
=
\frac{2p_i+1}{p_i(p_i+1)^2}
=
O(p_i^{-2}).
\]
由于 \(\sum_i p_i^{-2}<\infty\)，故
\[
\sum_{i=1}^{N}\frac{1}{p_i+1}
=
\sum_{i=1}^{N}\frac1{p_i}+O(1),
\qquad
\sum_{i=1}^{N}\frac{p_i}{(p_i+1)^2}
=
\sum_{i=1}^{N}\frac1{p_i}+O(1).
\]
再由素数倒数和的 Mertens 渐近，
\[
\sum_{i=1}^{N}\frac1{p_i}=\log\log p_N+O(1),
\]
即得所示 \(\mu_N,\sigma_N^2\) 的表达式与发散结论。
最后，由 \(\sigma_N^2=\mathrm{Var}(\omega_N)+O(1)\) 且 \(\sigma_N^2\to\infty\)，两种标准化只相差 \(1+o(1)\) 因子，从而极限正态性保持。证毕。
\end{proof}

\begin{proposition}[素数寄存器与 Lee--Yang 谱的 Newton--Maclaurin 对应及不等式链]\label{prop:xi-prime-hardcore-newton-maclaurin}
\leanverified{paper\_xi\_prime\_hardcore\_newton\_maclaurin}
设 \(m_N=\deg P_N=\lfloor(N+1)/2\rfloor\)，并写
\[
P_N(t)=\sum_{k=0}^{m_N}a_{N,k}\,t^k.
\]
则：
\begin{enumerate}
  \item 系数具有素数寄存器刻画
  \[
  a_{N,k}
  =\sum_{\substack{S\subseteq\{1,\dots,N\}\\ S\ \mathrm{无相邻},\ |S|=k}}
  \ \prod_{i\in S}\frac1{p_i}.
  \]
  \item 若 \(P_N(t)=\prod_{j=1}^{m_N}(1+t/\alpha_{N,j})\)（\(\alpha_{N,j}>0\)；定理 \ref{con:xi-zg-weighted-path-leyang-poisson-binomial}），则
  \[
  a_{N,k}=e_k\!\left(\frac{1}{\alpha_{N,1}},\dots,\frac{1}{\alpha_{N,m_N}}\right),
  \]
  其中 \(e_k\) 为第 \(k\) 个初等对称多项式；并且
  \[
  \sum_{j=1}^{m_N}\frac{1}{\alpha_{N,j}}=\sum_{i=1}^{N}\frac{1}{p_i}.
  \]
  \item 序列
  \[
  \left(\frac{a_{N,k}}{\binom{m_N}{k}}\right)^{1/k}\qquad(k=1,\dots,m_N)
  \]
  随 \(k\) 单调不增；并且 \(a_{N,k}\) 满足严格 Newton 不等式 \(a_{N,k}^2>a_{N,k-1}a_{N,k+1}\)（\(1\le k\le m_N-1\)）。
\end{enumerate}
\end{proposition}
\begin{proof}
前两条由 \(P_N\) 的独立集展开与 Lee--Yang 实根分解展开逐项比较得到。
第三条由正数列的 Maclaurin 不等式与实负互异零点多项式的严格 Newton 不等式推出。证毕。
\end{proof}


\begin{theorem}[Lee--Yang 单位圆提升、Joukowsky 椭圆输运与 \texorpdfstring{$\delta_{\mathrm{ZG}}$}{deltaZG} 的谱行列式表达]\label{con:xi-zg-unit-circle-joukowsky-spectral-determinant}
在定理 \ref{con:xi-zg-weighted-path-leyang-poisson-binomial} 的记号下：
\begin{enumerate}
  \item 设 \(\chi_{N+1}(x)=\det(xI-A_{N+1})\)，\(\Phi_{N+1}(x):=\chi_{N+1}(x/\sqrt2)\)。
  由于 \(\|A_{N+1}\|_{2}\le 2\max_k w_k=\sqrt2\)，\(\Phi_{N+1}\) 全部根落在 \([-2,2]\)。
  因此
  \[
  \mathcal L_{N+1}(z)
  :=z^{N+1}\Phi_{N+1}(z+z^{-1})
  =z^{N+1}\chi_{N+1}\!\Bigl(\frac{z+z^{-1}}{\sqrt2}\Bigr)
  \]
  的全部零点位于单位圆 \(|z|=1\)。
  \item 对任意 \(r>1\)，令 \(J_r(z):=rz+r^{-1}z^{-1}\)。
  若 \(z_j\) 为 \(\mathcal L_{N+1}\) 的零点，则 \(w_j:=J_r(z_j)\) 落在椭圆 \(J_r(S^1)\) 上。
  \item 取 \(x=i\) 于恒等式
  \(
  \chi_{N+1}(x)=x^{N+1}P_N(-x^{-2})
  \)
  得
  \[
  P_N(1)=|\chi_{N+1}(i)|=|\det(iI-A_{N+1})|.
  \]
  因而
  \[
  \delta_{\mathrm{ZG}}
  =\lim_{N\to\infty}
  \Bigl(\prod_{k=1}^N\Bigl(1-\frac1{p_k}\Bigr)\Bigr)\,|\det(iI-A_{N+1})|.
  \]
  又
  \[
  |\det(iI-A_{N+1})|^2
  =\det\!\bigl((iI-A_{N+1})(-iI-A_{N+1})\bigr)
  =\det(I+A_{N+1}^2),
  \]
  故
  \[
  P_N(1)=\det(I+A_{N+1}^2)^{1/2},
  \quad
  \delta_{\mathrm{ZG}}
  =\lim_{N\to\infty}
  \Bigl(\prod_{k=1}^N\Bigl(1-\frac1{p_k}\Bigr)\Bigr)\det(I+A_{N+1}^2)^{1/2}.
  \]
\end{enumerate}
\end{theorem}
\begin{proof}
第一条由谱半径界和参数化 \(x=z+z^{-1}=2\cos\theta\) 直接得到。
第二条是 \(J_r\) 对单位圆的标准像性质。
第三条由 \(x=i\) 代入与行列式恒等式逐步化简得出。证毕。
\end{proof}

\begin{corollary}[\texorpdfstring{$s=1$}{s=1} 处单极点与 Dirichlet 留数常数]\label{cor:xi-zg-dirichlet-residue}
令 \(P(s)=\sum_{k\ge 1}p_k^{-s}\) 为 prime zeta，并令 \(G(s)\) 为定理 \ref{thm:xi-zg-dirichlet-log-renormalization} 的正则因子。
采用 prime zeta 在 \(s\to 1^+\) 的经典展开
\[
P(s)=\log\frac1{s-1}+B_{\mathsf P}+o(1),
\]
其中 \(B_{\mathsf P}\in\RR\) 为常数（可由 Mertens 定理推导；参见 \cite{Apostol1976,Titchmarsh1986RiemannZeta}）。
则 \(F(s)=\sum_{n\in\mathrm{ZG}}n^{-s}\) 在 \(s=1\) 处具有单极点，且其留数为
\[
\lim_{s\to1^+}(s-1)F(s)=\exp\!\bigl(B_{\mathsf P}+G(1)\bigr)\in(0,\infty).
\]
并且该留数等于 \(\mathrm{ZG}\) 的 Dirichlet 密度；结合定理 \ref{con:xi-zg-density-closed-form-tail-bound} 的自然密度存在性，两者一致。
\end{corollary}

\paragraph{Abel 边界律与前缀递推证书}

\begin{theorem}[ZG 的 Dirichlet 级数在 \texorpdfstring{$s=1$}{s=1} 的 Abel 型边界律、调和主项与收敛临界线]\label{thm:xi-zg-abel-residue-log-density}
沿用本段记号，记
\[
A_{\mathrm{ZG}}(x):=\#\{n\le x:\ n\in\mathrm{ZG}\},
\qquad
D_{\mathrm{ZG}}(s):=\sum_{n\in\mathrm{ZG}}n^{-s}
\qquad (\Re(s)>1),
\]
以及
\[
\delta_{\mathrm{ZG}}:=\lim_{x\to\infty}\frac{A_{\mathrm{ZG}}(x)}{x}.
\]
则 \(\delta_{\mathrm{ZG}}\in(0,1)\)，并且成立：
\begin{enumerate}
  \item Abel 型边界律
  \[
  \lim_{\sigma\downarrow 1}(\sigma-1)D_{\mathrm{ZG}}(\sigma)=\delta_{\mathrm{ZG}}.
  \]
  \item 调和和主项
  \[
  \sum_{\substack{n\le X\\ n\in\mathrm{ZG}}}\frac{1}{n}
  =\delta_{\mathrm{ZG}}\log X+o(\log X)
  \qquad (X\to\infty).
  \]
  \item 绝对收敛临界线满足
  \[
  \sigma_c(D_{\mathrm{ZG}})=1.
  \]
\end{enumerate}
\end{theorem}
\begin{proof}
由定理 \ref{con:xi-zg-density-closed-form-tail-bound}，自然密度存在且
\[
\delta_{\mathrm{ZG}}\in(0,1).
\]
记
\[
A_{\mathrm{ZG}}(x)=\delta_{\mathrm{ZG}}\,x+r(x),
\qquad
r(x)=o(x)
\qquad (x\to\infty).
\]

对任意 \(\sigma>1\)，Abel--Stieltjes 分部积分给出
\[
D_{\mathrm{ZG}}(\sigma)
=
\sigma\int_1^\infty A_{\mathrm{ZG}}(x)x^{-\sigma-1}\,\dd x.
\]
代入 \(A_{\mathrm{ZG}}(x)=\delta_{\mathrm{ZG}}x+r(x)\) 后得到
\[
D_{\mathrm{ZG}}(\sigma)
=
\frac{\sigma\,\delta_{\mathrm{ZG}}}{\sigma-1}
+
\sigma\int_1^\infty r(x)x^{-\sigma-1}\,\dd x.
\]
因此只需证明
\[
\lim_{\sigma\downarrow1}
(\sigma-1)\int_1^\infty r(x)x^{-\sigma-1}\,\dd x=0.
\]
任取 \(\varepsilon>0\)。
由 \(r(x)=o(x)\)，可取 \(X_0\ge 1\) 使得 \(|r(x)|\le \varepsilon x\) 对所有 \(x\ge X_0\) 成立。
于是
\[
\begin{aligned}
(\sigma-1)\left|\sigma\int_1^\infty r(x)x^{-\sigma-1}\,\dd x\right|
&\le
(\sigma-1)\sigma\int_1^{X_0}|r(x)|x^{-\sigma-1}\,\dd x \\
&\qquad
+(\sigma-1)\sigma\int_{X_0}^{\infty}\varepsilon x\cdot x^{-\sigma-1}\,\dd x \\
&\le
(\sigma-1)C_{X_0}
+\sigma\varepsilon X_0^{1-\sigma},
\end{aligned}
\]
其中 \(C_{X_0}<\infty\) 与 \(\sigma\) 无关。
令 \(\sigma\downarrow1\)，上式上极限不超过 \(\varepsilon\)；再令 \(\varepsilon\downarrow0\)，即得
\[
\lim_{\sigma\downarrow 1}(\sigma-1)D_{\mathrm{ZG}}(\sigma)=\delta_{\mathrm{ZG}}.
\]

再看第二条。
对和式作 Abel 分部积分，
\[
\sum_{\substack{n\le X\\ n\in\mathrm{ZG}}}\frac1n
=
\frac{A_{\mathrm{ZG}}(X)}{X}+\int_1^X \frac{A_{\mathrm{ZG}}(t)}{t^2}\,\dd t.
\]
代入 \(A_{\mathrm{ZG}}(t)=\delta_{\mathrm{ZG}}t+r(t)\) 得
\[
\sum_{\substack{n\le X\\ n\in\mathrm{ZG}}}\frac1n
=
\delta_{\mathrm{ZG}}
+\frac{r(X)}{X}
+\delta_{\mathrm{ZG}}\log X
+\int_1^X \frac{r(t)}{t^2}\,\dd t.
\]
由 \(r(t)=o(t)\)，对任意 \(\varepsilon>0\) 仍可取 \(T_0\) 使 \(|r(t)|\le \varepsilon t\) 对 \(t\ge T_0\) 成立，从而
\[
\left|\int_1^X \frac{r(t)}{t^2}\,\dd t\right|
\le C_{T_0}+\varepsilon\log X,
\]
其中 \(C_{T_0}\) 与 \(X\) 无关。
因此
\[
\sum_{\substack{n\le X\\ n\in\mathrm{ZG}}}\frac1n
=
\delta_{\mathrm{ZG}}\log X+o(\log X).
\]
最后，由 \(\mathrm{ZG}\subset \NN\) 可知对任意 \(\sigma>1\) 都有
\[
0\le D_{\mathrm{ZG}}(\sigma)\le \zeta(\sigma)<\infty,
\]
故 \(D_{\mathrm{ZG}}\) 在半平面 \(\Re(s)>1\) 上绝对收敛，从而 \(\sigma_c(D_{\mathrm{ZG}})\le 1\)。另一方面，第二条等价于
\[
\sum_{\substack{n\le X\\ n\in\mathrm{ZG}}}\frac1n
=
\delta_{\mathrm{ZG}}\log X+o(\log X)\to\infty,
\]
从而级数 \(D_{\mathrm{ZG}}(1)=\sum_{n\in\mathrm{ZG}}n^{-1}\) 发散。于是其收敛 abscissa 恰为 \(1\)。
\end{proof}

\begin{corollary}[ZG 的 Dirichlet 级数的绝对收敛临界线恰为 \texorpdfstring{$\Re(s)=1$}{Re(s)=1}]\label{cor:xi-zg-absolute-convergence-critical-line}
在定理 \ref{thm:xi-zg-abel-residue-log-density} 的设定下，\(D_{\mathrm{ZG}}(s)\) 在 \(\Re(s)>1\) 上绝对收敛，在 \(s=1\) 发散，并满足
\[
D_{\mathrm{ZG}}(\sigma)\sim \frac{\delta_{\mathrm{ZG}}}{\sigma-1}
\qquad (\sigma\downarrow 1).
\]
因而 \(\Re(s)=1\) 正是其绝对收敛临界线。
\end{corollary}
\begin{proof}
由定理 \ref{thm:xi-zg-abel-residue-log-density} 的第一条，
\[
(\sigma-1)D_{\mathrm{ZG}}(\sigma)\longrightarrow \delta_{\mathrm{ZG}}>0
\qquad (\sigma\downarrow 1),
\]
故
\[
D_{\mathrm{ZG}}(\sigma)
=
\frac{\delta_{\mathrm{ZG}}}{\sigma-1}
+o\!\bigl((\sigma-1)^{-1}\bigr)
\qquad (\sigma\downarrow 1).
\]
特别地，\(D_{\mathrm{ZG}}(\sigma)\to+\infty\)。再结合同一定理的第三条，便得到所述绝对收敛临界线结论。证毕。
\end{proof}

\begin{proposition}[Zeckendorf--Gödel 密度的归一化前缀递推、单调收敛与尾误差证书]\label{prop:xi-zg-prefix-recursion-monotone-tail-certificate}
沿用定理 \ref{con:xi-zg-density-closed-form-tail-bound} 与
\ref{con:xi-zg-weighted-path-leyang-poisson-binomial} 的记号。
定义数列
\[
(a_0,b_0)=(1,0),
\]
并对 \(N\ge 1\) 递推地令
\[
a_N=(a_{N-1}+b_{N-1})\frac{p_N}{p_N+1},
\qquad
b_N=\frac{a_{N-1}}{p_N+1}.
\]
再记
\[
u_N:=a_N+b_N,
\qquad
\widetilde\delta_N:=\frac{u_N}{\zeta(2)}.
\]
则成立：
\begin{enumerate}
  \item 对每个 \(N\ge 0\)，有
  \[
  u_N=
  \frac{P_N(1)}{\prod_{k=1}^{N}\left(1+\frac1{p_k}\right)},
  \]
  其中约定空积等于 \(1\)。
  因而
  \[
  \delta_{\mathrm{ZG}}
  =
  \frac{1}{\zeta(2)}\lim_{N\to\infty}u_N.
  \]
  \item \(u_0=u_1=1\)，并且数列 \((\widetilde\delta_N)_{N\ge 1}\) 严格单调下降，满足
  \[
  \widetilde\delta_N\downarrow \delta_{\mathrm{ZG}}.
  \]
  \item 对每个 \(N\ge 1\)，有显式尾误差界
  \[
  0\le \widetilde\delta_N-\delta_{\mathrm{ZG}}
  \le
  \widetilde\delta_N
  \sum_{j=N}^{\infty}\frac{1}{(p_j+1)(p_{j+1}+1)}.
  \]
\end{enumerate}
\end{proposition}
\begin{proof}
对 \(N\ge 1\)，定义两类部分配分函数
\[
A_N:=
\sum_{\substack{S\subseteq\{1,\dots,N\}\\ S\ \mathrm{无相邻}\\ N\notin S}}
\prod_{k\in S}\frac1{p_k},
\qquad
B_N:=
\sum_{\substack{S\subseteq\{1,\dots,N\}\\ S\ \mathrm{无相邻}\\ N\in S}}
\prod_{k\in S}\frac1{p_k}.
\]
则
\[
P_N(1)=A_N+B_N.
\]
并且由末位是否被选中，直接得到
\[
A_N=P_{N-1}(1),
\qquad
B_N=\frac{1}{p_N}A_{N-1}.
\]
把它们按
\[
a_N:=\frac{A_N}{\prod_{k=1}^{N}\left(1+\frac1{p_k}\right)},
\qquad
b_N:=\frac{B_N}{\prod_{k=1}^{N}\left(1+\frac1{p_k}\right)}
\]
归一化，便得到所述递推关系
\[
a_N=(a_{N-1}+b_{N-1})\frac{p_N}{p_N+1},
\qquad
b_N=\frac{a_{N-1}}{p_N+1}.
\]
于是
\[
u_N=a_N+b_N
=
\frac{P_N(1)}{\prod_{k=1}^{N}\left(1+\frac1{p_k}\right)}.
\]
这就证明了第一条中的第一式。

另一方面，由定理 \ref{con:xi-zg-density-closed-form-tail-bound}，
\[
\delta_N
=
\left(\prod_{k=1}^{N}\left(1-\frac1{p_k}\right)\right)P_N(1).
\]
将上一式代入并整理得
\[
\delta_N
=
\left(\prod_{k=1}^{N}\left(1-\frac1{p_k^2}\right)\right)u_N.
\]
由于 \(\delta_N\to \delta_{\mathrm{ZG}}\) 且
\[
\prod_{k=1}^{N}\left(1-\frac1{p_k^2}\right)\longrightarrow \frac1{\zeta(2)},
\]
只要 \(u_N\) 收敛，便有
\[
\delta_{\mathrm{ZG}}=\frac{1}{\zeta(2)}\lim_{N\to\infty}u_N.
\]

现证单调性。
由递推直接计算
\[
u_N
=
u_{N-1}-\frac{b_{N-1}}{p_N+1}.
\]
因此 \(u_0=u_1=1\)。
而对 \(N\ge 2\)，由 \(a_{N-2}>0\) 可得
\[
b_{N-1}=\frac{a_{N-2}}{p_{N-1}+1}>0,
\]
故
\[
u_N<u_{N-1}\qquad (N\ge 2).
\]
于是 \((u_N)_{N\ge 1}\) 严格递减且下有界于 \(0\)，故极限
\(
u_\infty:=\lim_{N\to\infty}u_N
\)
存在。
由前述关系，便得到
\[
\widetilde\delta_N=\frac{u_N}{\zeta(2)}\downarrow \frac{u_\infty}{\zeta(2)}=\delta_{\mathrm{ZG}}.
\]

最后证明尾误差界。
由
\[
u_N-u_{N+1}=\frac{b_N}{p_{N+1}+1}
\]
以及
\[
b_N=\frac{a_{N-1}}{p_N+1},
\qquad
u_N=a_{N-1}+\frac{p_N}{p_N+1}b_{N-1}\ge a_{N-1},
\]
可得
\[
u_N-u_{N+1}
\le
\frac{u_N}{(p_N+1)(p_{N+1}+1)}.
\]
而由于 \((u_N)\) 递减，对任意 \(j\ge N\) 亦有
\[
u_j-u_{j+1}
\le
\frac{u_j}{(p_j+1)(p_{j+1}+1)}
\le
\frac{u_N}{(p_j+1)(p_{j+1}+1)}.
\]
故
\[
u_N-u_\infty
=
\sum_{j=N}^{\infty}(u_j-u_{j+1})
\le
u_N\sum_{j=N}^{\infty}\frac{1}{(p_j+1)(p_{j+1}+1)}.
\]
两边乘以 \(\zeta(2)^{-1}\) 即得
\[
0\le \widetilde\delta_N-\delta_{\mathrm{ZG}}
\le
\widetilde\delta_N\sum_{j=N}^{\infty}\frac{1}{(p_j+1)(p_{j+1}+1)}.
\]
证毕。
\end{proof}


\begin{proposition}[临界重整化常数的可计算误差界]\label{prop:xi-zg-holographic-constant-tail-bound}
\leanverified{paper\_xi\_zg\_holographic\_constant\_tail\_bound}
取 \(s=1\)，并令 \(\Delta_N(1)\) 与 \(G(1)=\sum_{N\ge 2}\Delta_N(1)\) 如定理 \ref{thm:xi-zg-dirichlet-log-renormalization}。
定义截断
\[
G^{(N)}(1):=\sum_{k=2}^{N}\Delta_k(1).
\]
则存在绝对常数 \(C>0\) 使对所有 \(N\ge 3\) 有
\[
\bigl|G(1)-G^{(N)}(1)\bigr|
\le
C\sum_{k>N}\Bigl(\frac1{p_k^2}+\frac1{p_kp_{k-1}}\Bigr),
\]
从而正则因子 \(\exp(G(1))\) 可被一条显式可求和尾项在任意精度下有效逼近。
\end{proposition}

\begin{proof}
由定理 \ref{thm:xi-zg-dirichlet-log-renormalization} 的一致估计，在 \(s=1\) 时存在常数 \(C\) 使
\(
|\Delta_k(1)|
\le C\bigl(p_k^{-2}+(p_kp_{k-1})^{-1}\bigr)
\)。
对尾和求并界即得。证毕。
\end{proof}

\begin{theorem}[\(\mathrm{ZG}\) 的对数 Mertens 定律]\label{thm:xi-zg-log-mertens}
设
\[
H(x):=\sum_{\substack{n\le x\\ n\in\mathrm{ZG}}}\frac1n.
\]
则极限存在并满足
\[
\lim_{x\to\infty}\frac{H(x)}{\log x}=\delta_{\mathrm{ZG}},
\]
并且常数项可由 \(F(s)=\sum_{n\in\mathrm{ZG}}n^{-s}\) 在 \(s=1\) 的 Laurent 常数项显式识别：存在 \(C_{\mathrm{ZG}}\in\RR\) 使
\[
H(x)=\delta_{\mathrm{ZG}}\log x+C_{\mathrm{ZG}}+o(1)\qquad(x\to\infty).
\]
其中
\[
C_{\mathrm{ZG}}
=\frac{\mathfrak C_{\mathrm{HC}}}{\zeta(2)}
\left(\gamma-\frac{2\zeta'(2)}{\zeta(2)}+\kappa_{\mathrm{HC}}\right)
=\delta_{\mathrm{ZG}}
\left(\gamma-\frac{2\zeta'(2)}{\zeta(2)}+\kappa_{\mathrm{HC}}\right).
\]
\end{theorem}

\begin{proof}
由推论 \ref{cor:xi-zg-dirichlet-residue} 或推论 \ref{cor:xi-zg-hardcore-constant-residue}，\(F(s)=\sum_{n\in\mathrm{ZG}}n^{-s}\) 在 \(s=1\) 处的留数为 \(\delta_{\mathrm{ZG}}\)。
对非负系数 Dirichlet 级数应用标准 Tauberian 定理可得
\(
H(x)=\delta_{\mathrm{ZG}}\log x+C_{\mathrm{ZG}}+o(1)
\)，其中 \(C_{\mathrm{ZG}}\) 等于 \(F\) 在 \(s=1\) 的 Laurent 常数项。
而命题 \ref{prop:xi-zg-hardcore-laurent} 给出了该 Laurent 常数项的显式表达式，从而得到所示恒等式。证毕。
\end{proof}

\begin{conjecture}[正则因子的全半平面延拓与 Lee--Yang 型零点边界]\label{conj:xi-zg-leyang-godel-transport-boundary}
令 \(\mathcal R(s):=\exp(G(s))\)（定理 \ref{thm:xi-zg-dirichlet-log-renormalization}），其在 \(\Re(s)>\tfrac12\) 全纯且无零。
猜想 \(\mathcal R\) 可延拓为 \(\Re(s)>0\) 上的全纯无零函数；从而 \(F(s)=\exp(P(s))\mathcal R(s)\) 在 \(\Re(s)>0\) 的全部奇性结构由 prime zeta \(P(s)\) 的奇性传输唯一决定。
进一步地，对充分大的截断 \(N\)，将 \(F_N(s)\) 视为有限体积分配函数，则其归一化零点云在 \(N\to\infty\) 的极限中形成一条由 Joukowsky 椭圆输运诱导的边界曲线，并以 \(\Re(s)=\tfrac12\) 为临界对称轴，与 \(\mathcal R\) 的零点真空相容。
\end{conjecture}

\begin{conjecture}[弱局域约束下的 Erd\H{o}s--Kac 型高斯极限]\label{conj:xi-zg-erdos-kac-under-local-hardcore}
记
\[
\omega(n):=\#\{p:\ p\mid n\},
\qquad
\mathcal Z(X):=\{n\le X:\ n\in\mathrm{ZG}\}.
\]
猜想存在常数 \(c_0\in\RR\)，使得在 \(\mathcal Z(X)\) 上均匀抽样的 \(n\) 满足
\[
\frac{\omega(n)-\log\log X-c_0}{\sqrt{\log\log X}}
\Longrightarrow
\mathcal N(0,1)
\qquad(X\to\infty).
\]
该猜想表明：相邻素数不可同现的局域硬核约束在中心极限定律尺度上仅导致有限平移，而不改变标准化极限分布。
\end{conjecture}

\begin{remark}[可复现核验脚本]\label{rem:xi-zg-density-leyang-spectral-audit-script}
定理 \ref{con:xi-zg-density-closed-form-tail-bound}--\ref{con:xi-zg-unit-circle-joukowsky-spectral-determinant}
中的闭式恒等式与有限窗数值核验由脚本
\texttt{scripts/exp\_xi\_prime\_register\_zg\_density\_leyang\_spectrum\_audit.py}
给出，对应证书文件为
\texttt{artifacts/export/xi\_prime\_register\_zg\_density\_leyang\_spectrum\_audit.json}。
\par
定理 \ref{thm:xi-zg-dirichlet-log-renormalization} 与命题 \ref{prop:xi-zg-holographic-constant-tail-bound}
在 \(s=1\) 处的递推与尾项模板核验由脚本
\texttt{scripts/exp\_xi\_prime\_register\_zg\_dirichlet\_log\_renormalization\_audit.py}
给出，对应证书文件为
\texttt{artifacts/export/xi\_prime\_register\_zg\_dirichlet\_log\_renormalization\_audit.json}。
\par
定理 \ref{thm:xi-zg-zeta-factorization}--命题 \ref{prop:xi-zg-hardcore-laurent}
中 \(\mathfrak C_{\mathrm{HC}}\) 与 \(\kappa_{\mathrm{HC}}\) 的数值核验由脚本
\texttt{scripts/exp\_xi\_prime\_register\_zg\_hardcore\_zeta\_factorization\_constant\_audit.py}
给出，对应证书文件为
\texttt{artifacts/export/xi\_prime\_register\_zg\_hardcore\_zeta\_factorization\_constant\_audit.json}。
\end{remark}


\begin{theorem}[行列式素因子分解给出的椭圆面积守恒律]\label{thm:xi-terminal-ellipse-area-euler-prime-law}
\leanverified{paper\_xi\_terminal\_ellipse\_area\_euler\_prime\_law}
设 \(A\in\mathrm{GL}_2(\QQ)\)，单位圆盘 \(D:=\{x\in\RR^2:\|x\|\le 1\}\)，椭圆盘 \(E_A:=A(D)\)。
则
\[
\mathrm{Area}(E_A)=\pi\,|\det A|.
\]
并且若
\(
\det A=\pm\prod_p p^{v_p(\det A)}
\)
（仅有限个 \(v_p\neq 0\)），则
\[
\mathrm{Area}(E_A)=\pi\prod_p p^{v_p(\det A)}.
\]
\end{theorem}
\begin{proof}
二维 Lebesgue 测度在可逆线性变换下按雅可比因子 \( |\det A| \) 缩放，故
\(
\mathrm{Area}(A(D))=|\det A|\,\mathrm{Area}(D)=\pi|\det A|
\)。
其余结论由有理数乘法群的素因子分解直接得到。证毕。
\end{proof}

\begin{theorem}[有理椭圆合同类的二元不变量完备性]\label{thm:xi-terminal-rational-ellipse-two-invariant-completeness}
\leanverified{paper\_xi\_terminal\_rational\_ellipse\_two\_invariant\_completeness}
设 \(A\in\mathrm{GL}_2(\QQ)\)，记 \(E_A:=A(S^1)\subset\RR^2\)。
定义
\[
\Delta_1(A):=(\det A)^2\in\QQ_{>0},
\qquad
\Delta_2(A):=\mathrm{tr}(A^\top A)\in\QQ_{>0}.
\]
则 \(E_A\) 的欧氏合同类由 \((\Delta_1(A),\Delta_2(A))\) 完全决定。
若 \(\sigma_1\ge \sigma_2>0\) 为 \(A\) 的奇异值，则
\[
\sigma_1^2,\sigma_2^2\ \text{是}\ t^2-\Delta_2(A)t+\Delta_1(A)=0\ \text{的两根},
\]
因而
\[
\sigma_{1,2}
=\sqrt{\frac{\Delta_2(A)\pm\sqrt{\Delta_2(A)^2-4\Delta_1(A)}}{2}}.
\]
特别地，对任意 \(A,B\in\mathrm{GL}_2(\QQ)\),
\[
E_A\cong E_B\ (\text{欧氏合同})
\Longleftrightarrow
\bigl(\Delta_1(A),\Delta_2(A)\bigr)=\bigl(\Delta_1(B),\Delta_2(B)\bigr).
\]
\end{theorem}
\begin{proof}
极分解 \(A=UP\)（\(U\) 正交，\(P=(A^\top A)^{1/2}\) 对称正定）给出
\(
E_A=U(P(S^1))
\)；
故合同类由 \(P\) 决定。
\(P\) 的特征值为 \(\sigma_1,\sigma_2\)，即椭圆半轴长。
又
\[
\sigma_1^2+\sigma_2^2=\mathrm{tr}(A^\top A)=\Delta_2(A),
\qquad
\sigma_1^2\sigma_2^2=\det(A^\top A)=(\det A)^2=\Delta_1(A),
\]
因此 \((\sigma_1,\sigma_2)\) 由 \((\Delta_1,\Delta_2)\) 唯一恢复，反之亦然。证毕。
\end{proof}

% AUTO-GENERATED by scripts/exp_xi_prime_register_semidirect_two_shadow_audit.py
\begin{table}[H]
\centering
\scriptsize
\setlength{\tabcolsep}{6pt}
\caption{有理矩阵样本的椭圆二元不变量审计：$\Delta_1=(\det A)^2$ 与 $\Delta_2=\mathrm{tr}(A^\top A)$。}
\label{tab:xi_elliptic_two_shadow_invariants_audit}
\begin{tabular}{l l l l r r}
\toprule
id & $\det A$ & $\Delta_1(A)$ & $\Delta_2(A)$ & $\sigma_1$ (num.) & $\sigma_2$ (num.)\\
\midrule
$A_1$ & $209/120$ & $43681/14400$ & $14701/3600$ & 1.762791958 & 0.988016004\\
$B_1=A_1Q_{90}$ & $209/120$ & $43681/14400$ & $14701/3600$ & 1.762791958 & 0.988016004\\
$A_2$ & $163/70$ & $26569/4900$ & $233209/44100$ & 1.974001870 & 1.179619667\\
\bottomrule
\end{tabular}
\end{table}


\begin{theorem}[有限素部格与实部椭圆影的二影刚性]\label{thm:xi-terminal-finite-archimedean-two-shadow-rigidity}
定义 prime-register 纤维
\[
\widehat{\ZZ}:=\varprojlim_m \ZZ/m\ZZ\cong\prod_p \ZZ_p.
\]
对 \(A\in\mathrm{GL}_2(\QQ)\)，记有限素部格影
\[
\Lambda_A:=A\,\widehat{\ZZ}^{\,2}\subset \mathbb A_f^{\,2},
\]
并记实部椭圆影 \(E_A:=A(S^1)\subset\RR^2\)。
若 \(A,B\in\mathrm{GL}_2(\QQ)\) 满足
\[
\Lambda_A=\Lambda_B,
\qquad
E_A=E_B,
\]
则存在 \(Q\in O_2(\ZZ)\) 使 \(A=BQ\)。
若再施加定向保持规范（等价于 \(\det(B^{-1}A)=1\)），则 \(Q\in SO_2(\ZZ)\)，因此仅余 \(4\) 个离散候选；
进一步施加任一破缺该四元对称的规范锚定后，\(Q\) 唯一，从而 \(A\) 由 \((\Lambda_A,E_A)\) 唯一确定。
\end{theorem}
\begin{proof}
令 \(C:=B^{-1}A\)。
由 \(\Lambda_A=\Lambda_B\) 得
\(
C\,\widehat{\ZZ}^{\,2}=\widehat{\ZZ}^{\,2}
\)，故 \(C\in\mathrm{GL}_2(\widehat{\ZZ})\)。

由 \(E_A=E_B\) 得同一椭圆二次型：
\[
x^\top(AA^\top)^{-1}x=1
\Longleftrightarrow
x^\top(BB^\top)^{-1}x=1,
\]
从而 \(AA^\top=BB^\top\)。
左乘 \(B^{-1}\)、右乘 \((B^{-1})^\top\) 得
\[
CC^\top=I,
\]
即 \(C\in O_2(\QQ)\)。

同时 \(C\in\mathrm{GL}_2(\widehat{\ZZ})\) 蕴含其每个有理条目在所有 \(\ZZ_p\) 中整，故条目属于
\(
\QQ\cap\bigcap_p\ZZ_p=\ZZ
\)；
即 \(C\in\mathrm{GL}_2(\ZZ)\)。
综上 \(C\in O_2(\ZZ)\)，即 \(A=BQ\)（\(Q:=C\in O_2(\ZZ)\)）。

当 \(\det Q=1\) 时 \(Q\in SO_2(\ZZ)\)。
二维整正交群中 \(|SO_2(\ZZ)|=4\)，故仅有四个定向保持候选；对任意固定规范锚定（例如长轴正向锚定）即可消去残余离散歧义并得到唯一性。证毕。
\end{proof}

\paragraph{全对称 Belyi 覆盖、判别式闭式与多临界 Puiseux 缩放}
本段将一参数代数族
\(
\Pi_r(\lambda,y)=(\lambda-1)^r+(y-1)\lambda^{r-1}
\)
组织为三点分歧覆盖，并给出其判别式闭式、Puiseux 多临界指数与 \(r=3\) 的显式 Lee--Yang 区间模型。

\begin{theorem}[全对称 Belyi 覆盖与满对称单值群]\label{thm:xi-terminal-belyi-full-symmetric-cover}
设 \(r\ge 2\) 为整数，并定义
\[
\Pi_r(\lambda,y):=(\lambda-1)^r+(y-1)\lambda^{r-1}\in\ZZ[\lambda,y],
\qquad
f_r(\lambda):=1-\frac{(\lambda-1)^r}{\lambda^{r-1}}.
\]
则成立：
\begin{enumerate}
  \item \(f_r:\PP^1_\lambda\to\PP^1_y\) 为次数 \(r\) 的三点分歧覆盖；其有限分歧值恰为
  \[
  y=1,\qquad y=y_r:=1+\frac{r^r}{(r-1)^{r-1}},
  \]
  并与 \(y=\infty\) 共同构成全部分歧值。
  \item 三处分歧值的分歧型分别为
  \[
  y=1:\ (r),\qquad
  y=\infty:\ (r-1)(1),\qquad
  y=y_r:\ (2)(1^{r-2}).
  \]
  \item 覆盖的几何单值群为
  \[
  \mathrm{Mon}_{\mathrm{geo}}(f_r)\cong S_r.
  \]
  因而 \(\Pi_r(\lambda,y)\) 在 \(\CC(y)[\lambda]\) 上不可约。
\end{enumerate}
\end{theorem}
\begin{proof}
写
\[
f_r(\lambda)=\frac{\lambda^{r-1}-(\lambda-1)^r}{\lambda^{r-1}}=\frac{P_r(\lambda)}{Q_r(\lambda)},
\]
其中 \(\deg P_r=r\)、\(\deg Q_r=r-1\)，且 \(P_r(0)\neq 0\)，故 \(\deg f_r=r\)。

直接计算得
\[
f_r'(\lambda)=-\frac{(\lambda-1)^{r-1}(\lambda+r-1)}{\lambda^r}.
\]
因此有限临界点为 \(\lambda=1\)（阶 \(r-1\)）与 \(\lambda=1-r\)（简单临界点）：
在 \(\lambda=1\) 附近令 \(\mu=\lambda-1\)，则
\[
f_r(1+\mu)-1=-\frac{\mu^r}{(1+\mu)^{r-1}}=-\mu^r+O(\mu^{r+1}),
\]
故 \(y=1\) 的分歧型为 \((r)\)。
在 \(\lambda=1-r\) 处，
\[
y_r=f_r(1-r)=1+\frac{r^r}{(r-1)^{r-1}},
\]
且临界点简单，故对应换位型 \((2)(1^{r-2})\)。

对 \(y=\infty\)，\(\lambda=0\) 为阶 \(r-1\) 极点，\(\lambda=\infty\) 为简单极点（因 \(f_r(\lambda)\sim-\lambda\)），故分歧型为 \((r-1)(1)\)。
因此仅有三处分歧值。

由分歧型可选取分支循环 \(\sigma_1,\sigma_{y_r},\sigma_\infty\in S_r\) 满足
\[
\sigma_1\sim (r),\qquad
\sigma_{y_r}\sim (2)(1^{r-2}),\qquad
\sigma_\infty\sim (r-1)(1),\qquad
\sigma_1\sigma_{y_r}\sigma_\infty=1.
\]
群 \(\langle \sigma_1,\sigma_{y_r}\rangle\) 含 \(r\)-循环与换位，故为 \(S_r\)。
因此几何单值群为 \(S_r\)。
单值作用传递即等价于 \(\Pi_r\) 在 \(\CC(y)[\lambda]\) 上不可约。证毕。
\end{proof}

\begin{theorem}[判别式的闭式分解与双有限分歧点]\label{thm:xi-terminal-belyi-discriminant-closed-form}
\leanverified{paper\_xi\_terminal\_belyi\_discriminant\_closed\_form}
在定理 \ref{thm:xi-terminal-belyi-full-symmetric-cover} 的记号下，存在符号常数 \(\varepsilon_r\in\{\pm1\}\) 使
\[
\mathrm{Disc}_{\lambda}(\Pi_r)(y)
=\varepsilon_r\,(y-1)^{r-1}
\Bigl((r-1)^{r-1}y-\bigl((r-1)^{r-1}+r^r\bigr)\Bigr).
\]
特别地
\[
\mathrm{Disc}_{\lambda}(\Pi_r)(y)=0
\iff
y\in\left\{1,\ 1+\frac{r^r}{(r-1)^{r-1}}\right\},
\]
且在 \(y=1\) 的消失阶为 \(r-1\)，在 \(y=y_r\) 的消失阶为 \(1\)。
\end{theorem}
\begin{proof}
由判别式恒等式
\(
\mathrm{Disc}(P)=(-1)^{r(r-1)/2}\mathrm{Res}(P,P')
\)，只需判定 \(P=\Pi_r(\cdot,y)\) 与 \(P'\) 的公共根位置与重数。

若 \(y=1\)，则 \(P(\lambda)=(\lambda-1)^r\)，故判别式在 \(y=1\) 至少有 \((r-1)\) 重零点。
若 \(y\neq1\)，由
\[
P(\lambda)=0,\qquad
P'(\lambda)=r(\lambda-1)^{r-1}+(y-1)(r-1)\lambda^{r-2}=0
\]
消去 \((\lambda-1)^{r-1}\) 可得唯一候选 \(\lambda=1-r\)，再代回得 \(y=y_r\)。
且该临界点简单，故对应判别式一次消失。
因此判别式除符号常数外恰为所示两因子乘积。证毕。
\end{proof}

% AUTO-GENERATED by scripts/exp_xi_belyi_full_symmetric_puiseux_audit.py
\begin{equation}\label{eq:xi_terminal_belyi_family_discriminant_closed}
\mathrm{Disc}_{\lambda}(\Pi_r)(y)=\varepsilon_r\,(y-1)^{r-1}\Bigl((r-1)^{r-1}y-\bigl((r-1)^{r-1}+r^r\bigr)\Bigr),\quad \varepsilon_r\in\{\pm1\}.
\end{equation}

\begin{equation}\label{eq:xi_terminal_belyi_r3_recurrence_discriminant}
\begin{aligned}
\Pi_3(\lambda,y)&=(\lambda-1)^3+(y-1)\lambda^2=\lambda^3+(y-4)\lambda^2+3\lambda-1,\\
Z_{n+3}(y)&+(y-4)Z_{n+2}(y)+3Z_{n+1}(y)-Z_n(y)=0,\\
Z_0(y)&=3,\qquad Z_1(y)=4 - y,\qquad Z_2(y)=y^{2} - 8 y + 10,\\
\mathrm{Disc}_{\lambda}(\Pi_3)(y)&=(y-1)^2(4y-31).
\end{aligned}
\end{equation}

% AUTO-GENERATED by scripts/exp_xi_belyi_full_symmetric_puiseux_audit.py
\begin{table}[H]
\centering
\scriptsize
\setlength{\tabcolsep}{5pt}
\caption{Belyi 三点分歧族 \(\Pi_r\)（\(2\le r\le 9\)）的判别式因子、局部分歧与置换生成审计。}
\label{tab:xi_belyi_branch_discriminant_group_audit_r2_r9}
\begin{tabular}{r l r c c c r c}
\toprule
$r$ & $y_r$ & $\varepsilon_r$ & Disc & $e_{y_r}=2$ & $\sigma_\infty$ type & $|\langle\sigma_1,\sigma_{y_r}\rangle|$ & $=r!$\\
\midrule
2 & $5$ & 1 & 1 & 1 & $(1)(1)$ & 2 & 1\\
3 & $31/4$ & 1 & 1 & 1 & $(2)(1)$ & 6 & 1\\
4 & $283/27$ & -1 & 1 & 1 & $(3)(1)$ & 24 & 1\\
5 & $3381/256$ & -1 & 1 & 1 & $(4)(1)$ & 120 & 1\\
6 & $49781/3125$ & 1 & 1 & 1 & $(5)(1)$ & 720 & 1\\
7 & $870199/46656$ & 1 & 1 & 1 & $(6)(1)$ & 5040 & 1\\
8 & $17600759/823543$ & -1 & 1 & 1 & $(7)(1)$ & 40320 & 1\\
9 & $404197705/16777216$ & -1 & 1 & 1 & $(8)(1)$ & 362880 & 1\\
\bottomrule
\end{tabular}
\end{table}


\paragraph{一类六次 \texorpdfstring{$j$}{j}-覆叠的判别式指纹、分歧值与 \texorpdfstring{$S_6$}{S6} 正则性}
定义有理函数
\[
j(a):=-\frac{4096\,(a^2-a+1)^3}{(a-1)^2(16a^3-13a^2-78a+43)}\ \in\ \QQ(a),
\]
并对超越参数 \(t\) 令
\[
P_t(a):=-4096\,(a^2-a+1)^3-t\,(a-1)^2(16a^3-13a^2-78a+43)\ \in\ \ZZ[t][a],
\]
即 \(P_t(a)=0\) 等价于 \(j(a)=t\) 的交叉相乘形式。记
\(
\Delta(t):=\Disc_a(P_t)\in\ZZ[t]
\)。

\begin{theorem}[判别式全因子分解、分歧值代数指纹与满对称单值群]\label{thm:xi-j-sextic-discriminant-s6-regular}
以上记号下成立：
\begin{enumerate}
  \item 判别式存在完全显式分解
  \[
  \boxed{\;
  \Delta(t)=2^{43}\,79^3\cdot t^{4}\,(t-1728)^{3}\,\bigl(t^2+1862t+161792\bigr).\;}
  \]
  因而 \(j:\PP^1_a\to\PP^1_t\) 的有限分歧值恰为
  \[
  t\in\left\{0,\ 1728,\ -931\pm 89\sqrt{89}\right\}.
  \]
  \item 令 \(K=\QQ(t)\)，并令 \(L\) 为 \(P_t(a)\) 在 \(K\) 上的分裂域，则该扩张为正则且
  \[
  \boxed{\;\Gal(L/K)\ \cong\ S_6.\;}
  \]
  \item \(j(a)\) 在 \(\overline{\QQ}\) 上不可分解为非平凡复合 \(j=F\circ G\)（\(\deg F,\deg G\ge2\)），并且满足 \(j(\phi(a))\equiv j(a)\) 的 \(\phi\in\mathrm{PGL}_2(\overline{\QQ})\) 只能为恒等映射。
  \item \(L/K\) 的唯一二次子域由判别式平方根生成：
  \[
  \boxed{\;L^{A_6}=\QQ\bigl(t,\sqrt{\Delta(t)}\bigr).\;}
  \]
  在平方类 \(\QQ(t)^\times/(\QQ(t)^\times)^2\) 中，
  \[
  \sqrt{\Delta(t)}\ \sim\ \sqrt{158\,(t-1728)\,\bigl(t^2+1862t+161792\bigr)},
  \]
  从而除 \(t\in\{0,1728,\infty\}\) 外的代数分歧由二次域 \(\QQ(\sqrt{89})\) 控制。
\end{enumerate}
\end{theorem}
\begin{proof}
第 1 条等价于对线性笔 \(P_t=f+t g\) 消去公共根条件 \(P_t=\partial_aP_t=0\)，可由判别式或结式计算给出；二次因子判别式为 \(1862^2-4\cdot161792=4\cdot 89^3\)，故其根为 \(-931\pm 89\sqrt{89}\)。
\par
对第 2 条，正则性由 \(j\) 为 \(\PP^1\) 上有理覆叠且常数域不扩张得到。取 \(t_0=12345\)，则 \(P_{t_0}\in\ZZ[a]\) 为可分离六次多项式；对不整除 \(\Disc(P_{t_0})\) 的素数 \(p\)，其模 \(p\) 因子型给出 Frobenius 的循环类型。直接检验可取 \(p=13\) 使 \(P_{t_0}\) 在 \(\FF_{13}[a]\) 上不可约（给出 \(6\)-循环），并可取 \(p=149\) 使其因子型为 \((2)(1)(1)(1)(1)\)（给出换位）；含 \(6\)-循环与换位的传递子群只能为 \(S_6\)，故某一特化的 Galois 群为 \(S_6\)。由 Hilbert 不可约性，函数域 Galois 群亦为 \(S_6\)。
\par
第 3 条：非平凡分解等价于存在非平凡中间覆叠，亦即单值作用非本原；但 \(S_6\) 的自然作用本原，故不可分解。甲板变换群同构于单值群在自然置换表示中的中心化子；\(S_6\) 的中心化子平凡，故仅有恒等自同构。
\par
第 4 条：\(S_6\) 的唯一指数 \(2\) 正规子群为 \(A_6\)，对应的二次子域由符号特征所给；对任意六次可分离多项式，该二次子域由 \(\sqrt{\Delta(t)}\) 生成。代入第 1 条分解并在平方类中消去平方因子 \(t^4\) 与常数平方部分即得所示平方类代表。证毕。
\end{proof}

\paragraph{判别式二次子域的椭圆化、属 \(5\) 正规化与 \(A_6\) 基变压缩}
本段固定定理 \ref{thm:xi-j-sextic-discriminant-s6-regular} 的六次有理覆叠 \(j:\PP^1_a\to\PP^1_t\) 与其判别式因子
\[
f(t):=(t-1728)\bigl(t^2+1862t+161792\bigr)\in\ZZ[t].
\]

\begin{theorem}[判别式二次子域的椭圆化与 \(2\)-挠点分裂域]\label{thm:xi-j-sextic-discriminant-quadratic-subfield-ellipticization}
沿用定理 \ref{thm:xi-j-sextic-discriminant-s6-regular} 的记号。
令
\[
E_0:\quad y^2=f(t)=(t-1728)\bigl(t^2+1862t+161792\bigr)
=t^3+134t^2-3055744t-279576576,
\]
并令
\[
E_\Delta:\quad Y^2=158\,f(t).
\]
则：
\begin{enumerate}
  \item 二次子域 \(\QQ(t,\sqrt{\Delta(t)})\) 与 \(E_\Delta\) 的函数域同构：
  \[
  \boxed{\ \QQ\bigl(t,\sqrt{\Delta(t)}\bigr)\ \cong\ \QQ(E_\Delta)\ }.
  \]
  \item \(E_0\) 是 \(E_\Delta\) 的二次扭曲；二者的 \(2\)-挠点的 \(t\)-坐标集合完全一致，恰为 \(f(t)=0\) 的根。
  \item \(E_0(\QQ)\) 恰有一个非平凡 \(\QQ\)-有理 \(2\)-挠点 \((1728,0)\)，并且
  \[
  E_0(\QQ)[2]\cong \ZZ/2\ZZ,\qquad
  E_0\bigl(\QQ(\sqrt{89})\bigr)[2]\cong (\ZZ/2\ZZ)^2,
  \]
  且额外两枚有限 \(2\)-挠点的 \(t\)-坐标为
  \[
  \boxed{\ t=-931\pm 89\sqrt{89}\ }.
  \]
  \item 若以 Weierstrass 形式 \(E_0:y^2=t^3+134t^2-3055744t-279576576\) 表示，则其判别式满足
  \[
  \boxed{\ \Delta(E_0)=2^{20}\cdot 89^{3}\cdot 223^{4}\ }.
  \]
\end{enumerate}
\end{theorem}
\begin{proof}
由定理 \ref{thm:xi-j-sextic-discriminant-s6-regular}，
\[
\Delta(t)=2^{43}\,79^3\cdot t^{4}\,(t-1728)^{3}\,\bigl(t^2+1862t+161792\bigr).
\]
因此在 \(\QQ(t)^\times/(\QQ(t)^\times)^2\) 中有
\[
\sqrt{\Delta(t)}\ \sim\ \sqrt{(2\cdot 79)\,(t-1728)\bigl(t^2+1862t+161792\bigr)}
=\sqrt{158\,f(t)}.
\]
从而 \(\QQ(t,\sqrt{\Delta(t)})=\QQ(t,\sqrt{158\,f(t)})\cong \QQ(E_\Delta)\)，得第 (1) 条。
第 (2) 条由定义直接成立：把方程 \(y^2=f(t)\) 乘以常数 \(158\) 得到 \(E_\Delta\)，扭曲不改变根集合 \(f(t)=0\)，故 \(2\)-挠点的 \(t\)-坐标不变。
\par
对第 (3) 条，\(2\)-挠点等价于右端三次多项式的根（令 \(y=0\)）。
由于 \(f(t)=(t-1728)(t^2+1862t+161792)\)，显然存在且仅存在一个 \(\QQ\)-有理根 \(t=1728\)，对应唯一非平凡 \(\QQ\)-有理 \(2\)-挠点。
二次因子判别式为
\[
1862^2-4\cdot 161792=4\cdot 89^3,
\]
其平方自由部分为 \(89\)，故其两根属于且仅属于 \(\QQ(\sqrt{89})\)，并且显式为 \(-931\pm 89\sqrt{89}\)。
\par
第 (4) 条为对该整系数 Weierstrass 模型的判别式直接计算与整数因子分解。证毕。
\end{proof}

\begin{theorem}[判别式椭圆曲线的有理 \(2\)-挠点平移与显式 Möbius 反演]\label{thm:xi-j-discriminant-elliptic-2torsion-mobius}
沿用定理 \ref{thm:xi-j-sextic-discriminant-quadratic-subfield-ellipticization} 的记号，并在 \(E_0\) 上记其有理 \(2\)-挠点
\[
Q:=(1728,0)\in E_0(\QQ).
\]
令 \(\tau_Q:E_0\to E_0\) 为按群律加 \(Q\) 的平移自同构。
令 \(u:=t-1728\)，则 \(E_0\) 可改写为
\[
E_0:\quad y^2=u(u^2+Au+B),
\qquad
A:=5318,\ \ B:=6365312,
\]
并且 \(\tau_Q\) 在仿射坐标上满足完全显式的变换律
\[
\boxed{\;
\tau_Q:\ (u,y)\longmapsto \left(\frac{B}{u},\ -\frac{B}{u^2}\,y\right). \;}
\]
等价地，在 \(t\)-坐标上存在 Möbius 反演 \(\phi\in\mathrm{PGL}_2(\QQ)\) 使得
\[
\boxed{\;
\tau_Q:\ (t,y)\longmapsto\left(\phi(t),\ -\frac{B}{(t-1728)^2}\,y\right),\qquad
\phi(t):=1728+\frac{B}{t-1728}=\frac{1728t+3379328}{t-1728}. \;}
\]
并且：
\begin{enumerate}
  \item \(\tau_Q\) 为阶 \(2\) 自同构，且在 \(E_0\) 上无不动点；
  \item \(\phi\) 交换 \(t=1728\) 与 \(t=\infty\)，并精确交换二次分歧值
  \[
  t=-931\pm 89\sqrt{89},
  \qquad\text{即}\qquad
  \phi(-931+89\sqrt{89})=-931-89\sqrt{89}.
  \]
\end{enumerate}
\leanverified{paper\_xi\_j\_discriminant\_elliptic\_2torsion\_mobius}
\end{theorem}
\begin{proof}
把 \(t=u+1728\) 代回 \(E_0\) 的方程并展开即得
\(
y^2=u(u^2+Au+B)
\)
其中 \(A=5318,B=6365312\)，并且 \(Q\) 对应 \((u,y)=(0,0)\)。
对一般 Weierstrass 形式 \(y^2=x(x^2+Ax+B)\) 上的 \(2\)-挠点 \((0,0)\)，由群律加法公式可验证
\[
x(P+(0,0))=\frac{B}{x(P)},\qquad
y(P+(0,0))=-\frac{B}{x(P)^2}\,y(P),
\]
从而得到所示 \((u,y)\) 变换律；换回 \(t=u+1728\) 即得 \(\phi\) 与 \(y\) 的变换律。
\par
由 \(\phi(\phi(t))=t\) 立得 \(\tau_Q^2=\mathrm{id}\)。
又椭圆曲线上的平移 \(P\mapsto P+Q\) 若有不动点，则必有 \(P=P+Q\)，从而 \(Q=O\)，与 \(Q\neq O\) 矛盾，故无不动点。
\par
最后，定理 \ref{thm:xi-j-sextic-discriminant-quadratic-subfield-ellipticization} 已给出额外两枚 \(2\)-挠点的 \(t\)-坐标为 \(t=-931\pm 89\sqrt{89}\)。
由于 \(\tau_Q\) 置换非平凡 \(2\)-挠点集合且 \(\phi\) 为阶 \(2\) 的 Möbius 变换，故其在这两点上要么逐点固定、要么交换。
而 \(\phi(t)=t\) 的不动点需满足 \((t-1728)^2=B\)，与上述二点的最小多项式 \(t^2+1862t+161792\) 不相容，因而必为交换。证毕。
\end{proof}

\begin{proposition}[Möbius 反演的刚性刻画与二次分歧因子的协变律]\label{prop:xi-j-discriminant-mobius-rigidity-covariance}
沿用定理 \ref{thm:xi-j-discriminant-elliptic-2torsion-mobius} 的记号，并令
\(
Q(t):=t^2+1862t+161792
\)。
则：
\begin{enumerate}
  \item \(\phi\) 是 \(\mathrm{PGL}_2(\QQ)\) 中唯一满足
  \[
  \phi(1728)=\infty,\qquad \phi(\infty)=1728,\qquad
  \phi(-931+89\sqrt{89})=-931-89\sqrt{89}
  \]
  的变换；
  \item 成立严格协变恒等式
  \[
  \boxed{\ \phi(t)-1728=\frac{B}{t-1728}\ },
  \qquad
  \boxed{\ Q(\phi(t))=\frac{B\,Q(t)}{(t-1728)^2}\ }.
  \]
  因而 \(\phi\) 在平方类意义下保持 \((t-1728)Q(t)\) 所确定的二次扩张结构，并与定理 \ref{thm:xi-j-discriminant-elliptic-2torsion-mobius} 的 \(y\)-变换共同给出 \(E_0\) 的自同构。
\end{enumerate}
\leanverified{paper\_xi\_j\_discriminant\_mobius\_rigidity\_covariance}
\end{proposition}
\begin{proof}
第 (1) 条：Möbius 变换由三个点的像唯一确定。
由 \(\phi(1728)=\infty\) 可写 \(\phi(t)=(\alpha t+\beta)/(t-1728)\)；
由 \(\phi(\infty)=1728\) 得 \(\alpha=1728\)；
再由 \(\phi(-931+89\sqrt{89})=-931-89\sqrt{89}\) 解得 \(\beta=3379328\)，从而唯一性成立。
\par
第 (2) 条第一式由定义直接给出。
第二式为代数恒等式：把 \(\phi(t)=1728+\frac{B}{t-1728}\) 代入 \(Q(\cdot)\) 并整理即可得到
\(
Q(\phi(t))(t-1728)^2=BQ(t)
\)。
最后一句由 \(f(t)=(t-1728)Q(t)\) 与上式立刻推出。证毕。
\end{proof}

\begin{theorem}[核为 \(\langle Q\rangle\) 的 \(2\)-同源与商曲线的 Vélu 显式公式]\label{thm:xi-j-discriminant-2isogeny-velu}
\leanverified{paper\_xi\_j\_discriminant\_2isogeny\_velu}
在定理 \ref{thm:xi-j-discriminant-elliptic-2torsion-mobius} 的 \((u,y)\) 坐标下，令 \(Q=(0,0)\)。
则存在次数 \(2\) 的同源
\[
\pi:E_0\longrightarrow E':=E_0/\langle Q\rangle
\]
其在仿射坐标上由 Vélu 公式给出为
\[
\boxed{\ x'=u+\frac{B}{u},\qquad y'=y\left(1-\frac{B}{u^2}\right)\ },
\]
并且目标椭圆曲线可取为 Weierstrass 模型
\[
\boxed{\ E':\quad y'^2=x'\bigl(x'^2-2Ax'+(A^2-4B)\bigr)\ }.
\]
此外 \(\ker(\pi)=\{O,Q\}\)，从而 \(\deg(\pi)=2\)。
\end{theorem}
\begin{proof}
对核 \(G=\{O,Q\}\) 的 Vélu 公式给出
\[
x'(P)=x(P)+\sum_{R\in G\setminus\{O\}}\bigl(x(P+R)-x(R)\bigr),\qquad
y'(P)=y(P)+\sum_{R\in G\setminus\{O\}}\bigl(y(P+R)-y(R)\bigr).
\]
此处 \(x(Q)=0,y(Q)=0\)，并且由定理 \ref{thm:xi-j-discriminant-elliptic-2torsion-mobius} 的加法计算 \(x(P+Q)=B/x(P)\)、\(y(P+Q)=-B\,y(P)/x(P)^2\)，代回即得
\(
x'=u+B/u
\)
与
\(
y'=y(1-B/u^2)
\)。
将 \((x',y')\) 代入并化简可验证其满足所示 \(E'\) 方程（这也是 Vélu 公式的标准输出形式：对 \(y^2=x^3+Ax^2+Bx\)，核为 \((0,0)\) 的 \(2\)-同源商曲线为 \(y^2=x^3-2Ax^2+(A^2-4B)x\)）。
核与次数为椭圆曲线同源的基本性质。证毕。
\end{proof}

\begin{corollary}[平移自同构的固定子域与无分歧商的生成元]\label{cor:xi-j-discriminant-translation-fixed-field-quotient}
\leanverified{paper\_xi\_j\_discriminant\_translation\_fixed\_field\_quotient}
沿用定理 \ref{thm:xi-j-discriminant-elliptic-2torsion-mobius} 与定理 \ref{thm:xi-j-discriminant-2isogeny-velu} 的记号。
则函数域固定子域满足
\[
\boxed{\ \QQ(E_0)^{\langle\tau_Q\rangle}=\QQ(E')\ },
\]
并且 \(\QQ(E_0)^{\langle\tau_Q\rangle}\) 由
\(
x'=u+B/u
\) 与
\(
y'=y(1-B/u^2)
\)
生成，其唯一代数关系为 \(E'\) 的 Weierstrass 方程。
\end{corollary}
\begin{proof}
由定理 \ref{thm:xi-j-discriminant-elliptic-2torsion-mobius}，\(\tau_Q\) 为无不动点的阶 \(2\) 自同构，因此商曲线 \(E_0/\langle\tau_Q\rangle\) 仍为属 \(1\) 曲线，并与 \(E_0/\langle Q\rangle=E'\) 同构。
另一方面，\(\tau_Q\) 作用下 \(u\mapsto B/u\) 使得 \(x'=u+B/u\) 不变，且
\(
y(1-B/u^2)
\)
同样不变，故 \(\QQ(x',y')\subseteq \QQ(E_0)^{\langle\tau_Q\rangle}\)。
由定理 \ref{thm:xi-j-discriminant-2isogeny-velu}，\(\pi\) 为次数 \(2\) 映射，故 \([\QQ(E_0):\QQ(x',y')]=2\)，从而固定子域只能是 \(\QQ(x',y')\)。证毕。
\end{proof}

\begin{theorem}[\(2\)-同源对的判别式互换律与 \(j\)-不变量闭式]\label{thm:xi-j-discriminant-2isogeny-discriminant-j-swap}
令 \(k\) 为特征不为 \(2\) 的域，并考虑含显式 \(2\)-挠点的椭圆曲线族
\[
E_{A,B}:\quad y^2=x(x^2+Ax+B),
\]
以及按核 \(\langle(0,0)\rangle\) 商得到的 \(2\)-同源曲线
\[
E'_{A,B}:\quad y^2=x\bigl(x^2-2Ax+(A^2-4B)\bigr).
\]
则其 Weierstrass 判别式满足
\[
\boxed{\ \Delta(E_{A,B})=2^4\,B^2\,(A^2-4B)\ },
\qquad
\boxed{\ \Delta(E'_{A,B})=2^8\,B\,(A^2-4B)^2\ }.
\]
并且 \(j\)-不变量具有闭式
\[
\boxed{\ j(E_{A,B})=256\,\frac{(A^2-3B)^3}{B^2(A^2-4B)}\ },
\qquad
\boxed{\ j(E'_{A,B})=16\,\frac{(A^2+12B)^3}{B(A^2-4B)^2}\ }.
\]
因此在所有奇素数处，\(B\) 与 \(A^2-4B\) 的 \(p\)-进赋值在同源对的判别式指数中以 \((2,1)\) 与 \((1,2)\) 的方式互换出现。
\end{theorem}
\begin{proof}
对 \(a_1=a_3=0\) 的 Weierstrass 形式 \(y^2=x^3+a_2x^2+a_4x+a_6\)（此处 \(a_6=0\)），有标准不变量
\[
c_4=16(a_2^2-3a_4),\qquad
\Delta=16a_4^2(a_2^2-4a_4),
\qquad
j=\frac{c_4^3}{\Delta}.
\]
对 \(E_{A,B}\) 取 \(a_2=A,a_4=B\) 得 \(\Delta(E_{A,B})=2^4B^2(A^2-4B)\) 与 \(j(E_{A,B})\) 的闭式。
对 \(E'_{A,B}\) 取 \(a_2=-2A,a_4=A^2-4B\) 得 \(\Delta(E'_{A,B})=2^8B(A^2-4B)^2\) 与 \(j(E'_{A,B})\) 的闭式。
最后一条为对上述两式在奇素数处的指数读取。证毕。
\end{proof}

\paragraph{倍映射在 \texorpdfstring{$t$}{t}-线上诱导的 Latt\`es 自映射：闭式、临界分解与 Möbius 对称}
在判别式椭圆曲线 \(E_0:y^2=f(t)\) 的 \(t\)-坐标中，倍映射 \([2]\) 诱导一个次数 \(4\) 的有理自映射。
其后临界集由 \(2\)-挠点的 \(t\)-坐标完全控制，并且 \(E_0(\QQ)[2]\) 的平移在 \(t\)-线上实现为 Klein 四元反演族。

\begin{theorem}[倍映射诱导的 \texorpdfstring{$t$}{t}-线 Latt\`es 映射的闭式与后临界集]\label{thm:xi-j-sextic-elliptic-lattes-t-doubling-pcf}
取定理 \ref{thm:xi-j-sextic-discriminant-quadratic-subfield-ellipticization} 中的椭圆曲线
\[
E_0:\quad y^2=(t-1728)\bigl(t^2+1862t+161792\bigr)
=t^3+134t^2-3055744t-279576576.
\]
令 \(\mathcal L:\PP^1_t\dashrightarrow \PP^1_t\) 为由倍映射 \([2]:E_0\to E_0\) 在 \(t\)-坐标下诱导的有理自映射，即对一般点 \(P\in E_0\) 定义
\(
\mathcal L(t(P)):=t([2]P)
\)。
则 \(\deg(\mathcal L)=4\)，并且具有完全显式表达式
\[
\boxed{\;
\mathcal L(t)=
\frac{t^{4}+6111488\,t^{2}+2236612608\,t+9487424438272}
{4\,(t-1728)\,(t^{2}+1862t+161792)}\; }.
\]
其导数满足可约分解
\[
\boxed{\;
\mathcal L'(t)=
\frac{(t^{2}-3456t-3379328)\,
(t^{4}+3724t^{3}+970752t^{2}+10348004864t-8393927966720)}
{4\,(t-1728)^{2}\,(t^{2}+1862t+161792)^{2}}\; }.
\]
从而 \(\mathcal L\) 的有限临界值集合严格等于
\[
\boxed{\ \{1728,\ \alpha,\ \beta\}\ },
\qquad
\alpha,\beta\ \text{为}\ t^{2}+1862t+161792=0\ \text{的两根}.
\]
并且后临界集合为
\[
\boxed{\ \mathrm{PostCrit}(\mathcal L)=\{\infty,\,1728,\,\alpha,\,\beta\}\ },
\]
且所有临界轨道在至多两步内进入不动点 \(\infty\)。
\end{theorem}
\begin{proof}
将 \(E_0\) 写成 Weierstrass 形式 \(y^2=t^3+a_2t^2+a_4t+a_6\)，其中
\(
a_2=134,\ a_4=-3055744,\ a_6=-279576576
\)。
对一般点 \(P=(t,y)\)，切线斜率为
\[
m=\frac{3t^{2}+2a_2t+a_4}{2y}=\frac{3t^{2}+268t-3055744}{2y},
\]
倍点公式给出
\[
t([2]P)=m^{2}-a_2-2t.
\]
将 \(y^2=(t-1728)(t^2+1862t+161792)\) 代入并化简，得到所示 \(\mathcal L(t)\)。
对 \(\mathcal L\) 求导并因式分解得到 \(\mathcal L'(t)\) 的分解式。
\par
二次因子 \(t^{2}-3456t-3379328\) 的任一根 \(\xi\) 满足 \(\xi^{2}=3456\xi+3379328\)，代回 \(\mathcal L(t)\) 可直接化简得 \(\mathcal L(\xi)=1728\)，故 \(1728\) 为临界值。
对四次因子 \(g_4(t)=t^{4}+3724t^{3}+970752t^{2}+10348004864t-8393927966720\)，
以消去恒等式
\[
\mathrm{Res}_{t}\!\bigl(g_4(t),\,\mathrm{num}(\mathcal L)(t)-z\,\mathrm{den}(\mathcal L)(t)\bigr)
=\mathrm{const}\cdot(z^{2}+1862z+161792)^{2}
\]
刻画其像：任意 \(g_4(\eta)=0\) 的根 \(\eta\) 的像必落在 \(\{\alpha,\beta\}\)。
由于 \(\mathcal L\) 在 \(t=1728,\alpha,\beta\) 处有极点，故 \(\mathcal L(1728)=\mathcal L(\alpha)=\mathcal L(\beta)=\infty\)。
又由 \(\deg\mathrm{num}=4\)、\(\deg\mathrm{den}=3\) 得 \(\mathcal L(t)\sim \frac14 t\)（\(t\to\infty\)），从而 \(\mathcal L(\infty)=\infty\)。
综上得到后临界集合与“至多两步入 \(\infty\)”的断言。证毕。
\end{proof}
\leanverified{paper\_xi\_j\_sextic\_elliptic\_lattes\_t\_doubling\_pcf}

\begin{theorem}[\(\QQ\)-有理 \(2\)-挠平移诱导的降阶因子化与 \texorpdfstring{$\sqrt{89}$}{sqrt(89)} 临界值]\label{thm:xi-j-sextic-elliptic-lattes-degree-drop-by-2torsion}
\leanverified{paper\_xi\_j\_sextic\_elliptic\_lattes\_degree\_drop\_by\_2torsion}
仍取 \(E_0\) 与 \(\mathcal L\) 如定理 \ref{thm:xi-j-sextic-elliptic-lattes-t-doubling-pcf}。
记 \(\QQ\)-有理 \(2\)-挠点 \(Q=(1728,0)\in E_0(\QQ)\)，并令 \(\tau_Q:E_0\to E_0\) 为平移 \(P\mapsto P+Q\)。
则存在 involution \(\phi\in\mathrm{PGL}_2(\QQ)\) 使得对一般点 \(P\) 有 \(t(\tau_Q(P))=\phi(t(P))\)，并且
\[
\boxed{\ \phi(t)=1728+\frac{B}{t-1728}\ },
\qquad
B:=(1728-\alpha)(1728-\beta)=6365312.
\]
进一步，倍映射满足强不变性
\[
\boxed{\ \mathcal L\circ \phi=\mathcal L\ }.
\]
令 \(u:=t-1728\) 并定义 \(\phi\)-不变函数
\[
\pi(t):=u+\frac{B}{u},
\qquad (\deg\pi=2),
\]
则存在唯一的次数 \(2\) 有理函数 \(F\) 使得
\[
\boxed{\ \mathcal L = F\circ \pi\ },
\qquad
\boxed{\ F(x)=\frac{x^{2}+6912x+11296768}{4(x+5318)}\ }.
\]
并且
\[
\mathrm{Crit}(F)=\{-5318\pm 178\sqrt{89}\},
\qquad
F\bigl(\mathrm{Crit}(F)\bigr)=\{-931\pm 89\sqrt{89}\}=\{\alpha,\beta\}.
\]
\end{theorem}
\begin{proof}
由于 \(Q\in E_0[2]\)，对任意 \(P\in E_0\) 有
\(
[2](P+Q)=[2]P
\)。
取 \(t\)-坐标即得 \(\mathcal L(t(P+Q))=\mathcal L(t(P))\)，从而 \(\mathcal L\circ\phi=\mathcal L\)。
又 \(E_0\) 的 \(2\)-挠点的 \(t\)-坐标为 \(\{1728,\alpha,\beta\}\)，而 \(t=1728\) 对应的平移在 \(\PP^1_t\) 上为以 \(1728\) 为不动点的分式线性反演；
由对称性与归一化条件（交换 \(\infty\leftrightarrow 1728\) 且保持交比）得到所示 \(\phi\)，其中
\(
B=(1728-\alpha)(1728-\beta)
\)。
\par
\(\phi\) 的不变有理函数域为 \(\QQ(\pi(t))\)，故 \(\mathcal L\) 必经由商映射 \(\pi\) 因子化：存在 \(F\) 使 \(\mathcal L=F\circ\pi\)。
将 \(t=u+1728\) 并代入 \(\pi=u+B/u\)，对 \(\mathcal L(u+1728)\) 作消元化简得到所示 \(F\)；
其导数与判别式给出临界点为 \(-5318\pm 178\sqrt{89}\)，代回 \(F\) 得临界值为 \(-931\pm 89\sqrt{89}\)。
证毕。
\end{proof}

\begin{theorem}[\texorpdfstring{$E_0[2]$}{E0[2]} 的 Klein 平移群在 \texorpdfstring{$t$}{t}-线上实现为 Möbius 反演族]\label{thm:xi-j-sextic-elliptic-lattes-klein-mobius}
\leanverified{paper\_xi\_j\_sextic\_elliptic\_lattes\_klein\_mobius}
令 \(e_1:=1728\)，并令 \(e_2,e_3\) 为二次方程 \(t^{2}+1862t+161792=0\) 的两根（即 \(e_{2,3}=-931\pm 89\sqrt{89}\)）。
对 \(\{i,j,k\}=\{1,2,3\}\) 定义
\[
\phi_i(t):=e_i+\frac{(e_i-e_j)(e_i-e_k)}{t-e_i}.
\]
则：
\begin{enumerate}
  \item \(\phi_i\) 为二阶 Möbius 变换，且在四点集 \(\{\infty,e_1,e_2,e_3\}\) 上的置换为
  \[
  \phi_i:\ (\infty\ e_i)(e_j\ e_k),
  \]
  从而 \(\langle \phi_1,\phi_2,\phi_3\rangle\cong (\ZZ/2\ZZ)^2\) 且 \(\phi_i\circ\phi_j=\phi_k\)（\(i\neq j\)）。
  \item 该 Klein 四元群等同于 \(E_0[2]\) 的平移群在 \(t\)-线上诱导的作用：\(\phi_i\) 对应平移 \(\tau_{(e_i,0)}\)。
  \item 对每个 \(i\in\{1,2,3\}\) 有
  \[
  \boxed{\ \mathcal L\circ \phi_i=\mathcal L\ },
  \]
  因而 \(\mathrm{Aut}_{\overline{\QQ}}(\mathcal L)\) 至少包含该 Klein 四元群；其中在 \(\QQ\) 上可见的 Möbius 对称子群为 \(\langle \phi_1\rangle\)。
\end{enumerate}
\end{theorem}
\begin{proof}
对 (1)，直接代入验证 \(\phi_i(e_i)=\infty\)、\(\phi_i(\infty)=e_i\) 且 \(\phi_i(e_j)=e_k\)。
由此得到 \(\phi_i^2=\mathrm{id}\) 与乘法关系。
\par
对 (2)，在模型 \(y^2=(t-e_1)(t-e_2)(t-e_3)\) 上，\(2\)-挠点平移在 \(t\)-坐标下的作用为上述反演式，亦可由群律加法公式逐项推导。
\par
对 (3)，对任意 \(T\in E_0[2]\) 都有 \([2](P+T)=[2]P\)，取 \(t\)-坐标即得 \(\mathcal L\circ\phi_i=\mathcal L\)。证毕。
\end{proof}

\begin{theorem}[三分歧值的平方分解：\texorpdfstring{$\mathcal L(t)-v$}{L(t)-v} 的二次平方表示]\label{thm:xi-j-sextic-elliptic-lattes-three-branch-square-factorization}
令
\[
\alpha:=-931-89\sqrt{89},\qquad \beta:=-931+89\sqrt{89},
\qquad
Q(t):=t^2+1862t+161792=(t-\alpha)(t-\beta),
\]
并令 \(\mathcal L\) 为定理 \ref{thm:xi-j-sextic-elliptic-lattes-t-doubling-pcf} 的四次 Latt\`es 自映射。
则在 \(\QQ(\sqrt{89})(t)\) 中成立恒等式
\[
\boxed{\;
\mathcal L(t)-1728=\frac{A(t)^2}{4\,(t-1728)\,Q(t)}\;},
\qquad
A(t):=t^2-3456t-3379328,
\]
\[
\boxed{\;
\mathcal L(t)-\alpha=\frac{B_{\alpha}(t)^2}{4\,(t-1728)\,Q(t)}\;},
\qquad
B_{\alpha}(t):=t^2+(1862+178\sqrt{89})t+161792-307584\sqrt{89},
\]
\[
\boxed{\;
\mathcal L(t)-\beta=\frac{B_{\beta}(t)^2}{4\,(t-1728)\,Q(t)}\;},
\qquad
B_{\beta}(t):=t^2+(1862-178\sqrt{89})t+161792+307584\sqrt{89}.
\]
其中 \(B_{\beta}\) 为 \(B_{\alpha}\) 在 \(\Gal(\QQ(\sqrt{89})/\QQ)\) 下的共轭。
\leanverified{paper\_xi\_j\_sextic\_elliptic\_lattes\_three\_branch\_square\_factorization}
\end{theorem}
\begin{proof}
由定理 \ref{thm:xi-j-sextic-elliptic-lattes-t-doubling-pcf}，可写
\(
\mathcal L(t)=N(t)/D(t)
\)
其中
\(
D(t)=4(t-1728)Q(t)
\)
且 \(N(t)\in\QQ[t]\) 为四次多项式。
三式分别等价于
\[
N(t)-1728D(t)=A(t)^2,\qquad
N(t)-\alpha D(t)=B_{\alpha}(t)^2,\qquad
N(t)-\beta D(t)=B_{\beta}(t)^2,
\]
它们均为 \(\QQ(\sqrt{89})[t]\) 中的多项式恒等式；清分母后直接展开并整理即可验证成立。
最后一句由 \(B_{\beta}\) 的系数是 \(B_{\alpha}\) 的 \(\sqrt{89}\)-共轭得到。证毕。
\end{proof}

\begin{corollary}[分歧型 \texorpdfstring{$(2^2,2^2,2^2)$}{(2^2,2^2,2^2)} 与 \texorpdfstring{$V_4$}{V4}-Galois 性]\label{cor:xi-j-sextic-elliptic-lattes-v4-galois-branching}
在代数闭域 \(\overline{\QQ}\) 上，把 \(\mathcal L\) 视为覆盖
\[
\mathcal L:\PP^1_t\longrightarrow \PP^1_{z},\qquad z=\mathcal L(t).
\]
则：
\begin{enumerate}
\normalfont
  \item 分歧值严格等于 \(\{1728,\alpha,\beta\}\)；
  \item 对每个 \(v\in\{1728,\alpha,\beta\}\)，纤维 \(\mathcal L^{-1}(v)\) 中恰有两点分歧，且分歧指数均为 \(2\)（即每个分歧纤维的循环型为 \((2,2)\)）；
  \item 扩张 \(\overline{\QQ}(t)/\overline{\QQ}(z)\) 为 Galois 扩张，且
  \[
  \boxed{\ \Gal\bigl(\overline{\QQ}(t)/\overline{\QQ}(z)\bigr)\ \cong\ (\ZZ/2\ZZ)^2\ }.
  \]
  其甲板变换群由定理 \ref{thm:xi-j-sextic-elliptic-lattes-klein-mobius} 的三条 Möbius 反演 \(\phi_i\) 生成。
\end{enumerate}
\leanverified{paper\_xi\_j\_sextic\_elliptic\_lattes\_v4\_galois\_branching}
\end{corollary}
\begin{proof}
由定理 \ref{thm:xi-j-sextic-elliptic-lattes-three-branch-square-factorization}，对每个 \(v\in\{1728,\alpha,\beta\}\)，函数 \(\mathcal L(t)-v\) 的零点均为二重零，且由对应二次多项式的两根给出；因此每个分歧纤维由两点组成且分歧指数均为 \(2\)，得到第 (2) 条。
由 Riemann--Hurwitz 公式，四次有理映射的总分歧量为 \(2\deg(\mathcal L)-2=6\)，而三条分歧纤维共贡献 \(3\times 2=6\)，故无其他分歧值，从而得到第 (1) 条。
\par
由定理 \ref{thm:xi-j-sextic-elliptic-lattes-klein-mobius}，\(\mathrm{Aut}_{\overline{\QQ}}(\mathcal L)\) 至少包含阶为 \(4\) 的 Klein 四元群。
另一方面，任意次数 \(4\) 覆盖的甲板变换群阶至多为 \(4\)，故该 Klein 群即为全部甲板变换群，扩张遂为 Galois，且伽罗瓦群同构于 \((\ZZ/2\ZZ)^2\)。证毕。
\end{proof}

\begin{theorem}[特征二次微分的本征拉回：\texorpdfstring{$\mathcal L^\ast q=4q$}{L^*q=4q}]\label{thm:xi-j-sextic-elliptic-lattes-quadratic-differential-eigen}
令
\[
f(t):=(t-1728)\,Q(t)=(t-1728)(t-\alpha)(t-\beta),
\qquad
q(t):=\frac{(\dd t)^2}{f(t)}.
\]
则 \(q\) 为 \(\PP^1_t\) 上的非零亚纯二次微分，其极点至多为单极并且仅可能出现在 \(\{\infty,1728,\alpha,\beta\}\)。
并且成立严格本征关系
\[
\boxed{\ \mathcal L^{\ast}q=4q\ },
\qquad\text{等价地}\qquad
\boxed{\ (\mathcal L'(t))^{2}\,f(t)=4\,f(\mathcal L(t)).\ }
\]
\end{theorem}
\begin{proof}
在椭圆曲线 \(E_0:y^2=f(t)\) 上，\(\omega:=\dd t/y\) 为全纯 \(1\)-形式。
倍映射满足经典恒等式 \([2]^{\ast}\omega=2\omega\)。
另一方面，\(q=\omega^2\) 在超椭圆对合 \(y\mapsto -y\) 下不变，因而可下推到 \(\PP^1_t\simeq E_0/\{\pm 1\}\) 上；
而 \(\mathcal L\) 正是 \([2]\) 在该商上的诱导映射。
因此在 \(\PP^1_t\) 上有
\(
\mathcal L^{\ast}q=([2]^{\ast}\omega)^2=(2\omega)^2=4q
\)。
把 \(q=(\dd t)^2/f(t)\) 代回即可得到所示函数恒等式。证毕。
\end{proof}

\begin{corollary}[迭代导数的望远镜恒等式与周期乘子刚性]\label{cor:xi-j-sextic-elliptic-lattes-multiplier-rigidity}
\leanverified{paper\_xi\_j\_sextic\_elliptic\_lattes\_multiplier\_rigidity}
对任意整数 \(n\ge 1\)，有
\[
\boxed{\ \bigl((\mathcal L^{\circ n})'(t)\bigr)^2
=4^n\,\frac{f(\mathcal L^{\circ n}(t))}{f(t)}\ }.
\]
从而若 \(t_0\in\overline{\QQ}\setminus\{1728,\alpha,\beta\}\) 为周期点（\(\mathcal L^{\circ n}(t_0)=t_0\)），则
\[
\boxed{\ \bigl((\mathcal L^{\circ n})'(t_0)\bigr)^2=4^n\ },
\qquad\text{因而}\qquad
\boxed{\ (\mathcal L^{\circ n})'(t_0)=\pm 2^n\ }.
\]
进一步，取 \(P\in E_0(\overline{\QQ})\) 使 \(t(P)=t_0\)，则
\[
(\mathcal L^{\circ n})'(t_0)=+2^n\iff [2^n]P=P,
\qquad
(\mathcal L^{\circ n})'(t_0)=-2^n\iff [2^n]P=-P.
\]
\end{corollary}
\begin{proof}
由定理 \ref{thm:xi-j-sextic-elliptic-lattes-quadratic-differential-eigen} 迭代得到 \((\mathcal L^{\circ n})^\ast q=4^n q\)，展开即得第一式。
若 \(\mathcal L^{\circ n}(t_0)=t_0\) 且 \(t_0\notin\{1728,\alpha,\beta\}\)，则 \(f(t_0)\neq 0\) 且代入可消去得到 \(((\mathcal L^{\circ n})'(t_0))^2=4^n\)。
\par
对符号判别：令 \(\omega=\dd t/y\)。
由 \([2^n]^\ast\omega=2^n\omega\)，在点 \(P\) 处写出
\[
\frac{\dd(\mathcal L^{\circ n}(t))}{y([2^n]P)}=\frac{2^n\,\dd t}{y(P)}.
\]
注意 \(y([2^n]P)=y(P)\) 当且仅当 \([2^n]P=P\)，而 \(y([2^n]P)=-y(P)\) 当且仅当 \([2^n]P=-P\)。
于是得到 \((\mathcal L^{\circ n})'(t_0)=2^n\,y([2^n]P)/y(P)\in\{\pm 2^n\}\)，并与所述等价条件一致。证毕。
\end{proof}
\input{subsubsec__xi-j-sextic-elliptic-lattes-exact-dynamics}

\begin{corollary}[无抛物周期与 Artin--Mazur \texorpdfstring{$\zeta$}{zeta} 的闭式]\label{cor:xi-j-sextic-elliptic-lattes-artin-mazur-zeta}
\(\mathcal L\) 不存在抛物周期点（即不存在乘子为 \(1\) 的周期点）。
因此对每个 \(n\ge 1\)，方程 \(\mathcal L^{\circ n}(t)=t\) 的全部解均为单根，并且
\[
\boxed{\ \card{\mathrm{Fix}(\mathcal L^{\circ n})}=4^n+1\ }.
\]
从而其 Artin--Mazur 动力学 \(\zeta\)-函数
\[
\zeta_{\mathcal L}(z):=\exp\!\left(\sum_{n\ge 1}\frac{\card{\mathrm{Fix}(\mathcal L^{\circ n})}}{n}z^n\right)
\]
满足
\[
\boxed{\ \zeta_{\mathcal L}(z)=\frac{1}{(1-z)(1-4z)}\ }.
\]
\end{corollary}
\begin{proof}
由推论 \ref{cor:xi-j-sextic-elliptic-lattes-multiplier-rigidity}，任意有限周期点的乘子均为 \(\pm 2^n\neq 1\)。
另一方面，由 \(\mathcal L(t)\sim \tfrac14 t\)（\(t\to\infty\)）知 \(\infty\) 为不动点，且在局部坐标 \(s=1/t\) 下有 \(\mathcal L(s)=4s+O(s^2)\)，从而 \((\mathcal L^{\circ n})'(\infty)=4^n\neq 1\)。
故无抛物周期点。
\par
由定理 \ref{thm:xi-j-sextic-elliptic-lattes-fixedpoints-odd-torsion}，
\(\card{\mathrm{Fix}(\mathcal L^{\circ n})}=4^n+1\)。
代入定义并用 \(\sum_{n\ge1}z^n/n=-\log(1-z)\)、\(\sum_{n\ge1}(4z)^n/n=-\log(1-4z)\) 即得闭式。证毕。
\end{proof}

\begin{proposition}[临界点的算术分裂：\texorpdfstring{$\sqrt{2}$}{sqrt2} 与 \texorpdfstring{$\sqrt{89}$}{sqrt89} 的两支来源]\label{prop:xi-j-sextic-elliptic-lattes-critical-arithmetic-splitting}
令 \(\mathrm{Crit}(\mathcal L)\subset\PP^1(\overline{\QQ})\) 为 \(\mathcal L\) 的临界点集合。
则：
\begin{enumerate}
\normalfont
  \item 映到 \(1728\) 的两临界点为
  \[
  \boxed{\ t=1728\pm \sqrt{B}=1728\pm 1784\sqrt2\ }\in \QQ(\sqrt2),
  \]
  它们恰为 \(A(t)=t^2-3456t-3379328\) 的两根；
  \item 映到 \(\alpha\) 与 \(\beta\) 的四临界点分别由 \(B_{\alpha}(t)\) 与 \(B_{\beta}(t)\) 的零点给出；
  因而该部分的系数域为 \(\QQ(\sqrt{89})\)，并在 \(\Gal(\QQ(\sqrt{89})/\QQ)\) 作用下成对交换。
\end{enumerate}
\end{proposition}
\begin{proof}
由定理 \ref{thm:xi-j-sextic-elliptic-lattes-t-doubling-pcf} 的导数分解，\(\mathcal L^{-1}(1728)\) 的分歧点正是二次因子 \(A(t)\) 的零点。
又 \(A(t)=(t-1728)^2-B\) 且 \(B=6365312=2\cdot 1784^2\)，故得到第 (1) 条。
第 (2) 条由定理 \ref{thm:xi-j-sextic-elliptic-lattes-three-branch-square-factorization} 直接推出：\(\mathcal L(t)-\alpha\) 与 \(\mathcal L(t)-\beta\) 的零点分别为 \(B_{\alpha}\)、\(B_{\beta}\) 的二重零点。证毕。
\end{proof}

\begin{theorem}[在 \texorpdfstring{$\QQ(\sqrt{89})$}{Q(sqrt89)} 上的 Belyi 归一化与 passport]\label{thm:xi-j-sextic-elliptic-lattes-belyi-normalization}
令 \(K:=\QQ(\sqrt{89})\)，并在 \(K\) 上定义 Möbius 变换
\[
\psi(z):=\frac{1728-\beta}{1728-\alpha}\cdot\frac{z-\alpha}{z-\beta},
\]
则 \(\psi(\alpha)=0\)、\(\psi(1728)=1\)、\(\psi(\beta)=\infty\)。
令
\[
\mathcal B(t):=\psi(\mathcal L(t))\in K(t).
\]
则 \(\mathcal B:\PP^1_t\to\PP^1\) 为定义在 \(K\) 上的 Belyi 映射，并且其在 \(0,1,\infty\) 上的分歧型均为 \((2,2)\)，亦即其 passport 为
\[
\boxed{\ (2^2,\ 2^2,\ 2^2)\ }.
\]
此外，\(\mathcal B\) 为 Galois 覆盖，其甲板变换群同构于 \((\ZZ/2\ZZ)^2\)（与推论 \ref{cor:xi-j-sextic-elliptic-lattes-v4-galois-branching} 的群一致）。
\end{theorem}
\begin{proof}
由推论 \ref{cor:xi-j-sextic-elliptic-lattes-v4-galois-branching}，\(\mathcal L\) 的分歧值为 \(\{1728,\alpha,\beta\}\)，且三处分歧纤维循环型皆为 \((2,2)\)。
将值域以 \(\psi\) 归一化后，分歧值变为 \(\{1,0,\infty\}\)，从而 \(\mathcal B\) 为 Belyi 映射且 passport 如所示。
甲板变换群由 \(\mathcal L\) 的甲板变换群通过共轭识别得到，故仍同构于 \((\ZZ/2\ZZ)^2\)。证毕。
\end{proof}
\leanverified{paper\_xi\_j\_sextic\_elliptic\_lattes\_belyi\_normalization}

\endinput


\begin{theorem}[交错分歧的平方消去与属 \(5\) 超椭圆正规化]\label{thm:xi-j-sextic-fiber-product-genus5-hyperelliptic-model}
令 \(j(a)\) 如定理 \ref{thm:xi-j-sextic-discriminant-s6-regular}。
记
\[
r(a):=16a^3-13a^2-78a+43,
\qquad
H(a):=2048a^8-8752a^7+16455a^6-5030a^5-39667a^4+82096a^3-74559a^2+33406a-5869.
\]
并定义曲线
\[
\boxed{\ C:\quad Y^2=-2^{17}\,r(a)\,H(a)\ }.
\]
则：
\begin{enumerate}
  \item \(C\) 为光滑超椭圆曲线，且
  \[
  \boxed{\ g(C)=5\ }.
  \]
  \item 令 \(\pi:E_\Delta\to\PP^1_t\) 为定理 \ref{thm:xi-j-sextic-discriminant-quadratic-subfield-ellipticization} 中的二次覆叠（以 \(t\) 为坐标），则 \(C\) 与纤维积
  \[
  \PP^1_a\times_{\PP^1_t}E_\Delta
  \qquad(t=j(a))
  \]
  的正规化在双有理意义下等价。
\end{enumerate}
\end{theorem}
\begin{proof}
令
\[
s(a):=8a^3+15a^2-66a+35,\qquad p(a):=2a^2-9a-1.
\]
直接代入并约分可得两条恒等式：
\[
j(a)-1728=-\frac{64\,s(a)^2}{(a-1)^2\,r(a)},
\qquad
j(a)^2+1862j(a)+161792=\frac{2048\,p(a)^2\,H(a)}{(a-1)^4\,r(a)^2}.
\]
从而
\[
158\,(j(a)-1728)\bigl(j(a)^2+1862j(a)+161792\bigr)
\;=\;
-\frac{2^{17}\cdot 79\, s(a)^2p(a)^2\,H(a)}{(a-1)^6\,r(a)^3}.
\]
在函数域层面，把右端乘以平方因子 \(\bigl((a-1)^3r(a)s(a)p(a)\bigr)^{-2}\) 不改变二次扩张，于是得到与之同一二次扩张的模型
\[
Y^2=-2^{17}\,r(a)\,H(a),
\]
即为所定义的 \(C\)。
\par
又因 \(\gcd(r,H)=1\) 且 \(rH\) 平方自由（等价于分支多项式无重根），并且 \(\deg(rH)=11\)，故 \(C\) 为属
\[
g(C)=\left\lfloor\frac{11-1}{2}\right\rfloor=5
\]
的光滑超椭圆曲线。
最后，按构造 \(E_\Delta\) 的函数域为 \(\QQ(t,\sqrt{158f(t)})\)；把 \(t=j(a)\) 代入并消去平方因子后得到的正规化恰为上述 \(C\)，从而与纤维积正规化双有理等价。证毕。
\end{proof}

\begin{theorem}[属 \(5\) 超椭圆曲线到判别式椭圆曲线的显式六次覆叠]\label{thm:xi-j-sextic-genus5-to-elliptic-explicit-degree6-cover}
沿用定理 \ref{thm:xi-j-sextic-fiber-product-genus5-hyperelliptic-model} 的记号，并令
\[
E_0:\quad y^2=f(t)=(t-1728)\bigl(t^2+1862t+161792\bigr).
\]
定义多项式
\[
s(a):=8a^3+15a^2-66a+35,\qquad p(a):=2a^2-9a-1.
\]
则存在一个在 \(\QQ\) 上定义的非平凡态射
\[
g_0:C\longrightarrow E_0
\]
由下式完全显式给出：
\[
\boxed{\quad t=j(a),\qquad y=\frac{s(a)p(a)}{(a-1)^3\,r(a)^2}\,Y.\quad}
\]
并且 \(g_0\) 为次数 \(6\) 的有限覆叠，即 \(\deg(g_0)=6\)。
\end{theorem}
\begin{proof}
由定理 \ref{thm:xi-j-sextic-fiber-product-genus5-hyperelliptic-model} 证明中给出的恒等式
\[
j(a)-1728=-\frac{64\,s(a)^2}{(a-1)^2\,r(a)},
\qquad
j(a)^2+1862j(a)+161792=\frac{2048\,p(a)^2\,H(a)}{(a-1)^4\,r(a)^2}
\]
可得
\[
\bigl(j(a)-1728\bigr)\bigl(j(a)^2+1862j(a)+161792\bigr)
\,=\,
-\frac{2^{17}\,s(a)^2p(a)^2\,H(a)}{(a-1)^6\,r(a)^3}.
\]
另一方面，曲线 \(C\) 满足 \(Y^2=-2^{17}r(a)H(a)\)，从而
\[
\left(\frac{s(a)p(a)}{(a-1)^3\,r(a)^2}Y\right)^2
=
-\frac{2^{17}\,s(a)^2p(a)^2\,H(a)}{(a-1)^6\,r(a)^3}
=
f\bigl(j(a)\bigr).
\]
于是 \((t,y)=(j(a),\frac{s(a)p(a)}{(a-1)^3r(a)^2}Y)\) 确满足 \(E_0\) 的方程，故 \(g_0\) 的值落于 \(E_0\)。

对次数断言，\(t=j(a)\) 作为从 \(\PP^1_a\) 到 \(\PP^1_t\) 的有理函数具有次数 \(6\)（见定理 \ref{thm:xi-j-sextic-discriminant-s6-regular}），故在函数域层面
\(
[\QQ(a):\QQ(t)]=6
\)。
又 \(\QQ(C)=\QQ(a,Y)\) 为 \(\QQ(a)\) 的二次扩张，而 \(\QQ(E_0)=\QQ(t,y)\) 为 \(\QQ(t)\) 的二次扩张。
由上式 \(y\in \QQ(a,Y)\) 得到嵌入 \(\QQ(E_0)\hookrightarrow \QQ(C)\)，并且
\[
[\QQ(C):\QQ(t)]=2\,[\QQ(a):\QQ(t)]=12.
\]
因此
\(
[\QQ(C):\QQ(E_0)]=12/2=6
\)，从而 \(\deg(g_0)=6\)。证毕。
\end{proof}

\begin{theorem}[不变量微分的拉回消去]\label{thm:xi-j-sextic-invariant-differential-pullback-quartic-factor}
令
\[
\omega_{E_0}:=\frac{\dd t}{y}\in H^0(E_0,\Omega^1_{E_0})
\]
为 \(E_0\) 的标准不变量微分。
则对定理 \ref{thm:xi-j-sextic-genus5-to-elliptic-explicit-degree6-cover} 的态射 \(g_0:C\to E_0\)，在 \(C\) 上有完全显式的恒等式
\[
\boxed{\quad g_0^*(\omega_{E_0})=-4096\,(a^2-a+1)^2\,\frac{\dd a}{Y}.\quad}
\]
\end{theorem}
\begin{proof}
由 \(t=j(a)\) 得 \(\dd t=j'(a)\dd a\)。
另一方面，定理 \ref{thm:xi-j-sextic-genus5-to-elliptic-explicit-degree6-cover} 给出
\(
y=\frac{s(a)p(a)}{(a-1)^3r(a)^2}Y
\)，从而
\[
g_0^*\!\left(\frac{\dd t}{y}\right)
=
j'(a)\,\frac{(a-1)^3r(a)^2}{s(a)p(a)}\,\frac{\dd a}{Y}.
\]
对 \(j(a)=-4096\frac{(a^2-a+1)^3}{(a-1)^2r(a)}\) 直接求导并化简（可符号化简逐式复核）得到恒等式
\[
j'(a)=-4096\,(a^2-a+1)^2\,\frac{s(a)p(a)}{(a-1)^3\,r(a)^2}.
\]
代入即得所示拉回公式。证毕。
\end{proof}

\begin{corollary}[分歧除子与有效 theta 特征]\label{cor:xi-j-sextic-ramification-divisor-theta-characteristic}
令 \(\infty\in C\) 为超椭圆模型的无穷远点（因 \(\deg(rH)=11\) 为奇数而唯一）。
令 \(\alpha,\beta\) 为二次方程 \(a^2-a+1=0\) 的两根，并记其在 \(C\) 上的两点纤维为
\[
P_\alpha^\pm=(\alpha,\pm Y_\alpha),\qquad P_\beta^\pm=(\beta,\pm Y_\beta).
\]
则 \(g_0:C\to E_0\) 的分歧除子满足
\[
\boxed{\quad R_{g_0}=2\bigl(P_\alpha^++P_\alpha^-+P_\beta^++P_\beta^-\bigr),\qquad \deg R_{g_0}=8.\quad}
\]
并且该除子类与正则类一致：
\[
K_C\sim R_{g_0}.
\]
因此
\[
\boxed{\quad D:=P_\alpha^++P_\alpha^-+P_\beta^++P_\beta^- \ \text{给出一个有效 theta 特征：}\ 2D\sim K_C.\quad}
\]
\end{corollary}
\begin{proof}
在椭圆曲线 \(E_0\) 上，\(\omega_{E_0}\) 无零无极点，故 \(\Div(\omega_{E_0})=0\)。
一般覆叠理论给出
\(
\Div(g_0^*\omega_{E_0})=R_{g_0}
\)。
由定理 \ref{thm:xi-j-sextic-invariant-differential-pullback-quartic-factor}，
\[
\Div(g_0^*\omega_{E_0})
=
\Div\!\bigl((a^2-a+1)^2\bigr)+\Div\!\left(\frac{\dd a}{Y}\right).
\]
对超椭圆曲线 \(Y^2=f(a)\) 且 \(\deg f=2g+1\) 的标准事实给出
\(
\Div(\dd a/Y)=(2g-2)\,\infty
\)；此处 \(g=5\)，故 \(\Div(\dd a/Y)=8\,\infty\)。
又函数 \(a\) 在 \(\infty\) 处具有二阶极点，故 \((a^2-a+1)^2\) 在 \(\infty\) 处具有八阶极点，并在 \(a=\alpha,\beta\) 处分别给出二重零点，从而
\[
\Div\!\bigl((a^2-a+1)^2\bigr)
=
2\bigl(P_\alpha^++P_\alpha^-+P_\beta^++P_\beta^-\bigr)-8\,\infty.
\]
相加得到
\(
\Div(g_0^*\omega_{E_0})=2(P_\alpha^++P_\alpha^-+P_\beta^++P_\beta^-)
\)，即为 \(R_{g_0}\)。
又 \(K_{E_0}\sim 0\)，由 Riemann--Hurwitz 得 \(K_C\sim g_0^*K_{E_0}+R_{g_0}\sim R_{g_0}\)。证毕。
\end{proof}

\begin{proposition}[Jacobian 的显式椭圆因子与 \(\QQ\)-有理幂等元]\label{prop:xi-j-sextic-jacobian-elliptic-factor-idempotent}
设 \(J(C)\) 为 \(C\) 的 Jacobian。令
\[
g_{0*}:J(C)\to E_0,
\qquad
g_0^*:E_0\to J(C)
\]
分别为由 \(g_0\) 诱导的 push-forward 与 pull-back。
则同态满足
\[
\boxed{\quad g_{0*}\circ g_0^*=[6]_{E_0}.\quad}
\]
从而
\[
\boxed{\quad e:=\frac{1}{6}\,g_0^*\circ g_{0*}\ \in\ \End^0(J(C))\ \text{为幂等元： }e^2=e.\quad}
\]
并由此得到在 \(\QQ\) 上定义的同源分解
\[
\boxed{\quad J(C)\ \sim_{\QQ}\ g_0^*(E_0)\times P,\qquad P:=\ker(g_{0*})^0,\qquad \dim P=4.\quad}
\]
\end{proposition}
\begin{proof}
对任意曲线间有限态射 \(h\) 都有 \(h_*\circ h^*=[\deg h]\)。
应用于 \(h=g_0\) 且 \(\deg(g_0)=6\)，得到 \(g_{0*}g_0^*=[6]\)。
于是
\[
e^2=\frac{1}{36}g_0^*g_{0*}g_0^*g_{0*}
=\frac{1}{36}g_0^*[6]g_{0*}
=\frac{1}{6}g_0^*g_{0*}=e.
\]
由幂等元给出的像与核的同源分裂得到所示分解；维数断言由 \(\dim J(C)=g(C)=5\) 与 \(\dim E_0=1\) 得出。证毕。
\end{proof}

\begin{proposition}[良好约化处的 zeta 因子整除]\label{prop:xi-j-sextic-zeta-factor-divisibility}
令 \(S\) 为 \(C\) 与 \(E_0\) 的潜在坏素数集合（可取分支多项式判别式与椭圆判别式的素因子并集）。
对任意素数 \(\ell\notin S\)，在 \(\FF_\ell\) 上将 \((C,E_0,g_0)\) 约化得到 \((C_\ell,E_{0,\ell},g_{0,\ell})\)。
则在 \(\ell\)-进 \'{e}tale 上同调层面，拉回
\[
g_{0,\ell}^*:H^1_{\mathrm{\acute et}}(E_{0,\overline{\FF}_\ell},\QQ_\ell)\hookrightarrow H^1_{\mathrm{\acute et}}(C_{\overline{\FF}_\ell},\QQ_\ell)
\]
为单射且与 Frobenius 交换。
因此 \(C_\ell\) 的 Frobenius 特征多项式包含 \(E_{0,\ell}\) 的特征多项式因子；等价地，\(\zeta(C_\ell,T)\) 的分子多项式在 \(\ZZ[T]\) 中被 \(\zeta(E_{0,\ell},T)\) 的分子多项式整除。
\end{proposition}
\begin{proof}
在良好约化处，态射 \(g_{0,\ell}\) 诱导的 \(g_{0,\ell*}\circ g_{0,\ell}^*=[6]\) 在 \'{e}tale 上同调中仍成立。
从而 \(g_{0,\ell}^*\) 在 \(H^1\) 上为单射，并与 Frobenius 作用交换。
线性代数直接推出特征多项式的整除性。证毕。
\end{proof}

\begin{proposition}[分歧点的定义域与有理分歧除子]\label{prop:xi-j-sextic-branchpoints-field-and-rational-branch-divisor}
令 \(O_{E_0}\) 为 \(E_0\) 的无穷远点。
则 \(g_0:C\to E_0\) 的分歧点在 \(E_0\) 上恰为
\[
Q_\pm=(0,\pm\sqrt{f(0)})=(0,\pm 768\sqrt{-474}),
\]
并且
\(
\QQ(Q_\pm)=\QQ(\sqrt{-474})
\)，且 \(Q_-=\overline{Q_+}\)。
尽管两点各自不在 \(\QQ\) 上有理，分歧除子
\[
B:=Q_++Q_-
\]
却在 \(\QQ\) 上有理，且满足线性等价
\[
\boxed{\quad B\sim 2O_{E_0}.\quad}
\]
此外，推论 \ref{cor:xi-j-sextic-ramification-divisor-theta-characteristic} 中的四个分歧点按符号分裂：对每个 \(\alpha\in\{\text{\(a^2-a+1=0\) 的根}\}\)，两点 \(P_\alpha^\pm\) 分别映到 \(Q_\pm\)。
\end{proposition}
\begin{proof}
由定理 \ref{thm:xi-j-sextic-genus5-to-elliptic-explicit-degree6-cover} 的显式公式，
当 \(a^2-a+1=0\) 时有 \(t=j(a)=0\)，且 \(y=\frac{s(a)p(a)}{(a-1)^3r(a)^2}Y\) 随 \(Y\mapsto -Y\) 变号，故两点 \(P_\alpha^\pm\) 映到 \(t=0\) 纤维上的两点 \(Q_\pm\)。
直接计算
\(
f(0)=-279576576=-768^2\cdot 474
\)，得 \(Q_\pm\) 的定义域断言。
\par
函数 \(t\) 作为 \(E_0\) 的坐标函数在 \(O_{E_0}\) 处有二阶极点，且在 \(t=0\) 处有两简单零点，因此
\[
\Div(t)=Q_++Q_- - 2O_{E_0},
\]
从而 \(B\sim 2O_{E_0}\) 且 \(B\) 在 \(\QQ\) 上有理。证毕。
\end{proof}

\begin{conjecture}[四维因子的实乘结构与良好素数处分解型]\label{conj:xi-j-sextic-prym-real-multiplication-sqrt89}
令 \(P:=\ker(g_{0*})^0\) 为命题 \ref{prop:xi-j-sextic-jacobian-elliptic-factor-idempotent} 的四维因子。
猜想 \(\End^0_{\QQ}(P)\) 含有实二次域 \(\QQ(\sqrt{89})\) 的嵌入。
并且对任意良好约化素数 \(\ell\notin S\)，若 \(\bigl(\frac{89}{\ell}\bigr)=-1\)（\(\ell\) 在 \(\QQ(\sqrt{89})\) 中惰性），则 \(P_\ell\) 的 Frobenius 特征多项式在 \(\ZZ[T]\) 上通常不可约；
若 \(\bigl(\frac{89}{\ell}\bigr)=+1\)（\(\ell\) 分裂），则该特征多项式通常分裂为两个次数 \(4\) 的互为 Galois 共轭因子。
\end{conjecture}

\begin{theorem}[二次基变的奇置换消去：几何单值群从 \(S_6\) 压缩到 \(A_6\)，且仅余两处分歧]\label{thm:xi-j-sextic-basechange-compress-s6-to-a6-two-branchpoints}
令 \(f:\PP^1_a\to\PP^1_t\) 为 \(t=j(a)\) 给出的次数 \(6\) 覆叠，其几何单值群为 \(S_6\)（定理 \ref{thm:xi-j-sextic-discriminant-s6-regular}）。
令 \(\pi:E_\Delta\to\PP^1_t\) 为 \(E_\Delta:Y^2=158f(t)\) 给出的二次覆叠，并令 \(C\) 为定理 \ref{thm:xi-j-sextic-fiber-product-genus5-hyperelliptic-model} 的正规化纤维积。
则诱导得到次数 \(6\) 映射
\[
g_\Delta:C\to E_\Delta.
\]
并且：
\begin{enumerate}
  \item \(g_\Delta\) 的分歧点在 \(E_\Delta\) 上恰有两点，互为椭圆反演；在每个分歧点之上的置换型为 \((3)(3)\)，每点的分歧贡献为 \(4\)。
  \item 因而 \(g(C)=5\)（与定理 \ref{thm:xi-j-sextic-fiber-product-genus5-hyperelliptic-model} 的属数一致）。
  \item \(g_\Delta\) 的几何单值群包含于 \(A_6\)，且其 \(E_\Delta\) 上的 Galois 闭包的 Galois 群为 \(A_6\)。
\end{enumerate}
\end{theorem}
\begin{proof}
由定理 \ref{thm:xi-j-sextic-discriminant-s6-regular} 的判别式指数，
次数 \(6\) 覆叠 \(f\) 在五个分歧值处的总分歧量分别为：
\[
t=0:\ 4,\qquad
t=1728:\ 3,\qquad
t=-931\pm 89\sqrt{89}:\ 1,\qquad
t=\infty:\ 1.
\]
这与置换型
\[
t=0:\ (3)(3),\quad
t=1728:\ (2)(2)(2),\quad
t=-931\pm 89\sqrt{89}:\ (2),\quad
t=\infty:\ (2)
\]
相容，其中除 \(t=0\) 外均为奇置换。
\par
二次基变 \(\pi\) 的分歧恰发生在 \(t=1728\)、\(t=-931\pm 89\sqrt{89}\) 与 \(\infty\)。
在这些点上，小回路在基变后提升为绕两圈的回路，从而局部单值置换被替换为原置换的平方。
对换位与 \((2)(2)(2)\) 皆有平方为恒等，故上述四处分歧在拉回后被消去。
相反，\(t=0\) 处 \(\pi\) 不分歧，而局部单值为阶 \(3\) 的偶置换 \((3)(3)\)，故该分歧保留；
又因 \(\pi\) 在 \(t=0\) 处上方有两点，遂 \(g_\Delta\) 的分歧点恰为这两点且置换型仍为 \((3)(3)\)，得第 (1) 条。
\par
第 (2) 条由 Riemann--Hurwitz：
由于 \(g(E_\Delta)=1\)，并且两处分歧各贡献 \(4\)，故
\[
2g(C)-2=8\quad\Rightarrow\quad g(C)=5.
\]
\par
第 (3) 条：二次基变对应于把 \(\pi_1(\PP^1_t\setminus\{\text{分歧值}\})\) 限制到符号同态 \(\mathrm{sgn}:S_6\to\{\pm 1\}\) 的核，
等价于剔除奇单值贡献；因此拉回后的几何单值群包含于 \(A_6\)。
又因原群为 \(S_6\)，且 \(A_6\) 为其唯一指数 \(2\) 正规子群，故 \(E_\Delta\) 上的 Galois 闭包群为 \(A_6\)。证毕。
\end{proof}

\begin{proposition}[一组完全显式的分歧循环：一个 Nielsen 类代表]\label{prop:xi-j-sextic-explicit-branch-cycles-nielsen-rep}
存在对次数 \(6\) 纤维的标号，使得绕五个分歧值 \(t=0,1728,-931\pm 89\sqrt{89},\infty\) 的几何单值置换可取为下列五元组（右作用、从左到右相乘）：
\[
\boxed{
\begin{aligned}
\sigma_0&=(1,3,4)(2,6,5),&
\sigma_{1728}&=(1,6)(2,5)(3,4),\\
\sigma_{+}&=(1,5),&
\sigma_{-}&=(1,6),&
\sigma_{\infty}&=(1,3),
\end{aligned}}
\qquad
\boxed{\ \sigma_0\,\sigma_{1728}\,\sigma_{+}\,\sigma_{-}\,\sigma_{\infty}=1\ }.
\]
并且它们生成 \(S_6\)。
\end{proposition}
\begin{proof}
乘积恒等式可直接在 \(S_6\) 中验证。
对生成性，令
\(
c:=(1,4,5,2,3,6)
\)
为群 \(G:=\langle\sigma_0,\sigma_{1728},\sigma_{+},\sigma_{-},\sigma_{\infty}\rangle\) 中的一个 \(6\)-循环（由生成元显式相乘得到）。
取换位 \(\tau:=\sigma_{\infty}=(1,3)\in G\)。
则 \(c^k\tau c^{-k}\in G\) 给出若干换位；再结合 \(\sigma_-=(1,6)\) 可得到一族换位，其支撑边集合在顶点 \(\{1,\dots,6\}\) 上连通。
由“连通图的边换位生成整个对称群”得 \(G=S_6\)。证毕。
\end{proof}

\begin{corollary}[\(A_6\) 正则闭包曲线的属数与两处分歧]\label{cor:xi-j-sextic-a6-galois-closure-genus-241}
令 \(Y\to E_\Delta\) 为定理 \ref{thm:xi-j-sextic-basechange-compress-s6-to-a6-two-branchpoints} 中 \(g_\Delta:C\to E_\Delta\) 的 Galois 闭包。
则 \(\mathrm{Gal}(Y/E_\Delta)\cong A_6\)，并且 \(Y\to E_\Delta\) 仅在 \(E_\Delta\) 上那两点处分歧，每处分歧惯性群阶为 \(3\)。
因此
\[
\boxed{\ g(Y)=241\ }.
\]
\end{corollary}
\begin{proof}
由于 \(g(E_\Delta)=1\)，对 \(A_6\) 正则覆叠应用 Riemann--Hurwitz：
\[
2g(Y)-2
=\sum_{\text{branch }Q}|A_6|\Bigl(1-\frac1{e_Q}\Bigr).
\]
此处 \(|A_6|=360\)，分歧点数为 \(2\)，且 \(e_Q=3\)，故
\[
2g(Y)-2=2\cdot 360\cdot\frac{2}{3}=480,
\]
从而 \(g(Y)=241\)。证毕。
\end{proof}

\begin{theorem}[超椭圆分支多项式的判别式与坏素数集合]\label{thm:xi-j-sextic-hyperelliptic-branch-polynomial-discriminant-factorization}
\leanverified{paper\_xi\_j\_sextic\_hyperelliptic\_branch\_polynomial\_discriminant\_factorization}
沿用定理 \ref{thm:xi-j-sextic-fiber-product-genus5-hyperelliptic-model} 的 \(r(a)\)、\(H(a)\)，令
\[
P(a):=r(a)H(a)\in\ZZ[a].
\]
则其判别式与结式满足完全显式分解：
\[
\boxed{\ \Disc(r)=2^{6}\cdot 79^{3}\ },
\qquad
\boxed{\ \Disc(H)=-2^{30}\cdot 79^{9}\cdot 89^{10}\cdot 223^{3}\ },
\]
\[
\boxed{\ \Res(r,H)=2^{27}\cdot 79^{8}\ },
\qquad
\boxed{\ \Disc(P)=-2^{90}\cdot 79^{28}\cdot 89^{10}\cdot 223^{3}\ }.
\]
从而曲线 \(C:Y^2=-2^{17}P(a)\) 的坏素数集合精确包含于 \(\{2,79,89,223\}\)；
并且对任意素数 \(p\nmid 2\cdot 79\cdot 89\cdot 223\)，其分支多项式在 \(\FF_p\) 上仍平方自由，故 \(C\) 在 \(p\) 处具有良好约化（以“分支多项式平方自由”意义）。
\end{theorem}
\begin{proof}
四个分解均可由判别式与结式的标准公式在整数环上直接计算并因子分解得到；相应的可复现证书由脚本 \texttt{scripts/exp\_xi\_j\_sextic\_ellipticization\_a6\_hyperelliptic\_audit.py} 给出（见注记 \ref{rem:xi-j-sextic-ellipticization-audit-path}）。
\par
良好约化断言来自一般事实：若 \(P(a)\) 在 \(\FF_p\) 上平方自由，则二次覆盖 \(Y^2=P(a)\) 的分支在模 \(p\) 下保持分离，从而给出光滑约化；平方自由性在 \(p\) 处失败当且仅当 \(p\mid \Disc(P)\)。证毕。
\end{proof}

\begin{remark}[可复现证书路径]\label{rem:xi-j-sextic-ellipticization-audit-path}
本段的判别式/结式分解、椭圆判别式与分歧循环核验由脚本
\texttt{scripts/exp\_xi\_j\_sextic\_ellipticization\_a6\_hyperelliptic\_audit.py}
生成，输出文件为
\texttt{artifacts/export/xi\_j\_sextic\_ellipticization\_a6\_hyperelliptic\_audit.json}。
\end{remark}

\endinput


\begin{proposition}[两处素数模因子型的充分性：\texorpdfstring{$6$}{6}-循环与换位强制 \texorpdfstring{$S_6$}{S6}]\label{prop:xi-j-sextic-two-primes-force-s6}
取 \(t_0\in\ZZ\) 使 \(P_{t_0}(a)\) 可分离。若存在互异素数 \(p,q\nmid \Disc(P_{t_0})\) 满足
\[
P_{t_0}(a)\bmod p\ \text{在}\ \FF_p[a]\ \text{上不可约},\qquad
P_{t_0}(a)\bmod q\ \text{的因子型为}\ (2)(1)(1)(1)(1),
\]
则 \(\Gal(P_{t_0}/\QQ)\cong S_6\)。
\end{proposition}
\leanverified{paper\_xi\_j\_sextic\_two\_primes\_force\_s6}
\begin{proof}
由 Dedekind 定理（不整除判别式时）模素数因子型给出 Frobenius 循环类型；不可约给出 \(6\)-循环，因子型 \((2)(1^4)\) 给出换位。含 \(6\)-循环与换位的传递子群必为 \(S_6\)。证毕。
\end{proof}

\begin{theorem}[多临界 Puiseux 自由能与 \(n^{-r}\) 零点缩放极限]\label{thm:xi-terminal-belyi-puiseux-scaling}
\leanverified{paper\_xi\_terminal\_belyi\_puiseux\_scaling}
设 \(\lambda_1(y),\dots,\lambda_r(y)\) 为 \(\Pi_r(\lambda,y)=0\) 的根（按重数），并定义
\[
Z_n^{(r)}(y):=\sum_{j=1}^{r}\lambda_j(y)^n,\qquad n\in\NN.
\]
则：
\begin{enumerate}
  \item \(Z_n^{(r)}(y)\in\ZZ[y]\)，且 \(\deg Z_n^{(r)}=n\)；并满足阶 \(r\) 线性递推（系数由 \(\Pi_r\) 的 \(\lambda\)-系数给定，且在 \(y\) 中至多一次）。
  \item 在谱半径不并列区域，极限自由能
  \[
  F_r(y):=\lim_{n\to\infty}\frac{1}{n}\log|Z_n^{(r)}(y)|
  \]
  存在且
  \(
  F_r(y)=\log\rho(y)
  \),
  其中 \(\rho(y):=\max_j|\lambda_j(y)|\)。
  \item 在多临界点 \(y=1\) 附近，存在 \(\kappa_r\neq0\)（\(\kappa_r^r=-1\)）与 \(r\) 个分支，使
  \[
  \lambda_j(y)
  =1+\omega_r^j\kappa_r\,(y-1)^{1/r}
  +O\!\bigl(|y-1|^{2/r}\bigr),
  \qquad
  \omega_r=e^{2\pi i/r},
  \]
  从而边缘指数为 \(1/r\)。
  \item 对任意固定 \(t\in\CC\)，取 \(y=1+t^r/n^r\)，则有局部缩放
  \[
  Z_n^{(r)}\!\left(1+\frac{t^r}{n^r}\right)
  =\sum_{k=0}^{r-1}\exp\!\bigl(\omega_r^k\kappa_r t\bigr)+O(n^{-1}),
  \]
  误差对 \(t\) 的紧集一致。
  等价地，经重参数 \(t\mapsto\kappa_r^{-1}t\) 后，极限核为
  \[
  \Phi_r(t):=\sum_{k=0}^{r-1}\exp(\omega_r^k t).
  \]
\end{enumerate}
\end{theorem}
\begin{proof}
第 1 条由 Newton 幂和恒等式直接成立。
第 2 条在主导根唯一时由标准主模分解
\(
Z_n^{(r)}=\lambda_\ast^n(1+o(1))
\)
得到。

对第 3 条，令 \(\mu=\lambda-1,\ \delta=y-1\)，则
\[
\Pi_r(1+\mu,1+\delta)=\mu^r+\delta(1+\mu)^{r-1}=0.
\]
故主导平衡为 \(\mu^r\sim-\delta\)，由 Puiseux 定理得
\(
\mu=\omega_r^j\kappa_r\delta^{1/r}+O(\delta^{2/r})
\)。
于是边缘指数为 \(1/r\)。

第 4 条由上述展开代入 \(\delta=t^r/n^r\) 并对 \(\lambda_j^n\) 取对数展开得到：
\[
\lambda_j^n
=\exp\!\Bigl(\omega_r^j\kappa_r t+O(n^{-1})\Bigr),
\]
求和即得。证毕。
\end{proof}

% AUTO-GENERATED by scripts/exp_xi_belyi_full_symmetric_puiseux_audit.py
\begin{table}[H]
\centering
\scriptsize
\setlength{\tabcolsep}{6pt}
\caption{局部缩放 \(y=1+t^r/n^r\) 下的数值审计：\(Z_n^{(r)}\) 与 \(\sum_{k=0}^{r-1}\exp(\omega_r^k\kappa_r t)\) 的最大误差（\(\kappa_r^r=-1\)）。}
\label{tab:xi_belyi_scaling_limit_audit}
\begin{tabular}{r r l l}
\toprule
$r$ & $n$ & max error & $n\cdot$max error\\
\midrule
3 & 120 & $0.00026738388$ & $0.032086065$\\
3 & 240 & $0.00013369047$ & $0.032085714$\\
4 & 120 & $0.00010725135$ & $0.012870162$\\
4 & 240 & $5.3477609e-5$ & $0.012834626$\\
5 & 120 & $2.1660776e-5$ & $0.0025992932$\\
5 & 240 & $1.074835e-5$ & $0.0025796041$\\
6 & 120 & $2.9283989e-6$ & $0.00035140787$\\
6 & 240 & $1.4430866e-6$ & $0.00034634079$\\
\bottomrule
\end{tabular}
\end{table}


\begin{proposition}[\(r=3\) 的显式 Lee--Yang 区间与三次尖点指数]\label{prop:xi-terminal-belyi-r3-explicit-leyang-interval}
\leanverified{paper\_xi\_terminal\_belyi\_r3\_explicit\_leyang\_interval}
当 \(r=3\) 时，
\[
\Pi_3(\lambda,y)=(\lambda-1)^3+(y-1)\lambda^2
=\lambda^3+(y-4)\lambda^2+3\lambda-1.
\]
记 \(Z_n(y):=\sum_{j=1}^{3}\lambda_j(y)^n\)，则：
\begin{enumerate}
  \item \(Z_n(y)\in\ZZ[y]\)、\(\deg Z_n=n\)，且
  \[
  Z_{n+3}+(y-4)Z_{n+2}+3Z_{n+1}-Z_n=0,
  \]
  并有初值
  \(
  Z_0=3,\ Z_1=4-y,\ Z_2=y^2-8y+10
  \)（见式 \eqref{eq:xi_terminal_belyi_r3_recurrence_discriminant}）。
  \item 判别式分解为
  \[
  \mathrm{Disc}_{\lambda}(\Pi_3)(y)=(y-1)^2(4y-31).
  \]
  因此仅有两处有限分歧值 \(y=1\) 与 \(y=31/4\)。
  \item 对任意实 \(y\in(1,31/4)\)，\(\Pi_3(\cdot,y)\) 恰有一实根 \(\lambda_0(y)\in(0,1)\) 与一对共轭根
  \(
  \lambda_\pm(y)=\rho(y)e^{\pm i\theta(y)}
  \)（\(\rho(y)>1\)），并满足
  \[
  \lambda_0(y)=\rho(y)^{-2}.
  \]
  因而实轴上的谱半径并列区间恰为 \([1,31/4]\)，其 \(y=1\) 端点边缘指数为 \(1/3\)，零点逼近速率为 \(\Theta(n^{-3})\)。
\end{enumerate}
\end{proposition}
\begin{proof}
第 1 条由定理 \ref{thm:xi-terminal-belyi-puiseux-scaling} 的一般幂和递推在 \(r=3\) 的专门化给出。
第 2 条由三次判别式公式直接代入
\(
a=y-4,\ b=3,\ c=-1
\)
得到。

对第 3 条，\(\mathrm{Disc}<0\) 在 \((1,31/4)\) 上成立，故恰有一实根与一对共轭非实根。
又
\[
\Pi_3(0,y)=-1<0,\qquad
\Pi_3(1,y)=y-1>0,
\]
故唯一实根 \(\lambda_0(y)\in(0,1)\)。
由 Vieta，
\(
\lambda_0\lambda_+\lambda_-=1
\)，
写 \(\lambda_\pm=\rho e^{\pm i\theta}\) 得
\(
\lambda_0\rho^2=1
\)，即 \(\lambda_0=\rho^{-2}\) 且 \(\rho>1\)。
因此区间内由共轭对给出谱半径并列；区间外判别式非负，谱半径为单根。
端点指数与缩放率由定理 \ref{thm:xi-terminal-belyi-puiseux-scaling} 的 \(r=3\) 情形立即得到。证毕。
\end{proof}

\begin{corollary}[可调多临界指数、代数端点显式性与最大 Galois 复杂度]\label{cor:xi-terminal-belyi-adjustable-exponent-maximal-galois}
对每个 \(r\ge2\)，族 \(\Pi_r\) 同时满足：
\begin{enumerate}
  \item 多临界点固定于 \(y=1\)，边缘指数严格为 \(1/r\)，局部缩放核由 \(\Phi_r\) 显式给出；
  \item 唯一另一有限端点 \(y_r=1+\frac{r^r}{(r-1)^{r-1}}\) 可由判别式闭式直接读出；
  \item 几何单值群为 \(S_r\)，故在同阶参数族中实现满对称置换复杂度。
\end{enumerate}
\end{corollary}
\begin{proof}
由定理 \ref{thm:xi-terminal-belyi-full-symmetric-cover}、定理 \ref{thm:xi-terminal-belyi-discriminant-closed-form} 与定理 \ref{thm:xi-terminal-belyi-puiseux-scaling} 逐项合并即得。证毕。
\end{proof}

\begin{conjecture}[有限态代数谱配分族的整数 Puiseux 指数原理]\label{conj:xi-terminal-algebraic-spectrum-integer-puiseux-principle}
设 \(\{Z_n(y)\}_{n\ge0}\subset\ZZ[y]\) 来自有限维转移算子或代数谱曲线：
存在 \(\Pi(\lambda,y)\in\ZZ[\lambda,y]\)，使
\[
Z_n(y)=\sum_j c_j(y)\lambda_j(y)^n,
\]
其中 \(\lambda_j\) 为 \(\Pi(\lambda,y)=0\) 的代数分支。
猜想任一自由能非解析点 \(y_c\) 均由主导根的有限重合并（重数 \(m\ge2\)）触发，且边缘指数恒为 \(1/m\)。
进一步，在无额外对称约束的一般位置下，几何单值群应达到满对称群 \(S_{\deg_\lambda\Pi}\)。
若该猜想成立，则临界定位与指数读取可统一归约为 \(\Pi\) 的判别式/子结式链的符号判别与主导重根阶的抽取。
\end{conjecture}

\paragraph{次数 \texorpdfstring{$11$}{11} 的满对称单值化：循环三重归一化与 \texorpdfstring{$S_{11}$}{S11}}
定义二元多项式
\[
E_t(\mu):=
2t^3(\mu^2+1)^2(\mu^2-\mu-1)^2+\mu^5(3\mu^2-2\mu-1)^3\in\ZZ[t,\mu].
\]

\begin{theorem}[一族 \texorpdfstring{$11$}{11} 次特化多项式的全对称 Galois 群]\label{thm:xi-degree11-Et-specialization-galois-S11}
取 \(t=3\)，记 \(f(\mu):=E_3(\mu)\in\ZZ[\mu]\)。
则 \(f\) 在 \(\QQ\) 上不可约，并且
\[
\Gal\bigl(f/\QQ\bigr)\cong S_{11}.
\]
\leanverified{paper\_xi\_degree11\_generic\_galois\_S11\_seeds}
\leanverified{paper\_xi\_degree11\_generic\_galois\_S11\_package}
\leanverified{paper\_xi\_degree11\_et\_specialization\_galois\_s11}
\end{theorem}
\begin{proof}
令 \(n=\deg f=11\)。
\par
\emph{不可约性。}
在素数 \(31\) 下，将 \(f\) 约化并按首项单位化，可得
\[
\overline f(\mu)\equiv \mu^{11}-2\mu^{10}-10\mu^9-5\mu^8-11\mu^7-12\mu^6+4\mu^5+4\mu^3+2\mu^2+4\mu+2\pmod{31}.
\]
该多项式在 \(\FF_{31}[\mu]\) 中不可约（可由有限域因式分解直接核验）。
故由 Gauss 引理与模素不可约判别，\(f\) 在 \(\QQ[\mu]\) 上不可约，因而 \(G:=\Gal(f/\QQ)\subset S_{11}\) 在根集上传递。
\par
\emph{存在 \texorpdfstring{$7$}{7}-循环。}
在素数 \(7\) 下，\(f\) 的约化分解为一次、三次与七次不可约因子之积：
\[
f(\mu)\equiv (\mu-3)\,(\mu^3+\mu^2-\mu+3)\,(\mu^7+2\mu^5+2\mu^3+2\mu-1)\pmod 7.
\]
并且判别式满足
\[
\Disc(f)=2^{22}\cdot 3^{37}\cdot 5^{8}\cdot 23\cdot 2741\cdot 5153\cdot 1315229,
\]
从而 \(7\nmid \Disc(f)\)。
因此 Dedekind 定理给出 Frobenius 置换在根上的循环型与上述分解型一致，
从而 \(G\) 含有循环型 \((7)(3)(1)\) 的元，特别含有一个 \(7\)-循环。
\par
\emph{非交错性。}
\(\Disc(f)\) 非平方数（例如其中素因子 \(23\) 的指数为 \(1\)），故 \(G\nsubseteq A_{11}\)。
\par
\emph{Jordan 逼近。}
由于 \(n=11\) 为素数，传递性蕴含 \(G\) 为本原群。
又 \(G\) 含有 \(p\)-循环（\(p=7\le n-3\)），由 Jordan 定理知 \(A_{11}\subseteq G\)。
结合 \(G\nsubseteq A_{11}\)，得到 \(G=S_{11}\)。证毕。
\end{proof}

\begin{corollary}[泛型 \texorpdfstring{$t$}{t} 的全对称性与不可解性]\label{cor:xi-degree11-Et-generic-galois-S11-hilbert}
令 \(K=\QQ(t)\)，并将 \(E_t(\mu)\) 视作 \(K[\mu]\) 中多项式。
则
\[
\Gal\bigl(E_t(\mu)/\QQ(t)\bigr)\cong S_{11}.
\]
此外，除去一个薄集外，对所有 \(a\in\QQ\) 的可分特化 \(E_a(\mu)\in\QQ[\mu]\) 均满足
\[
\Gal\bigl(E_a(\mu)/\QQ\bigr)\cong S_{11},
\]
从而存在无穷多个整数 \(a\in\ZZ\) 使上述结论成立。
特别地，方程 \(E_t(\mu)=0\) 在 \(\QQ(t)\) 上一般不可由根式求解。
\leanverified{paper\_xi\_degree11\_generic\_galois\_S11\_seeds}
\end{corollary}
\begin{proof}
由定理 \ref{thm:xi-degree11-Et-specialization-galois-S11}，在特化 \(t=3\) 处的伽罗瓦群已为 \(S_{11}\)。
特化同态的基本性质表明：在可分特化点处，特化伽罗瓦群嵌入泛型伽罗瓦群。
因此泛型伽罗瓦群包含 \(S_{11}\)，只能等于 \(S_{11}\)。
\par
Hilbert 不可约定理断言：除去一个薄集外，\(\QQ\) 上的特化保持与泛型相同的伽罗瓦群；
从而得到对几乎所有 \(a\in\QQ\) 的结论，并推出存在无穷多个整数特化仍为 \(S_{11}\)。
最后，由经典伽罗瓦理论，根式可解当且仅当伽罗瓦群可解；而 \(S_{11}\) 不可解，故一般不可由根式求解。证毕。
\end{proof}

\begin{theorem}[判别轨迹的循环三重归一化与属数 \texorpdfstring{$3$}{3}]\label{thm:xi-degree11-cyclic-cubic-normalization-genus3}
设
\[
F(\mu):=-\frac{2\mu(\mu^2-\mu-1)^2}{\mu^2+1}\in\QQ(\mu),
\]
并令函数域扩张
\[
\QQ(\mu)\subset \QQ(\mu,u),\qquad u^3=F(\mu).
\]
则存在唯一（至同构）光滑射影曲线 \(\mathcal C\) 使其函数域为 \(\QQ(\mathcal C)=\QQ(\mu,u)\)。
映射
\[
\pi:\mathcal C\to\PP^1,\qquad (\mu,u)\mapsto \mu
\]
为次数 \(3\) 的循环覆盖，且恰有五个分歧点
\[
\mu=0,\quad \mu=\frac{1\pm\sqrt5}{2},\quad \mu=\pm i,
\]
其分歧指数均为 \(3\)。
因此
\[
g(\mathcal C)=3,
\]
并且 \(\Aut(\mathcal C)\) 含有阶 \(3\) 元 \(u\mapsto \omega u\)（\(\omega^3=1\)）。
\end{theorem}
\begin{proof}
由 \(u^3=F(\mu)\) 可知 \(\pi\) 为 Kummer 型次数 \(3\) 扩张，其 Galois 群由 \(u\mapsto \omega u\) 生成，故为循环覆盖。
分歧由 \(F\) 在 \(\PP^1_\mu\) 上的零极点阶决定：凡零极点阶不被 \(3\) 整除者产生全分歧 \(e=3\)。
\par
由表达式可见：
\(\mu=0\) 为零点阶 \(1\)；
\(\mu=\frac{1\pm\sqrt5}{2}\) 为零点阶 \(2\)；
\(\mu=\pm i\) 为极点阶 \(1\)；
而 \(\mu=\infty\) 处有 \(F(\mu)\sim -2\mu^3\)，零极点阶为 \(3\)，可被 \(3\) 整除，故不分歧。
因此恰有五个全分歧点，各贡献 \(e-1=2\)。
由 Riemann--Hurwitz 公式
\[
2g(\mathcal C)-2=3\cdot(-2)+5\cdot 2
\]
得 \(g(\mathcal C)=3\)。证毕。
\end{proof}

\begin{theorem}[\texorpdfstring{$\mathcal C$}{C} 上自然参数的极点除子与 \texorpdfstring{$\deg(t)=11$}{deg(t)=11}]\label{thm:xi-degree11-parameter-t-divisor-degree}
在定理 \ref{thm:xi-degree11-cyclic-cubic-normalization-genus3} 的 \(\mathcal C\) 上定义亚纯函数
\[
t:=\frac{X(\mu)}{u},
\qquad
X(\mu):=\frac{\mu^2(3\mu^2-2\mu-1)}{\mu^2+1}.
\]
则映射 \(t:\mathcal C\to\PP^1\) 的次数为 \(\deg(t)=11\)，并且其极点除子满足
\[
\div_\infty(t)
=2P_{\alpha_+}+2P_{\alpha_-}+2P_{i}+2P_{-i}+P_{\infty,1}+P_{\infty,2}+P_{\infty,3},
\]
其中 \(\alpha_\pm=\frac{1\pm\sqrt5}{2}\)，\(P_{\alpha_\pm},P_{\pm i}\) 为 \(\pi^{-1}(\alpha_\pm),\pi^{-1}(\pm i)\) 中唯一点（全分歧），而 \(P_{\infty,1},P_{\infty,2},P_{\infty,3}\) 为 \(\pi^{-1}(\infty)\) 中三点（不分歧）。
\end{theorem}
\begin{proof}
在 \(\mu=\alpha_\pm\) 处，\(F(\mu)\sim c(\mu-\alpha_\pm)^2\)。
取局部参数 \(s\) 使 \(\mu-\alpha_\pm=s^3\)，则 \(u\sim c^{1/3}s^2\)，故 \(v_{P_{\alpha_\pm}}(u)=2\)。
又 \(X(\alpha_\pm)\neq 0\)，故 \(v_{P_{\alpha_\pm}}(X)=0\)，从而 \(v_{P_{\alpha_\pm}}(t)=-2\)，给出极点阶 \(2\)。
\par
在 \(\mu=\pm i\) 处，\(\mu^2+1\) 为一阶零，故 \(F\) 为一阶极点并全分歧。
取 \(\mu\mp i=s^3\)，则 \(u\sim c^{1/3}s^{-1}\)，故 \(v_{P_{\pm i}}(u)=-1\)。
另一方面，\(X\) 在 \(\mu=\pm i\) 处为一次极点，在全分歧点上变为三次极点，即 \(v_{P_{\pm i}}(X)=-3\)。
因此 \(v_{P_{\pm i}}(t)=-3-(-1)=-2\)，亦为二阶极点。
\par
在 \(\mu=\infty\) 上方的三点 \(P_{\infty,j}\) 处，\(\pi\) 不分歧且 \(F(\mu)\sim -2\mu^3\)，故 \(u\sim c\mu\)，从而 \(v_{P_{\infty,j}}(u)=-1\)。
又 \(X(\mu)\sim 3\mu^2\)，故 \(v_{P_{\infty,j}}(X)=-2\)，于是 \(v_{P_{\infty,j}}(t)=-1\)，即每点给出一个一次极。
\par
综上得到所述极点除子，其次数为 \(2+2+2+2+1+1+1=11\)，从而 \(\deg(t)=11\)。证毕。
\end{proof}

\begin{corollary}[次数 \texorpdfstring{$11$}{11} 覆盖的满对称算术单值化]\label{cor:xi-degree11-cover-arithmetic-monodromy-S11}
\leanpartial{paper\_xi\_degree11\_cover\_arithmetic\_monodromy\_S11}{The requested Lean signature is false as stated: taking \(G\) to be the trivial group contradicts \(\mathrm{Nonempty}(G \simeq \mathrm{Perm}(\mathrm{Fin}\ 11))\).}
在上述设定下，函数 \(t:\mathcal C\to\PP^1\) 的伽罗瓦闭包之算术单值化群同构于 \(S_{11}\)。
等价地，\(\QQ(\mathcal C)/\QQ(t)\) 的伽罗瓦闭包伽罗瓦群为 \(S_{11}\)。
\end{corollary}
\begin{proof}
由 \(t=X(\mu)/u\) 与 \(u^3=F(\mu)\) 得 \((X(\mu)/t)^3=F(\mu)\)，清分母即得 \(\mu\) 满足方程 \(E_t(\mu)=0\)。
因此 \(\QQ(\mathcal C)=\QQ(t,\mu)\)，并且 \([\QQ(\mathcal C):\QQ(t)]=\deg(t)=11\)（定理 \ref{thm:xi-degree11-parameter-t-divisor-degree}）。
由推论 \ref{cor:xi-degree11-Et-generic-galois-S11-hilbert}，\(\Gal(E_t/\QQ(t))\cong S_{11}\)，从而对应的次数 \(11\) 覆盖之伽罗瓦闭包算术单值化群即为 \(S_{11}\)。证毕。
\end{proof}

\begin{remark}[可复现的有限域因式分解与判别式审计]\label{rem:xi-degree11-Et-audit-script}
定理 \ref{thm:xi-degree11-Et-specialization-galois-S11} 中对模素不可约性、模素分解型与判别式奇偶平方性的核验，可通过符号计算脚本一致复现。
一个实现见 \texttt{scripts/exp\_xi\_degree11\_Et\_S11\_audit.py}。
\end{remark}

\paragraph{平方根谱碰撞与前导零点的 \texorpdfstring{$n^{-2}$}{n^{-2}} 量子化}
令 \(U\subset\CC\) 为开集，\(A:U\to\Mat_{d}(\CC)\) 为解析矩阵族（\(d\ge 2\)）。
取非零线性泛函 \(\ell\in(\CC^d)^\vee\) 与非零向量 \(r\in\CC^d\)，并定义有限尺度配分函数
\[
Z_n(u):=\ell\!\bigl(A(u)^n r\bigr),\qquad n\in\NN.
\]
记特征多项式
\[
P(u,\lambda):=\det(\lambda I-A(u)).
\]

\begin{theorem}[平方根谱碰撞下 Lee--Yang 前导零点的 \texorpdfstring{$n^{-2}$}{n^{-2}} 量子化定标]\label{thm:xi-leyang-square-root-collision-leading-zeros-n2}
\leanverified{paper\_xi\_leyang\_square\_root\_collision\_leading\_zeros\_n2}
设存在 \((u_c,\lambda_c)\in U\times\CC\) 满足：
\begin{enumerate}
  \item \(P(u_c,\lambda_c)=0\) 且 \(\partial_\lambda P(u_c,\lambda_c)=0\)；
  \item \(\partial_{\lambda\lambda}P(u_c,\lambda_c)\neq 0\) 且 \(\partial_u P(u_c,\lambda_c)\neq 0\)；
  \item \(A(u_c)\) 的其余全部特征值 \(\{\lambda_j(u_c)\}_{j\ge 3}\) 满足 \(|\lambda_j(u_c)|<|\lambda_c|\)；
  \item \(\ell,r\) 在该二重特征值对应的二维广义谱簇上诱导的两支谱投影系数不同时为零。
\end{enumerate}
则存在 \(\varepsilon>0\) 及解析函数 \(c_\pm(t)\)（满足 \(c_\pm(0)\neq 0\)），以及定义于 \(|t|<\varepsilon\) 的两支特征值分支 \(\lambda_\pm(t)\)，使得在二重覆盖参数化
\[
u=u_c+t^2
\]
下成立展开
\[
\lambda_\pm(t)=\lambda_c\pm \alpha\,t+O(t^2),
\qquad
\alpha^2=-\frac{2\,\partial_u P(u_c,\lambda_c)}{\partial_{\lambda\lambda}P(u_c,\lambda_c)}\neq 0,
\]
并且存在常数 \(\eta\in(0,|\lambda_c|)\) 使对 \(|t|<\varepsilon\) 一致有
\[
Z_n(u_c+t^2)=c_+(t)\lambda_+(t)^n+c_-(t)\lambda_-(t)^n+O(\eta^n)
\qquad(n\to\infty).
\]
令
\[
\kappa:=-\frac{c_-(0)}{c_+(0)}\in\CC^\times,
\]
并固定一条对数分支 \(\log\) 在某邻域内解析。
则存在常数 \(C>0\) 使得对每个固定整数 \(k\in\ZZ\)，当 \(n\) 足够大时，方程 \(Z_n(u)=0\) 在缩放窗 \(|u-u_c|\le Cn^{-2}\) 内存在唯一零点 \(u_{n,k}\)，且满足
\[
u_{n,k}
=
u_c+\frac{\lambda_c^{2}}{4\alpha^{2}n^{2}}
\bigl(\log\kappa+2\pi i k\bigr)^{2}
+O(n^{-3}),
\qquad(n\to\infty).
\]
\end{theorem}
\begin{proof}
由 (1)(2) 与 Weierstrass 制备定理，\(P(u,\lambda)\) 在 \((u_c,\lambda_c)\) 的邻域内可分解为
\[
P(u,\lambda)=\Bigl((\lambda-\lambda_c)^2+a_1(u)(\lambda-\lambda_c)+a_0(u)\Bigr)\,Q(u,\lambda),
\qquad Q(u_c,\lambda_c)\neq 0,
\]
其中 \(a_0(u_c)=a_1(u_c)=0\)，且 \(a_0'(u_c)\neq 0\)（横截性）。
因此存在二重覆盖 \(u=u_c+t^2\)，使得该二次因子的两根成为解析分支 \(\lambda_\pm(t)\)，并由二次项比较得到
\[
(\lambda-\lambda_c)^2
=
-\frac{2\,\partial_u P(u_c,\lambda_c)}{\partial_{\lambda\lambda}P(u_c,\lambda_c)}(u-u_c)
+O\bigl(|u-u_c|^{3/2}\bigr),
\]
从而 \(\lambda_\pm(t)=\lambda_c\pm\alpha t+O(t^2)\) 且 \(\alpha^2\) 如述。

由解析谱投影与 (3) 的严格谱隙，可将 \(A(u_c+t^2)^n\) 分解为两主支谱投影项与余项：
\[
A(u_c+t^2)^n=\lambda_+(t)^n\Pi_+(t)+\lambda_-(t)^n\Pi_-(t)+R_n(t),
\]
其中 \(\Pi_\pm(t)\) 解析，且 \(\|R_n(t)\|=O(\eta^n)\) 对 \(|t|<\varepsilon\) 一致成立。
左右夹以 \(\ell(\cdot r)\) 得
\(
Z_n(u_c+t^2)=c_+(t)\lambda_+(t)^n+c_-(t)\lambda_-(t)^n+O(\eta^n)
\)，其中 \(c_\pm(t):=\ell(\Pi_\pm(t)r)\)；
由 (4) 得 \(c_\pm(0)\neq 0\)。

令 \(Z_n(u_c+t^2)=0\)。
将余项并入误差后，有
\[
\Bigl(\frac{\lambda_+(t)}{\lambda_-(t)}\Bigr)^n=-\frac{c_-(t)}{c_+(t)}+O\bigl((\eta/|\lambda_c|)^n\bigr).
\]
在 \(|t|<\varepsilon\) 内取对数并吸收分支不定性 \(2\pi i k\) 得
\[
n\log\frac{\lambda_+(t)}{\lambda_-(t)}
=
\log\kappa+2\pi i k+O(t)+O\bigl((\eta/|\lambda_c|)^n\bigr).
\]
由 \(\lambda_\pm(t)=\lambda_c\pm\alpha t+O(t^2)\) 得
\(
\log\frac{\lambda_+(t)}{\lambda_-(t)}=\frac{2\alpha}{\lambda_c}t+O(t^2)
\)。
因此
\[
t=\frac{\lambda_c}{2\alpha n}\bigl(\log\kappa+2\pi i k\bigr)+O(n^{-2}),
\]
进而
\(
u-u_c=t^2=\frac{\lambda_c^2}{4\alpha^2n^2}(\log\kappa+2\pi i k)^2+O(n^{-3})
\)。
唯一性可由上述对数方程在 \(|t|\le c/n\) 区域的收缩估计得到。证毕。
\end{proof}

\begin{corollary}[对称可见性下的两零点外推与幅度消去]\label{cor:xi-leyang-two-leading-zeros-extrapolate-uc}
在定理 \ref{thm:xi-leyang-square-root-collision-leading-zeros-n2} 的假设下，若进一步有 \(c_+(0)=c_-(0)\)，则 \(\kappa=-1\)。
取对数分支满足 \(\log(-1)=i\pi\)，并令
\[
C:=\frac{\pi^2\lambda_c^2}{4\alpha^2}\neq 0.
\]
则当 \(n\to\infty\) 时，靠近 \(u_c\) 的两枚最小尺度零点（对应 \(k=0,1\)）满足
\[
u_{n,0}=u_c-\frac{C}{n^2}+O(n^{-3}),
\qquad
u_{n,1}=u_c-\frac{9C}{n^2}+O(n^{-3}).
\]
从而两零点外推
\[
\widehat u_c(n):=\frac{9u_{n,0}-u_{n,1}}{8}
=u_c+O(n^{-3}),
\]
并且幅度可由
\[
\widehat C(n):=\frac{n^2}{8}\bigl(u_{n,0}-u_{n,1}\bigr)=C+O(n^{-1})
\]
一致估计。
\par
同样地，幅度亦可由两零点外推出的消元式直接读出：
\[
\widehat C_{\mathrm{elim}}(n):=-n^2\bigl(u_{n,0}-\widehat u_c(n)\bigr)=C+O(n^{-1}).
\]
\leanverified{paper\_xi\_leyang\_two\_leading\_zeros\_extrapolate\_uc}
\end{corollary}
\begin{proof}
由定理 \ref{thm:xi-leyang-square-root-collision-leading-zeros-n2} 的量子化公式代入 \(\log\kappa=\log(-1)=i\pi\)，得到
\[
u_{n,k}
=u_c+\frac{\lambda_c^2}{4\alpha^2n^2}\bigl(i\pi+2\pi i k\bigr)^2+O(n^{-3})
=u_c-\frac{\pi^2\lambda_c^2}{4\alpha^2n^2}(2k+1)^2+O(n^{-3}).
\]
取 \(k=0,1\) 即得前两式；线性消元得到 \(\widehat u_c(n)\) 的表达式，\(\widehat C(n)\) 由差分读出。证毕。
\end{proof}

\begin{theorem}[三零点端点外推的 \texorpdfstring{$O(n^{-6})$}{O(n^-6)} 误差改进]\label{thm:xi-leyang-three-leading-zeros-extrapolate-uc-n6}
\leanverified{paper\_xi\_leyang\_three\_leading\_zeros\_extrapolate\_uc\_n6}
设在某端点机制下，存在常数 \(u_c\in\CC\)、\(C\neq 0\)、\(D\in\CC\)，使得对固定的 \(k\in\{0,1,2\}\) 有渐近展开
\[
u_{n,k}
=
u_c
-\frac{C}{n^{2}}(2k+1)^{2}
-\frac{D}{n^{4}}(2k+1)^{4}
O(n^{-6})
\qquad(n\to\infty).
\]
定义三零点外推量
\[
\boxed{\
\widehat u_{c}^{(n)}
:=\frac{75}{64}\,u_{n,0}-\frac{25}{128}\,u_{n,1}+\frac{3}{128}\,u_{n,2}.
}
\]
则成立严格误差改进
\[
\boxed{\ \widehat u_{c}^{(n)}=u_{c}+O(n^{-6}).\ }
\]
\end{theorem}
\begin{proof}
令 \(a_k:=(2k+1)^2\)，则 \(a_0=1,a_1=9,a_2=25\)。
取系数 \(\alpha_0,\alpha_1,\alpha_2\) 满足
\[
\alpha_0+\alpha_1+\alpha_2=1,\qquad
\alpha_0a_0+\alpha_1a_1+\alpha_2a_2=0,\qquad
\alpha_0a_0^2+\alpha_1a_1^2+\alpha_2a_2^2=0.
\]
该 Vandermonde 线性系统唯一解为
\(
(\alpha_0,\alpha_1,\alpha_2)=(75/64,\,-25/128,\,3/128)
\)。
将展开式代入 \(\sum_{k=0}^2\alpha_k u_{n,k}\)，由约束恰消去 \(n^{-2}\) 与 \(n^{-4}\) 两阶项，余项阶为 \(O(n^{-6})\)。证毕。
\end{proof}

\begin{theorem}[\texorpdfstring{$(M+1)$}{M+1} 点端点外推的闭式权重与最优误差阶]\label{thm:xi-leyang-mplus1-point-extrapolate-uc-optimal}
\leanverified{paper\_xi\_leyang\_mplus1\_point\_extrapolate\_uc\_optimal}
在与定理 \ref{thm:xi-leyang-three-leading-zeros-extrapolate-uc-n6} 同型的端点机制下，固定整数 \(M\ge 1\)，并令
\[
a_k:=(2k+1)^2\qquad (k=0,1,\dots,M).
\]
定义 Lagrange 权重
\[
\boxed{\
\alpha_{k}^{(M)}
:=
\prod_{\substack{0\le i\le M\\ i\ne k}}
\frac{-a_{i}}{a_{k}-a_{i}}
\qquad(k=0,\dots,M),
}
\]
以及外推估计量
\[
\boxed{\
\widehat u_{c,M}^{(n)}:=\sum_{k=0}^{M}\alpha_{k}^{(M)}\,u_{n,k}.
}
\]
若存在常数 \(C_1,\dots,C_{M+1}\in\CC\) 使得对每个固定 \(k\in\{0,1,\dots,M\}\) 有
\[
u_{n,k}
=
u_c+\sum_{j=1}^{M+1}\frac{C_j}{n^{2j}}\,a_k^{j}+O(n^{-2M-4}),
\qquad(n\to\infty),
\]
则
\[
\boxed{\ \widehat u_{c,M}^{(n)}=u_c+O(n^{-2M-2}).\ }
\]
并且在“仅使用 \(\{u_{n,0},\dots,u_{n,M}\}\) 并要求消去前 \(M\) 个偶阶项”的信息模型下，上述权重是唯一解，从而达到该模型内的最优误差阶。
\end{theorem}
\begin{proof}
在假设展开中，截断到 \(n^{-2M-2}\) 的部分在变量 \(a_k\) 上是次数至多 \(M+1\) 的多项式，并且 \(u_c\) 对应于 \(a=0\) 处的取值。
因此 \(\widehat u_{c,M}^{(n)}\) 等价于对节点 \(\{a_0,\dots,a_M\}\) 的 Lagrange 插值多项式在 \(a=0\) 的取值，其权重即为 \(L_k(0)\)，从而给出闭式 \(\alpha_k^{(M)}\)。
插值误差由最高项 \(a^{M+1}\) 控制，对应尺度 \(n^{-2M-2}\)，得到误差阶。
唯一性来自 Vandermonde 矩阵非退化。证毕。
\end{proof}

\begin{corollary}[由三零点同时恢复主标度常数与次阶常数]\label{cor:xi-leyang-recover-C-D-from-three-leading-zeros}
\leanverified{paper\_xi\_leyang\_recover\_C\_D\_from\_three\_leading\_zeros}
在定理 \ref{thm:xi-leyang-three-leading-zeros-extrapolate-uc-n6} 的假设下，令 \(\widehat u_c^{(n)}\) 为其三零点外推量，并定义
\[
B_k:=n^{2}\bigl(\widehat u_{c}^{(n)}-u_{n,k}\bigr)\qquad(k=0,1).
\]
则常数 \(C,D\) 可由显式公式一致估计：
\[
\boxed{\
\widehat C^{(n)}:=\frac{81B_{0}-B_{1}}{72}
\ =\ C+O(n^{-4}),\ }
\qquad
\boxed{\
\widehat D^{(n)}:=\frac{n^{2}(B_{1}-9B_{0})}{72}
\ =\ D+O(n^{-2}).\ }
\]
\end{corollary}
\begin{proof}
由定理 \ref{thm:xi-leyang-three-leading-zeros-extrapolate-uc-n6}，
\(
\widehat u_c^{(n)}=u_c+O(n^{-6})
\)。
结合 \(a_0=1,a_1=9\) 与展开式可得
\[
B_k=C\,a_k+\frac{D}{n^2}\,a_k^2+O(n^{-4})\qquad(k=0,1).
\]
对 \(k=0,1\) 作线性消元：\(\widehat C^{(n)}\) 消去 \(D/n^2\) 项而达到 \(O(n^{-4})\)，\(\widehat D^{(n)}\) 由差分提取 \(D\) 并得到 \(O(n^{-2})\)。证毕。
\end{proof}

\begin{proposition}[平方根分歧点的等模曲线切线方向]\label{prop:xi-leyang-equimodular-tangent-square-root}
在定理 \ref{thm:xi-leyang-square-root-collision-leading-zeros-n2} 的记号下，定义局部等模集合
\[
\Sigma:=\{u\in U:\ |\lambda_+(u)|=|\lambda_-(u)|\},
\]
其中 \(\lambda_\pm\) 取定理中的两支 Puiseux 解析延拓。
则在足够小邻域内，\(\Sigma\) 由恰好两条穿过 \(u_c\) 的实解析弧组成；
其在 \(u\)-平面中穿过 \(u_c\) 的切线方向唯一确定为
\[
\arg(u-u_c)\equiv \pi-2\arg\bigl(\lambda_c\overline{\alpha}\bigr)\pmod{\pi}.
\]
更精确地，在均匀化坐标 \(u=u_c+t^2\) 下，等模条件的一阶近似等价于
\[
\Re\bigl(\lambda_c\overline{\alpha}\,t\bigr)=0,
\]
从而在 \(t\)-平面给出两条互逆直线，其经 \(t\mapsto t^2\) 推送为上述两条切线方向。
\end{proposition}
\begin{proof}
由 \(\lambda_\pm(t)=\lambda_c\pm\alpha t+O(t^2)\) 得
\[
|\lambda_+(t)|^2-|\lambda_-(t)|^2
=|\lambda_c+\alpha t|^2-|\lambda_c-\alpha t|^2+O(|t|^2)
=4\,\Re(\lambda_c\overline{\alpha}\,t)+O(|t|^2).
\]
因此 \(|\lambda_+|=|\lambda_-|\) 在一阶上等价于 \(\Re(\lambda_c\overline{\alpha}\,t)=0\)，给出 \(t\)-平面中过原点的两条直线。
由 \(u-u_c=t^2\) 切向角加倍，得到所示同余式；高阶项由实解析隐函数定理给出两条实解析弧。证毕。
\end{proof}

\begin{theorem}[多项式传递矩阵的临界点代数性与判别式复杂度上界]\label{thm:xi-polynomial-transfer-matrix-branchpoint-discriminant-degree}
设 \(A(u)\in\Mat_d(\ZZ[u])\) 为多项式矩阵，且每个矩阵元 \(A_{ij}(u)\) 的次数不超过 \(k\ge 1\)。
令
\[
P(u,\lambda):=\det(\lambda I-A(u))\in\ZZ[u,\lambda],
\qquad
D(u):=\mathrm{Disc}_\lambda P(u,\lambda)\in\ZZ[u].
\]
则成立：
\begin{enumerate}
  \item \(D(u)\) 非恒零，且
  \[
  \deg_u D\le (2d-1)\deg_u P\le (2d-1)dk.
  \]
  \item 任一谱分歧点 \(u_c\)（即存在 \(\lambda\) 使 \(P(u_c,\lambda)=\partial_\lambda P(u_c,\lambda)=0\)）必为代数数，且
  \[
  [\QQ(u_c):\QQ]\le \deg_u D \le (2d-1)dk.
  \]
  \item 若某分歧点 \(u_c\) 为 \(D\) 的单根（\(D(u_c)=0\) 且 \(D'(u_c)\neq 0\)），并且在该点存在严格谱隙 \(|\lambda_j(u_c)|<|\lambda_c|\)（除去发生碰撞的二重根 \(\lambda_c\)），则定理 \ref{thm:xi-leyang-square-root-collision-leading-zeros-n2} 的平方根碰撞情形在该点为代数可审计：可由判别式单根性与谱隙验证推出 \(n^{-2}\) 量子化定标。
\end{enumerate}
\end{theorem}
\begin{proof}
由 \(D(u)=\mathrm{Disc}_\lambda(P)=(-1)^{d(d-1)/2}\mathrm{Res}_\lambda(P,\partial_\lambda P)\)。
对任意 \(F,G\in\ZZ[u,\lambda]\)，其关于 \(\lambda\) 的结式满足次数估计
\[
\deg_u\mathrm{Res}_\lambda(F,G)
\le (\deg_\lambda F)(\deg_u G)+(\deg_\lambda G)(\deg_u F).
\]
取 \(F=P\)、\(G=\partial_\lambda P\)，则 \(\deg_\lambda P=d\)、\(\deg_\lambda(\partial_\lambda P)=d-1\)，且 \(\deg_u(\partial_\lambda P)\le \deg_u P\)，从而
\[
\deg_u D\le d\,\deg_u P+(d-1)\deg_u P=(2d-1)\deg_u P.
\]
又 \(P(u,\lambda)=\det(\lambda I-A(u))\) 的每一项为 \(d\) 个矩阵元之积，故 \(\deg_u P\le dk\)，得第 1 条。

第 2 条由分歧点满足 \(D(u_c)=0\) 得出：\(u_c\) 为整系数多项式 \(D\) 的零点，因而为代数数且次数不超过 \(\deg_u D\)。

对第 3 条，\(D'(u_c)\neq 0\) 等价于 \(P(u,\lambda)\) 在 \((u_c,\lambda_c)\) 发生单纯二重根碰撞（局部仅两根以平方根方式分歧），与定理 \ref{thm:xi-leyang-square-root-collision-leading-zeros-n2} 的 (1)(2) 同型；
严格谱隙与可见性条件给出两主支主导，从而可应用定理 \ref{thm:xi-leyang-square-root-collision-leading-zeros-n2} 得到 \(n^{-2}\) 量子化。证毕。
\end{proof}

\begin{conjecture}[平方根端点下的两零点比值刚性与超收敛外推]\label{conj:xi-leyang-square-root-endpoint-two-zero-ratio}
设一族有限尺度配分函数 \(Z_n(u)\) 来自有限维解析传递算子或传递矩阵（例如 \(Z_n(u)=\ell(A(u)^n r)\)），并且其自由能的主要解析延拓障碍由某临界端点 \(u_c\) 处两条主导特征值的单纯平方根碰撞控制（判别式在 \(u_c\) 为单根，且存在严格谱隙与可见性）。
则存在编号约定使得靠近 \(u_c\) 的前两枚零点满足
\[
u_{n,1}-u_c=9\,(u_{n,0}-u_c)+o(n^{-2}),
\]
并且由两零点组合
\(
\widehat u_c(n)=(9u_{n,0}-u_{n,1})/8
\)
给出的端点外推误差为 \(o(n^{-2})\)。
若该刚性在数值数据中成立，则可据此反推临界端点处主导谱碰撞为单纯平方根类型，并据判别式与谱隙对端点作代数定位。
\end{conjecture}

\begin{theorem}[\texorpdfstring{$k$}{k}-重主导谱碰撞的 Mittag--Leffler 普遍临界核]\label{thm:xi-mittag-leffler-universal-kfold-jordan-collision}
\leanverified{paper\_xi\_mittag\_leffler\_universal\_kfold\_jordan\_collision}
令 \(M(t)\in M_d(\CC)\) 为 \(t\) 在 \(0\) 的邻域内解析的矩阵族。
设 \(M(0)\) 的谱满足：
\begin{enumerate}
  \item 存在 \(\eta>0\) 使得除特征值 \(1\) 外，其余全部特征值 \(\lambda\) 满足 \(|\lambda|\le 1-\eta\)；
  \item 特征值 \(1\) 的代数重数为 \(k\)，几何重数为 \(1\)（对应单 Jordan 块）；
  \item 在适当选择的主导 Jordan 基下，限制到主导广义特征子空间的特征多项式满足一参数非退化展开
  \[
  \chi(\lambda,t)
  =(\lambda-1)^k-\mathbf c\,t
  +O\!\left(t^2+t(\lambda-1)+(\lambda-1)^{k+1}\right),
  \qquad
  \mathbf c\neq 0.
  \]
\end{enumerate}
取向量 \(u,v\in\CC^d\) 使得其在主导 Jordan 链上的最高阶配对不为零（等价地：在主导子空间标准形中，双线性型 \(u^\top(\cdot)\,v\) 不消去 \((1,k)\) 矩阵元）。
定义标量序列
\[
Z_n(t):=u^\top M(t)^{\,n+k-1}v,\qquad n\ge 1.
\]
则存在常数 \(\mathbf C\neq 0\)（由 \(u,v\) 与主导 Jordan 链的配对给出）使得对任意紧集 \(K\subset\CC\) 有一致收敛
\[
\boxed{\;
n^{-(k-1)}\,Z_n\!\left(-\frac{s}{n^k}\right)\ \longrightarrow\
\mathbf C\,E_{k,k}\!\left(-\mathbf c\,s\right)
\qquad (n\to\infty,\ s\in K)\; }.
\]
其中二参数 Mittag--Leffler 函数定义为
\[
E_{k,k}(z):=\sum_{r\ge 0}\frac{z^r}{\Gamma(kr+k)}.
\]
\end{theorem}
\begin{proof}
由谱隙假设，取围绕主导谱簇的 Riesz 投影 \(P(t)\)，则存在常数 \(C>0\) 与 \(0<\vartheta<1\) 使
\[
M(t)^n=\bigl(P(t)M(t)P(t)\bigr)^n+R_n(t),
\qquad
\|R_n(t)\|\le C\,\vartheta^n
\]
在 \(t\) 的充分小邻域内一致成立。
故在 \(n^{-(k-1)}\) 的尺度下，\(R_n(t)\) 的贡献指数小，可忽略。
\par
记 \(A(t):=P(t)M(t)P(t)\) 并把其视为作用在 \(k\)-维主导广义特征子空间上的算子。
由单 Jordan 块的一参数解析变形，可取解析可逆变换 \(S(t)\) 使
\[
S(t)^{-1}A(t)S(t)=J_k+\tau(t)E_{k1},
\qquad
\tau(t)=\mathbf c\,t+O(t^2),
\]
其中 \(J_k\) 为特征值 \(1\) 的 \(k\times k\) Jordan 块，\(E_{k1}\) 为仅 \((k,1)\) 位置为 \(1\) 的矩阵。
并且由边界配对的非消去假设，存在 \(\mathbf C\neq 0\) 使得
\[
Z_n(t)=\mathbf C\cdot \bigl(J_k+\tau(t)E_{k1}\bigr)^{n+k-1}_{1k}
 +O(n^{k-2})+O(\vartheta^n)
\]
在 \(t\) 的充分小邻域内成立。
\par
对矩阵族 \(J_k+\tau E_{k1}\)，其 \((1,k)\) 元满足阶为 \(k\) 的常系数线性递推，其特征多项式为 \((\lambda-1)^k-\tau\)。
由此得到临界缩放 \(t=-s/n^k\) 下的统一极限核为 Mittag--Leffler 函数：
\[
n^{-(k-1)}\bigl(J_k+\tau(t)E_{k1}\bigr)^{n+k-1}_{1k}
\longrightarrow
E_{k,k}\!\left(-\mathbf c\,s\right),
\]
且误差对 \(s\) 的紧集一致。
合并上述各步即得结论。证毕。
\end{proof}

\begin{corollary}[\texorpdfstring{$k$}{k}-阶缩放下零点的普适量子化]\label{cor:xi-mittag-leffler-kfold-zero-quantization}
\leanverified{paper\_xi\_mittag\_leffler\_kfold\_zero\_quantization}
在定理 \ref{thm:xi-mittag-leffler-universal-kfold-jordan-collision} 的条件下，
设 \(t_{n,j}\) 为 \(Z_n(t)\) 的零点序列且满足 \(n^k t_{n,j}\to -s_j\)。
则必有
\[
E_{k,k}(-\mathbf c\,s_j)=0.
\]
反之，若 \(s\) 是函数 \(E_{k,k}(-\mathbf c\,\cdot)\) 的简单零点，则存在零点族 \(t_n(s)\) 使
\[
t_n(s)=-\frac{s}{n^k}+o(n^{-k}),
\qquad
Z_n\bigl(t_n(s)\bigr)=0.
\]
\end{corollary}
\begin{proof}
定理 \ref{thm:xi-mittag-leffler-universal-kfold-jordan-collision} 给出
\(
f_n(s):=n^{-(k-1)}Z_n(-s/n^k)
\)
在紧集上一致收敛到 \(\mathbf C\,E_{k,k}(-\mathbf c\,s)\)（\(\mathbf C\neq 0\)）。
由 Hurwitz 定理，任何零点极限必须落在极限函数的零点集合上；而简单零点在一致收敛下稳定延拓，给出所述反向存在性。证毕。
\end{proof}

\endinput


\paragraph{有限阶累积量的代数零点反演：Hankel 铅笔、收敛阶与稳定性}
本段把 Lee--Yang 零点的反演组织为一个纯代数问题：给定 \(\log Z\) 在原点的有限阶导数（等价的累积量/对数导数泰勒系数），以 Hankel 铅笔的谱读出最近零点族，并给出可证的指数收敛、锐性与扰动稳定性判据（参见 \cite{LeeYang1952LatticeGasIsing,FlindtGarrahan2013Trajectory} 的背景脉络）。

\begin{definition}[对数累积量与规范化幂和]\label{def:xi-leyang-cumulant-power-sums}
设 \(Z(z)\in\RR[z]\) 为多项式且 \(Z(0)\neq 0\)。记其非零零点（按重数计）为 \(\{\zeta_j\}_{j=1}^N\subset\CC\setminus\{0\}\)，并定义对数累积量
\[
\kappa_n:=\left.\frac{\dd^n}{\dd z^n}\log Z(z)\right|_{z=0}\qquad(n\ge 1).
\]
定义规范化幂和
\begin{equation}\label{eq:xi-leyang-power-sum-sn}
s_n:=-\frac{\kappa_n}{(n-1)!}=\sum_{j=1}^N \zeta_j^{-n}\qquad(n\ge 1),
\end{equation}
并记 \(\xi_j:=\zeta_j^{-1}\)。
\end{definition}

\begin{remark}[零点--累积量幂和恒等式]\label{rem:xi-leyang-zero-cumulant-identity}
在上述设定下有因子分解
\[
Z(z)=Z(0)\prod_{j=1}^N\left(1-\frac{z}{\zeta_j}\right),
\]
从而在 \(z=0\) 的邻域
\[
\log Z(z)=\log Z(0)-\sum_{n\ge 1}\frac{z^n}{n}\sum_{j=1}^N\zeta_j^{-n}.
\]
逐项求导并在 \(z=0\) 取值给出
\(
\kappa_n=-(n-1)!\sum_j\zeta_j^{-n}
\)，即 \eqref{eq:xi-leyang-power-sum-sn}。
\end{remark}

\begin{theorem}[支配零点的 Hankel 铅笔谱逼近与指数收敛]\label{thm:xi-leyang-hankel-pencil-dominant-zeros}
在定义 \ref{def:xi-leyang-cumulant-power-sums} 的设定下，将零点按模长排序 \(|\zeta_1|\le\cdots\le|\zeta_N|\)，并取整数 \(r\ge 1\) 使
\[
|\zeta_r|<|\zeta_{r+1}|\qquad (|\zeta_{N+1}|:=+\infty).
\]
令 \(\xi_j=\zeta_j^{-1}\)，则
\(
|\xi_1|\ge\cdots\ge|\xi_r|>|\xi_{r+1}|
\)。
对任意 \(n\ge 1\) 与固定 \(r\) 定义 \(r\times r\) Hankel 矩阵
\[
H_0^{(r)}(n):=\bigl[s_{n+i+j}\bigr]_{i,j=0}^{r-1},
\qquad
H_1^{(r)}(n):=\bigl[s_{n+i+j+1}\bigr]_{i,j=0}^{r-1},
\]
并在 \(H_0^{(r)}(n)\) 可逆时定义 Hankel 铅笔
\[
\mathcal{S}_r(n):=\bigl(H_0^{(r)}(n)\bigr)^{-1}H_1^{(r)}(n).
\]
记间隙比
\(
q:=|\xi_{r+1}|/|\xi_r|\in(0,1)
\)。
则存在常数 \(C>0\) 与 \(n_0\in\NN\)（仅依赖于 \(\{\xi_j\}_{j=1}^{r+1}\)）使得对所有 \(n\ge n_0\)：
\begin{enumerate}
  \item \(H_0^{(r)}(n)\) 可逆；
  \item 若以 \(\{\widehat\xi_j^{(r)}(n)\}_{j=1}^r\) 表示 \(\mathcal{S}_r(n)\) 的特征值（计代数重数），则可重排使
  \begin{equation}\label{eq:xi-leyang-hankel-pencil-eig-error}
  \max_{1\le j\le r}\abs{\widehat\xi_j^{(r)}(n)-\xi_j}\ \le\ C\,|\xi_1|\,q^{\,n}.
  \end{equation}
\end{enumerate}
因此 \(\widehat\zeta_j^{(r)}(n):=1/\widehat\xi_j^{(r)}(n)\) 给出最近 \(r\) 个零点的指数级逼近，其收敛率由模长间隙 \(q\) 精确控制。
\end{theorem}
\begin{proof}
令 Vandermonde 型矩阵 \(V_r\in\CC^{r\times r}\) 与对角矩阵 \(D_n,\Lambda\) 为
\[
(V_r)_{i,j}:=\xi_j^{\,i}\quad(0\le i\le r-1,\,1\le j\le r),\qquad
D_n:=\mathrm{diag}(\xi_1^n,\dots,\xi_r^n),\qquad
\Lambda:=\mathrm{diag}(\xi_1,\dots,\xi_r).
\]
按定义 \(s_m=\sum_{j=1}^N\xi_j^m\)，可分解为
\[
H_0^{(r)}(n)=V_rD_nV_r^{\mathsf T}+E_0(n),\qquad
H_1^{(r)}(n)=V_rD_n\Lambda V_r^{\mathsf T}+E_1(n),
\]
其中误差矩阵 \(E_0(n),E_1(n)\) 仅由 \(j\ge r+1\) 的项贡献。由 \(|\xi_j|^{n+i+j}\le |\xi_{r+1}|^{n}|\xi_{r+1}|^{i+j}\) 得
\[
\norm{E_\ell(n)}\le C_\ell |\xi_{r+1}|^{\,n}\qquad(\ell=0,1),
\]
其中 \(C_\ell\) 与 \(n\) 无关。
另一方面，\(V_r\) 可逆且 \(D_n\) 可逆，故 \(V_rD_nV_r^{\mathsf T}\) 可逆，且存在常数 \(C'>0\) 使
\[
\norm{(V_rD_nV_r^{\mathsf T})^{-1}}\le C'|\xi_r|^{-n}.
\]
因此
\[
\norm{(V_rD_nV_r^{\mathsf T})^{-1}E_0(n)}\le C''\left(\frac{|\xi_{r+1}|}{|\xi_r|}\right)^{n}=C''q^{\,n}.
\]
取 \(n_0\) 充分大使上式 \(\le \tfrac12\)，则 \(H_0^{(r)}(n)\) 可逆且
\(
\norm{(H_0^{(r)}(n))^{-1}}\le 2\norm{(V_rD_nV_r^{\mathsf T})^{-1}}
\)。
进一步，
\[
\mathcal{S}_r(n)
=(V_rD_nV_r^{\mathsf T})^{-1}(V_rD_n\Lambda V_r^{\mathsf T})
O(q^{\,n})
=V_r^{-{\mathsf T}}\Lambda V_r^{\mathsf T}+O(q^{\,n}).
\]
其中 \(V_r^{-{\mathsf T}}\Lambda V_r^{\mathsf T}\) 与 \(\Lambda\) 相似，特征值恰为 \(\{\xi_1,\dots,\xi_r\}\)。
对矩阵扰动应用 Bauer--Fike 型特征值扰动界即可得到 \eqref{eq:xi-leyang-hankel-pencil-eig-error}。证毕。
\end{proof}

\begin{corollary}[最近共轭对的二阶 Hankel 判别式估计器]\label{cor:xi-leyang-nearest-conjugate-pair-discriminant-estimator}
在定理 \ref{thm:xi-leyang-hankel-pencil-dominant-zeros} 的设定下取 \(r=2\)，并进一步假设最近两零点构成非实共轭对
\[
|\zeta_1|=|\zeta_2|=:R_1<R_2:=\min_{j\ge 3}|\zeta_j|,\qquad \zeta_2=\overline{\zeta_1},\qquad \zeta_1\notin\RR.
\]
对 \(n\ge 1\) 定义二阶 Hankel 判别式
\[
\Delta_n:=\det\begin{pmatrix}
s_n & s_{n+1}\\
s_{n+1} & s_{n+2}
\end{pmatrix}
=s_ns_{n+2}-s_{n+1}^2,
\]
并在 \(\Delta_n\neq 0\) 时定义
\[
A_n:=\frac{s_{n+3}s_n-s_{n+2}s_{n+1}}{\Delta_n},\qquad
B_n:=\frac{s_{n+3}s_{n+1}-s_{n+2}^2}{\Delta_n}.
\]
令 \(\widehat\xi_\pm^{(n)}\) 为二次多项式 \(q_n(\lambda)=\lambda^2-A_n\lambda+B_n\) 的两根，并令 \(\widehat\zeta_\pm^{(n)}:=1/\widehat\xi_\pm^{(n)}\)。
则存在常数 \(C>0\) 与 \(n_0\)（仅依赖于 \(\{\zeta_j\}\)）使对所有 \(n\ge n_0\) 成立：在适当重排下
\[
\max_{\sigma\in\{+,-\}}\abs{\widehat\zeta_\sigma^{(n)}-\zeta_\sigma}
\ \le\ C\,R_1\left(\frac{R_1}{R_2}\right)^{n}.
\]
\end{corollary}
\begin{proof}
当 \(r=2\) 时，\(\mathcal{S}_2(n)=(H_0^{(2)}(n))^{-1}H_1^{(2)}(n)\) 的特征多项式恰为
\(
\lambda^2-A_n\lambda+B_n
\)，其根为 \(\widehat\xi_\pm^{(n)}\)。
由定理 \ref{thm:xi-leyang-hankel-pencil-dominant-zeros} 的特征值收敛界（此时 \(q=R_1/R_2\)）得到
\(\widehat\xi_\pm^{(n)}\to \xi_{1,2}\) 且误差为 \(O(q^n)\)。
再由映射 \(\xi\mapsto 1/\xi\) 在 \(|\xi|=R_1^{-1}\) 的邻域上 Lipschitz（常数 \(\asymp R_1^2\)）得到对应的 \(\zeta\)-平面误差界。证毕。
\end{proof}

\begin{proposition}[指数收敛率的一般锐性]\label{prop:xi-leyang-exponential-rate-sharpness}
在推论 \ref{cor:xi-leyang-nearest-conjugate-pair-discriminant-estimator} 的设定下，进一步假设存在第二个最近共轭对
\[
|\zeta_3|=|\zeta_4|=R_2,\qquad \zeta_4=\overline{\zeta_3},
\]
且其余零点模长严格大于 \(R_2\)。
则存在常数 \(c\neq 0\) 与 \(n_1\in\NN\) 使对所有 \(n\ge n_1\) 有
\[
\max_{\sigma\in\{+,-\}}\abs{\widehat\zeta_\sigma^{(n)}-\zeta_\sigma}
\ \ge\ |c|\,R_1\left(\frac{R_1}{R_2}\right)^{n}.
\]
因此在仅由模长间隙 \(R_1<R_2\) 所刻画的类上，收敛率 \((R_1/R_2)^n\) 不能被统一改进。
\end{proposition}
\begin{proof}[证明要点]
将幂和分解为两对共轭主项与更远余项：
\(
s_n=(\xi_1^n+\overline{\xi_1}^n)+(\xi_3^n+\overline{\xi_3}^n)+\widetilde R_n
\)，其中 \(|\widetilde R_n|=o(|\xi_3|^n)\)。
把 \(A_n,B_n\) 视为 \((s_n,s_{n+1},s_{n+2},s_{n+3})\) 的有理函数，并按小参数
\(
\varepsilon_n:=(|\xi_3|/|\xi_1|)^n=(R_1/R_2)^n
\)
作一阶展开，可得在一般位置（相位非共振的开稠密条件）下
\[
A_n=(\xi_1+\overline{\xi_1})+\alpha\,\varepsilon_n+o(\varepsilon_n),\qquad
B_n=\xi_1\overline{\xi_1}+\beta\,\varepsilon_n+o(\varepsilon_n),
\]
其中 \((\alpha,\beta)\neq (0,0)\)。
二次多项式根对系数的一阶扰动给出根位移 \(\asymp \varepsilon_n\)，从而经 \(\xi\mapsto 1/\xi\) 得到 \(\zeta\)-平面误差下界 \(\asymp R_1\varepsilon_n\)。证毕。
\end{proof}

\begin{theorem}[Hankel 奇异值的对数斜率识别支配模长]\label{thm:xi-leyang-hankel-singular-slope-detects-moduli}
在定理 \ref{thm:xi-leyang-hankel-pencil-dominant-zeros} 的设定下，固定任意 \(k\ge r\)，并令
\[
H_0^{(k)}(n):=\bigl[s_{n+i+j}\bigr]_{i,j=0}^{k-1}.
\]
记其奇异值按降序为 \(\sigma_1(n)\ge\cdots\ge\sigma_k(n)\ge 0\)。
则对每个 \(1\le j\le r\) 存在常数 \(c_j,C_j>0\) 与 \(n_j\) 使对所有 \(n\ge n_j\) 有
\[
c_j\,|\xi_j|^{\,n}\ \le\ \sigma_j(n)\ \le\ C_j\,|\xi_j|^{\,n}.
\]
因此
\[
\lim_{n\to\infty}\frac{1}{n}\log\sigma_j(n)=\log|\xi_j|=-\log|\zeta_j|\qquad(1\le j\le r).
\]
\end{theorem}
\begin{proof}[证明要点]
将 \(H_0^{(k)}(n)\) 分解为秩 \(r\) 的主项与余项：
\(
H_0^{(k)}(n)=V_k D_n^{(r)} V_k^{\mathsf T}+\widetilde E(n)
\)，其中 \((V_k)_{i,j}=\xi_j^i\)（\(0\le i\le k-1,\,1\le j\le r\)），\(D_n^{(r)}=\mathrm{diag}(\xi_1^n,\dots,\xi_r^n)\)，且 \(\norm{\widetilde E(n)}\le C|\xi_{r+1}|^{\,n}\)。
对主项应用奇异值的 min--max 表达与 \(V_k\) 的左右有界性得到 \(\sigma_j(V_k D_n^{(r)} V_k^{\mathsf T})\asymp |\xi_j|^n\)；
再用 Weyl 型奇异值扰动不等式吸收余项，即得双边界与极限。证毕。
\end{proof}

\begin{theorem}[累积量噪声扰动下的零点重构稳定性]\label{thm:xi-leyang-reconstruction-stability-under-noise}
在定理 \ref{thm:xi-leyang-hankel-pencil-dominant-zeros} 的设定下，设观测到的幂和为 \(\widetilde s_n=s_n+\delta_n\)，并由此构造 \(\widetilde H_0^{(r)}(n),\widetilde H_1^{(r)}(n)\) 与
\(
\widetilde{\mathcal S}_r(n):=(\widetilde H_0^{(r)}(n))^{-1}\widetilde H_1^{(r)}(n)
\)。
若对某个 \(n\ge n_0\) 有
\begin{equation}\label{eq:xi-leyang-stability-invertibility-condition}
\norm{(H_0^{(r)}(n))^{-1}}\cdot \norm{\widetilde H_0^{(r)}(n)-H_0^{(r)}(n)}\le \frac12,
\end{equation}
则：
\begin{enumerate}
\item \(\widetilde H_0^{(r)}(n)\) 可逆，且
\begin{align}\label{eq:xi-leyang-stability-pencil-bound}
\norm{\widetilde{\mathcal S}_r(n)-\mathcal S_r(n)}
&\le
4\norm{(H_0^{(r)}(n))^{-1}}
\Bigl(\norm{\widetilde H_1^{(r)}(n)-H_1^{(r)}(n)}
\nonumber\\
&\hspace{3.7em}
+\norm{\mathcal S_r(n)}\cdot\norm{\widetilde H_0^{(r)}(n)-H_0^{(r)}(n)}\Bigr).
\end{align}
\item 若 \(\mathcal S_r(n)\) 可对角化，\(\mathcal S_r(n)=PDP^{-1}\)，则任意 \(\widetilde{\mathcal S}_r(n)\) 的特征值 \(\widetilde\xi\) 满足
\[
\mathrm{dist}\!\bigl(\widetilde\xi,\sigma(\mathcal S_r(n))\bigr)\ \le\ \kappa(P)\,\norm{\widetilde{\mathcal S}_r(n)-\mathcal S_r(n)},
\qquad
\kappa(P):=\norm{P}\,\norm{P^{-1}}.
\]
从而在 \(|\xi_j|\) 有下界的区域内，零点估计 \(\widetilde\zeta=1/\widetilde\xi\) 的误差亦可由同一右端给出显式控制。
\end{enumerate}
\end{theorem}
\begin{proof}
条件 \eqref{eq:xi-leyang-stability-invertibility-condition} 蕴含
\(
\widetilde H_0^{(r)}(n)=H_0^{(r)}(n)\bigl(I+(H_0^{(r)}(n))^{-1}(\widetilde H_0^{(r)}(n)-H_0^{(r)}(n))\bigr)
\)
中括号可逆，且
\(
\norm{(\widetilde H_0^{(r)}(n))^{-1}}\le 2\norm{(H_0^{(r)}(n))^{-1}}
\)。
再由恒等式
\[
\widetilde{\mathcal S}-\mathcal S
=(\widetilde H_0^{-1}-H_0^{-1})H_1+\widetilde H_0^{-1}(\widetilde H_1-H_1),
\qquad
\widetilde H_0^{-1}-H_0^{-1}=\widetilde H_0^{-1}(H_0-\widetilde H_0)H_0^{-1},
\]
直接得到 \eqref{eq:xi-leyang-stability-pencil-bound}。
第二条为 Bauer--Fike 定理的直接推论。证毕。
\end{proof}

\endinput

\paragraph{外二次谱与局部可识别性：二阶 Hankel--Wronskian、判别式平方与六窗消参证书}
沿用结论 \ref{con:xi-terminal-zm-order4-rational} 的双变量配分函数 \(Z_m(y)\) 及其最小湮灭多项式
\(\Pi(\lambda,y)=\lambda^4-\lambda^3-(2y+1)\lambda^2+\lambda+y(y+1)\)。
并沿用结论 \ref{con:xi-terminal-zm-discriminant-leyang} 的 Lee--Yang 三次因子
\[
P_{\mathrm{LY}}(y):=256y^3+411y^2+165y+32.
\]
\par
本段将二阶 Hankel 行列式视为外二次表示的 Pl\"ucker 坐标，并据此给出其闭递推、外二次判别式的完全平方分解、黄金比共振点处的降阶回文递推，以及仅用六窗局部数据恢复参数与消去参数的代数证书。
\par

\begin{conclusion}[外二次谱控制的二阶 Hankel--Wronskian 六阶闭递推]\label{con:xi-terminal-zm-exterior-square-hankel-wronskian-order6}
定义二阶 Hankel--Wronskian
\[
\Delta^{(2)}_m(y):=Z_m(y)Z_{m+3}(y)-Z_{m+1}(y)Z_{m+2}(y)\qquad(m\ge 0),
\]
并在固定 \(y\) 下简记 \(\Delta^{(2)}_m:=\Delta^{(2)}_m(y)\)。
则序列 \(\bigl(\Delta^{(2)}_m\bigr)_{m\ge 0}\) 满足一条六阶齐次线性递推：
\[
\boxed{\;
\begin{aligned}
&\Delta^{(2)}_{m+6}+(2y+1)\Delta^{(2)}_{m+5}-(y^2+y+1)\Delta^{(2)}_{m+4}-(4y^3+7y^2+3y+1)\Delta^{(2)}_{m+3}\\
&\qquad-(y^4+2y^3+2y^2+y)\Delta^{(2)}_{m+2}+(2y^5+5y^4+4y^3+y^2)\Delta^{(2)}_{m+1}+y^3(y+1)^3\Delta^{(2)}_{m}=0.
\end{aligned}\;}
\]
等价地，其特征多项式可取为
\[
\chi^{(2)}(\mu;y)
=\mu^6+(2y+1)\mu^5-(y^2+y+1)\mu^4-(4y^3+7y^2+3y+1)\mu^3-(y^4+2y^3+2y^2+y)\mu^2
\]
\[
\qquad\qquad\qquad\qquad\qquad\quad
+(2y^5+5y^4+4y^3+y^2)\mu+y^3(y+1)^3.
\]
若 \(\lambda_1,\dots,\lambda_4\) 为 \(\Pi(\lambda,y)=0\) 的根（在代数闭包中），则 \(\chi^{(2)}(\mu;y)\) 的根集合可表述为
\[
\{\lambda_i\lambda_j:\ 1\le i<j\le 4\}.
\]
\end{conclusion}
\begin{proof}[证明要点]
在代数闭包中，若 \(\Pi(\lambda,y)\) 的四根互异，则存在系数 \(c_1,\dots,c_4\) 使
\(
Z_m(y)=\sum_{i=1}^4 c_i\lambda_i^m
\)。
将该表达代入
\(
\Delta^{(2)}_m
=\det\!\bigl(\begin{smallmatrix}Z_m&Z_{m+1}\\ Z_{m+2}&Z_{m+3}\end{smallmatrix}\bigr)
\)
并展开行列式，可得
\(
\Delta^{(2)}_m=\sum_{1\le i<j\le 4}c_{ij}(\lambda_i\lambda_j)^m
\)
（反对称性消去 \(i=j\) 项）。
因此 \(\Delta^{(2)}_m\) 的指数模态由六个乘积 \(\lambda_i\lambda_j\) 控制，递推的特征多项式为
\(
\prod_{i<j}(\mu-\lambda_i\lambda_j)
\)，即为 \(\chi^{(2)}(\mu;y)\)。
将该乘积在 \(\QQ(y)\) 中展开得到所示显式系数。
由于两边系数均为 \(\ZZ[y]\) 中多项式，上述恒等式由稠密参数集推出对全体 \(y\) 成立。
显式展开与核验可由脚本 \texttt{scripts/exp\_fold\_zm\_bivariate\_partition\_audit.py} 复核。证毕。
\end{proof}

\begin{conclusion}[外二次判别式的完全平方分解与二阶共振临界集]\label{con:xi-terminal-zm-exterior-square-discriminant-square}
将 \(\chi^{(2)}(\mu;y)\) 视为关于 \(\mu\) 的六次多项式，则其判别式在 \(\QQ(y)\) 中具有完全平方分解：
\[
\boxed{\;
\mathrm{Disc}_{\mu}\bigl(\chi^{(2)}(\mu;y)\bigr)
=y^8(y-1)^2(y+1)^6(y^2+y-1)^2P_{\mathrm{LY}}(y)^2
=\Bigl(y^4(y-1)(y+1)^3(y^2+y-1)P_{\mathrm{LY}}(y)\Bigr)^2.
\;}
\]
从而其判别式支撑（等价地，单生成阶 \(A[\mu]\) 的有限判别式支撑，\(A:=\QQ[y]\)）包含于
\[
\boxed{\;y\in\{0,1,-1\}\ \cup\ \{y^2+y-1=0\}\ \cup\ \{P_{\mathrm{LY}}(y)=0\}.}
\]
并且由于判别式为平方，\(\mathrm{Gal}\bigl(\chi^{(2)}/\QQ(y)\bigr)\subseteq A_6\)。
\par
其中 \(y=-1\) 与 \(y^2+y-1=0\) 的出现属于单生成模型 \(A[\mu]\subset B\)（\(B\) 为 \(A\) 在 \(L=\QQ(y)(\mu)\) 中的整闭包）的指数理想缺陷；其并不对应扩张 \(L/\QQ(y)\) 的真实有限分歧。
真实有限分歧支撑精确等于 \(y=0\)、\(y=1\) 与 \(P_{\mathrm{LY}}(y)=0\)（见定理 \ref{thm:xi-terminal-zm-exterior-square-finite-ramification-index-ideal}）。
\end{conclusion}
\begin{proof}[证明要点]
由结论 \ref{con:xi-terminal-zm-exterior-square-hankel-wronskian-order6} 的显式表达，对 \(\chi^{(2)}(\mu;y)\) 作 \(\mu\)-判别式并在 \(\QQ[y]\) 中因式分解即可得到所示完全平方分解。
判别式在基域中为平方当且仅当伽罗瓦群包含于交错群，应用于次数 \(6\) 即得结论。
上述恒等式可由脚本 \texttt{scripts/exp\_fold\_zm\_bivariate\_partition\_audit.py} 复核。证毕。
\end{proof}

\leanverified{paper\_xi\_terminal\_zm\_exterior\_square\_finite\_ramification\_index\_ideal}
\begin{theorem}[外二次六次原始元的有限分歧精确定位与指数理想闭式公式]\label{thm:xi-terminal-zm-exterior-square-finite-ramification-index-ideal}
设 \(K=\QQ(y)\)，\(A=\QQ[y]\)，并令 \(\chi^{(2)}(\mu;y)\in A[\mu]\) 如结论 \ref{con:xi-terminal-zm-exterior-square-hankel-wronskian-order6}。
取 \(\mu\) 为其任一根并置 \(L:=K(\mu)\)。
记 \(B\) 为 \(A\) 在 \(L\) 中的整闭包。
则成立如下三点：
\begin{enumerate}
  \item \(\chi^{(2)}\) 的伽罗瓦闭包与谱四次 \(\Pi(\lambda,y)\) 的分裂域一致；
  且 \(\mathrm{Gal}(L/K)\cong S_4\) 并通过四元集合之 \(2\)-子集置换表示嵌入 \(A_6\)
  （见结论 \ref{con:xi-terminal-zm-exterior-square-galois-s4}）。
  \item \(B/A\) 的真实有限分歧支撑精确等于
  \[
  \{y=0\}\ \cup\ \{y=1\}\ \cup\ \{P_{\mathrm{LY}}(y)=0\},
  \]
  并且其有限判别式理想为主理想
  \[
  \mathfrak D_{B/A,\mathrm{fin}}=\bigl(y(y-1)P_{\mathrm{LY}}(y)\bigr)^2.
  \]
  \item 单生成阶 \(A[\mu]\subset B\) 的指数理想满足闭式恒等式
  \[
  [B:A]_{\mathrm{fin}}=\bigl(y^3(y+1)^3(y^2+y-1)\bigr),
  \]
  从而在有限素处有判别式指数公式
  \[
  \Disc_A(A[\mu])_{\mathrm{fin}}
  =\mathfrak D_{B/A,\mathrm{fin}}\cdot [B:A]_{\mathrm{fin}}^{2}.
  \]
  特别地，\(\mathrm{Disc}_{\mu}\bigl(\chi^{(2)}(\mu;y)\bigr)\) 中出现的 \(y+1\) 与 \(y^2+y-1\) 因子在有限处完全由指数缺陷贡献，而非真实分歧。
\end{enumerate}
上述结论的分歧--指数分解结构可参照内部审计记录 \cite{Internal2026ExteriorSquareFiniteRamificationIndexIdealAudit}。
\end{theorem}
\begin{proof}[证明要点]
第 \(1\) 点已由结论 \ref{con:xi-terminal-zm-exterior-square-galois-s4} 给出。
\par
对第 \(2\) 点，谱四次覆盖的 \(S_4\)-伽罗瓦闭包在 \(y\)-直线上的有限分歧值由 \(y=0,1\) 与 \(P_{\mathrm{LY}}(y)=0\) 给出（见结论 \ref{con:xi-terminal-zm-leyang-ramification-critical-values}），且相应惰性生成元可取为换位。
外二次扩张 \(L/K\) 对应 \(S_4\) 在 \(2\)-子集上的次数 \(6\) 置换表示；在该表示下换位诱导为循环型 \(2^2\cdot 1^2\)，故每个有限分歧值对不同的总贡献为 \(2\)，从而
\(
\mathfrak D_{B/A,\mathrm{fin}}=\bigl(y(y-1)P_{\mathrm{LY}}(y)\bigr)^2
\)。
\par
对第 \(3\) 点，\(\chi^{(2)}(\mu;y)\) 关于 \(\mu\) 首一且可分，故
\(
\Disc_A(A[\mu])=\mathrm{Disc}_{\mu}\bigl(\chi^{(2)}(\mu;y)\bigr)\in A
\)，并由 Dedekind 域上的经典指数公式
\(
\Disc_A(A[\mu])=\Disc_A(B)\,[B:A]^2
\)
比较有限素处指数即得
\(
[B:A]_{\mathrm{fin}}=(y^3(y+1)^3(y^2+y-1))
\)。
由此亦可见 \(y=-1\) 与 \(y^2+y-1=0\) 处仅为单生成阶的非极大性，而非 \(L/K\) 的真实有限分歧。证毕。
\end{proof}

\leanverified{paper\_xi\_terminal\_zm\_exterior\_square\_curve\_genus2}
\begin{theorem}[外二次谱曲线的正规化亏格]\label{thm:xi-terminal-zm-exterior-square-curve-genus2}
令 \(C^{\wedge 2}\) 为仿射曲线 \(\chi^{(2)}(\mu;y)=0\) 的正规化之光滑射影模型，并视为到 \(\PP^1_y\) 的次数 \(6\) 覆叠。
则
\[
g(C^{\wedge 2})=2.
\]
\end{theorem}
\begin{proof}[证明要点]
由定理 \ref{thm:xi-terminal-zm-exterior-square-finite-ramification-index-ideal}，有限分歧值为五点 \(y=0,1\) 与 \(P_{\mathrm{LY}}(y)=0\) 的三根，且各点惰性诱导为循环型 \(2^2\cdot 1^2\)，故每点分歧贡献为 \(2\)。
另一方面，无穷远处由谱四次的四循环惰性诱导（见结论 \ref{con:xi-terminal-zm-s4-splitting-curve-genus16}），其在 \(2\)-子集表示下循环型为 \(4\cdot 2\)，故无穷远分歧贡献为 \(4\)。
对次数 \(6\) 的覆盖应用 Riemann--Hurwitz 公式得
\[
2g(C^{\wedge 2})-2=-12+\bigl(4+5\cdot 2\bigr)=2,
\]
从而 \(g(C^{\wedge 2})=2\)。证毕。
\end{proof}

\begin{corollary}[外二次属二曲线 Jacobian 的椭圆分裂]\label{cor:xi-terminal-zm-exterior-square-jacobian-splitting}
\leanverified{paper\_xi\_terminal\_zm\_exterior\_square\_jacobian\_splitting}
在泛型参数下，\(C^{\wedge 2}\) 的 Jacobian 在 \(\QQ(y)\) 上同源分解为两条椭圆曲线之积。
更精确地，若 \(\widetilde{X}\to\PP^1_y\) 为谱四次的 \(S_4\)-伽罗瓦闭包，且取配对稳定子 \(V_4^{\mathrm{pair}}=\langle(12),(34)\rangle\le S_4\)，则 \(C^{\wedge 2}\) 与 \(\widetilde{X}/V_4^{\mathrm{pair}}\) 双有理等价，且
\[
J(C^{\wedge 2})\ \sim\ E_{\mathrm{res}}\times E,
\]
其中 \(E=\widetilde{X}/S_3\) 与 \(E_{\mathrm{res}}=\widetilde{X}/D_8\) 为两条椭圆曲线（见结论 \ref{con:xi-terminal-zm-s4-jacobian-group-algebra-prym-e3-b3}）。
\end{corollary}
\begin{proof}[证明要点]
由结论 \ref{con:xi-terminal-zm-exterior-square-galois-s4}，扩张 \(L/K\) 的单值化由 \(S_4\) 在 \(2\)-子集上的置换表示给出；
因此其对应的次数 \(6\) 中间域可识别为分裂域曲线 \(\widetilde{X}\) 对 \(2\)-子集稳定子的商，即 \(\widetilde{X}/V_4^{\mathrm{pair}}\)（参见结论 \ref{con:xi-terminal-zm-s4-v4pair-explicit-model} 中对该商的显式模型）。
于是 \(C^{\wedge 2}\) 与 \(\widetilde{X}/V_4^{\mathrm{pair}}\) 具有同一函数域，从而双有理等价。
最后，结论 \ref{con:xi-terminal-zm-s4-jacobian-group-algebra-prym-e3-b3} 给出 \(J(\widetilde{X}/V_4^{\mathrm{pair}})\sim E_{\mathrm{res}}\times E\)，遂得到所述分解。证毕。
\end{proof}

\begin{conclusion}[黄金比参数处的 \texorpdfstring{$(-1)$}{-1}-倒易对称与外二次谱塌缩]\label{con:xi-terminal-zm-minus-one-reciprocity-exterior-square}
对所有 \(y\) 恒成立恒等式
\[
\Pi(\lambda,y)-\lambda^{4}\Pi(-1/\lambda,y)=-(\lambda^{4}-1)(y^{2}+y-1).
\]
因此当 \(y^{2}+y-1=0\) 时，谱四次满足 \((-1)\)-倒易性
\[
\Pi(\lambda,y)=\lambda^{4}\Pi(-1/\lambda,y),
\]
并诱导根集合上的对合 \(\lambda\mapsto -1/\lambda\)。
从而外二次多项式满足
\[
(\mu+1)^{2}\ \bigm|\ \chi^{(2)}(\mu;y).
\]
\end{conclusion}
\begin{proof}[证明要点]
差值恒等式由对 \(\Pi(\lambda,y)\) 直接代入并化简得到。
当 \(y^{2}+y-1=0\) 时，\(\lambda\mapsto -1/\lambda\) 保持方程 \(\Pi(\lambda,y)=0\) 不变，故根集在该对合下成对出现；从而存在至少两组 \(i<j\) 使 \(\lambda_i\lambda_j=-1\)，于是 \(\chi^{(2)}\) 在 \(\mu=-1\) 处至少二重零点，即 \((\mu+1)^2\mid \chi^{(2)}(\mu;y)\)。证毕。
\end{proof}

\begin{conclusion}[\texorpdfstring{$\mu=-1$}{mu=-1} 的黄金比共振判别与二重根结构]\label{con:xi-terminal-zm-exterior-square-mu-minus-one-resonance}
对所有 \(y\) 有显式恒等式
\[
\boxed{\;\chi^{(2)}(-1;y)=y(y-1)(y^2+y-1)^2.\;}
\]
因此 \(\mu=-1\) 为 \(\chi^{(2)}(\mu;y)\) 的根当且仅当
\(
y\in\{0,1\}
\)
或
\(
y^2+y-1=0
\)。
并且在 \(y^2+y-1=0\) 且 \(y\notin\{0,1\}\) 时，\(\mu=-1\) 为 \(\chi^{(2)}(\mu;y)\) 的二重根。
\end{conclusion}
\begin{proof}[证明要点]
直接代入 \(\mu=-1\) 并在 \(\ZZ[y]\) 中因式分解得到恒等式。
再对 \(\partial_\mu\chi^{(2)}(\mu;y)\) 代入 \(\mu=-1\) 可见其亦被 \(y^2+y-1\) 整除，从而在 \(y^2+y-1=0\) 时至少二重；
结合结论 \ref{con:xi-terminal-zm-exterior-square-discriminant-square} 的判别式分解可知该处重根阶数恰为 \(2\)。证毕。
\end{proof}

\begin{conclusion}[黄金比参数处 \texorpdfstring{$\Delta^{(2)}_m$}{Delta(2)} 的最小递推降阶与回文五阶关系]\label{con:xi-terminal-zm-exterior-square-hankel-wronskian-order5-golden}
若 \(y\) 满足 \(y^2+y-1=0\)，则 \(\chi^{(2)}(\mu;y)\) 含因子 \((\mu+1)^2\)，并且 \(\bigl(\Delta^{(2)}_m(y)\bigr)_{m\ge 0}\) 的最小湮灭多项式为 \(\chi^{(2)}(\mu;y)\) 的平方自由部分；因而递推阶数由 \(6\) 降至 \(5\)。
在该约束下可取如下回文五阶递推：
\[
\boxed{\;
\Delta^{(2)}_{m+5}+2y\,\Delta^{(2)}_{m+4}-2(y+1)\Delta^{(2)}_{m+3}-2(y+1)\Delta^{(2)}_{m+2}+2y\,\Delta^{(2)}_{m+1}+\Delta^{(2)}_{m}=0.
\;}
\]
\end{conclusion}
\begin{proof}[证明要点]
由结论 \ref{con:xi-terminal-zm-exterior-square-mu-minus-one-resonance} 得在 \(y^2+y-1=0\) 时 \(\mu=-1\) 为二重根，故 \((\mu+1)^2\mid \chi^{(2)}(\mu;y)\)。
另一方面，结论 \ref{con:xi-terminal-zm-discriminant-leyang} 给出 \(\mathrm{Disc}_{\lambda}(\Pi)=-y(y-1)P_{\mathrm{LY}}(y)\)，而 \(y^2+y-1\) 与 \(y(y-1)P_{\mathrm{LY}}(y)\) 在 \(\QQ[y]\) 中互素，故在 \(y^2+y-1=0\) 下 \(\Pi(\lambda,y)\) 具有四个互异根；
因此其伴随矩阵在代数闭包中可对角化，而外二次表示在同一特征向量楔积基下亦对角化，故最小多项式为特征多项式的平方自由部分。
将 \(y^2+y-1=0\) 代入 \(\chi^{(2)}(\mu;y)\) 并约去一个 \((\mu+1)\) 因子，化简得到所示回文五阶递推。证毕。
\end{proof}

\begin{conclusion}[任意六窗局部数据的有理可识别性]\label{con:xi-terminal-zm-six-window-rational-identifiability}
令 \(m\ge 0\)，并设
\[
U_m:=Z_{m+4}-Z_{m+3}+Z_{m+1},\qquad
V_m:=Z_{m+5}-Z_{m+4}+Z_{m+2}.
\]
若 \(\Delta^{(2)}_m\neq 0\)，则参数 \(y\) 可由任意六个连续观测值 \(\{Z_m,\dots,Z_{m+5}\}\) 以有理式唯一恢复：
\[
\boxed{\;
y=\frac12\left(\frac{Z_mV_m-Z_{m+1}U_m}{\Delta^{(2)}_m}-1\right).
\;}
\]
\end{conclusion}
\begin{proof}[证明要点]
由结论 \ref{con:xi-terminal-zm-order4-rational} 的四阶递推在指标 \(m+4\)、\(m+5\) 处可写为二元线性方程组：
\[
U_m=(2y+1)Z_{m+2}-y(y+1)Z_m,\qquad
V_m=(2y+1)Z_{m+3}-y(y+1)Z_{m+1}.
\]
对未知量 \(A:=2y+1\)、\(B:=y(y+1)\) 用 Cramer 法则求解，系数行列式恰为
\(\Delta^{(2)}_m=Z_mZ_{m+3}-Z_{m+1}Z_{m+2}\)。
由此得到
\(
A=(Z_mV_m-Z_{m+1}U_m)/\Delta^{(2)}_m
\)，进而 \(y=(A-1)/2\)。证毕。
\end{proof}

\begin{conclusion}[消去参数的六窗四次 Pl\"ucker 型恒等式]\label{con:xi-terminal-zm-six-window-parameter-free-quartic}
在结论 \ref{con:xi-terminal-zm-six-window-rational-identifiability} 的记号下，不显式引入参数 \(y\)，六窗数据 \(\{Z_m,\dots,Z_{m+5}\}\) 必满足如下参数自由四次恒等式：
\[
\boxed{\;
\bigl(Z_mV_m-Z_{m+1}U_m\bigr)^2-\bigl(\Delta^{(2)}_m\bigr)^2-4\,\Delta^{(2)}_m\bigl(Z_{m+2}V_m-Z_{m+3}U_m\bigr)=0.
\;}
\]
\end{conclusion}
\begin{proof}[证明要点]
由上一结论同理可得
\[
A=\frac{Z_mV_m-Z_{m+1}U_m}{\Delta^{(2)}_m},
\qquad
B=\frac{Z_{m+2}V_m-Z_{m+3}U_m}{\Delta^{(2)}_m}.
\]
而 \((A,B)\) 来自同一参数 \(y\)，故恒满足代数约束 \(4B=A^2-1\)（由 \(A=2y+1\)、\(B=y(y+1)\) 直接计算）。
将 \(A,B\) 的有理表达代入并清分母，即得所示四次恒等式。证毕。
\end{proof}

\begin{conclusion}[外二次六次多项式的伽罗瓦群刻画与信息完备性]\label{con:xi-terminal-zm-exterior-square-galois-s4}
令 \(\chi^{(2)}(\mu;y)\) 如结论 \ref{con:xi-terminal-zm-exterior-square-hankel-wronskian-order6}。
则当 \(y\) 不落入结论 \ref{con:xi-terminal-zm-exterior-square-discriminant-square} 的分歧集合时，成立：
\begin{enumerate}
  \item \(\chi^{(2)}(\mu;y)\) 的分裂域与 \(\Pi(\lambda,y)\) 的分裂域相同；
  \item \(\mathrm{Gal}\bigl(\chi^{(2)}/\QQ(y)\bigr)\cong S_4\)，并以 \(S_4\) 作用于四元集合之 \(2\)-子集的标准方式忠实嵌入 \(A_6\subset S_6\)。
\end{enumerate}
特别地，结论 \ref{con:xi-terminal-zm-exterior-square-discriminant-square} 中判别式为平方与该嵌入的全偶置换性一致。
\end{conclusion}
\begin{proof}[证明要点]
设 \(\lambda_1,\dots,\lambda_4\) 为 \(\Pi(\lambda,y)\) 的根，则由外二次构造，\(\chi^{(2)}\) 的根为 \(\lambda_i\lambda_j\)（\(i<j\)），故其分裂域包含于 \(\Pi\) 的分裂域。
反之，若某自同构固定全部 \(\lambda_i\lambda_j\)，则其在 \(2\)-子集集合 \(\binom{\{1,2,3,4\}}{2}\) 上的诱导作用为恒等；
该作用为忠实，故该自同构必为恒等，从而两分裂域相同。
另一方面，由定理 \ref{thm:xi-terminal-zm-leyang-elliptic-structure} 知 \(\mathrm{Gal}(\Pi/\QQ(y))\cong S_4\)，于是 \(\mathrm{Gal}(\chi^{(2)}/\QQ(y))\cong S_4\)。
再由结论 \ref{con:xi-terminal-zm-exterior-square-discriminant-square} 的判别式平方分解可知该像包含于 \(A_6\)。
该嵌入的全偶性亦可由生成元循环型直接检验，例如换位 \((12)\) 在 \(2\)-子集上诱导为 \((1^2\,2^2)\)，而 \(4\)-循环 \((1234)\) 诱导为 \((4\cdot 2)\)，二者均为偶置换。证毕。
\end{proof}

\begin{conclusion}[谱对相位共振的结果式证书与 Chebyshev 压缩：等模条件与判别式的统一消元]\label{con:xi-terminal-zm-phase-resonance-resultant-chebyshev}
沿用结论 \ref{con:xi-terminal-zm-order4-rational} 的谱四次 \(\Pi(\lambda,y)\)。
对任意 \(u\in\CC^{\ast}\) 定义二参数消元多项式
\[
\mathcal{R}(u,y):=\operatorname{Res}_{\lambda}\bigl(\Pi(\lambda,y),\Pi(u\lambda,y)\bigr)\in\ZZ[u,y].
\]
则 \(\mathcal{R}(u,y)\) 在 \(\ZZ[u,y]\) 中严格分解为
\[
\mathcal{R}(u,y)=y(y+1)(u-1)^4\,F(u,y),
\]
从而归一化结果式可取为
\[
\mathcal F(u,y):=\frac{\mathcal{R}(u,y)}{y(y+1)(u-1)^4}=F(u,y).
\]
其中 \(F(u,y)\) 在 \(u\) 上为 \(12\) 次回文多项式、在 \(y\) 上为 \(6\) 次多项式，并可作相位压缩：令 \(s:=u+u^{-1}\)，则存在 \(s\) 的六次多项式 \(G(s,y)\in\ZZ[s,y]\) 使得
\[
F(u,y)=u^6G(s,y).
\]
上述 \(F\) 与 \(G\) 的显式系数与退化恒等式由下式给出。
% Auto-generated by scripts/exp_fold_zm_phase_resonance_resultant_chebyshev_audit.py
\[
\Pi(\lambda,y)=\lambda^{4}-\lambda^{3}-(2y+1)\lambda^{2}+\lambda+y(y+1),\qquad \mathcal{R}(u,y):=\operatorname{Res}_{\lambda}\bigl(\Pi(\lambda,y),\Pi(u\lambda,y)\bigr).
\]
\[
\mathcal{R}(u,y)=y(y+1)(u-1)^{4}\,F(u,y).
\]
\[
\begin{aligned}
F(u,y)=&\ y^{3} \left(y + 1\right)^{3}(u^{12}+1)
+\ y^{2} \left(y + 1\right)^{2} \left(2 y + 1\right)^{2}(u^{11}+u)
+\ y \left(y + 1\right) \left(2 y^{2} - 1\right) \left(y^{2} + y + 1\right)(u^{10}+u^{2})\\
&-\ \left(12 y^{6} + 50 y^{5} + 63 y^{4} + 38 y^{3} + 19 y^{2} + 6 y + 1\right)(u^{9}+u^{3})
-\ \left(17 y^{6} + 87 y^{5} + 96 y^{4} + 33 y^{3} + 16 y^{2} + 7 y + 2\right)(u^{8}+u^{4})\\
&+\ \left(8 y^{6} - 26 y^{5} + 11 y^{4} + 94 y^{3} + 51 y^{2} + 14 y + 1\right)(u^{7}+u^{5})
+\ \left(28 y^{6} + 32 y^{5} + 103 y^{4} + 186 y^{3} + 103 y^{2} + 32 y + 4\right)u^{6}.
\end{aligned}
\]
\[
F(u,y)=u^{6}G(s,y),\qquad s:=u+u^{-1}.
\]
\[
\begin{aligned}
G(s,y)=&\ \left(y^{6} + 3 y^{5} + 3 y^{4} + y^{3}\right)s^{6}
+\left(4 y^{6} + 12 y^{5} + 13 y^{4} + 6 y^{3} + y^{2}\right)s^{5}\\
&+\left(- 4 y^{6} - 14 y^{5} - 15 y^{4} - 6 y^{3} - 2 y^{2} - y\right)s^{4}
+\left(- 32 y^{6} - 110 y^{5} - 128 y^{4} - 68 y^{3} - 24 y^{2} - 6 y - 1\right)s^{3}\\
&+\left(- 16 y^{6} - 76 y^{5} - 81 y^{4} - 24 y^{3} - 8 y^{2} - 3 y - 2\right)s^{2}
+\left(64 y^{6} + 184 y^{5} + 265 y^{4} + 238 y^{3} + 113 y^{2} + 32 y + 4\right)s\\
&+\left(64 y^{6} + 208 y^{5} + 295 y^{4} + 250 y^{3} + 131 y^{2} + 44 y + 8\right).
\end{aligned}
\]
\[
G(2,y)=-y(y-1)\bigl(256y^{3}+411y^{2}+165y+32\bigr).
\]

\par
当 \(u=e^{i\theta}\) 时 \(s=2\cos\theta\)，从而“存在一对谱根满足比值 \(u\)”（因而模长相等）的消元条件可降维为代数曲线 \(G(2\cos\theta,y)=0\)。
并且在 \(u=1\)（即 \(s=2\)）处出现严格退化 \(G(2,y)=-y(y-1)P_{\mathrm{LY}}(y)\)，从而该二参数证书在单一框架下同时回收了判别式零轨迹（含 \(y_{\mathrm{LY}}\)）并编码全部等模谱对的相位共振集合。
\end{conclusion}
\begin{proof}[证明要点]
对变量 \(\lambda\) 作结果式消元并在 \(\ZZ[u,y]\) 中因式分解得到 \(\mathcal{R}(u,y)=y(y+1)(u-1)^4F(u,y)\)。
回文化与 \(s=u+u^{-1}\) 的压缩来自 \(u\mapsto u^{-1}\) 的互反对称；退化式 \(G(2,y)\) 由代入 \(s=2\) 后在 \(\ZZ[y]\) 中因式分解得到。
上述恒等式与系数可由脚本 \texttt{scripts/exp\_fold\_zm\_phase\_resonance\_resultant\_chebyshev\_audit.py} 复核。证毕。
\end{proof}

\begin{conclusion}[相位共振判别式的闭式分解与平方结构]\label{con:xi-terminal-zm-phase-compression-disc-square}
沿用结论 \ref{con:xi-terminal-zm-phase-resonance-resultant-chebyshev} 的记号：
\[
F(u,y)=u^{6}G(u+u^{-1},y),\qquad s=u+u^{-1},
\]
并记
\[
P_{\mathrm{LY}}(y):=256y^{3}+411y^{2}+165y+32.
\]
则关于 \(G\) 与 \(F\) 的端点值、判别式闭式分解与平方结构由下式统一给出：
% Auto-generated by scripts/exp_fold_zm_phase_compression_discriminant_square_audit.py
\[
G(2,y)=-y(y-1)P_{\mathrm{LY}}(y),\qquad G(-2,y)=y^{2}(y-1)^{2}.
\]
\[
\mathrm{Disc}_{s}\bigl(G(s,y)\bigr)=y^{10}(y-1)^{4}(y+1)^{6}(y^{2}+y-1)^{4}P_{\mathrm{LY}}(y)^{2}P_{9}(y)^{2}.
\]
\[
\mathrm{Disc}_{s}\bigl(G(s,y)\bigr)=\Delta_{G}(y)^{2},\qquad \Delta_{G}(y)=y^{5}(y-1)^{2}(y+1)^{3}(y^{2}+y-1)^{2}P_{\mathrm{LY}}(y)P_{9}(y).
\]
\[
P_{9}(y)=256 y^{9} + 1219 y^{8} + 2542 y^{7} + 3090 y^{6} + 2446 y^{5} + 1315 y^{4} + 478 y^{3} + 112 y^{2} + 16 y + 1.
\]
\[
\mathrm{Disc}_{u}\bigl(F(u,y)\bigr)=-y^{23}(y-1)^{11}(y+1)^{12}(y^{2}+y-1)^{8}P_{\mathrm{LY}}(y)^{5}P_{9}(y)^{4}.
\]
\[
\mathrm{Disc}_{u}\bigl(F(u,y)\bigr)=\mathrm{Disc}_{s}\bigl(G(s,y)\bigr)^{2}\,G(2,y)\,G(-2,y).
\]
\[
\mathrm{Disc}_{u}\bigl(F(u,y)\bigr)=-y(y-1)P_{\mathrm{LY}}(y)\,\Omega(y)^{2},\quad \Omega(y)=y^{11}(y-1)^{5}(y+1)^{6}(y^{2}+y-1)^{4}P_{\mathrm{LY}}(y)^{2}P_{9}(y)^{2}.
\]
\[
\gcd\Bigl(\operatorname{Res}_{s}\bigl(G,\partial_{s}G\bigr),\ \operatorname{Res}_{s}\bigl(G,\partial_{y}G\bigr)\Bigr)=y^{10}(y-1)^{3}(y+1)^{6}(y^{2}+y-1)^{4}P_{9}(y)^{2}.
\]
\[
G\!\left(s,\frac{-1-\sqrt5}{2}\right)=\bigl(s^{2}+(1-2\sqrt5)s+7\bigr)\bigl(s^{2}+(2+\sqrt5)s-1+2\sqrt5\bigr)^{2},
\]
\[
G\!\left(s,\frac{-1+\sqrt5}{2}\right)=\bigl(s^{2}+(1+2\sqrt5)s+7\bigr)\bigl(s^{2}+(2-\sqrt5)s-1-2\sqrt5\bigr)^{2}.
\]
\[
P_{\mathrm{LY}}(y_{0})=0\ \Longrightarrow\ G(s,y_{0})=y_{0}^{3}(y_{0}+1)^{3}(s-2)L(s)Q(s)^{2},
\]
\[
\widetilde{G}(s,y):=\frac{G(s,y)}{y^{3}(y+1)^{3}}\ \Longrightarrow\ \widetilde{G}(s,y_{0})=(s-2)L(s)Q(s)^{2}.
\]
\[
L(s)=s+\frac{-1280y_{0}^{2}+2041y_{0}+6583}{3168},\quad Q(s)=s^{2}+\frac{6400y_{0}^{2}-10205y_{0}-1235}{3168}\,s+\frac{16640y_{0}^{2}-26533y_{0}-9547}{1584}.
\]
\[
G(s,0)=-(s-2)(s+2)^{2},\qquad G(s,1)=(s-2)(s+2)^{2}(2s+5)^{2}(2s-5).
\]

特别地，\(\mathrm{Disc}_{s}(G)\) 在 \(\ZZ[y]\) 中为完全平方，且满足桥接恒等式
\[
\mathrm{Disc}_{u}\bigl(F(u,y)\bigr)=\mathrm{Disc}_{s}\bigl(G(s,y)\bigr)^2\,G(2,y)\,G(-2,y).
\]
\end{conclusion}
\begin{proof}[证明要点]
记 \(d:=\deg_s G=6\)，并在代数闭包上写
\[
G(s,y)=a(y)\prod_{i=1}^{d}(s-s_i),\qquad a(y)=y^{3}(y+1)^{3}.
\]
由 \(s=u+u^{-1}\) 得
\[
F(u,y)=u^{d}G(u+u^{-1},y)=a(y)\prod_{i=1}^{d}\bigl(u^{2}-s_i u+1\bigr).
\]
对二次因子 \(f_i(u):=u^2-s_i u+1\) 使用判别式与结式乘法恒等式
\[
\mathrm{Disc}\!\Bigl(\prod_i f_i\Bigr)=\prod_i \mathrm{Disc}(f_i)\cdot \prod_{i<j}\mathrm{Res}(f_i,f_j)^2
\]
并代入
\(
\mathrm{Disc}_u(f_i)=s_i^2-4,\ 
\mathrm{Res}_u(f_i,f_j)=(s_i-s_j)^2
\)，
再结合
\[
\prod_i(s_i^2-4)=\frac{G(2,y)G(-2,y)}{a(y)^2},\qquad
\prod_{i<j}(s_i-s_j)^4=\frac{\mathrm{Disc}_s(G)^2}{a(y)^{4d-4}},
\]
即可得到
\(
\mathrm{Disc}_u(F)=\mathrm{Disc}_s(G)^2\,G(2,y)\,G(-2,y)
\)，其中首项因子 \(a(y)\) 完全消去。
将显式 \(G(s,y)\) 代入并在 \(\ZZ[y]\) 中因式分解，得到所列闭式。
上述恒等式可由脚本 \texttt{scripts/exp\_fold\_zm\_phase\_compression\_discriminant\_square\_audit.py} 复核。证毕。
\end{proof}

\begin{theorem}[Lee--Yang 三次因子根处的双重碰撞与局部单值群型]\label{thm:xi-terminal-zm-leyang-cubic-double-collision-monodromy-type}
\leanverified{paper\_xi\_terminal\_zm\_leyang\_cubic\_double\_collision\_monodromy\_type}
在结论 \ref{con:xi-terminal-zm-phase-compression-disc-square} 的设定下，取 \(y_0\) 为 \(P_{\mathrm{LY}}(y)\) 的任一根。
设 \(\pi_{\mathrm{ph}}:X_{\mathrm{ph}}\to\PP^1_y\) 为平面曲线 \(G(s,y)=0\) 的射影闭包之正规化在 \(y\)-坐标上的投影（为六叶覆叠，见结论 \ref{con:xi-terminal-zm-phase-compression-galois-a6}）。
则在分歧点 \(y_0\) 处，\(\pi_{\mathrm{ph}}\) 的局部单值生成元在六张纸上的循环型为
\[
(2)(2)(1)(1)\in \mathfrak S_6,
\]
因而属于交错群 \(\mathfrak A_6\)。
并且 \(\mathrm{Disc}_{s}(G(s,y))\) 在 \(y=y_0\) 处的消失阶恰为 \(2\)。
\end{theorem}

\begin{proof}
由结论 \ref{con:xi-terminal-zm-phase-compression-disc-square} 的显式分解（式 \(\widetilde G(s,y)=G(s,y)/(y^3(y+1)^3)\) 的专化分解），在 \(y=y_0\) 处存在一次多项式 \(L\in\QQ(y_0)[s]\) 与二次多项式 \(Q\in\QQ(y_0)[s]\)（在 \(\overline{\QQ}\) 上有两互异根）使得
\[
G(s,y_0)=y_0^3(y_0+1)^3\,(s-2)\,L(s)\,Q(s)^2.
\]
因此专化多项式 \(G(\cdot,y_0)\) 在 \(s\)-坐标上恰有两个互异二重根（来自 \(Q(s)^2\)）与两个单根（来自 \((s-2)L(s)\)）。
对一参数代数族的局部解析分支，围绕 \(y_0\) 的小回路将每个二重根邻域内的两条分支互换（给出一换位），并且两处二重根互不相交，故局部单值为两不交换位之积，循环型为 \((2)(2)(1)(1)\)。
其为偶置换，从而落在 \(\mathfrak A_6\)。
\par
另一方面，\(\mathrm{Disc}_s(G)\) 在 \(\ZZ[y]\) 中含因子 \(P_{\mathrm{LY}}(y)^2\)（结论 \ref{con:xi-terminal-zm-phase-compression-disc-square}），故在 \(y=y_0\) 处至少二阶消失；
而上述专化仅出现两处简单二重根且其余根单纯，故判别式消失阶恰为 \(2\)。证毕。
\end{proof}

\leanverified{paper\_xi\_terminal\_zm\_leyang\_u\_root\_pinning\_double}
\begin{corollary}[Lee--Yang 临界根处的 \(u\)-根钉扎与互反二重根对]\label{cor:xi-terminal-zm-leyang-u-root-pinning-double}
在定理 \ref{thm:xi-terminal-zm-leyang-cubic-double-collision-monodromy-type} 的设定下，令 \(Q(s)=(s-\alpha)(s-\beta)\)（\(\alpha\neq\beta\)）为其在 \(\overline{\QQ}\) 上的分解。
则在 \(\overline{\QQ}[u]\) 中，\(F(u,y_0)=u^6G(u+u^{-1},y_0)\) 具有分解
\[
F(u,y_0)=u^{6}\cdot y_0^3(y_0+1)^3\cdot (u-1)^2\cdot \bigl(u^2-\sigma u+1\bigr)\cdot (u^2-\alpha u+1)^2\cdot (u^2-\beta u+1)^2,
\]
其中 \(\sigma\in\overline{\QQ}\) 为线性因子 \(L(s)=s-\sigma\) 的根。
因此：
\begin{enumerate}
  \item \(u=1\) 为 \(F(u,y_0)\) 的二重根；
  \item \(F(u,y_0)\) 额外包含两对互反二重根 \(\{u_\alpha,u_\alpha^{-1}\}\)、\(\{u_\beta,u_\beta^{-1}\}\)，分别为二次方程 \(u^2-\alpha u+1=0\)、\(u^2-\beta u+1=0\) 的根，并均以重数 \(2\) 出现。
\end{enumerate}
\end{corollary}

\begin{proof}
由 \(F(u,y)=u^6G(u+u^{-1},y)\) 与定理 \ref{thm:xi-terminal-zm-leyang-cubic-double-collision-monodromy-type} 中的专化分解，
每个 \(s\)-因子 \(s-s_0\) 在 \(u\)-坐标下对应二次因子 \(u^2-s_0u+1\)（参见结论 \ref{con:xi-terminal-zm-phase-compression-disc-square} 的一般分解式）。
当 \(s_0=2\) 时 \(u^2-2u+1=(u-1)^2\)，给出 \(u=1\) 的二重根；
当 \(s_0=\alpha,\beta\) 为二重根时，相应二次因子在 \(F(u,y_0)\) 中以平方出现，从而得到两对互反二重根。
线性因子 \(L(s)\) 贡献的 \(u^2-\sigma u+1\) 为一对互反单根，不影响上述二重性断言。证毕。
\end{proof}

\begin{theorem}[导子与单值群型给出的局部退化分层]\label{thm:xi-langlands-leyang-local-degeneration-stratification}
\leanverified{paper\_xi\_langlands\_leyang\_local\_degeneration\_stratification}
在 \(p_7\) 的 \(S_5\)-Artin 层（定理 \ref{thm:xi-p7-s5-odd-ramification-discriminant-and-conductor}--\ref{thm:xi-p7-s5-dyadic-c3-inertia-quadratic-functors-conductors}）与 Chebyshev--Lee--Yang 相位层（结论 \ref{con:xi-terminal-zm-phase-resonance-resultant-chebyshev}--推论 \ref{cor:xi-terminal-zm-leyang-u-root-pinning-double}）的共同设定下：
\begin{enumerate}
  \item 在奇分歧素数 \(q\) 处，\(\rho_4,\rho_5,\rho_6\) 的惰性限制具有共同三维不变子空间（定理 \ref{thm:xi-p7-s5-q-inertia-vacuum3-schur-functor-conductors}），从而任意 Schur 函子 \(\operatorname{Sym}^r,\bigwedge^r\) 的 \(q\)-导子指数可由有限组合计数给出（\eqref{eq:xi-p7-q-schur-symr-conductor-closed}--\eqref{eq:xi-p7-q-schur-wedger-conductor-closed}）。
  \item 在 Lee--Yang 三次因子根 \(y_0\) 处，相位压缩六叶覆叠 \(X_{\mathrm{ph}}\to\PP^1_y\) 发生双重碰撞，其局部单值循环型为 \((2)(2)(1)(1)\)（定理 \ref{thm:xi-terminal-zm-leyang-cubic-double-collision-monodromy-type}），并在 \(u\)-平面上表现为 \(u=1\) 与两对互反二重根的钉扎（推论 \ref{cor:xi-terminal-zm-leyang-u-root-pinning-double}）。
\end{enumerate}
\end{theorem}

\begin{proof}
第一条为定理 \ref{thm:xi-p7-s5-q-inertia-vacuum3-schur-functor-conductors} 的直接重述。
第二条由定理 \ref{thm:xi-terminal-zm-leyang-cubic-double-collision-monodromy-type} 与推论 \ref{cor:xi-terminal-zm-leyang-u-root-pinning-double} 合并得到。证毕。
\end{proof}

\begin{conclusion}[Chebyshev 压缩下的二次判别式类严格不变性]\label{con:xi-terminal-zm-phase-compression-disc-squareclass-invariance}
在函数域平方类群 \(\QQ(y)^{\times}/(\QQ(y)^{\times})^{2}\) 中成立
\[
\bigl[\mathrm{Disc}_{u}\bigl(F(u,y)\bigr)\bigr]
=\bigl[-y(y-1)P_{\mathrm{LY}}(y)\bigr]
=\bigl[\mathrm{Disc}_{\lambda}\bigl(\Pi(\lambda,y)\bigr)\bigr].
\]
等价地，\(\mathrm{Disc}_{u}(F)\) 的平方自由核恰为 \(-y(y-1)P_{\mathrm{LY}}(y)\)。
\end{conclusion}
\begin{proof}[证明要点]
由结论 \ref{con:xi-terminal-zm-phase-compression-disc-square} 的显式分解，
\[
\mathrm{Disc}_{u}(F)
=-y(y-1)P_{\mathrm{LY}}(y)\,\Omega(y)^2
\]
其中 \(\Omega(y)\in\QQ(y)\)。
因此其平方类仅由 \(-y(y-1)P_{\mathrm{LY}}(y)\) 决定。
再由结论 \ref{con:xi-terminal-zm-discriminant-leyang}
\(
\mathrm{Disc}_{\lambda}(\Pi)=-y(y-1)P_{\mathrm{LY}}(y)
\)
立即得到两者平方类相等。证毕。
\end{proof}

\begin{theorem}[判别式平方类锁定的同一二次中心特征]\label{thm:xi-terminal-zm-leyang-phase-resonance-discriminant-center-character}
\leanverified{paper\_xi\_terminal\_zm\_leyang\_phase\_resonance\_discriminant\_center\_character}
令 \(K:=\QQ(y)\)，并沿用结论 \ref{con:xi-terminal-zm-phase-resonance-resultant-chebyshev} 的
\(\Pi(\lambda,y)\in K[\lambda]\) 与 \(F(u,y)\in K[u]\)。
设 \(\Pi\) 与 \(F\) 在 \(K\) 上可分。
则由结论 \ref{con:xi-terminal-zm-phase-compression-disc-squareclass-invariance} 的平方类恒等式可知二次扩张严格同一：
\[
K\!\left(\sqrt{\mathrm{Disc}_{u}\bigl(F(u,y)\bigr)}\right)
\;=\;
K\!\left(\sqrt{\mathrm{Disc}_{\lambda}\bigl(\Pi(\lambda,y)\bigr)}\right)
\;=\;
K\!\left(\sqrt{-y(y-1)P_{\mathrm{LY}}(y)}\right).
\]
并且令 \(L_\Pi/K\)、\(L_F/K\) 分别为 \(\Pi\)、\(F\) 的分裂域，令
\[
\chi_\Pi:\mathrm{Gal}(L_\Pi/K)\to\{\pm 1\},\qquad
\chi_F:\mathrm{Gal}(L_F/K)\to\{\pm 1\}
\]
为其在根集置换表示上的行列式（等价地符号角色），则二者诱导的二次子扩张同一：
\[
L_\Pi^{\ker\chi_\Pi}=K\!\left(\sqrt{\mathrm{Disc}_{\lambda}(\Pi)}\right),
\qquad
L_F^{\ker\chi_F}=K\!\left(\sqrt{\mathrm{Disc}_{u}(F)}\right).
\]
从而 \(\chi_\Pi\) 与 \(\chi_F\) 的二次中心特征由同一平方类 \(-y(y-1)P_{\mathrm{LY}}(y)\) 统一刻画。
\end{theorem}
\begin{proof}
结论 \ref{con:xi-terminal-zm-phase-compression-disc-squareclass-invariance} 给出
\(
[\mathrm{Disc}_{u}(F)]=[\mathrm{Disc}_{\lambda}(\Pi)]
\)
在 \(K^\times/(K^\times)^2\) 中成立。
因此存在 \(\Omega\in K^\times\) 使得
\(
\mathrm{Disc}_{u}(F)=\Omega^2\,\mathrm{Disc}_{\lambda}(\Pi)
\)，从而
\(
K(\sqrt{\mathrm{Disc}_{u}(F)})=K(\sqrt{\mathrm{Disc}_{\lambda}(\Pi)})
\)。
再由结论 \ref{con:xi-terminal-zm-discriminant-leyang} 的判别式分解
\(
\mathrm{Disc}_{\lambda}(\Pi)=-y(y-1)P_{\mathrm{LY}}(y)
\)
即得首个等式链。
\par
对可分次数 \(n\) 多项式 \(f\in K[T]\)，其分裂域 \(L_f/K\) 上根集的置换表示给出一条符号角色
\(\chi_f:\mathrm{Gal}(L_f/K)\to\{\pm 1\}\)。
判别式的平方类刻画交错子群的固定子域：有标准等式
\[
L_f^{\ker\chi_f}=K\!\left(\sqrt{\mathrm{Disc}(f)}\right).
\]
将 \(f=\Pi\) 与 \(f=F\) 代入并结合已证的二次扩张同一性，即得结论。证毕。
\end{proof}

\begin{theorem}[良好素数处的 Frobenius 因子数奇偶同调]\label{thm:xi-terminal-zm-leyang-phase-resonance-frobenius-parity}
\leanverified{paper\_xi\_terminal\_zm\_leyang\_phase\_resonance\_frobenius\_parity}
令 \(p\) 为奇素数，取 \(y_0\in\FF_p\) 使得专化多项式
\(\Pi(\lambda,y_0)\in\FF_p[\lambda]\) 与 \(F(u,y_0)\in\FF_p[u]\) 均可分。
记
\[
r_{4}(p;y_0):=\#\{\text{\(\Pi(\lambda,y_0)\) 在 \(\FF_p[\lambda]\) 中的不可约因子}\},
\qquad
r_{12}(p;y_0):=\#\{\text{\(F(u,y_0)\) 在 \(\FF_p[u]\) 中的不可约因子}\}.
\]
则有同余关系
\[
r_{4}(p;y_0)\equiv r_{12}(p;y_0)\pmod 2.
\]
\end{theorem}
\begin{proof}
由定理 \ref{thm:xi-terminal-zm-leyang-phase-resonance-discriminant-center-character}，
\(\mathrm{Disc}_{u}(F)\) 与 \(\mathrm{Disc}_{\lambda}(\Pi)\) 在 \(K=\QQ(y)\) 的平方类群中相等；
在满足可分性且 \(p\neq 2\) 的良好约化点 \((p,y_0)\) 上，该平方类恒等式在
\(\FF_p^\times/(\FF_p^\times)^2\) 中保持，从而勒让德符号一致：
\[
\left(\frac{\mathrm{Disc}_{\lambda}(\Pi(\cdot,y_0))}{p}\right)
=
\left(\frac{\mathrm{Disc}_{u}(F(\cdot,y_0))}{p}\right).
\]
另一方面，Stickelberger 定理给出：对任意首一可分多项式 \(f\in\FF_p[T]\)（\(\deg f=n\)）记其不可约因子数为 \(r\)，则
\[
\left(\frac{\mathrm{Disc}(f)}{p}\right)=(-1)^{n-r}.
\]
将其分别用于 \(f=\Pi(\lambda,y_0)\)（\(n=4\)）与 \(f=F(u,y_0)\)（\(n=12\)），并注意 \(4,12\) 为偶数，得
\[
(-1)^{r_4(p;y_0)}
=(-1)^{4-r_4(p;y_0)}
=\left(\frac{\mathrm{Disc}_{\lambda}(\Pi(\cdot,y_0))}{p}\right)
=\left(\frac{\mathrm{Disc}_{u}(F(\cdot,y_0))}{p}\right)
=(-1)^{12-r_{12}(p;y_0)}
=(-1)^{r_{12}(p;y_0)}.
\]
故 \(r_{4}(p;y_0)\equiv r_{12}(p;y_0)\pmod 2\)。证毕。
\end{proof}

\begin{conclusion}[压缩六次多项式的二子集表示伽罗瓦像]\label{con:xi-terminal-zm-phase-compression-galois-a6}
设 \(k:=\QQ(y)\)，\(\Pi(\lambda,y)\in k[\lambda]\) 为谱四次多项式，且 \(\mathrm{Gal}(\Pi/k)\cong S_4\)；
记其分裂域为 \(L/k\)，根为 \(\lambda_1,\dots,\lambda_4\)。
定义
\[
s_{ij}:=\frac{\lambda_i}{\lambda_j}+\frac{\lambda_j}{\lambda_i}\in L\qquad (1\le i<j\le4).
\]
则 \(G(s,y)\) 的六个根恰为 \(\{s_{ij}\}_{i<j}\)，并且 \(G\) 在 \(k[s]\) 上不可约、\([k(s_{12}):k]=6\)。
此外，\(G\) 的分裂域等于 \(L\)，且
\[
\mathrm{Gal}\bigl(G/k\bigr)\cong S_4
\]
其在六个根上的作用同构于 \(S_4\) 在 \(\binom{\{1,2,3,4\}}{2}\) 上的自然作用。
该像包含于 \(A_6\)，故 \(\mathrm{Disc}_s(G)\) 为平方，与结论 \ref{con:xi-terminal-zm-phase-compression-disc-square} 一致。
\end{conclusion}
\begin{proof}[证明要点]
令 \(u_{ij}:=\lambda_i/\lambda_j\)。
由结果式构造
\(
F(u,y)=\operatorname{Res}_{\lambda}(\Pi(\lambda,y),\Pi(u\lambda,y))/\bigl(y(y+1)(u-1)^4\bigr)
\)
可知 \(F(u_{ij},y)=0\)，于是
\[
0=F(u_{ij},y)=u_{ij}^{6}G(u_{ij}+u_{ij}^{-1},y)=u_{ij}^{6}G(s_{ij},y),
\]
故每个 \(s_{ij}\) 为 \(G\) 的根。
由于 \(\deg_s G=6\) 且 \(\{s_{ij}\}_{i<j}\) 恰有 \(6\) 个，根集即为全部 \(s_{ij}\)。
对 \(s_{12}\) 而言，其在 \(S_4\) 下稳定子为集合稳定子 \(\mathrm{Stab}(\{1,2\})\)，阶为 \(4\)，故轨道大小为 \(24/4=6\)；
因而 \(s_{12}\) 在 \(k\) 上的最小多项式次数为 \(6\)，即 \(G\) 不可约且 \([k(s_{12}):k]=6\)。
再者，所有根 \(s_{ij}\in L\)，故 \(G\) 的分裂域 \(M\subseteq L\)。
而 \(S_4\) 在二子集上的作用是忠实作用，给出阶 \(24\) 的像；由 \(M\subseteq L\) 及 \([L:k]=24\) 得 \(M=L\)。
最后，\(S_4\) 的换位在二子集上的诱导置换是两个互换之积，为偶置换，故像落在 \(A_6\)。证毕。
\end{proof}

\begin{conclusion}[反相位点 \texorpdfstring{$u=-1$}{u=-1} 的刚性]\label{con:xi-terminal-zm-phase-compression-antiphase-rigidity}
对任意参数 \(y\) 有
\[
G(-2,y)=y^2(y-1)^2.
\]
因此 \(u=-1\)（即 \(s=-2\)）是方程 \(F(u,y)=0\) 的根当且仅当 \(y\in\{0,1\}\)。
特别地，Lee--Yang 条件 \(P_{\mathrm{LY}}(y)=0\) 不产生 \(u=-1\) 的相位共振根。
\end{conclusion}
\begin{proof}[证明要点]
由结论 \ref{con:xi-terminal-zm-phase-compression-disc-square}，\(G(-2,y)=y^2(y-1)^2\)。
又 \(F(-1,y)=(-1)^6G(-2,y)=G(-2,y)\)，故 \(F(-1,y)=0\iff y\in\{0,1\}\)。
当 \(P_{\mathrm{LY}}(y)=0\) 时 \(y\neq 0,1\)，从而不能触发 \(u=-1\) 根。证毕。
\end{proof}

\begin{conclusion}[黄金比例判别式分量的显式平方塌缩]\label{con:xi-terminal-zm-phase-compression-golden-square-collapse}
令
\[
\alpha:=\frac{-1-\sqrt5}{2},\qquad
\beta:=\frac{-1+\sqrt5}{2},
\]
则在 \(\QQ(\sqrt5)\) 上有精确因式分解
\[
G(s,\alpha)=\bigl(s^2+(1-2\sqrt5)s+7\bigr)\bigl(s^2+(2+\sqrt5)s-1+2\sqrt5\bigr)^2,
\]
\[
G(s,\beta)=\bigl(s^2+(1+2\sqrt5)s+7\bigr)\bigl(s^2+(2-\sqrt5)s-1-2\sqrt5\bigr)^2.
\]
\end{conclusion}
\begin{proof}[证明要点]
将 \(y=\alpha,\beta\) 代入显式 \(G(s,y)\)，并在 \(\QQ(\sqrt5)[s]\) 中进行代数因式分解即可得到两式。
该分解由脚本 \texttt{scripts/exp\_fold\_zm\_phase\_compression\_discriminant\_square\_audit.py} 直接核验。证毕。
\end{proof}

\begin{conclusion}[Lee--Yang 边界处的二重共振对机制]\label{con:xi-terminal-zm-phase-compression-leyang-double-resonance-pairs}
设 \(y_0\) 满足 \(P_{\mathrm{LY}}(y_0)=0\)，并记 \(k_0:=\QQ(y_0)\)。
则在 \(k_0[s]\) 中有
\[
G(s,y_0)=y_0^3(y_0+1)^3\,(s-2)\,L(s)\,Q(s)^2,
\]
其中
\[
L(s)=s+\frac{-1280y_0^2+2041y_0+6583}{3168},
\]
\[
Q(s)=s^2+\frac{6400y_0^2-10205y_0-1235}{3168}\,s+\frac{16640y_0^2-26533y_0-9547}{1584}.
\]
等价地，对首项归一化多项式 \(\widetilde G(s,y):=G(s,y)/(y^3(y+1)^3)\) 有
\[
\widetilde G(s,y_0)=(s-2)L(s)Q(s)^2.
\]
因此在 Lee--Yang 边界上，压缩六次多项式呈现 \((1,1,2,2)\) 根型。
\end{conclusion}
\begin{proof}[证明要点]
由结论 \ref{con:xi-terminal-zm-phase-compression-disc-square} 的端点恒等式，
\(
G(2,y)=-y(y-1)P_{\mathrm{LY}}(y)
\)，故 \(P_{\mathrm{LY}}(y_0)=0\Rightarrow s=2\) 为根。
将 \(y\) 约化到商域 \(\QQ[y]/(P_{\mathrm{LY}}(y))\) 后对 \(G\) 进行精确因式分解，可得所示
\((s-2)LQ^2\) 结构（对 \(G\) 本身仅差首项系数 \(y_0^3(y_0+1)^3\)）。
于是重根结构为一个二次因子的平方与两个一次因子，根型即 \((1,1,2,2)\)。
上述等式由脚本 \texttt{scripts/exp\_fold\_zm\_phase\_compression\_discriminant\_square\_audit.py} 核验。证毕。
\end{proof}

\begin{conclusion}[相位压缩曲线的正规化：属 \(2\) 与分歧签名由二子集表示完全决定]\label{con:xi-terminal-zm-phase-compression-curve-genus2}
令 \(X_{\mathrm{ph}}\) 为平面曲线 \(G(s,y)=0\) 的射影闭包之正规化，并令
\[
\pi:X_{\mathrm{ph}}\longrightarrow \PP^1_y
\]
为由 \(y\) 坐标诱导的有限态射。
则 \(\deg(\pi)=6\)，且 \(\pi\) 的分歧仅可能发生在
\(
y\in\{0,1,P_{\mathrm{LY}}(y)=0,\infty\}
\)
之上；其局部分歧型为：在每个有限分歧点上惯性置换在六张面上的循环型为 \((1^2\,2^2)\)，而在 \(y=\infty\) 上循环型为 \((4\cdot 2)\)。
此外，结论 \ref{con:xi-terminal-zm-phase-compression-disc-square} 的判别式分解在 \(y=-1\)、\(y^{2}+y-1=0\) 与 \(P_9(y)=0\) 处出现额外退化因子；这些对应于仿射平面模型的奇异纤维（结论 \ref{con:xi-terminal-zm-phase-compression-affine-singular-fibers}），而不产生新的域分歧点。
因此
\[
g(X_{\mathrm{ph}})=2.
\]
\end{conclusion}
\begin{proof}[证明要点]
由结论 \ref{con:xi-terminal-zm-phase-compression-galois-a6}，\(\QQ(X_{\mathrm{ph}})=\QQ(y,s_{12})\) 为 \(\Pi(\lambda,y)\) 的分裂域 \(L/\QQ(y)\) 的中间子扩张；
因此 \(X_{\mathrm{ph}}\to\PP^1_y\) 不可能在 \(L/\QQ(y)\) 的分歧集合之外引入新的域分歧。
又谱四次覆盖的 Hurwitz 分歧签名为 \((2,2,2,2,2,4)\)（结论 \ref{con:xi-terminal-zm-leyang-ramification-critical-values}），
将其中换位与四循环通过“对 \(2\)-子集的作用”诱导到次数 \(6\) 的表示，得到所述循环型 \((1^2\,2^2)\) 与 \((4\cdot 2)\)。
\par
对 \(\pi\) 应用 Riemann--Hurwitz 公式：
五个有限分歧点各贡献 \((2-1)+(2-1)=2\)，总计 \(10\)；
\(\infty\) 处贡献 \((4-1)+(2-1)=4\)，总分歧贡献 \(14\)。
于是
\(
2g(X_{\mathrm{ph}})-2=6(-2)+14
\),
从而 \(g(X_{\mathrm{ph}})=2\)。证毕。
\end{proof}

\begin{conclusion}[到 \texorpdfstring{$D_4$}{D4}--resolvent 椭圆曲线的分歧二重覆盖与 Jacobian 椭圆分解]\label{con:xi-terminal-zm-phase-compression-jacobian-splitting}
令 \(H:=\langle(12),(34)\rangle\le S_4\)，并令 \(D_4:=N_{S_4}(H)\) 为其正规化子（阶 \(8\) 的二面体 Sylow \(2\)-子群）。
取谱四次覆盖的 \(S_4\) 伽罗瓦闭包曲线 \(X_e\to\PP^1_y\)（结论 \ref{con:xi-terminal-zm-s4-d4fixedfield-resolvent-elliptic}），并记
\[
X_{\mathrm{ph}}:=X_e/H,\qquad R:=X_e/D_4.
\]
则存在自然次数 \(2\) 的态射
\[
\varphi:X_{\mathrm{ph}}\longrightarrow R
\]
使得 \(\pi\) 分解为 \(X_{\mathrm{ph}}\xrightarrow{\varphi}R\to\PP^1_y\)。
并且 \(\varphi\) 的分歧除子次数恰为 \(2\)，从而在 \(\QQ\) 上成立同源分解
\[
J(X_{\mathrm{ph}})\ \sim_{\QQ}\ R\times \operatorname{Prym}(\varphi),
\]
其中 \(\operatorname{Prym}(\varphi)\) 为一条定义于 \(\QQ\) 的椭圆曲线。
\end{conclusion}
\begin{proof}[证明要点]
由于 \(H\triangleleft D_4\) 且 \([D_4:H]=2\)，商映射 \(X_e/H\to X_e/D_4\) 给出 \(\varphi\)，其次数为 \(2\)。
由结论 \ref{con:xi-terminal-zm-phase-compression-curve-genus2} 得 \(g(X_{\mathrm{ph}})=2\)，而 \(R\) 为椭圆曲线（结论 \ref{con:xi-terminal-zm-s4-d4fixedfield-resolvent-elliptic}），故 \(g(R)=1\)。
Riemann--Hurwitz 对二重覆盖给出
\(
2g(X_{\mathrm{ph}})-2=2(2g(R)-2)+\deg\mathrm{Ram}(\varphi)
\),
从而 \(\deg\mathrm{Ram}(\varphi)=2\)。
最后，对任意二重覆盖 \(\varphi:C\to E\)（\(E\) 椭圆曲线），Prym 分解给出 \(J(C)\sim E\times \operatorname{Prym}(\varphi)\)，且 \(\dim\operatorname{Prym}=g(C)-g(E)=1\)，故 Prym 为椭圆曲线。证毕。
\end{proof}

\begin{conclusion}[平面相位压缩模型的奇异纤维投影：由 \texorpdfstring{$y=-1$}{y=-1}、黄金共轭与 \texorpdfstring{$P_9$}{P9} 控制]\label{con:xi-terminal-zm-phase-compression-affine-singular-fibers}
令 \(C_{\mathrm{aff}}\subset\mathbb A^2_{s,y}\) 为仿射曲线 \(G(s,y)=0\)。
则其奇异点 \((s_0,y_0)\) 必满足
\[
y_0(y_0-1)(y_0+1)(y_0^{2}+y_0-1)P_{9}(y_0)=0,
\]
其中 \(P_9\) 由结论 \ref{con:xi-terminal-zm-phase-compression-disc-square} 给出。
更精确地：
\begin{itemize}
  \item 在 \(y=0\) 与 \(y=1\) 上，唯一有限奇异点分别为 \((s,y)=(-2,0)\) 与 \((s,y)=(-2,1)\)。
  \item 在 \(y^{2}+y-1=0\) 的每个根上，存在恰好两个有限奇异点；其 \(s\)-坐标在 \(\QQ(\sqrt5)\) 中可取为
  \[
  s=\frac{\sqrt5-2}{2}\ \pm\ \frac12\sqrt{13+4\sqrt5},
  \]
  另一根对应其 \(\sqrt5\)-共轭。
  \item 在 \(P_9(y)=0\) 的每个根上，\(C_{\mathrm{aff}}\) 存在有限奇异点；其发生机制为在该参数处 \(G(\cdot,y)\) 出现二重 \(s\)-根。
\end{itemize}
\end{conclusion}
\begin{proof}[证明要点]
仿射奇异点当且仅当 \(G=\partial_sG=\partial_yG=0\)。
消去 \(s\) 并记
\[
D_1(y):=\operatorname{Res}_{s}\bigl(G,\partial_sG\bigr),\qquad
D_2(y):=\operatorname{Res}_{s}\bigl(G,\partial_yG\bigr).
\]
若存在奇异点，则必有 \(D_1(y)=D_2(y)=0\)。
由结论 \ref{con:xi-terminal-zm-phase-compression-disc-square} 中给出的
\(
\gcd(D_1,D_2)=y^{10}(y-1)^3(y+1)^6(y^2+y-1)^4P_9(y)^2
\)
可得 \(y_0\) 的必要条件即为所述零点集合。
\par
对 \(y=0,1\) 的断言由结论 \ref{con:xi-terminal-zm-phase-compression-disc-square} 中的分解
\(
G(s,0)=-(s-2)(s+2)^2
\)、\(
G(s,1)=(s-2)(s+2)^2(2s+5)^2(2s-5)
\)
直接检验三方程同解点得到。
对 \(y^{2}+y-1=0\) 的情形，在系数域 \(\QQ(\sqrt5)\) 上将 \(y\) 特化到任一根 \(y_0\) 后，
可计算得到
\(
\gcd_{\,\QQ(\sqrt5)[s]}\bigl(G(\cdot,y_0),\partial_sG(\cdot,y_0)\bigr)
\)
为一条可分二次多项式；例如当 \(y_0=\frac{-1+\sqrt5}{2}\) 时该因子可取为
\(
s^{2}-(\sqrt5-2)s-(1+2\sqrt5)
\)，另一根对应其 \(\sqrt5\)-共轭。
并且该二次因子整除 \(\partial_yG(\cdot,y_0)\)；
从而在该 \(y_0\) 上恰存在两条有限奇异分支，给出“恰好两个”有限奇异点。
所示 \(s\)-坐标即该二次多项式的求根公式展开。
对 \(P_9(y)=0\) 的情形，由 \(\gcd(D_1,D_2)\) 含因子 \(P_9(y)^2\) 可知每个 \(P_9\) 的根 \(y_0\) 处必存在 \(s_0\) 使 \(G=\partial_sG=\partial_yG=0\)，从而出现有限奇异点。
上述恒等式与投影必要条件可由脚本 \texttt{scripts/exp\_fold\_zm\_phase\_compression\_discriminant\_square\_audit.py} 复核。证毕。
\end{proof}

\begin{conclusion}[三重等模点的九次整系数刻画与算术不变量]\label{con:xi-terminal-zm-triple-equimodular-q9}
沿用结论 \ref{con:xi-terminal-zm-order4-rational} 的记号，令
\[
\Pi(\lambda,y):=\lambda^4-\lambda^3-(2y+1)\lambda^2+\lambda+y(y+1).
\]
则存在唯一参数
\[
y_\diamond\in(-1,0)
\]
使得 \(\Pi(\lambda,y_\diamond)\) 的谱半径由三个不同根同时取得：一条负实根 \(\lambda_-(y_\diamond)<0\) 与一对共轭复根 \(\lambda_c(y_\diamond),\overline{\lambda_c(y_\diamond)}\)，满足
\[
|\lambda_c(y_\diamond)|=-\lambda_-(y_\diamond)=:\rho_\diamond,
\qquad
|\lambda_+(y_\diamond)|<\rho_\diamond,
\]
其中 \(\lambda_+(y_\diamond)>0\) 为另一条实根。

该参数由如下九次整系数多项式在 \((-1,0)\) 上的唯一实根给出：
\[
Q(y)=256 y^9 + 1219 y^8 + 2542 y^7 + 3090 y^6 + 2446 y^5 + 1315 y^4 + 478 y^3 + 112 y^2 + 16 y + 1.
\]
并且 \(Q\) 在 \(\QQ[y]\) 上不可约，且
\[
\operatorname{Disc}(Q)=2^{28}\cdot 3^{9}\cdot 37\cdot 43^2\cdot 211\cdot 82301.
\]

数值上，
\[
y_\diamond\approx -0.160091463528933,\qquad
\rho_\diamond\approx 0.956702551378141,
\]
\[
\lambda_-(y_\diamond)\approx-0.956702551378141,\qquad
\lambda_c(y_\diamond)\approx 0.901572819032290\pm 0.320072216531856\,i,
\]
\[
\lambda_+(y_\diamond)\approx 0.153556913313561.
\]
\end{conclusion}
\begin{proof}[证明要点]
在 \(y\in(-1,0)\) 且 \(\operatorname{Disc}_{\lambda}(\Pi)\neq 0\) 的区域，\(\Pi\) 有两条实根与一对共轭复根且均为单根。
设四根记为 \((\lambda,b,\mu,\overline\mu)\)（\(\lambda<0<b\)），并令 \(u:=y(y+1)\)。
等模条件
\(
|\mu|=|\lambda|=-\lambda
\)
等价于 \(\mu\overline\mu=\lambda^2\)，由韦达关系得
\[
b=\frac{u}{\lambda^3},
\qquad
\sigma_3=\lambda b(\mu+\overline\mu)+b\mu\overline\mu=-1.
\]
消去 \(\mu+\overline\mu\) 后可得必要且充分的代数条件
\[
S(\lambda,y):=\lambda^8+\lambda^5+u\lambda^3-u^2=0.
\]
于是三重等模点满足联立方程
\[
\Pi(\lambda,y)=0,\qquad S(\lambda,y)=0,\qquad \lambda<0,\ y\in(-1,0).
\]
对 \(\lambda\) 取结果式：
\[
R(y):=\operatorname{Res}_{\lambda}\bigl(\Pi(\lambda,y),S(\lambda,y)\bigr)
=y(y-1)(y+1)\bigl(y^3+3y^2+y+1\bigr)^2Q(y).
\]
在 \((-1,0)\) 内可排除 \(y=0,\pm 1\) 与三次因子实根（该实根位于 \((- \infty,-1)\)），故三重等模条件在该区间内等价于 \(Q(y)=0\)。
由 Sturm 链可知 \(Q\) 在 \((-1,0)\) 上恰有一个实根，即 \(y_\diamond\)。
再由谱分支连续性与数值核验可得 \(|\lambda_+(y_\diamond)|<\rho_\diamond\)，故谱半径恰由三根同时取得。
不可约性与判别式可由 \(\QQ[y]\) 上因式分解及判别式计算直接验证。证毕。
\end{proof}

\begin{conclusion}[振荡主导区间上的实零点量子化律与指数误差]\label{con:xi-terminal-zm-leyang-real-zero-quantization}
沿用结论 \ref{con:xi-terminal-zm-triple-equimodular-q9} 的记号，取任意紧区间
\[
I=[a,b]\subset(-\infty,y_\diamond),
\]
并令 \(\lambda_c(y)\) 为上半平面中模最大的特征根分支，写作
\[
\lambda_c(y)=\rho(y)e^{i\theta(y)},\qquad \rho(y)>0,\ \theta(y)\in(0,\pi).
\]
则存在 \(\eta\in(0,1)\) 与在 \(I\) 上解析且处处非零的函数 \(A(y)\)、\(\phi(y)\)，使得一致地
\[
Z_m(y)
=
2\,\rho(y)^m\,|A(y)|\cos\!\bigl(m\theta(y)+\phi(y)\bigr)
+ O\!\bigl((\eta\rho(y))^m\bigr)
\qquad (m\to\infty,\ y\in I).
\]

若进一步假设 \(\theta'(y)\neq 0\) 于 \(I\) 上成立，则对充分大 \(m\)，\(Z_m\) 在 \(I\) 内存在一族互异实零点 \(\{y_{m,k}\}\)，并满足
\[
m\theta(y_{m,k})+\phi(y_{m,k})
=
\left(k+\frac12\right)\pi + O(\eta^m).
\]
这些零点均为单根，且计数满足
\[
\#\{y_{m,k}\in I\}
=
\frac{m}{\pi}\bigl(\theta(b)-\theta(a)\bigr)+O(1).
\]
对应经验测度弱收敛到绝对连续测度
\[
\dd \nu_I(y)=\frac{1}{\pi}\,|\theta'(y)|\,\dd y.
\]
\end{conclusion}
\begin{proof}[证明要点]
由结论 \ref{con:xi-terminal-zm-order4-rational}，
\[
\sum_{m\ge 0}Z_m(y)t^m=\frac{N(t,y)}{D(t,y)},
\]
\[
N(t,y)=1+yt-yt^2-y^2t^3,\qquad
D(t,y)=1-t-(2y+1)t^2+t^3+y(y+1)t^4.
\]
在判别式非零处，\(D\) 的根 \(t_j(y)\) 均为单根，令 \(\lambda_j(y):=1/t_j(y)\)，由留数展开
\[
Z_m(y)=\sum_{j=1}^4 c_j(y)\lambda_j(y)^m,\qquad
c_j(y)=-\lambda_j(y)\frac{N(1/\lambda_j(y),y)}{\partial_t D(1/\lambda_j(y),y)}.
\]
并且
\[
\operatorname{Res}_{t}(N,D)=-y^4(y-1)^3,
\]
故在 \(I\subset(-\infty,y_\diamond)\) 上（尤其 \(y\neq 0,1\)）各 \(c_j\) 解析且不消失。
由 \(y_\diamond\) 的定义，\(I\) 上主导谱为 \(\lambda_c,\overline{\lambda_c}\)，且与其余两根存在一致谱隙：存在 \(\eta\in(0,1)\) 使
\[
\max_{j\notin\{c,\overline c\}}|\lambda_j(y)|\le \eta\,|\lambda_c(y)|,\qquad y\in I.
\]
因此
\[
Z_m(y)=c_c(y)\lambda_c(y)^m+\overline{c_c(y)\lambda_c(y)^m}+O((\eta\rho(y))^m),
\]
写 \(c_c=|A|e^{i\phi}\) 即得主式。

若 \(\theta'(y)\neq 0\)，则
\(
g_m(y):=m\theta(y)+\phi(y)
\)
在大 \(m\) 时严格单调，主项 \(\cos g_m\) 的符号翻转与小误差叠加给出唯一零点。
对零点处导数使用
\[
\frac{\dd}{\dd y}\cos g_m(y)=-g_m'(y)\sin g_m(y),
\]
并由 \(g_m'(y)=m\theta'(y)+O(1)\) 得单根性。
量子化关系由将误差吸收入相位获得；计数与极限测度由 Riemann 和逼近得到。证毕。
\end{proof}

\begin{conclusion}[零点密度的显式代数微分表达与椭圆谱几何解释]\label{con:xi-terminal-zm-density-elliptic-dlog}
在结论 \ref{con:xi-terminal-zm-leyang-real-zero-quantization} 的设定下，
\[
\theta'(y)
=
\Im\!\left(\frac{\lambda_c'(y)}{\lambda_c(y)}\right)
=
\Im\!\left(
\frac{2\lambda_c(y)^2-(2y+1)}
{\lambda_c(y)\bigl(4\lambda_c(y)^3-3\lambda_c(y)^2-(4y+2)\lambda_c(y)+1\bigr)}
\right).
\]
因此零点极限测度密度可写为
\[
\dd\nu_I(y)
=
\frac{1}{\pi}
\left|
\Im\!\left(
\frac{2\lambda_c(y)^2-(2y+1)}
{\lambda_c(y)\bigl(4\lambda_c(y)^3-3\lambda_c(y)^2-(4y+2)\lambda_c(y)+1\bigr)}
\right)
\right|
\dd y.
\]

进一步，在谱曲线 \(\mathcal C:\Pi(\lambda,y)=0\) 上定义亚纯 \(1\)-形式
\[
\omega:=\dd\log\lambda
=
-\frac{\partial_y\Pi(\lambda,y)}{\lambda\,\partial_\lambda\Pi(\lambda,y)}\,\dd y.
\]
则沿主导复支有
\[
\dd\theta=\Im(\omega),\qquad
\nu_I=\frac{1}{\pi}\,(\pi_y)_*\!\bigl(|\Im(\omega)|\bigr).
\]
\end{conclusion}
\begin{proof}[证明要点]
由隐函数关系 \(\Pi(\lambda_c(y),y)=0\) 可得
\[
\partial_\lambda\Pi(\lambda_c(y),y)\lambda_c'(y)+\partial_y\Pi(\lambda_c(y),y)=0,
\]
\[
\frac{\lambda_c'(y)}{\lambda_c(y)}
=
-\frac{\partial_y\Pi(\lambda_c(y),y)}{\lambda_c(y)\partial_\lambda\Pi(\lambda_c(y),y)}.
\]
代入
\[
\partial_y\Pi(\lambda,y)=-2\lambda^2+2y+1,\qquad
\partial_\lambda\Pi(\lambda,y)=4\lambda^3-3\lambda^2-(4y+2)\lambda+1
\]
即得 \(\theta'(y)\) 的显式公式。
再由结论 \ref{con:xi-terminal-zm-leyang-real-zero-quantization} 的
\(
\dd\nu_I=(1/\pi)|\theta'(y)|\dd y
\)
得到密度表达。
最后 \(\dd\theta=\Im(\dd\log\lambda)=\Im(\omega)\)，故推前测度表述成立。证毕。
\end{proof}

\begin{conclusion}[\texorpdfstring{$m^{-1}$}{1/m} 局部缩放下的三相干涉极限方程]\label{con:xi-terminal-zm-triple-equimodular-local-scaling}
令 \(y_\diamond\) 为结论 \ref{con:xi-terminal-zm-triple-equimodular-q9} 的三重等模点，并记
\[
\lambda_-(y_\diamond)=-\rho_\diamond,\qquad
\lambda_c(y_\diamond)=\rho_\diamond e^{i\theta_\diamond}.
\]
在谱展开
\(
Z_m(y)=\sum_{j=1}^4 c_j(y)\lambda_j(y)^m
\)
中，定义
\[
c_-^\diamond:=c_-(y_\diamond)\in\RR\setminus\{0\},\qquad
c_c^\diamond:=c_c(y_\diamond)\in\CC\setminus\{0\},
\]
\[
\alpha:=
\Re\!\left(\frac{\lambda_c'}{\lambda_c}-\frac{\lambda_-'}{\lambda_-}\right)\Big|_{y=y_\diamond},
\qquad
\beta:=
\Im\!\left(\frac{\lambda_c'}{\lambda_c}\right)\Big|_{y=y_\diamond},
\qquad
\phi_\diamond:=\arg(c_c^\diamond).
\]
对任意满足 \(m_n\theta_\diamond\equiv\delta\pmod{2\pi}\) 的整数子列 \(\{m_n\}\)（并可固定奇偶类），在任意紧致 \(\xi\)-区间上一致有
\[
\rho_\diamond^{-m_n}
e^{-\Re(\lambda_-'/\lambda_-)|_{y=y_\diamond}\,\xi}
Z_{m_n}\!\left(y_\diamond+\frac{\xi}{m_n}\right)
\longrightarrow
(-1)^{m_n}c_-^\diamond
+2|c_c^\diamond|e^{\alpha\xi}
\cos\!\bigl(\delta+\beta\xi+\phi_\diamond\bigr).
\]
因此，窗口 \(|y-y_\diamond|=O(m^{-1})\) 内零点在缩放变量 \(\xi\) 下收敛到方程
\[
(-1)^m c_-^\diamond
+2|c_c^\diamond|e^{\alpha\xi}
\cos\!\bigl(\delta+\beta\xi+\phi_\diamond\bigr)=0
\]
的解集。

数值上
\[
\alpha\approx -0.1820719187487,\qquad
\beta\approx -1.454672700610.
\]
\end{conclusion}
\begin{proof}[证明要点]
在 \(y_\diamond\) 的小邻域内，四条特征根保持单纯，故 \(\lambda_j(y)\)、\(c_j(y)\) 解析。
令 \(y=y_\diamond+\xi/m\)，则
\[
\lambda_j(y)
=
\lambda_j(y_\diamond)\exp\!\left(
\frac{\lambda_j'(y_\diamond)}{\lambda_j(y_\diamond)}\frac{\xi}{m}
+O(m^{-2})
\right),
\]
\[
\lambda_j(y)^m
=
\lambda_j(y_\diamond)^m
\exp\!\left(
\frac{\lambda_j'(y_\diamond)}{\lambda_j(y_\diamond)}\xi
+O(m^{-1})
\right),
\]
且 \(c_j(y)=c_j(y_\diamond)+O(m^{-1})\)。
代回谱展开后，第四根满足 \(|\lambda_+(y_\diamond)|<\rho_\diamond\)，故在 \(\rho_\diamond^{-m}\) 归一化下指数消失。
其余三根中，\(\lambda_-=-\rho_\diamond\) 贡献 \((-1)^m\) 相位，而 \(\lambda_c,\overline{\lambda_c}\) 合并为余弦项并带指数倾斜 \(e^{\alpha\xi}\) 与线性相位 \(\beta\xi\)。
沿 \(m_n\theta_\diamond\to\delta\ (\mathrm{mod}\ 2\pi)\) 的子列即得一致极限。
零点极限方程由极限函数的过零集合直接给出。证毕。
\end{proof}


% (User input window)

\begin{conclusion}[\texorpdfstring{$(q=2)$}{(q=2)} 的重量—碰撞配分函数：六阶常系数递推与有理生成函数]\label{con:xi-terminal-zm2-order6-rational}
令
\(
X_m=\{x\in\{0,1\}^m:\ x_ix_{i+1}=0\}
\)，
并记 Fold 纤维多重度 \(d_m(x)=|\Fold_m^{-1}(x)|\)。
定义二阶重量—碰撞配分函数
\[
Z_m^{(2)}(y):=\sum_{x\in X_m} y^{|x|_1}\,d_m(x)^2\in\ZZ[y].
\]
则对一切 \(m\ge 7\) 成立严格六阶递推
\[
\begin{aligned}
Z_m^{(2)}(y)=&\ Z_{m-1}^{(2)}(y)+(2+3y)Z_{m-2}^{(2)}(y)+(y-2)Z_{m-3}^{(2)}(y)\\
&\ +(-1-2y-3y^2)Z_{m-4}^{(2)}(y)+(1-y)Z_{m-5}^{(2)}(y)+(y+y^3)Z_{m-6}^{(2)}(y),
\end{aligned}
\]
并且其常规生成函数为显式有理函数
\[
\boxed{\;
\sum_{m\ge 0} Z_m^{(2)}(y)t^m
=\frac{1+yt+(1-2y)t^2+(y-2y^2)t^3+(y^2-y)t^4+(y^3-y^2)t^5}{1-t-(2+3y)t^2+(2-y)t^3+(1+2y+3y^2)t^4-(1-y)t^5-(y+y^3)t^6}\; }.
\]
其中可取初值
\[
Z_0^{(2)}(y)=1,\ 
Z_1^{(2)}(y)=y+1,\ 
Z_2^{(2)}(y)=2y+4,\ 
Z_3^{(2)}(y)=y^2+9y+4,
\]
\[
Z_4^{(2)}(y)=6y^2+21y+9,\ 
Z_5^{(2)}(y)=y^3+31y^2+47y+9,\ 
Z_6^{(2)}(y)=15y^3+96y^2+93y+16.
\]
\end{conclusion}
\begin{proof}[证明要点]
与结论 \ref{con:xi-terminal-zm-order4-rational} 同理，
\(Z_m^{(2)}(y)\) 是有限状态规约器的二阶碰撞计数在重量变量 \(y^{|x|_1}\) 下的扭曲求和，
故其 \(t\)-生成函数为有理函数，约分后分母次数等于最小线性实现的维数，从而诱导常系数递推。
分子分母的显式系数可由有限窗口精确枚举与有理重建（Pad\'e/Prony）恢复，并可由脚本
\texttt{scripts/exp\_fold\_zq\_weight\_collision\_partition\_audit.py} 一键复核。证毕。
\end{proof}

\begin{conclusion}[\texorpdfstring{$(q=3)$}{(q=3)} 的重量—碰撞配分函数：八阶常系数递推与有理生成函数]\label{con:xi-terminal-zm3-order8-rational}
定义三阶重量—碰撞配分函数
\[
Z_m^{(3)}(y):=\sum_{x\in X_m} y^{|x|_1}\,d_m(x)^3\in\ZZ[y].
\]
则对一切 \(m\ge 9\) 成立严格八阶递推
\[
\begin{aligned}
Z_m^{(3)}(y)=&\ Z_{m-1}^{(3)}(y)+(3+4y)Z_{m-2}^{(3)}(y)+(4y-3)Z_{m-3}^{(3)}(y)\\
&\ +(-3+y-6y^2)Z_{m-4}^{(3)}(y)+(3-6y+y^2)Z_{m-5}^{(3)}(y)\\
&\ +(1+2y-5y^2+4y^3)Z_{m-6}^{(3)}(y)+(-1+2y-y^2)Z_{m-7}^{(3)}(y)\\
&\ +(-y^4+y^3+y^2-y)Z_{m-8}^{(3)}(y),
\end{aligned}
\]
并且生成函数可写为
\[
\sum_{m\ge 0} Z_m^{(3)}(y)t^m=\frac{N_3(t,y)}{D_3(t,y)},
\]
其中
\[
N_3(t,y)=1+yt+(4-3y)t^2+(4y-3y^2)t^3+(1-6y+3y^2)t^4+(y-6y^2+3y^3)t^5-y(y-1)^2t^6-y^2(y-1)^2t^7,
\]
并且
\[
\begin{aligned}
D_3(t,y)=&\ 1-t-(3+4y)t^2-(4y-3)t^3+(6y^2-y+3)t^4-(y^2-6y+3)t^5\\
&\ -(4y^3-5y^2+2y+1)t^6+(y-1)^2t^7+(y^4-y^3-y^2+y)t^8.
\end{aligned}
\]
其中可取初值
\[
Z_0^{(3)}(y)=1,\ 
Z_1^{(3)}(y)=y+1,\ 
Z_2^{(3)}(y)=2y+8,\ 
Z_3^{(3)}(y)=y^2+17y+8,
\]
\[
Z_4^{(3)}(y)=10y^2+51y+27,\ 
Z_5^{(3)}(y)=y^3+79y^2+153y+27,
\]
\[
Z_6^{(3)}(y)=37y^3+324y^2+395y+64,\ 
Z_7^{(3)}(y)=y^4+298y^3+1304y^2+837y+64,\ 
Z_8^{(3)}(y)=101y^4+1755y^3+3965y^2+1822y+125.
\]
\end{conclusion}
\begin{proof}[证明要点]
同上：有限状态碰撞计数在重量变量下保持有限维线性实现，
从而给出有理生成函数与常系数递推；显式系数与初值由有限数据重建与枚举核验得到，
并由脚本 \texttt{scripts/exp\_fold\_zq\_weight\_collision\_partition\_audit.py} 输出可审计证书。证毕。
\end{proof}

\begin{conclusion}[\texorpdfstring{$(q=4)$}{(q=4)} 的重量—碰撞配分函数：十阶常系数递推与核多项式的在位回收]\label{con:xi-terminal-zm4-order10-rational}
定义四阶重量—碰撞配分函数
\[
Z_m^{(4)}(y):=\sum_{x\in X_m} y^{|x|_1}\,d_m(x)^4\in\ZZ[y].
\]
则对一切 \(m\ge 11\) 成立严格十阶递推
\[
\begin{aligned}
Z_m^{(4)}(y)=&\ Z_{m-1}^{(4)}(y)+(4+5y)Z_{m-2}^{(4)}(y)+(11y-4)Z_{m-3}^{(4)}(y)\\
&\ +(-10y^2+19y-6)Z_{m-4}^{(4)}(y)+(11y^2-18y+6)Z_{m-5}^{(4)}(y)\\
&\ +(10y^3-27y^2+2y+4)Z_{m-6}^{(4)}(y)+(y^3-13y^2+9y-4)Z_{m-7}^{(4)}(y)\\
&\ +(-5y^4+5y^3+2y^2-3y-1)Z_{m-8}^{(4)}(y)+(-y^3+2y^2-2y+1)Z_{m-9}^{(4)}(y)\\
&\ +(y^5-y^4+2y^3-y^2+y)Z_{m-10}^{(4)}(y),
\end{aligned}
\]
并且生成函数 \(F_4(t,y)=\sum_{m\ge 0}Z_m^{(4)}(y)t^m\) 为有理函数 \(N_4(t,y)/D_4(t,y)\)，其中分母由递推系数给出
\[
D_4(t,y)=1-\sum_{j=1}^{10}a_j(y)t^j
\]
（\(a_j(y)\) 为上式右端的十个系数多项式），而分子可取显式整系数形式
\[
\begin{aligned}
N_4(t,y)=&\ 1+yt+(11-4y)t^2+(11y-4y^2)t^3+(11-18y+6y^2)t^4\\
&\ +(11y-18y^2+6y^3)t^5+(1-13y+9y^2-4y^3)t^6+(y-13y^2+9y^3-4y^4)t^7\\
&\ +y(y-1)(y^2-y+1)t^8+y^2(y-1)(y^2-y+1)t^9.
\end{aligned}
\]
特别地，取 \(y=1\)（忘却重量层，仅保留四阶碰撞矩）时，分母在 \(t\)-变量上精确析出不可约核因子
\[
2t^5-2t^4-7t^2-2t+1,
\]
并且该因子在 \(F_4(t,1)\) 中保持在位（不被分子约去），从而以代数不可约方式回收为四阶碰撞核的本征结构常数。
\end{conclusion}
\begin{proof}[证明要点]
有理性与递推由有限状态碰撞核的有限维线性实现给出；
分子分母的显式形式及 \(y=1\) 的核因子析出可由脚本
\texttt{scripts/exp\_fold\_zq\_weight\_collision\_partition\_audit.py} 在 \(\ZZ[y,t]\) 中精确化简并输出证书。证毕。
\end{proof}

\paragraph{四阶 moment-kernel 的分裂域与符号二次脊柱}
令
\[
P^{\mathrm{mom}}_{4}(t):=2t^{5}-2t^{4}-7t^{2}-2t+1\in\ZZ[t],
\qquad
F_4(t,1)=\frac{7t^2+1}{P^{\mathrm{mom}}_{4}(t)}.
\]

\begin{proposition}[四阶 moment-kernel 的满对称伽罗瓦性]\label{prop:xi-terminal-zm4-moment-kernel-s5}
\leanverified{paper\_xi\_terminal\_zm4\_moment\_kernel\_s5}
\(P^{\mathrm{mom}}_{4}\) 在 \(\QQ[t]\) 上不可约。记其分裂域为
\(
K:=\mathrm{Spl}\bigl(P^{\mathrm{mom}}_{4}\bigr)
\)，
则
\[
\Gal(K/\QQ)\cong S_{5},
\qquad
\Disc_{t}\!\bigl(P^{\mathrm{mom}}_{4}\bigr)=-2^{4}\cdot 985219.
\]
\end{proposition}
\begin{proof}
置
\(
p_7(x):=x^{5}-2x^{4}-7x^{3}-2x+2
\)。
直接验算得互反关系
\[
p_7(x)=x^{5}P^{\mathrm{mom}}_{4}(1/x).
\]
因此两者根集互为倒数，分裂域相同，且判别式相等。
由命题 \ref{prop:pom-a4-galois-s5} 得 \(\Gal(\mathrm{Spl}(p_7)/\QQ)\cong S_5\) 且 \(\Disc_x(p_7)=-2^4\cdot 985219\)，从而结论成立。证毕。
\end{proof}

\begin{proposition}[唯一二次子域与分歧刚性]\label{prop:xi-terminal-zm4-moment-kernel-unique-quadratic}
\leanverified{paper\_xi\_terminal\_zm4\_moment\_kernel\_unique\_quadratic}
在命题 \ref{prop:xi-terminal-zm4-moment-kernel-s5} 的记号下，
\(
K/\QQ
\)
的唯一非平凡正规子扩张为
\[
K^{A_5}=\QQ\!\left(\sqrt{\Disc_{t}(P^{\mathrm{mom}}_{4})}\right)=\QQ(\sqrt{-985219}).
\]
并且若 \(L/\QQ\) 为伽罗瓦扩张且在素数 \(985219\) 处不分歧，则 \(K\cap L=\QQ\)。
\end{proposition}
\begin{proof}
由 \(\Gal(K/\QQ)\cong S_5\) 及 \(A_5\triangleleft S_5\) 的唯一指数 \(2\) 正规性，\(K\) 的非平凡正规子扩张只能为 \(K^{A_5}\)，而符号表示对应的二次子域由判别式平方根生成，故第一条成立。
\par
设 \(L/\QQ\) 伽罗瓦，则交域 \(E:=K\cap L\) 亦为 \(\QQ\) 上伽罗瓦扩张，且 \(E\subseteq K\)。
由第一条，\(E\in\{\QQ,\QQ(\sqrt{-985219})\}\)。
若 \(E=\QQ(\sqrt{-985219})\)，则 \(L\) 必包含在 \(985219\) 处分歧的二次域，矛盾；故 \(E=\QQ\)。证毕。
\end{proof}

\begin{corollary}[三通道 Chebotarev 乘积律：\texorpdfstring{$S_5\times S_4\times S_3$}{S5xS4xS3} 的联合分布]\label{cor:xi-zm4-moment-kernel-leyang-s4-s3-s5-chebotarev-triple-product}
令 \(K=\mathrm{Spl}(P_4^{\mathrm{mom}})\)。
令
\(
L_{\mathrm{LY}}:=\mathrm{Spl}(P_{\mathrm{LY}})
\)
为 Lee--Yang 三次因子分裂域（结论 \ref{con:xi-terminal-zm-leyang-cubic-arithmetic}，\(\Gal(L_{\mathrm{LY}}/\QQ)\cong S_3\)）。
取有理参数 \(y_0=a/b\in\QQ\)（\(\gcd(a,b)=1,\ b>0\)），并令
\(
P_{a,b}(\lambda)=b^{2}\Pi(\lambda,a/b)\in\ZZ[\lambda]
\)
如结论 \ref{con:xi-terminal-zm-specialization-discriminant-ab}。
设其分裂域 \(L_4:=\mathrm{Spl}(P_{a,b})\) 满足 \(\Gal(L_4/\QQ)\cong S_4\)，并且
\[
985219\nmid \Disc_{\lambda}(P_{a,b}),
\qquad
\QQ\!\left(\sqrt{\Disc_{\lambda}(P_{a,b})}\right)\neq \QQ(\sqrt{-111}).
\]
则
\[
\Gal(KL_4L_{\mathrm{LY}}/\QQ)\ \cong\ S_{5}\times S_{4}\times S_{3}.
\]
因此对除去有限坏素集合后的素数 \(p\)，三域中 Frobenius 共轭类的联合分布严格乘积化：任取共轭类
\(
C_5\subset S_5
\)、\(
C_4\subset S_4
\)、\(
C_3\subset S_3
\)，满足
\[
\mathrm{Frob}_p(K)\in C_5,\ \mathrm{Frob}_p(L_4)\in C_4,\ \mathrm{Frob}_p(L_{\mathrm{LY}})\in C_3
\]
的素数集合具有自然密度
\[
\delta
=\frac{|C_5|}{|S_5|}\cdot\frac{|C_4|}{|S_4|}\cdot\frac{|C_3|}{|S_3|}.
\]
\end{corollary}
\begin{proof}
由命题 \ref{prop:xi-terminal-zm4-moment-kernel-unique-quadratic}，由于 \(985219\nmid \Disc_{\lambda}(P_{a,b})\)，得 \(K\cap L_4=\QQ\)；同理 \(K\cap L_{\mathrm{LY}}=\QQ\)。
另一方面，\(L_4/\QQ\) 为 \(S_4\)-伽罗瓦扩张，其唯一二次子域为 \(\QQ(\sqrt{\Disc_{\lambda}(P_{a,b})})\)；
而 \(L_{\mathrm{LY}}/\QQ\) 的唯一二次子域为 \(\QQ(\sqrt{-111})\)（结论 \ref{con:xi-terminal-zm-leyang-cubic-arithmetic}）。
由假设 \(\QQ(\sqrt{\Disc_{\lambda}(P_{a,b})})\neq \QQ(\sqrt{-111})\)，推出 \(L_4\cap L_{\mathrm{LY}}=\QQ\)。
于是三者两两线性互素，从而
\(
\Gal(KL_4L_{\mathrm{LY}}/\QQ)\cong \Gal(K/\QQ)\times \Gal(L_4/\QQ)\times \Gal(L_{\mathrm{LY}}/\QQ)
\)
并得到直积同构。
对不分歧素数应用 Chebotarev 密度定理，直积群上的共轭类为笛卡尔积，密度即为共轭类比例之乘积。证毕。
\end{proof}

\begin{theorem}[四阶碰撞矩的主极点常数与对数增长率]\label{thm:xi-zm4-moment-kernel-residue-asymptotics}
\leanverified{paper\_xi\_zm4\_moment\_kernel\_residue\_asymptotics}
令
\(
F_4(t,1)=\sum_{m\ge 0}a_m t^m
\)
如上，并令 \(r\in(0,1)\) 为 \(P^{\mathrm{mom}}_{4}\) 在正实轴上的最小模根，置 \(\rho:=r^{-1}\)。
则 \(t=r\) 为 \(F_4(t,1)\) 的简单极点，且存在常数 \(C>0\) 与 \(0<\alpha<\rho\) 使得
\[
a_m=C\,\rho^{m}+O(\alpha^{m})\qquad(m\to\infty),
\qquad
C=-\frac{7r^{2}+1}{r\,(P^{\mathrm{mom}}_{4})'(r)}.
\]
并且对数增长率极限存在并等于
\[
\kappa_{4}:=\lim_{m\to\infty}\frac1m\log a_m=\log\rho.
\]
\end{theorem}
\begin{proof}
由于 \(P^{\mathrm{mom}}_{4}\) 可分离且 \(r\) 为其根，\(F_4(t,1)\) 在 \(t=r\) 处为单极，其留数由
\(
\mathrm{Res}_{t=r}\frac{7t^2+1}{P_4^{\mathrm{mom}}(t)}=\frac{7r^2+1}{(P_4^{\mathrm{mom}})'(r)}
\)
给出。
按部分分式分解（或 Cauchy 系数公式拾取主极点留数）可得主项系数
\(
-[t^m](t-r)^{-1}=r^{-m-1}
\)，从而得到所示 \(C\) 的表达。
\par
由命题 \ref{prop:pom-s4-recurrence} 给出的五态非负整数 realization，\(F_4(t,1)\) 的极点集合与该矩阵 \(A_4\) 的特征值集合互为倒数。
由 Perron--Frobenius，Perron 根 \(\rho=\rho(A_4)\) 为单根且严格支配其余特征值模长；因此除 \(t=r=\rho^{-1}\) 外的极点模长均严格大于 \(r\)。
取 \(\alpha:=\max_{\mu\in\sigma(A_4)\setminus\{\rho\}}|\mu|\)，则 \(\alpha<\rho\)（其数值证书见注 \ref{rem:pom-collision-rh-break-a4-spectrum}），并给出误差项 \(O(\alpha^m)\)。
最后，由 \(a_m=C\rho^m(1+o(1))\) 直接推出 \(\frac1m\log a_m\to \log\rho\)。证毕。
\end{proof}


\begin{remark}[\texorpdfstring{$y=1$}{y=1} 的幽灵模态全消与 moment-kernel 的严格同态化]\label{rem:xi-terminal-zq-ghost-cancellation-y1}
记 \(F_q(t,y)=\sum_{m\ge 0}Z_m^{(q)}(y)t^m=N_q(t,y)/D_q(t,y)\)。
在 \(y=1\) 处出现统一的“幽灵模态”严格约去：
\begin{enumerate}
  \item \(q=2\)：\(N_2(t,1)=-(t-1)(t+1)^2\)，且 \(D_2(t,1)=-(t-1)(t+1)^2(2t^3-2t^2-2t+1)\)，故
  \[
  F_2(t,1)=\frac{1}{2t^3-2t^2-2t+1}.
  \]
  \item \(q=3\)：\(N_3(t,1)=-(t-1)(t+1)^2(2t^2+1)\)，且 \(D_3(t,1)=-(t-1)(t+1)^2(2t^3-4t^2-2t+1)\)，故
  \[
  F_3(t,1)=\frac{2t^2+1}{2t^3-4t^2-2t+1}.
  \]
  \item \(q=4\)：\(N_4(t,1)=-(t-1)(t+1)^2(t^2+1)(7t^2+1)\)，且
  \(D_4(t,1)=-(t-1)(t+1)^2(t^2+1)(2t^5-2t^4-7t^2-2t+1)\)，故
  \[
  F_4(t,1)=\frac{7t^2+1}{2t^5-2t^4-7t^2-2t+1}.
  \]
\end{enumerate}
因此在 \(y=1\) 的“碰撞投影”下，所有额外的 \(\pm1\) 与 \(\pm i\) 幽灵模态均被分子严格抵消，
留下的唯一动力学分母恰为 moment-kernel 的核多项式；该约去机制把重量细分与碰撞核谱系在生成函数层面严格同态化。
\end{remark}

\begin{conclusion}[谱分歧判别式的结构分解：普适因子与割圆因子的出现]\label{con:xi-terminal-zq-discriminant-factorization}
令 \(\Pi_q(\lambda,y)\) 表示 \(q=2,3,4\) 的递推特征多项式（即 \(D_q(t,y)\) 在 \(t=1/\lambda\) 下清分母后得到的 \(\lambda\)-多项式）。
则其关于 \(\lambda\) 的判别式在 \(\ZZ[y]\) 中具有严格分解：
\begin{enumerate}
  \item \(q=2\)：
  \[
  \mathrm{Disc}_{\lambda}(\Pi_2)=4\,y^{3}(y-1)\,P_{9}(y),
  \]
  其中
  \[
  \begin{aligned}
  P_{9}(y)=&\ 186624y^9-695632y^8+980411y^7-1121640y^6+1021286y^5\\
  &-458803y^4+127602y^3-35316y^2-10260y+6912.
  \end{aligned}
  \]
  \item \(q=3\)：
  \[
  \mathrm{Disc}_{\lambda}(\Pi_3)=-16\,y^{5}(y-1)^3\,P_{17}(y),
  \]
  其中 \(P_{17}(y)\in\ZZ[y]\) 为次数 \(17\) 多项式，且具有唯一负实根 \(y\approx-1.07316827\)。
  \item \(q=4\)：
  \[
  \mathrm{Disc}_{\lambda}(\Pi_4)=9216\,y^{7}(y-1)(y^2-y+1)\,R_{31}(y),
  \]
  其中 \(R_{31}(y)\in\ZZ[y]\) 为次数 \(31\) 多项式。
\end{enumerate}
因此除去普适因子 \(y^{2q-1}\) 与 \(y=1\) 的投影退化外，
\(q=4\) 首次强制出现割圆因子 \(y^2-y+1\)（即 \(6\) 次单位根），
表明四阶重量细分在谱分歧层面已不可约地感知到单位圆对称通道的额外算术结构。
\end{conclusion}
\begin{proof}[证明要点]
将 \(\Pi_q(\lambda,y)\) 视为 \(\lambda\)-多项式并计算判别式，在 \(\ZZ[y]\) 中作因式分解即可得到分解结构；
剩余因子 \(P_{17},R_{31}\) 的显式系数可由同一分解过程唯一确定，并由脚本
\texttt{scripts/exp\_fold\_zq\_weight\_collision\_partition\_audit.py} 导出为可审计工件。证毕。
\end{proof}

\begin{conclusion}[零重量纤维的完全闭式：\texorpdfstring{$d_m(0^m)=\lfloor (m+2)/2\rfloor$}{d_m(0^m)=floor((m+2)/2)}]\label{con:xi-terminal-zero-weight-fiber-closed}
记 \(0^m\in X_m\) 为全零稳定类型。
则其纤维大小满足对一切 \(m\ge 1\) 的精确公式
\[
d_m(0^m)=\Bigl\lfloor\frac{m+2}{2}\Bigr\rfloor.
\]
从而对任意整数 \(q\ge 1\) 有
\[
Z_m^{(q)}(0)=d_m(0^m)^q=\Bigl\lfloor\frac{m+2}{2}\Bigr\rfloor^q,
\]
给出重量细分在 \(k=0\) 层上的完全可解子扇区；该层的增长为多项式阶而非指数阶，对应一个严格的“边界零熵”极端相。
\end{conclusion}
\begin{proof}
由 Fold 的定义（见脚本 \texttt{scripts/common\_phi\_fold.py} 的 Zeckendorf 正规形口径），
\(\Fold_m(a)=0^m\) 当且仅当微态数值 \(N(a)\) 的 Zeckendorf 数码在位置 \(1,\dots,m\) 全为 \(0\)；
等价地 \(N(a)\in\{0,F_{m+2}\}\)。
其中 \(N(a)=0\) 仅对应全零微态，贡献 \(1\)。
另一方面，设 \(r_m\) 为用权集合 \(\{F_2,\dots,F_{m+1}\}\) 表示 \(F_{m+2}\) 的 \(0/1\) 子集表示数。
由于 \(\sum_{i=2}^{m}F_i=F_{m+2}-1\)，任意表示必须包含 \(F_{m+1}\)，故 \(r_m\) 等于用 \(\{F_2,\dots,F_m\}\) 表示 \(F_m\) 的表示数。
对后者再作“最大项是否取用”的分解：若取用 \(F_m\) 得到一项表示；否则由于 \(\sum_{i=2}^{m-2}F_i=F_m-1\)，表示必包含 \(F_{m-1}\)，余项为 \(F_{m-2}\)，从而表示数满足递推 \(r_m=1+r_{m-2}\)（\(m\ge 3\)），且 \(r_1=0,r_2=1\)。
解得 \(r_m=\lfloor m/2\rfloor\)，因此
\(
d_m(0^m)=1+r_m=\lfloor (m+2)/2\rfloor
\)。
代回 \(Z_m^{(q)}(0)\) 即得结论。证毕。
\end{proof}

\begin{conclusion}[最大重量层的稳定词分类与纤维谱闭式：\texorpdfstring{$|x|_1=\lceil m/2\rceil$}{|x|=ceil(m/2)}]\label{con:xi-terminal-max-weight-layer-fiber-spectrum-closed}
定义最大重量层
\[
X_m^{\max}:=\Bigl\{x\in X_m:\ |x|_1=\Bigl\lceil\frac{m}{2}\Bigr\rceil\Bigr\}.
\]
则：
\begin{enumerate}
  \item 若 \(m=2k+1\) 为奇数，则 \(X_{2k+1}^{\max}=\{(10)^k1\}\)，并且
  \[
  d_{2k+1}\bigl((10)^k1\bigr)=1.
  \]
  \item 若 \(m=2k\) 为偶数，则 \(|X_{2k}^{\max}|=k+1\)，且可取显式参数化
  \[
  x^{(r)}:=(10)^r(01)^{k-r}\qquad(0\le r\le k),
  \]
  其中 \((10)^r\) 表示长度 \(2r\) 的交替前缀 \(10\cdots10\)。
  该层上的纤维多重度满足
  \[
  d_{2k}\bigl(x^{(0)}\bigr)=1,\qquad
  d_{2k}\bigl(x^{(r)}\bigr)=k-r+1\ \ (1\le r\le k),
  \]
  等价地，
  \[
  \{d_{2k}(x):x\in X_{2k}^{\max}\}=\{1,1,2,3,\dots,k\}.
  \]
\end{enumerate}
\end{conclusion}
\begin{proof}
关于 \(X_m^{\max}\) 的分类仅依赖合法语法约束 \(x_ix_{i+1}=0\)：为达到 \(\lceil m/2\rceil\) 的最大 \(1\)-计数，除边界外每个 \(1\) 必被 \(0\) 隔开，故词必须为交替结构；当 \(m\) 为奇数时交替词唯一，偶数时恰允许出现一次 \(00\) 缺陷，从而得到上述 \(x^{(r)}\) 族。
\par
纤维多重度闭式使用 Zeckendorf 值函数 \(V_m(x)=\sum_{i=1}^m x_iF_{i+1}\) 与微态值 \(N(\omega)=\sum_{i=1}^m\omega_iF_{i+1}\)（见脚本 \texttt{scripts/common\_phi\_fold.py} 的口径）。当 \(x\in X_m\) 时，
\(\Fold_m(\omega)=x\) 当且仅当 \(N(\omega)\) 的 Zeckendorf 数码在长度 \(m\) 上与 \(x\) 相同；因此在本结论所述的最大重量词上（均满足 \(V_m(x)<F_{m+2}\)），纤维大小等于
\[
d_m(x)=\#\Bigl\{\omega\in\{0,1\}^m:\ N(\omega)=V_m(x)\Bigr\}.
\]
下证偶数情形的闭式；奇数情形可由同一“表示唯一性”论证推出 \(d_{2k+1}((10)^k1)=1\)。
\par
令 \(m=2k\) 且 \(0\le r\le k\)。由 \(x^{(r)}\) 的交替结构与两类求和恒等式
\[
\sum_{j=1}^{r}F_{2j}=F_{2r+1}-1,\qquad
\sum_{j=r+1}^{k}F_{2j+1}=F_{2k+2}-F_{2r+2},
\]
得到其 Zeckendorf 值的闭式
\[
V_{2k}\bigl(x^{(r)}\bigr)=\sum_{j=1}^{r}F_{2j}+\sum_{j=r+1}^{k}F_{2j+1}
=F_{2k+2}-F_{2r}-1.
\]
因此，若令
\[
\mathcal{R}_{k,r}:=\Bigl\{S\subseteq\{2,3,\dots,2k+1\}:\ \sum_{n\in S}F_n=F_{2k+2}-F_{2r}-1\Bigr\},
\]
则 \(d_{2k}(x^{(r)})=\card{\mathcal{R}_{k,r}}\)。
\par
\noindent\textbf{(1) \(r=0\) 时的唯一性。}
此时目标和为 \(F_{2k+2}-1\)。若不取最大项 \(F_{2k+1}\)，则最大可得和为
\(\sum_{n=2}^{2k}F_n=F_{2k+2}-2\)，不足以达到目标；故每个表示都必须取用 \(F_{2k+1}\)。
从而表示集合与 \(\mathcal{R}_{k-1,0}\) 一一对应，且以 \(k=1\) 为初值知 \(\card{\mathcal{R}_{k,0}}=1\)。
\par
\noindent\textbf{(2) \(1\le r\le k\) 时的递推与闭式。}
把 \(\mathcal{R}_{k,r}\) 按是否取用最大项 \(F_{2k+1}\) 分解。若 \(F_{2k+1}\in S\)，则余项必须表示
\[
F_{2k+2}-F_{2r}-1-F_{2k+1}=F_{2k}-F_{2r}-1,
\]
并且由于 \(F_{2k}-F_{2r}-1<F_{2k}\)，该表示不可能再取用 \(F_{2k}\)，从而与 \(\mathcal{R}_{k-1,r}\)（在权集合 \(\{2,\dots,2k-1\}\) 上）严格一一对应。
\par
若 \(F_{2k+1}\notin S\)，则 \(S\subseteq\{2,\dots,2k\}\)，而此时目标和相对于全取用之最大和的缺口为
\[
\Bigl(\sum_{n=2}^{2k}F_n\Bigr)-\bigl(F_{2k+2}-F_{2r}-1\bigr)=F_{2r}-1.
\]
因此不取用 \(F_{2k+1}\) 的表示与对缺口 \(F_{2r}-1\) 的“删去子集”一一对应。又由贪心强制律：
若 \(T\subseteq\{2,\dots,2r-1\}\) 满足 \(\sum_{n\in T}F_n=F_{2r}-1\)，则必须有 \(F_{2r-1}\in T\)，并递归推出
\(
T=\{3,5,\dots,2r-1\}
\) 唯一。
从而 \(\mathcal{R}_{k,r}\) 中恰存在唯一一个不含 \(F_{2k+1}\) 的表示。
\par
综上，对 \(1\le r\le k\) 有递推
\(
\card{\mathcal{R}_{k,r}}=\card{\mathcal{R}_{k-1,r}}+1
\)
并以边界初值 \(\card{\mathcal{R}_{r,r}}=1\)（对应 \(F_{2r+1}-1\) 的唯一表示）解得
\(
\card{\mathcal{R}_{k,r}}=k-r+1
\)。
这与 \(d_{2k}(x^{(r)})=\card{\mathcal{R}_{k,r}}\) 合并即得结论。
\end{proof}

\begin{conclusion}[重量—碰撞配分函数的顶层系数：最大重量层强制 Faulhaber 幂和]\label{con:xi-terminal-zq-top-coefficient-faulhaber}
对整数 \(q\ge 1\) 定义重量—碰撞配分函数
\[
Z_m^{(q)}(y):=\sum_{x\in X_m}y^{|x|_1}\,d_m(x)^q\in\ZZ[y].
\]
则其顶层系数（即 \(y^{\lceil m/2\rceil}\) 的系数）具有完全闭式：
\begin{enumerate}
  \item 若 \(m=2k+1\)，则 \([y^{k+1}]Z_{2k+1}^{(q)}(y)=1\)。
  \item 若 \(m=2k\)，则
  \[
  [y^{k}]Z_{2k}^{(q)}(y)=1+\sum_{j=1}^{k}j^{q}.
  \]
\end{enumerate}
\end{conclusion}
\begin{proof}
顶层系数等于最大重量层上幂和 \(\sum_{x\in X_m^{\max}}d_m(x)^q\)。由结论 \ref{con:xi-terminal-max-weight-layer-fiber-spectrum-closed} 的谱 \(\{1\}\)（奇数）与 \(\{1,1,2,\dots,k\}\)（偶数）直接得到两式。证毕。
\end{proof}

\begin{conclusion}[顶层扇区的生成函数：分母阶为 \texorpdfstring{$(1-t)^{q+2}$}{(1-t)^{q+2}}]\label{con:xi-terminal-max-weight-top-sector-ogf}
令
\(
b_k^{(q)}:=[y^{k}]Z_{2k}^{(q)}(y)=1+\sum_{j=1}^{k}j^{q}
\)
（\(q\ge 1\) 固定）。
则其普通生成函数满足严格恒等式
\[
\sum_{k\ge 0} b_k^{(q)}t^{k}
=\frac{1}{1-t}+\frac{1}{1-t}\sum_{j\ge 1}j^{q}t^{j}.
\]
从而 \(\sum_{k\ge 0} b_k^{(q)}t^{k}\) 为有理函数，且在 \(t=1\) 处具有阶恰为 \(q+2\) 的极点（因此分母可取为 \((1-t)^{q+2}\)）。
\end{conclusion}
\begin{proof}
由 \(b_k^{(q)}=1+\sum_{j=1}^k j^q\) 作求和交换：
\[
\sum_{k\ge 0}b_k^{(q)}t^k=\sum_{k\ge 0}t^k+\sum_{k\ge 0}\sum_{j=1}^{k}j^q t^k
=\frac{1}{1-t}+\sum_{j\ge 1}j^q\sum_{k\ge j}t^k,
\]
而 \(\sum_{k\ge j}t^k=t^j/(1-t)\) 得所示恒等式。
又 \(\sum_{j\ge 1}j^q t^j\) 为负整数阶 polylog，等价于一个以 \((1-t)^{-(q+1)}\) 为分母的有理函数；
乘以 \((1-t)^{-1}\) 后得到分母阶至多 \(q+2\)。
另一方面 \(b_k^{(q)}\sim k^{q+1}/(q+1)\)，故生成函数在 \(t=1\) 处的极点阶不能小于 \(q+2\)，从而阶恰为 \(q+2\)。证毕。
\end{proof}

\begin{conclusion}[硬边界事件的精确概率：\texorpdfstring{$H_m=\lceil m/2\rceil$}{H_m=ceil(m/2)} 的多项式前因子]\label{con:xi-terminal-hard-boundary-event-probability}
令 \(w_m(x)=d_m(x)/2^m\) 为均匀微态推前分布，取 \(X\sim w_m\)，并记 \(H_m:=|X|_1\)。
则最大重量事件 \(H_m=\lceil m/2\rceil\) 的概率具有闭式：
\begin{enumerate}
  \item \(m=2k+1\)：\(\mathbb{P}(H_m=k+1)=2^{-(2k+1)}\)；
  \item \(m=2k\)：\(\mathbb{P}(H_m=k)=\bigl(1+\tfrac{k(k+1)}{2}\bigr)\,2^{-2k}\)。
\end{enumerate}
并且当 \(k\to\infty\) 时
\[
\log \mathbb{P}(H_{2k}=k)= -2k\log 2 +2\log k +O(1),
\]
表明该硬边界处的指数速率固定为 \(\log 2\)，而前因子具有可计算的多项式阶 \(k^2\) 修正。
\end{conclusion}
\begin{proof}
由定义 \(\mathbb{P}(H_m=\ell)=2^{-m}\sum_{|x|_1=\ell}d_m(x)\)，右端即为 \(Z_m(y)\) 的 \(y^\ell\) 系数除以 \(2^m\)。
对 \(\ell=\lceil m/2\rceil\)，该系数等于结论 \ref{con:xi-terminal-zq-top-coefficient-faulhaber} 在 \(q=1\) 的顶层系数，代回即得两条闭式；渐近式由 \(\log(1+k(k+1)/2)=2\log k+O(1)\) 得到。证毕。
\end{proof}

\begin{conclusion}[最大重量层上的 \texorpdfstring{$q$}{q}-倾斜极限：\texorpdfstring{$J_k/k\Rightarrow \mathrm{Beta}(q+1,1)$}{Jk/k -> Beta(q+1,1)}]\label{con:xi-terminal-max-weight-escort-beta-limit}
在偶数长度 \(m=2k\) 上，考虑最大重量层条件下的 \(q\)-倾斜测度
\[
\mathbb{P}_{k,q}(x\in\cdot):=\frac{\sum_{x\in X_{2k}^{\max}} d_{2k}(x)^{q}\,\mathbf 1_{\{x\in\cdot\}}}{\sum_{x\in X_{2k}^{\max}} d_{2k}(x)^{q}}
\qquad(q\ge 1).
\]
令 \(X\sim\mathbb{P}_{k,q}\)，并设 \(J_k:=d_{2k}(X)\)。
则由结论 \ref{con:xi-terminal-max-weight-layer-fiber-spectrum-closed} 的谱 \(\{1,1,2,\dots,k\}\) 有
\[
\mathbb{P}(J_k=j)=\frac{j^{q}}{1+\sum_{i=1}^{k} i^{q}}\ (2\le j\le k),\qquad
\mathbb{P}(J_k=1)=\frac{2}{1+\sum_{i=1}^{k} i^{q}}.
\]
并且在尺度极限下有弱收敛
\[
\frac{J_k}{k}\ \Rightarrow\ \mathrm{Beta}(q+1,1),
\qquad \text{密度为 }(q+1)x^{q}\mathbf 1_{[0,1]}(x).
\]
\end{conclusion}
\begin{proof}
概率质量函数由谱的计数（\(j=1\) 出现两次，其余 \(j=2,\dots,k\) 各出现一次）直接得到。
对弱收敛，令 \(f\) 为 \([0,1]\) 上连续函数，则
\[
\mathbb{E}\Bigl[f\Bigl(\frac{J_k}{k}\Bigr)\Bigr]
=\frac{2\cdot f(1/k)\cdot 1^{q}+\sum_{j=2}^{k} f(j/k)\,j^{q}}{1+\sum_{i=1}^{k} i^{q}}
=\frac{\sum_{j=1}^{k} f(j/k)\,j^{q}}{\sum_{i=1}^{k} i^{q}}+o(1),
\]
其中 \(o(1)\) 来自 \(j=1\) 的双重计数在归一化下的可忽略性。
将分子分母同除以 \(k^{q+1}\)，并把和式视为 Riemann 和，即得极限
\(
\int_{0}^{1}f(x)x^{q}\,\dd x\,/\,\int_0^1 x^{q}\dd x
\)
对应的 \(\mathrm{Beta}(q+1,1)\) 分布。证毕。
\end{proof}

\begin{conclusion}[缺陷位置的尺度极限：\texorpdfstring{$R_k/k\Rightarrow \mathrm{Beta}(1,q+1)$}{Rk/k -> Beta(1,q+1)} 与重心比例]\label{con:xi-terminal-max-weight-defect-localization}
沿结论 \ref{con:xi-terminal-max-weight-layer-fiber-spectrum-closed} 的参数化 \(x^{(r)}=(10)^r(01)^{k-r}\)，令 \(R_k\) 为随机缺陷参数（即 \(X=x^{(R_k)}\) 的指标）。
则在结论 \ref{con:xi-terminal-max-weight-escort-beta-limit} 的 \(\mathbb{P}_{k,q}\) 下，有
\[
\frac{R_k}{k}\ \Rightarrow\ \mathrm{Beta}(1,q+1),
\qquad \text{密度为 }(q+1)(1-x)^q\mathbf 1_{[0,1]}(x).
\]
特别地其重心比例满足
\[
\frac{1}{k}\mathbb{E}[R_k]\to \frac{1}{q+2},
\qquad
\frac{1}{k}\mathbb{E}[d_{2k}(X)]\to \frac{q+1}{q+2}.
\]
\end{conclusion}
\begin{proof}
由结论 \ref{con:xi-terminal-max-weight-layer-fiber-spectrum-closed}，在 \(r\ge 1\) 上有确定关系
\(
J_k=d_{2k}(x^{(R_k)})=k-R_k+1
\)。
并且
\(
\mathbb{P}_{k,q}(R_k=0)=\frac{1}{1+\sum_{i=1}^{k}i^q}\to 0
\)。
因此对任意连续函数 \(f:[0,1]\to\RR\)，可作变量代换 \(j=k-r+1\) 得
\[
\mathbb{E}_{k,q}\Bigl[f\Bigl(\frac{R_k}{k}\Bigr)\Bigr]
=\frac{f(0)+\sum_{r=1}^{k} f(r/k)\,(k-r+1)^q}{1+\sum_{i=1}^{k}i^q}
=\frac{\sum_{j=1}^{k} f\!\bigl(1-j/k+1/k\bigr)\,j^q}{\sum_{i=1}^{k}i^q}+o(1).
\]
将分子分母同除以 \(k^{q+1}\) 并取 Riemann 和极限，得到
\[
\lim_{k\to\infty}\mathbb{E}_{k,q}\Bigl[f\Bigl(\frac{R_k}{k}\Bigr)\Bigr]
=\frac{\int_{0}^{1} f(1-x)\,x^q\,\dd x}{\int_{0}^{1}x^q\,\dd x}
=\int_{0}^{1} f(u)\,(q+1)(1-u)^q\,\dd u,
\]
即 \(R_k/k\Rightarrow \mathrm{Beta}(1,q+1)\)。
由于 \(R_k/k\) 与 \(J_k/k\) 均一致有界于 \([0,1]\)，重心比例由弱收敛与 Beta 分布均值公式直接推出。证毕。
\end{proof}

\begin{conclusion}[\texorpdfstring{$t=\pm1$}{t=±1} 直线的零轨迹：重量参数的代数曲线量子化]\label{con:xi-terminal-zq-line-tpm1-zeros}
设 \(D_q(t,y)\) 为 \(F_q(t,y)=N_q(t,y)/D_q(t,y)\) 的分母多项式。
则在 \(q=2,3,4\) 上均有严格等式
\[
D_q(1,y)=D_q(-1,y),
\]
并且其零点集合在 \(y\)-平面上量子化为有限代数集合：
\[
D_2(\pm1,y)=-y(y-1)(y-2),
\]
\[
D_3(\pm1,y)=y(y-1)(y^2-4y+6),
\]
\[
D_4(\pm1,y)=-y(y-1)(y^3-5y^2+12y-24).
\]
因此，“出现特征模态 \(\lambda=\pm1\)”仅在 \(y\) 落入上述代数集合时发生。
并且在上述三式中，除去共同因子 \(y(y-1)\) 后出现的附加因子在常数项处满足
\[
(y-2)\big|_{y=0}=-2,\qquad
(y^2-4y+6)\big|_{y=0}=6,\qquad
(y^3-5y^2+12y-24)\big|_{y=0}=-24,
\]
即呈现统一的阶乘尺度律 \(({-}1)^{q+1}q!\)（对 \(q=2,3,4\)）。
\end{conclusion}
\begin{proof}[证明要点]
把 \(t=\pm1\) 代入各 \(D_q(t,y)\) 并在 \(\ZZ[y]\) 中因式分解即可；
等式 \(D_q(1,y)=D_q(-1,y)\) 反映分母在 \(t\mapsto -t\) 的偶性结构在该三阶族中的稳定性。证毕。
\end{proof}



\begin{conclusion}[符号纤维配分函数的三阶递推：Tribonacci 常数给出主振荡模的模平方]\label{con:xi-terminal-signed-fiber-partition-recurrence}
令 \(\Fold_m:\Omega_m=\{0,1\}^m\to X_m\) 为标准折叠映射，并令纤维多重度 \(d_m(x):=|\Fold_m^{-1}(x)|\)。
定义符号纤维配分函数
\[
B_m^{\mathrm{sgn}}:=\sum_{x\in X_m}(-1)^{|x|_1}\,d_m(x)
=\sum_{a\in\Omega_m}(-1)^{|\Fold_m(a)|_1}.
\]
由结论 \ref{con:xi-terminal-zm-order4-rational} 取 \(y=-1\) 得 \(B_m^{\mathrm{sgn}}=Z_m(-1)\)，
并且 \(y(y+1)=0\) 使四阶可积递推在该点退化为三阶递推（见结论 \ref{con:xi-terminal-zm-hankel-rank4}）。
则 \(B_m^{\mathrm{sgn}}\) 满足三阶线性递推
\[
B_m^{\mathrm{sgn}}=B_{m-1}^{\mathrm{sgn}}-B_{m-2}^{\mathrm{sgn}}-B_{m-3}^{\mathrm{sgn}}\qquad(m\ge 4),
\]
并具有初值 \(B_1^{\mathrm{sgn}}=0,B_2^{\mathrm{sgn}}=0,B_3^{\mathrm{sgn}}=-2\)。
相应特征多项式为
\[
P(\lambda):=\lambda^3-\lambda^2+\lambda+1.
\]
设 \(\alpha,\overline{\alpha}\) 为其非实共轭根，并令 \(\tau\) 为 Tribonacci 常数（即 \(\tau^3-\tau^2-\tau-1=0\) 的唯一正实根）。
则
\[
\abs{\alpha}^2=\tau,
\qquad
\rho_2:=\abs{\alpha}=\sqrt{\tau}=1.3562030656262951\ldots .
\]
写 \(\alpha=\rho_2 e^{i\vartheta}\)（\(\vartheta\in(0,\pi)\)），则相位满足
\[
\cos\vartheta=\frac{\tau+1}{2\tau^{3/2}}.
\]
因此存在常数 \(C\neq 0\) 与 \(\psi\in\RR\) 使得
\[
B_m^{\mathrm{sgn}}=C\,\tau^{m/2}\cos(m\vartheta+\psi)+O(\tau^{-m})\qquad(m\to\infty),
\]
从而该符号通道在指数尺度上以 \(\tau^{m/2}\) 量级增长并携带非平凡旋转相位；该振荡模态与正配分矩族 \(\{S_q(m)\}_q\) 完全正交，构成可审计的频域指纹。
\end{conclusion}
\begin{proof}[证明]
由结论 \ref{con:xi-terminal-zm-order4-rational} 的四阶递推在 \(y=-1\) 处取值，并注意 \(y(y+1)=0\) 消去四阶项，
得到所示三阶递推；初值由 \(m\le 3\) 的直接核算给出。
特征多项式为 \(P(\lambda)=\lambda^3-\lambda^2+\lambda+1\)。
令 \(r\) 为其唯一实根，则 \(r=-1/\tau\)：代入 \(\lambda=-1/t\) 得 \(P(-1/t)=-(t^3-t^2-t-1)/t^3\)，故 \(t=\tau\)。
由根乘积 \(r\,\alpha\,\overline{\alpha}=-1\) 得 \(\abs{\alpha}^2=\tau\)，从而 \(\rho_2=\sqrt{\tau}\)；
又由根和 \(r+\alpha+\overline{\alpha}=1\) 得 \(2\rho_2\cos\vartheta=1-r\)，即 \(\cos\vartheta=(\tau+1)/(2\tau^{3/2})\)。
一般项的主振荡展开由根分解直接推出。证毕。
\end{proof}

\begin{remark}[判别式与二次分辨率子域]\label{rem:xi-terminal-signed-fiber-partition-discriminant}
三次特征多项式 \(P(\lambda)=\lambda^3-\lambda^2+\lambda+1\) 的判别式为 \(-44\)，
故其分裂域的唯一二次子域为 \(\QQ(\sqrt{-11})\)。
该虚二次域为 Heegner 数 \(11\) 对应的类数 \(1\) 情形，因此奇偶扭曲通道的最小算术外延被约束在 \(\QQ(\sqrt{-11})\) 上。
\end{remark}

\begin{conclusion}[奇偶偏差的指数衰减：\texorpdfstring{$\kappa_2=\log(2/\sqrt{\tau})$}{kappa2=log(2/sqrt(tau))}]\label{con:xi-terminal-parity-gap-kappa2}
令 \(w_m(x):=d_m(x)/2^m\) 为均匀微态推前分布，并令 \(X\sim w_m\)。
记 \(H_m:=|X|_1\) 为稳定类型的 \(1\)-计数随机变量。
则奇偶偏差满足严格恒等式
\[
\boxed{\;
\mathbb P(H_m\equiv 0\!\!\!\pmod 2)-\mathbb P(H_m\equiv 1\!\!\!\pmod 2)
=\EE\bigl[(-1)^{H_m}\bigr]
=\frac{B_m^{\mathrm{sgn}}}{2^m}\; }.
\]
特别地存在常数 \(\kappa_2>0\) 使得
\[
\abs{\mathbb P(H_m\equiv 0\!\!\!\pmod 2)-\mathbb P(H_m\equiv 1\!\!\!\pmod 2)}
=\exp\!\bigl(-m(\kappa_2+o(1))\bigr),
\qquad
\kappa_2=\log\!\Bigl(\frac{2}{\rho_2}\Bigr)=\log\!\Bigl(\frac{2}{\sqrt{\tau}}\Bigr)\approx 0.3884582488.
\]
\end{conclusion}
\begin{proof}
由定义
\(
\EE[(-1)^{H_m}]=\sum_{x\in X_m}(-1)^{|x|_1}w_m(x)=2^{-m}B_m^{\mathrm{sgn}}
\)
得到恒等式。
再由结论 \ref{con:xi-terminal-signed-fiber-partition-recurrence} 的显式展开可得
\(
\abs{B_m^{\mathrm{sgn}}}=\exp(m\log\rho_2+o(m))
\)，
从而
\(
\abs{B_m^{\mathrm{sgn}}}/2^m=\exp(-m(\log(2/\rho_2)+o(1)))
\)，并由 \(\rho_2=\sqrt{\tau}\) 得 \(\kappa_2\) 的闭式与数值。证毕。
\end{proof}

\begin{conclusion}[模 \(3\) Fourier 通道的 \(8\) 阶可逆递推：自反特征多项式]\label{con:xi-terminal-mod3-fourier-order8-recurrence}
令 \(\omega=e^{2\pi i/3}\)，并定义根单位扭曲配分
\[
C_m^{(3)}:=\sum_{x\in X_m}\omega^{|x|_1}\,d_m(x)
=2^m\,\EE\!\left[\omega^{H_m}\right].
\]
则整数序列
\[
a_m:=2\Re C_m^{(3)},\qquad b_m:=\frac{2}{\sqrt3}\Im C_m^{(3)}
\]
满足同一条 \(8\) 阶齐次整系数递推
\[
\boxed{\;
s_m
=2s_{m-1}-s_{m-2}-2s_{m-3}+s_{m-4}-2s_{m-5}-s_{m-6}+2s_{m-7}-s_{m-8}\;}
\]
其特征多项式为
\[
\Pi^{(3)}(\lambda)=\lambda^8-2\lambda^7+\lambda^6+2\lambda^5-\lambda^4+2\lambda^3+\lambda^2-2\lambda+1,
\]
并满足严格自反性 \(\lambda^8\Pi^{(3)}(1/\lambda)=\Pi^{(3)}(\lambda)\)。
令 \(\rho_3\) 为 \(\Pi^{(3)}\) 的谱半径，则 \(\rho_3\approx 1.7030650701\)，从而存在常数 \(c_3>0\) 使
\[
\abs{\EE\!\left[\omega^{H_m}\right]}\le c_3\Bigl(\frac{\rho_3}{2}\Bigr)^m,
\qquad
\abs{\EE\!\left[\omega^{2H_m}\right]}\le c_3\Bigl(\frac{\rho_3}{2}\Bigr)^m.
\]
\end{conclusion}
\begin{proof}[证明要点]
与结论 \ref{con:xi-terminal-signed-fiber-partition-recurrence} 同理，
\(C_m^{(3)}\) 是有限状态带权转移矩阵的矩阵系数，故满足某个有限阶常系数线性递推；
该递推可由有限窗口的精确枚举与最小递推恢复得到，并在本实例中稳定化为所示 \(8\) 阶形式。
自反性可由结论 \ref{con:xi-terminal-mod3-integral-symplectic} 的整辛结构推出（或直接检验系数回文性）。
谱半径 \(\rho_3\) 控制 \(\abs{C_m^{(3)}}\) 的指数级增长，从而给出 \(\abs{\EE[\omega^{H_m}]}\) 的衰减界。证毕。
\end{proof}

\begin{conclusion}[模 \(4\) Fourier 通道的 \(8\) 阶递推与更慢的混合]\label{con:xi-terminal-mod4-fourier-order8-recurrence}
令 \(i=\sqrt{-1}\)，并定义
\[
C_m^{(4)}:=\sum_{x\in X_m} i^{|x|_1}\,d_m(x)
=2^m\,\EE\!\left[i^{H_m}\right].
\]
则 \(\Re C_m^{(4)}\) 与 \(\Im C_m^{(4)}\) 同样满足一条 \(8\) 阶齐次整系数递推，其特征多项式为
\[
\Pi^{(4)}(\lambda)=\lambda^8-2\lambda^7-\lambda^6+4\lambda^5+\lambda^4-\lambda^2-2\lambda+2.
\]
令 \(\rho_4\) 为 \(\Pi^{(4)}\) 的谱半径，则 \(\rho_4\approx 1.8315468016\)，从而存在常数 \(c_4>0\) 使
\[
\abs{\EE\!\left[i^{H_m}\right]}\le c_4\Bigl(\frac{\rho_4}{2}\Bigr)^m,
\qquad
\abs{\EE\!\left[(-i)^{H_m}\right]}\le c_4\Bigl(\frac{\rho_4}{2}\Bigr)^m.
\]
该通道的衰减率严格慢于奇偶通道（因 \(\rho_4>\rho_2\)）。
\end{conclusion}
\begin{proof}[证明要点]
同上：根单位扭曲配分是有限维带权转移矩阵的矩阵系数，故满足有限阶整系数递推；
本实例中其最小递推稳定化为阶 \(8\)。
谱半径 \(\rho_4\) 给出 \(\abs{C_m^{(4)}}\) 的指数界，从而得到非平凡 Fourier 系数的衰减率。证毕。
\end{proof}

\begin{conclusion}[同余类的指数趋于均匀：由非平凡谱半径给出全变差界]\label{con:xi-terminal-congruence-mixing-tv-bound}
对 \(k\in\{2,3,4\}\)，记 \(\mu_{m,k}\) 为 \(H_m\bmod k\) 的分布，并记其 Fourier 系数
\[
\widehat{\mu}_{m,k}(j):=\EE\!\left[e^{2\pi i j H_m/k}\right]\qquad(1\le j\le k-1).
\]
令 \(\rho_2=\sqrt{\tau}\)，并令 \(\rho_3,\rho_4\) 如结论 \ref{con:xi-terminal-mod3-fourier-order8-recurrence} 与结论 \ref{con:xi-terminal-mod4-fourier-order8-recurrence}。
则存在常数 \(c_k>0\) 使得对所有 \(m\ge 1\) 与所有 \(1\le j\le k-1\) 有
\[
\abs{\widehat{\mu}_{m,k}(j)}\le c_k\Bigl(\frac{\rho_k}{2}\Bigr)^m.
\]
于是到均匀分布 \(U_k\) 的全变差距离满足
\[
\boxed{\;
\norm{\mu_{m,k}-U_k}_{\mathrm{TV}}
\le \frac12\sum_{j=1}^{k-1}\abs{\widehat{\mu}_{m,k}(j)}
\le C_k\Bigl(\frac{\rho_k}{2}\Bigr)^m\;}
\]
其中常数 \(C_k>0\) 仅依赖于 \(k\) 与有限初值数据。
\end{conclusion}
\begin{proof}
由有限阿贝尔群 \(\ZZ/k\ZZ\) 上的 Fourier 反演与三角不等式得到
\(
\norm{\mu_{m,k}-U_k}_{\mathrm{TV}}\le \frac12\sum_{j\ne 0}\abs{\widehat{\mu}_{m,k}(j)}
\)。
而当 \(k=2,3,4\) 时，各非平凡系数分别由
\(\EE[(-1)^{H_m}]\)、\(\EE[\omega^{H_m}]\)、\(\EE[i^{H_m}]\) 与其共轭给出，
其指数衰减由结论 \ref{con:xi-terminal-parity-gap-kappa2}、
\ref{con:xi-terminal-mod3-fourier-order8-recurrence}、
\ref{con:xi-terminal-mod4-fourier-order8-recurrence} 给出。证毕。
\end{proof}

\begin{conclusion}[同余混合的严格层级：模数加深导致主 Fourier 缺口收缩]\label{con:xi-terminal-congruence-mixing-hierarchy}
上述三通道满足严格序
\[
\rho_2<\rho_3<\rho_4<2,
\]
从而对应的主 Fourier 缺口（混合指数）严格递减：
\[
\log\Bigl(\frac{2}{\rho_2}\Bigr)\approx 0.3884582488\;>\;
\log\Bigl(\frac{2}{\rho_3}\Bigr)\approx 0.1607175705\;>\;
\log\Bigl(\frac{2}{\rho_4}\Bigr)\approx 0.0879863239.
\]
因此，Fold 的均匀微态推前在同余可观测量上呈现刚性的谱层级：奇偶类最先均匀化，其后为模 \(3\)，模 \(4\) 最慢。
\end{conclusion}

\begin{conclusion}[模 \(3\) 通道的整辛结构：伴随矩阵保辛并强制谱互反成对]\label{con:xi-terminal-mod3-integral-symplectic}
令 \(s_m\) 表示结论 \ref{con:xi-terminal-mod3-fourier-order8-recurrence} 的整数序列（可取 \(s_m=a_m\) 或 \(s_m=b_m\)），并定义其 companion 矩阵
\[
C=
\begin{pmatrix}
2&-1&-2&1&-2&-1&2&-1\\
1&0&0&0&0&0&0&0\\
0&1&0&0&0&0&0&0\\
0&0&1&0&0&0&0&0\\
0&0&0&1&0&0&0&0\\
0&0&0&0&1&0&0&0\\
0&0&0&0&0&1&0&0\\
0&0&0&0&0&0&1&0
\end{pmatrix},
\qquad \det C=1.
\]
则存在一条整的满秩反对称形式 \(\Omega\in \Mat_{8}(\ZZ)\)（\(\det\Omega=1\)）使得
\[
C^{\mathsf T}\Omega C=\Omega.
\]
可取
\(
\Omega=\begin{psmallmatrix}0&J\\ -J^{\mathsf T}&0\end{psmallmatrix}
\)
其中
\[
J=
\begin{pmatrix}
1&0&0&0\\
-2&1&0&0\\
1&-2&1&0\\
2&1&-2&1
\end{pmatrix}.
\]
因此 \(C\) 为某个整辛群 \(\mathrm{Sp}(8,\ZZ;\Omega)\) 的元素；其特征值必在 \(\lambda\mapsto 1/\lambda\) 下成对，
从而特征多项式 \(\Pi^{(3)}\) 必为自反多项式（与结论 \ref{con:xi-terminal-mod3-fourier-order8-recurrence} 一致）。
在生成函数层面，这等价于相应分母多项式满足严格的 \(z\mapsto 1/z\) 反演对称。
\end{conclusion}
\begin{proof}[证明要点]
直接代入所示 \(\Omega\) 计算可得 \(C^{\mathsf T}\Omega C=\Omega\)。
保辛性推出 \(C\) 的谱在取逆下成对，故特征多项式系数回文，即 \(\lambda^8\Pi^{(3)}(1/\lambda)=\Pi^{(3)}(\lambda)\)。证毕。
\end{proof}

\begin{conclusion}[顶端三层纤维谱的稳定相：缺口 Fibonacci 封口与偶数支达到者计数平台]\label{con:xi-terminal-top3-fiber-spectrum-stable-multiplicity-phase}
记 \(D_m^{(1)}>\!D_m^{(2)}>\!D_m^{(3)}\) 为 Fold 纤维多重度谱的前三个不同取值（定义 \ref{def:pom-top-fiber-spectrum}），并令 \(F_1=F_2=1\) 为 Fibonacci 数列。
则当 \(m=2k\ge 12\) 时，
\[
D_{2k}^{(1)}-D_{2k}^{(2)}=F_{k-4},
\qquad
D_{2k}^{(1)}-D_{2k}^{(3)}=F_{k-3},
\]
并且归一化缺口在 \(k\to\infty\) 时收敛到常数比例
\[
\lim_{k\to\infty}\frac{D_{2k}^{(1)}-D_{2k}^{(2)}}{D_{2k}^{(1)}}=\varphi^{-6},
\qquad
\lim_{k\to\infty}\frac{D_{2k}^{(1)}-D_{2k}^{(3)}}{D_{2k}^{(1)}}=\varphi^{-5}.
\]
此外，偶数支在阈值之后进入稳定多重度相：对一切偶数 \(m\ge 18\)，达到 \(D_m^{(1)},D_m^{(2)},D_m^{(3)}\) 的稳定类型数分别稳定为
\[
\#\{x:d_m(x)=D_m^{(1)}\}=2,\qquad
\#\{x:d_m(x)=D_m^{(2)}\}=6,\qquad
\#\{x:d_m(x)=D_m^{(3)}\}=4.
\]
\end{conclusion}
\begin{proof}[证明]
Fibonacci 缺口闭式由定理 \ref{thm:pom-second-max-fiber-closed-form} 与定理 \ref{thm:pom-third-max-fiber-even-closed-form}（并见注 \ref{rem:pom-top-fiber-spectrum-ladder}）直接给出。
由 \(D_{2k}^{(1)}=F_{k+2}\)（定理 \ref{thm:pom-max-fiber}）与 Binet 公式 \(F_n\sim \varphi^n/\sqrt5\) 得到归一化缺口的极限比例常数。
\par
达到者计数平台可由同一 residue DP/精确枚举口径逐 \(m\) 核验：在可审计窗口内其数值稳定出现在偶数支 \(m\ge 18\)，并与最大达到者稳定化定理 \ref{thm:pom-max-achievers-phase-stabilization}（偶数支为 \(2\)）相容。证毕。
\end{proof}

\begin{conclusion}[\texorpdfstring{$S_q$}{Sq}-对称压缩等于对称幂表示：\texorpdfstring{$B_q=\mathrm{Sym}^q(M)$}{Bq=Sym^q(M)} 且幺模可逆并给出显式逆矩阵]\label{con:xi-terminal-disjointness-bq-sympower-unimodular-inverse}
令 \(M=\bigl(\begin{smallmatrix}1&1\\[2pt]1&0\end{smallmatrix}\bigr)\)，并令 \(B_q\) 为不交图 \(\mathcal D_q\) 的 \(S_q\)-对称压缩算子（命题 \ref{prop:pom-disjointness-symmetric-compression-bq}）。
则在齐次多项式基
\(
e_j:=x^{q-j}y^j\ (0\le j\le q)
\)
下，\(B_q\) 严格等于对称幂表示 \(\mathrm{Sym}^q(M)\) 的矩阵：
\[
\boxed{\;B_q=\mathrm{Sym}^q(M)\;}
\qquad\text{且}\qquad
B_q(i,j)=\binom{q-j}{i}.
\]
进一步，令 \(\mathbf{P}_q\) 为 Pascal 矩阵 \(\mathbf{P}_q(i,j)=\binom{j}{i}\)（\(i\le j\)，否则为 \(0\)），并令 \(\mathbf{J}_q\) 为列反转置换矩阵 \(\mathbf{J}_q(i,j)=\mathbf{1}_{i=q-j}\)。
则有分解
\[
B_q=\mathbf{P}_q\,\mathbf{J}_q,
\]
从而 \(B_q\) 为整矩阵幺模（unimodular），并且
\[
\det(B_q)=\det(\mathbf{J}_q)=(-1)^{q(q+1)/2}.
\]
其逆矩阵亦完全显式：对 \(0\le i,j\le q\)，若 \(i+j<q\) 则 \((B_q^{-1})(i,j)=0\)；
若 \(i+j\ge q\) 则
\[
(B_q^{-1})(i,j)=(-1)^{i+j-q}\binom{j}{q-i}.
\]
\end{conclusion}
\begin{proof}[证明]
由命题 \ref{prop:pom-disjointness-symmetric-compression-bq} 已知 \(B_q(i,j)=\binom{q-j}{i}\)。
另一方面，对称幂表示满足：对任意 \(a,b,c,d\)，矩阵 \(\bigl(\begin{smallmatrix}a&b\\ c&d\end{smallmatrix}\bigr)\) 在基 \(e_j\) 上的作用为
\[
e_j(x,y)\longmapsto (ax+by)^{q-j}(cx+dy)^j.
\]
取 \(M=\bigl(\begin{smallmatrix}1&1\\ 1&0\end{smallmatrix}\bigr)\) 得
\[
e_j\longmapsto (x+y)^{q-j}x^j=\sum_{i=0}^{q-j}\binom{q-j}{i}x^{q-i}y^i,
\]
故 \(\mathrm{Sym}^q(M)(i,j)=\binom{q-j}{i}=B_q(i,j)\)，从而 \(B_q=\mathrm{Sym}^q(M)\)。
\par
分解 \(B_q=\mathbf{P}_q\mathbf{J}_q\) 由
\((\mathbf{P}_q\mathbf{J}_q)(i,j)=\mathbf{P}_q(i,q-j)=\binom{q-j}{i}\)
立即得到。
\(\mathbf{P}_q\) 为下三角且对角元全为 \(1\)，故 \(\det(\mathbf{P}_q)=1\)，从而 \(\det(B_q)=\det(\mathbf{J}_q)=(-1)^{q(q+1)/2}\)。
显式逆矩阵由 \(\mathbf{P}_q^{-1}\) 的标准闭式与 \(\mathbf{J}_q^{-1}=\mathbf{J}_q\) 合成得到，并可直接逐项核验 \(B_qB_q^{-1}=I\)。证毕。
\end{proof}

\begin{conclusion}[\texorpdfstring{$B_q^n$}{Bq^n} 的全显式闭式：所有层间传输系数为 Fibonacci 幂的非负整系数组合]\label{con:xi-terminal-bq-power-fibonacci-nonnegative-sum}
对任意整数 \(n\ge 1\)，有严格恒等式
\[
B_q^{\,n}=\mathrm{Sym}^q(M^n),
\qquad
M^n=\begin{pmatrix}F_{n+1}&F_n\\[2pt]F_n&F_{n-1}\end{pmatrix}.
\]
因此在同一基下，对一切 \(0\le i,j\le q\) 有完全显式条目公式
\[
\boxed{\;
\bigl(B_q^{\,n}\bigr)(i,j)
=\sum_{v=0}^{\min(i,j)}
\binom{q-j}{i-v}\binom{j}{v}\,
F_{n+1}^{\,q-j-i+v}\,
F_{n}^{\,i+j-2v}\,
F_{n-1}^{\,v}\; }.
\]
\end{conclusion}
\begin{proof}[证明]
由结论 \ref{con:xi-terminal-disjointness-bq-sympower-unimodular-inverse}，\(B_q=\mathrm{Sym}^q(M)\)。
对称幂表示是一个表示函子，故 \(\mathrm{Sym}^q(M^n)=\mathrm{Sym}^q(M)^n=B_q^{\,n}\)。
又由 Fibonacci 矩阵恒等式 \(M^n=\bigl(\begin{smallmatrix}F_{n+1}&F_n\\ F_n&F_{n-1}\end{smallmatrix}\bigr)\) 得到所示条目展开：
对基向量 \(e_j\) 作用为
\[
e_j\longmapsto (F_{n+1}x+F_ny)^{q-j}(F_nx+F_{n-1}y)^j,
\]
分别作二项式展开并收集 \(y^i\) 的系数即得求和式；各项系数均为非负整数，且仅含 Fibonacci 幂单项式。证毕。
\end{proof}

\begin{conclusion}[边界层无混合定律：\texorpdfstring{$B_q^n$}{Bq^n} 的第一行与最后一列直接重建 Fibonacci 幂谱]\label{con:xi-terminal-bq-boundary-nonmixing-law}
在结论 \ref{con:xi-terminal-bq-power-fibonacci-nonnegative-sum} 的条目公式中，边界层出现完全退化的无混合结构：
对任意 \(0\le j\le q\) 有
\[
\bigl(B_q^{\,n}\bigr)(0,j)=F_{n+1}^{\,q-j}F_n^{\,j},
\]
并且对任意 \(0\le i\le q\) 有
\[
\bigl(B_q^{\,n}\bigr)(i,q)=\binom{q}{i}F_n^{\,q-i}F_{n-1}^{\,i}.
\]
因此 \(B_q^{\,n}\) 在两条边界上逐点编码了 Fibonacci 的幂次谱（并带有精确二项式权），可作为对称幂动力的最短结构基准。
\end{conclusion}
\begin{proof}[证明]
在结论 \ref{con:xi-terminal-bq-power-fibonacci-nonnegative-sum} 的求和式中取 \(i=0\) 时只有 \(v=0\) 可能，立即得到第一式。
取 \(j=q\) 时 \(\binom{q-j}{i-v}=\binom{0}{i-v}\) 强制 \(v=i\)，代回即得第二式。证毕。
\end{proof}

